\IfFileExists{stacks-project-book.cls}{%
\documentclass{stacks-project-book}
}{%
\documentclass{amsbook}
}

\usepackage{verbatim}
\newenvironment{reference}{\comment}{\endcomment}
\newenvironment{slogan}{\comment}{\endcomment}
\newenvironment{history}{\comment}{\endcomment}

\usepackage[all]{xy}

\xyoption{2cell}
\UseAllTwocells

\usepackage{multicol}

\usepackage{lmodern}
\usepackage[T1]{fontenc}
\usepackage[french]{babel}
\AtBeginDocument{\shorthandoff{?}}


\usepackage{hyperref}


\theoremstyle{plain}
\newtheorem{theorem}[subsection]{Théorème}
\newtheorem{proposition}[subsection]{Proposition}
\newtheorem{lemma}[subsection]{Lemme}

\theoremstyle{definition}
\newtheorem{definition}[subsection]{Définition}
\newtheorem{example}[subsection]{Exemple}
\newtheorem{exercise}[subsection]{Exercice}
\newtheorem{situation}[subsection]{Situation}

\theoremstyle{remark}
\newtheorem{remark}[subsection]{Remarque}
\newtheorem{remarks}[subsection]{Remarques}

\renewcommand{\proofname}{Démonstration}

\numberwithin{equation}{subsection}

\def\lim{\mathop{\mathrm{lim}}\nolimits}
\def\colim{\mathop{\mathrm{colim}}\nolimits}
\def\Spec{\mathop{\mathrm{Spec}}}
\def\Hom{\mathop{\mathrm{Hom}}\nolimits}
\def\Ext{\mathop{\mathrm{Ext}}\nolimits}
\def\SheafHom{\mathop{\mathcal{H}\!\mathit{om}}\nolimits}
\def\SheafExt{\mathop{\mathcal{E}\!\mathit{xt}}\nolimits}
\def\Sch{\mathit{Sch}}
\def\Mor{\mathop{\mathrm{Mor}}\nolimits}
\def\Ob{\mathop{\mathrm{Ob}}\nolimits}
\def\Sh{\mathop{\mathit{Sh}}\nolimits}
\def\NL{\mathop{N\!L}\nolimits}
\def\CH{\mathop{\mathrm{CH}}\nolimits}
\def\proetale{{pro\text{-}\acute{e}tale}}
\def\etale{{\acute{e}tale}}
\def\QCoh{\mathit{QCoh}}
\def\Ker{\mathop{\mathrm{Ker}}}
\def\Im{\mathop{\mathrm{Im}}}
\def\Coker{\mathop{\mathrm{Coker}}}
\def\Coim{\mathop{\mathrm{Coim}}}

\DeclareMathSymbol{\boxtimes}{\mathbin}{AMSa}{"02}

\def\QCohstack{\mathcal{QC}\!\mathit{oh}}
\def\Cohstack{\mathcal{C}\!\mathit{oh}}
\def\Spacesstack{\mathcal{S}\!\mathit{paces}}
\def\Quotfunctor{\mathrm{Quot}}
\def\Hilbfunctor{\mathrm{Hilb}}
\def\Curvesstack{\mathcal{C}\!\mathit{urves}}
\def\Polarizedstack{\mathcal{P}\!\mathit{olarized}}
\def\Complexesstack{\mathcal{C}\!\mathit{omplexes}}
\def\Pic{\mathop{\mathrm{Pic}}\nolimits}
\def\Picardstack{\mathcal{P}\!\mathit{ic}}
\def\Picardfunctor{\mathrm{Pic}}
\def\Deformationcategory{\mathcal{D}\!\mathit{ef}}
\begin{document}
\title{Le projet Stacks --- édition française, partie 10}
\maketitle
\tableofcontents

% OK, on commence ici.
%

\chapter{Amorçage}


\phantomsection
\label{bootstrap-section-phantom}





\section{Introduction}
\label{bootstrap-section-introduction}

\noindent
Dans ce chapitre, nous utilisons les résultats des sections précédentes pour
donner des critères assurant qu'un préfaisceau d'ensembles sur la catégorie des
schémas est un espace algébrique. Une partie de ces résultats provient des
travaux d'Artin, voir \cite{ArtinI}, \cite{ArtinII},
\cite{Artin-Theorem-Representability},
\cite{Artin-Construction-Techniques},
\cite{Artin-Algebraic-Spaces},
\cite{Artin-Algebraic-Approximation},
\cite{Artin-Implicit-Function},
et \cite{ArtinVersal}.
Notre méthode consistera toutefois à employer autant que possible des arguments
semblables à ceux de l'article de Keel et Mori, voir
\cite{K-M}.

\section{Conventions}
\label{bootstrap-section-conventions}

\noindent
Nous supposons une fois pour toutes que tous les schémas considérés appartiennent
à un grand site fppf $\Sch_{fppf}$. De même, tout anneau $A$ considéré
est tel que $\Spec(A)$ est (isomorphe à) un
objet de ce grand site.

\medskip\noindent
Soient $S$ un schéma et $X$ un espace algébrique sur $S$.
Dans ce chapitre et dans les suivants, nous noterons $X \times_S X$
le produit de $X$ avec lui-même (dans la catégorie des espaces
algébriques sur $S$), plutôt que $X \times X$.




\section{Morphismes représentables par des espaces algébriques}
\label{bootstrap-section-morphism-representable-by-spaces}

\noindent
Nous définissons ici la notion de préfaisceau relativement représentable
par des espaces algébriques au-dessus d'un autre, puis en établissons quelques propriétés.

\begin{definition}
\label{bootstrap-definition-morphism-representable-by-spaces}
Soit $S$ un schéma appartenant à $\Sch_{fppf}$.
Soient $F$, $G$ des préfaisceaux sur $\Sch_{fppf}/S$.
On dit qu'un morphisme $a : F \to G$ est
{\it représentable par des espaces algébriques}
si, pour tout $U \in \Ob((\Sch/S)_{fppf})$ et
tout $\xi : U \to G$, le produit fibré $U \times_{\xi, G} F$
est un espace algébrique.
\end{definition}

\noindent
Voici une vérification élémentaire.

\begin{lemma}
\label{bootstrap-lemma-morphism-spaces-is-representable-by-spaces}
Soit $S$ un schéma.
Soit $f : X \to Y$ un morphisme d'espaces algébriques sur $S$.
Alors $f$ est représentable par des espaces algébriques.
\end{lemma}

\begin{proof}
C'est formel. Cela repose sur le fait que
la catégorie des espaces algébriques sur $S$ admet des produits fibrés, voir
Espaces, lemme \ref{spaces-lemma-fibre-product-spaces}.
\end{proof}

\begin{lemma}
\label{bootstrap-lemma-base-change-transformation}
\begin{slogan}
Tout changement de base d'un morphisme de préfaisceaux représentable
par des espaces algébriques est représentable par des espaces algébriques.
\end{slogan}
Soit $S$ un schéma. Soit
$$
\xymatrix{
G' \times_G F \ar[r] \ar[d]^{a'} & F \ar[d]^a \\
G' \ar[r] & G
}
$$
soit un carré cartésien de préfaisceaux sur $(\Sch/S)_{fppf}$.
Si $a$ est représentable par des espaces algébriques, alors $a'$ l'est aussi.
\end{lemma}

\begin{proof}
Omis. Indication : c'est formel.
\end{proof}

\begin{lemma}
\label{bootstrap-lemma-representable-by-spaces-transformation-to-sheaf}
Soit $S$ un schéma appartenant à $\Sch_{fppf}$.
Soient $F, G : (\Sch/S)_{fppf}^{opp} \to \textit{Ensembles}$.
Soit $a : F \to G$ représentable par des espaces algébriques.
Si $G$ est un faisceau, alors $F$ en est un aussi.
\end{lemma}

\begin{proof}
(Même démonstration que celle de
Espaces, lemme \ref{spaces-lemma-representable-transformation-to-sheaf}.)
Soit $\{\varphi_i : T_i \to T\}$ un recouvrement du site
$(\Sch/S)_{fppf}$.
Soient $s_i \in F(T_i)$ satisfaisant à la condition de faisceau.
Alors les $\sigma_i = a(s_i) \in G(T_i)$ satisfont eux aussi à la condition
de faisceau. Il existe donc un unique $\sigma \in G(T)$ tel
que $\sigma_i = \sigma|_{T_i}$. Par hypothèse,
$F' = h_T \times_{\sigma, G, a} F$ est un faisceau.
Notons que les $(\varphi_i, s_i) \in F'(T_i)$ satisfont eux aussi à la
condition de faisceau et proviennent donc d'un unique
$(\text{id}_T, s) \in F'(T)$. Manifestement, $s$ est la section de
$F$ cherchée.
\end{proof}

\begin{lemma}
\label{bootstrap-lemma-representable-by-spaces-transformation-diagonal}
Soit $S$ un schéma appartenant à $\Sch_{fppf}$.
Soient $F, G : (\Sch/S)_{fppf}^{opp} \to \textit{Ensembles}$.
Soit $a : F \to G$ représentable par des espaces algébriques.
Alors $\Delta_{F/G} : F \to F \times_G F$ est représentable par
des espaces algébriques.
\end{lemma}

\begin{proof}
(Même démonstration que celle de
Espaces, lemme \ref{spaces-lemma-representable-transformation-diagonal}.)
Soit $U$ un schéma. Soit $\xi = (\xi_1, \xi_2) \in (F \times_G F)(U)$.
Posons $\xi' = a(\xi_1) = a(\xi_2) \in G(U)$.
Par hypothèse, il existe un espace algébrique $V$ et un morphisme $V \to U$
qui représentent le produit fibré $U \times_{\xi', G} F$.
En particulier, les éléments $\xi_1, \xi_2$ donnent des morphismes
$f_1, f_2 : U \to V$ au-dessus de $U$. Comme $V$ représente le
produit fibré $U \times_{\xi', G} F$ et que
$\xi' = a \circ \xi_1 = a \circ \xi_2$,
on voit que, si $g : U' \to U$ est un morphisme, alors
$$
g^*\xi_1 = g^*\xi_2
\Leftrightarrow
f_1 \circ g = f_2 \circ g.
$$
Autrement dit, $U \times_{\xi, F \times_G F} F$
est représenté par $V \times_{\Delta, V \times V, (f_1, f_2)} U$,
qui est un espace algébrique.
\end{proof}

\noindent
La démonstration du
lemme \ref{bootstrap-lemma-representable-by-spaces-over-space}
ci-dessous est en fait quelque peu délicate. En effet,
nous ne pouvons employer l'argument de la démonstration de
Espaces, lemme \ref{spaces-lemma-representable-over-space},
car nous ne savons pas encore qu'un composé de transformations
représentables par des espaces algébriques est représentable par des espaces
algébriques. Nous utiliserons précisément ce lemme pour établir cet énoncé.

\begin{lemma}
\label{bootstrap-lemma-representable-by-spaces-over-space}
Soit $S$ un schéma appartenant à $\Sch_{fppf}$.
Soient $F, G : (\Sch/S)_{fppf}^{opp} \to \textit{Ensembles}$.
Soit $a : F \to G$ représentable par des espaces algébriques.
Si $G$ est un espace algébrique, alors $F$ en est un aussi.
\end{lemma}

\begin{proof}
Nous avons vu au
lemme \ref{bootstrap-lemma-representable-by-spaces-transformation-to-sheaf}
que $F$ est un faisceau.

\medskip\noindent
Soit $U$ un schéma et soit $U \to G$ un morphisme étale surjectif.
Alors $U \times_G F$ est un espace algébrique. Soit $W$ un schéma
et soit $W \to U \times_G F$ un morphisme étale surjectif.

\medskip\noindent
Affirmons d'abord que $W \to F$ est représentable.
Pour le voir, soient $X$ un schéma et $X \to F$ un morphisme.
Alors
$$
W \times_F X = W \times_{U \times_G F} U \times_G F \times_F X
= W \times_{U \times_G F} (U \times_G X)
$$
Comme $U \times_G F$ et $G$ sont tous deux des espaces algébriques, on voit
qu'il s'agit d'un schéma.

\medskip\noindent
Affirmons ensuite que $W \to F$ est étale et surjectif (cela a désormais
un sens puisque nous savons qu'il est représentable). Cela résulte de la
formule précédente : $W \to U \times_G F$ et $U \to G$
sont tous deux étales et surjectifs, donc
$W \times_{U \times_G F} (U \times_G X) \to U \times_G X$ et
$U \times_G X \to X$ sont étales et surjectifs, et le composé de
morphismes étales surjectifs est étale et surjectif.

\medskip\noindent
Posons $R = W \times_F W$. D'après ce qui précède, $R$ est un schéma et
les projections $t, s : R \to W$
sont étales. Il est clair que $R$ est une relation d'équivalence et que
$W \to F$ est un épimorphisme de faisceaux. Ainsi, $R$ est une relation
d'équivalence étale et $F = W/R$. Par conséquent, $F$ est un espace algébrique d'après
Espaces,
théorème \ref{spaces-theorem-presentation}.
\end{proof}

\begin{lemma}
\label{bootstrap-lemma-representable-by-spaces}
Soit $S$ un schéma.
Soit $a : F \to G$ un morphisme de préfaisceaux sur $(\Sch/S)_{fppf}$.
Supposons $a : F \to G$ représentable par des espaces algébriques.
Si $X$ est un espace algébrique sur $S$ et si $X \to G$ est un morphisme de préfaisceaux,
alors $X \times_G F$ est un espace algébrique.
\end{lemma}

\begin{proof}
D'après le lemme \ref{bootstrap-lemma-base-change-transformation}, la transformation
$X \times_G F \to X$ est représentable par des espaces algébriques. Sa source est donc
un espace algébrique d'après le
lemme \ref{bootstrap-lemma-representable-by-spaces-over-space}.
\end{proof}

\begin{lemma}
\label{bootstrap-lemma-composition-transformation}
Soit $S$ un schéma.
Soient
$$
\xymatrix{
F \ar[r]^a & G \ar[r]^b & H
}
$$
des morphismes de préfaisceaux sur $(\Sch/S)_{fppf}$.
Si $a$ et $b$ sont représentables par des espaces algébriques, alors
$b \circ a$ l'est aussi.
\end{lemma}

\begin{proof}
Soit $T$ un schéma sur $S$, et soit $T \to H$ un morphisme.
Par hypothèse, $T \times_H G$ est un espace algébrique. Le
lemme \ref{bootstrap-lemma-representable-by-spaces}
montre donc que $T \times_H F = (T \times_H G) \times_G F$ est lui aussi
un espace algébrique.
\end{proof}

\begin{lemma}
\label{bootstrap-lemma-product-transformations}
Soit $S$ un schéma.
Soient $F_i, G_i : (\Sch/S)_{fppf}^{opp} \to \textit{Ensembles}$, $i = 1, 2$.
Soient $a_i : F_i \to G_i$, $i = 1, 2$,
représentables par des espaces algébriques.
Alors
$$
a_1 \times a_2 : F_1 \times F_2 \longrightarrow G_1 \times G_2
$$
est représentable par des espaces algébriques.
\end{lemma}

\begin{proof}
Écrivons $a_1 \times a_2$ comme le composé
$F_1 \times F_2 \to G_1 \times F_2 \to G_1 \times G_2$.
La première flèche est le changement de base de $a_1$ par le morphisme
$G_1 \times F_2 \to G_1$, et la seconde
est le changement de base de $a_2$ par le morphisme
$G_1 \times G_2 \to G_2$. Ce lemme est donc une conséquence formelle des
lemmes \ref{bootstrap-lemma-composition-transformation}
et \ref{bootstrap-lemma-base-change-transformation}.
\end{proof}

\begin{lemma}
\label{bootstrap-lemma-representable-by-spaces-permanence}
Soit $S$ un schéma. Soient $a : F \to G$ et $b : G \to H$ des
transformations de foncteurs $(\Sch/S)_{fppf}^{opp} \to \textit{Ensembles}$.
Supposons que
\begin{enumerate}
\item $\Delta : G \to G \times_H G$ soit représentable
par des espaces algébriques, et que
\item $b \circ a : F \to H$ soit représentable par des espaces algébriques.
\end{enumerate}
Alors $a$ est représentable par des espaces algébriques.
\end{lemma}

\begin{proof}
Soit $U$ un schéma sur $S$ et soit $\xi \in G(U)$. Alors
$$
U \times_{\xi, G, a} F =
(U \times_{b(\xi), H, b \circ a} F) \times_{(\xi, a), (G \times_H G), \Delta} G
$$
Le résultat découle donc du lemme \ref{bootstrap-lemma-representable-by-spaces}.
\end{proof}

\begin{lemma}
\label{bootstrap-lemma-glueing-sheaves}
Soit $S \in \Ob(\Sch_{fppf})$. Soit $F$ un préfaisceau d'ensembles sur
$(\Sch/S)_{fppf}$. Supposons que
\begin{enumerate}
\item $F$ soit un faisceau pour la topologie de Zariski sur $(\Sch/S)_{fppf}$,
\item il existe un ensemble d'indices $I$ et des sous-foncteurs $F_i \subset F$ tels que
\begin{enumerate}
\item chaque $F_i$ soit un faisceau fppf,
\item chaque $F_i \to F$ soit représentable par des espaces algébriques,
\item $\coprod F_i \to F$ devienne surjectif après faisceautisation fppf.
\end{enumerate}
\end{enumerate}
Alors $F$ est un faisceau fppf.
\end{lemma}

\begin{proof}
Soient $T \in \Ob((\Sch/S)_{fppf})$ et $s \in F(T)$. D'après (2)(c),
il existe un recouvrement fppf $\{T_j \to T\}$ tel que
$s|_{T_j}$ soit une section de $F_{\alpha(j)}$ pour un certain $\alpha(j) \in I$.
Soit $W_j \subset T$ l'image de $T_j \to T$ ;
c'est un sous-schéma ouvert d'après Morphismes, lemme \ref{morphisms-lemma-fppf-open}.
D'après (2)(b),
$F_{\alpha(j)} \times_{F, s|_{W_j}} W_j \to W_j$ est un monomorphisme
d'espaces algébriques par lequel $T_j$ se factorise. Puisque $\{T_j \to W_j\}$
est un recouvrement fppf, on en déduit que
$F_{\alpha(j)} \times_{F, s|_{W_j}} W_j = W_j$, in other words
$s|_{W_j} \in F_{\alpha(j)}(W_j)$. On en conclut donc que
$\coprod F_i \to F$ est surjectif pour la topologie de Zariski.

\medskip\noindent
Soit $\{T_j \to T\}$ un recouvrement fppf de $(\Sch/S)_{fppf}$.
Soient $s, s' \in F(T)$ tels que $s|_{T_j} = s'|_{T_j}$ pour tout $j$.
Montrons que $s, s'$ sont égaux. Comme $F$ est un faisceau de Zariski d'après (1),
nous pouvons raisonner localement pour la topologie de Zariski sur $T$. D'après l'alinéa précédent,
nous pouvons supposer qu'il existe $i$ tel que $s \in F_i(T)$. Alors
$s'|_{T_j}$ est une section de $F_i$. D'après (2)(b),
$F_{i} \times_{F, s'} T \to T$ est un monomorphisme d'espaces algébriques
par lequel tous les $T_j$ se factorisent. On en conclut que
$s' \in F_i(T)$. Comme $F_i$ est un faisceau pour la topologie fppf,
on obtient $s = s'$.

\medskip\noindent
Soit $\{T_j \to T\}$ un recouvrement fppf de $(\Sch/S)_{fppf}$ et soient
$s_j \in F(T_j)$ tels que
$s_j|_{T_j \times_T T_{j'}} = s_{j'}|_{T_j \times_T T_{j'}}$. D'après l'hypothèse
(2)(c), nous pouvons raffiner le recouvrement et supposer que $s_j \in F_{\alpha(j)}(T_j)$
pour un certain $\alpha(j) \in I$. Soit $W_j \subset T$ l'image de $T_j \to T$,
qui est un sous-schéma ouvert d'après Morphismes, lemme \ref{morphisms-lemma-fppf-open}.
Alors $\{T_j \to W_j\}$ est un recouvrement fppf. Comme $F_{\alpha(j)}$ est un
sous-préfaisceau de $F$, les deux restrictions de $s_j$ à
$T_j \times_{W_j} T_j$ coïncident comme éléments de
$F_{\alpha(j)}(T_j \times_{W_j} T_j)$. La condition de faisceau pour
$F_{\alpha(j)}$ implique donc l'existence d'un $s'_j \in F_{\alpha(j)}(W_j)$
dont la restriction à $T_j$ est $s_j$. Pour toute paire d'indices
$j$ et $j'$, les sections $s'_j|_{W_j \cap W_{j'}}$ et
$s'_{j'}|_{W_j \cap W_{j'}}$ de $F$ coïncident d'après l'alinéa
précédent. Le fait que $F$ soit un faisceau de Zariski achève
la démonstration.
\end{proof}





\section{Propriétés des morphismes de préfaisceaux représentables par des espaces algébriques}
\label{bootstrap-section-representable-by-spaces-properties}

\noindent
Voici la définition qui convient.

\begin{definition}
\label{bootstrap-definition-property-transformation}
Soit $S$ un schéma. Soit $a : F \to G$ un morphisme de préfaisceaux sur
$(\Sch/S)_{fppf}$ représentable par des espaces algébriques.
Soit $\mathcal{P}$ une propriété des morphismes d'espaces algébriques qui
\begin{enumerate}
\item est stable par tout changement de base, et
\item est locale sur la base pour la topologie fppf, voir
Descente sur les espaces,
définition \ref{spaces-descent-definition-property-morphisms-local}.
\end{enumerate}
Dans ce cas, on dit que $a$ possède {\it la propriété $\mathcal{P}$} si, pour tout
schéma $U$ et tout $\xi : U \to G$, le morphisme d'espaces algébriques ainsi obtenu
$U \times_G F \to U$ possède la propriété $\mathcal{P}$.
\end{definition}

\noindent
Il importe de noter que nous n'utiliserons cette définition que pour
des propriétés de morphismes stables par changement de base et
locales sur la base pour la topologie fppf. Ce n'est
pas que la définition n'ait autrement aucun sens ; c'est plutôt que
nous pouvons vouloir donner une autre définition, mieux
adaptée à la propriété envisagée.

\medskip\noindent
La définition précédente s'applique\footnote{La stabilité par changement
de base résulte de
Morphismes d'espaces, lemmes
\ref{spaces-morphisms-lemma-base-change-surjective},
\ref{spaces-morphisms-lemma-base-change-quasi-compact},
\ref{spaces-morphisms-lemma-base-change-etale},
\ref{spaces-morphisms-lemma-base-change-smooth},
\ref{spaces-morphisms-lemma-base-change-flat},
\ref{spaces-morphisms-lemma-base-change-separated},
\ref{spaces-morphisms-lemma-base-change-finite-type},
\ref{spaces-morphisms-lemma-base-change-quasi-finite},
\ref{spaces-morphisms-lemma-base-change-finite-presentation},
\ref{spaces-morphisms-lemma-base-change-affine},
\ref{spaces-morphisms-lemma-base-change-proper}, et
Espaces, lemme
\ref{spaces-lemma-base-change-immersions}.
Le caractère local sur la base pour la topologie fppf résulte de
Descente sur les espaces, lemmes
\ref{spaces-descent-lemma-descending-property-surjective},
\ref{spaces-descent-lemma-descending-property-quasi-compact},
\ref{spaces-descent-lemma-descending-property-etale},
\ref{spaces-descent-lemma-descending-property-smooth},
\ref{spaces-descent-lemma-descending-property-flat},
\ref{spaces-descent-lemma-descending-property-separated},
\ref{spaces-descent-lemma-descending-property-finite-type},
\ref{spaces-descent-lemma-descending-property-quasi-finite},
\ref{spaces-descent-lemma-descending-property-locally-finite-presentation},
\ref{spaces-descent-lemma-descending-property-affine},
\ref{spaces-descent-lemma-descending-property-proper}, et
\ref{spaces-descent-lemma-descending-property-closed-immersion}.
}
par exemple aux propriétés d'être
« surjectif »,
« quasi-compact »,
« étale »,
« lisse »,
« plat »,
« séparé »,
« (localement) de type fini »,
« (localement) quasi-fini »,
« (localement) de présentation finie »,
« affine »,
« propre » et
« une immersion fermée ».
Autrement dit, $a$ est
{\it surjectif}
(resp.\ {\it quasi-compact},
{\it étale},
{\it lisse},
{\it plat},
{\it séparé},
{\it (localement) de type fini},
{\it (localement) quasi-fini},
{\it (localement) de présentation finie},
{\it propre},
{\it une immersion fermée})
si, pour tout schéma $T$ et tout morphisme $\xi : T \to G$,
le morphisme d'espaces algébriques $T \times_{\xi, G} F \to T$ est
surjectif
(resp.\ quasi-compact,
étale,
plat,
séparé,
(localement) de type fini,
(localement) quasi-fini,
(localement) de présentation finie,
propre,
une immersion fermée).

\medskip\noindent
Vérifions ensuite la compatibilité avec les notions déjà définies. D'après le
lemme \ref{bootstrap-lemma-morphism-spaces-is-representable-by-spaces},
tout morphisme entre espaces algébriques sur $S$ est représentable par
des espaces algébriques. D'autre part, d'après
Morphismes d'espaces,
lemme \ref{spaces-morphisms-lemma-surjective-local}
(resp.\ \ref{spaces-morphisms-lemma-quasi-compact-local},
\ref{spaces-morphisms-lemma-etale-local},
\ref{spaces-morphisms-lemma-smooth-local},
\ref{spaces-morphisms-lemma-flat-local},
\ref{spaces-morphisms-lemma-separated-local},
\ref{spaces-morphisms-lemma-finite-type-local},
\ref{spaces-morphisms-lemma-quasi-finite-local},
\ref{spaces-morphisms-lemma-finite-presentation-local},
\ref{spaces-morphisms-lemma-affine-local},
\ref{spaces-morphisms-lemma-proper-local},
\ref{spaces-morphisms-lemma-closed-immersion-local})
la définition précédente de
surjectif
(resp.\ quasi-compact,
étale,
lisse,
plat,
séparé,
(localement) de type fini,
(localement) quasi-fini,
(localement) de présentation finie,
affine,
propre,
immersion fermée)
coïncide avec la définition déjà donnée pour les morphismes
d'espaces algébriques.

\medskip\noindent
Voici maintenant quelques lemmes formels.

\begin{lemma}
\label{bootstrap-lemma-base-change-transformation-property}
Soit $S$ un schéma.
Soit $\mathcal{P}$ une propriété comme dans la
définition \ref{bootstrap-definition-property-transformation}.
Soit
$$
\xymatrix{
G' \times_G F \ar[r] \ar[d]^{a'} & F \ar[d]^a \\
G' \ar[r] & G
}
$$
un carré cartésien de préfaisceaux sur $(\Sch/S)_{fppf}$.
Si $a$ est représentable par des espaces algébriques et possède $\mathcal{P}$,
alors $a'$ aussi.
\end{lemma}

\begin{proof}
Omis. Indication : c'est formel.
\end{proof}

\begin{lemma}
\label{bootstrap-lemma-composition-transformation-property}
Soit $S$ un schéma.
Soit $\mathcal{P}$ une propriété comme dans la
définition \ref{bootstrap-definition-property-transformation},
et supposons $\mathcal{P}$ stable par composition.
Soient
$$
\xymatrix{
F \ar[r]^a & G \ar[r]^b & H
}
$$
des morphismes de préfaisceaux sur $(\Sch/S)_{fppf}$.
Si $a$ et $b$ sont représentables par des espaces algébriques et possèdent
$\mathcal{P}$, alors $b \circ a$ la possède aussi.
\end{lemma}

\begin{proof}
Omis. Indication : voir le
lemme \ref{bootstrap-lemma-composition-transformation}
et utiliser la stabilité par composition.
\end{proof}

\begin{lemma}
\label{bootstrap-lemma-product-transformations-property}
Soit $S$ un schéma.
Soient $F_i, G_i : (\Sch/S)_{fppf}^{opp} \to \textit{Ensembles}$,
$i = 1, 2$.
Soient $a_i : F_i \to G_i$, $i = 1, 2$, représentables par des espaces algébriques.
Soit $\mathcal{P}$ une propriété comme dans la
définition \ref{bootstrap-definition-property-transformation},
stable par composition.
Si $a_1$ et $a_2$ possèdent la propriété $\mathcal{P}$, alors
$a_1 \times a_2 : F_1 \times F_2 \longrightarrow G_1 \times G_2$ la possède aussi.
\end{lemma}

\begin{proof}
Notons que l'énoncé a un sens d'après le
lemme \ref{bootstrap-lemma-product-transformations}.
Démonstration omise.
\end{proof}

\begin{lemma}
\label{bootstrap-lemma-transformations-property-implication}
Soit $S$ un schéma.
Soient $F, G : (\Sch/S)_{fppf}^{opp} \to \textit{Ensembles}$.
Soit $a : F \to G$ une transformation de foncteurs représentable par
des espaces algébriques.
Soient $\mathcal{P}$, $\mathcal{P}'$ des propriétés comme dans la
définition \ref{bootstrap-definition-property-transformation}.
Supposons que, pour tout morphisme $f : X \to Y$ d'espaces algébriques sur $S$,
on ait $\mathcal{P}(f) \Rightarrow \mathcal{P}'(f)$.
Si $a$ possède la propriété $\mathcal{P}$, alors
$a$ possède la propriété $\mathcal{P}'$.
\end{lemma}

\begin{proof}
C'est formel.
\end{proof}

\begin{lemma}
\label{bootstrap-lemma-surjective-flat-locally-finite-presentation}
Soit $S$ un schéma.
Soient $F, G : (\Sch/S)_{fppf}^{opp} \to \textit{Ensembles}$ des faisceaux.
Soit $a : F \to G$ représentable par des espaces algébriques, plat,
localement de présentation finie et surjectif.
Alors $a : F \to G$ est surjectif comme morphisme de faisceaux.
\end{lemma}

\begin{proof}
Soit $T$ un schéma sur $S$ et soit $g : T \to G$ un point à valeurs dans
$T$ de $G$. Par hypothèse, $T' = F \times_G T$ est un espace algébrique et
le morphisme $T' \to T$ est un morphisme d'espaces algébriques plat,
localement de présentation finie et surjectif.
Soit $U \to T'$ un morphisme étale surjectif, où $U$ est un schéma.
Par définition des morphismes plats d'espaces algébriques,
le morphisme de schémas $U \to T$ est alors plat. Il en va de même pour la
propriété « localement de présentation finie ». Le morphisme $U \to T$ est aussi surjectif,
voir
Morphismes d'espaces, lemme \ref{spaces-morphisms-lemma-surjective-local}.
Ainsi, $\{U \to T\}$ est un recouvrement fppf tel
que $g|_U \in G(U)$ provienne d'un élément de $F(U)$, à savoir du
morphisme $U \to T' \to F$. Cela prouve que le morphisme est surjectif comme
morphisme de faisceaux, voir
Sites, définition \ref{sites-definition-sheaves-injective-surjective}.
\end{proof}




\section{Amorçage de la diagonale}
\label{bootstrap-section-bootstrap-diagonal}

\noindent
Dans cette section, nous prouvons que la diagonale d'un faisceau $F$ sur
$(\Sch/S)_{fppf}$ est représentable dès qu'il existe
un « recouvrement fppf » de $F$ par un schéma ou par un espace algébrique, voir le
lemme \ref{bootstrap-lemma-bootstrap-diagonal}.

\begin{lemma}
\label{bootstrap-lemma-representable-diagonal}
\begin{slogan}
La diagonale d'un préfaisceau est représentable par des espaces algébriques si et seulement si
tout morphisme d'un schéma vers ce préfaisceau est représentable par des espaces algébriques.
\end{slogan}
Soit $S$ un schéma.
Soit $F$ un préfaisceau sur $(\Sch/S)_{fppf}$.
Les conditions suivantes sont équivalentes :
\begin{enumerate}
\item $\Delta_F : F \to F \times F$ est représentable par des espaces algébriques ;
\item pour tout schéma $T$, tout morphisme $T \to F$ est représentable par des espaces
algébriques ;
\item pour tout espace algébrique $X$, tout morphisme $X \to F$ est représentable
par des espaces algébriques.
\end{enumerate}
\end{lemma}

\begin{proof}
Supposons (1). Soit $X \to F$ comme en (3). Soit $T$ un schéma et soit
$T \to F$ un morphisme. On a alors
$$
T \times_F X = (T \times_S X) \times_{F \times F, \Delta} F
$$
qui est un espace algébrique d'après le
lemme \ref{bootstrap-lemma-representable-by-spaces}
et (1). Ainsi, $X \to F$ est représentable, c'est-à-dire que (3) est vérifiée.
L'implication (3) $\Rightarrow$ (2) est immédiate. Supposons (2).
Soit $T$ un schéma et soit $(a, b) : T \to F \times F$ un morphisme.
Alors
$$
F \times_{\Delta_F, F \times F} T =
(T \times_{a, F, b} T) \times_{T \times T, \Delta_T} T
$$
qui est un espace algébrique par hypothèse. Ainsi, $\Delta_F$ est
représentable par des espaces algébriques, c'est-à-dire que (1) est vérifiée.
\end{proof}

\noindent
En particulier, si $F$ est un préfaisceau satisfaisant aux conditions équivalentes du
lemme, alors, pour tout morphisme $X \to F$ où $X$ est un espace algébrique,
il est légitime de dire que $X \to F$ est surjectif (resp.\ étale, plat,
localement de présentation finie) en appliquant la
définition \ref{bootstrap-definition-property-transformation}.

\medskip\noindent
Avant de procéder à l'amorçage proprement dit, établissons un petit lemme utile.

\begin{lemma}
\label{bootstrap-lemma-after-fppf-sep-lqf}
Soit $S$ un schéma.
Soit
$$
\xymatrix{
E \ar[r]_a \ar[d]_f & F \ar[d]^g \\
H \ar[r]^b & G
}
$$
un diagramme cartésien de faisceaux sur $(\Sch/S)_{fppf}$, de sorte que
$E = H \times_G F$. Si
\begin{enumerate}
\item $g$ est représentable par des espaces algébriques, surjectif, plat et
localement de présentation finie ;
\item $a$ est représentable par des espaces algébriques, séparé et
localement quasi-fini,
\end{enumerate}
alors $b$ est représentable (par des schémas), séparé et
localement quasi-fini.
\end{lemma}

\begin{proof}
Soit $T$ un schéma et soit $T \to G$ un morphisme.
Il faut montrer que $T \times_G H$ est un schéma et que
le morphisme $T \times_G H \to T$ est séparé et
localement quasi-fini. Nous pouvons donc effectuer sur tout le diagramme le changement de base vers $T$
et supposer que $G$ est un schéma. Alors $F$ est un espace algébrique.
Soit $U$ un schéma et soit $U \to F$ un morphisme étale surjectif.
Alors $U \to F$ est représentable, surjectif, plat et
localement de présentation finie d'après
Morphismes d'espaces,
lemmes \ref{spaces-morphisms-lemma-etale-flat} et
\ref{spaces-morphisms-lemma-etale-locally-finite-presentation}.
Le
lemme \ref{bootstrap-lemma-composition-transformation}
montre que $U \to G$ est lui aussi surjectif, plat et localement de présentation finie.
Notons que le changement de base $E \times_F U \to U$ de $a$ reste
séparé et localement quasi-fini (d'après le
lemme \ref{bootstrap-lemma-base-change-transformation-property}). Nous pouvons donc
remplacer la partie supérieure du diagramme du lemme par
$E \times_F U \to U$. Autrement dit, nous pouvons supposer que
$F \to G$ est un morphisme de schémas surjectif et plat,
localement de présentation finie.
En particulier, $\{F \to G\}$ est un recouvrement fppf de schémas.
D'après
Morphismes d'espaces, proposition
\ref{spaces-morphisms-proposition-locally-quasi-finite-separated-over-scheme}
on en conclut que $E$ est aussi un schéma.
D'après
Descente, lemme \ref{descent-lemma-descent-data-sheaves},
l'égalité $E = H \times_G F$ signifie que l'on obtient une donnée de descente
sur $E$ relativement au recouvrement fppf $\{F \to G\}$.
D'après
Compléments sur les morphismes, lemme
\ref{more-morphisms-lemma-separated-locally-quasi-finite-morphisms-fppf-descend}
Cette donnée de descente est effective.
En appliquant de nouveau
Descente, lemme \ref{descent-lemma-descent-data-sheaves},
on en déduit que $H$ est un schéma.
Enfin, d'après
Descente, lemmes \ref{descent-lemma-descending-property-separated} et
\ref{descent-lemma-descending-property-quasi-finite}
il s'ensuit que $b$ est séparé et localement quasi-fini.
\end{proof}

\noindent
Voici le résultat annoncé par le titre de la section.

\begin{lemma}
\label{bootstrap-lemma-bootstrap-diagonal}
Soit $S$ un schéma.
Soit $F : (\Sch/S)_{fppf}^{opp} \to \textit{Ensembles}$ un foncteur.
Supposons que
\begin{enumerate}
\item le préfaisceau $F$ soit un faisceau ;
\item il existe un espace algébrique $X$ et un morphisme $X \to F$
représentable par des espaces algébriques, surjectif, plat et
localement de présentation finie.
\end{enumerate}
Alors $\Delta_F$ est représentable (par des schémas).
\end{lemma}

\begin{proof}
Soit $U \to X$ un morphisme étale surjectif d'un schéma vers $X$.
Alors $U \to X$ est représentable, surjectif, plat et
localement de présentation finie d'après
Morphismes d'espaces,
lemmes \ref{spaces-morphisms-lemma-etale-flat} et
\ref{spaces-morphisms-lemma-etale-locally-finite-presentation}.
Le
lemme \ref{bootstrap-lemma-composition-transformation-property}
montre que le composé $U \to F$ est lui aussi représentable par des espaces algébriques,
surjectif, plat et localement de présentation finie.
Ainsi, $R = U \times_F U$ est un espace algébrique, voir le
lemme \ref{bootstrap-lemma-representable-by-spaces}.
Le morphisme d'espaces algébriques $R \to U \times_S U$ est
un monomorphisme, donc il est séparé (puisque la diagonale d'un monomorphisme
est un isomorphisme, voir
Morphismes d'espaces,
lemme \ref{spaces-morphisms-lemma-monomorphism}).
Comme $U \to F$ est localement de présentation finie, les deux
morphismes $R \to U$ le sont aussi, voir le
lemme \ref{bootstrap-lemma-base-change-transformation-property}.
Ainsi, $R \to U \times_S U$ est localement de type fini (utiliser
Morphismes d'espaces,
lemmes \ref{spaces-morphisms-lemma-finite-presentation-finite-type} et
\ref{spaces-morphisms-lemma-permanence-finite-type}).
Au total,
$R \to U \times_S U$ est un monomorphisme localement de type
fini, donc un morphisme séparé et localement quasi-fini, voir
Morphismes d'espaces, lemme
\ref{spaces-morphisms-lemma-monomorphism-loc-finite-type-loc-quasi-finite}.

\medskip\noindent
Nous sommes maintenant en mesure de prouver que $\Delta_F$ est représentable.
Soit $T$ un schéma et soit $(a, b) : T \to F \times F$ un morphisme.
Posons
$$
T' = (U \times_S U) \times_{F \times F} T.
$$
Notons que $U \times_S U \to F \times F$ est
représentable par des espaces algébriques, surjectif, plat et
localement de présentation finie d'après le
lemme \ref{bootstrap-lemma-product-transformations-property}.
Ainsi, $T'$ est un espace algébrique et le morphisme de projection
$T' \to T$ est surjectif, plat et localement de présentation finie.
Considérons $Z = T \times_{F \times F} F$ (c'est un faisceau) et
$$
Z' = T' \times_{U \times_S U} R
= T' \times_T Z.
$$
On voit que $Z'$ est un espace algébrique et que
$Z' \to T'$ est séparé et localement quasi-fini d'après la
discussion du premier alinéa de la démonstration, qui a établi que $R$ est
un espace algébrique et que le
morphisme $R \to U \times_S U$ possède ces propriétés.
Nous pouvons donc appliquer le
lemme \ref{bootstrap-lemma-after-fppf-sep-lqf}
au diagramme
$$
\xymatrix{
Z' \ar[r] \ar[d] & T' \ar[d] \\
Z \ar[r] & T
}
$$
ce qui permet de conclure.
\end{proof}

\noindent
Voici une variante du résultat précédent.

\begin{lemma}
\label{bootstrap-lemma-bootstrap-locally-quasi-finite}
Soit $S$ un schéma. Soit $F : (\Sch/S)_{fppf}^{opp} \to \textit{Ensembles}$ un
foncteur. Soit $X$ un schéma et soit $X \to F$ représentable par des espaces
algébriques et localement quasi-fini. Alors $X \to F$ est représentable
(par des schémas).
\end{lemma}

\begin{proof}
Soit $T$ un schéma et soit $T \to F$ un morphisme. Il faut montrer que
l'espace algébrique $X \times_F T$ est représentable par un schéma. Considérons
le morphisme
$$
X \times_F T  \longrightarrow X \times_{\Spec(\mathbf{Z})} T
$$
Comme $X \times_F T \to T$ est localement quasi-fini, la flèche affichée
l'est aussi (Morphismes d'espaces, lemme
\ref{spaces-morphisms-lemma-permanence-quasi-finite}).
D'autre part, la flèche affichée est un monomorphisme,
donc elle est séparée (Morphismes d'espaces, lemme
\ref{spaces-morphisms-lemma-monomorphism-separated}).
Ainsi, $X \times_F T$ est un schéma d'après Morphismes d'espaces, proposition
\ref{spaces-morphisms-proposition-locally-quasi-finite-separated-over-scheme}.
\end{proof}











\section{Amorçage}
\label{bootstrap-section-bootstrap}

\noindent
Prévenons d'emblée le lecteur que le résultat de cette section sera
supplanté par le théorème plus fort
\ref{bootstrap-theorem-final-bootstrap}.
En revanche, le théorème de cette section est nettement plus facile à
démontrer et donne déjà une bonne compréhension du mécanisme en jeu,
notamment aux lecteurs qui s'intéressent surtout aux champs de
Deligne--Mumford.

\medskip\noindent
Dans
Espaces, section \ref{spaces-section-algebraic-spaces},
nous avons défini un espace algébrique comme un faisceau pour la topologie fppf
dont la diagonale est représentable et vers lequel il existe un morphisme étale
surjectif provenant d'un schéma. Dans cette section, nous montrons qu'un
faisceau pour la topologie fppf dont la diagonale est représentable par des espaces
algébriques et qui admet un recouvrement étale surjectif par un espace algébrique
est lui aussi un espace algébrique.
Autrement dit, la catégorie des espaces algébriques s'obtient en élargissant la
catégorie des schémas par les faisceaux fppf $F$ qui possèdent une diagonale
représentable et un recouvrement étale par un schéma. Le
résultat de cette section affirme que recommencer ce procédé à partir de
la catégorie des espaces algébriques ne conduit pas à une nouvelle catégorie.

\medskip\noindent
Une autre motivation des résultats de cette section est qu'ils permettront
plus loin d'assurer qu'un champ de Deligne--Mumford dont le champ d'inertie est
trivial est équivalent à un espace algébrique, voir
Champs algébriques, lemme \ref{algebraic-lemma-algebraic-stack-no-automorphisms}.

\medskip\noindent
Voici le résultat principal de cette section (comme indiqué plus haut, il
sera supplanté par le théorème plus fort
\ref{bootstrap-theorem-final-bootstrap}).

\begin{theorem}
\label{bootstrap-theorem-bootstrap}
Soit $S$ un schéma.
Soit $F : (\Sch/S)_{fppf}^{opp} \to \textit{Ensembles}$ un foncteur.
Supposons que
\begin{enumerate}
\item le préfaisceau $F$ soit un faisceau ;
\item le morphisme diagonal $F  \to F \times F$ soit représentable par
des espaces algébriques ;
\item il existe un espace algébrique $X$
et un morphisme $X \to F$ surjectif et étale.
\end{enumerate}
ou bien supposons que
\begin{enumerate}
\item[(a)] le préfaisceau $F$ soit un faisceau ;
\item[(b)] il existe un espace algébrique $X$ et un morphisme $X \to F$
représentable par des espaces algébriques, surjectif et étale.
\end{enumerate}
Alors $F$ est un espace algébrique.
\end{theorem}

\begin{proof}
Nous utiliserons sans plus les mentionner les remarques qui suivent immédiatement la
définition \ref{bootstrap-definition-property-transformation}.

\medskip\noindent
Supposons (1), (2) et (3), et soit $X \to F$ comme en (3).
D'après le lemme \ref{bootstrap-lemma-representable-diagonal}, le morphisme
$X \to F$ est représentable par des espaces algébriques. Ainsi,
(a) et (b) sont vérifiées.

\medskip\noindent
Supposons (a) et (b), et soit $X \to F$ comme en (b).
Soit $U \to X$ un morphisme étale surjectif d'un schéma vers $X$.
D'après le lemme \ref{bootstrap-lemma-composition-transformation}, la transformation
$U \to F$ est représentable par des espaces algébriques, surjective et étale.
Pour prouver que $F$ est un espace algébrique, il suffit donc de montrer que
$\Delta_F$ est représentable (Espaces, définition
\ref{spaces-definition-algebraic-space}). Cela résulte immédiatement du
lemme \ref{bootstrap-lemma-bootstrap-diagonal}.
Nous pouvons aussi éviter ce lemme et montrer directement que $F$
est un espace algébrique, comme dans l'alinéa suivant.

\medskip\noindent
En effet, soient $U$ un schéma et $U \to F$ un morphisme représentable par des espaces
algébriques, surjectif et étale. Considérons le produit fibré $R = U \times_F U$.
Les deux projections $R \to U$ sont représentables par des espaces algébriques, surjectives
et étales (lemme \ref{bootstrap-lemma-base-change-transformation-property}).
En particulier, $R$ est un espace algébrique d'après le
lemme \ref{bootstrap-lemma-representable-by-spaces-over-space}.
Le morphisme d'espaces algébriques $R \to U \times_S U$ est un monomorphisme,
donc il est séparé (puisque la diagonale d'un monomorphisme est un isomorphisme).
Comme $R \to U$ est étale, $R \to U$ est localement quasi-fini, voir
Morphismes d'espaces,
lemme \ref{spaces-morphisms-lemma-etale-locally-quasi-finite}.
On en déduit que $R \to U \times_S U$ est lui aussi
localement quasi-fini d'après
Morphismes d'espaces,
lemme \ref{spaces-morphisms-lemma-permanence-quasi-finite}.
Ainsi,
Morphismes d'espaces, proposition
\ref{spaces-morphisms-proposition-locally-quasi-finite-separated-over-scheme}
s'applique et $R$ est un schéma. D'après le
lemme \ref{bootstrap-lemma-surjective-flat-locally-finite-presentation},
le morphisme $U \to F$ est un épimorphisme de faisceaux. Ainsi $F = U/R$.
On en conclut que $F$ est un espace algébrique d'après
Espaces, théorème \ref{spaces-theorem-presentation}.
\end{proof}









\section{Recherche d'ouverts}
\label{bootstrap-section-finding-opens}


\medskip\noindent
Commençons par établir un lemme qui améliore et généralise légèrement
Espaces, lemme \ref{spaces-lemma-finding-opens},
aux faisceaux quotients associés à des groupoïdes.

\begin{lemma}
\label{bootstrap-lemma-better-finding-opens}
Soit $S$ un schéma.
Soit $(U, R, s, t, c)$ un schéma en groupoïdes sur $S$.
Soit $g : U' \to U$ un morphisme.
Supposons que
\begin{enumerate}
\item le composé
$$
\xymatrix{
U' \times_{g, U, t} R \ar[r]_-{\text{pr}_1} \ar@/^3ex/[rr]^h
& R \ar[r]_s & U
}
$$
a pour image un ouvert $W \subset U$ ;
\item le morphisme induit $h : U' \times_{g, U, t} R \to W$
définit un épimorphisme de faisceaux pour la topologie fppf.
\end{enumerate}
Soit $R' = R|_{U'}$ la restriction de $R$ à $U'$. Alors le morphisme
de faisceaux quotients
$$
U'/R' \to U/R
$$
pour la topologie fppf est représentable et est une immersion ouverte.
\end{lemma}

\begin{proof}
Notons que $W$ est un sous-schéma ouvert $R$-invariant de $U$.
En effet, l'ensemble des points de $W$ est l'ensemble
des points de $U$ équivalents, au sens de
Groupoïdes,
lemme \ref{groupoids-lemma-pre-equivalence-equivalence-relation-points},
à un point de $g(U') \subset U$ (le lemme s'applique puisque $j : R \to U \times_S U$
est une pré-relation d'équivalence d'après
Groupoïdes, lemme \ref{groupoids-lemma-groupoid-pre-equivalence}).
De plus, $g : U' \to U$ se factorise par $W$.
Soit $R|_W$ la restriction de $R$ à $W$.
Il s'ensuit que $R'$ est aussi la restriction de $R|_W$ à $U'$.
On peut donc factoriser le morphisme de faisceaux du lemme sous la forme
$$
U'/R' \longrightarrow W/R|_W \longrightarrow U/R
$$
D'après Groupoïdes, lemme \ref{groupoids-lemma-quotient-groupoid-restrict},
la première flèche est un isomorphisme de faisceaux.
Il suffit donc de démontrer le lemme lorsque $g$ est l'immersion
d'un ouvert $R$-invariant dans $U$.

\medskip\noindent
Supposons que $U' \subset U$ soit un ouvert $R$-invariant et que $g$ soit le
morphisme d'inclusion. Posons $F = U/R$ et $F' = U'/R'$. D'après
Groupoïdes,
lemme \ref{groupoids-lemma-quotient-pre-equivalence-relation-restrict}
ou \ref{groupoids-lemma-quotient-groupoid-restrict},
le morphisme $F' \to F$ est injectif. Soit $\xi \in F(T)$.
Il faut montrer que $T \times_{\xi, F} F'$ est représentable
par un sous-schéma ouvert de $T$.
Il existe un recouvrement fppf $\{f_i : T_i \to T\}$ tel que
$\xi|_{T_i}$ soit l'image par $U \to U/R$ d'un morphisme $a_i : T_i \to U$.
Posons $V_i = a_i^{-1}(U')$.
Affirmons que $V_i \times_T T_j = T_i \times_T V_j$ comme sous-schémas ouverts
de $T_i \times_T T_j$.

\medskip\noindent
Comme $a_i \circ \text{pr}_0$ et $a_j \circ \text{pr}_1$ sont des morphismes
$T_i \times_T T_j \to U$ ayant tous deux pour image la section
$\xi|_{T_i \times_T T_j} \in F(T_i \times_T T_j)$, il existe
un recouvrement fppf $\{f_{ijk} : T_{ijk} \to T_i \times_T T_j\}$ et des morphismes
$r_{ijk} : T_{ijk} \to R$ tels que
$$
a_i \circ \text{pr}_0 \circ f_{ijk} = s \circ r_{ijk},
\quad
a_j \circ \text{pr}_1 \circ f_{ijk} = t \circ r_{ijk},
$$
voir
Groupoïdes, lemme \ref{groupoids-lemma-quotient-pre-equivalence}.
Puisque $U'$ est $R$-invariant, on a $s^{-1}(U') = t^{-1}(U')$, donc
$f_{ijk}^{-1}(V_i \times_T T_j) = f_{ijk}^{-1}(T_i \times_T V_j)$.
Comme $\{f_{ijk}\}$ est surjectif, l'affirmation précédente en résulte.
Ainsi, d'après
Descente, lemme \ref{descent-lemma-open-fpqc-covering},
il existe un sous-schéma ouvert $V \subset T$ tel que
$f_i^{-1}(V) = V_i$. Affirmons que $V$ représente $T \times_{\xi, F} F'$.

\medskip\noindent
Dans un premier temps, montrons que $\xi|_V$ appartient à $F'(V) \subset F(V)$.
En effet, la famille de morphismes $\{V_i \to V\}$ est un recouvrement fppf
et, par construction, $\xi|_{V_i} \in F'(V_i)$.
La propriété de faisceau de $F'$ donne donc $\xi|_V \in F'(V)$.
Enfin, soient $T' \to T$ un morphisme de schémas et
$\xi|_{T'} \in F'(T')$. Pour achever la démonstration, il faut montrer que
$T' \to T$ se factorise par $V$.
Il existe un recouvrement fppf $\{T'_j \to T'\}_{j \in J}$ et des morphismes
$b_j : T'_j \to U'$ tels que $\xi|_{T'_j}$ soit l'image par
$U' \to U/R$ de $b_j$. Il suffit manifestement de montrer que les composés
$T'_j \to T$ se factorisent par $V$. Nous pouvons donc supposer que $\xi|_{T'}$
est l'image d'un morphisme $b : T' \to U'$. Il suffit alors de montrer
que $T'\times_T T_i \to T_i$ se factorise par $V_i$. Sur le schéma
$T' \times_T T_i$, la restriction de $\xi$ est l'image de deux
éléments de $(U/R)(T' \times_T T_i)$, à savoir $a_i \circ \text{pr}_1$ et
$b \circ \text{pr}_0$, dont le second se factorise par l'ouvert $R$-invariant
$U'$. Ainsi, d'après
Groupoïdes, lemme \ref{groupoids-lemma-quotient-pre-equivalence},
il existe un recouvrement $\{h_k : Z_k \to T' \times_T T_i\}$ et des morphismes
$r_k : Z_k \to R$ tels que $a_i \circ \text{pr}_1 \circ h_k = s \circ r_k$
et $b \circ \text{pr}_0 \circ h_k = t \circ r_k$. Comme $U'$ est un ouvert
$R$-invariant et que l'image de $b$ est contenue dans $U'$, l'image de chaque
$a_i \circ \text{pr}_1 \circ h_k$ est contenue dans $U'$. Il s'ensuit que
l'image de $T' \times_T T_i \to T_i$ est contenue dans $V_i$, par définition de $V_i$,
ce qui achève la démonstration.
\end{proof}












\section{Découpage des relations d'équivalence}
\label{bootstrap-section-slicing}

\noindent
Dans cette section, nous expliquons comment « améliorer » une relation
d'équivalence donnée en la découpant. Il ne s'agit pas du type de « découpage
étale » auquel on pourrait être habitué, mais d'un découpage beaucoup plus grossier.


\begin{lemma}
\label{bootstrap-lemma-slice-equivalence-relation}
Soit $S$ un schéma.
Soit $j : R \to U \times_S U$ une relation d'équivalence sur des schémas au-dessus de $S$.
Supposons que $s, t : R \to U$ soient plats et localement de présentation finie.
Alors il existe une relation d'équivalence $j' : R' \to U'\times_S U'$
sur des schémas au-dessus de $S$, et un isomorphisme
$$
U'/R' \longrightarrow U/R
$$
induit par un morphisme $U' \to U$ qui envoie $R'$ dans $R$, tel que
$s', t' : R \to U$ soient plats, localement de présentation finie
et localement quasi-finis.
\end{lemma}

\begin{proof}
Nous démontrerons ce lemme en plusieurs étapes. Nous utiliserons sans plus
le mentionner le fait qu'une relation d'équivalence donne lieu à un schéma en
groupoïdes et que la restriction d'une relation d'équivalence est encore une relation
d'équivalence, voir
Groupoïdes, lemmes
\ref{groupoids-lemma-restrict-relation},
\ref{groupoids-lemma-equivalence-groupoid}, et
\ref{groupoids-lemma-restrict-groupoid-relation}.

\medskip\noindent
Étape 1 : nous pouvons supposer que $s, t : R \to U$ sont des morphismes
localement de présentation finie et de Cohen--Macaulay. En effet, comme dans
Compléments sur les groupoïdes, lemme \ref{more-groupoids-lemma-make-CM},
soit $g : U' \to U$ le sous-schéma ouvert tel que
$t^{-1}(U') \subset R$ soit l'ouvert maximal au-dessus duquel $s : R \to U$ est
de Cohen--Macaulay, et notons $R'$ la restriction de $R$ à $U'$.
D'après le lemme cité, le morphisme
$$
\xymatrix{
t^{-1}(U') \ar@{=}[r] &
U' \times_{g, U, t} R \ar[r]_-{\text{pr}_1} \ar@/^3ex/[rr]^h &
R \ar[r]_s &
U
}
$$
est surjectif. Comme $h$ est plat et localement de présentation finie,
$\{h\}$ est un recouvrement fppf. Ainsi, d'après
Groupoïdes, lemme \ref{groupoids-lemma-quotient-groupoid-restrict},
$U'/R' \to U/R$ est un isomorphisme. Par construction de $U'$,
$s', t'$ sont de Cohen--Macaulay et localement de présentation finie.

\medskip\noindent
Étape 2. Supposons $s, t$ de Cohen--Macaulay et localement de présentation finie.
Soit $u \in U$ un point de type fini. D'après
Compléments sur les groupoïdes, lemme \ref{more-groupoids-lemma-max-slice-quasi-finite},
il existe un schéma affine $U'$ et un morphisme $g : U' \to U$ tels que
\begin{enumerate}
\item $g$ est une immersion ;
\item $u \in U'$ ;
\item $g$ est localement de présentation finie ;
\item $h$ est plat, localement de présentation finie et localement quasi-fini ;
\item les morphismes $s', t' : R' \to U'$ sont plats, localement de
présentation finie et localement quasi-finis.
\end{enumerate}
Nous avons utilisé ici les notations introduites dans
Compléments sur les groupoïdes, situation \ref{more-groupoids-situation-slice}.

\medskip\noindent
Étape 3. Pour chaque point $u \in U$ de type fini,
choisissons $g_u : U'_u \to U$ comme à
l'étape 2 et notons $R'_u$ la restriction de $R$ à $U'_u$.
Notons $h_u = s \circ \text{pr}_1 : U'_u \times_{g_u, U, t} R \to U$. Posons
$U' = \coprod_{u \in U} U'_u$ et $g = \coprod g_u$. Soit $R'$ la
restriction de $R$ à $U'$ comme ci-dessus. Affirmons que
le couple $(U', g)$ convient\footnote{Il faut ici vérifier que $U'$ n'est pas
trop grand, c'est-à-dire qu'il est isomorphe à un objet de la catégorie
$\Sch_{fppf}$, voir la
section \ref{bootstrap-section-conventions}.
Il s'agit d'une question purement ensembliste ; utilisons la notion de taille d'un
schéma introduite dans
Ensembles, section \ref{sets-section-categories-schemes}.
Notons que la taille de chaque $U'_u$ est au plus celle de $U$
et que le cardinal de l'ensemble d'indices est au plus le cardinal de
$|U|$, lui-même borné par la taille de $U$. Ainsi, $U'$ est isomorphe
à un objet de $\Sch_{fppf}$ d'après
Ensembles, lemme \ref{sets-lemma-what-is-in-it}, partie (6).}.
Notons que
\begin{align*}
R' = &
\coprod\nolimits_{u_1, u_2 \in U}
(U'_{u_1} \times_{g_{u_1}, U, t} R)
\times_R
(R \times_{s, U, g_{u_2}} U'_{u_2}) \\
= &
\coprod\nolimits_{u_1, u_2 \in U}
(U'_{u_1} \times_{g_{u_1}, U, t} R) \times_{h_{u_1}, U, g_{u_2}} U'_{u_2}
\end{align*}
Ainsi, la projection $s' : R' \to U' = \coprod U'_{u_2}$
est plate, localement de
présentation finie et localement quasi-finie, car elle est un changement de base de $\coprod h_{u_1}$.
Enfin, par construction, le morphisme
$h : U' \times_{g, U, t} R \to U$ est égal à $\coprod h_u$ ; son
image contient donc tous les points de type fini de $U$.
Comme chaque $h_u$ est plat et localement de présentation finie, on en déduit que
$h$ est plat et localement de présentation finie.
En particulier, l'image de $h$ est ouverte (voir
Morphismes, lemme \ref{morphisms-lemma-fppf-open})
et, puisque l'ensemble des points de type fini est dense (voir
Morphismes, lemme \ref{morphisms-lemma-enough-finite-type-points}),
on en conclut que l'image de $h$ est $U$. Il s'ensuit que
$\{h\}$ est un recouvrement fppf. D'après
Groupoïdes, lemme \ref{groupoids-lemma-quotient-groupoid-restrict},
cela signifie que $U'/R' \to U/R$ est un isomorphisme.
La démonstration du lemme est achevée.
\end{proof}










\section{Quotient par un sous-groupoïde}
\label{bootstrap-section-dividing}

\noindent
Il nous faut encore un lemme avant de pouvoir effectuer l'amorçage final.
Examinons d'abord ce qui se passe pour les groupoïdes « ordinaires », avant
d'aborder la version schématique.

\medskip\noindent
Soit $\mathcal{C}$ un groupoïde, voir
Catégories, définition \ref{categories-definition-groupoid}.
Comme expliqué dans
Groupoïdes, section \ref{groupoids-section-groupoids},
il correspond à un quintuple $(\text{Ob}, \text{Arrows}, s, t, c)$.
Supposons donné un sous-ensemble $P \subset \text{Arrows}$ tel que
$(\text{Ob}, P, s|_P, t|_P, c|_P)$ soit lui aussi un groupoïde et que
$P$ ne contienne aucun automorphisme non trivial. On peut alors construire
le groupoïde quotient
$(\overline{\text{Ob}}, \overline{\text{Arrows}}, \overline{s},
\overline{t}, \overline{c})$
de la manière suivante :
\begin{enumerate}
\item $\overline{\text{Ob}} = \text{Ob}/P$
est l'ensemble des classes de $P$-isomorphie ;
\item $\overline{\text{Arrows}} = P\backslash \text{Arrows}/P$
est l'ensemble des flèches de $\mathcal{C}$ modulo précomposition et
postcomposition par des flèches de $P$ ;
\item les applications source et but
$\overline{s}, \overline{t} : P\backslash \text{Arrows}/P \to \text{Ob}/P$
sont induites par $s, t$ ;
\item la composition est définie par la règle
$\overline{c}(\overline{a}, \overline{b}) = \overline{c(a, b)}$,
qui est bien définie.
\end{enumerate}
En fait, le groupoïde initial
$(\text{Ob}, \text{Arrows}, s, t, c)$ est canoniquement
isomorphe à la restriction (voir la discussion dans
Groupoïdes, section \ref{groupoids-section-restrict-groupoid})
du groupoïde
$(\overline{\text{Ob}}, \overline{\text{Arrows}}, \overline{s},
\overline{t}, \overline{c})$ par l'application quotient
$g : \text{Ob} \to \overline{\text{Ob}}$. Rappelons que cela signifie
que
$$
\text{Arrows} =
\text{Ob}
\times_{g, \overline{\text{Ob}}, \overline{t}}
\overline{\text{Arrows}}
\times_{\overline{s}, \overline{\text{Ob}}, g}
\text{Ob}
$$
ce qui est vrai puisque $P$ ne possède aucun automorphisme non trivial.
Nous omettons les détails.

\medskip\noindent
Le lemme suivant vaut dans une généralité bien plus grande, mais c'est
la version que nous utiliserons dans la démonstration de l'amorçage final (après quoi
nous pourrons démontrer plus facilement les versions plus générales de ce lemme).

\begin{lemma}
\label{bootstrap-lemma-divide-subgroupoid}
Soit $S$ un schéma.
Soit $(U, R, s, t, c)$ un schéma en groupoïdes sur $S$.
Soit $P \to R$ un monomorphisme de schémas. Supposons que
\begin{enumerate}
\item $(U, P, s|_P, t|_P, c|_{P \times_{s, U, t}P})$ soit un schéma en groupoïdes ;
\item $s|_P, t|_P : P \to U$ soient finis localement libres ;
\item $j|_P : P \to U \times_S U$ soit un monomorphisme ;
\item $U$ soit affine ;
\item $j : R \to U \times_S U$ soit séparé et localement quasi-fini.
\end{enumerate}
Alors $U/P$ est représentable par un schéma affine $\overline{U}$, le
morphisme quotient $U \to \overline{U}$ est fini localement libre et
$P = U \times_{\overline{U}} U$. De plus, $R$ est la restriction d'un
schéma en groupoïdes
$(\overline{U}, \overline{R}, \overline{s}, \overline{t}, \overline{c})$
sur $\overline{U}$ par le morphisme quotient $U \to \overline{U}$.
\end{lemma}

\begin{proof}
Les conditions (1), (2), (3) et (4), ainsi que
Groupoïdes, proposition \ref{groupoids-proposition-finite-flat-equivalence},
impliquent l'existence du schéma affine $\overline{U}$ représentant $U/P$ ;
le morphisme $U \to \overline{U}$ est fini localement libre et
$P = U \times_{\overline{U}} U$. L'identification
$P = U \times_{\overline{U}} U$ est telle que $t|_P = \text{pr}_0$ et
$s|_P = \text{pr}_1$, et que la composition est égale à
$\text{pr}_{02} : U \times_{\overline{U}} U \times_{\overline{U}} U
\to U \times_{\overline{U}} U$.
Un produit de morphismes finis localement libres est fini localement libre (voir
Espaces, lemme \ref{spaces-lemma-product-representable-transformations-property}
et
Morphismes, lemmes \ref{morphisms-lemma-base-change-finite-locally-free} et
\ref{morphisms-lemma-composition-finite-locally-free}).
Pour obtenir $\overline{R}$, nous allons descendre
le schéma $R$ par le morphisme fini localement libre
$U \times_S U \to \overline{U} \times_S \overline{U}$.
Notons en effet que
$$
(U \times_S U)
\times_{(\overline{U} \times_S \overline{U})}
(U \times_S U)
=
P \times_S P
$$
d'après ce qui précède. Ainsi, se donner une donnée de descente (voir
Descente, définition \ref{descent-definition-descent-datum})
pour $R / U \times_S U / \overline{U} \times_S \overline{U}$
revient à se donner un isomorphisme
$$
\varphi :
R \times_{(U \times_S U), t \times t} (P \times_S P)
\longrightarrow
(P \times_S P) \times_{s \times s, (U \times_S U)} R
$$
au-dessus de $P \times_S P$ satisfaisant à une condition de cocycle. Nous définissons $\varphi$
sur les points à valeurs dans $T$ par la règle
$$
\varphi : (r, (p, p')) \longmapsto ((p, p'), p^{-1} \circ r \circ p')
$$
où la composition est prise dans la catégorie groupoïde
$(U(T), R(T), s, t, c)$.
Cela a un sens, car pour que $(r, (p, p'))$ soit un point à valeurs dans $T$
de la source de $\varphi$, il faut que $t(r) = t(p)$
et $s(r) = t(p')$. Notons que ce morphisme est un isomorphisme
dont l'inverse est donné par
$((p, p'), r') \mapsto (p \circ r' \circ (p')^{-1}, (p, p'))$.
Pour vérifier la condition de cocycle, il faut établir que
$\varphi_{02} = \varphi_{12} \circ \varphi_{01}$
comme morphismes au-dessus de
$$
(U \times_S U)
\times_{(\overline{U} \times_S \overline{U})} (U \times_S U)
\times_{(\overline{U} \times_S \overline{U})} (U \times_S U) =
(P \times_S P) \times_{s \times s, (U \times_S U), t \times t} (P \times_S P)
$$
Un calcul explicite donne
$$
\begin{matrix}
\varphi_{02} & (r, (p_1, p_1'), (p_2, p_2')) & \mapsto &
((p_1, p_1'), (p_2, p_2'),
(p_1 \circ p_2)^{-1} \circ r \circ (p_1' \circ p_2')) \\
\varphi_{01} & (r, (p_1, p_1'), (p_2, p_2')) & \mapsto &
((p_1, p_1'), p_1^{-1} \circ r \circ p_1', (p_2, p_2')) \\
\varphi_{12} & ((p_1, p_1'), r, (p_2, p_2')) & \mapsto &
((p_1, p_1'), (p_2, p_2'), p_2^{-1} \circ r \circ p_2')
\end{matrix}
$$
(avec des notations évidentes), ce qui donne le résultat voulu.
Comme $j$ est séparé et localement quasi-fini d'après (5), nous pouvons appliquer
Compléments sur les morphismes, lemme
\ref{more-morphisms-lemma-separated-locally-quasi-finite-morphisms-fppf-descend}
pour obtenir un schéma $\overline{R} \to \overline{U} \times_S \overline{U}$
et un isomorphisme
$$
R \to \overline{R} \times_{(\overline{U} \times_S \overline{U})} (U \times_S U)
$$
qui identifie la donnée de descente $\varphi$ à la donnée de descente
canonique sur
$\overline{R} \times_{(\overline{U} \times_S \overline{U})} (U \times_S U)$,
voir
Descente, définition \ref{descent-definition-effective}.

\medskip\noindent
Comme $U \times_S U \to \overline{U} \times_S \overline{U}$ est fini
localement libre, on en déduit que $R \to \overline{R}$ est fini localement libre
comme changement de base. Ainsi, $R \to \overline{R}$ est surjectif comme morphisme de
faisceaux sur $(\Sch/S)_{fppf}$.
Notre choix de $\varphi$ implique que, pour des points à valeurs dans $T$, $r, r' \in R(T)$,
ceux-ci ont la même image dans $\overline{R}$ si et seulement si
$p^{-1} \circ r \circ p'$ pour certains $p, p' \in P(T)$. Ainsi,
$\overline{R}$ représente le faisceau
$$
T \longmapsto  \overline{R(T)} = P(T)\backslash R(T)/P(T)
$$
avec les notations de la discussion précédant le lemme.
Nous pouvons donc définir la structure de groupoïde sur
$(\overline{U} = U/P, \overline{R} = P\backslash R/P)$ exactement comme dans
la discussion du cas d'un groupoïde « ordinaire ».
Il s'ensuit que $(U, R, s, t, c)$ est l'image inverse de
cette structure de groupoïde par le morphisme $U \to \overline{U}$.
Cela achève la démonstration.
\end{proof}















\section{Amorçage final}
\label{bootstrap-section-final-bootstrap}

\noindent
Le résultat suivant va nettement plus loin que les précédents.

\begin{theorem}
\label{bootstrap-theorem-final-bootstrap}
Soit $S$ un schéma.
Soit $F : (\Sch/S)_{fppf}^{opp} \to \textit{Ensembles}$ un foncteur.
Chacune des conditions suivantes implique que $F$ est un espace algébrique :
\begin{enumerate}
\item $F = U/R$, où $(U, R, s, t, c)$ est un groupoïde en espaces algébriques
sur $S$ tel que $s, t$ soient plats et localement de présentation finie, et
$j = (t, s) : R \to U \times_S U$ soit une relation d'équivalence ;
\item $F = U/R$, où $(U, R, s, t, c)$ est un schéma en groupoïdes
sur $S$ tel que $s, t$ soient plats et localement de présentation finie, et
$j = (t, s) : R \to U \times_S U$ soit une relation d'équivalence ;
\item $F$ est un faisceau et il existe un espace algébrique $U$ et un morphisme
$U \to F$ représentable par des espaces algébriques,
surjectif, plat et localement de présentation finie ;
\item $F$ est un faisceau et il existe un schéma $U$ et un morphisme
$U \to F$ représentable par des espaces algébriques ou par des schémas,
surjectif, plat et localement de présentation finie ;
\item $F$ est un faisceau, $\Delta_F$ est représentable par des espaces algébriques,
et il existe un espace algébrique $U$ et un morphisme $U \to F$
surjectif, plat et localement de présentation finie ;
\item $F$ est un faisceau, $\Delta_F$ est représentable,
et il existe un schéma $U$ et un morphisme $U \to F$
surjectif, plat et localement de présentation finie.
\end{enumerate}
\end{theorem}

\begin{proof}
Observations immédiates : (6) est un cas particulier de (5) et
(4) est un cas particulier de (3).
Montrons d'abord que les cas (5) et (3) se ramènent au cas (1).
En effet, le lemme d'amorçage de la diagonale
\ref{bootstrap-lemma-bootstrap-diagonal}
montre que (3) implique (5). Dans le cas (5), posons $R = U \times_F U$ ; c'est
un espace algébrique par hypothèse. De plus, toujours par hypothèse, les deux
projections $s, t : R \to U$ sont surjectives, plates et localement de
présentation finie. Le morphisme $j : R \to U \times_S U$ est manifestement une
relation d'équivalence. D'après le
lemme \ref{bootstrap-lemma-surjective-flat-locally-finite-presentation},
le morphisme $U \to F$ est un épimorphisme de faisceaux. Ainsi $F = U/R$,
ce qui nous ramène au cas (1).

\medskip\noindent
Montrons ensuite que (1) se ramène à (2).
Soit $(U, R, s, t, c)$ un groupoïde en espaces algébriques
sur $S$ tel que $s, t$ soient plats et localement de présentation finie, et que
$j = (t, s) : R \to U \times_S U$ soit une relation d'équivalence.
Choisissons un schéma $U'$ et un morphisme étale surjectif $U' \to U$.
Soit $R' = R|_{U'}$ la restriction de $R$ à $U'$. D'après
Groupoïdes dans les espaces,
lemme \ref{spaces-groupoids-lemma-quotient-pre-equivalence-relation-restrict},
on a $U/R = U'/R'$. Comme $s', t' : R' \to U'$ sont eux aussi
plats et localement de présentation finie (voir
Compléments sur les groupoïdes dans les espaces,
lemme \ref{spaces-more-groupoids-lemma-restrict-preserves-type}),
nous sommes ramenés au cas où $U$ est un schéma.
Puisque $j$ est une relation d'équivalence, $j$ est un monomorphisme.
Comme $s : R \to U$ est localement de présentation finie,
$j : R \to U \times_S U$ est localement de type fini, voir
Morphismes d'espaces, lemme \ref{spaces-morphisms-lemma-permanence-finite-type}.
D'après
Morphismes d'espaces, lemme
\ref{spaces-morphisms-lemma-monomorphism-loc-finite-type-loc-quasi-finite}
$j$ est localement quasi-fini et séparé.
Ainsi, si $U$ est un schéma, alors $R$ est un schéma d'après
Morphismes d'espaces, proposition
\ref{spaces-morphisms-proposition-locally-quasi-finite-separated-over-scheme}.
Nous sommes donc ramenés à démontrer le théorème dans le cas (2).

\medskip\noindent
Supposons $F = U/R$, où $(U, R, s, t, c)$ est un schéma en groupoïdes
sur $S$ tel que $s, t$ soient plats et localement de présentation finie, et que
$j = (t, s) : R \to U \times_S U$ soit une relation d'équivalence. D'après le
lemme \ref{bootstrap-lemma-slice-equivalence-relation},
nous sommes ramenés au cas où $s, t$ sont plats,
localement de présentation finie et localement quasi-finis.
Soit $U = \bigcup_{i \in I} U_i$ un recouvrement ouvert affine
(où l'ensemble d'indices $I$ est de cardinal $\leq$ à la taille de $U$, afin d'éviter
plus loin des problèmes ensemblistes --- la plupart des lecteurs peuvent ignorer cette remarque).
Soit $(U_i, R_i, s_i, t_i, c_i)$ la restriction de $R$
à $U_i$. Il est clair que $s_i, t_i$ restent plats, localement de présentation
finie et localement quasi-finis, puisque $R_i$ est le sous-schéma ouvert
$s^{-1}(U_i) \cap t^{-1}(U_i)$ de $R$
et que $s_i, t_i$ sont les restrictions de $s, t$ à cet ouvert. D'après le
lemme \ref{bootstrap-lemma-better-finding-opens}
(ou, plus simplement,
Espaces, lemme \ref{spaces-lemma-finding-opens}),
le morphisme $U_i/R_i \to U/R$ est représentable par des immersions ouvertes.
Ainsi, si l'on peut montrer que $F_i = U_i/R_i$ est un espace algébrique, alors
$\coprod_{i \in I} F_i$ est un espace algébrique d'après
Espaces, lemme \ref{spaces-lemma-coproduct-algebraic-spaces}.
Comme $U = \bigcup U_i$ est un recouvrement ouvert, il est clair que
$\coprod F_i \to F$ est surjectif. Il s'ensuit que
$U/R$ est un espace algébrique, d'après
Espaces, lemme \ref{spaces-lemma-glueing-algebraic-spaces}.
Nous sommes ainsi ramenés au cas où $U$ est affine, $s, t$ sont plats,
localement de présentation finie et localement quasi-finis, et où
$j$ est une relation d'équivalence.

\medskip\noindent
Supposons que $(U, R, s, t, c)$ soit un schéma en groupoïdes sur $S$,
avec $U$ affine, tel que $s, t$ soient plats, localement de présentation finie
et localement quasi-finis, et que $j$ soit une relation d'équivalence.
Choisissons $u \in U$. Appliquons
Compléments sur les groupoïdes dans les espaces,
lemme \ref{spaces-more-groupoids-lemma-quasi-splitting-affine-scheme},
à $u \in U, R, s, t, c$. On obtient un schéma affine $U'$, un morphisme étale
$g : U' \to U$ et un point $u' \in U'$ tel que $\kappa(u) = \kappa(u')$
et que la restriction $R' = R|_{U'}$ soit quasi-scindée au-dessus de $u'$.
Notons que l'image $g(U')$ est ouverte, puisque $g$ est étale, et qu'elle contient $u$.
En appliquant le lemme à plusieurs reprises, on peut donc trouver un nombre fini de
points $u_i \in U$, $i = 1, \ldots, n$,
de schémas affines $U'_i$, de morphismes étales $g_i : U_i' \to U$ et de points
$u'_i \in U'_i$ tels que $g(u'_i) = u_i$, avec (a) chaque
restriction $R'_i$ quasi-scindée au-dessus d'un point de $U'_i$, et
(b) $U = \bigcup_{i = 1, \ldots, n} g_i(U'_i)$.
Reprenons maintenant la dernière partie de l'argument de l'alinéa précédent :
d'après le
lemme \ref{bootstrap-lemma-better-finding-opens}
(ou, plus simplement,
Espaces, lemme \ref{spaces-lemma-finding-opens}),
le morphisme $U'_i/R'_i \to U/R$ est représentable par des immersions ouvertes.
Si l'on peut montrer que $F_i = U'_i/R'_i$ est un espace algébrique, alors
$\coprod_{i \in I} F_i$ est un espace algébrique d'après
Espaces, lemme \ref{spaces-lemma-coproduct-algebraic-spaces}.
Comme $\{g_i : U'_i \to U\}$ est un recouvrement étale,
il est clair que $\coprod F_i \to F$ est surjectif. Il s'ensuit que
$U/R$ est un espace algébrique, d'après
Espaces, lemme \ref{spaces-lemma-glueing-algebraic-spaces}.
Nous sommes ainsi ramenés au cas où $U$ est affine, $s, t$ sont plats,
localement de présentation finie et localement quasi-finis,
$j$ est une relation d'équivalence et $R$ est quasi-scindé au-dessus de $u$ pour un certain
$u \in U$.

\medskip\noindent
Supposons que $(U, R, s, t, c)$ soit un schéma en groupoïdes sur $S$,
avec $U$ affine et $u \in U$, tel que $s, t$ soient plats, localement
de présentation finie et localement quasi-finis, que
$j = (t, s) : R \to U \times_S U$ soit une relation d'équivalence
et que $R$ soit quasi-scindé au-dessus de $u$. Soit $P \subset R$ un quasi-scindage
de $R$ au-dessus de $u$. D'après le
lemme \ref{bootstrap-lemma-divide-subgroupoid},
$(U, R, s, t, c)$ est la restriction d'un groupoïde
$(\overline{U}, \overline{R}, \overline{s}, \overline{t}, \overline{c})$
par un morphisme fini localement libre surjectif $U \to \overline{U}$ tel que
$P = U \times_{\overline{U}} U$. Notons que $s$ admet la factorisation
$$
R = U \times_{\overline{U}, \overline{t}} \overline{R}
\times_{\overline{s}, \overline{U}} U
\xrightarrow{\text{pr}_{23}}
\overline{R} \times_{\overline{s}, \overline{U}} U
\xrightarrow{\text{pr}_2} U
$$
Le morphisme $\text{pr}_2$ est le changement de base de $\overline{s}$, et
$\text{pr}_{23}$ est un changement de base du morphisme fini localement libre
surjectif $U \to \overline{U}$. Comme $s$ est plat, localement
de présentation finie et localement quasi-fini, et comme $\text{pr}_{23}$
est fini localement libre surjectif (comme changement de base d'un tel morphisme),
on en déduit que $\text{pr}_2$ est plat, localement
de présentation finie et localement quasi-fini d'après
Descente, lemmes
\ref{descent-lemma-flat-fpqc-local-source} et
\ref{descent-lemma-locally-finite-presentation-fppf-local-source}, ainsi que
Morphismes, lemme \ref{morphisms-lemma-quasi-finite-local-source}.
Comme $\text{pr}_2$ est le changement de base du morphisme
$\overline{s}$ par $U \to \overline{U}$ et que $\{U \to \overline{U}\}$
est un recouvrement fppf, on en conclut que $\overline{s}$ est
plat, localement de présentation finie et localement quasi-fini, voir
Descente, lemmes \ref{descent-lemma-descending-property-flat},
\ref{descent-lemma-descending-property-locally-finite-presentation} et
\ref{descent-lemma-descending-property-quasi-finite}. Il en va de même
pour $\overline{t}$. Considérons le diagramme commutatif
$$
\xymatrix{
U \times_{\overline{U}} U \ar@{=}[r] \ar[rd] & P \ar[r] \ar[d] & R \ar[d] \\
& \overline{U} \ar[r]^{\overline{e}} & \overline{R}
}
$$
Un fait général sur les restrictions affirme que les quatre sommets extérieurs
forment un diagramme cartésien. L'égalité montre que le carré intérieur est
cartésien. Comme $P$ est ouvert dans $R$ (par définition d'un quasi-scindage),
on en conclut que $\overline{e}$ est une immersion ouverte d'après
Descente, lemme \ref{descent-lemma-descending-property-open-immersion}.
En appliquant
Groupoïdes,
lemme \ref{groupoids-lemma-quotient-pre-equivalence-relation-restrict},
on obtient $U/R = \overline{U}/\overline{R}$. Nous sommes donc ramenés
au cas où $(U, R, s, t, c)$ est un schéma en groupoïdes sur $S$,
avec $U$ affine et $u \in U$, tel que $s, t$ soient plats, localement
de présentation finie et localement quasi-finis, que
$j = (t, s) : R \to U \times_S U$ soit une relation d'équivalence
et que $e : U \to R$ soit une immersion ouverte !

\medskip\noindent
Mais bien entendu, si $e$ est une immersion ouverte et si
$s, t$ sont plats et localement de présentation finie,
alors les morphismes $t, s$ sont étales.
On peut par exemple le voir en appliquant
Compléments sur les groupoïdes, lemme \ref{more-groupoids-lemma-sheaf-differentials},
qui montre que $\Omega_{R/U} = 0$, ce qui implique à son tour
que $s, t : R \to U$ sont G-non ramifiés (voir
Morphismes, lemme \ref{morphisms-lemma-unramified-omega-zero}),
et implique donc que $s, t$ sont étales (voir
Morphismes, lemme \ref{morphisms-lemma-flat-unramified-etale}).
Enfin, si $s, t$ sont étales, alors $U/R$ est un espace
algébrique d'après
Espaces, théorème \ref{spaces-theorem-presentation}.
\end{proof}





\section{Applications}
\label{bootstrap-section-applications}

\noindent
Comme première application, nous obtenons le fait fondamental suivant :
$$
\fbox{Un faisceau qui est localement un espace algébrique pour la topologie fppf est un espace algébrique.}
$$
C'est le contenu du lemme suivant.
Notons que l'hypothèse (2) équivaut à la condition selon laquelle
$F|_{(\Sch/S_i)_{fppf}}$ est un espace algébrique, voir
Espaces, lemme \ref{spaces-lemma-rephrase}.
L'hypothèse (3) est une condition ensembliste que peuvent ignorer
ceux qui ne se soucient pas des questions ensemblistes.

\begin{lemma}
\label{bootstrap-lemma-locally-algebraic-space}
\begin{slogan}
La définition d'un espace algébrique est locale pour la topologie fppf.
\end{slogan}
Soit $S$ un schéma.
Soit $F : (\Sch/S)_{fppf}^{opp} \to \textit{Ensembles}$ un foncteur.
Soit $\{S_i \to S\}_{i \in I}$ un recouvrement de $(\Sch/S)_{fppf}$.
Supposons que
\begin{enumerate}
\item $F$ soit un faisceau ;
\item chaque $F_i = h_{S_i} \times F$ soit un espace algébrique ;
\item $\coprod_{i \in I} F_i$ soit un espace algébrique (voir
Espaces, lemme \ref{spaces-lemma-coproduct-algebraic-spaces}).
\end{enumerate}
Alors $F$ est un espace algébrique.
\end{lemma}

\begin{proof}
Considérons le morphisme $\coprod F_i \to F$. C'est le changement de base
de $\coprod S_i \to S$ par $F \to S$. Il est donc représentable,
localement de présentation finie, plat et surjectif, d'après notre définition
d'un recouvrement fppf et le
lemme \ref{bootstrap-lemma-base-change-transformation-property}.
Ainsi, le
théorème \ref{bootstrap-theorem-final-bootstrap}
s'applique et montre que $F$ est un espace algébrique.
\end{proof}

\noindent
Voici un cas particulier du lemme \ref{bootstrap-lemma-locally-algebraic-space}
dans lequel les questions ensemblistes ne se posent pas.

\begin{lemma}
\label{bootstrap-lemma-locally-algebraic-space-finite-type}
Soit $S$ un schéma.
Soit $F : (\Sch/S)_{fppf}^{opp} \to \textit{Ensembles}$ un foncteur.
Soit $\{S_i \to S\}_{i \in I}$ un recouvrement de $(\Sch/S)_{fppf}$.
Supposons que
\begin{enumerate}
\item $F$ soit un faisceau ;
\item chaque $F_i = h_{S_i} \times F$ soit un espace algébrique ;
\item les morphismes $F_i \to S_i$ soient de type fini.
\end{enumerate}
Alors $F$ est un espace algébrique.
\end{lemma}

\begin{proof}
Nous allons utiliser le
lemme \ref{bootstrap-lemma-locally-algebraic-space}
ci-dessus. Pour cela, nous montrerons que l'hypothèse selon laquelle
$F_i$ est de type fini sur $S_i$ entraîne que la condition ensembliste
du lemme est satisfaite (quitte à raffiner légèrement le
recouvrement donné de $S$).
Nous suggérons au lecteur de passer la suite de la démonstration.

\medskip\noindent
Si $S'_i \to S_i$ est un morphisme de schémas, alors
$$
h_{S'_i} \times F =
h_{S'_i} \times_{h_{S_i}} h_{S_i} \times F =
h_{S'_i} \times_{h_{S_i}} F_i
$$
est un espace algébrique de type fini sur $S'_i$, voir
Espaces, lemme \ref{spaces-lemma-fibre-product-spaces},
et
Morphismes d'espaces,
lemme \ref{spaces-morphisms-lemma-base-change-finite-type}.
Nous pouvons donc raffiner le recouvrement donné. Après ce raffinement, nous pouvons supposer :
(a) chaque $S_i$ est affine ; (b) le cardinal de $I$ est au plus
le cardinal de l'ensemble des points de $S$. (En effet, pour recouvrir
tout $S$, il suffit que chaque point appartienne à l'image de $S_i \to S$
pour un certain $i$.)

\medskip\noindent
Comme chaque $S_i$ est affine et chaque $F_i$ de type fini sur $S_i$,
on en déduit que $F_i$ est quasi-compact. Ainsi, d'après
Propriétés des espaces,
lemme \ref{spaces-properties-lemma-quasi-compact-affine-cover},
il existe un $U_i \in \Ob((\Sch/S)_{fppf})$ affine
et un morphisme étale surjectif $U_i \to F_i$. Le fait que
$F_i \to S_i$ soit localement de type fini implique alors que
$U_i \to S_i$ est localement de type fini, et en particulier que
$U_i \to S$ est localement de type fini. D'après
Ensembles, lemme \ref{sets-lemma-bound-finite-type},
on a $\text{size}(U_i) \leq \text{size}(S)$.
Comme on a aussi $|I| \leq \text{size}(S)$,
$\coprod_{i \in I} U_i$ est isomorphe à un objet de
$(\Sch/S)_{fppf}$ d'après
Ensembles, lemme \ref{sets-lemma-bound-size},
et la construction de $\Sch$. Il s'ensuit que
$\coprod F_i$ est un espace algébrique d'après
Espaces, lemme \ref{spaces-lemma-coproduct-algebraic-spaces},
ce qui conclut.
\end{proof}

\noindent
Comme deuxième application, nous obtenons
$$
\fbox{Toute donnée de descente fppf pour des espaces algébriques est effective.}
$$
Cet énoncé vaut sous réserve de difficultés ensemblistes ; nous donnons
le lemme suivant comme résultat type.

\begin{lemma}
\label{bootstrap-lemma-descend-algebraic-space}
\begin{slogan}
Les données de descente fppf pour les espaces algébriques sont effectives.
\end{slogan}
Soit $S$ un schéma. Soit $\{X_i \to X\}_{i \in I}$ un recouvrement
fppf d'espaces algébriques sur $S$.
\begin{enumerate}
\item Si $I$ est dénombrable\footnote{La restriction de dénombrabilité peut être
ignorée par ceux que les questions ensemblistes ne préoccupent pas. On peut admettre
ici des ensembles d'indices plus grands si l'on peut borner la taille des espaces
algébriques que l'on descend. Voir par exemple le
lemme \ref{bootstrap-lemma-locally-algebraic-space-finite-type}.}, alors toute
donnée de descente d'espaces algébriques relativement à $\{X_i \to X\}$ est effective.
\item Toute donnée de descente $(Y_i, \varphi_{ij})$ relativement à
$\{X_i \to X\}_{i \in I}$ (Descente sur les espaces, définition
\ref{spaces-descent-definition-descent-datum-for-family-of-morphisms})
pour laquelle $Y_i \to X_i$ est de type fini
est effective.
\end{enumerate}
\end{lemma}

\begin{proof}
Démonstration de (1). D'après
Descente sur les espaces, lemme \ref{spaces-descent-lemma-descent-data-sheaves},
cela revient à affirmer qu'un faisceau fppf $F$
muni d'un morphisme $F \to X$ est un espace algébrique dès que
chaque $F \times_X X_i$ est un espace algébrique.
La restriction sur le cardinal de $I$ implique que
les coproduits d'espaces algébriques indexés par $I$ sont des espaces algébriques, voir
Espaces, lemme \ref{spaces-lemma-coproduct-algebraic-spaces},
et
Ensembles, lemme \ref{sets-lemma-what-is-in-it}.
Le morphisme
$$
\coprod F \times_X X_i \longrightarrow F
$$
est représentable par des espaces algébriques (comme changement de base de
$\coprod X_i \to X$, voir le lemme \ref{bootstrap-lemma-base-change-transformation}),
et surjectif, plat et localement de présentation finie
(comme changement de base de $\coprod X_i \to X$, voir le
lemme \ref{bootstrap-lemma-base-change-transformation-property}).
La partie (1) résulte donc du théorème \ref{bootstrap-theorem-final-bootstrap}.

\medskip\noindent
Démonstration de (2). Appliquons d'abord
Descente sur les espaces, lemme \ref{spaces-descent-lemma-descent-data-sheaves},
pour obtenir un faisceau fppf $F$ muni d'un morphisme $F \to X$
tel que $F \times_X X_i = Y_i$ pour tout $i \in I$.
Il s'agit de montrer que $F$ est un espace algébrique.
Choisissons un schéma $U$ et un morphisme étale surjectif $U \to X$.
Alors $F' = U \times_X F \to F$ est représentable, surjectif et étale
comme changement de base de $U \to X$.
D'après le théorème \ref{bootstrap-theorem-final-bootstrap}, il suffit de montrer
que $F' = U \times_X F$ est un espace algébrique.
On peut choisir un recouvrement fppf $\{U_j \to U\}_{j \in J}$,
où $U_j$ est un schéma raffinant le recouvrement fppf
$\{X_i \times_X U \to U\}_{i \in I}$, voir
Topologies sur les espaces, lemme
\ref{spaces-topologies-lemma-refine-fppf-schemes}.
On obtient ainsi une application $a : J \to I$ et, pour chaque $j$,
un morphisme $U_j \to X_{a(j)}$ au-dessus de $X$.
Alors $U_j \times_U F' = U_j \times_{X_{a(j)}} Y_{a(j)}$
est de type fini sur $U_j$. Par conséquent, $F'$ est un espace
algébrique d'après le lemme \ref{bootstrap-lemma-locally-algebraic-space-finite-type}.
\end{proof}

\noindent
Voici une application d'un autre type.

\begin{lemma}
\label{bootstrap-lemma-representable-by-spaces-cover}
Soit $S$ un schéma. Soient $a : F \to G$ et $b : G \to H$ des
transformations de foncteurs $(\Sch/S)_{fppf}^{opp} \to \textit{Ensembles}$.
Supposons que
\begin{enumerate}
\item $F, G, H$ soient des faisceaux ;
\item $a : F \to G$ soit représentable par des espaces algébriques, plat,
localement de présentation finie et surjectif ;
\item $b \circ a : F \to H$ soit représentable par des espaces algébriques.
\end{enumerate}
Alors $b$ est représentable par des espaces algébriques.
\end{lemma}

\begin{proof}
Soit $U$ un schéma sur $S$ et soit $\xi \in H(U)$. Il faut montrer que
$U \times_{\xi, H} G$ est un espace algébrique. Or, nous savons
que $U \times_{\xi, H} F$ est un espace algébrique et que
$U \times_{\xi, H} F \to U \times_{\xi, H} G$ est représentable par
des espaces algébriques, plat, localement de présentation finie et surjectif,
comme changement de base du morphisme $a$ (voir le
lemme \ref{bootstrap-lemma-base-change-transformation-property}).
Le résultat découle donc du théorème \ref{bootstrap-theorem-final-bootstrap}.
\end{proof}

\begin{lemma}
\label{bootstrap-lemma-quotient-stack-isom}
Supposons que $B \to S$ et $(U, R, s, t, c)$ soient comme dans
Groupoïdes dans les espaces,
définition \ref{spaces-groupoids-definition-quotient-stack} (1).
Pour tout schéma $T$ sur $S$ et tous objets $x, y$ de $[U/R]$ au-dessus de $T$,
le faisceau $\mathit{Isom}(x, y)$ sur $(\Sch/T)_{fppf}$
est un espace algébrique.
\end{lemma}

\begin{proof}
D'après
Groupoïdes dans les espaces,
lemme \ref{spaces-groupoids-lemma-quotient-stack-isom},
il existe un recouvrement fppf $\{T_i \to T\}_{i \in I}$
tel que $\mathit{Isom}(x, y)|_{(\Sch/T_i)_{fppf}}$
soit un espace algébrique pour chaque $i$. D'après
Espaces, lemme \ref{spaces-lemma-rephrase},
cela signifie que chaque $F_i = h_{S_i} \times \mathit{Isom}(x, y)$
est un espace algébrique.
Pour démontrer le lemme, il ne reste donc qu'à vérifier la condition ensembliste
du lemme \ref{bootstrap-lemma-locally-algebraic-space} ci-dessus, à savoir que
$\coprod F_i$ est un espace algébrique. Pour cela, nous utilisons
Espaces, lemme \ref{spaces-lemma-coproduct-algebraic-spaces},
qui demande de montrer que $I$ et les $F_i$ ne sont pas « trop grands ».
Nous suggérons au lecteur de passer la suite de la démonstration.

\medskip\noindent
Choisissons $U' \in \Ob(\Sch/S)_{fppf}$ et un morphisme
étale surjectif $U' \to U$. Soit $R'$ la restriction de $R$ à $U'$.
Comme $[U/R] = [U'/R']$, nous pouvons, après avoir remplacé $U$ par $U'$,
supposer que $U$ est un schéma. (Cette étape sert à faire en sorte que les
produits fibrés ci-dessous soient pris au-dessus d'un schéma.)

\medskip\noindent
Notons que, si l'on raffine le recouvrement $\{T_i \to T\}$, chaque
$F_i$ reste un espace algébrique.
Nous pouvons donc supposer que chaque $T_i$ est affine. Comme
$T_i \to T$ est localement de présentation finie, il s'ensuit que
$\text{size}(T_i) \leq \text{size}(T)$, voir
Ensembles, lemme \ref{sets-lemma-bound-finite-type}.
Nous pouvons aussi supposer que le cardinal de l'ensemble d'indices $I$ est au plus le
cardinal de l'ensemble des points de $T$, car, pour obtenir un
recouvrement, il suffit que chaque point de $T$ appartienne à l'image.
Ainsi $|I| \leq \text{size}(T)$.
Choisissons $W \in \Ob((\Sch/S)_{fppf})$
et un morphisme étale surjectif $W \to R$. Notons que, dans la démonstration de
Groupoïdes dans les espaces,
lemme \ref{spaces-groupoids-lemma-quotient-stack-isom},
nous avons montré que $F_i$ est représenté par
$T_i \times_{(y_i, x_i), U \times_B U} R$ pour certains
$x_i, y_i : T_i \to U$. On voit alors que
$V_i = T_i \times_{(y_i, x_i), U \times_B U} W$ est un
schéma muni d'un morphisme étale surjectif $V_i \to F_i$.
D'après
Ensembles, lemme \ref{sets-lemma-bound-size-fibre-product},
on a
$$
\text{size}(V_i) \leq \max\{\text{size}(T_i), \text{size}(W)\}
\leq \max\{\text{size}(T), \text{size}(W)\}
$$
Ainsi, d'après
Ensembles, lemme \ref{sets-lemma-bound-size},
on en déduit que
$$
\text{size}(\coprod\nolimits_{i \in I} V_i)
\leq \max\{|I|, \text{size}(T), \text{size}(W)\}.
$$
La construction de $\Sch$ montre donc
que $\coprod_{i \in I} V_i$ est isomorphe à un objet
$V$ de $(\Sch/S)_{fppf}$. Cela vérifie
l'hypothèse de
Espaces, lemme \ref{spaces-lemma-coproduct-algebraic-spaces},
ce qui conclut.
\end{proof}

\begin{lemma}
\label{bootstrap-lemma-covering-quotient}
Soit $S$ un schéma. Considérons un espace algébrique $F$ de la forme $F = U/R$,
où $(U, R, s, t, c)$ est un groupoïde en espaces algébriques
sur $S$ tel que $s, t$ soient plats et localement de présentation finie, et que
$j = (t, s) : R \to U \times_S U$ soit une relation d'équivalence.
Alors $U \to F$ est surjectif, plat et localement de présentation finie.
\end{lemma}

\begin{proof}
Ce résultat est presque, mais pas tout à fait, immédiat. En effet, d'après
Groupoïdes dans les espaces, lemme
\ref{spaces-groupoids-lemma-quotient-pre-equivalence}
et puisque $j$ est un monomorphisme, on a $R = U \times_F U$.
Choisissons un schéma $W$ et un morphisme étale surjectif $W \to F$.
Comme $U \to F$ est un épimorphisme de faisceaux, il existe un recouvrement fppf
$\{W_i \to W\}$ et des morphismes $W_i \to U$ relevant les morphismes $W_i \to F$.
Alors
$$
W_i \times_F U = W_i \times_U U \times_F U = W_i \times_{U, t} R
$$
et la projection $W_i \times_F U \to W_i$ est le changement de base de
$t : R \to U$, donc elle est plate et localement de présentation finie, voir
Morphismes d'espaces, lemmes
\ref{spaces-morphisms-lemma-base-change-flat} et
\ref{spaces-morphisms-lemma-base-change-finite-presentation}.
Ainsi, d'après
Descente sur les espaces, lemmes
\ref{spaces-descent-lemma-descending-property-flat} et
\ref{spaces-descent-lemma-descending-property-locally-finite-presentation}
on voit que $U \to F$ est plat et localement de présentation finie.
Il est surjectif d'après
Espaces, remarque \ref{spaces-remark-warning}.
\end{proof}

\begin{lemma}
\label{bootstrap-lemma-quotient-free-action}
Soit $S$ un schéma. Soit $X \to B$ un morphisme d'espaces algébriques sur
$S$. Soit $G$ un espace algébrique en groupes sur $B$ et soit
$a : G \times_B X \to X$ une action de $G$ sur $X$ au-dessus de $B$.
Si
\begin{enumerate}
\item $a$ est une action libre ;
\item $G \to B$ est plat et localement de présentation finie,
\end{enumerate}
alors $X/G$ (voir
Groupoïdes dans les espaces, définition
\ref{spaces-groupoids-definition-quotient-sheaf})
est un espace algébrique, le morphisme $X \to X/G$ est
surjectif, plat et localement de présentation finie, et
$X$ est un torseur fppf sous $G$ au-dessus de $X/G$.
\end{lemma}

\begin{proof}
Le fait que $X/G$ soit un espace algébrique résulte immédiatement du
théorème \ref{bootstrap-theorem-final-bootstrap}
et des définitions. En effet, $X/G = X/R$, où $R = G \times_B X$.
Les morphismes $s, t : G \times_B X \to X$ sont plats et localement de
présentation finie (c'est clair pour $s$, qui est un changement de base de $G \to B$,
et cela en résulte pour $t$ par symétrie, en utilisant l'inverse), et le morphisme
$j : G \times_B X \to X \times_B X$ est un monomorphisme d'après
Groupoïdes dans les espaces, lemme \ref{spaces-groupoids-lemma-free-action},
puisque l'action est libre. Le morphisme $X \to X/G$ est surjectif, plat et
localement de présentation finie d'après le lemme \ref{bootstrap-lemma-covering-quotient}.
Pour voir que $X \to X/G$ est un torseur fppf sous $G$
(Groupoïdes dans les espaces, définition
\ref{spaces-groupoids-definition-principal-homogeneous-space})
il faut montrer que $G \times_S X \to X \times_{X/G} X$
est un isomorphisme et que $X \to X/G$ admet localement des sections pour la topologie fppf.
La seconde assertion résulte des propriétés déjà établies de $X \to X/G$.
Le morphisme $G \times_S X \to X \times_{X/G} X$ est injectif (comme morphisme de
faisceaux fppf), puisque l'action est libre. Enfin, il est aussi surjectif
comme morphisme de faisceaux d'après Groupoïdes dans les espaces, lemme
\ref{spaces-groupoids-lemma-quotient-pre-equivalence}.
Cela achève la démonstration.
\end{proof}

\begin{lemma}
\label{bootstrap-lemma-descent-torsor}
Soit $\{S_i \to S\}_{i \in I}$ un recouvrement de $(\Sch/S)_{fppf}$.
Soit $G$ un espace algébrique en groupes sur $S$, et notons
$G_i = G_{S_i}$ ses changements de base. Supposons donnés
\begin{enumerate}
\item pour chaque $i \in I$, un torseur fppf sous $G_i$, noté $X_i$, au-dessus de $S_i$ ;
\item pour chaque $i, j \in I$, un isomorphisme $G_{S_i \times_S S_j}$-équivariant
$\varphi_{ij} : X_i \times_S S_j \to S_i \times_S X_j$ satisfaisant à la condition
de cocycle sur chaque $S_i \times_S S_j \times_S S_j$.
\end{enumerate}
Alors il existe un torseur fppf sous $G$, noté $X$, au-dessus de $S$
dont le changement de base à $S_i$ est isomorphe à $X_i$, de telle sorte que l'on
retrouve la donnée de descente $\varphi_{ij}$.
\end{lemma}

\begin{proof}
Nous pouvons considérer $X_i$ comme un faisceau sur $(\Sch/S_i)_{fppf}$, voir
Espaces, section \ref{spaces-section-change-base-scheme}.
D'après
Sites, section \ref{sites-section-glueing-sheaves},
la donnée de descente $(X_i, \varphi_{ij})$ est effective au sens où
il existe un unique faisceau $X$ sur $(\Sch/S)_{fppf}$ qui
redonne les espaces algébriques $X_i$ après restriction à
$(\Sch/S_i)_{fppf}$. On a donc
$X_i = h_{S_i} \times X$. D'après le
lemme \ref{bootstrap-lemma-locally-algebraic-space},
$X$ est un espace algébrique, sous réserve de vérifier que $\coprod X_i$
est un espace algébrique, ce que nous ferons à la fin de la démonstration.
Par l'équivalence de catégories de
Sites, lemme \ref{sites-lemma-mapping-property-glue},
les morphismes d'action $G_i \times_{S_i} X_i \to X_i$
se recollent en un morphisme $a : G \times_S X \to X$.
Il faut maintenant montrer que $a$ est une action, que $X$
est un pseudo-torseur et qu'il est localement trivial pour la topologie fppf (voir
Groupoïdes dans les espaces,
définition \ref{spaces-groupoids-definition-principal-homogeneous-space}).
Ces propriétés se vérifient localement pour la topologie fppf et
découlent donc des propriétés correspondantes des actions
$G_i \times_{S_i} X_i \to X_i$. Le lemme est ainsi démontré.

\medskip\noindent
Nous suggérons au lecteur de passer la suite de la démonstration, qui est purement
ensembliste. Choisissons des recouvrements $\{S_{ij} \to S_j\}_{j \in J_i}$ de
$(\Sch/S)_{fppf}$
qui trivialisent les torseurs sous $G_i$ notés $X_i$ (c'est possible par hypothèse et par
Topologies, lemme \ref{topologies-lemma-fppf-induced}, partie (1)).
Alors $\{S_{ij} \to S\}_{i \in I, j \in J_i}$ est un recouvrement de
$(\Sch/S)_{fppf}$ ; nous pouvons donc supposer que chaque $X_i$
est le torseur trivial ! Bien entendu, nous pouvons encore raffiner le recouvrement ;
nous pouvons donc supposer que chaque $S_i$ est affine et que le cardinal de l'ensemble
d'indices $I$ est borné par le cardinal de l'ensemble des points
de $S$. Choisissons $U \in \Ob((\Sch/S)_{fppf})$ et un morphisme
étale surjectif $U \to G$. Alors $U_i = U \times_S S_i$ est muni
d'un morphisme étale surjectif vers $X_i \cong G_i$. D'après
Ensembles, lemme \ref{sets-lemma-bound-size-fibre-product},
on a $\text{size}(U_i) \leq \max\{\text{size}(U), \text{size}(S_i)\}$. D'après
Ensembles, lemme \ref{sets-lemma-bound-finite-type},
on a $\text{size}(S_i) \leq \text{size}(S)$.
Ainsi,
$\text{size}(U_i) \leq \max\{\text{size}(U), \text{size}(S)\}$
pour tout $i \in I$. En combinant cela avec la borne sur $|I|$ obtenue plus haut,
Ensembles, lemme \ref{sets-lemma-bound-size},
donne $\text{size}(\coprod U_i) \leq \max\{\text{size}(U), \text{size}(S)\}$.
Par conséquent,
Espaces, lemme \ref{spaces-lemma-coproduct-algebraic-spaces},
s'applique et montre que $\coprod X_i$ est un espace algébrique, comme
il fallait le prouver.
\end{proof}




\section{Espaces algébriques pour la topologie étale}
\label{bootstrap-section-spaces-etale}

\noindent
Soit $S$ un schéma. Au lieu de travailler avec des faisceaux sur
le grand site fppf $(\Sch/S)_{fppf}$, nous pourrions travailler avec des faisceaux
sur le grand site étale $(\Sch/S)_\etale$. Tous les résultats de
Espaces algébriques, sections \ref{spaces-section-representable} et
\ref{spaces-section-representable-properties}
ont un sens pour les faisceaux sur $(\Sch/S)_\etale$.
On obtient ainsi une seconde notion d'espace algébrique en travaillant pour la
topologie étale. Cette notion est a priori plus faible que celle introduite dans
Espaces algébriques, définition \ref{spaces-definition-algebraic-space},
puisqu'un faisceau pour la topologie fppf est assurément un faisceau pour la
topologie étale. Les deux notions sont toutefois équivalentes, comme le montre le
lemme suivant.

\begin{lemma}
\label{bootstrap-lemma-spaces-etale}
La catégorie sous-jacente commune à $\Sch_{fppf}$ et $\Sch_\etale$ sera notée
$\Sch_\alpha$
(voir Topologies, remarque \ref{topologies-remark-choice-sites}).
Soit $S$ un objet de $\Sch_\alpha$. Soit
$$
F : (\Sch_\alpha/S)^{opp} \longrightarrow \textit{Ensembles}
$$
un préfaisceau possédant les propriétés suivantes :
\begin{enumerate}
\item $F$ est un faisceau pour la topologie étale ;
\item la diagonale $\Delta : F \to F \times F$ est représentable ;
\item il existe $U \in \Ob(\Sch_\alpha/S)$
et un morphisme $U \to F$ surjectif et étale.
\end{enumerate}
Alors $F$ est un espace algébrique au sens de
Espaces algébriques, définition \ref{spaces-definition-algebraic-space}.
\end{lemma}

\begin{proof}
Notons que les propriétés (2) et (3) du lemme et les propriétés correspondantes
(2) et (3) de
Espaces algébriques, définition \ref{spaces-definition-algebraic-space},
ne dépendent pas de la topologie. En effet, ces propriétés
ne font intervenir que la notion de produit fibré de préfaisceaux, les morphismes de
préfaisceaux, la notion de transformation représentable de foncteurs
et ce que signifie pour une telle transformation d'être surjective et étale.
Il ne reste donc qu'à prouver qu'un faisceau étale $F$ possédant les propriétés
(2) et (3) est aussi un faisceau fppf.

\medskip\noindent
Pour cela, posons $R = U \times_F U$. D'après (2), le préfaisceau $R$ est représenté
par un schéma et, d'après (3), les projections $R \to U$ sont étales. Ainsi,
$j : R \to U \times_S U$ est une relation d'équivalence étale. De plus,
$U \to F$ identifie $F$ au quotient de $U$ par $R$ pour la
topologie étale : (a) si $T \to F$ est un morphisme, alors $\{T \times_F U \to T\}$
est un recouvrement étale, donc $U \to F$ est un épimorphisme de faisceaux pour la
topologie étale ; (b) si $a, b : T \to U$ ont pour image la même section de $F$,
alors $(a, b) : T \to R$, donc $a$ et $b$ ont la même image dans le quotient
de $U$ par $R$ pour la topologie étale. Notons ensuite $U/R$ le faisceau quotient
pour la topologie fppf ; c'est un espace algébrique d'après
Espaces, théorème \ref{spaces-theorem-presentation}.
Nous avons donc des morphismes (transformations de foncteurs)
$$
U \to F \to U/R.
$$
D'après le
théorème \ref{spaces-theorem-presentation} d'Espaces cité plus haut,
le composé est représentable, surjectif et étale. Ainsi, pour tout
schéma $T$ et tout morphisme $T \to U/R$, le produit fibré $V = T \times_{U/R} U$
est un schéma étale et surjectif sur $T$. Autrement dit, $\{V \to U\}$
est un recouvrement étale. Cela prouve que $U \to U/R$ est surjectif comme
morphisme de faisceaux pour la topologie étale. Il s'ensuit que
$F \to U/R$ est surjectif comme morphisme de faisceaux pour la topologie étale.
D'autre part, le morphisme $F \to U/R$ est injectif (comme morphisme de préfaisceaux),
puisque $R = U \times_{U/R} U$, de nouveau d'après
Espaces, théorème \ref{spaces-theorem-presentation}.
Ainsi, $F \to U/R$ est un isomorphisme de faisceaux étales, voir
Sites, lemme \ref{sites-lemma-mono-epi-sheaves},
ce qui achève la démonstration.
\end{proof}

\noindent
Il existe aussi un analogue de
Espaces, lemme \ref{spaces-lemma-etale-locally-representable-gives-space}.

\begin{lemma}
\label{bootstrap-lemma-spaces-etale-locally-representable}
La catégorie sous-jacente commune à $\Sch_{fppf}$ et $\Sch_\etale$ sera notée
$\Sch_\alpha$
(voir Topologies, remarque \ref{topologies-remark-choice-sites}).
Soit $S$ un objet de $\Sch_\alpha$. Soit
$$
F : (\Sch_\alpha/S)^{opp} \longrightarrow \textit{Ensembles}
$$
un préfaisceau possédant les propriétés suivantes :
\begin{enumerate}
\item $F$ est un faisceau pour la topologie étale ;
\item il existe un espace algébrique $U$ sur $S$
et un morphisme $U \to F$ représentable par
des espaces algébriques, surjectif et étale.
\end{enumerate}
Alors $F$ est un espace algébrique au sens de
Espaces algébriques, définition \ref{spaces-definition-algebraic-space}.
\end{lemma}

\begin{proof}
Posons $R = U \times_F U$. C'est un espace algébrique puisque $U \to F$ est supposé
représentable par des espaces algébriques. Les projections $s, t : R \to U$ sont des
morphismes étales d'espaces algébriques puisque $U \to F$ est supposé étale.
Le morphisme $j = (t, s) : R \to U \times_S U$ est un monomorphisme et une
relation d'équivalence puisque $R = U \times_F U$. D'après le
théorème \ref{bootstrap-theorem-final-bootstrap},
le faisceau quotient fppf $F' = U/R$ est un espace algébrique.
Le morphisme $U \to F'$ est surjectif, plat et localement de présentation
finie d'après le lemme \ref{bootstrap-lemma-covering-quotient}.
Le morphisme $R \to U \times_{F'} U$ est surjectif comme morphisme de faisceaux
fppf d'après Groupoïdes dans les espaces, lemme
\ref{spaces-groupoids-lemma-quotient-pre-equivalence}
et, puisque $j$ est un monomorphisme, c'est un isomorphisme.
Ainsi, le changement de base de $U \to F'$ par $U \to F'$ est étale,
et l'on en déduit que $U \to F'$ est étale d'après
Descente sur les espaces, lemme \ref{spaces-descent-lemma-descending-property-etale}.
Le morphisme $U \to F'$ est donc surjectif comme morphisme de faisceaux étales.
Cela signifie que $F'$ est égal au faisceau quotient $U/R$
pour la topologie étale (petite vérification omise). On obtient donc
une factorisation canonique $U \to F' \to F$, et $F' \to F$ est un morphisme
injectif de faisceaux. D'autre part, $U \to F$ est surjectif comme morphisme
de faisceaux étales, et il en va donc de même de $F' \to F$. Ainsi $F' = F$,
ce qui achève la démonstration.
\end{proof}

\noindent
En fait, il suffit de disposer d'un recouvrement lisse par un schéma et de supposer
que la diagonale est représentable par des espaces algébriques.

\begin{lemma}
\label{bootstrap-lemma-spaces-etale-smooth-cover}
La catégorie sous-jacente commune à $\Sch_{fppf}$
et $\Sch_\etale$ sera notée $\Sch_\alpha$ (voir
Topologies, remarque \ref{topologies-remark-choice-sites}). Soit $S$ un objet
de $\Sch_\alpha$. Soit
$$
F : (\Sch_\alpha/S)^{opp} \longrightarrow \textit{Ensembles}
$$
un préfaisceau possédant les propriétés suivantes :
\begin{enumerate}
\item $F$ est un faisceau pour la topologie étale ;
\item la diagonale $\Delta : F \to F \times F$ est représentable
par des espaces algébriques ;
\item il existe $U \in \Ob(\Sch_\alpha/S)$
et un morphisme $U \to F$ surjectif et lisse.
\end{enumerate}
Alors $F$ est un espace algébrique au sens de
Espaces algébriques, définition \ref{spaces-definition-algebraic-space}.
\end{lemma}

\begin{proof}
La démonstration suit celle du lemme \ref{bootstrap-lemma-spaces-etale}. Posons
$R = U \times_F U$. D'après (2), le préfaisceau $R$ est un espace algébrique et, d'après (3),
les projections $R \to U$ sont lisses et surjectives. Notons $(U, R, s, t, c)$
le groupoïde associé à la relation d'équivalence $j : R \to U \times_S U$
(voir Groupoïdes dans les espaces, lemme
\ref{spaces-groupoids-lemma-equivalence-groupoid}).
D'après le théorème \ref{bootstrap-theorem-final-bootstrap}, $X = U/R$ (quotient
pour la topologie fppf) est un espace algébrique. Comme la topologie lisse
et la topologie étale ont les mêmes faisceaux (d'après
Compléments sur les morphismes, lemme \ref{more-morphisms-lemma-etale-dominates-smooth}),
le morphisme $U \to F$ identifie $F$ au quotient de
$U$ par $R$ pour la topologie lisse (détails omis).
Nous avons donc des morphismes (transformations de foncteurs)
$$
U \to F \to X.
$$
D'après le lemme \ref{bootstrap-lemma-covering-quotient}, $U \to X$ est
surjectif, plat et localement de présentation finie. D'après
Groupoïdes dans les espaces, lemme
\ref{spaces-groupoids-lemma-quotient-pre-equivalence}
(et puisque $j$ est un monomorphisme), on a $R = U \times_X U$. D'après
Descente sur les espaces, lemme \ref{spaces-descent-lemma-descending-property-smooth},
on en déduit que $U \to X$ est lisse et surjectif (puisque les projections
$R \to U$ sont lisses et surjectives et que $\{U \to X\}$ est un recouvrement
fppf). Ainsi, pour tout schéma $T$ et tout morphisme $T \to X$, le produit fibré
$T \times_X U$ est un espace algébrique lisse et surjectif sur $T$.
Choisissons un schéma $V$ et un morphisme étale surjectif $V \to T \times_X U$.
Alors $\{V \to T\}$ est un recouvrement lisse tel que $V \to T \to X$
se relève en un morphisme $V \to U$. Cela prouve que
$U \to X$ est surjectif comme morphisme de faisceaux pour la topologie lisse.
Il s'ensuit que $F \to X$ est surjectif comme morphisme de faisceaux pour la topologie
lisse. D'autre part, le morphisme $F \to X$ est injectif (comme morphisme
de préfaisceaux), puisque $R = U \times_X U$.
Ainsi, $F \to X$ est un isomorphisme de faisceaux lisses ($=$ étales),
voir Sites, lemme \ref{sites-lemma-mono-epi-sheaves},
ce qui achève la démonstration.
\end{proof}

\noindent
Enfin, voici l'analogue de
Espaces, lemme \ref{spaces-lemma-etale-locally-representable-gives-space},
où l'espace est recouvert par un morphisme lisse.

\begin{lemma}
\label{bootstrap-lemma-spaces-smooth-locally-representable}
La catégorie sous-jacente commune à $\Sch_{fppf}$ et $\Sch_\etale$ sera notée
$\Sch_\alpha$
(voir Topologies, remarque \ref{topologies-remark-choice-sites}).
Soit $S$ un objet de $\Sch_\alpha$. Soit
$$
F : (\Sch_\alpha/S)^{opp} \longrightarrow \textit{Ensembles}
$$
un préfaisceau possédant les propriétés suivantes :
\begin{enumerate}
\item $F$ est un faisceau pour la topologie étale ;
\item il existe un espace algébrique $U$ sur $S$
et un morphisme $U \to F$ représentable par
des espaces algébriques, surjectif et lisse.
\end{enumerate}
Alors $F$ est un espace algébrique au sens de
Espaces algébriques, définition \ref{spaces-definition-algebraic-space}.
\end{lemma}

\begin{proof}
La démonstration est identique à celle du
lemme \ref{bootstrap-lemma-spaces-etale-locally-representable}.
Posons $R = U \times_F U$. C'est un espace algébrique puisque $U \to F$ est supposé
représentable par des espaces algébriques. Les projections $s, t : R \to U$ sont des
morphismes lisses d'espaces algébriques puisque $U \to F$ est supposé lisse.
Le morphisme $j = (t, s) : R \to U \times_S U$ est un monomorphisme et une
relation d'équivalence puisque $R = U \times_F U$. D'après le
théorème \ref{bootstrap-theorem-final-bootstrap},
le faisceau quotient fppf $F' = U/R$ est un espace algébrique.
Le morphisme $U \to F'$ est surjectif, plat et localement de présentation
finie d'après le lemme \ref{bootstrap-lemma-covering-quotient}.
Le morphisme $R \to U \times_{F'} U$ est surjectif comme morphisme de faisceaux
fppf d'après Groupoïdes dans les espaces, lemme
\ref{spaces-groupoids-lemma-quotient-pre-equivalence}
et, puisque $j$ est un monomorphisme, c'est un isomorphisme.
Ainsi, le changement de base de $U \to F'$ par $U \to F'$ est lisse,
et l'on en déduit que $U \to F'$ est lisse d'après
Descente sur les espaces, lemme \ref{spaces-descent-lemma-descending-property-smooth}.
Le morphisme $U \to F'$ est donc surjectif comme morphisme de faisceaux étales (car
la topologie lisse et la topologie étale ont les mêmes faisceaux d'après
Compléments sur les morphismes, lemme \ref{more-morphisms-lemma-etale-dominates-smooth}).
Cela signifie que $F'$ est égal au faisceau quotient $U/R$
pour la topologie étale (petite vérification omise). On obtient donc
une factorisation canonique $U \to F' \to F$, et $F' \to F$ est un morphisme
injectif de faisceaux. D'autre part, $U \to F$ est surjectif comme morphisme
de faisceaux étales (puisque les topologies lisse et étale ont les mêmes
faisceaux), et il en va donc de même de $F' \to F$. Ainsi $F' = F$,
ce qui achève la démonstration.
\end{proof}











% OK, commencez ici.
%

\chapter{Sommes amalgamées d'espaces algébriques}



\phantomsection
\label{spaces-pushouts-section-phantom}


\section{Introduction}
\label{spaces-pushouts-section-introduction}

\noindent
Le but de ce chapitre est d'étudier les sommes amalgamées dans la catégorie des espaces algébriques, sous diverses hypothèses. Une construction assez générale est donnée dans \cite{Temkin-Tyomkin} : l'un des morphismes est affine et l'autre est une immersion fermée. Nous étudions un cas particulier de cette construction dans la section \ref{spaces-pushouts-section-pushouts} où nous supposons que l'un des morphismes est affine et l'autre est un épaississement, situation qui apparaît souvent en théorie des déformations.

\medskip\noindent
Dans les sections \ref{spaces-pushouts-section-formal-glueing} et
\ref{spaces-pushouts-section-formal-glueing-spaces} nous discutons des diagrammes
$$
\xymatrix{
f^{-1}(X \setminus Z) \ar[r] \ar[d] & Y \ar[d]^f \\
X \setminus Z \ar[r] & X
}
$$
où $f$ est un morphisme quasi-compact et quasi-séparé d'espaces algébriques, $Z \to X$ est une immersion fermée de présentation finie, le morphisme $f^{-1}(Z) \to Z$ est un isomorphisme, et $f$ est plat le long de $f^{-1}(Z)$. Dans cette situation, nous recollons des modules quasi-cohérents sur $X \setminus Z$ et $Y$
(dans la section \ref{spaces-pushouts-section-formal-glueing}) à des modules quasi-cohérents sur $X$
et nous recollons des espaces algébriques sur $X \setminus Z$ et sur $Y$
(dans la section \ref{spaces-pushouts-section-formal-glueing-spaces}) à des espaces algébriques sur $X$.

\medskip\noindent
Dans la section \ref{spaces-pushouts-section-coequalizer-glue}, nous discutons de la manière dont les morphismes birationnels propres des espaces algébriques noethériens donnent lieu à des diagrammes de coégalisation dans la catégorie des espaces algébriques, en un sens précisé plus loin.

\medskip\noindent
Dans la section \ref{spaces-pushouts-section-compactifications}, nous utilisons la construction
des carrés distingués élémentaires
dans la section \ref{spaces-pushouts-section-elementary-dsitnguished-squares}
pour prouver le théorème de Nagata sur les compactifications dans le cadre
des espaces algébriques.






\section{Conventions}
\label{spaces-pushouts-section-conventions}

\noindent
Notre convention permanente est que tous les schémas sont contenus dans
un grand site fppf $\Sch_{fppf}$. Tous les anneaux $A$ considérés
ont la propriété que $\Spec(A)$ est (isomorphe) à un
objet de ce grand site.

\medskip\noindent
Soit $S$ un schéma et soit $X$ un espace algébrique sur $S$.
Dans ce chapitre et le suivant, nous écrirons $X \times_S X$
pour le produit de $X$ avec lui-même (dans la catégorie des espaces algébriques sur $S$), au lieu de $X \times X$.




\section{Limites inductives des espaces algébriques}
\label{spaces-pushouts-section-pushouts-generalities}

\noindent
Nous discutons brièvement des limites inductives des espaces algébriques. Soit $S$ un schéma.
Soit $\mathcal{I} \to (\Sch/S)_{fppf}$, $i \mapsto X_i$
un diagramme (voir Catégories, section \ref{categories-section-limits}).
Pour chaque $i$, nous pouvons considérer le petit site \'etale $X_{i, \etale}$
dont les objets sont des schémas \'etales sur $X_i$, voir
Propriétés des espaces, section \ref{spaces-properties-section-etale-site}.
Pour chaque morphisme $i \to j$ de $\mathcal{I}$, nous avons le morphisme
$X_i \to X_j$ et donc un foncteur de changement de base
$X_{j, \etale} \to X_{i, \etale}$.
Ainsi, nous obtenons un pseudo-foncteur de $\mathcal{I}^{opp}$ dans
la $2$-catégorie des catégories. Désignons
$$
\lim_i X_{i, \etale}
$$
la $2$-limite (voir la référence future à insérer ici). Qu'est-ce que cela signifie
concrètement ? Un objet de cette limite est un système de morphismes \'etales
$U_i \to X_i$ sur $\mathcal{I}$ tel que pour chaque $i \to j$ dans $\mathcal{I}$ le diagramme
$$
\xymatrix{
U_i \ar[r] \ar[d] & U_j \ar[d] \\
X_i \ar[r] & X_j
}
$$
est cartésien. Les morphismes entre les objets sont définis de manière évidente.
Supposons que $f_i : X_i \to T$ soit une famille de morphismes telle que
pour chaque $i \to j$, la composition $X_i \to X_j \to T$ soit égale à $f_i$.
Alors nous obtenons un foncteur $T_\etale \to \lim X_{i, \etale}$.
Avec ces notations, nous pouvons formuler notre lemme.

\begin{lemma}
\label{spaces-pushouts-lemma-colimit-agrees}
Soit $S$ un schéma. Soit $\mathcal{I} \to (\Sch/S)_{fppf}$, $i \mapsto X_i$
un diagramme de schémas sur $S$ comme ci-dessus. Supposons que
\begin{enumerate}
\item $X = \colim X_i$ existe dans la catégorie des schémas,
\item $\coprod X_i \to X$ est surjectif,
\item si $U \to X$ est \'etale et $U_i = X_i \times_X U$, alors
$U = \colim U_i$ dans la catégorie des schémas, et
\item chaque objet $(U_i \to X_i)$ de $\lim X_{i, \etale}$
avec $U_i \to X_i$ séparé est dans l'image essentielle du
foncteur $X_\etale \to \lim X_{i, \etale}$.
\end{enumerate}
Alors $X = \colim X_i$ également dans la catégorie des espaces algébriques sur $S$.
\end{lemma}

\begin{proof}
Soit $Z$ un espace algébrique sur $S$. Supposons que $f_i : X_i \to Z$ soit
une famille de morphismes telle que pour chaque $i \to j$ la composition
$X_i \to X_j \to Z$ soit égale à $f_i$. Nous devons construire un morphisme
d'espaces algébriques $f : X \to Z$ tel que nous puissions retrouver $f_i$ comme
la composition $X_i \to X \to Z$. Soit $W \to Z$ un morphisme
\'etale surjectif d'un schéma vers $Z$. Nous pouvons supposer que $W$ est une
union disjointe d'affines et en particulier nous pouvons supposer que
$W \to Z$ est séparé. Pour chaque $i$, posons
$U_i = W \times_{Z, f_i} X_i$ et notons $h_i : U_i \to W$ la projection.
Alors $U_i \to X_i$ forme un objet de $\lim X_{i, \etale}$
avec $U_i \to X_i$ séparé. Par
l'hypothèse (4), nous pouvons trouver un morphisme \'etale $U \to X$ et des isomorphismes (fonctoriels) $U_i = X_i \times_X U$. D'après l'hypothèse (3), il existe un morphisme $h : U \to W$ tel que les compositions $U_i \to U \to W$ soient $h_i$. Soit $g : U \to Z$ la composée de $h$ avec le morphisme $W \to Z$. Pour terminer la preuve, nous devons montrer que $g : U \to Z$ provient d'un morphisme $X \to Z$. Pour ce faire, considérons le morphisme $(h, h) : U \times_X U \to W \times_S W$. En composant avec $U_i \times_{X_i} U_i \to U \times_X U$, nous obtenons $(h_i, h_i)$ qui se factorise par $W \times_Z W$. Comme $U \times_X U$ est la limite inductive des schémas $U_i \times_{X_i} U_i$ par (3), nous voyons que $(h, h)$ se factorise par $W \times_Z W$. Ainsi, les deux compositions $U \times_X U \to U \to W \to Z$ sont égales. Comme chaque $U_i \to X_i$ est surjectif et que, d'après l'hypothèse (2), $U \to X$ est surjectif, le fait que $Z$ soit un faisceau pour la topologie \'etale montre que $g : U \to Z$ provient d'un morphisme $f : X \to Z$ comme voulu.
\end{proof}

\noindent
La propriété, pour un cocône, d'être une limite inductive se vérifie localement pour la topologie fpqc sur le cocône.

\begin{lemma}
\label{spaces-pushouts-lemma-pushout-fpqc-local}
Soit $S$ un schéma. Soit $B$ un espace algébrique sur $S$.
Soit $\mathcal{I} \to (\Sch/S)_{fppf}$, $i \mapsto X_i$
un diagramme d'espaces algébriques sur $B$. Soit $(X, X_i \to X)$
un cocône pour le diagramme dans la catégorie des espaces algébriques sur $B$
(Catégories, Remarque \ref{categories-remark-cones-and-cocones}).
S'il existe un recouvrement fpqc $\{U_a \to X\}_{a \in A}$ tel que
\begin{enumerate}
\item pour tout $a \in A$ nous avons
$U_a = \colim X_i \times_X U_a$
dans la catégorie des espaces algébriques sur $B$, et
\item pour tout $a, b \in A$ nous avons
$U_a \times_X U_b = \colim X_i \times_X U_a \times_X U_b$
dans la catégorie des espaces algébriques sur $B$,
\end{enumerate}
alors $X = \colim X_i$ dans la catégorie des espaces algébriques sur $B$.
\end{lemma}

\begin{proof}
En effet, pour un espace algébrique $Y$ sur $B$, un morphisme $X \to Y$ sur $B$
est la même chose qu'une collection de morphismes $U_a \to Y$ qui coïncident sur
les intersections $U_a \times_X U_b$ pour tout $a, b \in A$, voir
Descente sur les espaces, Lemme
\ref{spaces-descent-lemma-fpqc-universal-effective-epimorphisms}.
\end{proof}

\noindent
Nous allons trouver une généralisation partielle commune des
Lemmes \ref{spaces-pushouts-lemma-colimit-agrees} et \ref{spaces-pushouts-lemma-pushout-fpqc-local}
qui peut en particulier être utilisée pour réduire une construction de limite inductive à une
sous-catégorie de la catégorie de tous les espaces algébriques.

\medskip\noindent
Soit $S$ un schéma et soit $B$ un espace algébrique sur $S$. Soit
$\mathcal{I}$ une catégorie d'indices et soit $i \mapsto X_i$
un diagramme dans la catégorie des espaces algébriques sur $B$,
voir Catégories, section \ref{categories-section-limits}.
Pour chaque $i$ nous pouvons considérer le petit site \'etale
$X_{i, spaces, \etale}$
dont les objets sont les espaces algébriques \'etales sur $X_i$, voir
Propriétés des espaces, section \ref{spaces-properties-section-etale-site}.
Pour chaque morphisme $i \to j$ de $\mathcal{I}$ nous avons le morphisme
$X_i \to X_j$ et donc un foncteur d'image inverse
$X_{j, spaces, \etale} \to X_{i, spaces, \etale}$.
Ainsi nous obtenons un pseudo-foncteur de $\mathcal{I}^{opp}$ vers
la $2$-catégorie des catégories. Notons
$$
\lim_i X_{i, spaces, \etale}
$$
la $2$-limite (voir la référence future à insérer ici). Que signifie ceci
concrètement ? Un objet de cette limite est un diagramme $i \mapsto (U_i \to X_i)$ dans la catégorie des flèches des espaces algébriques sur $B$ tel que pour chaque $i \to j$ dans $\mathcal{I}$ le diagramme $$
\xymatrix{
U_i \ar[r] \ar[d] & U_j \ar[d] \\
X_i \ar[r] & X_j
}
$$ soit cartésien. Les morphismes entre les objets sont définis de manière évidente. Supposons que $f_i : X_i \to Z$ soit une famille de morphismes d'espaces algébriques sur $B$ telle que, pour chaque $i \to j$, la composée $X_i \to X_j \to Z$ soit égale à $f_i$. Alors nous obtenons un foncteur $Z_{spaces, \etale} \to \lim X_{i, spaces, \etale}$. Avec ces notations, nous pouvons formuler le lemme suivant.

\begin{lemma}
\label{spaces-pushouts-lemma-colimit-check-etale-locally}
Soit $S$ un schéma. Soit $B$ un espace algébrique sur $S$.
Soit $\mathcal{I} \to (\Sch/S)_{fppf}$, $i \mapsto X_i$
un diagramme d'espaces algébriques sur $B$. Soit $(X, X_i \to X)$
un cocône pour le diagramme dans la catégorie des espaces algébriques sur $B$
(Catégories, Remarque \ref{categories-remark-cones-and-cocones}).
Supposons que
\begin{enumerate}
\item le foncteur de changement de base
$X_{spaces, \'etale} \to \lim X_{i, spaces, \etale}$,
envoyant $U$ à $U_i = X_i \times_X U$ est une équivalence,
\item si l'on se donne
\begin{enumerate}
\item $B'$ affine et \'etale sur $B$,
\item $Z$ un schéma affine sur $B'$,
\item $U \to X \times_B B'$ un morphisme \'etale d'espaces algébriques
avec $U$ affine,
\item $f_i : U_i \to Z$ un cocône sur $B'$ du diagramme
$i \mapsto U_i = U \times_X X_i$,
\end{enumerate}
il existe un morphisme unique $f : U \to Z$ sur $B'$
tel que $f_i$ soit égal à la composition $U_i \to U \to Z$.
\end{enumerate}
Alors $X = \colim X_i$ dans la catégorie de tous les espaces algébriques sur $B$.
\end{lemma}

\begin{proof}
Dans ce paragraphe, nous réduisons au cas où $B$ est un schéma affine.
Soit $B' \to B$ un morphisme \'etale d'espaces algébriques.
Remarquons que les conditions (1) et (2) sont préservées si nous remplaçons
$B$, $X_i$, $X$ par $B'$, $X_i \times_B B'$, $X \times_B B'$.
Soit $\{B_a \to B\}_{a \in A}$ un recouvrement \'etale avec $B_a$
affine, voir Propriétés des espaces, Lemme
\ref{spaces-properties-lemma-cover-by-union-affines}.
Pour $a \in A$, notons $X_a$, $X_{a, i}$ les changements de base de $X$ et le
diagramme vers $B_a$. Pour $a, b \in A$, notons
$X_{a, b}$ et $X_{a, b, i}$ les changements de base de $X$ et le
diagramme vers $B_a \times_B B_b$.
D'après le Lemme \ref{spaces-pushouts-lemma-pushout-fpqc-local}
il suffit de prouver que $X_a = \colim X_{a, i}$ et
$X_{a, b} = \colim X_{a, b, i}$.
Cela nous réduit au cas où $B = B_a$ (un schéma affine) ou $B = B_a \times_B B_b$ (un schéma séparé). En répétant l'argument une fois de plus, nous concluons que nous pouvons supposer que $B$ est un schéma affine (cela utilise le fait que l'intersection d'ouverts affines dans un schéma séparé est affine).

\medskip\noindent
Supposons que $B$ soit un schéma affine. Soit $Z$ un espace algébrique sur $B$.
Nous devons montrer que
$$
\Mor_B(X, Z) \longrightarrow \lim \Mor_B(X_i, Z)
$$
est une bijection.

\medskip\noindent
Preuve de l'injectivité. Soit $f, g : X \to Z$ des morphismes tels que les composées $f_i, g_i : X_i \to Z$ soient les mêmes pour tous les $i$. Choisissons un schéma affine $Z'$ et un morphisme \'etale $Z' \to Z$. Par Propriétés des espaces, Lemme \ref{spaces-properties-lemma-cover-by-union-affines}, nous savons que nous pouvons recouvrir $Z$ par de tels schémas affines. Posons $U = X \times_{f, Z} Z'$ et $U' = X \times_{g, Z} Z'$ et notons $p : U \to X$ et $p' : U' \to X$ les projections. Comme $f_i = g_i$ pour tous les $i$, nous voyons que $$
U_i = X_i \times_{f_i, Z} Z' = X_i \times_{g_i, Z} Z' = U'_i
$$ est compatible avec les morphismes de transition. Par (1) il existe un unique isomorphisme $\epsilon : U \to U'$ comme espaces algébriques sur $X$, c'est-à-dire avec $p = p' \circ \epsilon$ qui est compatible avec les identifications affichées. Choisissons un recouvrement \'etale $\{h_a : U_a \to U\}$ avec $U_a$ affine. Par (2) nous voyons que $f \circ p \circ h_a = g \circ p' \circ \epsilon \circ h_a =
g \circ p \circ h_a$. Puisque $\{h_a : U_a \to U\}$ est un recouvrement \'etale, nous concluons $f \circ p = g \circ p$. Puisque la collection de morphismes $p : U \to X$ que nous obtenons de cette manière est un recouvrement \'etale, nous concluons que $f = g$.

\medskip\noindent
Preuve de la surjectivité. Soit $f_i : X_i \to Z$ un élément du membre de droite de la flèche affichée dans le premier paragraphe de la preuve. Il suffit de trouver un recouvrement \'etale $\{U_c \to X\}_{c \in C}$ tel que les familles $f_{c, i} \in \lim_i \Mor_B(X_i \times_X U_c, Z)$ proviennent de morphismes $f_c : U_c \to Z$. En effet, par l'unicité prouvée ci-dessus, les morphismes $f_c$ coïncideront sur $U_c \times_X U_b$ et se recolleront donc en le morphisme souhaité $f : X \to Z$. Pour construire ce recouvrement, choisissons d'abord un recouvrement \'etale $\{g_a : Z_a \to Z\}_{a \in A}$ où chaque $Z_a$ est affine. Posons ensuite $U_{a, i} = X_i \times_{f_i, Z} Z_a$. Par (1), il existe des espaces algébriques $U_a$, \'etales sur $X$, tels que $U_{a, i} = X_i \times_X U_a$. Choisissons enfin un recouvrement \'etale $\{U_{a, b} \to U_a\}_{b \in B_a}$ avec $U_{a, b}$ affine et considérons les morphismes $$
U_{a, b, i} = X_i \times_X U_{a, b} \to
X_i \times_X U_a = X_i \times_{f_i, Z} Z_a \to Z_a
$$. Par (2), nous obtenons des morphismes $f_{a, b} : U_{a, b} \to Z_a$ compatibles avec ceux-ci. Posons $C = \coprod_{a \in A} B_a$ et, si $c \in C$ correspond à $b \in B_a$, posons $U_c = U_{a, b}$ et $f_c = g_a \circ f_{a, b} : U_c \to Z$. Cela conclut la preuve.
\end{proof}

\noindent
Voici une application de ces idées pour réduire le cas général au cas des espaces algébriques séparés.

\begin{lemma}
\label{spaces-pushouts-lemma-colimit-separated-enough}
Soit $S$ un schéma. Soit $B$ un espace algébrique sur $S$.
Soit $\mathcal{I} \to (\Sch/S)_{fppf}$, $i \mapsto X_i$
un diagramme d'espaces algébriques sur $B$. Supposons que
\begin{enumerate}
\item chaque $X_i$ est séparé sur $B$,
\item $X = \colim X_i$ existe dans la catégorie des
espaces algébriques séparés sur $B$,
\item $\coprod X_i \to X$ est surjectif,
\item si $U \to X$ est un morphisme \'etale séparé d'espaces algébriques et
$U_i = X_i \times_X U$, alors $U = \colim U_i$ dans
la catégorie des espaces algébriques séparés sur $B$, et
\item chaque objet $(U_i \to X_i)$ de $\lim X_{i, spaces, \etale}$
avec $U_i \to X_i$ séparé est de la forme $U_i = X_i \times_X U$
pour un certain morphisme \'etale séparé d'espaces algébriques $U \to X$.
\end{enumerate}
Alors $X = \colim X_i$ dans la catégorie de tous les espaces algébriques sur $B$.
\end{lemma}

\begin{proof}
Nous encourageons le lecteur à consulter plutôt
le Lemme \ref{spaces-pushouts-lemma-colimit-check-etale-locally}
et sa démonstration.

\medskip\noindent
Soit $Z$ un espace algébrique sur $B$. Supposons que $f_i : X_i \to Z$ soit
une famille de morphismes telle que pour chaque $i \to j$, la composition
$X_i \to X_j \to Z$ soit égale à $f_i$. Nous devons construire un morphisme
d'espaces algébriques $f : X \to Z$ sur $B$ de sorte que nous puissions retrouver $f_i$ comme
la composition $X_i \to X \to Z$. Soit $W \to Z$ un morphisme surjectif
\'etale d'un schéma vers $Z$. Nous pouvons supposer que $W$ est une
union disjointe d'affines et en particulier nous pouvons supposer que
$W \to Z$ est séparé et que $W$ est séparé sur $B$. Pour chaque $i$, posons
$U_i = W \times_{Z, f_i} X_i$ et notons $h_i : U_i \to W$ la projection.
Alors $U_i \to X_i$ forme un objet de $\lim X_{i, spaces, \etale}$
avec $U_i \to X_i$ séparé. Par l'hypothèse (5) nous pouvons trouver un morphisme \'etale séparé $U \to X$ d'espaces algébriques et des isomorphismes (fonctoriels) $U_i = X_i \times_X U$.  
D'après l'hypothèse (4), il existe un morphisme $h : U \to W$ sur $B$ tel que les compositions $U_i \to U \to W$ soient $h_i$.  
Soit $g : U \to Z$ la composée de $h$ avec le morphisme $W \to Z$. Pour finir la démonstration, nous devons montrer que $g : U \to Z$ provient d'un morphisme $X \to Z$. Pour ce faire, considérons le morphisme $(h, h) : U \times_X U \to W \times_S W$.  
En composant avec $U_i \times_{X_i} U_i \to U \times_X U$ nous obtenons $(h_i, h_i)$ qui se factorise par $W \times_Z W$. Comme $U \times_X U$ est la limite inductive des espaces algébriques $U_i \times_{X_i} U_i$ dans la catégorie des espaces algébriques séparés sur $B$, d'après (4), nous voyons que $(h, h)$ se factorise par $W \times_Z W$. Ainsi, les deux compositions $U \times_X U \to U \to W \to Z$ sont égales. Comme chaque $U_i \to X_i$ est surjectif et que, d'après l'hypothèse (2), $U \to X$ est surjectif, le fait que $Z$ soit un faisceau pour la topologie \'etale montre que $g : U \to Z$ provient d'un morphisme $f : X \to Z$ comme voulu. 
\end{proof}





\section{Descente des faisceaux \'etales}
\label{spaces-pushouts-section-glueing-etale}

\noindent
Cette section est l'analogue pour les espaces algébriques de
Cohomologie \'etale, section \ref{etale-cohomology-section-glueing-etale}.

\medskip\noindent
Afin d'exprimer commodément nos résultats, nous avons besoin de quelques notations.
Soit $S$ un schéma.
Soit $\mathcal{U} = \{f_i : X_i \to X\}$
une famille de morphismes d'espaces algébriques sur $S$ de même but.
Une {\it donnée de descente} pour les faisceaux \'etales par rapport à
$\mathcal{U}$ est une famille
$((\mathcal{F}_i)_{i \in I}, (\varphi_{ij})_{i, j \in I})$ où
\begin{enumerate}
\item $\mathcal{F}_i$ est dans $\Sh(X_{i, \etale})$, et
\item $\varphi_{ij} :
\text{pr}_{0, small}^{-1} \mathcal{F}_i
\longrightarrow
\text{pr}_{1, small}^{-1} \mathcal{F}_j$
est un isomorphisme dans $\Sh((X_i \times_X X_j)_\etale)$
\end{enumerate}
et telle que la {\it condition de cocycle} soit satisfaite : les diagrammes
$$
\xymatrix{
\text{pr}_{0, small}^{-1}\mathcal{F}_i
\ar[dr]_{\text{pr}_{02, small}^{-1}\varphi_{ik}}
\ar[rr]^{\text{pr}_{01, small}^{-1}\varphi_{ij}} & &
\text{pr}_{1, small}^{-1}\mathcal{F}_j
\ar[dl]^{\text{pr}_{12, small}^{-1}\varphi_{jk}} \\
& \text{pr}_{2, small}^{-1}\mathcal{F}_k
}
$$
commutent dans $\Sh((X_i \times_X X_j \times_X X_k)_\etale)$.
Il existe une notion évidente de {\it morphismes de données de descente}
et nous obtenons une catégorie de données de descente.
Une donnée de descente
$((\mathcal{F}_i)_{i \in I}, (\varphi_{ij})_{i, j \in I})$
est dite {\it effective} s'il existe un objet
$\mathcal{F}$ de $\Sh(X_\etale)$ et des isomorphismes
$\varphi_i : f_{i, small}^{-1} \mathcal{F} \to \mathcal{F}_i$
dans $\Sh(X_{i, \etale})$ compatibles avec les $\varphi_{ij}$, c'est-à-dire
$$
\varphi_{ij} =
\text{pr}_{1, small}^{-1} (\varphi_j) \circ
\text{pr}_{0, small}^{-1} (\varphi_i^{-1})
$$
Une autre façon de le dire est la suivante. Étant donné un objet $\mathcal{F}$
de $\Sh(X_\etale)$, nous obtenons la {\it donnée de descente canonique}
$(f_{i, small}^{-1}\mathcal{F}, c_{ij})$ où $c_{ij}$
est l'isomorphisme canonique
$$
c_{ij} : \text{pr}_{0, small}^{-1} f_{i, small}^{-1}\mathcal{F}
\longrightarrow
\text{pr}_{1, small}^{-1} f_{j, small}^{-1}\mathcal{F}
$$
La donnée de descente
$((\mathcal{F}_i)_{i \in I}, (\varphi_{ij})_{i, j \in I})$
est effective si et seulement si elle est isomorphe à la donnée de descente canonique associée à un objet $\mathcal{F}$ dans $\Sh(X_\etale)$.

\medskip\noindent
Si la famille consiste en un seul morphisme $\{X \to Y\}$,
alors nous considérons une donnée de descente comme une paire $(\mathcal{F}, \varphi)$
où $\mathcal{F}$ est un objet de $\Sh(X_\etale)$ et
$\varphi$ est un isomorphisme
$$
\text{pr}_{0, small}^{-1} \mathcal{F}
\longrightarrow
\text{pr}_{1, small}^{-1} \mathcal{F}
$$
dans $\Sh((X \times_Y X)_\etale)$ tel que la condition de cocycle soit satisfaite :
$$
\xymatrix{
\text{pr}_{0, small}^{-1}\mathcal{F}
\ar[dr]_{\text{pr}_{02, small}^{-1}\varphi}
\ar[rr]^{\text{pr}_{01, small}^{-1}\varphi} & &
\text{pr}_{1, small}^{-1}\mathcal{F}
\ar[dl]^{\text{pr}_{12, small}^{-1}\varphi} \\
& \text{pr}_{2, small}^{-1}\mathcal{F}
}
$$
commute dans $\Sh((X \times_Y X \times_Y X)_\etale)$.
Il existe une notion de morphismes de données de descente et d'effectivité
exactement comme auparavant.

\begin{lemma}
\label{spaces-pushouts-lemma-glue-etale-sheaf-etale}
Soit $S$ un schéma. Soit $\{f_i : X_i \to X\}$ un recouvrement \'etale d'espaces algébriques. Le foncteur
$$
\Sh(X_\etale)
\longrightarrow
\text{données de descente pour les faisceaux \'etales relativement à }\{f_i : X_i \to X\}
$$
est une équivalence de catégories.
\end{lemma}

\begin{proof}
Dans Propriétés des espaces, section \ref{spaces-properties-section-etale-site}
nous avons défini un site $X_{spaces, \etale}$ dont les objets
sont des espaces algébriques \'etales sur $X$ avec des recouvrements \'etales.
De plus, nous avons des identifications
$\Sh(X_\etale) = \Sh(X_{spaces, \etale})$ compatibles
avec les morphismes d'espaces algébriques, c'est-à-dire compatibles avec
les images directes et les images inverses. Par conséquent, l'énoncé du lemme découle
de la discussion beaucoup plus générale dans
Sites, section \ref{sites-section-glueing-sheaves}.
\end{proof}

\begin{lemma}
\label{spaces-pushouts-lemma-reduce-to-scheme-base}
Soit $S$ un schéma. Soit $f : X \to Y$ un morphisme d'espaces algébriques sur $S$. Soit $\{Y_i \to Y\}_{i \in I}$ un recouvrement \'etale d'espaces algébriques. Si pour chaque $i \in I$ le foncteur $$
\Sh(Y_{i, \etale})
\longrightarrow
\text{données de descente pour les faisceaux \'etales relativement à }\{X \times_Y Y_i \to Y_i\}
$$ est une équivalence de catégories et pour chaque $i, j \in I$ le foncteur $$
\Sh((Y_i \times_Y Y_j)_\etale)
\longrightarrow
\text{données de descente pour les faisceaux \'etales relativement à }
\{X \times_Y Y_i \times_Y Y_j \to Y_i \times_Y Y_j\}
$$ est une équivalence de catégories, alors $$
\Sh(Y_\etale)
\longrightarrow
\text{données de descente pour les faisceaux \'etales relativement à }\{X \to Y\}
$$ est une équivalence de catégories.
\end{lemma}

\begin{proof}
Conséquence formelle du Lemme \ref{spaces-pushouts-lemma-glue-etale-sheaf-etale}
et des définitions.
\end{proof}

\begin{lemma}
\label{spaces-pushouts-lemma-representable-case}
Soit $S$ un schéma. Soit $f : X \to Y$ un morphisme d'espaces algébriques sur $S$. Supposons que $f$ soit représentable (par des schémas) et que $f$ possède l'une des propriétés suivantes :
surjectif et entier,
surjectif et propre, ou
surjectif, plat et localement de présentation finie.
Alors
$$
\Sh(Y_\etale)
\longrightarrow
\text{données de descente pour les faisceaux \'etales relativement à }\{X \to Y\}
$$
est une équivalence de catégories.
\end{lemma}

\begin{proof}
Chacune des propriétés des morphismes d'espaces algébriques
mentionnées dans l'énoncé du lemme est préservée par
un changement de base arbitraire, voir les listes dans
Espaces, section \ref{spaces-section-lists}.
Ainsi, nous pouvons appliquer le Lemme \ref{spaces-pushouts-lemma-reduce-to-scheme-base}
pour voir que nous pouvons travailler \'etale localement sur $Y$.
De cette manière, nous réduisons au cas où $Y$ est un schéma ;
certains détails sont omis. Dans ce cas, $X$ est également un schéma
et le résultat découle de Cohomologie \'etale, Lemme
\ref{etale-cohomology-lemma-glue-etale-sheaf-integral-surjective},
\ref{etale-cohomology-lemma-glue-etale-sheaf-proper-surjective}, ou
\ref{etale-cohomology-lemma-glue-etale-sheaf-fppf-cover}.
\end{proof}

\begin{lemma}
\label{spaces-pushouts-lemma-reduce-to-scheme-source}
Soit $S$ un schéma. Soit $f : X \to Y$ un morphisme d'espaces algébriques sur $S$. Soit $\pi : X' \to X$ un morphisme d'espaces algébriques. Supposons
\begin{enumerate}
\item $f \circ \pi$ est représentable (par des schémas),
\item $f \circ \pi$ possède l'une des propriétés suivantes :
surjectif et intégral,
surjectif et propre, ou
surjectif, plat et localement de présentation finie.
\end{enumerate}
Alors
$$
\Sh(Y_\etale)
\longrightarrow
\text{données de descente pour les faisceaux \'etales relativement à }\{X \to Y\}
$$
est une équivalence de catégories.
\end{lemma}

\begin{proof}
Conséquence formelle du Lemme \ref{spaces-pushouts-lemma-representable-case}
et de Champs, Lemme \ref{stacks-lemma-compare-descent-condition}.
\end{proof}

\begin{lemma}
\label{spaces-pushouts-lemma-glue-etale-sheaf-proper-surjective}
Soit $S$ un schéma. Soit $f : X \to Y$ un morphisme d'espaces algébriques sur $S$ qui possède l'une des propriétés suivantes : surjectif et intégral, surjectif et propre, ou surjectif, plat et localement de présentation finie. Alors le foncteur $$
\Sh(Y_\etale)
\longrightarrow
\text{données de descente pour les faisceaux \'etales relativement à }\{X \to Y\}
$$ est une équivalence de catégories.
\end{lemma}

\begin{proof}
Remarquons que le changement de base d'un morphisme propre et surjectif est
propre et surjectif, voir
Morphismes d'espaces, Lemme \ref{spaces-morphisms-lemma-base-change-proper}
et \ref{spaces-morphisms-lemma-base-change-surjective}.
Ainsi, d'après le Lemme \ref{spaces-pushouts-lemma-reduce-to-scheme-base}
nous pouvons travailler \'etale localement sur $Y$. Ainsi
nous nous ramenons au cas où $Y$ est un schéma affine ; quelques détails sont omis.

\medskip\noindent
Supposons que $Y$ soit affine. D'après le Lemme \ref{spaces-pushouts-lemma-reduce-to-scheme-source},
il suffit de trouver un morphisme $X' \to X$, où $X'$ est un schéma, tel que $X' \to Y$ soit surjectif et entier, surjectif et propre, ou surjectif, plat et localement de présentation finie.

\medskip\noindent
Dans le cas où $X \to Y$ est entier et surjectif, nous pouvons prendre $X' = X$,
puisqu'un morphisme entier est représentable.

\medskip\noindent
Si $f$ est propre et surjectif, alors l'espace algébrique
$X$ est quasi-compact et séparé, voir
Morphismes d'espaces, section \ref{spaces-morphisms-section-quasi-compact} et
Lemme \ref{spaces-morphisms-lemma-separated-over-separated}.
Choisissons un schéma $X'$ et un morphisme fini surjectif $X' \to X$, voir
Limites d'espaces, Proposition
\ref{spaces-limits-proposition-there-is-a-scheme-finite-over}.
Alors $X' \to Y$ est surjectif et propre.

\medskip\noindent
Enfin, si $X \to Y$ est surjectif et plat et localement de présentation finie, alors nous pouvons prendre un recouvrement \'etale affine $\{U_i \to X\}$ et définir $X'$ comme la réunion disjointe $\coprod U_i$.
\end{proof}

\begin{lemma}
\label{spaces-pushouts-lemma-glue-etale-sheaf-fppf}
Soit $S$ un schéma.
Soit $\{f_i : X_i \to X\}$ un recouvrement fppf d'espaces algébriques sur $S$.
Le foncteur
$$
\Sh(X_\etale)
\longrightarrow
\text{données de descente pour les faisceaux \'etales relativement à }\{f_i : X_i \to X\}
$$
est une équivalence de catégories.
\end{lemma}

\begin{proof}
Nous avons le Lemme \ref{spaces-pushouts-lemma-glue-etale-sheaf-proper-surjective}
pour le morphisme $f : \coprod X_i \to X$.
Alors un argument formel montre que les données de descente pour $f$
reviennent aux données de descente pour le recouvrement ; comparer
avec Descente, Lemme \ref{descent-lemma-family-is-one}.
Détails omis.
\end{proof}

\begin{lemma}
\label{spaces-pushouts-lemma-glue-etale-sheaf-modification}
Soit $S$ un schéma. Soit $f : Y' \to Y$ un morphisme propre d'espaces algébriques sur $S$. Soit $i : Z \to Y$ une immersion fermée. Posons $E = Z \times_Y Y'$. Considérons le diagramme
$$
\xymatrix{
E \ar[d]_g \ar[r]_j & Y' \ar[d]^f \\
Z \ar[r]^i & Y
}
$$
Si $f$ est un isomorphisme sur $Y \setminus Z$, alors le foncteur $$
\Sh(Y_\etale)
\longrightarrow
\Sh(Y'_\etale) \times_{\Sh(E_\etale)} \Sh(Z_\etale)
$$ est une équivalence de catégories.
\end{lemma}

\begin{proof}
Observons que $X = Y' \coprod Z \to Y$ est un morphisme propre surjectif.
Ainsi, il suffit de construire une équivalence de catégories
$$
\Sh(Y'_\etale) \times_{\Sh(E_\etale)} \Sh(Z_\etale)
\longrightarrow
\text{données de descente pour les faisceaux \'etales relativement à }\{X \to Y\}
$$
compatible avec les foncteurs d'image inverse issus de $Y$
car ensuite nous pouvons utiliser le Lemme \ref{spaces-pushouts-lemma-glue-etale-sheaf-proper-surjective}
pour conclure. Soit donc $(\mathcal{G}', \mathcal{G}, \alpha)$ un
objet de $\Sh(Y'_\etale) \times_{\Sh(E_\etale)} \Sh(Z_\etale)$ avec
les notations de Catégories, Exemple
\ref{categories-example-2-fibre-product-categories}.
Nous pouvons alors considérer le faisceau $\mathcal{F}$ sur $X$ défini
en prenant $\mathcal{G}'$ sur la composante $Y'$ et $\mathcal{G}$
sur la composante $Z$. Nous avons
$$
X \times_Y X = Y' \times_Y Y' \amalg
Y' \times_Y Z \amalg Z \times_Y Y' \amalg Z \times_Y Z =
Y' \times_Y Y' \amalg E \amalg E \amalg Z
$$
Les isomorphismes des deux images inverses de $\mathcal{F}$ sur cet espace algébrique sont évidents sur les composantes $E$, $E$, $Z$. La partie intéressante de la preuve est de trouver un isomorphisme
$\text{pr}_{0, small}^{-1}\mathcal{G}' \to
\text{pr}_{1, small}^{-1}\mathcal{G}'$
sur $Y' \times_Y Y'$ satisfaisant la condition de cocycle. Cependant, l'hypothèse que $Y' \to Y$ soit un isomorphisme sur $Y \setminus Z$ implique que $$
h : Y' \coprod E \times_Z E \longrightarrow Y' \times_Y Y'
$$ est un morphisme propre surjectif. (C'est en fait un morphisme fini car il s'agit de l'union disjointe de deux immersions fermées.) Par conséquent, il suffit de construire un isomorphisme des images inverses de $\text{pr}_{0, small}^{-1}\mathcal{G}'$ et $\text{pr}_{1, small}^{-1}\mathcal{G}'$ par $h_{small}$ satisfaisant une certaine condition de cocycle. Pour la diagonale, il est clair comment procéder. Sur $E \times_Z E$, nous utilisons le fait que les images inverses des deux faisceaux s'identifient à l'image inverse de $\mathcal{G}$ par le morphisme $E \times_Z E \to Z$. Nous omettons les détails. 
\end{proof}






\section{Descente des morphismes \'etales d'espaces algébriques}
\label{spaces-pushouts-section-descending-etale}

\noindent
Dans cette section, nous combinons les résultats de recollement pour les faisceaux \'etales donnés dans la section \ref{spaces-pushouts-section-glueing-etale} avec la flexibilité des espaces algébriques afin d'obtenir certaines affirmations de descente pour les morphismes \'etales d'espaces algébriques.

\begin{lemma}
\label{spaces-pushouts-lemma-descend-etale-proper-surjective}
Soit $S$ un schéma. Soit $f : X \to Y$ un morphisme propre et surjectif
d'espaces algébriques sur $S$. Toute donnée de descente $(U/X, \varphi)$ relative à $f$
(Descente sur les espaces, Définition \ref{spaces-descent-definition-descent-datum})
avec $U$ \'etale sur $X$ est effective
(Descente sur les espaces, Définition \ref{spaces-descent-definition-effective}).
Plus précisément, il existe un morphisme \'etale $V \to Y$ d'espaces algébriques
dont la donnée de descente canonique correspondante est isomorphe à $(U/X, \varphi)$.
\end{lemma}

\begin{proof}
Rappelons que $U$ donne lieu à un faisceau représentable
$\mathcal{F} = h_U$ dans $\Sh(X_{spaces, \etale}) = \Sh(X_\etale)$, voir
Propriétés des espaces, section \ref{spaces-properties-section-etale-site}.
La donnée de descente sur $U$ relative à $f$
donne exactement une donnée de descente $(\mathcal{F}, \varphi)$
pour les faisceaux \'etales relativement à $\{X \to Y\}$.
D'après le Lemme \ref{spaces-pushouts-lemma-glue-etale-sheaf-proper-surjective}
cette donnée de descente est effective.
Soit $\mathcal{G}$ le faisceau correspondant sur $Y_\etale$.
D'après Propriétés des espaces, Lemme \ref{spaces-properties-lemma-sheaf-gives-space}
nous obtenons un morphisme \'etale $V \to Y$ d'espaces algébriques
correspondant à $\mathcal{G}$ ; nous omettons la vérification de
la condition ensembliste\footnote{Cela découle du fait que $\mathcal{F}$ satisfait la condition correspondante.}.
L'isomorphisme donné $\mathcal{F} \to f_{small}^{-1}\mathcal{G}$
correspond à un isomorphisme $U \to V \times_Y X$ compatible avec la donnée de descente.
\end{proof}

\begin{lemma}
\label{spaces-pushouts-lemma-glue-etale-space-modification}
Soit $S$ un schéma. Soit $f : Y' \to Y$ un morphisme propre d'espaces algébriques sur $S$. Soit $i : Z \to Y$ une immersion fermée. Posons $E = Z \times_Y Y'$. Considérons le diagramme
$$
\xymatrix{
E \ar[d]_g \ar[r]_j & Y' \ar[d]^f \\
Z \ar[r]^i & Y
}
$$
Si $f$ est un isomorphisme sur $Y \setminus Z$, alors le foncteur $$
Y_{spaces, \etale}
\longrightarrow
Y'_{spaces, \etale} \times_{E_{spaces, \etale}} Z_{spaces, \etale}
$$ est une équivalence de catégories.
\end{lemma}

\begin{proof}
Soit $(V' \to Y', W \to Z, \alpha)$ un objet du membre de droite.
Rappelons que $V'$ et $W$ donnent lieu respectivement à des faisceaux représentables
$\mathcal{G}' = h_{V'}$ dans $\Sh(Y'_{spaces, \etale}) = \Sh(Y'_\etale)$,
et $\mathcal{G} = h_W$
dans $\Sh(Z_{spaces, \etale}) = \Sh(Z_\etale)$, voir
Propriétés des espaces, section \ref{spaces-properties-section-etale-site}.
L'isomorphisme $\alpha : V' \times_{Y'} E \to W \times_Z E$
détermine un isomorphisme
$j_{small}^{-1}\mathcal{G}' \to g_{small}^{-1}\mathcal{G}$ de
faisceaux sur $E$.
Par le Lemme \ref{spaces-pushouts-lemma-glue-etale-sheaf-modification},
nous obtenons un unique faisceau $\mathcal{F}$ sur $Y$ dont les images inverses s'identifient
à $\mathcal{G}'$ et $\mathcal{G}$ de manière compatible avec l'isomorphisme.
D'après Propriétés des espaces, Lemme \ref{spaces-properties-lemma-sheaf-gives-space},
nous obtenons un morphisme \'etale $V \to Y$ d'espaces algébriques
correspondant à $\mathcal{F}$ ; nous omettons la vérification de
la condition ensembliste\footnote{Cela découle du fait que $\mathcal{G}$ et $\mathcal{G}'$ satisfont la condition correspondante.}.
Les isomorphismes donnés $\mathcal{G}' \to f_{small}^{-1}\mathcal{F}$
et $\mathcal{G} \to i_{small}^{-1}\mathcal{F}$
correspondent aux isomorphismes $V' \to V \times_Y Y'$
et $W \to V \times_Y Z$ compatibles
avec $\alpha$ comme voulu.
\end{proof}









\section{Sommes amalgamées le long d'épaississements et de morphismes affines}
\label{spaces-pushouts-section-pushouts}

\noindent
Cette section est l'analogue de
Compléments sur les morphismes, section \ref{more-morphisms-section-pushouts}.

\begin{lemma}
\label{spaces-pushouts-lemma-pushout-along-thickening-schemes}
Soit $S$ un schéma. Soit $X \to X'$ un épaississement de schémas sur $S$ et soit $X \to Y$ un morphisme affine de schémas sur $S$.
Soit $Y' = Y \amalg_X X'$ la somme amalgamée dans la catégorie des schémas (voir
Compléments sur les morphismes, Lemme \ref{more-morphisms-lemma-pushout-along-thickening}).
Alors $Y'$ est aussi une somme amalgamée dans la catégorie des espaces algébriques sur $S$.
\end{lemma}

\begin{proof}
Ceci est une conséquence immédiate du Lemme \ref{spaces-pushouts-lemma-colimit-agrees} et
Compléments sur les morphismes, Lemmes
\ref{more-morphisms-lemma-pushout-along-thickening},
\ref{more-morphisms-lemma-equivalence-categories-schemes-over-pushout}, et
\ref{more-morphisms-lemma-equivalence-categories-schemes-over-pushout-flat}.
\end{proof}

\begin{lemma}
\label{spaces-pushouts-lemma-pushout-along-thickening}
Soit $S$ un schéma. Soit $X \to X'$ un épaississement d'espaces algébriques sur $S$ et soit $X \to Y$ un morphisme affine d'espaces algébriques sur $S$.
Alors il existe une somme amalgamée
$$
\xymatrix{
X \ar[r] \ar[d]_f
&
X' \ar[d]^{f'}
\\
Y \ar[r]
&
Y \amalg_X X'
}
$$
dans la catégorie des espaces algébriques sur $S$. De plus, $Y' = Y \amalg_X X'$
est un épaississement de $Y$ et
$$
\mathcal{O}_{Y'} = \mathcal{O}_Y \times_{f_*\mathcal{O}_X} f'_*\mathcal{O}_{X'}
$$
comme faisceaux sur $Y_\etale = (Y')_\etale$.
\end{lemma}

\begin{proof}
Choisissons un schéma $V$ et un morphisme \'etale surjectif $V \to Y$.
Définissons $U = V \times_Y X$. Il s'agit d'un schéma affine sur $V$ avec un morphisme \'etale surjectif $U \to X$. D'après Compléments sur les morphismes d'espaces, Lemme \ref{spaces-more-morphisms-lemma-thickening-equivalence}
il existe un $U' \to X'$ \'etale surjectif avec $U = U' \times_{X'} X$. En particulier, le morphisme de schémas $U \to U'$ est aussi un épaississement. Appliquons Compléments sur les morphismes, Lemme \ref{more-morphisms-lemma-pushout-along-thickening}
pour obtenir une somme amalgamée $V' = V \amalg_U U'$ dans la catégorie des schémas.

\medskip\noindent
Nous répétons cette procédure pour construire une somme amalgamée
$$
\xymatrix{
U \times_X U \ar[d] \ar[r] & U' \times_{X'} U' \ar[d] \\
V \times_Y V \ar[r] & R'
}
$$
dans la catégorie des schémas. Considérons les morphismes
$$
U \times_X U \to U \to V',\quad
U' \times_{X'} U' \to U' \to V',\quad
V \times_Y V \to V \to V'
$$
où nous utilisons la première projection dans chaque cas. Il est clair que ceux-ci se recollent pour donner un morphisme $t' : R' \to V'$ qui est \'etale d'après
Compléments sur les morphismes, Lemme
\ref{more-morphisms-lemma-equivalence-categories-schemes-over-pushout-flat}.
De même, nous obtenons $s' : R' \to V'$ \'etale.
Le morphisme $j' = (t', s') : R' \to V' \times_S V'$ est non ramifié
(puisque $t'$ est \'etale) et un monomorphisme lorsqu'on le restreint au sous-schéma fermé $V \times_Y V \subset R'$. Comme $V \times_Y V \subset R'$ est
un épaississement, il s'ensuit que $j'$ est également un monomorphisme. Enfin, $j'$
définit également une relation d'équivalence, car nous pouvons utiliser la fonctorialité des sommes amalgamées
de schémas pour construire un morphisme $c' : R' \times_{s', V', t'} R' \to R'$
(détails omis). À ce stade, nous posons $Y' = V'/R'$, voir
Espaces, Théorème \ref{spaces-theorem-presentation}.

\medskip\noindent
Nous avons des morphismes $X' = U'/U' \times_{X'} U' \to V'/R' = Y'$ et
$Y = V/V \times_Y V \to V'/R' = Y'$.
Par construction, ceux-ci s'insèrent dans le diagramme commutatif
$$
\xymatrix{
X \ar[r] \ar[d]_f & X' \ar[d]^{f'} \\
Y \ar[r] & Y'
}
$$
Puisque $Y \to Y'$ est un épaississement, nous avons
$Y_\etale = (Y')_\etale$, voir Compléments sur les morphismes d'espaces,
Lemme \ref{spaces-more-morphisms-lemma-thickening-equivalence}.
La commutativité du diagramme fournit un morphisme de faisceaux
$$
\mathcal{O}_{Y'}
\longrightarrow
\mathcal{O}_Y \times_{f_*\mathcal{O}_X} f'_*\mathcal{O}_{X'}
$$
sur ce site. D'après Compléments sur les morphismes, Lemme
\ref{more-morphisms-lemma-pushout-along-thickening}
ce morphisme est un isomorphisme lorsqu'on le restreint au
schéma $V'$, donc c'est un isomorphisme.

\medskip\noindent
Pour terminer la preuve, nous montrons que le diagramme ci-dessus est une somme amalgamée dans la catégorie des espaces algébriques. Pour voir cela, soit $Z$ un espace algébrique et soient $a' : X' \to Z$ et $b : Y \to Z$ des morphismes d'espaces algébriques. Par le Lemme \ref{spaces-pushouts-lemma-pushout-along-thickening-schemes}, nous obtenons un morphisme unique $h : V' \to Z$ s'insérant dans les diagrammes commutatifs $$
\vcenter{
\xymatrix{
U' \ar[d] \ar[r] & V' \ar[d]^h \\
X' \ar[r]^{a'} & Z
}
}
\quad\text{et}\quad
\vcenter{
\xymatrix{
V \ar[r] \ar[d] & V' \ar[d]^h \\
Y \ar[r]^b & Z
}
}
$$. L'unicité montre que $h \circ t' = h \circ s'$. Ainsi, $h$ se factorise de manière unique sous la forme $V' \to Y' \to Z$, ce qui conclut.
\end{proof}

\noindent
Dans le lemme suivant, nous utilisons le produit fibré de catégories au sens de
Catégories, Exemple \ref{categories-example-2-fibre-product-categories}.

\begin{lemma}
\label{spaces-pushouts-lemma-categories-spaces-over-pushout}
Soit $S$ un schéma de base. Soit $X \to X'$ un épaississement d'espaces algébriques sur $S$ et soit $X \to Y$ un morphisme affine d'espaces algébriques sur $S$.
Soit $Y' = Y \amalg_X X'$ la somme amalgamée (voir
Lemme \ref{spaces-pushouts-lemma-pushout-along-thickening}). Le changement de base donne un foncteur
$$
F :
(\textit{Espaces}/Y')
\longrightarrow
(\textit{Espaces}/Y) \times_{(\textit{Espaces}/Y')} (\textit{Espaces}/X')
$$
défini par $V' \longmapsto (V' \times_{Y'} Y, V' \times_{Y'} X', 1)$, qui
envoie $(\Sch/Y')$ dans $(\Sch/Y) \times_{(\Sch/Y')} (\Sch/X')$.
Le foncteur $F$ a un adjoint à gauche
$$
G :
(\textit{Espaces}/Y) \times_{(\textit{Espaces}/Y')} (\textit{Espaces}/X')
\longrightarrow
(\textit{Espaces}/Y')
$$
qui envoie le triple $(V, U', \varphi)$ vers la somme amalgamée
$V \amalg_{(V \times_Y X)} U'$ dans la catégorie des espaces algébriques sur $S$.
Le foncteur $G$ envoie $(\Sch/Y) \times_{(\Sch/Y')} (\Sch/X')$ dans $(\Sch/Y')$.
\end{lemma}

\begin{proof} La preuve est complètement formelle. Puisque les morphismes $X \to X'$ et $X \to Y$ sont représentables, il est clair que $F$ envoie $(\Sch/Y')$ dans $(\Sch/Y) \times_{(\Sch/Y')} (\Sch/X')$.

\medskip\noindent
Construisons $G$. Soit $(V, U', \varphi)$ un objet de la catégorie du produit fibré. Posons $U = U' \times_{X'} X$. Notons que $U \to U'$ est un épaississement. Puisque $\varphi : V \times_Y X \to U' \times_{X'} X = U$ est un isomorphisme, nous avons un morphisme $U \to V$ sur $X \to Y$ qui identifie $U$ avec le produit fibré $X \times_Y V$. En particulier, $U \to V$ est affine, voir Morphismes d'espaces, Lemme \ref{spaces-morphisms-lemma-base-change-affine}. 
Ainsi, nous pouvons appliquer le Lemme \ref{spaces-pushouts-lemma-pushout-along-thickening} pour obtenir une somme amalgamée $V' = V \amalg_U U'$. La propriété cocartésienne de $V'$ et les morphismes $V \to Y$ et $U' \to X'$, qui coïncident sur $U$ comme morphismes vers $Y'$, définissent un morphisme $V' \to Y'$. En posant $G(V, U', \varphi) = V'$, on obtient le foncteur $G$.

\medskip\noindent
Si $(V, U', \varphi)$ est un objet de $(\Sch/Y) \times_{(\Sch/Y')} (\Sch/X')$
alors $U = U' \times_{X'} X$ est aussi un schéma et nous pouvons former la somme amalgamée
$V' = V \amalg_U U'$ dans la catégorie des schémas par
Compléments sur les morphismes, Lemme \ref{more-morphisms-lemma-pushout-along-thickening}.
D'après le Lemme \ref{spaces-pushouts-lemma-pushout-along-thickening-schemes}
c'est aussi une somme amalgamée dans la catégorie des espaces algébriques, donc
$G$ envoie $(\Sch/Y) \times_{(\Sch/Y')} (\Sch/X')$ dans $(\Sch/Y')$.

\medskip\noindent
Prouvons que $G$ est un adjoint à gauche de $F$. Soit $Z$ un espace algébrique sur $Y'$. Il nous faut montrer que
$$
\Mor(V', Z) = \Mor((V, U', \varphi), F(Z))
$$
où les ensembles de morphismes sont pris dans leurs catégories respectives.
Soit $g' : V' \to Z$ un morphisme. Notons $\tilde g$ et $\tilde f'$ les composées de $g'$ avec les morphismes $V \to V'$ et $U' \to V'$, respectivement.
Effectuons les changements de base de $\tilde g$ et $\tilde f'$ le long de $Y \to Y'$ et $X' \to Y'$, respectivement, pour obtenir des morphismes $g : V \to Z \times_{Y'} Y$ et $f' : U' \to Z \times_{Y'} X'$. Alors $(g, f')$ est un élément du membre de droite de l'équation ci-dessus (détails omis).
Inversement, supposons que $(g, f') : (V, U', \varphi) \to F(Z)$ soit un élément du membre de droite. Nous pouvons considérer les composées $\tilde g : V \to Z$ et $\tilde f' : U' \to Z$ de $g$ avec $Z \times_{Y'} Y \to Z$ et de $f'$ avec $Z \times_{Y'} X' \to Z$, respectivement. Alors $\tilde g$ et $\tilde f'$ coïncident comme morphismes de $U$ vers $Z$. Par la propriété universelle du carré cocartésien, nous obtenons un morphisme $g' : V' \to Z$, c'est-à-dire un élément du membre de gauche. Nous omettons la vérification que ces constructions sont mutuellement inverses. 
\end{proof}

\begin{lemma}
\label{spaces-pushouts-lemma-diagram}
Soit $S$ un schéma. Soit
$$
\xymatrix{
A \ar[r] \ar[d] & C \ar[d] \ar[r] & E \ar[d] \\
B \ar[r] & D \ar[r] & F
}
$$
un diagramme commutatif d'espaces algébriques sur $S$.
Supposons que $A, B, C, D$ et $A, B, E, F$ forment des carrés cartésiens
et que $B \to D$ soit \'etale surjectif.
Alors $C, D, E, F$ est un carré cartésien.
\end{lemma}

\begin{proof}
Ceci est formel.
\end{proof}

\begin{lemma}
\label{spaces-pushouts-lemma-equivalence-categories-spaces-over-pushout}
Dans le cas du Lemme \ref{spaces-pushouts-lemma-categories-spaces-over-pushout}
le foncteur $F \circ G$ est isomorphe au foncteur identité.
\end{lemma}

\begin{proof}
Nous allons prouver que $F \circ G$ est isomorphe à l'identité en
réduisant cela à l'énoncé correspondant de
Compléments sur les morphismes, Lemme
\ref{more-morphisms-lemma-equivalence-categories-schemes-over-pushout}.

\medskip\noindent
Choisissons un schéma $Y_1$ et un morphisme \'etale surjectif $Y_1 \to Y$. Posons $X_1 = Y_1 \times_Y X$. C'est un schéma affine sur $Y_1$ muni d'un morphisme \'etale surjectif $X_1 \to X$. D'après Compléments sur les morphismes d'espaces, Lemme \ref{spaces-more-morphisms-lemma-thickening-equivalence}, il existe un morphisme \'etale surjectif $X'_1 \to X'$ tel que $X_1 = X_1' \times_{X'} X$. En particulier, le morphisme de schémas $X_1 \to X_1'$ est également un épaississement. Appliquons Compléments sur les morphismes, Lemme \ref{more-morphisms-lemma-pushout-along-thickening}, pour obtenir une somme amalgamée $Y_1' = Y_1 \amalg_{X_1} X_1'$ dans la catégorie des schémas. Dans la preuve du Lemme \ref{spaces-pushouts-lemma-pushout-along-thickening}, nous avons construit $Y'$ comme quotient d'une relation d'équivalence \'etale sur $Y_1'$, d'où un diagramme commutatif 
\begin{equation}
\label{spaces-pushouts-equation-cube}
\vcenter{
\xymatrix{
& X \ar[rr] \ar'[d][dd] & & X' \ar[dd] \\
X_1 \ar[rr] \ar[dd] \ar[ru] & & X_1' \ar[dd] \ar[ru] & \\
& Y \ar'[r][rr] & & Y' \\
Y_1 \ar[rr] \ar[ru] & & Y_1' \ar[ru]
}
}
\end{equation} où tous les carrés, sauf les carrés avant et arrière, sont cartésiens (les carrés avant et arrière sont des sommes amalgamées), et où les flèches vers le nord-est
sont \'etales et surjectives. Notons $F_1$, $G_1$ les
foncteurs construits dans
Compléments sur les morphismes, Lemme
\ref{more-morphisms-lemma-equivalence-categories-schemes-over-pushout}
pour la face avant. Alors le diagramme de catégories
$$
\xymatrix{
(\Sch/Y_1') \ar@<-1ex>[r]_-{F_1} \ar[d] &
(\Sch/Y_1) \times_{(\Sch/Y_1')} (\Sch/X_1') \ar[d] \ar@<-1ex>[l]_-{G_1} \\
(\textit{Espaces}/Y') \ar@<-1ex>[r]_-F &
(\textit{Espaces}/Y) \times_{(\textit{Espaces}/Y')} (\textit{Espaces}/X')
\ar@<-1ex>[l]_-G
}
$$
est commutatif par des considérations simples sur les foncteurs de changement de base
et parce que les sommes amalgamées dans les schémas coïncident avec les sommes amalgamées dans
les espaces du Lemme \ref{spaces-pushouts-lemma-pushout-along-thickening-schemes}.

\medskip\noindent
Soit $(V, U', \varphi)$ un objet de
$(\textit{Espaces}/Y) \times_{(\textit{Espaces}/Y')} (\textit{Espaces}/X')$.
Notons $U = U' \times_{X'} X$ de sorte que $G(V, U', \varphi) = V \amalg_U U'$.
Choisissons un schéma $V_1$ et un morphisme \'etale surjectif
$V_1 \to Y_1 \times_Y V$. Posons $U_1 = V_1 \times_Y X$. Alors
$$
U_1 = V_1 \times_Y X
\longrightarrow
(Y_1 \times_Y V) \times_Y X =
X_1 \times_Y V = X_1 \times_X X \times_Y V = X_1 \times_X U
$$
est lui aussi \'etale et surjectif. Par
Compléments sur les morphismes d'espaces, Lemme
\ref{spaces-more-morphisms-lemma-thickening-equivalence}
il existe un épaississement $U_1 \to U_1'$ et un morphisme \'etale surjectif
$U_1' \to X_1' \times_{X'} U'$ dont le changement de base à $X_1 \times_X U$ est le
morphisme affiché. À ce stade $(V_1, U'_1, \varphi_1)$ est un objet de
$(\Sch/Y_1) \times_{(\Sch/Y_1')} (\Sch/X_1')$. Dans la preuve du
Lemme \ref{spaces-pushouts-lemma-pushout-along-thickening} nous avons construit
$G(V, U', \varphi) = V \amalg_U U'$ comme un quotient d'une relation
d'équivalence \'etale sur $G_1(V_1, U_1', \varphi_1) = V_1 \amalg_{U_1} U_1'$
de sorte que nous obtenons un diagramme commutatif
\begin{equation}
\label{spaces-pushouts-equation-cube-over}
\vcenter{
\xymatrix{
& U \ar[rr] \ar'[d][dd] & & U' \ar[dd] \\
U_1 \ar[rr] \ar[dd] \ar[ru] & & U_1' \ar[dd] \ar[ru] & \\
& V \ar'[r][rr] & & G(V, U', \varphi) \\
V_1 \ar[rr] \ar[ru] & & G_1(V_1, U_1', \varphi_1) \ar[ru]
}
}
\end{equation}
où tous les carrés sauf les carrés avant et arrière sont cartésiens (les carrés avant et arrière sont des sommes amalgamées), et où les flèches vers le nord-est sont \'etales et surjectives. En particulier, $$
G_1(V_1, U_1', \varphi_1) \to G(V, U', \varphi)
$$ est \'etale et surjectif.

\medskip\noindent
Enfin, démontrons le lemme. Nous devons montrer que le morphisme d'adjonction $(V, U', \varphi) \to F(G(V, U', \varphi))$ est un isomorphisme. Nous savons que $(V_1, U_1', \varphi_1) \to F_1(G_1(V_1, U_1', \varphi_1))$ est un isomorphisme d'après Compléments sur les morphismes, Lemme \ref{more-morphisms-lemma-equivalence-categories-schemes-over-pushout}. Rappelons que $F$ et $F_1$ sont donnés par changement de base. En utilisant les propriétés de (\ref{spaces-pushouts-equation-cube-over}) et le Lemme \ref{spaces-pushouts-lemma-diagram}, nous voyons que $V \to G(V, U', \varphi) \times_{Y'} Y$ et $U' \to G(V, U', \varphi) \times_{Y'} X'$ sont des isomorphismes, c'est-à-dire que $(V, U', \varphi) \to F(G(V, U', \varphi))$ est un isomorphisme. 
\end{proof}

\begin{lemma}
\label{spaces-pushouts-lemma-space-over-pushout-flat-modules}
Soit $S$ un schéma de base.
Soit $X \to X'$ un épaississement d'espaces algébriques sur $S$
et soit $X \to Y$ un morphisme affine d'espaces algébriques sur $S$.
Soit $Y' = Y \amalg_X X'$ la somme amalgamée
(voir le Lemme \ref{spaces-pushouts-lemma-pushout-along-thickening}).
Soit $V' \to Y'$ un morphisme d'espaces algébriques sur $S$. Posons
$V = Y \times_{Y'} V'$, $U' = X' \times_{Y'} V'$, et $U = X \times_{Y'} V'$.
Il existe une équivalence de catégories entre
\begin{enumerate}
\item $\mathcal{O}_{V'}$-modules quasi-cohérents plats sur $Y'$, et
\item la catégorie des triplets $(\mathcal{G}, \mathcal{F}', \varphi)$ où
\begin{enumerate}
\item $\mathcal{G}$ est un $\mathcal{O}_V$-module quasi-cohérent plat sur $Y$,
\item $\mathcal{F}'$ est un $\mathcal{O}_{U'}$-module quasi-cohérent plat
sur $X'$, et
\item $\varphi : (U \to V)^*\mathcal{G} \to (U \to U')^*\mathcal{F}'$ est un isomorphisme de $\mathcal{O}_U$-modules.
\end{enumerate}
\end{enumerate}
L'équivalence envoie $\mathcal{G}'$ sur
$((V \to V')^*\mathcal{G}', (U' \to V')^*\mathcal{G}', can)$.
Supposons que $\mathcal{G}'$ corresponde au triple
$(\mathcal{G}, \mathcal{F}', \varphi)$. Alors
\begin{enumerate}
\item[(a)] $\mathcal{G}'$ est un $\mathcal{O}_{V'}$-module de type fini si et seulement si $\mathcal{G}$ et $\mathcal{F}'$ sont des $\mathcal{O}_V$- et $\mathcal{O}_{U'}$-modules de type fini.
\item[(b)] si $V' \to Y'$ est localement de présentation finie, alors $\mathcal{G}'$ est un $\mathcal{O}_{V'}$-module de présentation finie si et seulement si $\mathcal{G}$ et $\mathcal{F}'$ sont des $\mathcal{O}_V$- et $\mathcal{O}_{U'}$-modules de présentation finie.
\end{enumerate}
\end{lemma}

\begin{proof}
Un foncteur quasi-inverse associe au triple
$(\mathcal{G}, \mathcal{F}', \varphi)$ le produit fibré
$$
(V \to V')_*\mathcal{G}
\times_{(U \to V')_*\mathcal{F}}
(U' \to V')_*\mathcal{F}'
$$
où $\mathcal{F} = (U \to U')^*\mathcal{F}'$. Cette construction convient car, sur
les objets affines \'etales sur $V'$ et sur $Y'$, nous retrouvons l'équivalence de
Compléments d'algèbre, Lemme
\ref{more-algebra-lemma-relative-flat-module-over-fibre-product}.
Détails omis.

\medskip\noindent
Les parties (a) et (b) se réduisent par la localisation \'etale
(Propriétés des espaces, section
\ref{spaces-properties-section-properties-modules})
au cas où $V'$ et $Y'$ sont affines, auquel cas le résultat
suit de
Compléments d'algèbre, Lemmes
\ref{more-algebra-lemma-relative-finite-module-over-fibre-product} et
\ref{more-algebra-lemma-relative-finitely-presented-module-over-fibre-product}.
\end{proof}


\begin{lemma}
\label{spaces-pushouts-lemma-equivalence-categories-spaces-pushout-flat}
Dans la situation du
Lemme \ref{spaces-pushouts-lemma-equivalence-categories-spaces-over-pushout},
Si $V' = G(V, U', \varphi)$ pour un certain triple $(V, U', \varphi)$, alors
\begin{enumerate}
\item $V' \to Y'$ est localement de type fini si et seulement si $V \to Y$ et
$U' \to X'$ sont localement de type fini,
\item $V' \to Y'$ est plat si et seulement si $V \to Y$ et $U' \to X'$ sont plats,
\item $V' \to Y'$ est plat et localement de présentation finie si et seulement si
$V \to Y$ et $U' \to X'$ sont plats et localement de présentation finie,
\item $V' \to Y'$ est lisse si et seulement si $V \to Y$ et $U' \to X'$ sont lisses,
\item $V' \to Y'$ est \'etale si et seulement si $V \to Y$ et $U' \to X'$
sont \'etales, et
\item ajouter plus ici si nécessaire.
\end{enumerate}
Si $W'$ est plat sur $Y'$, alors le morphisme d'adjonction $G(F(W')) \to W'$ est un isomorphisme. Par conséquent, $F$ et $G$ définissent des foncteurs quasi-inverses mutuels entre la catégorie des espaces algébriques plats sur $Y'$ et la catégorie des triples $(V, U', \varphi)$ avec $V \to Y$ et $U' \to X'$ plats. 
\end{lemma}

\begin{proof}
Choisissons un diagramme (\ref{spaces-pushouts-equation-cube}) comme dans la démonstration du
Lemme \ref{spaces-pushouts-lemma-equivalence-categories-spaces-over-pushout}.

\medskip\noindent
Preuve de (1) -- (5). Soit $(V, U', \varphi)$ un objet de
$(\textit{Espaces}/Y) \times_{(\textit{Espaces}/Y')} (\textit{Espaces}/X')$.
Construisons un diagramme (\ref{spaces-pushouts-equation-cube-over}) comme dans la preuve du
Lemme \ref{spaces-pushouts-lemma-equivalence-categories-spaces-over-pushout}.
Alors le changement de base de $G(V, U', \varphi) \to Y'$ à
$Y'_1$ est $G_1(V_1, U_1', \varphi_1) \to Y_1'$. Par conséquent, (1) -- (5)
suivent immédiatement des énoncés correspondants de Compléments sur les morphismes, Lemme
\ref{more-morphisms-lemma-equivalence-categories-schemes-over-pushout-flat}
pour les schémas.

\medskip\noindent
Supposons que $W' \to Y'$ soit plat. Choisissons un schéma $W'_1$ et un morphisme \'etale surjectif $W'_1 \to Y_1' \times_{Y'} W'$. Observons que $W'_1 \to W'$ est \'etale et surjectif comme composée de morphismes \'etales et surjectifs. Nous savons que $G_1(F_1(W_1')) \to W_1'$ est un isomorphisme par Compléments sur les morphismes, Lemme \ref{more-morphisms-lemma-equivalence-categories-schemes-over-pushout-flat}, appliqué à $W'_1$ sur $Y'_1$ et à la face avant du diagramme (avec les foncteurs $G_1$ et $F_1$ comme dans la preuve du Lemme \ref{spaces-pushouts-lemma-equivalence-categories-spaces-over-pushout}).
Alors la construction de $G(F(W'))$ (comme somme amalgamée, c'est-à-dire par la construction du Lemme \ref{spaces-pushouts-lemma-pushout-along-thickening}) montre que $G_1(F_1(W'_1)) \to G(F(W'))$ est \'etale et surjectif. Nous en concluons que $G(F(W')) \to W'$ est \'etale, voir par exemple Propriétés des espaces, Lemme \ref{spaces-properties-lemma-etale-local}.
Mais $G(F(W')) \to W'$ induit un isomorphisme entre les espaces algébriques réduits associés (par construction), donc c'est un isomorphisme.
\end{proof}







\section{Sommes amalgamées le long d'immersions fermées et de morphismes entiers}
\label{spaces-pushouts-section-pushouts-II}

\noindent
Cette section est l'analogue de
Compléments sur les morphismes, section \ref{more-morphisms-section-pushouts-II}.

\begin{lemma}
\label{spaces-pushouts-lemma-pushout-along-closed-immersion-and-integral}
Dans Compléments sur les morphismes, Situation
\ref{more-morphisms-situation-pushout-along-closed-immersion-and-integral}
soit $Y \amalg_Z X$ la somme amalgamée dans la catégorie des schémas
(Compléments sur les morphismes, Proposition
\ref{more-morphisms-proposition-pushout-along-closed-immersion-and-integral}).
Alors $Y \amalg_Z X$
est également une somme amalgamée dans la catégorie des espaces algébriques sur $S$.
\end{lemma}

\begin{proof}
Ceci est une conséquence du Lemme \ref{spaces-pushouts-lemma-colimit-agrees}, de la proposition mentionnée dans le lemme, ainsi que de Compléments sur les morphismes, Lemmes \ref{more-morphisms-lemma-pushout-functor} et \ref{more-morphisms-lemma-pushout-functor-equivalence-flat}.
Les conditions (1) et (2) du Lemme \ref{spaces-pushouts-lemma-colimit-agrees} suivent immédiatement. Pour voir (3) et (4), remarquons qu'un morphisme \'etale est localement quasi-fini et utilisons le fait que l'équivalence de catégories de Compléments sur les morphismes, Lemme \ref{more-morphisms-lemma-pushout-functor-equivalence-flat}, est construite à l'aide de la construction de la somme amalgamée de Compléments sur les morphismes, Lemme \ref{more-morphisms-lemma-pushout-functor}.
Détails mineurs omis.
\end{proof}








\section{Sommes amalgamées et catégories dérivées}
\label{spaces-pushouts-section-pushouts-derived}

\noindent
Dans cette section, nous discutons du comportement de la catégorie dérivée des modules lors des sommes amalgamées.

\begin{lemma}
\label{spaces-pushouts-lemma-pushout-along-thickening-derived}
Soit $S$ un schéma. Considérons une somme amalgamée
$$
\xymatrix{
X \ar[r]_i \ar[d]_f & X' \ar[d]^{f'}
\\
Y \ar[r]^j & Y'
}
$$
dans la catégorie des espaces algébriques sur $S$
comme dans le Lemme \ref{spaces-pushouts-lemma-pushout-along-thickening}.
Supposons que $i$ soit un épaississement. Alors l'image essentielle du foncteur\footnote{Tous les foncteurs sont donnés par l'image inverse dérivée.}
$$
D(\mathcal{O}_{Y'}) \longrightarrow
D(\mathcal{O}_Y) \times_{D(\mathcal{O}_X)} D(\mathcal{O}_{X'})
$$
contient tout triplet $(M, K', \alpha)$ où $M \in D(\mathcal{O}_Y)$
et $K' \in D(\mathcal{O}_{X'})$ sont pseudo-cohérents.
\end{lemma}

\begin{proof}
Soit $(M, K', \alpha)$ un objet de la cible du foncteur
du lemme. Ici $\alpha : Lf^*M \to Li^*K'$
est un isomorphisme adjoint à un morphisme $\beta : M \to Rf_*Li^*K'$.
Ainsi nous obtenons des morphismes
$$
Rj_*M \xrightarrow{Rj_*\beta}
Rj_*Rf_*Li^*K' = Rf'_*Ri_*Li^*K' \leftarrow Rf'_*K'
$$
où la flèche pointant vers la gauche provient de $K' \to Ri_*Li^*K'$.
Choisissons un triangle distingué
$$
M' \to Rj_*M \oplus Rf'_*K' \to Rj_*Rf_*Li^*K' \to M'[1]
$$
dans $D(\mathcal{O}_{Y'})$. La première flèche définit des morphismes canoniques
$Lj^*M' \to M$ et $L(f')^*M' \to K'$ compatibles avec $\alpha$.
Ainsi, il suffit de montrer que les morphismes
$Lj^*M' \to M$ et $L(f')^*M' \to K'$ sont des isomorphismes.
Cette assertion se vérifie localement pour la topologie \'etale sur $Y'$, donc nous pouvons
supposer que $Y'$ est affine.

\medskip\noindent
Supposons que $Y'$ soit affine et que $M \in D(\mathcal{O}_Y)$
et $K' \in D(\mathcal{O}_{X'})$ soient pseudo-cohérents.
Supposons que notre somme amalgamée corresponde au produit fibré
$$
\xymatrix{
B & B' \ar[l] \\
A \ar[u] & A' \ar[l] \ar[u]
}
$$
d'anneaux où $B' \to B$ est surjectif avec un noyau localement nilpotent $I$
(et donc $A' \to A$ est également surjectif avec un noyau localement nilpotent $I$).
Les hypothèses sur $M$ et $K'$ impliquent que $M$ provient d'un objet pseudo-cohérent de $D(A)$ et que $K'$ provient d'un objet pseudo-cohérent de $D(B')$, voir
Catégories dérivées des espaces, Lemmes
\ref{spaces-perfect-lemma-pseudo-coherent},
\ref{spaces-perfect-lemma-derived-quasi-coherent-small-etale-site}, et
\ref{spaces-perfect-lemma-descend-pseudo-coherent}
et
Catégories dérivées des schémas, Lemme
\ref{perfect-lemma-affine-compare-bounded} et
\ref{perfect-lemma-pseudo-coherent-affine}.
De plus, les foncteurs image directe et image inverse dérivée coïncident avec les
opérations correspondantes sur les catégories dérivées de modules, voir Catégories dérivées des espaces, Remarque
\ref{spaces-perfect-remark-match-total-direct-images}
et
Catégories dérivées des schémas, Lemmes
\ref{perfect-lemma-quasi-coherence-pushforward} et
\ref{perfect-lemma-quasi-coherence-pullback}.
Cela nous ramène à l'énoncé formulé dans le paragraphe suivant.
(Pour vérifier que ces références montrent que l'objet $M'$ appartient à
$D_\QCoh(\mathcal{O}_{Y'})$, on utilise le fait qu'il s'agit d'une
sous-catégorie triangulée de $D(\mathcal{O}_{Y'})$.)

\medskip\noindent
Étant donné un diagramme d'anneaux comme ci-dessus et un triple
$(M, K', \alpha)$ où $M \in D(A)$, $K' \in D(B')$ sont
pseudo-cohérents et
$\alpha : M \otimes_A^\mathbf{L} B \to K' \otimes_{B'}^\mathbf{L} B$
est un isomorphisme, supposons que nous ayons un triangle distingué
$$
M' \to M \oplus K' \to K' \otimes_{B'}^\mathbf{L} B \to M'[1]
$$
dans $D(A')$. Il s'agit de montrer que les morphismes induits
$M' \otimes_{A'}^\mathbf{L} A \to M$ et
$M' \otimes_{A'}^\mathbf{L} B' \to K'$ sont des isomorphismes.
Pour ce faire, choisissons un complexe borné supérieurement
$E^\bullet$ de $A$-modules libres finis représentant $M$.
Puisque $(B', I)$ est une paire hensélienne
(Compléments d'algèbre, Lemme \ref{more-algebra-lemma-locally-nilpotent-henselian})
avec $B = B'/I$, nous pouvons appliquer Compléments d'algèbre, Lemme
\ref{more-algebra-lemma-lift-complex-finite-projectives}
pour voir qu'il existe un complexe borné supérieurement $P^\bullet$
de $B'$-modules libres tel que $\alpha$ soit représenté
par un isomorphisme $E^\bullet \otimes_A B \cong P^\bullet \otimes_{B'} B$.
Ensuite, nous pouvons considérer la suite exacte courte
$$
0 \to L^\bullet \to
E^\bullet \oplus P^\bullet \to P^\bullet \otimes_{B'} B \to 0
$$
de complexes de $A'$-modules.  
Compléments d'algèbre, Lemme  
\ref{more-algebra-lemma-finitely-presented-module-over-fibre-product}  
implique que $L^\bullet$ est un complexe borné supérieurement de  
$A'$-modules projectifs finis  
(en fait, il est assez facile de montrer directement que $L^n$ est libre fini  
dans notre cas) et que nous avons  
$L^\bullet \otimes_{A'} A = E^\bullet$ et  
$L^\bullet \otimes_{A'} B' = P^\bullet$.  
La suite exacte courte donne un triangle distingué  
$$
L^\bullet \to M \oplus K' \to K' \otimes_{B'}^\mathbf{L} B \to (L^\bullet)[1]
$$  
dans $D(A')$ (Catégories dérivées, section  
\ref{derived-section-canonical-delta-functor}) qui est isomorphe  
au triangle distingué donné par les propriétés générales des  
catégories triangulées (Catégories dérivées, section  
\ref{derived-section-elementary-results}). En d'autres termes, $L^\bullet$  
représente $M'$ de manière compatible avec les morphismes donnés. Ainsi, les morphismes  
$M' \otimes_{A'}^\mathbf{L} A \to M$ et  
$M' \otimes_{A'}^\mathbf{L} B' \to K'$ sont  
des isomorphismes parce que nous venons de voir que l'assertion correspondante est vraie pour $L^\bullet$.  
\end{proof}







\section{Construction de carrés distingués élémentaires}
\label{spaces-pushouts-section-elementary-dsitnguished-squares}

\noindent
Les carrés distingués élémentaires ont été définis dans
Catégories dérivées des espaces, section \ref{spaces-perfect-section-induction}.

\begin{lemma}
\label{spaces-pushouts-lemma-elementary-distinguished-square-pushout}
Soit $S$ un schéma. Soit $(U \subset W, f : V \to W)$ un carré distingué élémentaire. Alors
$$
\xymatrix{
U \times_W V \ar[r] \ar[d] &
V \ar[d]^f \\
U \ar[r] & W
}
$$
est un carré cocartésien dans la catégorie des espaces algébriques sur $S$.
\end{lemma}

\begin{proof}
Observons que $U \amalg V \to W$ est un morphisme \'etale surjectif.
Le produit fibré
$$
(U \amalg V) \times_W (U \amalg V)
$$
est la réunion disjointe de quatre composantes, à savoir
$U = U \times_W U$, $U \times_W V$, $V \times_W U$,
et $V \times_W V$.
Il existe un morphisme \'etale surjectif
$$
V \amalg (U \times_W V) \times_U (U \times_W V) \longrightarrow V \times_W V
$$
car $f$ induit un isomorphisme sur $W \setminus U$
(faisant partie de la définition d'un carré distingué élémentaire).
Soit $B$ un espace algébrique sur $S$ et soit
$g : V \to B$ et $h : U \to B$ des morphismes sur
$S$ qui coïncident après restriction à $U \times_W V$.
Alors la description de
$(U \amalg V) \times_W (U \amalg V)$ donnée ci-dessus
montre que $h \amalg g : U \amalg V \to B$
égalise les deux projections. Comme $B$ est un faisceau
pour la topologie \'etale, nous obtenons une factorisation unique
de $h \amalg g$ par $W$ comme voulu.
\end{proof}

\begin{lemma}
\label{spaces-pushouts-lemma-construct-elementary-distinguished-square}
Soit $S$ un schéma. Soient $V$, $U$ des espaces algébriques sur $S$.
Soit $V' \subset V$ un sous-espace ouvert et soit $f' : V' \to U$ un
morphisme \'etale séparé d'espaces algébriques sur $S$.
Alors il existe une somme amalgamée
$$
\xymatrix{
V' \ar[r] \ar[d] &
V \ar[d]^f \\
U \ar[r] & W
}
$$
dans la catégorie des espaces algébriques sur $S$ et de plus
$(U \subset W, f : V \to W)$ est un carré distingué élémentaire.
\end{lemma}

\begin{proof}
Nous allons construire $W$ comme le quotient d'une relation d'équivalence \'etale $R$ sur $U \amalg V$. Un tel quotient est un espace algébrique, par exemple par Amorçage, Théorème \ref{bootstrap-theorem-final-bootstrap}.
De plus, la démonstration du Lemme \ref{spaces-pushouts-lemma-elementary-distinguished-square-pushout} nous indique de prendre $$
R = U \amalg V' \amalg V' \amalg V \amalg
(V' \times_U V' \setminus \Delta_{V'/U}(V'))
$$.
Puisque nous avons supposé que $V' \to U$ est séparé, l'image de $\Delta_{V'/U}$ est fermée et donc le complémentaire est un sous-espace ouvert. Le morphisme $j : R \to (U \amalg V) \times_S (U \amalg V)$ est donné par $$
u,\ v',\ v',\ v,\ (v'_1, v'_2) \mapsto
(u, u),\ (f'(v'), v'),\ (v', f'(v')),\ (v, v),\ (v'_1, v'_2)
$$ avec une notation évidente. On vérifie immédiatement qu'il s'agit d'un monomorphisme et d'une relation d'équivalence, et que les morphismes induits $s, t : R \to U \amalg V$ sont \'etales. Soit $W = (U \amalg V)/R$ l'espace algébrique quotient.
Nous obtenons un diagramme commutatif comme dans l'énoncé du lemme.
Pour terminer la démonstration, il suffit de montrer que ce diagramme est un carré distingué élémentaire, car alors le
Lemme \ref{spaces-pushouts-lemma-elementary-distinguished-square-pushout} implique qu'il s'agit d'une somme amalgamée.  
Ainsi, nous devons montrer que $U \to W$ est ouvert et que  
$f$ est \'etale et induit un isomorphisme sur $W \setminus U$.  
Cela découle du choix de $R$ ; nous omettons les détails.  
\end{proof}






\section{Recollement formel de modules quasi-cohérents}
\label{spaces-pushouts-section-formal-glueing}

\noindent
Cette section est l'analogue de
Compléments d'algèbre, section \ref{more-algebra-section-formal-glueing}.
Dans le cas des morphismes de schémas, le résultat se trouve dans
l'article de Joyet \cite{Joyet} ; c'est un bon point de départ pour la lecture.
Pour une discussion des applications aux problèmes de descente pour les champs, voir
l'article de Moret-Bailly \cite{MB}. Dans le cas d'un morphisme affine de schémas, on trouve un énoncé dans l'appendice de l'article
\cite{Ferrand-Raynaud}, mais il est nécessaire d'ajouter l'hypothèse
que le sous-schéma fermé est défini par un idéal de type fini (comme dans l'article de Joyet), sans quoi l'énoncé est faux.
Une généralisation de ce matériel aux catégories dérivées supérieures, avec d'éventuelles applications à des situations non plates,
se trouve dans \cite[section 5]{Bhatt-Algebraize}.

\medskip\noindent
Nous commençons par un lemme sur les faisceaux abéliens supportés sur des sous-ensembles fermés.

\begin{lemma}
\label{spaces-pushouts-lemma-stalk-pushforward-with-support}
Soit $S$ un schéma. Soit $f : Y \to X$ un morphisme d'espaces algébriques
sur $S$. Soit $Z \subset X$ un sous-espace fermé tel que $f^{-1}Z \to Z$ soit
intégral et universellement injectif. Soit $\overline{y}$ un point géométrique
de $Y$ et $\overline{x} = f(\overline{y})$. Nous avons
$$
(Rf_*Q)_{\overline{x}} = Q_{\overline{y}}
$$
dans $D(\textit{Ab})$ pour tout objet $Q$ de $D(Y_\etale)$ supporté
sur $|f^{-1}Z|$.
\end{lemma}

\begin{proof}
Considérons le diagramme commutatif des espaces algébriques
$$
\xymatrix{
f^{-1}Z \ar[r]_{i'} \ar[d]_{f'} & Y \ar[d]_f \\
Z \ar[r]^i & X
}
$$
D'après Cohomologie des espaces, Lemme
\ref{spaces-cohomology-lemma-complexes-with-support-on-closed} nous pouvons écrire
$Q = Ri'_*K'$ pour un certain objet $K'$ de $D(f^{-1}Z_\etale)$.
D'après Morphismes d'espaces, Lemme
\ref{spaces-morphisms-lemma-integral-universally-injective-push-pull}
nous avons $K' = (f')^{-1}K$ avec $K = Rf'_*K'$.
Alors nous avons $Rf_*Q = Rf_*Ri'_*K' = Ri_*Rf'_*K' = Ri_*K$.
Soit $\overline{z}$ le point géométrique de $Z$ correspondant à $\overline{x}$ et soit $\overline{z}'$ le point géométrique de $f^{-1}Z$ correspondant à $\overline{y}$. Nous obtenons alors le résultat du lemme comme suit
$$
Q_{\overline{y}} = (Ri'_*K')_{\overline{y}} = K'_{\overline{z}'} =
(f')^{-1}K_{\overline{z}'} = K_{\overline{z}} = Ri_*K_{\overline{x}} =
Rf_*Q_{\overline{x}}
$$
L'égalité du milieu tient grâce à la description de la fibre d'une image inverse donnée dans Propriétés des espaces, Lemme \ref{spaces-properties-lemma-stalk-pullback}.
\end{proof}

\begin{lemma}
\label{spaces-pushouts-lemma-stalk-formal-glueing}
Soit $S$ un schéma. Soit $f : Y \to X$ un morphisme d'espaces algébriques sur $S$. Soit $Z \subset X$ un sous-espace fermé tel que $f^{-1}Z \to Z$ soit intégral et universellement injectif. Soit $\overline{y}$ un point géométrique de $Y$ et $\overline{x} = f(\overline{y})$. Soit $\mathcal{G}$ un faisceau abélien sur $Y$. Alors le morphisme entre les complexes à deux termes $$
\left(f_*\mathcal{G}_{\overline{x}} \to
(f \circ j')_*(\mathcal{G}|_V)_{\overline{x}}\right)
\longrightarrow
\left(\mathcal{G}_{\overline{y}} \to j'_*(\mathcal{G}|_V)_{\overline{y}}\right)
$$ induit un isomorphisme sur les noyaux et une injection sur les conoyaux. Ici $V = Y \setminus f^{-1}Z$ et $j' : V \to Y$ est l'inclusion.
\end{lemma}

\begin{proof}
Choisissons un triangle distingué
$$
\mathcal{G} \to Rj'_*\mathcal{G}|_V \to Q \to \mathcal{G}[1]
$$
dans $D(Y_\etale)$. Les faisceaux de cohomologie de $Q$
sont supportés sur $|f^{-1}Z|$. Nous appliquons $Rf_*$ et nous obtenons
$$
Rf_*\mathcal{G} \to Rf_*Rj'_*\mathcal{G}|_V \to Rf_*Q
\to Rf_*\mathcal{G}[1]
$$
En prenant les fibres en $\overline{x}$ nous obtenons une suite exacte
$$
0 \to
(R^{-1}f_*Q)_{\overline{x}} \to
f_*\mathcal{G}_{\overline{x}} \to
(f \circ j')_*(\mathcal{G}|_V)_{\overline{x}} \to
(R^0f_*Q)_{\overline{x}}
$$
Nous pouvons comparer ceci avec la suite exacte
$$
0 \to
H^{-1}(Q)_{\overline{y}} \to
\mathcal{G}_{\overline{y}} \to
j'_*(\mathcal{G}|_V)_{\overline{y}} \to
H^0(Q)_{\overline{y}}
$$
Ainsi, nous voyons que le lemme en découle car
$Q_{\overline{y}} = Rf_*Q_{\overline{x}}$ par
Lemme \ref{spaces-pushouts-lemma-stalk-pushforward-with-support}.
\end{proof}

\begin{lemma}
\label{spaces-pushouts-lemma-stalk-of-pushforward}
Soit $S$ un schéma. Soit $X$ un espace algébrique sur $S$.
Soit $f : Y \to X$ un morphisme quasi-compact et quasi-séparé.
Soit $\overline{x}$ un point géométrique de $X$ et soit
$\Spec(\mathcal{O}_{X, \overline{x}}) \to X$
le morphisme canonique. Pour un module quasi-cohérent
$\mathcal{G}$ sur $Y$, nous avons
$$
f_*\mathcal{G}_{\overline{x}} =
\Gamma(Y \times_X \Spec(\mathcal{O}_{X, \overline{x}}), p^*\mathcal{G})
$$
où $p : Y \times_X \Spec(\mathcal{O}_{X, \overline{x}}) \to Y$
est la projection.
\end{lemma}

\begin{proof}
Observons que $f_*\mathcal{G}_{\overline{x}} =
\Gamma(\Spec(\mathcal{O}_{X, \overline{x}}), h^*f_*\mathcal{G})$
où $h : \Spec(\mathcal{O}_{X, \overline{x}}) \to X$.
Par conséquent, le résultat est vrai car $h$ est plat, si bien que
Cohomologie des espaces, Lemme
\ref{spaces-cohomology-lemma-flat-base-change-cohomology}
s'applique.
\end{proof}

\begin{lemma}
\label{spaces-pushouts-lemma-stalk-of-module-with-support}
Soit $S$ un schéma. Soit $X$ un espace algébrique sur $S$.
Soit $i : Z \to X$ une immersion fermée de présentation finie.
Soit $Q \in D_\QCoh(\mathcal{O}_X)$ supporté sur $|Z|$.
Soit $\overline{x}$ un point géométrique de $X$ et soit
$I_{\overline{x}} \subset \mathcal{O}_{X, \overline{x}}$ le germe du
faisceau d'idéaux de $Z$. Alors les modules de cohomologie
$H^n(Q_{\overline{x}})$ sont constitués d'éléments annulés par une puissance de $I_{\overline{x}}$
(voir Compléments d'algèbre, Définition
\ref{more-algebra-definition-f-power-torsion}).
\end{lemma}

\begin{proof}
Choisissons un schéma affine $U$ et un morphisme \'etale $U \to X$ tel que $\overline{x}$ se relève en un point géométrique $\overline{u}$ de $U$. Ensuite, nous pouvons remplacer $X$ par $U$, $Z$ par $U \times_X Z$, $Q$ par la restriction $Q|_U$, et $\overline{x}$ par $\overline{u}$. Ainsi, nous pouvons supposer que $X = \Spec(A)$ est affine. Soit $I \subset A$ l'idéal définissant $Z$. Puisque $i : Z \to X$ est de présentation finie, l'idéal $I = (f_1, \ldots, f_r)$ est de type fini. L'objet $Q$ provient d'un complexe de $A$-modules $M^\bullet$, voir Catégories dérivées des espaces, Lemme \ref{spaces-perfect-lemma-derived-quasi-coherent-small-etale-site} et Catégories dérivées des schémas, Lemme \ref{perfect-lemma-affine-compare-bounded}. Puisque les faisceaux de cohomologie de $Q$ sont supportés sur $Z$, nous voyons que la localisation $M^\bullet_f$ est acyclique pour chaque $f \in I$. Prenons $x \in H^p(M^\bullet)$. D'après ce qui précède, nous pouvons trouver un entier $n_i$ tel que $f_i^{n_i} x = 0$ dans $H^p(M^\bullet)$ pour chaque $i$. Ensuite, en posant $n = \sum n_i$, nous voyons que $I^n$ annule $x$. Ainsi, tout élément de $H^p(M^\bullet)$ est annulé par une puissance de $I$. Comme l'homomorphisme d'anneaux $A \to \mathcal{O}_{X, \overline{x}}$ est plat et comme $I_{\overline{x}} = I\mathcal{O}_{X, \overline{x}}$, nous concluons. 
\end{proof}

\begin{lemma}
\label{spaces-pushouts-lemma-formal-glueing-on-closed}
Soit $S$ un schéma. Soit $f : Y \to X$ un morphisme d'espaces algébriques sur $S$. Soit $Z \subset X$ un sous-espace fermé. Supposons que $f^{-1}Z \to Z$ soit un isomorphisme et que $f$ soit plat en chaque point de $f^{-1}Z$. Pour tout $Q$ dans $D_\QCoh(\mathcal{O}_Y)$ supporté sur $|f^{-1}Z|$, nous avons $Lf^*Rf_*Q = Q$.
\end{lemma}

\begin{proof}
Nous montrons que le morphisme canonique $Lf^*Rf_*Q \to Q$ est un isomorphisme
en vérifiant sur les germes en $\overline{y}$. Si $\overline{y}$ n'est pas
dans $f^{-1}Z$, alors les deux côtés sont nuls et le résultat est vrai.
Supposons que l'image $\overline{x}$ de $\overline{y}$ soit dans $Z$.
D'après le Lemme \ref{spaces-pushouts-lemma-stalk-pushforward-with-support} nous avons
$Rf_*Q_{\overline{x}} = Q_{\overline{y}}$ et puisque $f$ est plat
en $\overline{y}$, nous voyons que
$$
(Lf^*Rf_*Q)_{\overline{y}} =
(Rf_*Q)_{\overline{x}}
\otimes_{\mathcal{O}_{X, \overline{x}}}
\mathcal{O}_{Y, \overline{y}} =
Q_{\overline{y}} \otimes_{\mathcal{O}_{X, \overline{x}}}
\mathcal{O}_{Y, \overline{y}}
$$
Ainsi, nous devons vérifier que le morphisme canonique
$$
Q_{\overline{y}} \otimes_{\mathcal{O}_{X, \overline{x}}}
\mathcal{O}_{Y, \overline{y}}
\longrightarrow Q_{\overline{y}}
$$
est un isomorphisme dans la catégorie dérivée. Soit
$I_{\overline{x}} \subset \mathcal{O}_{X, \overline{x}}$ le
germe du faisceau d'idéaux définissant $Z$. Puisque $Z \to X$ est localement de
présentation finie, cet idéal est de type fini et les
groupes de cohomologie de $Q_{\overline{y}}$
sont constitués d'éléments annulés par une puissance de $I_{\overline{y}} = I_{\overline{x}}\mathcal{O}_{Y, \overline{y}}$ d'après le Lemme \ref{spaces-pushouts-lemma-stalk-of-module-with-support} appliqué à $Q$ sur $Y$. Il s'ensuit qu'ils sont également constitués d'éléments annulés par une puissance de $I_{\overline{x}}$. 
L'homomorphisme d'anneaux 
$\mathcal{O}_{X, \overline{x}} \to \mathcal{O}_{Y, \overline{y}}$ 
est plat et induit un isomorphisme après passage aux quotients par 
$I_{\overline{x}}$ et $I_{\overline{y}}$ car nous avons supposé que 
$f^{-1}Z \to Z$ est un isomorphisme. Par conséquent, nous voyons que 
les modules de cohomologie de 
$Q_{\overline{y}} \otimes_{\mathcal{O}_{X, \overline{x}}}
\mathcal{O}_{Y, \overline{y}}$ 
sont égaux aux modules de cohomologie de $Q_{\overline{y}}$ d'après 
Compléments d'algèbre, Lemme \ref{more-algebra-lemma-neighbourhood-isomorphism} 
ce qui termine la démonstration. 
\end{proof}

\begin{situation}
\label{spaces-pushouts-situation-formal-glueing}
Ici $S$ est un schéma de base, $f : Y \to X$ est un morphisme quasi-compact et quasi-séparé d'espaces algébriques sur $S$, et $Z \to X$ est une immersion fermée de présentation finie. Nous supposons que $f^{-1}(Z) \to Z$ est un isomorphisme et que $f$ est plat en tout point $y \in |f^{-1}Z|$. Nous posons $U = X \setminus Z$ et $V = Y \setminus f^{-1}(Z)$.
Considérons le diagramme
$$
\xymatrix{
V \ar[r]_{j'} \ar[d]_{f|_V} & Y \ar[d]^f \\
U \ar[r]^j & X
}
$$
\end{situation}

\noindent
Dans la situation \ref{spaces-pushouts-situation-formal-glueing}, nous définissons
$\textit{QCoh}(Y \to X, Z)$ comme la catégorie des
triplets $(\mathcal{H}, \mathcal{G}, \varphi)$ où
$\mathcal{H}$ est un faisceau quasi-cohérent de
$\mathcal{O}_U$-modules, $\mathcal{G}$ est un faisceau quasi-cohérent
de $\mathcal{O}_Y$-modules, et
$\varphi : f^*\mathcal{H} \to \mathcal{G}|_V$ est un isomorphisme
de $\mathcal{O}_V$-modules. Il existe un foncteur canonique
\begin{equation}
\label{spaces-pushouts-equation-formal-glueing-modules}
\QCoh(\mathcal{O}_X) \longrightarrow \textit{QCoh}(Y \to X, Z)
\end{equation}
qui associe $\mathcal{F}$ au système
$(\mathcal{F}|_U, f^*\mathcal{F}, can)$.
Par analogie avec la preuve donnée dans le cas affine, nous construisons
un foncteur dans la direction opposée. À un objet
$(\mathcal{H}, \mathcal{G}, \varphi)$, nous associons le $\mathcal{O}_X$-module
\begin{equation}
\label{spaces-pushouts-equation-reverse}
\Ker(j_*\mathcal{H} \oplus f_*\mathcal{G} \to (f \circ j')_*\mathcal{G}|_V)
\end{equation}.
Remarquons que $j$ et $j'$ sont des morphismes quasi-compacts car
$Z \to X$ est de présentation finie. Par conséquent, $f_*$, $j_*$ et $(f \circ j')_*$
transforment les modules quasi-cohérents en modules quasi-cohérents (Morphismes d'espaces, Lemme \ref{spaces-morphisms-lemma-pushforward}). Ainsi, le module (\ref{spaces-pushouts-equation-reverse}) est quasi-cohérent.

\begin{lemma}
\label{spaces-pushouts-lemma-adjoint}
Dans la situation \ref{spaces-pushouts-situation-formal-glueing}.
Le foncteur (\ref{spaces-pushouts-equation-reverse}) est adjoint à droite du
foncteur (\ref{spaces-pushouts-equation-formal-glueing-modules}).
\end{lemma}

\begin{proof}
Cela découle facilement de l'adjonction de $f^*$ à $f_*$
et de $j^*$ à $j_*$. Détails omis.
\end{proof}

\begin{lemma}
\label{spaces-pushouts-lemma-reverse-commutes-with-flat-base-change}
Dans la situation \ref{spaces-pushouts-situation-formal-glueing}.
Soit $X' \to X$ un morphisme plat d'espaces algébriques.
Posons $Z' = X' \times_X Z$ et $Y' = X' \times_X Y$.
Les images inverses $\QCoh(\mathcal{O}_X) \to \QCoh(\mathcal{O}_{X'})$
et $\QCoh(Y \to X, Z) \to \QCoh(Y' \to X', Z')$ sont compatibles
avec les foncteurs (\ref{spaces-pushouts-equation-reverse}) et
(\ref{spaces-pushouts-equation-formal-glueing-modules}).
\end{lemma}

\begin{proof}
Ceci est vrai parce que les images inverses commutent entre elles et parce que l'image inverse par un morphisme plat commute à l'image directe le long de morphismes quasi-compacts et quasi-séparés, voir Cohomologie des espaces, Lemme \ref{spaces-cohomology-lemma-flat-base-change-cohomology}.
\end{proof}

\begin{proposition}
\label{spaces-pushouts-proposition-formal-glueing-modules}
Dans la situation \ref{spaces-pushouts-situation-formal-glueing}, le foncteur
(\ref{spaces-pushouts-equation-formal-glueing-modules}) est une équivalence
avec un quasi-inverse donné par (\ref{spaces-pushouts-equation-reverse}).
\end{proposition}

\begin{proof}
Nous traitons d'abord le cas particulier où $X$ et $Y$ sont des schémas affines
et où le morphisme $f$ est plat. Écrivons $X = \Spec(R)$ et $Y = \Spec(S)$.
Alors $f$ correspond à un homomorphisme d'anneaux plat $R \to S$. De plus, $Z \subset X$
est découpé par un idéal de type fini $I \subset R$. Choisissons des générateurs $f_1, \ldots, f_t \in I$.
D'après la description des modules quasi-cohérents
en termes de modules
(Schémas, section \ref{schemes-section-quasi-coherent-affine}),
nous voyons que la catégorie $\textit{QCoh}(Y \to X, Z)$
est canoniquement équivalente à la catégorie
$\text{Glue}(R \to S, f_1, \ldots, f_t)$
de Compléments d'algèbre, Remarque \ref{more-algebra-remark-glueing-data}
de telle sorte que les foncteurs
(\ref{spaces-pushouts-equation-formal-glueing-modules}) et (\ref{spaces-pushouts-equation-reverse})
correspondent aux foncteurs $\text{Can}$ et $H^0$.
Ainsi, le résultat suit de
Compléments d'algèbre, Proposition \ref{more-algebra-proposition-equivalence}
dans ce cas.

\medskip\noindent
Nous revenons au cas général.
Soit $\mathcal{F}$ un module quasi-cohérent sur $X$.
Nous allons montrer que
$$
\alpha :
\mathcal{F}
\longrightarrow
\Ker\left(j_*\mathcal{F}|_U \oplus f_*f^*\mathcal{F} \to
(f \circ j')_*f^*\mathcal{F}|_V\right)
$$
est un isomorphisme. Soit $(\mathcal{H}, \mathcal{G}, \varphi)$
un objet de $\QCoh(Y \to X, Z)$. Nous allons montrer que
$$
\beta :
f^*\Ker\left(
j_*\mathcal{H} \oplus f_*\mathcal{G} \to (f \circ j')_*\mathcal{G}|_V
\right)
\longrightarrow
\mathcal{G}
$$
et
$$
\gamma :
j^*\Ker\left(
j_*\mathcal{H} \oplus f_*\mathcal{G} \to (f \circ j')_*\mathcal{G}|_V
\right)
\longrightarrow
\mathcal{H}
$$
sont des isomorphismes. Pour voir que ces affirmations sont vraies, il suffit de
examiner les germes. Soit $\overline{y}$ un point géométrique de $Y$ s'envoyant
sur le point géométrique $\overline{x}$ de $X$.

\medskip\noindent
Fixons un objet $(\mathcal{H}, \mathcal{G}, \varphi)$ de $\QCoh(Y \to X, Z)$.
Par le Lemme \ref{spaces-pushouts-lemma-stalk-formal-glueing}
et une vérification de diagramme (omise), le morphisme canonique
$$
\Ker(j_*\mathcal{H} \oplus f_*\mathcal{G} \to
(f \circ j')_*\mathcal{G}|_V)_{\overline{x}}
\longrightarrow
\Ker(
j_*\mathcal{H}_{\overline{x}} \oplus \mathcal{G}_{\overline{y}}
\to
j'_*\mathcal{G}_{\overline{y}}
)
$$
est un isomorphisme.

\medskip\noindent
En particulier, si $\overline{y}$ est un point géométrique de $V$, alors
nous voyons que $j'_*\mathcal{G}_{\overline{y}} = \mathcal{G}_{\overline{y}}$
et donc que ce noyau est égal à $\mathcal{H}_{\overline{x}}$.
Cela implique facilement que $\alpha_{\overline{x}}$, $\gamma_{\overline{x}}$,
et $\beta_{\overline{y}}$ sont des isomorphismes dans ce cas.

\medskip\noindent
Ensuite, supposons que $\overline{y}$ soit un point de $f^{-1}Z$.
Soient $I_{\overline{x}} \subset \mathcal{O}_{X, \overline{x}}$ et
$I_{\overline{y}} \subset \mathcal{O}_{Y, \overline{y}}$ les germes des idéaux
définissant respectivement $Z$ et $f^{-1}Z$.
Alors $I_{\overline{x}}$ est un idéal de type fini,
$I_{\overline{y}} = I_{\overline{x}}\mathcal{O}_{Y, \overline{y}}$,
et $\mathcal{O}_{X, \overline{x}} \to \mathcal{O}_{Y, \overline{y}}$
est un homomorphisme local plat induisant un isomorphisme
$\mathcal{O}_{X, \overline{x}}/I_{\overline{x}} =
\mathcal{O}_{Y, \overline{y}}/I_{\overline{y}}$.
À ce stade, nous pouvons conclure à l'aide du diagramme de catégories
$$
\xymatrix{
\QCoh(\mathcal{O}_X) \ar[r]_-{(\ref{spaces-pushouts-equation-formal-glueing-modules})} \ar[d] &
\QCoh(Y \to X, Z) \ar[d] \ar@/_2pc/[l]^{(\ref{spaces-pushouts-equation-reverse})} \\
\text{Mod}_{\mathcal{O}_{X, \overline{x}}} \ar[r]^-{\text{Can}} &
\text{Glue}(\mathcal{O}_{X, \overline{x}} \to \mathcal{O}_{Y, \overline{y}},
f_1, \ldots, f_t) \ar@/^2pc/[l]_{H^0}
}
$$
En effet, comme dans le premier paragraphe de la démonstration, nous identifions
$$
\text{Glue}(\mathcal{O}_{X, \overline{x}} \to \mathcal{O}_{Y, \overline{y}},
f_1, \ldots, f_t)
=
\QCoh(\Spec(\mathcal{O}_{Y, \overline{y}}) \to
\Spec(\mathcal{O}_{X, \overline{x}}), V(I_{\overline{x}}))
$$
Le foncteur vertical droit est donné par l'image inverse, et il est clair que
le carré intérieur est commutatif. Notre calcul du germe du
noyau dans le troisième paragraphe de la démonstration combiné avec
le Lemme \ref{spaces-pushouts-lemma-stalk-of-pushforward} implique que
le carré extérieur (en utilisant les flèches courbes) commute. Nous concluons alors en appliquant le cas d'un morphisme plat de schémas affines traité dans le premier paragraphe de la démonstration. 
\end{proof}

\begin{lemma}
\label{spaces-pushouts-lemma-derived-equivalent}
Dans la situation \ref{spaces-pushouts-situation-formal-glueing}, le foncteur
$Rf_*$ induit une équivalence entre $D_{\QCoh, |f^{-1}Z|}(\mathcal{O}_Y)$
et $D_{\QCoh, |Z|}(\mathcal{O}_X)$ avec quasi-inverse donné par
$Lf^*$.
\end{lemma}

\begin{proof}  
Puisque $f$ est quasi-compact et quasi-séparé, nous voyons que $Rf_*$  
définit un foncteur de $D_{\QCoh, |f^{-1}Z|}(\mathcal{O}_Y)$  
vers $D_{\QCoh, |Z|}(\mathcal{O}_X)$, voir  
Catégories dérivées des espaces, Lemme  
\ref{spaces-perfect-lemma-quasi-coherence-direct-image}.  
D'après Catégories dérivées des espaces, Lemme  
\ref{spaces-perfect-lemma-quasi-coherence-pullback}  
nous voyons que $Lf^*$ transforme $D_{\QCoh, |Z|}(\mathcal{O}_X)$  
en $D_{\QCoh, |f^{-1}Z|}(\mathcal{O}_Y)$.  
Dans le Lemme \ref{spaces-pushouts-lemma-formal-glueing-on-closed}, nous avons vu que  
$Lf^*Rf_*Q = Q$ pour $Q$ dans $D_{\QCoh, |f^{-1}Z|}(\mathcal{O}_Y)$.  
Par l'énoncé dual de Catégories dérivées, Lemme  
\ref{derived-lemma-fully-faithful-adjoint-kernel-zero}  
pour terminer la preuve, il suffit de montrer que $Lf^*K = 0$  
implique $K = 0$ pour $K$ dans $D_{\QCoh, |Z|}(\mathcal{O}_X)$.  
Ceci découle du fait que $f$ est plat en tous les points de  
$f^{-1}Z$ et que $f^{-1}Z \to Z$ est surjectif.  
\end{proof}

\begin{lemma}
\label{spaces-pushouts-lemma-dominate-by-fpqc-covering}
Dans la situation \ref{spaces-pushouts-situation-formal-glueing}, il existe un
recouvrement fpqc $\{X_i \to X\}_{i \in I}$ affinant la
famille $\{U \to X, Y \to X\}$.
\end{lemma}

\begin{proof}
Pour la définition et les propriétés générales des recouvrements fpqc, nous renvoyons à
Topologies, section \ref{topologies-section-fpqc}. En particulier, nous pouvons
d'abord choisir un recouvrement \'etale $\{X_i \to X\}$ avec $X_i$ affine et en
changeant $Y$, $Z$ et $U$ de base par chaque $X_i$, nous nous ramenons au cas où
$X$ est affine. Dans ce cas, $U$ est quasi-compact et donc une union finie
$U = U_1 \cup \ldots \cup U_n$ d'ouverts affines.
Alors $Z$ est quasi-compact donc aussi $f^{-1}Z$ est quasi-compact.
Ainsi nous pouvons choisir un schéma affine $W$ et un morphisme \'etale
$h : W \to Y$ tel que $h^{-1}f^{-1}Z \to f^{-1}Z$ soit surjectif.
Écrivons $W = \Spec(B)$ et $h^{-1}f^{-1}Z = V(J)$ où $J \subset B$
est un idéal de type fini.
Par Cohomologie pro-\'etale, Lemme \ref{proetale-lemma-localization}, il existe une localisation $B \to B'$ telle que les points de $\Spec(B')$ correspondent exactement aux points de $W = \Spec(B)$ se spécialisant en $h^{-1}f^{-1}Z = V(J)$. Il s'ensuit que le morphisme composé $\Spec(B') \to \Spec(B) = W \to Y \to X$ est plat car, par hypothèse, $f : Y \to X$ est plat en tous les points de $f^{-1}Z$. Ensuite, $\{\Spec(B') \to X, U_1 \to X, \ldots, U_n \to X\}$ est un recouvrement fpqc d'après Topologies, Lemme \ref{topologies-lemma-recognize-fpqc-covering}. 
\end{proof}




\section{Recollement formel d'espaces algébriques}
\label{spaces-pushouts-section-formal-glueing-spaces}

\noindent
Dans la situation \ref{spaces-pushouts-situation-formal-glueing}, nous considérons la catégorie
$\textit{Espaces}(Y \to X, Z)$
des diagrammes commutatifs d'espaces algébriques sur $S$ de la forme
$$
\xymatrix{
U' \ar[d] & V' \ar[l] \ar[d] \ar[r] & Y' \ar[d] \\
U & V \ar[l] \ar[r] & Y
}
$$
où les deux carrés sont cartésiens. Il existe un foncteur canonique
\begin{equation}
\label{spaces-pushouts-equation-formal-glueing-spaces}
\textit{Espaces}/X \longrightarrow \textit{Espaces}(Y \to X, Z)
\end{equation}
qui envoie $X' \to X$ sur les morphismes
$U \times_X X' \leftarrow V \times_X X' \rightarrow Y \times_X X'$.

\begin{lemma}
\label{spaces-pushouts-lemma-equivalence-on-affine}
Dans la situation \ref{spaces-pushouts-situation-formal-glueing}, le foncteur
(\ref{spaces-pushouts-equation-formal-glueing-spaces}) se restreint à une
équivalence
\begin{enumerate}
\item de la catégorie des espaces algébriques affines sur $X$
vers la sous-catégorie pleine de $\textit{Espaces}(Y \to X, Z)$ consistant
en $(U' \leftarrow V' \rightarrow Y')$ avec $U' \to U$, $V' \to V$,
et $Y' \to Y$ affines,
\item de la catégorie des immersions fermées $X' \to X$
vers la sous-catégorie pleine de $\textit{Espaces}(Y \to X, Z)$ consistant
en $(U' \leftarrow V' \rightarrow Y')$ avec $U' \to U$, $V' \to V$,
et $Y' \to Y$ immersions fermées, et
\item même énoncé que dans (2) pour les morphismes finis.
\end{enumerate}
\end{lemma}

\begin{proof}
La catégorie des espaces algébriques affines sur $X$ est équivalente à la
catégorie des faisceaux quasi-cohérents $\mathcal{A}$ de $\mathcal{O}_X$-algèbres.
La sous-catégorie pleine de $\textit{Espaces}(Y \to X, Z)$ consistant en
$(U' \leftarrow V' \rightarrow Y')$ avec $U' \to U$, $V' \to V$,
et $Y' \to Y$ affines est équivalente à la catégorie des
objets en algèbres de $\QCoh(Y \to X, Z)$. Dans les deux cas, cela découle
de Morphismes d'espaces, Lemme
\ref{spaces-morphisms-lemma-affine-equivalence-algebras}
avec un quasi-inverse donné par la construction du spectre relatif
(Morphismes d'espaces, Définition
\ref{spaces-morphisms-definition-relative-spec})
qui commute à tout changement de base. Ainsi, la partie (1) du
lemme découle de la Proposition \ref{spaces-pushouts-proposition-formal-glueing-modules}.

\medskip\noindent
La pleine fidélité dans la partie (2) découle de la partie (1). Pour la surjectivité essentielle, on réduit par la partie (1) à prouver que $X' \to X$ est une immersion fermée si et seulement si à la fois $U \times_X X' \to U$ et $Y \times_X X' \to Y$ sont des immersions fermées. Par le Lemme \ref{spaces-pushouts-lemma-dominate-by-fpqc-covering}, $\{U \to X, Y \to X\}$ peut être raffiné par un recouvrement fpqc. Par conséquent, le résultat découle de Descente sur les espaces, Lemme \ref{spaces-descent-lemma-descending-property-closed-immersion}.

\medskip\noindent
Pour (3), le même argument que pour (2), joint à
Descente sur les espaces, Lemme
\ref{spaces-descent-lemma-descending-property-finite}.
\end{proof}

\begin{lemma}
\label{spaces-pushouts-lemma-reflects-isomorphisms}
Dans la situation \ref{spaces-pushouts-situation-formal-glueing}, le foncteur
(\ref{spaces-pushouts-equation-formal-glueing-spaces}) reflète les isomorphismes.
\end{lemma}

\begin{proof}
Par un argument formel de changement de base, cela se réduit à l'assertion suivante : un morphisme $a : X' \to X$ d'espaces algébriques pour lequel $U \times_X X' \to U$ et $Y \times_X X' \to Y$ sont des isomorphismes est lui-même un isomorphisme. La famille $\{U \to X, Y \to X\}$ peut être affinée par un recouvrement fpqc d'après le Lemme \ref{spaces-pushouts-lemma-dominate-by-fpqc-covering}. 
Par conséquent, le résultat découle de
Descente sur les espaces, Lemme \ref{spaces-descent-lemma-descending-property-isomorphism}.
\end{proof}

\begin{lemma}
\label{spaces-pushouts-lemma-fully-faithful-on-separated}
Dans la situation \ref{spaces-pushouts-situation-formal-glueing}, le foncteur
(\ref{spaces-pushouts-equation-formal-glueing-spaces}) est pleinement fidèle
sur les espaces algébriques séparés sur $X$. Plus précisément, il induit
une bijection
$$
\Mor_X(X'_1, X'_2)
\longrightarrow
\Mor_{\textit{Espaces}(Y \to X, Z)}(F(X'_1), F(X'_2))
$$
lorsque $X'_2 \to X$ est séparé.
\end{lemma}

\begin{proof}
Puisque $X'_2 \to X$ est séparé, le graphe $i : X'_1 \to X'_1 \times_X X'_2$
d'un morphisme $X'_1 \to X'_2$ sur $X$ est une immersion fermée, voir
Morphismes d'espaces, Lemme \ref{spaces-morphisms-lemma-semi-diagonal}.
De plus, une immersion fermée $i : T \to X'_1 \times_X X'_2$ est le graphe d'un
morphisme si et seulement si $\text{pr}_1 \circ i$ est un isomorphisme.
Il en est de même pour
\begin{enumerate}
\item le graphe d'un morphisme $U \times_X X'_1 \to U \times_X X'_2$ sur $U$,
\item le graphe d'un morphisme $V \times_X X'_1 \to V \times_X X'_2$ sur $V$,
et
\item le graphe d'un morphisme $Y \times_X X'_1 \to Y \times_X X'_2$ sur $Y$.
\end{enumerate}
De plus, si des morphismes comme dans (1), (2), (3) se recollent pour former un
morphisme dans la catégorie $\textit{Espaces}(Y \to X, Z)$, alors ces
graphes se recollent pour donner un objet de
$\textit{Espaces}(Y \times_X (X'_1 \times_X X'_2) \to X'_1 \times_X X'_2,
Z \times_X (X'_1 \times_X X'_2))$
dont les trois morphismes sont des immersions fermées. La preuve est terminée en appliquant les Lemmes \ref{spaces-pushouts-lemma-equivalence-on-affine} et \ref{spaces-pushouts-lemma-reflects-isomorphisms}. 
\end{proof}









\section{Recollement et le théorème de Beauville-Laszlo}
\label{spaces-pushouts-section-glueing-beauville-laszlo}

\noindent
Soit $R \to R'$ un homomorphisme d'anneaux et soit $f \in R$ un élément tel que
$$
0 \to R \to R_f \oplus R' \to R'_f \to 0
$$
soit une suite exacte courte. Cela implique que $R/f^nR \cong R'/f^nR'$
pour tout $n$ et que $(R \to R', f)$ est une paire de recollement au sens de
Compléments d'algèbre, section \ref{more-algebra-section-beauville-laszlo}.
Posons $X = \Spec(R)$, $U = \Spec(R_f)$, $X' = \Spec(R')$ et
$U' = \Spec(R'_f)$. Considérons le diagramme
$$
\xymatrix{
U' \ar[r] \ar[d] & X' \ar[d] \\
U \ar[r] & X
}
$$
Dans cette situation, nous pouvons considérer la catégorie
$\textit{Espaces}(U \leftarrow U' \to X')$ dont les objets
sont des diagrammes commutatifs
$$
\xymatrix{
V \ar[d] & V' \ar[l] \ar[d] \ar[r] & Y' \ar[d] \\
U & U' \ar[l] \ar[r] & X'
}
$$
d'espaces algébriques avec les deux carrés cartésiens et dont les morphismes
sont définis de manière évidente. Un objet de cette catégorie sera
noté $(V, V', Y')$ en omettant les flèches.
Il existe un foncteur
\begin{equation}
\label{spaces-pushouts-equation-beauville-laszlo-glueing-spaces}
\textit{Espaces}/X
\longrightarrow
\textit{Espaces}(U \leftarrow U' \to X')
\end{equation}
donné par le changement de base : $Y \mapsto (U \times_X Y, U' \times_X Y, X' \times_X Y)$.

\medskip\noindent
Nous avons vu dans
Compléments d'algèbre, section \ref{more-algebra-section-beauville-laszlo}
que les données de recollement ne permettent pas de retrouver tout $R$-module $M$ à partir de ses
données de recollement. De même, le foncteur
(\ref{spaces-pushouts-equation-beauville-laszlo-glueing-spaces})
ne sera pas pleinement fidèle sur la catégorie de tous les espaces sur $X$.
Afin de dégager une sous-catégorie appropriée d'espaces algébriques sur $X$, nous avons besoin d'un lemme.

\begin{lemma}
\label{spaces-pushouts-lemma-glueable}
Soit $(R \to R', f)$ une paire de recollement, voir ci-dessus. Soit $Y$ un espace algébrique sur $X$. Les énoncés suivants sont équivalents
\begin{enumerate}
\item il existe un recouvrement \'etale $\{Y_i \to Y\}_{i \in I}$
avec $Y_i$ affine et $\Gamma(Y_i, \mathcal{O}_{Y_i})$
recollable comme $R$-module,
\item pour tout morphisme \'etale $W \to Y$ avec $W$ affine
$\Gamma(W, \mathcal{O}_W)$ est un $R$-module recollable.
\end{enumerate}
\end{lemma}

\begin{proof}
Il est immédiat que (2) implique (1). Supposons que $\{Y_i \to Y\}$
soit comme dans (1) et soit $W \to Y$ comme dans (2). Alors
$\{Y_i \times_Y W \to W\}_{i \in I}$ est un recouvrement \'etale,
que nous pouvons affiner par un recouvrement \'etale
$\{W_j \to W\}_{j = 1, \ldots, m}$ avec $W_j$ affine
(Topologies, Lemme \ref{topologies-lemma-etale-affine}).
Ainsi, pour terminer la démonstration, il suffit de montrer
les trois assertions algébriques suivantes :
\begin{enumerate}
\item si $R \to A \to B$ sont des homomorphismes d'anneaux avec $A \to B$ \'etale
et $A$ est recollable en tant que $R$-module, alors $B$ est recollable en tant que
$R$-module,
\item les produits finis de $R$-modules recollables sont recollables,
\item si $R \to A \to B$ sont des homomorphismes d'anneaux avec $A \to B$ \'etale et fidèlement plat
et $B$ est recollable en tant que $R$-module, alors $A$ est recollable en tant que $R$-module.  
\end{enumerate}
  
En effet, le premier énoncé implique que $\Gamma(W_j, \mathcal{O}_{W_j})$ est un $R$-module recollable, le deuxième que $\prod \Gamma(W_j, \mathcal{O}_{W_j})$ est un $R$-module recollable, et le troisième que $\Gamma(W, \mathcal{O}_W)$ est un $R$-module recollable.

\medskip\noindent
Considérons un homomorphisme \'etale de $R$-algèbres $A \to B$. Posons $A' = A \otimes_R R'$ et $B' = B \otimes_R R' = A' \otimes_A B$.
Les affirmations (1) et (3) découlent alors des faits suivants :
(a) $A$, respectivement $B$, est recollable si et seulement si la suite $$
0 \to A \to A_f \oplus A' \to A'_f \to 0,
\quad\text{respectivement}\quad
0 \to B \to B_f \oplus B' \to B'_f \to 0,
$$
est exacte, (b) la seconde suite s'obtient en appliquant à la première le foncteur $- \otimes_A B$, et (c) la platitude, respectivement la fidèle platitude, de $A \to B$. Nous omettons la démonstration de (2).
\end{proof}

\noindent
Soit $(R \to R', f)$ une paire de recollement, voir ci-dessus.
Nous dirons qu'un espace algébrique $Y$ sur $X = \Spec(R)$
est {\it recollable pour $(R \to R', f)$}
si les conditions équivalentes du Lemme \ref{spaces-pushouts-lemma-glueable}
sont satisfaites.

\begin{lemma}
\label{spaces-pushouts-lemma-glueing-affines}
Soit $(R \to R', f)$ une paire de recollement, voir ci-dessus.
Le foncteur (\ref{spaces-pushouts-equation-beauville-laszlo-glueing-spaces})
se restreint à une équivalence entre la catégorie des objets affines
$Y/X$ qui sont recollables pour $(R \to R', f)$ et la
sous-catégorie pleine des objets $(V, V', Y')$ de
$\textit{Espaces}(U \leftarrow U' \to X')$
avec $V$, $V'$, $Y'$ affines.
\end{lemma}

\begin{proof}
Soit $(V, V', Y')$ un objet de
$\textit{Espaces}(U \leftarrow U' \to X')$
avec $V$, $V'$, $Y'$ affines.
Écrivons $V = \Spec(A_1)$ et $Y' = \Spec(A')$. D'après notre définition de la
catégorie $\textit{Espaces}(U \leftarrow U' \to X')$, nous trouvons que
$V'$ est le spectre de $A_1 \otimes_{R_f} R'_f = A_1 \otimes_R R'$
et le spectre de $A'_f$. Ainsi, nous obtenons un isomorphisme
$\varphi : A'_f \to A_1 \otimes_R R'$ de $R'_f$-algèbres.
D'après Compléments d'algèbre, Théorème \ref{more-algebra-theorem-BL-glueing}
il existe un unique $R$-module recollable $A$ et des isomorphismes
$A_f \to A_1$ et $A \otimes_R R' \to A'$ de modules compatibles avec
$\varphi$. Puisque la suite
$$
0 \to A \to A_1 \oplus A' \to A'_f \to 0
$$
est exacte courte, les multiplications sur $A_1$ et $A'$ définissent
une structure de $R$-algèbre unique sur $A$ telle que les morphismes $A \to A_1$
et $A \to A'$ sont des homomorphismes d'anneaux. Nous omettons la vérification que cette construction définit un quasi-inverse au foncteur
(\ref{spaces-pushouts-equation-beauville-laszlo-glueing-spaces})
restreint aux sous-catégories mentionnées dans l'énoncé du lemme.
\end{proof}

\begin{lemma}
\label{spaces-pushouts-lemma-glueing-affines-etale}
Soit $P$ l'une des propriétés suivantes des morphismes :
``fini'', ``immersion fermée'', ``plat'', ``de type fini'',
``plat et de présentation finie'', ``\'etale''.
Sous l'équivalence du Lemme \ref{spaces-pushouts-lemma-glueing-affines}
les morphismes ayant $P$ correspondent à des morphismes de triples
dont les composantes ont $P$.
\end{lemma}

\begin{proof}
Soit $P'$ l'une des propriétés suivantes des homomorphismes d'anneaux :
``fini'', ``surjectif'', ``plat'', ``de type fini'',
``plat et de présentation finie'', ``\'etale''.
En termes algébriques, l'énoncé signifie ce qui suit :
Si $A \to B$ est un homomorphisme de $R$-algèbres et que $A$ et $B$
sont recollables pour $(R \to R', f)$, alors
$A_f \to B_f$ et $A \otimes_R R' \to B \otimes_R R'$ ont $P'$
si et seulement si $A \to B$ a $P'$.

\medskip\noindent
Par Compléments d'algèbre, Lemmes \ref{more-algebra-lemma-faithful-descent}
et \ref{more-algebra-lemma-BL-flat} l'énoncé algébrique
est vrai lorsque $P'$ est « fini » ou « plat ».

\medskip\noindent
Si $A_f \to B_f$ et $A \otimes_R R' \to B \otimes_R R'$ sont surjectifs,
alors le conoyau $N = B/A$ est un $R$-module avec $N_f = 0$ et $N \otimes_R R' = 0$ et
par conséquent est nul d'après Compléments d'algèbre, Lemme
\ref{more-algebra-lemma-BL-faithful}. Ainsi $A \to B$ est surjectif.

\medskip\noindent
Si $A_f \to B_f$ et $A \otimes_R R' \to B \otimes_R R'$ sont de type fini,
alors nous pouvons choisir un homomorphisme de $A$-algèbres $A[x_1, \ldots, x_n] \to B$
tel que $A_f[x_1, \ldots, x_n] \to B_f$ et
$(A \otimes_R R')[x_1, \ldots, x_n] \to B \otimes_R R'$ sont surjectifs
(détail mineur omis). Nous en concluons que $A[x_1, \ldots, x_n] \to B$
est surjectif d'après le résultat précédent. Ainsi $A \to B$ est de type fini.

\medskip\noindent
Si $A_f \to B_f$ et $A \otimes_R R' \to B \otimes_R R'$ sont plats et de présentation finie, alors nous savons que $A \to B$ est plat et de type fini grâce à ce que nous avons déjà démontré. Choisissons une surjection $A[x_1, \ldots, x_n] \to B$ et notons $I$ son noyau. Par la platitude de $B$ sur $A$, nous voyons que $I_f$ est le noyau de $A_f[x_1, \ldots, x_n] \to B_f$ et que $I \otimes_R R'$ est le noyau de $A \otimes_R R'[x_1, \ldots, x_n] \to B \otimes_R R'$. Ainsi, $I_f$ est un $A_f[x_1, \ldots, x_n]$-module de type fini et $I \otimes_R R'$ est un $(A \otimes_R R')[x_1, \ldots, x_n]$-module de type fini. D'après Compléments d'algèbre, Lemme \ref{more-algebra-lemma-faithful-descent}, appliqué à $I$ considéré comme module sur $A[x_1, \ldots, x_n]$, nous concluons que $I$ est un idéal de type fini, puis que $A \to B$ est plat et de présentation finie.

\medskip\noindent
Si $A_f \to B_f$ et $A \otimes_R R' \to B \otimes_R R'$ sont \'etales,
alors nous savons que $A \to B$ est plat et de présentation finie d'après ce que
nous avons déjà montré. Puisque les fibres de $\Spec(B) \to \Spec(A)$
sont isomorphes aux fibres de $\Spec(B_f) \to \Spec(A_f)$ ou
de $\Spec(B/fB) \to \Spec(A/fA)$, nous concluons que $A \to B$ est non ramifié,
voir Morphismes, Lemmes \ref{morphisms-lemma-unramified-over-field}
et \ref{morphisms-lemma-unramified-etale-fibres}.
Nous concluons que $A \to B$ est \'etale d'après
Morphismes, Lemme \ref{morphisms-lemma-flat-unramified-etale} par exemple.
\end{proof}

\begin{lemma}
\label{spaces-pushouts-lemma-glueing-f}
Soit $(R \to R', f)$ une paire de recollement, voir ci-dessus.
Le foncteur (\ref{spaces-pushouts-equation-beauville-laszlo-glueing-spaces})
est fidèle sur la sous-catégorie pleine des
espaces algébriques $Y/X$ recollables pour $(R \to R', f)$.
\end{lemma}

\begin{proof}
Soient $f, g : Y \to Z$ deux morphismes d'espaces algébriques sur $X$
avec $Y$ et $Z$ recollables pour $(R \to R', f)$ de telle sorte que $f$ et $g$ soient envoyés
sur le même morphisme dans la catégorie $\textit{Espaces}(U \leftarrow U' \to X')$.
Nous devons montrer que l'égaliseur $E \to Y$ de $f$ et $g$ est un isomorphisme.
En travaillant \'etale localement sur $Y$, nous pouvons supposer que $Y$ est un schéma affine.
Alors $E$ est un schéma et le morphisme $E \to Y$ est un monomorphisme
et localement quasi-fini, voir Morphismes d'espaces, Lemme
\ref{spaces-morphisms-lemma-properties-diagonal}.
De plus, le changement de base de $E \to Y$ vers $U$ et vers $X'$ est
un isomorphisme. Comme l'espace topologique sous-jacent à $Y$ est la réunion disjointe de l'ouvert affine
$V = U \times_X Y$ et du fermé affine $V(f) \times_X Y$, nous concluons que $E$ est l'union disjointe de leurs images inverses isomorphes.  
Il s'ensuit en particulier que $E$ est quasi-compact.  
Par le théorème principal de Zariski (Compléments sur les morphismes, Lemme \ref{more-morphisms-lemma-quasi-finite-separated-pass-through-finite})  
nous concluons que $E$ est quasi-affine.  
Posons $B = \Gamma(E, \mathcal{O}_E)$ et $A = \Gamma(Y, \mathcal{O}_Y)$  
de sorte que nous ayons un homomorphisme de $R$-algèbres $A \to B$.  
Puisque $E \to Y$ devient un isomorphisme après changement de base vers $U$ et $X'$,  
nous obtenons des morphismes d'anneaux $B \to A_f$ et $B \to A \otimes_R R'$  
qui coïncident comme morphismes vers $A \otimes_R R'_f$. Puisque $A$ est recollable  
par rapport à $(R \to R', f)$, nous obtenons un morphisme d'anneaux $B \to A$ qui est l'inverse à gauche  
du morphisme $A \to B$. Le morphisme correspondant $Y = \Spec(A) \to \Spec(B)$  
a son image ensembliste contenue dans le sous-schéma ouvert $E \subset \Spec(B)$ car  
cela est vrai après changement de base vers $U$ et $X'$. Ainsi, nous obtenons un morphisme $Y \to E$ sur $Y$. Puisque $E \to Y$ est un monomorphisme, nous en concluons que $Y \to E$ est un isomorphisme comme voulu. 
\end{proof}

\begin{lemma}
\label{spaces-pushouts-lemma-glueing-ff}
Soit $(R \to R', f)$ une paire de recollement, voir ci-dessus.
Le foncteur (\ref{spaces-pushouts-equation-beauville-laszlo-glueing-spaces})
est pleinement fidèle sur la sous-catégorie pleine des
espaces algébriques $Y/X$ qui sont (a) recollables pour $(R \to R', f)$
et (b) ont une diagonale affine $Y \to Y \times_X Y$.
\end{lemma}

\begin{proof}
Soient $Y, Z$ deux espaces algébriques sur $X$ qui sont tous deux recollables pour
$(R \to R', f)$ et supposons que la diagonale de $Z$ soit affine. Soient
$a : U \times_X Y \to U \times_X Z$ sur $U$ et
$b : X' \times_X Y \to X' \times_X Z$ sur $X'$ deux morphismes
d'espaces algébriques
qui induisent le même morphisme $c : U' \times_X Y \to U' \times_X Z$
sur $U'$.
Nous voulons construire un morphisme $f : Y \to Z$ sur $X$
qui induise les morphismes $a$ et $b$ après changement de base à $U$ et à $X'$.
Par la fidélité du Lemme \ref{spaces-pushouts-lemma-glueing-f}, il suffit de construire
le morphisme $f$ \'etale localement sur $Y$ (détails omis).
Ainsi, nous pouvons supposer que $Y$ est affine.

\medskip\noindent
Soit $y \in |Y|$ un point. Si $y$ s'envoie dans l'ouvert $U \subset X$,
alors $U \times_X Y$ est un ouvert de $Y$ sur lequel le morphisme $f$
est défini (nous pouvons simplement prendre $a$).
Ainsi nous pouvons supposer que $y$ s'envoie dans le sous-ensemble fermé
$V(f)$ de $X$. Puisque $R/fR = R'/fR'$ il existe un point unique
$y' \in |X' \times_X Y|$ s'envoyant sur $y$. On note
$z' = b(y') \in |X' \times_X Z|$ et $z \in |Z|$ les images de $y'$.
Choisissons un voisinage \'etale $(W, w) \to (Z, z)$
avec $W$ affine. Remarquons que
$$
(U \times_X W) \times_{U \times_X Z, a} (U \times_X Y),\quad
(U' \times_X W) \times_{U' \times_X Z, c} (U' \times_X Y),
$$
et
$$
(X' \times_X W) \times_{X' \times_X Z, b} (X' \times_X Y)
$$
forment un objet de $\textit{Espaces}(U \leftarrow U' \to X')$
dont les composantes sont affines (c'est ici que nous utilisons le fait que $Z$ a une diagonale affine).
Ainsi, par le Lemme \ref{spaces-pushouts-lemma-glueing-affines}
il existe un schéma affine unique $V$ recollable pour $(R \to R', f)$ tel que
$$
(U \times_X V, U' \times_X V, X' \times_X V)
$$
est le triple affiché ci-dessus. Par la pleine fidélité dans le cas affine (Lemme \ref{spaces-pushouts-lemma-glueing-affines}), nous obtenons deux morphismes uniques $V \to W$ et $V \to Y$ coïncidant avec les morphismes donnés par la première et la seconde projection sur $U$ et $X'$ dans la construction ci-dessus. Par le Lemme \ref{spaces-pushouts-lemma-glueing-affines-etale}, le morphisme $V \to Y$ est \'etale. Pour terminer la preuve, il suffit de montrer qu'il existe un point $v \in |V|$ se projetant sur $y$ (car alors $f$ est défini sur un voisinage \'etale de $y$, à savoir $V$). Il existe un point unique $w' \in |X' \times_X W|$ s'envoyant sur $w$. Par unicité, $w'$ s'envoie sur $z'$ par le morphisme $|X' \times_X W| \to |X' \times_X Z|$. Considérons alors le diagramme cartésien $$
\xymatrix{
X' \times_X V \ar[r] \ar[d] & X' \times_X W \ar[d] \\
X' \times_X Y \ar[r] & X' \times_X Z
}
$$
pour voir qu'il existe un point $v' \in |X' \times_X V|$ s'envoyant sur $y'$ et $w'$, voir Propriétés des espaces, Lemme \ref{spaces-properties-lemma-points-cartesian}. Bien sûr, l'image $v$ de $v'$ dans $|V|$ s'envoie bien sur $y$ et la preuve est complète. 
\end{proof}

\begin{lemma}
\label{spaces-pushouts-lemma-glueing-quasi-affines}
Soit $(R \to R', f)$ une paire de recollement, voir ci-dessus. Tout objet
$(V, V', Y')$ de $\textit{Espaces}(U \leftarrow U' \to X')$
avec $V$, $V'$, $Y'$ quasi-affine est isomorphe à
l'image sous le foncteur (\ref{spaces-pushouts-equation-beauville-laszlo-glueing-spaces})
d'un espace algébrique séparé $Y$ sur $X$.
\end{lemma}

\begin{proof}
Choisissons $n'$, $T' \to Y'$ et $n_1$, $T_1 \to V$ comme dans
Propriétés, Lemme \ref{properties-lemma-quasi-affine-presentation}.
Considérons le diagramme
$$
\xymatrix{
& &
T_1 \times_V V' \times_Y T' \ar[ld] \ar[rd] \\
T_1 \ar[d] &
T_1 \times_V V' \ar[l] \ar[dr] & &
V' \times_{Y'} T' \ar[r] \ar[dl] &
T' \ar[d] \\
V & &
V' \ar[rr] \ar[ll] & &
Y'
}
$$
Observons que $T_1 \times_V V'$ et $V' \times_{Y'} T'$
sont affines (les morphismes $V' \to V$ et $V' \to Y'$
sont affines comme changements de base des morphismes affines $U' \to U$
et $U' \to X'$). Par construction, nous voyons que
$$
\mathbf{A}^{n'}_{T_1 \times_V V'} \cong
T_1 \times_V V' \times_{Y'} T' \cong
\mathbf{A}^{n_1}_{V' \times_{Y'} T'}
$$
En d'autres termes, les schémas affines $\mathbf{A}^{n'}_{T_1}$
et $\mathbf{A}^{n_1}_{T'}$ font partie d'un triple formant un objet affine de
$\textit{Espaces}(U \leftarrow U' \to X')$.
D'après le Lemme \ref{spaces-pushouts-lemma-glueing-affines}
il existe un morphisme de schémas affines $T \to X$
et des isomorphismes $U \times_X T \cong \mathbf{A}^{n'}_{T_1}$
et $X' \times_X T \cong \mathbf{A}^{n_1}_{T'}$ compatibles
avec les isomorphismes affichés ci-dessus.
Ces isomorphismes produisent des morphismes
$$
U \times_X T \longrightarrow V
\quad\text{et}\quad
X' \times_X T \longrightarrow Y'
$$
satisfaisant la propriété de
Propriétés, Lemme \ref{properties-lemma-quasi-affine-presentation}, avec $n = n' + n_1$, et définissant en outre un morphisme du triple $(U \times_X T, U' \times_X T, X' \times_X T)$ vers notre triple $(V, V', Y')$ dans la catégorie $\textit{Espaces}(U \leftarrow U' \to X')$.

\medskip\noindent
Par le Lemme \ref{spaces-pushouts-lemma-glueing-affines} il existe un schéma affine $W$ dont
l'image dans $\textit{Espaces}(U \leftarrow U' \to X')$ est isomorphe au
triple
$$
((U \times_X T) \times_V (U \times_X T),
(U' \times_X T) \times_{V'} (U' \times_X T),
(X' \times_X T) \times_{Y'} (X' \times_X T))
$$
Par la pleine fidélité de cette construction, nous obtenons
deux morphismes $p_0, p_1 : W \to T$ dont les changements de base
à $U, U', X'$ sont les morphismes de projection.
Par le Lemme \ref{spaces-pushouts-lemma-glueing-affines-etale}
les morphismes $p_0, p_1$ sont plats et de présentation finie et
le morphisme $(p_0, p_1) : W \to T \times_X T$ est une immersion fermée.
En fait, $W \to T \times_X T$ définit une relation d'équivalence : par les lemmes
démontrés ci-dessus, nous pouvons vérifier la symétrie, la réflexivité et la transitivité
après changement de base à $U$ et $X'$, car elles y sont évidentes (détails omis).
Ainsi, le faisceau quotient
$$
Y = T/W
$$
est un espace algébrique par exemple par
Amorçage, Théorème \ref{bootstrap-theorem-final-bootstrap}. Par construction, $Y/X$ est envoyé sur le triple $(V, V', Y')$. Le changement de base de la diagonale $\Delta : Y \to Y \times_X Y$ par le morphisme plat, quasi-compact et surjectif $T \times_X T \to Y \times_X Y$ est l'immersion fermée $W \to T \times_X T$. Ainsi $\Delta$ est une immersion fermée par Descente sur les espaces, Lemme \ref{spaces-descent-lemma-descending-property-closed-immersion}. Il s'ensuit que l'espace algébrique $Y$ est séparé et la preuve est complète. 
\end{proof}








\section{Coégalisateurs et recollement}
\label{spaces-pushouts-section-coequalizer-glue}

\noindent
Soit $X$ un espace algébrique noethérien et $Z \to X$ une immersion fermée.
Soit $X' \to X$ l'éclatement dans $Z$. Dans cette section, nous montrons que
$X$ peut être reconstitué à partir de $X'$, $Z_n$ et des données de recollement, où $Z_n$
est le $n$-ième voisinage infinitésimal de $Z$ dans $X$.

\begin{lemma}
\label{spaces-pushouts-lemma-coequalizer}
Soit $S$ un schéma. Soit
$$
g : Y \longrightarrow X
$$
un morphisme d'espaces algébriques sur $S$. Supposons que $X$ soit localement noethérien,
et que $g$ soit propre. Soit $R = Y \times_X Y$ avec des morphismes de projection
$t, s : R \to Y$. Il existe un coégalisateur $X'$ de $s, t : R \to Y$
dans la catégorie des espaces algébriques sur $S$. De plus
\begin{enumerate}
\item Le morphisme $X' \to X$ est fini.
\item Le morphisme $Y \to X'$ est propre.
\item Le morphisme $Y \to X'$ est surjectif.
\item Le morphisme $X' \to X$ est universellement injectif.
\item Si $g$ est surjectif, le morphisme $X' \to X$
est un homéomorphisme universel.
\end{enumerate}
\end{lemma}

\begin{proof}
Notons $h : R \to X$ la composée de $s$ (ou de $t$)
avec $g$. Alors $h$ est propre par Morphismes d'espaces, Lemmes
\ref{spaces-morphisms-lemma-base-change-proper} et
\ref{spaces-morphisms-lemma-composition-proper}.
Les faisceaux
$$
g_*\mathcal{O}_Y
\quad\text{et}\quad
h_*\mathcal{O}_R
$$
sont des $\mathcal{O}_X$-algèbres cohérentes par Cohomologie des espaces, Lemme
\ref{spaces-cohomology-lemma-proper-pushforward-coherent}.
Les $X$-morphismes $s$, $t$ induisent des homomorphismes de $\mathcal{O}_X$-algèbres
$s^\sharp, t^\sharp$ de la première à la seconde.
Posons
$$
\mathcal{A} = \text{Égalisateur}\left(s^\sharp, t^\sharp :
g_*\mathcal{O}_Y \longrightarrow h_*\mathcal{O}_R\right)
$$
Alors $\mathcal{A}$ est une $\mathcal{O}_X$-algèbre cohérente et nous pouvons
définir
$$
X' = \underline{\Spec}_X(\mathcal{A})
$$
comme dans Morphismes d'espaces, Définition
\ref{spaces-morphisms-definition-relative-spec}.
Par Morphismes d'espaces, Remarque
\ref{spaces-morphisms-remark-factorization-quasi-compact-quasi-separated}
et par fonctorialité de la construction $\underline{\Spec}$
il existe une factorisation
$$
Y \longrightarrow X' \longrightarrow X
$$
et le morphisme $g' : Y \to X'$ égalise $s$ et $t$.

\medskip\noindent
Avant de montrer que $X'$ est le coégalisateur de $s$ et $t$, nous montrons que $Y \to X'$ et $X' \to X$ ont les propriétés souhaitées. Puisque $\mathcal{A}$ est un $\mathcal{O}_X$-module cohérent, il est clair que $X' \to X$ est un morphisme fini d'espaces algébriques. Cela prouve (1).
Le morphisme $Y \to X'$ est propre d'après Morphismes d'espaces, Lemme \ref{spaces-morphisms-lemma-universally-closed-permanence}. Cela prouve (2).
Notons $Y \to Y' \to X$, avec $Y' = \underline{\Spec}_X(g_*\mathcal{O}_Y)$, la factorisation de Stein de $g$ ; voir Compléments sur les morphismes d'espaces, Théorème \ref{spaces-more-morphisms-theorem-stein-factorization-Noetherian}.
Bien sûr, nous obtenons des morphismes $Y \to Y' \to X' \to X$ correspondant aux morphismes étudiés ci-dessus.
Puisque $\mathcal{O}_{X'} \subset g_*\mathcal{O}_Y$ est une extension finie, nous voyons que $Y' \to X'$ est fini et surjectif.
Quelques détails omis ; indice : utiliser Algèbre, Lemme \ref{algebra-lemma-integral-overring-surjective}, et réduire au cas affine par localisation \'etale. Puisque $Y \to Y'$ est surjectif (avec des fibres géométriquement connexes), nous concluons que $Y \to X'$ est surjectif. Cela prouve (3). Pour montrer que $X' \to X$ est universellement injectif, nous devons montrer que $X' \to X' \times_X X'$ est surjectif, voir Morphismes d'espaces, Définition \ref{spaces-morphisms-definition-universally-injective} et Lemme \ref{spaces-morphisms-lemma-universally-injective}. Puisque $Y \to X'$ est surjectif (voir ci-dessus) et que les changements de base et les compositions de morphismes surjectifs sont surjectifs d'après Morphismes d'espaces, Lemme \ref{spaces-morphisms-lemma-base-change-surjective} et \ref{spaces-morphisms-lemma-composition-surjective}, nous voyons que $Y \times_X Y \to X' \times_X X'$ est surjectif. Cependant, puisque $Y \to X'$ égalise $s$ et $t$, nous voyons que $Y \times_X Y \to X' \times_X X'$ se factorise par $X' \to X' \times_X X'$ et nous concluons que ce dernier morphisme est surjectif. Cela prouve (4). Enfin, si $g$ est surjectif, alors puisque $g$ se factorise par $X' \to X$, nous voyons que $X' \to X$ est surjectif. Comme un morphisme surjectif, universellement injectif et fini est un homéomorphisme universel (car il est universellement bijectif et universellement fermé), cela prouve (5).

\medskip\noindent
Dans le reste de la preuve, nous montrons que $Y \to X'$ est le coégalisateur de $s$ et $t$ dans la catégorie des espaces algébriques sur $S$. Observons que $X'$ est localement noethérien (Morphismes d'espaces, Lemme \ref{spaces-morphisms-lemma-locally-finite-type-locally-noetherian}). Observons aussi que $Y \times_{X'} Y \to Y \times_X Y$ est un isomorphisme car $Y \to X'$ égalise $s$ et $t$ (c'est une assertion catégorique). Ainsi, afin de prouver que $Y \to X'$ est le coégalisateur de $s$ et $t$, nous pouvons supposer, et supposons, que $X = X'$. En d'autres termes, $\mathcal{O}_X$ est l'égaliseur des morphismes $s^\sharp, t^\sharp : g_*\mathcal{O}_Y \to h_*\mathcal{O}_R$.

\medskip\noindent
Soit $X_1 \to X$ un morphisme plat d'espaces algébriques sur $S$ avec $X_1$ localement noethérien. Notons $g_1 : Y_1 \to X_1$, $h_1 : R_1 \to X_1$ et $s_1, t_1 : R_1 \to Y_1$ les changements de base de $g, h, s, t$ à $X_1$.
Bien sûr, $g_1$ est propre et $R_1 = Y_1 \times_{X_1} Y_1$.
Par le théorème de changement de base plat pour les images directes de modules quasi-cohérents, Cohomologie des espaces, Lemme \ref{spaces-cohomology-lemma-flat-base-change-cohomology}, nous voyons que $\mathcal{O}_{X_1}$ est l'égaliseur des morphismes $s_1^\sharp, t_1^\sharp : g_{1, *}\mathcal{O}_{Y_1} \to
h_{1, *}\mathcal{O}_{R_1}$. Ainsi, toutes les hypothèses que nous avons sont préservées par ce changement de base.

\medskip\noindent
À ce stade, nous allons vérifier les conditions (1) et (2) du
Lemme \ref{spaces-pushouts-lemma-colimit-check-etale-locally}. La condition (1)
suit du Lemme \ref{spaces-pushouts-lemma-descend-etale-proper-surjective}
et du fait que $g$ est propre et surjectif (parce que $X = X'$).
Pour vérifier la condition (2), compte tenu des remarques sur le changement de base ci-dessus,
nous revenons à l'énoncé discuté et prouvé dans le paragraphe suivant.

\medskip\noindent
Supposons que $S = \Spec(A)$ soit un schéma affine, que $X = X'$ soit un schéma affine, et que $Z$ soit un schéma affine sur $S$. Nous devons montrer que $$
\Mor_S(X, Z) \longrightarrow
\text{Égalisateur}(s, t : \Mor_S(Y, Z) \to \Mor_S(R, Z))
$$
est bijectif. Cela est toutefois clair puisque $X = X'$, ce qui implique que $\mathcal{O}_X$ est l'égaliseur des morphismes $s^\sharp, t^\sharp : g_*\mathcal{O}_Y \to h_*\mathcal{O}_R$, et entraîne à son tour $$
\Gamma(X, \mathcal{O}_X) =
\text{Égalisateur}\left(
s^\sharp, t^\sharp : \Gamma(Y, \mathcal{O}_Y) \to
\Gamma(R, \mathcal{O}_R)
\right)
$$. En effet, nous avons $$
\Mor_S(X, Z) = \Hom_A(\Gamma(Z, \mathcal{O}_Z), \Gamma(X, \mathcal{O}_X))
$$ et de même pour $Y$ et $R$, voir Propriétés des espaces, Lemme \ref{spaces-properties-lemma-morphism-to-affine-scheme}. 
\end{proof}

\noindent
Nous travaillerons dans la situation suivante.

\begin{situation}
\label{spaces-pushouts-situation-coequalizer-glue}
Soit $S$ un schéma. Soit $X$ un espace algébrique localement noethérien sur $S$.
Soit $Z \to X$ une immersion fermée et soit $U \subset X$ le sous-espace ouvert complémentaire. Enfin, soit $f : X' \to X$ un morphisme propre d'espaces algébriques tel que $f^{-1}(U) \to U$ soit un isomorphisme.
\end{situation}

\begin{lemma}
\label{spaces-pushouts-lemma-coequalizer-glue}
Dans la situation \ref{spaces-pushouts-situation-coequalizer-glue}, soit $Y = X' \amalg Z$ et
$R = Y \times_X Y$ avec des projections $t, s : R \to Y$. Il existe un
coégalisateur $X_1$ de $s, t : R \to Y$ dans la catégorie des espaces algébriques
sur $S$. Le morphisme $X_1 \to X$ est un homéomorphisme universel fini,
un isomorphisme sur $U$, et $Z \to X$ admet un relèvement à $X_1$.
\end{lemma}

\begin{proof}
L'existence de $X_1$ et le fait que $X_1 \to X$ est un homéomorphisme universel fini sont un cas particulier du Lemme \ref{spaces-pushouts-lemma-coequalizer}.
La formation de $X_1$ commute avec la localisation \'etale sur $X$ (voir la preuve du Lemme \ref{spaces-pushouts-lemma-coequalizer}).
Ainsi, le morphisme $X_1 \to X$ est un isomorphisme sur $U$.
Il résulte immédiatement de la construction que $Z \to X$ admet un relèvement à $X_1$.
\end{proof}

\noindent
Dans la situation \ref{spaces-pushouts-situation-coequalizer-glue} pour $n \geq 1$, soit
$Z_n \subset X$ le voisinage infinitésimal d'ordre $n$ de $Z$ dans $X$, c'est-à-dire le sous-espace fermé défini par la $n$-ième puissance du faisceau d'idéaux coupant $Z$. Considérons $Y_n = X' \amalg Z_n$
et $R_n = Y_n \times_X Y_n$ et le coégalisateur
$$
\xymatrix{
R_n \ar@<1ex>[r] \ar@<-1ex>[r] & Y_n \ar[r] & X_n \ar[r] & X
}
$$
comme dans le Lemme \ref{spaces-pushouts-lemma-coequalizer-glue}. Les morphismes $Y_n \to Y_{n + 1}$
et $R_n \to R_{n + 1}$ induisent des morphismes
\begin{equation}
\label{spaces-pushouts-equation-system-coequalizers}
X_1 \to X_2 \to X_3 \to \ldots \to X
\end{equation}
Chacun de ces morphismes est un homéomorphisme universel puisque les morphismes
$X_n \to X$ sont des homéomorphismes universels.

\begin{lemma}
\label{spaces-pushouts-lemma-essentially-constant}
Dans la situation \ref{spaces-pushouts-situation-coequalizer-glue}, supposons que $X$ soit quasi-compact.
Dans (\ref{spaces-pushouts-equation-system-coequalizers}), pour tout $n$ suffisamment grand, il existe un entier $m$ tel que $X_n \to X_{n + m}$ se factorise par une immersion fermée $X \to X_{n + m}$.
\end{lemma}

\begin{proof}
Examinons de plus près la construction de $X_n$
et comment elle change lorsque nous augmentons $n$. Nous avons
$X_n = \underline{\Spec}(\mathcal{A}_n)$
où $\mathcal{A}_n$ est l'égaliseur de $s_n^\sharp$ et $t_n^\sharp$
allant de $g_{n, *}\mathcal{O}_{Y_n}$ à $h_{n, *}\mathcal{O}_{R_n}$.
Ici $g_n : Y_n = X' \amalg Z_n \to X$ et $h_n : R_n = Y_n \times_X Y_n \to X$
sont les morphismes donnés. Soit $\mathcal{I} \subset \mathcal{O}_X$ le
faisceau cohérent d'idéaux correspondant à $Z$. Alors
$$
g_{n, *}\mathcal{O}_{Y_n} =
f_*\mathcal{O}_{X'} \times \mathcal{O}_X/\mathcal{I}^n
$$
De même, nous avons une décomposition
$$
R_n = X' \times_X X' \amalg X' \times_X Z_n \amalg
Z_n \times_X X' \amalg Z_n \times_X Z_n
$$
Comme $Z_n \to X$ est un monomorphisme, nous voyons que
$X' \times_X Z_n = Z_n \times_X X'$ et que cette identification
est compatible avec les deux morphismes vers $X$, avec les deux morphismes vers
$X'$, et avec les deux morphismes vers $Z_n$.
Notons $f_n : X' \times_X Z_n \to X$ le morphisme vers $X$.
Posons
$$
\mathcal{A} = \text{Égalisateur}(
\xymatrix{
f_*\mathcal{O}_{X'} \ar@<1ex>[r] \ar@<-1ex>[r] &
(f \times f)_*\mathcal{O}_{X' \times_X X'}
}
)
$$
Par les remarques ci-dessus, nous trouvons que
$$
\mathcal{A}_n =
\text{Égalisateur}(
\xymatrix{
\mathcal{A} \times \mathcal{O}_X/\mathcal{I}^n \ar@<1ex>[r] \ar@<-1ex>[r] &
f_{n, *}\mathcal{O}_{X' \times_X Z_n}
}
)
$$
Nous avons des morphismes canoniques
$$
\mathcal{O}_X \to \ldots \to \mathcal{A}_3 \to \mathcal{A}_2 \to \mathcal{A}_1
$$
de $\mathcal{O}_X$-algèbres cohérentes. L'énoncé du lemme signifie que
pour $n$ suffisamment grand, il existe un $m \geq 0$ tel que l'image de
$\mathcal{A}_{n + m} \to \mathcal{A}_n$ soit isomorphe à $\mathcal{O}_X$.
Cette assertion se vérifie localement pour la topologie \'etale sur $X$. Ainsi, d'après Propriétés des espaces,
Lemme \ref{spaces-properties-lemma-quasi-compact-affine-cover}
nous pouvons supposer que $X$ est un schéma noethérien affine.

\medskip\noindent
Puisque $X_n \to X$ est un isomorphisme sur $U$, nous voyons que le noyau de $\mathcal{O}_X \to \mathcal{A}_n$ est supporté sur $|Z|$.
Puisque $X$ est noethérien, la suite des noyaux $\mathcal{J}_n = \Ker(\mathcal{O}_X \to \mathcal{A}_n)$ se stabilise (Cohomologie des espaces, Lemme \ref{spaces-cohomology-lemma-acc-coherent}).
Écrivons $\mathcal{J}_{n_0} = \mathcal{J}_{n_0 + 1} = \ldots = \mathcal{J}$.
D'après Cohomologie des espaces, Lemme \ref{spaces-cohomology-lemma-power-ideal-kills-sheaf} nous trouvons que $\mathcal{I}^t \mathcal{J} = 0$ pour un entier $t \geq 0$.
D'autre part, il existe un homomorphisme de $\mathcal{O}_X$-algèbres $\mathcal{A}_n \to \mathcal{O}_X/\mathcal{I}^n$ et donc $\mathcal{J} \subset \mathcal{I}^n$ pour tout $n$.
Par Artin-Rees (Cohomologie des espaces, Lemme \ref{spaces-cohomology-lemma-Artin-Rees}), nous trouvons que $\mathcal{J} \cap \mathcal{I}^n \subset \mathcal{I}^{n - c}\mathcal{J}$ pour un entier $c \geq 0$ et tout $n \gg 0$. Nous concluons que $\mathcal{J} = 0$.

\medskip\noindent
Choisissons $n \geq n_0$ comme dans le paragraphe précédent. Ensuite
$\mathcal{O}_X \to \mathcal{A}_n$ est injectif. Par conséquent, il suffit maintenant de trouver $m \geq 0$ tel que l'image de
$\mathcal{A}_{n + m} \to \mathcal{A}_n$ soit égale à l'image de $\mathcal{O}_X$. Observons que $\mathcal{A}_n$
s'insère dans une suite exacte courte
$$
0 \to \Ker(\mathcal{A} \to f_{n, *}\mathcal{O}_{X' \times_X Z_n})
\to \mathcal{A}_n \to \mathcal{O}_X/\mathcal{I}^n \to 0
$$
et de même pour $\mathcal{A}_{n + m}$. Par conséquent, il suffit de montrer
$$
\Ker(\mathcal{A} \to f_{n + m, *}\mathcal{O}_{X' \times_X Z_{n + m}})
\subset
\Im(\mathcal{I}^n \to \mathcal{A})
$$
pour un entier $m \geq 0$. Pour ce faire, nous pouvons travailler \'etale localement sur
$X$ et puisque $X$ est noethérien, nous pouvons supposer que $X$ est
un schéma affine noethérien. Écrivons $X = \Spec(R)$ et $\mathcal{I}$
correspond à l'idéal $I \subset R$. Soit $\mathcal{A} = \widetilde{A}$
pour une $R$-algèbre finie $A$. Soit $f_*\mathcal{O}_{X'} = \widetilde{B}$
pour une $R$-algèbre finie $B$. Alors $R \to A \subset B$ et ces homomorphismes deviennent des isomorphismes en inversant n'importe quel élément de $I$.

\medskip\noindent
Notons que $f_{n, *}\mathcal{O}_{X' \times_X Z_n}$
est égal à $f_*(\mathcal{O}_{X'}/I^n\mathcal{O}_{X'})$
dans la notation utilisée dans Cohomologie des espaces, section
\ref{spaces-cohomology-section-theorem-formal-functions}.
D'après la Cohomologie des espaces, Lemme
\ref{spaces-cohomology-lemma-ML-cohomology-powers-ideal}
nous voyons qu'il existe un entier $c \geq 0$ tel que
$$
\Ker(B \to \Gamma(X, f_*(\mathcal{O}_{X'}/I^{n + m + c}\mathcal{O}_{X'})))
$$
est contenu dans $I^{n + m}B$. D'autre part, comme $R \to B$ est
fini et un isomorphisme après inversion de tout élément de $I$
nous voyons que $I^{n + m}B \subset \Im(I^n \to B)$ pour $m$ assez grand
(peut être choisi indépendamment de $n$). Cela termine la preuve puisque $A \subset B$.
\end{proof}

\begin{remark}
\label{spaces-pushouts-remark-essentially-constant}
La signification du Lemme \ref{spaces-pushouts-lemma-essentially-constant}
est que le système $X_1 \to X_2 \to X_3 \to \ldots$ est essentiellement
constant de valeur $X$. Voir Catégories, Définition
\ref{categories-definition-essentially-constant-diagram}.
\end{remark}






\section{Compactifications}
\label{spaces-pushouts-section-compactifications}

\noindent
Cette section est l'analogue de Compléments sur la platitude, section
\ref{flat-section-compactifications}. Le théorème de cette
section est le théorème principal de \cite{CLO}.

\medskip\noindent
Soit $B$ un espace algébrique quasi-compact et quasi-séparé sur un schéma de base $S$. Nous dirons qu'un espace algébrique $X$ sur $B$ {\it a une compactification sur $B$} ou {\it est compactifiable sur $B$} s'il existe une immersion ouverte quasi-compacte $X \to \overline{X}$ dans un espace algébrique $\overline{X}$ propre sur $B$. Si $X$ a une compactification sur $B$, alors $X \to B$ est séparé et de type fini. Le théorème principal de cette section est que la réciproque est également vraie.

\begin{lemma}
\label{spaces-pushouts-lemma-check-separated}
Soit $S$ un schéma. Soit $X \to Y$ un morphisme d'espaces algébriques sur $S$.
Si $(U \subset X, f : V \to X)$ est un carré distingué élémentaire
tel que $U \to Y$ et $V \to Y$ soient séparés et que
$U \times_X V \to U \times_Y V$ soit fermé, alors $X \to Y$ est séparé.
\end{lemma}

\begin{proof}
Nous devons vérifier que $\Delta : X \to X \times_Y X$ est une immersion fermée.
Il existe un recouvrement \'etale de $X \times_Y X$ donné par les quatre composantes
$U \times_Y U$, $U \times_Y V$, $V \times_Y U$ et $V \times_Y V$.
Observons que
$(U \times_Y U) \times_{(X \times_Y X), \Delta} X  = U$,
$(U \times_Y V) \times_{(X \times_Y X), \Delta} X = U \times_X V$,
$(V \times_Y U) \times_{(X \times_Y X), \Delta} X = V \times_X U$ et
$(V \times_Y V) \times_{(X \times_Y X), \Delta} X = V$.
Ainsi, les hypothèses du lemme
nous disent exactement que $\Delta$ est une immersion fermée.
\end{proof}

\begin{lemma}
\label{spaces-pushouts-lemma-separate-disjoint-locally-closed-by-blowing-up}
Soit $S$ un schéma.
Soit $X$ un espace algébrique quasi-compact et quasi-séparé sur $S$.
Soit $U \subset X$ un ouvert quasi-compact.
\begin{enumerate}
\item Si $Z_1, Z_2 \subset X$ sont des sous-espaces fermés de présentation finie tels que $Z_1 \cap Z_2 \cap U = \emptyset$, alors
il existe un éclatement $U$-admissible $X' \to X$
tel que les transformées strictes de $Z_1$ et $Z_2$ soient disjointes.
\item Si $T_1, T_2 \subset |U|$ sont des sous-ensembles constructibles fermés disjoints, alors il existe un éclatement $U$-admissible $X' \to X$
tel que les adhérences de $T_1$ et $T_2$ soient disjointes.
\end{enumerate}
\end{lemma}

\begin{proof}
Preuve de (1). L'hypothèse que $Z_i \to X$ est de présentation finie signifie que le faisceau d'idéaux quasi-cohérent $\mathcal{I}_i$ de $Z_i$ est de type fini, voir 
Morphismes d'espaces, Lemme
\ref{spaces-morphisms-lemma-closed-immersion-finite-presentation}.
Notons $Z \subset X$ le sous-espace fermé défini par le produit $\mathcal{I}_1 \mathcal{I}_2$.
Observons que $Z \cap U$ est la réunion disjointe de $Z_1 \cap U$ et $Z_2 \cap U$. D'après Diviseurs sur les espaces, Lemme
\ref{spaces-divisors-lemma-separate-disjoint-opens-by-blowing-up}
il existe un éclatement $U \cap Z$-admissible $Z' \to Z$ tel que les transformées strictes de $Z_1$ et $Z_2$ soient disjointes.
Notons $Y \subset Z$ le centre de cet éclatement.
Alors $Y \to X$ est une immersion fermée de présentation finie comme composée de $Y \to Z$ et $Z \to X$ (Diviseurs sur les espaces, Définition
\ref{spaces-divisors-definition-admissible-blowup} et
Morphismes d'espaces, Lemme
\ref{spaces-morphisms-lemma-composition-finite-presentation}). Ainsi, l'éclatement $X' \to X$ de $Y$ est un éclatement $U$-admissible. D'après les propriétés générales des transformées strictes, les transformées strictes de $Z_1$ et $Z_2$ par rapport à $X' \to X$ sont les mêmes que leurs transformées strictes par rapport à $Z' \to Z$, voir Diviseurs sur les espaces, Lemme \ref{spaces-divisors-lemma-strict-transform}. Ainsi, (1) est démontré.

\medskip\noindent
Preuve de (2). Par Limites d'espaces, Lemme
\ref{spaces-limits-lemma-quasi-coherent-finite-type-ideals}
il existe un faisceau quasi-cohérent d'idéaux de type fini
$\mathcal{J}_i \subset \mathcal{O}_U$ tel que
$T_i = V(\mathcal{J}_i)$ (au sens ensembliste).
Par Limites d'espaces, Lemme \ref{spaces-limits-lemma-extend}
il existe un faisceau quasi-cohérent d'idéaux de type fini $\mathcal{I}_i \subset \mathcal{O}_X$
dont la restriction à $U$ est $\mathcal{J}_i$.
Appliquons le résultat de la partie (1) aux sous-espaces fermés $Z_i = V(\mathcal{I}_i)$ pour conclure.
\end{proof}

\begin{lemma}
\label{spaces-pushouts-lemma-blowup-iso-along}
Soit $S$ un schéma. Soit $f : X \to Y$ un morphisme propre d'espaces algébriques quasi-compacts et quasi-séparés sur $S$. Soit $V \subset Y$ un ouvert quasi-compact et $U = f^{-1}(V)$. Soit $T \subset |V|$ un sous-ensemble fermé tel que $f|_U : U \to V$ soit un isomorphisme sur un voisinage ouvert de $T$ dans $V$. Alors il existe un éclatement $V$-admissible $Y' \to Y$ tel que la transformée stricte $f' : X' \to Y'$ de $f$ soit un isomorphisme sur un voisinage ouvert de l'adhérence de $T$ dans $|Y'|$.
\end{lemma}

\begin{proof}
Soit $T' \subset |V|$ le complément de l'ouvert maximal sur lequel
$f|_U$ est un isomorphisme. Alors $T', T$ sont fermés dans $|V|$ et
$T \cap T' = \emptyset$. Comme $|V|$ est un espace topologique spectral
(Propriétés des espaces, Lemme
\ref{spaces-properties-lemma-quasi-compact-quasi-separated-spectral})
nous pouvons trouver des sous-ensembles constructibles fermés $T_c, T'_c$ de $|V|$
avec $T \subset T_c$, $T' \subset T'_c$ tels que
$T_c \cap T'_c = \emptyset$ (choisir un ouvert quasi-compact
$W$ de $|V|$ qui contient $T'$ et ne rencontre pas $T$
et poser $T_c = |V| \setminus W$, puis choisir un ouvert quasi-compact
$W'$ de $|V|$ qui contient $T_c$ et ne rencontre pas $T'$
et poser $T'_c = |V| \setminus W'$).
D'après le Lemme \ref{spaces-pushouts-lemma-separate-disjoint-locally-closed-by-blowing-up}
nous pouvons, après avoir remplacé $Y$ par un éclatement $V$-admissible, supposer que $T_c$ et $T'_c$ ont des adhérences disjointes dans $|Y|$.  
Soit $Y_0$ le sous-espace ouvert de $Y$ correspondant à l'ouvert $|Y| \setminus \overline{T}'_c$ et posons $V_0 = V \cap Y_0$, $U_0 = U \times_V V_0$ et $X_0 = X \times_Y Y_0$.  
Comme $U_0 \to V_0$ est un isomorphisme, nous pouvons trouver un éclatement $V_0$-admissible $Y'_0 \to Y_0$ tel que la transformée stricte $X'_0$ de $X_0$ s'envoie isomorphiquement sur $Y'_0$, voir Compléments sur les morphismes d'espaces, Lemme \ref{spaces-more-morphisms-lemma-zariski-after-blowup}.  
D'après Diviseurs sur les espaces, Lemme \ref{spaces-divisors-lemma-extend-admissible-blowups} il existe un éclatement $V$-admissible $Y' \to Y$ dont la restriction à $Y_0$ est $Y'_0 \to Y_0$. Si $f' : X' \to Y'$ désigne la transformée stricte de $f$, alors l'assertion voulue est vraie puisque $f'$ se restreint à un isomorphisme sur $Y'_0$.  
\end{proof}

\begin{lemma}
\label{spaces-pushouts-lemma-blowup-etale-along}
Soit $S$ un schéma. Considérons un diagramme
$$
\xymatrix{
X \ar[d]_f & U \ar[l] \ar[d]_{f|_U} & A \ar[d] \ar[l] \\
Y & V \ar[l] & B \ar[l]
}
$$
d'espaces algébriques quasi-compacts et quasi-séparés sur $S$.
Supposons que
\begin{enumerate}
\item $f$ soit propre,
\item $V$ soit un ouvert quasi-compact de $Y$, $U = f^{-1}(V)$,
\item $B \subset V$ et $A \subset U$ soient des sous-espaces fermés,
\item $f|_A : A \to B$ soit un isomorphisme, et que
$f$ soit \'etale en chaque point de $A$.
\end{enumerate}
Alors il existe un éclatement $V$-admissible $Y' \to Y$ tel que la transformée stricte
$f' : X' \to Y'$ ait la propriété suivante : pour chaque point géométrique
$\overline{a}$ de l'adhérence de $|A|$ dans $|X'|$
il existe un quotient $\mathcal{O}_{X', \overline{a}} \to \mathcal{O}$
tel que $\mathcal{O}_{Y', f'(\overline{a})} \to \mathcal{O}$
soit fini et plat.
\end{lemma}

\noindent
Comme le montre la preuve, un résultat plus fort est vrai, mais cet énoncé, déjà assez long, suffira par la suite.

\begin{proof}
Soit $T' \subset |U|$ le complément de l'ouvert maximal sur lequel
$f|_U$ est \'etale. Alors $T'$ est fermé dans $|U|$ et disjoint de $|A|$.
Comme $|U|$ est un espace topologique spectral (Propriétés des espaces, Lemme
\ref{spaces-properties-lemma-quasi-compact-quasi-separated-spectral})
nous pouvons trouver des sous-ensembles constructibles fermés $T_c, T'_c$ de $|U|$
avec $|A| \subset T_c$, $T' \subset T'_c$ tels que
$T_c \cap T'_c = \emptyset$ (voir la preuve du Lemme \ref{spaces-pushouts-lemma-blowup-iso-along}).
D'après le Lemme \ref{spaces-pushouts-lemma-separate-disjoint-locally-closed-by-blowing-up}
il existe un éclatement $U$-admissible $X_1 \to X$ tel que
$T_c$ et $T'_c$ aient des adhérences disjointes dans $|X_1|$.
Soit $X_{1, 0}$ le sous-espace ouvert de $X_1$ correspondant à l'ouvert
$|X_1| \setminus \overline{T}'_c$ et posons $U_0 = U \cap X_{1, 0}$.
Remarquons que l'image schématique $\overline{A}_1 \subset X_1$ de $A$ est contenue dans $X_{1, 0}$ par construction.

\medskip\noindent
Après avoir remplacé $Y$ par un éclatement $V$-admissible et pris les transformées strictes, nous pouvons supposer que $X_{1, 0} \to Y$ est plat, quasi-fini et de présentation finie, voir
Compléments sur les morphismes d'espaces, Lemmes
\ref{spaces-more-morphisms-lemma-flat-after-blowing-up} et
\ref{spaces-more-morphisms-lemma-flat-fp-dimension-over-dense-open}.
Considérons le diagramme commutatif
$$
\vcenter{
\xymatrix{
X_1 \ar[rr] \ar[rd] & & X \ar[ld] \\
& Y
}
}
\quad\text{et le diagramme}\quad
\vcenter{
\xymatrix{
\overline{A}_1 \ar[rr] \ar[rd] & & \overline{A} \ar[ld] \\
& \overline{B}
}
}
$$
d'images schématiques. Le morphisme $\overline{A}_1 \to \overline{A}$
est surjectif : comme il est propre, l'image schématique de $\overline{A}_1 \to \overline{A}$ est $\overline{A}$, et nous pouvons alors appliquer Morphismes d'espaces, Lemme
\ref{spaces-morphisms-lemma-scheme-theoretic-image-is-proper}.
L'énoncé sur les anneaux locaux strictement henséliens s'obtient
en relevant le point géométrique $\overline{a}$
en un point géométrique $\overline{a}_1$ de $\overline{A}_1$ et en posant
$\mathcal{O} = \mathcal{O}_{X_1, \overline{a}_1}$. En effet, puisque
$X_1 \to Y$ est plat et quasi-fini sur
$X_{1, 0} \supset \overline{A}_1$, l'homomorphisme
$\mathcal{O}_{Y', f'(\overline{a})} \to \mathcal{O}_{X_1, \overline{a}_1}$ est fini et plat, voir Algèbre, Lemme \ref{algebra-lemma-quasi-finite-strict-henselization} et \ref{algebra-lemma-characterize-henselian}. 
\end{proof}

\begin{lemma}
\label{spaces-pushouts-lemma-replaced-by-strict-transform}
Soit $S$ un schéma. Soient $X \to B$ et $Y \to B$ des morphismes d'espaces algébriques sur $S$. Soit $U \subset X$ un sous-espace ouvert. Soit $V \to X \times_B Y$ un morphisme quasi-compact dont la composée avec la première projection a son image dans $U$. Soit $Z \subset X \times_B Y$ l'image schématique de $V \to X \times_B Y$. Soit $X' \to X$ un éclatement $U$-admissible. Alors l'image schématique de $V \to X' \times_B Y$ est la transformée stricte de $Z$ par rapport à cet éclatement.
\end{lemma}

\begin{proof}
Notons $Z' \to Z$ la transformée stricte. Le morphisme $Z' \to X'$
induit un morphisme $Z' \to X' \times_B Y$ qui est une immersion fermée
(car $Z'$ est un sous-espace fermé de $X' \times_X Z$ par définition).
Ainsi, pour terminer la preuve, il suffit de montrer que l'image schématique $Z''$ de $V \to Z'$ est $Z'$. Observons que $Z'' \subset Z'$
est un sous-espace fermé tel que $V \to Z'$ se factorise par $Z''$.
Comme $V \to X \times_B Y$ et $V \to X' \times_B Y$ sont tous deux
quasi-compacts (pour le second, cela découle de Morphismes d'espaces, Lemme
\ref{spaces-morphisms-lemma-quasi-compact-permanence}
et du fait que $X' \times_B Y \to X \times_B Y$ est séparé
comme un changement de base d'un morphisme propre), d'après Morphismes d'espaces, Lemme
\ref{spaces-morphisms-lemma-quasi-compact-scheme-theoretic-image}
nous voyons que $Z \cap (U \times_B Y) = Z'' \cap (U \times_B Y)$.
Ainsi, le morphisme d'inclusion $Z'' \to Z'$ est un isomorphisme loin du diviseur exceptionnel $E$ de $Z' \to Z$. Cependant, le faisceau structural de $Z'$ n'a pas de sections non nulles supportées sur $E$ (par définition des transformées strictes) et nous concluons que la surjection $\mathcal{O}_{Z'} \to \mathcal{O}_{Z''}$ doit être un isomorphisme. 
\end{proof}

\begin{lemma}
\label{spaces-pushouts-lemma-compactification-dominates}
Soit $S$ un schéma. Soit $B$ un espace algébrique quasi-compact et quasi-séparé sur $S$. Soit $U$ un espace algébrique de type fini et séparé sur $B$. Soit $V \to U$ un morphisme \'etale. Si $V$ possède une compactification $V \subset Y$ sur $B$, alors il existe un éclatement $V$-admissible $Y' \to Y$ et un ouvert $V \subset V' \subset Y'$ tel que $V \to U$ s'étende en un morphisme propre $V' \to U$.
\end{lemma}

\begin{proof}
Considérons l'image schématique $Z \subset Y \times_B U$
du morphisme ``diagonal'' $V \to Y \times_B U$. Si nous remplaçons
$Y$ par un éclatement $V$-admissible, alors $Z$ est remplacé par
la transformée stricte par rapport à cet éclatement, voir
Lemme \ref{spaces-pushouts-lemma-replaced-by-strict-transform}. Ainsi, d'après
Compléments sur les morphismes d'espaces, Lemme
\ref{spaces-more-morphisms-lemma-zariski-after-blowup},
nous pouvons supposer que $Z \to Y$
est une immersion ouverte. Si $V' \subset Y$ désigne l'image, alors nous
voyons que le morphisme induit $V' \to U$ est propre car la
projection $Y \times_B U \to U$ est propre et $V' \cong Z$
est un sous-espace fermé de $Y \times_B U$.
\end{proof}

\noindent
Le lemme suivant est formulé pour des espaces algébriques séparés de type fini sur un espace algébrique de type fini sur $\mathbf{Z}$. La version pour les espaces algébriques quasi-compacts et quasi-séparés est également vraie (avec essentiellement la même démonstration), mais sera trivialement impliquée par le théorème principal de cette section. Nous encourageons vivement le lecteur à lire d'abord la démonstration de ce lemme dans le cas des schémas.

\begin{lemma}
\label{spaces-pushouts-lemma-two-compactifications}
Soit $B$ un espace algébrique de type fini sur $\mathbf{Z}$.
Soit $U$ un espace algébrique de type fini et séparé sur $B$.
Soit $(U_2 \subset U, f : U_1 \to U)$ un carré distingué élémentaire. Supposons que $U_1$ et $U_2$ possèdent des compactifications sur $B$ et que $U_1 \times_U U_2 \to U$ ait une image dense.
Alors $U$ a une compactification sur $B$.
\end{lemma}

\begin{proof}
Choisissons une compactification $U_i \subset X_i$ sur $B$ pour $i = 1, 2$. Nous pouvons supposer que $U_i$ est schématiquement dense dans $X_i$. Nous pouvons supposer qu'il existe un ouvert $V_i \subset X_i$ et un morphisme propre $\psi_i : V_i \to U$ prolongeant $U_i \to U$, voir le Lemme \ref{spaces-pushouts-lemma-compactification-dominates}. Considérons le diagramme
$$
\xymatrix{
U_i \ar[r] \ar[d] & V_i \ar[r] \ar[dl]^{\psi_i} & X_i \\
U
}
$$
Notons $Z_1 \subset U$ le sous-espace fermé réduit correspondant au sous-ensemble fermé $|U| \setminus |U_2|$. Rappelons que $f^{-1}Z_1$ est un sous-espace fermé de $U_1$ s'envoyant isomorphiquement sur $Z_1$. Notons $Z_2 \subset U$ le sous-espace fermé réduit correspondant au sous-ensemble fermé $|U| \setminus \Im(|f|) =
|U_2| \setminus \Im(|U_1 \times_U U_2| \to |U_2|)$. Ainsi nous avons
$$
U = U_2 \amalg Z_1 = Z_2 \amalg \Im(f) =
Z_2 \amalg \Im(U_1 \times_U U_2 \to U_2) \amalg Z_1
$$
au sens ensembliste. Notons $Z_{i, i} \subset V_i$ l'image inverse de $Z_i$ par $\psi_i$. Observons que $\psi_2$ est un isomorphisme sur un voisinage ouvert de $Z_2$. Observons que $Z_{1, 1} = \psi_1^{-1}Z_1 = f^{-1}Z_1 \amalg T$ pour un certain sous-espace fermé $T \subset V_1$ disjoint de $f^{-1}Z_1$ et que $\psi_1$ est en outre \'etale le long de $f^{-1}Z_1$.  
Notons $Z_{i, j} \subset V_i$ l'image inverse de $Z_j$ par $\psi_i$.  
Observons que $\psi_i : Z_{i, j} \to Z_j$ est un morphisme propre.  
Puisque $Z_i$ et $Z_j$ sont des sous-espaces fermés disjoints de $U$, on voit que $Z_{i, i}$ et $Z_{i, j}$ sont des sous-espaces fermés disjoints de $V_i$.

\medskip\noindent
Notons $\overline{Z}_{i, i}$ et $\overline{Z}_{i, j}$ les images schématiques de $Z_{i, i}$ et $Z_{i, j}$ dans $X_i$.
Rappelons que $|Z_{i, j}|$ est dense dans $|\overline{Z}_{i, j}|$, voir
Morphismes d'espaces, Lemme
\ref{spaces-morphisms-lemma-quasi-compact-immersion}.
Après avoir remplacé $X_i$ par un éclatement $V_i$-admissible, nous pouvons supposer que $\overline{Z}_{i, i}$ et $\overline{Z}_{i, j}$ sont disjoints, voir
Lemme \ref{spaces-pushouts-lemma-separate-disjoint-locally-closed-by-blowing-up}.
Nous supposons que cela vaut pour $X_1$ et $X_2$.
Remarquons que cette propriété est préservée si nous remplaçons $X_i$ par un autre éclatement $V_i$-admissible. Par conséquent, nous pouvons remplacer $X_1$ par un autre éclatement $V_1$-admissible et supposer que $|\overline{Z}_{1, 1}|$ est la réunion disjointe des adhérences de $|T|$ et $|f^{-1}Z_1|$ dans $|X_1|$.

\medskip\noindent
Soit $V_{12} = V_1 \times_U V_2$. Nous avons une immersion
$V_{12} \to X_1 \times_B X_2$ qui est la composée de l'immersion fermée $V_{12} = V_1 \times_U V_2 \to V_1 \times_B V_2$
(Morphismes d'espaces, Lemme \ref{spaces-morphisms-lemma-fibre-product-after-map})
et de l'immersion ouverte $V_1 \times_B V_2 \to X_1 \times_B X_2$.
Soit $X_{12} \subset X_1 \times_B X_2$ l'image schématique de $V_{12} \to X_1 \times_B X_2$. Les morphismes de projection
$$
p_1 : X_{12} \to X_1
\quad\text{et}\quad
p_2 : X_{12} \to X_2
$$
sont propres car $X_1$ et $X_2$ sont propres sur $B$. Si nous remplaçons $X_1$ par un
éclatement $V_1$-admissible, alors $X_{12}$ est remplacé par
la transformée stricte par rapport à cet éclatement, voir
Lemme \ref{spaces-pushouts-lemma-replaced-by-strict-transform}.

\medskip\noindent  
Notons $\psi : V_{12} \to U$ les composées  
$\psi = \psi_1 \circ p_1|_{V_{12}} = \psi_2 \circ p_2|_{V_{12}}$.  
Considérons le sous-espace fermé  
$$
Z_{12, 2} =
(p_1|_{V_{12}})^{-1}Z_{1, 2} =
(p_2|_{V_{12}})^{-1}Z_{2, 2} =
\psi^{-1}Z_2 \subset V_{12}
$$.  
Le morphisme $p_1|_{V_{12}} : V_{12} \to V_1$ est un isomorphisme  
sur un voisinage ouvert de $Z_{1, 2}$ parce que $\psi_2 : V_2 \to U$  
est un isomorphisme sur un voisinage ouvert de $Z_2$ et  
$V_{12} = V_1 \times_U V_2$. Par le Lemme \ref{spaces-pushouts-lemma-blowup-iso-along}  
il existe un éclatement $V_1$-admissible $X_1' \to X_1$  
tel que la transformée stricte $p'_1 : X'_{12} \to X'_1$  
de $p_1$ soit un isomorphisme sur un voisinage ouvert de  
l'adhérence de $|Z_{1, 2}|$ dans $|X'_1|$.  
Après avoir remplacé $X_1$ par $X'_1$ et $X_{12}$ par $X'_{12}$  
nous pouvons supposer que $p_1$ est un isomorphisme sur un voisinage ouvert de $|\overline{Z}_{1, 2}|$.

\medskip\noindent
Le résultat du paragraphe précédent nous dit que
$$
X_{12} \cap (\overline{Z}_{1, 2} \times_B \overline{Z}_{2, 1}) = \emptyset
$$
où l'intersection est prise dans $X_1 \times_B X_2$. En effet, l'image inverse
$p_1^{-1}\overline{Z}_{1, 2}$ dans $X_{12}$ s'envoie isomorphiquement
à $\overline{Z}_{1, 2}$. En particulier, nous voyons que $|Z_{12, 2}|$
est dense dans $|p_1^{-1}\overline{Z}_{1, 2}|$. Ainsi, $p_2$ envoie
$|p_1^{-1}\overline{Z}_{1, 2}|$ dans $|\overline{Z}_{2, 2}|$.
Puisque $|\overline{Z}_{2, 2}| \cap |\overline{Z}_{2, 1}| = \emptyset$
nous concluons.

\medskip\noindent
Il s'avère que nous devons effectuer un éclatement supplémentaire avant de pouvoir conclure l'argument. En effet, soit $V_2 \subset W_2 \subset X_2$ le sous-espace ouvert dont l'espace topologique sous-jacent est $$
|W_2| =
|V_2| \cup (|X_2| \setminus |\overline{Z}_{2, 1}|) =
|X_2| \setminus \left(|\overline{Z}_{2, 1}| \setminus |Z_{2, 1}|\right)
$$. Puisque $p_2(p_1^{-1}\overline{Z}_{1, 2})$ est contenu dans $W_2$ (voir ci-dessus), nous voyons que remplacer $X_2$ par un éclatement $W_2$-admissible et $X_{12}$ par la transformée stricte correspondante préserve la propriété de $p_1$ d'être un isomorphisme sur un voisinage ouvert de $\overline{Z}_{1, 2}$. Puisque $\overline{Z}_{2, 1} \cap W_2 = \overline{Z}_{2, 1} \cap V_2 = Z_{2, 1}$, nous voyons que $Z_{2, 1}$ est un sous-espace fermé de $W_2$ et $V_2$. Remarquons que $V_{12} = V_1 \times_U V_2 = p_1^{-1}(V_1) = p_2^{-1}(V_2)$ comme sous-espaces ouverts de $X_{12}$ puisqu'il s'agit du plus grand sous-espace ouvert de $X_{12}$ sur lequel le morphisme $\psi : V_{12} \to U$ se prolonge ; détails omis\footnote{En effet, $V_1 \times_U V_2$ est propre sur $U$ donc si $\psi$ s'étend à un ouvert plus grand de $X_{12}$, alors $V_1 \times_U V_2$ serait fermé dans cet ouvert par Morphismes d'espaces, Lemme \ref{spaces-morphisms-lemma-universally-closed-permanence}. Ensuite, nous obtenons l'égalité car $V_{12} \subset X_{12}$ est dense.}. Nous avons les égalités suivantes de sous-espaces fermés de $V_{12}$ : $$
p_2^{-1}Z_{2, 1} = p_2^{-1} \psi_2^{-1} Z_1 =
p_1^{-1} \psi_1^{-1} Z_1= p_1^{-1}Z_{1, 1} =
p_1^{-1}f^{-1}Z_1 \amalg p_1^{-1}T
$$. Ici et ci-dessous, nous utilisons le léger abus de notation consistant à écrire $p_2$ au lieu de la restriction de $p_2$ à $V_{12}$, etc. Puisque $p_2^{-1}(Z_{2, 1})$ est un sous-espace fermé de $p_2^{-1}(W_2)$ et que $Z_{2, 1}$ est un sous-espace fermé de $W_2$, nous concluons que $p_1^{-1}f^{-1}Z_1$ est également un sous-espace fermé de $p_2^{-1}(W_2)$. Enfin, le morphisme $p_2 : X_{12} \to X_2$ est \'etale aux points de $p_1^{-1}f^{-1}Z_1$ car $\psi_1$ est \'etale le long de $f^{-1}Z_1$ et $V_{12} = V_1 \times_U V_2$. Ainsi, nous pouvons appliquer le Lemme \ref{spaces-pushouts-lemma-blowup-etale-along} au morphisme $p_2 : X_{12} \to X_2$, à l'ouvert $W_2$, au sous-espace fermé $Z_{2, 1} \subset W_2$ et au sous-espace fermé $p_1^{-1}f^{-1}Z_1 \subset p_2^{-1}(W_2)$. Par conséquent, après avoir remplacé $X_2$ par un éclatement $W_2$-admissible et $X_{12}$ par la transformée stricte correspondante, nous obtenons pour chaque point géométrique $\overline{y}$ de l'adhérence de $|p_1^{-1}f^{-1}Z_1|$ un homomorphisme local d'anneaux $\mathcal{O}_{X_{12}, \overline{y}} \to \mathcal{O}$ tel que $\mathcal{O}_{X_2, p_2(\overline{y})} \to \mathcal{O}$ soit fini et plat.

\medskip\noindent
Considérons l'espace algébrique
$$
W_2 = U \coprod\nolimits_{U_2} (X_2 \setminus \overline{Z}_{2, 1}),
$$
et, avec $T \subset V_1$ comme dans le premier paragraphe, l'espace algébrique
$$
W_1 = U \coprod\nolimits_{U_1}
(X_1 \setminus \overline{Z}_{1, 2} \cup \overline{T}),
$$
Ces deux espaces sont obtenus par somme amalgamée, voir
Lemme \ref{spaces-pushouts-lemma-construct-elementary-distinguished-square}.
Appliquons le Lemme \ref{spaces-pushouts-lemma-check-separated}
pour voir que $W_i \to B$ est séparé. Tout d'abord,
$U \to B$ et $X_i \to B$ sont séparés. Vérifions que l'immersion quasi-compacte $U_i \to U \times_B (X_i \setminus \overline{Z}_{i, j})$
est fermée en utilisant le critère valuatif, voir Morphismes d'espaces, Lemme
\ref{spaces-morphisms-lemma-quasi-compact-existence-universally-closed}.
Choisissons un anneau de valuation $A$ sur $B$ avec corps des fractions $K$ et
morphismes compatibles $(u, x_i) : \Spec(A) \to U \times_B X_i$ et
$u_i : \Spec(K) \to U_i$. Puisque $\psi_i$ est propre, nous
pouvons trouver un unique $v_i : \Spec(A) \to V_i$ compatible avec
$u$ et $u_i$. Puisque $X_i$ est propre sur $B$ nous voyons que $x_i = v_i$. Si $v_i$ ne se factorise pas par $U_i \subset V_i$, alors nous concluons que $x_i$ envoie le point fermé de $\Spec(A)$ dans $Z_{i, j}$ ou $T$ lorsque $i = 1$. Cela termine la démonstration car nous avons supprimé $\overline{Z}_{i, j}$ et $\overline{T}$ dans la construction de $W_i$.

\medskip\noindent
D'autre part, pour tout anneau de valuation $A$ sur $B$ avec
corps des fractions $K$ et tout morphisme
$$
\gamma : \Spec(K) \to \Im(U_1 \times_U U_2 \to U)
$$
sur $B$, nous affirmons qu'après avoir remplacé $A$ par une extension d'anneaux de valuation, il existe un $i$ et un prolongement de $\gamma$ en un morphisme $h_i : \Spec(A) \to W_i$. En effet, nous étendons d'abord $\gamma$ à un morphisme $g_2 : \Spec(A) \to X_2$ en utilisant le critère valuatif de propreté. Si l'image de $g_2$ ne rencontre pas $\overline{Z}_{2, 1}$, alors nous obtenons notre morphisme vers $W_2$.
Sinon, notons $\overline{z} \in \overline{Z}_{2, 1}$ un point géométrique situé au-dessus de l'image par $g_2$ du point fermé. Nous pouvons le relever en un point géométrique $\overline{y}$ de $X_{12}$
dans l'adhérence de $|p_1^{-1}f^{-1}Z_1|$ car l'application continue $|p_1^{-1}f^{-1}Z_1| \to |\overline{Z}_{2, 1}|$ est fermée, avec une image contenant l'ouvert dense $|Z_{2, 1}|$. Après avoir remplacé $A$ par sa hensélisation stricte (Compléments d'algèbre, Lemme \ref{more-algebra-lemma-henselization-valuation-ring}), nous obtenons le diagramme suivant $$
\xymatrix{
A \ar@{..>}[rr] & & A' \\
\mathcal{O}_{X_2, \overline{z}} \ar[r] \ar[u] &
\mathcal{O}_{X_{12}, \overline{y}} \ar[r] &
\mathcal{O} \ar@{..>}[u]
}
$$ où $\mathcal{O}_{X_{12}, \overline{y}} \to \mathcal{O}$ est l'homomorphisme que nous avons trouvé au cinquième paragraphe de la preuve. Puisque la composée horizontale est finie et plate, nous pouvons trouver une extension d'anneaux de valuation $A'/A$ et une flèche pointillée rendant le diagramme commutatif. Après avoir remplacé $A$ par $A'$, cela signifie que nous obtenons un relèvement $g_{12} : \Spec(A) \to X_{12}$ dont le point fermé s'envoie dans l'adhérence de $|p_1^{-1}f^{-1}Z_1|$. Alors $g_1 = p_1 \circ g_{12} : \Spec(A) \to X_1$ est un morphisme dont le point fermé s'envoie dans l'adhérence de $|f^{-1}Z_1|$. Puisque l'adhérence de $|f^{-1}Z_1|$ est disjointe de l'adhérence de $|T|$ et contenue dans $|\overline{Z}_{1, 1}|$ qui est disjoint de $|\overline{Z}_{1, 2}|$, nous en concluons que $g_1$ définit un morphisme $h_1 : \Spec(A) \to W_1$ comme voulu.

\medskip\noindent
Considérons un diagramme
$$
\xymatrix{
W_1' \ar[d] \ar[r] & W & W_2' \ar[l] \ar[d] \\
W_1 & U \ar[l] \ar[lu] \ar[u] \ar[ru] \ar[r] & W_2
}
$$
comme dans Compléments sur les morphismes d'espaces, Lemme
\ref{spaces-more-morphisms-lemma-find-common-blowups}.
D'après le paragraphe précédent, pour tout diagramme dont les flèches pleines forment le diagramme commutatif
$$
\xymatrix{
\Spec(K) \ar[r]_\gamma  \ar[d] & W \ar[d] \\
\Spec(A) \ar@{..>}[ru] \ar[r] & B
}
$$
où $\Im(\gamma) \subset \Im(U_1 \times_U U_2 \to U)$,
il existe un $i$ et une extension $h_i : \Spec(A) \to W_i$ de $\gamma$
après avoir éventuellement remplacé $A$ par une extension d'anneaux de valuation.
En utilisant le critère valuatif de propreté pour $W'_i \to W_i$,
nous pouvons alors relever $h_i$ en $h'_i : \Spec(A) \to W'_i$.
Ainsi, la flèche en pointillés dans le diagramme existe après
avoir éventuellement étendu $A$. Puisque $W$ est séparé sur $B$,
nous voyons que le choix de l'extension n'est pas nécessaire et que la flèche est également unique, voir
Morphismes d'espaces, Lemmes \ref{spaces-morphisms-lemma-usual-enough}
et \ref{spaces-morphisms-lemma-separated-implies-valuative}.
Enfin, l'existence de la flèche en pointillés
implique que $W \to B$ est universellement fermé par Morphismes d'espaces, Lemme \ref{spaces-morphisms-lemma-refined-valuative-criterion-universally-closed}. Comme $W \to B$ est déjà de type fini et séparé, le résultat s'ensuit. 
\end{proof}

\begin{lemma}
\label{spaces-pushouts-lemma-filter-Noetherian-space}
Soit $S$ un schéma. Soit $X$ un espace algébrique noethérien sur $S$.
Soit $U \subset X$ un sous-espace ouvert dense, proprement contenu dans l'espace ambiant. Alors il existe un
schéma affine $V$ et un morphisme \'etale $V \to X$ tel que
\begin{enumerate}
\item le sous-espace ouvert $W = U \cup \Im(V \to X)$ est strictement plus grand
que $U$,
\item $(U \subset W, V \to W)$ est un carré distingué, et
\item $U \times_W V \to U$ a une image dense.
\end{enumerate}
\end{lemma}

\begin{proof}
Choisissons une stratification
$$
\emptyset = U_{n + 1} \subset
U_n \subset U_{n - 1} \subset \ldots \subset U_1 = X
$$
et des morphismes $f_p : V_p \to U_p$ comme dans Espaces décents, Lemme
\ref{decent-spaces-lemma-filter-quasi-compact-quasi-separated}.
Soit $p$ le plus petit entier tel que $U_p \not \subset U$
(cela est possible car $U \not = X$). Choisissons un ouvert affine $V \subset V_p$
tel que le morphisme \'etale $f = f_p|_V : V \to X$ ne se factorise pas par $U$.
Considérons l'ouvert $W = U \cup \Im(V \to X)$ et notons $Z \subset W$ le sous-espace fermé réduit tel que $|Z| = |W| \setminus |U|$.
Alors $f^{-1}Z \to Z$ est un isomorphisme car nous avons la
propriété correspondante pour le morphisme $f_p$, voir le lemme cité ci-dessus.
Ainsi $(U \subset W, f : V \to W)$ est un carré distingué.
Il se peut que l'ouvert $I = \Im(U \times_W V \to U)$
ne soit pas dense dans $U$. L'espace algébrique $U' \subset U$ dont l'ensemble
sous-jacent est $|U| \setminus \overline{|I|}$ est noethérien, et donc nous pouvons trouver un sous-schéma ouvert dense $U'' \subset U'$, voir par exemple Propriétés des espaces, Proposition \ref{spaces-properties-proposition-locally-quasi-separated-open-dense-scheme}. Ensuite, nous pouvons trouver un ouvert affine dense $U''' \subset U''$, voir Propriétés, Lemme \ref{properties-lemma-Noetherian-irreducible-components} et \ref{properties-lemma-maximal-points-affine}. Après avoir remplacé $f$ par $V \amalg U''' \to X$, tout est clair. 
\end{proof}

\begin{theorem}
\label{spaces-pushouts-theorem-nagata}
\begin{reference}
\cite{CLO}
\end{reference}
Soit $S$ un schéma. Soit $B$ un espace algébrique quasi-compact et
quasi-séparé sur $S$. Soit $X \to B$ un morphisme séparé de type fini.
Alors $X$ admet une compactification sur $B$.
\end{theorem}

\begin{proof}
Nous réduisons d'abord au cas noethérien. Nous conseillons vivement au lecteur de passer ce paragraphe. Tout d'abord, nous pouvons remplacer $S$ par $\Spec(\mathbf{Z})$ ; voir Espaces, section \ref{spaces-section-change-base-scheme}, et Propriétés des espaces, Définition \ref{spaces-properties-definition-separated}. Il existe une immersion fermée $X \to X'$ avec $X' \to B$ de présentation finie et séparée ; voir Limites d'espaces, Proposition \ref{spaces-limits-proposition-separated-closed-in-finite-presentation}. Si nous trouvons une compactification de $X'$ sur $B$, alors l'adhérence schématique de $X$ dans celle-ci donnera une compactification de $X$ sur $B$. Ainsi, nous pouvons supposer que $X \to B$ est séparé et de présentation finie. Nous pouvons écrire $B = \lim B_i$ comme une limite projective d'un système filtrant d'espaces algébriques noethériens de type fini sur $\Spec(\mathbf{Z})$, à morphismes de transition affines ; voir Limites des espaces, Proposition \ref{spaces-limits-proposition-approximate}.  
Nous pouvons choisir un $i$ et un morphisme $X_i \to B_i$ de présentation finie dont le changement de base à $B$ est $X \to B$, voir Limites des espaces, Lemme \ref{spaces-limits-lemma-descend-finite-presentation}.  
Après avoir remplacé $i$ par un indice plus grand, nous pouvons supposer que $X_i \to B_i$ est séparé, voir Limites des espaces, Lemme \ref{spaces-limits-lemma-descend-separated-morphism}.  
Si nous pouvons trouver une compactification de $X_i$ sur $B_i$, alors son changement de base à $B$ sera une compactification de $X$ sur $B$.  
Cela nous ramène au cas discuté dans le paragraphe suivant.

\medskip\noindent
Supposons que $B$ soit de type fini sur $\mathbf{Z}$ en plus d'être quasi-compact et quasi-séparé. Soit $U \to X$ un morphisme \'etale d'espaces algébriques tel que $U$ possède une compactification $Y$ sur $\Spec(\mathbf{Z})$. Le morphisme $$
U \longrightarrow B \times_{\Spec(\mathbf{Z})} Y
$$ est séparé et quasi-fini d'après Morphismes d'espaces, Lemme \ref{spaces-morphisms-lemma-monomorphism-loc-finite-type-loc-quasi-finite} (le morphisme affiché se factorise en une immersion, donc c'est un monomorphisme). Par conséquent, d'après le théorème principal de Zariski (Compléments sur les morphismes d'espaces, Lemme \ref{spaces-more-morphisms-lemma-quasi-finite-separated-pass-through-finite}), il existe une immersion ouverte de $U$ dans un espace algébrique $Y'$ fini sur $B \times_{\Spec(\mathbf{Z})} Y$. Alors $Y' \to B$ est propre comme composée $Y' \to B \times_{\Spec(\mathbf{Z})} Y \to B$ de deux morphismes propres (utiliser Morphismes d'espaces, Lemmes \ref{spaces-morphisms-lemma-finite-proper}, \ref{spaces-morphisms-lemma-composition-proper} et \ref{spaces-morphisms-lemma-base-change-proper}). Nous concluons que $U$ a une compactification sur $B$.

\medskip\noindent
Il existe un sous-espace ouvert et dense $U \subset X$ qui est un schéma
(Propriétés des espaces, Proposition \ref{spaces-properties-proposition-locally-quasi-separated-open-dense-scheme}).
En fait, nous pouvons choisir $U$ de façon que ce soit un schéma affine
(Propriétés, Lemme \ref{properties-lemma-Noetherian-irreducible-components} et \ref{properties-lemma-maximal-points-affine}).
Ainsi, $U$ possède une compactification sur $\Spec(\mathbf{Z})$;
cela se démontre facilement directement mais découle aussi du
théorème pour les schémas, voir
Compléments sur la platitude, Théorème \ref{flat-theorem-nagata}.
D'après le paragraphe précédent, $U$ possède une compactification sur $B$.
Par induction noethérienne, nous pouvons trouver un sous-espace ouvert et dense maximal
$U \subset X$ qui possède une compactification sur $B$. Nous allons montrer
que l'hypothèse que $U \not = X$ mène à une contradiction.
En effet, d'après le Lemme \ref{spaces-pushouts-lemma-filter-Noetherian-space}
nous pouvons trouver un ouvert strictement plus grand $U \subset W \subset X$ et un carré distingué $(U \subset W, f : V \to W)$ avec $V$ affine et $U \times_W V$ ayant une image dense dans $U$. Puisque $V$ est affine, comme précédemment, il possède une compactification sur $B$. Ainsi, le Lemme \ref{spaces-pushouts-lemma-two-compactifications} s'applique pour montrer que $W$ possède une compactification sur $B$, ce qui est la contradiction désirée. 
\end{proof}









% OK, start here.
%

\chapter{Groupes de Chow des espaces}


\phantomsection
\label{spaces-chow-section-phantom}




\section{Introduction}
\label{spaces-chow-section-introduction}

\noindent
Dans ce chapitre, nous étudions d'abord les groupes de Chow des espaces
algébriques. Après les avoir définis, nous définissons les classes de Chern des
fibrés vectoriels comme des opérateurs sur ces groupes de Chow. La stratégie
est exactement la même que dans le cas des schémas. Nous invitons donc le
lecteur à consulter l'introduction du chapitre correspondant pour les schémas
(Homologie de Chow, section \ref{chow-section-introduction}).

\medskip\noindent
Quelques articles sur le sujet : \cite{edidin-graham} et \cite{kresch_cycle}.



\section{Cadre}
\label{spaces-chow-section-setup}

\noindent
Nous fixons d'abord la catégorie des espaces algébriques que nous considérerons.
Il faut garder à l'esprit, tout au long de ce chapitre, que
« décent $+$ localement noethérien » équivaut à « quasi-séparé $+$ localement
noethérien » d'après Espaces décents, lemme
\ref{decent-spaces-lemma-locally-Noetherian-decent-quasi-separated}.

\begin{situation}
\label{spaces-chow-situation-setup}
Ici, $S$ est un schéma et $B$ un espace algébrique sur $S$. Nous supposons
que $B$ est quasi-séparé, localement noethérien et universellement caténaire
(Espaces décents, définition
\ref{decent-spaces-definition-universally-catenary}). De plus, nous supposons
donnée une fonction de dimension $\delta : |B| \longrightarrow
\mathbf{Z}$. On dit que $X/B$ est {\it bon} si $X$ est un espace algébrique
sur $B$ dont le morphisme de structure $f : X \to B$ est quasi-séparé et
localement de type fini. Dans ce cas, nous définissons
$$
\delta = \delta_{X/B} : |X| \longrightarrow \mathbf{Z}
$$
comme l'application qui à $x$ associe $\delta(f(x))$ plus le degré de
transcendance de $x/f(x)$ (Morphismes d'espaces, définition
\ref{spaces-morphisms-definition-dimension-fibre}). Il s'agit d'une fonction
de dimension d'après Compléments sur les morphismes d'espaces, lemme
\ref{spaces-more-morphisms-lemma-universally-catenary-dimension-function}.
\end{situation}

\noindent
Un cas particulier est celui où $S = B$ est un schéma et $(S, \delta)$ vérifie
les hypothèses d'Homologie de Chow, situation \ref{chow-situation-setup}.
Ainsi, $B$ peut être le spectre d'un corps (Homologie de Chow, exemple
\ref{chow-example-field}) ou $B = \Spec(\mathbf{Z})$ (Homologie de Chow, exemple
\ref{chow-example-domain-dimension-1}).

\medskip\noindent
De nombreux lemmes, propositions, théorèmes et définitions concernant les
espaces algébriques sont plus faciles à établir dans le cadre de la situation
\ref{spaces-chow-situation-setup} car les espaces algébriques avec lesquels nous
travaillons sont quasi-séparés (et donc a fortiori décents) et localement
noethériens. Nous le signalerons au fil du chapitre par des remarques comme
les suivantes.

\begin{remark}
\label{spaces-chow-remark-sober}
Dans la situation \ref{spaces-chow-situation-setup}, si $X/B$ est bon, alors $|X|$ est un
espace topologique sobre. Voir Propriétés des espaces, lemme
\ref{spaces-properties-lemma-quasi-separated-sober} ou Espaces décents,
proposition \ref{decent-spaces-proposition-reasonable-sober}. Nous
l'utiliserons sans autre mention pour choisir des points génériques de
sous-ensembles fermés irréductibles de $|X|$.
\end{remark}

\begin{remark}
\label{spaces-chow-remark-integral}
Dans la situation \ref{spaces-chow-situation-setup}, si $X/B$ est bon, alors $X$ est
intègre (Espaces sur les corps, définition
\ref{spaces-over-fields-definition-integral-algebraic-space}) si et seulement
si $X$ est réduit et $|X|$ est irréductible. De plus, pour tout point $\xi \in
|X|$, il existe un sous-espace fermé intègre unique $Z \subset X$ tel que
$\xi$ soit le point générique du sous-ensemble fermé $|Z| \subset |X|$ ; voir
Espaces sur les corps, lemme
\ref{spaces-over-fields-lemma-decent-irreducible-closed}.
\end{remark}

\noindent
Si $B$ est jacobsonien et que $\delta$ envoie les points fermés sur zéro, alors
$\delta$ est la fonction qui envoie un point à la dimension de sa fermeture.

\begin{lemma}
\label{spaces-chow-lemma-delta-is-dimension}
Dans la situation \ref{spaces-chow-situation-setup}, supposons que $B$ soit jacobsonien et que
$\delta(b) = 0$ pour chaque point fermé $b$ de $|B|$. Soit $X/B$ bon. Si
$Z \subset X$ est un sous-espace fermé intègre de point générique $\xi \in
|Z|$, alors les entiers suivants sont égaux :
\begin{enumerate}
\item $\delta(\xi) = \delta_{X/B}(\xi)$,
\item $\dim(|Z|)$,
\item $\text{codim}(\{z\}, |Z|)$ pour $z \in |Z|$ fermé,
\item la dimension de l'anneau local de $Z$ à $z$ pour $z \in |Z|$ fermé, et
\item $\dim(\mathcal{O}_{Z, \overline{z}})$ pour $z \in |Z|$ fermé.
\end{enumerate}
\end{lemma}

\begin{proof}
Soient $X$, $Z$, $\xi$ comme dans le lemme. Puisque $X$ est localement de type
fini sur $B$, l'espace $X$ est jacobsonien ; voir Espaces décents, lemme
\ref{decent-spaces-lemma-Jacobson-universally-Jacobson}. Ainsi
$X_{\text{ft-pts}} \subset |X|$ est l'ensemble des points fermés d'après
Espaces décents, lemme \ref{decent-spaces-lemma-decent-Jacobson-ft-pts}. Étant donné
une chaîne $T_0 \supset \ldots \supset T_e$ de sous-ensembles fermés
irréductibles de $|Z|$, l'ensemble $T_e \cap X_{\text{ft-pts}}$ est non vide
d'après Morphismes d'espaces, lemme
\ref{spaces-morphisms-lemma-enough-finite-type-points}. Ainsi on peut toujours
supposer qu'une telle chaîne se termine par $T_e = \{z\}$ pour un point fermé
$z \in |Z|$. Il s'ensuit que $\dim(Z) = \sup_z \text{codim}(\{z\}, |Z|)$, où
$z$ court sur les points fermés de $|Z|$. Nous avons $\text{codim}(\{z\}, Z) =
\delta(\xi) - \delta(z)$ d'après Topologie, lemme
\ref{topology-lemma-dimension-function-catenary}. Par Morphismes d'espaces,
lemme \ref{spaces-morphisms-lemma-finite-type-points-morphism}, l'image de $z$
est un point de type fini de $B$, c'est-à-dire un point fermé de $|B|$. Par
Morphismes d'espaces, lemme
\ref{spaces-morphisms-lemma-jacobson-finite-type-points}, le degré de
transcendance de $z/b$ est $0$. Par hypothèse, nous concluons que
$\delta(z) = \delta(b) = 0$. On obtient donc l'égalité
$$
\dim(|Z|) = \text{codim}(\{z\}, Z) = \delta(\xi)
$$
pour tout point fermé $z \in |Z|$. Enfin, $\text{codim}(\{z\}, Z)$ est égal à
la dimension de l'anneau local de $Z$ en $z$ d'après Espaces décents, lemme
\ref{decent-spaces-lemma-codimension-local-ring}, laquelle est à son tour égale
à $\dim(\mathcal{O}_{Z, \overline{z}})$ d'après Propriétés des espaces, lemme
\ref{spaces-properties-lemma-dimension-local-ring}.
\end{proof}

\noindent
Dans la situation du lemme précédent, la valeur de $\delta$ au point
générique d'un sous-ensemble fermé irréductible est la dimension de ce
sous-ensemble. Cela motive la définition suivante.

\begin{definition}
\label{spaces-chow-definition-delta-dimension}
Dans la situation \ref{spaces-chow-situation-setup}, pour tout $X/B$ bon et tout
sous-ensemble fermé irréductible $T \subset |X|$, on définit
$$
\dim_\delta(T) = \delta(\xi)
$$
où $\xi \in T$ est le point générique de $T$. Nous appelons ce nombre la {\it
$\delta$-dimension de $T$}. Si $T \subset |X|$ est un sous-ensemble fermé,
alors nous définissons $\dim_\delta(T)$ comme la borne supérieure des
$\delta$-dimensions des composantes irréductibles de $T$. Si $Z$ est un sous-espace fermé
de $X$, alors nous définissons $\dim_\delta(Z) = \dim_\delta(|Z|)$.
\end{definition}

\noindent
Bien sûr, cela signifie simplement que $\dim_\delta(T) = \sup \{\delta(t) \mid
t \in T\}$.







\section{Cycles}
\label{spaces-chow-section-cycles}

\noindent
Cette section est l'analogue d'Homologie de Chow, section
\ref{chow-section-cycles}.

\medskip\noindent
Puisque nous ne supposons pas nos espaces quasi-compacts, il faut être un peu
prudent dans la définition des cycles. Nous devons autoriser les sommes infinies,
car une fonction rationnelle peut, par exemple, avoir une infinité de pôles.
En revanche, si $X$ est quasi-compact, alors un cycle est
une somme finie comme d'habitude.

\begin{definition}
\label{spaces-chow-definition-cycles}
Dans la situation \ref{spaces-chow-situation-setup}, soit $X/B$ bon. Soit $k
\in \mathbf{Z}$.
\begin{enumerate}
\item Un {\it cycle sur $X$} est une somme formelle
$$
\alpha = \sum n_Z [Z]
$$
où la somme s'étend sur les sous-espaces fermés intègres $Z \subset X$,
chaque $n_Z \in \mathbf{Z}$ et $\{|Z|; n_Z \not = 0\}$ est une famille
localement finie de sous-ensembles de $|X|$ (Topologie, définition
\ref{topology-definition-locally-finite}).
\item Un {\it $k$-cycle} sur $X$ est un cycle
$$
\alpha = \sum n_Z [Z]
$$
où $n_Z \not = 0 \Rightarrow \dim_\delta(Z) = k$.
\item Le groupe abélien de tous les $k$-cycles sur $X$ est noté $Z_k(X)$.
\end{enumerate}
\end{definition}

\noindent
En d'autres termes, un $k$-cycle sur $X$ est une combinaison formelle
$\mathbf{Z}$-linéaire localement finie de sous-espaces fermés intègres (remarque
\ref{spaces-chow-remark-integral}) de $\delta$-dimension $k$. L'addition des $k$-cycles
$\alpha = \sum n_Z[Z]$ et $\beta = \sum m_Z[Z]$ est donné par
$$
\alpha + \beta = \sum (n_Z + m_Z)[Z],
$$
c'est-à-dire en additionnant les coefficients.




\section{Multiplicités}
\label{spaces-chow-section-multiplicities}

\noindent
Cette section rassemble quelques résultats simples sur les longueurs et les
multiplicités.

\begin{lemma}
\label{spaces-chow-lemma-length}
Soit $S$ un schéma et soit $X$ un espace algébrique sur $S$. Soit $\mathcal{F}$
un $\mathcal{O}_X$-module quasi-cohérent. Soit $x \in |X|$. Soit $d \in
\{0, 1, 2, \ldots, \infty\}$. Les assertions suivantes sont équivalentes :
\begin{enumerate}
\item
$\text{length}_{\mathcal{O}_{X, \overline{x}}} \mathcal{F}_{\overline{x}} = d$
\item il existe un morphisme \'etale $U \to X$, avec $U$ un schéma, et un
point $u \in U$ d'image $x$, tels que $\text{length}_{\mathcal{O}_{U, u}}
(\mathcal{F}|_U)_u = d$
\item pour tout morphisme \'etale $U \to X$, avec $U$ un schéma, et tout
point $u \in U$ d'image $x$, on a $\text{length}_{\mathcal{O}_{U, u}}
(\mathcal{F}|_U)_u = d$
\end{enumerate}
\end{lemma}

\begin{proof}
Soient $U \to X$ et $u \in U$ comme en (2) ou (3). On sait alors que
$\mathcal{O}_{X, \overline{x}}$ est l'hensélisation stricte de
$\mathcal{O}_{U, u}$ et que
$$
\mathcal{F}_{\overline{x}} =
(\mathcal{F}|_U)_u \otimes_{\mathcal{O}_{U, u}} \mathcal{O}_{X, \overline{x}}
$$
Voir Propriétés des espaces, lemmes
\ref{spaces-properties-lemma-describe-etale-local-ring} et
\ref{spaces-properties-lemma-stalk-quasi-coherent}. Il s'ensuit, par Algèbre,
lemme \ref{algebra-lemma-pullback-module}, par la platitude de
$\mathcal{O}_{U, u} \to \mathcal{O}_{X, \overline{x}}$ et par le fait que
$\mathcal{O}_{X, \overline{x}}/\mathfrak m_u\mathcal{O}_{X,
\overline{x}}$ est le corps résiduel de $\mathcal{O}_{X, \overline{x}}$, que
les longueurs sont égales. Ces propriétés des hensélisations strictes se
trouvent dans Compléments d'algèbre, lemme
\ref{more-algebra-lemma-dumb-properties-henselization}.
\end{proof}

\begin{definition}
\label{spaces-chow-definition-length-at-x}
Soit $S$ un schéma et soit $X$ un espace algébrique sur $S$. Soit $\mathcal{F}$
un $\mathcal{O}_X$-module quasi-cohérent. Soit $x \in |X|$. Soit $d \in
\{0, 1, 2, \ldots, \infty\}$. Nous disons que {\it $\mathcal{F}$ est de longueur
$d$ en $x$} si les conditions équivalentes du lemme \ref{spaces-chow-lemma-length} sont
satisfaites.
\end{definition}

\begin{lemma}
\label{spaces-chow-lemma-length-closed-immersion}
Soit $S$ un schéma. Soit $i : Y \to X$ une immersion fermée d'espaces
algébriques sur $S$. Soit $\mathcal{G}$ un $\mathcal{O}_Y$-module
quasi-cohérent. Soit $y \in |Y|$ d'image $x \in |X|$. Soit $d \in \{0,
1, 2, \ldots, \infty\}$. Les assertions suivantes sont équivalentes :
\begin{enumerate}
\item $\mathcal{G}$ a une longueur $d$ à $y$, et
\item $i_*\mathcal{G}$ a une longueur $d$ à $x$.
\end{enumerate}
\end{lemma}

\begin{proof}
Choisissons un morphisme \'etale $f : U \to X$, où $U$ est un schéma, et un
point $u \in U$ d'image $x$. Posons $V = Y \times_X U$. Notons $g : V \to Y$
et $j : V \to U$ les projections. Alors $j : V \to U$ est une immersion fermée
et il existe un unique point $v \in V$ d'images $y \in |Y|$ et $u \in U$
(voir Propriétés des espaces, lemme
\ref{spaces-properties-lemma-points-cartesian}, et Espaces, lemme
\ref{spaces-lemma-base-change-immersions}). Nous avons $j_*(\mathcal{G}|_V) =
(i_*\mathcal{G})|_U$ en tant que modules sur le schéma $V$, et $j_*$ est
l'image directe « usuelle » des modules pour le morphisme de schémas $j$ ;
voir la discussion qui entoure Cohomologie des espaces, équation
(\ref{spaces-cohomology-equation-representable-higher-direct-image}). On se
ramène ainsi au cas des schémas : si $i : Y \to X$ est une immersion fermée de
schémas, alors
$$
(i_*\mathcal{G})_x = \mathcal{G}_y
$$
en tant que modules sur $\mathcal{O}_{X, x}$, la structure de module au membre
de droite étant donnée par la surjection $i_y^\sharp : \mathcal{O}_{X, x} \to
\mathcal{O}_{Y, y}$. L'égalité résulte donc d'Algèbre, lemme
\ref{algebra-lemma-length-independent}.
\end{proof}

\begin{lemma}
\label{spaces-chow-lemma-length-finite}
Soit $S$ un schéma et soit $X$ un espace algébrique localement noethérien sur
$S$. Soit $\mathcal{F}$ un $\mathcal{O}_X$-module cohérent. Soit $x \in |X|$.
Les assertions suivantes sont équivalentes :
\begin{enumerate}
\item il existe un morphisme \'etale $U \to X$, avec $U$ un schéma, et un point
$u \in U$ d'image $x$, tel que $u$ soit un point générique d'une composante
irréductible de $\text{Supp}(\mathcal{F}|_U)$,
\item pour tout morphisme \'etale $U \to X$, avec $U$ un schéma, et tout point
$u \in U$ d'image $x$, le point $u$ est générique dans une composante irréductible
de $\text{Supp}(\mathcal{F}|_U)$,
\item la longueur de $\mathcal{F}$ à $x$ est finie et non nulle.
\end{enumerate}
Si $X$ est décent (de manière équivalente, quasi-séparé), ces assertions sont
également équivalentes à
\begin{enumerate}
\item[(4)] $x$ est un point générique d'une composante irréductible de
$\text{Supp}(\mathcal{F})$.
\end{enumerate}
\end{lemma}

\begin{proof}
Supposons que $f : U \to X$ soit un morphisme \'etale, avec $U$ un schéma, et
que $u \in U$ ait pour image $x$. Alors $\mathcal{F}|_U = f^*\mathcal{F}$ est
un $\mathcal{O}_U$-module cohérent sur le schéma localement noethérien $U$ ; en
particulier, $(\mathcal{F}|_U)_u$ est un $\mathcal{O}_{U, u}$-module fini ; voir
Cohomologie des espaces, lemme
\ref{spaces-cohomology-lemma-coherent-Noetherian} et Cohomologie des schémas,
lemme \ref{coherent-lemma-coherent-Noetherian}. Rappelons que le support de
$\mathcal{F}|_U$ est un sous-ensemble fermé de $U$ (Morphismes, lemme
\ref{morphisms-lemma-support-finite-type}) et que le support de
$(\mathcal{F}|_U)_u$ est l'image réciproque du support de $\mathcal{F}|_U$ par le
morphisme $\Spec(\mathcal{O}_{U, u}) \to U$. Ainsi $u$ est un point générique
d'une composante irréductible de $\text{Supp}(\mathcal{F}|_U)$ si et seulement
si le support de $(\mathcal{F}|_U)_u$ est égal à l'idéal maximal de
$\mathcal{O}_{U, u}$. L'équivalence de (1), (2) et (3) résulte alors d'Algèbre,
lemme \ref{algebra-lemma-support-point}.

\medskip\noindent
Si $X$ est décent, choisissons un morphisme \'etale $f : U \to X$ et un point
$u \in U$ d'image $x$. L'image réciproque du support de $\mathcal{F}$ est le
support de $\mathcal{F}|_U$ ; voir Morphismes d'espaces, lemme
\ref{spaces-morphisms-lemma-support-finite-type}. De plus, les spécialisations
$x' \leadsto x$ dans $|X|$ se relèvent en des spécialisations $u' \leadsto u$
dans $U$ et toute spécialisation non triviale $u' \leadsto u$ dans $U$
correspond à une spécialisation non triviale $f(u') \leadsto f(u)$ dans $|X|$,
voir Espaces décents, lemmes \ref{decent-spaces-lemma-decent-specialization}
et \ref{decent-spaces-lemma-decent-no-specializations-map-to-same-point}. En
utilisant le fait que $|X|$ et $U$ sont des espaces topologiques sobres
(Espaces décents, proposition \ref{decent-spaces-proposition-reasonable-sober}
et Schémas, lemme \ref{schemes-lemma-scheme-sober}), nous concluons que $x$
est un point générique du support de $\mathcal{F}$ si et seulement si $u$ est
un point générique du support de $\mathcal{F}|_U$. Nous concluons que (4) est
équivalent à (1).

\medskip\noindent
L'assertion entre parenthèses résulte d'Espaces décents, lemme
\ref{decent-spaces-lemma-locally-Noetherian-decent-quasi-separated}.
\end{proof}

\begin{lemma}
\label{spaces-chow-lemma-point-of-max-dimension}
Dans la situation \ref{spaces-chow-situation-setup}, soit $X/B$ bon. Soit $T
\subset |X|$ un sous-ensemble fermé et $t \in T$. Si $\dim_\delta(T) \leq k$
et $\delta(t) = k$, alors $t$ est un point générique d'une composante
irréductible de $T$.
\end{lemma}

\begin{proof}
Nous savons que $t$ appartient à une composante irréductible $T' \subset T$.
Soit $t' \in T'$ le point générique. Alors $k \geq \delta(t') \geq \delta(t)$.
Puisque $\delta$ est une fonction de dimension, nous voyons que $t = t'$.
\end{proof}



\section{Cycle associé à un sous-espace fermé}
\label{spaces-chow-section-cycle-of-closed-subscheme}

\noindent
Cette section est l'analogue d'Homologie de Chow, section
\ref{chow-section-cycle-of-closed-subscheme}.

\begin{remark}
\label{spaces-chow-remark-irreducible-component}
Dans la situation \ref{spaces-chow-situation-setup}, soit $X/B$ bon. Soit $Y
\subset X$ un sous-espace fermé. D'après les remarques \ref{spaces-chow-remark-sober} et
\ref{spaces-chow-remark-integral}, il existe des bijections (c'est-à-dire des
correspondances $1$-à-$1$) entre les ensembles suivants :
\begin{enumerate}
\item les composantes irréductibles $T$ de $|Y|$,
\item les points génériques des composantes irréductibles de $|Y|$, et
\item les sous-espaces fermés intègres $Z \subset Y$ tels que $|Z|$ soit une
composante irréductible de $|Y|$.
\end{enumerate}
Dans ce chapitre, nous appelons tout $Z$ comme en (3) une {\it composante
irréductible de $Y$}, et nous appelons $\xi \in |Z|$ son {\it point générique}.
\end{remark}

\begin{definition}
\label{spaces-chow-definition-cycle-associated-to-closed-subscheme}
Dans la situation \ref{spaces-chow-situation-setup}, soit $X/B$ bon. Soit $Y
\subset X$ un sous-espace fermé.
\begin{enumerate}
\item Pour une composante irréductible $Z \subset Y$ de point générique $\xi$,
la longueur de $\mathcal{O}_Y$ en $\xi$ (définition
\ref{spaces-chow-definition-length-at-x}) est appelée la {\it multiplicité de $Z$ dans
$Y$}. Par le lemme \ref{spaces-chow-lemma-length-finite} appliqué à $\mathcal{O}_Y$ sur
$Y$, c'est un entier positif.
\item Supposons $\dim_\delta(Y) \leq k$. Le {\it $k$-cycle associé à $Y$} est
$$
[Y]_k = \sum m_{Z, Y}[Z]
$$
où la somme porte sur les composantes irréductibles $Z$ de $Y$ de
$\delta$-dimension $k$ et $m_{Z, Y}$ est la multiplicité de $Z$ dans $Y$. Il
s'agit d'un $k$-cycle d'après Espaces sur les corps, lemme
\ref{spaces-over-fields-lemma-components-locally-finite}.
\end{enumerate}
\end{definition}

\noindent
Il importe de noter que nous ne définissons $[Y]_k$ que si la
$\delta$-dimension de $Y$ ne dépasse pas $k$. En d'autres termes, par
convention, l'écriture $[Y]_k$ implique donc que $\dim_\delta(Y) \leq
k$.







\section{Cycle associé à un faisceau cohérent}
\label{spaces-chow-section-cycle-of-coherent-sheaf}

\noindent
Cette section est l'analogue d'Homologie de Chow, section
\ref{chow-section-cycle-of-coherent-sheaf}.

\begin{definition}
\label{spaces-chow-definition-cycle-associated-to-coherent-sheaf}
Dans la situation \ref{spaces-chow-situation-setup}, soit $X/B$ bon. Soient
$\mathcal{F}$ un $\mathcal{O}_X$-module cohérent.
\begin{enumerate}
\item Pour un sous-espace fermé intègre $Z \subset X$ de point générique
$\xi$ tel que $|Z|$ soit une composante irréductible de
$\text{Supp}(\mathcal{F})$, la longueur de $\mathcal{F}$ en $\xi$ (définition
\ref{spaces-chow-definition-length-at-x}) est appelée la {\it multiplicité de $Z$ dans
$\mathcal{F}$}. D'après le lemme \ref{spaces-chow-lemma-length-finite}, il s'agit d'un
entier positif.
\item Supposons $\dim_\delta(\text{Supp}(\mathcal{F})) \leq k$. Le {\it
$k$-cycle associé à $\mathcal{F}$} est
$$
[\mathcal{F}]_k = \sum m_{Z, \mathcal{F}}[Z]
$$
où la somme porte sur les sous-espaces fermés intègres $Z \subset X$
correspondant aux composantes irréductibles de $\text{Supp}(\mathcal{F})$ de
$\delta$-dimension $k$, et $m_{Z, \mathcal{F}}$ est la multiplicité de $Z$
dans $\mathcal{F}$. Il s'agit d'un $k$-cycle d'après Espaces sur les corps, lemme
\ref{spaces-over-fields-lemma-components-locally-finite}.
\end{enumerate}
\end{definition}

\noindent
Il importe de noter que nous ne définissons $[\mathcal{F}]_k$ que si
$\mathcal{F}$ est cohérent et que la $\delta$-dimension de
$\text{Supp}(\mathcal{F})$ ne dépasse pas $k$. En d'autres termes, par
convention, l'écriture $[\mathcal{F}]_k$ implique donc que
$\mathcal{F}$ est cohérent sur $X$ et $\dim_\delta(\text{Supp}(\mathcal{F}))
\leq k$.

\begin{lemma}
\label{spaces-chow-lemma-reformulate-coeff-coherent}
Dans la situation \ref{spaces-chow-situation-setup}, soit $X/B$ bon. Soit
$\mathcal{F}$ un $\mathcal{O}_X$-module cohérent avec
$\dim_\delta(\text{Supp}(\mathcal{F})) \leq k$. Soit $Z$ un sous-espace fermé
intègre de $X$ tel que $\dim_\delta(Z) = k$. Soit $\xi \in |Z|$ son point
générique. Alors le coefficient de $Z$ dans $[\mathcal{F}]_k$ est la
longueur de $\mathcal{F}$ en $\xi$.
\end{lemma}

\begin{proof}
Observons que $|Z|$ est une composante irréductible de
$\text{Supp}(\mathcal{F})$ si et seulement si $\xi \in
\text{Supp}(\mathcal{F})$ ; voir le lemme \ref{spaces-chow-lemma-point-of-max-dimension}.
De plus, la longueur de $\mathcal{F}$ en $\xi$ est nulle si $\xi \not \in
\text{Supp}(\mathcal{F})$. Le résultat découle alors de la définition
\ref{spaces-chow-definition-cycle-associated-to-coherent-sheaf}.
\end{proof}

\begin{lemma}
\label{spaces-chow-lemma-cycle-closed-coherent}
Dans la situation \ref{spaces-chow-situation-setup}, soit $X/B$ bon. Soit $Y
\subset X$ un sous-espace fermé. Si $\dim_\delta(Y) \leq k$, alors $[Y]_k =
[i_*\mathcal{O}_Y]_k$ où $i : Y \to X$ est le morphisme d'inclusion.
\end{lemma}

\begin{proof}
Soit $Z$ un sous-espace fermé intègre de $X$ tel que $\dim_\delta(Z) = k$. Si $Z
\not \subset Y$, le coefficient de $Z$ est nul dans $[Y]_k$ comme dans
$[i_*\mathcal{O}_Y]_k$. Si $Z \subset Y$, alors le point générique de $Z$ peut
être considéré comme un point $y \in |Y|$ d'image $x \in |X|$. Alors le
coefficient de $Z$ dans $[Y]_k$ est la longueur de $\mathcal{O}_Y$ en $y$ et
le coefficient de $Z$ dans $[i_*\mathcal{O}_Y]_k$ est la longueur de
$i_*\mathcal{O}_Y$ en $x$. L'égalité des coefficients résulte donc du lemme
\ref{spaces-chow-lemma-length-closed-immersion}.
\end{proof}

\begin{lemma}
\label{spaces-chow-lemma-additivity-sheaf-cycle}
Dans la situation \ref{spaces-chow-situation-setup}, soit $X/B$ bon. Soit $0 \to
\mathcal{F} \to \mathcal{G} \to \mathcal{H} \to 0$ une suite exacte courte
de $\mathcal{O}_X$-modules cohérents. Supposons que la $\delta$-dimension de
chacun des supports de $\mathcal{F}$, $\mathcal{G}$ et $\mathcal{H}$ soit
$\leq k$. Alors
$[\mathcal{G}]_k = [\mathcal{F}]_k + [\mathcal{H}]_k$.
\end{lemma}

\begin{proof}
Soit $Z$ un sous-espace fermé intègre de $X$ tel que $\dim_\delta(Z) = k$. Il
suffit de montrer que les coefficients de $Z$ dans $[\mathcal{G}]_k$,
$[\mathcal{F}]_k$ et $[\mathcal{H}]_k$ satisfont l'additivité correspondante.
Par le lemme \ref{spaces-chow-lemma-reformulate-coeff-coherent}, il suffit de montrer
$$
\text{the length of }\mathcal{G}\text{ at }x =
\text{the length of }\mathcal{F}\text{ at }x +
\text{the length of }\mathcal{H}\text{ at }x
$$
pour tout $x \in |X|$. D'après la définition \ref{spaces-chow-definition-length-at-x},
cela résulte immédiatement de l'additivité des longueurs ; voir Algèbre, lemme
\ref{algebra-lemma-length-additive}.
\end{proof}





\section{Préliminaires à l'image directe propre}
\label{spaces-chow-section-preparation-pushforward}

\noindent
Cette section est l'analogue de l'Homologie de Chow, section
\ref{chow-section-preparation-pushforward}.

\begin{lemma}
\label{spaces-chow-lemma-proper-image}
Dans la situation \ref{spaces-chow-situation-setup}, soient $X,Y/B$ bons et soit $f : X
\to Y$ un morphisme au-dessus de $B$. Si $Z \subset X$ est un sous-espace fermé
intègre, alors il existe un unique sous-espace fermé intègre $Z' \subset Y$
tel qu'il existe un diagramme commutatif
$$
\xymatrix{
Z \ar[r] \ar[d] & X \ar[d]^f \\
Z' \ar[r] & Y
}
$$
avec $Z \to Z'$ dominant. Si $f$ est propre, alors $Z \to Z'$ est propre et
surjectif.
\end{lemma}

\begin{proof}
Soit $\xi \in |Z|$ le point générique. Soit $Z' \subset Y$ le sous-espace
fermé intègre dont le point générique est $\xi' = f(\xi)$ ; voir la remarque
\ref{spaces-chow-remark-integral}. Puisque $\xi \in |f^{-1}(Z')| = |f|^{-1}(|Z'|)$ d'après
Propriétés des espaces, lemme \ref{spaces-properties-lemma-points-cartesian},
et puisque $Z$ est le sous-espace réduit tel que $|Z| =
\overline{\{\xi\}}$, nous voyons que $Z \subset f^{-1}(Z')$ en tant que
sous-espaces fermés de $X$ (voir Propriétés des espaces, lemme
\ref{spaces-properties-lemma-map-into-reduction}). On obtient ainsi notre
morphisme $Z \to Z'$. Ce morphisme est dominant puisque le point générique de
$Z$ s'envoie sur le point générique de $Z'$. L'unicité de $Z'$ est claire. Si
$f$ est propre, alors $Z \to Y$ est propre en tant que
composition de morphismes propres (Morphismes d'espaces, lemmes
\ref{spaces-morphisms-lemma-base-change-proper} et
\ref{spaces-morphisms-lemma-closed-immersion-proper}). Il en résulte que $Z \to
Z'$ est propre d'après Morphismes d'espaces, lemme
\ref{spaces-morphisms-lemma-universally-closed-permanence}. La surjectivité en
résulte, puisque l'image d'un morphisme propre est fermée.
\end{proof}

\begin{remark}
\label{spaces-chow-remark-residue-field}
Dans la situation \ref{spaces-chow-situation-setup}, soit $X/B$ bon. Tout $x \in
|X|$ peut être représenté par un monomorphisme (unique) $\Spec(k) \to X$, où
$k$ est un corps ; voir Espaces décents, lemme
\ref{decent-spaces-lemma-decent-points-monomorphism}. Alors $k$ est le {\it
corps résiduel de} $x$ et se note $\kappa(x)$. Rappelons que $X$ possède un
sous-schéma ouvert dense $U \subset X$ (Propriétés des espaces, proposition
\ref{spaces-properties-proposition-locally-quasi-separated-open-dense-scheme}).
Si $x \in U$, alors $\kappa(x)$ coïncide avec le corps résiduel de $x$
dans le schéma $U$. Voir Espaces décents, section
\ref{decent-spaces-section-residue-fields-henselian-local-rings}.
\end{remark}

\begin{remark}
\label{spaces-chow-remark-function-field}
Dans la situation \ref{spaces-chow-situation-setup}, soit $X/B$ bon. Supposons que
$X$ soit intègre. Dans ce cas, le {\it corps des fonctions} $R(X)$ de $X$ est
défini et est égal au corps résiduel de $X$ en son point générique. Voir
Espaces sur les corps, définition
\ref{spaces-over-fields-definition-function-field}. En combinant cela avec la
remarque \ref{spaces-chow-remark-integral}, nous voyons que, pour tout $x \in X$, le
corps résiduel $\kappa(x)$ est le corps des fonctions de l'unique sous-espace fermé
intègre unique $Z \subset X$ dont le point générique est $x$.
\end{remark}

\begin{lemma}
\label{spaces-chow-lemma-equal-dimension}
Dans la situation \ref{spaces-chow-situation-setup}, soient $X, Y/B$ bons et soit $f : X
\to Y$ un morphisme au-dessus de $B$. Supposons $X$ et $Y$ intègres et
$\dim_\delta(X) = \dim_\delta(Y)$. Alors, ou bien $f$ se factorise par un
sous-espace fermé strict de $Y$, ou bien $f$ est dominant et l'extension de
corps de fonctions $R(X) / R(Y)$ est finie.
\end{lemma}

\begin{proof}
Par le lemme \ref{spaces-chow-lemma-proper-image}, il existe un sous-espace fermé intègre
unique $Z \subset Y$ tel que $f$ se factorise par un morphisme dominant $X
\to Z$. Alors $Z = Y$ si et seulement si $\dim_\delta(Z) = \dim_\delta(Y)$.
D'autre part, par notre construction des fonctions de dimension (voir la
situation \ref{spaces-chow-situation-setup}), nous avons $\dim_\delta(X) =
\dim_\delta(Z) + r$, où $r$ est le degré de transcendance de l'extension
$R(X)/R(Z)$. Le lemme résulte alors d'Espaces sur les corps, lemme
\ref{spaces-over-fields-lemma-finite-degree}.
\end{proof}

\begin{lemma}
\label{spaces-chow-lemma-quasi-compact-locally-finite}
Dans la situation \ref{spaces-chow-situation-setup}, soient $X, Y/B$ bons. Soit $f :
X \to Y$ un morphisme sur $B$. Supposons que $f$ soit quasi-compact et que
$\{T_i\}_{i \in I}$ soit une famille localement finie de parties fermées de
$|X|$. Alors $\{\overline{|f|(T_i)}\}_{i \in I}$ est une famille localement
finie de parties fermées de $|Y|$.
\end{lemma}

\begin{proof}
Soit $V \subset |Y|$ une partie ouverte quasi-compacte. Alors $|f|^{-1}(V)
\subset |X|$ est quasi-compact par Morphismes d'espaces, Lemme
\ref{spaces-morphisms-lemma-quasi-compact-is-quasi-compact}. L'ensemble $\{i
\in I : T_i \cap |f|^{-1}(V) \not = \emptyset \}$ est donc fini par un simple
argument topologique que nous omettons. Or cet ensemble est aussi égal à
$$
\{i \in I : |f|(T_i) \cap V \not = \emptyset \} =
\{i \in I : \overline{|f|(T_i)} \cap V \not = \emptyset \}
$$
ce qui démontre le lemme.
\end{proof}










\section{Image directe propre}
\label{spaces-chow-section-proper-pushforward}

\noindent
Cette section est l'analogue de l'Homologie de Chow, section
\ref{chow-section-proper-pushforward}.

\begin{definition}
\label{spaces-chow-definition-proper-pushforward}
Dans la situation \ref{spaces-chow-situation-setup}, soient $X, Y/B$ bons. Soit $f : X
\to Y$ un morphisme au-dessus de $B$. Supposons $f$ propre.
\begin{enumerate}
\item Soit $Z \subset X$ un sous-espace fermé intègre avec $\dim_\delta(Z) =
k$. Soit $Z' \subset Y$ l'image de $Z$ comme dans le lemme
\ref{spaces-chow-lemma-proper-image}. Nous définissons
$$
f_*[Z] =
\left\{
\begin{matrix}
0 & \text{if} & \dim_\delta(Z')< k, \\
\deg(Z/Z') [Z'] & \text{if} & \dim_\delta(Z') = k.
\end{matrix}
\right.
$$
Le degré de $Z$ sur $Z'$ est défini et fini si $\dim_\delta(Z') =
\dim_\delta(Z)$ d'après le lemme \ref{spaces-chow-lemma-equal-dimension} et Espaces sur
les corps, définition \ref{spaces-over-fields-definition-degree}.
\item Soit $\alpha = \sum n_Z [Z]$ un $k$-cycle sur $X$. L'{\it image directe}
de $\alpha$ est la somme
$$
f_* \alpha = \sum n_Z f_*[Z]
$$
où chaque $f_*[Z]$ est défini comme ci-dessus. Cette somme est localement finie
d'après le lemme \ref{spaces-chow-lemma-quasi-compact-locally-finite} ci-dessus.
\end{enumerate}
\end{definition}

\noindent
Par définition, l'image directe propre des cycles
$$
f_* : Z_k(X) \longrightarrow Z_k(Y)
$$
est un homomorphisme de groupes abéliens. Elle fait de $X \mapsto Z_k(X)$ un
foncteur covariant sur la catégorie dont les objets sont les espaces
algébriques bons au-dessus de $B$ et dont les morphismes sont les morphismes
propres au-dessus de $B$.

\begin{lemma}
\label{spaces-chow-lemma-compose-pushforward}
Dans la situation \ref{spaces-chow-situation-setup}, soient $X, Y, Z/B$ bons. Soient $f
: X \to Y$ et $g : Y \to Z$ des morphismes propres au-dessus de $B$. Alors $g_* \circ
f_* = (g \circ f)_*$ comme applications $Z_k(X) \to Z_k(Z)$.
\end{lemma}

\begin{proof}
Soit $W \subset X$ un sous-espace fermé intègre de dimension $k$. Considérons
les sous-espaces fermés intègres $W' \subset Y$ et $W'' \subset Z$ obtenus en
appliquant le lemme \ref{spaces-chow-lemma-proper-image} à $f$ et $W$, puis à $g$ et $W'$.
Alors $W \to W'$ et $W' \to W''$ sont surjectifs et propres. Il faut montrer
que $g_*(f_*[W]) = (f \circ g)_*[W]$. Si $\dim_\delta(W'') <
k$, alors les deux côtés sont nuls. Si $\dim_\delta(W'') = k$, alors nous
voyons que $W \to W'$ et $W' \to W''$ satisfont tous deux aux hypothèses du
lemme \ref{spaces-chow-lemma-equal-dimension}. Ainsi
$$
g_*(f_*[W]) = \deg(W/W')\deg(W'/W'')[W''],
\quad
(f \circ g)_*[W] = \deg(W/W'')[W''].
$$
La conclusion résulte alors d'Espaces sur les corps, lemme
\ref{spaces-over-fields-lemma-degree-composition}.
\end{proof}

\begin{lemma}
\label{spaces-chow-lemma-cycle-push-sheaf}
Dans la situation \ref{spaces-chow-situation-setup}, soit $f : X \to Y$ un morphisme
propre d'espaces algébriques bons au-dessus de $B$.
\begin{enumerate}
\item Soit $Z \subset X$ un sous-espace fermé avec $\dim_\delta(Z) \leq k$.
Alors
$$
f_*[Z]_k = [f_*{\mathcal O}_Z]_k.
$$
\item Soit $\mathcal{F}$ un faisceau cohérent sur $X$ tel que
$\dim_\delta(\text{Supp}(\mathcal{F})) \leq k$. Alors
$$
f_*[\mathcal{F}]_k = [f_*{\mathcal F}]_k.
$$
\end{enumerate}
Remarquons que l'énoncé a un sens puisque $f_*\mathcal{F}$ et $f_*\mathcal{O}_Z$
sont des $\mathcal{O}_Y$-modules cohérents d'après Cohomologie des espaces, lemme
\ref{spaces-cohomology-lemma-proper-pushforward-coherent}.
\end{lemma}

\begin{proof}
La partie (1) découle de (2) et du lemme \ref{spaces-chow-lemma-cycle-closed-coherent}.
Soit $\mathcal{F}$ un faisceau cohérent sur $X$. Supposons que
$\dim_\delta(\text{Supp}(\mathcal{F})) \leq k$. D'après Cohomologie des espaces,
lemme \ref{spaces-cohomology-lemma-coherent-support-closed}, il existe une
immersion fermée $i : Z \to X$ et un $\mathcal{O}_Z$-module cohérent
$\mathcal{G}$ tel que $i_*\mathcal{G} \cong \mathcal{F}$ et tel que le support
de $\mathcal{F}$ soit $Z$. Soit $Z' \subset Y$ l'image schématique de
$f|_Z : Z \to Y$ ; voir Morphismes d'espaces, définition
\ref{spaces-morphisms-definition-scheme-theoretic-image}. Considérons le
diagramme commutatif
$$
\xymatrix{
Z \ar[r]_i \ar[d]_{f|_Z} &
X \ar[d]^f \\
Z' \ar[r]^{i'} & Y
}
$$
d'espaces algébriques au-dessus de $B$. Remarquons que $f|_Z$ est surjectif
(cela résulte de Morphismes d'espaces, lemme
\ref{spaces-morphisms-lemma-quasi-compact-scheme-theoretic-image}, et du fait
que $|f|$ est fermé) et propre (cela résulte de Morphismes d'espaces, lemmes
\ref{spaces-morphisms-lemma-base-change-proper},
\ref{spaces-morphisms-lemma-closed-immersion-proper}, et
\ref{spaces-morphisms-lemma-universally-closed-permanence}). En parcourant le
diagramme par les deux chemins, nous obtenons $f_*\mathcal{F} =
f_*i_*\mathcal{G} = i'_*(f|_Z)_*\mathcal{G}$. Supposons le résultat établi pour
les immersions fermées et pour $f|_Z$. Alors
$$
f_*[\mathcal{F}]_k = f_*i_*[\mathcal{G}]_k
= (i')_*(f|_Z)_*[\mathcal{G}]_k =
(i')_*[(f|_Z)_*\mathcal{G}]_k =
[(i')_*(f|_Z)_*\mathcal{G}]_k = [f_*\mathcal{F}]_k
$$
comme voulu. Le cas d'une immersion fermée découle du lemme
\ref{spaces-chow-lemma-length-closed-immersion} et des définitions. Nous nous sommes donc
réduits au cas où $\dim_\delta(X) \leq k$ et $f : X \to Y$ est propre et
surjectif.

\medskip\noindent
Supposons que $\dim_\delta(X) \leq k$ et que $f : X \to Y$ soit propre et
surjectif. Pour toute composante irréductible $Z \subset Y$ de point
générique $\eta$, il existe un point $\xi \in X$ tel que $f(\xi) = \eta$. Ainsi
$\delta(\eta) \leq \delta(\xi) \leq k$. On voit donc que dans les expressions
$$
f_*[\mathcal{F}]_k = \sum n_Z[Z],
\quad
\text{et}
\quad
[f_*\mathcal{F}]_k = \sum m_Z[Z].
$$
dès que $n_Z \not = 0$, ou que $m_Z \not = 0$, le sous-espace fermé
intègre $Z$ est en fait une composante irréductible de $Y$ de
$\delta$-dimension $k$ (voir le lemme \ref{spaces-chow-lemma-point-of-max-dimension}).
Choisissons un tel sous-espace fermé intègre $Z \subset Y$ et notons $\eta$ son
point générique. Remarquons que, pour tout $\xi \in X$ tel que $f(\xi) = \eta$,
nous avons
$\delta(\xi) \geq k$ et donc $\xi$ est également un point générique d'une
composante irréductible de $X$ de $\delta$-dimension $k$ (voir le lemme
\ref{spaces-chow-lemma-point-of-max-dimension}). D'après Espaces sur les corps, lemme
\ref{spaces-over-fields-lemma-finite-in-codim-1}, il existe un sous-espace
ouvert $\eta \in V \subset Y$ tel que $f^{-1}(V) \to V$ soit fini. Puisque
$\eta$ est un point générique d'une composante irréductible de $|Y|$, nous
pouvons supposer que $V$ est un schéma affine ; voir Propriétés des espaces,
proposition
\ref{spaces-properties-proposition-locally-quasi-separated-open-dense-scheme}.
En remplaçant $Y$ par $V$ et $X$ par $f^{-1}(V)$, nous nous réduisons au cas où
$Y$ est affine et où $f$ est fini. En particulier, $X$ et $Y$ sont des schémas,
et nous sommes ramenés au résultat correspondant pour les schémas ; voir
Homologie de Chow, lemme \ref{chow-lemma-cycle-push-sheaf} (appliqué avec $S = Y$).
\end{proof}













\section{Préliminaires à l'image réciproque plate}
\label{spaces-chow-section-preparation-flat-pullback}

\noindent
Cette section est l'analogue de l'Homologie de Chow, section
\ref{chow-section-preparation-flat-pullback}.

\medskip\noindent
Rappelons qu'un morphisme d'espaces algébriques est dit de dimension relative
$r$ si, localement pour la topologie \'etale sur la source et sur le but, on
obtient un morphisme de schémas de dimension relative $r$. La définition
précise est équivalente, quoique légèrement différente en réalité ; voir
Morphismes d'espaces, définition
\ref{spaces-morphisms-definition-relative-dimension}.

\begin{lemma}
\label{spaces-chow-lemma-flat-inverse-image-dimension}
Dans la situation \ref{spaces-chow-situation-setup}, soient $X, Y/B$ bons. Soit $f :
X \to Y$ un morphisme au-dessus de $B$. Supposons que $f$ soit plat de
dimension relative $r$. Pour toute partie fermée $T \subset |Y|$, nous avons
$$
\dim_\delta(|f|^{-1}(T)) = \dim_\delta(T) + r.
$$
pourvu que $|f|^{-1}(T)$ ne soit pas vide. Si $Z \subset Y$ est un
sous-schéma fermé intègre et que $Z' \subset f^{-1}(Z)$ est une composante
irréductible, alors $Z'$ domine $Z$ et $\dim_\delta(Z') = \dim_\delta(Z) + r$.
\end{lemma}

\begin{proof}
Puisque la $\delta$-dimension d'une partie fermée est la borne supérieure des
$\delta$-dimensions de ses composantes irréductibles, il suffit de démontrer
la dernière assertion. Nous pouvons remplacer $Y$ par le sous-schéma fermé intègre
$Z$ et $X$ par $f^{-1}(Z) = Z \times_Y X$. Nous pouvons donc supposer que $Z =
Y$ est intègre et que $f$ est un morphisme plat de dimension relative $r$.
Puisque $Y$ est localement noethérien, le morphisme $f$, qui est localement de
type fini, est en fait localement de présentation finie. Le lemme
\ref{spaces-morphisms-lemma-fppf-open} de Morphismes d'espaces s'applique donc,
et $f$ est ouvert. Soit $\xi \in X$ un point générique d'une composante
irréductible de $X$. Comme $f$ est ouvert, $f(\xi)$ est le point
générique $\eta$ de $Z = Y$. Ainsi $Z'$ domine $Z = Y$. Enfin, la proposition
\ref{spaces-properties-proposition-locally-quasi-separated-open-dense-scheme}
de Propriétés des espaces montre que $\xi$ et $\eta$ appartiennent aux lieux
schématiques de $X$ et de $Y$.
Puisque $\xi$ est un point générique de $X$, nous voyons que $\mathcal{O}_{X,
\xi} = \mathcal{O}_{X_\eta, \xi}$ n'a qu'un seul idéal premier et a donc la
dimension $0$ (on peut utiliser les anneaux locaux usuels, car $\xi$ et
$\eta$ appartiennent aux lieux schématiques de $X$ et de $Y$). Le lemme
\ref{spaces-morphisms-lemma-dimension-fibre-at-a-point} de Morphismes d'espaces
(et la définition des morphismes d'une dimension relative donnée) montre alors que
le degré de transcendance de $\kappa(\xi)$ sur $\kappa(\eta)$ est $r$. En
d'autres termes, $\delta(\xi) = \delta(\eta) + r$, comme voulu.
\end{proof}

\noindent
Voici le lemme que nous utiliserons pour montrer que l'image réciproque plate
d'une famille localement finie de sous-schémas fermés est localement finie.

\begin{lemma}
\label{spaces-chow-lemma-inverse-image-locally-finite}
Dans la situation \ref{spaces-chow-situation-setup}, soient $X, Y/B$ bons. Soit $f :
X \to Y$ un morphisme au-dessus de $B$. Supposons que $\{T_i\}_{i \in I}$ soit
une famille localement finie de parties fermées de $|Y|$. Alors
$\{|f|^{-1}(T_i)\}_{i \in I}$ est une famille localement finie de
parties fermées de $X$.
\end{lemma}

\begin{proof}
Soit $U \subset |X|$ une partie ouverte quasi-compacte. Puisque l'image
$|f|(U) \subset |Y|$ est une partie quasi-compacte, il existe une partie $V
\subset |Y|$ ouverte quasi-compacte telle que $|f|(U) \subset V$. Remarquons que
$$
\{i \in I : |f|^{-1}(T_i) \cap U \not = \emptyset \}
\subset
\{i \in I : T_i \cap V \not = \emptyset \}.
$$
Le membre de droite étant fini par hypothèse, on conclut.
\end{proof}

















\section{Image réciproque plate}
\label{spaces-chow-section-flat-pullback}

\noindent
Cette section est l'analogue de l'Homologie de Chow, section
\ref{chow-section-flat-pullback}.

\medskip\noindent
Soit $S$ un schéma et soit $f : X \to Y$ un morphisme d'espaces algébriques
au-dessus de $S$. Soit $Z \subset Y$ un sous-espace fermé. Dans ce chapitre, nous
utiliserons parfois la terminologie {\it image réciproque schématique}
pour l'image réciproque $f^{-1}(Z)$ de $Z$ construite dans Morphismes des
espaces, définition
\ref{spaces-morphisms-definition-inverse-image-closed-subspace}. L'image
réciproque schématique est le produit fibré
$$
\xymatrix{
f^{-1}(Z) \ar[r] \ar[d] & X \ar[d] \\
Z \ar[r] & Y
}
$$
Si $\mathcal{I} \subset \mathcal{O}_Y$ est le faisceau quasi-cohérent d'idéaux
correspondant à $Z$ dans $Y$, alors $f^{-1}(\mathcal{I})\mathcal{O}_X$ est le
faisceau quasi-cohérent d'idéaux correspondant à $f^{-1}(Z)$ dans $X$.

\begin{definition}
\label{spaces-chow-definition-flat-pullback}
Dans la situation \ref{spaces-chow-situation-setup}, soient $X, Y/B$ bons. Soit $f :
X \to Y$ un morphisme au-dessus de $B$. Supposons que $f$ soit plat de dimension
relative $r$.
\begin{enumerate}
\item Soit $Z \subset Y$ un sous-espace fermé intègre de $\delta$-dimension
$k$. Nous définissons $f^*[Z]$ comme le $(k+r)$-cycle sur $X$ associé à
l'image réciproque schématique
$$
f^*[Z] = [f^{-1}(Z)]_{k+r}.
$$
Cela a du sens puisque $\dim_\delta(f^{-1}(Z)) = k + r$ par le lemme
\ref{spaces-chow-lemma-flat-inverse-image-dimension}.
\item Soit $\alpha = \sum n_i [Z_i]$ un $k$-cycle sur $Y$. L'{\it image
réciproque plate de $\alpha$ par $f$} est la somme
$$
f^* \alpha = \sum n_i f^*[Z_i]
$$
où chaque $f^*[Z_i]$ est défini comme ci-dessus. La somme est localement finie
par le lemme \ref{spaces-chow-lemma-inverse-image-locally-finite}.
\item On note $f^* : Z_k(Y) \to Z_{k + r}(X)$ l'homomorphisme de groupes abéliens
ainsi obtenu.
\end{enumerate}
\end{definition}

\noindent
Une immersion ouverte est plate. C'est un cas particulier important, quoique
trivial, de morphisme plat. Si $U \subset X$ est ouvert, l'image réciproque par
$j : U \to X$ d'un cycle est parfois appelée la {\it restriction} du cycle à
$U$. Remarquons que, dans ce cas, les applications
$$
j^* : Z_k(X) \longrightarrow Z_k(U)
$$
sont toutes {\it surjectives}. En effet, étant donné un sous-espace fermé
intègre $Z' \subset U$, nous pouvons prendre pour $Z$ l'adhérence de
$Z'$ dans $X$ et la considérer comme un sous-espace fermé réduit de $X$ (voir
Propriétés des espaces, définition
\ref{spaces-properties-definition-reduced-induced-space}). On a alors clairement
$Z \cap U = Z'$, autrement dit $j^*[Z] = [Z']$, d'où la surjectivité. En fait,
on a un peu mieux.

\begin{lemma}
\label{spaces-chow-lemma-exact-sequence-open}
Dans la situation \ref{spaces-chow-situation-setup}, soit $X/B$ bon. Soit $U
\subset X$ un sous-espace ouvert. Soit $Y$ le sous-espace fermé réduit de $X$
avec $|Y| = |X| \setminus |U|$ et notons $i : Y \to X$ le morphisme
d'inclusion. Pour chaque $k \in \mathbf{Z}$, la séquence
$$
\xymatrix{
Z_k(Y) \ar[r]^{i_*} & Z_k(X) \ar[r]^{j^*} & Z_k(U) \ar[r] & 0
}
$$
est une suite exacte de groupes abéliens.
\end{lemma}

\begin{proof}
Nous avons vu ci-dessus que $j^*$ est surjectif. Supposons d'abord que $X$
soit quasi-compact. Alors $Z_k(X)$ est un $\mathbf{Z}$-module libre de base
donnée par les éléments $[Z]$ où $Z \subset X$ est fermé intègre de
$\delta$-dimension $k$. Un tel élément de base s'envoie soit sur l'élément de
base $[Z \cap U]$ de $Z_k(U)$, soit sur zéro si $Z \subset Y$. Le lemme est donc
clair dans ce cas. Le cas général est similaire et la preuve est omise.
\end{proof}

\begin{lemma}
\label{spaces-chow-lemma-etale-pullback}
Dans la situation \ref{spaces-chow-situation-setup}, soit $f : X \to Y$ un morphisme
\'etale de bons espaces algébriques au-dessus de $B$. Si $Z \subset Y$ est un
sous-espace fermé intègre, alors $f^*[Z] = \sum [Z']$ où la somme est sur les
composantes irréductibles (remarque \ref{spaces-chow-remark-irreducible-component}) de
$f^{-1}(Z)$.
\end{lemma}

\begin{proof}
La signification du lemme est que le coefficient de $[Z']$ est $1$. Cela
découle du fait que $f^{-1}(Z)$ est un espace algébrique réduit car il est
\'etale sur l'espace algébrique intègre $Z$.
\end{proof}

\begin{lemma}
\label{spaces-chow-lemma-compose-flat-pullback}
Dans la situation \ref{spaces-chow-situation-setup}, soient $X, Y, Z/B$ bons. Soient $f
: X \to Y$ et $g : Y \to Z$ des morphismes plats de dimensions relatives $r$
et $s$ au-dessus de $B$. Alors $g \circ f$ est plat de dimension relative $r + s$ et
$$
f^* \circ g^* = (g \circ f)^*
$$
comme applications $Z_k(Z) \to Z_{k + r + s}(X)$.
\end{lemma}

\begin{proof}
La composition est plate de dimension relative $r + s$ d'après Morphismes
d'espaces, lemmes
\ref{spaces-morphisms-lemma-dimension-fibre-at-a-point-additive} et
\ref{spaces-morphisms-lemma-composition-flat}. Supposons que
\begin{enumerate}
\item $A \subset Z$ est un sous-espace fermé intègre de $\delta$-dimension
$k$,
\item $A' \subset Y$ est un sous-espace fermé intègre de $\delta$-dimension
$k + s$ tel que $A' \subset g^{-1}(A)$, et
\item $A'' \subset Y$ est un sous-espace fermé intègre de $\delta$-dimension
$k + s + r$ tel que $A'' \subset f^{-1}(W')$.
\end{enumerate}
Nous devons montrer que le coefficient $n$ de $[A'']$ dans $(g \circ f)^*[A]$
est égal au coefficient $m$ de $[A'']$ dans $f^*(g^*[A])$. On peut
choisir un diagramme commutatif
$$
\xymatrix{
U \ar[d] \ar[r] & V \ar[d] \ar[r] & W \ar[d] \\
X \ar[r] & Y \ar[r] & Z
}
$$
où $U, V, W$ sont des schémas, les flèches verticales sont \'etales, et il
existe des points $u \in U$, $v \in V$, $w \in W$ tels que $u \mapsto v
\mapsto w$ et tels que $u, v, w$ s'envoient sur les points génériques de $A'',
A', A$. (Détails omis.) Nous avons alors des homomorphismes plats d'anneaux
locaux $\mathcal{O}_{W, w} \to \mathcal{O}_{V, v}$, $\mathcal{O}_{V, v} \to
\mathcal{O}_{U, u}$, et en utilisant à plusieurs reprises le lemme
\ref{spaces-chow-lemma-length}, nous trouvons
$$
n = \text{length}_{\mathcal{O}_{U, u}}(
\mathcal{O}_{U, u}/\mathfrak m_w\mathcal{O}_{U, u})
$$
et
$$
m =
\text{length}_{\mathcal{O}_{V, v}}(
\mathcal{O}_{V, v}/\mathfrak m_w\mathcal{O}_{V, v})
\text{length}_{\mathcal{O}_{U, u}}(
\mathcal{O}_{U, u}/\mathfrak m_v\mathcal{O}_{U, u})
$$
L'égalité découle donc du lemme d'Algèbre
\ref{algebra-lemma-pullback-transitive}.
\end{proof}

\begin{lemma}
\label{spaces-chow-lemma-pullback-coherent}
Dans la situation \ref{spaces-chow-situation-setup}, soient $X, Y/B$ bons. Soit $f :
X \to Y$ un morphisme plat de dimension relative $r$.
\begin{enumerate}
\item Soit $Z \subset Y$ un sous-espace fermé avec $\dim_\delta(Z) \leq k$.
Alors $\dim_\delta(f^{-1}(Z)) \leq k + r$ et $[f^{-1}(Z)]_{k +
r} = f^*[Z]_k$ dans $Z_{k + r}(X)$.
\item Soit $\mathcal{F}$ un faisceau cohérent sur $Y$ avec
$\dim_\delta(\text{Supp}(\mathcal{F})) \leq k$. Nous avons alors
$\dim_\delta(\text{Supp}(f^*\mathcal{F})) \leq k + r$ et
$$
f^*[{\mathcal F}]_k = [f^*{\mathcal F}]_{k+r}
$$
dans $Z_{k + r}(X)$.
\end{enumerate}
\end{lemma}

\begin{proof}
La partie (1) découle de la partie (2) par le lemme
\ref{spaces-chow-lemma-cycle-closed-coherent} et le fait que $f^*\mathcal{O}_Z =
\mathcal{O}_{f^{-1}(Z)}$.

\medskip\noindent
Preuve de (2). Comme $X$, $Y$ sont localement noethériens, nous pouvons
appliquer le lemme
\ref{spaces-cohomology-lemma-coherent-Noetherian} de Cohomologie des espaces :
$\mathcal{F}$ est de type fini, donc $f^*\mathcal{F}$ est de type fini (Modules
sur les sites, lemme \ref{sites-modules-lemma-local-pullback}), et par suite
$f^*\mathcal{F}$ est cohérent (de nouveau Cohomologie des espaces, lemme
\ref{spaces-cohomology-lemma-coherent-Noetherian}). Le lemme a donc un sens.
Soit $W \subset Y$ un sous-espace fermé intègre de $\delta$-dimension $k$, et
soit $W' \subset X$ un sous-espace fermé intègre de dimension $k + r$
s'envoyant dans $W$ par $f$. Nous devons montrer que le
coefficient $n$ de $[W']$ dans $f^*[{\mathcal F}]_k$ est égal au
coefficient $m$ de $[W']$ dans $[f^*{\mathcal F}]_{k+r}$. On peut choisir un
diagramme commutatif
$$
\xymatrix{
U \ar[d] \ar[r] & V \ar[d] \\
X \ar[r] & Y
}
$$
où $U, V$ sont des schémas, les flèches verticales sont \'etales, et il existe
des points $u \in U$, $v \in V$ tels que $u \mapsto v$ et tels que $u, v$
s'envoient sur les points génériques de $W', W$. (Détails omis.) Considérons le
germe $M = (\mathcal{F}|_V)_v$ comme un $\mathcal{O}_{V, v}$-module.
(Remarquons que $M$ est de longueur finie d'après nos hypothèses de dimension,
mais nous n'avons en fait pas besoin de le vérifier ; voir le lemme
\ref{spaces-chow-lemma-length-finite}.) Nous avons $(f^*\mathcal{F}|_U)_u =
\mathcal{O}_{U, u} \otimes_{\mathcal{O}_{V, v}} M$. Ainsi on voit que
$$
n = \text{length}_{\mathcal{O}_{U, u}}
(\mathcal{O}_{U, u} \otimes_{\mathcal{O}_{V, v}} M)
\quad
\text{et}
\quad
m = \text{length}_{\mathcal{O}_{V, v}}(M)
\text{length}_{\mathcal{O}_{V, v}}(
\mathcal{O}_{U, u}/\mathfrak m_v \mathcal{O}_{U, u})
$$
Ainsi, l'égalité découle du lemme d'Algèbre
\ref{algebra-lemma-pullback-module}.
\end{proof}









\section{Image directe et image réciproque}
\label{spaces-chow-section-push-pull}

\noindent
Cette section est l'analogue de l'Homologie de Chow, section
\ref{chow-section-push-pull}.

\medskip\noindent
Dans cette section, nous vérifions que l'image directe propre et l'image
réciproque plate sont compatibles lorsque cela a un sens. D'après le travail que
nous avons effectué ci-dessus, cela est une conséquence de la cohomologie et
du changement de base.

\begin{lemma}
\label{spaces-chow-lemma-flat-pullback-proper-pushforward}
Dans la situation \ref{spaces-chow-situation-setup}, soit
$$
\xymatrix{
X' \ar[r]_{g'} \ar[d]_{f'} & X \ar[d]^f \\
Y' \ar[r]^g & Y
}
$$
un diagramme cartésien de bons espaces algébriques au-dessus de $B$.
Supposons $f : X \to Y$ propre et $g : Y' \to Y$ plat de dimension
relative $r$. Alors également $f'$ est propre et $g'$ est plat de dimension
relative $r$. Pour tout $k$-cycle $\alpha$ sur $X$, nous avons
$$
g^*f_*\alpha = f'_*(g')^*\alpha
$$
dans $Z_{k + r}(Y')$.
\end{lemma}

\begin{proof}
L'assertion selon laquelle $f'$ est propre découle de Morphismes d'espaces,
lemme \ref{spaces-morphisms-lemma-base-change-proper}. Celle selon laquelle
$g'$ est plat de dimension relative $r$ découle de Morphismes d'espaces, lemmes
\ref{spaces-morphisms-lemma-dimension-fibre-after-base-change} et
\ref{spaces-morphisms-lemma-base-change-flat}. Il suffit de prouver l'égalité
des cycles lorsque $\alpha = [W]$ pour un sous-espace fermé intègre $W
\subset X$ de $\delta$-dimension $k$. Remarquons que, dans ce cas, $\alpha
= [\mathcal{O}_W]_k$ ; voir le lemme \ref{spaces-chow-lemma-cycle-closed-coherent}. D'après les
lemmes \ref{spaces-chow-lemma-cycle-push-sheaf} et \ref{spaces-chow-lemma-pullback-coherent} il suffit
donc de montrer que $f'_*(g')^*\mathcal{O}_W$ est isomorphe à
$g^*f_*\mathcal{O}_W$. Cela découle de la cohomologie et du changement de
base ; voir Cohomologie des espaces, lemme
\ref{spaces-cohomology-lemma-flat-base-change-cohomology}.
\end{proof}

\begin{lemma}
\label{spaces-chow-lemma-finite-flat}
Dans la situation \ref{spaces-chow-situation-setup}, soient $X, Y/B$ bons. Soit $f :
X \to Y$ un morphisme fini localement libre de degré $d$ (voir Morphismes des
espaces, définition \ref{spaces-morphisms-definition-finite-locally-free}).
Alors $f$ est à la fois propre et plat de dimension relative $0$, et
$$
f_*f^*\alpha = d\alpha
$$
pour tout $\alpha \in Z_k(Y)$.
\end{lemma}

\begin{proof}
Un morphisme fini localement libre est plat et fini d'après Morphismes
d'espaces, lemme \ref{spaces-morphisms-lemma-finite-flat}, et un morphisme
fini est propre d'après Morphismes d'espaces, lemme
\ref{spaces-morphisms-lemma-finite-proper}. Nous omettons de vérifier qu'un
morphisme fini est de dimension relative $0$. La formule a donc un sens. Pour
la démontrer, soit $Z \subset Y$ un sous-schéma fermé intègre de
$\delta$-dimension $k$. Il suffit de démontrer la formule pour
$\alpha = [Z]$. Puisque le changement de base d'un morphisme fini localement
libre est fini localement libre (Morphismes des espaces, lemme
\ref{spaces-morphisms-lemma-base-change-finite-locally-free}), on voit que
$f_*f^*\mathcal{O}_Z$ est un faisceau localement libre de rang fini $d$ sur
$Z$. Il est donc clair que $f_*f^*\mathcal{O}_Z$ est de longueur $d$ au point
générique de $Z$. Par conséquent,
$$
f_*f^*[Z] = f_*f^*[\mathcal{O}_Z]_k =
[f_*f^*\mathcal{O}_Z]_k = d[Z]
$$
où nous avons utilisé les lemmes \ref{spaces-chow-lemma-pullback-coherent} et
\ref{spaces-chow-lemma-cycle-push-sheaf}.
\end{proof}










\section{Préliminaires aux diviseurs principaux}
\label{spaces-chow-section-preparation-principal-divisors}

\noindent
Cette section est l'analogue de l'Homologie de Chow, section
\ref{chow-section-preparation-principal-divisors}. Certains éléments de cette
section recoupent en partie la discussion d'Espaces sur les corps, section
\ref{spaces-over-fields-section-Weil-divisors}.

\begin{lemma}
\label{spaces-chow-lemma-divisor-delta-dimension}
Dans la situation \ref{spaces-chow-situation-setup}, soit $X/B$ bon. Supposons que
$X$ soit intègre.
\begin{enumerate}
\item Si $Z \subset X$ est un sous-espace fermé intègre, alors les assertions
suivantes sont équivalentes :
\begin{enumerate}
\item $Z$ est un diviseur premier,
\item $|Z|$ est de codimension $1$ dans $|X|$, et
\item $\dim_\delta(Z) = \dim_\delta(X) - 1$.
\end{enumerate}
\item Si $Z$ est une composante irréductible d'un diviseur de Cartier effectif
sur $X$, alors $\dim_\delta(Z) = \dim_\delta(X) - 1$.
\end{enumerate}
\end{lemma}

\begin{proof}
La partie (1) découle de la définition d'un diviseur premier (Espaces sur les
corps, définition \ref{spaces-over-fields-definition-Weil-divisor}), d'Espaces
décents, lemme \ref{decent-spaces-lemma-codimension-local-ring}, et de la
définition d'une fonction de dimension (Topologie, définition
\ref{topology-definition-dimension-function}).

\medskip\noindent
Soit $D \subset X$ un diviseur de Cartier effectif. Soit $Z \subset D$
une composante irréductible et soit $\xi \in |Z|$ le point générique.
Choisissons un voisinage \'etale $(U, u) \to (X, \xi)$ où $U = \Spec(A)$ et où
$D \times_X U$ est défini par un non-diviseur de zéro $f \in A$ ; voir
Diviseurs sur les espaces, lemme
\ref{spaces-divisors-lemma-characterize-effective-Cartier-divisor}. Alors $u$
est le point générique de $V(f)$ d'après Espaces décents, lemme
\ref{decent-spaces-lemma-decent-generic-points}. Par conséquent,
$\mathcal{O}_{U, u}$ est de dimension $1$ d'après le théorème de l'idéal
principal de Krull (Algèbre, lemme \ref{algebra-lemma-minimal-over-1}). Ainsi,
$\xi$ est un point de codimension $1$ sur $X$ et $Z$ est un diviseur premier,
comme voulu.
\end{proof}












\section{Diviseurs principaux}
\label{spaces-chow-section-principal-divisors}

\noindent
Cette section est l'analogue de l'Homologie de Chow, section
\ref{chow-section-principal-divisors}. La définition suivante est l'analogue
d'Espaces sur les corps, définition
\ref{spaces-over-fields-definition-principal-divisor}, dans notre cadre.

\begin{definition}
\label{spaces-chow-definition-principal-divisor}
Dans la situation \ref{spaces-chow-situation-setup}, soit $X/B$ bon. Supposons que
$X$ soit intègre et que $\dim_\delta(X) = n$. Soit $f \in R(X)^*$. Le
{\it diviseur principal associé à $f$} est le $(n - 1)$-cycle
$$
\text{div}(f) = \text{div}_X(f) = \sum \text{ord}_Z(f) [Z]
$$
défini dans Espaces sur les corps, définition
\ref{spaces-over-fields-definition-principal-divisor}. Cela a un sens car
les diviseurs premiers sont de $\delta$-dimension $n - 1$ par le lemme
\ref{spaces-chow-lemma-divisor-delta-dimension}.
\end{definition}

\noindent
Dans la situation de la définition, pour $f, g \in R(X)^*$, nous avons
$$
\text{div}_X(fg) = \text{div}_X(f) + \text{div}_X(g)
$$
dans $Z_{n - 1}(X)$. Voir Espaces sur les corps, lemme
\ref{spaces-over-fields-lemma-div-additive}. Le lemme suivant nous permettra
de ramener les assertions sur les diviseurs principaux au cas des schémas.

\begin{lemma}
\label{spaces-chow-lemma-etale-pullback-principal-divisor}
Dans la situation \ref{spaces-chow-situation-setup}, soit $f : X \to Y$ un morphisme
\'etale de bons espaces algébriques au-dessus de $B$. Supposons que $Y$ soit intègre.
Soit $g \in R(Y)^*$. Comme cycles sur $X$ nous avons
$$
f^*(\text{div}_Y(g)) =
\sum\nolimits_{X'} (X' \to X)_*\text{div}_{X'}(g \circ f|_{X'})
$$
où la somme s'étend sur les composantes irréductibles de $X$ (remarque
\ref{spaces-chow-remark-irreducible-component}).
\end{lemma}

\begin{proof}
L'application $|X| \to |Y|$ est ouverte. L'ensemble des composantes irréductibles
de $|X|$ est localement fini dans $|X|$. Nous concluons que $f|_{X'} : X' \to
Y$ est dominant pour toute composante irréductible $X' \subset X$. Ainsi $g
\circ f|_{X'}$ est défini (Morphismes des espaces, section
\ref{spaces-morphisms-section-rational-maps}), donc $\text{div}_{X'}(g \circ
f|_{X'})$ est défini. De plus, la somme est localement finie et on trouve que
le membre de droite est bien un cycle sur $X$. Le membre de gauche est défini par
la définition \ref{spaces-chow-definition-flat-pullback} et le fait qu'un morphisme
\'etale est plat de dimension relative $0$.

\medskip\noindent
Puisque $f$ est \'etale, nous voyons que $\delta_X(x) = \delta_y(f(x))$ pour
tout $x \in |X|$. Ainsi, si $\dim_\delta(Y) = n$, alors $\dim_\delta(X') =
n$ pour toute composante irréductible $X'$ de $X$ (puisque les points
génériques de $X$ s'envoient sur le point générique de $Y$ ; voir ci-dessus).
Les deux membres sont donc des $(n - 1)$-cycles.

\medskip\noindent
Soit $Z \subset X$ un sous-espace fermé intègre avec $\dim_\delta(Z) = n -
1$. Pour démontrer l'égalité, il faut montrer que les coefficients de $Z$
dans les deux membres sont égaux. Soit $Z' \subset Y$ le sous-espace fermé intègre construit
dans le lemme \ref{spaces-chow-lemma-proper-image}. Alors $\dim_\delta(Z') = n - 1$ aussi.
Soit $\xi \in |Z|$ le point générique. Alors $\xi' = f(\xi) \in |Z'|$ est le
point générique. Considérons le diagramme commutatif
$$
\xymatrix{
\Spec(\mathcal{O}_{X, \xi}^h) \ar[r] \ar[d] & X \ar[d] \\
\Spec(\mathcal{O}_{Y, \xi'}^h) \ar[r] & Y
}
$$
d'Espaces décents, remarque
\ref{decent-spaces-remark-functoriality-henselian-local-ring}. Nous devons
être légèrement prudents car les anneaux locaux noéthériens réduits
$\mathcal{O}_{X, \xi}^h$ et $\mathcal{O}_{Y, \xi'}^h$ ne sont pas
nécessairement intègres. Nous travaillons donc avec des anneaux totaux de
fractions $Q(-)$ plutôt que des corps des fractions. Par définition, pour
obtenir le coefficient de $Z'$ dans $\text{div}_Y(g)$ nous écrivons l'image de
$g$ dans $Q(\mathcal{O}_{Y, \xi'}^h)$ comme $a/b$ avec $a, b \in
\mathcal{O}_{Y, \xi'}^h$ non-diviseurs de zéro et nous prenons
$$
\text{ord}_{Z'}(g) =
\text{length}_{\mathcal{O}_{Y, \xi'}^h}
(\mathcal{O}_{Y, \xi'}^h/a \mathcal{O}_{Y, \xi'}^h) -
\text{length}_{\mathcal{O}_{Y, \xi'}^h}
(\mathcal{O}_{Y, \xi'}^h/b \mathcal{O}_{Y, \xi'}^h)
$$
Remarquons que le coefficient de $Z$ dans $f^*\text{div}_Y(G)$ est ce même
entier ; voir le lemme \ref{spaces-chow-lemma-etale-pullback}. Supposons que $\xi \in X'$.
On peut alors considérer les applications
$$
\mathcal{O}_{Y, \xi'}^h \to
\mathcal{O}_{X, \xi}^h \to
\mathcal{O}_{X', \xi}^h
$$
La première flèche est plate et la seconde est un homomorphisme surjectif
d'anneaux locaux noethériens réduits de dimension $1$. Par conséquent, ces
deux homomorphismes envoient les non-diviseurs de zéro sur des non-diviseurs
de zéro, et le coefficient de $Z'$ dans $\text{div}_{X'}(g \circ f|_{X'})$
est
$$
\text{ord}_Z(g \circ f|_{X'}) =
\text{length}_{\mathcal{O}_{X', \xi}^h}
(\mathcal{O}_{X', \xi}^h/a \mathcal{O}_{X', \xi}^h) -
\text{length}_{\mathcal{O}_{Y, \xi'}^h}
(\mathcal{O}_{X', \xi}^h/b \mathcal{O}_{X', \xi}^h)
$$
par la même prescription que ci-dessus. Il suffit donc de montrer
$$
\text{length}_{\mathcal{O}_{Y, \xi'}^h}
(\mathcal{O}_{Y, \xi'}^h/a \mathcal{O}_{Y, \xi'}^h) =
\sum\nolimits_{\xi \in |X'|}
\text{length}_{\mathcal{O}_{X', \xi}^h}
(\mathcal{O}_{X', \xi}^h/a \mathcal{O}_{X', \xi}^h)
$$
Tout d'abord, puisque l'homomorphisme d'anneaux $\mathcal{O}_{Y, \xi'}^h \to
\mathcal{O}_{X, \xi}^h$ est plate et non ramifiée, nous avons
$$
\text{length}_{\mathcal{O}_{Y, \xi'}^h}
(\mathcal{O}_{Y, \xi'}^h/a \mathcal{O}_{Y, \xi'}^h) =
\text{length}_{\mathcal{O}_{X, \xi}^h}
(\mathcal{O}_{X, \xi}^h/a \mathcal{O}_{X, \xi}^h)
$$
par Algèbre, lemme \ref{algebra-lemma-pullback-module}. Soient $\mathfrak q_1,
\ldots, \mathfrak q_t$ les idéaux premiers non maximaux de $\mathcal{O}_{X,
\xi}^h$, et posons $R_j = \mathcal{O}_{X, \xi}^h/\mathfrak q_j$. Pour $X'$
comme ci-dessus, notons $J(X') \subset \{1, \ldots, t\}$ l'ensemble des
indices tels que $\mathfrak q_j$ corresponde à un point de $X'$, c'est-à-dire
tels que, par l'homomorphisme surjectif $\mathcal{O}_{X, \xi}^h \to
\mathcal{O}_{X', \xi}$, l'idéal premier $\mathfrak q_j$ corresponde à un idéal
premier de $\mathcal{O}_{X', \xi}$. D'après Homologie de Chow, lemme
\ref{chow-lemma-additivity-divisors-restricted}, on obtient
$$
\text{length}_{\mathcal{O}_{X, \xi}^h}
(\mathcal{O}_{X, \xi}^h/a \mathcal{O}_{X, \xi}^h) =
\sum\nolimits_j \text{length}_{R_j}(R_j/a R_j)
$$
et
$$
\text{length}_{\mathcal{O}_{X', \xi}^h}
(\mathcal{O}_{X', \xi}^h/a \mathcal{O}_{X', \xi}^h) =
\sum\nolimits_{j \in J(X')} \text{length}_{R_j}(R_j/a R_j)
$$
Le résultat du lemme en découle, car $\{1, \ldots, t\}$ est l'union
disjointe des ensembles $J(X')$ : chaque point de codimension $0$ sur $X$
appartient à un unique $X'$.
\end{proof}







\section{Diviseurs principaux et image directe}
\label{spaces-chow-section-two-fun}

\noindent
Cette section est l'analogue de l'Homologie de Chow, section
\ref{chow-section-two-fun}.

\begin{lemma}
\label{spaces-chow-lemma-proper-pushforward-alteration}
Dans la situation \ref{spaces-chow-situation-setup}, soient $X, Y/B$ bons. Supposons
que $X$, $Y$ soient intègres et que $n = \dim_\delta(X) = \dim_\delta(Y)$.
Soit $p : X \to Y$ un morphisme propre dominant. Soit $f \in R(X)^*$.
Posons
$$
g = \text{Nm}_{R(X)/R(Y)}(f).
$$
Nous avons alors $p_*\text{div}(f) = \text{div}(g)$.
\end{lemma}

\begin{proof}
Nous allons le déduire du cas des schémas par localisation \'etale. Soit $Z
\subset Y$ un sous-espace fermé intègre de $\delta$-dimension $n - 1$.
Nous voulons montrer que les coefficients de $[Z]$ dans $p_*\text{div}(f)$ et
$\text{div}(g)$ sont égaux. Appliquons le lemme
\ref{spaces-over-fields-lemma-finite-in-codim-1} d'Espaces sur les corps au
morphisme $p : X \to Y$ et au point générique $\xi \in |Z|$. Nous pouvons
alors remplacer $Y$ par un sous-espace ouvert contenant $\xi$ et supposer que
$p : X \to Y$ est fini. Choisissons un voisinage \'etale $(V, v) \to (Y, \xi)$
où $V$ est un schéma affine. D'après le lemme \ref{spaces-chow-lemma-etale-pullback}, il
suffit de démontrer l'égalité des cycles après image réciproque sur $V$. Posons
$U = V \times_Y X$ et
considérez le diagramme commutatif
$$
\xymatrix{
U \ar[r]_a \ar[d]_{p'} & X \ar[d]^p \\
V \ar[r]^b & Y
}
$$
Soient $V_j \subset V$, $j = 1, \ldots, m$, les composantes irréductibles de
$V$. Pour chaque $i$, soient $U_{j, i}$, $i = 1, \ldots, n_j$, les composantes
irréductibles de $U$ dominant $V_j$. Notons $p'_{j, i} : U_{j, i} \to V_j$ la
restriction de $p' : U \to V$. D'après le cas des schémas (Homologie de Chow,
lemme \ref{chow-lemma-proper-pushforward-alteration}), on a
$$
p'_{j, i, *}\text{div}_{U_{j, i}}(f_{j, i}) = \text{div}_{V_j}(g_{j, i})
$$
où $f_{j, i}$ est la restriction de $f$ à $U_{j, i}$ et $g_{j, i}$ est la
norme de $f_{j, i}$ pour l'extension finie $R(U_{j, i})/R(V_j)$. Nous
avons
\begin{align*}
b^* p_*\text{div}_X(f)
& =
p'_* a^* \text{div}_X(f) \\
& =
p'_*\left(\sum\nolimits_{j, i}
(U_{j, i} \to U)_*\text{div}_{U_{j, i}}(f_{j, i})\right) \\
& =
\sum\nolimits_{j, i}
(V_j \to V)_*p'_{j, i, *}\text{div}_{U_{j, i}}(f_{j, i}) \\
& =
\sum\nolimits_j
(V_j \to V)_*\left(\sum\nolimits_i \text{div}_{V_j}(g_{j, i})\right) \\
& =
\sum\nolimits_j (V_j \to V)_*\text{div}_{V_j}(\prod\nolimits_i g_{j, i})
\end{align*}
par les lemmes \ref{spaces-chow-lemma-flat-pullback-proper-pushforward},
\ref{spaces-chow-lemma-etale-pullback-principal-divisor} et
\ref{spaces-chow-lemma-compose-pushforward}. Pour terminer la preuve, en utilisant à
nouveau le lemme \ref{spaces-chow-lemma-etale-pullback-principal-divisor}, il suffit de
montrer que
$$
g \circ b|_{V_j} = \prod\nolimits_i g_{j, i}
$$
comme éléments du corps des fonctions de $V_j$. En termes de corps, il s'agit
de l'énoncé suivant : soit $L/K$ une extension finie. Soit $M/K$ une extension
finie séparable. Écrivons $M \otimes_K L = \prod M_i$. Alors, pour $t \in L$
dont les images sont $t_i \in M_i$, l'image de $\text{Norm}_{L/K}(t)$ dans $M$ est
$\prod \text{Norm}_{M_i/M}(t_i)$. Nous omettons la preuve.
\end{proof}







\section{Équivalence rationnelle}
\label{spaces-chow-section-rational-equivalence}

\noindent
Cette section est l'analogue de l'Homologie de Chow, section
\ref{chow-section-rational-equivalence}. Dans cette section, nous définissons
l'{\it équivalence rationnelle} sur les $k$-cycles. Nous autoriserons des
sommes localement finies d'images de diviseurs principaux (par immersions
fermées). Cela conduit à des phénomènes assez étranges (voir les exemples du
chapitre sur les schémas). Toutefois, sans les autoriser, nous ne savons pas
montrer que l'intersection avec les classes de Chern de
fibrés en droites passe au quotient par l'équivalence rationnelle.

\begin{definition}
\label{spaces-chow-definition-rational-equivalence}
Dans la situation \ref{spaces-chow-situation-setup}, soit $X/B$ bon. Soit $k
\in \mathbf{Z}$.
\begin{enumerate}
\item Étant donné une famille localement finie $\{W_j \subset X\}$ de
sous-espaces fermés intègres telle que $\dim_\delta(W_j) = k + 1$, et un
élément $f_j \in R(W_j)^*$ quelconque, nous pouvons considérer
$$
\sum (i_j)_*\text{div}(f_j) \in Z_k(X)
$$
où $i_j : W_j \to X$ est le morphisme d'inclusion. Cela a du sens puisque le
morphisme $\coprod i_j : \coprod W_j \to X$ est propre.
\item On dit que $\alpha \in Z_k(X)$ est {\it rationnellement équivalent à
zéro} si $\alpha$ est un cycle de la forme affichée ci-dessus.
\item On dit que $\alpha, \beta \in Z_k(X)$ sont {\it rationnellement
équivalents} et on écrit $\alpha \sim_{rat} \beta$ si $\alpha - \beta$ est
rationnellement équivalent à zéro.
\item Nous définissons
$$
\CH_k(X) = Z_k(X) / \sim_{rat}
$$
comme le {\it groupe de Chow des $k$-cycles sur $X$}. On l'appelle parfois le
{\it groupe de Chow des $k$-cycles modulo l'équivalence rationnelle sur $X$}.
\end{enumerate}
\end{definition}

\noindent
Il existe de nombreuses autres relations d'équivalence intéressantes.
L'équivalence rationnelle est la plus grossière de toutes. Un lemme très
simple mais important est le suivant.

\begin{lemma}
\label{spaces-chow-lemma-restrict-to-open}
Dans la situation \ref{spaces-chow-situation-setup}, soit $X/B$ bon. Soit $U
\subset X$ un sous-espace ouvert. Soit $Y$ le sous-espace fermé réduit de $X$
avec $|Y| = |X| \setminus |U|$ et notons $i : Y \to X$ le morphisme
d'inclusion. Soit $k \in \mathbf{Z}$. Supposons que $\alpha, \beta \in
Z_k(X)$. Si $\alpha|_U \sim_{rat} \beta|_U$ alors il existe un cycle $\gamma
\in Z_k(Y)$ tel que
$$
\alpha \sim_{rat} \beta + i_*\gamma.
$$
Autrement dit, la séquence
$$
\xymatrix{
\CH_k(Y) \ar[r]^{i_*} & \CH_k(X) \ar[r]^{j^*} & \CH_k(U) \ar[r] & 0
}
$$
est une suite exacte de groupes abéliens.
\end{lemma}

\begin{proof}
Soit $\{W_j\}_{j \in J}$ une famille localement finie de sous-espaces
fermés intègres de $U$ de $\delta$-dimension $k + 1$, et soient $f_j \in
R(W_j)^*$ des éléments tels que $(\alpha - \beta)|_U = \sum
(i_j)_*\text{div}(f_j)$ comme dans la définition. Soit $W_j' \subset X$ le
sous-espace fermé intègre correspondant de $X$, c'est-à-dire ayant le même
point générique que $W_j$. Supposons que $V \subset X$ soit un ouvert
quasi-compact. Alors $V \cap U$ est également un ouvert quasi-compact de $U$,
car $V$ est noethérien. Ainsi, l'ensemble $\{j \in J \mid W_j \cap V \not =
\emptyset\} = \{j \in J \mid W'_j \cap V \not = \emptyset\}$ est fini puisque
$\{W_j\}$ est localement finie. Autrement dit, $\{W'_j\}$ est également
localement finie. Puisque $R(W_j) = R(W'_j)$, on voit que
$$
\alpha - \beta - \sum (i'_j)_*\text{div}(f_j)
$$
est un cycle sur $X$ dont la restriction à $U$ est nulle. Le lemme découle
alors du lemme \ref{spaces-chow-lemma-exact-sequence-open}.
\end{proof}

\begin{remark}
\label{spaces-chow-remark-infinite-sums-rational-equivalences}
Dans la situation \ref{spaces-chow-situation-setup}, soit $X/B$ bon. Supposons que
soient données des familles infinies $\alpha_i, \beta_i \in Z_k(X)$,
$i \in I$, de $k$-cycles sur $X$. Supposons que les supports de $\alpha_i$ et
$\beta_i$ forment des familles localement finies de parties fermées de $X$, de
sorte que $\sum \alpha_i$ et $\sum \beta_i$ soient définies comme cycles. De
plus, supposons que $\alpha_i \sim_{rat} \beta_i$ pour tout $i$. Il n'est alors
pas clair que $\sum \alpha_i \sim_{rat} \sum \beta_i$. En effet, le problème
est que les équivalences rationnelles peuvent être données
par des familles localement finies $\{W_{i, j}, f_{i, j} \in R(W_{i,
j})^*\}_{j \in J_i}$ mais que l'union $\{W_{i, j}\}_{i \in I, j\in J_i}$ peut
ne pas être localement finie.

\medskip\noindent
En pratique, dans de nombreux cas, on dispose d'une famille localement finie de
parties fermées $\{T_i\}_{i \in I}$ de $|X|$ telle que $\alpha_i,
\beta_i$ soient supportés sur $T_i$ et que $\alpha_i \sim_{rat} \beta_i$ «
sur » $T_i$. Plus précisément, les familles $\{W_{i, j}, f_{i, j} \in R(W_{i,
j})^*\}_{j \in J_i}$ sont constituées de sous-espaces fermés intègres $W_{i,
j}$ tels que $|W_{i, j}| \subset T_i$. Dans ce cas, on a bien $\sum \alpha_i
\sim_{rat} \sum \beta_i$ sur $X$, simplement parce que la famille
$\{W_{i, j}\}_{i \in I, j\in J_i}$ est alors automatiquement localement finie.
\end{remark}











\section{Équivalence rationnelle, image directe et image réciproque}
\label{spaces-chow-section-properties-rational-equivalence}

\noindent
Cette section est l'analogue de l'Homologie de Chow, section
\ref{chow-section-properties-rational-equivalence}. Dans cette section, nous
montrons que l'image réciproque plate et l'image directe propre sont
compatibles avec l'équivalence rationnelle.

\begin{lemma}
\label{spaces-chow-lemma-prepare-flat-pullback-rational-equivalence}
Dans la situation \ref{spaces-chow-situation-setup}, soient $X, Y/B$ bons. Supposons
que $Y$ soit intègre et que $\dim_\delta(Y) = k$. Soit $f : X \to Y$ un morphisme
plat de dimension relative $r$. Alors pour $g \in R(Y)^*$ nous avons
$$
f^*\text{div}_Y(g) =
\sum m_{X', X} (X' \to X)_*\text{div}_{X'}(g \circ f|_{X'})
$$
comme $(k + r - 1)$-cycles sur $X$, où la somme est indexée par les composantes
irréductibles $X'$ de $X$, et où $m_{X', X}$ est la multiplicité de $X'$ dans $X$.
\end{lemma}

\begin{proof}
Remarquons que toute composante irréductible de $X$ domine $Y$ (lemme
\ref{spaces-chow-lemma-flat-inverse-image-dimension}) et donc la composition $g \circ
f|_{X'}$ est définie (Morphismes des espaces, section
\ref{spaces-morphisms-section-rational-maps}). Nous réduirons cela au cas des
schémas. Choisissons un schéma $V$ et un morphisme \'etale surjectif $V \to Y$.
Choisissons un schéma $U$ et un morphisme \'etale surjectif $U \to V \times_Y
X$. Considérons le diagramme $$
\xymatrix{
U \ar[r]_a \ar[d]_h & X \ar[d]^f \\
V \ar[r]^b & Y
}
$$
Puisque $a$ est surjectif et \'etale, il découle du lemme
\ref{spaces-chow-lemma-etale-pullback} qu'il suffit de prouver l'égalité des cycles après
image réciproque par $a$. On peut utiliser le lemme
\ref{spaces-chow-lemma-etale-pullback-principal-divisor} pour écrire
$$
b^*\text{div}_Y(g) = \sum (V' \to V)_*\text{div}_{V'}(g \circ b|_{V'})
$$
où la somme est indexée par les composantes irréductibles $V'$ de $V$. En utilisant le
lemme \ref{spaces-chow-lemma-flat-pullback-proper-pushforward}, on trouve
$$
h^*b^*\text{div}_Y(g) =
\sum (V' \times_V U \to U)_*(h')^*\text{div}_{V'}(g \circ b|_{V'})
$$
où $h' : V' \times_V U \to V'$ est la projection. On peut appliquer le lemme
dans le cas des schémas (Homologie de Chow, lemme
\ref{chow-lemma-prepare-flat-pullback-rational-equivalence}) au morphisme $h'
: V' \times_V U \to V'$ pour voir que l'on a
$$
(h')^*\text{div}_{V'}(g \circ b|_{V'}) =
\sum
m_{U', V' \times_V U}
(U' \to V' \times_V U)_*\text{div}_{U'}(g \circ b|_{V'} \circ h'|_{U'})
$$
où la somme est indexée par les composantes irréductibles $U'$ de $V' \times_V U$.
Chaque $U'$ apparaissant dans cette somme est une composante irréductible de
$U$ et inversement chaque composante irréductible $U'$ de $U$ est une
composante irréductible de $V' \times_V U$ pour une composante irréductible
unique $V' \subset V$. Étant donné une composante irréductible $U' \subset U$,
notons $\overline{a(U')} \subset X$ l'« image » dans $X$ (lemme
\ref{spaces-chow-lemma-proper-image}) ; c'est une composante irréductible de $X$ par
exemple d'après le lemme \ref{spaces-chow-lemma-flat-inverse-image-dimension}. La multiplicité
$m_{U', V' \times_V U}$ est égale à la multiplicité $m_{\overline{a(U')}, X}$.
Cela découle de l'égalité $h^*a^*[Y] = b^*f^*[Y]$ (lemme
\ref{spaces-chow-lemma-compose-flat-pullback}), des définitions et du lemme
\ref{spaces-chow-lemma-etale-pullback}. En combinant tout ce que nous venons de dire, nous
obtenons
$$
a^*f^*\text{div}_Y(g) =
h^*b^*\text{div}_Y(g) =
\sum m_{\overline{a(U')}, X}
(U' \to U)_*\text{div}_{U'}(g \circ (f \circ a)|_{U'})
$$
À présent, analysons ce que devient le membre de droite de la formule de
l'énoncé après image réciproque par $a$. Tout d'abord, le lemme
\ref{spaces-chow-lemma-flat-pullback-proper-pushforward} donne
$$
a^*\sum m_{X', X} (X' \to X)_*\text{div}_{X'}(g \circ f|_{X'}) =
\sum m_{X', X} (X' \times_X U \to U)_*(a')^*\text{div}_{X'}(g \circ f|_{X'})
$$
où $a' : X' \times_X U \to X'$ est la projection. Par le lemme
\ref{spaces-chow-lemma-etale-pullback-principal-divisor}, on obtient
$$
(a')^*\text{div}_{X'}(g \circ f|_{X'}) =
\sum (U' \to X' \times_X U)_*\text{div}_{U'}(g \circ (f \circ a)|_{U'})
$$
où la somme est indexée par les composantes irréductibles $U'$ de $X' \times_X U$. Ces
$U'$ sont des composantes irréductibles de $U$ et sont en fait exactement les
composantes irréductibles de $U$ telles que $\overline{a(U')} = X'$. La comparaison
avec ce qui précède permet de conclure.
\end{proof}

\begin{lemma}
\label{spaces-chow-lemma-flat-pullback-rational-equivalence}
Dans la situation \ref{spaces-chow-situation-setup}, soient $X, Y/B$ bons. Soit $f :
X \to Y$ un morphisme plat de dimension relative $r$. Soit donnée une
équivalence rationnelle $\alpha \sim_{rat} \beta$ entre deux $k$-cycles sur
$Y$. Alors $f^*\alpha
\sim_{rat} f^*\beta$ comme $(k + r)$-cycles sur $X$.
\end{lemma}

\begin{proof}
Que faut-il montrer ? Supposons que l'on se donne une famille
$$
i_j : W_j \longrightarrow Y
$$
d'immersions fermées, chaque $W_j$ étant intègre de $\delta$-dimension $k +
1$, ainsi que des fonctions rationnelles $g_j \in R(W_j)^*$. Supposons en outre que la
famille $\{|i_j|(|W_j|)\}_{j \in J}$ soit localement finie dans $|Y|$. Il
faut alors montrer que
$$
f^*(\sum i_{j, *}\text{div}(g_j)) = \sum f^*i_{j, *}\text{div}(g_j)
$$
est rationnellement équivalent à zéro sur $X$. La somme de droite a un sens
d'après le lemme \ref{spaces-chow-lemma-inverse-image-locally-finite}.

\medskip\noindent
Considérons les produits fibrés
$$
i'_j : W'_j = W_j \times_Y X \longrightarrow X.
$$
et notons $f_j : W'_j \to W_j$ la première projection. Par le lemme
\ref{spaces-chow-lemma-flat-pullback-proper-pushforward}, nous pouvons écrire la somme
ci-dessus sous la forme
$$
\sum i'_{j, *}(f_j^*\text{div}(g_j))
$$
Par le lemme \ref{spaces-chow-lemma-prepare-flat-pullback-rational-equivalence}, on voit
que chaque $f_j^*\text{div}(g_j)$ est rationnellement équivalent à zéro sur
$W'_j$. Par conséquent, chaque $i'_{j, *}(f_j^*\text{div}(g_j))$ est
rationnellement équivalent à zéro. Il en va de même pour la somme affichée par
la discussion dans la remarque
\ref{spaces-chow-remark-infinite-sums-rational-equivalences}.
\end{proof}

\begin{lemma}
\label{spaces-chow-lemma-proper-pushforward-rational-equivalence}
Dans la situation \ref{spaces-chow-situation-setup}, soient $X, Y/B$ bons. Soit $p :
X \to Y$ un morphisme propre. Supposons que $\alpha, \beta \in Z_k(X)$ soient
rationnellement équivalents. Alors $p_*\alpha$ est rationnellement équivalent
à $p_*\beta$.
\end{lemma}

\begin{proof}
Que faut-il montrer ? Supposons que l'on se donne une famille
$$
i_j : W_j \longrightarrow X
$$
d'immersions fermées, chaque $W_j$ étant intègre de $\delta$-dimension $k +
1$, ainsi que des fonctions rationnelles $f_j \in R(W_j)^*$. Supposons en outre
que la famille $\{i_j(W_j)\}_{j \in J}$ soit localement finie dans $X$. Il faut
alors montrer que
$$
p_*\left(\sum i_{j, *}\text{div}(f_j)\right)
$$
est rationnellement équivalent à zéro sur $X$.

\medskip\noindent
Remarquons que la somme est égale à
$$
\sum p_*i_{j, *}\text{div}(f_j).
$$
Soit $W'_j \subset Y$ le sous-espace fermé intègre qui est l'image de $p
\circ i_j$ ; voir le lemme \ref{spaces-chow-lemma-proper-image}. La famille $\{W'_j\}$ est
localement finie dans $Y$ par le lemme
\ref{spaces-chow-lemma-quasi-compact-locally-finite}. Il suffit donc de montrer, pour un
$j$ donné, soit que $p_*i_{j, *}\text{div}(f_j) = 0$, soit que ce cycle soit égal à
$i'_{j, *}\text{div}(g_j)$ pour un certain $g_j \in R(W'_j)^*$.

\medskip\noindent
Les arguments ci-dessus nous ramènent donc au cas d'un seul sous-espace fermé
intègre $W \subset X$ de $\delta$-dimension $k + 1$. Soit $f \in
R(W)^*$. Soit $W' = p(W)$ comme ci-dessus. On obtient un diagramme commutatif
de morphismes
$$
\xymatrix{
W \ar[r]_i \ar[d]_{p'} & X \ar[d]^p \\
W' \ar[r]^{i'} & Y
}
$$
Remarquons que $p_*i_*\text{div}(f) = i'_*(p')_*\text{div}(f)$ d'après le lemme
\ref{spaces-chow-lemma-compose-pushforward}. Comme expliqué ci-dessus, nous devons montrer
que $(p')_*\text{div}(f)$ est le diviseur d'une fonction rationnelle sur $W'$
ou zéro. Il y a trois cas à distinguer.

\medskip\noindent
Le cas $\dim_\delta(W') < k$. Dans ce cas automatiquement $(p')_*\text{div}(f)
= 0$ et il n'y a rien à prouver.

\medskip\noindent
Le cas $\dim_\delta(W') = k$. Montrons que $(p')_*\text{div}(f) = 0$ dans ce
cas. Puisque $(p')_*\text{div}(f)$ est un $k$-cycle, nous voyons que
$(p')_*\text{div}(f) = n[W']$ pour un certain $n \in \mathbf{Z}$. Afin de
prouver que $n = 0$, on peut remplacer $W'$ par un sous-espace ouvert non vide.
En particulier, nous pouvons supposer, et supposons, que $W'$ est
un schéma. Soit $\eta \in W'$ le point générique. Soit $K = \kappa(\eta) =
R(W')$ le corps des fonctions. Considérons le diagramme de changement de base
$$
\xymatrix{
W_\eta \ar[r] \ar[d]_c & W \ar[d]^{p'} \\
\Spec(K) \ar[r]^\eta & W'
}
$$
Remarquons que $c$ est propre. De plus, $|W_\eta|$ est de dimension $1$ : le
lemme \ref{decent-spaces-lemma-topology-fibre} d'Espaces décents identifie
$|W_\eta|$ à la partie de $|W|$ formée des points s'envoyant sur $\eta$ ; or,
puisque $\dim_\delta(W) = k + 1$ et $\delta(\eta) = k$, les points de
$W_\eta$ ont nécessairement une valeur $\delta$ égale à $k$ ou à $k + 1$.
Ainsi, les anneaux locaux sont de dimension $\leq 1$ d'après Espaces décents, lemme
\ref{decent-spaces-lemma-codimension-local-ring}. Par Espaces sur les corps,
lemme \ref{spaces-over-fields-lemma-codim-1-point-in-schematic-locus}, on voit
que $W_\eta$ est un schéma. Puisque $\Spec(K)$ est la limite des
sous-schémas ouverts affines non vides de $W'$, nous concluons que nous
pouvons supposer que $W$ est un schéma d'après Limites d'espaces, lemme
\ref{spaces-limits-lemma-limit-is-scheme}. Enfin, le cas des schémas
(Homologie de Chow, lemme
\ref{chow-lemma-proper-pushforward-rational-equivalence}) donne $n = 0$.

\medskip\noindent
Le cas $\dim_\delta(W') = k + 1$. Dans ce cas, le lemme
\ref{spaces-chow-lemma-proper-pushforward-alteration} s'applique, et on voit qu'en effet
$p'_*\text{div}(f) = \text{div}(g)$ pour un certain $g \in R(W')^*$, comme
voulu.
\end{proof}












\section{Le diviseur associé à un faisceau inversible}
\label{spaces-chow-section-divisor-invertible-sheaf}

\noindent
Cette section est l'analogue de l'Homologie de Chow, section
\ref{chow-section-divisor-invertible-sheaf}. La définition suivante est
l'analogue d'Espaces sur les corps, définition
\ref{spaces-over-fields-definition-divisor-invertible-sheaf}, dans notre cadre.

\begin{definition}
\label{spaces-chow-definition-divisor-invertible-sheaf}
Dans la situation \ref{spaces-chow-situation-setup}, soit $X/B$ bon. Supposons que
$X$ soit intègre et que $n = \dim_\delta(X)$. Soit $\mathcal{L}$ un
$\mathcal{O}_X$-module inversible.
\begin{enumerate}
\item Pour toute section méromorphe non nulle $s$ de $\mathcal{L}$, nous
définissons le {\it diviseur de Weil associé à $s$} comme le $(n - 1)$-cycle
$$
\text{div}_\mathcal{L}(s) =
\sum \text{ord}_{Z, \mathcal{L}}(s) [Z]
$$
défini dans Espaces sur les corps, définition
\ref{spaces-over-fields-definition-divisor-invertible-sheaf}. Cela a un sens
car les diviseurs de Weil sont de $\delta$-dimension $n - 1$ par le lemme
\ref{spaces-chow-lemma-divisor-delta-dimension}.
\item Nous définissons le {\it diviseur de Weil associé à $\mathcal{L}$} comme
$$
c_1(\mathcal{L}) \cap [X] =
\text{class of }\text{div}_\mathcal{L}(s) \in \CH_{n - 1}(X)
$$
où $s$ est une section méromorphe non nulle quelconque de $\mathcal{L}$ sur
$X$. Cette définition ne dépend pas du choix, d'après Espaces sur les corps, lemme
\ref{spaces-over-fields-lemma-divisor-meromorphic-well-defined}.
\end{enumerate}
\end{definition}

\noindent
Le schéma des zéros d'une section non nulle est un diviseur de Cartier effectif
dont la classe de Weil donne le diviseur de Weil associé au module inversible.

\begin{lemma}
\label{spaces-chow-lemma-compute-c1}
Dans la situation \ref{spaces-chow-situation-setup}, soit $X/B$ bon. Supposons que
$X$ soit intègre et que $n = \dim_\delta(X)$. Soit $\mathcal{L}$ un
$\mathcal{O}_X$-module inversible. Soit $s \in \Gamma(X, \mathcal{L})$ une
section globale non nulle. Alors
$$
\text{div}_\mathcal{L}(s) = [Z(s)]_{n - 1}
$$
dans $Z_{n - 1}(X)$ et
$$
c_1(\mathcal{L}) \cap [X] = [Z(s)]_{n - 1}
$$
dans $\CH_{n - 1}(X)$.
\end{lemma}

\begin{proof}
Soit $Z \subset X$ un sous-espace fermé intègre de $\delta$-dimension $n -
1$. Soit $\xi \in |Z|$ son point générique. Pour prouver la première égalité,
comparons les coefficients de $Z$ dans les deux membres. Choisissons un
voisinage \'etale élémentaire $(U, u) \to (X, \xi)$ ; voir Espaces décents, section
\ref{decent-spaces-section-residue-fields-henselian-local-rings}, et rappelons
que $\mathcal{O}_{X, \xi}^h = \mathcal{O}_{U, u}^h$ dans ce cas. Après avoir
remplacé $U$ par un voisinage ouvert de $u$, nous pouvons supposer qu'il existe
une section banalisante $s_U$ de $\mathcal{L}|_U$. Écrivons $s|_U = f s_U$ pour
un certain $f \in \Gamma(U, \mathcal{O}_U)$. Alors $Z \times_X U$ est égal à
$V(f)$ comme sous-schéma fermé de $U$ ; voir Diviseurs sur les espaces,
définition \ref{spaces-divisors-definition-zero-scheme-s}. Comme dans Espaces
sur les corps, section \ref{spaces-over-fields-section-c1}, notons
$\mathcal{L}_\xi$ l'image réciproque de $\mathcal{L}$ par le morphisme canonique
$c_\xi : \Spec(\mathcal{O}_{X, \xi}^h) \to X$. Notons $s_\xi$ l'image réciproque de
$s_U$ ; c'est une banalisation de $\mathcal{L}_\xi$. Alors
$c_\xi^*(s) = fs_\xi$. Le coefficient de $Z$ dans $[Z(s)]_{n - 1}$ est par
définition
$$
\text{length}_{\mathcal{O}_{U, u}}(\mathcal{O}_{U, u}/f\mathcal{O}_{U, u})
$$
Puisque $\mathcal{O}_{U, u} \to \mathcal{O}_{X, \xi}^h$ est plat et identifie
les corps résiduels, cela est égal à
$$
\text{length}_{\mathcal{O}_{X, \xi}^h}
(\mathcal{O}_{X, \xi}^h/f\mathcal{O}_{X, \xi}^h)
$$
par Algèbre, lemme \ref{algebra-lemma-pullback-module}. Cette dernière quantité
est égale à $\text{ord}_{Z, \mathcal{L}}(s)$ par Espaces sur les corps,
définition \ref{spaces-over-fields-definition-order-vanishing-meromorphic},
c'est-à-dire au coefficient de $Z$ dans $\text{div}_\mathcal{L}(s)$ comme
voulu.
\end{proof}

\begin{lemma}
\label{spaces-chow-lemma-Gm-torsor}
Dans la situation \ref{spaces-chow-situation-setup}, soit $X/B$ bon. Soit
$\mathcal{L}$ un $\mathcal{O}_X$-module inversible. Le morphisme
$$
q :
T = \underline{\Spec}\left(
\bigoplus\nolimits_{n \in \mathbf{Z}} \mathcal{L}^{\otimes n}\right)
\longrightarrow
X
$$
a les propriétés suivantes :
\begin{enumerate}
\item $q$ est surjectif, lisse, affine, de dimension relative $1$,
\item il existe un isomorphisme $\alpha : q^*\mathcal{L} \cong \mathcal{O}_T$,
\item la formation de $(q : T \to X, \alpha)$ commute aux changements de
base,
\item $q^* : Z_k(X) \to Z_{k + 1}(T)$ est injectif,
\item si $Z \subset X$ est un sous-espace fermé intègre, alors $q^{-1}(Z)
\subset T$ est un sous-espace fermé intègre,
\item si $Z \subset X$ est un sous-espace fermé de $X$ de $\delta$-dimension
$\leq k$, alors $q^{-1}(Z)$ est un sous-espace fermé de $T$ de
$\delta$-dimension $\leq k + 1$ et $q^*[Z]_k = [q^{-1}(Z)]_{k + 1}$,
\item si $\xi' \in |T|$ est le point générique de la fibre de $|T| \to |X|$
sur $\xi$, alors l'homomorphisme d'anneaux $\mathcal{O}_{X, \xi}^h \to
\mathcal{O}_{T, \xi'}^h$ est plat, nous avons $\mathfrak m_{\xi'}^h =
\mathfrak m_\xi^h \mathcal{O}_{T, \xi'}^h$, et l'extension de corps résiduels
est purement transcendante, de degré de transcendance $1$, et
\item ajouter ici d'autres propriétés au besoin.
\end{enumerate}
\end{lemma}

\begin{proof}
Soit $U \to X$ un morphisme \'etale tel que $\mathcal{L}|_U$ soit trivial.
Alors $T \times_X U \to U$ est isomorphe au morphisme de projection
$\mathbf{G}_m \times U \to U$, où $\mathbf{G}_m$ est le schéma en groupes
multiplicatif ; voir Groupoïdes, exemple
\ref{groupoids-example-multiplicative-group}. Ainsi (1) est clair.

\medskip\noindent
Pour (2), remarquons que $q_*q^*\mathcal{L} = \bigoplus_{n \in \mathbf{Z}}
\mathcal{L}^{\otimes n + 1}$. Il existe donc un isomorphisme évident
$q_*q^*\mathcal{L} \to q_*\mathcal{O}_T$ de $q_*\mathcal{O}_T$-modules. D'après
Morphismes d'espaces, lemme
\ref{spaces-morphisms-lemma-affine-equivalence-modules}, cela détermine un
isomorphisme $q^*\mathcal{L} \to \mathcal{O}_T$.

\medskip\noindent
La partie (3) résulte du fait que la formation du spectre relatif commute à
tout changement de base ; il en va clairement de même pour
l'isomorphisme $\alpha$.

\medskip\noindent
La partie (4) découle immédiatement de (1) et des définitions.

\medskip\noindent
La partie (5) découle du fait que si $Z$ est un espace algébrique intègre,
alors $\mathbf{G}_m \times Z$ est un espace algébrique intègre.

\medskip\noindent
La partie (6) découle de la conservation des longueurs : si $(A,
\mathfrak m)$ est un anneau local, si $B = A[x]_{\mathfrak m A[x]}$ et si $M$
est un $A$-module, alors $\text{length}_A(M) = \text{length}_B(M \otimes_A
B)$. Cela implique que, si $\mathcal{F}$ est un $\mathcal{O}_X$-module cohérent,
si $\xi \in |X|$ et si $\xi' \in |T|$ est le point générique de la fibre
au-dessus de $\xi$, alors la longueur de $\mathcal{F}$ en $\xi$ est égale à celle de
$q^*\mathcal{F}$ en $\xi'$. En remontant les définitions, on obtient (6), et
même davantage.

\medskip\noindent
L'homomorphisme de la partie (7) provient d'Espaces décents, remarque
\ref{decent-spaces-remark-functoriality-henselian-local-ring}. Cependant, dans
notre cas, nous avons
$$
\Spec(\mathcal{O}_{X, \xi}^h) \times_X T =
\mathbf{G}_m \times \Spec(\mathcal{O}_{X, \xi}^h) =
\Spec(\mathcal{O}_{X, \xi}^h[t, t^{-1}])
$$
et $\xi'$ correspond au point générique de la fibre spéciale de ce morphisme au-dessus de
$\Spec(\mathcal{O}_{X, \xi}^h)$. Ainsi $\mathcal{O}_{T, \xi'}^h$ est
l'hensélisation de la localisation de $\mathcal{O}_{X, \xi}^h[t, t^{-1}]$ en
l'idéal premier correspondant. La partie (7) en découle par quelques résultats
d'algèbre commutative ; détails omis.
\end{proof}

\begin{lemma}
\label{spaces-chow-lemma-Gm-torsor-divisor-meromorphic-section}
Dans la situation \ref{spaces-chow-situation-setup}, soit $X/B$ bon. Soit
$\mathcal{L}$ un $\mathcal{O}_X$-module inversible. Supposons que $X$ soit
intègre. Soit $s$ une section méromorphe non nulle de $\mathcal{L}$. Soit $q
: T \to X$ le morphisme du lemme \ref{spaces-chow-lemma-Gm-torsor}. Alors
$$
q^*\text{div}_\mathcal{L}(s) = \text{div}_T(q^*(s))
$$
où nous considérons l'image réciproque $q^*(s)$ comme une fonction méromorphe non
nulle sur $T$ au moyen de l'isomorphisme $q^*\mathcal{L} \to \mathcal{O}_T$.
\end{lemma}

\begin{proof}
Remarquons que $\text{div}_T(q^*(s)) = \text{div}_{\mathcal{O}_T}(q^*(s))$ par
la compatibilité entre les constructions données dans Espaces sur les corps,
sections \ref{spaces-over-fields-section-Weil-divisors} et
\ref{spaces-over-fields-section-c1}. Dans cette preuve, nous établirons
l'égalité avec $\text{div}_{\mathcal{O}_T}(q^*(s))$. Nous utiliserons
toutes les propriétés de $q : T \to X$ indiquées dans le lemme
\ref{spaces-chow-lemma-Gm-torsor} sans autre mention. Soit $Z \subset T$ un diviseur
premier. Ou bien $Z \to X$ est dominant, ou bien $Z = q^{-1}(Z')$ pour un
diviseur premier $Z' \subset X$. Si $Z \to X$ est dominant, alors le
coefficient de $Z$ dans chacun des deux membres de l'égalité du lemme est nul. On peut
donc supposer $Z = q^{-1}(Z')$ où $Z' \subset X$ est un diviseur premier. Soit
$\xi' \in |Z'|$ et $\xi \in |Z|$ les points génériques. On obtient alors un
diagramme commutatif
$$
\xymatrix{
\Spec(\mathcal{O}_{T, \xi}^h) \ar[r]_-{c_\xi} \ar[d]_h & T \ar[d]^q \\
\Spec(\mathcal{O}_{X, \xi'}^h) \ar[r]^-{c_{\xi'}} & X
}
$$
voir Espaces décents, remarque
\ref{decent-spaces-remark-functoriality-henselian-local-ring}. Choisissons une
banalisation $s_{\xi'}$ de $\mathcal{L}_{\xi'} = c_{\xi'}^*\mathcal{L}$.
Nous pouvons alors utiliser l'image réciproque $s_\xi$ de $s_{\xi'}$ par $h$
comme banalisation de $\mathcal{L}_\xi = c_\xi^* q^*\mathcal{L}$. Écrivons
$s/s_{\xi'} = a/b$ avec $a, b \in \mathcal{O}_{X, \xi'}$ non-diviseurs de
zéro. Par définition, le coefficient de $Z'$ dans $\text{div}_\mathcal{L}(s)$
est
$$
\text{length}_{\mathcal{O}_{X, \xi'}^h}(
\mathcal{O}_{X, \xi'}^h/a \mathcal{O}_{X, \xi'}^h)
-
\text{length}_{\mathcal{O}_{X, \xi'}^h}(
\mathcal{O}_{X, \xi'}^h/b \mathcal{O}_{X, \xi'}^h)
$$
Puisque $Z = q^{-1}(Z')$, c'est aussi le coefficient de $Z$ dans
$q^*\text{div}_\mathcal{L}(s)$. Étant donné que $\mathcal{O}_{X, \xi'}^h \to
\mathcal{O}_{T, \xi}^h$ est plat, les éléments $a, b$ sont envoyés sur des
non-diviseurs de zéro dans $\mathcal{O}_{T, \xi}^h$. Ainsi $q^*(s)/s_\xi = a/b$
dans l'anneau total des fractions de $\mathcal{O}_{T, \xi}^h$. Par définition, le
coefficient de $Z$ dans $\text{div}_T(q^*(s))$ est
$$
\text{length}_{\mathcal{O}_{T, \xi}^h}(
\mathcal{O}_{T, \xi}^h/a \mathcal{O}_{T, \xi}^h)
-
\text{length}_{\mathcal{O}_{T, \xi}^h}(
\mathcal{O}_{T, \xi}^h/b \mathcal{O}_{T, \xi}^h)
$$
La preuve est achevée, car ces longueurs sont les mêmes que précédemment
d'après Algèbre, lemme \ref{algebra-lemma-pullback-module}, et l'égalité
$\mathfrak m_\xi^h = \mathfrak m_{\xi'}^h\mathcal{O}_{T, \xi}^h$ établie dans
le lemme \ref{spaces-chow-lemma-Gm-torsor}.
\end{proof}









\section{Intersection avec un faisceau inversible}
\label{spaces-chow-section-intersecting-with-divisors}

\noindent
Cette section est l'analogue de l'Homologie de Chow, section
\ref{chow-section-intersecting-with-divisors}. Dans cette section, nous
étudions la construction suivante.

\begin{definition}
\label{spaces-chow-definition-cap-c1}
Dans la situation \ref{spaces-chow-situation-setup}, soit $X/B$ bon. Soit
$\mathcal{L}$ un $\mathcal{O}_X$-module inversible. On définit, pour tout
entier $k$, une opération
$$
c_1(\mathcal{L}) \cap - :
Z_{k + 1}(X) \to \CH_k(X)
$$
appelée {\it intersection avec la première classe de Chern de $\mathcal{L}$}.
\begin{enumerate}
\item Étant donné un sous-espace fermé intègre $i : W \to X$ avec
$\dim_\delta(W) = k + 1$, nous définissons
$$
c_1(\mathcal{L}) \cap [W] = i_*(c_1({i^*\mathcal{L}}) \cap [W])
$$
où le membre de droite est défini dans la définition
\ref{spaces-chow-definition-divisor-invertible-sheaf}.
\item Pour un $(k + 1)$-cycle général $\alpha = \sum n_i [W_i]$, nous
définissons
$$
c_1(\mathcal{L}) \cap \alpha = \sum n_i c_1(\mathcal{L}) \cap [W_i]
$$
\end{enumerate}
\end{definition}

\noindent
Écrivons chaque $c_1(\mathcal{L}) \cap W_i = \sum_j n_{i, j} [Z_{i, j}]$, où
$\{Z_{i, j}\}_j$ est une famille localement finie de sous-espaces fermés intègres
de $W_i$. Puisque $\{W_i\}$ est une famille localement finie de
sous-espaces fermés intègres de $X$, il s'ensuit facilement que $\{Z_{i,
j}\}_{i, j}$ est une famille localement finie de sous-espaces fermés de
$X$. $c_1(\mathcal{L}) \cap \alpha = \sum n_in_{i, j}[Z_{i, j}]$ est donc un
cycle. Une autre interprétation, souvent plus commode, consiste à observer
que le morphisme $\coprod W_i \to X$ est propre. Par conséquent,
$c_1(\mathcal{L}) \cap \alpha$ peut être considéré comme l'image directe d'une
classe dans $\CH_k(\coprod W_i) = \prod \CH_k(W_i)$. Cela explique aussi
pourquoi le résultat est bien défini modulo l'équivalence rationnelle sur $X$.

\medskip\noindent
L'objectif principal des prochaines sections est de montrer que l'intersection
par $c_1(\mathcal{L})$ passe au quotient par l'équivalence rationnelle. Cela
n'est pas immédiat.

\begin{lemma}
\label{spaces-chow-lemma-c1-cap-additive}
Dans la situation \ref{spaces-chow-situation-setup}, soit $X/B$ bon. Soient
$\mathcal{L}$, $\mathcal{N}$ des faisceaux inversibles sur $X$. Alors
$$
c_1(\mathcal{L}) \cap \alpha  + c_1(\mathcal{N}) \cap \alpha =
c_1(\mathcal{L} \otimes_{\mathcal{O}_X} \mathcal{N}) \cap \alpha
$$
dans $\CH_k(X)$ pour tout $\alpha \in Z_{k - 1}(X)$. De plus,
$c_1(\mathcal{O}_X) \cap \alpha = 0$ pour tout $\alpha$.
\end{lemma}

\begin{proof}
L'additivité découle directement d'Espaces sur les corps, lemme
\ref{spaces-over-fields-lemma-c1-additive} et des définitions. Pour voir que
$c_1(\mathcal{O}_X) \cap \alpha = 0$, considérons la section $1 \in \Gamma(X,
\mathcal{O}_X)$. Celle-ci se restreint en une section partout non nulle sur tout
sous-espace fermé intègre $W \subset X$. D'où $c_1(\mathcal{O}_X) \cap [W] =
0$, comme voulu.
\end{proof}

\noindent
Rappelons que $Z(s) \subset X$ désigne le schéma des zéros d'une section globale $s$
d'un faisceau inversible sur un espace algébrique $X$, voir Diviseurs sur les
espaces, définition \ref{spaces-divisors-definition-zero-scheme-s}.

\begin{lemma}
\label{spaces-chow-lemma-prepare-geometric-cap}
Dans la situation \ref{spaces-chow-situation-setup}, soit $Y/B$ bon. Soit
$\mathcal{L}$ un $\mathcal{O}_Y$-module inversible. Soit $s \in \Gamma(Y,
\mathcal{L})$ une section régulière et supposons $\dim_\delta(Y) \leq k + 1$.
Écrivons $[Y]_{k + 1} = \sum n_i[Y_i]$ où $Y_i \subset Y$ sont les composantes
irréductibles de $Y$ de $\delta$-dimension $k + 1$. Posons $s_i = s|_{Y_i}
\in \Gamma(Y_i, \mathcal{L}|_{Y_i})$. Alors
\begin{equation}
\label{spaces-chow-equation-equal-as-cycles}
[Z(s)]_k =  \sum n_i[Z(s_i)]_k
\end{equation}
comme $k$-cycles sur $Y$.
\end{lemma}

\begin{proof}
Soit $\varphi : V \to Y$ un morphisme \'etale surjectif, où $V$ est un schéma.
Il suffit de démontrer l'égalité après image réciproque par $\varphi$ ; voir le lemme
\ref{spaces-chow-lemma-etale-pullback}. Ce même lemme nous dit que $\varphi^*[Y_i] =
[\varphi^{-1}(Y_i)] = \sum [V_{i, j}]$ où $V_{i, j}$ sont les composantes
irréductibles de $V$ situées au-dessus de $Y_i$. Par conséquent, si nous
appliquons d'abord le cas des schémas (Homologie de Chow, lemme
\ref{chow-lemma-prepare-geometric-cap}) à $\varphi^*s_i$ sur $Y_i \times_Y V$,
nous trouvons que $\varphi^*[Z(s_i)]_k = [Z(\varphi^*s_i)] = \sum [Z(s_{i,
j})]_k$ où $s_{i, j}$ est l'image réciproque de $s$ sur $V_{i, j}$. En
appliquant le cas des schémas à $\varphi^*s$, on obtient
$$
\varphi^*[Z(s)]_k =
[Z(\varphi^*s)]_k =
\sum n_i[Z(s_{i, j})]_k
$$
par notre remarque sur les multiplicités ci-dessus. En combinant tout ce qui
précède, la preuve est complète.
\end{proof}

\noindent
Le lemme suivant permet de calculer le produit d'intersection de la classe
$c_1$ d'un faisceau inversible avec le cycle associé à un sous-schéma fermé.
Rappelons que $Z(s) \subset X$ désigne le schéma des zéros d'une section globale $s$
d'un faisceau inversible sur un schéma $X$ ; voir Diviseurs, définition
\ref{divisors-definition-zero-scheme-s}.

\begin{lemma}
\label{spaces-chow-lemma-geometric-cap}
Dans la situation \ref{spaces-chow-situation-setup}, soit $X/B$ bon. Soit
$\mathcal{L}$ un $\mathcal{O}_X$-module inversible. Soit $Y \subset X$ un
sous-schéma fermé avec $\dim_\delta(Y) \leq k + 1$, et soit $s \in \Gamma(Y,
\mathcal{L}|_Y)$ une section régulière. Alors
$$
c_1(\mathcal{L}) \cap [Y]_{k + 1} = [Z(s)]_k
$$
dans $\CH_k(X)$.
\end{lemma}

\begin{proof}
Écrivons
$$
[Y]_{k + 1} = \sum n_i[Y_i]
$$
où $Y_i \subset Y$ sont les composantes irréductibles de $Y$ de $\delta$-dimension $k + 1$ et $n_i > 0$. Par hypothèse, la restriction $s_i = s|_{Y_i}
\in \Gamma(Y_i, \mathcal{L}|_{Y_i})$ n'est pas nulle et constitue donc une
section régulière. D'après le lemme \ref{spaces-chow-lemma-compute-c1},
$[Z(s_i)]_k$ représente $c_1(\mathcal{L}|_{Y_i})$. Par définition,
$$
c_1(\mathcal{L}) \cap [Y]_{k + 1} = \sum n_i[Z(s_i)]_k
$$
Le résultat découle donc du lemme \ref{spaces-chow-lemma-prepare-geometric-cap}.
\end{proof}









\section{Intersection avec un faisceau inversible, image directe et image réciproque}
\label{spaces-chow-section-intersecting-with-divisors-push-pull}

\noindent
Cette section est l'analogue de l'Homologie de Chow, section
\ref{chow-section-intersecting-with-divisors-push-pull}. Dans cette section,
nous prouvons que l'opération $c_1(\mathcal{L}) \cap -$ commute avec
l'image réciproque plate et l'image directe propre.

\begin{lemma}
\label{spaces-chow-lemma-prepare-flat-pullback-cap-c1}
Dans la situation \ref{spaces-chow-situation-setup}, soient $X, Y/B$ bons. Soit $f :
X \to Y$ un morphisme plat de dimension relative $r$. Soit $\mathcal{L}$ un
faisceau inversible sur $Y$. Supposons $Y$ intègre et posons $n =
\dim_\delta(Y)$. Soit $s$ une section méromorphe non nulle de $\mathcal{L}$.
Alors
$$
f^*\text{div}_\mathcal{L}(s) = \sum n_i\text{div}_{f^*\mathcal{L}|_{X_i}}(s_i)
$$
dans $Z_{n + r - 1}(X)$. Ici, la somme porte sur les composantes irréductibles
$X_i \subset X$ de $\delta$-dimension $n + r$, la section $s_i =
f|_{X_i}^*(s)$ est l'image réciproque de $s$ et $n_i = m_{X_i, X}$ est la multiplicité
de $X_i$ dans $X$.
\end{lemma}

\begin{proof}
Par une astuce, nous allons déduire cet énoncé du lemme
\ref{spaces-chow-lemma-prepare-flat-pullback-rational-equivalence}. (Une alternative est
de refaire la preuve de ce lemme dans le cadre de sections méromorphes de
modules inversibles.) Plus précisément, soit $q : T \to Y$ le morphisme du lemme
\ref{spaces-chow-lemma-Gm-torsor} construit en utilisant $\mathcal{L}$ sur $Y$. Nous
utiliserons toutes les propriétés de $T$ énoncées dans ce lemme. Considérons le
diagramme cartésien
$$
\xymatrix{
T' \ar[r]_{q'} \ar[d]_h & X \ar[d]^f \\
T \ar[r]^q & Y
}
$$
Alors $q' : T' \to X$ est le morphisme construit en utilisant $f^*\mathcal{L}$
sur $X$. Il suffit alors de démontrer
$$
(q')^*f^*\text{div}_\mathcal{L}(s) =
\sum n_i (q')^*\text{div}_{f^*\mathcal{L}|_{X_i}}(s_i)
$$
Remarquons que $T'_i = q^{-1}(X_i)$ sont les composantes irréductibles de $T'$
et que $n_i$ est la multiplicité de $T'_i$ dans $T'$. Le membre de gauche est
égal à
$$
h^*q^*\text{div}_\mathcal{L}(s) = h^*\text{div}_T(q^*(s))
$$
par le lemme \ref{spaces-chow-lemma-Gm-torsor-divisor-meromorphic-section} (et le lemme
\ref{spaces-chow-lemma-compose-flat-pullback}). D'autre part, si $q'_i : T'_i \to
X_i$ désigne la restriction de $q'$, nous constatons que le lemme
\ref{spaces-chow-lemma-Gm-torsor-divisor-meromorphic-section} nous dit également que le
membre de droite est égal à
$$
\sum n_i \text{div}_{T_i}((q'_i)^*(s_i))
$$
Dans ces deux formules, les expressions $q^*(s)$ et $(q'_i)^*(s_i)$
représentent les fonctions rationnelles correspondant aux sections méromorphes
de $q^*\mathcal{L}$ et de $(q'_i)^*f^*\mathcal{L}|_{X_i}$ obtenues par image
réciproque ; l'identification se fait au moyen de l'isomorphisme
$\alpha : q^*\mathcal{L} \to \mathcal{O}_T$ et de ses images réciproques sur les
espaces au-dessus de $T$. Avec cette convention, il est clair que
$(q'_i)^*(s_i)$ est la composition de la fonction rationnelle $q^*(s)$ sur $T$
et du morphisme $h|_{T'_i} : T'_i \to T$. Ainsi le lemme
\ref{spaces-chow-lemma-prepare-flat-pullback-rational-equivalence} dit exactement que
$$
h^*\text{div}_T(q^*(s)) = \sum n_i \text{div}_{T_i}((q'_i)^*(s_i))
$$
comme voulu.
\end{proof}

\begin{lemma}
\label{spaces-chow-lemma-flat-pullback-cap-c1}
Dans la situation \ref{spaces-chow-situation-setup}, soient $X, Y/B$ bons. Soit $f :
X \to Y$ un morphisme plat de dimension relative $r$. Soit $\mathcal{L}$ un
faisceau inversible sur $Y$. Soit $\alpha$ un $k$-cycle sur $Y$. Alors
$$
f^*(c_1(\mathcal{L}) \cap \alpha) = c_1(f^*\mathcal{L}) \cap f^*\alpha
$$
dans $\CH_{k + r - 1}(X)$.
\end{lemma}

\begin{proof}
Écrivons $\alpha = \sum n_i[W_i]$. Nous allons montrer que
$$
f^*(c_1(\mathcal{L}) \cap [W_i]) = c_1(f^*\mathcal{L}) \cap f^*[W_i]
$$
dans $\CH_{k + r - 1}(X)$ en produisant une équivalence rationnelle sur le
sous-espace fermé $f^{-1}(W_i)$ de $X$. D'après la discussion de la remarque
\ref{spaces-chow-remark-infinite-sums-rational-equivalences}, cela prouvera que l'égalité
du lemme est vraie.

\medskip\noindent
Soit $W \subset Y$ un sous-espace fermé intègre de $\delta$-dimension $k$.
Considérons le sous-espace fermé $W' = f^{-1}(W) = W \times_Y X$, ce qui donne
le diagramme cartésien
$$
\xymatrix{
W' \ar[r] \ar[d]_h & X \ar[d]^f \\
W \ar[r] & Y
}
$$
Nous devons montrer que $f^*(c_1(\mathcal{L}) \cap [W]) = c_1(f^*\mathcal{L})
\cap f^*[W]$. Choisissons une section méromorphe non nulle $s$ de
$\mathcal{L}|_W$. Soient $W'_i \subset W'$ les composantes irréductibles de
$\delta$-dimension $k + r$. Écrivons $[W']_{k + r} = \sum n_i[W'_i]$ avec $n_i$
la multiplicité de $W'_i$ dans $W'$ selon la définition. Donc $f^*[W] = \sum
n_i[W'_i]$ dans $Z_{k + r}(X)$. Puisque chaque $W'_i \to W$ est dominant, nous
voyons que $s_i = s|_{W'_i}$ est une section méromorphe non nulle pour tout
$i$. Par le lemme \ref{spaces-chow-lemma-prepare-flat-pullback-cap-c1}, nous avons
l'égalité de cycles suivante
$$
h^*\text{div}_{\mathcal{L}|_W}(s) =
\sum n_i\text{div}_{f^*\mathcal{L}|_{W'_i}}(s_i)
$$
dans $Z_{k + r - 1}(W')$. Ceci achève la preuve, puisque le membre de gauche est un
cycle sur $W'$ dont l'image directe est $f^*(c_1(\mathcal{L}) \cap [W])$ dans $\CH_{k +
r - 1}(X)$ et le membre de droite est un cycle sur $W'$ dont l'image directe est
$c_1(f^*\mathcal{L}) \cap f^*[W]$ dans $\CH_{k + r - 1}(X)$.
\end{proof}

\begin{lemma}
\label{spaces-chow-lemma-equal-c1-as-cycles}
Dans la situation \ref{spaces-chow-situation-setup}, soient $X, Y/B$ bons. Soit $f :
X \to Y$ un morphisme propre. Soit $\mathcal{L}$ un faisceau inversible sur $Y$.
Supposons $X$, $Y$ intègres, $f$ dominant et $\dim_\delta(X) =
\dim_\delta(Y)$. Soit $s$ une section méromorphe non nulle, la notation $s$ désignant une section de
$\mathcal{L}$ sur $Y$. Alors
$$
f_*\left(\text{div}_{f^*\mathcal{L}}(f^*s)\right) =
[R(X) : R(Y)]\text{div}_\mathcal{L}(s).
$$
Il s'agit d'une égalité de cycles sur $Y$. En particulier,
$$
f_*(c_1(f^*\mathcal{L}) \cap [X]) = c_1(\mathcal{L}) \cap f_*[Y].
$$
\end{lemma}

\begin{proof}
La dernière égalité découle de la première puisque $f_*[X] = [R(X) :
R(Y)][Y]$ par définition. Démontrons la première égalité. Soit $q : T \to Y$ le
morphisme du lemme \ref{spaces-chow-lemma-Gm-torsor} construit en utilisant $\mathcal{L}$
sur $Y$. Nous utiliserons toutes les propriétés de $T$ énoncées dans ce lemme.
Considérons le diagramme cartésien
$$
\xymatrix{
T' \ar[r]_{q'} \ar[d]_h & X \ar[d]^f \\
T \ar[r]^q & Y
}
$$
Alors $q' : T' \to X$ est le morphisme construit en utilisant $f^*\mathcal{L}$
sur $X$. Il suffit de démontrer l'égalité après image réciproque sur $T'$. L'image
réciproque du membre de gauche vaut
\begin{align*}
q^*f_*\left(\text{div}_{f^*\mathcal{L}}(f^*s)\right)
& =
h_*(q')^*\text{div}_{f^*\mathcal{L}}(f^*s) \\
& =
h_*\text{div}_{(q')^*f^*\mathcal{L}}((q')^*f^*s) \\
& =
h_*\text{div}_{h^*q^*\mathcal{L}}(h^*q^*s)
\end{align*}
La première égalité découle du lemme \ref{spaces-chow-lemma-flat-pullback-proper-pushforward}.
La deuxième découle du lemme \ref{spaces-chow-lemma-prepare-flat-pullback-cap-c1}, puisque
$T'$ est intègre. La troisième résulte de la commutativité du diagramme affiché.
L'image réciproque du membre de droite est
$$
[R(X) : R(Y)]q^*\text{div}_\mathcal{L}(s) =
[R(T') : R(T)]\text{div}_{q^*\mathcal{L}}(q^*s)
$$
Cela découle du lemme \ref{spaces-chow-lemma-prepare-flat-pullback-cap-c1}, du fait que
$T$ est intègre, et de l'égalité $[R(T') : R(T)] = [R(X) : R(Y)]$ dont nous
omettons la preuve (cela découle du lemme
\ref{spaces-chow-lemma-flat-pullback-proper-pushforward}, mais ce serait une manière
inutilement détournée de démontrer l'égalité). Il suffit donc de démontrer le lemme pour $h : T' \to T$,
le module inversible $q^\mathcal{L}$ et la section $q^*s$. Puisque
$q^*\mathcal{L} = \mathcal{O}_T$, nous nous ramenons au cas où $\mathcal{L} \cong
\mathcal{O}$ discuté dans le paragraphe suivant.

\medskip\noindent
Supposons que $\mathcal{L} = \mathcal{O}_Y$. Dans ce cas $s$ correspond à une
fonction rationnelle $g \in R(Y)$. En utilisant le plongement $R(Y) \subset
R(X)$, nous pouvons considérer $g$ comme une fonction rationnelle sur $X$ et il
s'agit simplement de démontrer
$$
f_*\left(\text{div}_X(g)\right) = [R(X) : R(Y)]\text{div}_Y(g).
$$
En comparant avec le résultat du lemme
\ref{spaces-chow-lemma-proper-pushforward-alteration}, nous voyons que cela est vrai
puisque $\text{Nm}_{R(X)/R(Y)}(g) = g^{[R(X) : R(Y)]}$ pour $g \in R(Y)^*$.
\end{proof}

\begin{lemma}
\label{spaces-chow-lemma-pushforward-cap-c1}
Dans la situation \ref{spaces-chow-situation-setup}, soient $X, Y/B$ bons. Soit $p :
X \to Y$ un morphisme propre. Soit $\alpha \in Z_{k + 1}(X)$. Soit
$\mathcal{L}$ un faisceau inversible sur $Y$. Alors
$$
p_*(c_1(p^*\mathcal{L}) \cap \alpha) = c_1(\mathcal{L}) \cap p_*\alpha
$$
dans $\CH_k(Y)$.
\end{lemma}

\begin{proof}
Supposons que $p$ ait la propriété que, pour tout sous-espace fermé intègre
$W \subset X$, l'application $p|_W : W \to Y$ est une immersion fermée.
Alors, par définition de l'intersection avec $c_1(\mathcal{L})$, le lemme est
valable.

\medskip\noindent
Nous allons utiliser cette remarque pour nous ramener à un cas particulier. Plus précisément,
écrivons $\alpha = \sum n_i[W_i]$ avec $n_i \not = 0$ et $W_i$ distincts deux à
deux. Soit $W'_i \subset Y$ l'« image » de $W_i$ comme dans le lemme
\ref{spaces-chow-lemma-proper-image}. Considérons le diagramme
$$
\xymatrix{
X' = \coprod W_i \ar[r]_-q \ar[d]_{p'} & X \ar[d]^p \\
Y' = \coprod W'_i \ar[r]^-{q'} & Y.
}
$$
Puisque $\{W_i\}$ est localement fini sur $X$ et que $p$ est propre, on voit
que $\{W'_i\}$ est localement fini sur $Y$ et que $q, q', p'$ sont eux aussi
des morphismes propres. Nous pouvons également considérer $\sum n_i[W_i]$
comme un $k$-cycle $\alpha' \in Z_k(X')$. Clairement, $q_*\alpha' = \alpha$.
Nous avons $q_*(c_1(q^*p^*\mathcal{L}) \cap \alpha') = c_1(p^*\mathcal{L})
\cap q_*\alpha'$ et $(q')_*(c_1((q')^*\mathcal{L}) \cap p'_*\alpha') =
c_1(\mathcal{L}) \cap q'_*p'_*\alpha'$ par la remarque initiale de la preuve.
Il suffit donc de démontrer le lemme pour le morphisme $p'$ et le cycle $\sum
n_i[W_i]$. Cela signifie que nous pouvons supposer $X$, $Y$
intègres, $f : X \to Y$ dominant et $\alpha = [X]$. Dans ce cas, le résultat
découle du lemme \ref{spaces-chow-lemma-equal-c1-as-cycles}.
\end{proof}













\section{La formule clé}
\label{spaces-chow-section-key}

\noindent
Cette section est l'analogue de l'Homologie de Chow, section
\ref{chow-section-key}. Nous recommandons vivement au lecteur de lire d'abord
la preuve dans ce cas.

\medskip\noindent
Dans la situation \ref{spaces-chow-situation-setup}, soit $X/B$ bon. Supposons
$X$ intègre et $\dim_\delta(X) = n$. Soient $\mathcal{L}$ et
$\mathcal{N}$ des $\mathcal{O}_X$-modules inversibles. Soit $s$ une section
méromorphe non nulle de $\mathcal{L}$ et soit $t$ une section méromorphe non
nulle de $\mathcal{N}$. Soit $Z \subset X$ un diviseur premier ayant pour point
générique $\xi \in |Z|$. Considérons le morphisme
$$
c_\xi : \Spec(\mathcal{O}_{X, \xi}^h) \longrightarrow X
$$
utilisé dans Espaces sur les corps, section \ref{spaces-over-fields-section-c1}.
Nous notons $\mathcal{L}_\xi$ et $\mathcal{N}_\xi$ les images réciproques de
$\mathcal{L}$ et $\mathcal{N}$ par $c_\xi$ ; nous considérons souvent
$\mathcal{L}_\xi$ et $\mathcal{N}_\xi$ comme
les modules libres de rang $1$ sur $\mathcal{O}_{X, \xi}^h$ qu'ils définissent.
Notons que l'image réciproque de $s$, resp.\ $t$, est une section méromorphe régulière de $\mathcal{L}_\xi$,
resp.\ $\mathcal{N}_\xi$.

\medskip\noindent
Soit $Z_i \subset X$, $i \in I$, une famille localement finie de diviseurs
premiers ayant la propriété suivante : si $Z \not \in \{Z_i\}$, alors $s$ est
un générateur de $\mathcal{L}_\xi$ et $t$ est un générateur de
$\mathcal{N}_\xi$. Une telle famille existe d'après le lemme d'Espaces sur les corps
\ref{spaces-over-fields-lemma-divisor-meromorphic-locally-finite}. Alors
$$
\text{div}_\mathcal{L}(s) = \sum \text{ord}_{Z_i, \mathcal{L}}(s) [Z_i]
$$
et de même
$$
\text{div}_\mathcal{N}(t) = \sum \text{ord}_{Z_i, \mathcal{N}}(t) [Z_i]
$$
En explicitant davantage les définitions, choisissons pour tout $i$ des générateurs
$s_i \in \mathcal{L}_{\xi_i}$ et $t_i \in \mathcal{N}_{\xi_i}$,
où $\xi_i$ est le point générique de $Z_i$. Nous pouvons alors écrire
$$
s = f_i s_i
\quad\text{et}\quad
t = g_i t_i
$$
où $f_i, g_i$ sont des éléments inversibles de l'anneau total des fractions
$Q(\mathcal{O}_{X, \xi_i}^h)$. Posons, pour abréger, $B_i = \mathcal{O}_{X,
\xi_i}^h$. Notons
$$
\text{ord}_{B_i} : Q(B_i)^* \longrightarrow \mathbf{Z},\quad
a/b \longmapsto
\text{length}_{B_i}(B_i/aB_i) - \text{length}_{B_i}(B_i/bB_i)
$$
Autrement dit, nous étendons provisoirement la définition
\ref{algebra-definition-ord} d'Algèbre à ces anneaux locaux noethériens réduits
de dimension $1$. Alors, par définition,
$$
\text{ord}_{Z_i, \mathcal{L}}(s) = \text{ord}_{B_i}(f_i)
\quad\text{et}\quad
\text{ord}_{Z_i, \mathcal{N}}(t) = \text{ord}_{B_i}(g_i)
$$
Puisque $\xi_i$ est le point générique de $Z_i$, nous voyons que le corps
résiduel $\kappa(\xi_i)$ est le corps des fonctions de $Z_i$. De plus,
$\kappa(\xi_i)$ est le corps résiduel de $B_i$ ; voir le lemme d'Espaces décents
\ref{decent-spaces-lemma-residue-field-henselian-local-ring}. Puisque $t_i$
est un générateur de $\mathcal{N}_{\xi_i}$, on voit que son image dans la fibre
$\mathcal{N}_{\xi_i} \otimes_{B_i} \kappa(\xi_i)$ est une section méromorphe
non nulle de $\mathcal{N}|_{Z_i}$. Nous noterons cette image $t_i|_{Z_i}$. Il
résulte de nos définitions que
$$
c_1(\mathcal{N}) \cap \text{div}_\mathcal{L}(s) =
\sum \text{ord}_{B_i}(f_i)
(Z_i \to X)_*\text{div}_{\mathcal{N}|_{Z_i}}(t_i|_{Z_i})
$$
et de même
$$
c_1(\mathcal{L}) \cap \text{div}_\mathcal{N}(t) =
\sum \text{ord}_{B_i}(g_i)
(Z_i \to X)_*\text{div}_{\mathcal{L}|_{Z_i}}(s_i|_{Z_i})
$$
dans $\CH_{n - 2}(X)$. Nous allons construire une équivalence rationnelle entre
ces deux cycles. Pour cela, considérons le symbole modéré
$$
\partial_{B_i}(f_i, g_i) \in \kappa(\xi_i)^* = R(Z_i)^*
$$
voir la section \ref{chow-section-tame-symbol} de l'Homologie de Chow.

\begin{lemma}[Formule clé]
\label{spaces-chow-lemma-key-formula}
Dans la situation ci-dessus, le cycle
$$
\sum
(Z_i \to X)_*\left(
\text{ord}_{B_i}(f_i) \text{div}_{\mathcal{N}|_{Z_i}}(t_i|_{Z_i}) -
\text{ord}_{B_i}(g_i) \text{div}_{\mathcal{L}|_{Z_i}}(s_i|_{Z_i}) \right)
$$
est égal au cycle
$$
\sum (Z_i \to X)_*\text{div}(\partial_{B_i}(f_i, g_i))
$$
\end{lemma}

\begin{proof}
Voici la stratégie de la preuve : nous nous ramenons d'abord au cas où $\mathcal{L}$ et
$\mathcal{N}$ sont des modules inversibles triviaux, puis nous modifions nos choix
de banalisations locales et, enfin, nous utilisons la localisation \'etale pour
nous ramener au cas des schémas\footnote{Il est possible qu'une preuve plus courte
s'obtienne en appliquant immédiatement la localisation \'etale.}.

\medskip\noindent
Première étape. Soit $q : T \to X$ le morphisme construit dans le lemme
\ref{spaces-chow-lemma-Gm-torsor}. Nous utiliserons toutes les propriétés énoncées dans ce
lemme sans plus ample mention. En particulier, il suffit de montrer que les cycles
deviennent égaux après passage aux images réciproques par $q$. Désignons par $s'$ et $t'$ les images réciproques de $s$ et
$t$, vues comme sections méromorphes de $q^*\mathcal{L}$ et $q^*\mathcal{N}$.
Notons $Z'_i = q^{-1}(Z_i)$, notons $\xi'_i \in |Z'_i|$ le point générique,
notons $B'_i = \mathcal{O}_{T, \xi'_i}^h$, notons $\mathcal{L}_{\xi'_i}$ et
$\mathcal{N}_{\xi'_i}$ les images réciproques de $\mathcal{L}$ et $\mathcal{N}$ sur
$\Spec(B'_i)$. Rappelons que nous avons des diagrammes commutatifs
$$
\xymatrix{
\Spec(B'_i) \ar[r]_-{c_{\xi'_i}} \ar[d] & T \ar[d]^q \\
\Spec(B_i) \ar[r]^-{c_{\xi_i}} & X
}
$$
voir la remarque d'Espaces décents
\ref{decent-spaces-remark-functoriality-henselian-local-ring}. Notons $s'_i$
et $t'_i$ les images réciproques de $s_i$ et $t_i$ ; ce sont des générateurs de
$\mathcal{L}_{\xi'_i}$ et $\mathcal{N}_{\xi'_i}$. Alors
$$
s' = f'_i s'_i
\quad\text{et}\quad
t' = g'_i t'_i
$$
où $f'_i$ et $g'_i$ sont les images de $f_i, g_i$ par l'application $Q(B_i) \to
Q(B'_i)$ induite par $B_i \to B'_i$. D'après le lemme d'Algèbre
\ref{algebra-lemma-pullback-module}, nous avons
$$
\text{ord}_{B_i}(f_i) = \text{ord}_{B'_i}(f'_i)
\quad\text{et}\quad
\text{ord}_{B_i}(g_i) = \text{ord}_{B'_i}(g'_i)
$$
Par le lemme \ref{spaces-chow-lemma-prepare-flat-pullback-cap-c1} appliqué au morphisme $q : Z'_i \to
Z_i$, nous avons
$$
q^*\text{div}_{\mathcal{N}|_{Z_i}}(t_i|_{Z_i}) =
\text{div}_{q^*\mathcal{N}|_{Z'_i}}(t'_i|_{Z'_i})
\quad\text{et}\quad
q^*\text{div}_{\mathcal{L}|_{Z_i}}(s_i|_{Z_i}) =
\text{div}_{q^*\mathcal{L}|_{Z'_i}}(s'_i|_{Z'_i})
$$
Cela montre déjà que l'image réciproque du premier cycle de l'énoncé du lemme
est le cycle correspondant aux données $s', t', Z'_i, s'_i, t'_i$. Pour établir
la même assertion pour le second, remarquons que le lemme de l'Homologie de Chow
\ref{chow-lemma-tame-symbol-formally-smooth} donne
$$
\partial_{B_i}(f_i, g_i) \mapsto
\partial_{B'_i}(f'_i, g'_i)
\quad\text{via}\quad
\kappa(\xi_i) \to \kappa(\xi'_i)
$$
Le même lemme qu'auparavant montre donc que
$$
q^*\text{div}(\partial_{B_i}(f_i, g_i)) =
\text{div}(\partial_{B'_i}(f'_i, g'_i))
$$
Puisque $q^*\mathcal{L} \cong \mathcal{O}_T$, nous constatons qu'il suffit de
démontrer l'égalité dans le cas où $\mathcal{L}$ est trivial. En échangeant les
rôles de $\mathcal{L}$ et $\mathcal{N}$, nous voyons que nous pouvons de même
supposer que $\mathcal{N}$ est trivial. Ceci achève la preuve de la
première étape.

\medskip\noindent
Deuxième étape. Supposons $\mathcal{L} = \mathcal{O}_X$ et $\mathcal{N} =
\mathcal{O}_X$. Désignons par $1$ la section qui banalise $\mathcal{L}$. Alors
$s_i = u \cdot 1$ pour une certaine unité $u \in B_i$. Examinons ce qui se passe
si nous remplaçons $s_i$ par $1$. Alors, $f_i$ est remplacé par $u f_i$.
Ainsi la première partie de la première expression du lemme est inchangée et
au second terme nous ajoutons
$$
\text{ord}_{B_i}(g_i)\text{div}(u|_{Z_i})
$$
où $u|_{Z_i}$ est l'image de $u$ dans le corps résiduel, d'après le lemme d'Espaces sur les corps \ref{spaces-over-fields-lemma-divisor-meromorphic-well-defined},
et à la deuxième expression nous ajoutons
$$
\text{div}(\partial_{B_i}(u, g_i))
$$
par bilinéarité du symbole modéré. Ces termes sont égaux en vertu de la
propriété du symbole modéré donnée par l'équation de l'Homologie de Chow
(\ref{chow-item-normalization}).

\medskip\noindent
Soit $Y \subset X$ un sous-espace fermé intègre avec $\dim_\delta(Y) = n -
2$. Pour montrer que les coefficients de $Y$ dans les deux cycles du lemme sont
égaux, nous pouvons remplacer $s_i$ par $1$ comme au paragraphe
précédent. De même, nous pouvons remplacer $t_i$ par $1$. Puisqu'il n'y a qu'un nombre fini de $Z_i$ tels
que $Y \subset Z_i$, nous pouvons supposer $s_i = 1$ et $t_i = 1$ pour tous
ces $Z_i$.

\medskip\noindent
Troisième étape. Nous montrons ici que les coefficients de $Y$ dans les cycles
du lemme coïncident pour un sous-espace fermé intègre $Y$ avec
$\dim_\delta(Y) = n - 2$, dans la situation où, de plus, $\mathcal{L} = \mathcal{O}_X$ et
$\mathcal{N} = \mathcal{O}_X$, et où $s_i = 1$ et $t_i = 1$ pour tous les $Z_i$ tels
que $Y \subset Z_i$. Quitte à remplacer $X$ par un sous-espace ouvert plus
petit, nous pouvons supposer que $s_i$ et $t_i$ sont égaux à $1$ pour
tout $i$. Dans ce cas, le premier cycle est nul. Il reste à montrer
que le coefficient de $Y$ dans le deuxième cycle est également nul.

\medskip\noindent
Tout d'abord, puisque $\mathcal{L} = \mathcal{O}_X$ et $\mathcal{N} =
\mathcal{O}_X$, nous identifions $s, t$ à des
fonctions rationnelles $f, g$ sur $X$. Puisque $s_i$ et $t_i$ sont égaux à
$1$, nous voyons que $f_i$, resp.\ $g_i$, est l'image de $f$, resp.\ $g$, dans
$Q(B_i)$ pour tous les $i$. Soit $\zeta \in |Y|$ le point générique.
Choisissons un voisinage \'etale
$$
(U, u) \longrightarrow (X, \zeta)
$$
et notons $Y' = \overline{\{u\}} \subset U$. Puisqu'un morphisme \'etale
est plat, nous pouvons prendre les images réciproques de $f$ et $g$, qui sont des fonctions méromorphes
régulières sur $U$ que nous noterons encore $f$ et $g$. Pour tout
diviseur premier $Y \subset Z \subset X$, le schéma $Z \times_X U$ est une
réunion de diviseurs premiers de $U$. Réciproquement, étant donné un diviseur
premier $Y' \subset Z' \subset U$, il existe un diviseur premier $Y \subset Z
\subset X$ tel que $Z'$ soit une composante de $Z \times_X U$. Pour une
telle paire $(Z, Z')$, l'homomorphisme d'anneaux
$$
\mathcal{O}_{X, \xi}^h \to \mathcal{O}_{U, \xi'}^h
$$
est \'etale (il est même fini \'etale). On obtient donc
$$
\partial_{\mathcal{O}_{X, \xi}^h}(f, g) \mapsto
\partial_{\mathcal{O}_{U, \xi'}^h}(f, g)
\quad\text{via}\quad
\kappa(\xi) \to \kappa(\xi')
$$
en vertu du lemme de l'Homologie de Chow \ref{chow-lemma-tame-symbol-formally-smooth}. Ainsi
le lemme \ref{spaces-chow-lemma-etale-pullback-principal-divisor} s'applique et montre
$$
(Z \times_X U \to Z)^*\text{div}_Z(\partial_{\mathcal{O}_{X, \xi}^h}(f, g))
=
\sum\nolimits_{Z' \subset Z \times_X U}
\text{div}_{Z'}(\partial_{\mathcal{O}_{U, \xi'}^h}(f, g))
$$
Puisque l'image réciproque plate commute à l'image directe par les immersions
fermées (lemme \ref{spaces-chow-lemma-flat-pullback-proper-pushforward}), nous voyons
qu'il suffit de démontrer que le coefficient de $Y'$ dans
$$
\sum\nolimits_{Z' \subset U} 
(Z' \to U)_*\text{div}_{Z'}(\partial_{\mathcal{O}_{U, \xi'}^h}(f, g))
$$
est nul.

\medskip\noindent
Soit $A = \mathcal{O}_{U, u}$. Alors $f, g \in Q(A)^*$. Écrivons
$f = a/b$ et $g = c/d$, où $a, b, c, d \in A$ sont des non-diviseurs de zéro.
Le coefficient de $Y'$ dans l'expression ci-dessus est
$$
\sum\nolimits_{\mathfrak q \subset A\text{ height }1}
\text{ord}_{A/\mathfrak q}(\partial_{A_\mathfrak q}(f, g))
$$
Par bilinéarité de $\partial_A$, il suffit de démontrer que
$$
\sum\nolimits_{\mathfrak q \subset A\text{ height }1}
\text{ord}_{A/\mathfrak q}(\partial_{A_\mathfrak q}(a, c))
$$
est nul. Il en est de même pour les autres paires $(a, d)$, $(b, c)$ et $(b, d)$. Cela
résulte du lemme de l'Homologie de Chow \ref{chow-lemma-key-nonzerodivisors}.
\end{proof}











\section{Intersection avec un faisceau inversible et équivalence rationnelle}
\label{spaces-chow-section-commutativity}

\noindent
Cette section est l'analogue de l'Homologie de Chow, section
\ref{chow-section-commutativity}. En appliquant le lemme clé, nous obtenons
les propriétés fondamentales de l'intersection avec des faisceaux inversibles.
En particulier, nous verrons que $c_1(\mathcal{L}) \cap -$ passe au quotient par
l'équivalence rationnelle et que les opérations associées à différents faisceaux inversibles
commutent.

\begin{lemma}
\label{spaces-chow-lemma-commutativity-on-integral}
Dans la situation \ref{spaces-chow-situation-setup}, soit $X/B$ bon. Supposons
$X$ intègre et $\dim_\delta(X) = n$. Soient $\mathcal{L}$ et $\mathcal{N}$ des faisceaux
inversibles sur $X$. Choisissons une section méromorphe non nulle $s$ de
$\mathcal{L}$ et une section méromorphe non nulle $t$ de $\mathcal{N}$. Posons
$\alpha = \text{div}_\mathcal{L}(s)$ et $\beta = \text{div}_\mathcal{N}(t)$.
Alors
$$
c_1(\mathcal{N}) \cap \alpha
=
c_1(\mathcal{L}) \cap \beta
$$
dans $\CH_{n - 2}(X)$.
\end{lemma}

\begin{proof}
Cela résulte immédiatement du lemme clé \ref{spaces-chow-lemma-key-formula} et de la discussion qui le
précède.
\end{proof}

\begin{lemma}
\label{spaces-chow-lemma-factors}
Dans la situation \ref{spaces-chow-situation-setup}, soit $X/B$ bon. Soit
$\mathcal{L}$ un faisceau inversible sur $X$. L'opération $\alpha \mapsto c_1(\mathcal{L})
\cap \alpha$ passe au quotient par l'équivalence rationnelle pour donner une opération
$$
c_1(\mathcal{L}) \cap - : \CH_{k + 1}(X) \to \CH_k(X)
$$
\end{lemma}

\begin{proof}
Soit $\alpha \in Z_{k + 1}(X)$ et $\alpha \sim_{rat} 0$. Nous devons montrer
que $c_1(\mathcal{L}) \cap \alpha$, tel qu'il est défini dans la définition
\ref{spaces-chow-definition-cap-c1}, est nul. D'après la définition
\ref{spaces-chow-definition-rational-equivalence}, il existe une famille localement finie
$\{W_j\}$ de sous-espaces fermés intègres avec $\dim_\delta(W_j) = k + 2$ et
des fonctions rationnelles $f_j \in R(W_j)^*$ telles que
$$
\alpha = \sum (i_j)_*\text{div}_{W_j}(f_j)
$$
Notons que $p : \coprod W_j \to X$ est un morphisme propre, et donc $\alpha =
p_*\alpha'$ où $\alpha' \in Z_{k + 1}(\coprod W_j)$ est la somme des diviseurs
principaux $\text{div}_{W_j}(f_j)$. Par le lemme
\ref{spaces-chow-lemma-pushforward-cap-c1}, nous avons $c_1(\mathcal{L}) \cap \alpha =
p_*(c_1(p^*\mathcal{L}) \cap \alpha')$. Il suffit donc de montrer que chacun des cycles
$c_1(\mathcal{L}|_{W_j}) \cap \text{div}_{W_j}(f_j)$ est nul. En d'autres
termes, nous pouvons supposer que $X$ est intègre et $\alpha =
\text{div}_X(f)$ pour un certain $f \in R(X)^*$.

\medskip\noindent
Supposons $X$ intègre et $\alpha = \text{div}_X(f)$ pour
un certain $f \in R(X)^*$. Nous pouvons considérer $f$ comme une section
méromorphe régulière du faisceau inversible $\mathcal{N} = \mathcal{O}_X$.
Choisissons une section méromorphe $s$ de $\mathcal{L}$ et posons $\beta =
\text{div}_\mathcal{L}(s)$. Par le lemme \ref{spaces-chow-lemma-commutativity-on-integral}
nous concluons que
$$
c_1(\mathcal{L}) \cap \alpha = c_1(\mathcal{O}_X) \cap \beta.
$$
Cependant, le lemme \ref{spaces-chow-lemma-c1-cap-additive} montre que le membre
de droite est nul dans $\CH_k(X)$, comme voulu.
\end{proof}

\noindent
Dans la situation \ref{spaces-chow-situation-setup}, soit $X/B$ bon. Soit
$\mathcal{L}$ un faisceau inversible sur $X$. Nous noterons
$$
c_1(\mathcal{L})^s \cap - : \CH_{k + s}(X) \to \CH_k(X)
$$
l'opération $c_1(\mathcal{L}) \cap - $. Cette notation a un sens d'après le lemme
\ref{spaces-chow-lemma-factors}. Nous noterons $c_1(\mathcal{L}^s \cap -$ l'itérée
$s$-ième de cette opération pour tout $s \geq 0$.

\begin{lemma}
\label{spaces-chow-lemma-cap-commutative}
Dans la situation \ref{spaces-chow-situation-setup}, soit $X/B$ bon. Soient
$\mathcal{L}$ et $\mathcal{N}$ des faisceaux inversibles sur $X$. Pour tout $\alpha \in \CH_{k
+ 2}(X)$, nous avons
$$
c_1(\mathcal{L}) \cap c_1(\mathcal{N}) \cap \alpha
=
c_1(\mathcal{N}) \cap c_1(\mathcal{L}) \cap \alpha
$$
comme éléments de $\CH_k(X)$.
\end{lemma}

\begin{proof}
Écrivons $\alpha = \sum m_j[Z_j]$ pour une famille localement finie de
sous-espaces fermés intègres $Z_j \subset X$ avec $\dim_\delta(Z_j) = k + 2$.
Considérons le morphisme propre $p : \coprod Z_j \to X$. Posons
$\alpha' = \sum m_j[Z_j]$ comme un $(k + 2)$-cycle sur $\coprod Z_j$. Après
plusieurs applications du lemme \ref{spaces-chow-lemma-pushforward-cap-c1}, nous voyons que
$c_1(\mathcal{L}) \cap c_1(\mathcal{N}) \cap \alpha = p_*(c_1(p^*\mathcal{L})
\cap c_1(p^*\mathcal{N}) \cap \alpha')$ et $c_1(\mathcal{N}) \cap
c_1(\mathcal{L}) \cap \alpha = p_*(c_1(p^*\mathcal{N}) \cap
c_1(p^*\mathcal{L}) \cap \alpha')$. Il suffit donc de démontrer la formule dans
le cas où $X$ est intègre et $\alpha = [X]$. Dans ce cas, le résultat
découle du lemme \ref{spaces-chow-lemma-commutativity-on-integral} et des définitions.
\end{proof}

















\section{Intersection avec les diviseurs de Cartier effectifs}
\label{spaces-chow-section-intersecting-effective-Cartier}

\noindent
Cette section est l'analogue de l'Homologie de Chow, section
\ref{chow-section-intersecting-effective-Cartier}. Nous renvoyons à
l'introduction de cette section pour les motivations.

\medskip\noindent
Rappelons que les diviseurs de Cartier effectifs correspondent de façon $1$-à-$1$ aux
classes d'isomorphisme de couples $(\mathcal{L}, s)$ où $\mathcal{L}$ est un
faisceau inversible et $s$ une section globale ; voir le lemme de Diviseurs sur les
espaces \ref{spaces-divisors-lemma-characterize-OD}. Si $D$ correspond
à $(\mathcal{L}, s)$, alors $\mathcal{L} = \mathcal{O}_X(D)$. Gardons
ce fait à l'esprit au cours de cette section.

\begin{definition}
\label{spaces-chow-definition-gysin-homomorphism}
Dans la situation \ref{spaces-chow-situation-setup}, soit $X/B$ bon. Soit
$(\mathcal{L}, s)$ un couple constitué d'un faisceau inversible et d'une section
globale $s \in \Gamma(X, \mathcal{L})$. Soit $D = Z(s)$ le schéma des zéros
de $s$, et notons $i : D \to X$ l'immersion fermée. Nous définissons, pour
tout entier $k$, un {\it homomorphisme de Gysin raffiné}
$$
i^* : Z_{k + 1}(X) \to \CH_k(D).
$$
suivant les règles ci-dessous :
\begin{enumerate}
\item Étant donné un sous-espace fermé intègre $W \subset X$ avec
$\dim_\delta(W) = k + 1$, nous posons
\begin{enumerate}
\item si $W \not \subset D$, alors $i^*[W] = [D \cap W]_k$ comme
$k$-cycle sur $D$, et
\item si $W \subset D$, alors $i^*[W] = i'_*(c_1(\mathcal{L}|_W) \cap [W])$,
où $i' : W \to D$ est l'immersion fermée induite.
\end{enumerate}
\item Pour un $(k + 1)$-cycle général $\alpha = \sum n_j[W_j]$, nous
posons
$$
i^*\alpha = \sum n_j i^*[W_j]
$$
\item Si $D$ est un diviseur de Cartier effectif, nous notons $D \cdot
\alpha = i_*i^*\alpha$ l'image directe de la classe, considérée comme une classe sur $X$.
\end{enumerate}
\end{definition}

\noindent
En fait, comme nous le verrons plus tard, cet homomorphisme de Gysin $i^*$
peut être considéré comme un exemple d'image réciproque non plate. Ainsi, nous
appellerons parfois de manière informelle la classe $i^*\alpha$ l'{\it
image réciproque} de la classe $\alpha$.

\begin{remark}
\label{spaces-chow-remark-on-cycles}
Soient $S$, $B$, $X$, $\mathcal{L}$, $s$, $i : D \to X$ les données de la définition
\ref{spaces-chow-definition-gysin-homomorphism} et supposons que $\mathcal{L}|_D \cong
\mathcal{O}_D$. Dans ce cas, nous pouvons définir une application canonique $i^* :
Z_{k + 1}(X) \to Z_k(D)$ au niveau des cycles, en imposant $i^*[W] = 0$ lorsque
$W \subset D$. Cette possibilité sera utile plus loin.
\end{remark}

\begin{remark}
\label{spaces-chow-remark-pullback-pairs}
Soit $f : X' \to X$ un morphisme de bons espaces algébriques au-dessus de $B$ comme
dans la situation \ref{spaces-chow-situation-setup}. Soit $(\mathcal{L}, s, i : D \to X)$
un triplet comme dans la définition \ref{spaces-chow-definition-gysin-homomorphism}.
Posons $\mathcal{L}' = f^*\mathcal{L}$, $s' = f^*s$ et
$D' = X' \times_X D = Z(s')$. Nous obtenons un diagramme commutatif
$$
\xymatrix{
D' \ar[d]_g \ar[r]_{i'} & X' \ar[d]^f \\
D \ar[r]^i & X
}
$$
et nous pouvons demander diverses compatibilités entre $i^*$ et $(i')^*$.
\end{remark}

\begin{lemma}
\label{spaces-chow-lemma-support-cap-effective-Cartier}
Dans la situation \ref{spaces-chow-situation-setup}, soit $X/B$ bon. Soit
$(\mathcal{L}, s, i : D \to X)$ comme dans la définition
\ref{spaces-chow-definition-gysin-homomorphism}. Soit $\alpha$ un $(k + 1)$-cycle sur $X$.
Alors $i_*i^*\alpha = c_1(\mathcal{L}) \cap \alpha$ dans $\CH_k(X)$. En
particulier, si $D$ est un diviseur de Cartier effectif, alors $D \cdot \alpha =
c_1(\mathcal{O}_X(D)) \cap \alpha$.
\end{lemma}

\begin{proof}
Écrivons $\alpha = \sum n_j[W_j]$, où $i_j : W_j \to X$ sont les immersions de sous-espaces
fermés intègres avec $\dim_\delta(W_j) = k$. Puisque $D$ est le schéma des zéros
de $s$, nous voyons que $D \cap W_j$ est le schéma des zéros de
la restriction $s|_{W_j}$. Par conséquent, pour tout $j$ tel que $W_j \not
\subset D$, nous avons $c_1(\mathcal{L}) \cap [W_j] = [D \cap W_j]_k$ par le
lemme \ref{spaces-chow-lemma-geometric-cap}. Nous avons donc
$$
c_1(\mathcal{L}) \cap \alpha
=
\sum\nolimits_{W_j \not \subset D} n_j[D \cap W_j]_k
+
\sum\nolimits_{W_j \subset D}
n_j i_{j, *}(c_1(\mathcal{L})|_{W_j}) \cap [W_j])
$$
dans $\CH_k(X)$ d'après la définition \ref{spaces-chow-definition-cap-c1}. Le membre de droite
coïncide, terme à terme, avec l'image directe de la classe $i^*\alpha$ sur $D$
définie dans la définition \ref{spaces-chow-definition-gysin-homomorphism}. Le résultat en découle.
\end{proof}

\begin{lemma}
\label{spaces-chow-lemma-closed-in-X-gysin}
Dans la situation \ref{spaces-chow-situation-setup}, soit $f : X' \to X$ un morphisme
propre entre de bons espaces algébriques au-dessus de $B$. Soit $(\mathcal{L}, s, i : D \to
X)$ comme dans la définition \ref{spaces-chow-definition-gysin-homomorphism}. Formons le
diagramme
$$
\xymatrix{
D' \ar[d]_g \ar[r]_{i'} & X' \ar[d]^f \\
D \ar[r]^i & X
}
$$
comme dans la remarque \ref{spaces-chow-remark-pullback-pairs}. Pour tout $(k + 1)$-cycle
$\alpha'$ sur $X'$, nous avons $i^*f_*\alpha' = g_*(i')^*\alpha'$ dans
$\CH_k(D)$ (ce qui a un sens puisque $f_*$ est défini au niveau des cycles).
\end{lemma}

\begin{proof}
Supposons $\alpha = [W']$ pour un sous-espace fermé intègre $W' \subset X'$.
Soit $W \subset X$ l'« image » de $W'$ comme dans le lemme
\ref{spaces-chow-lemma-proper-image}. Si $W' \not \subset D'$, alors $W \not
\subset D$ et nous voyons que
$$
[W' \cap D']_k = \text{div}_{\mathcal{L}'|_{W'}}({s'|_{W'}})
\quad\text{et}\quad
[W \cap D]_k = \text{div}_{\mathcal{L}|_W}(s|_W)
$$
L'image par $f_*$ du premier cycle est donc égale au second, d'après le lemme
\ref{spaces-chow-lemma-equal-c1-as-cycles}.
L'égalité vaut ainsi au niveau des cycles.
Si $W' \subset D'$, alors $W \subset D$ et
$f_*(c_1(\mathcal{L}|_{W'}) \cap [W'])$ est égal à $c_1(\mathcal{L}|_W) \cap
[W]$ dans $\CH_k(W)$ par la deuxième assertion du lemme
\ref{spaces-chow-lemma-equal-c1-as-cycles}. Par la remarque
\ref{spaces-chow-remark-infinite-sums-rational-equivalences}, le résultat en découle pour
$\alpha'$ quelconque.
\end{proof}

\begin{lemma}
\label{spaces-chow-lemma-gysin-flat-pullback}
Dans la situation \ref{spaces-chow-situation-setup}, soit $f : X' \to X$ un morphisme plat
de dimension relative $r$ entre de bons espaces algébriques au-dessus de $B$. Soit
$(\mathcal{L}, s, i : D \to X)$ comme dans la définition
\ref{spaces-chow-definition-gysin-homomorphism}. Considérons le diagramme
$$
\xymatrix{
D' \ar[d]_g \ar[r]_{i'} & X' \ar[d]^f \\
D \ar[r]^i & X
}
$$
comme dans la remarque \ref{spaces-chow-remark-pullback-pairs}. Pour tout $(k + 1)$-cycle
$\alpha$ sur $X$, nous avons $(i')^*f^*\alpha = g^*i^*\alpha'$ dans $\CH_{k +
r}(D)$ (ce qui a un sens puisque $f^*$ est défini au niveau des cycles).
\end{lemma}

\begin{proof}
Supposons $\alpha = [W]$ pour un sous-espace fermé intègre $W \subset X$.
Soit $W' = f^{-1}(W) \subset X'$. Si $W \not \subset D$, alors $W'
\not \subset D'$ et nous voyons que
$$
W' \cap D' = g^{-1}(W \cap D)
$$
comme sous-espaces fermés de $D'$. L'égalité vaut donc au niveau des cycles ;
voir le lemme \ref{spaces-chow-lemma-pullback-coherent}. Si $W
\subset D$, alors $W' \subset D'$ et $W' = g^{-1}(W)$, avec $[W']_{k + 1 + r} =
g^*[W]$, et l'égalité vaut dans $\CH_{k + r}(D')$ par le lemme
\ref{spaces-chow-lemma-flat-pullback-cap-c1}. Par la remarque
\ref{spaces-chow-remark-infinite-sums-rational-equivalences}, le résultat en découle pour
$\alpha'$ quelconque.
\end{proof}

\begin{lemma}
\label{spaces-chow-lemma-easy-gysin}
Dans la situation \ref{spaces-chow-situation-setup}, soit $X/B$ bon. Soit
$(\mathcal{L}, s, i : D \to X)$ comme dans la définition
\ref{spaces-chow-definition-gysin-homomorphism}. Soit $Z \subset X$ un sous-schéma fermé
tel que $\dim_\delta(Z) \leq k + 1$ et tel que $D \cap Z$ soit un diviseur de
Cartier effectif sur $Z$. Alors $i^*([Z]_{k + 1}) = [D \cap Z]_k$.
\end{lemma}

\begin{proof}
L'hypothèse signifie que $s|_Z$ est une section régulière de $\mathcal{L}|_Z$.
Ainsi, $D \cap Z = Z(s)$ et nous obtenons
$$
[D \cap Z]_k = \sum n_i [Z(s_i)]_k
$$
comme égalité de cycles, où $s_i = s|_{Z_i}$, les $Z_i$ sont les composantes irréductibles
de $\delta$-dimension $k + 1$ et $[Z]_{k + 1} = \sum n_i[Z_i]$ ; voir le lemme
\ref{spaces-chow-lemma-prepare-geometric-cap}. Nous avons $D \cap Z_i = Z(s_i)$. La
comparaison avec la définition de l'application de Gysin achève la preuve.
\end{proof}













\section{Homomorphismes de Gysin}
\label{spaces-chow-section-gysin}

\noindent
Cette section est l'analogue de l'Homologie de Chow, section
\ref{chow-section-gysin}. Dans cette section, nous utilisons la formule clé
pour montrer que l'homomorphisme de Gysin passe au quotient par l'équivalence rationnelle.

\begin{lemma}
\label{spaces-chow-lemma-gysin-factors-general}
Dans la situation \ref{spaces-chow-situation-setup}, soit $X/B$ bon. Supposons
$X$ intègre et $n = \dim_\delta(X)$. Soit $i : D \to X$ un
diviseur de Cartier effectif. Soit $\mathcal{N}$ un $\mathcal{O}_X$-module
inversible et soit $t$ une section méromorphe non nulle de $\mathcal{N}$. Alors
$i^*\text{div}_\mathcal{N}(t) = c_1(\mathcal{N}) \cap [D]_{n - 1}$ dans
$\CH_{n - 2}(D)$.
\end{lemma}

\begin{proof}
Écrivons $\text{div}_\mathcal{N}(t) = \sum \text{ord}_{Z_i,
\mathcal{N}}(t)[Z_i]$ pour des sous-espaces fermés intègres $Z_i \subset
X$ de $\delta$-dimension $n - 1$. Nous pouvons supposer que la famille $\{Z_i\}$
est localement finie, que $t \in \Gamma(U, \mathcal{N}|_U)$ est un générateur
où $U = X \setminus \bigcup Z_i$, et que chaque composante irréductible de $D$
est l'un des $Z_i$ ; voir les lemmes d'Espaces sur les corps
\ref{spaces-over-fields-lemma-components-locally-finite},
\ref{spaces-over-fields-lemma-divisor-locally-finite} et
\ref{spaces-over-fields-lemma-divisor-meromorphic-locally-finite}.

\medskip\noindent
Posons $\mathcal{L} = \mathcal{O}_X(D)$. Notons $s \in \Gamma(X,
\mathcal{O}_X(D)) = \Gamma(X, \mathcal{L})$ la section canonique. Nous
appliquerons la discussion de la section \ref{spaces-chow-section-key} à notre situation
actuelle. Pour tout $i$, soit $\xi_i \in |Z_i|$ son point générique. Soit
$B_i = \mathcal{O}_{X, \xi_i}^h$. Pour tout $i$, choisissons des
générateurs $s_i$ de $\mathcal{L}_{\xi_i}$ et $t_i$ de $\mathcal{N}_{\xi_i}$
sur $B_i$, en imposant $s_i = s$ si
$Z_i \not \subset D$. Écrivons $s = f_i s_i$ et $t = g_i t_i$ avec $f_i, g_i
\in B_i$. Alors $\text{ord}_{Z_i, \mathcal{N}}(t) = \text{ord}_{B_i}(g_i)$.
D'autre part, nous avons $f_i \in B_i$ et
$$
[D]_{n - 1} = \sum \text{ord}_{B_i}(f_i)[Z_i]
$$
grâce à notre choix des $s_i$. Nous affirmons que
$$
i^*\text{div}_\mathcal{N}(t) =
\sum \text{ord}_{B_i}(g_i) \text{div}_{\mathcal{L}|_{Z_i}}(s_i|_{Z_i})
$$
comme égalité de cycles. Plus précisément, le membre de droite est un cycle
représentant le membre de gauche. En effet, cela résulte de notre formule
pour $\text{div}_\mathcal{N}(t)$ et le fait que
$\text{div}_{\mathcal{L}|_{Z_i}}(s_i|_{Z_i}) = [Z(s_i|_{Z_i})]_{n - 2} = [Z_i
\cap D]_{n - 2}$ lorsque $Z_i \not \subset D$ car dans ce cas $s_i|_{Z_i} =
s|_{Z_i}$ est une section régulière ; voir le lemme \ref{spaces-chow-lemma-compute-c1}. De
même,
$$
c_1(\mathcal{N}) \cap [D]_{n - 1} =
\sum \text{ord}_{B_i}(f_i) \text{div}_{\mathcal{N}|_{Z_i}}(t_i|_{Z_i})
$$
La formule clé (lemme \ref{spaces-chow-lemma-key-formula}) donne l'égalité
$$
\sum \left(
\text{ord}_{B_i}(f_i) \text{div}_{\mathcal{N}|_{Z_i}}(t_i|_{Z_i}) -
\text{ord}_{B_i}(g_i) \text{div}_{\mathcal{L}|_{Z_i}}(s_i|_{Z_i}) \right) =
\sum \text{div}_{Z_i}(\partial_{B_i}(f_i, g_i))
$$
de cycles. Si $Z_i \not \subset D$, alors $f_i = 1$ et donc
$\text{div}_{Z_i}(\partial_{B_i}(f_i, g_i)) = 0$. Nous obtenons ainsi une
équivalence rationnelle entre nos cycles spécifiques représentant
$i^*\text{div}_\mathcal{N}(t)$ et $c_1(\mathcal{N}) \cap [D]_{n - 1}$ sur $D$.
Ceci achève la preuve.
\end{proof}

\begin{lemma}
\label{spaces-chow-lemma-gysin-factors}
Dans la situation \ref{spaces-chow-situation-setup}, soit $X/B$ bon. Soit
$(\mathcal{L}, s, i : D \to X)$ comme dans la définition
\ref{spaces-chow-definition-gysin-homomorphism}. L'homomorphisme de Gysin passe au quotient par
l'équivalence rationnelle et donne une application $i^* : \CH_{k + 1}(X) \to
\CH_k(D)$.
\end{lemma}

\begin{proof}
Soit $\alpha \in Z_{k + 1}(X)$ et supposons que $\alpha \sim_{rat} 0$. Cela
signifie qu'il existe une famille localement finie de sous-espaces fermés
intègres $W_j \subset X$ de $\delta$-dimension $k + 2$ et des fonctions
$f_j \in R(W_j)^*$ telles que $\alpha = \sum i_{j, *}\text{div}_{W_j}(f_j)$. Posons $X' =
\coprod W_i$ et considérons le diagramme
$$
\xymatrix{
D' \ar[d]_q \ar[r]_{i'} & X' \ar[d]^p \\
D \ar[r]^i & X
}
$$
de la remarque \ref{spaces-chow-remark-pullback-pairs}. Puisque $X' \to X$ est propre,
nous voyons que $i^*p_* = q_*(i')^*$ par le lemme
\ref{spaces-chow-lemma-closed-in-X-gysin}. Comme nous savons que $q_*$ passe au quotient par
l'équivalence rationnelle (lemme
\ref{spaces-chow-lemma-proper-pushforward-rational-equivalence}), il suffit de démontrer le
résultat pour $\alpha' = \sum \text{div}_{W_j}(f_j)$ sur $X'$. Cela
nous ramène au cas où $X$ est intègre et $\alpha = \text{div}(f)$ pour
un certain $f \in R(X)^*$.

\medskip\noindent
Supposons $X$ intègre et $\alpha = \text{div}(f)$ pour
un certain $f \in R(X)^*$. Si $X = D$, alors nous voyons que $i^*\alpha$ est
égal à $c_1(\mathcal{L}) \cap \alpha$. Ceci est rationnellement équivalent à
zéro par le lemme \ref{spaces-chow-lemma-factors}. Si $D \not = X$, alors nous voyons que
$i^*\text{div}_X(f)$ est égal à $c_1(\mathcal{O}_D) \cap [D]_{n - 1}$ dans
$\CH_k(D)$ par le lemme \ref{spaces-chow-lemma-gysin-factors-general}. Bien entendu,
l'intersection avec $c_1(\mathcal{O}_D)$ est l'application nulle.
\end{proof}

\begin{lemma}
\label{spaces-chow-lemma-gysin-commutes-cap-c1}
Dans la situation \ref{spaces-chow-situation-setup}, soit $X/B$ bon. Soit
$(\mathcal{L}, s, i : D \to X)$ un triplet comme dans la définition
\ref{spaces-chow-definition-gysin-homomorphism}. Soit $\mathcal{N}$ un
$\mathcal{O}_X$-module inversible. Alors $i^*(c_1(\mathcal{N}) \cap \alpha) =
c_1(i^*\mathcal{N}) \cap i^*\alpha$ dans $\CH_{k - 2}(D)$ pour tout
$\alpha \in \CH_k(Z)$.
\end{lemma}

\begin{proof}
Avec exactement la même preuve que dans le lemme \ref{spaces-chow-lemma-gysin-factors},
cela découle des lemmes \ref{spaces-chow-lemma-pushforward-cap-c1},
\ref{spaces-chow-lemma-cap-commutative} et \ref{spaces-chow-lemma-gysin-factors-general}.
\end{proof}

\begin{lemma}
\label{spaces-chow-lemma-gysin-commutes-gysin}
Dans la situation \ref{spaces-chow-situation-setup}, soit $X/B$ bon. Soient
$(\mathcal{L}, s, i : D \to X)$ et $(\mathcal{L}', s', i' : D' \to X)$ deux
triplets comme dans la définition \ref{spaces-chow-definition-gysin-homomorphism}. Alors le
diagramme
$$
\xymatrix{
\CH_k(X) \ar[r]_{i^*} \ar[d]_{(i')^*} & \CH_{k - 1}(D) \ar[d] \\
\CH_{k - 1}(D') \ar[r] & \CH_{k - 2}(D \cap D')
}
$$
est commutatif, chaque flèche étant un homomorphisme de Gysin.
\end{lemma}

\begin{proof}
Désignons par $j : D \cap D' \to D$ et $j' : D \cap D' \to D'$ les immersions fermées
correspondant à $(\mathcal{L}|_{D'}, s|_{D'}$ et $(\mathcal{L}'_D, s|_D)$.
Nous devons montrer que $(j')^*i^*\alpha = j^* (i')^*\alpha$ pour tout
$\alpha \in \CH_k(X)$. Soit $W \subset X$ un sous-schéma fermé intègre de
dimension $k$. Nous établirons l'égalité dans le cas $\alpha = [W]$. Le cas
général découlera alors de l'observation de la remarque
\ref{spaces-chow-remark-infinite-sums-rational-equivalences} (et de la forme particulière de
notre équivalence rationnelle produite ci-dessous). Nous déduirons l'égalité
pour $\alpha = [W]$ de la formule clé.

\medskip\noindent
Prenons pour $\sigma$ une section méromorphe non nulle de
$\mathcal{L}|_W$, égale à $s|_W$ si $W \not
\subset D$. Prenons pour $\sigma'$ une section méromorphe non nulle de
$\mathcal{L}'|_W$, égale à $s'|_W$ si $W \not
\subset D'$. Écrivons
$$
\text{div}_{\mathcal{L}|_W}(\sigma) =
\sum \text{ord}_{Z_i, \mathcal{L}|_W}(\sigma)[Z_i] = \sum n_i[Z_i]
$$
et de même
$$
\text{div}_{\mathcal{L}'|_W}(\sigma') =
\sum \text{ord}_{Z_i, \mathcal{L}'|_W}(\sigma')[Z_i] = \sum n'_i[Z_i]
$$
comme dans la discussion de la section \ref{spaces-chow-section-key}. Nous voyons alors
que $Z_i \subset D$ si $n_i \not = 0$ et $Z'_i \subset D'$ si $n'_i \not = 0$.
Pour tout $i$, soit $\xi_i \in |Z_i|$ le point générique. Comme dans la
section \ref{spaces-chow-section-key}, choisissons pour tout $i$ un élément $\sigma_i \in
\mathcal{L}_{\xi_i}$, resp.\ $\sigma'_i \in \mathcal{L}'_{\xi_i}$, qui soit un générateur
sur $B_i = \mathcal{O}_{W, \xi_i}^h$ et qui soit égal à l'image de $s$, resp.\
$s'$, si $Z_i \not \subset D$, resp.\ $Z_i \not \subset D'$. Écrivons $\sigma =
f_i \sigma_i$ et $\sigma' = f'_i\sigma'_i$ de sorte que $n_i =
\text{ord}_{B_i}(f_i)$ et $n'_i = \text{ord}_{B_i}(f'_i)$. Il résulte de nos définitions que
$$
(j')^*i^*[W] =
\sum \text{ord}_{B_i}(f_i) \text{div}_{\mathcal{L}'|_{Z_i}}(\sigma'_i|_{Z_i})
$$
comme égalité de cycles, et
$$
j^*(i')^*[W] =
\sum \text{ord}_{B_i}(f'_i) \text{div}_{\mathcal{L}|_{Z_i}}(\sigma_i|_{Z_i})
$$
La formule clé (lemme \ref{spaces-chow-lemma-key-formula}) donne alors l'égalité
$$
\sum \left(
\text{ord}_{B_i}(f_i) \text{div}_{\mathcal{L}'|_{Z_i}}(\sigma'_i|_{Z_i}) -
\text{ord}_{B_i}(f'_i) \text{div}_{\mathcal{L}|_{Z_i}}(\sigma_i|_{Z_i})
\right) =
\sum \text{div}_{Z_i}(\partial_{B_i}(f_i, f'_i))
$$
de cycles. Notons que $\text{div}_{Z_i}(\partial_{B_i}(f_i, f'_i)) = 0$ si
$Z_i \not \subset D \cap D'$, car dans ce cas soit $f_i = 1$, soit $f'_i = 1$. Nous
obtenons ainsi une équivalence rationnelle entre nos cycles spécifiques
représentant $(j')^*i^*[W]$ et $j^*(i')^*[W]$ sur $D \cap D' \cap W$.
\end{proof}










\section{Diviseurs de Cartier effectifs relatifs}
\label{spaces-chow-section-relative-effective-cartier}

\noindent
Cette section est l'analogue de l'Homologie de Chow, section
\ref{chow-section-relative-effective-cartier}. Les diviseurs de Cartier effectifs
relatifs sont définis dans la section de Diviseurs sur les espaces
\ref{spaces-divisors-section-effective-Cartier-morphisms}. Pour développer les
résultats de base sur les classes de Chern de fibrés vectoriels, nous n'avons
besoin que du cas où le schéma ambiant et le diviseur de Cartier effectif sont
plats sur la base.

\begin{lemma}
\label{spaces-chow-lemma-relative-effective-cartier}
Dans la situation \ref{spaces-chow-situation-setup}, soient $X, Y/B$ bons. Soit $p : X
\to Y$ un morphisme plat de dimension relative $r$. Soit $i : D \to X$ un
diviseur de Cartier effectif relatif (Diviseurs sur les espaces, définition
\ref{spaces-divisors-definition-relative-effective-Cartier-divisor}). Soit
$\mathcal{L} = \mathcal{O}_X(D)$. Pour tout $\alpha \in \CH_{k + 1}(Y)$, nous
avons
$$
i^*p^*\alpha = (p|_D)^*\alpha
$$
dans $\CH_{k + r}(D)$ et
$$
c_1(\mathcal{L}) \cap p^*\alpha = i_* ((p|_D)^*\alpha)
$$
dans $\CH_{k + r}(X)$.
\end{lemma}

\begin{proof}
Soit $W \subset Y$ un sous-espace fermé intègre de $\delta$-dimension $k +
1$. Le lemme de Diviseurs sur les espaces
\ref{spaces-divisors-lemma-relative-Cartier} montre que $D \cap p^{-1}W$
est un diviseur de Cartier effectif sur $p^{-1}W$. Par le lemme
\ref{spaces-chow-lemma-easy-gysin}, nous obtenons la première égalité dans
$$
i^*[p^{-1}W]_{k + r + 1} =
[D \cap p^{-1}W]_{k + r} =
[(p|_D)^{-1}(W)]_{k + r}.
$$
et la seconde vient de ce que $D \cap p^{-1}(W) = (p|_D)^{-1}(W)$ comme espaces
algébriques. Puisque, par définition, $p^*[W] = [p^{-1}W]_{k + r + 1}$, nous
voyons que $i^*p^*[W] = (p|_D)^*[W]$ comme égalité de cycles. Si $\alpha = \sum
m_j[W_j]$ est un $k + 1$-cycle général, alors nous obtenons $i^*\alpha = \sum
m_j i^*p^*[W_j] = \sum m_j(p|_D)^*[W_j]$ comme égalité de cycles. Cela démontre
la première égalité. Pour déduire la seconde de la première, appliquons le lemme
\ref{spaces-chow-lemma-support-cap-effective-Cartier}.
\end{proof}












\section{Fibrés affines}
\label{spaces-chow-section-affine-vector}

\noindent
Cette section est l'analogue de l'Homologie de Chow, section
\ref{chow-section-affine-vector}. Pour un fibré affine, l'application d'image réciproque
est surjective sur les groupes de Chow.

\begin{lemma}
\label{spaces-chow-lemma-pullback-affine-fibres-surjective}
Dans la situation \ref{spaces-chow-situation-setup}, soient $X, Y/B$ bons. Soit $f :
X \to Y$ un morphisme plat quasi-compact au-dessus de $B$, de dimension relative $r$.
Supposons que, pour tout $y \in Y$, nous ayons $X_y \cong
\mathbf{A}^r_{\kappa(y)}$. Alors $f^* : \CH_k(Y) \to \CH_{k + r}(X)$ est
surjectif pour tout $k \in \mathbf{Z}$.
\end{lemma}

\begin{proof}
Soit $\alpha \in \CH_{k + r}(X)$. Écrivons $\alpha = \sum m_j[W_j]$ avec
$m_j \not = 0$ et $W_j$ des sous-espaces fermés intègres distincts deux à deux
de $\delta$-dimension $k + r$. Alors la famille $\{W_j\}$ est localement
finie sur $X$. Soit $Z_j \subset Y$ le sous-espace fermé intègre tel que
le morphisme obtenu $W_j \to Z_j$ soit dominant, comme dans le lemme
\ref{spaces-chow-lemma-proper-image}. Pour tout $V \subset Y$ ouvert quasi-compact, nous
voyons que $f^{-1}(V) \cap W_j$ n'est non vide que pour un nombre fini de $j$.
Par conséquent, les $Z_j$, fermetures des images, forment une famille
localement finie de sous-espaces fermés intègres de $Y$.

\medskip\noindent
Considérons les diagrammes cartésiens
$$
\xymatrix{
f^{-1}(Z_j) \ar[r] \ar[d]_{f_j} & X \ar[d]^f \\
Z_j \ar[r] & Y
}
$$
Supposons que $[W_j] \in Z_{k + r}(f^{-1}(Z_j))$ soit rationnellement
équivalent à $f_j^*\beta_j$ pour un $k$-cycle $\beta_j \in \CH_k(Z_j)$.
Alors $\beta = \sum m_j \beta_j$ sera un $k$-cycle sur $Y$ et $f^*\beta = \sum
m_j f_j^*\beta_j$ sera rationnellement équivalent à $\alpha$ (voir la remarque
\ref{spaces-chow-remark-infinite-sums-rational-equivalences}). Cela nous ramène au cas où $Y$ est
intègre et où $\alpha = [W]$ pour un sous-schéma fermé intègre de $X$ dominant
$Y$. En particulier, nous pouvons supposer que $d = \dim_\delta(Y) < \infty$.

\medskip\noindent
Nous pouvons donc raisonner par récurrence sur $d = \dim_\delta(Y)$. Si $d < k$,
alors $\CH_{k + r}(X) = 0$ et le lemme est valable ; c'est le cas initial de
la récurrence. Soit $V \subset Y$ un ouvert non vide. Supposons que nous
puissions montrer que $\alpha|_{f^{-1}(V)} = f^*\beta$ pour un certain $\beta
\in Z_k(V)$. Le lemme \ref{spaces-chow-lemma-exact-sequence-open} montre que
$\beta = \beta'|_V$ pour un certain $\beta' \in Z_k(Y)$. La suite exacte
$\CH_k(f^{-1}(Y \setminus V)) \to \CH_k(X) \to \CH_k(f^{-1}(V))$ du lemme
\ref{spaces-chow-lemma-restrict-to-open} montre que $\alpha - f^*\beta'$ provient d'un
cycle $\alpha' \in \CH_{k + r}(f^{-1}(Y \setminus V))$. Puisque $\dim_\delta(Y
\setminus V) < d$, nous concluons par récurrence sur $d$.

\medskip\noindent
En particulier, quitte à remplacer $Y$ par un ouvert convenable, nous pouvons
supposer que $Y$ est un schéma ayant pour point générique $\eta$. L'isomorphisme
$Y_\eta \cong \mathbf{A}^r_\eta$ s'étend à un isomorphisme sur un $V \subset
Y$ ouvert non vide ; voir le lemme de Limites d'espaces
\ref{spaces-limits-lemma-descend-finite-presentation}. Cela nous ramène au cas
des schémas, traité dans le lemme de l'Homologie de Chow
\ref{chow-lemma-pullback-affine-fibres-surjective}.
\end{proof}

\begin{lemma}
\label{spaces-chow-lemma-linebundle}
Dans la situation \ref{spaces-chow-situation-setup}, soit $X/B$ bon. Soit
$\mathcal{L}$ un $\mathcal{O}_X$-module inversible. Soit
$$
p :
L = \underline{\Spec}(\text{Sym}^*(\mathcal{L}))
\longrightarrow
X
$$
le fibré vectoriel associé au-dessus de $X$. Alors $p^* : \CH_k(X) \to \CH_{k +
1}(L)$ est un isomorphisme pour tout $k$.
\end{lemma}

\begin{proof}
Pour la surjectivité, voir le lemme
\ref{spaces-chow-lemma-pullback-affine-fibres-surjective}. Soit $o : X \to L$ la section
nulle de $L \to X$, c'est-à-dire le morphisme correspondant à la surjection
$\text{Sym}^*(\mathcal{L}) \to \mathcal{O}_X$ qui envoie $\mathcal{L}^{\otimes
n}$ sur zéro pour tout $n > 0$. Alors $p \circ o = \text{id}_X$ et $o(X)$
est un diviseur de Cartier effectif sur $L$. Ainsi, d'après le lemme
\ref{spaces-chow-lemma-relative-effective-cartier}, nous avons $o^* \circ p^* =
\text{id}$ et nous concluons que $p^*$ est également injectif.
\end{proof}























\section{Théorie bivariante de l'intersection}
\label{spaces-chow-section-bivariant}

\noindent
Cette section est l'analogue de l'Homologie de Chow, section
\ref{chow-section-bivariant}. Afin de traiter convenablement les classes de Chern
supérieures des fibrés vectoriels, nous introduisons la notion suivante, à la
suite de \cite{FM}. Il résulte de \cite[théorème 17.1]{F} que notre définition
coïncide avec celle de \cite{F}, sous réserve de la différence entre les
contextes considérés.

\begin{definition}
\label{spaces-chow-definition-bivariant-class}
\begin{reference}
Comparer avec \cite[Definition 17.1]{F}
\end{reference}
Dans la situation \ref{spaces-chow-situation-setup}, soit $f : X \to Y$ un morphisme de
bons espaces algébriques au-dessus de $B$. Soit $p \in \mathbf{Z}$. Une
{\it classe bivariante $c$ de degré $p$ pour $f$} est la donnée d'une règle qui
associe à tout morphisme $Y' \to Y$ de bons espaces algébriques au-dessus de $B$ et à
tout $k$ une application
$$
c \cap - : \CH_k(Y') \longrightarrow \CH_{k - p}(X')
$$
où $X' = Y' \times_Y X$, et qui satisfait aux conditions suivantes :
\begin{enumerate}
\item si $Y'' \to Y'$ est un morphisme propre, alors $c \cap (Y'' \to
Y')_*\alpha'' = (X'' \to X')_*(c \cap \alpha'')$ pour tout $\alpha''$ sur
$Y''$,
\item si $Y'' \to Y'$ est un morphisme de bons espaces algébriques au-dessus de $B$ qui est
plat de dimension relative $r$, alors $c \cap (Y'' \to Y')^*\alpha' = (X'' \to
X')^*(c \cap \alpha')$ pour tout $\alpha'$ sur $Y'$,
\item si $(\mathcal{L}', s', i' : D' \to Y')$ est comme dans la définition
\ref{spaces-chow-definition-gysin-homomorphism}, dont l'image réciproque est $(\mathcal{N}', t', j' : E'
\to X')$ sur $X'$, alors nous avons $c \cap (i')^*\alpha' = (j')^*(c \cap
\alpha')$ pour tout $\alpha'$ sur $Y'$.
\end{enumerate}
L'ensemble de toutes les classes bivariantes de degré $p$ pour $f$ est noté
$A^p(X \to Y)$.
\end{definition}

\noindent
Dans la situation \ref{spaces-chow-situation-setup}, soient $X \to Y$ et $Y \to Z$ des
morphismes de bons espaces algébriques au-dessus de $B$. Soit $p \in \mathbf{Z}$. Il
est clair que $A^p(X \to Y)$ est un groupe abélien. De plus, il est clair que
nous avons une composition bilinéaire
$$
A^p(X \to Y) \times A^q(Y \to Z) \to A^{p + q}(X \to Z)
$$
qui est associative. Nous nous intéresserons surtout à $A^p(X) = A^p(X \to X)$,
notation qui désignera toujours les classes de cohomologie bivariantes pour
$\text{id}_X$. C'est en effet là que vivront les classes de Chern.

\begin{definition}
\label{spaces-chow-definition-chow-cohomology}
Dans la situation \ref{spaces-chow-situation-setup}, soit $X/B$ bon. La
{\it cohomologie de Chow} de $X$ est la $\mathbf{Z}$-algèbre graduée $A^*(X)$ dont
la composante de degré $p$ est $A^p(X \to X)$.
\end{definition}

\noindent
Attention : il n'est pas clair que la structure de $\mathbf{Z}$-algèbre sur
$A^*(X)$ soit commutative, mais nous verrons que les classes de Chern vivent
en son centre.

\begin{remark}
\label{spaces-chow-remark-pullback-cohomology}
Dans la situation \ref{spaces-chow-situation-setup}, soit $f : X \to Y$ un morphisme de
bons espaces algébriques au-dessus de $B$. Il existe alors un homomorphisme canonique de
$\mathbf{Z}$-algèbres $A^*(Y) \to A^*(X)$. Plus précisément, étant donné $c \in
A^p(Y)$ et $X' \to X$, nous pouvons définir $f^*c$ par
l'application $c \cap - : \CH_k(X') \to \CH_{k - p}(X')$ qui est donnée en
considérant $X'$ comme un espace algébrique au-dessus de $Y$.
\end{remark}

\begin{lemma}
\label{spaces-chow-lemma-cap-c1-bivariant}
Dans la situation \ref{spaces-chow-situation-setup}, soit $X/B$ bon. Soit
$\mathcal{L}$ un $\mathcal{O}_X$-module inversible. Alors la règle qui
associe à $f : X' \to X$ $c_1(f^*\mathcal{L}) \cap - : \CH_k(X') \to \CH_{k -
1}(X')$ est une classe bivariante de degré $1$.
\end{lemma}

\begin{proof}
Cela découle des lemmes \ref{spaces-chow-lemma-factors}, \ref{spaces-chow-lemma-pushforward-cap-c1},
\ref{spaces-chow-lemma-flat-pullback-cap-c1} et \ref{spaces-chow-lemma-gysin-commutes-cap-c1}.
\end{proof}

\begin{lemma}
\label{spaces-chow-lemma-flat-pullback-bivariant}
Dans la situation \ref{spaces-chow-situation-setup}, soit $f : X \to Y$ un morphisme de
bons espaces algébriques au-dessus de $B$ qui est plat de dimension relative $r$.
Alors la règle qui associe à $Y' \to Y$ $(f')^* : \CH_k(Y') \to \CH_{k +
r}(X')$ où $X' = X \times_Y Y'$ est une classe bivariante de degré $-r$.
\end{lemma}

\begin{proof}
Cela découle des lemmes \ref{spaces-chow-lemma-flat-pullback-rational-equivalence},
\ref{spaces-chow-lemma-compose-flat-pullback},
\ref{spaces-chow-lemma-flat-pullback-proper-pushforward} et
\ref{spaces-chow-lemma-gysin-flat-pullback}.
\end{proof}

\begin{lemma}
\label{spaces-chow-lemma-gysin-bivariant}
Dans la situation \ref{spaces-chow-situation-setup}, soit $X/B$ bon. Soit
$(\mathcal{L}, s, i : D \to X)$ un triplet comme dans la définition
\ref{spaces-chow-definition-gysin-homomorphism}. Alors la règle qui associe à $f : X'
\to X$ $(i')^* : \CH_k(X') \to \CH_{k - 1}(D')$ où $D' = D \times_X X'$ est
une classe bivariante de degré $1$.
\end{lemma}

\begin{proof}
Cela découle des lemmes \ref{spaces-chow-lemma-gysin-factors},
\ref{spaces-chow-lemma-closed-in-X-gysin}, \ref{spaces-chow-lemma-gysin-flat-pullback} et
\ref{spaces-chow-lemma-gysin-commutes-gysin}.
\end{proof}

\begin{lemma}
\label{spaces-chow-lemma-push-proper-bivariant}
Dans la situation \ref{spaces-chow-situation-setup}, soient $f : X \to Y$ et $g : Y \to Z$
des morphismes de bons espaces algébriques au-dessus de $B$. Soit $c \in A^p(X \to Z)$
et supposons $f$ propre. Alors la règle qui associe à $X' \to X$
$\alpha \longmapsto f_*(c \cap \alpha)$ est une classe bivariante de degré
$p$.
\end{lemma}

\begin{proof}
Cela découle des lemmes \ref{spaces-chow-lemma-compose-pushforward},
\ref{spaces-chow-lemma-flat-pullback-proper-pushforward} et \ref{spaces-chow-lemma-closed-in-X-gysin}.
\end{proof}

\noindent
Ici, nous voyons que $c_1(\mathcal{L})$ est au centre de $A^*(X)$.

\begin{lemma}
\label{spaces-chow-lemma-c1-center}
Dans la situation \ref{spaces-chow-situation-setup}, soit $X/B$ bon. Soit
$\mathcal{L}$ un $\mathcal{O}_X$-module inversible. Alors $c_1(\mathcal{L})
\in A^1(X)$ commute avec tout élément $c \in A^p(X)$.
\end{lemma}

\begin{proof}
Soit $p : L \to X$ comme dans le lemme \ref{spaces-chow-lemma-linebundle} et soit $o : X
\to L$ la section nulle. Remarquons que $p^*\mathcal{L}^{\otimes -1}$ possède une
section canonique dont le schéma des zéros est exactement le diviseur de
Cartier effectif $o(X)$. Soit $\alpha \in \CH_k(X)$. Nous avons alors
$$
p^*(c_1(\mathcal{L}^{\otimes -1}) \cap \alpha) =
c_1(p^*\mathcal{L}^{\otimes -1}) \cap p^*\alpha =
o_* o^* p^*\alpha
$$
par les lemmes \ref{spaces-chow-lemma-flat-pullback-cap-c1} et
\ref{spaces-chow-lemma-relative-effective-cartier}. Puisque $c$ est une classe bivariante,
nous avons
\begin{align*}
p^*(c \cap c_1(\mathcal{L}^{\otimes -1}) \cap \alpha)
& =
c \cap p^*(c_1(\mathcal{L}^{\otimes -1}) \cap \alpha) \\
& =
c \cap o_* o^* p^*\alpha \\
& =
o_* o^* p^*(c \cap \alpha) \\
& =
p^*(c_1(\mathcal{L}^{\otimes -1}) \cap c \cap \alpha)
\end{align*}
(la dernière égalité résulte de ce qui précède, appliqué à $c \cap \alpha$). Puisque
$p^*$ est injectif d'après un lemme cité plus haut, nous obtenons que
$c_1(\mathcal{L}^{\otimes -1})$ est au centre de $A^*(X)$. Cela prouve le
lemme.
\end{proof}

\noindent
Voici un critère pour savoir quand une classe bivariante est nulle.

\begin{lemma}
\label{spaces-chow-lemma-bivariant-zero}
Dans la situation \ref{spaces-chow-situation-setup}, soit $X/B$ bon. Soit $c
\in A^p(X)$. Alors $c$ est nul si et seulement si $c \cap [Y] = 0$ dans
$\CH_*(Y)$ pour tout espace algébrique intègre $Y$ localement de type fini
sur $X$.
\end{lemma}

\begin{proof}
Le sens direct est clair. Réciproquement, supposons que $c \cap [Y] = 0$ dans
$\CH_*(Y)$ pour tout espace algébrique intègre $Y$ localement de type fini
sur $X$. Soit $X' \to X$ localement de type fini. Soit $\alpha \in
\CH_k(X')$. Écrivons $\alpha = \sum n_i [Y_i]$, les $Y_i \subset X'$ formant une
famille localement finie de sous-schémas fermés intègres de $\delta$-dimension $k$. Alors $\alpha$ est l'image directe du cycle
$\alpha' = \sum n_i[Y_i]$ sur $X'' = \coprod Y_i$ par le morphisme propre
$X'' \to X'$. Par les propriétés des classes bivariantes, il suffit de prouver
que $c \cap \alpha' = 0$ dans $\CH_{k - p}(X'')$. Nous avons $\CH_{k - p}(X'')
= \prod \CH_{k - p}(Y_i)$, comme il résulte immédiatement des définitions. Les applications
de projection $\CH_{k - p}(X'') \to \CH_{k - p}(Y_i)$ sont données par image réciproque
plate. Puisque l'intersection avec $c$ commute avec l'image réciproque plate,
nous voyons qu'il suffit de montrer que $c \cap [Y_i]$ est nul dans $\CH_{k -
p}(Y_i)$, ce qui est vrai par hypothèse.
\end{proof}















\section{Formule du fibré projectif}
\label{spaces-chow-section-projective-space-bundle-formula}

\noindent
Dans la situation \ref{spaces-chow-situation-setup}, soit $X/B$ bon. Considérons un
$\mathcal{O}_X$-module localement libre de type fini $\mathcal{E}$, de rang
$r$. Par convention, le {\it fibré projectif associé à $\mathcal{E}$} est le
morphisme
$$
\xymatrix{
\mathbf{P}(\mathcal{E}) =
\underline{\text{Proj}}_X(\text{Sym}^*(\mathcal{E}))
\ar[r]^-\pi
& X
}
$$
au-dessus de $X$, muni de $\mathcal{O}_{\mathbf{P}(\mathcal{E})}(1)$ normalisé
par l'égalité $\pi_*(\mathcal{O}_{\mathbf{P}(\mathcal{E})}(1)) = \mathcal{E}$.
En particulier, il existe une surjection $\pi^*\mathcal{E} \to
\mathcal{O}_{\mathbf{P}(\mathcal{E})}(1)$. Nous dirons informellement « soit
$(\pi : P \to X, \mathcal{O}_P(1))$ le fibré projectif associé à
$\mathcal{E}$ » pour désigner la situation où $P = \mathbf{P}(\mathcal{E})$ et
$\mathcal{O}_P(1) = \mathcal{O}_{\mathbf{P}(\mathcal{E})}(1)$.

\begin{lemma}
\label{spaces-chow-lemma-cap-projective-bundle}
Dans la situation \ref{spaces-chow-situation-setup}, soit $X/B$ bon. Soit $\mathcal{E}$ un
$\mathcal{O}_X$-module localement libre de type fini ; ce module
$\mathcal{E}$ est de rang $r$. Soit $(\pi : P \to X, \mathcal{O}_P(1))$ le
fibré projectif associé à $\mathcal{E}$. Pour tout $\alpha \in \CH_k(X)$,
l'élément
$$
\pi_*\left(
c_1(\mathcal{O}_P(1))^s \cap \pi^*\alpha
\right)
\in
\CH_{k + r - 1 - s}(X)
$$
est $0$ si $s < r - 1$ et égal à $\alpha$ si $s = r - 1$.
\end{lemma}

\begin{proof}
Soit $Z \subset X$ un sous-espace fermé intègre de $\delta$-dimension $k$.
Nous démontrerons le lemme pour $\alpha = [Z]$. Nous omettons l'argument qui
déduit le cas général de ce cas particulier ; indication : procéder comme dans
la remarque \ref{spaces-chow-remark-infinite-sums-rational-equivalences}.

\medskip\noindent
Soit $P_Z = P \times_X Z$ le changement de base ; bien entendu, $\pi_Z : P_Z
\to Z$ est le fibré projectif associé à $\mathcal{E}|_Z$, et l'image réciproque
de $\mathcal{O}_P(1)$ est le module inversible correspondant sur $P_Z$. Puisque
$c_1(\mathcal{O}_P(1) \cap -$ et $\pi^*$ sont des classes bivariantes par les
lemmes \ref{spaces-chow-lemma-cap-c1-bivariant} et \ref{spaces-chow-lemma-flat-pullback-bivariant},
il vient
$$
\pi_*\left(
c_1(\mathcal{O}_P(1))^s \cap \pi^*[Z]
\right)
=
(Z \to X)_*\pi_{Z, *}\left(
c_1(\mathcal{O}_{P_Z}(1))^s \cap \pi_Z^*[Z]
\right)
$$
Il suffit donc de prouver le lemme dans le cas où $X$ est intègre et $\alpha
= [X]$.

\medskip\noindent
Supposons $X$ intègre, $\dim_\delta(X) = k$ et $\alpha = [X]$. Remarquons que
$\pi^*[X] = [P]$, puisque $P$ est intègre de $\delta$-dimension $r - 1$. Si
$s < r - 1$, alors, par construction, $c_1(\mathcal{O}_P(1))^s \cap [P]$ est
un $(k + r - 1 - s)$-cycle. Son image directe est donc nulle pour des raisons
de dimension.

\medskip\noindent
Soit $s = r - 1$. L'argument précédent donne
$\pi_*(c_1(\mathcal{O}_P(1))^s \cap [P]) = n [X]$ pour un certain $n \in
\mathbf{Z}$. Nous voulons montrer que $n = 1$. Pour les mêmes raisons de
dimension, il suffit de le faire après remplacement de $X$ par un ouvert
dense. Nous pouvons donc supposer que $X$ est un schéma, et le résultat découle
alors du lemme \ref{chow-lemma-cap-projective-bundle} de l'Homologie de Chow.
\end{proof}

\begin{lemma}[Formule du fibré projectif]
\label{spaces-chow-lemma-chow-ring-projective-bundle}
Soit $(S, \delta)$ comme dans la situation \ref{spaces-chow-situation-setup}. Soit $X$
localement de type fini sur $S$. Soit $\mathcal{E}$ un $\mathcal{O}_X$-module
localement libre de type fini ; ce module $\mathcal{E}$ est de rang $r$. Soit
$(\pi : P \to X, \mathcal{O}_P(1))$ le fibré projectif associé à
$\mathcal{E}$. L'application
$$
\bigoplus\nolimits_{i = 0}^{r - 1}
\CH_{k + i}(X)
\longrightarrow
\CH_{k + r - 1}(P),
$$
$$
(\alpha_0, \ldots, \alpha_{r-1})
\longmapsto
\pi^*\alpha_0 +
c_1(\mathcal{O}_P(1)) \cap \pi^*\alpha_1
+ \ldots +
c_1(\mathcal{O}_P(1))^{r - 1} \cap \pi^*\alpha_{r-1}
$$
est un isomorphisme.
\end{lemma}

\begin{proof}
Fixons $k \in \mathbf{Z}$. Montrons d'abord que l'application est injective.
Supposons que $(\alpha_0, \ldots, \alpha_{r - 1})$ soit un élément du membre de
gauche dont l'image est nulle. Le lemme \ref{spaces-chow-lemma-cap-projective-bundle} donne
$$
0 = \pi_*(\pi^*\alpha_0 +
c_1(\mathcal{O}_P(1)) \cap \pi^*\alpha_1
+ \ldots +
c_1(\mathcal{O}_P(1))^{r - 1} \cap \pi^*\alpha_{r-1})
= \alpha_{r - 1}
$$
Ensuite,
$$
0 = \pi_*(c_1(\mathcal{O}_P(1)) \cap (\pi^*\alpha_0 +
c_1(\mathcal{O}_P(1)) \cap \pi^*\alpha_1
+ \ldots +
c_1(\mathcal{O}_P(1))^{r - 2} \cap \pi^*\alpha_{r - 2}))
= \alpha_{r - 2}
$$
et ainsi de suite. L'application est donc injective.

\medskip\noindent
Pour démontrer la surjectivité, nous raisonnerons exactement comme dans la
preuve du lemme \ref{spaces-chow-lemma-pullback-affine-fibres-surjective}, afin de nous
ramener au cas des schémas. Le lecteur est invité à omettre cette preuve.

\medskip\noindent
Soit $\beta \in \CH_{k + r - 1}(P)$. Écrivons $\beta = \sum m_j[W_j]$ avec
$m_j \not = 0$ et $W_j$ des sous-espaces fermés intègres distincts par paires
et de $\delta$-dimension $k + r$. La famille $\{W_j\}$ est alors localement
finie dans $P$. Soit $Z_j \subset X$ l'« image » de $W_j$ comme dans le lemme
\ref{spaces-chow-lemma-proper-image}. Pour toute partie ouverte quasi-compacte $U \subset
X$, l'intersection $\pi^{-1}(U) \cap W_j$ n'est non vide que pour un nombre
fini de $j$. Par conséquent, les images $Z_j$ forment une famille localement
finie de sous-espaces fermés intègres de $X$.

\medskip\noindent
Considérons les diagrammes cartésiens
$$
\xymatrix{
P_j \ar[r] \ar[d]_{\pi_j} & P \ar[d]^\pi \\
Z_j \ar[r] & X
}
$$
Supposons que $[W_j] \in Z_{k + r - 1}(P_j)$ soit rationnellement équivalent à
$$
\pi_j^*\alpha_{j, 0} +
c_1(\mathcal{O}(1)) \cap \pi_j^*\alpha_{j, 1} +
\ldots +
c_1(\mathcal{O}(1))^{r - 1} \cap \pi_j^*\alpha_{j, r - 1}
$$
pour certains $(k + i)$-cycles $\alpha_{j, i} \in \CH_{k + i}(Z_j)$. Alors
$\alpha_i = \sum m_j \beta_{j, i}$ sera un $(k + i)$-cycle sur $X$, et
$$
\pi^*\alpha_0 +
c_1(\mathcal{O}(1)) \cap \pi^*\alpha_1 +
\ldots +
c_1(\mathcal{O}(1))^{r - 1} \cap \pi^*\alpha_{r - 1}
$$
sera rationnellement équivalent à $\beta$ (voir la remarque
\ref{spaces-chow-remark-infinite-sums-rational-equivalences}). Cela nous réduit au cas $X$
intègre et $\alpha = [W]$ pour un sous-schéma fermé intègre de $P$ dominant
$X$. En particulier, nous pouvons supposer que $d = \dim_\delta(X) < \infty$.

\medskip\noindent
Nous pouvons donc raisonner par récurrence sur $d = \dim_\delta(X)$. Si $d <
k$, alors $\CH_{k + r - 1}(X) = 0$ et le lemme est établi ; c'est le cas
initial de la récurrence. Considérons une partie ouverte non vide $U \subset
X$. Supposons que
l'on puisse montrer que
$$
\beta|_{\pi^{-1}(U)} =
\pi^*\alpha_0 +
c_1(\mathcal{O}(1)) \cap \pi^*\alpha_1 +
\ldots +
c_1(\mathcal{O}(1))^{r - 1} \cap \pi^*\alpha_{r - 1}
$$
pour certains $\alpha_i \in Z_{k + i}(U)$. Le lemme
\ref{spaces-chow-lemma-exact-sequence-open} donne $\alpha_i = \alpha'_i|_U$ pour certains
$\alpha'_i \in Z_{k + i}(X)$. Par les suites exactes $\CH_{k +
i}(\pi^{-1}(X \setminus U)) \to \CH_{k + i}(P) \to \CH_{k + i}(\pi^{-1}(U))$
du lemme \ref{spaces-chow-lemma-restrict-to-open}, il vient que
$$
\beta -
\left(\pi^*\alpha'_0 +
c_1(\mathcal{O}(1)) \cap \pi^*\alpha'_1 +
\ldots +
c_1(\mathcal{O}(1))^{r - 1} \cap \pi^*\alpha'_{r - 1}\right)
$$
provient d'un cycle $\beta' \in \CH_{k + r}(\pi^{-1}(X \setminus U))$. Puisque
$\dim_\delta(X \setminus U) < d$, on conclut par récurrence sur $d$.

\medskip\noindent
En particulier, après remplacement de $X$ par un ouvert convenable, nous
pouvons supposer que $X$ est un schéma ; le problème est alors ramené au lemme
\ref{chow-lemma-chow-ring-projective-bundle} de l'Homologie de Chow.
\end{proof}

\begin{lemma}
\label{spaces-chow-lemma-vectorbundle}
Dans la situation \ref{spaces-chow-situation-setup}, soit $X/B$ bon. Soit $\mathcal{E}$ un
faisceau localement libre de type fini et de rang $r$ sur $X$. Soit
$$
p :
E = \underline{\Spec}(\text{Sym}^*(\mathcal{E}))
\longrightarrow
X
$$
le fibré vectoriel associé au-dessus de $X$. Alors $p^* : \CH_k(X) \to
\CH_{k + r}(E)$ est un isomorphisme pour tout $k$.
\end{lemma}

\begin{proof}
(Pour le cas des fibrés en droites, voir le lemme \ref{spaces-chow-lemma-linebundle}.) La
surjectivité résulte du lemme \ref{spaces-chow-lemma-pullback-affine-fibres-surjective}.
Soit $(\pi  : P \to X, \mathcal{O}_P(1))$ le fibré en espaces projectifs
associé au faisceau localement libre de type fini $\mathcal{E} \oplus
\mathcal{O}_X$. Soit $s \in \Gamma(P, \mathcal{O}_P(1))$ la section globale
correspondant à $(0, 1) \in \Gamma(X, \mathcal{E} \oplus \mathcal{O}_X)$.
Soit $D = Z(s) \subset P$. Remarquons que $(\pi|_D : D \to X ,
\mathcal{O}_P(1)|_D)$ est le fibré en espaces projectifs associé à
$\mathcal{E}$. Posons $\pi_D = \pi|_D$ et $\mathcal{O}_D(1) =
\mathcal{O}_P(1)|_D$. De plus, $D$ est un diviseur de Cartier effectif sur $P$.
Ainsi, $\mathcal{O}_P(D) = \mathcal{O}_P(1)$ (voir Diviseurs sur les espaces,
lemme \ref{spaces-divisors-lemma-characterize-OD}). Il existe aussi un
isomorphisme $E \cong P \setminus D$. Notons $j : E \to P$ l'immersion ouverte
correspondante. Pour l'injectivité, utilisons le fait que les éléments du noyau
de
$$
j^* :
\CH_{k + r}(P)
\longrightarrow
\CH_{k + r}(E)
$$
sont les cycles à support dans le diviseur de Cartier effectif $D$ ; voir le
lemme \ref{spaces-chow-lemma-restrict-to-open}. Ainsi, si $p^*\alpha = 0$, alors
$\pi^*\alpha = i_*\beta$ pour un certain $\beta \in \CH_{k + r}(D)$. Le lemme
\ref{spaces-chow-lemma-chow-ring-projective-bundle} permet d'écrire
$$
\beta = \pi_D^*\beta_0 +
\ldots + c_1(\mathcal{O}_D(1))^{r - 1} \cap \pi_D^* \beta_{r - 1}.
$$
pour certains $\beta_i \in \CH_{k + i}(X)$. Les lemmes
\ref{spaces-chow-lemma-relative-effective-cartier} et \ref{spaces-chow-lemma-pushforward-cap-c1}
donnent alors
$$
\pi^*\alpha = i_*\beta =
c_1(\mathcal{O}_P(1)) \cap \pi^*\beta_0 +
\ldots +
c_1(\mathcal{O}_D(1))^r \cap \pi^*\beta_{r - 1}.
$$
Puisque le rang de $\mathcal{E} \oplus \mathcal{O}_X$ est $r + 1$, cela
contredit le lemme \ref{spaces-chow-lemma-pushforward-cap-c1}, à moins que tous les
$\alpha$ et tous les $\beta_i$ ne soient nuls.
\end{proof}








\section{Les classes de Chern d'un fibré vectoriel}
\label{spaces-chow-section-chern-classes-vector-bundles}

\noindent
Cette section est l'analogue des sections
\ref{chow-section-chern-classes-vector-bundles} et
\ref{chow-section-intersecting-chern-classes} de l'Homologie de Chow.
Cependant, contrairement à ce qui y est fait, nous définissons directement les
classes de Chern d'un fibré vectoriel comme des classes bivariantes. Cela
économise un travail considérable.

\begin{lemma}
\label{spaces-chow-lemma-segre-classes}
Dans la situation \ref{spaces-chow-situation-setup}, soit $X/B$ bon. Soit $\mathcal{E}$ un
faisceau localement libre de type fini et de rang $r$ sur $X$. Soit
$(\pi : P \to X, \mathcal{O}_P(1))$ le fibré en espaces projectifs associé à
$\mathcal{E}$. Pour tout morphisme $X' \to X$ de bons espaces algébriques
au-dessus de $B$, il existe des applications uniques
$$
c_i(\mathcal{E}) \cap - : \CH_k(X') \longrightarrow \CH_{k - i}(X'),\quad
i = 0, \ldots, r
$$
telles que, pour $\alpha \in \CH_k(X')$, nous ayons $c_0(\mathcal{E}) \cap
\alpha = \alpha$ et
$$
\sum\nolimits_{i = 0, \ldots, r}
(-1)^i c_1(\mathcal{O}_{P'}(1))^i \cap
(\pi')^*\left(c_{r - i}(\mathcal{E}) \cap \alpha\right) = 0
$$
où $\pi' : P' \to X'$ est le changement de base de $\pi$. De plus, ces applications
définissent une classe bivariante $c_i(\mathcal{E})$ de degré $i$ sur $X$.
\end{lemma}

\begin{proof}
L'existence et l'unicité des applications $c_i(\mathcal{E}) \cap -$ résultent
immédiatement du lemme \ref{spaces-chow-lemma-chow-ring-projective-bundle} et de la
description donnée de $c_0(\mathcal{E})$. Pour tout $i \in \mathbf{Z}$, la
règle qui associe à tout morphisme $X' \to X$ de bons espaces algébriques
au-dessus de $B$ l'application
$$
t_i(\mathcal{E}) \cap - :
\CH_k(X') \longrightarrow \CH_{k - i}(X'),\quad
\alpha \longmapsto
\pi'_*(c_1(\mathcal{O}_{P'}(1))^{r - 1 + i} \cap (\pi')^*\alpha)
$$
est une classe bivariante\footnote{À des signes près, ce sont les classes de
Segre de $\mathcal{E}$.}, par les lemmes \ref{spaces-chow-lemma-cap-c1-bivariant},
\ref{spaces-chow-lemma-flat-pullback-bivariant} et \ref{spaces-chow-lemma-push-proper-bivariant}. Par le
lemme \ref{spaces-chow-lemma-cap-projective-bundle}, nous avons $t_i(\mathcal{E}) = 0$ pour
$i < 0$ et $t_0(\mathcal{E}) = 1$. En appliquant l'image directe à l'équation
de l'énoncé du lemme, le lemme \ref{spaces-chow-lemma-cap-projective-bundle} donne
$$
(-1)^r t_1(\mathcal{E}) + (-1)^{r - 1}c_1(\mathcal{E}) = 0
$$
En particulier, $c_1(\mathcal{E})$ est une classe bivariante. Si nous
multiplions l'équation de l'énoncé du lemme par
$c_1(\mathcal{O}_{P'}(1))$ et prenons l'image directe du résultat sur $X'$, il
vient
$$
(-1)^r t_2(\mathcal{E}) +
(-1)^{r - 1} t_1(\mathcal{E}) \cap c_1(\mathcal{E}) +
(-1)^{r - 2} c_2(\mathcal{E}) = 0
$$
Comme précédemment, nous concluons que $c_2(\mathcal{E})$ est une classe
bivariante. Et ainsi de suite.
\end{proof}

\begin{definition}
\label{spaces-chow-definition-chern-classes}
Dans la situation \ref{spaces-chow-situation-setup}, soit $X/B$ bon. Soit $\mathcal{E}$ un
faisceau localement libre de type fini et de rang $r$ sur $X$. Pour $i = 0,
\ldots, r$, la {\it $i$-ième classe de Chern de $\mathcal{E}$} est la classe
bivariante $c_i(\mathcal{E}) \in A^i(X)$ de degré $i$ construite dans le lemme
\ref{spaces-chow-lemma-segre-classes}. La {\it classe totale de Chern de $\mathcal{E}$} est
la somme formelle
$$
c(\mathcal{E}) =
c_0(\mathcal{E}) + c_1(\mathcal{E}) + \ldots + c_r(\mathcal{E})
$$
qui est considérée comme une classe bivariante non homogène sur $X$.
\end{definition}

\noindent
Par commodité, nous posons souvent $c_i(\mathcal{E}) = 0$ pour
$i > r$ et $i < 0$. Par définition, nous avons $c_0(\mathcal{E}) = 1 \in
A^0(X)$. Voici une vérification élémentaire.

\begin{lemma}
\label{spaces-chow-lemma-first-chern-class}
Dans la situation \ref{spaces-chow-situation-setup}, soit $X/B$ bon. Soit $\mathcal{L}$ un
$\mathcal{O}_X$-module inversible. La première classe de Chern de
$\mathcal{L}$ sur $X$, définie dans la définition
\ref{spaces-chow-definition-chern-classes}, est égale à la classe bivariante du lemme
\ref{spaces-chow-lemma-cap-c1-bivariant}.
\end{lemma}

\begin{proof}
En effet, dans ce cas $P = \mathbf{P}(\mathcal{L}) = X$ et $\mathcal{O}_P(1) =
\mathcal{L}$ par notre normalisation du fibré projectif ; voir la section
\ref{spaces-chow-section-projective-space-bundle-formula}. L'équation du lemme
\ref{spaces-chow-lemma-segre-classes} s'écrit donc
$$
(-1)^0 c_1(\mathcal{L})^0 \cap c^{new}_1(\mathcal{L}) \cap \alpha +
(-1)^1 c_1(\mathcal{L})^1 \cap c^{new}_0(\mathcal{L}) \cap \alpha = 0
$$
où $c_i^{new}(\mathcal{L})$ est défini comme dans la définition
\ref{spaces-chow-definition-chern-classes}. Puisque $c_0^{new}(\mathcal{L}) = 1$ et
$c_1(\mathcal{L})^0 = 1$, la conclusion suit.
\end{proof}

\noindent
Nous voyons ensuite que les classes de Chern appartiennent au centre de
l'anneau de cohomologie de Chow bivariante $A^*(X)$.

\begin{lemma}
\label{spaces-chow-lemma-cap-commutative-chern}
Dans la situation \ref{spaces-chow-situation-setup}, soit $X/B$ bon. Soit $\mathcal{E}$ un
$\mathcal{O}_X$-module localement libre de rang $r$. Alors
$c_j(\mathcal{L}) \in A^j(X)$ commute avec tout élément $c \in A^p(X)$. En
particulier, si $\mathcal{F}$ est un second $\mathcal{O}_X$-module localement
libre sur $X$, de rang $s$, alors
$$
c_i(\mathcal{E}) \cap c_j(\mathcal{F}) \cap \alpha
=
c_j(\mathcal{F}) \cap c_i(\mathcal{E}) \cap \alpha
$$
comme éléments de $\CH_{k - i - j}(X)$ pour tout $\alpha \in \CH_k(X)$.
\end{lemma}

\begin{proof}
Soit $X' \to X$ un morphisme de bons espaces algébriques au-dessus de $B$.
Soit $\alpha \in \CH_k(X')$. Écrivons $\alpha_j = c_j(\mathcal{E}) \cap
\alpha$, de sorte que $\alpha_0 = \alpha$. Le lemme
\ref{spaces-chow-lemma-segre-classes} donne
$$
\sum\nolimits_{i = 0}^r
(-1)^i c_1(\mathcal{O}_{P'}(1))^i \cap
(\pi')^*(\alpha_{r - i}) = 0
$$
dans le groupe de Chow du fibré projectif $(\pi' : P' \to X',
\mathcal{O}_{P'}(1))$ associé à $(X' \to X)^*\mathcal{E}$. En appliquant $c
\cap -$, puis en utilisant le lemme \ref{spaces-chow-lemma-c1-center} et les propriétés des
classes bivariantes, nous obtenons
$$
\sum\nolimits_{i = 0}^r
(-1)^i c_1(\mathcal{O}_{P'}(1))^i \cap
\pi^*(c \cap \alpha_{r - i}) = 0
$$
dans le groupe de Chow de $P'$. L'unicité affirmée dans le lemme
\ref{spaces-chow-lemma-segre-classes} montre donc que $c \cap \alpha_j$ est égal à
$c_j(\mathcal{E}) \cap (c \cap \alpha)$. Cela prouve le lemme.
\end{proof}

\begin{remark}
\label{spaces-chow-remark-extend-to-finite-locally-free}
Dans la situation \ref{spaces-chow-situation-setup}, soit $X/B$ bon. Soit $\mathcal{E}$ un
$\mathcal{O}_X$-module localement libre de type fini. Si le rang de
$\mathcal{E}$ n'est pas constant, nous pouvons encore définir les classes de
Chern de $\mathcal{E}$. En effet, dans ce cas, on peut écrire
$$
X = X_0 \amalg X_1 \amalg X_2 \amalg \ldots
$$
où $X_r \subset X$ est le sous-espace ouvert et fermé sur lequel le rang de
$\mathcal{E}$ est $r$. Si $X' \to X$ est un morphisme de bons espaces
algébriques au-dessus de $B$, l'image réciproque fournit une décomposition
correspondante de $X'$, et les définitions donnent
$$
\CH_*(X') = \prod\nolimits_{r \geq 0} \CH_*(X'_r)
$$
Nous définissons alors $c_i(\mathcal{E})$ comme la classe bivariante qui
préserve ces décompositions en produits directs et agit sur les facteurs par
les opérations déjà définies $c_i(\mathcal{E}|_{X_r}) \cap -$. Remarquons que,
dans ce cadre, $c_i(\mathcal{E})$ peut être non nul pour une infinité de $i$.
\end{remark}












\section{Relations polynomiales entre les classes de Chern}
\label{spaces-chow-section-relations-chern-classes}

\noindent
Dans la situation \ref{spaces-chow-situation-setup}, soit $X/B$ bon. Soient
$\mathcal{E}_i$ les membres d'une famille finie de $\mathcal{O}_X$-modules
localement libres de type fini. Le lemme \ref{spaces-chow-lemma-cap-commutative-chern}
montre que les classes de Chern
$$
c_j(\mathcal{E}_i) \in A^*(X)
$$
engendrent une sous-$\mathbf{Z}$-algèbre commutative, et même centrale, de la
cohomologie de Chow $A^*(X)$. Il est donc bien défini qu'un polynôme en ces
classes de Chern soit nul, ou que deux tels polynômes soient égaux. À titre
d'exemple, dire que $c_1(\mathcal{E}_1)^5 +
c_2(\mathcal{E}_2)c_3(\mathcal{E}_3) = 0$ signifie que les opérations
$$
\CH_k(Y) \longrightarrow \CH_{k - 5}(Y), \quad
\alpha \longmapsto
c_1(\mathcal{E}_1)^5 \cap \alpha +
c_2(\mathcal{E}_2) \cap c_3(\mathcal{E}_3) \cap \alpha
$$
sont nulles pour tout morphisme $f : Y \to X$ de bons espaces algébriques
au-dessus de $B$. Par le lemme \ref{spaces-chow-lemma-bivariant-zero}, cette condition
équivaut à demander que, pour tout morphisme $f : Y \to X$, où $Y$ est un
espace algébrique intègre localement de type fini sur $X$, le cycle
$$
c_1(\mathcal{E}_1)^5 \cap [Y] +
c_2(\mathcal{E}_2) \cap c_3(\mathcal{E}_3) \cap [Y]
$$
est nul dans $\CH_{\dim(Y) - 5}(Y)$.

\medskip\noindent
Un exemple concret est la relation
$$
c_1(\mathcal{L} \otimes_{\mathcal{O}_X} \mathcal{N})
=
c_1(\mathcal{L}) + c_1(\mathcal{N})
$$
établie dans le lemme \ref{spaces-chow-lemma-c1-cap-additive}. Plus généralement, voici ce
qui se passe lorsqu'on tensorise un faisceau localement libre quelconque par un
faisceau inversible.

\begin{lemma}
\label{spaces-chow-lemma-chern-classes-E-tensor-L}
Dans la situation \ref{spaces-chow-situation-setup}, soit $X/B$ bon. Soit $\mathcal{E}$ un
faisceau localement libre de type fini et de rang $r$ sur $X$. Soit
$\mathcal{L}$ un faisceau inversible sur $X$. Alors
\begin{equation}
\label{spaces-chow-equation-twist}
c_i({\mathcal E} \otimes {\mathcal L})
=
\sum\nolimits_{j = 0}^i
\binom{r - i + j}{j} c_{i - j}({\mathcal E}) c_1({\mathcal L})^j
\end{equation}
dans $A^*(X)$.
\end{lemma}

\begin{proof}
La preuve est identique à celle du lemme
\ref{chow-lemma-chern-classes-E-tensor-L} de l'Homologie de Chow, les lemmes
\ref{spaces-chow-lemma-bivariant-zero} et \ref{spaces-chow-lemma-segre-classes} remplaçant ceux qui y
sont utilisés.
\end{proof}












\section{Additivité des classes de Chern}
\label{spaces-chow-section-additivity-chern-classes}

\noindent
Cette section est l'analogue de la section
\ref{chow-section-additivity-chern-classes} de l'Homologie de Chow.

\begin{lemma}
\label{spaces-chow-lemma-get-rid-of-trivial-subbundle}
Dans la situation \ref{spaces-chow-situation-setup}, soit $X/B$ bon. Soient
$\mathcal{E}$, $\mathcal{F}$ des faisceaux localement libres de type fini sur
$X$, de rangs respectifs $r$, $r - 1$, qui s'insèrent dans une suite exacte
courte
$$
0 \to \mathcal{O}_X \to \mathcal{E} \to \mathcal{F} \to 0
$$
Alors
$$
c_r(\mathcal{E}) = 0, \quad
c_j(\mathcal{E}) = c_j(\mathcal{F}), \quad j = 0, \ldots, r - 1
$$
dans $A^*(X)$.
\end{lemma}

\begin{proof}
La preuve est identique à celle du lemme
\ref{chow-lemma-get-rid-of-trivial-subbundle} de l'Homologie de Chow, les
lemmes \ref{spaces-chow-lemma-bivariant-zero},
\ref{spaces-chow-lemma-relative-effective-cartier}, \ref{spaces-chow-lemma-pushforward-cap-c1} et
\ref{spaces-chow-lemma-segre-classes} remplaçant ceux qui y sont utilisés.
\end{proof}

\begin{lemma}
\label{spaces-chow-lemma-additivity-invertible-subsheaf}
Dans la situation \ref{spaces-chow-situation-setup}, soit $X/B$ bon. Soient
$\mathcal{E}$, $\mathcal{F}$ des faisceaux localement libres de type fini sur
$X$, de rangs respectifs $r$, $r - 1$, qui s'insèrent dans une suite exacte
courte
$$
0 \to \mathcal{L} \to \mathcal{E} \to \mathcal{F} \to 0
$$
où $\mathcal{L}$ est un faisceau inversible. Alors
$$
c(\mathcal{E}) = c(\mathcal{L}) c(\mathcal{F})
$$
dans $A^*(X)$.
\end{lemma}

\begin{proof}
La preuve est identique à celle du lemme
\ref{chow-lemma-additivity-invertible-subsheaf} de l'Homologie de Chow, les
lemmes \ref{spaces-chow-lemma-get-rid-of-trivial-subbundle} et
\ref{spaces-chow-lemma-chern-classes-E-tensor-L} remplaçant ceux qui y sont utilisés.
\end{proof}

\begin{lemma}
\label{spaces-chow-lemma-additivity-chern-classes}
Dans la situation \ref{spaces-chow-situation-setup}, soit $X/B$ bon. Supposons que
$\mathcal{E}$ s'insère dans une suite exacte
$$
0
\to
\mathcal{E}_1
\to
\mathcal{E}
\to
\mathcal{E}_2
\to
0
$$
de faisceaux localement libres de type fini $\mathcal{E}_i$, de rang $r_i$.
Les classes totales de Chern vérifient
$$
c(\mathcal{E}) = c(\mathcal{E}_1) c(\mathcal{E}_2)
$$
dans $A^*(X)$.
\end{lemma}

\begin{proof}
La preuve est identique à celle du lemme
\ref{chow-lemma-additivity-chern-classes} de l'Homologie de Chow, les lemmes
\ref{spaces-chow-lemma-bivariant-zero}, \ref{spaces-chow-lemma-additivity-invertible-subsheaf} et
\ref{spaces-chow-lemma-segre-classes} remplaçant ceux qui y sont utilisés.
\end{proof}

\begin{lemma}
\label{spaces-chow-lemma-chern-filter-by-linebundles}
Dans la situation \ref{spaces-chow-situation-setup}, soit $X/B$ bon. Soient
${\mathcal L}_i$, $i = 1, \ldots, r$, des $\mathcal{O}_X$-modules inversibles.
Soit $\mathcal{E}$ un $\mathcal{O}_X$-module localement libre, muni d'une
filtration
$$
0 = \mathcal{E}_0 \subset \mathcal{E}_1 \subset \mathcal{E}_2
\subset \ldots \subset \mathcal{E}_r = \mathcal{E}
$$
tel que $\mathcal{E}_i/\mathcal{E}_{i - 1} \cong \mathcal{L}_i$. Posons
$c_1({\mathcal L}_i) = x_i$. Alors
$$
c(\mathcal{E})
=
\prod\nolimits_{i = 1}^r (1 + x_i)
$$
dans $A^*(X)$.
\end{lemma}

\begin{proof}
Appliquons le lemme \ref{spaces-chow-lemma-additivity-invertible-subsheaf} et raisonnons
par récurrence.
\end{proof}



















\section{Le principe de scindage}
\label{spaces-chow-section-splitting-principle}

\noindent
Cette section est l'analogue de la section
\ref{chow-section-additivity-chern-classes} de l'Homologie de Chow.

\begin{lemma}
\label{spaces-chow-lemma-splitting-principle}
Dans la situation \ref{spaces-chow-situation-setup}, soit $X/B$ bon. Soient
$\mathcal{E}_i$ les membres d'une famille finie de $\mathcal{O}_X$-modules
localement libres, de rang $r_i$. Il existe un morphisme projectif et plat
$\pi : P \to X$, de dimension relative $d$, tel que
\begin{enumerate}
\item pour tout morphisme $f : Y \to X$ de bons espaces algébriques au-dessus de $B$,
l'application $\pi_Y^* : \CH_*(Y) \to \CH_{* + d}(Y \times_X P)$ est
injective, et
\item chaque $\pi^*\mathcal{E}_i$ possède une filtration dont les quotients
successifs $\mathcal{L}_{i, 1}, \ldots, \mathcal{L}_{i, r_i}$ sont des
${\mathcal O}_P$-modules inversibles.
\end{enumerate}
\end{lemma}

\begin{proof}
Raisonnons par récurrence sur l'entier $r = \sum r_i$. Si $r = 0$, nous
pouvons prendre $\pi = \text{id}_X$. Si $r_i = 1$ pour tout $i$, nous pouvons
encore prendre $\pi = \text{id}_X$. Supposons que $r_{i_0} > 1$ pour un certain
$i_0$. Soit $(\pi : P \to X, \mathcal{O}_P(1))$ le fibré projectif associé à
$\mathcal{E}_{i_0}$. L'homomorphisme canonique $\pi^*\mathcal{E}_{i_0} \to
\mathcal{O}_P(1)$ est surjectif ; son noyau $\mathcal{E}'_{i_0}$ est donc
localement libre de type fini et de rang $r_{i_0} - 1$. Remarquons que
$\pi_Y^*$ est injectif pour tout morphisme $f : Y \to X$ de bons espaces
algébriques au-dessus de $B$ ; voir le lemme
\ref{spaces-chow-lemma-chow-ring-projective-bundle}. Il suffit donc de démontrer le lemme
pour $P$ et les faisceaux localement libres $\pi^*\mathcal{E}_i$. Cependant,
puisque nous disposons du sous-fibré $\mathcal{E}_{i_0} \subset
\pi^*\mathcal{E}_{i_0}$ à quotient inversible, il suffit maintenant de traiter
la famille $\{\mathcal{E}_i\}_{i \not = i_0} \cup
\{\mathcal{E}'_{i_0}\}$. Cela diminue $r$ de $1$, et l'hypothèse de récurrence
permet de conclure.
\end{proof}

\noindent
Plutôt que d'expliquer ce qu'affirme le principe de scindage, utilisons-le
dans la preuve de quelques lemmes.

\begin{lemma}
\label{spaces-chow-lemma-chern-classes-dual}
Dans la situation \ref{spaces-chow-situation-setup}, soit $X/B$ bon. Soit $\mathcal{E}$ un
$\mathcal{O}_X$-module localement libre de type fini, dont le dual est
$\mathcal{E}^\vee$. Alors
$$
c_i(\mathcal{E}^\vee) = (-1)^i c_i(\mathcal{E})
$$
dans $A^i(X)$.
\end{lemma}

\begin{proof}
Choisissons un morphisme $\pi : P \to X$ comme dans le lemme
\ref{spaces-chow-lemma-splitting-principle}. Puisque $\pi^*$ est injectif après tout
changement de base, il suffit d'établir la relation entre les classes de Chern
de $\mathcal{E}$ et de $\mathcal{E}^\vee$ après image réciproque sur $P$. Nous
pouvons donc supposer qu'il existe des $\mathcal{O}_X$-modules inversibles
${\mathcal L}_i$, $i = 1, \ldots, r$, et une filtration
$$
0 = \mathcal{E}_0 \subset \mathcal{E}_1 \subset \mathcal{E}_2
\subset \ldots \subset \mathcal{E}_r = \mathcal{E}
$$
telle que $\mathcal{E}_i/\mathcal{E}_{i - 1} \cong \mathcal{L}_i$. On obtient
alors la filtration duale
$$
0 = \mathcal{E}_r^\perp \subset \mathcal{E}_1^\perp \subset \mathcal{E}_2^\perp
\subset \ldots \subset \mathcal{E}_0^\perp = \mathcal{E}^\vee
$$
telle que $\mathcal{E}_{i - 1}^\perp/\mathcal{E}_i^\perp \cong
\mathcal{L}_i^{\otimes -1}$. Posons $x_i = c_1(\mathcal{L}_i)$. Alors
$c_1(\mathcal{L}_i^{\otimes -1}) = - x_i$ par le lemme
\ref{spaces-chow-lemma-c1-cap-additive}. Par le lemme
\ref{spaces-chow-lemma-chern-filter-by-linebundles}, nous avons
$$
c(\mathcal{E}) = \prod\nolimits_{i = 1}^r (1 + x_i)
\quad\text{et}\quad
c(\mathcal{E}^\vee) = \prod\nolimits_{i = 1}^r (1 - x_i)
$$
dans $A^*(X)$. Le résultat découle d'un calcul formel que nous omettons.
\end{proof}

\begin{lemma}
\label{spaces-chow-lemma-chern-classes-tensor-product}
Dans la situation \ref{spaces-chow-situation-setup}, soit $X/B$ bon. Soient
$\mathcal{E}$ et $\mathcal{F}$ des $\mathcal{O}_X$-modules localement libres de
type fini, de rangs respectifs $r$ et $s$. Alors
$$
c_1(\mathcal{E} \otimes \mathcal{F})
=
r c_1(\mathcal{F}) + s c_1(\mathcal{E})
$$
$$
c_2(\mathcal{E} \otimes \mathcal{F})
=
r^2 c_2(\mathcal{F}) +
rs c_1(\mathcal{F})c_1(\mathcal{E}) +
s^2 c_2(\mathcal{E})
$$
et ainsi de suite (voir preuve).
\end{lemma}

\begin{proof}
En raisonnant exactement comme dans la preuve du lemme
\ref{spaces-chow-lemma-chern-classes-dual}, nous pouvons supposer donnés des
$\mathcal{O}_X$-modules inversibles ${\mathcal L}_i$, $i = 1, \ldots, r$, des
${\mathcal N}_i$, $i = 1, \ldots, s$, et des filtrations
$$
0 = \mathcal{E}_0 \subset \mathcal{E}_1 \subset \mathcal{E}_2
\subset \ldots \subset \mathcal{E}_r = \mathcal{E}
\quad\text{et}\quad
0 = \mathcal{F}_0 \subset \mathcal{F}_1 \subset \mathcal{F}_2
\subset \ldots \subset \mathcal{F}_s = \mathcal{F}
$$
telles que $\mathcal{E}_i/\mathcal{E}_{i - 1} \cong \mathcal{L}_i$ et que
$\mathcal{F}_j/\mathcal{F}_{j - 1} \cong \mathcal{N}_j$. En ordonnant
lexicographiquement les couples $(i, j)$, on obtient une filtration
$$
0 \subset \ldots \subset
\mathcal{E}_i \otimes \mathcal{F}_j
+
\mathcal{E}_{i - 1} \otimes \mathcal{F}
\subset \ldots \subset \mathcal{E} \otimes \mathcal{F}
$$
dont les quotients successifs sont
$$
\mathcal{L}_1 \otimes \mathcal{N}_1,
\mathcal{L}_1 \otimes \mathcal{N}_2,
\ldots,
\mathcal{L}_1 \otimes \mathcal{N}_s,
\mathcal{L}_2 \otimes \mathcal{N}_1,
\ldots,
\mathcal{L}_r \otimes \mathcal{N}_s
$$
Par le lemme \ref{spaces-chow-lemma-chern-filter-by-linebundles}, nous avons
$$
c(\mathcal{E}) = \prod (1 + x_i),
\quad
c(\mathcal{F}) = \prod (1 + y_j),
\quad\text{et}\quad
c(\mathcal{F}) = \prod (1 + x_i + y_j),
$$
dans $A^*(X)$. Le résultat découle d'un calcul formel que nous omettons.
\end{proof}











\section{Degrés des zéro-cycles}
\label{spaces-chow-section-degree-zero-cycles}

\noindent
Cette section est l'analogue de la section
\ref{chow-section-degree-zero-cycles} de l'Homologie de Chow. Commençons par
définir le degré d'un zéro-cycle sur un espace algébrique propre au-dessus d'un
corps.

\begin{definition}
\label{spaces-chow-definition-degree-zero-cycle}
Soit $k$ un corps. Soit $p : X \to \Spec(k)$ un morphisme propre d'espaces
algébriques. Le {\it degré d'un zéro-cycle} sur $X$ est donné par l'image
directe propre
$$
p_* : \CH_0(X) \longrightarrow \CH_0(\Spec(k)) \longrightarrow \mathbf{Z}
$$
(lemme \ref{spaces-chow-lemma-proper-pushforward-rational-equivalence}), composée avec
l'isomorphisme naturel $\CH_0(\Spec(k)) \to \mathbf{Z}$ qui envoie $[\Spec(k)]$
à $1$. Notation : $\deg(\alpha)$.
\end{definition}

\noindent
Expliquons cela plus en détail.

\begin{lemma}
\label{spaces-chow-lemma-spell-out-degree-zero-cycle}
Soit $k$ un corps. Soit $X$ un espace algébrique propre sur $k$. Soit $\alpha
= \sum n_i[Z_i]$ dans $Z_0(X)$. Alors
$$
\deg(\alpha) = \sum n_i\deg(Z_i)
$$
où $\deg(Z_i)$ est le degré de $Z_i \to \Spec(k)$, c'est-à-dire $\deg(Z_i) =
\dim_k \Gamma(Z_i, \mathcal{O}_{Z_i})$.
\end{lemma}

\begin{proof}
C'est la définition de l'image directe propre (définition
\ref{spaces-chow-definition-proper-pushforward}).
\end{proof}

\begin{lemma}
\label{spaces-chow-lemma-degrees-and-numerical-intersections}
Soit $k$ un corps. Soit $X$ un espace algébrique propre sur $k$. Soit $Z
\subset X$ un sous-espace fermé de dimension $d$. Soient $\mathcal{L}_1,
\ldots, \mathcal{L}_d$ des $\mathcal{O}_X$-modules inversibles. Alors
$$
(\mathcal{L}_1 \cdots \mathcal{L}_d \cdot Z) =
\deg(
c_1(\mathcal{L}_1) \cap \ldots \cap c_1(\mathcal{L}_1) \cap [Z]_d)
$$
où le membre de gauche est défini dans Espaces sur les corps, définition
\ref{spaces-over-fields-definition-intersection-number}.
\end{lemma}

\begin{proof}
Soient $Z_i \subset Z$, $i = 1, \ldots, t$, les composantes irréductibles de
dimension $d$. Soit $m_i$ la multiplicité de $Z_i$ dans $Z$. Alors $[Z]_d =
\sum m_i[Z_i]$, et $c_1(\mathcal{L}_1) \cap \ldots \cap c_1(\mathcal{L}_d)
\cap [Z]_d$ est la somme des cycles $m_i c_1(\mathcal{L}_1) \cap \ldots \cap
c_1(\mathcal{L}_d) \cap [Z_i]$. Comme on dispose d'une décomposition analogue
pour $(\mathcal{L}_1 \cdots \mathcal{L}_d \cdot Z)$ d'après le lemme
\ref{spaces-over-fields-lemma-numerical-polynomial-leading-term} d'Espaces sur
les corps, il suffit de démontrer le lemme lorsque $Z = X$ est un espace
algébrique intègre propre sur $k$.

\medskip\noindent
D'après le lemme de Chow, il existe un morphisme propre $f : X' \to X$ qui est
un isomorphisme au-dessus d'une partie ouverte dense $U \subset X$ et tel que
$X'$ soit un schéma ; voir le lemme
\ref{spaces-more-morphisms-lemma-chow-noetherian-separated} de Compléments sur
les morphismes d'espaces. Alors $X'$ est un schéma propre sur $k$. Après avoir
remplacé $X'$ par l'adhérence schématique de $f^{-1}(U)$, nous pouvons supposer
que $X'$ est intègre. Alors
$$
(f^*\mathcal{L}_1 \cdots f^*\mathcal{L}_d \cdot X') =
(\mathcal{L}_1 \cdots \mathcal{L}_d \cdot X)
$$
par le lemme \ref{spaces-over-fields-lemma-intersection-number-and-pullback}
d'Espaces sur les corps, et nous avons
$$
f_*(c_1(f^*\mathcal{L}_1) \cap \ldots \cap c_1(f^*\mathcal{L}_d) \cap [Y]) =
c_1(\mathcal{L}_1) \cap \ldots \cap c_1(\mathcal{L}_d) \cap [X]
$$
par le lemme \ref{spaces-chow-lemma-pushforward-cap-c1}. Nous pouvons donc remplacer $X$
par $X'$ et supposer que $X$ est un schéma propre sur $k$. Ce cas est établi
dans le lemme \ref{chow-lemma-degrees-and-numerical-intersections} de
l'Homologie de Chow.
\end{proof}



















% OK, start here.
%

\chapter{Quotients de groupoïdes}



\phantomsection
\label{groupoids-quotients-section-phantom}


\section{Introduction}
\label{groupoids-quotients-section-introduction}

\noindent
Ce chapitre est consacré à des généralités sur les groupoïdes et leurs
quotients (dans la mesure où ceux-ci existent).
Il existe une abondante littérature sur ce sujet ; voir par exemple
\cite{GIT}, \cite{seshadri_quotients}, \cite{KollarQuotients},
\cite{K-M}, \cite{KollarFinite}, et bien d'autres références.






\section{Conventions et notation}
\label{groupoids-quotients-section-conventions-notation}

\noindent
Dans ce chapitre, les conventions et la notation sont celles introduites dans
Groupoïdes in Espaces, sections \ref{spaces-groupoids-section-conventions}
et \ref{spaces-groupoids-section-notation}.


\section{Morphismes invariants}
\label{groupoids-quotients-section-invariant}

\begin{definition}
\label{groupoids-quotients-definition-invariant}
Soit $S$ un schéma et soit $B$ un espace algébrique sur $S$.
Soit $j = (t, s) : R \to U \times_B U$ une prérelation d'espaces
algébriques sur $B$. On dit qu'un morphisme $\phi : U \to X$ d'espaces algébriques
sur $B$ est {\it $R$-invariant} si le diagramme
$$
\xymatrix{
R \ar[r]_s \ar[d]_t & U \ar[d]^\phi \\
U \ar[r]^\phi & X
}
$$
est commutatif. Si $j : R \to U \times_B U$ provient de l'action
d'un espace algébrique en groupes $G$ sur $U$ au-dessus de $B$, comme dans
Groupoïdes in Espaces, lemme \ref{spaces-groupoids-lemma-groupoid-from-action},
on dit alors que $\phi$ est {\it $G$-invariant}.
\end{definition}

\noindent
Autrement dit, un morphisme $U \to X$ est $R$-invariant s'il égalise
$s$ et $t$. On peut reformuler ceci à l'aide des faisceaux quotients
associés, comme suit.

\begin{lemma}
\label{groupoids-quotients-lemma-invariant}
Soit $S$ un schéma et soit $B$ un espace algébrique sur $S$.
Soit $j = (t, s) : R \to U \times_B U$ une prérelation d'espaces
algébriques sur $B$. Un morphisme d'espaces algébriques $\phi : U \to X$ est
$R$-invariant si et seulement s'il se factorise sous la forme
$U \to U/R \to X$.
\end{lemma}

\begin{proof}
Cela résulte immédiatement de la définition du faisceau quotient donnée dans
Groupoïdes in Espaces, section \ref{spaces-groupoids-section-quotient-sheaves}.
\end{proof}

\begin{lemma}
\label{groupoids-quotients-lemma-base-change-on-invariant}
Soit $S$ un schéma et soit $B$ un espace algébrique sur $S$.
Soit $j = (t, s) : R \to U \times_B U$ une prérelation d'espaces
algébriques sur $B$. Soit $U \to X$ un morphisme $R$-invariant d'espaces
algébriques sur $B$. Soit $X' \to X$ un morphisme quelconque d'espaces algébriques.
\begin{enumerate}
\item En posant $U' = X' \times_X U$, $R' = X' \times_X R$, on obtient
une prérelation $j' : R' \to U' \times_B U'$.
\item Si $j$ est une relation, alors $j'$ est une relation.
\item Si $j$ est une prérelation d'équivalence, alors $j'$ est une
prérelation d'équivalence.
\item Si $j$ est une relation d'équivalence, alors $j'$ est une relation
d'équivalence.
\item Si $j$ provient d'un groupoïde en espaces algébriques
$(U, R, s, t, c)$ sur $B$, alors
\begin{enumerate}
\item $(U, R, s, t, c)$ est un groupoïde en espaces algébriques sur $X$, et
\item $j'$ provient du changement de base $(U', R', s', t', c')$
de ce groupoïde à $X'$ ; voir
Groupoïdes in Espaces, lemme
\ref{spaces-groupoids-lemma-base-change-groupoid}.
\end{enumerate}
\item Si $j$ provient de l'action d'un espace algébrique en groupes $G/B$ sur $U$,
comme dans Groupoïdes in Espaces, lemme
\ref{spaces-groupoids-lemma-groupoid-from-action},
alors $j'$ provient de l'action induite de $G$ sur $U'$.
\end{enumerate}
\end{lemma}

\begin{proof}
Omise. Indication : combiner le point de vue fonctoriel avec le diagramme
$$
\xymatrix{
R' = X' \times_X R \ar[dd] \ar[rr] \ar[rd] & &
X' \times_X U = U' \ar'[d][dd] \ar[rd] \\
& R \ar[dd] \ar[rr] & & U \ar[dd] \\
U' = X' \times_X U \ar'[r][rr] \ar[rd] & & X' \ar[rd] \\
& U \ar[rr] & & X
}
$$
\end{proof}

\begin{definition}
\label{groupoids-quotients-definition-base-change}
Dans la situation du lemme \ref{groupoids-quotients-lemma-base-change-on-invariant},
on appelle $j' : R' \to U' \times_B U'$ le {\it changement de base} de la prérelation
$j$ à $X'$. On dit que c'est un {\it changement de base plat} si $X' \to X$ est un
morphisme plat d'espaces algébriques.
\end{definition}

\noindent
Ce type de changement de base se comporte bien avec la formation des faisceaux
quotients et des champs quotients.

\begin{lemma}
\label{groupoids-quotients-lemma-base-change-quotient-sheaf}
Dans la situation du lemme \ref{groupoids-quotients-lemma-base-change-on-invariant},
il existe un isomorphisme de faisceaux
$$
U'/R' = X' \times_X U/R
$$
Pour la construction des faisceaux quotients, voir
Groupoïdes in Espaces, section \ref{spaces-groupoids-section-quotient-sheaves}.
\end{lemma}

\begin{proof}
Puisque $U \to X$ est $R$-invariant, il est clair que l'application
$U \to X$ se factorise par le faisceau quotient $U/R$.
Rappelons que, par définition,
$$
\xymatrix{
R \ar@<1ex>[r] \ar@<-1ex>[r] &
U \ar[r] &
U/R
}
$$
est un diagramme coégalisateur dans la catégorie $\Sh$ des faisceaux d'ensembles sur
$(\Sch/S)_{fppf}$. En fait, c'est un diagramme coégalisateur dans la
catégorie au-dessus $\Sh/X$. Puisque le foncteur de changement de base
$X' \times_X - : \Sh/X \to \Sh/X'$ est exact (ce qui est vrai dans tout topos),
on conclut.
\end{proof}

\begin{lemma}
\label{groupoids-quotients-lemma-base-change-quotient-stack}
Soit $S$ un schéma. Soit $B$ un espace algébrique sur $S$.
Soit $(U, R, s, t, c)$ un groupoïde en espaces algébriques sur $B$.
Soit $U \to X$ un morphisme $R$-invariant d'espaces algébriques sur
$B$. Soit $g : X' \to X$ un morphisme d'espaces algébriques sur $B$
et soit $(U', R', s', t', c')$ le changement de base défini dans le
lemme \ref{groupoids-quotients-lemma-base-change-on-invariant}. Alors
$$
\xymatrix{
[U'/R'] \ar[r] \ar[d] & [U/R] \ar[d] \\
\mathcal{S}_{X'} \ar[r] & \mathcal{S}_X
}
$$
est un $2$-produit fibré de champs en groupoïdes sur $(\Sch/S)_{fppf}$.
Pour la construction des champs quotients et des morphismes de ce
diagramme, voir
Groupoïdes in Espaces, section \ref{spaces-groupoids-section-stacks}.
\end{lemma}

\begin{proof}
Nous allons le démontrer en utilisant la description explicite
des champs quotients donnée dans
Groupoïdes in Espaces, lemme
\ref{spaces-groupoids-lemma-quotient-stack-objects}.
Nous encourageons néanmoins vivement le lecteur à trouver sa propre démonstration.
Tout d'abord, on peut considérer $(U, R, s, t, c)$ comme un groupoïde en
espaces algébriques sur $X$, d'où une application
$f : [U/R] \to \mathcal{S}_X$ ; voir
Groupoïdes in Espaces, lemme \ref{spaces-groupoids-lemma-quotient-stack-arrows}.
De même, on a $f' : [U'/R'] \to X'$.

\medskip\noindent
Un objet du $2$-produit fibré
$\mathcal{S}_{X'} \times_{\mathcal{S}_X} [U/R]$ au-dessus d'un schéma $T$ sur $S$
revient à un morphisme $x' : T \to X'$ et à un objet $y$ de $[U/R]$ au-dessus de $T$
tels que la composée $g \circ x'$ soit égale à $f(y)$.
Cela a un sens, car les objets de $\mathcal{S}_X$ au-dessus de $T$
sont les morphismes $T \to X$. D'après Groupoïdes in Espaces, lemme
\ref{spaces-groupoids-lemma-quotient-stack-objects},
on peut supposer que $y$ est donné par une donnée de descente $[U/R]$, $(u_i, r_{ij})$,
relative à un recouvrement fppf $\{T_i \to T\}$.
L'égalité $g \circ x' = f(y)$ signifie que les diagrammes
$$
\vcenter{
\xymatrix{
T_i \ar[rr]_{u_i} \ar[d] & & U \ar[d] \\
T \ar[r]^{x'} & X' \ar[r]^g & X
}
}
\quad\text{et}\quad
\vcenter{
\xymatrix{
T_i \times_T T_j \ar[rr]_{r_{ij}} \ar[d] & & R \ar[d] \\
T \ar[r]^{x'} & X' \ar[r]^g & X
}
}
$$
sont commutatifs.

\medskip\noindent
D'autre part, d'après Groupoïdes in Espaces, lemme
\ref{spaces-groupoids-lemma-quotient-stack-objects}, un objet $y'$ de $[U'/R']$
au-dessus d'un schéma $T$ sur $S$
est donné par une donnée de descente $[U'/R']$, $(u'_i, r'_{ij})$,
relative à un recouvrement fppf $\{T_i \to T\}$.
En posant $f'(y') = x' : T \to X'$, on voit que
les diagrammes
$$
\vcenter{
\xymatrix{
T_i \ar[r]_{u'_i} \ar[d] & U' \ar[d] \\
T \ar[r]^{x'} & X'
}
}
\quad\text{et}\quad
\vcenter{
\xymatrix{
T_i \times_T T_j \ar[r]_{r'_{ij}} \ar[d] & U' \ar[d] \\
T \ar[r]^{x'} & X'
}
}
$$
sont commutatifs.

\medskip\noindent
Ces notations étant posées, on définit un foncteur
$$
[U'/R'] \longrightarrow \mathcal{S}_{X'} \times_{\mathcal{S}_X} [U/R]
$$
en envoyant $y' = (u'_i, r'_{ij})$, comme ci-dessus, sur l'objet
$(x', (u_i, r_{ij}))$, où $x' = f'(y')$, où
$u_i$ est la composée $T_i \to U' \to U$, et où
$r_{ij}$ est la composée $T_i \times_T T_j \to R' \to R$.
Réciproquement, étant donné un objet $(x', (u_i, r_{ij})$
du membre de droite, on peut l'envoyer sur l'objet
$((x', u_i), (x', r_{ij}))$ du membre de gauche.
Nous omettons d'indiquer ce qu'il faut faire des morphismes (cela fonctionne
exactement de la même manière).
\end{proof}





\section{Quotients catégoriques}
\label{groupoids-quotients-section-categorical}

\noindent
C'est le type de quotient le plus élémentaire que l'on puisse considérer.

\begin{definition}
\label{groupoids-quotients-definition-categorical}
Soit $S$ un schéma et soit $B$ un espace algébrique sur $S$.
Soit $j = (t, s) : R \to U \times_B U$ une prérelation en espaces algébriques
sur $B$.
\begin{enumerate}
\item On dit qu'un morphisme $\phi : U \to X$ d'espaces algébriques sur $B$
est un {\it quotient catégorique} s'il est $R$-invariant et si,
pour tout morphisme $R$-invariant $\psi : U \to Y$ d'espaces algébriques sur $B$,
il existe un unique morphisme $\chi : X \to Y$ tel que
$\psi = \phi \circ \chi$.
\item Soit $\mathcal{C}$ une sous-catégorie pleine de la catégorie des espaces
algébriques sur $B$. Supposons que $U$, $R$ soient des objets de $\mathcal{C}$.
Dans cette situation, on dit
qu'un morphisme $\phi : U \to X$ d'espaces algébriques sur $B$
est un {\it quotient catégorique dans $\mathcal{C}$}
si $X \in \Ob(\mathcal{C})$, si $\phi$ est $R$-invariant,
et si, pour tout morphisme $R$-invariant
$\psi : U \to Y$ avec $Y \in \Ob(\mathcal{C})$,
il existe un unique morphisme $\chi : X \to Y$ tel
que $\psi = \phi \circ \chi$.
\item Si $B = S$ et si $\mathcal{C}$ est la catégorie des schémas sur $S$,
on dit alors que $U \to X$ est un
{\it quotient catégorique dans la catégorie des schémas}, ou simplement un
{\it quotient catégorique en schémas}.
\end{enumerate}
\end{definition}

\noindent
On distingue souvent une catégorie $\mathcal{C}$ d'espaces algébriques sur $B$
par un axiome de séparation ; voir
l'exemple \ref{groupoids-quotients-example-categories}
pour quelques cas usuels.
Remarquons que $\phi : U \to X$ est un quotient catégorique si et seulement
si $U \to X$ est un coégalisateur des
morphismes $t, s : R \to U$ dans la catégorie. On en déduit immédiatement
le lemme suivant.

\begin{lemma}
\label{groupoids-quotients-lemma-categorical-unique}
Soit $S$ un schéma et soit $B$ un espace algébrique sur $S$.
Soit $j : R \to U \times_B U$ une prérelation en espaces algébriques sur $B$.
S'il existe un quotient catégorique dans la catégorie des espaces algébriques
sur $B$, il est unique à isomorphisme unique près.
Il en va de même pour les quotients catégoriques dans les sous-catégories pleines de
$\textit{Espaces}/B$.
\end{lemma}

\begin{proof}
Voir Catégories, section \ref{categories-section-coequalizers}.
\end{proof}

\begin{example}
\label{groupoids-quotients-example-categories}
Soit $S$ un schéma et soit $B$ un espace algébrique sur $S$.
Voici quelques exemples usuels de catégories $\mathcal{C}$
qui se présentent fréquemment lorsqu'on applique la
définition \ref{groupoids-quotients-definition-categorical} :
\begin{enumerate}
\item $\mathcal{C}$ est la catégorie de tous les espaces algébriques sur $B$ ;
\item $B$ est séparé et $\mathcal{C}$ est la catégorie de tous les espaces
algébriques séparés sur $B$ ;
\item $B$ est quasi-séparé et $\mathcal{C}$ est la catégorie de tous les espaces
algébriques quasi-séparés sur $B$ ;
\item $B$ est localement séparé et $\mathcal{C}$ est la catégorie de tous les espaces
algébriques localement séparés sur $B$ ;
\item $B$ est décent et $\mathcal{C}$ est la catégorie de tous les espaces algébriques
décents sur $B$ ; enfin,
\item $S = B$ et $\mathcal{C}$ est la catégorie des schémas sur $S$.
\end{enumerate}
Dans ces cas, si $\phi : U \to X$ est un quotient catégorique, on dit que
$U \to X$ est
(1) un {\it quotient catégorique},
(2) un {\it quotient catégorique dans les espaces algébriques séparés},
(3) un {\it quotient catégorique dans les espaces algébriques quasi-séparés},
(4) un {\it quotient catégorique dans les espaces algébriques localement séparés},
(5) un {\it quotient catégorique dans les espaces algébriques décents},
(6) un {\it quotient catégorique en schémas}.
\end{example}

\begin{definition}
\label{groupoids-quotients-definition-universal-categorical}
Soit $S$ un schéma et soit $B$ un espace algébrique sur $S$.
Soit $\mathcal{C}$ une sous-catégorie pleine de la catégorie des espaces
algébriques sur $B$, stable par produits fibrés.
Soit $j = (t, s) : R \to U \times_B U$ une prérelation dans
$\mathcal{C}$, et soit $U \to X$ un morphisme $R$-invariant avec
$X \in \Ob(\mathcal{C})$.
\begin{enumerate}
\item On dit que $U \to X$ est un {\it quotient catégorique universel}
dans $\mathcal{C}$ si, pour tout morphisme $X' \to X$ dans $\mathcal{C}$,
le morphisme $U' = X' \times_X U \to X'$ est le quotient catégorique dans
$\mathcal{C}$ du changement de base $j' : R' \to U'$ de $j$.
\item On dit que $U \to X$ est un {\it quotient catégorique uniforme}
dans $\mathcal{C}$ si, pour tout morphisme plat $X' \to X$ dans $\mathcal{C}$,
le morphisme $U' = X' \times_X U \to X'$ est le quotient catégorique dans
$\mathcal{C}$ du changement de base $j' : R' \to U'$ de $j$.
\end{enumerate}
\end{definition}

\begin{lemma}
\label{groupoids-quotients-lemma-categorical-reduced}
Dans la situation de la
définition \ref{groupoids-quotients-definition-categorical},
si $\phi : U \to X$ est un quotient catégorique et si $U$ est réduit,
alors $X$ est réduit. Le même énoncé vaut pour les quotients catégoriques dans
les catégories d'espaces $\mathcal{C}$ énumérées dans
l'exemple \ref{groupoids-quotients-example-categories}.
\end{lemma}

\begin{proof}
Soit $X_{red}$ la réduction de l'espace algébrique $X$.
Puisque $U$ est réduit, le morphisme $\phi : U \to X$ se factorise par
$i : X_{red} \to X$ (Propriétés des espaces algébriques, lemme
\ref{spaces-properties-lemma-map-into-reduction}). Notons ce morphisme
$\phi_{red} : U \to X_{red}$. Comme $\phi \circ s = \phi \circ t$, on
voit aussi que $\phi_{red} \circ s = \phi_{red} \circ t$ (car
$i : X_{red} \to X$ est un monomorphisme). La propriété universelle
de $\phi$ donne donc un morphisme $\chi : X \to X_{red}$ tel que
$\phi_{red} = \phi \circ \chi$. Par unicité, on voit que
$i \circ \chi = \text{id}_X$ et $\chi \circ i = \text{id}_{X_{red}}$.
Ainsi, $i$ est un isomorphisme et $X$ est réduit.

\medskip\noindent
Pour voir que cet argument fonctionne dans une catégorie $\mathcal{C}$, il
suffit de montrer que la réduction d'un objet de $\mathcal{C}$
est un objet de $\mathcal{C}$. Nous omettons de vérifier que c'est
le cas pour chacun des exemples usuels.
\end{proof}





\section{Quotients comme espaces d'orbites}
\label{groupoids-quotients-section-orbits}

\noindent
Soit $j = (t, s) : R \to U \times_B U$ une prérelation.
Si $j$ est une prérelation d'équivalence, alors, de façon informelle,
les « orbites » de $R$ dans $U$
sont les sous-ensembles $t(s^{-1}(\{u\}))$ de $U$. Cependant, si $j$ n'est qu'une
prérelation, il faut prendre la relation d'équivalence engendrée
par $R$.

\begin{definition}
\label{groupoids-quotients-definition-orbit}
Soit $S$ un schéma et soit $B$ un espace algébrique sur $S$.
Soit $j : R \to U \times_B U$ une prérelation sur $B$.
Si $u \in |U|$, alors l'{\it orbite}, ou plus précisément l'
{\it $R$-orbite} de $u$, est
$$
O_u =
\left\{
u' \in |U|\ :
\begin{matrix}
\exists n \geq 1, \ \exists u_0, \ldots, u_n \in |U|\text{ tels que }
u_0 = u \text{ et } u_n = u' \\
\text{et, pour tout }i \in \{0, \ldots, n - 1\},\text{ soit }
u_i = u_{i + 1}\text{ or } \\
\exists r \in |R|, \ s(r) = u_i, t(r) = u_{i + 1}
\text{ or } \\
\exists r \in |R|, \ t(r) = u_i, s(r) = u_{i + 1}
\end{matrix}
\right\}
$$
\end{definition}

\noindent
Il est clair que ce sont les classes d'équivalence d'une relation d'équivalence,
c'est-à-dire que $u' \in O_u$ si et seulement si $u \in O_{u'}$. Le lemme suivant
est une reformulation de
Groupoïdes in Espaces,
lemme \ref{spaces-groupoids-lemma-pre-equivalence-equivalence-relation-points}.

\begin{lemma}
\label{groupoids-quotients-lemma-pre-equivalence-equivalence-relation-points}
Soit $B \to S$ comme dans la section \ref{groupoids-quotients-section-conventions-notation}.
Soit $j : R \to U \times_B U$ une prérelation d'équivalence
d'espaces algébriques sur $B$. Alors
$$
O_u =
\{u' \in |U| \text{ tel qu'il existe } r \in |R|, \ s(r) = u, \ t(r) = u'\}.
$$
\end{lemma}

\begin{proof}
D'après le lemme précité de
Groupoïdes in Espaces,
lemme \ref{spaces-groupoids-lemma-pre-equivalence-equivalence-relation-points},
les orbites $O_u$ définies dans le lemme donnent une décomposition de
$|U|$ en réunion disjointe. Elles sont donc égales aux
orbites définies dans la définition \ref{groupoids-quotients-definition-orbit}.
\end{proof}

\begin{lemma}
\label{groupoids-quotients-lemma-invariant-map-constant-on-orbit}
Dans la situation de la définition \ref{groupoids-quotients-definition-orbit},
soit $\phi : U \to X$ un morphisme $R$-invariant d'espaces algébriques sur
$B$. Alors $|\phi| : |U| \to |X|$ est constant sur les orbites.
\end{lemma}

\begin{proof}
Il suffit de montrer que $\phi(u) = \phi(u')$
pour tous $u, u' \in |U|$ tels qu'il existe
$r \in |R|$ vérifiant $s(r) = u$ et $t(r) = u'$.
Cela est clair puisque $\phi$ égalise $s$ et $t$.
\end{proof}

\noindent
L'utilisation des orbites $O_u \subset |U|$ pour caractériser les propriétés
des applications quotients présente plusieurs difficultés. En voici une.
Supposons que $\Spec(k) \to B$
soit un point géométrique de $B$. Considérons l'application canonique
$$
U(k) \longrightarrow |U|.
$$
En général, les classes d'équivalence
de la relation d'équivalence engendrée par $j(R(k)) \subset U(k) \times U(k)$
ne sont pas les images réciproques des orbites $O_u \subset |U|$.
Un exemple élémentaire consiste à prendre $S = B = \Spec(\mathbf{Z})$,
$U = R = \Spec(k)$ avec $s = t = \text{id}_k$. Alors $|U| = |R|$ se réduit
à un point, mais $U(k)/R(k)$ est énorme.
Un exemple plus intéressant consiste à prendre $S = B = \Spec(\mathbf{Q})$,
à choisir des corps de nombres $K \subset L$, puis à poser $U = \Spec(L)$
et $R = \Spec(L \otimes_K L)$, avec les applications évidentes $s, t : R \to U$.
Dans ce cas, $|U|$ n'a toujours qu'un point, mais le quotient
$$
U(k)/R(k) = \Hom(K, k)
$$
contient plus d'un élément. Ces deux exemples montrent
que, si $U \to X$ est une application $R$-invariante et si l'on veut qu'elle
« sépare les orbites », on obtient une notion beaucoup plus forte et intéressante en
considérant les applications induites $U(k) \to X(k)$ et en exigeant que
ces applications séparent les orbites.

\medskip\noindent
Cette formulation présente encore une difficulté. Supposons en effet que
$S = B = \Spec(\mathbf{R})$,
$U = \Spec(\mathbf{C})$, et
$R = \Spec(\mathbf{C}) \amalg \Spec(K)$
pour une extension de corps $\sigma : \mathbf{C} \to K$.
Les applications $s, t$ sont données par l'identité sur la composante $\Spec(\mathbf{C})$,
mais, sur la seconde composante, elles sont données par $\sigma, \sigma \circ \tau$, où
$\tau$ est la conjugaison complexe. Si
$K$ est une extension non triviale de $\mathbf{C}$, les deux points
$1, \tau \in U(\mathbf{C})$ ne sont pas équivalents modulo
$j(R(\mathbf{C}))$. Toutefois, après avoir choisi une extension $\mathbf{C} \subset \Omega$
de cardinal suffisamment grand (par exemple supérieur au cardinal
de $K$), les images de $1, \tau \in U(\mathbf{C})$ dans
$U(\Omega)$ deviennent équivalentes ! Intuitivement, il semble clair que
cela se produit soit parce que $s, t : R \to U$ ne sont pas localement de type fini,
soit parce que le cardinal du corps $k$ n'est pas assez grand.

\medskip\noindent
Compte tenu de ce qui précède, posons la définition suivante.

\begin{definition}
\label{groupoids-quotients-definition-geometric-orbits}
Soit $S$ un schéma et soit $B$ un espace algébrique sur $S$.
Soit $j : R \to U \times_B U$ une prérelation sur $B$.
Soit $\Spec(k) \to B$ un point géométrique de $B$.
\begin{enumerate}
\item On dit que $\overline{u}, \overline{u}' \in U(k)$ sont
{\it faiblement $R$-équivalents} s'ils appartiennent à la même classe d'équivalence
pour la relation d'équivalence engendrée par la relation
$j(R(k)) \subset U(k) \times U(k)$.
\item On dit que $\overline{u}, \overline{u}' \in U(k)$ sont
{\it $R$-équivalents} si, pour un surcorps $k \subset \Omega$,
leurs images dans $U(\Omega)$ sont faiblement $R$-équivalentes.
\item L'{\it orbite faible}, ou plus précisément l'{\it $R$-orbite faible},
de $\overline{u} \in U(k)$ est l'ensemble de tous les
éléments de $U(k)$ faiblement $R$-équivalents à $\overline{u}$.
\item L'{\it orbite}, ou plus précisément l'{\it $R$-orbite},
de $\overline{u} \in U(k)$ est l'ensemble de tous les
éléments de $U(k)$ qui sont $R$-équivalents à $\overline{u}$.
\end{enumerate}
\end{definition}

\noindent
On verra que, dans les cas favorables, orbites et orbites faibles coïncident ; voir
le lemme \ref{groupoids-quotients-lemma-geometric-orbits}. Le lemme suivant illustre
la différence dans le cas particulier d'une prérelation d'équivalence.

\begin{lemma}
\label{groupoids-quotients-lemma-weak-orbit-pre-equivalence}
Soit $S$ un schéma et soit $B$ un espace algébrique sur $S$.
Soit $\Spec(k) \to B$ un point géométrique de $B$.
Soit $j : R \to U \times_B U$ une prérelation d'équivalence sur $B$.
Dans ce cas, l'orbite faible de $\overline{u} \in U(k)$ est simplement
$$
\{
\overline{u}' \in U(k)
\text{ tel que }
\exists \overline{r} \in R(k),
\ s(\overline{r}) = \overline{u},
\ t(\overline{r}) = \overline{u}'
\}
$$
et l'orbite de $\overline{u} \in U(k)$ est
$$
\{
\overline{u}' \in U(k) :
\exists\text{ field extension }K/k, \ \exists\ r \in R(K),
\ s(r) = \overline{u}, \ t(r) = \overline{u}'\}
$$
\end{lemma}

\begin{proof}
Cela résulte du fait que, par définition d'une prérelation d'équivalence, l'image
$j(R(k)) \subset U(k) \times U(k)$ est une relation d'équivalence.
\end{proof}

\noindent
Décrivons le procédé qui transforme toute prérelation en une
prérelation d'équivalence. Nous utiliserons les morphismes
\begin{equation}
\label{groupoids-quotients-equation-list}
\begin{matrix}
j_{diag} &
: &
U &
\longrightarrow &
U \times_B U, &
u &
\longmapsto &
(u, u) \\
j_{flip} &
: &
R &
\longrightarrow &
U \times_B U, &
r &
\longmapsto &
(s(r), t(r)) \\
j_{comp} &
: &
R \times_{s, U, t} R &
\longrightarrow &
U \times_B U, &
(r, r') &
\longmapsto &
(t(r), s(r'))
\end{matrix}
\end{equation}
On définit $j_1 = (t_1, s_1) : R_1 \to U \times_B U$ comme étant le morphisme
$$
j \amalg j_{diag} \amalg j_{flip} :
R \amalg U \amalg R
\longrightarrow
U \times_B U
$$
avec les notations de
l'équation (\ref{groupoids-quotients-equation-list}).
Pour $n > 1$, on pose
$$
j_n = (t_n, s_n) :
R_n = R_1 \times_{s_1, U, t_{n - 1}} R_{n - 1} \longrightarrow U \times_B U
$$
où $t_n$ provient de $t_1$ précomposé avec la projection sur $R_1$, et
$s_n$ de $s_{n - 1}$ précomposé avec la projection sur $R_{n - 1}$.
Enfin, on note
$$
j_\infty = (t_\infty, s_\infty) :
R_\infty = \coprod\nolimits_{n \geq 1} R_n
\longrightarrow
U \times_B U.
$$

\begin{lemma}
\label{groupoids-quotients-lemma-make-pre-equivalence}
Soit $S$ un schéma et soit $B$ un espace algébrique sur $S$.
Soit $j : R \to U \times_B U$ une prérelation sur $B$.
Alors $j_\infty : R_\infty \to U \times_B U$ est une
prérelation d'équivalence sur $B$. De plus,
\begin{enumerate}
\item $\phi : U \to X$ est $R$-invariant si et seulement s'il est
$R_\infty$-invariant ;
\item l'application canonique de faisceaux quotients $U/R \to U/R_\infty$ (voir
Groupoïdes in Espaces, section \ref{spaces-groupoids-section-quotient-sheaves})
est un isomorphisme ;
\item les $R$-orbites faibles coïncident avec les $R_\infty$-orbites faibles ;
\item les $R$-orbites coïncident avec les $R_\infty$-orbites ;
\item si $s, t$ sont localement de type fini, alors $s_\infty$, $t_\infty$
sont localement de type fini ;
\item ajouter ici d'autres propriétés au besoin.
\end{enumerate}
\end{lemma}

\begin{proof}
Omise. Indication pour (5) : toute propriété de $s, t$ stable par composition,
stable par changement de base et locale pour la topologie de Zariski sur la source
est héritée par $s_\infty, t_\infty$.
\end{proof}

\begin{lemma}
\label{groupoids-quotients-lemma-geometric-orbits}
Soit $S$ un schéma et soit $B$ un espace algébrique sur $S$.
Soit $j : R \to U \times_B U$ une prérelation sur $B$.
Soit $\Spec(k) \to B$ un point géométrique de $B$.
\begin{enumerate}
\item Si $s, t : R \to U$ sont localement de type fini,
alors l'équivalence $R$-faible sur $U(k)$ coïncide avec l'$R$-équivalence, et
les $R$-orbites faibles coïncident avec les $R$-orbites dans $U(k)$.
\item Si $k$ est de cardinal suffisamment grand, alors l'équivalence $R$-faible
sur $U(k)$ coïncide avec l'$R$-équivalence, et les $R$-orbites faibles
coïncident avec les $R$-orbites dans $U(k)$.
\end{enumerate}
\end{lemma}

\begin{proof}
Démontrons d'abord (1). Supposons $s, t$ localement de type fini. D'après le
lemme \ref{groupoids-quotients-lemma-make-pre-equivalence},
on peut supposer que $R$ est une prérelation d'équivalence.
Soit $k$ un corps algébriquement clos au-dessus de $B$.
Supposons que $\overline{u}, \overline{u}' \in U(k)$ soient $R$-équivalents.
Alors, pour une extension de corps $\Omega/k$, il existe
un point $\overline{r} \in R(\Omega)$ qui s'envoie sur
$(\overline{u}, \overline{u}') \in (U \times_B U)(\Omega)$ ; voir le
lemme \ref{groupoids-quotients-lemma-weak-orbit-pre-equivalence}.
Ainsi,
$$
Z = R \times_{j, U \times_B U, (\overline{u}, \overline{u}')} \Spec(k)
$$
est non vide. Comme $s$ est localement de type fini,
$j$ l'est aussi ; voir
Morphismes d'espaces algébriques, lemme \ref{spaces-morphisms-lemma-permanence-finite-type}.
Il s'ensuit que $Z$ est un espace algébrique non vide, localement de type fini
sur le corps algébriquement clos $k$ (utiliser
Morphismes d'espaces algébriques,
lemme \ref{spaces-morphisms-lemma-base-change-finite-type}).
Par conséquent, $Z$ possède un point à valeurs dans $k$ ; voir
Morphismes d'espaces algébriques, lemme
\ref{spaces-morphisms-lemma-locally-finite-type-surjective-geometric-points}.
Il existe donc $\overline{r} \in R(k)$ tel que
$j(\overline{r}) = (\overline{u}, \overline{u}')$, et l'on conclut comme voulu que
$\overline{u}, \overline{u}'$ sont $R$-équivalents.

\medskip\noindent
La démonstration de (2) est la même, si ce n'est qu'elle utilise
Morphismes d'espaces algébriques, lemme
\ref{spaces-morphisms-lemma-large-enough}
au lieu de
Morphismes d'espaces algébriques, lemme
\ref{spaces-morphisms-lemma-locally-finite-type-surjective-geometric-points}.
Cela montre que l'assertion vaut dès que $|k| > \lambda(R)$, où
$\lambda(R)$ est introduit juste avant
Morphismes d'espaces algébriques, lemme
\ref{spaces-morphisms-lemma-locally-finite-type-surjective-geometric-points}.
\end{proof}

\noindent
Dans la définition suivante, dire que « $k$ est un corps
au-dessus de $B$ » signifie que $\Spec(k)$ est muni d'un morphisme
$\Spec(k) \to B$.

\begin{definition}
\label{groupoids-quotients-definition-set-theoretically-invariant}
Soit $S$ un schéma et soit $B$ un espace algébrique sur $S$.
Soit $j : R \to U \times_B U$ une prérelation sur $B$.
\begin{enumerate}
\item On dit que $\phi : U \to X$ est {\it ensemblistement $R$-invariant}
si et seulement si l'application $U(k) \to X(k)$ égalise les deux applications
$s, t : R(k) \to U(k)$ pour tout corps algébriquement clos $k$
au-dessus de $B$.
\item On dit que $\phi : U \to X$ {\it sépare les orbites}, ou
{\it sépare les $R$-orbites}, s'il est ensemblistement $R$-invariant et si
$\phi(\overline{u}) = \phi(\overline{u}')$ dans $X(k)$ implique que
$\overline{u}, \overline{u}' \in U(k)$ appartiennent à la même orbite,
pour tout corps algébriquement clos $k$ au-dessus de $B$.
\end{enumerate}
\end{definition}

\noindent
Dans
l'exemple \ref{groupoids-quotients-example-not-invariant},
on montre que l'invariance ensembliste est une notion « trop faible » dans
la catégorie des espaces algébriques. Une reformulation plus géométrique
de l'invariance ensembliste et de la séparation des orbites est donnée au
lemme \ref{groupoids-quotients-lemma-separates-orbits}.

\begin{lemma}
\label{groupoids-quotients-lemma-set-theoretic-invariant}
Dans la situation de la définition \ref{groupoids-quotients-definition-set-theoretically-invariant},
un morphisme $\phi : U \to X$ est ensemblistement $R$-invariant si et
seulement si, pour tout corps algébriquement clos $k$ au-dessus de $B$, l'application
$U(k) \to X(k)$ est constante sur les orbites.
\end{lemma}

\begin{proof}
Cela résulte du fait que la condition est imposée pour tous les corps
algébriquement clos au-dessus de $B$.
\end{proof}

\begin{lemma}
\label{groupoids-quotients-lemma-invariant-set-theoretically-invariant}
Dans la situation de la définition \ref{groupoids-quotients-definition-set-theoretically-invariant},
un morphisme invariant est ensemblistement invariant.
\end{lemma}

\begin{proof}
Cela découle immédiatement des définitions.
\end{proof}

\begin{lemma}
\label{groupoids-quotients-lemma-set-theoretically-invariant-invariant-when-reduced}
Dans la situation de la définition \ref{groupoids-quotients-definition-set-theoretically-invariant},
soit $\phi : U \to X$ un morphisme d'espaces algébriques sur $B$.
Supposons que
\begin{enumerate}
\item $\phi$ soit ensemblistement $R$-invariant ;
\item $R$ soit réduit ; et
\item $X$ soit localement séparé sur $B$.
\end{enumerate}
Alors $\phi$ est $R$-invariant.
\end{lemma}

\begin{proof}
Considérons l'égalisateur
$$
Z = R \times_{(\phi, \phi) \circ j, X \times_B X, \Delta_{X/B}} X
$$
comme espace algébrique. Alors $Z \to R$ est une immersion en vertu de (3).
D'après (1), $|Z| \to |R|$ est surjectif. Il s'ensuit que
$Z \to R$ est une immersion fermée bijective (utiliser
Schémas, lemme \ref{schemes-lemma-immersion-when-closed})
et, par (2), on conclut que $Z = R$.
\end{proof}

\begin{example}
\label{groupoids-quotients-example-not-invariant}
Il existe des espaces algébriques réduits quasi-séparés $X$, $Y$ et deux
morphismes $a, b : Y \to X$ qui coïncident sur tous les points à valeurs dans $k$ sans être
égaux. Pour obtenir un exemple, prendre $Y = \Spec(k[[x]])$ et
$$
X = \mathbf{A}^1_k \Big/ \big(\Delta \amalg \{(x, -x) \mid x \not = 0\}\big)
$$
l'espace algébrique de
Espaces, exemple \ref{spaces-example-affine-line-involution}.
Les deux morphismes $a, b : Y \to X$
proviennent des deux applications $x \mapsto x$ et $x \mapsto -x$
de $Y$ vers $\mathbf{A}^1_k = \Spec(k[x])$. Au point générique,
les deux applications sont égales car, sur l'ouvert $x \not = 0$ de l'espace
$X$, les fonctions $x$ et $-x$ sont égales. Au point fermé,
les applications sont manifestement égales. On a aussi $a \not = b$.
Cela implique que le
lemme \ref{groupoids-quotients-lemma-set-theoretically-invariant-invariant-when-reduced}
cesse d'être vrai si l'on remplace (3) par l'hypothèse que $X$
est quasi-séparé. En effet, considérons le diagramme
$$
\xymatrix{
Y \ar[d]_{-1} \ar[r]_1 & Y \ar[d]^a \\
Y \ar[r]^a & X
}
$$
alors la composée $a \circ (-1) = b$. On peut donc poser $R = Y$,
$U = Y$, $s = 1$, $t = -1$, $\phi = a$ pour obtenir un exemple de morphisme
ensemblistement invariant qui n'est pas invariant.
\end{example}

\noindent
L'exemple précédent est instructif car l'application $Y \to X$ sépare même
les orbites. Il montre que, dans la catégorie des espaces algébriques, il existe tout
simplement trop de morphismes ensemblistement invariants. Définissons maintenant
ce que signifie pour $R$ d'être une relation d'équivalence ensembliste, en gardant
à l'esprit qu'il faut autoriser des extensions de corps pour que la définition
soit correcte.

\begin{definition}
\label{groupoids-quotients-definition-set-theoretic-equivalence}
Soit $S$ un schéma et soit $B$ un espace algébrique sur $S$.
Soit $j : R \to U \times_B U$ une prérelation sur $B$.
\begin{enumerate}
\item On dit que $j$ est une {\it prérelation d'équivalence ensembliste} si,
pour tout corps algébriquement clos $k$ au-dessus de $B$, la relation
$\sim_R$ sur $U(k)$ définie par
$$
\overline{u} \sim_R \overline{u}'
\Leftrightarrow
\begin{matrix}
\exists\text{ field extension }K/k, \ \exists\ r \in R(K), \\
s(r) = \overline{u}, \ t(r) = \overline{u}'
\end{matrix}
$$
est une relation d'équivalence.
\item On dit que $j$ est une {\it relation d'équivalence ensembliste}
si $j$ est universellement injectif et est une prérelation d'équivalence
ensembliste.
\end{enumerate}
\end{definition}

\noindent
Reformulons cette définition en termes plus géométriques.

\begin{lemma}
\label{groupoids-quotients-lemma-set-theoretic-pre-equivalence-geometric}
Dans la situation de la définition \ref{groupoids-quotients-definition-set-theoretic-equivalence},
les conditions suivantes sont équivalentes :
\begin{enumerate}
\item Le morphisme $j$ est une prérelation d'équivalence ensembliste.
\item Le sous-ensemble $j(|R|) \subset |U \times_B U|$ contient l'image de
$|j'|$ pour chacun des morphismes $j'$ de l'équation (\ref{groupoids-quotients-equation-list}).
\item Pour tout corps algébriquement clos $k$ au-dessus de $B$ et de cardinal suffisamment grand,
le sous-ensemble $j(R(k)) \subset U(k) \times U(k)$ est une relation
d'équivalence.
\end{enumerate}
Si $s, t$ sont localement de type fini, ces conditions équivalent aussi à
\begin{enumerate}
\item[(4)] Pour tout corps algébriquement clos $k$ au-dessus de $B$,
le sous-ensemble $j(R(k)) \subset U(k) \times U(k)$ est une relation d'équivalence.
\end{enumerate}
\end{lemma}

\begin{proof}
Supposons (2). Soit $k$ un corps algébriquement clos au-dessus de $B$.
Montrons que $\sim_R$ est une relation d'équivalence.
Supposons que $\overline{u}_i : \Spec(k) \to U$, $i = 1, 2$,
soient des points à valeurs dans $k$ de $U$. Supposons que $(\overline{u}_1, \overline{u}_2)$
soit l'image d'un point à valeurs dans $K$, $r \in R(K)$. Considérons le
diagramme commutatif en traits pleins
$$
\xymatrix{
\Spec(K') \ar@{..>}[r] \ar@{..>}[d]
&
\Spec(k) \ar[d]_{(\overline{u}_2, \overline{u}_1)} &
\Spec(K) \ar[d] \ar[l] \\
R \ar[r]^-j &
U \times_B U &
R \ar[l]_-{j_{flip}}
}
$$
On note aussi $r \in |R|$ l'image de $r$.
Par hypothèse, l'image de $|j_{flip}|$ est contenue dans l'image de
$|j|$ ; autrement dit, il existe $r' \in |R|$ tel que
$|j|(r') = |j_{flip}|(r)$. Mais remarquons que $(\overline{u}_2, \overline{u}_1)$
appartient à la classe d'équivalence qui définit $|j|(r')$ (par commutativité
de la partie en traits pleins du diagramme). Il existe donc une extension de corps
$K'/k$ et un morphisme $r' : \Spec(K) \to R$
(encore noté abusivement $r'$) tels que
$j \circ r' = (\overline{u}_2, \overline{u}_1) \circ i$,
où $i : \Spec(K') \to \Spec(K)$ est l'application évidente.
Autrement dit, la partie pointillée du diagramme est commutative.
Cela prouve que $\sim_R$ est une relation symétrique sur $U(k)$.
De même, le fait que l'image de $|j_{diag}|$ soit contenue
dans l'image de $|j|$ montre que $\sim_R$ est réflexive (détails omis).

\medskip\noindent
Pour montrer que $\sim_R$ est transitive, supposons donnés
$\overline{u}_i : \Spec(k) \to U$, $i = 1, 2, 3$,
des extensions de corps $K_i/k$ et des points
$r_i : \Spec(K_i) \to R$, $i = 1, 2$, tels que
$j(r_1) = (\overline{u}_1, \overline{u}_2)$ et
$j(r_1) = (\overline{u}_2, \overline{u}_3)$. On peut alors choisir un
diagramme commutatif de corps
$$
\xymatrix{
K & K_2 \ar[l] \\
K_1 \ar[u] & k \ar[l] \ar[u]
}
$$
et considérer $r_1, r_2 \in R(K)$. Considérons le
diagramme commutatif en traits pleins
$$
\xymatrix{
\Spec(K') \ar@{..>}[r] \ar@{..>}[d]
&
\Spec(k) \ar[d]_{(\overline{u}_1, \overline{u}_3)} &
\Spec(K) \ar[d]^{(r_1, r_2)} \ar[l]
\\
R \ar[r]^-j &
U \times_B U &
R \times_{s, U, t} R \ar[l]_-{j_{comp}}
}
$$
Le même raisonnement que dans la première partie de la démonstration, en utilisant
cette fois que $|j_{comp}|((r_1, r_2))$ appartient à l'image de $|j|$,
montre qu'il existe un corps $K'$ et des flèches pointillées qui rendent
le diagramme commutatif. Cela prouve que $\sim_R$ est transitive et achève
la preuve de l'implication (2) vers (1).

\medskip\noindent
Supposons (1) et soit $k$ un corps algébriquement clos au-dessus de $B$, dont le
cardinal est supérieur à $\lambda(R)$ ; voir
Morphismes d'espaces algébriques, lemme \ref{spaces-morphisms-lemma-large-enough}.
Supposons que $\overline{u} \sim_R \overline{u}'$, avec
$\overline{u}, \overline{u}' \in U(k)$. Par hypothèse, il existe
un point de $|R|$ qui s'envoie sur $(\overline{u}, \overline{u}') \in |U \times_B U|$.
Ainsi, d'après
Morphismes d'espaces algébriques, lemme \ref{spaces-morphisms-lemma-large-enough},
il existe $\overline{r} \in R(k)$ tel que
$j(\overline{r}) = (\overline{u}, \overline{u}')$. On voit ainsi
que (1) implique (3).

\medskip\noindent
Supposons (3). Montrons que
$\Im(|j_{comp}|) \subset \Im(|j|)$. Choisissons un point quelconque
$c \in |R \times_{s, U, t} R|$. On peut le représenter par un morphisme
$\overline{c} : \Spec(k) \to R \times_{s, U, t} R$, le corps $k$ au-dessus de $B$
étant de cardinal suffisamment grand. Par hypothèse,
$j_{comp}(\overline{c}) \in U(k) \times U(k) = (U \times_B U)(k)$
est aussi l'image $j(\overline{r})$ d'un certain $\overline{r} \in R(k)$.
Par conséquent, $j_{comp}(c) = j(r)$ dans $|U \times_B U|$, comme voulu (où
$r \in |R|$ est la classe d'équivalence de $\overline{r}$). Le même argument
montre aussi que $\Im(|j_{diag}|) \subset \Im(|j|)$ et
$\Im(|j_{flip}|) \subset \Im(|j|)$ (détails omis).
On voit ainsi que (3) implique (2). Nous avons donc montré que
(1), (2) et (3) sont toutes équivalentes.

\medskip\noindent
Il est clair que (4) implique (3) (sans aucune hypothèse sur $s$, $t$).
Pour achever la démonstration du lemme, montrons que (1) implique (4) si $s, t$
sont localement de type fini. Soit donc $k$ un corps algébriquement clos
au-dessus de $B$. Supposons que $\overline{u} \sim_R \overline{u}'$, avec
$\overline{u}, \overline{u}' \in U(k)$. Par hypothèse, l'espace algébrique
$Z = R \times_{j, U \times_B U, (\overline{u}, \overline{u}')} \Spec(k)$
est non vide. D'autre part, puisque $j = (t, s)$ est localement de type fini,
le morphisme $Z \to \Spec(k)$ est lui aussi localement de type fini (utiliser
Morphismes d'espaces algébriques, lemmes \ref{spaces-morphisms-lemma-permanence-finite-type}
et \ref{spaces-morphisms-lemma-base-change-finite-type}).
Ainsi, $Z$ possède un point à valeurs dans $k$ d'après
Morphismes d'espaces algébriques, lemme
\ref{spaces-morphisms-lemma-locally-finite-type-surjective-geometric-points},
et l'on conclut que $(\overline{u}, \overline{u}') \in j(R(k))$,
comme voulu. Cela achève la démonstration du lemme.
\end{proof}

\begin{lemma}
\label{groupoids-quotients-lemma-set-theoretic-equivalence-geometric}
Dans la situation de la définition \ref{groupoids-quotients-definition-set-theoretic-equivalence},
les conditions suivantes sont équivalentes :
\begin{enumerate}
\item Le morphisme $j$ est une relation d'équivalence ensembliste.
\item Le morphisme $j$ est universellement injectif et
$j(|R|) \subset |U \times_B U|$ contient l'image de
$|j'|$ pour chacun des morphismes $j'$ de l'équation (\ref{groupoids-quotients-equation-list}).
\item Pour tout corps algébriquement clos $k$ au-dessus de $B$ et de cardinal suffisamment grand,
l'application $j : R(k) \to U(k) \times U(k)$ est injective et
son image est une relation d'équivalence.
\end{enumerate}
Si $j$ est décent, localement séparé ou quasi-séparé,
ces conditions équivalent aussi à
\begin{enumerate}
\item[(4)] Pour tout corps algébriquement clos $k$ au-dessus de $B$,
l'application $j : R(k) \to U(k) \times U(k)$ est injective et son image
est une relation d'équivalence.
\end{enumerate}
\end{lemma}

\begin{proof}
Les implications (1) $\Rightarrow$ (2) et (2) $\Rightarrow$ (3) résultent du
lemme \ref{groupoids-quotients-lemma-set-theoretic-pre-equivalence-geometric}
et des définitions. Le même lemme montre que (3) implique que
$j$ est une prérelation d'équivalence ensembliste. Mais la condition
(3) implique aussi, bien entendu, que $j$ est universellement injectif ; voir
Morphismes d'espaces algébriques, lemme \ref{spaces-morphisms-lemma-universally-injective} ;
ainsi, $j$ est bien une relation d'équivalence ensembliste.
Nous savons donc que (1), (2) et (3) sont toutes équivalentes.

\medskip\noindent
La condition (4) implique (3) sans autre hypothèse sur $j$. Supposons que $j$
soit décent, localement séparé ou quasi-séparé, et que les conditions équivalentes
(1), (2), (3) soient satisfaites. D'après
Compléments sur les morphismes d'espaces algébriques,
lemme \ref{spaces-more-morphisms-lemma-when-universally-injective-radicial},
$j$ est radiciel.
Soit $k$ un corps algébriquement clos quelconque au-dessus de $B$. Soient
$\overline{u}, \overline{u}' \in U(k)$ tels que
$\overline{u} \sim_R \overline{u}'$. On voit que
$R \times_{U \times_B U, (\overline{u}, \overline{u}')} \Spec(k)$
est non vide. Comme $j$ est radiciel, sa réduction est donc le spectre d'un
corps purement inséparable sur $k$. Puisque $k = \overline{k}$, on voit que
c'est le spectre de $k$. Il existe donc un point $\overline{r} \in R(k)$
tel que $t(\overline{r}) = \overline{u}$ et $s(\overline{r}) = \overline{u}'$,
comme voulu.
\end{proof}

\begin{lemma}
\label{groupoids-quotients-lemma-set-theoretic-equivalence}
Soit $S$ un schéma et soit $B$ un espace algébrique sur $S$.
Soit $j : R \to U \times_B U$ une prérelation sur $B$.
\begin{enumerate}
\item Si $j$ est une prérelation d'équivalence, alors $j$ est une
prérelation d'équivalence ensembliste. C'est en particulier le cas
lorsque $j$ provient d'un groupoïde en espaces algébriques ou de
l'action d'un espace algébrique en groupes sur $U$.
\item Si $j$ est une relation d'équivalence, alors $j$ est une
relation d'équivalence ensembliste.
\end{enumerate}
\end{lemma}

\begin{proof}
Omise.
\end{proof}

\begin{lemma}
\label{groupoids-quotients-lemma-separates-orbits}
Soit $B \to S$ comme dans la section \ref{groupoids-quotients-section-conventions-notation}.
Soit $j : R \to U \times_B U$ une prérelation.
Soit $\phi : U \to X$ un morphisme d'espaces algébriques sur $B$.
Considérons le diagramme
$$
\xymatrix{
(U \times_X U) \times_{(U \times_B U)} R \ar[d]^q \ar[r]_-p & R \ar[d]^j \\
U \times_X U \ar[r]^c & U \times_B U
}
$$
Alors :
\begin{enumerate}
\item Le morphisme $\phi$ est ensemblistement invariant si et seulement
si $p$ est surjectif.
\item Si $j$ est une prérelation d'équivalence ensembliste, alors
$\phi$ sépare les orbites si et seulement si $p$ et $q$ sont surjectifs.
\item Si $p$ et $q$ sont surjectifs, alors $j$ est une prérelation
d'équivalence ensembliste (et $\phi$ sépare les orbites).
\item Si $\phi$ est $R$-invariant et si $j$ est une prérelation d'équivalence
ensembliste, alors $\phi$ sépare les orbites si et seulement si le morphisme induit
$R \to U \times_X U$ est surjectif.
\end{enumerate}
\end{lemma}

\begin{proof}
Supposons que $\phi$ soit ensemblistement invariant. Cela signifie que, pour tout
corps algébriquement clos $k$ au-dessus de $B$ et tout $\overline{r} \in R(k)$,
on a $\phi(s(\overline{r})) = \phi(t(\overline{r}))$. Ainsi,
$((\phi(t(\overline{r})), \phi(s(\overline{r}))), \overline{r})$
définit un point du produit fibré qui s'envoie sur $\overline{r}$ par
$p$. Cela montre que $p$ est surjectif. Réciproquement, supposons $p$
surjectif. Choisissons $\overline{r} \in R(k)$. Puisque $p$ est surjectif, on
peut trouver une extension de corps $K/k$ et un point à valeurs dans $K$
$\tilde r$ du produit fibré tel que $p(\tilde r) = \overline{r}$.
Alors $q(\tilde r) \in U \times_X U$ s'envoie sur
$(t(\overline{r}), s(\overline{r}))$ dans $U \times_B U$, et l'on conclut
que $\phi(s(\overline{r})) = \phi(t(\overline{r}))$. Cela prouve
que $\phi$ est ensemblistement invariant.

\medskip\noindent
Les démonstrations de (2), (3) et (4) sont omises. Indication : supposer que $k$ est un
corps algébriquement clos au-dessus de $B$, de grand cardinal. Considérer le
diagramme d'ensembles associé
$$
\xymatrix{
(U(k) \times_{X(k)} U(k)) \times_{U(k) \times U(k)} R(k) \ar[d]^q \ar[r]_-p &
R(k) \ar[d]^j \\
U(k) \times_{X(k)} U(k) \ar[r]^c & U(k) \times U(k)
}
$$
Grâce aux lemmes précédents, les équivalences formulées en (2), (3) et (4) deviennent
des questions ensemblistes relatives au diagramme que l'on vient d'afficher, en utilisant
le fait que la surjectivité se traduit par la surjectivité sur les points à valeurs dans $k$, d'après
Morphismes d'espaces algébriques, lemme \ref{spaces-morphisms-lemma-large-enough}.
\end{proof}

\noindent
Comme nous avons vu plus haut que la notion de morphisme ensemblistement
invariant est assez faible dans la catégorie des espaces algébriques,
nous définissons comme suit un espace d'orbites pour une prérelation.

\begin{definition}
\label{groupoids-quotients-definition-orbit-space}
Soit $B \to S$ comme dans la section \ref{groupoids-quotients-section-conventions-notation}.
Soit $j : R \to U \times_B U$ une prérelation.
On dit que $\phi : U \to X$ est un {\it espace d'orbites pour $R$} si
\begin{enumerate}
\item $\phi$ est $R$-invariant ;
\item $\phi$ sépare les $R$-orbites ; et
\item $\phi$ est surjectif.
\end{enumerate}
\end{definition}

\noindent
La définition de la séparation des $R$-orbites fait intervenir des
points à valeurs dans des corps algébriquement clos. Mais, comme on l'a vu,
cela correspond dans de nombreux cas à la surjectivité de certains
morphismes d'espaces algébriques canoniquement associés.
Résumons la discussion précédente dans la caractérisation suivante
des espaces d'orbites.

\begin{lemma}
\label{groupoids-quotients-lemma-orbit-space}
Soit $B \to S$ comme dans la section \ref{groupoids-quotients-section-conventions-notation}.
Soit $j : R \to U \times_B U$ une prérelation d'équivalence
ensembliste. Un morphisme $\phi : U \to X$ est un espace d'orbites pour $R$ si et seulement si
\begin{enumerate}
\item $\phi \circ s = \phi \circ t$, c'est-à-dire que $\phi$ est invariant ;
\item le morphisme induit $(t, s) : R \to U \times_X U$ est surjectif ; et
\item le morphisme $\phi : U \to X$ est surjectif.
\end{enumerate}
Cette caractérisation s'applique par exemple si $j$ est une prérelation d'équivalence,
s'il provient d'un groupoïde en espaces algébriques sur $B$, ou s'il provient de l'action
d'un espace algébrique en groupes sur $B$ sur $U$.
\end{lemma}

\begin{proof}
Cela résulte immédiatement du lemme \ref{groupoids-quotients-lemma-separates-orbits}, partie (4).
\end{proof}

\noindent
Dans le lemme suivant, il ne suffit (probablement) pas de supposer seulement que
les morphismes $s, t$ sont localement de type fini. En effet,
il peut arriver qu'une application $\phi : U \to X$ soit un espace d'orbites sans être
localement de type fini. Dans ce cas, $U(k) \to X(k)$ peut ne pas
être surjective pour tout corps algébriquement clos $k$ au-dessus de $B$.

\begin{lemma}
\label{groupoids-quotients-lemma-orbit-space-locally-finite-type-over-base}
Soit $B \to S$ comme dans la section \ref{groupoids-quotients-section-conventions-notation}.
Soit $j = (t, s) : R \to U \times_B U$ une prérelation.
Supposons que $R, U$ soient localement de type fini sur $B$.
Soit $\phi : U \to X$ un morphisme $R$-invariant d'espaces algébriques sur $B$.
Alors $\phi$ est un espace d'orbites pour $R$ si et seulement si l'application naturelle
$$
U(k)/\big(\text{equivalence relation generated by }j(R(k))\big)
\longrightarrow
X(k)
$$
est bijective pour tout corps algébriquement clos $k$ au-dessus de $B$.
\end{lemma}

\begin{proof}
Remarquons que, puisque $U$, $R$ sont localement de type fini sur $B$, tous les
morphismes $s, t, j, \phi$ sont localement de type fini ; voir
Morphismes d'espaces algébriques, lemme \ref{spaces-morphisms-lemma-permanence-finite-type}.
Nous utiliserons aussi sans autre mention
Morphismes d'espaces algébriques, lemme
\ref{spaces-morphisms-lemma-locally-finite-type-surjective-geometric-points}.
Supposons que $\phi$ soit un espace d'orbites. Soit $k$ un corps algébriquement clos
quelconque au-dessus de $B$. Soit $\overline{x} \in X(k)$. Considérons
$U \times_{\phi, X, \overline{x}} \Spec(k)$.
C'est un espace algébrique non vide,
localement de type fini sur $k$. Il possède donc un point à valeurs dans $k$.
Cela montre que l'application affichée dans le lemme est surjective.
Supposons que $\overline{u}, \overline{u}' \in U(k)$ aient la même
image dans $X(k)$. D'après la
définition \ref{groupoids-quotients-definition-set-theoretically-invariant},
cela signifie que $\overline{u}, \overline{u}'$ sont dans la même
$R$-orbite. D'après le lemme \ref{groupoids-quotients-lemma-geometric-orbits}, cela signifie qu'ils
sont équivalents pour la relation d'équivalence engendrée par
$j(R(k))$. Ainsi, le morphisme affiché est injectif.

\medskip\noindent
Réciproquement, supposons que l'application affichée soit bijective pour tout corps
algébriquement clos $k$ au-dessus de $B$. Cette condition implique clairement que $\phi$
est surjectif. Nous avons déjà supposé que $\phi$ est $R$-invariant.
Enfin, l'injectivité de toutes les applications affichées implique que
$\phi$ sépare les orbites. Par conséquent, $\phi$ est un espace d'orbites.
\end{proof}










\section{Quotients grossiers}
\label{groupoids-quotients-section-coarse}

\noindent
Nous ajoutons seulement cette notion ici afin de pouvoir dire plus tard que les quotients
grossiers correspondent aux espaces de modules grossiers (ou aux schémas de modules).

\begin{definition}
\label{groupoids-quotients-definition-coarse}
Soit $S$ un schéma et soit $B$ un espace algébrique sur $S$.
Soit $j : R \to U \times_B U$ une prérelation.
Un morphisme $\phi : U \to X$ d'espaces algébriques sur $B$
est appelé {\it quotient grossier} si
\begin{enumerate}
\item $\phi$ est un quotient catégorique ; et
\item $\phi$ est un espace d'orbites.
\end{enumerate}
Si $S = B$ et si $U$, $R$ sont tous deux des schémas, on dit qu'un morphisme de schémas
$\phi : U \to X$ est un {\it quotient grossier en schémas} si
\begin{enumerate}
\item $\phi$ est un quotient catégorique en schémas ; et
\item $\phi$ est un espace d'orbites.
\end{enumerate}
\end{definition}

\noindent
Dans de nombreuses situations, les espaces algébriques $R$ et $U$ sont localement de type fini
sur $B$, et la condition d'espace d'orbites signifie simplement que
$$
U(k)/\big(\text{equivalence relation generated by }j(R(k))\big)
\cong
X(k)
$$
pour tout corps algébriquement clos $k$. Voir le
lemme \ref{groupoids-quotients-lemma-orbit-space-locally-finite-type-over-base}.
Si $j$ est en outre une prérelation d'équivalence (ensembliste), alors la condition
équivaut simplement à ce que $U(k)/j(R(k)) \to X(k)$ soit bijective pour tout
corps algébriquement clos $k$.









\section{Propriétés topologiques}
\label{groupoids-quotients-section-topological}

\noindent
Soit $S$ un schéma et soit $B$ un espace algébrique sur $S$.
Soit $j : R \to U \times_B U$ une prérelation.
On dit qu'un sous-ensemble $T \subset |U|$ est {\it $R$-invariant} si
$s^{-1}(T) = t^{-1}(T)$ comme sous-ensembles de $|R|$.
Remarquons que, si $T$ est fermé, il peut arriver que
le sous-espace fermé réduit correspondant de $U$ ne soit pas $R$-invariant
(au sens de
Groupoïdes in Espaces, définition
\ref{spaces-groupoids-definition-invariant-open}),
car les images réciproques $s^{-1}(T)$, $t^{-1}(T)$ peuvent ne pas être réduites.
Voici quelques conditions que l'on peut considérer pour un
morphisme invariant $\phi : U \to X$.

\begin{definition}
\label{groupoids-quotients-definition-topological}
Soit $S$ un schéma et soit $B$ un espace algébrique sur $S$.
Soit $j : R \to U \times_B U$ une prérelation.
Soit $\phi : U \to X$ un morphisme $R$-invariant d'espaces algébriques sur $B$.
\begin{enumerate}
\item
\label{groupoids-quotients-item-submersive}
Le morphisme $\phi$ est submersif.
\item
\label{groupoids-quotients-item-invariant-closed}
Pour tout sous-ensemble fermé $R$-invariant $Z \subset |U|$, l'image
$\phi(Z)$ est fermée dans $|X|$.
\item
\label{groupoids-quotients-item-intersect-invariant-closed}
La condition (\ref{groupoids-quotients-item-invariant-closed}) est satisfaite et, pour toute paire de
sous-ensembles fermés $R$-invariants $Z_1, Z_2 \subset |U|$, on a
$$
\phi(Z_1 \cap Z_2) = \phi(Z_1) \cap \phi(Z_2)
$$
\item Le morphisme $(t, s) : R \to U \times_X U$ est universellement submersif.
\label{groupoids-quotients-item-strong}
\end{enumerate}
Pour chacune de ces propriétés, on peut aussi exiger qu'elle subsiste après tout
changement de base plat, ou après tout changement de base ; voir la
définition \ref{groupoids-quotients-definition-base-change}. On dit alors que la condition
(\ref{groupoids-quotients-item-submersive}),
(\ref{groupoids-quotients-item-invariant-closed}),
(\ref{groupoids-quotients-item-intersect-invariant-closed}) ou
(\ref{groupoids-quotients-item-strong}) est satisfaite {\it uniformément} ou {\it universellement}.
\end{definition}









\section{Fonctions invariantes}
\label{groupoids-quotients-section-functions}

\noindent
Dans certains cas, il est commode de déterminer précisément le faisceau structural
d'un quotient en exigeant que toute fonction invariante soit une section locale
du faisceau structural du quotient.

\begin{definition}
\label{groupoids-quotients-definition-functions}
Soit $S$ un schéma et soit $B$ un espace algébrique sur $S$.
Soit $j : R \to U \times_B U$ une prérelation.
Soit $\phi : U \to X$ un morphisme $R$-invariant.
Notons $\phi' = \phi \circ s = \phi \circ t : R \to X$.
\begin{enumerate}
\item On note $(\phi_*\mathcal{O}_U)^R$ la sous-$\mathcal{O}_X$-algèbre
de $\phi_*\mathcal{O}_U$ qui est le noyau de la double flèche
$$
\xymatrix{
\phi_*\mathcal{O}_U
\ar@<1ex>[rr]^{\phi_*s^\sharp}
\ar@<-1ex>[rr]_{\phi_*t^\sharp}
& &
\phi'_*\mathcal{O}_R
}
$$
sur $X_\etale$. On l'appelle parfois le
{\it faisceau des fonctions $R$-invariantes sur $X$}.
\item On dit que {\it les fonctions sur $X$ sont les fonctions $R$-invariantes sur
$U$} si l'application naturelle $\mathcal{O}_X \to (\phi_*\mathcal{O}_U)^R$
est un isomorphisme.
\end{enumerate}
\end{definition}

\noindent
On peut bien sûr exiger que cette propriété subsiste après tout
changement de base (plat ou quelconque), ce qui donne une notion (uniforme ou) universelle.
Cette condition est souvent ajoutée
à d'autres afin d'obtenir un quotient davantage déterminé. Une grande partie de
la motivation de toute la théorie provient bien sûr du cas particulier suivant :
$U = \Spec(A)$ est un schéma affine sur un corps
$S = B = \Spec(k)$ et $R = G \times U$,
où $G$ est un schéma en groupes affine sur $k$. Dans ce cas,
on peut choisir comme quotient
$$
X = \Spec(A^G)
$$
de sorte qu'au moins la condition de la définition précédente soit satisfaite.
Bien que cette construction soit agréable, elle ne donne souvent pas le bon
quotient ; par exemple, si $U = \text{GL}_{n, k}$ et si $G$ est le groupe des
matrices triangulaires supérieures, elle donne $X = \Spec(k)$, alors qu'il
existe un bien meilleur quotient (à savoir la variété de drapeaux).









\section{Bons quotients}
\label{groupoids-quotients-section-good}

\noindent
La définition suivante est particulièrement utile lorsqu'on prend des quotients
par des actions de groupes.

\begin{definition}
\label{groupoids-quotients-definition-good}
Soit $S$ un schéma et soit $B$ un espace algébrique sur $S$.
Soit $j : R \to U \times_B U$ une prérelation.
Un morphisme $\phi : U \to X$ d'espaces algébriques sur $B$
est appelé {\it bon quotient} si
\begin{enumerate}
\item $\phi$ est invariant ;
\item $\phi$ est affine ;
\item $\phi$ est surjectif ;
\item la condition (\ref{groupoids-quotients-item-intersect-invariant-closed}) est satisfaite universellement ; et
\item les fonctions sur $X$ sont les fonctions $R$-invariantes sur $U$.
\end{enumerate}
\end{definition}

\noindent
Dans \cite{seshadri_quotients}, Seshadri donne presque la même définition,
mais, à la place de (4), il exige simplement que la condition
(\ref{groupoids-quotients-item-intersect-invariant-closed}) soit satisfaite, sans exiger
qu'elle le soit universellement.








\section{Quotients géométriques}
\label{groupoids-quotients-section-geometric}

\noindent
C'est la définition de Mumford d'un quotient géométrique (du moins celle
de la première édition de GIT ; pour autant que nous puissions en juger, les éditions ultérieures
ont remplacé « universellement submersif » par « submersif »).

\begin{definition}
\label{groupoids-quotients-definition-geometric}
Soit $S$ un schéma et soit $B$ un espace algébrique sur $S$.
Soit $j : R \to U \times_B U$ une prérelation.
Un morphisme $\phi : U \to X$ d'espaces algébriques sur $B$
est appelé {\it quotient géométrique} si
\begin{enumerate}
\item $\phi$ est un espace d'orbites ;
\item la condition (\ref{groupoids-quotients-item-submersive}) est satisfaite universellement, c'est-à-dire
que $\phi$ est universellement submersif ; et
\item les fonctions sur $X$ sont les fonctions $R$-invariantes sur $U$.
\end{enumerate}
\end{definition}















% OK, start here.
%

\chapter{Compléments sur la cohomologie des espaces}


\phantomsection
\label{spaces-more-cohomology-section-phantom}





\section{Introduction}
\label{spaces-more-cohomology-section-introduction}

\noindent
Nous poursuivons dans ce chapitre l'étude commencée dans
Cohomologie des espaces, section \ref{spaces-cohomology-section-introduction}.
On peut aussi voir ce chapitre comme l'analogue, pour les espaces algébriques,
du chapitre consacré à la cohomologie étale des schémas; voir
Cohomologie étale, section \ref{etale-cohomology-section-introduction}.

\medskip\noindent
En fait, ce chapitre vise principalement à transposer dans le langage des
espaces algébriques des résultats déjà démontrés pour les schémas.
Certains de nos résultats se trouvent dans \cite{Kn}.





\section{Conventions}
\label{spaces-more-cohomology-section-conventions}

\noindent
Nous supposons une fois pour toutes que tous les schémas appartiennent à
un grand site fppf $\Sch_{fppf}$. Tous les anneaux $A$ considérés
ont en outre la propriété que $\Spec(A)$ est (isomorphe à) un
objet de ce grand site.

\medskip\noindent
Soient $S$ un schéma et $X$ un espace algébrique sur $S$.
Dans ce chapitre et dans les suivants, nous noterons $X \times_S X$
le produit de $X$ par lui-même (dans la catégorie des espaces
algébriques sur $S$), plutôt que $X \times X$.







\section{Transfert des résultats établis pour les schémas}
\label{spaces-more-cohomology-section-api}

\noindent
Nous expliquons brièvement dans cette section comment des résultats concernant
les schémas entraînent des résultats pour les espaces algébriques
(représentables) et les morphismes (représentables) d'espaces algébriques.
On dispose d'énoncés plus forts pour les modules quasi-cohérents
(car la cohomologie étale d'un module quasi-cohérent sur un schéma
coïncide avec sa cohomologie de Zariski); cette question a déjà été
traitée dans Cohomologie des espaces, section
\ref{spaces-cohomology-section-higher-direct-image}.

\medskip\noindent
Soient $S$ un schéma et $X$ un espace algébrique sur $S$.
Supposons maintenant que $X$ soit représentable par le schéma $X_0$
(notation peu commode mais temporaire; nous disons habituellement simplement
que « $X$ est un schéma »). Dans ce cas, $X$ et $X_0$ ont le même petit
site étale:
$$
X_\etale = (X_0)_\etale
$$
Ceci est signalé dans Propriétés des espaces, section
\ref{spaces-properties-section-etale-site}.
De plus, si $f : X \to Y$ est un morphisme d'espaces algébriques représentables
sur $S$ et si $f_0 : X_0 \to Y_0$ est un morphisme de schémas
représentant $f$, alors les morphismes induits entre les petits topos
étales coïncident:
$$
\xymatrix{
\Sh(X_\etale) \ar[rr]_{f_{small}} \ar@{=}[d] & &
\Sh(Y_\etale) \ar@{=}[d] \\
\Sh((X_0)_\etale) \ar[rr]^{(f_0)_{small}} & &
\Sh((Y_0)_\etale)
}
$$
Voir Propriétés des espaces, lemme
\ref{spaces-properties-lemma-functoriality-etale-site} et
Topologies, lemme \ref{topologies-lemma-morphism-big-small-etale}.

\medskip\noindent
Il n'y a donc absolument aucune différence entre la cohomologie étale
d'un schéma et celle de l'espace algébrique correspondant.
Il en va de même des images directes supérieures par des morphismes de schémas.
En effet, si $f : X \to Y$ est un morphisme d'espaces algébriques sur $S$
qui est représentable (par des schémas), les images directes supérieures
$R^if_*\mathcal{F}$ d'un faisceau $\mathcal{F}$ sur $X_\etale$
peuvent se calculer localement pour la topologie étale sur $Y$ (Cohomologie
sur les sites, lemme \ref{sites-cohomology-lemma-higher-direct-images});
ceci ramène donc souvent les calculs et les démonstrations au cas où
$Y$ et $X$ sont des schémas.

\medskip\noindent
Nous utiliserons ce qui précède sans autre mention dans ce chapitre.
Le même fait vaut pour les autres topologies; nous l'énonçons ici
explicitement sous forme d'un lemme pour la cohomologie.

\begin{lemma}
\label{spaces-more-cohomology-lemma-compare-cohomology-other-topologies}
Soit $S$ un schéma. Soit $\tau \in \{\etale, fppf, ph\}$ (en ajouter ici).
Le foncteur d'inclusion
$$
(\Sch/S)_\tau \longrightarrow (\textit{Espaces}/S)_\tau
$$
est un foncteur cocontinu spécial
(Sites, définition \ref{sites-definition-special-cocontinuous-functor})
et identifie par conséquent les topos.
\end{lemma}

\begin{proof}
Les conditions de Sites, lemme \ref{sites-lemma-equivalence},
se vérifient immédiatement, puisque notre foncteur est pleinement fidèle
et que tout espace algébrique admet un recouvrement étale par des schémas.
\end{proof}







\section{Changement de base propre}
\label{spaces-more-cohomology-section-proper-base-change}

\noindent
Le théorème de changement de base propre pour les espaces algébriques se déduit,
au prix d'un peu de travail, du théorème de changement de base propre
pour les schémas et du lemme de Chow.

\begin{lemma}
\label{spaces-more-cohomology-lemma-surjective-proper}
Soit $S$ un schéma. Soit $f : Y \to X$ un morphisme propre surjectif
d'espaces algébriques sur $S$. Soit $\mathcal{F}$ un faisceau sur $X_\etale$.
Alors $\mathcal{F} \to f_*f^{-1}\mathcal{F}$ est injectif et son
image est l'égalisateur des deux applications
$f_*f^{-1}\mathcal{F} \to g_*g^{-1}\mathcal{F}$, où
$g$ est le morphisme structural $g : Y \times_X Y \to X$.
\end{lemma}

\begin{proof}
Pour tout morphisme surjectif $f : Y \to X$ d'espaces algébriques sur $S$,
l'application $\mathcal{F} \to f_*f^{-1}\mathcal{F}$ est injective.
En effet, si $\overline{x}$ est un point géométrique de $X$, choisissons
un point géométrique $\overline{y}$ de $Y$ situé au-dessus de $\overline{x}$
et considérons
$$
\mathcal{F}_{\overline{x}} \to
(f_*f^{-1}\mathcal{F})_{\overline{x}} \to
(f^{-1}\mathcal{F})_{\overline{y}} = \mathcal{F}_{\overline{x}}
$$
Voir Propriétés des espaces, lemme \ref{spaces-properties-lemma-stalk-pullback},
pour la dernière égalité.

\medskip\noindent
Le second énoncé est local sur $X$ pour la topologie étale; nous pouvons donc
supposer, et nous supposons, que $Y$ est un schéma affine.

\medskip\noindent
Choisissons un morphisme propre surjectif $Z \to Y$, où $Z$ est un schéma; voir
Cohomologie des espaces, lemme \ref{spaces-cohomology-lemma-weak-chow}.
Le résultat pour $Z \to X$ entraîne celui pour $Y \to X$.
Comme $Z \to X$ est un morphisme propre surjectif de schémas,
et par conséquent un recouvrement ph
(Topologies, lemme \ref{topologies-lemma-surjective-proper-ph}),
le résultat pour $Z \to X$ découle de Cohomologie étale, lemme
\ref{etale-cohomology-lemma-describe-pullback-pi-ph}
(il est en fait, en un certain sens, équivalent à ce lemme).
\end{proof}

\begin{lemma}
\label{spaces-more-cohomology-lemma-h0-proper-over-henselian-pair}
Soit $(A, I)$ un couple hensélien. Soit $X$ un espace algébrique sur $A$
tel que le morphisme structural $f : X \to \Spec(A)$ soit propre.
Soit $i : X_0 \to X$ l'inclusion de $X \times_{\Spec(A)} \Spec(A/I)$.
Pour tout faisceau $\mathcal{F}$ sur $X_\etale$, on a
$\Gamma(X, \mathcal{F}) = \Gamma(X_0, i^{-1}\mathcal{F})$.
\end{lemma}

\begin{proof}
Choisissons un morphisme propre surjectif $Y \to X$, où $Y$ est un schéma; voir
Cohomologie des espaces, lemme \ref{spaces-cohomology-lemma-weak-chow}.
Considérons le diagramme
$$
\xymatrix{
\Gamma(X_0, \mathcal{F}_0) \ar[r] &
\Gamma(Y_0, \mathcal{G}_0) \ar@<1ex>[r] \ar@<-1ex>[r] &
\Gamma((Y \times_X Y)_0, \mathcal{H}_0) \\
\Gamma(X, \mathcal{F}) \ar[r] \ar[u] &
\Gamma(Y, \mathcal{G}) \ar@<1ex>[r] \ar@<-1ex>[r] \ar[u] &
\Gamma(Y \times_X Y, \mathcal{H}) \ar[u]
}
$$
Ici $\mathcal{G}$, resp.\ $\mathcal{H}$, est l'image réciproque de
$\mathcal{F}$ sur $Y$, resp.\ sur $Y \times_X Y$, et l'indice $0$
indique le changement de base à $\Spec(A/I)$. D'après le cas des schémas
(Cohomologie étale, lemme
\ref{etale-cohomology-lemma-h0-proper-over-henselian-pair})
nous voyons que les flèches verticales du milieu et de droite sont bijectives.
Le lemme \ref{spaces-more-cohomology-lemma-surjective-proper} montre qu'il en va de même de celle de gauche.
\end{proof}

\begin{lemma}
\label{spaces-more-cohomology-lemma-h0-proper-over-henselian-local}
Soit $A$ un anneau local hensélien. Soit $X$ un espace algébrique
sur $A$ tel que $f : X \to \Spec(A)$ soit un morphisme propre.
Soit $X_0 \subset X$ la fibre de $f$ au-dessus du point fermé.
Pour tout faisceau $\mathcal{F}$ sur $X_\etale$, on a
$\Gamma(X, \mathcal{F}) = \Gamma(X_0, \mathcal{F}|_{X_0})$.
\end{lemma}

\begin{proof}
C'est un cas particulier du lemme \ref{spaces-more-cohomology-lemma-h0-proper-over-henselian-pair}.
\end{proof}

\begin{lemma}
\label{spaces-more-cohomology-lemma-proper-base-change-f-star}
Soit $S$ un schéma. Soient $f : X \to Y$ et $g : Y' \to Y$
des morphismes d'espaces algébriques sur $S$. Supposons $f$ propre.
Posons $X' = Y' \times_Y X$, de projections $f' : X' \to Y'$ et $g' : X' \to X$.
Soit $\mathcal{F}$ un faisceau quelconque sur $X_\etale$. Alors
$g^{-1}f_*\mathcal{F} = f'_*(g')^{-1}\mathcal{F}$.
\end{lemma}

\begin{proof}
La question est locale pour la topologie étale sur $Y'$. Choisissons un schéma
$V$ et un morphisme étale surjectif $V \to Y$. Choisissons un schéma $V'$ et un
morphisme étale surjectif $V' \to V \times_Y Y'$. Nous pouvons alors remplacer
$Y'$ par $V'$ et $Y$ par $V$, donc supposer que $Y$ et $Y'$ sont des schémas.
Nous pouvons ensuite travailler localement pour la topologie de Zariski sur
$Y$ et $Y'$, et donc supposer que $Y$ et $Y'$ sont des schémas affines.

\medskip\noindent
Supposons que $Y$ et $Y'$ soient des schémas affines. Choisissons un morphisme
propre surjectif $h_1 : X_1 \to X$, où $X_1$ est un schéma; voir
Cohomologie des espaces, lemme \ref{spaces-cohomology-lemma-weak-chow}.
Posons $X_2 = X_1 \times_X X_1$ et notons
$h_2 : X_2 \to X$ le morphisme structural. Observons qu'il s'agit d'un schéma.
D'après le cas des schémas
(Cohomologie étale, lemme
\ref{etale-cohomology-lemma-proper-base-change-f-star})
nous savons que le lemme est vrai pour les diagrammes cartésiens
$$
\vcenter{
\xymatrix{
X'_1 \ar[r] \ar[d] & X_1 \ar[d] \\
Y' \ar[r] & Y
}
}
\quad\text{et}\quad
\vcenter{
\xymatrix{
X'_2 \ar[r] \ar[d] & X_2 \ar[d] \\
Y' \ar[r] & Y
}
}
$$
et les faisceaux $\mathcal{F}_i = (X_i \to X)^{-1}\mathcal{F}$.
D'après le lemme \ref{spaces-more-cohomology-lemma-surjective-proper}, nous avons une suite exacte
$0 \to \mathcal{F} \to h_{1, *}\mathcal{F}_1 \to h_{2, *}\mathcal{F}_2$
et de même pour $(g')^{-1}\mathcal{F}$, puisque
$X'_2 = X'_1 \times_{X'} X'_1$. Nous en concluons que le
lemme est vrai (certains détails sont omis).
\end{proof}

\noindent
Soit $S$ un schéma.
Soit $f : Y \to X$ un morphisme d'espaces algébriques sur $S$. Soit
$\overline{x} : \Spec(k) \to S$ un point géométrique. La fibre
de $f$ en $\overline{x}$ est l'espace algébrique
$Y_{\overline{x}} = \Spec(k) \times_{\overline{x}, X} Y$ sur $\Spec(k)$.
Si $\mathcal{F}$ est un faisceau sur $Y_\etale$, notons
$\mathcal{F}_{\overline{x}} = p^{-1}\mathcal{F}$
l'image réciproque de $\mathcal{F}$ sur $(Y_{\overline{x}})_\etale$.
Ici $p : Y_{\overline{x}} \to Y$ est la projection.
Dans la suite, nous considérerons l'ensemble
$\Gamma(Y_{\overline{x}}, \mathcal{F}_{\overline{x}})$.

\begin{lemma}
\label{spaces-more-cohomology-lemma-proper-pushforward-stalk}
Soit $S$ un schéma.
Soit $f : Y \to X$ un morphisme propre d'espaces algébriques sur $S$. Soit
$\overline{x} \to X$ un point géométrique.
Pour tout faisceau $\mathcal{F}$ sur $Y_\etale$,
l'application canonique
$$
(f_*\mathcal{F})_{\overline{x}} \longrightarrow
\Gamma(Y_{\overline{x}}, \mathcal{F}_{\overline{x}})
$$
est bijective.
\end{lemma}

\begin{proof}
C'est un cas particulier du lemme \ref{spaces-more-cohomology-lemma-proper-base-change-f-star}.
\end{proof}

\begin{theorem}
\label{spaces-more-cohomology-theorem-proper-base-change}
Soit $S$ un schéma. Soit
$$
\xymatrix{
X' \ar[r]_{g'} \ar[d]_{f'} & X \ar[d]^f \\
Y' \ar[r]^g & Y
}
$$
un carré cartésien d'espaces algébriques sur $S$.
Supposons $f$ propre.
Soit $\mathcal{F}$ un faisceau abélien de torsion sur $X_\etale$.
Alors le morphisme de changement de base
$$
g^{-1}Rf_*\mathcal{F} \longrightarrow Rf'_*(g')^{-1}\mathcal{F}
$$
est un isomorphisme.
\end{theorem}

\begin{proof}
Cette démonstration reprend certains des arguments de la démonstration du
théorème de changement de base propre pour les schémas. Voir
Cohomologie étale, section \ref{etale-cohomology-section-proper-base-change},
pour davantage de détails.

\medskip\noindent
L'énoncé est local pour la topologie étale sur $Y'$ et sur $Y$; nous pouvons
donc supposer que $Y$ et $Y'$ sont tous deux des schémas affines. Observons que
ceci démontre en particulier le théorème lorsque $f$ est représentable
(nous l'utiliserons plus bas).

\medskip\noindent
Pour tout $n \geq 1$, soit $\mathcal{F}[n]$ le sous-faisceau des sections
de $\mathcal{F}$ annulées par $n$. Alors
$\mathcal{F} = \colim \mathcal{F}[n]$. D'après 
Cohomologie des espaces, lemme \ref{spaces-cohomology-lemma-colimit-cohomology},
les foncteurs $g^{-1}R^pf_*$ et $R^pf'_*(g')^{-1}$ commutent
aux limites inductives filtrantes. Il suffit donc de démontrer le théorème
lorsque $\mathcal{F}$ est annulée par $n$.

\medskip\noindent
Soit $\mathcal{F} \to \mathcal{I}^\bullet$ une résolution par des
faisceaux injectifs de $\mathbf{Z}/n\mathbf{Z}$-modules.
Observons que
$g^{-1}f_*\mathcal{I}^\bullet = f'_*(g')^{-1}\mathcal{I}^\bullet$
d'après le lemme \ref{spaces-more-cohomology-lemma-proper-base-change-f-star}.
En appliquant le lemme d'acyclicité de Leray
(Catégories dérivées, lemme \ref{derived-lemma-leray-acyclicity}),
nous concluons qu'il suffit de démontrer
$R^pf'_*(g')^{-1}\mathcal{I}^m = 0$ pour $p > 0$ et $m \in \mathbf{Z}$.

\medskip\noindent
Choisissons un morphisme propre surjectif
$h : Z \to X$, où $Z$ est un schéma; voir
Cohomologie des espaces, lemme \ref{spaces-cohomology-lemma-weak-chow}.
Choisissons une application injective $h^{-1}\mathcal{I}^m \to \mathcal{J}$,
où $\mathcal{J}$ est un faisceau injectif de
$\mathbf{Z}/n\mathbf{Z}$-modules sur $Z_\etale$.
Comme $h$ est surjectif, l'application $\mathcal{I}^m \to h_*\mathcal{J}$
est injective (voir le lemme \ref{spaces-more-cohomology-lemma-surjective-proper}).
Puisque $\mathcal{I}^m$ est injectif, nous voyons que $\mathcal{I}^m$
est un facteur direct de
$h_*\mathcal{J}$. Il suffit donc de démontrer l'annulation voulue
pour $h_*\mathcal{J}$.

\medskip\noindent
Notons $h'$ son changement de base par $g$ et
$g'' : Z' \to Z$ la projection. Il existe une suite spectrale
$$
E_2^{p, q} = R^pf'_* R^qh'_* (g'')^{-1}\mathcal{J}
$$
qui aboutit à $R^{p + q}(f' \circ h')_*(g'')^{-1}\mathcal{J}$.
Comme $h$ et $f \circ h$ sont représentables (par des schémas),
nous savons que le résultat voulu vaut pour eux. La suite spectrale montre donc
que $E_2^{p, q} = 0$ pour $q > 0$ et que
$R^{p + q}(f' \circ h')_*(g'')^{-1}\mathcal{J} = 0$
pour $p + q > 0$. Il s'ensuit que $E_2^{p, 0} = 0$ pour $p > 0$.
Or
$$
E_2^{p, 0} = R^pf'_* h'_* (g'')^{-1}\mathcal{J} =
R^pf'_* (g')^{-1}h_*\mathcal{J}
$$
d'après le lemme \ref{spaces-more-cohomology-lemma-proper-base-change-f-star}. Ceci achève la démonstration.
\end{proof}

\begin{lemma}
\label{spaces-more-cohomology-lemma-proper-base-change}
Soit $S$ un schéma. Soit
$$
\xymatrix{
X' \ar[r]_{g'} \ar[d]_{f'} & X \ar[d]^f \\
Y' \ar[r]^g & Y
}
$$
un carré cartésien d'espaces algébriques sur $S$. Supposons $f$ propre.
Soit $E \in D^+(X_\etale)$ à faisceaux de cohomologie de torsion.
Alors le morphisme de changement de base $g^{-1}Rf_*E \to Rf'_*(g')^{-1}E$
est un isomorphisme.
\end{lemma}

\begin{proof}
C'est une conséquence immédiate du théorème de changement de base propre
(théorème \ref{spaces-more-cohomology-theorem-proper-base-change}), en utilisant les suites
spectrales
$$
E_2^{p, q} = R^pf_*H^q(E)
\quad\text{et}\quad
{E'}_2^{p, q} = R^pf'_*(g')^{-1}H^q(E)
$$
qui aboutissent à $R^nf_*E$ et $R^nf'_*(g')^{-1}E$.
Ces suites spectrales sont construites dans Catégories dérivées, lemme
\ref{derived-lemma-two-ss-complex-functor}.
Certains détails sont omis.
\end{proof}

\begin{lemma}
\label{spaces-more-cohomology-lemma-proper-base-change-stalk}
Soit $S$ un schéma.
Soit $f : X \to Y$ un morphisme propre d'espaces algébriques.
Soit $\overline{y} \to Y$ un point géométrique.
\begin{enumerate}
\item Pour un faisceau abélien de torsion $\mathcal{F}$ sur $X_\etale$, on a
$(R^nf_*\mathcal{F})_{\overline{y}} =
H^n_\etale(X_{\overline{y}}, \mathcal{F}_{\overline{y}})$.
\item Pour $E \in D^+(X_\etale)$ à faisceaux de cohomologie de torsion, on a
$(R^nf_*E)_{\overline{y}} = H^n_\etale(X_{\overline{y}}, E_{\overline{y}})$.
\end{enumerate}
\end{lemma}

\begin{proof}
Dans l'énoncé, $\mathcal{F}_{\overline{y}}$ désigne l'image réciproque
de $\mathcal{F}$ sur $X_{\overline{y}} = \overline{y} \times_Y X$.
Comme l'image réciproque par $\overline{y} \to Y$ donne la fibre de
$\mathcal{F}$, la première assertion du lemme est un cas particulier du
théorème \ref{spaces-more-cohomology-theorem-proper-base-change}. La seconde est un cas particulier
du lemme \ref{spaces-more-cohomology-lemma-proper-base-change}.
\end{proof}

\begin{lemma}
\label{spaces-more-cohomology-lemma-base-change-separably-closed}
Soit $k'/k$ une extension de corps séparablement clos.
Soit $X$ un espace algébrique propre sur $k$.
Soit $\mathcal{F}$ un faisceau abélien de torsion sur $X$.
Alors l'application $H^q_\etale(X, \mathcal{F}) \to
H^q_\etale(X_{k'}, \mathcal{F}|_{X_{k'}})$ est un isomorphisme
pour $q \geq 0$.
\end{lemma}

\begin{proof}
C'est un cas particulier du théorème \ref{spaces-more-cohomology-theorem-proper-base-change}.
\end{proof}










\section{Comparaison des grands et petits topos}
\label{spaces-more-cohomology-section-compare}

\noindent
Soient $S$ un schéma et $X$ un espace algébrique sur $S$.
Dans Topologies sur les espaces, lemme
\ref{spaces-topologies-lemma-at-the-bottom-etale}
nous avons introduit les morphismes de comparaison
$\pi_X : (\textit{Espaces}/X)_\etale \to X_{spaces, \etale}$ et
$i_X : \Sh(X_\etale) \to \Sh((\textit{Espaces}/X)_\etale)$
avec $\pi_X \circ i_X = \text{id}$ comme morphismes de topos et
$\pi_{X, *} = i_X^{-1}$.
Plus généralement, si $f : Y \to X$ est un objet de $(\textit{Espaces}/X)_\etale$,
il existe un morphisme
$i_f : \Sh(Y_\etale) \to \Sh((\textit{Espaces}/X)_\etale)$
tel que $f_{small} = \pi_X \circ i_f$; voir
Topologies sur les espaces, lemmes \ref{spaces-topologies-lemma-put-in-T-etale} et
\ref{spaces-topologies-lemma-morphism-big-small-etale}. Dans
Topologies sur les espaces, remarque
\ref{spaces-topologies-remark-change-topologies-ringed}
nous avons étendu ceux-ci en un morphisme de sites annelés
$$
\pi_X :
((\textit{Espaces}/X)_\etale, \mathcal{O})
\to
(X_{spaces, \etale}, \mathcal{O}_X)
$$
et en des morphismes de topos annelés
$$
i_X :
(\Sh(X_\etale), \mathcal{O}_X)
\to
(\Sh((\textit{Espaces}/X)_\etale), \mathcal{O})
$$
et
$$
i_f :
(\Sh(Y_\etale), \mathcal{O}_Y)
\to
(\Sh((\textit{Espaces}/X)_\etale, \mathcal{O}))
$$
Notons que la restriction $i_X^{-1} = \pi_{X, *}$ (voir
Topologies, définition \ref{topologies-definition-restriction-small-etale})
transforme $\mathcal{O}$ en $\mathcal{O}_X$.
De même, $i_f^{-1}$ transforme $\mathcal{O}$ en $\mathcal{O}_Y$.
Voir Topologies sur les espaces, remarque
\ref{spaces-topologies-remark-change-topologies-ringed}.
Ainsi $i_X^*\mathcal{F} = i_X^{-1}\mathcal{F}$ et
$i_f^*\mathcal{F} = i_f^{-1}\mathcal{F}$ pour tout $\mathcal{O}$-module
$\mathcal{F}$ sur $(\textit{Espaces}/X)_\etale$. En particulier, $i_X^*$ et
$i_f^*$ sont des foncteurs exacts. Le foncteur $i_X^*$ est souvent noté
$\mathcal{F} \mapsto \mathcal{F}|_{X_\etale}$ (ce qui n'entre pas
en conflit avec la notation de Topologies sur les espaces, définition
\ref{spaces-topologies-definition-restriction-small-etale}).

\begin{lemma}
\label{spaces-more-cohomology-lemma-describe-pullback}
Soit $S$ un schéma. Soit $X$ un espace algébrique sur $S$.
Soit $\mathcal{F}$ un faisceau sur $X_\etale$. Alors
$\pi_X^{-1}\mathcal{F}$ est donné par la règle
$$
(\pi_X^{-1}\mathcal{F})(Y) = \Gamma(Y_\etale, f_{small}^{-1}\mathcal{F})
$$
pour $f : Y \to X$ dans $(\textit{Espaces}/X)_\etale$.
De plus, $\pi_Y^{-1}\mathcal{F}$ satisfait à la condition de faisceau pour les
recouvrements lisses, syntomiques, fppf, fpqc et ph.
\end{lemma}

\begin{proof}
Comme l'image réciproque est transitive et que $f_{small} = \pi_X \circ i_f$
(voir ci-dessus), nous avons
$i_f^{-1} \pi_X^{-1}\mathcal{F} = f_{small}^{-1}\mathcal{F}$.
Ceci donne la description de $\pi_X^{-1}$ annoncée dans le lemme.

\medskip\noindent
Pour démontrer que $\pi_X^{-1}\mathcal{F}$ est un faisceau pour la topologie ph,
il suffit, d'après Topologies sur les espaces, lemme
\ref{spaces-topologies-lemma-characterize-sheaf}
de montrer que, pour tout morphisme propre surjectif
$V \to U$ d'espaces algébriques sur $X$, l'ensemble
$(\pi_X^{-1}\mathcal{F})(U)$ est l'égalisateur des deux applications
$(\pi_X^{-1}\mathcal{F})(V) \to (\pi_X^{-1}\mathcal{F})(V \times_U V)$.
C'est précisément le lemme \ref{spaces-more-cohomology-lemma-surjective-proper}.

\medskip\noindent
Le cas des recouvrements lisses, syntomiques et fppf se déduit de celui
des recouvrements ph par Topologies sur les espaces, lemme
\ref{spaces-topologies-lemma-zariski-etale-smooth-syntomic-fppf-ph}.

\medskip\noindent
Soit $\mathcal{U} = \{U_i \to U\}_{i \in I}$ un recouvrement fpqc d'espaces
algébriques sur $X$. Soient $s_i \in (\pi_X^{-1}\mathcal{F})(U_i)$ des sections
qui coïncident sur $U_i \times_U U_j$. Il faut démontrer qu'il existe un unique
$s \in (\pi_X^{-1}\mathcal{F})(U)$ qui se restreint en $s_i$ sur $U_i$.
Cas I: $U$ et les $U_i$ sont des schémas. Ce cas découle de
Cohomologie étale, lemme \ref{etale-cohomology-lemma-describe-pullback}.
Cas II: $U$ est un schéma. Choisissons des morphismes étales surjectifs
$T_i \to U_i$, où les $T_i$ sont des schémas. Alors
$\mathcal{T} = \{T_i \to U\}$ est un recouvrement fpqc par des schémas et,
d'après le cas I, le résultat vaut pour $\mathcal{T}$. Nous omettons de vérifier
que ceci entraîne le résultat pour $\mathcal{U}$.
Cas III: cas général. Soit $W \to U$ un morphisme étale surjectif, où $W$ est
un schéma. Alors $\mathcal{W} = \{U_i \times_U W \to W\}$ est un recouvrement
fpqc (par des espaces algébriques) du schéma $W$.
D'après le cas II, le résultat vaut pour $\mathcal{W}$. Nous omettons de vérifier
que ceci entraîne le résultat pour $\mathcal{U}$.
\end{proof}

\begin{lemma}
\label{spaces-more-cohomology-lemma-compare-injectives}
Soit $S$ un schéma.
Soit $Y \to X$ un morphisme de $(\textit{Espaces}/S)_\etale$.
\begin{enumerate}
\item Si $\mathcal{I}$ est injectif dans
$\textit{Ab}((\textit{Espaces}/X)_\etale)$, alors
\begin{enumerate}
\item $i_f^{-1}\mathcal{I}$ est injectif dans $\textit{Ab}(Y_\etale)$,
\item $\mathcal{I}|_{X_\etale}$ est injectif dans $\textit{Ab}(X_\etale)$,
\end{enumerate}
\item Si $\mathcal{I}^\bullet$ est un complexe K-injectif
dans $\textit{Ab}((\textit{Espaces}/X)_\etale)$, alors
\begin{enumerate}
\item $i_f^{-1}\mathcal{I}^\bullet$ est un complexe K-injectif dans
$\textit{Ab}(Y_\etale)$,
\item $\mathcal{I}^\bullet|_{X_\etale}$ est un complexe K-injectif dans
$\textit{Ab}(X_\etale)$,
\end{enumerate}
\end{enumerate}
Les énoncés correspondants ne valent pas pour les modules.
\end{lemma}

\begin{proof}
Les assertions (1)(b) et (2)(b) découlent formellement du fait que le foncteur
de restriction $\pi_{X, *} = i_X^{-1}$ est adjoint à droite du foncteur exact
$\pi_X^{-1}$; voir Homologie, lemme
\ref{homology-lemma-adjoint-preserve-injectives}, et Catégories dérivées,
lemme \ref{derived-lemma-adjoint-preserve-K-injectives}.

\medskip\noindent
Les assertions (1)(a) et (2)(a) peuvent se voir de deux façons.
Première démonstration: on utilise que $i_f^{-1}$ est adjoint à droite du
foncteur exact $i_{f, !}$. Ce foncteur est construit, pour les faisceaux
d'ensembles, dans Topologies, lemme \ref{topologies-lemma-put-in-T-etale},
et, pour les faisceaux abéliens, dans Modules sur les sites, lemme
\ref{sites-modules-lemma-g-shriek-adjoint}. Il est démontré dans Modules sur
les sites, lemme \ref{sites-modules-lemma-exactness-lower-shriek}, qu'il est exact.
Seconde démonstration: on utilise l'égalité $i_f = i_Y \circ f_{big}$,
établie dans Topologies, lemme
\ref{topologies-lemma-morphism-big-small-etale}. Comme $f_{big}$ est une
localisation, l'image réciproque par ce morphisme préserve les injectifs et les
K-injectifs; voir Cohomologie sur les sites, lemmes
\ref{sites-cohomology-lemma-cohomology-of-open} et
\ref{sites-cohomology-lemma-restrict-K-injective-to-open}.
On conclut alors en appliquant au foncteur $i_Y^{-1}$ les assertions (1)(b)
et (2)(b), déjà démontrées.

\medskip\noindent
Pour un contre-exemple dans le cas des modules, nous renvoyons à
Cohomologie étale, lemme \ref{etale-cohomology-lemma-compare-injectives}.
\end{proof}

\noindent
Soit $S$ un schéma.
Soit $f : Y \to X$ un morphisme d'espaces algébriques sur $S$.
Le diagramme commutatif de Topologies sur les espaces, lemme
\ref{spaces-topologies-lemma-morphism-big-small-etale} (3)
donne un diagramme commutatif de sites annelés
$$
\xymatrix{
(Y_{spaces, \etale}, \mathcal{O}_Y) \ar[d]_{f_{spaces, \etale}} &
((\textit{Espaces}/Y)_\etale, \mathcal{O}) \ar[d]^{f_{big}} \ar[l]^{\pi_Y} \\
(X_{spaces, \etale}, \mathcal{O}_X) &
((\textit{Espaces}/X)_\etale, \mathcal{O}) \ar[l]_{\pi_X}
}
$$
comme on le voit aisément en explicitant les définitions de
$f_{small}^\sharp$, $f_{big}^\sharp$, $\pi_X^\sharp$ et $\pi_Y^\sharp$.
En particulier, cela signifie que
\begin{equation}
\label{spaces-more-cohomology-equation-compare-big-small}
(f_{big, *}\mathcal{F})|_{X_\etale} =
f_{small, *}(\mathcal{F}|_{Y_\etale})
\end{equation}
pour tout faisceau $\mathcal{F}$ sur $(\textit{Espaces}/Y)_\etale$; si
$\mathcal{F}$ est un faisceau de $\mathcal{O}$-modules, alors
(\ref{spaces-more-cohomology-equation-compare-big-small})
est un isomorphisme de $\mathcal{O}_X$-modules sur $X_\etale$.

\begin{lemma}
\label{spaces-more-cohomology-lemma-compare-higher-direct-image}
Soit $S$ un schéma.
Soit $f : Y \to X$ un morphisme d'espaces algébriques sur $S$.
\begin{enumerate}
\item Pour $K$ dans $D((\textit{Espaces}/Y)_\etale)$, on a
$
(Rf_{big, *}K)|_{X_\etale} = Rf_{small, *}(K|_{Y_\etale})
$
dans $D(X_\etale)$.
\item Pour $K$ dans $D((\textit{Espaces}/Y)_\etale, \mathcal{O})$, on a
$
(Rf_{big, *}K)|_{X_\etale} = Rf_{small, *}(K|_{Y_\etale})
$
dans $D(\textit{Mod}(X_\etale, \mathcal{O}_X))$.
\end{enumerate}
Plus généralement, soit $g : X' \to X$ un objet de
$(\textit{Espaces}/X)_\etale$. Considérons le produit fibré
$$
\xymatrix{
Y' \ar[r]_{g'} \ar[d]_{f'} & Y \ar[d]^f \\
X' \ar[r]^g & X
}
$$
Alors
\begin{enumerate}
\item[(3)] Pour $K$ dans $D((\textit{Espaces}/Y)_\etale)$, on a
$i_g^{-1}(Rf_{big, *}K) = Rf'_{small, *}(i_{g'}^{-1}K)$
dans $D(X'_\etale)$.
\item[(4)] Pour $K$ dans $D((\textit{Espaces}/Y)_\etale, \mathcal{O})$, on a
$i_g^*(Rf_{big, *}K) = Rf'_{small, *}(i_{g'}^*K)$
dans $D(\textit{Mod}(X'_\etale, \mathcal{O}_{X'}))$.
\item[(5)] Pour $K$ dans $D((\textit{Espaces}/Y)_\etale)$, on a
$g_{big}^{-1}(Rf_{big, *}K) = Rf'_{big, *}((g'_{big})^{-1}K)$
dans $D((\textit{Espaces}/X')_\etale)$.
\item[(6)] Pour $K$ dans $D((\textit{Espaces}/Y)_\etale, \mathcal{O})$, on a
$g_{big}^*(Rf_{big, *}K) = Rf'_{big, *}((g'_{big})^*K)$
dans $D(\textit{Mod}(X'_\etale, \mathcal{O}_{X'}))$.
\end{enumerate}
\end{lemma}

\begin{proof}
L'assertion (1) découle du lemme \ref{spaces-more-cohomology-lemma-compare-injectives}
et de (\ref{spaces-more-cohomology-equation-compare-big-small}),
en choisissant un complexe K-injectif de faisceaux abéliens représentant $K$.

\medskip\noindent
L'assertion (3) découle du lemme \ref{spaces-more-cohomology-lemma-compare-injectives}
et de Topologies, lemme
\ref{topologies-lemma-morphism-big-small-cartesian-diagram-etale}
en choisissant un complexe K-injectif de faisceaux abéliens représentant $K$.

\medskip\noindent
L'assertion (5) est Cohomologie sur les sites, lemme
\ref{sites-cohomology-lemma-localize-cartesian-square}.

\medskip\noindent
L'assertion (6) est Cohomologie sur les sites, lemme
\ref{sites-cohomology-lemma-localize-cartesian-square-modules}.

\medskip\noindent
L'assertion (2) se démontre comme suit. Nous avons vu ci-dessus que
$\pi_X \circ f_{big} = f_{small} \circ \pi_Y$ comme morphismes de sites
annelés. Nous obtenons donc
$R\pi_{X, *} \circ Rf_{big, *} = Rf_{small, *} \circ R\pi_{Y, *}$
d'après Cohomologie sur les sites, lemme
\ref{sites-cohomology-lemma-derived-pushforward-composition}.
Comme les foncteurs de restriction $\pi_{X, *}$ et $\pi_{Y, *}$
sont exacts, on conclut.

\medskip\noindent
L'assertion (4) découle des assertions (6) et (2), cette dernière étant
appliquée à $f' : Y' \to X'$.
\end{proof}

\noindent
Soit $S$ un schéma. Soit $X$ un espace algébrique sur $S$.
Soit $\mathcal{H}$ un faisceau abélien sur
$(\textit{Espaces}/X)_\etale$. Rappelons que $H^n_\etale(U, \mathcal{H})$
désigne la cohomologie de $\mathcal{H}$ sur un objet
$U$ de $(\textit{Espaces}/X)_\etale$.

\begin{lemma}
\label{spaces-more-cohomology-lemma-compare-cohomology}
Soit $S$ un schéma.
Soit $f : Y \to X$ un morphisme d'espaces algébriques sur $S$. Alors
\begin{enumerate}
\item Pour $K$ dans $D(X_\etale)$, on a
$H^n_\etale(X, \pi_X^{-1}K) = H^n(X_\etale, K)$.
\item Pour $K$ dans $D(X_\etale, \mathcal{O}_X)$, on a
$H^n_\etale(X, L\pi_X^*K) = H^n(X_\etale, K)$.
\item Pour $K$ dans $D(X_\etale)$, on a
$H^n_\etale(Y, \pi_X^{-1}K) = H^n(Y_\etale, f_{small}^{-1}K)$.
\item Pour $K$ dans $D(X_\etale, \mathcal{O}_X)$, on a
$H^n_\etale(Y, L\pi_X^*K) = H^n(Y_\etale, Lf_{small}^*K)$.
\item Pour $M$ dans $D((\textit{Espaces}/X)_\etale)$, on a
$H^n_\etale(Y, M) = H^n(Y_\etale, i_f^{-1}M)$.
\item Pour $M$ dans $D((\textit{Espaces}/X)_\etale, \mathcal{O})$, on a
$H^n_\etale(Y, M) = H^n(Y_\etale, i_f^*M)$.
\end{enumerate}
\end{lemma}

\begin{proof}
Pour démontrer (5), représentons $M$ par un complexe K-injectif de faisceaux
abéliens, appliquons le lemme \ref{spaces-more-cohomology-lemma-compare-injectives} et explicitons
les définitions. L'assertion (3) en découle puisque
$i_f^{-1}\pi_X^{-1} = f_{small}^{-1}$. L'assertion (1) est un cas
particulier de (3).

\medskip\noindent
L'assertion (6) découle du très général lemme
\ref{sites-cohomology-lemma-pullback-same-cohomology} de Cohomologie sur les
sites. Ensuite, l'assertion (4) découle de l'égalité
$Lf_{small}^* = i_f^* \circ L\pi_X^*$. L'assertion (2) est un cas particulier
de (4).
\end{proof}

\begin{lemma}
\label{spaces-more-cohomology-lemma-cohomological-descent-etale}
Soit $S$ un schéma. Soit $X$ un espace algébrique sur $S$.
Pour $K \in D(X_\etale)$, l'application
$$
K \longrightarrow R\pi_{X, *}\pi_X^{-1}K
$$
est un isomorphisme, où
$\pi_X : \Sh((\textit{Espaces}/X)_\etale) \to \Sh(X_\etale)$ est défini ci-dessus.
\end{lemma}

\begin{proof}
Ceci est vrai parce que $\pi_X^{-1}$ et $\pi_{X, *} = i_X^{-1}$ sont tous deux
des foncteurs exacts et que le composé $\pi_{X, *} \circ \pi_X^{-1}$
est le foncteur identité.
\end{proof}

\begin{lemma}
\label{spaces-more-cohomology-lemma-compare-higher-direct-image-proper}
Soit $S$ un schéma.
Soit $f : Y \to X$ un morphisme propre d'espaces algébriques sur $S$.
Alors on a
\begin{enumerate}
\item $\pi_X^{-1} \circ f_{small, *} = f_{big, *} \circ \pi_Y^{-1}$
comme foncteurs $\Sh(Y_\etale) \to \Sh((\textit{Espaces}/X)_\etale)$,
\item $\pi_X^{-1}Rf_{small, *}K = Rf_{big, *}\pi_Y^{-1}K$
pour $K$ dans $D^+(Y_\etale)$ dont les faisceaux de cohomologie sont de torsion, et
\item $\pi_X^{-1}Rf_{small, *}K = Rf_{big, *}\pi_Y^{-1}K$
pour tout $K$ dans $D(Y_\etale)$ si $f$ est fini.
\end{enumerate}
\end{lemma}

\begin{proof}
Démonstration de (1). Soit $\mathcal{F}$ un faisceau sur $Y_\etale$.
Soit $g : X' \to X$ un objet de $(\textit{Espaces}/X)_\etale$.
Considérons le produit fibré
$$
\xymatrix{
Y' \ar[r]_{f'} \ar[d]_{g'} & X' \ar[d]^g \\
Y \ar[r]^f & X
}
$$
On a alors
$$
(f_{big, *}\pi_Y^{-1}\mathcal{F})(X') =
(\pi_Y^{-1}\mathcal{F})(Y') =
((g'_{small})^{-1}\mathcal{F})(Y')  =
(f'_{small, *}(g'_{small})^{-1}\mathcal{F})(X')
$$
la seconde égalité résultant du lemme \ref{spaces-more-cohomology-lemma-describe-pullback}.
D'autre part,
$$
(\pi_X^{-1}f_{small, *}\mathcal{F})(X') =
(g_{small}^{-1}f_{small, *}\mathcal{F})(X')
$$
de nouveau d'après le lemme \ref{spaces-more-cohomology-lemma-describe-pullback}.
Le changement de base propre pour les faisceaux d'ensembles
(lemme \ref{spaces-more-cohomology-lemma-proper-base-change-f-star}) montre donc que les deux
ensembles sont canoniquement isomorphes. Cet isomorphisme est compatible
avec les applications de restriction et définit un isomorphisme
$\pi_X^{-1}f_{small, *}\mathcal{F} = f_{big, *}\pi_Y^{-1}\mathcal{F}$.
On obtient ainsi un isomorphisme de foncteurs
$\pi_X^{-1} \circ f_{small, *} = f_{big, *} \circ \pi_Y^{-1}$.

\medskip\noindent
Démonstration de (2). Il existe un morphisme canonique de changement de base
$\pi_X^{-1}Rf_{small, *}K \to Rf_{big, *}\pi_Y^{-1}K$
pour tout $K$ dans $D(Y_\etale)$; voir Cohomologie sur les sites, remarque
\ref{sites-cohomology-remark-base-change}.
Pour démontrer que c'est un isomorphisme, il suffit de montrer que l'image
réciproque du morphisme de changement de base par
$i_g : \Sh(X'_\etale) \to \Sh((\Sch/X)_\etale)$ est un isomorphisme pour tout
objet $g : X' \to X$ de $(\Sch/X)_\etale$.
Soient $T', g', f'$ comme au paragraphe précédent.
L'image réciproque du morphisme de changement de base est
\begin{align*}
g_{small}^{-1}Rf_{small, *}K
& =
i_g^{-1}\pi_X^{-1}Rf_{small, *}K \\
& \to
i_g^{-1}Rf_{big, *}\pi_Y^{-1}K \\
& =
Rf'_{small, *}(i_{g'}^{-1}\pi_Y^{-1}K) \\
& =
Rf'_{small, *}((g'_{small})^{-1}K)
\end{align*}
où nous avons utilisé $\pi_X \circ i_g = g_{small}$,
$\pi_Y \circ i_{g'} = g'_{small}$ et le lemme
\ref{spaces-more-cohomology-lemma-compare-higher-direct-image}.
Cette application est un isomorphisme d'après le théorème de changement de base
propre (lemme \ref{spaces-more-cohomology-lemma-proper-base-change}), pourvu que $K$ soit borné
inférieurement et que les faisceaux de cohomologie de $K$ soient de torsion.

\medskip\noindent
Démonstration de (3). Si $f$ est fini, les foncteurs
$f_{small, *}$ et $f_{big, *}$ sont exacts. Pour $f_{small}$, ceci découle de
Cohomologie des espaces, lemme
\ref{spaces-cohomology-lemma-finite-higher-direct-image-zero}.
Comme tout changement de base $f'$ de $f$ est encore fini, le lemme
\ref{spaces-more-cohomology-lemma-compare-higher-direct-image}, assertion (3), montre que
$f_{big, *}$ est lui aussi exact (puisque les foncteurs dérivés supérieurs
sont nuls). Ce cas découle donc de l'assertion (1).
\end{proof}







\section{Comparaison des topologies fppf et étale}
\label{spaces-more-cohomology-section-fppf-etale}

\noindent
Cette section est l'analogue de Cohomologie étale, section
\ref{etale-cohomology-section-fppf-etale}.

\medskip\noindent
Soient $S$ un schéma et $X$ un espace algébrique sur $S$.
Sur la catégorie $\textit{Espaces}/X$, considérons les topologies fppf
et étale. Le foncteur identité
$(\textit{Espaces}/X)_\etale \to (\textit{Espaces}/X)_{fppf}$
est continu et définit un morphisme de sites
$$
\epsilon_X :
(\textit{Espaces}/X)_{fppf} \longrightarrow (\textit{Espaces}/X)_\etale
$$
par application de Sites, proposition \ref{sites-proposition-get-morphism}.
Notons que $\epsilon_{X, *}$ est le foncteur identité sur les préfaisceaux
sous-jacents et que $\epsilon_X^{-1}$ associe à un faisceau étale son
faisceau associé pour la topologie fppf.
Considérons le morphisme de sites
$$
\pi_X : (\textit{Espaces}/X)_\etale \longrightarrow X_{spaces, \etale}
$$
qui compare les grands et petits sites étales; voir la section \ref{spaces-more-cohomology-section-compare}.
Le composé détermine un morphisme de sites
$$
a_X = \pi_X \circ \epsilon_X :
(\textit{Espaces}/X)_{fppf}
\longrightarrow
X_{spaces, \etale}
$$
Si $\mathcal{H}$ est un faisceau abélien sur $(\textit{Espaces}/X)_{fppf}$,
nous noterons $H^n_{fppf}(U, \mathcal{H})$ la cohomologie de
$\mathcal{H}$ sur un objet $U$ de $(\textit{Espaces}/X)_{fppf}$.

\begin{lemma}
\label{spaces-more-cohomology-lemma-comparison-fppf-etale}
Soit $S$ un schéma. Soit $X$ un espace algébrique sur $S$.
\begin{enumerate}
\item Pour $\mathcal{F} \in \Sh(X_\etale)$, on a
$\epsilon_{X, *}a_X^{-1}\mathcal{F} = \pi_X^{-1}\mathcal{F}$
et $a_{X, *}a_X^{-1}\mathcal{F} = \mathcal{F}$.
\item Pour $\mathcal{F} \in \textit{Ab}(X_\etale)$, on a
$R^i\epsilon_{X, *}(a_X^{-1}\mathcal{F}) = 0$ pour $i > 0$.
\end{enumerate}
\end{lemma}

\begin{proof}
On a $a_X^{-1}\mathcal{F} = \epsilon_X^{-1} \pi_X^{-1}\mathcal{F}$.
D'après le lemme \ref{spaces-more-cohomology-lemma-describe-pullback}, le faisceau étale
$\pi_X^{-1}\mathcal{F}$ est un faisceau pour la topologie fppf; il est donc
égal à $a_X^{-1}\mathcal{F}$ (puisque l'image réciproque par $\epsilon_X$
est donnée par le faisceau associé pour la topologie fppf).
Rappelons en outre que $\epsilon_{X, *}$ est le foncteur identité
sur les préfaisceaux sous-jacents.
L'assertion (1) résulte alors immédiatement de la description explicite de
$\pi_X^{-1}$ donnée au lemme \ref{spaces-more-cohomology-lemma-describe-pullback}.

\medskip\noindent
Nous démontrerons (2) en nous ramenant au cas des schémas; voir l'assertion (1)
de Cohomologie étale, lemme
\ref{etale-cohomology-lemma-V-C-all-n-etale-fppf}.
Cette méthode « fonctionne clairement », puisque tout espace algébrique est
localement un schéma pour la topologie étale. Les détails sont donnés ci-dessous,
mais nous invitons vivement le lecteur à omettre la démonstration.

\medskip\noindent
Pour un faisceau abélien $\mathcal{H}$ sur $(\textit{Espaces}/X)_{fppf}$,
l'image directe supérieure $R^p\epsilon_{X, *}\mathcal{H}$ est le faisceau
associé au préfaisceau $U \mapsto H^p_{fppf}(U, \mathcal{H})$
sur $(\textit{Espaces}/X)_\etale$. Voir Cohomologie sur les sites, lemme
\ref{sites-cohomology-lemma-higher-direct-images}.
Comme tout objet de $(\textit{Espaces}/X)_\etale$ admet un recouvrement par des
schémas, il suffit de démontrer que, pour tout schéma $U/X$ et tout
$\xi \in H^p_{fppf}(U, a_X^{-1}\mathcal{F})$, il existe un recouvrement étale
$\{U_i \to U\}$ tel que la restriction de $\xi$ à chaque $U_i$ soit nulle.
On a
\begin{align*}
H^p_{fppf}(U, a_X^{-1}\mathcal{F})
& =
H^p((\textit{Espaces}/U)_{fppf}, (a_X^{-1}\mathcal{F})|_{\textit{Espaces}/U}) \\
& =
H^p((\Sch/U)_{fppf}, (a_X^{-1}\mathcal{F})|_{\Sch/U})
\end{align*}
où la seconde identification est le lemme
\ref{spaces-more-cohomology-lemma-compare-cohomology-other-topologies}, et la première un fait général
sur la restriction (Cohomologie sur les sites, lemme
\ref{sites-cohomology-lemma-cohomology-of-open}).
En comparant le premier paragraphe et le résultat correspondant dans le cas
des schémas (Cohomologie étale, lemme
\ref{etale-cohomology-lemma-describe-pullback-pi-fppf})
nous concluons que le faisceau $(a_X^{-1}\mathcal{F})|_{\Sch/U}$
coïncide avec l'image réciproque par la « version pour les schémas de $a_U$ ».
Nous pouvons donc trouver un recouvrement étale $\{U_i \to U\}$ tel que notre
classe s'annule dans
$H^p((\Sch/U_i)_{fppf}, (a_X^{-1}\mathcal{F})|_{\Sch/U_i})$
pour tout $i$; voir Cohomologie étale, lemme
\ref{etale-cohomology-lemma-V-C-all-n-etale-fppf}
(l'énoncé précis à utiliser ici est que $V_n$ vaut pour tout $n$, ce qui est
l'assertion (2) dans le cas des schémas).
En revenant en arrière (au moyen des mêmes formules que ci-dessus, appliquées
cette fois à $U_i$), nous concluons que la restriction de $\xi$ à $U_i$
est nulle, comme souhaité.
\end{proof}

\noindent
Le travail difficile effectué dans le cas des schémas montre maintenant que
les cohomologies étale et fppf coïncident pour les faisceaux provenant du
petit site étale.

\begin{lemma}
\label{spaces-more-cohomology-lemma-cohomological-descent-etale-fppf}
Soit $S$ un schéma. Soit $X$ un espace algébrique sur $S$.
Pour $K \in D^+(X_\etale)$, les applications
$$
\pi_X^{-1}K \longrightarrow R\epsilon_{X, *}a_X^{-1}K
\quad\text{et}\quad
K \longrightarrow Ra_{X, *}a_X^{-1}K
$$
sont des isomorphismes, où
$a_X : \Sh((\textit{Espaces}/X)_{fppf}) \to \Sh(X_\etale)$ est défini ci-dessus.
\end{lemma}

\begin{proof}
Nous démontrons seulement le second énoncé; le premier est plus facile et se
démontre exactement de la même manière. On se ramène immédiatement au cas où
$K$ est donné par un seul faisceau abélien. En effet, représentons $K$ par un
complexe borné inférieurement $\mathcal{F}^\bullet$. Le cas d'un faisceau donne
$\mathcal{F}^n = a_{X, *} a_X^{-1} \mathcal{F}^n$
et les faisceaux $R^qa_{X, *}a_X^{-1}\mathcal{F}^n$ sont nuls pour $q > 0$.
Le lemme d'acyclicité de Leray (Catégories dérivées, lemme
\ref{derived-lemma-leray-acyclicity}), appliqué à
$a_X^{-1}\mathcal{F}^\bullet$ et au foncteur $a_{X, *}$, permet de conclure.
Supposons désormais $K = \mathcal{F}$.

\medskip\noindent
Le lemme \ref{spaces-more-cohomology-lemma-comparison-fppf-etale} donne
$a_{X, *}a_X^{-1}\mathcal{F} = \mathcal{F}$. Il suffit donc de montrer que
$R^qa_{X, *}a_X^{-1}\mathcal{F} = 0$ pour $q > 0$.
Pour cela, nous pouvons utiliser $a_X = \epsilon_X \circ \pi_X$ et la suite
spectrale de Leray (Cohomologie sur les sites, lemme
\ref{sites-cohomology-lemma-relative-Leray}).
D'après le lemme \ref{spaces-more-cohomology-lemma-comparison-fppf-etale}, on a
$R^i\epsilon_{X, *}(a_X^{-1}\mathcal{F}) = 0$ pour $i > 0$.
On a
$\epsilon_{X, *}a_X^{-1}\mathcal{F} = \pi_X^{-1}\mathcal{F}$
et, d'après le lemme \ref{spaces-more-cohomology-lemma-cohomological-descent-etale},
$R^j\pi_{X, *}(\pi_X^{-1}\mathcal{F}) = 0$ pour $j > 0$.
Ceci achève la démonstration.
\end{proof}

\begin{lemma}
\label{spaces-more-cohomology-lemma-compare-cohomology-etale-fppf}
Soient $S$ un schéma et $X$ un espace algébrique sur $S$.
Avec $a_X : \Sh((\textit{Espaces}/X)_{fppf}) \to \Sh(X_\etale)$
défini ci-dessus:
\begin{enumerate}
\item $H^q(X_\etale, \mathcal{F}) = H^q_{fppf}(X, a_X^{-1}\mathcal{F})$
pour un faisceau abélien $\mathcal{F}$ sur $X_\etale$,
\item $H^q(X_\etale, K) = H^q_{fppf}(X, a_X^{-1}K)$ pour $K \in D^+(X_\etale)$.
\end{enumerate}
Exemple: si $A$ est un groupe abélien, alors
$H^q_\etale(X, \underline{A}) = H^q_{fppf}(X, \underline{A})$.
\end{lemma}

\begin{proof}
Ceci découle du lemme \ref{spaces-more-cohomology-lemma-cohomological-descent-etale-fppf}, d'après
Cohomologie sur les sites, remarque \ref{sites-cohomology-remark-before-Leray}.
\end{proof}

\begin{lemma}
\label{spaces-more-cohomology-lemma-push-pull-fppf-etale}
Soit $S$ un schéma.
Soit $f : X \to Y$ un morphisme d'espaces algébriques sur $S$.
Il existe alors des diagrammes commutatifs de topos
$$
\xymatrix{
\Sh((\textit{Espaces}/X)_{fppf}) \ar[rr]_{f_{big, fppf}} \ar[d]_{\epsilon_X} & &
\Sh((\textit{Espaces}/Y)_{fppf}) \ar[d]^{\epsilon_Y} \\
\Sh((\textit{Espaces}/X)_\etale) \ar[rr]^{f_{big, \etale}} & &
\Sh((\textit{Espaces}/Y)_\etale)
}
$$
et
$$
\xymatrix{
\Sh((\textit{Espaces}/X)_{fppf}) \ar[rr]_{f_{big, fppf}} \ar[d]_{a_X} & &
\Sh((\textit{Espaces}/Y)_{fppf}) \ar[d]^{a_Y} \\
\Sh(X_\etale) \ar[rr]^{f_{small}} & &
\Sh(Y_\etale)
}
$$
avec $a_X = \pi_X \circ \epsilon_X$ et $a_Y = \pi_X \circ \epsilon_X$.
\end{lemma}

\begin{proof}
Ceci découle immédiatement des définitions des morphismes en jeu; voir
Topologies sur les espaces, section \ref{spaces-topologies-section-fppf},
et la section \ref{spaces-more-cohomology-section-compare}.
\end{proof}

\begin{lemma}
\label{spaces-more-cohomology-lemma-proper-push-pull-fppf-etale}
Dans le lemme \ref{spaces-more-cohomology-lemma-push-pull-fppf-etale}, si $f$ est propre, alors on a
\begin{enumerate}
\item $a_Y^{-1} \circ f_{small, *} = f_{big, fppf, *} \circ a_X^{-1}$, et
\item
$a_Y^{-1}(Rf_{small, *}K) = Rf_{big, fppf, *}(a_X^{-1}K)$
pour $K$ dans $D^+(X_\etale)$ à faisceaux de cohomologie de torsion.
\end{enumerate}
\end{lemma}

\begin{proof}
Démonstration de (1). On pourrait la donner en répétant celle du lemme
\ref{spaces-more-cohomology-lemma-compare-higher-direct-image-proper}, assertion (1); nous allons
plutôt en déduire le résultat. Comme $\epsilon_{Y, *}$ est le foncteur identité
sur les préfaisceaux sous-jacents, il reflète les isomorphismes. Le lemme
\ref{spaces-more-cohomology-lemma-comparison-fppf-etale} montre que
$\epsilon_{Y, *} \circ a_Y^{-1} = \pi_Y^{-1}$, et de même pour $X$.
Pour montrer que l'application canonique
$a_Y^{-1}f_{small, *}\mathcal{F} \to f_{big, fppf, *}a_X^{-1}\mathcal{F}$
est un isomorphisme, il suffit de montrer que
\begin{align*}
\pi_Y^{-1}f_{small, *}\mathcal{F}
& =
\epsilon_{Y, *}a_Y^{-1}f_{small, *}\mathcal{F} \\
& \to 
\epsilon_{Y, *}f_{big, fppf, *}a_X^{-1}\mathcal{F} \\
& =
f_{big, \etale, *} \epsilon_{X, *}a_X^{-1}\mathcal{F} \\
& =
f_{big, \etale, *}\pi_X^{-1}\mathcal{F}
\end{align*}
est un isomorphisme. C'est l'assertion (1) du lemme
\ref{spaces-more-cohomology-lemma-compare-higher-direct-image-proper}.

\medskip\noindent
Pour démontrer (2), nous utilisons
\begin{align*}
R\epsilon_{Y, *}Rf_{big, fppf, *}a_X^{-1}K
& =
Rf_{big, \etale, *}R\epsilon_{X, *}a_X^{-1}K \\
& =
Rf_{big, \etale, *}\pi_X^{-1}K \\
& =
\pi_Y^{-1}Rf_{small, *}K \\
& =
R\epsilon_{Y, *} a_Y^{-1}Rf_{small, *}K
\end{align*}
La première égalité résulte du diagramme commutatif du lemme
\ref{spaces-more-cohomology-lemma-push-pull-fppf-etale} et de Cohomologie sur les sites, lemme
\ref{sites-cohomology-lemma-derived-pushforward-composition}.
La deuxième est le lemme \ref{spaces-more-cohomology-lemma-cohomological-descent-etale-fppf}; la
troisième est l'assertion (2) du lemme
\ref{spaces-more-cohomology-lemma-compare-higher-direct-image-proper}; la quatrième est de nouveau
le lemme \ref{spaces-more-cohomology-lemma-cohomological-descent-etale-fppf}.
Ainsi le morphisme de changement de base
$a_Y^{-1}(Rf_{small, *}K) \to Rf_{big, fppf, *}(a_X^{-1}K)$
induit un isomorphisme
$$
R\epsilon_{Y, *}a_Y^{-1}Rf_{small, *}K \to
R\epsilon_{Y, *}Rf_{big, fppf, *}a_X^{-1}K
$$
La remarque suivante achève la démonstration: soit
$\alpha : a_Y^{-1}L \to M$, avec $L$ dans $D^+(Y_\etale)$ et $M$ dans
$D^+((\textit{Espaces}/Y)_{fppf})$, telle que $R\epsilon_{Y, *}\alpha$ soit un
isomorphisme. Alors cette application est un isomorphisme. En effet, montrons par
récurrence sur $i$ que $H^i(\alpha)$ est un isomorphisme. Ceci est vrai pour
tout $i$ suffisamment petit. Si cela vaut pour $i \leq i_0$, alors
$R^j\epsilon_{Y, *}H^i(M) = 0$ pour $j > 0$ et $i \leq i_0$ d'après le lemme
\ref{spaces-more-cohomology-lemma-comparison-fppf-etale}, puisque
$H^i(M) = a_Y^{-1}H^i(L)$ dans cet intervalle. Un argument de suite spectrale
donne alors
$\epsilon_{Y, *}H^{i_0 + 1}(M) = H^{i_0 + 1}(R\epsilon_{Y, *}M)$.
Par conséquent, $\epsilon_{Y, *}H^{i_0 + 1}(M) =
\pi_Y^{-1}H^{i_0 + 1}(L) =
\epsilon_{Y, *}a_Y^{-1}H^{i_0 + 1}(L)$. Il s'ensuit que
$H^{i_0 + 1}(\alpha)$ est un isomorphisme (car $\epsilon_{Y, *}$ reflète les
isomorphismes, puisqu'il est l'identité sur les préfaisceaux sous-jacents),
comme voulu.
\end{proof}

\begin{lemma}
\label{spaces-more-cohomology-lemma-finite-push-pull-fppf-etale}
Dans le lemme \ref{spaces-more-cohomology-lemma-push-pull-fppf-etale}, si $f$ est fini, alors
$a_Y^{-1}(Rf_{small, *}K) = Rf_{big, fppf, *}(a_X^{-1}K)$
pour $K$ dans $D^+(X_\etale)$.
\end{lemma}

\begin{proof}
Soit $V \to Y$ un morphisme étale surjectif, où $V$ est un schéma.
Il suffit de montrer que le morphisme de changement de base est un isomorphisme
après restriction à $V$. Nous pouvons donc supposer que $Y$ est un schéma.
Comme le morphisme est fini, donc représentable, nous pouvons alors supposer
que $X$ et $Y$ sont tous deux des schémas. Dans ce cas, le résultat découle du
cas des schémas (Cohomologie étale, lemme
\ref{etale-cohomology-lemma-V-C-all-n-etale-fppf}, assertion (2)), au moyen de
la comparaison des topos exposée à la section \ref{spaces-more-cohomology-section-api} et, en
particulier, dans le lemme \ref{spaces-more-cohomology-lemma-compare-cohomology-other-topologies}.
Certains détails sont omis.
\end{proof}

\begin{lemma}
\label{spaces-more-cohomology-lemma-descent-sheaf-fppf-etale}
Dans le lemme \ref{spaces-more-cohomology-lemma-push-pull-fppf-etale}, supposons que
$f$ soit plat, localement de présentation finie et surjectif.
Alors le foncteur
$$
\Sh(Y_\etale) \longrightarrow
\left\{
(\mathcal{G}, \mathcal{H}, \alpha)
\middle|
\begin{matrix}
\mathcal{G} \in \Sh(X_\etale),\ \mathcal{H} \in \Sh((\Sch/Y)_{fppf}), \\
\alpha : a_X^{-1}\mathcal{G} \to f_{big, fppf}^{-1}\mathcal{H}
\text{ un isomorphisme}
\end{matrix}
\right\}
$$
qui envoie $\mathcal{F}$ sur
$(f_{small}^{-1}\mathcal{F}, a_Y^{-1}\mathcal{F}, can)$ est une équivalence.
\end{lemma}

\begin{proof}
Le foncteur $a_X^{-1}$ est pleinement fidèle (car
$a_{X, *}a_X^{-1} = \text{id}$ d'après le lemme
\ref{spaces-more-cohomology-lemma-comparison-fppf-etale}). Le foncteur d'oubli
$(\mathcal{G}, \mathcal{H}, \alpha) \mapsto \mathcal{H}$ identifie donc la
catégorie des triplets à une sous-catégorie pleine de $\Sh((\Sch/Y)_{fppf})$.
De plus, le foncteur $a_Y^{-1}$ est pleinement fidèle; le foncteur du lemme
l'est donc lui aussi.

\medskip\noindent
Supposons donné un recouvrement étale $\{Y_i \to Y\}$.
Soit $f_i : X_i \to Y_i$ le changement de base de $f$.
Notons $f_{ij} = f_i \times f_j : X_i \times_X X_j \to Y_i \times_Y Y_j$.
Affirmation: si le lemme est vrai pour $f_i$ et $f_{ij}$ pour tous $i,j$, il
est vrai pour $f$. En effet, le recouvrement étale donné détermine un
recouvrement étale de l'objet final de chacun des quatre sites
$Y_\etale, X_\etale, (\Sch/Y)_{fppf}, (\Sch/X)_{fppf}$.
Dans chacun des quatre cas, la catégorie des faisceaux est donc équivalente à
la catégorie des données de recollement pour ce recouvrement (Sites, lemme
\ref{sites-lemma-mapping-property-glue}). Un immense diagramme commutatif de
catégories achève alors la démonstration de l'affirmation. Nous omettons les
détails. L'affirmation montre que nous pouvons travailler localement pour la
topologie étale sur $Y$. En particulier, nous pouvons supposer que $Y$ est un
schéma.

\medskip\noindent
Supposons que $Y$ soit un schéma. Choisissons un schéma $X'$ et un morphisme
étale surjectif $s : X' \to X$. Posons $f' = f \circ s : X' \to Y$ et
observons que $f'$ est surjectif, localement de présentation finie et plat.
Affirmation: si le lemme est vrai pour $f'$, il est vrai pour $f$.
En effet, à partir d'un triplet $(\mathcal{G}, \mathcal{H}, \alpha)$ pour $f$,
on obtient par image réciproque par $s$ un triplet
$(s_{small}^{-1}\mathcal{G}, \mathcal{H}, s_{big, fppf}^{-1}\alpha)$ pour $f'$.
Une solution de ce triplet fournit un faisceau $\mathcal{F}$ sur $Y_\etale$
tel que $a_Y^{-1}\mathcal{F} = \mathcal{H}$. D'après le premier paragraphe de
la démonstration, cela signifie que le triplet appartient à l'image essentielle.
Nous sommes ainsi ramenés au cas où $X$ et $Y$ sont tous deux des schémas.
Ce cas découle de Cohomologie étale, lemme
\ref{etale-cohomology-lemma-descent-sheaf-fppf-etale}
au moyen de la discussion de la section \ref{spaces-more-cohomology-section-api}, et en particulier
du lemme \ref{spaces-more-cohomology-lemma-compare-cohomology-other-topologies}.
\end{proof}



\section{Comparaison des topologies fppf et étale: modules}
\label{spaces-more-cohomology-section-fppf-etale-modules}

\noindent
Nous poursuivons l'étude de la section \ref{spaces-more-cohomology-section-fppf-etale}, en examinant
brièvement ici ce qui se passe pour les faisceaux de modules.

\medskip\noindent
Soient $S$ un schéma et $X$ un espace algébrique sur $S$.
Les morphismes de sites $\epsilon_X$, $\pi_X$ et leur composé $a_X$, introduits
à la section \ref{spaces-more-cohomology-section-fppf-etale}, se prolongent naturellement en des
morphismes de sites annelés. Le premier s'écrit
$$
\epsilon_X :
((\textit{Espaces}/X)_{fppf}, \mathcal{O})
\longrightarrow
((\textit{Espaces}/X)_\etale, \mathcal{O})
$$
Notons que nous pouvons employer le même symbole pour le faisceau structural,
car ces faisceaux ont effectivement le même préfaisceau sous-jacent. Le second est
$$
\pi_X :
((\textit{Espaces}/X)_\etale, \mathcal{O})
\longrightarrow
(X_\etale, \mathcal{O}_X)
$$
Le troisième est le morphisme
$$
a_X :
((\textit{Espaces}/X)_{fppf}, \mathcal{O})
\longrightarrow
(X_\etale, \mathcal{O}_X)
$$
Rappelons ce que nous savons déjà des modules quasi-cohérents sur ces sites.

\begin{lemma}
\label{spaces-more-cohomology-lemma-review-quasi-coherent}
Soit $S$ un schéma. Soit $X$ un espace algébrique sur $S$.
Soit $\mathcal{F}$ un $\mathcal{O}_X$-module quasi-cohérent.
\begin{enumerate}
\item La règle
$$
\mathcal{F}^a : (\textit{Espaces}/X)_\etale \longrightarrow \textit{Ab},\quad
(f : Y \to X) \longmapsto \Gamma(Y, f^*\mathcal{F})
$$
satisfait à la condition de faisceau pour les recouvrements fpqc et, a fortiori,
pour les recouvrements fppf et étales,
\item $\mathcal{F}^a = \pi_X^*\mathcal{F}$ sur $(\textit{Espaces}/X)_\etale$,
\item $\mathcal{F}^a = a_X^*\mathcal{F}$ sur $(\textit{Espaces}/X)_{fppf}$,
\item la règle $\mathcal{F} \mapsto \mathcal{F}^a$ définit une équivalence entre
les $\mathcal{O}_X$-modules quasi-cohérents et les modules quasi-cohérents sur
$((\textit{Espaces}/X)_\etale, \mathcal{O})$,
\item la règle $\mathcal{F} \mapsto \mathcal{F}^a$ définit une équivalence entre
les $\mathcal{O}_X$-modules quasi-cohérents et les modules quasi-cohérents sur
$((\textit{Espaces}/X)_{fppf}, \mathcal{O})$,
\item on a $\epsilon_{X, *}a_X^*\mathcal{F} = \pi_X^*\mathcal{F}$
et $a_{X, *}a_X^*\mathcal{F} = \mathcal{F}$,
\item on a $R^i\epsilon_{X, *}(a_X^*\mathcal{F}) = 0$
et $R^ia_{X, *}(a_X^*\mathcal{F}) = 0$ pour $i > 0$.
\end{enumerate}
\end{lemma}

\begin{proof}
L'assertion (1) est une conséquence de la descente fppf des modules
quasi-cohérents. Supposons en effet que $\{f_i : U_i \to U\}$ soit un
recouvrement fpqc dans $(\textit{Espaces}/X)_\etale$. Notons $g : U \to X$ le
morphisme structural. Supposons donnée une famille de sections
$s_i \in \Gamma(U_i , f_i^*g^*\mathcal{F})$ telle que
$s_i|_{U_i \times_U U_j} = s_j|_{U_i \times_U U_j}$. Il faut trouver la section
correspondante $s \in \Gamma(U, g^*\mathcal{F})$. Nous pouvons interpréter
les $s_i$ comme une famille d'applications
$\varphi_i : f_i^*\mathcal{O}_U = \mathcal{O}_{U_i} \to f_i^*g^*\mathcal{F}$
compatibles avec les données de descente canoniques associées aux faisceaux
quasi-cohérents $\mathcal{O}_U$ et $g^*\mathcal{F}$ sur $U$.
Par conséquent, d'après Descente sur les espaces, proposition
\ref{spaces-descent-proposition-fpqc-descent-quasi-coherent}
nous pouvons les descendre (de manière unique) en une application
$\mathcal{O}_U \to g^*\mathcal{F}$, qui fournit notre section $s$.

\medskip\noindent
Nous déduirons les assertions (2) à (7) de l'énoncé correspondant pour les
schémas. Choisissons un recouvrement étale $\{X_i \to X\}_{i \in I}$, où
chaque $X_i$ est un schéma. Observons que $X_i \times_X X_j$ est également un
schéma. Ce recouvrement induit un recouvrement de l'objet final de chacun des
trois sites
$(\textit{Espaces}/X)_{fppf}$, $(\textit{Espaces}/X)_\etale$ et $X_\etale$.
Ainsi, la catégorie des faisceaux sur chacun de ces sites est équivalente à
celle des données de descente pour ces recouvrements; voir Sites, lemme
\ref{sites-lemma-mapping-property-glue}. Les assertions (2) et (3) sont locales
(grâce à l'énoncé de recollement). La quasi-cohérence est une propriété locale;
les assertions (4) et (5) sont donc locales. Les assertions (6) et (7) le sont
manifestement. Il suffit par conséquent de démontrer les assertions (2) à (7)
du lemme lorsque $X$ est un schéma.

\medskip\noindent
Supposons que $X$ soit un schéma. Les plongements
$(\Sch/X)_\etale \subset (\textit{Espaces}/X)_\etale$ et
$(\Sch/X)_{fppf} \subset (\textit{Espaces}/X)_{fppf}$
déterminent des équivalences de topos annelés d'après le lemme
\ref{spaces-more-cohomology-lemma-compare-cohomology-other-topologies}.
Nous concluons que les assertions (2) à (7) découlent du cas des schémas,
Cohomologie étale, lemme
\ref{etale-cohomology-lemma-review-quasi-coherent}.
Pour transporter la propriété de quasi-cohérence par cette équivalence, on
utilise le fait que la quasi-cohérence est une propriété intrinsèque des
modules, comme expliqué dans Modules sur les sites, section
\ref{sites-modules-section-local}. Quelques détails mineurs sont omis.
\end{proof}

\begin{lemma}
\label{spaces-more-cohomology-lemma-cohomological-descent-etale-fppf-modules}
Soit $S$ un schéma. Soit $X$ un espace algébrique sur $S$.
Pour $\mathcal{F}$ un $\mathcal{O}_X$-module quasi-cohérent, les applications
$$
\pi_X^*\mathcal{F} \longrightarrow R\epsilon_{X, *}(a_X^*\mathcal{F})
\quad\text{et}\quad
\mathcal{F} \longrightarrow Ra_{X, *}(a_X^*\mathcal{F})
$$
sont des isomorphismes.
\end{lemma}

\begin{proof}
C'est une conséquence immédiate des assertions (6) et (7) du lemme
\ref{spaces-more-cohomology-lemma-review-quasi-coherent}.
\end{proof}

\begin{lemma}
\label{spaces-more-cohomology-lemma-vanishing-adequate}
Soit $S$ un schéma. Soit $X$ un espace algébrique sur $S$.
Soit $\mathcal{F}_1 \to \mathcal{F}_2 \to \mathcal{F}_3$
un complexe de $\mathcal{O}_X$-modules quasi-cohérents.
Posons
$$
\mathcal{H}_\etale =
\Ker(\pi_X^*\mathcal{F}_2 \to \pi_X^*\mathcal{F}_3)/
\Im(\pi_X^*\mathcal{F}_1 \to \pi_X^*\mathcal{F}_2)
$$
sur $(\textit{Espaces}/X)_\etale$, et posons
$$
\mathcal{H}_{fppf} =
\Ker(a_X^*\mathcal{F}_2 \to a_X^*\mathcal{F}_3)/
\Im(a_X^*\mathcal{F}_1 \to a_X^*\mathcal{F}_2)
$$
sur $(\textit{Espaces}/X)_{fppf}$.
Alors $\mathcal{H}_\etale = \epsilon_{X, *}\mathcal{H}_{fppf}$ et
$$
H^p_\etale(U, \mathcal{H}_\etale) = H^p_{fppf}(U, \mathcal{H}_{fppf}) = 0
$$
pour $p > 0$ et tout objet affine $U$ de $(\textit{Espaces}/X)_\etale$.
\end{lemma}

\noindent
Il y a mieux: les modules sur $(\textit{Espaces}/X)_{fppf}$ qui, localement pour
la topologie fppf, sont de la forme considérée dans le lemme sont appelés
modules adéquats. Ils forment une sous-catégorie de Serre faible de la catégorie
de tous les $\mathcal{O}$-modules, et leur cohomologie est étudiée dans Modules
adéquats, section \ref{adequate-section-adequate}.

\begin{proof}
Pour tout objet $f : U \to X$ de $(\textit{Espaces}/X)_\etale$, considérons la
restriction $\mathcal{H}_\etale|_{U_\etale}$ de $\mathcal{H}_\etale$ à
$U_\etale$ par le foncteur $i_f^* = i_f^{-1}$ étudié à la section
\ref{spaces-more-cohomology-section-compare}. Le faisceau $\mathcal{H}_\etale|_{U_\etale}$ est égal à
l'homologie du complexe $f^*\mathcal{F}_\bullet$ en degré $1$.
Ceci résulte de l'égalité $i_f \circ \pi_X = f$ comme morphismes de sites
annelés $U_\etale \to X_\etale$. En particulier,
$\mathcal{H}_\etale|_{U_\etale}$ est un $\mathcal{O}_U$-module quasi-cohérent.
Soit ensuite $g : V \to U$ un morphisme plat dans
$(\textit{Espaces}/X)_\etale$. Comme
$$
i_{f \circ g}^* \circ \pi_X^* = (f \circ g)^* = g^* \circ f^*
$$
comme morphismes de sites $V_\etale \to X_\etale$, et comme $g$ est plat, donc
$g^*$ exact, nous obtenons
$$
\mathcal{H}_\etale|_{V_\etale} =
g^*\left(\mathcal{H}_\etale|_{U_\etale}\right)
$$
Ces préparatifs achevés, nous pouvons démontrer le lemme.

\medskip\noindent
Soit $\mathcal{U} = \{g_i : U_i \to U\}_{i \in I}$ un recouvrement fppf, avec
$f : U \to X$ comme ci-dessus. La propriété de faisceau vaut pour
$\mathcal{H}_\etale$ et le recouvrement $\mathcal{U}$, d'après l'assertion (1)
du lemme \ref{spaces-more-cohomology-lemma-review-quasi-coherent}, appliquée à
$\mathcal{H}_\etale|_{U_\etale}$, et d'après ce qui précède.
Ainsi $\mathcal{H}_\etale$ est déjà un faisceau fppf, ce qui signifie que
$\mathcal{H}_{fppf}$ est égal à $\mathcal{H}_\etale$ comme préfaisceau.
En particulier,
$\mathcal{H}_\etale = \epsilon_{X, *}\mathcal{H}_{fppf}$.

\medskip\noindent
Enfin, pour démontrer l'annulation, nous utilisons Cohomologie sur les sites,
lemme \ref{sites-cohomology-lemma-cech-vanish-collection}. Soit $\mathcal{B}$
la classe des objets affines de $(\textit{Espaces}/X)_{fppf}$ et soit
$\text{Cov}$ l'ensemble des recouvrements fppf finis
$\mathcal{U} = \{U_i \to U\}_{i = 1, \ldots, n}$ où $U$ et les $U_i$ sont affines.
On a
$$
{\check H}^p(\mathcal{U}, \mathcal{H}_\etale) =
{\check H}^p(\mathcal{U}, \left(\mathcal{H}_\etale|_{U_\etale}\right)^a)
$$
car, d'après le premier paragraphe, les valeurs de $\mathcal{H}_\etale$ sur les
schémas affines $U_{i_0} \times_U \ldots \times_U U_{i_p}$, plats sur $U$,
coïncident avec les valeurs de l'image réciproque du module quasi-cohérent
$\mathcal{H}_\etale|_{U_\etale}$. L'annulation résulte alors de Descente, lemme
\ref{descent-lemma-standard-covering-Cech-quasi-coherent}.
Ceci achève la démonstration.
\end{proof}

\begin{lemma}
\label{spaces-more-cohomology-lemma-cohomological-descent-etale-fppf-modules-unbounded}
Soit $S$ un schéma. Soit $X$ un espace algébrique sur $S$.
Pour $K \in D_\QCoh(\mathcal{O}_X)$, les applications
$$
L\pi_X^*K \longrightarrow R\epsilon_{X, *}(La_X^*K)
\quad\text{et}\quad
K \longrightarrow Ra_{X, *}(La_X^*K)
$$
sont des isomorphismes. Ici,
$a_X : \Sh((\textit{Espaces}/X)_{fppf}) \to \Sh(X_\etale)$ est défini ci-dessus.
\end{lemma}

\begin{proof}
La question est locale pour la topologie étale sur $X$; nous pouvons donc
supposer que $X$ est affine, disons $X = \Spec(A)$. Alors
$D_\QCoh(\mathcal{O}_X) = D(A)$ d'après Catégories dérivées des espaces, lemme
\ref{spaces-perfect-lemma-derived-quasi-coherent-small-etale-site}
et Catégories dérivées des schémas, lemme
\ref{perfect-lemma-affine-compare-bounded}.
Nous pouvons donc choisir un complexe K-plat de $A$-modules $K^\bullet$ dont
le complexe correspondant $\mathcal{K}^\bullet$ de $\mathcal{O}_X$-modules
quasi-cohérents représente $K$.
Nous affirmons que $\mathcal{K}^\bullet$ est un complexe K-plat de
$\mathcal{O}_X$-modules.

\medskip\noindent
Démonstration de l'affirmation. D'après Catégories dérivées des schémas, lemme
\ref{perfect-lemma-affine-K-flat}
le complexe $\widetilde{K}^\bullet$ est K-plat sur le schéma
$(\Spec(A), \mathcal{O}_{\Spec(A)})$.
Ensuite, notons que $\mathcal{K}^\bullet = \epsilon^*\widetilde{K}^\bullet$,
où $\epsilon$ est défini comme dans Catégories dérivées des espaces, lemme
\ref{spaces-perfect-lemma-derived-quasi-coherent-small-etale-site}
; le complexe $\mathcal{K}^\bullet$ est donc K-plat d'après Cohomologie sur les
sites, lemme \ref{sites-cohomology-lemma-pullback-K-flat-points}, et le fait que
le site étale d'un schéma possède suffisamment de points (Cohomologie étale,
remarques \ref{etale-cohomology-remarks-enough-points}).

\medskip\noindent
D'après l'affirmation, on a
$La_X^*K = a_X^*\mathcal{K}^\bullet$ et
$L\pi_X^*K = \pi_X^*\mathcal{K}^\bullet$.
Comme la première partie de la démonstration montre que l'image réciproque
$a_X^*\mathcal{K}^n$ du module quasi-cohérent est acyclique pour
$\epsilon_{X, *}$, resp.\ $a_{X, *}$, le lemme d'acyclicité de Leray ne
permet-il pas aussitôt de conclure? En fait, non, car ce lemme ne s'applique
qu'aux complexes bornés inférieurement. Nous montrerons toutefois au paragraphe
suivant que le résultat découle bien du cas borné inférieurement, parce que
notre complexe est la limite dérivée de complexes bornés inférieurement de
modules quasi-cohérents.

\medskip\noindent
Les faisceaux de cohomologie de $\pi_X^*\mathcal{K}^\bullet$ et de
$a_X^*\mathcal{K}^\bullet$ ont des groupes de cohomologie supérieure nuls sur
les objets affines de $(\textit{Espaces}/X)_\etale$, d'après le lemme
\ref{spaces-more-cohomology-lemma-vanishing-adequate}. Par conséquent, on a
$$
L\pi_X^*K = R\lim \tau_{\geq -n}(L\pi_X^*K)
\quad\text{et}\quad
La_X^*K = R\lim \tau_{\geq -n}(La_X^*K)
$$
d'après Cohomologie sur les sites, lemme
\ref{sites-cohomology-lemma-is-limit-dimension}.

\medskip\noindent
Démonstration de $L\pi_X^*K = R\epsilon_{X, *}(La_X^*\mathcal{F})$.
D'après ce qui précède, on a
$$
R\epsilon_{X, *}La_X^*K =
R\lim R\epsilon_{X, *}(\tau_{\geq -n}(La_X^*K))
$$
d'après Cohomologie sur les sites, lemme
\ref{sites-cohomology-lemma-Rf-commutes-with-Rlim}.
Notons que $\tau_{\geq -n}(La_X^*K)$ est représenté par
$\tau_{\geq -n}(a_X^*\mathcal{K}^\bullet)$, qui peut différer de
$a_X^*(\tau_{\geq -n}\mathcal{K}^\bullet)$.
Mais les systèmes
$$
\{\tau_{\geq -n}(a_X^*\mathcal{K}^\bullet)\}_{n \geq 1}
\quad\text{et}\quad
\{a_X^*(\tau_{\geq -n}\mathcal{K}^\bullet)\}_{n \geq 1}
$$
sont manifestement isomorphes comme pro-systèmes.
D'après le lemme d'acyclicité de Leray (Catégories dérivées, lemme
\ref{derived-lemma-leray-acyclicity}) et la première partie du lemme, on a
$$
R\epsilon_{X, *}(a_X^*(\tau_{\geq -n}\mathcal{K}^\bullet)) =
\pi_X^*(\tau_{\geq -n}\mathcal{K}^\bullet)
$$
Nous pouvons ensuite utiliser le fait que les systèmes
$$
\{\tau_{\geq -n}(\pi_X^*\mathcal{K}^\bullet)\}_{n \geq 1}
\quad\text{et}\quad
\{\pi_X^*(\tau_{\geq -n}\mathcal{K}^\bullet)\}_{n \geq 1}
$$
sont isomorphes comme pro-systèmes. Enfin, regroupons toutes les égalités:
\begin{align*}
R\epsilon_{X, *}La_X^*K
& =
R\epsilon_{X, *} (R\lim \tau_{\geq -n}(La_X^*K)) \\
& =
R\lim R\epsilon_{X, *}(\tau_{\geq -n}(La_X^*K)) \\
& =
R\lim R\epsilon_{X, *}(\tau_{\geq -n}(a_X^*\mathcal{K}^\bullet)) \\
& =
R\lim R\epsilon_{X, *}(a_X^*(\tau_{\geq -n}\mathcal{K}^\bullet)) \\
& =
R\lim \pi_X^*(\tau_{\geq -n}\mathcal{K}^\bullet) \\
& =
R\lim \tau_{\geq -n}(\pi_X^*\mathcal{K}^\bullet) \\
& =
R\lim \tau_{\geq -n}(L\pi_X^*K) \\
& =
L\pi_X^*K
\end{align*}
Dans les quatrième et sixième égalités, nous avons utilisé que des pro-systèmes
isomorphes ont le même $R\lim$ (un petit détail est omis).
On peut éviter cette étape en utilisant les résultats plus précis sur la
cohomologie des termes du complexe $\tau_{\geq -n}a_X^*\mathcal{K}^\bullet$
démontrés au lemme \ref{spaces-more-cohomology-lemma-vanishing-adequate}; ils donnent directement
$R\epsilon_{X, *}(\tau_{\geq -n}(a_X^*\mathcal{K}^\bullet)) =
\tau_{\geq -n}(\pi_X^*\mathcal{K}^\bullet)$.

\medskip\noindent
L'égalité $K = Ra_{X, *}(La_X^*\mathcal{F})$ se démontre exactement de la même
manière, en utilisant à la dernière étape l'égalité
$K = R\lim \tau_{\geq -n}K$, donnée par Catégories dérivées des espaces,
lemme \ref{spaces-perfect-lemma-nice-K-injective}.
\end{proof}








\section{Comparaison des topologies ph et étale}
\label{spaces-more-cohomology-section-ph-etale}

\noindent
Cette section est l'analogue de Cohomologie étale, section
\ref{etale-cohomology-section-ph-etale}.

\medskip\noindent
Soient $S$ un schéma et $X$ un espace algébrique sur $S$.
Sur la catégorie $\textit{Espaces}/X$, considérons les topologies ph et étale.
Le foncteur identité
$(\textit{Espaces}/X)_\etale \to (\textit{Espaces}/X)_{ph}$
est continu, puisque tout recouvrement étale est un recouvrement ph d'après
Topologies sur les espaces, lemme
\ref{spaces-topologies-lemma-zariski-etale-smooth-syntomic-fppf-ph}.
Il définit donc un morphisme de sites
$$
\epsilon_X :
(\textit{Espaces}/X)_{ph} \longrightarrow (\textit{Espaces}/X)_\etale
$$
par application de Sites, proposition \ref{sites-proposition-get-morphism}.
Notons que $\epsilon_{X, *}$ est le foncteur identité sur les préfaisceaux
sous-jacents et que $\epsilon_X^{-1}$ associe à un faisceau étale son faisceau
associé pour la topologie ph.
Considérons le morphisme de sites
$$
\pi_X : (\textit{Espaces}/X)_\etale \longrightarrow X_{spaces, \etale}
$$
qui compare les grands et petits sites étales; voir la section \ref{spaces-more-cohomology-section-compare}.
Le composé détermine un morphisme de sites
$$
a_X = \pi_X \circ \epsilon_X :
(\textit{Espaces}/X)_{ph}
\longrightarrow
X_{spaces, \etale}
$$
Si $\mathcal{H}$ est un faisceau abélien sur $(\textit{Espaces}/X)_{ph}$,
nous noterons $H^n_{ph}(U, \mathcal{H})$ la cohomologie de
$\mathcal{H}$ sur un objet $U$ de $(\textit{Espaces}/X)_{ph}$.

\begin{lemma}
\label{spaces-more-cohomology-lemma-comparison-ph-etale}
Soit $S$ un schéma. Soit $X$ un espace algébrique sur $S$.
\begin{enumerate}
\item Pour $\mathcal{F} \in \Sh(X_\etale)$, on a
$\epsilon_{X, *}a_X^{-1}\mathcal{F} = \pi_X^{-1}\mathcal{F}$
et $a_{X, *}a_X^{-1}\mathcal{F} = \mathcal{F}$.
\item Pour $\mathcal{F} \in \textit{Ab}(X_\etale)$ de torsion, on a
$R^i\epsilon_{X, *}(a_X^{-1}\mathcal{F}) = 0$ pour $i > 0$.
\end{enumerate}
\end{lemma}

\begin{proof}
On a $a_X^{-1}\mathcal{F} = \epsilon_X^{-1} \pi_X^{-1}\mathcal{F}$.
D'après le lemme \ref{spaces-more-cohomology-lemma-describe-pullback}, le faisceau étale
$\pi_X^{-1}\mathcal{F}$ est un faisceau pour la topologie ph; il est donc égal
à $a_X^{-1}\mathcal{F}$ (puisque l'image réciproque par $\epsilon_X$ est donnée
par le faisceau associé pour la topologie ph). Rappelons en outre que
$\epsilon_{X, *}$ est le foncteur identité sur les préfaisceaux sous-jacents.
L'assertion (1) résulte alors immédiatement de la description explicite de
$\pi_X^{-1}$ donnée au lemme \ref{spaces-more-cohomology-lemma-describe-pullback}.

\medskip\noindent
Nous démontrerons (2) en nous ramenant au cas des schémas; voir l'assertion (1)
de Cohomologie étale, lemme
\ref{etale-cohomology-lemma-V-C-all-n-etale-ph}.
Cette méthode « fonctionne clairement », puisque tout espace algébrique est
localement un schéma pour la topologie étale. Les détails sont donnés ci-dessous,
mais nous invitons vivement le lecteur à omettre la démonstration.

\medskip\noindent
Pour un faisceau abélien $\mathcal{H}$ sur $(\textit{Espaces}/X)_{ph}$,
l'image directe supérieure $R^p\epsilon_{X, *}\mathcal{H}$ est le faisceau
associé au préfaisceau $U \mapsto H^p_{ph}(U, \mathcal{H})$
sur $(\textit{Espaces}/X)_\etale$. Voir Cohomologie sur les sites, lemme
\ref{sites-cohomology-lemma-higher-direct-images}.
Comme tout objet de $(\textit{Espaces}/X)_\etale$ admet un recouvrement par des
schémas, il suffit de démontrer que, pour tout schéma $U/X$ et tout
$\xi \in H^p_{ph}(U, a_X^{-1}\mathcal{F})$, il existe un recouvrement étale
$\{U_i \to U\}$ tel que la restriction de $\xi$ à chaque $U_i$ soit nulle.
On a
\begin{align*}
H^p_{ph}(U, a_X^{-1}\mathcal{F})
& =
H^p((\textit{Espaces}/U)_{ph}, (a_X^{-1}\mathcal{F})|_{\textit{Espaces}/U}) \\
& =
H^p((\Sch/U)_{ph}, (a_X^{-1}\mathcal{F})|_{\Sch/U})
\end{align*}
où la seconde identification est le lemme
\ref{spaces-more-cohomology-lemma-compare-cohomology-other-topologies}, et la première un fait général
sur la restriction (Cohomologie sur les sites, lemme
\ref{sites-cohomology-lemma-cohomology-of-open}).
En comparant le premier paragraphe et le résultat correspondant dans le cas
des schémas (Cohomologie étale, lemme
\ref{etale-cohomology-lemma-describe-pullback-pi-ph})
nous concluons que le faisceau $(a_X^{-1}\mathcal{F})|_{\Sch/U}$
coïncide avec l'image réciproque par la « version pour les schémas de $a_U$ ».
Nous pouvons donc trouver un recouvrement étale $\{U_i \to U\}$ tel que notre
classe s'annule dans
$H^p((\Sch/U_i)_{ph}, (a_X^{-1}\mathcal{F})|_{\Sch/U_i})$
pour tout $i$; voir Cohomologie étale, lemme
\ref{etale-cohomology-lemma-V-C-all-n-etale-ph}
(l'énoncé précis à utiliser ici est que $V_n$ vaut pour tout $n$, ce qui est
l'assertion (2) dans le cas des schémas).
En revenant en arrière (au moyen des mêmes formules que ci-dessus, appliquées
cette fois à $U_i$), nous concluons que la restriction de $\xi$ à $U_i$
est nulle, comme souhaité.
\end{proof}

\noindent
Le travail difficile effectué dans le cas des schémas montre maintenant que
les cohomologies étale et ph coïncident pour les faisceaux abéliens de torsion
provenant du petit site étale.

\begin{lemma}
\label{spaces-more-cohomology-lemma-cohomological-descent-etale-ph}
Soit $S$ un schéma. Soit $X$ un espace algébrique sur $S$.
Pour $K \in D^+(X_\etale)$ à faisceaux de cohomologie de torsion, les applications
$$
\pi_X^{-1}K \longrightarrow R\epsilon_{X, *}a_X^{-1}K
\quad\text{et}\quad
K \longrightarrow Ra_{X, *}a_X^{-1}K
$$
sont des isomorphismes, où
$a_X : \Sh((\textit{Espaces}/X)_{ph}) \to \Sh(X_\etale)$ est défini ci-dessus.
\end{lemma}

\begin{proof}
Nous démontrons seulement le second énoncé; le premier est plus facile et se
démontre exactement de la même manière. On se ramène au cas où $K$ est donné
par un seul faisceau abélien de torsion. En effet, représentons $K$ par un
complexe borné inférieurement $\mathcal{F}^\bullet$ de faisceaux abéliens de
torsion. C'est possible d'après Cohomologie sur les sites, lemme
\ref{sites-cohomology-lemma-torsion}.
Le cas d'un faisceau donne
$\mathcal{F}^n = a_{X, *} a_X^{-1} \mathcal{F}^n$
et les faisceaux $R^qa_{X, *}a_X^{-1}\mathcal{F}^n$ sont nuls pour $q > 0$.
Le lemme d'acyclicité de Leray (Catégories dérivées, lemme
\ref{derived-lemma-leray-acyclicity}), appliqué à
$a_X^{-1}\mathcal{F}^\bullet$ et au foncteur $a_{X, *}$, permet de conclure.
Supposons désormais $K = \mathcal{F}$, où $\mathcal{F}$ est un faisceau
abélien de torsion.

\medskip\noindent
Le lemme \ref{spaces-more-cohomology-lemma-comparison-ph-etale} donne
$a_{X, *}a_X^{-1}\mathcal{F} = \mathcal{F}$. Il suffit donc de montrer que
$R^qa_{X, *}a_X^{-1}\mathcal{F} = 0$ pour $q > 0$.
Pour cela, nous pouvons utiliser $a_X = \epsilon_X \circ \pi_X$ et la suite
spectrale de Leray (Cohomologie sur les sites, lemme
\ref{sites-cohomology-lemma-relative-Leray}).
D'après le lemme \ref{spaces-more-cohomology-lemma-comparison-ph-etale}, on a
$R^i\epsilon_{X, *}(a_X^{-1}\mathcal{F}) = 0$ pour $i > 0$.
On a
$\epsilon_{X, *}a_X^{-1}\mathcal{F} = \pi_X^{-1}\mathcal{F}$
et, d'après le lemme \ref{spaces-more-cohomology-lemma-cohomological-descent-etale},
$R^j\pi_{X, *}(\pi_X^{-1}\mathcal{F}) = 0$ pour $j > 0$.
Ceci achève la démonstration.
\end{proof}

\begin{lemma}
\label{spaces-more-cohomology-lemma-compare-cohomology-etale-ph}
Soient $S$ un schéma et $X$ un espace algébrique sur $S$.
Avec $a_X : \Sh((\textit{Espaces}/X)_{ph}) \to \Sh(X_\etale)$
défini ci-dessus:
\begin{enumerate}
\item $H^q(X_\etale, \mathcal{F}) = H^q_{ph}(X, a_X^{-1}\mathcal{F})$
pour un faisceau abélien de torsion $\mathcal{F}$ sur $X_\etale$,
\item $H^q(X_\etale, K) = H^q_{ph}(X, a_X^{-1}K)$ pour $K \in D^+(X_\etale)$
à faisceaux de cohomologie de torsion.
\end{enumerate}
Exemple: si $A$ est un groupe abélien de torsion, alors
$H^q_\etale(X, \underline{A}) = H^q_{ph}(X, \underline{A})$.
\end{lemma}

\begin{proof}
Ceci découle du lemme \ref{spaces-more-cohomology-lemma-cohomological-descent-etale-ph}, d'après
Cohomologie sur les sites, remarque \ref{sites-cohomology-remark-before-Leray}.
\end{proof}

\begin{lemma}
\label{spaces-more-cohomology-lemma-push-pull-ph-etale}
Soit $S$ un schéma.
Soit $f : X \to Y$ un morphisme d'espaces algébriques sur $S$.
Il existe alors des diagrammes commutatifs de topos
$$
\xymatrix{
\Sh((\textit{Espaces}/X)_{ph}) \ar[rr]_{f_{big, ph}} \ar[d]_{\epsilon_X} & &
\Sh((\textit{Espaces}/Y)_{ph}) \ar[d]^{\epsilon_Y} \\
\Sh((\textit{Espaces}/X)_\etale) \ar[rr]^{f_{big, \etale}} & &
\Sh((\textit{Espaces}/Y)_\etale)
}
$$
et
$$
\xymatrix{
\Sh((\textit{Espaces}/X)_{ph}) \ar[rr]_{f_{big, ph}} \ar[d]_{a_X} & &
\Sh((\textit{Espaces}/Y)_{ph}) \ar[d]^{a_Y} \\
\Sh(X_\etale) \ar[rr]^{f_{small}} & &
\Sh(Y_\etale)
}
$$
avec $a_X = \pi_X \circ \epsilon_X$ et $a_Y = \pi_X \circ \epsilon_X$.
\end{lemma}

\begin{proof}
Ceci découle immédiatement des définitions des morphismes en jeu; voir
Topologies sur les espaces, section \ref{spaces-topologies-section-ph},
et la section \ref{spaces-more-cohomology-section-compare}.
\end{proof}

\begin{lemma}
\label{spaces-more-cohomology-lemma-proper-push-pull-ph-etale}
Dans le lemme \ref{spaces-more-cohomology-lemma-push-pull-ph-etale}, si $f$ est propre, alors on a
\begin{enumerate}
\item $a_Y^{-1} \circ f_{small, *} = f_{big, ph, *} \circ a_X^{-1}$, et
\item
$a_Y^{-1}(Rf_{small, *}K) = Rf_{big, ph, *}(a_X^{-1}K)$
pour $K$ dans $D^+(X_\etale)$ à faisceaux de cohomologie de torsion.
\end{enumerate}
\end{lemma}

\begin{proof}
Démonstration de (1). On pourrait la donner en répétant celle du lemme
\ref{spaces-more-cohomology-lemma-compare-higher-direct-image-proper}, assertion (1); nous allons
plutôt en déduire le résultat. Comme $\epsilon_{Y, *}$ est le foncteur identité
sur les préfaisceaux sous-jacents, il reflète les isomorphismes. Le lemme
\ref{spaces-more-cohomology-lemma-comparison-ph-etale} montre que
$\epsilon_{Y, *} \circ a_Y^{-1} = \pi_Y^{-1}$, et de même pour $X$.
Pour montrer que l'application canonique
$a_Y^{-1}f_{small, *}\mathcal{F} \to f_{big, ph, *}a_X^{-1}\mathcal{F}$
est un isomorphisme, il suffit de montrer que
\begin{align*}
\pi_Y^{-1}f_{small, *}\mathcal{F}
& =
\epsilon_{Y, *}a_Y^{-1}f_{small, *}\mathcal{F} \\
& \to 
\epsilon_{Y, *}f_{big, ph, *}a_X^{-1}\mathcal{F} \\
& =
f_{big, \etale, *} \epsilon_{X, *}a_X^{-1}\mathcal{F} \\
& =
f_{big, \etale, *}\pi_X^{-1}\mathcal{F}
\end{align*}
est un isomorphisme. C'est l'assertion (1) du lemme
\ref{spaces-more-cohomology-lemma-compare-higher-direct-image-proper}.

\medskip\noindent
Pour démontrer (2), nous utilisons
\begin{align*}
R\epsilon_{Y, *}Rf_{big, ph, *}a_X^{-1}K
& =
Rf_{big, \etale, *}R\epsilon_{X, *}a_X^{-1}K \\
& =
Rf_{big, \etale, *}\pi_X^{-1}K \\
& =
\pi_Y^{-1}Rf_{small, *}K \\
& =
R\epsilon_{Y, *} a_Y^{-1}Rf_{small, *}K
\end{align*}
La première égalité résulte du diagramme commutatif du lemme
\ref{spaces-more-cohomology-lemma-push-pull-ph-etale} et de Cohomologie sur les sites, lemme
\ref{sites-cohomology-lemma-derived-pushforward-composition}.
La deuxième est le lemme \ref{spaces-more-cohomology-lemma-cohomological-descent-etale-ph}; la
troisième est l'assertion (2) du lemme
\ref{spaces-more-cohomology-lemma-compare-higher-direct-image-proper}; la quatrième est de nouveau
le lemme \ref{spaces-more-cohomology-lemma-cohomological-descent-etale-ph}.
Ainsi le morphisme de changement de base
$a_Y^{-1}(Rf_{small, *}K) \to Rf_{big, ph, *}(a_X^{-1}K)$
induit un isomorphisme
$$
R\epsilon_{Y, *}a_Y^{-1}Rf_{small, *}K \to
R\epsilon_{Y, *}Rf_{big, ph, *}a_X^{-1}K
$$
La remarque suivante achève la démonstration: soit
$\alpha : a_Y^{-1}L \to M$, avec $L$ dans $D^+(Y_\etale)$ à faisceaux de
cohomologie de torsion et $M$ dans $D^+((\textit{Espaces}/Y)_{ph})$.
Si $R\epsilon_{Y, *}\alpha$ est un isomorphisme, alors $\alpha$ est un
isomorphisme. En effet, montrons par récurrence sur $i$ que $H^i(\alpha)$ est un
isomorphisme. Ceci est vrai pour tout $i$ suffisamment petit. Si cela vaut pour
$i \leq i_0$, alors $R^j\epsilon_{Y, *}H^i(M) = 0$ pour $j > 0$ et
$i \leq i_0$ d'après le lemme \ref{spaces-more-cohomology-lemma-comparison-ph-etale}, puisque
$H^i(M) = a_Y^{-1}H^i(L)$ dans cet intervalle. Un argument de suite spectrale
donne alors
$\epsilon_{Y, *}H^{i_0 + 1}(M) = H^{i_0 + 1}(R\epsilon_{Y, *}M)$.
Par conséquent, $\epsilon_{Y, *}H^{i_0 + 1}(M) =
\pi_Y^{-1}H^{i_0 + 1}(L) =
\epsilon_{Y, *}a_Y^{-1}H^{i_0 + 1}(L)$. Il s'ensuit que
$H^{i_0 + 1}(\alpha)$ est un isomorphisme (car $\epsilon_{Y, *}$ reflète les
isomorphismes, puisqu'il est l'identité sur les préfaisceaux sous-jacents),
comme voulu.
\end{proof}

















% Bien, commençons ici.
%

\chapter{Espaces simpliciaux}



\phantomsection
\label{spaces-simplicial-section-phantom}


\section{Introduction}
\label{spaces-simplicial-section-introduction}

\noindent
Ce chapitre développe une théorie des espaces topologiques simpliciaux, des
espaces annelés simpliciaux, des schémas simpliciaux et des espaces
algébriques simpliciaux. Cette théorie permet parfois de démontrer des principes
du local au global qui paraissent difficiles à établir autrement.
On trouvera quelques exemples d'applications dans les articles
\cite{faltings_finiteness}, \cite{Kiehl} et \cite{HodgeIII}.

\medskip\noindent
Nous supposons dans tout ce chapitre que le lecteur connaît les notions et
résultats fondamentaux du chapitre Méthodes simpliciales, voir
Méthodes simpliciales, section \ref{simplicial-section-introduction}.
En particulier, nous continuons à écrire $X$, et non $X_\bullet$,
pour un objet simplicial.








\section{Espaces topologiques simpliciaux}
\label{spaces-simplicial-section-simplicial-top}

\noindent
Un {\it espace simplicial} est un objet simplicial de la catégorie des
espaces topologiques, dont les morphismes sont les applications continues.
(Nous emploierons ``espace algébrique simplicial'' pour désigner les objets
simpliciaux de la catégorie des espaces algébriques.)
On peut représenter un espace simplicial $X$ comme suit :
$$
\xymatrix{
X_2
\ar@<2ex>[r]
\ar@<0ex>[r]
\ar@<-2ex>[r]
&
X_1
\ar@<1ex>[r]
\ar@<-1ex>[r]
\ar@<1ex>[l]
\ar@<-1ex>[l]
&
X_0
\ar@<0ex>[l]
}
$$
Il y a ici deux morphismes $d^1_0, d^1_1 : X_1 \to X_0$
et un unique morphisme $s^0_0 : X_0 \to X_1$, etc.
Il importe de garder à l'esprit que $d^n_i : X_n \to X_{n - 1}$
doit être considéré comme une ``projection qui oublie la
$i$-ième coordonnée'', et $s^n_j : X_n \to X_{n + 1}$ comme l'application
diagonale qui répète la $j$-ième coordonnée.

\medskip\noindent
Soit $X$ un espace simplicial. Nous lui associons un site
$X_{Zar}$\footnote{Cette notation est analogue à celle de
Sites, exemple \ref{sites-example-site-topological}
et de
Topologies, définition \ref{topologies-definition-big-small-Zariski}.}
à $X$ de la manière suivante.
\begin{enumerate}
\item Un objet de $X_{Zar}$ est un ouvert $U$ de $X_n$ pour un certain $n$,
\item un morphisme $U \to V$ de $X_{Zar}$ est donné par un
$\varphi : [m] \to [n]$, où $n, m$ sont tels que
$U \subset X_n$, $V \subset X_m$, et où $\varphi$ vérifie
$X(\varphi)(U) \subset V$, et
\item un recouvrement $\{U_i \to U\}$ dans $X_{Zar}$ signifie
que $U, U_i \subset X_n$ sont ouverts, que les applications $U_i \to U$ sont
données par $\text{id} : [n] \to [n]$, et que $U = \bigcup U_i$.
\end{enumerate}
Remarquons en particulier que, si $U \to V$ est un morphisme de $X_{Zar}$
donné par $\varphi$, alors $X(\varphi) : X_n \to X_m$ induit bien
une application continue $U \to V$ d'espaces topologiques.

\noindent
Il est clair que ce qui précède est un cas particulier d'une construction
qui associe un site à tout diagramme d'espaces topologiques. Énonçons
le lemme obligé.

\begin{lemma}
\label{spaces-simplicial-lemma-simplicial-site}
Soit $X$ un espace simplicial. Alors $X_{Zar}$, tel qu'il est
défini ci-dessus, est un site.
\end{lemma}

\begin{proof}
Omis.
\end{proof}

\noindent
Soit $X$ un espace simplicial. Soit $\mathcal{F}$ un faisceau sur $X_{Zar}$.
Il résulte clairement de la définition des recouvrements que la restriction
de $\mathcal{F}$ aux ouverts de $X_n$ définit un faisceau $\mathcal{F}_n$
sur l'espace topologique $X_n$. Pour tout $\varphi : [m] \to [n]$, les
applications de restriction de $\mathcal{F}$ pour les couples $U \subset X_n$, $V \subset X_m$
tels que $X(\varphi)(U) \subset V$ définissent un $X(\varphi)$-morphisme
$\mathcal{F}(\varphi) : \mathcal{F}_m \to \mathcal{F}_n$, voir
Faisceaux, définition \ref{sheaves-definition-f-map}.
De plus, étant donnés $\varphi : [m] \to [n]$ et $\psi : [l] \to [m]$,
on a
$$
\mathcal{F}(\varphi) \circ \mathcal{F}(\psi) =
\mathcal{F}(\varphi \circ \psi)
$$
(dans le membre de gauche, on utilise la composition des $f$-morphismes, voir
Faisceaux, définition \ref{sheaves-definition-composition-f-maps}).
La réciproque est manifestement vraie elle aussi : si l'on dispose d'un système
$(\{\mathcal{F}_n\}_{n \geq 0},
\{\mathcal{F}(\varphi)\}_{\varphi \in \text{Flèches}(\Delta)})$
comme ci-dessus, satisfaisant les égalités affichées,
on obtient alors un faisceau sur $X_{Zar}$.

\begin{lemma}
\label{spaces-simplicial-lemma-describe-sheaves-simplicial-site}
Soit $X$ un espace simplicial. Il existe une équivalence de
catégories entre
\begin{enumerate}
\item $\Sh(X_{Zar})$, et
\item la catégorie des systèmes $(\mathcal{F}_n, \mathcal{F}(\varphi))$
décrits ci-dessus.
\end{enumerate}
\end{lemma}

\begin{proof}
Voir la discussion ci-dessus.
\end{proof}

\begin{lemma}
\label{spaces-simplicial-lemma-simplicial-space-site-functorial}
Soit $f : Y \to X$ un morphisme d'espaces simpliciaux.
Alors le foncteur $u : X_{Zar} \to Y_{Zar}$
qui associe à l'ouvert $U \subset X_n$ l'ouvert
$f_n^{-1}(U) \subset Y_n$ définit un morphisme de sites
$f_{Zar} : Y_{Zar} \to X_{Zar}$.
\end{lemma}

\begin{proof}
Il est clair que $u$ est un foncteur continu. On obtient donc les foncteurs
$f_{Zar, *} = u^s$ et $f_{Zar}^{-1} = u_s$, voir
Sites, section \ref{sites-section-morphism-sites}.
Pour voir que l'on obtient un morphisme de sites, il faut montrer
que $u_s$ est exact. Nous utiliserons
Sites, lemme \ref{sites-lemma-directed-morphism} à cette fin.
Soit $V \subset Y_n$ un ouvert. La catégorie
$\mathcal{I}_V^u$ (voir Sites, section \ref{sites-section-functoriality-PSh})
est formée des couples $(U, \varphi)$ où
$\varphi : [m] \to [n]$ et où $U \subset X_m$ est un ouvert tel que
$Y(\varphi)(V) \subset f_m^{-1}(U)$. De plus, un morphisme
$(U, \varphi) \to (U', \varphi')$ est donné par un
$\psi : [m'] \to [m]$ tel que $X(\psi)(U) \subset U'$
et $\varphi \circ \psi = \varphi'$.
Il s'agit de montrer que $\mathcal{I}_V^u$ est cofiltrante.

\medskip\noindent
Vérifions les conditions de
Catégories, définition \ref{categories-definition-codirected}.
La condition (1) est satisfaite parce que $(X_n, \text{id}_{[n]})$ est un objet.
Soit $(U, \varphi)$ un objet. La condition
$Y(\varphi)(V) \subset f_m^{-1}(U)$ équivaut à
$V \subset f_n^{-1}(X(\varphi)^{-1}(U))$. On obtient donc un morphisme
$(X(\varphi)^{-1}(U), \text{id}_{[n]}) \to (U, \varphi)$ en posant
$\psi = \varphi$. De plus, étant donné un couple d'objets
de la forme $(U, \text{id}_{[n]})$ et $(U', \text{id}_{[n]})$,
on voit qu'il existe un objet, à savoir $(U \cap U', \text{id}_{[n]})$,
qui se projette sur chacun d'eux. La condition (2) est donc satisfaite.
Pour vérifier la condition (3), supposons donnés deux morphismes
$a, a': (U, \varphi) \to (U', \varphi')$ définis par $\psi, \psi' : [m'] \to [m]$.
Leur précomposition par le morphisme
$(X(\varphi)^{-1}(U), \text{id}_{[n]}) \to (U, \varphi)$ défini
par $\varphi$ égalise alors $a, a'$, puisque
$\varphi \circ \psi = \varphi' = \varphi \circ \psi'$.
Ceci achève la démonstration.
\end{proof}

\begin{lemma}
\label{spaces-simplicial-lemma-describe-functoriality}
Soit $f : Y \to X$ un morphisme d'espaces simpliciaux. En termes de la
description des faisceaux donnée dans le
lemme \ref{spaces-simplicial-lemma-describe-sheaves-simplicial-site}, le
morphisme $f_{Zar}$ du lemme \ref{spaces-simplicial-lemma-simplicial-space-site-functorial}
se décrit comme suit.
\begin{enumerate}
\item Si $\mathcal{G}$ est un faisceau sur $Y$, alors
$(f_{Zar, *}\mathcal{G})_n = f_{n, *}\mathcal{G}_n$.
\item Si $\mathcal{F}$ est un faisceau sur $X$, alors
$(f_{Zar}^{-1}\mathcal{F})_n = f_n^{-1}\mathcal{F}_n$.
\end{enumerate}
\end{lemma}

\begin{proof}
La première assertion résulte immédiatement des définitions. Pour la seconde,
remarquons que, dans la démonstration du
lemme \ref{spaces-simplicial-lemma-simplicial-space-site-functorial},
nous avons montré que, pour un ouvert $V \subset Y_n$, la catégorie
$(\mathcal{I}_V^u)^{opp}$ contient comme sous-catégorie cofinale
la catégorie des ouverts $U \subset X_n$ tels que $f_n^{-1}(U) \supset V$,
les morphismes étant les inclusions. On voit donc que la restriction
de $u_p\mathcal{F}$ aux ouverts de $Y_n$ est le préfaisceau
$f_{n, p}\mathcal{F}_n$ défini dans
Faisceaux, lemme \ref{sheaves-lemma-pullback-presheaves}.
Puisque $f_{Zar}^{-1}\mathcal{F} = u_s\mathcal{F}$ est le faisceau associé à
$u_p\mathcal{F}$, que sa formation ne fait intervenir que les recouvrements et
que les recouvrements dans $Y_{Zar}$ ne font intervenir que des inclusions
entre ouverts d'un même $Y_n$; le résultat découle du fait que $f_n^{-1}\mathcal{F}_n$
est le faisceau associé à $f_{n, p}\mathcal{F}_n$, voir
Faisceaux, section \ref{sheaves-section-presheaves-functorial}.
\end{proof}

\noindent
Soit $X$ un espace topologique. Dans
Sites, exemple \ref{sites-example-site-topological},
nous avons noté $X_{Zar}$ le site dont les objets sont les ouverts de $X$,
les morphismes les inclusions et les recouvrements les recouvrements ouverts.
Nous identifions le topos $\Sh(X_{Zar})$ à la catégorie
des faisceaux sur $X$.

\begin{lemma}
\label{spaces-simplicial-lemma-restriction-to-components}
Soit $X$ un espace simplicial. Le foncteur
$X_{n, Zar} \to X_{Zar}$, $U \mapsto U$, est continu
et cocontinu. Le morphisme de topos associé
$g_n : \Sh(X_n) \to \Sh(X_{Zar})$ vérifie les propriétés suivantes :
\begin{enumerate}
\item $g_n^{-1}$ associe au faisceau $\mathcal{F}$ sur $X$
le faisceau $\mathcal{F}_n$ sur $X_n$,
\item $g_n^{-1} : \Sh(X_{Zar}) \to \Sh(X_n)$ admet un adjoint à gauche $g^{Sh}_{n!}$,
\item $g^{Sh}_{n!}$ commute aux limites finies connexes,
\item $g_n^{-1} : \textit{Ab}(X_{Zar}) \to \textit{Ab}(X_n)$
admet un adjoint à gauche $g_{n!}$, et
\item $g_{n!}$ est exact.
\end{enumerate}
\end{lemma}

\begin{proof}
Outre les propriétés de notre foncteur mentionnées dans l'énoncé,
la catégorie $X_{n, Zar}$ possède des produits fibrés et des égalisateurs,
et le foncteur commute à ceux-ci (attention, $X_{Zar}$ ne possède pas
tous les produits fibrés). Le lemme résulte donc de la discussion de
Sites, sections \ref{sites-section-cocontinuous-functors} et
\ref{sites-section-cocontinuous-morphism-topoi}
et
Modules sur les sites, section \ref{sites-modules-section-exactness-lower-shriek}.
Plus précisément,
Sites, lemmes \ref{sites-lemma-cocontinuous-morphism-topoi},
\ref{sites-lemma-when-shriek} et
\ref{sites-lemma-preserve-equalizers},
ainsi que de
Modules sur les sites, lemmes
\ref{sites-modules-lemma-g-shriek-adjoint} et
\ref{sites-modules-lemma-exactness-lower-shriek}.
\end{proof}

\begin{lemma}
\label{spaces-simplicial-lemma-restriction-injective-to-component}
Soit $X$ un espace simplicial. Si $\mathcal{I}$ est un faisceau abélien
injectif sur $X_{Zar}$, alors $\mathcal{I}_n$ est un faisceau abélien injectif
sur $X_n$.
\end{lemma}

\begin{proof}
Ceci résulte de
Homologie, lemme \ref{homology-lemma-adjoint-preserve-injectives}
et du
lemme \ref{spaces-simplicial-lemma-restriction-to-components}.
\end{proof}

\begin{lemma}
\label{spaces-simplicial-lemma-restriction-to-components-functorial}
Soit $f : Y \to X$ un morphisme d'espaces simpliciaux. Alors
$$
\xymatrix{
\Sh(Y_n) \ar[d] \ar[r]_{f_n} & \Sh(X_n) \ar[d] \\
\Sh(Y_{Zar}) \ar[r]^{f_{Zar}} & \Sh(X_{Zar})
}
$$
est un diagramme commutatif de topos.
\end{lemma}

\begin{proof}
Cela résulte directement de la description des foncteurs image inverse dans les
lemmes \ref{spaces-simplicial-lemma-describe-functoriality} et
\ref{spaces-simplicial-lemma-restriction-to-components}.
\end{proof}

\begin{lemma}
\label{spaces-simplicial-lemma-augmentation}
Soit $Y$ un espace simplicial et soit $a : Y \to X$ une augmentation
(Méthodes simpliciales, définition \ref{simplicial-definition-augmentation}).
Soient $a_n : Y_n \to X$ les morphismes d'espaces topologiques correspondants.
Il existe un morphisme canonique de topos
$$
a : \Sh(Y_{Zar}) \to \Sh(X)
$$
qui possède les propriétés suivantes :
\begin{enumerate}
\item $a^{-1}\mathcal{F}$ est le faisceau dont la restriction est $a_n^{-1}\mathcal{F}$
sur $Y_n$,
\item $a_m \circ Y(\varphi) = a_n$ pour tout $\varphi : [m] \to [n]$,
\item $a \circ g_n = a_n$ comme morphismes de topos, où
$g_n$ est défini au lemme \ref{spaces-simplicial-lemma-restriction-to-components},
\item $a_*\mathcal{G}$, pour $\mathcal{G} \in \Sh(Y_{Zar})$,
est l'égalisateur des deux applications
$a_{0, *}\mathcal{G}_0 \to a_{1, *}\mathcal{G}_1$.
\end{enumerate}
\end{lemma}

\begin{proof}
La partie (2) vaut pour les augmentations d'objets simpliciaux de toute catégorie.
Ainsi, $Y(\varphi)^{-1} a_m^{-1} \mathcal{F} = a_n^{-1}\mathcal{F}$,
ce qui définit un $Y(\varphi)$-morphisme de $a_m^{-1}\mathcal{F}$
vers $a_n^{-1}\mathcal{F}$.
On peut donc prendre (1) pour définition de $a^{-1}\mathcal{F}$ (en utilisant le
lemme \ref{spaces-simplicial-lemma-describe-sheaves-simplicial-site}) et
(4) pour définition de $a_*$. Si ceci définit un morphisme de topos,
alors la partie (3) en résulte, car on aura $g_n^{-1} \circ a^{-1} = a_n^{-1}$
par construction. Pour vérifier que $a$ est un morphisme de topos, il faut montrer
que $a^{-1}$ est adjoint à gauche de $a_*$ et
que $a^{-1}$ est exact. Ce dernier fait résulte immédiatement de l'exactitude des
foncteurs $a_n^{-1}$.

\medskip\noindent
Soit $\mathcal{F}$ un objet de $\Sh(X)$ et soit $\mathcal{G}$
un objet de $\Sh(Y_{Zar})$. Étant donné
$\beta : a^{-1}\mathcal{F} \to \mathcal{G}$, considérons ses
composantes $\beta_n : a_n^{-1}\mathcal{F} \to \mathcal{G}_n$.
Ces applications sont adjointes aux applications
$\beta_n : \mathcal{F} \to a_{n, *}\mathcal{G}_n$.
La compatibilité à la structure simpliciale montre que
$\beta_0$ se factorise par $a_*\mathcal{G}$.
Réciproquement, supposons donnée une application $\alpha : \mathcal{F} \to a_*\mathcal{G}$.
Pour tout $n$, choisissons un $\varphi : [0] \to [n]$. On peut alors considérer
la composée
$$
\mathcal{F} \xrightarrow{\alpha} a_*\mathcal{G}
\to a_{0, *}\mathcal{G}_0 \xrightarrow{\mathcal{G}(\varphi)}
a_{n, *}\mathcal{G}_n
$$
Ses adjointes sont des applications $a_n^{-1}\mathcal{F} \to \mathcal{G}_n$
qui définissent un morphisme de faisceaux $a^{-1}\mathcal{F} \to \mathcal{G}$.
Nous omettons la vérification que les constructions ci-dessus définissent
des bijections inverses l'une de l'autre
$$
\Mor_{\Sh(Y_{Zar})}(a^{-1}\mathcal{F}, \mathcal{G}) =
\Mor_{\Sh(X)}(\mathcal{F}, a_*\mathcal{G})
$$
Ceci achève la démonstration. Une observation intéressante est que
ce morphisme de topos ne correspond à aucun foncteur géométrique évident
entre les sites qui définissent les topos.
\end{proof}

\begin{lemma}
\label{spaces-simplicial-lemma-simplicial-resolution-Z}
Soit $X$ un espace topologique simplicial. Le complexe de
préfaisceaux abéliens sur $X_{Zar}$
$$
\ldots \to \mathbf{Z}_{X_2} \to \mathbf{Z}_{X_1} \to \mathbf{Z}_{X_0}
$$
dont le différentiel est $\sum (-1)^i d^n_i$ est une résolution
du préfaisceau constant $\mathbf{Z}$.
\end{lemma}

\begin{proof}
Soit $U \subset X_m$ un objet de $X_{Zar}$. La valeur
du complexe ci-dessus sur $U$ est le complexe de groupes abéliens
$$
\ldots \to
\mathbf{Z}[\Mor_\Delta([2], [m])] \to
\mathbf{Z}[\Mor_\Delta([1], [m])] \to
\mathbf{Z}[\Mor_\Delta([0], [m])]
$$
Autrement dit, c'est le complexe associé au
groupe abélien libre engendré par l'ensemble simplicial $\Delta[m]$, voir
Méthodes simpliciales, exemple \ref{simplicial-example-simplex-simplicial-set}.
Puisque $\Delta[m]$ est homotopiquement équivalent à $\Delta[0]$, voir
Méthodes simpliciales, exemple \ref{simplicial-example-simplex-contractible},
et puisque ``prendre les groupes abéliens libres'' est un foncteur,
on voit que le complexe ci-dessus est homotopiquement équivalent au
groupe abélien libre engendré par $\Delta[0]$
(Méthodes simpliciales, remarque \ref{simplicial-remark-homotopy-better} et
lemme \ref{simplicial-lemma-homotopy-equivalence-s-N}).
Ce complexe est acyclique en degrés strictement positifs
et égal à $\mathbf{Z}$ en degré $0$.
\end{proof}

\begin{lemma}
\label{spaces-simplicial-lemma-simplicial-sheaf-cohomology}
Soit $X$ un espace topologique simplicial. Soit $\mathcal{F}$ un faisceau abélien
sur $X$. Il existe une suite spectrale $(E_r, d_r)_{r \geq 0}$ telle que
$$
E_1^{p, q} = H^q(X_p, \mathcal{F}_p)
$$
et qui aboutit à $H^{p + q}(X_{Zar}, \mathcal{F})$.
Cette suite spectrale est fonctorielle en $\mathcal{F}$.
\end{lemma}

\begin{proof}
Soit $\mathcal{F} \to \mathcal{I}^\bullet$ une résolution injective.
Considérons le complexe double de termes
$$
A^{p, q} = \mathcal{I}^q(X_p)
$$
et dont le premier différentiel est donné par la somme alternée des morphismes
$d^{p + 1}_i$ induits $\mathcal{I}_p^q \to \mathcal{I}_{p + 1}^q$, voir
le lemme \ref{spaces-simplicial-lemma-describe-sheaves-simplicial-site}. Remarquons que
$$
A^{p, q} = \Gamma(X_p, \mathcal{I}_p^q) =
\Mor_{\textit{PSh}}(h_{X_p}, \mathcal{I}^q) =
\Mor_{\textit{PAb}}(\mathbf{Z}_{X_p}, \mathcal{I}^q)
$$
Il résulte donc du lemme \ref{spaces-simplicial-lemma-simplicial-resolution-Z} et de
Cohomologie sur les sites, lemme
\ref{sites-cohomology-lemma-injective-abelian-sheaf-injective-presheaf}
que les lignes du complexe double sont exactes en degrés strictement positifs et
ont pour valeur $\Gamma(X_{Zar}, \mathcal{I}^q)$ en degré $0$.
D'autre part, puisque la restriction est exacte
(lemme \ref{spaces-simplicial-lemma-restriction-to-components}),
l'application
$$
\mathcal{F}_p \to \mathcal{I}_p^\bullet
$$
est une résolution. Les faisceaux $\mathcal{I}_p^q$ sont des faisceaux
abéliens injectifs sur $X_p$
(lemme \ref{spaces-simplicial-lemma-restriction-injective-to-component}).
La cohomologie des colonnes calcule donc les groupes
$H^q(X_p, \mathcal{F}_p)$. On conclut en appliquant
Homologie, lemmes \ref{homology-lemma-first-quadrant-ss} et
\ref{homology-lemma-double-complex-gives-resolution}.
\end{proof}

\begin{lemma}
\label{spaces-simplicial-lemma-augmentation-pushforward-higher-direct-image}
Soit $X$ un espace simplicial et soit $a : X \to Y$
une augmentation. Soit $\mathcal{F}$ un faisceau abélien
sur $X_{Zar}$. Alors $R^na_*\mathcal{F}$ est le faisceau associé
au préfaisceau
$$
V \longmapsto H^n((X \times_Y V)_{Zar}, \mathcal{F}|_{(X \times_Y V)_{Zar}})
$$
\end{lemma}

\begin{proof}
C'est l'analogue de
Cohomologie, lemme \ref{cohomology-lemma-describe-higher-direct-images} ou de
Cohomologie sur les sites, lemme \ref{sites-cohomology-lemma-higher-direct-images},
et nous conseillons vivement au lecteur d'omettre la démonstration.
En choisissant une résolution injective de $\mathcal{F}$ sur
$X_{Zar}$ et en appliquant les définitions, on voit qu'il suffit de montrer :
(1) que la restriction d'un faisceau abélien injectif
sur $X_{Zar}$ à $(X \times_Y V)_{Zar}$ est un faisceau abélien injectif, et
(2) que $a_*\mathcal{F}$ est égal à la règle
$$
V \longmapsto H^0((X \times_Y V)_{Zar}, \mathcal{F}|_{(X \times_Y V)_{Zar}})
$$
La partie (2) résulte des faits suivants :
\begin{enumerate}
\item[(2a)] $a_*\mathcal{F}$ est l'égalisateur des deux applications
$a_{0, *}\mathcal{F}_0 \to a_{1, *}\mathcal{F}_1$
d'après le lemme \ref{spaces-simplicial-lemma-augmentation},
\item[(2b)] $a_{0, *}\mathcal{F}_0(V) =
H^0(a_0^{-1}(V), \mathcal{F}_0)$ et
$a_{1, *}\mathcal{F}_1(V) = H^0(a_1^{-1}(V), \mathcal{F}_1)$,
\item[(2c)] $X_0 \times_Y V = a_0^{-1}(V)$ et $X_1 \times_Y V = a_1^{-1}(V)$,
\item[(2d)] $H^0((X \times_Y V)_{Zar}, \mathcal{F}|_{(X \times_Y V)_{Zar}})$
est l'égalisateur des deux applications
$H^0(X_0 \times_Y V, \mathcal{F}_0) \to H^0(X_1 \times_Y V, \mathcal{F}_1)$,
par exemple d'après le lemme \ref{spaces-simplicial-lemma-simplicial-sheaf-cohomology}.
\end{enumerate}
La partie (1) résulte de ce que l'on définit un adjoint à gauche exact
$j_! : \textit{Ab}((X \times_Y V)_{Zar}) \to \textit{Ab}(X_{Zar})$
(extension par zéro) de la restriction
$\textit{Ab}(X_{Zar}) \to \textit{Ab}((X \times_Y V)_{Zar})$,
et que l'on applique Homologie, lemme \ref{homology-lemma-adjoint-preserve-injectives}.
\end{proof}

\noindent
Soit $X$ un espace topologique. Notons $X_\bullet$ l'espace topologique
simplicial constant de valeur $X$. D'après le
lemme \ref{spaces-simplicial-lemma-describe-sheaves-simplicial-site},
se donner un faisceau sur $X_{\bullet, Zar}$ revient à se
donner un objet cosimplicial de la catégorie des faisceaux sur $X$.

\begin{lemma}
\label{spaces-simplicial-lemma-constant-simplicial-space}
Soit $X$ un espace topologique. Soit $X_\bullet$ l'espace topologique
simplicial constant de valeur $X$. Le foncteur
$$
X_{\bullet, Zar} \longrightarrow X_{Zar},\quad
U \longmapsto U
$$
est continu et cocontinu et définit un morphisme de
topos $g : \Sh(X_{\bullet, Zar}) \to \Sh(X)$, ainsi qu'un adjoint à gauche
$g_!$ de $g^{-1}$. On a
\begin{enumerate}
\item $g^{-1}$ associe à un faisceau sur $X$ le faisceau cosimplicial
constant sur $X$,
\item $g_!$ associe à un faisceau $\mathcal{F}$ sur $X_{\bullet, Zar}$ le
faisceau $\mathcal{F}_0$, et
\item $g_*$ associe à un faisceau $\mathcal{F}$ sur $X_{\bullet, Zar}$
l'égalisateur des deux applications $\mathcal{F}_0 \to \mathcal{F}_1$.
\end{enumerate}
\end{lemma}

\begin{proof}
Les assertions concernant le foncteur se vérifient immédiatement.
L'existence de $g$ et de $g_!$ résulte de
Sites, lemmes \ref{sites-lemma-cocontinuous-morphism-topoi} et
\ref{sites-lemma-when-shriek}. La description de
$g^{-1}$ résulte immédiatement de Sites, lemme \ref{sites-lemma-when-shriek}.
La description de $g_*$ et de $g_!$ résulte de ce que les foncteurs indiqués sont
respectivement adjoint à droite et adjoint à gauche de $g^{-1}$.
\end{proof}








\section{Sites simpliciaux et topos}
\label{spaces-simplicial-section-simplicial-sites}

\noindent
Il paraît naturel de définir un {\it site simplicial} comme un objet
simplicial de la (grande) catégorie dont les objets sont les sites
et dont les morphismes sont les morphismes de sites.
Voir Sites, définitions \ref{sites-definition-site} et
\ref{sites-definition-morphism-sites},
la composition des morphismes étant celle de
Sites, lemme \ref{sites-lemma-composition-morphisms-sites}.
Mais voici quelques variantes que l'on pourrait vouloir considérer :
(a) on pourrait plutôt travailler avec des foncteurs cocontinus
(voir Sites, sections \ref{sites-section-cocontinuous-functors} et
\ref{sites-section-cocontinuous-morphism-topoi}) entre sites,
(b) on pourrait travailler dans une $2$-catégorie convenable de sites où l'on introduit
la notion de $2$-morphisme entre morphismes de sites,
(c) on pourrait travailler dans une $2$-catégorie construite à partir de foncteurs
cocontinus. Plutôt que de choisir l'une de ces variantes comme définition,
nous développerons simplement la théorie au fur et à mesure des besoins.

\medskip\noindent
Un {\it topos simplicial} devrait sans doute être défini comme un
pseudo-foncteur de $\Delta^{opp}$ vers la $2$-catégorie des topos.
Voir Catégories, définition \ref{categories-definition-functor-into-2-category}
et Sites, section \ref{sites-section-topoi} et
\ref{sites-section-2-category}. Nous essaierons d'éviter, si possible,
de travailler avec un tel objet.

\medskip\noindent
{\bf Cas A.}
Soit $\mathcal{C}$ un objet simplicial de la catégorie dont les objets
sont les sites et dont les morphismes sont les morphismes de sites. Cela signifie que,
pour tout morphisme $\varphi : [m] \to [n]$ de $\Delta$, on dispose d'un morphisme
de sites $f_\varphi : \mathcal{C}_n \to \mathcal{C}_m$. Ce morphisme est
donné par un foncteur continu en sens opposé, que l'on notera
$u_\varphi : \mathcal{C}_m \to \mathcal{C}_n$.

\begin{lemma}
\label{spaces-simplicial-lemma-simplicial-site-site}
Soit $\mathcal{C}$ un objet simplicial de la catégorie des sites.
Avec les notations ci-dessus, construisons un site $\mathcal{C}_{total}$ comme suit.
\begin{enumerate}
\item Un objet de $\mathcal{C}_{total}$ est un objet $U$ de
$\mathcal{C}_n$ pour un certain $n$,
\item un morphisme $(\varphi, f) : U \to V$ de $\mathcal{C}_{total}$
est donné par une application $\varphi : [m] \to [n]$ avec
$U \in \Ob(\mathcal{C}_n)$, $V \in \Ob(\mathcal{C}_m)$,
et un morphisme $f : U \to u_\varphi(V)$ de $\mathcal{C}_n$, et
\item un recouvrement $\{(\text{id}, f_i) :  U_i \to U\}$ dans $\mathcal{C}_{total}$
est donné par un entier $n$ et un recouvrement $\{f_i : U_i \to U\}$
de $\mathcal{C}_n$.
\end{enumerate}
\end{lemma}

\begin{proof}
La composée de $(\varphi, f) : U \to V$ et de $(\psi, g) : V \to W$
est donnée par $(\varphi \circ \psi, u_\varphi(g) \circ f)$.
On utilise ici l'égalité $u_\varphi \circ u_\psi = u_{\varphi \circ \psi}$.

\medskip\noindent
Soit $\{(\text{id}, f_i) :  U_i \to U\}$ un recouvrement comme en (3),
et soit $(\varphi, g) : W \to U$ un morphisme avec
$W \in \Ob(\mathcal{C}_m)$. Nous affirmons que
$$
W \times_{(\varphi, g), U, (\text{id}, f_i)} U_i =
W \times_{g, u_\varphi(U), u_\varphi(f_i)} u_\varphi(U_i)
$$
dans la catégorie $\mathcal{C}_{total}$. Cela a un sens car, d'après notre
définition des morphismes de sites, les produits fibrés requis
existent dans $\mathcal{C}_m$, puisque $u_\varphi$ transforme les recouvrements en
recouvrements. Le même raisonnement démontre l'affirmation (détails omis).
On voit donc que la famille des recouvrements est stable par changement
de base. Les autres axiomes d'un site sont immédiats.
\end{proof}

\noindent
{\bf Cas B.}
Soit $\mathcal{C}$ un objet simplicial de la catégorie dont les objets sont les
sites et dont les morphismes sont les foncteurs cocontinus. Cela signifie que, pour
tout morphisme $\varphi : [m] \to [n]$ de $\Delta$, on dispose d'un foncteur cocontinu
noté $u_\varphi : \mathcal{C}_n \to \mathcal{C}_m$. Le morphisme de topos
associé est noté
$f_\varphi : \Sh(\mathcal{C}_n) \to \Sh(\mathcal{C}_m)$.

\begin{lemma}
\label{spaces-simplicial-lemma-simplicial-cocontinuous-site}
Soit $\mathcal{C}$ un objet simplicial de la catégorie dont les objets sont les
sites et dont les morphismes sont les foncteurs cocontinus. Avec les notations ci-dessus,
supposons que les foncteurs $u_\varphi : \mathcal{C}_n \to \mathcal{C}_m$
possèdent la propriété $P$ de Sites, remarque \ref{sites-remark-cartesian-cocontinuous}.
On peut alors construire un site $\mathcal{C}_{total}$ comme suit.
\begin{enumerate}
\item Un objet de $\mathcal{C}_{total}$ est un objet $U$ de
$\mathcal{C}_n$ pour un certain $n$,
\item un morphisme $(\varphi, f) : U \to V$ de $\mathcal{C}_{total}$
est donné par une application $\varphi : [m] \to [n]$ avec
$U \in \Ob(\mathcal{C}_n)$, $V \in \Ob(\mathcal{C}_m)$,
et un morphisme $f : u_\varphi(U) \to V$ de $\mathcal{C}_m$, et
\item un recouvrement $\{(\text{id}, f_i) :  U_i \to U\}$ dans $\mathcal{C}_{total}$
est donné par un entier $n$ et un recouvrement $\{f_i : U_i \to U\}$
de $\mathcal{C}_n$.
\end{enumerate}
\end{lemma}

\begin{proof}
La composée de $(\varphi, f) : U \to V$ et de $(\psi, g) : V \to W$
est donnée par $(\varphi \circ \psi, g \circ u_\psi(f))$.
On utilise ici l'égalité $u_\psi \circ u_\varphi = u_{\varphi \circ \psi}$.

\medskip\noindent
Soit $\{(\text{id}, f_i) :  U_i \to U\}$ un recouvrement comme en (3),
et soit $(\varphi, g) : W \to U$ un morphisme avec
$W \in \Ob(\mathcal{C}_m)$. Nous affirmons que
$$
W \times_{(\varphi, g), U, (\text{id}, f_i)} U_i =
W \times_{g, U, f_i} U_i
$$
dans la catégorie $\mathcal{C}_{total}$, où le membre de droite
est l'objet de $\mathcal{C}_m$ défini dans
Sites, remarque \ref{sites-remark-cartesian-cocontinuous},
qui existe d'après la propriété $P$. La compatibilité de ce type de produit fibré
avec la composition des foncteurs démontre l'affirmation (détails omis).
Comme la famille $\{W \times_{g, U, f_i} U_i \to W\}$ est un
recouvrement de $\mathcal{C}_m$ d'après la propriété $P$, on voit que
la famille des recouvrements est stable par changement
de base. Les autres axiomes d'un site sont immédiats.
\end{proof}

\begin{situation}
\label{spaces-simplicial-situation-simplicial-site}
On se trouve ici dans l'un des deux cas suivants :
\begin{enumerate}
\item[(A)] $\mathcal{C}$ est un objet simplicial de la catégorie dont les
objets sont les sites et dont les morphismes sont les morphismes de sites. Pour tout
morphisme $\varphi : [m] \to [n]$ de $\Delta$, on dispose d'un morphisme de sites
$f_\varphi : \mathcal{C}_n \to \mathcal{C}_m$ donné par un foncteur continu
$u_\varphi : \mathcal{C}_m \to \mathcal{C}_n$.
\item[(B)] $\mathcal{C}$ est un objet simplicial de la catégorie dont les
objets sont les sites et dont les morphismes sont des foncteurs cocontinus possédant la
propriété $P$ de Sites, remarque \ref{sites-remark-cartesian-cocontinuous}.
Pour tout morphisme $\varphi : [m] \to [n]$ de $\Delta$, on dispose d'un foncteur cocontinu
$u_\varphi : \mathcal{C}_n \to \mathcal{C}_m$ qui induit un
morphisme de topos $f_\varphi : \Sh(\mathcal{C}_n) \to \Sh(\mathcal{C}_m)$.
\end{enumerate}
Comme d'habitude, on notera $f_\varphi^{-1}$ et $f_{\varphi, *}$ les foncteurs
image inverse et image directe. Notons $\mathcal{C}_{total}$ le
site défini au
lemme \ref{spaces-simplicial-lemma-simplicial-site-site} (cas A) ou au
lemme \ref{spaces-simplicial-lemma-simplicial-cocontinuous-site} (cas B).
\end{situation}

\noindent
Soit $\mathcal{C}$ comme dans la situation \ref{spaces-simplicial-situation-simplicial-site}.
Soit $\mathcal{F}$ un faisceau sur $\mathcal{C}_{total}$.
Il résulte clairement de la définition des recouvrements que la restriction
de $\mathcal{F}$ aux objets de $\mathcal{C}_n$ définit un faisceau
$\mathcal{F}_n$ sur le site $\mathcal{C}_n$. Pour tout
$\varphi : [m] \to [n]$, les applications de restriction de $\mathcal{F}$
le long des morphismes $(\varphi, f) : U \to V$ avec
$U \in \Ob(\mathcal{C}_n)$ et $V \in \Ob(\mathcal{C}_m)$
définissent un élément $\mathcal{F}(\varphi)$ de
$$
\Mor_{\Sh(\mathcal{C}_m)}(\mathcal{F}_m, f_{\varphi, *}\mathcal{F}_n) =
\Mor_{\Sh(\mathcal{C}_n)}(f_\varphi^{-1}\mathcal{F}_m, \mathcal{F}_n)
$$
De plus, étant donnés $\varphi : [m] \to [n]$ et $\psi : [l] \to [m]$,
les diagrammes
$$
\vcenter{
\xymatrix{
\mathcal{F}_l \ar[rr]_{\mathcal{F}(\varphi \circ \psi)}
\ar[rd]_{\mathcal{F}(\psi)}
& & f_{\varphi \circ \psi, *} \mathcal{F}_n \\
& f_{\psi, *}\mathcal{F}_m \ar[ur]_{f_{\psi, *}\mathcal{F}(\varphi)}
}
}
\quad\text{et}\quad
\vcenter{
\xymatrix{
f_{\varphi \circ \psi}^{-1}\mathcal{F}_l
\ar[rr]_{\mathcal{F}(\varphi \circ \psi)}
\ar[rd]_{f_\varphi^{-1}\mathcal{F}(\psi)}
& & \mathcal{F}_n \\
& f_\varphi^{-1}\mathcal{F}_m \ar[ur]_{\mathcal{F}(\varphi)}
}
}
$$
sont commutatifs. La réciproque est manifestement vraie elle aussi : si l'on a un système
$(\{\mathcal{F}_n\}_{n \geq 0},
\{\mathcal{F}(\varphi)\}_{\varphi \in \text{Flèches}(\Delta)})$
satisfaisant les conditions de commutativité ci-dessus,
on obtient alors un faisceau sur $\mathcal{C}_{total}$.

\begin{lemma}
\label{spaces-simplicial-lemma-describe-sheaves-simplicial-site-site}
Dans la situation \ref{spaces-simplicial-situation-simplicial-site}, il existe une équivalence de
catégories entre
\begin{enumerate}
\item $\Sh(\mathcal{C}_{total})$, et
\item la catégorie des systèmes $(\mathcal{F}_n, \mathcal{F}(\varphi))$
décrits ci-dessus.
\end{enumerate}
En particulier, le topos $\Sh(\mathcal{C}_{total})$ ne dépend que des
topos $\Sh(\mathcal{C}_n)$ et des morphismes de topos $f_\varphi$.
\end{lemma}

\begin{proof}
Voir la discussion ci-dessus.
\end{proof}

\begin{lemma}
\label{spaces-simplicial-lemma-restriction-to-components-site}
Dans la situation \ref{spaces-simplicial-situation-simplicial-site}, le foncteur
$\mathcal{C}_n \to \mathcal{C}_{total}$, $U \mapsto U$, est continu
et cocontinu. Le morphisme de
topos associé $g_n : \Sh(\mathcal{C}_n) \to \Sh(\mathcal{C}_{total})$ vérifie
\begin{enumerate}
\item $g_n^{-1}$ associe au faisceau $\mathcal{F}$ sur $\mathcal{C}_{total}$
le faisceau $\mathcal{F}_n$ sur $\mathcal{C}_n$,
\item $g_n^{-1} : \Sh(\mathcal{C}_{total}) \to \Sh(\mathcal{C}_n)$
admet un adjoint à gauche $g^{Sh}_{n!}$,
\item pour $\mathcal{G}$ dans $\Sh(\mathcal{C}_n)$, la restriction de
$g_{n!}^{Sh}\mathcal{G}$ à $\mathcal{C}_m$ est
$\coprod\nolimits_{\varphi : [n] \to [m]} f_\varphi^{-1}\mathcal{G}$,
\item $g_{n!}^{Sh}$ commute aux limites finies connexes,
\item $g_n^{-1} : \textit{Ab}(\mathcal{C}_{total}) \to
\textit{Ab}(\mathcal{C}_n)$ admet un adjoint à gauche $g_{n!}$,
\item pour $\mathcal{G}$ dans $\textit{Ab}(\mathcal{C}_n)$, la restriction de
$g_{n!}\mathcal{G}$ à $\mathcal{C}_m$ est
$\bigoplus\nolimits_{\varphi : [n] \to [m]} f_\varphi^{-1}\mathcal{G}$, et
\item $g_{n!}$ est exact.
\end{enumerate}
\end{lemma}

\begin{proof}
Cas A. Si $\{U_i \to U\}_{i \in I}$ est un recouvrement dans $\mathcal{C}_n$,
alors son image $\{U_i \to U\}_{i \in I}$ est un recouvrement dans $\mathcal{C}_{total}$
par définition (lemme \ref{spaces-simplicial-lemma-simplicial-site-site}). Pour un morphisme
$V \to U$ de $\mathcal{C}_n$, le produit fibré
$V \times_U U_i$ dans $\mathcal{C}_n$ est aussi
le produit fibré dans $\mathcal{C}_{total}$ (d'après l'affirmation de la
démonstration du lemme \ref{spaces-simplicial-lemma-simplicial-site-site}).
Notre foncteur est donc continu. D'autre part, il
définit une bijection entre les recouvrements de $U$ dans $\mathcal{C}_n$
et les recouvrements de $U$ dans $\mathcal{C}_{total}$. Il est donc
assurément cocontinu.

\medskip\noindent
Cas B. Si $\{U_i \to U\}_{i \in I}$ est un recouvrement dans $\mathcal{C}_n$,
alors son image $\{U_i \to U\}_{i \in I}$ est un recouvrement dans $\mathcal{C}_{total}$
par définition (lemme \ref{spaces-simplicial-lemma-simplicial-cocontinuous-site}). Pour un morphisme
$V \to U$ de $\mathcal{C}_n$, le produit fibré
$V \times_U U_i$ dans $\mathcal{C}_n$ est aussi
le produit fibré dans $\mathcal{C}_{total}$ (d'après l'affirmation de la
démonstration du lemme \ref{spaces-simplicial-lemma-simplicial-cocontinuous-site}).
Notre foncteur est donc continu. D'autre part, il
définit une bijection entre les recouvrements de $U$ dans $\mathcal{C}_n$
et les recouvrements de $U$ dans $\mathcal{C}_{total}$. Il est donc
assurément cocontinu.

\medskip\noindent
À ce stade, la partie (1) et l'existence de $g^{Sh}_{n!}$ et de $g_{n!}$
dans les cas A et B résultent de
Sites, lemmes \ref{sites-lemma-cocontinuous-morphism-topoi} et
\ref{sites-lemma-when-shriek},
ainsi que de
Modules sur les sites, lemme \ref{sites-modules-lemma-g-shriek-adjoint}.

\medskip\noindent
Démonstration de (3). Soit $\mathcal{G}$ un faisceau sur $\mathcal{C}_n$.
Considérons le faisceau $\mathcal{H}$ sur $\mathcal{C}_{total}$
dont le terme de degré $m$ est le faisceau
$$
\mathcal{H}_m = \coprod\nolimits_{\varphi : [n] \to [m]}
f_\varphi^{-1}\mathcal{G}
$$
donné dans la partie (3) de l'énoncé du lemme.
Étant donnée une application $\psi : [m] \to [m']$, l'application
$\mathcal{H}(\psi) : f_\psi^{-1}\mathcal{H}_m \to \mathcal{H}_{m'}$
est donnée sur les composantes par les identifications
$$
f_\psi^{-1} f_\varphi^{-1} \mathcal{G} \to
f_{\psi \circ \varphi}^{-1}\mathcal{G}
$$
Observons qu'une application $\alpha : \mathcal{H} \to \mathcal{F}$
de faisceaux sur $\mathcal{C}_{total}$ fournit une application
$\mathcal{G} \to \mathcal{F}_n$
qui correspond à la restriction de $\alpha_n$ à la composante
$\mathcal{G}$ de $\mathcal{H}_n$. Réciproquement, étant donnée une application
$\beta : \mathcal{G} \to \mathcal{F}_n$ de faisceaux sur $\mathcal{C}_n$,
on peut définir
$\alpha : \mathcal{H} \to \mathcal{F}$ en prenant pour $\alpha_m$
l'application qui, sur les composantes
$$
f_\varphi^{-1}\mathcal{G} \to \mathcal{F}_m
$$
utilise les applications adjointes à $\mathcal{F}(\varphi) \circ f_\varphi^{-1}\beta$.
Nous omettons la vérification que ces deux constructions donnent
des applications inverses l'une de l'autre
$$
\Mor_{\Sh(\mathcal{C}_n)}(\mathcal{G}, \mathcal{F}_n) =
\Mor_{\Sh(\mathcal{C}_{total})}(\mathcal{H}, \mathcal{F})
$$
Ainsi, $\mathcal{H} = g^{Sh}_{n!}\mathcal{G}$, comme voulu.

\medskip\noindent
Démonstration de (4). Si $\mathcal{G}$ est un faisceau abélien sur $\mathcal{C}_n$,
on procède exactement de la même manière que ci-dessus, à ceci près que
l'on définit $\mathcal{H}$ comme le faisceau abélien sur $\mathcal{C}_{total}$
dont le terme de degré $m$ est le faisceau
$$
\bigoplus\nolimits_{\varphi : [n] \to [m]} f_\varphi^{-1}\mathcal{G}
$$
avec des applications de transition définies exactement comme ci-dessus. La bijection
$$
\Mor_{\textit{Ab}(\mathcal{C}_n)}(\mathcal{G}, \mathcal{F}_n) =
\Mor_{\textit{Ab}(\mathcal{C}_{total})}(\mathcal{H}, \mathcal{F})
$$
se démontre exactement comme ci-dessus.
Ainsi, $\mathcal{H} = g_{n!}\mathcal{G}$, comme voulu.

\medskip\noindent
Les propriétés d'exactitude de $g^{Sh}_{n!}$ et de $g_{n!}$ résultent
des formules données pour ces foncteurs.
\end{proof}

\begin{lemma}
\label{spaces-simplicial-lemma-restriction-injective-to-component-site}
\begin{slogan}
Un faisceau abélien injectif sur un site simplicial est injectif sur chaque composante
\end{slogan}
Dans la situation \ref{spaces-simplicial-situation-simplicial-site},
si $\mathcal{I}$ est injectif dans $\textit{Ab}(\mathcal{C}_{total})$,
alors $\mathcal{I}_n$ est injectif dans $\textit{Ab}(\mathcal{C}_n)$.
Si $\mathcal{I}^\bullet$ est un complexe K-injectif dans
$\textit{Ab}(\mathcal{C}_{total})$,
alors $\mathcal{I}_n^\bullet$ est K-injectif dans $\textit{Ab}(\mathcal{C}_n)$.
\end{lemma}

\begin{proof}
La première assertion résulte de
Homologie, lemme \ref{homology-lemma-adjoint-preserve-injectives}
et du
lemme \ref{spaces-simplicial-lemma-restriction-to-components-site}.
La seconde résulte de
Catégories dérivées, lemme \ref{derived-lemma-adjoint-preserve-K-injectives}
et du
lemme \ref{spaces-simplicial-lemma-restriction-to-components-site}.
\end{proof}







\section{Augmentations de sites simpliciaux}
\label{spaces-simplicial-section-augmentation-simplicial-sites}

\noindent
Nous poursuivons selon la méthode décrite dans la
section \ref{spaces-simplicial-section-simplicial-sites},
en explicitant le sens des augmentations dans les cas A et B
traités dans cette section.

\begin{remark}
\label{spaces-simplicial-remark-augmentation-site}
Dans la situation \ref{spaces-simplicial-situation-simplicial-site}, une
{\it augmentation $a_0$ vers un site $\mathcal{D}$} désignera
\begin{enumerate}
\item[(A)] un morphisme de sites $a_0 : \mathcal{C}_0 \to \mathcal{D}$
donné par un foncteur continu $u_0 : \mathcal{D} \to \mathcal{C}_0$
tel que, pour tous $\varphi, \psi : [0] \to [n]$, on ait
$u_\varphi \circ u_0 = u_\psi \circ u_0$.
\item[(B)] un morphisme $a_0 : \Sh(\mathcal{C}_0) \to \Sh(\mathcal{D})$
de topos donné par un foncteur cocontinu $u_0 : \mathcal{C}_0 \to \mathcal{D}$
tel que, pour tous $\varphi, \psi : [0] \to [n]$, on ait
$u_0 \circ u_\varphi = u_0 \circ u_\psi$.
\end{enumerate}
\end{remark}

\begin{lemma}
\label{spaces-simplicial-lemma-augmentation-site}
Dans la situation \ref{spaces-simplicial-situation-simplicial-site}, soit $a_0$ une
augmentation vers un site $\mathcal{D}$ au sens de la
remarque \ref{spaces-simplicial-remark-augmentation-site}. Alors $a_0$ induit
\begin{enumerate}
\item un morphisme de topos $a_n : \Sh(\mathcal{C}_n) \to \Sh(\mathcal{D})$
pour tout $n \geq 0$,
\item un morphisme de topos $a : \Sh(\mathcal{C}_{total}) \to \Sh(\mathcal{D})$
\end{enumerate}
tel que
\begin{enumerate}
\item pour tout $\varphi : [m] \to [n]$, on ait $a_m \circ f_\varphi = a_n$,
\item si $g_n : \Sh(\mathcal{C}_n) \to \Sh(\mathcal{C}_{total})$
est défini comme au lemme \ref{spaces-simplicial-lemma-restriction-to-components-site}, alors
$a \circ g_n = a_n$, et
\item $a_*\mathcal{F}$, pour $\mathcal{F} \in \Sh(\mathcal{C}_{total})$,
est l'égalisateur des deux applications
$a_{0, *}\mathcal{F}_0 \to a_{1, *}\mathcal{F}_1$.
\end{enumerate}
\end{lemma}

\begin{proof}
Cas A. Soit $u_n : \mathcal{D} \to \mathcal{C}_n$ la valeur commune
des foncteurs $u_\varphi \circ u_0$ pour $\varphi : [0] \to [n]$.
Alors $u_n$ correspond à un morphisme de sites
$a_n : \mathcal{C}_n \to \mathcal{D}$, voir
Sites, lemme \ref{sites-lemma-composition-morphisms-sites}.
Le même lemme montre que, pour tout $\varphi : [m] \to [n]$, on a
$a_m \circ f_\varphi = a_n$.

\medskip\noindent
Cas B. Soit $u_n : \mathcal{C}_n \to \mathcal{D}$ la valeur commune
des foncteurs $u_0 \circ u_\varphi$ pour $\varphi : [0] \to [n]$.
Alors $u_n$ est cocontinu et définit donc un morphisme de topos
$a_n : \Sh(\mathcal{C}_n) \to \Sh(\mathcal{D)}$, voir
Sites, lemme \ref{sites-lemma-composition-cocontinuous}.
Le même lemme montre que, pour tout $\varphi : [m] \to [n]$, on a
$a_m \circ f_\varphi = a_n$.

\medskip\noindent
Considérons le foncteur $a^{-1} : \Sh(\mathcal{D}) \to \Sh(\mathcal{C}_{total})$
qui, à un faisceau d'ensembles $\mathcal{G}$, associe le faisceau
$\mathcal{F} = a^{-1}\mathcal{G}$ dont les composantes sont $a_n^{-1}\mathcal{G}$
et dont les applications de transition $\mathcal{F}(\varphi)$ sont les identifications
$$
f_\varphi^{-1}\mathcal{F}_m =
f_\varphi^{-1} a_m^{-1}\mathcal{G} =
a_n^{-1}\mathcal{G} =
\mathcal{F}_n
$$
pour $\varphi : [m] \to [n]$, voir la description de
$\Sh(\mathcal{C}_{total})$ dans le
lemme \ref{spaces-simplicial-lemma-describe-sheaves-simplicial-site-site}.
Puisque les foncteurs $a_n^{-1}$ sont exacts, $a^{-1}$ est un foncteur exact.
Enfin, pour $a_* : \Sh(\mathcal{C}_{total}) \to \Sh(\mathcal{D})$,
on prend le foncteur qui, à un faisceau $\mathcal{F}$ sur $\Sh(\mathcal{D})$,
associe
$$
\xymatrix{
a_*\mathcal{F} \ar@{=}[r] &
\text{Égalisateur}(a_{0, *}\mathcal{F}_0
\ar@<1ex>[r] \ar@<-1ex>[r] &
a_{1, *}\mathcal{F}_1)
}
$$
Les deux applications proviennent ici des deux applications $\varphi : [0] \to [1]$
par
$$
a_{0, *}\mathcal{F}_0 \to
a_{0, *}f_{\varphi, *} f_\varphi^{-1}\mathcal{F}_0
\xrightarrow{\mathcal{F}(\varphi)}
a_{0, *}f_{\varphi, *} \mathcal{F}_1 = a_{1, *}\mathcal{F}_1
$$
où la première flèche provient de $1 \to f_{\varphi, *} f_\varphi^{-1}$.
Notons $\mathcal{G}_\bullet$ le faisceau cosimplicial constant
de valeur $\mathcal{G}$, et $a_{\bullet, *}\mathcal{F}$
le faisceau cosimplicial ayant $a_{n, *}\mathcal{F}_n$ en degré $n$.
Par l'adjonction usuelle pour les morphismes de topos $a_n$, on voit qu'une
application $a^{-1}\mathcal{G} \to \mathcal{F}$
revient à se donner une application
$$
\mathcal{G}_\bullet \longrightarrow a_{\bullet, *}\mathcal{F}
$$
de faisceaux cosimpliciaux.
Par dualité avec Méthodes simpliciales, lemme \ref{simplicial-lemma-augmentation-howto},
cela revient à se donner une application $\mathcal{G} \to a_*\mathcal{F}$.
Ainsi, $a^{-1}$ et $a_*$ sont des foncteurs adjoints, et l'on obtient
notre morphisme de topos $a$\footnote{Dans le cas B, le morphisme $a$
correspond au foncteur cocontinu
$\mathcal{C}_{total} \to \mathcal{D}$ qui envoie
$U$ dans $\mathcal{C}_n$ sur $u_n(U)$.}. Les égalités
$a \circ g_n = f_n$ résultent immédiatement des définitions.
\end{proof}



\section{Morphismes de sites simpliciaux}
\label{spaces-simplicial-section-morphism-simplicial-sites}

\noindent
Nous poursuivons de la manière décrite dans la
section \ref{spaces-simplicial-section-simplicial-sites},
en explicitant ce que signifient les morphismes de sites simpliciaux
dans les cas A et B traités dans cette section.

\begin{remark}
\label{spaces-simplicial-remark-morphism-simplicial-sites}
Soient $\mathcal{C}_n, f_\varphi, u_\varphi$ et
$\mathcal{C}'_n, f'_\varphi, u'_\varphi$ comme dans la
situation \ref{spaces-simplicial-situation-simplicial-site}. Par
{\it morphisme $h$ entre sites simpliciaux}, on entendra
\begin{enumerate}
\item[(A)] Des morphismes de sites
$h_n : \mathcal{C}_n \to \mathcal{C}'_n$
tels que $f'_\varphi \circ h_n = h_m \circ f_\varphi$
comme morphismes de sites pour tout $\varphi : [m] \to [n]$.
\item[(B)] Des foncteurs cocontinus
$v_n : \mathcal{C}_n \to \mathcal{C}'_n$
induisant des morphismes de topos $h_n : \Sh(\mathcal{C}_n) \to \Sh(\mathcal{C}'_n)$
tels que $u'_\varphi \circ v_n = v_m \circ u_\varphi$
comme foncteurs pour tout $\varphi : [m] \to [n]$.
\end{enumerate}
Dans les deux cas, on a
$f'_\varphi \circ h_n = h_m \circ f_\varphi$
comme morphismes de topos ; voir
Sites, lemme \ref{sites-lemma-composition-cocontinuous}
pour le cas B et Sites,
définition \ref{sites-definition-composition-morphisms-sites}
pour le cas A.
\end{remark}

\begin{lemma}
\label{spaces-simplicial-lemma-morphism-simplicial-sites}
Soient $\mathcal{C}_n, f_\varphi, u_\varphi$ et
$\mathcal{C}'_n, f'_\varphi, u'_\varphi$ comme dans la
situation \ref{spaces-simplicial-situation-simplicial-site}.
Soit $h$ un morphisme entre sites simpliciaux comme dans la
remarque \ref{spaces-simplicial-remark-morphism-simplicial-sites}.
Alors on obtient un morphisme de topos
$$
h_{total} : \Sh(\mathcal{C}_{total}) \to \Sh(\mathcal{C}'_{total})
$$
et des diagrammes commutatifs
$$
\xymatrix{
\Sh(\mathcal{C}_n) \ar[d]_{g_n} \ar[r]_{h_n} &
\Sh(\mathcal{C}'_n) \ar[d]^{g'_n} \\
\Sh(\mathcal{C}_{total}) \ar[r]^{h_{total}} &
\Sh(\mathcal{C}'_{total})
}
$$
De plus, on a $(g'_n)^{-1} \circ h_{total, *} = h_{n, *} \circ g_n^{-1}$.
\end{lemma}

\begin{proof}
Cas A. Disons que $h_n$ correspond au foncteur continu
$v_n : \mathcal{C}'_n \to \mathcal{C}_n$. On peut alors définir
un foncteur $v_{total} : \mathcal{C}'_{total} \to \mathcal{C}_{total}$
en utilisant $v_n$ en degré $n$. Il s'agit clairement d'un foncteur continu
(voir la définition des recouvrements dans le lemme \ref{spaces-simplicial-lemma-simplicial-site-site}).
Soient
$h_{total}^{-1} = v_{total, s} :
\Sh(\mathcal{C}'_{total}) \to \Sh(\mathcal{C}_{total})$ et
$h_{total, *} = v_{total}^s = v_{total}^p :
\Sh(\mathcal{C}_{total}) \to \Sh(\mathcal{C}'_{total})$
le couple de foncteurs adjoints construit et étudié dans
Sites, sections \ref{sites-section-continuous-functors} et
\ref{sites-section-morphism-sites}.
Pour voir que $h_{total}$ est un morphisme de topos,
il reste à vérifier que $h_{total}^{-1}$ est exact.
On observe d'abord que
$(g'_n)^{-1} \circ h_{total, *} = h_{n, *} \circ g_n^{-1}$ ;
cela résulte immédiatement du calcul des sections sur un objet $U$
de $\mathcal{C}'_n$. Ainsi, si l'on considère un faisceau $\mathcal{F}$
sur $\mathcal{C}_{total}$ comme un système $(\mathcal{F}_n, \mathcal{F}(\varphi))$
comme dans le lemme \ref{spaces-simplicial-lemma-describe-sheaves-simplicial-site-site}, alors
$h_{total, *}\mathcal{F}$ correspond
au système $(h_{n, *}\mathcal{F}_n, h_{n, *}\mathcal{F}(\varphi))$.
Manifestement, le foncteur
$(\mathcal{F}'_n, \mathcal{F}'(\varphi)) \to
(h_n^{-1}\mathcal{F}'_n, h_n^{-1}\mathcal{F}'(\varphi))$
est son adjoint à gauche. Par unicité des adjoints, on conclut que
$h_{total}^{-1}$ est donné sur les systèmes par cette règle. En particulier,
$h_{total}^{-1}$ est exact (d'après la description des faisceaux sur
$\mathcal{C}_{total}$ donnée dans le lemme et l'exactitude des
foncteurs $h_n^{-1}$), et l'on obtient notre morphisme de topos.
Enfin, on obtient également $g_n^{-1} \circ h_{total}^{-1} =
h_n^{-1} \circ (g'_n)^{-1}$, ce qui prouve que le
diagramme affiché dans le lemme est commutatif.

\medskip\noindent
Cas B. On dispose ici d'un foncteur
$v_{total} : \mathcal{C}_{total} \to \mathcal{C}'_{total}$
obtenu en utilisant $v_n$ en degré $n$. Il s'agit clairement d'un foncteur cocontinu
(voir la définition des recouvrements dans le lemme \ref{spaces-simplicial-lemma-simplicial-cocontinuous-site}).
Soit $h_{total}$ le morphisme de topos associé à $v_{total}$.
La commutativité du diagramme affiché dans le lemme résulte
immédiatement de Sites, lemme \ref{sites-lemma-composition-cocontinuous}.
En prenant les adjoints à gauche, la dernière égalité du lemme devient
$$
h_{total}^{-1} \circ (g'_n)^{Sh}_! = g^{Sh}_{n!} \circ h_n^{-1}
$$
Cela résulte immédiatement de la description explicite des foncteurs
$(g'_n)^{Sh}_!$ et $g^{Sh}_{n!}$ donnée dans le
lemme \ref{spaces-simplicial-lemma-restriction-to-components-site},
du fait que $h_n^{-1} \circ (f'_\varphi)^{-1} =
f_\varphi^{-1} \circ h_m^{-1}$ pour $\varphi : [m] \to [n]$, et
du fait que l'on sait déjà que $h_{total}^{-1}$ commute
aux restrictions aux composantes de degré $n$ des sites simpliciaux.
\end{proof}

\begin{lemma}
\label{spaces-simplicial-lemma-direct-image-morphism-simplicial-sites}
Avec les notations et les hypothèses du lemme \ref{spaces-simplicial-lemma-morphism-simplicial-sites},
pour $K \in D(\mathcal{C}_{total})$, on a
$(g'_n)^{-1}Rh_{total, *}K = Rh_{n, *}g_n^{-1}K$.
\end{lemma}

\begin{proof}
Soit $\mathcal{I}^\bullet$ un complexe K-injectif sur $\mathcal{C}_{total}$
représentant $K$. Alors $g_n^{-1}K$ est représenté par
$g_n^{-1}\mathcal{I}^\bullet = \mathcal{I}_n^\bullet$,
qui est K-injectif d'après le
lemme \ref{spaces-simplicial-lemma-restriction-injective-to-component-site}.
On a $(g'_n)^{-1}h_{total, *}\mathcal{I}^\bullet =
h_{n, *}g_n^{-1}\mathcal{I}_n^\bullet$ d'après le
lemme \ref{spaces-simplicial-lemma-morphism-simplicial-sites},
ce qui donne l'égalité voulue.
\end{proof}

\begin{remark}
\label{spaces-simplicial-remark-morphism-augmentation-simplicial-sites}
Soient $\mathcal{C}_n, f_\varphi, u_\varphi$ et
$\mathcal{C}'_n, f'_\varphi, u'_\varphi$ comme dans la
situation \ref{spaces-simplicial-situation-simplicial-site}.
Soit $a_0$, resp.\ $a'_0$, une augmentation
vers un site $\mathcal{D}$, resp.\ $\mathcal{D}'$,
comme dans la remarque \ref{spaces-simplicial-remark-augmentation-site}.
Soit $h$ un morphisme entre sites simpliciaux comme dans la
remarque \ref{spaces-simplicial-remark-morphism-simplicial-sites}.
On dit qu'un morphisme de topos $h_{-1} : \Sh(\mathcal{D}) \to \Sh(\mathcal{D}')$
est {\it compatible avec $h$, $a_0$, $a'_0$} si
\begin{enumerate}
\item[(A)] $h_{-1}$ provient d'un morphisme de sites
$h_{-1} : \mathcal{D} \to \mathcal{D}'$
tel que $a'_0 \circ h_0 = h_{-1} \circ a_0$
comme morphismes de sites.
\item[(B)] $h_{-1}$ provient d'un foncteur cocontinu
$v_{-1} : \mathcal{D} \to \mathcal{D}'$
tel que $u'_0 \circ v_0 = v_{-1} \circ u_0$
comme foncteurs.
\end{enumerate}
Dans les deux cas, on a $a'_0 \circ h_0 = h_{-1} \circ a_0$
comme morphismes de topos ; voir
Sites, lemme \ref{sites-lemma-composition-cocontinuous}
pour le cas B et Sites,
définition \ref{sites-definition-composition-morphisms-sites}
pour le cas A.
\end{remark}

\begin{lemma}
\label{spaces-simplicial-lemma-morphism-augmentation-simplicial-sites}
Soient $\mathcal{C}_n, f_\varphi, u_\varphi, \mathcal{D}, a_0$,
$\mathcal{C}'_n, f'_\varphi, u'_\varphi, \mathcal{D}', a'_0$, et
$h_n$, $n \geq -1$, comme dans la
remarque \ref{spaces-simplicial-remark-morphism-augmentation-simplicial-sites}.
On obtient alors un diagramme commutatif
$$
\xymatrix{
\Sh(\mathcal{C}_{total}) \ar[d]_a \ar[r]_{h_{total}} &
\Sh(\mathcal{C}'_{total}) \ar[d]^{a'} \\
\Sh(\mathcal{D}) \ar[r]^{h_{-1}} &
\Sh(\mathcal{D}')
}
$$
\end{lemma}

\begin{proof}
Le morphisme $h$ est défini dans le lemme \ref{spaces-simplicial-lemma-morphism-simplicial-sites}.
Les morphismes $a$ et $a'$ sont définis dans le lemme \ref{spaces-simplicial-lemma-augmentation-site}.
Il reste donc seulement à prouver la commutativité du diagramme.
Pour cela, on prouve que
$a^{-1} \circ h_{-1}^{-1} = h_{total}^{-1} \circ (a')^{-1}$.
D'après les diagrammes commutatifs du
lemme \ref{spaces-simplicial-lemma-morphism-simplicial-sites}
et la description de $\Sh(\mathcal{C}_{total})$
et $\Sh(\mathcal{C}'_{total})$ en termes de composantes
donnée dans le lemme \ref{spaces-simplicial-lemma-describe-sheaves-simplicial-site-site},
il suffit de montrer que
$$
\xymatrix{
\Sh(\mathcal{C}_n) \ar[d]_{a_n} \ar[r]_{h_n} &
\Sh(\mathcal{C}'_n) \ar[d]^{a'_n} \\
\Sh(\mathcal{D}) \ar[r]^{h_{-1}} &
\Sh(\mathcal{D}')
}
$$
est commutatif pour tout $n$. Cela résulte du cas $n = 0$
(qui est une hypothèse de la
remarque \ref{spaces-simplicial-remark-morphism-augmentation-simplicial-sites}) ;
pour $n > 0$, on choisit $\varphi : [0] \to [n]$,
et la commutativité requise résulte alors du cas $n = 0$,
des relations $a_n = a_0 \circ f_\varphi$
et $a'_n = a'_0 \circ f'_\varphi$,
ainsi que des relations de commutation
$f'_\varphi \circ h_n = h_0 \circ f_\varphi$.
\end{proof}




\section{Sites simpliciaux annelés}
\label{spaces-simplicial-section-simplicial-sites-modules}

\noindent
Munissons notre topos simplicial d'un faisceau d'anneaux.

\begin{lemma}
\label{spaces-simplicial-lemma-restriction-module-to-components-site}
Dans la situation \ref{spaces-simplicial-situation-simplicial-site}, soit $\mathcal{O}$
un faisceau d'anneaux sur $\mathcal{C}_{total}$.
Il existe un morphisme canonique de topos annelés
$g_n : (\Sh(\mathcal{C}_n), \mathcal{O}_n) \to
(\Sh(\mathcal{C}_{total}), \mathcal{O})$
qui coïncide, sur les topos sous-jacents, avec le morphisme $g_n$ du
lemme \ref{spaces-simplicial-lemma-restriction-to-components-site}.
Le foncteur
$g_n^* : \textit{Mod}(\mathcal{O}) \to \textit{Mod}(\mathcal{O}_n)$
admet un adjoint à gauche $g_{n!}$.
Pour $\mathcal{G}$ dans $\textit{Mod}(\mathcal{O}_n)$, la
restriction de $g_{n!}\mathcal{G}$ à $\mathcal{C}_m$ est
$$
\bigoplus\nolimits_{\varphi : [n] \to [m]} f_\varphi^*\mathcal{G}
$$
où $f_\varphi : (\Sh(\mathcal{C}_m), \mathcal{O}_m) \to
(\Sh(\mathcal{C}_n), \mathcal{O}_n)$ est le morphisme de topos annelés
qui coïncide sur les topos avec le morphisme $f_\varphi$ défini précédemment et
qui utilise l'application
$\mathcal{O}(\varphi) : f_\varphi^{-1}\mathcal{O}_n \to \mathcal{O}_m$
sur les faisceaux d'anneaux.
\end{lemma}

\begin{proof}
D'après le lemme \ref{spaces-simplicial-lemma-restriction-to-components-site}, on a
$g_n^{-1}\mathcal{O} = \mathcal{O}_n$, d'où notre
morphisme de topos annelés. D'après Modules sur les sites, lemme
\ref{sites-modules-lemma-lower-shriek-modules},
on obtient l'adjoint $g_{n!}$. Pour prouver la formule donnant $g_{n!}$,
on définit d'abord un faisceau de $\mathcal{O}$-modules $\mathcal{H}$
sur $\mathcal{C}_{total}$ dont la composante de degré $m$ est
le $\mathcal{O}_m$-module
$$
\mathcal{H}_m =
\bigoplus\nolimits_{\varphi : [n] \to [m]} f_\varphi^*\mathcal{G}
$$
Étant donnée une application $\psi : [m] \to [m']$, l'application
$\mathcal{H}(\psi) : f_\psi^{-1}\mathcal{H}_m \to \mathcal{H}_{m'}$
est donnée sur les composantes par
$$
f_\psi^{-1} f_\varphi^*\mathcal{G} \to
f_\psi^* f_\varphi^*\mathcal{G} \to
f_{\psi \circ \varphi}^*\mathcal{G}
$$
Puisque cette application $f_\psi^{-1}\mathcal{H}_m \to \mathcal{H}_{m'}$ est
semi-linéaire relativement à $\mathcal{O}(\psi) : f_\psi^{-1}\mathcal{O}_m \to \mathcal{O}_{m'}$,
elle définit bien un $\mathcal{O}$-module
(utiliser le lemme \ref{spaces-simplicial-lemma-describe-sheaves-simplicial-site-site}).
On prouve alors directement que
$$
\Mor_{\mathcal{O}_n}(\mathcal{G}, \mathcal{F}_n) =
\Mor_{\mathcal{O}}(\mathcal{H}, \mathcal{F})
$$
en procédant comme dans la preuve du lemme \ref{spaces-simplicial-lemma-restriction-to-components-site}.
Ainsi, $\mathcal{H} = g_{n!}\mathcal{G}$, comme voulu.
\end{proof}

\begin{lemma}
\label{spaces-simplicial-lemma-restriction-injective-to-component-limp}
Dans la situation \ref{spaces-simplicial-situation-simplicial-site},
soit $\mathcal{O}$ un faisceau d'anneaux sur $\mathcal{C}_{total}$.
Si $\mathcal{I}$ est injectif dans $\textit{Mod}(\mathcal{O})$, alors
$\mathcal{I}_n$ est un faisceau totalement acyclique sur $\mathcal{C}_n$.
\end{lemma}

\begin{proof}
Cela résulte de
Cohomologie sur les sites, lemme \ref{sites-cohomology-lemma-pullback-injective-limp},
appliqué au foncteur d'inclusion $\mathcal{C}_n \to \mathcal{C}_{total}$
et aux propriétés de celui-ci prouvées dans le lemme \ref{spaces-simplicial-lemma-restriction-to-components-site}.
\end{proof}

\begin{lemma}
\label{spaces-simplicial-lemma-exactness-g-shriek-modules}
Sous les hypothèses du
lemme \ref{spaces-simplicial-lemma-restriction-module-to-components-site}, le foncteur
$g_{n!} : \textit{Mod}(\mathcal{O}_n) \to \textit{Mod}(\mathcal{O})$
est exact si les applications $f_\varphi^{-1}\mathcal{O}_n \to \mathcal{O}_m$
sont plates pour tout $\varphi : [n] \to [m]$.
\end{lemma}

\begin{proof}
Rappelons que $g_{n!}\mathcal{G}$ est le $\mathcal{O}$-module
dont la composante de degré $m$ est le $\mathcal{O}_m$-module
$$
\bigoplus\nolimits_{\varphi : [n] \to [m]} f_\varphi^*\mathcal{G}
$$
Ici, le morphisme de topos annelés
$f_\varphi : (\Sh(\mathcal{C}_m), \mathcal{O}_m) \to
(\Sh(\mathcal{C}_n), \mathcal{O}_n)$ utilise l'application
$f_\varphi^{-1}\mathcal{O}_n \to \mathcal{O}_m$ de l'
énoncé du lemme. Si ces applications sont plates, alors
$f_\varphi^*$ est exact
(Modules sur les sites, lemme \ref{sites-modules-lemma-flat-pullback-exact}).
Par définition du site $\mathcal{C}_{total}$, on voit que ces
foncteurs possèdent les propriétés d'exactitude voulues, et l'on conclut.
\end{proof}

\begin{lemma}
\label{spaces-simplicial-lemma-restriction-injective-to-component-site-module}
Dans la situation \ref{spaces-simplicial-situation-simplicial-site},
soit $\mathcal{O}$ un faisceau d'anneaux sur $\mathcal{C}_{total}$
tel que $f_\varphi^{-1}\mathcal{O}_n \to \mathcal{O}_m$
soit plat pour tout $\varphi : [n] \to [m]$.
Si $\mathcal{I}$ est injectif dans $\textit{Mod}(\mathcal{O})$, alors
$\mathcal{I}_n$ est injectif dans $\textit{Mod}(\mathcal{O}_n)$.
\end{lemma}

\begin{proof}
Cela résulte du
lemme \ref{homology-lemma-adjoint-preserve-injectives} de Homologie
et du
lemme \ref{spaces-simplicial-lemma-exactness-g-shriek-modules}.
\end{proof}







\section{Morphismes de sites simpliciaux annelés}
\label{spaces-simplicial-section-morphism-simplicial-sites-modules}

\noindent
Nous poursuivons la discussion de la section \ref{spaces-simplicial-section-morphism-simplicial-sites}.

\begin{remark}
\label{spaces-simplicial-remark-morphism-simplicial-sites-modules}
Soient $\mathcal{C}_n, f_\varphi, u_\varphi$ et
$\mathcal{C}'_n, f'_\varphi, u'_\varphi$ comme dans la
situation \ref{spaces-simplicial-situation-simplicial-site}.
Soient $\mathcal{O}$ et $\mathcal{O}'$
des faisceaux d'anneaux sur $\mathcal{C}_{total}$ et $\mathcal{C}'_{total}$, respectivement.
On dira que $(h, h^\sharp)$ est un
{\it morphisme entre sites simpliciaux annelés}
si $h$ est un morphisme entre sites simpliciaux comme dans la
remarque \ref{spaces-simplicial-remark-morphism-simplicial-sites}
et si $h^\sharp : h_{total}^{-1}\mathcal{O}' \to \mathcal{O}$,
ou, de manière équivalente, $h^\sharp : \mathcal{O}' \to h_{total, *}\mathcal{O}$,
est un homomorphisme de faisceaux d'anneaux.
\end{remark}

\begin{lemma}
\label{spaces-simplicial-lemma-morphism-simplicial-sites-modules}
Soient $\mathcal{C}_n, f_\varphi, u_\varphi$ et
$\mathcal{C}'_n, f'_\varphi, u'_\varphi$ comme dans la
situation \ref{spaces-simplicial-situation-simplicial-site}.
Soient $\mathcal{O}$ et $\mathcal{O}'$
des faisceaux d'anneaux sur $\mathcal{C}_{total}$ et $\mathcal{C}'_{total}$, respectivement.
Soit $(h, h^\sharp)$ un morphisme entre sites simpliciaux annelés comme dans la
remarque \ref{spaces-simplicial-remark-morphism-simplicial-sites-modules}.
Alors on obtient un morphisme de topos annelés
$$
h_{total} :
(\Sh(\mathcal{C}_{total}), \mathcal{O})
\to
(\Sh(\mathcal{C}'_{total}), \mathcal{O}')
$$
et des diagrammes commutatifs
$$
\xymatrix{
(\Sh(\mathcal{C}_n), \mathcal{O}_n) \ar[d]_{g_n} \ar[r]_{h_n} &
(\Sh(\mathcal{C}'_n), \mathcal{O}'_n) \ar[d]^{g'_n} \\
(\Sh(\mathcal{C}_{total}), \mathcal{O}) \ar[r]^{h_{total}} &
(\Sh(\mathcal{C}'_{total}), \mathcal{O}')
}
$$
de topos annelés, où $g_n$ et $g'_n$ sont comme dans le
lemme \ref{spaces-simplicial-lemma-restriction-module-to-components-site}.
De plus, on a
$(g'_n)^* \circ h_{total, *} = h_{n, *} \circ g_n^*$
comme foncteurs $\textit{Mod}(\mathcal{O}) \to \textit{Mod}(\mathcal{O}'_n)$.
\end{lemma}

\begin{proof}
Cela résulte des
lemmes \ref{spaces-simplicial-lemma-morphism-simplicial-sites} et
\ref{spaces-simplicial-lemma-restriction-module-to-components-site},
en tenant compte des faisceaux d'anneaux.
Un petit point est que, pour définir $h_n$ comme morphisme
de topos annelés, on pose
$h_n^\sharp = g_n^{-1}h^\sharp :
g_n^{-1}h_{total}^{-1}\mathcal{O}' \to g_n^{-1}\mathcal{O}$,
ce qui a un sens puisque
$g_n^{-1}h_{total}^{-1}\mathcal{O}' = h_n^{-1}(g'_n)^{-1}\mathcal{O}' =
h_n^{-1}\mathcal{O}'_n$ et $g_n^{-1}\mathcal{O} = \mathcal{O}_n$.
Notons que $g_n^*\mathcal{F} = g_n^{-1}\mathcal{F}$
pour un faisceau de $\mathcal{O}$-modules $\mathcal{F}$,
et de même pour $g'_n$ ; cela aide à comprendre pourquoi
$(g'_n)^* \circ h_{total, *} = h_{n, *} \circ g_n^*$
résulte de l'énoncé correspondant du
lemme \ref{spaces-simplicial-lemma-morphism-simplicial-sites}.
\end{proof}

\begin{lemma}
\label{spaces-simplicial-lemma-direct-image-morphism-simplicial-sites-modules}
Avec les notations et les hypothèses du
lemme \ref{spaces-simplicial-lemma-morphism-simplicial-sites-modules},
pour $K \in D(\mathcal{O})$, on a
$(g'_n)^*Rh_{total, *}K = Rh_{n, *}g_n^*K$.
\end{lemma}

\begin{proof}
Rappelons que $g_n^* = g_n^{-1}$ parce que $g_n^{-1}\mathcal{O} = \mathcal{O}_n$
par la construction du lemme \ref{spaces-simplicial-lemma-restriction-module-to-components-site}.
En particulier, $g_n^*$ est exact, et $Lg_n^*$ s'obtient en appliquant $g_n^*$
à n'importe quel complexe de modules représentant l'objet considéré. De même pour $g'_n$.
Il existe un morphisme canonique de changement de base
$(g'_n)^*Rh_{total, *}K \to Rh_{n, *}g_n^*K$ ; voir
Cohomologie sur les sites, remarque \ref{sites-cohomology-remark-base-change}.
D'après Cohomologie sur les sites, lemme
\ref{sites-cohomology-lemma-modules-abelian-unbounded},
l'image de ce morphisme dans $D(\mathcal{C}'_n)$ est le morphisme
$(g'_n)^{-1}Rh_{total, *}K_{ab} \to Rh_{n, *}g_n^{-1}K_{ab}$,
où $K_{ab}$ est l'image de $K$ dans $D(\mathcal{C}_{total})$.
On a prouvé que celui-ci est un isomorphisme dans le
lemme \ref{spaces-simplicial-lemma-direct-image-morphism-simplicial-sites},
et le résultat en découle.
\end{proof}








\section{Cohomologie sur les sites simpliciaux}
\label{spaces-simplicial-section-cohomology-simplicial-sites}

\noindent
Soit $\mathcal{C}$ comme dans la situation \ref{spaces-simplicial-situation-simplicial-site}.
Dans l'énoncé des lemmes suivants, on notera
$g_n : \Sh(\mathcal{C}_n) \to \Sh(\mathcal{C}_{total})$ le
morphisme de topos du
lemme \ref{spaces-simplicial-lemma-restriction-to-components-site}. Si $\varphi : [m] \to [n]$
est un morphisme de $\Delta$, alors le diagramme de topos
$$
\xymatrix{
\Sh(\mathcal{C}_n) \ar[rd]_{g_n} \ar[rr]_{f_\varphi} & &
\Sh(\mathcal{C}_m) \ar[ld]^{g_m} \\
& \Sh(\mathcal{C}_{total})
}
$$
n'est pas commutatif, mais il existe un $2$-morphisme $g_n \to g_m \circ f_\varphi$
provenant des applications
$\mathcal{F}(\varphi) : f_\varphi^{-1}\mathcal{F}_m \to \mathcal{F}_n$.
Voir Sites, section \ref{sites-section-2-category}.

\begin{lemma}
\label{spaces-simplicial-lemma-simplicial-resolution-Z-site}
Dans la situation \ref{spaces-simplicial-situation-simplicial-site} et avec les notations ci-dessus,
il existe un complexe
$$
\ldots \to g_{2!}\mathbf{Z} \to g_{1!}\mathbf{Z} \to g_{0!}\mathbf{Z}
$$
de faisceaux abéliens sur $\mathcal{C}_{total}$ qui constitue une résolution du
faisceau constant de valeur $\mathbf{Z}$ sur $\mathcal{C}_{total}$.
\end{lemma}

\begin{proof}
On utilisera sans autre mention la description des foncteurs $g_{n!}$ donnée dans le
lemme \ref{spaces-simplicial-lemma-restriction-to-components-site}. Comme applications du complexe,
on prend $\sum (-1)^i d^n_i$, où
$d^n_i : g_{n!}\mathbf{Z} \to g_{n - 1!}\mathbf{Z}$ est l'
adjoint de l'application $\mathbf{Z} \to
\bigoplus_{[n - 1] \to [n]} \mathbf{Z} = g_n^{-1}g_{n - 1!}\mathbf{Z}$
correspondant au facteur étiqueté par $\delta^n_i : [n - 1] \to [n]$.
En appliquant $g_m^{-1}$ au complexe, on obtient le complexe
$$
\ldots \to
\bigoplus\nolimits_{\alpha \in \Mor_\Delta([2], [m])]} \mathbf{Z} \to
\bigoplus\nolimits_{\alpha \in \Mor_\Delta([1], [m])]} \mathbf{Z} \to
\bigoplus\nolimits_{\alpha \in \Mor_\Delta([0], [m])]} \mathbf{Z}
$$
sur $\mathcal{C}_m$.
Autrement dit, c'est le complexe associé au
faisceau abélien libre sur l'ensemble simplicial $\Delta[m]$ ; voir
Méthodes simpliciales, exemple \ref{simplicial-example-simplex-simplicial-set}.
Puisque $\Delta[m]$ est homotopiquement équivalent à $\Delta[0]$ ; voir
Méthodes simpliciales, exemple \ref{simplicial-example-simplex-contractible},
et puisque « prendre le faisceau abélien libre sur » est un foncteur,
on voit que le complexe ci-dessus est homotopiquement équivalent au
faisceau abélien libre sur $\Delta[0]$
(Méthodes simpliciales, remarque \ref{simplicial-remark-homotopy-better} et
lemme \ref{simplicial-lemma-homotopy-equivalence-s-N}).
Ce complexe est acyclique en degrés strictement positifs
et égal à $\mathbf{Z}$ en degré $0$.
\end{proof}

\begin{lemma}
\label{spaces-simplicial-lemma-cech-complex}
Dans la situation \ref{spaces-simplicial-situation-simplicial-site}, soit $\mathcal{F}$ un faisceau
abélien sur $\mathcal{C}_{total}$. Il existe un complexe canonique
$$
0 \to \Gamma(\mathcal{C}_{total}, \mathcal{F}) \to
\Gamma(\mathcal{C}_0, \mathcal{F}_0) \to
\Gamma(\mathcal{C}_1, \mathcal{F}_1) \to
\Gamma(\mathcal{C}_2, \mathcal{F}_2) \to \ldots
$$
qui est exact en degrés $-1, 0$ et partout exact
si $\mathcal{F}$ est injectif.
\end{lemma}

\begin{proof}
Observons que
$\Hom(\mathbf{Z}, \mathcal{F}) = \Gamma(\mathcal{C}_{total}, \mathcal{F})$
et que
$\Hom(g_{n!}\mathbf{Z}, \mathcal{F}) = \Gamma(\mathcal{C}_n, \mathcal{F}_n)$.
Ce lemme résulte donc immédiatement du
lemme \ref{spaces-simplicial-lemma-simplicial-resolution-Z-site}
et du fait que $\Hom(-, \mathcal{F})$ est exact si
$\mathcal{F}$ est injectif.
\end{proof}

\begin{lemma}
\label{spaces-simplicial-lemma-simplicial-sheaf-cohomology-site}
Dans la situation \ref{spaces-simplicial-situation-simplicial-site}, pour $K$ dans
$D^+(\mathcal{C}_{total})$, il existe une suite spectrale
$(E_r, d_r)_{r \geq 0}$ avec
$$
E_1^{p, q} = H^q(\mathcal{C}_p, K_p),\quad
d_1^{p, q} : E_1^{p, q} \to E_1^{p + 1, q}
$$
qui aboutit à $H^{p + q}(\mathcal{C}_{total}, K)$.
Cette suite spectrale est fonctorielle en $K$.
\end{lemma}

\begin{proof}
Soit $\mathcal{I}^\bullet$ un complexe borné inférieurement d'injectifs
représentant $K$. Considérons le complexe double dont les termes sont
$$
A^{p, q} = \Gamma(\mathcal{C}_p, \mathcal{I}^q_p)
$$
où les flèches horizontales proviennent du lemme \ref{spaces-simplicial-lemma-cech-complex}
et les flèches verticales des différentielles du
complexe $\mathcal{I}^\bullet$. Les lignes du complexe double sont exactes
en degrés strictement positifs et valent
$\Gamma(\mathcal{C}_{total}, \mathcal{I}^q)$ en degré $0$.
D'autre part, puisque la restriction à $\mathcal{C}_p$ est exacte
(lemme \ref{spaces-simplicial-lemma-restriction-to-components-site}),
le complexe $\mathcal{I}_p^\bullet$ représente $K_p$ dans
$D(\mathcal{C}_p)$. Les faisceaux $\mathcal{I}_p^q$ sont des
faisceaux abéliens injectifs sur $\mathcal{C}_p$
(lemme \ref{spaces-simplicial-lemma-restriction-injective-to-component-site}).
La cohomologie des colonnes calcule donc les groupes
$H^q(\mathcal{C}_p, K_p)$. On conclut en appliquant les
lemmes \ref{homology-lemma-first-quadrant-ss} et
\ref{homology-lemma-double-complex-gives-resolution} de Homologie.
\end{proof}

\begin{remark}
\label{spaces-simplicial-remark-simplicial-cohomology-site}
On fait les mêmes hypothèses et on emploie les mêmes notations que dans le
lemme \ref{spaces-simplicial-lemma-simplicial-sheaf-cohomology-site},
sauf que l'on ne suppose pas que $K$ dans $D(\mathcal{C}_{total})$
soit borné inférieurement. On affirme qu'il existe également dans ce cas une
suite spectrale naturelle. En effet, supposons que
$\mathcal{I}^\bullet$ soit un complexe K-injectif de faisceaux sur
$\mathcal{C}_{total}$, à termes injectifs, qui représente $K$.
On a
\begin{align*}
R\Gamma(\mathcal{C}_{total}, K)
& =
R\Hom(\mathbf{Z}, K) \\
& =
R\Hom(
\ldots \to g_{2!}\mathbf{Z} \to g_{1!}\mathbf{Z} \to g_{0!}\mathbf{Z}, K) \\
& =
\Gamma(\mathcal{C}_{total},
\SheafHom^\bullet(
\ldots \to g_{2!}\mathbf{Z} \to g_{1!}\mathbf{Z} \to g_{0!}\mathbf{Z},
\mathcal{I}^\bullet)) \\
& =
\text{Tot}_\pi(A^{\bullet, \bullet})
\end{align*}
où $A^{\bullet, \bullet}$ est le complexe double dont les termes sont
$A^{p, q} = \Gamma(\mathcal{C}_p, \mathcal{I}^q_p)$, et où $\text{Tot}_\pi$
désigne la totalisation par produits de ce complexe double.
En effet, la première égalité vaut dans tout site. La deuxième égalité
résulte du lemme \ref{spaces-simplicial-lemma-simplicial-resolution-Z-site}.
La troisième égalité vaut parce que $\mathcal{I}^\bullet$ est K-injectif ; voir
Cohomologie sur les sites, sections \ref{sites-cohomology-section-hom-complexes}
et \ref{sites-cohomology-section-internal-hom}.
La dernière égalité résulte de la construction de $\SheafHom^\bullet$
et du fait que $\Hom(g_{p!}\mathbf{Z}, \mathcal{I}^q) =
\Gamma(\mathcal{C}_p, \mathcal{I}^q_p)$.
On obtient alors notre suite spectrale en considérant
$\text{Tot}_\pi(A^{\bullet, \bullet})$ comme un
complexe filtré avec $F^i\text{Tot}^n_\pi(A^{\bullet, \bullet}) =
\prod_{p + q = n,\ p \geq i} A^{p, q}$. La suite spectrale
ainsi obtenue se comporte comme la suite spectrale $({}'E_r, {}'d_r)_{r \geq 0}$ de
Homologie, section \ref{homology-section-double-complex}
(où est traité le cas de la totalisation par sommes directes),
si ce n'est pour les questions de régularité, de bornitude, de convergence et
d'aboutissement. En particulier, on obtient $E_1^{p, q} = H^q(\mathcal{C}_p, K_p)$
comme dans le lemme \ref{spaces-simplicial-lemma-simplicial-sheaf-cohomology-site}.
\end{remark}

\begin{lemma}
\label{spaces-simplicial-lemma-simplicial-sheaf-cohomology-site-zero}
Dans la situation \ref{spaces-simplicial-situation-simplicial-site}, soit $K$ un objet de
$D(\mathcal{C}_{total})$.
\begin{enumerate}
\item Si $H^{-p}(\mathcal{C}_p, K_p) = 0$ pour tout $p \geq 0$, alors
$H^0(\mathcal{C}_{total}, K) = 0$.
\item Si $R\Gamma(\mathcal{C}_p, K_p) = 0$
pour tout $p \geq 0$, alors $R\Gamma(\mathcal{C}_{total}, K) = 0$.
\end{enumerate}
\end{lemma}

\begin{proof}
Avec les notations de la remarque \ref{spaces-simplicial-remark-simplicial-cohomology-site},
on voit que $R\Gamma(\mathcal{C}_{total}, K)$ est
représenté par $\text{Tot}_\pi(A^{\bullet, \bullet})$.
L'hypothèse de (1) nous dit que $H^{-p}(A^{p, \bullet}) = 0$.
L'annulation de (1) résulte donc du
lemme \ref{more-algebra-lemma-vanishing-coh-prod-totalization} de Compléments d'algèbre.
L'assertion (2) résulte de l'assertion (1) en prenant des décalés.
\end{proof}

\begin{lemma}
\label{spaces-simplicial-lemma-sanity-check}
Soit $\mathcal{C}$ comme dans la situation \ref{spaces-simplicial-situation-simplicial-site}.
Soit $U \in \Ob(\mathcal{C}_n)$. Soit
$\mathcal{F} \in \textit{Ab}(\mathcal{C}_{total})$.
Alors $H^p(U, \mathcal{F}) = H^p(U, g_n^{-1}\mathcal{F})$,
où, dans le membre de gauche, $U$ est considéré comme un objet de $\mathcal{C}_{total}$.
\end{lemma}

\begin{proof}
Notons que l'expression « $U$ considéré comme objet de $\mathcal{C}_{total}$ »
est expliquée par la construction de $\mathcal{C}_{total}$ dans le
lemme \ref{spaces-simplicial-lemma-simplicial-site-site} dans le cas (A) et le
lemme \ref{spaces-simplicial-lemma-simplicial-cocontinuous-site} dans le cas (B).
L'égalité résulte alors du
lemme \ref{spaces-simplicial-lemma-restriction-injective-to-component-site}
et de la définition de la cohomologie.
\end{proof}






\section{Cohomologie et augmentations de sites simpliciaux}
\label{spaces-simplicial-section-cohomology-augmentation-simplicial-sites}

\noindent
Considérons un site simplicial $\mathcal{C}$ comme dans la situation
\ref{spaces-simplicial-situation-simplicial-site}.
Soit $a_0$ une augmentation vers un site $\mathcal{D}$ comme dans la
remarque \ref{spaces-simplicial-remark-augmentation-site}.
Par le lemme \ref{spaces-simplicial-lemma-augmentation-site}, on obtient un morphisme de topoi
$$
a : \Sh(\mathcal{C}_{total}) \longrightarrow \Sh(\mathcal{D})
$$
ainsi que des morphismes de topoi
$g_n : \Sh(\mathcal{C}_n) \to \Sh(\mathcal{C}_{total})$
comme dans le lemme \ref{spaces-simplicial-lemma-restriction-to-components-site}.
Les compositions $a \circ g_n$ sont notées
$a_n : \Sh(\mathcal{C}_n) \to \Sh(\mathcal{D})$.
De plus, la structure simpliciale donne des morphismes
de topoi
$f_\varphi : \Sh(\mathcal{C}_n) \to \Sh(\mathcal{C}_m)$
tels que $a_n \circ f_\varphi = a_m$ pour tout $\varphi : [m] \to [n]$.

\begin{lemma}
\label{spaces-simplicial-lemma-simplicial-resolution-augmentation}
Dans la situation \ref{spaces-simplicial-situation-simplicial-site}, soit
$a_0$ une augmentation vers un site $\mathcal{D}$
comme dans la remarque \ref{spaces-simplicial-remark-augmentation-site}.
Pour tout faisceau abélien $\mathcal{G}$ sur $\mathcal{D}$,
il existe un complexe exact
$$
\ldots \to
g_{2!}(a_2^{-1}\mathcal{G}) \to
g_{1!}(a_1^{-1}\mathcal{G}) \to
g_{0!}(a_0^{-1}\mathcal{G}) \to
a^{-1}\mathcal{G} \to 0
$$
de faisceaux abéliens sur $\mathcal{C}_{total}$.
\end{lemma}

\begin{proof}
Nous conseillons au lecteur de lire d'abord la preuve du
lemme \ref{spaces-simplicial-lemma-simplicial-resolution-Z-site}.
Nous utiliserons le lemme \ref{spaces-simplicial-lemma-augmentation-site} et
la description des foncteurs $g_{n!}$ donnée dans le
lemme \ref{spaces-simplicial-lemma-restriction-to-components-site}, sans autre mention.
En particulier, $g_{n!}(a_n^{-1}\mathcal{G})$ est le
faisceau sur $\mathcal{C}_{total}$ dont la restriction à $\mathcal{C}_m$
est le faisceau
$$
\bigoplus\nolimits_{\varphi : [n] \to [m]} f_\varphi^{-1}a_n^{-1}\mathcal{G} =
\bigoplus\nolimits_{\varphi : [n] \to [m]} a_m^{-1}\mathcal{G}
$$
Comme morphismes du complexe, prenons $\sum (-1)^i d^n_i$, où
$d^n_i : g_{n!}(a_n^{-1}\mathcal{G}) \to g_{n - 1!}(a_{n - 1}^{-1}\mathcal{G})$
est adjoint au morphisme
$a_n^{-1}\mathcal{G} \to \bigoplus_{[n - 1] \to [n]} a_n^{-1}\mathcal{G} =
g_n^{-1}g_{n - 1!}(a_{n - 1}^{-1}\mathcal{G})$
correspondant au facteur étiqueté $\delta^n_i : [n - 1] \to [n]$.
Le morphisme $g_{0!}(a_0^{-1}\mathcal{G}) \to a^{-1}\mathcal{G}$ est adjoint
au morphisme identité de $a_0^{-1}\mathcal{G}$.
L'application de $g_m^{-1}$ au complexe de chaînes en degrés
$\ldots, 2, 1, 0$ donne le complexe
$$
\ldots \to
\bigoplus\nolimits_{\alpha \in \Mor_\Delta([2], [m])]} a_m^{-1}\mathcal{G} \to
\bigoplus\nolimits_{\alpha \in \Mor_\Delta([1], [m])]} a_m^{-1}\mathcal{G} \to
\bigoplus\nolimits_{\alpha \in \Mor_\Delta([0], [m])]} a_m^{-1}\mathcal{G}
$$
sur $\mathcal{C}_m$. Il est égal à $a_m^{-1}\mathcal{G}$
tensorisé sur le faisceau constant $\mathbf{Z}$ avec le complexe
$$
\ldots \to
\bigoplus\nolimits_{\alpha \in \Mor_\Delta([2], [m])]} \mathbf{Z} \to
\bigoplus\nolimits_{\alpha \in \Mor_\Delta([1], [m])]} \mathbf{Z} \to
\bigoplus\nolimits_{\alpha \in \Mor_\Delta([0], [m])]} \mathbf{Z}
$$
discuté dans la preuve du lemme \ref{spaces-simplicial-lemma-simplicial-resolution-Z-site}.
Nous y avons vu que ce complexe est homotopiquement équivalent à
$\mathbf{Z}$ placé en degré $0$, ce qui termine la preuve.
\end{proof}

\begin{lemma}
\label{spaces-simplicial-lemma-augmentation-cech-complex}
Dans la situation \ref{spaces-simplicial-situation-simplicial-site}, soit
$a_0$ une augmentation vers un site $\mathcal{D}$
comme dans la remarque \ref{spaces-simplicial-remark-augmentation-site}.
Pour un faisceau abélien $\mathcal{F}$ sur $\mathcal{C}_{total}$
il existe un complexe canonique
$$
0 \to a_*\mathcal{F} \to a_{0, *}\mathcal{F}_0 \to a_{1, *}\mathcal{F}_1 \to
a_{2, *}\mathcal{F}_2 \to \ldots
$$
sur $\mathcal{D}$ qui est exact en degrés $-1, 0$ et
exact partout si $\mathcal{F}$ est injectif.
\end{lemma}

\begin{proof}
Pour construire le complexe, il suffit, par le lemme de Yoneda, de construire pour tout
faisceau abélien $\mathcal{G}$ sur $\mathcal{D}$ un complexe
$$
0 \to \Hom(\mathcal{G}, a_*\mathcal{F}) \to
\Hom(\mathcal{G}, a_{0, *}\mathcal{F}_0) \to
\Hom(\mathcal{G}, a_{1, *}\mathcal{F}_1) \to \ldots
$$
fonctoriellement en $\mathcal{G}$. Pour cela, appliquons $\Hom(-, \mathcal{F})$
au complexe exact du lemme \ref{spaces-simplicial-lemma-simplicial-resolution-augmentation}
et utilisons l'adjonction entre image inverse et image directe.
Les propriétés d'exactitude en degrés $-1, 0$ découlent de
la construction, puisque $\Hom(-, \mathcal{F})$ est exact à gauche.
Si $\mathcal{F}$ est un faisceau abélien injectif, alors le complexe
est exact car $\Hom(-, \mathcal{F})$ est exact.
\end{proof}

\begin{lemma}
\label{spaces-simplicial-lemma-augmentation-spectral-sequence}
Dans la situation \ref{spaces-simplicial-situation-simplicial-site}, soit
$a_0$ une augmentation vers un site $\mathcal{D}$
comme dans la remarque \ref{spaces-simplicial-remark-augmentation-site}.
Pour tout $K$ dans $D^+(\mathcal{C}_{total})$, il existe une suite
spectrale
$(E_r, d_r)_{r \geq 0}$ avec
$$
E_1^{p, q} = R^qa_{p, *} K_p,\quad
d_1^{p, q} : E_1^{p, q} \to E_1^{p + 1, q}
$$
qui aboutit à $R^{p + q}a_*K$. Cette suite spectrale est fonctorielle en $K$.
\end{lemma}

\begin{proof}
Soit $\mathcal{I}^\bullet$ un complexe borné inférieurement d'injectifs
représentant $K$. Considérons le complexe double dont les termes sont
$$
A^{p, q} = a_{p, *}\mathcal{I}^q_p
$$
où les flèches horizontales proviennent du
lemme \ref{spaces-simplicial-lemma-augmentation-cech-complex}
et les flèches verticales des différentiels du complexe
$\mathcal{I}^\bullet$. Les lignes du double complexe sont exactes
en degrés positifs et s'évaluent à $a_*\mathcal{I}^q$ en degré $0$.
D'autre part, puisque la restriction à $\mathcal{C}_p$ est exacte
(lemme \ref{spaces-simplicial-lemma-restriction-to-components-site})
le complexe $\mathcal{I}_p^\bullet$ représente $K_p$ dans
$D(\mathcal{C}_p)$. Les faisceaux $\mathcal{I}_p^q$ sont des faisceaux abéliens injectifs
sur $\mathcal{C}_p$
(lemme \ref{spaces-simplicial-lemma-restriction-injective-to-component-site}).
La cohomologie des colonnes calcule donc $R^qa_{p, *}K_p$.
On conclut en appliquant les
lemmes \ref{homology-lemma-first-quadrant-ss} et
\ref{homology-lemma-double-complex-gives-resolution}.
\end{proof}





\section{Cohomologie sur les sites simpliciaux annelés}
\label{spaces-simplicial-section-cohomology-simplicial-sites-modules}

\noindent
Cette section est l'analogue de
la section \ref{spaces-simplicial-section-cohomology-simplicial-sites}
pour les faisceaux de modules.

\medskip\noindent
Dans la situation \ref{spaces-simplicial-situation-simplicial-site}, soit $\mathcal{O}$
un faisceau d'anneaux sur $\mathcal{C}_{total}$.
Dans les énoncés des lemmes suivants, nous noterons
$g_n : (\Sh(\mathcal{C}_n), \mathcal{O}_n) \to
(\Sh(\mathcal{C}_{total}), \mathcal{O})$
le morphisme de topos annelés du lemme
\ref{spaces-simplicial-lemma-restriction-module-to-components-site}.
Si $\varphi : [m] \to [n]$ est un morphisme de $\Delta$, alors le diagramme
de topoi annelés
$$
\xymatrix{
(\Sh(\mathcal{C}_n), \mathcal{O}_n) \ar[rd]_{g_n} \ar[rr]_{f_\varphi} & &
(\Sh(\mathcal{C}_m), \mathcal{O}_m) \ar[ld]^{g_m} \\
& (\Sh(\mathcal{C}_{total}), \mathcal{O})
}
$$
n'est pas commutatif, mais il existe un $2$-morphisme $g_n \to g_m \circ f_\varphi$
provenant des morphismes
$\mathcal{F}(\varphi) : f_\varphi^{-1}\mathcal{F}_m \to \mathcal{F}_n$.
Voir Sites, section \ref{sites-section-2-category}.

\begin{lemma}
\label{spaces-simplicial-lemma-simplicial-resolution-ringed}
Dans la situation \ref{spaces-simplicial-situation-simplicial-site}, soit $\mathcal{O}$
un faisceau d'anneaux sur $\mathcal{C}_{total}$. Il existe un complexe
$$
\ldots \to g_{2!}\mathcal{O}_2 \to g_{1!}\mathcal{O}_1 \to g_{0!}\mathcal{O}_0
$$
de $\mathcal{O}$-modules qui est une résolution de
$\mathcal{O}$.
Ici, $g_{n!}$ est comme dans le lemme \ref{spaces-simplicial-lemma-restriction-module-to-components-site}.
\end{lemma}

\begin{proof}
Nous utiliserons la description de $g_{n!}$ donnée dans le
lemme \ref{spaces-simplicial-lemma-restriction-to-components-site}.
Comme morphismes du complexe, prenons $\sum (-1)^i d^n_i$, où
$d^n_i : g_{n!}\mathcal{O}_n \to g_{n - 1!}\mathcal{O}_{n - 1}$
est adjoint au morphisme
$\mathcal{O}_n \to \bigoplus_{[n - 1] \to [n]} \mathcal{O}_n =
g_n^*g_{n - 1!}\mathcal{O}_{n - 1}$
correspondant au facteur étiqueté $\delta^n_i : [n - 1] \to [n]$.
Alors $g_m^{-1}$ appliqué au complexe donne le complexe
$$
\ldots \to
\bigoplus\nolimits_{\alpha \in \Mor_\Delta([2], [m])]} \mathcal{O}_m \to
\bigoplus\nolimits_{\alpha \in \Mor_\Delta([1], [m])]} \mathcal{O}_m \to
\bigoplus\nolimits_{\alpha \in \Mor_\Delta([0], [m])]} \mathcal{O}_m
$$
sur $\mathcal{C}_m$.
Autrement dit, il s'agit du complexe associé au
$\mathcal{O}_m$-module libre engendré par l'ensemble simplicial $\Delta[m]$; voir
Simplicial, exemple \ref{simplicial-example-simplex-simplicial-set}.
Puisque $\Delta[m]$ est homotopiquement équivalent à $\Delta[0]$, voir
Simplicial, exemple \ref{simplicial-example-simplex-contractible},
et puisque « prendre le faisceau abélien libre engendré par » est un foncteur,
on voit que le complexe ci-dessus est homotopiquement équivalent au
faisceau abélien libre engendré par $\Delta[0]$
(Simplicial, remarque \ref{simplicial-remark-homotopy-better} et
lemme \ref{simplicial-lemma-homotopy-equivalence-s-N}).
Ce complexe est acyclique en degrés positifs
et égal à $\mathcal{O}_m$ en degré $0$.
\end{proof}

\begin{lemma}
\label{spaces-simplicial-lemma-cech-complex-modules}
Dans la situation \ref{spaces-simplicial-situation-simplicial-site}, soit $\mathcal{O}$
un faisceau d'anneaux. Soit $\mathcal{F}$ un faisceau
de $\mathcal{O}$-modules. Il existe un complexe canonique
$$
0 \to \Gamma(\mathcal{C}_{total}, \mathcal{F}) \to
\Gamma(\mathcal{C}_0, \mathcal{F}_0) \to
\Gamma(\mathcal{C}_1, \mathcal{F}_1) \to
\Gamma(\mathcal{C}_2, \mathcal{F}_2) \to \ldots
$$
qui est exact en degrés $-1, 0$ et exact partout
si $\mathcal{F}$ est un $\mathcal{O}$-module injectif.
\end{lemma}

\begin{proof}
Observons que
$\Hom(\mathcal{O}, \mathcal{F}) = \Gamma(\mathcal{C}_{total}, \mathcal{F})$
et
$\Hom(g_{n!}\mathcal{O}_n, \mathcal{F}) = \Gamma(\mathcal{C}_n, \mathcal{F}_n)$.
Ce lemme est donc une conséquence immédiate du lemme
\ref{spaces-simplicial-lemma-simplicial-resolution-ringed}
et du fait que $\Hom(-, \mathcal{F})$ est exact si
$\mathcal{F}$ est injectif.
\end{proof}

\begin{lemma}
\label{spaces-simplicial-lemma-simplicial-module-cohomology-site}
Dans la situation \ref{spaces-simplicial-situation-simplicial-site}, soit $\mathcal{O}$
un faisceau d'anneaux. Pour $K$ dans $D^+(\mathcal{O})$
il existe une suite spectrale $(E_r, d_r)_{r \geq 0}$ telle que
$$
E_1^{p, q} = H^q(\mathcal{C}_p, K_p),\quad
d_1^{p, q} : E_1^{p, q} \to E_1^{p + 1, q}
$$
qui aboutit à $H^{p + q}(\mathcal{C}_{total}, K)$.
Cette suite spectrale est fonctorielle en $K$.
\end{lemma}

\begin{proof}
Soit $\mathcal{I}^\bullet$ un complexe borné inférieurement de
$\mathcal{O}$-modules injectifs représentant $K$. Considérons le complexe double de termes
$$
A^{p, q} = \Gamma(\mathcal{C}_p, \mathcal{I}^q_p)
$$
où les flèches horizontales proviennent du
lemme \ref{spaces-simplicial-lemma-cech-complex-modules}
et les flèches verticales des différentiels du complexe
$\mathcal{I}^\bullet$. Notons que
$\Gamma(\mathcal{D}, -) =
\Hom_{\mathcal{O}_\mathcal{D}}(\mathcal{O}_\mathcal{D}, -)$
sur $\textit{Mod}(\mathcal{O}_\mathcal{D})$. Le lemme
affirme donc que les lignes du complexe double sont exactes
en degrés positifs et donnent
$\Gamma(\mathcal{C}_{total}, \mathcal{I}^q)$ en degré $0$.
Ainsi le complexe total associé au double complexe
calcule $R\Gamma(\mathcal{C}_{total}, K)$ d'après
Homologie, lemme \ref{homology-lemma-double-complex-gives-resolution}.
D'un autre côté, puisque la restriction à $\mathcal{C}_p$ est exacte,
(lemme \ref{spaces-simplicial-lemma-restriction-to-components-site})
le complexe $\mathcal{I}_p^\bullet$ représente $K_p$ dans
$D(\mathcal{C}_p)$. Les faisceaux $\mathcal{I}_p^q$
sont totalement acycliques sur $\mathcal{C}_p$
(lemme \ref{spaces-simplicial-lemma-restriction-injective-to-component-limp}).
La cohomologie des colonnes calcule donc les groupes
$H^q(\mathcal{C}_p, K_p)$ par le lemme d'acyclicité de Leray
(Catégories dérivées, lemme \ref{derived-lemma-leray-acyclicity})
et
Cohomologie sur les sites, lemme \ref{sites-cohomology-lemma-limp-acyclic}.
On conclut en appliquant
Homologie, lemme \ref{homology-lemma-first-quadrant-ss}.
\end{proof}

\begin{lemma}
\label{spaces-simplicial-lemma-sanity-check-modules}
Dans la situation \ref{spaces-simplicial-situation-simplicial-site}, soit $\mathcal{O}$
un faisceau d'anneaux. Soient $U \in \Ob(\mathcal{C}_n)$ et
$\mathcal{F} \in \textit{Mod}(\mathcal{O})$.
Alors $H^p(U, \mathcal{F}) = H^p(U, g_n^*\mathcal{F})$
où, au membre de gauche, $U$ est considéré comme un objet de
$\mathcal{C}_{total}$.
\end{lemma}

\begin{proof}
Observons que « $U$ considéré comme objet de $\mathcal{C}_{total}$ »
est expliqué par la construction de $\mathcal{C}_{total}$ dans le
lemme \ref{spaces-simplicial-lemma-simplicial-site-site} dans le cas (A), et dans le
lemme \ref{spaces-simplicial-lemma-simplicial-cocontinuous-site} dans le cas (B).
Dans les deux cas, le foncteur $\mathcal{C}_n \to \mathcal{C}$
est continu et cocontinu, voir
le lemme \ref{spaces-simplicial-lemma-restriction-to-components-site}, et
$g_n^{-1}\mathcal{O} = \mathcal{O}_n$ par définition.
Le résultat est donc un cas particulier du lemme
de Cohomologie sur les sites
\ref{sites-cohomology-lemma-pullback-same-cohomology}.
\end{proof}






\section{Cohomologie et augmentations des sites simpliciaux annelés}
\label{spaces-simplicial-section-cohomology-augmentation-ringed-simplicial-sites}

\noindent
Cette section est l'analogue de
la section \ref{spaces-simplicial-section-cohomology-augmentation-simplicial-sites}
pour les faisceaux de modules.

\medskip\noindent
Considérons un site simplicial $\mathcal{C}$ comme dans
la situation \ref{spaces-simplicial-situation-simplicial-site}.
Soit $a_0$ une augmentation vers un site $\mathcal{D}$ comme dans
la remarque \ref{spaces-simplicial-remark-augmentation-site}.
Soit $\mathcal{O}$ un faisceau d'anneaux sur $\mathcal{C}_{total}$.
Soit $\mathcal{O}_\mathcal{D}$ un faisceau d'anneaux sur $\mathcal{D}$.
Supposons qu'on nous donne un morphisme
$$
a^\sharp : \mathcal{O}_\mathcal{D} \longrightarrow a_*\mathcal{O}
$$
où $a$ est comme dans le lemme \ref{spaces-simplicial-lemma-augmentation-site}.
Par conséquent, on obtient un morphisme de topoi annelés
$$
a :
(\Sh(\mathcal{C}_{total}), \mathcal{O})
\longrightarrow
(\Sh(\mathcal{D}), \mathcal{O}_\mathcal{D})
$$
Considérons le morphisme de topos annelés $g_n : (\Sh(\mathcal{C}_n), \mathcal{O}_n) \to
(\Sh(\mathcal{C}_{total}), \mathcal{O})$
du
lemme \ref{spaces-simplicial-lemma-restriction-module-to-components-site}. En formant
la composition $a_n = a \circ g_n$
(lemme \ref{spaces-simplicial-lemma-augmentation-site})
dans la catégorie des topoi annelés, on obtient
$$
a_n :
(\Sh(\mathcal{C}_n), \mathcal{O}_n)
\longrightarrow
(\Sh(\mathcal{D}), \mathcal{O}_\mathcal{D})
$$
À l'aide des morphismes de transition $f_\varphi^{-1}\mathcal{O}_m \to \mathcal{O}_n$
on obtient des morphismes de topoi annelés
$$
f_\varphi : (\Sh(\mathcal{C}_n), \mathcal{O}_n) \to
(\Sh(\mathcal{C}_m), \mathcal{O}_m)
$$
tels que $a_m \circ f_\varphi = a_n$ comme morphismes de topoi annelés
pour tout $\varphi : [m] \to [n]$.

\begin{lemma}
\label{spaces-simplicial-lemma-flat-augmentation-modules}
Avec les notations ci-dessus, le morphisme
$a : (\Sh(\mathcal{C}_{total}), \mathcal{O}) \to
(\Sh(\mathcal{D}), \mathcal{O}_\mathcal{D})$
est plat si et seulement si
$a_n : (\Sh(\mathcal{C}_n), \mathcal{O}_n) \to
(\Sh(\mathcal{D}), \mathcal{O}_\mathcal{D})$
est plat pour $n \geq 0$.
\end{lemma}

\begin{proof}
Puisque $g_n : (\Sh(\mathcal{C}_n), \mathcal{O}_n) \to
(\Sh(\mathcal{C}_{total}), \mathcal{O})$ est plat, on voit
que, si $a$ est plat, alors $a_n = a \circ g_n$ est plat comme
composé. Réciproquement, supposons $a_n$ plat pour tout $n$.
Il faut vérifier que $\mathcal{O}$ est plat comme faisceau de
$a^{-1}\mathcal{O}_\mathcal{D}$-modules. Soit $\mathcal{F} \to \mathcal{G}$
un morphisme injectif de $a^{-1}\mathcal{O}_\mathcal{D}$-modules.
Nous devons montrer que
$$
\mathcal{F} \otimes_{a^{-1}\mathcal{O}_\mathcal{D}} \mathcal{O}
\to
\mathcal{G} \otimes_{a^{-1}\mathcal{O}_\mathcal{D}} \mathcal{O}
$$
est injectif. On peut le vérifier sur $\mathcal{C}_n$, c'est-à-dire après
application de $g_n^{-1}$. Puisque $g_n^* = g_n^{-1}$ et que
$g_n^{-1}\mathcal{O} = \mathcal{O}_n$, on obtient
$$
g_n^{-1}\mathcal{F} \otimes_{g_n^{-1}a^{-1}\mathcal{O}_\mathcal{D}}
\mathcal{O}_n
\to
g_n^{-1}\mathcal{G} \otimes_{g_n^{-1}a^{-1}\mathcal{O}_\mathcal{D}}
\mathcal{O}_n
$$
qui est injectif car
$g_n^{-1}a^{-1}\mathcal{O}_\mathcal{D} = a_n^{-1}\mathcal{O}_\mathcal{D}$
et, par hypothèse, $a_n$ est plat.
\end{proof}

\begin{lemma}
\label{spaces-simplicial-lemma-simplicial-resolution-augmentation-modules}
Avec les notations ci-dessus, pour tout $\mathcal{O}_\mathcal{D}$-module $\mathcal{G}$,
il existe un complexe exact
$$
\ldots \to
g_{2!}(a_2^*\mathcal{G}) \to
g_{1!}(a_1^*\mathcal{G}) \to
g_{0!}(a_0^*\mathcal{G}) \to
a^*\mathcal{G} \to 0
$$
de faisceaux de $\mathcal{O}$-modules sur $\mathcal{C}_{total}$.
Ici, $g_{n!}$ est comme dans le lemme \ref{spaces-simplicial-lemma-restriction-module-to-components-site}.
\end{lemma}

\begin{proof}
Observons que $a^*\mathcal{G}$ est le $\mathcal{O}$-module sur
$\mathcal{C}_{total}$ dont la restriction à $\mathcal{C}_m$
est le $\mathcal{O}_m$-module $a_m^*\mathcal{G}$.
La description des foncteurs $g_{n!}$ sur les modules
dans le lemme \ref{spaces-simplicial-lemma-restriction-module-to-components-site}
montre que $g_{n!}(a_n^*\mathcal{G})$ est le
$\mathcal{O}$-module sur $\mathcal{C}_{total}$
dont la restriction à $\mathcal{C}_m$ est le $\mathcal{O}_m$-module
$$
\bigoplus\nolimits_{\varphi : [n] \to [m]} f_\varphi^*a_n^*\mathcal{G} =
\bigoplus\nolimits_{\varphi : [n] \to [m]} a_m^*\mathcal{G}
$$
Le reste de la démonstration est identique à celui du lemme
\ref{spaces-simplicial-lemma-simplicial-resolution-augmentation},
après remplacement de $a_m^{-1}\mathcal{G}$ par $a_m^*\mathcal{G}$.
\end{proof}

\begin{lemma}
\label{spaces-simplicial-lemma-augmentation-cech-complex-modules}
Avec les notations ci-dessus.
Pour un $\mathcal{O}$-module $\mathcal{F}$ sur $\mathcal{C}_{total}$
il existe un complexe canonique
$$
0 \to a_*\mathcal{F} \to a_{0, *}\mathcal{F}_0 \to a_{1, *}\mathcal{F}_1 \to
a_{2, *}\mathcal{F}_2 \to \ldots
$$
de $\mathcal{O}_\mathcal{D}$-modules, exact en degrés $-1, 0$.
Si $\mathcal{F}$ est un $\mathcal{O}$-module injectif, alors le complexe
est exact à tous les degrés et reste exact en appliquant le foncteur
$\Hom_{\mathcal{O}_\mathcal{D}}(\mathcal{G}, -)$ pour tout
$\mathcal{O}_\mathcal{D}$-module $\mathcal{G}$.
\end{lemma}

\begin{proof}
Pour construire le complexe, il suffit, par le lemme de Yoneda, de construire pour tout
$\mathcal{O}_\mathcal{D}$-module $\mathcal{G}$ sur $\mathcal{D}$
un complexe
$$
0 \to \Hom_{\mathcal{O}_\mathcal{D}}(\mathcal{G}, a_*\mathcal{F}) \to
\Hom_{\mathcal{O}_\mathcal{D}}(\mathcal{G}, a_{0, *}\mathcal{F}_0) \to
\Hom_{\mathcal{O}_\mathcal{D}}(\mathcal{G}, a_{1, *}\mathcal{F}_1) \to \ldots
$$
fonctoriellement en $\mathcal{G}$. Pour cela, appliquons
$\Hom_\mathcal{O}(-, \mathcal{F})$
au complexe exact du
lemme \ref{spaces-simplicial-lemma-simplicial-resolution-augmentation-modules}
et utilisons l'adjonction entre image inverse et image directe.
Les propriétés d'exactitude en degrés $-1, 0$ découlent de
la construction, puisque $\Hom_\mathcal{O}(-, \mathcal{F})$ est exact à gauche.
Si $\mathcal{F}$ est un $\mathcal{O}$-module injectif, alors le complexe
est exact car $\Hom_\mathcal{O}(-, \mathcal{F})$ est exact.
\end{proof}

\begin{lemma}
\label{spaces-simplicial-lemma-augmentation-spectral-sequence-modules}
Avec les notations ci-dessus, pour tout $K$ dans $D^+(\mathcal{O})$, il existe une suite
spectrale $(E_r, d_r)_{r \geq 0}$ dans $\textit{Mod}(\mathcal{O}_\mathcal{D})$
avec
$$
E_1^{p, q} = R^qa_{p, *} K_p
$$
qui aboutit à $R^{p + q}a_*K$. Cette suite spectrale est fonctorielle en $K$.
\end{lemma}

\begin{proof}
Soit $\mathcal{I}^\bullet$ un complexe borné inférieurement de
$\mathcal{O}$-modules injectifs représentant $K$. Considérons le complexe double de termes
$$
A^{p, q} = a_{p, *}\mathcal{I}^q_p
$$
où les flèches horizontales proviennent du
lemme \ref{spaces-simplicial-lemma-augmentation-cech-complex-modules}
et les flèches verticales des différentiels du complexe
$\mathcal{I}^\bullet$. Le lemme
affirme que les lignes du complexe double sont exactes
en degrés positifs et donnent
$a_*\mathcal{I}^q$ en degré $0$.
Ainsi le complexe total associé au double complexe
calcule $Ra_*K$ d'après
Homologie, lemme \ref{homology-lemma-double-complex-gives-resolution}.
D'autre part, puisque la restriction à $\mathcal{C}_p$ est exacte
(lemme \ref{spaces-simplicial-lemma-restriction-to-components-site})
le complexe $\mathcal{I}_p^\bullet$ représente $K_p$ dans
$D(\mathcal{C}_p)$. Les faisceaux $\mathcal{I}_p^q$
sont totalement acycliques sur $\mathcal{C}_p$
(lemme \ref{spaces-simplicial-lemma-restriction-injective-to-component-limp}).
La cohomologie des colonnes est donc donnée par les faisceaux
$R^qa_{p, *}K_p$ par le lemme d'acyclicité de Leray
(Catégories dérivées, lemme \ref{derived-lemma-leray-acyclicity})
et
Cohomologie sur les sites, lemme \ref{sites-cohomology-lemma-limp-acyclic}.
On conclut en appliquant
Homologie, lemme \ref{homology-lemma-first-quadrant-ss}.
\end{proof}





\section{Faisceaux et modules cartésiens}
\label{spaces-simplicial-section-cartesian}

\noindent
Voici la définition.

\begin{definition}
\label{spaces-simplicial-definition-cartesian-sheaf}
Dans la situation \ref{spaces-simplicial-situation-simplicial-site}.
\begin{enumerate}
\item Un faisceau $\mathcal{F}$ d'ensembles ou de groupes abéliens sur
$\mathcal{C}_{total}$ est {\it cartésien} si les morphismes
$\mathcal{F}(\varphi) : f_\varphi^{-1}\mathcal{F}_m \to \mathcal{F}_n$
sont des isomorphismes pour tout $\varphi : [m] \to [n]$.
\item Si $\mathcal{O}$ est un faisceau d'anneaux sur $\mathcal{C}_{total}$,
alors un faisceau $\mathcal{F}$ de $\mathcal{O}$-modules est
{\it cartésien} si les morphismes $f_\varphi^*\mathcal{F}_m \to \mathcal{F}_n$
sont des isomorphismes pour tout $\varphi : [m] \to [n]$.
\item Un objet $K$ de $D(\mathcal{C}_{total})$ est {\it cartésien} si les morphismes
$f_\varphi^{-1}K_m \to K_n$
sont des isomorphismes pour tout $\varphi : [m] \to [n]$.
\item Si $\mathcal{O}$ est un faisceau d'anneaux sur $\mathcal{C}_{total}$, alors
un objet $K$ de $D(\mathcal{O})$ est {\it cartésien} si les morphismes
$Lf_\varphi^*K_m \to K_n$
sont des isomorphismes pour tout $\varphi : [m] \to [n]$.
\end{enumerate}
\end{definition}

\noindent
Bien entendu, il existe une notion générale de section cartésienne d'une
catégorie fibrée, dont les notions ci-dessus ne sont que des exemples.
La propriété sur les images inverses doit uniquement être vérifiée pour les dégénérescences.

\begin{lemma}
\label{spaces-simplicial-lemma-check-cartesian-module}
Dans la situation \ref{spaces-simplicial-situation-simplicial-site}.
\begin{enumerate}
\item Un faisceau $\mathcal{F}$ d'ensembles ou de groupes abéliens est cartésien
si et seulement si les morphismes
$(f_{\delta^n_j})^{-1}\mathcal{F}_{n - 1} \to \mathcal{F}_n$
sont des isomorphismes.
\item Un objet $K$ de $D(\mathcal{C}_{total})$ est cartésien
si et seulement si les morphismes
$(f_{\delta^n_j})^{-1}K_{n - 1} \to K_n$
sont des isomorphismes.
\item Si $\mathcal{O}$ est un faisceau d'anneaux sur $\mathcal{C}_{total}$,
un faisceau $\mathcal{F}$ de $\mathcal{O}$-modules est cartésien
si et seulement si les morphismes
$(f_{\delta^n_j})^*\mathcal{F}_{n - 1} \to \mathcal{F}_n$
sont des isomorphismes.
\item Si $\mathcal{O}$ est un faisceau d'anneaux sur $\mathcal{C}_{total}$,
un objet $K$ de $D(\mathcal{O})$ est cartésien
si et seulement si les morphismes
$L(f_{\delta^n_j})^*K_{n - 1} \to K_n$
sont des isomorphismes.
\item Ajouter d'autres cas ici.
\end{enumerate}
\end{lemma}

\begin{proof}
Dans chaque cas, le point essentiel est que les foncteurs image inverse
se composent en un foncteur du même type; pour la partie (4), voir
Cohomologie sur les sites, lemme
\ref{sites-cohomology-lemma-derived-pullback-composition}.
Nous montrons comment fonctionne l'argument dans le cas (1) et omettons la preuve
dans les autres cas.
La catégorie $\Delta$ est engendrée par les morphismes
$\delta^n_j$ et $\sigma^n_j$; voir
Simplicial, lemme \ref{simplicial-lemma-face-degeneracy}.
Il suffit donc de vérifier que les morphismes
$(f_{\delta^n_j})^{-1}\mathcal{F}_{n - 1} \to \mathcal{F}_n$
et $(f_{\sigma^n_j})^{-1}\mathcal{F}_{n + 1} \to \mathcal{F}_n$ sont
des isomorphismes; voir
Simplicial, lemme \ref{simplicial-lemma-characterize-simplicial-object}
pour les notations. Puisque $\sigma^n_j \circ \delta_j^{n + 1} = \text{id}_{[n]}$
la composition
$$
\mathcal{F}_n =
(f_{\sigma^n_j})^{-1}
(f_{\delta_j^{n + 1}})^{-1}
\mathcal{F}_n \to
(f_{\sigma^n_j})^{-1}
\mathcal{F}_{n + 1} \to
\mathcal{F}_n
$$
est l'identité. Ainsi, le résultat pour $\delta^{n + 1}_j$ implique le résultat
pour $\sigma^n_j$.
\end{proof}

\begin{lemma}
\label{spaces-simplicial-lemma-augmentation-cartesian-module}
Dans la situation \ref{spaces-simplicial-situation-simplicial-site}, soit
$a_0$ une augmentation vers un site $\mathcal{D}$ comme dans la
remarque \ref{spaces-simplicial-remark-augmentation-site}.
\begin{enumerate}
\item L'image inverse $a^{-1}\mathcal{G}$ d'un faisceau d'ensembles ou de groupes abéliens
sur $\mathcal{D}$ est cartésienne.
\item L'image inverse $a^{-1}K$ d'un objet $K$ de $D(\mathcal{D})$
est cartésienne.
\end{enumerate}
Soit $\mathcal{O}$ un faisceau d'anneaux sur $\mathcal{C}_{total}$ et
$\mathcal{O}_\mathcal{D}$ un faisceau d'anneaux sur $\mathcal{D}$
et $a^\sharp : \mathcal{O}_\mathcal{D} \to a_*\mathcal{O}$ un
morphisme comme dans la section
\ref{spaces-simplicial-section-cohomology-augmentation-ringed-simplicial-sites}.
\begin{enumerate}
\item[(3)] L'image inverse $a^*\mathcal{F}$ d'un faisceau de
$\mathcal{O}_\mathcal{D}$-modules est cartésienne.
\item[(4)] L'image inverse dérivée $La^*K$ d'un objet
$K$ de $D(\mathcal{O}_\mathcal{D})$ est cartésienne.
\end{enumerate}
\end{lemma}

\begin{proof}
Cela découle immédiatement des identités
$a_m \circ f_\varphi = a_n$ pour tout $\varphi : [m] \to [n]$.
Voir le lemme \ref{spaces-simplicial-lemma-augmentation-site} et la discussion dans la section
\ref{spaces-simplicial-section-cohomology-augmentation-ringed-simplicial-sites}.
\end{proof}

\begin{lemma}
\label{spaces-simplicial-lemma-characterize-cartesian}
Dans la situation \ref{spaces-simplicial-situation-simplicial-site}.
La catégorie des faisceaux cartésiens d'ensembles (resp.\ de groupes abéliens)
est équivalente à la catégorie des couples $(\mathcal{F}, \alpha)$
où $\mathcal{F}$ est un faisceau d'ensembles (resp.\ de groupes abéliens)
sur $\mathcal{C}_0$ et
$$
\alpha :
(f_{\delta_1^1})^{-1}\mathcal{F}
\longrightarrow (f_{\delta_0^1})^{-1}\mathcal{F}
$$
est un isomorphisme de faisceaux d'ensembles (resp.\ de groupes abéliens)
sur $\mathcal{C}_1$ tel que
$(f_{\delta^2_1})^{-1}\alpha =
(f_{\delta^2_0})^{-1}\alpha \circ (f_{\delta^2_2})^{-1}\alpha$
comme morphismes de faisceaux sur $\mathcal{C}_2$.
\end{lemma}

\begin{proof}
Nous abrégeons
$d^n_j = f_{\delta^n_j} : \Sh(\mathcal{C}_n) \to \Sh(\mathcal{C}_{n - 1})$.
La condition imposée à $\alpha$ dans l'énoncé du lemme est bien définie, car
$$
d^1_1 \circ d^2_2 = d^1_1 \circ d^2_1, \quad
d^1_1 \circ d^2_0 = d^1_0 \circ d^2_2, \quad
d^1_0 \circ d^2_0 = d^1_0 \circ d^2_1
$$
comme morphismes de topoi $\Sh(\mathcal{C}_2) \to \Sh(\mathcal{C}_0)$, voir
Simplicial, remarque \ref{simplicial-remark-relations}. On peut donc
représenter ces morphismes comme suit
$$
\xymatrix{
& (d^2_0)^{-1}(d^1_1)^{-1}\mathcal{F} \ar[r]_-{(d^2_0)^{-1}\alpha} &
(d^2_0)^{-1}(d^1_0)^{-1}\mathcal{F} \ar@{=}[rd] & \\
(d^2_2)^{-1}(d^1_0)^{-1}\mathcal{F} \ar@{=}[ru] & & &
(d^2_1)^{-1}(d^1_0)^{-1}\mathcal{F} \\
& (d^2_2)^{-1}(d^1_1)^{-1}\mathcal{F} \ar[lu]^{(d^2_2)^{-1}\alpha} \ar@{=}[r] &
(d^2_1)^{-1}(d^1_1)^{-1}\mathcal{F} \ar[ru]_{(d^2_1)^{-1}\alpha}
}
$$
et la condition signifie que le diagramme est commutatif. Or,
étant donné un faisceau cartésien $\mathcal{G}$ d'ensembles (resp.\ de groupes abéliens)
sur $\mathcal{C}_{total}$,
on peut poser $\mathcal{F} = \mathcal{G}_0$ et prendre pour $\alpha$ le composé
$$
(d_1^1)^{-1}\mathcal{G}_0 \to \mathcal{G}_1
\leftarrow (d_1^0)^{-1}\mathcal{G}_0
$$
où les flèches sont inversibles puisque $\mathcal{G}$ est cartésien.
Pour prouver que ce foncteur
est une équivalence, construisons-en un quasi-inverse. La construction du
quasi-inverse est analogue à celle décrite dans
Descente, section \ref{descent-section-descent-modules}, à laquelle nous empruntons
les notations $\tau^n_i : [0] \to [n]$, $0 \mapsto i$ et
$\tau^n_{ij} : [1] \to [n]$, $0 \mapsto i$, $1 \mapsto j$.
Plus précisément, pour un couple $(\mathcal{F}, \alpha)$
comme dans le lemme, posons $\mathcal{G}_n = (f_{\tau^n_n})^{-1}\mathcal{F}$.
Pour $\varphi : [n] \to [m]$, définissons
$\mathcal{G}(\varphi) : (f_\varphi)^{-1}\mathcal{G}_n \to \mathcal{G}_m$
en utilisant
$$
\xymatrix{
(f_\varphi)^{-1}\mathcal{G}_n \ar@{=}[r] &
(f_\varphi)^{-1}(f_{\tau^n_n})^{-1}\mathcal{F} \ar@{=}[r] &
(f_{\tau^m_{\varphi(n)}})^{-1}\mathcal{F} \ar@{=}[r] &
(f_{\tau^m_{\varphi(n)m}})^{-1}(d^1_1)^{-1}\mathcal{F}
\ar[d]^{(f_{\tau^m_{\varphi(n)m}})^{-1}\alpha} \\
&
\mathcal{G}_m \ar@{=}[r] &
(f_{\tau^m_m})^{-1}\mathcal{F} \ar@{=}[r] &
(f_{\tau^m_{\varphi(n)m}})^{-1}(d^1_0)^{-1}\mathcal{F}
}
$$
Nous laissons au lecteur la vérification que la commutativité du diagramme
ci-dessus entraîne la compatibilité des morphismes avec la composition et définit donc un
faisceau sur $\mathcal{C}_{total}$; voir
le lemme \ref{spaces-simplicial-lemma-describe-sheaves-simplicial-site-site}.
Nous laissons également au lecteur le soin de vérifier
que les deux foncteurs sont quasi inverses l'un de l'autre.
\end{proof}

\begin{lemma}
\label{spaces-simplicial-lemma-characterize-cartesian-modules}
Dans la situation \ref{spaces-simplicial-situation-simplicial-site},
soit $\mathcal{O}$ un faisceau d'anneaux sur $\mathcal{C}_{total}$.
La catégorie des $\mathcal{O}$-modules cartésiens
est équivalente à la catégorie des couples $(\mathcal{F}, \alpha)$
où $\mathcal{F}$ est un $\mathcal{O}_0$-module
et
$$
\alpha :
(f_{\delta_1^1})^*\mathcal{F}
\longrightarrow (f_{\delta_0^1})^*\mathcal{F}
$$
est un isomorphisme de $\mathcal{O}_1$-modules tel que
$(f_{\delta^2_1})^*\alpha =
(f_{\delta^2_0})^*\alpha \circ (f_{\delta^2_2})^*\alpha$
comme morphismes de $\mathcal{O}_2$-modules.
\end{lemma}

\begin{proof}
La démonstration est identique à celle du
lemme \ref{spaces-simplicial-lemma-characterize-cartesian},
l'image inverse des faisceaux de groupes abéliens étant remplacée
par l'image inverse des modules.
\end{proof}

\begin{lemma}
\label{spaces-simplicial-lemma-Serre-subcat-cartesian-modules}
Dans la situation \ref{spaces-simplicial-situation-simplicial-site}.
\begin{enumerate}
\item La sous-catégorie pleine des faisceaux abéliens cartésiens est une
sous-catégorie de Serre faible de $\textit{Ab}(\mathcal{C}_{total})$.
Les colimites de systèmes de faisceaux abéliens cartésiens sont cartésiennes.
\item Soit $\mathcal{O}$ un faisceau d'anneaux sur $\mathcal{C}_{total}$
tel que les morphismes
$$
f_{\delta^n_j} : (\Sh(\mathcal{C}_n), \mathcal{O}_n)
\to (\Sh(\mathcal{C}_{n - 1}), \mathcal{O}_{n - 1})
$$
sont plats. La sous-catégorie pleine des $\mathcal{O}$-modules cartésiens est une
sous-catégorie de Serre faible de $\textit{Mod}(\mathcal{O})$.
Les colimites de systèmes de $\mathcal{O}$-modules cartésiens sont cartésiennes.
\end{enumerate}
\end{lemma}

\begin{proof}
Pour voir qu'on obtient en (1) une sous-catégorie de Serre faible,
vérifions les conditions énumérées dans
Homologie, lemme \ref{homology-lemma-characterize-weak-serre-subcategory}.
D'abord, si $\varphi : \mathcal{F} \to \mathcal{G}$ est un morphisme
entre faisceaux abéliens cartésiens, alors
$\Ker(\varphi)$ et $\Coker(\varphi)$ sont également cartésiens
car les foncteurs de restriction
$\Sh(\mathcal{C}_{total}) \to \Sh(\mathcal{C}_n)$
et les foncteurs $f_\varphi^{-1}$ sont exacts.
De même, si
$$
0 \to \mathcal{F} \to \mathcal{H} \to \mathcal{G} \to 0
$$
est une suite exacte courte de faisceaux abéliens sur $\mathcal{C}_{total}$
avec $\mathcal{F}$ et $\mathcal{G}$ cartésiens, il s'ensuit que
$\mathcal{H}$ est cartésien par le lemme des cinq. Pour les colimites,
utilisons le fait qu'elles commutent aux images inverses, puisque les foncteurs image inverse sont adjoints
à gauche. Dans le cas des modules,
on raisonne de même en utilisant l'exactitude de l'image inverse par un morphisme plat
(Modules sur les sites, lemme \ref{sites-modules-lemma-flat-pullback-exact})
et le fait qu'il suffit de vérifier la condition
pour $f_{\delta^n_j}$; voir le lemme \ref{spaces-simplicial-lemma-check-cartesian-module}.
\end{proof}

\begin{remark}[Avertissement]
\label{spaces-simplicial-remark-warning-cartesian-modules}
Malgré le lemme \ref{spaces-simplicial-lemma-Serre-subcat-cartesian-modules}, il
peut arriver que la catégorie des $\mathcal{O}$-modules cartésiens soit abélienne
sans être une sous-catégorie de Serre de $\textit{Mod}(\mathcal{O})$.
Plus précisément, supposons que nous sachions seulement que
$f_{\delta_1^1}$ et $f_{\delta_0^1}$ sont plats.
Il résulte alors facilement du lemme
\ref{spaces-simplicial-lemma-characterize-cartesian-modules}
que la catégorie des $\mathcal{O}$-modules cartésiens est abélienne.
Mais si $f_{\delta_0^2}$ n'est pas plat (par exemple),
il n'y a aucune raison pour que le foncteur d'inclusion
de la catégorie des $\mathcal{O}$-modules cartésiens
dans la catégorie de tous les $\mathcal{O}$-modules soit exact.
\end{remark}

\begin{lemma}
\label{spaces-simplicial-lemma-derived-cartesian-modules}
Dans la situation \ref{spaces-simplicial-situation-simplicial-site}.
\begin{enumerate}
\item Un objet $K$ de $D(\mathcal{C}_{total})$ est cartésien si et seulement
si $H^q(K)$ est un faisceau abélien cartésien pour tout $q$.
\item Soit $\mathcal{O}$ un faisceau
d'anneaux sur $\mathcal{C}_{total}$ tel que les morphismes
$f_{\delta^n_j} : (\Sh(\mathcal{C}_n), \mathcal{O}_n)
\to (\Sh(\mathcal{C}_{n - 1}), \mathcal{O}_{n - 1})$ sont plats.
Alors un objet $K$ de $D(\mathcal{O})$ est cartésien si et seulement
si $H^q(K)$ est un $\mathcal{O}$-module cartésien pour tout $q$.
\end{enumerate}
\end{lemma}

\begin{proof}
La partie (1) résulte de ce que les foncteurs image inverse $(f_\varphi)^{-1}$
sont exacts. La partie (2) découle de la caractérisation dans le lemme
\ref{spaces-simplicial-lemma-check-cartesian-module}
et du fait que $L(f_{\delta^n_j})^* = (f_{\delta^n_j})^*$
par planéité.
\end{proof}

\begin{lemma}
\label{spaces-simplicial-lemma-derived-cartesian-shriek}
Dans la situation \ref{spaces-simplicial-situation-simplicial-site}.
\begin{enumerate}
\item Un objet $K$ de $D(\mathcal{C}_{total})$ est cartésien si et seulement
le morphisme canonique
$$
g_{n!}K_n \longrightarrow
g_{n!}\mathbf{Z} \otimes^\mathbf{L}_\mathbf{Z} K
$$
est un isomorphisme pour tout $n$.
\item Soit $\mathcal{O}$ un faisceau d'anneaux sur $\mathcal{C}_{total}$
tel que les morphismes $f_\varphi^{-1}\mathcal{O}_n \to \mathcal{O}_m$
sont plats pour tout $\varphi : [n] \to [m]$. Alors un objet $K$ de
$D(\mathcal{O})$ est cartésien si et seulement si le morphisme canonique
$$
g_{n!}K_n \longrightarrow
g_{n!}\mathcal{O}_n \otimes^\mathbf{L}_\mathcal{O} K
$$
est un isomorphisme pour tout $n$.
\end{enumerate}
\end{lemma}

\begin{proof}
Pour (1). Puisque $g_{n!}$ est exact, il induit un foncteur
sur les catégories dérivées adjoint à $g_n^{-1}$.
Le morphisme est adjoint au morphisme
$K_n \to (g_n^{-1}g_{n!}\mathbf{Z}) \otimes^\mathbf{L}_\mathbf{Z} K_n$
correspondant à $\mathbf{Z} \to g_n^{-1}g_{n!}\mathbf{Z}$
qui à son tour est adjoint à
$\text{id} : g_{n!}\mathbf{Z} \to g_{n!}\mathbf{Z}$.
En utilisant la description de $g_{n!}$
donnée dans le lemme \ref{spaces-simplicial-lemma-restriction-to-components-site}
on voit que la restriction à $\mathcal{C}_m$ de ce morphisme
est
$$
\bigoplus\nolimits_{\varphi : [n] \to [m]} f_\varphi^{-1}K_n
\longrightarrow
\bigoplus\nolimits_{\varphi : [n] \to [m]} K_m
$$
L'assertion est donc claire.

\medskip\noindent
Pour (2). Puisque $g_{n!}$ est exact
(lemme \ref{spaces-simplicial-lemma-exactness-g-shriek-modules}), il induit un foncteur
sur les catégories dérivées adjoint à $g_n^*$ (également exact).
Le morphisme est adjoint au morphisme
$K_n \to (g_n^*g_{n!}\mathcal{O}_n) \otimes^\mathbf{L}_{\mathcal{O}_n} K_n$
correspondant à $\mathcal{O}_n \to g_n^*g_{n!}\mathcal{O}_n$
qui à son tour est adjoint à
$\text{id} : g_{n!}\mathcal{O}_n \to g_{n!}\mathcal{O}_n$.
En utilisant la description de $g_{n!}$
donnée dans le lemme \ref{spaces-simplicial-lemma-restriction-module-to-components-site}
on voit que la restriction à $\mathcal{C}_m$ de ce morphisme
est
$$
\bigoplus\nolimits_{\varphi : [n] \to [m]} f_\varphi^*K_n
\longrightarrow
\bigoplus\nolimits_{\varphi : [n] \to [m]}
f_\varphi^*\mathcal{O}_n \otimes_{\mathcal{O}_m} K_m =
\bigoplus\nolimits_{\varphi : [n] \to [m]} K_m
$$
L'assertion est donc claire.
\end{proof}

\begin{lemma}
\label{spaces-simplicial-lemma-quasi-coherent-sheaf}
Dans la situation \ref{spaces-simplicial-situation-simplicial-site},
soit $\mathcal{O}$ un faisceau d'anneaux sur $\mathcal{C}_{total}$.
Soit $\mathcal{F}$ un faisceau de $\mathcal{O}$-modules.
Alors $\mathcal{F}$ est quasi-cohérent au sens de
Modules sur les sites, définition \ref{sites-modules-definition-site-local},
si et seulement si $\mathcal{F}$ est cartésien
et si $\mathcal{F}_n$ est un $\mathcal{O}_n$-module quasi-cohérent pour tout $n$.
\end{lemma}

\begin{proof}
Supposons que $\mathcal{F}$ soit quasi-cohérent. Puisque les images inverses des modules quasi-cohérents
sont quasi-cohérents
(Modules sur les sites, lemme \ref{sites-modules-lemma-local-pullback}),
nous voyons que $\mathcal{F}_n$ est un module $\mathcal{O}_n$ quasi-cohérent
pour tout $n$. Pour montrer que $\mathcal{F}$ est cartésien, soit $U$
un objet de $\mathcal{C}_n$ pour un certain $n$. Regardons $U$
comme un objet de $\mathcal{C}_{total}$. Comme $\mathcal{F}$
est quasi-cohérent, il existe un recouvrement $\{U_i \to U\}$
et pour chaque $i$ une présentation
$$
\bigoplus\nolimits_{j \in J_i} \mathcal{O}_{\mathcal{C}_{total}/U_i} \to
\bigoplus\nolimits_{k \in K_i} \mathcal{O}_{\mathcal{C}_{total}/U_i} \to
\mathcal{F}|_{\mathcal{C}_{total}/U_i} \to 0
$$
Notons que $\{U_i \to U\}$ est un recouvrement de $\mathcal{C}_n$ par
la construction du site $\mathcal{C}_{total}$.
Soit ensuite $V$ un objet de $\mathcal{C}_m$ pour un certain $m$, et soit
$V \to U$ un morphisme de $\mathcal{C}_{total}$ au-dessus de
$\varphi : [n] \to [m]$. Les produits fibrés $V_i = V \times_U U_i$
existent et donnent un recouvrement induit $\{V_i \to V\}$ dans $\mathcal{C}_m$.
En restreignant la présentation ci-dessus aux sites
$\mathcal{C}_n/U_i$ et $\mathcal{C}_m/V_i$, on obtient
les présentations
$$
\bigoplus\nolimits_{j \in J_i} \mathcal{O}_{\mathcal{C}_m/U_i} \to
\bigoplus\nolimits_{k \in K_i} \mathcal{O}_{\mathcal{C}_m/U_i} \to
\mathcal{F}_n|_{\mathcal{C}_n/U_i} \to 0
$$
et
$$
\bigoplus\nolimits_{j \in J_i} \mathcal{O}_{\mathcal{C}_m/V_i} \to
\bigoplus\nolimits_{k \in K_i} \mathcal{O}_{\mathcal{C}_m/V_i} \to
\mathcal{F}_m|_{\mathcal{C}_m/V_i} \to 0
$$
Ces présentations sont compatibles avec le morphisme
$\mathcal{F}(\varphi) : f_\varphi^*\mathcal{F}_n \to \mathcal{F}_m$
(car ce morphisme est défini à l'aide des morphismes de restriction de $\mathcal{F}$
le long des morphismes de $\mathcal{C}_{total}$ recouvrant $\varphi$).
On en déduit que $\mathcal{F}(\varphi)|_{\mathcal{C}_m/V_i}$
est un isomorphisme. Comme $\{V_i \to V\}$ est un recouvrement,
$\mathcal{F}(\varphi)|_{\mathcal{C}_m/V}$ est donc un isomorphisme.
L'arbitraire de $V$ et de $U$ montre que $\mathcal{F}$ est cartésien.
(Dans le cas (A), on utilise Sites, lemme \ref{sites-lemma-morphism-of-sites-covering}.)

\medskip\noindent
Réciproquement, supposons $\mathcal{F}_n$
quasi-cohérent pour tout $n$ et que $\mathcal{F}$ est cartésien.
Alors, pour tout $n$ et tout objet $U$ de $\mathcal{C}_n$, nous
pouvons choisir un recouvrement $\{U_i \to U\}$ de $\mathcal{C}_n$
et pour chacun $i$ une présentation
$$
\bigoplus\nolimits_{j \in J_i} \mathcal{O}_{\mathcal{C}_m/U_i} \to
\bigoplus\nolimits_{k \in K_i} \mathcal{O}_{\mathcal{C}_m/U_i} \to
\mathcal{F}_n|_{\mathcal{C}_n/U_i} \to 0
$$
En prenant l'image inverse sur $\mathcal{C}_{total}/U_i$, on obtient les complexes
$$
\bigoplus\nolimits_{j \in J_i} \mathcal{O}_{\mathcal{C}_{total}/U_i} \to
\bigoplus\nolimits_{k \in K_i} \mathcal{O}_{\mathcal{C}_{total}/U_i} \to
\mathcal{F}|_{\mathcal{C}_{total}/U_i} \to 0
$$
de modules sur $\mathcal{C}_{total}/U_i$. La cartésianité de
$\mathcal{F}$ implique alors l'exactitude de ce complexe.
Nous omettons les détails.
\end{proof}











\section{Systèmes simpliciaux de la catégorie dérivée}
\label{spaces-simplicial-section-glueing}

\noindent
Dans cette section, nous allons prouver un cas particulier de
\cite[Proposition 3.2.9]{BBD} dans le cadre des catégories dérivées
de faisceaux abéliens. Le cas des modules
est abordé dans la section \ref{spaces-simplicial-section-glueing-modules}.

\begin{definition}
\label{spaces-simplicial-definition-cartesian-derived}
Dans la situation \ref{spaces-simplicial-situation-simplicial-site}. Un
{\it système simplicial de la catégorie dérivée}
se compose des données suivantes
\begin{enumerate}
\item pour tout $n$, un objet $K_n$ de $D(\mathcal{C}_n)$,
\item pour tout $\varphi : [m] \to [n]$, un morphisme
$K_\varphi : f_\varphi^{-1}K_m \to K_n$ dans $D(\mathcal{C}_n)$
\end{enumerate}
soumises à la condition
$$
K_{\varphi \circ \psi} = K_\varphi \circ f_\varphi^{-1}K_\psi :
f_{\varphi \circ \psi}^{-1}K_l = f_\varphi^{-1} f_\psi^{-1}K_l
\longrightarrow
K_n
$$
pour tout couple de morphismes $\varphi : [m] \to [n]$ et $\psi : [l] \to [m]$ de $\Delta$.
On dit que le système simplicial est {\it cartésien} si les morphismes $K_\varphi$
sont des isomorphismes pour tout $\varphi$.
Étant donné deux systèmes simpliciaux de la catégorie dérivée,
il existe une notion évidente de
{\it morphisme de systèmes simpliciaux de la catégorie dérivée}.
\end{definition}

\noindent
Nous avons intentionnellement donné à cette notion un nom ridiculement long.
Le but est de montrer qu'un système simplicial de la catégorie dérivée
provient d'un objet de $D(\mathcal{C}_{total})$ sous certaines
hypothèses.

\begin{lemma}
\label{spaces-simplicial-lemma-cartesian-objects-derived}
Dans la situation \ref{spaces-simplicial-situation-simplicial-site}.
Si $K \in D(\mathcal{C}_{total})$ est un objet,
alors $(K_n, K(\varphi))$ est un système simplicial de la catégorie dérivée.
Si $K$ est cartésien, le système l'est également.
\end{lemma}

\begin{proof}
C'est évident.
\end{proof}

\begin{lemma}
\label{spaces-simplicial-lemma-construct-cartesian-objects-derived}
Dans la situation \ref{spaces-simplicial-situation-simplicial-site},
supposons donnés $K_0 \in D(\mathcal{C}_0)$ et un isomorphisme
$$
\alpha :
f_{\delta_1^1}^{-1}K_0
\longrightarrow
f_{\delta_0^1}^{-1}K_0
$$
satisfaisant à la condition de cocycle. Posons
$\tau^n_i : [0] \to [n]$, $0 \mapsto i$ et
puis posons $K_n = f_{\tau^n_n}^{-1}K_0$. Les $K_n$
forment un système simplicial cartésien de la catégorie dérivée.
\end{lemma}

\begin{proof}
Comparer avec le lemme \ref{spaces-simplicial-lemma-characterize-cartesian}
et sa preuve (où la condition de cocycle est aussi explicitée).
La construction est analogue à celle décrite dans
Descente, section \ref{descent-section-descent-modules}, à laquelle
nous empruntons les notations $\tau^n_i : [0] \to [n]$, $0 \mapsto i$ et
$\tau^n_{ij} : [1] \to [n]$, $0 \mapsto i$, $1 \mapsto j$.
Pour $\varphi : [n] \to [m]$, définissons
$K_\varphi : f_\varphi^{-1}K_n \to K_m$
en utilisant
$$
\xymatrix{
f_\varphi^{-1}K_n \ar@{=}[r] &
f_\varphi^{-1} f_{\tau^n_n}^{-1}K_0 \ar@{=}[r] &
f_{\tau^m_{\varphi(n)}}^{-1}K_0 \ar@{=}[r] &
f_{\tau^m_{\varphi(n)m}}^{-1}f_{\delta^1_1}^{-1}K_0
\ar[d]_{f_{\tau^m_{\varphi(n)m}}^{-1}\alpha} \\
&
K_m \ar@{=}[r] &
f_{\tau^m_m}^{-1}K_0 \ar@{=}[r] &
f_{\tau^m_{\varphi(n)m}}^{-1}f_{\delta^1_0}^{-1}K_0
}
$$
Nous laissons au lecteur le soin de vérifier que la condition de cocycle
entraîne la compatibilité des morphismes avec la composition (dans leurs catégories dérivées
respectives) et définit donc un système simplicial
dans la catégorie dérivée.
\end{proof}

\begin{lemma}
\label{spaces-simplicial-lemma-abelian-postnikov}
Dans la situation \ref{spaces-simplicial-situation-simplicial-site}. Soit $K$
un objet de $D(\mathcal{C}_{total})$. Posons
$$
X_n = (g_{n!}\mathbf{Z})
\otimes^\mathbf{L}_\mathbf{Z} K
\quad\text{et}\quad
Y_n =
(g_{n!}\mathbf{Z} \to \ldots \to g_{0!}\mathbf{Z})[-n]
\otimes^\mathbf{L}_\mathbf{Z} K
$$
comme objets de $D(\mathcal{C}_{total})$, les morphismes étant
comme dans le lemme \ref{spaces-simplicial-lemma-simplicial-resolution-Z-site}.
Avec les morphismes canoniques évidents $Y_n \to X_n$ et
$Y_0 \to Y_1[1] \to Y_2[2] \to \ldots$, on a
\begin{enumerate}
\item les triangles distingués $Y_n \to X_n \to Y_{n - 1} \to Y_n[1]$
définissent un système de Postnikov
(Catégories dérivées, définition \ref{derived-definition-postnikov-system})
pour $\ldots \to X_2 \to X_1 \to X_0$,
\item $K = \text{hocolim} Y_n[n]$ dans $D(\mathcal{C}_{total})$.
\end{enumerate}
\end{lemma}

\begin{proof}
D'abord, si $K = \mathbf{Z}$, c'est la construction de
Catégories dérivées, exemple \ref{derived-example-key-postnikov},
appliquée au complexe
$$
\ldots \to
g_{2!}\mathbf{Z} \to
g_{1!}\mathbf{Z} \to
g_{0!}\mathbf{Z}
$$
dans $\textit{Ab}(\mathcal{C}_{total})$, compte tenu du fait que
ce complexe représente $K = \mathbf{Z}$ dans $D(\mathcal{C}_{total})$
par le lemme \ref{spaces-simplicial-lemma-simplicial-resolution-Z-site}.
Le cas général en découle, puisque le foncteur exact
$- \otimes^\mathbf{L}_\mathbf{Z} K$ transforme les systèmes de Postnikov en
systèmes de Postnikov et
que $- \otimes^\mathbf{L}_\mathbf{Z} K$ commute aux colimites homotopiques.
\end{proof}

\begin{lemma}
\label{spaces-simplicial-lemma-nullity-cartesian-objects-derived}
Dans la situation \ref{spaces-simplicial-situation-simplicial-site}.
Soient $K, K' \in D(\mathcal{C}_{total})$.
Supposons que
\begin{enumerate}
\item $K$ est cartésien,
\item $\Hom(K_i[i], K'_i) = 0$ pour $i > 0$ et
\item $\Hom(K_i[i + 1], K'_i) = 0$ pour $i \geq 0$.
\end{enumerate}
Alors tout morphisme $K \to K'$ induisant le morphisme nul $K_0 \to K'_0$ est nul.
\end{lemma}

\begin{proof}
Considérons les objets $X_n$ et le système de Postnikov $Y_n$
associé à $K$ dans le lemme \ref{spaces-simplicial-lemma-abelian-postnikov}.
Comme $K = \text{hocolim} Y_n[n]$, le morphisme $K \to K'$ fournit
une famille compatible de morphismes $Y_n[n] \to K'$.
D'après (1) et le lemme \ref{spaces-simplicial-lemma-derived-cartesian-shriek}, on a
$X_n = g_{n!}K_n$. Puisque $Y_0 = X_0$, la nullité de
$K_0 \to K'_0$ entraîne celle de $Y_0 \to K'$.
Supposons par récurrence que le morphisme $Y_n[n] \to K'$ soit nul
pour un certain $n \geq 0$. Le triangle distingué
$$
Y_n[n] \to Y_{n + 1}[n + 1] \to X_{n + 1}[n + 1] \to Y_n[n + 1]
$$
on déduit une suite exacte
$$
\Hom(X_{n + 1}[n + 1], K') \to
\Hom(Y_{n + 1}[n + 1], K') \to
\Hom(Y_n[n], K')
$$
Comme $X_{n + 1}[n + 1] = g_{n + 1!}K_{n + 1}[n + 1]$, le premier groupe est égal à
$$
\Hom(K_{n + 1}[n + 1], K'_{n + 1})
$$
qui est nul par l'hypothèse (2). Par récurrence, tous les morphismes
$Y_n[n] \to K'$ sont nuls. Considérons le triangle distingué qui définit
$$
\bigoplus Y_n[n] \to
\bigoplus Y_n[n] \to
K \to
(\bigoplus Y_n[n])[1]
$$
la colimite homotopique. En raisonnant comme ci-dessus, on voit qu'il suffit
de montrer que
$$
\Hom((\bigoplus Y_n[n])[1], K') = \prod \Hom(Y_n[n + 1], K')
$$
soit nul pour tout $n \geq 0$. Pour le vérifier, en raisonnant comme ci-dessus,
il suffit de montrer que
$$
\Hom(K_n[n + 1], K'_n)  = 0
$$
pour tout $n \geq 0$, ce qui découle de la condition (3).
\end{proof}

\begin{lemma}
\label{spaces-simplicial-lemma-hom-cartesian-objects-derived}
Dans la situation \ref{spaces-simplicial-situation-simplicial-site},
soient $K, K' \in D(\mathcal{C}_{total})$.
Supposons que
\begin{enumerate}
\item $K$ est cartésien,
\item $\Hom(K_i[i - 1], K'_i) = 0$ pour $i > 1$.
\end{enumerate}
Alors tout morphisme $\{K_n \to K'_n\}$ entre les systèmes simpliciaux associés
à $K$ et $K'$ provient d'un morphisme $K \to K'$ dans $D(\mathcal{C}_{total})$.
\end{lemma}

\begin{proof}
Soit $\{K_n \to K'_n\}_{n \geq 0}$
un morphisme de systèmes simpliciaux de la catégorie dérivée.
Considérons les objets $X_n$ et le système de Postnikov $Y_n$
associés au $K$ du lemme \ref{spaces-simplicial-lemma-abelian-postnikov}.
D'après (1) et le lemme \ref{spaces-simplicial-lemma-derived-cartesian-shriek}, on a
$X_n = g_{n!}K_n$. En particulier, le morphisme $K_0 \to K'_0$
induit un morphisme $X_0 \to K'$. Puisque $\{K_n \to K'_n\}$
est un morphisme de systèmes, un calcul (omis) montre que
la composition
$$
X_1 \to X_0 \to K'
$$
est nulle. Comme $Y_0 = X_0$ et que $Y_1$ s'insère dans un
triangle distingué
$$
Y_1 \to X_1 \to Y_0 \to Y_1[1]
$$
on en déduit l'existence d'un morphisme $Y_1[1] \to K'$ dont la composition
avec $X_0 = Y_0 \to Y_1[1]$ est le morphisme $X_0 \to K'$
ci-dessus. Supposons donné un morphisme $Y_n[n] \to K'$ pour $n \geq 1$.
À partir du triangle distingué
$$
X_{n + 1}[n] \to Y_n[n] \to Y_{n + 1}[n + 1] \to X_{n + 1}[n + 1]
$$
on obtient une suite exacte
$$
\Hom(Y_{n + 1}[n + 1], K') \to
\Hom(Y_n[n], K') \to
\Hom(X_{n + 1}[n], K')
$$
Comme $X_{n + 1}[n] = g_{n + 1!}K_{n + 1}[n]$, le dernier groupe est égal à
$$
\Hom(K_{n + 1}[n], K'_{n + 1})
$$
qui est nul par l'hypothèse (2). Par récurrence, on construit un système de morphismes
$Y_n[n] \to K'$ compatible avec les morphismes de transition et se réduisant
au morphisme donné sur $Y_0$. Cela fournit un morphisme
$$
\gamma :
K = \text{hocolim} Y_n[n]
\longrightarrow
K'
$$
Ce morphisme rend en particulier commutatif le diagramme
$$
\xymatrix{
X_0 \ar[rd] \ar[r] &
K \ar[d]^\gamma \\
& K'
}
$$
est commutatif. Par restriction à
$\mathcal{C}_0$, on en déduit que le morphisme $\gamma_0 : K_0 \to K'_0$
est le même que le premier morphisme $K_0 \to K'_0$ du morphisme
de systèmes simpliciaux. Comme $K$ est cartésien, on en déduit aisément que
$\{\gamma_n\}$ est le morphisme de systèmes simpliciaux dont on était parti.
\end{proof}

\begin{lemma}
\label{spaces-simplicial-lemma-cartesian-object-derived-from-simplicial}
Dans la situation \ref{spaces-simplicial-situation-simplicial-site}. Soit
$(K_n, K_\varphi)$ un système simplicial de la catégorie dérivée.
Supposons que
\begin{enumerate}
\item $(K_n, K_\varphi)$ est cartésien,
\item $\Hom(K_i[t], K_i) = 0$ pour $i \geq 0$ et $t > 0$.
\end{enumerate}
Il existe alors un objet cartésien $K$ de $D(\mathcal{C}_{total})$
dont le système simplicial associé est isomorphe à $(K_n, K_\varphi)$.
\end{lemma}

\begin{proof}
Posons $X_n = g_{n!}K_n$ dans $D(\mathcal{C}_{total})$. Pour tout $n \geq 1$,
on a
$$
\Hom(X_n, X_{n - 1}) =
\Hom(K_n, g_n^{-1}g_{n - 1!}K_{n - 1}) =
\bigoplus\nolimits_{\varphi : [n - 1] \to [n]}
\Hom(K_n, f_\varphi^{-1}K_{n - 1})
$$
On obtient ainsi un morphisme $X_n \to X_{n - 1}$ correspondant à la somme alternée
des morphismes
$K_\varphi^{-1} : K_n \to f_\varphi^{-1}K_{n - 1}$
où $\varphi$ parcourt $\delta^n_0, \ldots, \delta^n_n$.
Nous pouvons le faire car $K_\varphi$ est inversible par l'hypothèse (1).
Comparer avec la définition des morphismes
dans la preuve du lemme \ref{spaces-simplicial-lemma-simplicial-resolution-Z-site}.
On obtient un complexe
$$
\ldots \to X_2 \to X_1 \to X_0
$$
dans $D(\mathcal{C}_{total})$. Nous omettons le calcul qui montre
que les composés sont nuls. D'après
Catégories dérivées, lemme \ref{derived-lemma-existence-postnikov-system},
si l'on a
$$
\Hom(X_i[i - j - 2], X_j) = 0\text{ pour }i > j + 2
$$
on peut prolonger ce complexe en un système de Postnikov.
Le groupe est égal à
$$
\Hom(K_i[i - j - 2], g_i^{-1}g_{j!}K_j)
$$
En utilisant à nouveau la cartésianité de $(K_n, K_\varphi)$, on voit que
$g_i^{-1}g_{j!}K_j$ est isomorphe à une somme directe finie de copies de
$K_i$. Ce groupe s'annule donc par l'hypothèse (2).
Prenons le système de Postnikov défini par $Y_0 = X_0$ et les triangles distingués
$Y_n \to X_n \to Y_{n - 1} \to Y_n[1]$ pour $n \geq 1$.
Nous définissons
$$
K = \text{hocolim} Y_n[n]
$$
Pour terminer la preuve, nous devons montrer que $g_m^{-1}K$ est isomorphe
à $K_m$ pour tout $m$, de façon compatible avec les morphismes $K_\varphi$. Notons que
$$
g_m^{-1} K = \text{hocolim} g_m^{-1}Y_n[n]
$$
et que $g_m^{-1}Y_n[n]$ est un système Postnikov pour $g_m^{-1}X_n$.
Considérons les isomorphismes
$$
g_m^{-1}X_n =
\bigoplus\nolimits_{\varphi : [n] \to [m]} f_\varphi^{-1}K_n
\xrightarrow{\bigoplus K_\varphi}
\bigoplus\nolimits_{\varphi : [n] \to [m]} K_m
$$
Ces morphismes définissent un isomorphisme de complexes
$$
\xymatrix{
\ldots \ar[r] &
g_m^{-1}X_2 \ar[r] \ar[d] &
g_m^{-1}X_1 \ar[r] \ar[d] &
g_m^{-1}X_0 \ar[d] \\
\ldots \ar[r] &
\bigoplus\nolimits_{\varphi : [2] \to [m]} K_m \ar[r] &
\bigoplus\nolimits_{\varphi : [1] \to [m]} K_m \ar[r] &
\bigoplus\nolimits_{\varphi : [0] \to [m]} K_m
}
$$
dans $D(\mathcal{C}_m)$; les flèches de la rangée inférieure sont celles de
dans la preuve du lemme \ref{spaces-simplicial-lemma-simplicial-resolution-Z-site}.
Les carrés sont commutatifs par le choix des flèches du complexe
$\ldots \to X_2 \to X_1 \to X_0$ ; nous omettons le calcul.
Le complexe formé par la rangée inférieure admet la tour de Postnikov
$$
Y'_{m, n} =
\left(\bigoplus\nolimits_{\varphi : [n] \to [m]} \mathbf{Z} \to
\ldots \to
\bigoplus\nolimits_{\varphi : [0] \to [m]} \mathbf{Z}\right)[-n]
\otimes^\mathbf{L}_\mathbf{Z} K_m
$$
et l'on a $\text{hocolim} Y'_{m, n} = K_m$
(comparer avec la preuve du lemme \ref{spaces-simplicial-lemma-abelian-postnikov}
et Catégories dérivées, exemple \ref{derived-example-key-postnikov}).
La seconde partie de
Catégories dérivées, lemme \ref{derived-lemma-existence-postnikov-system},
permet de prolonger les morphismes verticaux du grand diagramme en un isomorphisme
de systèmes de Postnikov, pourvu que l'on ait
$$
\Hom(g_m^{-1}X_i[i - j - 1], \bigoplus\nolimits_{\varphi : [j] \to [m]} K_m)
= 0\text{ pour }i > j + 1
$$
Cette annulation découle de $\Hom(K_m[i - j - 1], K_m) = 0$ pour $i > j + 1$,
par l'hypothèse (2). Choisissons un isomorphisme
de systèmes de Postnikov $\gamma_{m, n} : g_m^{-1}Y_n \to Y'_{m, n}$
dans $D(\mathcal{C}_m)$. Par unicité des colimites homotopiques,
il existe un isomorphisme
$$
g_m^{-1} K = \text{hocolim} g_m^{-1}Y_n[n]
\xrightarrow{\gamma_m}
\text{hocolim} Y'_{m, n} = K_m
$$
compatible avec $\gamma_{m, n}$.

\medskip\noindent
Il reste à montrer que les morphismes $\gamma_m$ s'inscrivent dans les diagrammes commutatifs
$$
\xymatrix{
f_\varphi^{-1}g_m^{-1}K \ar[d]_{f_\varphi^{-1}\gamma_m} \ar[r]_{K(\varphi)} &
g_n^{-1}K \ar[d]^{\gamma_n} \\
f_\varphi^{-1}K_m \ar[r]^{K_\varphi} &
K_n
}
$$
pour tout $\varphi : [m] \to [n]$. Considérons le diagramme
$$
\xymatrix{
f_\varphi^{-1}(\bigoplus_{\psi : [0] \to [m]} f_\psi^{-1}K_0)
\ar@{=}[r] \ar[d]_{f_\varphi^{-1}(\bigoplus K_\psi)} &
f_\varphi^{-1}g_m^{-1}X_0 \ar[d] \ar[r]_{X_0(\varphi)} &
g_n^{-1}X_0 \ar[d] &
\bigoplus_{\chi : [0] \to [n]} f_\chi^{-1}K_0
\ar@{=}[l] \ar[d]^{\bigoplus K_\chi} \\
f_\varphi^{-1}(\bigoplus_{\psi : [0] \to [m]} K_m) \ar@{=}[d] &
f_\varphi^{-1}g_m^{-1}K \ar[d]_{f_\varphi^{-1}\gamma_m} \ar[r]_{K(\varphi)} &
g_n^{-1}K \ar[d]^{\gamma_n} &
\bigoplus_{\chi : [0] \to [n]} K_n \ar@{=}[d] \\
f_\varphi^{-1}Y'_{0, m} \ar[r] &
f_\varphi^{-1}K_m \ar[r]^{K_\varphi} &
K_n &
Y'_{0, n} \ar[l]
}
$$
Le carré du milieu supérieur est commutatif car $X_0 \to K$ est un morphisme
d'objets simpliciaux. Les rectangles de gauche et de droite sont
commutatifs, car $\gamma_m$ et $\gamma_n$ sont respectivement compatibles avec
$\gamma_{0, m}$ et $\gamma_{0, n}$, qui sont les flèches
$\bigoplus K_\psi$ et $\bigoplus K_\chi$ dans le diagramme.
Le rectangle extérieur du diagramme
est commutatif, car $(K_n, K_\varphi)$ est un système simplicial
et le morphisme $X_0(\varphi)$ est donné par les identifications évidentes
$f_\varphi^{-1}f_\psi^{-1}K_0 = f_{\varphi \circ \psi}^{-1}K_0$.
Notons que la flèche $\bigoplus_\psi K_m \to Y'_{0, m} \to K_m$
induit un isomorphisme sur chacun des facteurs directs
(en raison de notre construction explicite des systèmes de Postnikov
$Y'_{i, j}$ ci-dessus).
Ainsi, si l'on prend un facteur direct du
coin supérieur gauche, celui-ci s'envoie isomorphiquement sur
$f_\varphi^{-1}g_m^{-1}K$ car $\gamma_m$ est un isomorphisme.
En explicitant ce qui précède
et en ne considérant que ce facteur direct, on conclut que le carré
central inférieur est lui aussi commutatif. Cela termine la preuve.
\end{proof}









\section{Systèmes simpliciaux de la catégorie dérivée : modules}
\label{spaces-simplicial-section-glueing-modules}

\noindent
Dans cette section, nous allons prouver un cas particulier de
\cite[Proposition 3.2.9]{BBD} dans le cadre des catégories
dérivées de $\mathcal{O}$-modules. Le cas (un peu) plus simple des faisceaux abéliens
est discuté dans la section \ref{spaces-simplicial-section-glueing}.

\begin{definition}
\label{spaces-simplicial-definition-cartesian-derived-modules}
Dans la situation \ref{spaces-simplicial-situation-simplicial-site}. Soit $\mathcal{O}$
un faisceau d'anneaux sur $\mathcal{C}_{total}$. Un
{\it système simplicial de la catégorie dérivée de modules}
se compose des données suivantes
\begin{enumerate}
\item pour tout $n$, un objet $K_n$ de $D(\mathcal{O}_n)$,
\item pour tout $\varphi : [m] \to [n]$, un morphisme
$K_\varphi : Lf_\varphi^*K_m \to K_n$ dans $D(\mathcal{O}_n)$
\end{enumerate}
soumises à la condition
$$
K_{\varphi \circ \psi} = K_\varphi \circ Lf_\varphi^*K_\psi :
Lf_{\varphi \circ \psi}^*K_l = Lf_\varphi^* Lf_\psi^*K_l
\longrightarrow
K_n
$$
pour tout couple de morphismes $\varphi : [m] \to [n]$ et $\psi : [l] \to [m]$ de $\Delta$.
On dit que le système simplicial est {\it cartésien} si les morphismes $K_\varphi$
sont des isomorphismes pour tout $\varphi$.
Étant donnés deux systèmes simpliciaux de la catégorie dérivée,
il existe une notion évidente de
{\it morphisme de systèmes simpliciaux de la catégorie dérivée de modules}.
\end{definition}

\noindent
Nous avons intentionnellement donné à cette notion un nom ridiculement long.
Le but est de montrer qu'un système simplicial de la catégorie dérivée
de modules provient d'un objet de $D(\mathcal{O})$ sous certaines
hypothèses.

\begin{lemma}
\label{spaces-simplicial-lemma-cartesian-objects-derived-modules}
Dans la situation \ref{spaces-simplicial-situation-simplicial-site}, soit $\mathcal{O}$ un faisceau d'anneaux
sur $\mathcal{C}_{total}$.
Si $K \in D(\mathcal{O})$ est un objet, alors $(K_n, K(\varphi))$
est un système simplicial de la catégorie dérivée de modules.
Si $K$ est cartésien, le système l'est également.
\end{lemma}

\begin{proof}
Cela résulte immédiatement des définitions.
\end{proof}

\begin{lemma}
\label{spaces-simplicial-lemma-construct-cartesian-objects-derived-modules}
Dans la situation \ref{spaces-simplicial-situation-simplicial-site}, soit $\mathcal{O}$ un faisceau d'anneaux
sur $\mathcal{C}_{total}$.
Supposons donnés $K_0 \in D(\mathcal{O}_0)$ et un isomorphisme
$$
\alpha :
L(f_{\delta_1^1})^*K_0
\longrightarrow
L(f_{\delta_0^1})^*K_0
$$
satisfaisant à la condition de cocycle. Posons $\tau^n_i : [0] \to [n]$, $0 \mapsto i$
et posons $K_n = Lf_{\tau^n_n}^*K_0$. Les objets $K_n$ sont les termes d'un
système simplicial cartésien de la catégorie dérivée de modules.
\end{lemma}

\begin{proof}
Comparer avec les lemmes \ref{spaces-simplicial-lemma-construct-cartesian-objects-derived}
et \ref{spaces-simplicial-lemma-characterize-cartesian} et sa preuve
(également pour voir la condition du cocycle énoncée).
La construction est analogue à la construction discutée dans
Descente, section \ref{descent-section-descent-modules}, à laquelle nous empruntons
la notation $\tau^n_i : [0] \to [n]$, $0 \mapsto i$ et
$\tau^n_{ij} : [1] \to [n]$, $0 \mapsto i$, $1 \mapsto j$.
Pour $\varphi : [n] \to [m]$, définissons
$K_\varphi : L(f_\varphi)^*K_n \to K_m$
à l'aide du diagramme
$$
\xymatrix{
L(f_\varphi)^*K_n \ar@{=}[r] &
L(f_\varphi)^* L(f_{\tau^n_n})^*K_0 \ar@{=}[r] &
L(f_{\tau^m_{\varphi(n)}})^*K_0 \ar@{=}[r] &
L(f_{\tau^m_{\varphi(n)m}})^* L(f_{\delta^1_1})^*K_0
\ar[d]_{L(f_{\tau^m_{\varphi(n)m}})^*\alpha} \\
&
K_m \ar@{=}[r] &
L(f_{\tau^m_m})^*K_0 \ar@{=}[r] &
L(f_{\tau^m_{\varphi(n)m}})^* L(f_{\delta^1_0})^*K_0
}
$$
Nous laissons au lecteur le soin de vérifier que la condition de cocycle
entraîne la compatibilité des morphismes avec la composition (dans leurs catégories dérivées
respectives) et définit donc un système simplicial
de la catégorie dérivée de modules.
\end{proof}

\begin{lemma}
\label{spaces-simplicial-lemma-modules-postnikov}
Dans la situation \ref{spaces-simplicial-situation-simplicial-site}, soit $\mathcal{O}$
un faisceau d'anneaux sur $\mathcal{C}_{total}$. Soit $K$
un objet de $D(\mathcal{C}_{total})$. Posons
$$
X_n = (g_{n!}\mathcal{O}_n)
\otimes^\mathbf{L}_\mathcal{O} K
\quad\text{et}\quad
Y_n =
(g_{n!}\mathcal{O}_n \to \ldots \to g_{0!}\mathcal{O}_0)[-n]
\otimes^\mathbf{L}_\mathcal{O} K
$$
comme objets de $D(\mathcal{O})$, les morphismes étant
comme dans le lemme \ref{spaces-simplicial-lemma-simplicial-resolution-Z-site}.
Avec les morphismes canoniques évidents $Y_n \to X_n$ et
$Y_0 \to Y_1[1] \to Y_2[2] \to \ldots$, on a
\begin{enumerate}
\item les triangles distingués $Y_n \to X_n \to Y_{n - 1} \to Y_n[1]$
définissent un système de Postnikov
(Catégories dérivées, définition \ref{derived-definition-postnikov-system})
pour $\ldots \to X_2 \to X_1 \to X_0$,
\item $K = \text{hocolim} Y_n[n]$ dans $D(\mathcal{O})$.
\end{enumerate}
\end{lemma}

\begin{proof}
D'abord, si $K = \mathcal{O}$, c'est la construction de
Catégories dérivées, exemple \ref{derived-example-key-postnikov},
appliquée au complexe
$$
\ldots \to
g_{2!}\mathcal{O}_2 \to
g_{1!}\mathcal{O}_1 \to
g_{0!}\mathcal{O}_0
$$
dans $\textit{Ab}(\mathcal{C}_{total})$ combiné au fait que
ce complexe représente $K = \mathcal{O}$ dans $D(\mathcal{C}_{total})$
par le lemme \ref{spaces-simplicial-lemma-simplicial-resolution-ringed}.
Le cas général en découle, puisque le foncteur exact
$- \otimes^\mathbf{L}_\mathcal{O} K$ transforme les systèmes de Postnikov en
systèmes de Postnikov et
que $- \otimes^\mathbf{L}_\mathcal{O} K$ commute aux colimites homotopiques.
\end{proof}

\begin{lemma}
\label{spaces-simplicial-lemma-nullity-cartesian-modules-derived}
Dans la situation \ref{spaces-simplicial-situation-simplicial-site}, soit $\mathcal{O}$
un faisceau d'anneaux sur $\mathcal{C}_{total}$.
Soient $K, K' \in D(\mathcal{O})$.
Supposons que
\begin{enumerate}
\item $f_\varphi^{-1}\mathcal{O}_n \to \mathcal{O}_m$ est plat pour
$\varphi : [m] \to [n]$,
\item $K$ est cartésien,
\item $\Hom(K_i[i], K'_i) = 0$ pour $i > 0$ et
\item $\Hom(K_i[i + 1], K'_i) = 0$ pour $i \geq 0$.
\end{enumerate}
Alors tout morphisme $K \to K'$ induisant le morphisme nul $K_0 \to K'_0$ est nul.
\end{lemma}

\begin{proof}
La démonstration est identique à celle du
lemme \ref{spaces-simplicial-lemma-nullity-cartesian-objects-derived}, à ceci près qu'elle utilise le
lemme \ref{spaces-simplicial-lemma-modules-postnikov} au lieu du
lemme \ref{spaces-simplicial-lemma-abelian-postnikov}.
\end{proof}

\begin{lemma}
\label{spaces-simplicial-lemma-hom-cartesian-modules-derived}
Dans la situation \ref{spaces-simplicial-situation-simplicial-site}, soit $\mathcal{O}$
un faisceau d'anneaux sur $\mathcal{C}_{total}$.
Soient $K, K' \in D(\mathcal{O})$.
Supposons que
\begin{enumerate}
\item $f_\varphi^{-1}\mathcal{O}_n \to \mathcal{O}_m$ est plat pour
$\varphi : [m] \to [n]$,
\item $K$ est cartésien,
\item $\Hom(K_i[i - 1], K'_i) = 0$ pour $i > 1$.
\end{enumerate}
Alors tout morphisme $\{K_n \to K'_n\}$ entre les systèmes simpliciaux associés
à $K$ et $K'$ provient d'un morphisme $K \to K'$ dans $D(\mathcal{O})$.
\end{lemma}

\begin{proof}
La démonstration est identique à celle du
lemme \ref{spaces-simplicial-lemma-hom-cartesian-objects-derived}, à ceci près qu'elle utilise le
lemme \ref{spaces-simplicial-lemma-modules-postnikov} au lieu du
lemme \ref{spaces-simplicial-lemma-abelian-postnikov}.
\end{proof}

\begin{lemma}
\label{spaces-simplicial-lemma-cartesian-module-derived-from-simplicial}
Dans la situation \ref{spaces-simplicial-situation-simplicial-site}, soit $\mathcal{O}$
un faisceau d'anneaux sur $\mathcal{C}_{total}$. Soit
$(K_n, K_\varphi)$ un système simplicial de la catégorie dérivée
de modules. Supposons que
\begin{enumerate}
\item $f_\varphi^{-1}\mathcal{O}_n \to \mathcal{O}_m$ est plat pour
$\varphi : [m] \to [n]$,
\item $(K_n, K_\varphi)$ est cartésien,
\item $\Hom(K_i[t], K_i) = 0$ pour $i \geq 0$ et $t > 0$.
\end{enumerate}
Il existe alors un objet cartésien $K$ de $D(\mathcal{O})$
dont le système simplicial associé est isomorphe à $(K_n, K_\varphi)$.
\end{lemma}

\begin{proof}
La démonstration est identique à celle du
lemme \ref{spaces-simplicial-lemma-cartesian-object-derived-from-simplicial},
avec les modifications suivantes
\begin{enumerate}
\item utiliser $g_n^* = Lg_n^*$ partout au lieu de $g_n^{-1}$,
\item utiliser $f_\varphi^* = Lf_\varphi^*$ partout au lieu de $f_\varphi^{-1}$,
\item se référer au lemme \ref{spaces-simplicial-lemma-simplicial-resolution-ringed}
au lieu du lemme \ref{spaces-simplicial-lemma-simplicial-resolution-Z-site},
\item dans la construction de $Y'_{m, n}$ utiliser
$\mathcal{O}_m$ au lieu de $\mathbf{Z}$,
\item comparer avec la preuve du lemme \ref{spaces-simplicial-lemma-modules-postnikov}
plutôt que la preuve du lemme \ref{spaces-simplicial-lemma-abelian-postnikov}.
\end{enumerate}
Ceci termine la preuve.
\end{proof}







\section{Le site associé à un objet semi-représentable}
\label{spaces-simplicial-section-semi-representable}

\noindent
Soit $\mathcal{C}$ un site. Rappelons qu'un {\it objet semi-représentable}
de $\mathcal{C}$ est simplement une famille $\{U_i\}_{i \in I}$
d'objets de $\mathcal{C}$. Un
{\it morphisme $\{U_i\}_{i \in I} \to \{V_j\}_{j \in J}$ d'objets
semi-représentables} est donné par une application $\alpha : I \to J$
et pour tout $i \in I$ un morphisme $f_i : U_i \to V_{\alpha(i)}$
de $\mathcal{C}$.
La catégorie d'objets semi-représentables de $\mathcal{C}$
est notée $\text{SR}(\mathcal{C})$.
Voir Hyperrecouvrements, définition \ref{hypercovering-definition-SR},
et la section qui la contient pour davantage de précisions.

\medskip\noindent
Pour un objet semi-représentable $K = \{U_i\}_{i \in I}$ de $\mathcal{C}$,
on pose
$$
\mathcal{C}/K = \coprod\nolimits_{i \in I} \mathcal{C}/U_i
$$
l'union disjointe des localisations de $\mathcal{C}$ en $U_i$.
Cette catégorie porte une structure naturelle de site, dont les recouvrements sont
hérités de ceux des localisations $\mathcal{C}/U_i$.
Le site $\mathcal{C}/K$ s'appelle
{\it localisation de $\mathcal{C}$ en $K$}.
Observons qu'un faisceau sur $\mathcal{C}/K$ est la même chose qu'
une famille de faisceaux $\mathcal{F}_i$ sur $\mathcal{C}/U_i$, c'est-à-dire
$$
\Sh(\mathcal{C}/K) = \prod\nolimits_{i \in I} \Sh(\mathcal{C}/U_i)
$$
Cette observation est parfois utile pour comprendre la situation.

\medskip\noindent
Soient $\mathcal{C}$ un site et $K = \{U_i\}_{i \in I}$ un objet de
$\text{SR}(\mathcal{C})$. Il existe un foncteur de localisation
continu et cocontinu $j : \mathcal{C}/K \to \mathcal{C}$ qui est
le produit des foncteurs de localisation
$j_i : \mathcal{C}/V_i \to \mathcal{C}$.
On obtient les foncteurs $j_!$, $j^{-1}$, $j_*$ exactement
comme dans Sites, section \ref{sites-section-localize}.
En termes de la décomposition en produit
$\Sh(\mathcal{C}/K) = \prod\nolimits_{i \in I} \Sh(\mathcal{C}/U_i)$
nous avons
$$
\begin{matrix}
j_! & : &
(\mathcal{F}_i)_{i \in I} &
\longmapsto &
\coprod j_{i, !}\mathcal{F}_i \\
j^{-1} & : &
\mathcal{G} &
\longmapsto &
(j_i^{-1}\mathcal{G})_{i \in I} \\
j_* & : &
(\mathcal{F}_i)_{i \in I} &
\longmapsto &
\prod j_{i, *}\mathcal{F}_i
\end{matrix}
$$
ce que le lecteur vérifie aisément.

\medskip\noindent
Soit $f : K \to L$ un morphisme de $\text{SR}(\mathcal{C})$.
On obtient alors un foncteur continu et cocontinu
$$
v : \mathcal{C}/K \longrightarrow \mathcal{C}/L
$$
en appliquant aux composantes la construction de Sites, lemme \ref{sites-lemma-relocalize}.
aux composantes. Plus précisément, supposons $f = (\alpha, f_i)$
où $K = \{U_i\}_{i \in I}$, $L = \{V_j\}_{j \in J}$, $\alpha : I \to J$,
et $f_i : U_i \to V_{\alpha(i)}$. Le foncteur $v$ envoie la composante
$\mathcal{C}/U_i$ dans la composante $\mathcal{C}/V_{\alpha(i)}$
par la construction du lemme précité. En particulier,
on obtient un morphisme
$$
f : \Sh(\mathcal{C}/K) \to \Sh(\mathcal{C}/L)
$$
de topoi. En termes des décompositions en produits
$\Sh(\mathcal{C}/K) = \prod\nolimits_{i \in I} \Sh(\mathcal{C}/U_i)$ et
$\Sh(\mathcal{C}/L) = \prod\nolimits_{j \in J} \Sh(\mathcal{C}/V_j)$
le lecteur vérifiera que
$$
\begin{matrix}
f_! & : &
(\mathcal{F}_i)_{i \in I} &
\longmapsto &
(\coprod\nolimits_{i \in I, \alpha(i) = j} f_{i, !}\mathcal{F}_i)_{j \in J} \\
f^{-1} & : &
(\mathcal{G}_j)_{j \in J} &
\longmapsto &
(f_i^{-1}\mathcal{G}_{\alpha(i)})_{i \in I} \\
f_* & : &
(\mathcal{F}_i)_{i \in I} &
\longmapsto &
(\prod\nolimits_{i \in I, \alpha(i) = j} f_{i, *}\mathcal{F}_i)_{j \in J}
\end{matrix}
$$
où $f_i : \Sh(\mathcal{C}/U_i) \to \Sh(\mathcal{C}/V_{\alpha(i)})$
est le morphisme associé au foncteur de localisation
$\mathcal{C}/U_i \to \mathcal{C}/V_{\alpha(i)}$ correspondant à
$f_i : U_i \to V_{\alpha(i)}$.

\begin{lemma}
\label{spaces-simplicial-lemma-has-P}
Soit $\mathcal{C}$ un site.
\begin{enumerate}
\item Pour $K$ dans $\text{SR}(\mathcal{C})$, le foncteur
$j : \mathcal{C}/K \to \mathcal{C}$ est continu,
cocontinu et possède la propriété P de
Sites, remarque \ref{sites-remark-cartesian-cocontinuous}.
\item Pour $f : K \to L$ dans $\text{SR}(\mathcal{C})$,
le foncteur $v : \mathcal{C}/K \to \mathcal{C}/L$ (voir ci-dessus)
est continu, cocontinu et possède la propriété P de
Sites, remarque \ref{sites-remark-cartesian-cocontinuous}.
\end{enumerate}
\end{lemma}

\begin{proof}
Pour (2). Avec les notations de la discussion précédant le lemme,
les foncteurs de localisation $\mathcal{C}/U_i \to \mathcal{C}/V_{\alpha(i)}$
sont continus et cocontinus par
Sites, section \ref{sites-section-localize}
et possèdent la propriété $P$ d'après
Sites, remarque \ref{sites-remark-localization-cartesian-cocontinuous}.
On en déduit formellement que $v$ est continu, cocontinu et possède la propriété $P$.
Nous laissons au lecteur les détails, ainsi que la preuve de (1).
\end{proof}

\begin{lemma}
\label{spaces-simplicial-lemma-push-pull-localization}
Soit $\mathcal{C}$ un site et soit $K$ dans $\text{SR}(\mathcal{C})$.
Pour $\mathcal{F}$ dans $\Sh(\mathcal{C})$, on a
$$
j_*j^{-1}\mathcal{F} = \SheafHom(F(K)^\#, \mathcal{F})
$$
où $F$ est comme dans
Hyperrecouvrements, définition \ref{hypercovering-definition-SR-F}.
\end{lemma}

\begin{proof}
Écrivons $K = \{U_i\}_{i \in I}$.
En utilisant la description des foncteurs $j^{-1}$ et $j_*$
donnée ci-dessus, nous voyons que
$$
j_*j^{-1}\mathcal{F} =
\prod\nolimits_{i \in I} j_{i, *}(\mathcal{F}|_{\mathcal{C}/U_i}) =
\prod\nolimits_{i \in I} \SheafHom(h_{U_i}^\#, \mathcal{F})
$$
La seconde égalité résulte de Sites, lemme \ref{sites-lemma-hom-sheaf-hU}.
Puisque $F(K) = \coprod h_{U_i}$ dans $\textit{PSh}(\mathcal{C}$,
nous avons $F(K)^\# = \coprod h_{U_i}^\#$ dans $\Sh(\mathcal{C})$
et puisque $\SheafHom(-, \mathcal{F})$ transforme les coproduits en
produits (cela résulte immédiatement de la construction de
Sites, section \ref{sites-section-glueing-sheaves}), on conclut.
\end{proof}

\begin{lemma}
\label{spaces-simplicial-lemma-localize-compare}
Soit $\mathcal{C}$ un site.
\begin{enumerate}
\item Pour $K$ dans $\text{SR}(\mathcal{C})$, le foncteur $j_!$
donne une équivalence $\Sh(\mathcal{C}/K) \to \Sh(\mathcal{C})/F(K)^\#$
où $F$ est comme dans
Hyperrecouvrements, définition \ref{hypercovering-definition-SR-F}.
\item Le foncteur $j^{-1} : \Sh(\mathcal{C}) \to \Sh(\mathcal{C}/K)$
correspond, via l'identification de (1), à
$\mathcal{F} \mapsto (\mathcal{F} \times F(K)^\# \to F(K)^\#)$.
\item Pour $f : K \to L$ dans $\text{SR}(\mathcal{C})$, le foncteur
$f^{-1}$ correspond, via les identifications de (1), au foncteur
$\Sh(\mathcal{C})/F(L)^\# \to \Sh(\mathcal{C})/F(K)^\#$,
$(\mathcal{G} \to F(L)^\#) \mapsto
(\mathcal{G} \times_{F(L)^\#} F(K)^\# \to F(K)^\#)$.
\end{enumerate}
\end{lemma}

\begin{proof}
Observons que, si $K = \{U_i\}_{i \in I}$, la catégorie
$\Sh(\mathcal{C}/K)$ se décompose en produit des catégories
$\Sh(\mathcal{C}/U_i)$. Observons que
$F(K)^\# = \coprod_{i \in I} h_{U_i}^\#$ (coproduit en faisceaux).
$\Sh(\mathcal{C})/F(K)^\#$ est donc le produit des catégories
$\Sh(\mathcal{C})/h_{U_i}^\#$.
Ainsi, (1) et (2) résultent des assertions
correspondantes pour chaque $i$; voir
Sites, lemmes \ref{sites-lemma-essential-image-j-shriek} et
\ref{sites-lemma-compute-j-shriek-restrict}.
De même, si $L = \{V_j\}_{j \in J}$ et si $f$ est donné
par $\alpha : I \to J$ et $f_i : U_i \to V_{\alpha(i)}$,
alors on peut appliquer
Sites, lemme \ref{sites-lemma-relocalize-explicit}
à chacun des morphismes de relocalisation
$\mathcal{C}/U_i \to \mathcal{C}/V_{\alpha(i)}$
pour obtenir (3).
\end{proof}

\begin{lemma}
\label{spaces-simplicial-lemma-localize-injective}
Soit $\mathcal{C}$ un site et soit $K$ dans $\text{SR}(\mathcal{C})$.
Le foncteur $j^{-1}$ envoie les faisceaux abéliens injectifs sur des faisceaux
abéliens injectifs. De même, le foncteur $j^{-1}$ envoie les complexes K-injectifs
de faisceaux abéliens sur des complexes K-injectifs de faisceaux
abéliens.
\end{lemma}

\begin{proof}
La première assertion est la généralisation naturelle du lemme
de Cohomologie sur les sites
\ref{sites-cohomology-lemma-cohomology-of-open}
aux objets semi-représentables.
Elle découle en fait de ce lemme
par la décomposition en produit de $\Sh(\mathcal{C}/K)$
et la description du foncteur $j^{-1}$ donnée ci-dessus.
La seconde assertion est la généralisation naturelle du lemme
de Cohomologie sur les sites
\ref{sites-cohomology-lemma-restrict-K-injective-to-open}
et s'en déduit par la décomposition du topos en produit.

\medskip\noindent
Autre démonstration. Puisque $j$ induit une localisation de topos d'après la
partie (1) du lemme \ref{spaces-simplicial-lemma-localize-compare},
on le déduit aussi immédiatement de
Cohomologie sur les sites, lemmes \ref{sites-cohomology-lemma-cohomology-of-open}
et \ref{sites-cohomology-lemma-restrict-K-injective-to-open}
après agrandissement du site; comparer avec la preuve de
Cohomologie sur les sites, lemme
\ref{sites-cohomology-lemma-cohomology-on-sheaf-sets}
dans le cas des faisceaux injectifs.
\end{proof}

\begin{remark}[Variante au-dessus d'un objet]
\label{spaces-simplicial-remark-semi-representable-over-object}
Soit $\mathcal{C}$ un site et soit $X \in \Ob(\mathcal{C})$.
La catégorie $\text{SR}(\mathcal{C}, X)$
des {\it objets semi-représentables au-dessus de $X$}
est définie par la formule
$\text{SR}(\mathcal{C}, X) = \text{SR}(\mathcal{C}/X)$.
Voir Hyperrecouvrements, définition \ref{hypercovering-definition-SR}.
Nous pouvons donc appliquer la discussion ci-dessus au site
$\mathcal{C}/X$. En résumé, les constructions ci-dessus donnent
\begin{enumerate}
\item un site $\mathcal{C}/K$ pour $K$ dans $\text{SR}(\mathcal{C}, X)$,
\item une décomposition
$\Sh(\mathcal{C}/K) = \prod \Sh(\mathcal{C}/U_i)$ si $K = \{U_i/X\}$,
\item un foncteur de localisation $j : \mathcal{C}/K \to \mathcal{C}/X$,
\item un morphisme $f : \Sh(\mathcal{C}/K) \to \Sh(\mathcal{C}/L)$
pour $f : K \to L$ dans $\text{SR}(\mathcal{C}, X)$.
\end{enumerate}
Tous les résultats de cette section sont valables dans cette situation en remplaçant partout
$\mathcal{C}$ par $\mathcal{C}/X$.
\end{remark}

\begin{remark}[Variante annelée]
\label{spaces-simplicial-remark-semi-representable-ringed}
Soit $\mathcal{C}$ un site et soit $\mathcal{O}_\mathcal{C}$
un faisceau d'anneaux sur $\mathcal{C}$. Alors, pour tout
objet semi-représentable $K$ de $\mathcal{C}$, le site
$\mathcal{C}/K$ est annelé et muni du faisceau d'anneaux
$\mathcal{O}_K = j^{-1}\mathcal{O}_\mathcal{C}$.
Les constructions ci-dessus donnent
\begin{enumerate}
\item un site annelé $(\mathcal{C}/K, \mathcal{O}_K)$
pour $K$ dans $\text{SR}(\mathcal{C})$,
\item une décomposition
$\textit{Mod}(\mathcal{O}_K) =
\prod \textit{Mod}(\mathcal{O}_{U_i})$ si $K = \{U_i\}$,
\item un morphisme de localisation
$j : (\Sh(\mathcal{C}/K), \mathcal{O}_K) \to
(\Sh(\mathcal{C}), \mathcal{O}_\mathcal{C})$
de topoi annelés,
\item un morphisme
$f : (\Sh(\mathcal{C}/K), \mathcal{O}_K) \to
(\Sh(\mathcal{C}/L), \mathcal{O}_L)$ de topoi annelés
pour $f : K \to L$ dans $\text{SR}(\mathcal{C})$.
\end{enumerate}
La plupart des résultats ci-dessus sont valables dans ce cadre. Par exemple, le foncteur
$j^*$ a un adjoint exact à gauche
$$
j_! : \textit{Mod}(\mathcal{O}_K) \to \textit{Mod}(\mathcal{O}_\mathcal{C}),
$$
qui, en termes de la décomposition en produit donnée en (2), envoie
$(\mathcal{F}_i)_{i \in I}$ à $\bigoplus j_{i, !}\mathcal{F}_i$.
De même, étant donné $f : K \to L$ comme ci-dessus, le foncteur $f^*$ a
un adjoint exact à gauche
$f_! : \textit{Mod}(\mathcal{O}_K) \to \textit{Mod}(\mathcal{O}_L)$.
Ainsi, les foncteurs $j^*$ et $f^*$ sont exacts, c'est-à-dire que
$j$ et $f$ sont des morphismes plats de topos annelés (ce qui découle aussi
des égalités $\mathcal{O}_K = j^{-1}\mathcal{O}_\mathcal{C}$
et $\mathcal{O}_K = f^{-1}\mathcal{O}_L$).
\end{remark}

\begin{remark}[Variante annelée au-dessus d'un objet]
\label{spaces-simplicial-remark-semi-representable-ringed-over-object}
Soit $\mathcal{C}$ un site. Soit $\mathcal{O}_\mathcal{C}$
un faisceau d'anneaux sur $\mathcal{C}$. Soit $X \in \Ob(\mathcal{C})$
et notons $\mathcal{O}_X = \mathcal{O}_\mathcal{C}|_{\mathcal{C}/U}$.
On peut alors combiner les constructions données dans les
remarques \ref{spaces-simplicial-remark-semi-representable-over-object}
et \ref{spaces-simplicial-remark-semi-representable-ringed} pour obtenir
\begin{enumerate}
\item un site annelé $(\mathcal{C}/K, \mathcal{O}_K)$
pour $K$ dans $\text{SR}(\mathcal{C}, X)$,
\item une décomposition
$\textit{Mod}(\mathcal{O}_K) =
\prod \textit{Mod}(\mathcal{O}_{U_i})$ si $K = \{U_i\}$,
\item un morphisme de localisation
$j : (\Sh(\mathcal{C}/K), \mathcal{O}_K) \to
(\Sh(\mathcal{C}/X), \mathcal{O}_X)$
de topoi annelés,
\item un morphisme
$f : (\Sh(\mathcal{C}/K), \mathcal{O}_K) \to
(\Sh(\mathcal{C}/L), \mathcal{O}_L)$ de topoi annelés
pour $f : K \to L$ dans $\text{SR}(\mathcal{C}, X)$.
\end{enumerate}
Bien entendu, tous les résultats mentionnés dans
la remarque \ref{spaces-simplicial-remark-semi-representable-ringed}
sont également valables dans ce cadre.
\end{remark}





\section{Le site associé à un objet semi-représentable simplicial}
\label{spaces-simplicial-section-simplicial-semi-representable}

\noindent
Soient $\mathcal{C}$ un site et $K$ un objet simplicial de
$\text{SR}(\mathcal{C})$. Comme d'habitude, posons $K_n = K([n])$ et notons
$K(\varphi) : K_n \to K_m$ le morphisme associé à $\varphi : [m] \to [n]$.
D'après la construction de la
section \ref{spaces-simplicial-section-semi-representable}, on obtient un objet simplicial
$n \mapsto \mathcal{C}/K_n$ dans la catégorie dont les objets sont des sites et
dont les morphismes sont des foncteurs cocontinus. Autrement dit, on obtient
un diagramme comme dans le cas (B) de la section \ref{spaces-simplicial-section-simplicial-sites}.
Les foncteurs satisfont la propriété P par le lemme \ref{spaces-simplicial-lemma-has-P}.
Nous pouvons donc appliquer le lemme
\ref{spaces-simplicial-lemma-simplicial-cocontinuous-site}
pour obtenir un site $(\mathcal{C}/K)_{total}$.

\medskip\noindent
Nous pouvons décrire explicitement le site $(\mathcal{C}/K)_{total}$ comme suit.
Écrivons $K_n = \{U_{n, i}\}_{i \in I_n}$. Pour $\varphi : [m] \to [n]$
le morphisme $K(\varphi) : K_n \to K_m$ est donné par une application
$\alpha(\varphi) : I_n \to I_m$ et les morphismes
$f_{\varphi, i} : U_{n, i} \to U_{m, \alpha(\varphi)(i)}$ pour $i \in I_n$.
On a alors
\begin{enumerate}
\item un objet de $(\mathcal{C}/K)_{total}$
correspond à un objet $(U/U_{n, i})$
de $\mathcal{C}/U_{n, i}$ pour un certain $n$ et un certain $i \in I_n$;
\item un morphisme entre $U/U_{n, i}$ et $V/U_{m, j}$
est une paire $(\varphi, f)$ où $\varphi : [m] \to [n]$,
$j = \alpha(\varphi)(i)$ et $f : U \to V$ est un morphisme de
$\mathcal{C}$ tel que
$$
\vcenter{
\xymatrix{
U \ar[r]_f \ar[d] & V \ar[d] \\
U_{n, i} \ar[r]^-{f_{\varphi, i}} &
U_{m, j}
}
}
$$
est commutatif, et
\item les recouvrements de l'objet $U/U_{n, i}$ sont construits
à partir d'un recouvrement $\{f_j : U_j \to U\}$ dans $\mathcal{C}$, en déclarant
$\{(\text{id}, f_j) : U_j/U_{n, i} \to U/U_{n, i}\}$
recouvrement dans $(\mathcal{C}/K)_{total}$.
\end{enumerate}
Toute la théorie générale développée pour les sites simpliciaux s'applique à
$(\mathcal{C}/K)_{total}$. Observons que le foncteur oublieux évident
$$
j_{total} : (\mathcal{C}/K)_{total} \longrightarrow \mathcal{C}
$$
est continu et cocontinu. Le morphisme de topoi
associé provient d'une augmentation évidente.

\begin{lemma}
\label{spaces-simplicial-lemma-augmentation-simplicial-semi-representable}
Soient $\mathcal{C}$ un site et $K$ un objet simplicial de
$\text{SR}(\mathcal{C})$. Le foncteur de localisation
$j_0 : \mathcal{C}/K_0 \to \mathcal{C}$ définit une augmentation
$a_0 : \Sh(\mathcal{C}/K_0) \to \Sh(\mathcal{C})$, comme dans le cas (B) de
la remarque \ref{spaces-simplicial-remark-augmentation-site}.
Les morphismes de topoi correspondants
$$
a_n : \Sh(\mathcal{C}/K_n) \longrightarrow \Sh(\mathcal{C}),\quad
a : \Sh((\mathcal{C}/K)_{total}) \longrightarrow \Sh(\mathcal{C})
$$
du lemme \ref{spaces-simplicial-lemma-augmentation-site}
sont égaux aux morphismes de topos associés aux
foncteurs de localisation continus et cocontinus
$j_n : \mathcal{C}/K_n \to \mathcal{C}$ et
$j_{total} : (\mathcal{C}/K)_{total} \to \mathcal{C}$.
\end{lemma}

\begin{proof}
Cela résulte immédiatement des définitions.
Voir notamment la note de bas de page dans la preuve du lemme
\ref{spaces-simplicial-lemma-augmentation-site}
pour la relation entre $a$ et $j_{total}$.
\end{proof}

\begin{lemma}
\label{spaces-simplicial-lemma-comparison}
Avec les hypothèses et notations du
lemme \ref{spaces-simplicial-lemma-augmentation-simplicial-semi-representable},
nous avons les propriétés suivantes :
\begin{enumerate}
\item il existe un foncteur
$a^{Sh}_! : \Sh((\mathcal{C}/K)_{total}) \to \Sh(\mathcal{C})$
adjoint à gauche de $a^{-1} : \Sh(\mathcal{C}) \to \Sh((\mathcal{C}/K)_{total})$;
\item il existe un foncteur
$a_! : \textit{Ab}((\mathcal{C}/K)_{total}) \to \textit{Ab}(\mathcal{C})$
adjoint à gauche de
$a^{-1} : \textit{Ab}(\mathcal{C}) \to \textit{Ab}((\mathcal{C}/K)_{total})$,
\item le foncteur $a^{-1}$ associe à un faisceau
$\mathcal{F}$ dans $\Sh(\mathcal{C})$ le faisceau sur $(\mathcal{C}/K)_{total}$
qui en degré $n$ est égal à $a_n^{-1}\mathcal{F}$,
\item le foncteur $a_*$ associe à un faisceau abélien $\mathcal{G}$ sur
$\textit{Ab}((\mathcal{C}/K)_{total})$ l'égaliseur des deux morphismes
$j_{0, *}\mathcal{G}_0 \to j_{1, *}\mathcal{G}_1$,
\end{enumerate}
\end{lemma}

\begin{proof}
Les parties (3) et (4) valent pour toute augmentation d'un
site simplicial; voir le lemme \ref{spaces-simplicial-lemma-augmentation-site}.
Les parties (1) et (2) suivent de ce que $j_{total}$ est continu et cocontinu.
Le foncteur $a^{Sh}_!$ est construit dans
Sites, lemme \ref{sites-lemma-when-shriek},
et le foncteur $a_!$ est construit dans Modules
sur les sites, lemme
\ref{sites-modules-lemma-g-shriek-adjoint}.
\end{proof}

\begin{lemma}
\label{spaces-simplicial-lemma-sanity-check-simplicial-semi-representable}
Soient $\mathcal{C}$ un site et $K$ un objet simplicial de
$\text{SR}(\mathcal{C})$. Soit $U/U_{n, i}$ un objet de
$\mathcal{C}/K_n$. Soit
$\mathcal{F} \in \textit{Ab}((\mathcal{C}/K)_{total})$.
Alors
$$
H^p(U, \mathcal{F}) = H^p(U, \mathcal{F}_{n, i})
$$
où
\begin{enumerate}
\item au membre de gauche, $U$ est considéré comme un objet de
$\mathcal{C}_{total}$ et
\item au membre de droite, $\mathcal{F}_{n, i}$ est la $i$-ième composante
du faisceau $\mathcal{F}_n$ sur $\mathcal{C}/K_n$
dans la décomposition $\Sh(\mathcal{C}/K_n) = \prod \Sh(\mathcal{C}/U_{n, i})$
de la section \ref{spaces-simplicial-section-semi-representable}.
\end{enumerate}
\end{lemma}

\begin{proof}
Cela découle immédiatement du lemme \ref{spaces-simplicial-lemma-sanity-check}
et des décompositions en produits de la section \ref{spaces-simplicial-section-semi-representable}.
\end{proof}

\begin{remark}[Variante au-dessus d'un objet]
\label{spaces-simplicial-remark-augmentation-over-object}
Soient $\mathcal{C}$ un site et $X \in \Ob(\mathcal{C})$.
Rappelons que nous avons une catégorie
$\text{SR}(\mathcal{C}, X) = \text{SR}(\mathcal{C}/X)$
d'objets semi-représentables au-dessus de $X$;
voir la remarque \ref{spaces-simplicial-remark-semi-representable-over-object}.
Nous pouvons appliquer la discussion ci-dessus au site
$\mathcal{C}/X$. En bref, les constructions ci-dessus donnent
\begin{enumerate}
\item un site $(\mathcal{C}/K)_{total}$ pour un objet simplicial $K$
de $\text{SR}(\mathcal{C}, X)$,
\item un foncteur de localisation
$j_{total} : (\mathcal{C}/K)_{total} \to \mathcal{C}/X$,
\item des foncteurs de localisation $j_n : \mathcal{C}/K_n \to \mathcal{C}/X$,
\item un morphisme de topoi
$a : \Sh((\mathcal{C}/K)_{total}) \to \Sh(\mathcal{C}/X)$,
\item des morphismes de topoi
$a_n : \Sh(\mathcal{C}/K_n) \to \Sh(\mathcal{C}/X)$,
\item un foncteur
$a^{Sh}_! : \Sh((\mathcal{C}/K)_{total}) \to \Sh(\mathcal{C}/X)$
adjoint à gauche de $a^{-1}$, et
\item un foncteur
$a_! : \textit{Ab}((\mathcal{C}/K)_{total}) \to \textit{Ab}(\mathcal{C}/X)$
adjoint à gauche de $a^{-1}$.
\end{enumerate}
Tous les résultats de cette section valent dans ce cadre.
Pour le démontrer, on remplace partout
le site $\mathcal{C}$ par $\mathcal{C}/X$.
\end{remark}

\begin{remark}[Variante annelée]
\label{spaces-simplicial-remark-augmentation-ringed}
Soient $\mathcal{C}$ un site et $\mathcal{O}_\mathcal{C}$ un faisceau d'anneaux.
Étant donné un objet semi-représentable simplicial $K$ de $\mathcal{C}$,
nous définissons $\mathcal{O} = a^{-1}\mathcal{O}_\mathcal{C}$, où $a$
est comme dans les lemmes \ref{spaces-simplicial-lemma-augmentation-simplicial-semi-representable} et
\ref{spaces-simplicial-lemma-comparison}.
Les constructions ci-dessus, en tenant compte des faisceaux d'anneaux
comme dans la remarque \ref{spaces-simplicial-remark-semi-representable-ringed}, donnent
\begin{enumerate}
\item un site annelé $((\mathcal{C}/K)_{total}, \mathcal{O})$
pour un objet $K$ simplicial de $\text{SR}(\mathcal{C})$,
\item un morphisme de topoi annelés
$a : (\Sh((\mathcal{C}/K)_{total}), \mathcal{O}) \to
(\Sh(\mathcal{C}), \mathcal{O}_\mathcal{C})$,
\item des morphismes de topoi annelés
$a_n : (\Sh(\mathcal{C}/K_n), \mathcal{O}_n) \to
(\Sh(\mathcal{C}), \mathcal{O}_\mathcal{C})$,
\item un foncteur
$a_! : \textit{Mod}(\mathcal{O}) \to \textit{Mod}(\mathcal{O}_\mathcal{C})$
adjoint à gauche de $a^*$.
\end{enumerate}
Le foncteur $a_!$ existe (mais en général il n'est pas exact)
car $a^{-1}\mathcal{O}_\mathcal{C} = \mathcal{O}$
et l'on peut remplacer l'emploi du
lemme \ref{sites-modules-lemma-g-shriek-adjoint} de Modules sur les sites
dans la preuve du lemme \ref{spaces-simplicial-lemma-comparison}
par Modules sur les sites, lemme \ref{sites-modules-lemma-lower-shriek-modules}.
Comme indiqué dans la remarque \ref{spaces-simplicial-remark-semi-representable-ringed}
il existe des foncteurs exacts
$a_{n!} : \textit{Mod}(\mathcal{O}_n) \to
\textit{Mod}(\mathcal{O}_\mathcal{C})$
adjoints à gauche de $a_n^*$. Par conséquent, les morphismes $a$ et $a_n$ sont plats.
La remarque \ref{spaces-simplicial-remark-semi-representable-ringed}
implique que le morphisme de topoi annelés
$f_\varphi : (\Sh(\mathcal{C}/K_n), \mathcal{O}_n) \to
(\Sh(\mathcal{C}/K_m), \mathcal{O}_m)$
pour $\varphi : [m] \to [n]$ est plat, et il existe un foncteur exact
$f_{\varphi !} : \textit{Mod}(\mathcal{O}_n) \to \textit{Mod}(\mathcal{O}_m)$
adjoint à gauche de $f_\varphi^*$. Il en résulte que, pour
le morphisme plat de topos annelés
$g_n : (\Sh(\mathcal{C}/K_n), \mathcal{O}_n) \to
(\Sh((\mathcal{C}/K)_{total}), \mathcal{O})$
le foncteur $g_{n!} : \textit{Mod}(\mathcal{O}_n) \to
\textit{Mod}(\mathcal{O})$, adjoint à gauche de $g_n^*$, est exact; voir
le lemme \ref{spaces-simplicial-lemma-exactness-g-shriek-modules}.
\end{remark}

\begin{remark}[Variante annelée au-dessus d'un objet]
\label{spaces-simplicial-remark-augmentation-ringed-over-object}
Soient $\mathcal{C}$ un site et $\mathcal{O}_\mathcal{C}$ un faisceau d'anneaux.
Soit $X \in \Ob(\mathcal{C})$ et désignons
$\mathcal{O}_X = \mathcal{O}_\mathcal{C}|_{\mathcal{C}/X}$.
On peut alors combiner les constructions données dans les
remarques \ref{spaces-simplicial-remark-augmentation-over-object} et
\ref{spaces-simplicial-remark-augmentation-ringed}
pour obtenir
\begin{enumerate}
\item un site annelé $((\mathcal{C}/K)_{total}, \mathcal{O})$
pour un objet simplicial $K$ de $\text{SR}(\mathcal{C}, X)$,
\item un morphisme de topoi annelés
$a : (\Sh((\mathcal{C}/K)_{total}), \mathcal{O}) \to
(\Sh(\mathcal{C}/X), \mathcal{O}_X)$,
\item des morphismes de topoi annelés
$a_n : (\Sh(\mathcal{C}/K_n), \mathcal{O}_n) \to
(\Sh(\mathcal{C}/X), \mathcal{O}_X)$,
\item un foncteur
$a_! : \textit{Mod}(\mathcal{O}) \to \textit{Mod}(\mathcal{O}_X)$
adjoint à gauche de $a^*$.
\end{enumerate}
Bien entendu, tous les résultats mentionnés dans la remarque
\ref{spaces-simplicial-remark-augmentation-ringed}
valent aussi dans ce cadre.
\end{remark}






\section{Descente cohomologique pour les hyperrecouvrements}
\label{spaces-simplicial-section-cohomological-descent-hypercoverings}

\noindent
Soit $\mathcal{C}$ un site. Dans cette section, nous supposons que $\mathcal{C}$
admet des égalisateurs et des produits fibrés. Soit $K$ un hyperrecouvrement
au sens d'Hyperrecouvrements, définition
\ref{hypercovering-definition-hypercovering-variant}. Nous étudierons
l'augmentation
$$
a : \Sh((\mathcal{C}/K)_{total}) \longrightarrow \Sh(\mathcal{C})
$$
de la section \ref{spaces-simplicial-section-simplicial-semi-representable}.

\begin{lemma}
\label{spaces-simplicial-lemma-hypercovering-descent-sheaves}
Soit $\mathcal{C}$ un site admettant des égalisateurs et des produits fibrés.
Soit $K$ un hyperrecouvrement. Alors
\begin{enumerate}
\item $a^{-1} : \Sh(\mathcal{C}) \to \Sh((\mathcal{C}/K)_{total})$
est pleinement fidèle, d'image essentielle les faisceaux cartésiens d'ensembles,
\item $a^{-1} : \textit{Ab}(\mathcal{C}) \to
\textit{Ab}((\mathcal{C}/K)_{total})$
est pleinement fidèle, d'image essentielle les faisceaux cartésiens
de groupes abéliens.
\end{enumerate}
Dans les deux cas, $a_*$ est un foncteur quasi-inverse.
\end{lemma}

\begin{proof}
Le cas des faisceaux abéliens découle immédiatement du cas
de faisceaux d'ensembles, puisque le foncteur $a^{-1}$ commute aux produits.
Dans la suite de la preuve, nous travaillons avec des faisceaux d'ensembles.
Observons que $a^{-1}\mathcal{F}$ est cartésien pour
tout $\mathcal{F}$ dans $\Sh(\mathcal{C})$, d'après le
lemme \ref{spaces-simplicial-lemma-augmentation-cartesian-module}.
Il suffit de montrer que le morphisme d'adjonction
$\mathcal{F} \to a_*a^{-1}\mathcal{F}$
est un isomorphisme pour tout $\mathcal{F}$ dans $\Sh(\mathcal{C})$
et que, pour tout faisceau cartésien
$\mathcal{G}$ sur $(\mathcal{C}/K)_{total}$,
le morphisme d'adjonction
$a^{-1}a_*\mathcal{G} \to \mathcal{G}$ est un isomorphisme.

\medskip\noindent
Soit $\mathcal{F}$ un faisceau sur $\mathcal{C}$.
Rappelons que $a_*a^{-1}\mathcal{F}$ est l'égaliseur
des deux morphismes $a_{0, *}a_0^{-1}\mathcal{F} \to a_{1, *}a_1^{-1}\mathcal{F}$;
voir le lemme \ref{spaces-simplicial-lemma-comparison}.
D'après le lemme \ref{spaces-simplicial-lemma-push-pull-localization},
$$
a_{0, *}a_0^{-1}\mathcal{F} = \SheafHom(F(K_0)^\#, \mathcal{F})
\quad\text{et}\quad
a_{1, *}a_1^{-1}\mathcal{F} = \SheafHom(F(K_1)^\#, \mathcal{F})
$$
D'autre part, on sait que
$$
\xymatrix{
F(K_1)^\# \ar@<1ex>[r] \ar@<-1ex>[r] &
F(K_0)^\# \ar[r] & \text{objet final }*\text{ de }\Sh(\mathcal{C})
}
$$
est un diagramme coégalisateur dans les faisceaux d'ensembles, par définition d'
un hyperrecouvrement. Il suffit donc de prouver
que $\SheafHom(-, \mathcal{F})$ transforme les coégaliseurs
en égaliseurs, ce qui résulte immédiatement de la construction
dans Sites, section \ref{sites-section-glueing-sheaves}.

\medskip\noindent
Soit $\mathcal{G}$ un faisceau cartésien sur $(\mathcal{C}/K)_{total}$.
Montrons que $\mathcal{G} = a^{-1}\mathcal{F}$ pour un certain faisceau
$\mathcal{F}$ sur $\mathcal{C}$. Cela achèvera la preuve, car
on aura alors $a^{-1}a_*\mathcal{G} = a^{-1}a_*a^{-1}\mathcal{F} =
a^{-1}\mathcal{F} = \mathcal{G}$ par le résultat du paragraphe précédent.
Posons $\mathcal{K}_n = F(K_n)^\#$ pour $n \geq 0$. On dispose alors de morphismes de faisceaux
$$
\xymatrix{
\mathcal{K}_2
\ar@<1ex>[r]
\ar@<0ex>[r]
\ar@<-1ex>[r]
&
\mathcal{K}_1
\ar@<0.5ex>[r]
\ar@<-0.5ex>[r]
&
\mathcal{K}_0
}
$$
provenant du fait que $K$ est un objet semi-représentable simplicial.
Le fait que $K$ soit un hyperrecouvrement signifie que
$$
\mathcal{K}_1 \to \mathcal{K}_0 \times \mathcal{K}_0
\quad\text{et}\quad
\mathcal{K}_2 \to
\left(\text{cosk}_1(
\xymatrix{
\mathcal{K}_1
\ar@<0.5ex>[r]
\ar@<-0.5ex>[r]
&
\mathcal{K}_0 \ar[l]
})\right)_2
$$
sont des morphismes surjectifs de faisceaux. En utilisant la description des faisceaux cartésiens sur
$(\mathcal{C}/K)_{total}$ donnée dans le lemme \ref{spaces-simplicial-lemma-characterize-cartesian}
et en utilisant la description de $\Sh(\mathcal{C}/K_n)$ dans
le lemme \ref{spaces-simplicial-lemma-localize-compare},
on voit que notre problème peut se formuler entièrement\footnote{Bien que
la formulation précise n'a pas d'importance, nous l'expliquons :
le problème consiste à montrer qu'étant donné un objet
$\mathcal{G}_0/\mathcal{K}_0$ de $\Sh(\mathcal{C})/\mathcal{K}_0$
et un isomorphisme
$$
\alpha :
\mathcal{G}_0 \times_{\mathcal{K}_0, \mathcal{K}(\delta^1_1)} \mathcal{K}_1 \to
\mathcal{G}_0 \times_{\mathcal{K}_0, \mathcal{K}(\delta^1_0)} \mathcal{K}_1
$$
sur $\mathcal{K}_1$ satisfaisant une condition de cocycle dans
$\Sh(\mathcal{C})/\mathcal{K}_2$, il existe
$\mathcal{F}$ dans $\Sh(\mathcal{C})$ et un isomorphisme
$\mathcal{F} \times \mathcal{K}_0 \to \mathcal{G}_0$ sur $\mathcal{K}_0$
compatible avec $\alpha$.} en termes de
\begin{enumerate}
\item le topos $\Sh(\mathcal{C})$, et
\item l'objet simplicial $\mathcal{K}$ de $\Sh(\mathcal{C})$
dont les termes sont $\mathcal{K}_n$.
\end{enumerate}
Ainsi, après avoir remplacé $\mathcal{C}$ par un autre site $\mathcal{C}'$
comme dans Sites, lemme \ref{sites-lemma-topos-good-site}, on peut supposer
$\mathcal{C}$ admet toutes les limites finies,
la topologie sur $\mathcal{C}$ est sous-canonique,
une famille $\{V_j \to V\}$ de morphismes de $\mathcal{C}$
est un recouvrement si et seulement si $\coprod h_{V_j} \to V$ est surjectif, et
il existe un objet simplicial $U$ de $\mathcal{C}$
tel que $\mathcal{K}_n = h_{U_n}$ comme faisceaux simpliciaux.
En transportant la situation par ces équivalences, on peut supposer
$K_n = \{U_n\}$ pour tout $n$.

\medskip\noindent
Soit $X$ l'objet final de $\mathcal{C}$.
Alors $\{U_0 \to X\}$ est un recouvrement,
$\{U_1 \to U_0 \times U_0\}$ est un recouvrement et
$\{U_2 \to (\text{cosk}_1 \text{sk}_1 U)_2\}$ est un recouvrement.
Notons $d^n_i : U_n \to U_{n - 1}$ et
$s^n_j : U_n \to U_{n + 1}$ les morphismes correspondant
à $\delta^n_i$ et $\sigma^n_j$, comme dans
Simplicial, définition \ref{simplicial-definition-face-degeneracy}.
Par abus de notation, pour un morphisme
$c : V \to W$ de $\mathcal{C}$, on note le morphisme de topos
$c : \Sh(\mathcal{C}/V) \to \Sh(\mathcal{C}/W)$ de la même lettre.
Le faisceau $\mathcal{G}$ est alors déterminé par un faisceau $\mathcal{G}_0$
sur $\mathcal{C}/U_0$ et un isomorphisme
$\alpha : (d^1_1)^{-1}\mathcal{G}_0 \to (d^1_0)^{-1}\mathcal{G}_0$
satisfaisant à la condition de cocycle sur $\mathcal{C}/U_2$
formulée dans le lemme \ref{spaces-simplicial-lemma-characterize-cartesian}.
Puisque $\{U_2 \to (\text{cosk}_1 \text{sk}_1 U)_2\}$
est un recouvrement, le foncteur image inverse correspondant
sur les faisceaux est fidèle (nous omettons cette vérification élémentaire).
On peut donc remplacer $U$ par $\text{cosk}_1 \text{sk}_1 U$, car
cela remplace $U_2$ par $(\text{cosk}_1 \text{sk}_1 U)_2$ sans modifier
$U_1$ ni $U_0$. Alors
$$
(d^2_0, d^2_1, d^2_2) : U_2 \to U_1 \times U_1 \times U_1
$$
est un monomorphisme dont l'image sur les points à valeurs dans $T$ est
décrite dans Simplicial, lemme \ref{simplicial-lemma-work-out}.
En particulier, il existe un morphisme $c$ tel que le diagramme suivant soit commutatif :
$$
\xymatrix{
U_1 \times_{(d^1_1, d^1_0), U_0 \times U_0, (d^1_1, d^1_0)} U_1
\ar[d] \ar[rr]_c & & U_2 \ar[d] \\
U_1 \times U_1
\ar[rr]^{(\text{pr}_1, \text{pr}_2, s^0_0 \circ d^1_1 \circ \text{pr}_1)} & &
U_1 \times U_1 \times U_1
}
$$
car le parcours dans l'autre sens définit un point de $U_2$.
L'image inverse de la condition de cocycle vérifiée par $\alpha$ sur $U_2$
exprime que les images inverses de $\alpha$
par les deux projections sur
$U_1 \times_{(d^1_1, d^1_0), U_0 \times U_0, (d^1_1, d^1_0)} U_1$
coïncident, tandis que l'image inverse de $\alpha$ par
$s^0_0 \circ d^1_1 \circ \text{pr}_1$ est le morphisme identité
(en effet, l'image inverse de $\alpha$ par $s^0_0$ est l'identité).
D'après Sites, lemme \ref{sites-lemma-glue-maps},
on en déduit que $\alpha$ provient d'un isomorphisme
$$
\alpha' : \text{pr}_1^{-1}\mathcal{G}_0 \to \text{pr}_2^{-1}\mathcal{G}_0
$$
de faisceaux sur $\mathcal{C}/U_0 \times U_0$.
Enfin, le morphisme $U_2 \to U_0 \times U_0 \times U_0$
est surjectif sur les faisceaux associés, comme on le voit en utilisant la surjectivité
de $U_1 \to U_0 \times U_0$
et la description de $U_2$ indiquée ci-dessus. Par conséquent, $\alpha'$
satisfait la condition de cocycle sur $U_0 \times U_0 \times U_0$.
La preuve s'achève en appliquant
Sites, lemme \ref{sites-lemma-mapping-property-glue}
au recouvrement $\{U_0 \to X\}$.
\end{proof}

\begin{lemma}
\label{spaces-simplicial-lemma-hypercovering-cech-complex}
Soit $\mathcal{C}$ un site admettant des égalisateurs et des produits fibrés.
Soit $K$ un hyperrecouvrement. Le complexe de Čech
du lemme \ref{spaces-simplicial-lemma-augmentation-cech-complex}, associé à
$a^{-1}\mathcal{F}$
$$
a_{0, *}a_0^{-1}\mathcal{F} \to a_{1, *}a_1^{-1}\mathcal{F} \to
a_{2, *}a_2^{-1}\mathcal{F} \to \ldots
$$
s'identifie au complexe $\SheafHom(s(\mathbf{Z}_{F(K)}^\#), \mathcal{F})$.
Ici, $s(\mathbf{Z}_{F(K)}^\#)$ est le complexe défini dans
Hyperrecouvrements, définition \ref{hypercovering-definition-homology}.
\end{lemma}

\begin{proof}
D'après le lemme \ref{spaces-simplicial-lemma-push-pull-localization}, on a
$$
a_{n, *}a_n^{-1}\mathcal{F} = \SheafHom'(F(K_n)^\#, \mathcal{F})
$$
où $\SheafHom'$ est comme dans Sites, section \ref{sites-section-glueing-sheaves}.
Les morphismes de bord du complexe du
lemme \ref{spaces-simplicial-lemma-augmentation-cech-complex}
proviennent de la structure simpliciale.
L'égalité des complexes résulte donc
des identifications canoniques
$\SheafHom'(\mathcal{G}, \mathcal{F}) =
\SheafHom(\mathbf{Z}_\mathcal{G}, \mathcal{F})$ pour
$\mathcal{G}$ dans $\Sh(\mathcal{C})$.
\end{proof}

\begin{lemma}
\label{spaces-simplicial-lemma-hypercovering-descent-bounded-abelian}
Soit $\mathcal{C}$ un site admettant des égalisateurs et des produits fibrés.
Soit $K$ un hyperrecouvrement. Pour
$E \in D(\mathcal{C})$, le morphisme
$$
E \longrightarrow Ra_*a^{-1}E
$$
est un isomorphisme.
\end{lemma}

\begin{proof}
D'abord, soit $\mathcal{I}$ un faisceau abélien injectif sur $\mathcal{C}$.
La suite spectrale du lemme
\ref{spaces-simplicial-lemma-augmentation-spectral-sequence},
appliquée au faisceau $a^{-1}\mathcal{I}$, dégénère puisque
$(a^{-1}\mathcal{I})_p = a_p^{-1}\mathcal{I}$
est injectif par le lemme \ref{spaces-simplicial-lemma-localize-injective}.
Ainsi le complexe
$$
a_{0, *}a_0^{-1}\mathcal{I} \to
a_{1, *}a_1^{-1}\mathcal{I} \to
a_{2, *}a_2^{-1}\mathcal{I} \to\ldots
$$
calcule $Ra_*a^{-1}\mathcal{I}$. Par
le lemme \ref{spaces-simplicial-lemma-hypercovering-cech-complex},
ce complexe s'identifie à
$\SheafHom(s(\mathbf{Z}_{F(K)}^\#), \mathcal{I})$.
Puisque $K$ est un hyperrecouvrement,
$s(\mathbf{Z}_{F(K)}^\#)$ est exact en degrés $> 0$ par
Hyperrecouvrements, lemme \ref{hypercovering-lemma-acyclic-hypercover-sheaves}
appliqué au préfaisceau simplicial $F(K)$.
Puisque $\mathcal{I}$ est injectif, le foncteur $\SheafHom(-, \mathcal{I})$
est exact; on en déduit que
$\SheafHom(s(\mathbf{Z}_{F(K)}^\#), \mathcal{I})$
est exact en degrés positifs. Ainsi,
$R^pa_*a^{-1}\mathcal{I} = 0$ pour $p > 0$.
D'autre part, $\mathcal{I} = a_*a^{-1}\mathcal{I}$ d'après le
du lemme \ref{spaces-simplicial-lemma-hypercovering-descent-sheaves}.

\medskip\noindent
Cas borné. Soit $E \in D^+(\mathcal{C})$.
Choisissons un complexe borné inférieurement $\mathcal{I}^\bullet$ d'injectifs
représentant $E$. D'après le premier paragraphe et le
lemme d'acyclicité de Leray
(Catégories dérivées, lemme \ref{derived-lemma-leray-acyclicity}),
$Ra_*a^{-1}\mathcal{I}^\bullet$
est calculé par le complexe
$a_*a^{-1}\mathcal{I}^\bullet = \mathcal{I}^\bullet$
ce qui prouve le lemme dans ce cas.

\medskip\noindent
Cas non borné. Nous invitons le lecteur à sauter ce passage, puisque l'argument
est le même que ci-dessus, sauf que nous utilisons la représentation explicite
par doubles complexes pour contourner les problèmes de convergence.
Soit $E \in D(\mathcal{C})$.
Pour montrer que $E \to Ra_*a^{-1}E$ est un isomorphisme,
il suffit de montrer, pour tout objet $U$ de $\mathcal{C}$, que
$$
R\Gamma(U, E) = R\Gamma(U, Ra_*a^{-1}E)
$$
Calculons les deux membres et montrons que le morphisme $E \to Ra_*a^{-1}E$
induit un isomorphisme. Choisissons un complexe K-injectif
$\mathcal{I}^\bullet$ représentant $E$, puis un quasi-isomorphisme
$a^{-1}\mathcal{I}^\bullet \to \mathcal{J}^\bullet$
pour un complexe K-injectif $\mathcal{J}^\bullet$ sur
$(\mathcal{C}/K)_{total}$. Nous avons
$$
R\Gamma(U, E) = R\Hom(\mathbf{Z}_U^\#, E)
$$
et
$$
R\Gamma(U, Ra_*a^{-1}E) = R\Hom(\mathbf{Z}_U^\#, Ra_*a^{-1}E) =
R\Hom(a^{-1}\mathbf{Z}_U^\#, a^{-1}E)
$$
D'après le lemme \ref{spaces-simplicial-lemma-simplicial-resolution-augmentation},
nous avons un quasi-isomorphisme
$$
\Big(\ldots \to
g_{2!}(a_2^{-1}\mathbf{Z}_U^\#) \to
g_{1!}(a_1^{-1}\mathbf{Z}_U^\#) \to
g_{0!}(a_0^{-1}\mathbf{Z}_U^\#)\Big)
\longrightarrow
a^{-1}\mathbf{Z}_U^\#
$$
Par conséquent, $R\Hom(a^{-1}\mathbf{Z}_U^\#, a^{-1}E)$ est égal à
$$
R\Gamma((\mathcal{C}/K)_{total},
R\SheafHom(
\ldots \to
g_{2!}(a_2^{-1}\mathbf{Z}_U^\#) \to
g_{1!}(a_1^{-1}\mathbf{Z}_U^\#) \to
g_{0!}(a_0^{-1}\mathbf{Z}_U^\#),
\mathcal{J}^\bullet))
$$
D'après la construction de Cohomologie sur les sites, section
\ref{sites-cohomology-section-internal-hom}
et puisque $\mathcal{J}^\bullet$ est K-injectif, on voit que
l'objet précédent est représenté par le complexe de groupes abéliens dont les termes sont
$$
\prod\nolimits_{p + q = n}
\Hom(g_{p!}(a_p^{-1}\mathbf{Z}_U^\#), \mathcal{J}^q) =
\prod\nolimits_{p + q = n}
\Hom(a_p^{-1}\mathbf{Z}_U^\#, g_p^{-1}\mathcal{J}^q)
$$
Voir Cohomologie sur les sites, lemmes
\ref{sites-cohomology-lemma-RHom-into-K-injective} et
\ref{sites-cohomology-lemma-section-RHom-over-U} pour plus d'informations.
Ainsi, $R\Gamma(U, Ra_*a^{-1}E)$ est calculé par
le complexe total produit $\text{Tot}_\pi(B^{\bullet, \bullet})$
avec $B^{p, q} = \Hom(a_p^{-1}\mathbf{Z}_U^\#, g_p^{-1}\mathcal{J}^q)$.
Pour l'autre membre, raisonnons de la même manière. Observons d'abord que
$$
s(\mathbf{Z}_{F(K)}^\#) \longrightarrow \mathbf{Z}
$$
est un quasi-isomorphisme de complexes sur $\mathcal{C}$
d'après Hyperrecouvrements, lemme \ref{hypercovering-lemma-acyclic-hypercover-sheaves}.
Puisque $\mathbf{Z}_U^\#$ est un faisceau plat de $\mathbf{Z}$-modules,
on voit que
$$
s(\mathbf{Z}_{F(K)}^\#) \otimes_\mathbf{Z} \mathbf{Z}_U^\#
\longrightarrow
\mathbf{Z}_U^\#
$$
est un quasi-isomorphisme. Par conséquent,
$R\Hom(\mathbf{Z}_U^\#, E)$ est égal à
$$
R\Gamma(\mathcal{C}, R\SheafHom(
s(\mathbf{Z}_{F(K)}^\#) \otimes_\mathbf{Z} \mathbf{Z}_U^\#,
\mathcal{I}^\bullet))
$$
Par construction de $R\SheafHom$, et puisque $\mathcal{I}^\bullet$
est K-injectif, l'objet précédent est représenté par le complexe de groupes abéliens
dont les termes sont
$$
\prod\nolimits_{p + q = n}
\Hom(\mathbf{Z}^\#_{K_p} \otimes_\mathbf{Z} \mathbf{Z}_U^\#, \mathcal{I}^q)
=
\prod\nolimits_{p + q = n}
\Hom(a_p^{-1}\mathbf{Z}_U^\#, a_p^{-1}\mathcal{I}^q)
$$
L'égalité des termes découle du fait que
$\mathbf{Z}^\#_{K_p} \otimes_\mathbf{Z} \mathbf{Z}_U^\# =
a_{p!}a_p^{-1}\mathbf{Z}_U^\#$ d'après Modules sur les sites,
remarque \ref{sites-modules-remark-j-shriek-tensor}.
Ainsi, $R\Gamma(U, E)$ est calculé par
le complexe total produit $\text{Tot}_\pi(A^{\bullet, \bullet})$
avec $A^{p, q} = \Hom(a_p^{-1}\mathbf{Z}_U^\#, a_p^{-1}\mathcal{I}^q)$.

\medskip\noindent
Puisque $\mathcal{I}^\bullet$ est K-injectif, nous voyons que
$a_p^{-1}\mathcal{I}^\bullet$ est K-injectif, voir
le lemme \ref{spaces-simplicial-lemma-localize-injective}.
Puisque $\mathcal{J}^\bullet$ est K-injectif, nous voyons que
$g_p^{-1}\mathcal{J}^\bullet$ est K-injectif, voir
le lemme \ref{spaces-simplicial-lemma-restriction-injective-to-component-site}.
Ces deux complexes représentent l'objet $a_p^{-1}E$.
Ainsi, pour tout $p \geq 0$, le morphisme de complexes
$$
A^{p, \bullet} = \Hom(a_p^{-1}\mathbf{Z}_U^\#, a_p^{-1}\mathcal{I}^\bullet)
\longrightarrow
\Hom(a_p^{-1}\mathbf{Z}_U^\#, g_p^{-1}\mathcal{J}^\bullet) = B^{p, \bullet}
$$
obtenu en appliquant $g_p^{-1}$ au morphisme donné
$a^{-1}\mathcal{I}^\bullet \to \mathcal{J}^\bullet$
est un quasi-isomorphisme, puisque ces deux complexes calculent
$$
R\Hom(a_p^{-1}\mathbf{Z}_U^\#, a_p^{-1}E)
$$
D'après Compléments d'algèbre, lemme \ref{more-algebra-lemma-prod-qis-gives-qis},
nous concluons que la flèche verticale droite dans le diagramme commutatif
$$
\xymatrix{
R\Gamma(U, E) \ar[r] \ar[d] &
\text{Tot}_\pi(A^{\bullet, \bullet}) \ar[d] \\
R\Gamma(U, Ra_*a^{-1}E) \ar[r] &
\text{Tot}_\pi(B^{\bullet, \bullet})
}
$$
est un quasi-isomorphisme. Puisque nous avons vu plus haut que les flèches horizontales
sont des quasi-isomorphismes, la flèche verticale gauche l'est également.
\end{proof}

\begin{lemma}
\label{spaces-simplicial-lemma-compare-cohomology-hypercovering}
Soit $\mathcal{C}$ un site admettant des égalisateurs et des produits fibrés.
Soit $K$ un hyperrecouvrement.
Nous avons alors un isomorphisme canonique
$$
R\Gamma(\mathcal{C}, E) =
R\Gamma((\mathcal{C}/K)_{total}, a^{-1}E)
$$
pour $E \in D(\mathcal{C})$.
\end{lemma}

\begin{proof}
Cela découle du lemme \ref{spaces-simplicial-lemma-hypercovering-descent-bounded-abelian}
car $R\Gamma((\mathcal{C}/K)_{total}, -) =
R\Gamma(\mathcal{C}, -) \circ Ra_*$ par
Cohomologie sur les sites, remarque \ref{sites-cohomology-remark-before-Leray}.
\end{proof}

\begin{lemma}
\label{spaces-simplicial-lemma-hypercovering-equivalence-bounded}
Soit $\mathcal{C}$ un site admettant des égalisateurs et des produits fibrés.
Soit $K$ un hyperrecouvrement.
Soit $\mathcal{A} \subset \textit{Ab}((\mathcal{C}/K)_{total})$
la sous-catégorie de Serre faible des faisceaux abéliens cartésiens.
Alors le foncteur $a^{-1}$ définit une équivalence
$$
D^+(\mathcal{C}) \longrightarrow D_\mathcal{A}^+((\mathcal{C}/K)_{total})
$$
de quasi-inverse $Ra_*$.
\end{lemma}

\begin{proof}
Observons que $\mathcal{A}$ est une sous-catégorie de Serre faible d'après
le lemme \ref{spaces-simplicial-lemma-Serre-subcat-cartesian-modules}.
L'équivalence est une conséquence formelle des
résultats obtenus jusqu'ici. On utilise les lemmes
\ref{spaces-simplicial-lemma-hypercovering-descent-sheaves} et
\ref{spaces-simplicial-lemma-hypercovering-descent-bounded-abelian}, ainsi que
Cohomologie sur les sites, lemme \ref{sites-cohomology-lemma-equivalence-bounded}.
\end{proof}

\noindent
Nous invitons le lecteur à ignorer la remarque suivante.

\begin{remark}
\label{spaces-simplicial-remark-compare-cohomology-hypercovering-presheaf}
Soient $\mathcal{C}$ un site et $\mathcal{G}$ un préfaisceau d'ensembles sur
$\mathcal{C}$. Si $\mathcal{C}$ admet des égalisateurs et des produits fibrés,
nous avons défini la notion d'hyperrecouvrement de $\mathcal{G}$ dans
Hyperrecouvrements, définition \ref{hypercovering-definition-hypercovering-variant}.
Tous les résultats de cette section se transposent à
ce cadre.
Pour le voir,
considérons la localisation $\mathcal{C}/\mathcal{G}$
de $\mathcal{C}$ en $\mathcal{G}$, définie exactement comme dans
Sites, lemme \ref{sites-lemma-localize-topos-site}
(qui n'est énoncé que pour les faisceaux; le topos
$\Sh(\mathcal{C}/\mathcal{G})$ est égal à la localisation
du topos $\Sh(\mathcal{C})$ au faisceau $\mathcal{G}^\#$).
Le lecteur vérifie alors aisément que le site
$\mathcal{C}/\mathcal{G}$ admet des produits fibrés et des égalisateurs,
et qu'un hyperrecouvrement de $\mathcal{G}$ dans $\mathcal{C}$
équivaut à un hyperrecouvrement pour le site $\mathcal{C}/\mathcal{G}$.
Ainsi, en remplaçant le site $\mathcal{C}$ par $\mathcal{C}/\mathcal{G}$
dans les lemmes ci-dessus sur les hyperrecouvrements, on obtient les démonstrations des
résultats correspondants pour les hyperrecouvrements de $\mathcal{G}$.
Exemple : pour un hyperrecouvrement $K$ de $\mathcal{G}$, on a
$$
R\Gamma(\mathcal{C}/\mathcal{G}, E) =
R\Gamma((\mathcal{C}/K)_{total}, a^{-1}E)
$$
pour $E \in D^+(\mathcal{C}/\mathcal{G})$, où
$a : \Sh((\mathcal{C}/K)_{total}) \to \Sh(\mathcal{C}/\mathcal{G})$
est l'augmentation canonique. C'est un cas particulier du lemme
\ref{spaces-simplicial-lemma-compare-cohomology-hypercovering}.
Notons $R\Gamma(\mathcal{G}, -) : D(\mathcal{C}) \to D(\textit{Ab})$ le foncteur
défini comme le foncteur dérivé du foncteur
$H^0(\mathcal{G}, -) = H^0(\mathcal{G}^\#, -)$
étudié dans Hyperrecouvrements, section
\ref{hypercovering-section-hypercoverings-verdier}, et dans
Cohomologie sur les sites, section \ref{sites-cohomology-section-limp}.
Nous avons
$$
R\Gamma(\mathcal{G}, E) = R\Gamma(\mathcal{C}/\mathcal{G}, j^{-1}E)
$$
d'après l'analogue du lemme de Cohomologie sur les sites
\ref{sites-cohomology-lemma-cohomology-of-open}
pour le foncteur de localisation $j : \mathcal{C}/\mathcal{G} \to \mathcal{C}$.
En combinant ces égalités, on obtient
$$
R\Gamma(\mathcal{G}, E) =
R\Gamma((\mathcal{C}/K)_{total}, a^{-1}j^{-1}E) =
R\Gamma((\mathcal{C}/K)_{total}, g^{-1}E)
$$
pour $E \in D^+(\mathcal{C})$, où
$g : \Sh((\mathcal{C}/K)_{total}) \to \Sh(\mathcal{C})$
est la composition de $a$ et $j$.
\end{remark}







\section{Descente cohomologique pour les hyperrecouvrements : modules}
\label{spaces-simplicial-section-cohomological-descent-hypercoverings-modules}

\noindent
Soit $\mathcal{C}$ un site et soit $\mathcal{O}_\mathcal{C}$
un faisceau d'anneaux. Supposons que $\mathcal{C}$
admette des égalisateurs et des produits fibrés, et soit $K$ un hyperrecouvrement
au sens d'Hyperrecouvrements, définition
\ref{hypercovering-definition-hypercovering-variant}. Nous étudierons la descente cohomologique
pour l'augmentation
$$
a :
(\Sh((\mathcal{C}/K)_{total}), \mathcal{O})
\longrightarrow
(\Sh(\mathcal{C}), \mathcal{O}_\mathcal{C})
$$
de la remarque \ref{spaces-simplicial-remark-augmentation-ringed}.

\begin{lemma}
\label{spaces-simplicial-lemma-hypercovering-descent-modules}
Soit $\mathcal{C}$ un site admettant des égalisateurs et des produits fibrés.
Soit $\mathcal{O}_\mathcal{C}$ un faisceau d'anneaux.
Soit $K$ un hyperrecouvrement. Avec les notations ci-dessus,
$$
a^* : \textit{Mod}(\mathcal{O}_\mathcal{C}) \to \textit{Mod}(\mathcal{O})
$$
est pleinement fidèle, d'image essentielle les $\mathcal{O}$-modules cartésiens.
Le foncteur $a_*$ est un quasi-inverse.
\end{lemma}

\begin{proof}
Puisque $a^{-1}\mathcal{O}_\mathcal{C} = \mathcal{O}$, on a
$a^* = a^{-1}$. Le lemme résulte donc immédiatement
du lemme \ref{spaces-simplicial-lemma-hypercovering-descent-sheaves}.
\end{proof}

\begin{lemma}
\label{spaces-simplicial-lemma-hypercovering-descent-bounded-modules}
Soit $\mathcal{C}$ un site admettant des égalisateurs et des produits fibrés.
Soit $\mathcal{O}_\mathcal{C}$ un faisceau d'anneaux.
Soit $K$ un hyperrecouvrement. Pour
$E \in D(\mathcal{O}_\mathcal{C})$, le morphisme
$$
E \longrightarrow Ra_*La^*E
$$
est un isomorphisme.
\end{lemma}

\begin{proof}
Puisque $a^{-1}\mathcal{O}_\mathcal{C} = \mathcal{O}$, on a
$La^* = a^* = a^{-1}$. De plus, $Ra_*$ coïncide avec
$Ra_*$ sur les faisceaux abéliens; voir Cohomologie
sur les sites, lemme
\ref{sites-cohomology-lemma-modules-abelian-unbounded}.
Le lemme en résulte immédiatement,
du lemme \ref{spaces-simplicial-lemma-hypercovering-descent-bounded-abelian}.
\end{proof}

\begin{lemma}
\label{spaces-simplicial-lemma-compare-cohomology-hypercovering-modules}
Soit $\mathcal{C}$ un site admettant des égalisateurs et des produits fibrés.
Soit $\mathcal{O}_\mathcal{C}$ un faisceau d'anneaux.
Soit $K$ un hyperrecouvrement.
Nous avons alors un isomorphisme canonique
$$
R\Gamma(\mathcal{C}, E) =
R\Gamma((\mathcal{C}/K)_{total}, La^*E)
$$
pour $E \in D(\mathcal{O}_\mathcal{C})$.
\end{lemma}

\begin{proof}
Cela découle du lemme \ref{spaces-simplicial-lemma-hypercovering-descent-bounded-modules}
car $R\Gamma((\mathcal{C}/K)_{total}, -) =
R\Gamma(\mathcal{C}, -) \circ Ra_*$ par
Cohomologie sur les sites, remarque \ref{sites-cohomology-remark-before-Leray},
ou par
Cohomologie sur les sites, lemme \ref{sites-cohomology-lemma-Leray-unbounded}.
\end{proof}

\begin{lemma}
\label{spaces-simplicial-lemma-hypercovering-equivalence-bounded-modules}
Soit $\mathcal{C}$ un site admettant des égalisateurs et des produits fibrés.
Soit $\mathcal{O}_\mathcal{C}$ un faisceau d'anneaux.
Soit $K$ un hyperrecouvrement.
Soit $\mathcal{A} \subset \textit{Mod}(\mathcal{O})$
la sous-catégorie de Serre faible des $\mathcal{O}$-modules cartésiens.
Alors le foncteur $La^*$ définit une équivalence
$$
D^+(\mathcal{O}_\mathcal{C}) \longrightarrow D_\mathcal{A}^+(\mathcal{O})
$$
de quasi-inverse $Ra_*$.
\end{lemma}

\begin{proof}
Observons que $\mathcal{A}$ est une sous-catégorie de Serre faible d'après le
lemme \ref{spaces-simplicial-lemma-Serre-subcat-cartesian-modules}
(les hypothèses requises découlent de la discussion de la
remarque \ref{spaces-simplicial-remark-augmentation-ringed}).
L'équivalence est une conséquence formelle
des résultats obtenus jusqu'ici. On utilise les lemmes
\ref{spaces-simplicial-lemma-hypercovering-descent-modules},
\ref{spaces-simplicial-lemma-hypercovering-descent-bounded-modules}, ainsi que Cohomologie
sur les sites, lemme \ref{sites-cohomology-lemma-equivalence-bounded}.
\end{proof}






\section{Descente cohomologique pour les hyperrecouvrements d'un objet}
\label{spaces-simplicial-section-cohomological-descent-hypercoverings-X}

\noindent
Dans cette section, nous supposons que $\mathcal{C}$ admet des produits fibrés
et nous fixons $X \in \Ob(\mathcal{C})$. Soit $K$ un hyperrecouvrement de $X$
au sens de
Hyperrecouvrements, définition \ref{hypercovering-definition-hypercovering}.
Nous étudierons l'augmentation
$$
a : \Sh((\mathcal{C}/K)_{total}) \longrightarrow \Sh(\mathcal{C}/X)
$$
de la remarque \ref{spaces-simplicial-remark-augmentation-over-object}.
Observons que $\mathcal{C}/X$ est un site qui admet des égalisateurs
et des produits fibrés, et que $K$ est un hyperrecouvrement
du site $\mathcal{C}/X$\footnote{La réciproque peut être fausse :
si $K$ est un objet simplicial de
$\text{SR}(\mathcal{C}, X) = \text{SR}(\mathcal{C}/X)$
qui est un hyperrecouvrement du site $\mathcal{C}/X$ au sens de
Hyperrecouvrements, définition \ref{hypercovering-definition-hypercovering-variant},
alors $K$ ne définit pas nécessairement un hyperrecouvrement de $X$.
Par exemple, si $K_0 = \{U_{0, i}\}_{i \in I_0}$,
cette dernière condition garantit que
$\{U_{0, i} \to X\}$ est un recouvrement de $\mathcal{C}$
tandis que la première exige seulement que
$\coprod h_{U_{0, i}}^\# \to h_X^\#$ soit un morphisme surjectif
de faisceaux;} voir Hyperrecouvrements, lemme
\ref{hypercovering-lemma-hypercovering-F}.
Ainsi, tout résultat démontré pour les hyperrecouvrements
dans la section \ref{spaces-simplicial-section-cohomological-descent-hypercoverings}
se transpose immédiatement à la situation de la présente section.

\begin{lemma}
\label{spaces-simplicial-lemma-hypercovering-X-descent-sheaves}
Soit $\mathcal{C}$ un site avec des produits fibrés et $X \in \Ob(\mathcal{C})$.
Soit $K$ un hyperrecouvrement de $X$. Alors
\begin{enumerate}
\item $a^{-1} : \Sh(\mathcal{C}/X) \to \Sh((\mathcal{C}/K)_{total})$
est pleinement fidèle, d'image essentielle les faisceaux cartésiens d'ensembles,
\item $a^{-1} : \textit{Ab}(\mathcal{C}/X) \to
\textit{Ab}((\mathcal{C}/K)_{total})$
est pleinement fidèle, d'image essentielle les faisceaux cartésiens
de groupes abéliens.
\end{enumerate}
Dans les deux cas, $a_*$ est un foncteur quasi-inverse.
\end{lemma}

\begin{proof}
D'après les remarques \ref{spaces-simplicial-remark-semi-representable-over-object} et
\ref{spaces-simplicial-remark-augmentation-over-object}, ainsi que la discussion de
l'introduction de cette section,
le résultat découle du lemme \ref{spaces-simplicial-lemma-hypercovering-descent-sheaves}.
\end{proof}

\begin{lemma}
\label{spaces-simplicial-lemma-hypercovering-X-descent-bounded-abelian}
Soit $\mathcal{C}$ un site admettant des produits fibrés et soit $X \in \Ob(\mathcal{C})$.
Soit $K$ un hyperrecouvrement de $X$. Pour
$E \in D(\mathcal{C}/X)$, le morphisme
$$
E \longrightarrow Ra_*a^{-1}E
$$
est un isomorphisme.
\end{lemma}

\begin{proof}
D'après les remarques \ref{spaces-simplicial-remark-semi-representable-over-object} et
\ref{spaces-simplicial-remark-augmentation-over-object}, ainsi que la discussion de
l'introduction de cette section,
le résultat découle du lemme \ref{spaces-simplicial-lemma-hypercovering-descent-bounded-abelian}.
\end{proof}

\begin{lemma}
\label{spaces-simplicial-lemma-compare-cohomology-hypercovering-X}
Soit $\mathcal{C}$ un site avec des produits fibrés et $X \in \Ob(\mathcal{C})$.
Soit $K$ un hyperrecouvrement de $X$.
Nous avons alors un isomorphisme canonique
$$
R\Gamma(X, E) = R\Gamma((\mathcal{C}/K)_{total}, a^{-1}E)
$$
pour $E \in D(\mathcal{C}/X)$.
\end{lemma}

\begin{proof}
D'après les remarques \ref{spaces-simplicial-remark-semi-representable-over-object} et
\ref{spaces-simplicial-remark-augmentation-over-object},
le résultat découle du lemme \ref{spaces-simplicial-lemma-compare-cohomology-hypercovering}.
\end{proof}

\begin{lemma}
\label{spaces-simplicial-lemma-hypercovering-X-equivalence-bounded}
Soit $\mathcal{C}$ un site avec des produits fibrés et $X \in \Ob(\mathcal{C})$.
Soit $K$ un hyperrecouvrement de $X$.
Soit $\mathcal{A} \subset \textit{Ab}((\mathcal{C}/K)_{total})$
la sous-catégorie de Serre faible des faisceaux abéliens cartésiens.
Alors le foncteur $a^{-1}$ définit une équivalence
$$
D^+(\mathcal{C}/X) \longrightarrow D_\mathcal{A}^+((\mathcal{C}/K)_{total})
$$
de quasi-inverse $Ra_*$.
\end{lemma}

\begin{proof}
D'après les remarques \ref{spaces-simplicial-remark-semi-representable-over-object} et
\ref{spaces-simplicial-remark-augmentation-over-object},
le résultat découle du lemme \ref{spaces-simplicial-lemma-hypercovering-equivalence-bounded}.
\end{proof}








\section{Descente cohomologique pour les hyperrecouvrements d'un objet : modules}
\label{spaces-simplicial-section-cohomological-descent-hypercoverings-X-modules}

\noindent
Dans cette section, nous supposons que $\mathcal{C}$ admet des produits fibrés
et nous fixons $X \in \Ob(\mathcal{C})$. Soit $K$ un hyperrecouvrement de $X$
au sens de
Hyperrecouvrements, définition \ref{hypercovering-definition-hypercovering}.
Soit $\mathcal{O}_\mathcal{C}$ un faisceau d'anneaux sur $\mathcal{C}$.
Posons $\mathcal{O}_X = \mathcal{O}_\mathcal{C}|_{\mathcal{C}/X}$.
Nous étudierons l'augmentation
$$
a :
(\Sh((\mathcal{C}/K)_{total}), \mathcal{O})
\longrightarrow
(\Sh(\mathcal{C}/X), \mathcal{O}_X)
$$
de la remarque \ref{spaces-simplicial-remark-augmentation-ringed-over-object}.
Observons que $\mathcal{C}/X$ est un site qui admet des égalisateurs
et des produits fibrés, et que $K$ est un hyperrecouvrement
du site $\mathcal{C}/X$.
Par conséquent, les résultats de cette section découlent immédiatement
des résultats correspondants dans la section
\ref{spaces-simplicial-section-cohomological-descent-hypercoverings-modules}.

\begin{lemma}
\label{spaces-simplicial-lemma-hypercovering-X-descent-modules}
Soient $\mathcal{C}$ un site admettant des produits fibrés et $X \in \Ob(\mathcal{C})$.
Soit $\mathcal{O}_\mathcal{C}$ un faisceau d'anneaux.
Soit $K$ un hyperrecouvrement de $X$. Avec les notations ci-dessus,
$$
a^* : \textit{Mod}(\mathcal{O}_X) \to \textit{Mod}(\mathcal{O})
$$
est pleinement fidèle, d'image essentielle les $\mathcal{O}$-modules cartésiens.
Le foncteur $a_*$ est un quasi-inverse.
\end{lemma}

\begin{proof}
D'après les remarques \ref{spaces-simplicial-remark-semi-representable-ringed-over-object} et
\ref{spaces-simplicial-remark-augmentation-ringed-over-object}, ainsi que la discussion de
l'introduction de cette section,
le résultat découle du lemme \ref{spaces-simplicial-lemma-hypercovering-descent-modules}.
\end{proof}

\begin{lemma}
\label{spaces-simplicial-lemma-hypercovering-X-descent-bounded-modules}
Soient $\mathcal{C}$ un site admettant des produits fibrés et $X \in \Ob(\mathcal{C})$.
Soit $\mathcal{O}_\mathcal{C}$ un faisceau d'anneaux.
Soit $K$ un hyperrecouvrement de $X$. Pour
$E \in D(\mathcal{O}_X)$, le morphisme
$$
E \longrightarrow Ra_*La^*E
$$
est un isomorphisme.
\end{lemma}

\begin{proof}
D'après les remarques \ref{spaces-simplicial-remark-semi-representable-ringed-over-object} et
\ref{spaces-simplicial-remark-augmentation-ringed-over-object}, ainsi que la discussion de
l'introduction de cette section,
le résultat découle du lemme \ref{spaces-simplicial-lemma-hypercovering-descent-bounded-modules}.
\end{proof}

\begin{lemma}
\label{spaces-simplicial-lemma-compare-cohomology-hypercovering-X-modules}
Soient $\mathcal{C}$ un site admettant des produits fibrés et $X \in \Ob(\mathcal{C})$.
Soit $\mathcal{O}_\mathcal{C}$ un faisceau d'anneaux.
Soit $K$ un hyperrecouvrement de $X$.
Nous avons alors un isomorphisme canonique
$$
R\Gamma(X, E) = R\Gamma((\mathcal{C}/K)_{total}, La^*E)
$$
pour $E \in D(\mathcal{O}_\mathcal{C})$.
\end{lemma}

\begin{proof}
D'après les remarques \ref{spaces-simplicial-remark-semi-representable-ringed-over-object} et
\ref{spaces-simplicial-remark-augmentation-ringed-over-object}, ainsi que la discussion de
l'introduction de cette section,
le résultat découle du lemme \ref{spaces-simplicial-lemma-compare-cohomology-hypercovering-modules}.
\end{proof}

\begin{lemma}
\label{spaces-simplicial-lemma-hypercovering-X-equivalence-bounded-modules}
Soient $\mathcal{C}$ un site admettant des produits fibrés et $X \in \Ob(\mathcal{C})$.
Soit $\mathcal{O}_\mathcal{C}$ un faisceau d'anneaux.
Soit $K$ un hyperrecouvrement de $X$.
Soit $\mathcal{A} \subset \textit{Mod}(\mathcal{O})$
la sous-catégorie de Serre faible des $\mathcal{O}$-modules cartésiens.
Alors le foncteur $La^*$ définit une équivalence
$$
D^+(\mathcal{O}_X) \longrightarrow D_\mathcal{A}^+(\mathcal{O})
$$
de quasi-inverse $Ra_*$.
\end{lemma}

\begin{proof}
D'après les remarques \ref{spaces-simplicial-remark-semi-representable-ringed-over-object} et
\ref{spaces-simplicial-remark-augmentation-ringed-over-object}, ainsi que la discussion de
l'introduction de cette section,
le résultat découle du lemme \ref{spaces-simplicial-lemma-hypercovering-equivalence-bounded-modules}.
\end{proof}










\section{Hyperrecouvrement par un objet simplicial du site}
\label{spaces-simplicial-section-hypercovering}

\noindent
Soit $\mathcal{C}$ un site admettant des produits fibrés et
soit $X \in \Ob(\mathcal{C})$.
Dans cette section, nous expliquons les résultats de la section
\ref{spaces-simplicial-section-cohomological-descent-hypercoverings-X}
dans le cas où notre hyperrecouvrement est donné par
un objet simplicial du site.
Soit $U$ un objet simplicial de $\mathcal{C}$.
Comme d'habitude, on note $U_n = U([n])$ et $f_\varphi : U_n \to U_m$
le morphisme $f_\varphi = U(\varphi)$ correspondant à
$\varphi : [m] \to [n]$.
Supposons que nous ayons une augmentation
$$
a : U \to X
$$
On en déduit un site simplicial $(\mathcal{C}/U)_{total}$
et un morphisme d'augmentation
$$
a : \Sh((\mathcal{C}/U)_{total}) \longrightarrow \Sh(\mathcal{C}/X)
$$
Plus précisément, $U$ détermine un objet simplicial $K$ de
$\text{SR}(\mathcal{C}, X)$ dont le terme de degré $n$ est
$K_n = \{U_n \to X\}$; on peut donc appliquer les constructions de la
remarque \ref{spaces-simplicial-remark-augmentation-over-object}.
Plus précisément, un objet du site $(\mathcal{C}/U)_{total}$ est de la forme
$V/U_n$, et un morphisme $(\varphi, f) : V/U_n \to W/U_m$ est donné
par un morphisme $\varphi : [m] \to [n]$ dans $\Delta$ et un morphisme
$f : V \to W$ tel que le diagramme
$$
\xymatrix{
V \ar[r]_f \ar[d] & W \ar[d] \\
U_n \ar[r]^{f_\varphi} & U_m
}
$$
est commutatif. Le morphisme de topoi $a$ est défini par le foncteur cocontinu
$V/U_n \mapsto V/X$. Voilà toute la construction.

\medskip\noindent
Dans cette section, nous dirons que l'augmentation $a : U \to X$ est un
{\it hyperrecouvrement de $X$ dans $\mathcal{C}$} si les conditions suivantes sont satisfaites :
\begin{enumerate}
\item $\{U_0 \to X\}$ est un recouvrement de $\mathcal{C}$,
\item $\{U_1 \to U_0 \times_X U_0\}$ est un recouvrement de $\mathcal{C}$,
\item $\{U_{n + 1} \to (\text{cosk}_n\text{sk}_n U)_{n + 1}\}$
est un recouvrement de $\mathcal{C}$ pour $n \geq 1$.
\end{enumerate}
Cela équivaut à dire que $K$ (comme ci-dessus) est un hyperrecouvrement
de $X$; voir
Hyperrecouvrements, exemple \ref{hypercovering-example-hypercovering-in-C}.

\begin{lemma}
\label{spaces-simplicial-lemma-hypercovering-X-simple-descent-sheaves}
Soient $\mathcal{C}$ un site admettant des produits fibrés et $X \in \Ob(\mathcal{C})$.
Soit $a : U \to X$ un hyperrecouvrement de $X$ dans $\mathcal{C}$ comme défini ci-dessus.
Alors
\begin{enumerate}
\item $a^{-1} : \Sh(\mathcal{C}/X) \to \Sh((\mathcal{C}/U)_{total})$
est pleinement fidèle, d'image essentielle les faisceaux cartésiens d'ensembles,
\item $a^{-1} : \textit{Ab}(\mathcal{C}/X) \to
\textit{Ab}((\mathcal{C}/U)_{total})$
est pleinement fidèle, d'image essentielle les faisceaux cartésiens
de groupes abéliens.
\end{enumerate}
Dans les deux cas, $a_*$ est un foncteur quasi-inverse.
\end{lemma}

\begin{proof}
Il s'agit d'un cas particulier du lemme
\ref{spaces-simplicial-lemma-hypercovering-X-descent-sheaves}.
\end{proof}

\begin{lemma}
\label{spaces-simplicial-lemma-hypercovering-X-simple-descent-bounded-abelian}
Soient $\mathcal{C}$ un site admettant des produits fibrés et $X \in \Ob(\mathcal{C})$.
Soit $a : U \to X$ un hyperrecouvrement de $X$ dans $\mathcal{C}$ comme défini ci-dessus.
Pour $E \in D(\mathcal{C}/X)$, le morphisme
$$
E \longrightarrow Ra_*a^{-1}E
$$
est un isomorphisme.
\end{lemma}

\begin{proof}
Il s'agit d'un cas particulier du lemme
\ref{spaces-simplicial-lemma-hypercovering-X-descent-bounded-abelian}.
\end{proof}

\begin{lemma}
\label{spaces-simplicial-lemma-compare-cohomology-hypercovering-X-simple}
Soient $\mathcal{C}$ un site admettant des produits fibrés et $X \in \Ob(\mathcal{C})$.
Soit $a : U \to X$ un hyperrecouvrement de $X$ dans $\mathcal{C}$ comme défini ci-dessus.
Nous avons alors un isomorphisme canonique
$$
R\Gamma(X, E) = R\Gamma((\mathcal{C}/U)_{total}, a^{-1}E)
$$
pour $E \in D(\mathcal{C}/X)$.
\end{lemma}

\begin{proof}
Il s'agit d'un cas particulier du lemme
\ref{spaces-simplicial-lemma-compare-cohomology-hypercovering-X}.
\end{proof}

\begin{lemma}
\label{spaces-simplicial-lemma-hypercovering-X-simple-equivalence-bounded}
Soient $\mathcal{C}$ un site admettant des produits fibrés et $X \in \Ob(\mathcal{C})$.
Soit $a : U \to X$ un hyperrecouvrement de $X$ dans $\mathcal{C}$ comme défini ci-dessus.
Soit $\mathcal{A} \subset \textit{Ab}((\mathcal{C}/U)_{total})$
la sous-catégorie de Serre faible des faisceaux abéliens cartésiens.
Alors le foncteur $a^{-1}$ définit une équivalence
$$
D^+(\mathcal{C}/X) \longrightarrow D_\mathcal{A}^+((\mathcal{C}/U)_{total})
$$
de quasi-inverse $Ra_*$.
\end{lemma}

\begin{proof}
C'est un cas particulier du
lemme \ref{spaces-simplicial-lemma-hypercovering-X-equivalence-bounded}.
\end{proof}

\begin{lemma}
\label{spaces-simplicial-lemma-sr-when-fibre-products}
Soit $U$ un objet simplicial d'un site $\mathcal{C}$
admettant des produits fibrés.
\begin{enumerate}
\item $\mathcal{C}/U$ définit un objet simplicial
dans la catégorie dont les objets sont des sites et
dont les morphismes sont des morphismes de sites,
\item la construction du lemme \ref{spaces-simplicial-lemma-simplicial-site-site}
appliquée à la structure de (1)
reproduit le site $(\mathcal{C}/U)_{total}$ ci-dessus,
\item si $a : U \to X$ est une augmentation, alors
$a_0 : \mathcal{C}/U_0 \to \mathcal{C}/X$ est une augmentation
comme dans la partie (A) de la remarque \ref{spaces-simplicial-remark-augmentation-site} et donne le
même morphisme de topoi
$a : \Sh((\mathcal{C}/U)_{total}) \to \Sh(\mathcal{C}/X)$
identique à celui défini ci-dessus.
\end{enumerate}
\end{lemma}

\begin{proof}
Pour un morphisme d'objets $V \to W$ de $\mathcal{C}$, le morphisme de localisation
$j : \mathcal{C}/V \to \mathcal{C}/W$ est adjoint à gauche au
foncteur de changement de base $\mathcal{C}/W \to \mathcal{C}/V$.
Le foncteur de changement de base est continu et induit le même morphisme de topoi
que $j$. Voir
Sites, lemme \ref{sites-lemma-relocalize-given-fibre-products}.
Cela prouve (1).

\medskip\noindent
La partie (2) résulte de ce qu'un morphisme $V/U_n \to W/U_m$
de la catégorie construite
dans le lemme \ref{spaces-simplicial-lemma-simplicial-site-site}
est un morphisme $V \to W \times_{U_m, f_\varphi} U_n$ sur $U_n$
équivaut à un morphisme $f : V \to W$
sur le morphisme $f_\varphi : U_n \to U_m$, c'est-à-dire
la même chose qu'un morphisme dans la catégorie $(\mathcal{C}/U)_{total}$
défini ci-dessus. L'égalité des ensembles de recouvrements résulte
immédiatement de la définition.

\medskip\noindent
Nous omettons la preuve de (3).
\end{proof}







\section{Hyperrecouvrement par un objet simplicial du site : modules}
\label{spaces-simplicial-section-hypercovering-modules}

\noindent
Soient $\mathcal{C}$ un site admettant des produits fibrés et $X \in \Ob(\mathcal{C})$.
Soit $\mathcal{O}_\mathcal{C}$ un faisceau d'anneaux sur $\mathcal{C}$.
Soit $U \to X$ un hyperrecouvrement de $X$ dans $\mathcal{C}$, au sens de la
section \ref{spaces-simplicial-section-hypercovering}. Nous étudions ici l'augmentation
suivante :
$$
a :
(\Sh((\mathcal{C}/U)_{total}), \mathcal{O})
\longrightarrow
(\Sh(\mathcal{C}/X), \mathcal{O}_X)
$$
que l'on obtient en considérant $U$ comme un objet semi-représentable simplicial
de $\mathcal{C}/X$, dont le terme de degré $n$ est le singleton
$\{U_n/X\}$, puis en appliquant les constructions de la
remarque \ref{spaces-simplicial-remark-augmentation-ringed-over-object}.
Ainsi tous les résultats de cette section sont des conséquences immédiates
des résultats correspondants dans la section
\ref{spaces-simplicial-section-cohomological-descent-hypercoverings-X-modules}.

\begin{lemma}
\label{spaces-simplicial-lemma-hypercovering-X-simple-descent-modules}
Soient $\mathcal{C}$ un site admettant des produits fibrés et $X \in \Ob(\mathcal{C})$.
Soit $\mathcal{O}_\mathcal{C}$ un faisceau d'anneaux.
Soit $U$ un hyperrecouvrement de $X$ dans $\mathcal{C}$. Avec les notations ci-dessus,
$$
a^* : \textit{Mod}(\mathcal{O}_X) \to \textit{Mod}(\mathcal{O})
$$
est pleinement fidèle, d'image essentielle les $\mathcal{O}$-modules cartésiens.
Le foncteur $a_*$ est un quasi-inverse.
\end{lemma}

\begin{proof}
Il s'agit d'un cas particulier du lemme
\ref{spaces-simplicial-lemma-hypercovering-X-descent-modules}.
\end{proof}

\begin{lemma}
\label{spaces-simplicial-lemma-hypercovering-X-simple-descent-bounded-modules}
Soient $\mathcal{C}$ un site admettant des produits fibrés et $X \in \Ob(\mathcal{C})$.
Soit $\mathcal{O}_\mathcal{C}$ un faisceau d'anneaux.
Soit $U$ un hyperrecouvrement de $X$ dans $\mathcal{C}$. Pour
$E \in D(\mathcal{O}_X)$, le morphisme
$$
E \longrightarrow Ra_*La^*E
$$
est un isomorphisme.
\end{lemma}

\begin{proof}
Il s'agit d'un cas particulier du lemme
\ref{spaces-simplicial-lemma-hypercovering-X-descent-bounded-modules}.
\end{proof}

\begin{lemma}
\label{spaces-simplicial-lemma-compare-cohomology-hypercovering-X-simple-modules}
Soient $\mathcal{C}$ un site admettant des produits fibrés et $X \in \Ob(\mathcal{C})$.
Soit $\mathcal{O}_\mathcal{C}$ un faisceau d'anneaux.
Soit $U$ un hyperrecouvrement de $X$ dans $\mathcal{C}$.
Nous avons alors un isomorphisme canonique
$$
R\Gamma(X, E) = R\Gamma((\mathcal{C}/U)_{total}, La^*E)
$$
pour $E \in D(\mathcal{O}_\mathcal{C})$.
\end{lemma}

\begin{proof}
Il s'agit d'un cas particulier du lemme
\ref{spaces-simplicial-lemma-compare-cohomology-hypercovering-X-modules}.
\end{proof}

\begin{lemma}
\label{spaces-simplicial-lemma-hypercovering-X-simple-equivalence-bounded-modules}
Soient $\mathcal{C}$ un site admettant des produits fibrés et $X \in \Ob(\mathcal{C})$.
Soit $\mathcal{O}_\mathcal{C}$ un faisceau d'anneaux.
Soit $U$ un hyperrecouvrement de $X$ dans $\mathcal{C}$.
Soit $\mathcal{A} \subset \textit{Mod}(\mathcal{O})$
la sous-catégorie de Serre faible des $\mathcal{O}$-modules cartésiens.
Alors le foncteur $La^*$ définit une équivalence
$$
D^+(\mathcal{O}_X) \longrightarrow D_\mathcal{A}^+(\mathcal{O})
$$
de quasi-inverse $Ra_*$.
\end{lemma}

\begin{proof}
Il s'agit d'un cas particulier du lemme
\ref{spaces-simplicial-lemma-hypercovering-X-equivalence-bounded-modules}.
\end{proof}







\section{Descente cohomologique non bornée pour les hyperrecouvrements}
\label{spaces-simplicial-section-unbounded-cohomological-descent}

\noindent
Dans cette section, nous étudions la descente cohomologique non bornée.
Les résultats eux-mêmes seront des conséquences immédiates de
nos résultats sur la descente cohomologique bornée dans les
sections précédentes et des lemmes de Cohomologie sur les sites
\ref{sites-cohomology-lemma-equivalence-unbounded-one} et/ou
\ref{sites-cohomology-lemma-equivalence-unbounded-two} ; le vrai travail réside
dans la mise en place de la notation et le choix des hypothèses appropriées.
Cette discussion est motivée par celle de \cite{six-I},
bien que les détails diffèrent sensiblement.

\medskip\noindent
Soit $(\mathcal{C}, \mathcal{O}_\mathcal{C})$ un site annelé.
Supposons donnée, pour tout objet $U$ de $\mathcal{C}$,
une sous-catégorie de Serre faible $\mathcal{A}_U \subset \textit{Mod}(\mathcal{O}_U)$
satisfaisant aux propriétés suivantes :
\begin{enumerate}
\item
\label{spaces-simplicial-item-restriction}
pour tout morphisme $U \to V$ de $\mathcal{C}$, le foncteur de restriction
$\textit{Mod}(\mathcal{O}_V) \to \textit{Mod}(\mathcal{O}_U)$
envoie $\mathcal{A}_V$ dans $\mathcal{A}_U$,
\item
\label{spaces-simplicial-item-local}
pour tout recouvrement $\{U_i \to U\}_{i \in I}$ de $\mathcal{C}$,
un objet $\mathcal{F}$ de $\textit{Mod}(\mathcal{O}_U)$
appartient à $\mathcal{A}_U$ si et seulement si la restriction de
$\mathcal{F}$ à $\mathcal{C}/U_i$ appartient à $\mathcal{A}_{U_i}$
pour tout $i \in I$.
\item
\label{spaces-simplicial-item-bounded-dimension}
il existe un sous-ensemble $\mathcal{B} \subset \Ob(\mathcal{C})$
tel que
\begin{enumerate}
\item tout objet de $\mathcal{C}$ admet un recouvrement dont les membres
appartiennent à $\mathcal{B}$, et
\item pour tout $V \in \mathcal{B}$, il existe un entier $d_V$
et un système cofinal $\text{Cov}_V$ de recouvrements de $V$ tel
que
$$
H^p(V_i, \mathcal{F}) = 0 \text{ pour }
\{V_i \to V\} \in \text{Cov}_V,\ p > d_V, \text{ et }
\mathcal{F} \in \Ob(\mathcal{A}_V)
$$
\end{enumerate}
\end{enumerate}
Remarquons que cette propriété doit être satisfaite pour $\mathcal{F}$ dans
$\mathcal{A}_V$ et pas seulement pour les objets « globaux »
(elle est donc plus forte que la condition de
Cohomologie sur les sites, situation \ref{sites-cohomology-situation-olsson-laszlo}).
Dans cette situation, il existe une sous-catégorie de Serre faible
$\mathcal{A} \subset \textit{Mod}(\mathcal{O}_\mathcal{C})$
constituée des objets dont la restriction à $\mathcal{C}/U$
appartient à $\mathcal{A}_U$ pour tout $U \in \Ob(\mathcal{C})$.
De plus, il existe des catégories dérivées
$D_\mathcal{A}(\mathcal{O}_\mathcal{C})$ et
$D_{\mathcal{A}_U}(\mathcal{O}_U)$, et les foncteurs de restriction
envoient ces catégories les unes dans les autres.

\begin{example}
\label{spaces-simplicial-example-quasi-coherent-spaces-etale}
Soit $S$ un schéma et $X$ un espace algébrique sur $S$.
Soit $\mathcal{C} = X_{spaces, \etale}$ le site \'etale
sur la catégorie des espaces algébriques étales sur $X$; voir
Propriétés des espaces, définition
\ref{spaces-properties-definition-spaces-etale-site}.
Notons $\mathcal{O}_\mathcal{C}$ le faisceau structural, c'est-à-dire le faisceau
donné par la règle $U \mapsto \Gamma(U, \mathcal{O}_U)$.
Notons $\mathcal{A}_U$ la catégorie des $\mathcal{O}_U$-modules quasi-cohérents.
Soit $\mathcal{B} = \Ob(\mathcal{C})$ et, pour $V \in \mathcal{B}$,
posons $d_V = 0$ et notons $\text{Cov}_V$ l'ensemble des
recouvrements $\{V_i \to V\}$ tels que $V_i$ soit affine pour tout $i$.
Alors les hypothèses (1), (2) et (3) sont satisfaites.
Voir Propriétés des espaces, lemmes
\ref{spaces-properties-lemma-pullback-quasi-coherent} et
\ref{spaces-properties-lemma-properties-quasi-coherent}
pour les propriétés (1) et (2); l'annulation en (3) découle de
Cohomologie des schémas, lemme
\ref{coherent-lemma-quasi-coherent-affine-cohomology-zero}
et de la discussion de Cohomologie des espaces, section
\ref{spaces-cohomology-section-higher-direct-image}.
\end{example}

\begin{example}
\label{spaces-simplicial-example-etale}
Soit $S$ l'un des types de schémas suivants :
\begin{enumerate}
\item le spectre d'un corps fini,
\item le spectre d'un corps séparablement clos,
\item le spectre d'un anneau local noethérien strictement hensélien,
\item le spectre d'un anneau local noethérien hensélien
à corps résiduel fini,
\item ajouter d'autres exemples ici.
\end{enumerate}
Soit $\Lambda$ un anneau fini dont l'ordre est inversible sur $S$.
Soit $\mathcal{C} \subset (\Sch/S)_\etale$
la sous-catégorie pleine formée des schémas localement de type fini
sur $S$, munis de la topologie étale.
Soit $\mathcal{O}_\mathcal{C} = \underline{\Lambda}$ le
faisceau constant. Posons $\mathcal{A}_U = \textit{Mod}(\mathcal{O}_U)$;
autrement dit, on considère tous les faisceaux étales de $\Lambda$-modules.
Soit $\mathcal{B} \subset \Ob(\mathcal{C})$
l'ensemble des objets quasi-compacts. Pour $V \in \mathcal{B}$, posons
$$
d_V = 1 + 2\dim(S) +
\sup\nolimits_{v \in V}(\text{trdeg}_{\kappa(s)}(\kappa(v)) +
2 \dim \mathcal{O}_{V, v})
$$
et notons $\text{Cov}_V$ l'ensemble des recouvrements étales $\{V_i \to V\}$
tels que $V_i$ soit quasi-compact pour tout $i$.
Le choix de la borne $d_V$ provient du théorème de Gabber
sur la dimension cohomologique. Pour vérifier que la condition (3)
est satisfaite avec ce choix, on utilise
\cite[Exposé VIII-A, corollaire 1.2 et lemme 2.2]{Traveaux}
ainsi que des arguments élémentaires sur les dimensions cohomologiques des corps.
Nous ajoutons $1$ dans la formule parce que notre liste contient des cas où $S$
peut avoir un corps résiduel fini.
Nous reviendrons sur cet exemple plus tard (insérer la référence future).
\end{example}

\noindent
Soit $(\mathcal{C}, \mathcal{O}_\mathcal{C})$ un site annelé.
Supposons données des sous-catégories de Serre faibles
$\mathcal{A}_U \subset \textit{Mod}(\mathcal{O}_U)$
satisfassent à la condition (\ref{spaces-simplicial-item-restriction}).
Alors
\begin{enumerate}
\item pour tout objet semi-représentable $K = \{U_i\}_{i \in I}$,
nous obtenons une sous-catégorie de Serre faible
$\mathcal{A}_K \subset \textit{Mod}(\mathcal{O}_K)$
en prenant $\prod \mathcal{A}_{U_i} \subset 
\prod \textit{Mod}(\mathcal{O}_{U_i}) = \textit{Mod}(\mathcal{O}_K)$, et
\item étant donné un morphisme d'objets semi-représentables
$f : K \to L$, le foncteur image inverse
$f^* : \textit{Mod}(\mathcal{O}_L) \to \textit{Mod}(\mathcal{O}_L)$
envoie $\mathcal{A}_L$ dans $\mathcal{A}_K$.
\end{enumerate}
Voir la remarque \ref{spaces-simplicial-remark-semi-representable-ringed} pour les notations et les
explications. En particulier, pour un objet semi-représentable simplicial $K$,
l'expression « un objet $\mathcal{F}$ de
$\textit{Mod}(\mathcal{O})$, au sens de la remarque \ref{spaces-simplicial-remark-augmentation-ringed},
dont les restrictions $\mathcal{F}_n$ appartiennent à
$\mathcal{A}_{K_n}$ pour tout $n$ » est sans ambiguïté.

\begin{lemma}
\label{spaces-simplicial-lemma-hypercovering-equivalence-modules}
Soit $(\mathcal{C}, \mathcal{O}_\mathcal{C})$ un site annelé.
Supposons données des sous-catégories de Serre faibles
$\mathcal{A}_U \subset \textit{Mod}(\mathcal{O}_U)$
satisfaisant aux conditions (\ref{spaces-simplicial-item-restriction}),
(\ref{spaces-simplicial-item-local}) et (\ref{spaces-simplicial-item-bounded-dimension}) ci-dessus.
Supposons que $\mathcal{C}$ admette des égalisateurs et des produits fibrés, et
soit $K$ un hyperrecouvrement.
Considérons $((\mathcal{C}/K)_{total}, \mathcal{O})$ comme dans
la remarque \ref{spaces-simplicial-remark-augmentation-ringed}.
Soit $\mathcal{A}_{total} \subset \textit{Mod}(\mathcal{O})$
la sous-catégorie de Serre faible formée des $\mathcal{O}$-modules cartésiens
$\mathcal{F}$ tels que la restriction $\mathcal{F}_n$ appartienne à
$\mathcal{A}_{K_n}$ pour tout $n$ (au sens défini ci-dessus).
Alors le foncteur $La^*$ définit une équivalence
$$
D_\mathcal{A}(\mathcal{O}_\mathcal{C})
\longrightarrow
D_{\mathcal{A}_{total}}(\mathcal{O})
$$
de quasi-inverse $Ra_*$.
\end{lemma}

\begin{proof}
Les $\mathcal{O}$-modules cartésiens forment une sous-catégorie de Serre faible d'après le
lemme \ref{spaces-simplicial-lemma-Serre-subcat-cartesian-modules}
(les hypothèses requises découlent de la discussion de la
remarque \ref{spaces-simplicial-remark-augmentation-ringed}).
Puisque les foncteurs de restriction
$g_n^* : \textit{Mod}(\mathcal{O}) \to \textit{Mod}(\mathcal{O}_n)$
sont exacts, il s'ensuit que $\mathcal{A}_{total}$ est une
sous-catégorie de Serre faible.

\medskip\noindent
Montrons que $a^* : \mathcal{A} \to \mathcal{A}_{total}$
est une équivalence de catégories, de quasi-inverse $La_*$.
On sait déjà que $La_*a^*\mathcal{F} = \mathcal{F}$ par la
version bornée
(lemme \ref{spaces-simplicial-lemma-hypercovering-equivalence-bounded-modules}).
Il est clair que $a^*\mathcal{F}$ appartient à $\mathcal{A}_{total}$
pour tout $\mathcal{F}$ dans $\mathcal{A}$. Réciproquement, soit
$\mathcal{G} \in \mathcal{A}_{total}$. Comme $\mathcal{G}$
est cartésien, on a $\mathcal{G} = a^*\mathcal{F}$
pour un certain $\mathcal{O}_\mathcal{C}$-module $\mathcal{F}$, d'après le
lemme \ref{spaces-simplicial-lemma-hypercovering-descent-modules}.
Nous voulons montrer que $\mathcal{F}$ appartient à $\mathcal{A}$.
Fixons $U \in \Ob(\mathcal{C})$. Il faut montrer que la restriction
de $\mathcal{F}$ à $\mathcal{C}/U$ appartient à $\mathcal{A}_U$.
Comme d'habitude, écrivons $K_0 = \{U_{0, i}\}_{i \in I_0}$.
Puisque $K$ est un hyperrecouvrement, le morphisme $\coprod_{i \in I_0} h_{U_{0, i}} \to *$
devient surjectif après passage au faisceau associé. Il existe donc
un recouvrement $\{U_j \to U\}_{j \in J}$ et une application $\tau : J \to I_0$
et pour chaque $j \in J$ un morphisme $\varphi_j : U_j \to U_{0, \tau(j)}$.
Puisque $\mathcal{G}_0 = a_0^*\mathcal{F}$, on voit
que la restriction de $\mathcal{F}$ à $\mathcal{C}/U_j$
est égale à la restriction de la composante $\tau(j)$ de
$\mathcal{G}_0$ à $\mathcal{C}/U_j$ via le morphisme
$\varphi_j : U_j \to U_{0, \tau(i)}$. Par conséquent,
la condition (\ref{spaces-simplicial-item-restriction}) donne que $\mathcal{F}|_{\mathcal{C}/U_j}$
appartient à $\mathcal{A}_{U_j}$; la condition
(\ref{spaces-simplicial-item-local}) donne alors que $\mathcal{F}|_{\mathcal{C}/U}$
appartient à $\mathcal{A}_U$.

\medskip\noindent
En particulier, l'énoncé du lemme est bien défini.
Le lemme découle alors de Cohomologie sur les sites,
lemme \ref{sites-cohomology-lemma-equivalence-unbounded-one}.
L'hypothèse (1) est claire (voir la remarque \ref{spaces-simplicial-remark-augmentation-ringed}).
Les hypothèses (2) et (3) ont été démontrées au paragraphe précédent.
L'hypothèse (4) découle immédiatement de (\ref{spaces-simplicial-item-bounded-dimension}).
Pour l'hypothèse (5), soit $\mathcal{B}_{total}$ l'ensemble des objets
$U/U_{n, i}$ du site $(\mathcal{C}/K)_{total}$
tels que $U \in \mathcal{B}$, où $\mathcal{B}$ est comme dans
(\ref{spaces-simplicial-item-bounded-dimension}). Nous utilisons ici la description de
la description du site $(\mathcal{C}/K)_{total}$ donnée dans la section
\ref{spaces-simplicial-section-simplicial-semi-representable}.
De plus, posons $\text{Cov}_{U/U_{n, i}}$ égal à $\text{Cov}_U$
et $d_{U/U_{n, i}}$ égal à $d_U$, où $\text{Cov}_U$ et $d_U$
sont fournis par (\ref{spaces-simplicial-item-bounded-dimension}).
Avec ces choix, la condition (5) est satisfaite.
Cela découle immédiatement du lemme
\ref{spaces-simplicial-lemma-sanity-check-simplicial-semi-representable}
et du fait que $\mathcal{F} \in \mathcal{A}_{total}$
implique $\mathcal{F}_n \in \mathcal{A}_{K_n}$ et donc
$\mathcal{F}_{n, i} \in \mathcal{A}_{U_{n, i}}$.
(Le lecteur préoccupé par la différence entre la cohomologie
des faisceaux abéliens et la cohomologie
des faisceaux de modules peut consulter Cohomologie sur les sites, lemme
\ref{sites-cohomology-lemma-cohomology-modules-abelian-agree}.)
\end{proof}










\section{Recollement de complexes}
\label{spaces-simplicial-section-glueing-complexes}

\noindent
Cette section prolonge
Cohomologie, section \ref{cohomology-section-glueing-complexes}.
Le but est de prouver une légère généralisation de \cite[théorème 3.2.4]{BBD}.
Notre méthode différera légèrement, en ce qu'elle utilise
les résultats des sections \ref{spaces-simplicial-section-glueing} et
\ref{spaces-simplicial-section-glueing-modules}. Nous redémontrerons également la
version non bornée en suivant la démonstration de \cite{six-I}.

\medskip\noindent
Conseil au lecteur : nous suggérons de lire d'abord l'énoncé du
lemme \ref{spaces-simplicial-lemma-bbd-glueing-cech}, ainsi que sa seconde démonstration.

\medskip\noindent
Voici la situation qui nous intéresse.

\begin{situation}
\label{spaces-simplicial-situation-locally-given}
Soit $(\mathcal{C}, \mathcal{O}_\mathcal{C})$ un site annelé. On nous donne
\begin{enumerate}
\item une catégorie $\mathcal{B}$ et un foncteur
$u : \mathcal{B} \to \mathcal{C}$,
\item un objet $E_U$ dans $D(\mathcal{O}_{u(U)})$ pour $U \in \Ob(\mathcal{B})$,
\item un isomorphisme $\rho_a : E_U|_{\mathcal{C}/u(V)} \to E_V$ dans
$D(\mathcal{O}_{u(V)})$ pour $a : V \to U$ dans $\mathcal{B}$
\end{enumerate}
avec la compatibilité suivante : pour toutes flèches composables
$b : W \to V$ et $a : V \to U$ de $\mathcal{B}$, on a
$\rho_{a \circ b} = \rho_b \circ \rho_a|_{\mathcal{C}/u(W)}$.
\end{situation}

\noindent
Nous ne pourrons rien démontrer dans ce cadre sans faire davantage
d'hypothèses. Un cas intéressant est celui où $\mathcal{B}$ est une sous-catégorie pleine
telle que chaque objet de $\mathcal{C}$ a un recouvrement
dont les membres sont des objets de $\mathcal{B}$ (c'est le cas étudié
dans \cite{BBD}). Il importe ici d'autoriser aussi les situations où
ce n'est pas le cas; la principale variante est celle où l'on dispose d'un morphisme
de sites $f : \mathcal{C} \to \mathcal{D}$ et $\mathcal{B}$
est une sous-catégorie pleine de $\mathcal{D}$ telle que chaque objet de
$\mathcal{D}$ admet un recouvrement dont les membres sont des objets de $\mathcal{B}$.

\medskip\noindent
Dans la situation \ref{spaces-simplicial-situation-locally-given}, une {\it solution}
sera une paire $(E, \rho_U)$ où $E$ est un objet de
$D(\mathcal{O}_\mathcal{C})$
et où les morphismes $\rho_U : E|_{\mathcal{C}/u(U)} \to E_U$
pour $U \in \Ob(\mathcal{B})$
sont des isomorphismes tels que
l'on a $\rho_a \circ \rho_U|_{\mathcal{C}/u(V)} = \rho_V$
pour $a : V \to U$ dans $\mathcal{B}$.

\begin{lemma}
\label{spaces-simplicial-lemma-prepare-bbd-glueing}
Dans la situation \ref{spaces-simplicial-situation-locally-given}.
Supposons que les auto-Ext négatifs de $E_U$ dans $D(\mathcal{O}_{u(U)})$ soient nuls.
Soit $L$ un objet simplicial de $\text{SR}(\mathcal{B})$.
Considérons l'objet simplicial $K = u(L)$ de $\text{SR}(\mathcal{C})$
et considérons $((\mathcal{C}/K)_{total}, \mathcal{O})$ comme dans la
remarque \ref{spaces-simplicial-remark-augmentation-ringed}.
Il existe un objet cartésien $E$ de $D(\mathcal{O})$
tel que, si l'on écrit $L_n = \{U_{n, i}\}_{i \in I_n}$,
l'image de $E$ dans $D(\mathcal{O}_{\mathcal{C}/u(U_{n, i})})$
s'identifie de manière compatible à $E_{U_{n, i}}$ (voir la preuve).
De plus, $E$ est unique à isomorphisme unique près.
\end{lemma}

\begin{proof}
Rappelons que
$\Sh(\mathcal{C}/K_n) = \prod_{i \in I_n} \Sh(\mathcal{C}/u(U_{n, i}))$
et de même pour les catégories de modules. Cette décomposition en produit
se transporte également aux catégories dérivées de faisceaux de modules.
De plus, cette décomposition en produit est compatible avec
les morphismes de l'objet semi-représentable simplicial $K$.
Voir la section \ref{spaces-simplicial-section-semi-representable}.
Nous pouvons donc définir $E_n = \prod_{i \in I_n} E_{U_{n, i}}$
(produit « formel ») dans $D(\mathcal{O}_n)$.
En prenant les produits (formels) des morphismes $\rho_a$ de
la situation \ref{spaces-simplicial-situation-locally-given},
nous obtenons des isomorphismes $E_\varphi : f_\varphi^*E_n \to E_m$.
L'hypothèse sur les composés des morphismes $\rho_a$
implique immédiatement que $(E_n, E_\varphi)$
définit un système simplicial de la catégorie dérivée de modules
comme dans la définition \ref{spaces-simplicial-definition-cartesian-derived-modules}.
L'annulation des auto-Ext négatifs supposée dans le lemme entraîne que
$\Hom(E_n[t], E_n) = 0$ pour $n \geq 0$ et $t > 0$.
Par le
lemme \ref{spaces-simplicial-lemma-cartesian-module-derived-from-simplicial},
on obtient $E$.
L'unicité jusqu'à l'isomorphisme unique découle des lemmes
\ref{spaces-simplicial-lemma-nullity-cartesian-modules-derived} et
\ref{spaces-simplicial-lemma-hom-cartesian-modules-derived}.
\end{proof}

\begin{lemma}[Lemme de recollement BBD]
\label{spaces-simplicial-lemma-bbd-glueing}
Dans la situation \ref{spaces-simplicial-situation-locally-given}. Supposons que
\begin{enumerate}
\item $\mathcal{C}$ admet des égalisateurs et des produits fibrés,
\item il existe un morphisme de sites $f : \mathcal{C} \to \mathcal{D}$
donné par un foncteur continu $u : \mathcal{D} \to \mathcal{C}$
tel que
\begin{enumerate}
\item $\mathcal{D}$ admet des égalisateurs et des produits fibrés, et $u$
commute à leur formation,
\item $\mathcal{B}$ est une sous-catégorie pleine de $\mathcal{D}$
et $u : \mathcal{B} \to \mathcal{C}$ est la restriction de $u$,
\item tout objet de $\mathcal{D}$ admet un recouvrement dont les membres
sont des objets de $\mathcal{B}$,
\end{enumerate}
\item pour tout $U$ dans $\mathcal{B}$, tous les auto-Ext négatifs de $E_U$
dans $D(\mathcal{O}_{u(U)})$ sont nuls, et
\item il existe un $t \in \mathbf{Z}$ tel que $H^i(E_U) = 0$ pour $i < t$
et tout $U \in \Ob(\mathcal{B})$.
\end{enumerate}
Il existe alors une solution unique à isomorphisme unique près.
\end{lemma}

\begin{proof}
D'après Hyperrecouvrements, lemme \ref{hypercovering-lemma-hypercovering-site},
il existe un hyperrecouvrement $L$ pour le site $\mathcal{D}$ tel que
$L_n = \{U_{n, i}\}_{i \in I_n}$ avec $U_{i, n} \in \Ob(\mathcal{B})$.
Posons $K = u(L)$. Le lemme \ref{spaces-simplicial-lemma-prepare-bbd-glueing} donne
un objet cartésien $E$ de $D(\mathcal{O})$ sur le site
$(\mathcal{C}/K)_{total}$, dont la restriction est $E_{U_{n, i}}$ sur
$\mathcal{C}/u(U_{n, i})$, de façon compatible.
L'hypothèse sur $t$ implique que $E \in D^+(\mathcal{O})$.
D'après Hyperrecouvrements, lemme \ref{hypercovering-lemma-hypercovering-morphism-sites},
$K$ est lui aussi un hyperrecouvrement.
Le lemme \ref{spaces-simplicial-lemma-hypercovering-equivalence-bounded-modules}
donne $E = a^*F$ pour un certain $F$ dans $D^+(\mathcal{O}_\mathcal{C})$.

\medskip\noindent
Pour prouver que $F$ est une solution, nous utiliserons la construction de
$L_0$ et $L_1$ donnée dans la preuve de
Hyperrecouvrements, lemme \ref{hypercovering-lemma-hypercovering-site}.
(C'est un peu inélégant, mais il ne semble pas y avoir de moyen
tout à fait direct de l'éviter.)

\medskip\noindent
Plus précisément, $I_0 = \Ob(\mathcal{B})$ et donc
$L_0 = \{U\}_{U \in \Ob(\mathcal{B})}$.
La restriction de l'isomorphisme $a^*F \to E$ aux composantes
$\mathcal{C}/u(U)$ de $\mathcal{C}/K_0$ définit les isomorphismes
$\rho_U : F|_{\mathcal{C}/u(U)} \to E_U$ pour $U \in \Ob(\mathcal{B})$
par construction de $E$.

\medskip\noindent
Pour vérifier que les $\rho_U$ satisfont à la condition de compatibilité
avec les morphismes $\rho_a$ de la situation \ref{spaces-simplicial-situation-locally-given},
on utilise le fait que $I_1$ contient l'ensemble
$$
\Omega =
\{(U, V, W, a, b) \mid U, V, W \in \mathcal{B}, a : U \to V, b : U \to W\}
$$
et que, pour $i = (U, V, W, a, b)$ dans $\Omega$, on a
$U_{1, i} = U$. De plus, les composantes $f_{\delta^1_0, i}$ et
$f_{\delta^1_1, i}$ des deux morphismes $K_1 \to K_0$ sont les morphismes
$$
a : U \to V \quad\text{et}\quad b : U \to V
$$
La compatibilité du
lemme \ref{spaces-simplicial-lemma-prepare-bbd-glueing} donne donc
$$
\rho_a \circ \rho_V|_{\mathcal{C}/u(U)} = \rho_U
\quad\text{et}\quad
\rho_b \circ \rho_W|_{\mathcal{C}/u(U)} = \rho_U
$$
En prenant par exemple $i = (U, V, U, a, \text{id}_U) \in \Omega$, on obtient
la compatibilité voulue. L'unicité de $F$ découle
de celle de $E$ dans le lemme précédent (nous omettons ce détail).
\end{proof}

\begin{lemma}[Lemme de recollement BBD non borné]
\label{spaces-simplicial-lemma-bbd-unbounded-glueing}
Dans la situation \ref{spaces-simplicial-situation-locally-given}. Supposons que
\begin{enumerate}
\item $\mathcal{C}$ admet des égalisateurs et des produits fibrés,
\item il existe un morphisme de sites $f : \mathcal{C} \to \mathcal{D}$
donné par un foncteur continu $u : \mathcal{D} \to \mathcal{C}$
tel que
\begin{enumerate}
\item $\mathcal{D}$ admet des égalisateurs et des produits fibrés, et $u$
commute à leur formation,
\item $\mathcal{B}$ est une sous-catégorie pleine de $\mathcal{D}$
et $u : \mathcal{B} \to \mathcal{C}$ est la restriction de $u$,
\item chaque objet de $\mathcal{D}$ possède un recouvrement dont les membres
sont des objets de $\mathcal{B}$,
\end{enumerate}
\item tous les auto-Ext négatifs de $E_U$ dans $D(\mathcal{O}_{u(U)})$ sont nuls, et
\item il existe des sous-catégories de Serre faibles
$\mathcal{A}_U \subset \textit{Mod}(\mathcal{O}_U)$ pour tout
$U \in \Ob(\mathcal{C})$, satisfaisant aux conditions (\ref{spaces-simplicial-item-restriction}),
(\ref{spaces-simplicial-item-local}) et (\ref{spaces-simplicial-item-bounded-dimension}),
\item $E_U \in D_{\mathcal{A}_U}(\mathcal{O}_U)$.
\end{enumerate}
Il existe alors une solution unique à isomorphisme unique près.
\end{lemma}

\begin{proof}
La démonstration est {\bf exactement} la même que celle du
lemme \ref{spaces-simplicial-lemma-bbd-glueing}. La seule différence est que
$E$ est un objet de $D_{\mathcal{A}_{total}}(\mathcal{O})$
et l'on utilise donc le lemme \ref{spaces-simplicial-lemma-hypercovering-equivalence-modules}
pour obtenir $F$ avec $E = a^*F$
au lieu du lemme \ref{spaces-simplicial-lemma-hypercovering-equivalence-bounded-modules}.
\end{proof}

\noindent
Voici un exemple d'application de la théorie générale ci-dessus.

\begin{lemma}
\label{spaces-simplicial-lemma-bbd-glueing-cech}
\begin{reference}
Courriel de Martin Olsson daté du 9 septembre 2021.
\end{reference}
Soit $(\mathcal{C}, \mathcal{O}_\mathcal{C})$ un site annelé.
Supposons que $\mathcal{C}$ admette des produits fibrés.
Soit $\{U_i \to X\}_{i \in I}$ un recouvrement dans $\mathcal{C}$.
Pour $i \in I$, soit $E_i$ un objet de $D(\mathcal{O}_{U_i})$, et
pour $i, j \in I$, soit
$$
\rho_{ij} : E_i|_{\mathcal{C}/U_{ij}} \longrightarrow E_j|_{\mathcal{C}/U_{ij}}
$$
est un isomorphisme dans $D(\mathcal{O}_{U_{ij}})$, où
$U_{ij} = U_i \times_X U_j$. Supposons que
\begin{enumerate}
\item les $\rho_{ij}$ satisfassent à la condition de cocycle sur
$U_i \times_X U_j \times_X U_k$ pour tout $i, j, k \in I$,
\item $\SheafExt^p_{\mathcal{O}_{U_i}}(E_i, E_i) = 0$
pour tout $p < 0$ et tout $i \in I$, et
\item il existe un $t \in \mathbf{Z}$ tel que
$H^p(E_i) = 0$ pour $p < t$ et tout $i \in I$.
\end{enumerate}
Alors il existe un unique couple $(E, \rho_i)$, où $E$ est un objet de
$D(\mathcal{O}_X)$ et les $\rho_i : E|_{U_i} \to E_i$ sont des isomorphismes dans
$D(\mathcal{O}_{U_i})$ compatibles avec les $\rho_{ij}$.
\end{lemma}

\begin{proof}[Première démonstration]
Dans cette preuve, nous déduisons le résultat du lemme général
\ref{spaces-simplicial-lemma-bbd-glueing}. Nous invitons le lecteur à consulter plutôt
la seconde démonstration.

\medskip\noindent
Nous pouvons remplacer $\mathcal{C}$ par $\mathcal{C}/X$.
On peut donc supposer que $X$ est l'objet final de $\mathcal{C}$
et que $\mathcal{C}$ admet toutes les limites finies.

\medskip\noindent
Soit $\mathcal{B}$ la sous-catégorie pleine de $\mathcal{C}$
constituée des objets $U \in \Ob(\mathcal{C})$ pour lesquels il existe
un $i(U) \in I$ et un morphisme $a_U : U \to U_{i(U)}$.
Notons $E_U = a_U^*E_{i(U)}$ dans $D(\mathcal{O}_U)$,
image inverse (restriction) de $E_i$ par $a_U$.
Pour un morphisme $a : U \to U'$ de $\mathcal{B}$,
on obtient un morphisme
$(a_{U'} \circ a, a_U) : U \to U_{i(U')} \times_X U_{i(U)} = U_{i(U')i(U)}$
et donc un isomorphisme
$$
\rho_a :
a^*E_{U'} = a^*a_{U'}^*E_{i(U')}
\xrightarrow{(a_{U'} \circ a, a_U)^*\rho_{i(U')i(U)}}
a_{U}^*E_{i(U)} = E_{U}
$$
dans $D(\mathcal{O}_U)$. Les données $\mathcal{B}, E_U, \rho_a$
sont celles de la situation \ref{spaces-simplicial-situation-locally-given}; les
isomorphismes $\rho_a$ satisfont la condition de cocycle
précisément grâce à la condition (1) de l'énoncé du lemme
(détails omis).

\medskip\noindent
Nous allons appliquer le lemme \ref{spaces-simplicial-lemma-bbd-glueing} avec $\mathcal{B}$,
$E_U$, $\rho_a$ comme ci-dessus et avec $\mathcal{D} = \mathcal{C}$ et
$f : \mathcal{C} \to \mathcal{D}$ le morphisme identité.
Les hypothèses (1) et (2)(a) du lemme \ref{spaces-simplicial-lemma-bbd-glueing} ont été vérifiées ci-dessus.
L'hypothèse (2)(b) du lemme \ref{spaces-simplicial-lemma-bbd-glueing} est claire.
L'hypothèse (2)(c) du lemme \ref{spaces-simplicial-lemma-bbd-glueing} est satisfaite, car
$\{U_i \to X\}$ est un recouvrement\footnote{En fait, il suffirait
que le morphisme $\coprod_{i \in I} h_{U_i} \to h_X$ devienne
surjectif après passage au faisceau associé; le lemme resterait vrai dans ce cas
avec la même preuve.}.
L'hypothèse (3) du lemme \ref{spaces-simplicial-lemma-bbd-glueing} est satisfaite,
car nous avons supposé l'annulation de tous les faisceaux d'Ext négatifs
de $E_i$; cela implique, pour tout objet $U$ au-dessus de $U_i$, que
les auto-Ext négatifs de $E_i|_U$ sont nuls.
L'hypothèse (4) du lemme \ref{spaces-simplicial-lemma-bbd-glueing} est satisfaite, car
nous avons supposé que les faisceaux de cohomologie de chaque $E_i$ sont nuls
en degrés strictement inférieurs à $t$.

\medskip\noindent
Nous obtenons une solution unique $(E, \rho_U)$.
En posant $\rho_i = \rho_{U_i}$, on obtient le lemme.
\end{proof}

\begin{proof}[Seconde démonstration]
Nous esquissons une preuve plus directe.
Notons $K$ l'hyperrecouvrement de Čech de $X$ associé au
recouvrement $\{U_i \to X\}_{i \in I}$; voir
Hyperrecouvrements, exemple \ref{hypercovering-example-cech}.
Ainsi par exemple $K_0 = \{U_i \to X\}_{i \in I}$ et
$K_1 = \{U_i \times_X U_j \to X\}_{i, j \in I}$ et ainsi de suite.
Considérons $((\mathcal{C}/K)_{total}, \mathcal{O})$, $a$ et $a_n$ comme dans la
remarque \ref{spaces-simplicial-remark-augmentation-ringed-over-object}.
Les objets $E_i$ déterminent un objet $M_0$ dans
$D(\mathcal{O}_0) = \prod D(\mathcal{O}_{U_i})$.
De même, les isomorphismes $\rho_{ij}$ déterminent un isomorphisme
$$
\alpha :
L(f_{\delta_1^1})^*M_0
\longrightarrow
L(f_{\delta_0^1})^*M_0
$$
dans $D(\mathcal{O}_1)$ satisfaisant la condition de cocycle.
Par le lemme \ref{spaces-simplicial-lemma-construct-cartesian-objects-derived-modules}
on obtient un système simplicial cartésien $(M_n)$ de la catégorie dérivée.
L'annulation supposée des faisceaux d'Ext négatifs montre que
les auto-Ext négatifs des objets $M_n$ sont nuls.
On obtient donc un objet cartésien $M$ de $D(\mathcal{O})$
dont le système simplicial associé est isomorphe à $(M_n)$ d'après le
lemme \ref{spaces-simplicial-lemma-cartesian-module-derived-from-simplicial}.
Puisque les faisceaux de cohomologie de $M$ sont nuls en degrés $< t$,
le lemme \ref{spaces-simplicial-lemma-hypercovering-X-equivalence-bounded-modules}
donne $M = La^*E$ pour un certain $E$ dans $D(\mathcal{O}_X)$.
L'isomorphisme $La^*E \to M$ restreint à $\mathcal{C}/U_i$
induit les isomorphismes $\rho_i$. Nous omettons la vérification
de la compatibilité avec les isomorphismes $\rho_{ij}$.
\end{proof}












\section{Hyperrecouvrements propres en topologie}
\label{spaces-simplicial-section-proper-hypercovering}

\noindent
Travaillons dans la catégorie $\textit{LC}$ des espaces topologiques séparés et localement
quasi compacts, munie des applications continues; voir
Cohomologie sur les sites, section \ref{sites-cohomology-section-cohomology-LC}.
Soient $X$ un objet de $\textit{LC}$ et $U$ un objet
simplicial de $\textit{LC}$, muni d'une augmentation
$$
a : U \to X
$$
On dit que $U$ est un {\it hyperrecouvrement propre} de $X$ si
\begin{enumerate}
\item $U_0 \to X$ est un morphisme propre surjectif,
\item $U_1 \to U_0 \times_X U_0$ est un morphisme propre surjectif,
\item $U_{n + 1} \to (\text{cosk}_n\text{sk}_n U)_{n + 1}$
est un morphisme propre surjectif pour $n \geq 1$.
\end{enumerate}
La catégorie $\textit{LC}$ admet toutes les limites finies; les cosquelettes
utilisés dans la formulation ci-dessus existent donc.
$$
\fbox{Principe : les hyperrecouvrements propres permettent de calculer la cohomologie.}
$$
Une idée essentielle de la démonstration consiste à munir
$\textit{LC}$ d'une topologie plus fine que la topologie usuelle, telle que :
(a) tout morphisme propre surjectif définisse un recouvrement;
(b) la cohomologie, pour cette topologie plus fine, des faisceaux usuels
coïncide avec leur cohomologie usuelle. Les propriétés
(a) et (b) sont valables pour la topologie qc, voir Cohomologie
sur les sites, section \ref{sites-cohomology-section-cohomology-LC}.
Une fois (a) et (b) établis, le principe découle
des résultats précédents de ce chapitre.

\begin{lemma}
\label{spaces-simplicial-lemma-compare-simplicial-objects}
Soit $U$ un objet simplicial de $\textit{LC}$ et soit $a : U \to X$
une augmentation. Il existe un diagramme commutatif
$$
\xymatrix{
\Sh((\textit{LC}_{qc}/U)_{total}) \ar[r]_-h \ar[d]_{a_{qc}} &
\Sh(U_{Zar}) \ar[d]^a \\
\Sh(\textit{LC}_{qc}/X) \ar[r]^-{h_{-1}} &
\Sh(X)
}
$$
où la flèche verticale gauche est définie dans la section
\ref{spaces-simplicial-section-hypercovering}
et la flèche verticale droite est définie dans le lemme
\ref{spaces-simplicial-lemma-augmentation}.
\end{lemma}

\begin{proof}
Écrivons $\Sh(X) = \Sh(X_{Zar})$. Observons que
$(\textit{LC}_{qc}/U)_{total}$ et $U_{Zar}$ relèvent tous deux
du cas (A) de la situation \ref{spaces-simplicial-situation-simplicial-site}.
Cela résulte immédiatement de la construction de
$U_{Zar}$ dans la section \ref{spaces-simplicial-section-simplicial-top}
et découle du lemme \ref{spaces-simplicial-lemma-sr-when-fibre-products}
pour $(\textit{LC}_{qc}/U)_{total}$.
Considérons alors les foncteurs
$U_{n, Zar} \to \textit{LC}_{qc}/U_n$, $U \mapsto U/U_n$
et
$X_{Zar} \to \textit{LC}_{qc}/X$, $U \mapsto U/X$.
Ces foncteurs définissent des morphismes de sites, comme on l'a vu
dans Cohomologie sur les sites, section \ref{sites-cohomology-section-cohomology-LC}.
On obtient ainsi un morphisme de sites simpliciaux compatible avec les augmentations
comme dans la remarque \ref{spaces-simplicial-remark-morphism-augmentation-simplicial-sites}
et on peut appliquer le lemme
\ref{spaces-simplicial-lemma-morphism-augmentation-simplicial-sites} pour conclure.
\end{proof}

\begin{lemma}
\label{spaces-simplicial-lemma-descent-sheaves-for-proper-hypercovering}
Soit $U$ un objet simplicial de $\textit{LC}$, muni d'une augmentation $a : U \to X$.
Si $a : U \to X$ définit un hyperrecouvrement propre de $X$,
alors
$$
a^{-1} : \Sh(X) \to \Sh(U_{Zar})
\quad\text{et}\quad
a^{-1} : \textit{Ab}(X) \to \textit{Ab}(U_{Zar})
$$
sont pleinement fidèles, d'image essentielle les faisceaux cartésiens, et
admettent $a_*$ pour quasi-inverse. Ici, $a : \Sh(U_{Zar}) \to \Sh(X)$ est comme dans le lemme
\ref{spaces-simplicial-lemma-augmentation}.
\end{lemma}

\begin{proof}
Nous démontrerons l'énoncé pour les faisceaux d'ensembles. Ce sera une
conséquence presque formelle des résultats déjà établis.
Considérons le diagramme du lemme \ref{spaces-simplicial-lemma-compare-simplicial-objects}.
D'après Cohomologie sur les sites, lemme \ref{sites-cohomology-lemma-describe-pullback-pi},
le foncteur $(h_{-1})^{-1}$ est pleinement fidèle, de quasi-inverse $h_{-1, *}$.
Il en va de même pour les composantes $h_n$ de $h$.
Par la description des foncteurs $h^{-1}$ et $h_*$ du lemme
\ref{spaces-simplicial-lemma-morphism-simplicial-sites}
on conclut que $h^{-1}$ est pleinement fidèle, de quasi-inverse $h_*$.
Observons que $U$ est un hyperrecouvrement de $X$ dans $\textit{LC}_{qc}$
(au sens de la section \ref{spaces-simplicial-section-hypercovering}), d'après
Cohomologie sur les sites, lemme
\ref{sites-cohomology-lemma-proper-surjective-is-qc-covering}.
Par le lemme \ref{spaces-simplicial-lemma-hypercovering-X-simple-descent-sheaves}
on voit que $a_{qc}^{-1}$ est pleinement fidèle, de quasi-inverse $a_{qc, *}$
et d'image essentielle les faisceaux cartésiens sur
$(\textit{LC}_{qc}/U)_{total}$.
Une chasse formelle dans le diagramme montre alors que
$a^{-1}$ est pleinement fidèle.

\medskip\noindent
Enfin, supposons que $\mathcal{G}$ soit un faisceau cartésien sur $U_{Zar}$.
Alors $h^{-1}\mathcal{G}$ est un faisceau cartésien sur $\textit{LC}_{qc}/U$.
Ainsi, $h^{-1}\mathcal{G} = a_{qc}^{-1}\mathcal{H}$ pour un certain faisceau
$\mathcal{H}$ sur $\textit{LC}_{qc}/X$.
Calculons :
\begin{align*}
(h_{-1})^{-1}(a_*\mathcal{G})
& =
(h_{-1})^{-1}
\text{Eq}(
\xymatrix{
a_{0, *}\mathcal{G}_0
\ar@<1ex>[r] \ar@<-1ex>[r] &
a_{1, *}\mathcal{G}_1
}
) \\
& =
\text{Eq}(
\xymatrix{
(h_{-1})^{-1}a_{0, *}\mathcal{G}_0
\ar@<1ex>[r] \ar@<-1ex>[r] &
(h_{-1})^{-1}a_{1, *}\mathcal{G}_1
}
) \\
& =
\text{Eq}(
\xymatrix{
a_{qc, 0, *}h_0^{-1}\mathcal{G}_0
\ar@<1ex>[r] \ar@<-1ex>[r] &
a_{qc, 1, *}h_1^{-1}\mathcal{G}_1
}
) \\
& =
\text{Eq}(
\xymatrix{
a_{qc, 0, *}a_{qc, 0}^{-1}\mathcal{H}
\ar@<1ex>[r] \ar@<-1ex>[r] &
a_{qc, 1, *}a_{qc, 1}^{-1}\mathcal{H}
}
) \\
& =
a_{qc, *}a_{qc}^{-1}\mathcal{H} \\
& =
\mathcal{H}
\end{align*}
La première égalité découle du lemme \ref{spaces-simplicial-lemma-augmentation};
la deuxième, de l'exactitude du foncteur $(h_{-1})^{-1}$;
la troisième, de
Cohomologie sur les sites, lemme \ref{sites-cohomology-lemma-push-pull-LC}
(ici nous utilisons que $a_0 : U_0 \to X$ et $a_1: U_1 \to X$ sont propres),
la quatrième, de $a_{qc}^{-1}\mathcal{H} = h^{-1}\mathcal{G}$;
la cinquième, du lemme \ref{spaces-simplicial-lemma-augmentation-site}; enfin la sixième
a été établie ci-dessus. Puisque $a_{qc}^{-1}\mathcal{H} = h^{-1}\mathcal{G}$,
on en déduit que $h^{-1}\mathcal{G} \cong h^{-1}a^{-1}a_*\mathcal{G}$,
ce qui termine la preuve par la pleine fidélité de $h^{-1}$.
\end{proof}

\begin{lemma}
\label{spaces-simplicial-lemma-cohomological-descent-for-proper-hypercovering}
Soit $U$ un objet simplicial de $\textit{LC}$ et soit $a : U \to X$
une augmentation. Si $a : U \to X$ est un hyperrecouvrement propre de $X$,
alors pour $K \in D^+(X)$
$$
K \to Ra_*(a^{-1}K)
$$
est un isomorphisme, où $a : \Sh(U_{Zar}) \to \Sh(X)$ est comme dans le lemme
\ref{spaces-simplicial-lemma-augmentation}.
\end{lemma}

\begin{proof}
Considérons le diagramme du lemme \ref{spaces-simplicial-lemma-compare-simplicial-objects}.
Observons que $Rh_{n, *}h_n^{-1}$ est le foncteur identité
sur $D^+(U_n)$ d'après Cohomologie sur les sites, lemme
\ref{sites-cohomology-lemma-cohomological-descent-LC}.
Par conséquent, $Rh_*h^{-1}$ est le foncteur identité sur
$D^+(U_{Zar})$ par
le lemme \ref{spaces-simplicial-lemma-direct-image-morphism-simplicial-sites}.
Nous avons
\begin{align*}
Ra_*(a^{-1}K)
& =
Ra_*Rh_*h^{-1}a^{-1}K \\
& =
Rh_{-1, *}Ra_{qc, *}a_{qc}^{-1}(h_{-1})^{-1}K \\
& =
Rh_{-1, *}(h_{-1})^{-1}K \\
& =
K
\end{align*}
La première égalité découle de la discussion ci-dessus; la deuxième,
de la commutativité du diagramme du lemme
\ref{spaces-simplicial-lemma-compare-simplicial-objects}; la troisième, du lemme
\ref{spaces-simplicial-lemma-hypercovering-X-simple-descent-bounded-abelian}
($U$ est un hyperrecouvrement de $X$ dans $\textit{LC}_{qc}$ d'après Cohomologie
sur les sites, lemme
\ref{sites-cohomology-lemma-proper-surjective-is-qc-covering}),
et la dernière, du lemme déjà utilisé de Cohomologie sur les sites
\ref{sites-cohomology-lemma-cohomological-descent-LC}.
\end{proof}

\begin{lemma}
\label{spaces-simplicial-lemma-compute-via-proper-hypercovering}
Soit $U$ un objet simplicial de $\textit{LC}$ et soit $a : U \to X$
une augmentation. Si $U$ est un hyperrecouvrement propre de $X$, alors
$$
R\Gamma(X, K) = R\Gamma(U_{Zar}, a^{-1}K)
$$
pour $K \in D^+(X)$, où $a : \Sh(U_{Zar}) \to \Sh(X)$
est comme dans le lemme \ref{spaces-simplicial-lemma-augmentation}.
\end{lemma}

\begin{proof}
Cela découle du lemme
\ref{spaces-simplicial-lemma-cohomological-descent-for-proper-hypercovering}
car $R\Gamma(U_{Zar}, -) = R\Gamma(X, -) \circ Ra_*$ par
Cohomologie sur les sites, remarque \ref{sites-cohomology-remark-before-Leray}.
\end{proof}

\begin{lemma}
\label{spaces-simplicial-lemma-proper-hypercovering-equivalence-bounded}
Soit $U$ un objet simplicial de $\textit{LC}$ et soit $a : U \to X$
une augmentation.
Soit $\mathcal{A} \subset \textit{Ab}(U_{Zar})$
la sous-catégorie de Serre faible des faisceaux abéliens cartésiens.
Si $U$ est un hyperrecouvrement propre de $X$, alors
le foncteur $a^{-1}$ définit une équivalence
$$
D^+(X) \longrightarrow D_\mathcal{A}^+(U_{Zar})
$$
de quasi-inverse $Ra_*$, où $a : \Sh(U_{Zar}) \to \Sh(X)$
est comme dans le lemme \ref{spaces-simplicial-lemma-augmentation}.
\end{lemma}

\begin{proof}
Observons que $\mathcal{A}$ est une sous-catégorie de Serre faible d'après le
lemme \ref{spaces-simplicial-lemma-Serre-subcat-cartesian-modules}.
L'équivalence est une conséquence formelle
des résultats obtenus jusqu'ici. On utilise les
lemmes \ref{spaces-simplicial-lemma-descent-sheaves-for-proper-hypercovering},
\ref{spaces-simplicial-lemma-cohomological-descent-for-proper-hypercovering}, ainsi que
Cohomologie sur les sites, lemme \ref{sites-cohomology-lemma-equivalence-bounded}.
\end{proof}

\begin{lemma}
\label{spaces-simplicial-lemma-spectral-sequence-proper-hypercovering}
Soit $U$ un objet simplicial de $\textit{LC}$ et soit
$a : U \to X$ une augmentation. Soit $\mathcal{F}$ un faisceau abélien
sur $X$. Notons $\mathcal{F}_n$ son image inverse sur $U_n$.
Si $U$ est un hyperrecouvrement propre de $X$, alors
il existe une suite spectrale canonique
$$
E_1^{p, q} = H^q(U_p, \mathcal{F}_p)
$$
convergeant vers $H^{p + q}(X, \mathcal{F})$.
\end{lemma}

\begin{proof}
Conséquence immédiate des lemmes \ref{spaces-simplicial-lemma-compute-via-proper-hypercovering}
et \ref{spaces-simplicial-lemma-simplicial-sheaf-cohomology}.
\end{proof}




\section{Schémas simpliciaux}
\label{spaces-simplicial-section-simplicial}

\noindent
Un {\it schéma simplicial} est un objet simplicial dans la catégorie des schémas,
voir Simplicial, définition \ref{simplicial-definition-simplicial-object}.
Rappelons qu'un schéma simplicial ressemble à
$$
\xymatrix{
X_2
\ar@<2ex>[r]
\ar@<0ex>[r]
\ar@<-2ex>[r]
&
X_1
\ar@<1ex>[r]
\ar@<-1ex>[r]
\ar@<1ex>[l]
\ar@<-1ex>[l]
&
X_0
\ar@<0ex>[l]
}
$$
Il y a ici deux morphismes $d^1_0, d^1_1 : X_1 \to X_0$
et un seul morphisme $s^0_0 : X_0 \to X_1$, etc.
Ces morphismes satisfont notamment aux relations
$d^1_0 \circ s^0_0 = \text{id}_{X_0} = d^1_1 \circ s^0_0$, voir
Simplicial, lemme \ref{simplicial-lemma-characterize-simplicial-object}.
Il est utile de considérer $d^n_i : X_n \to X_{n - 1}$
comme la « projection qui oublie la $i$-ième coordonnée », et
$s^n_j : X_n \to X_{n + 1}$ comme le « morphisme diagonal qui répète
la $j$-ième coordonnée ».

\medskip\noindent
Un {\it morphisme de schémas simpliciaux} $h : X \to Y$ est la même chose
qu'un morphisme d'objets simpliciaux dans la catégorie des schémas,
voir Simplicial, définition \ref{simplicial-definition-simplicial-object}.
Ainsi, $h$ est constitué de morphismes de schémas $h_n : X_n \to Y_n$
tels que $h_{n - 1} \circ d^n_j = d^n_j \circ h_n$ et
$h_{n + 1} \circ s^n_j = s^n_j \circ h_n$ chaque fois que ces expressions sont définies.

\medskip\noindent
Une {\it augmentation} d'un schéma simplicial $X$ est un morphisme
de schémas $a_0 : X_0 \to S$ tel que $a_0 \circ d^1_0 = a_0 \circ d^1_1$.
Voir Simplicial, section \ref{simplicial-section-augmentation}.

\medskip\noindent
Soit $X$ un schéma simplicial. La construction de la section
\ref{spaces-simplicial-section-simplicial-top} appliquée à l'espace topologique simplicial
sous-jacent donne un site $X_{Zar}$.
En revanche, pour tout $n$, nous avons le petit site de Zariski
$X_{n, Zar}$ (Topologies, définition
\ref{topologies-definition-big-small-Zariski})
et pour tout morphisme $\varphi : [m] \to [n]$
on a un morphisme de sites
$f_\varphi = X(\varphi)_{small} : X_{n, Zar} \to X_{m, Zar}$,
associé au morphisme de schémas
$X(\varphi) : X_n \to X_m$ (Topologies, lemme
\ref{topologies-lemma-morphism-big-small}).
Cela donne un objet simplicial $\mathcal{C}$ dans la catégorie des sites.
Dans le lemme \ref{spaces-simplicial-lemma-simplicial-site-site}, on a construit le site
associé $\mathcal{C}_{total}$. Le foncteur qui associe son image à une immersion ouverte
définit une équivalence $\mathcal{C}_{total} \to X_{Zar}$ qui
identifie les faisceaux, c'est-à-dire $\Sh(\mathcal{C}_{total}) = \Sh(X_{Zar})$.
La différence entre $\mathcal{C}_{total}$ et $X_{Zar}$
est similaire à la différence entre le petit site Zariski $S_{Zar}$
et l'espace topologique sous-jacent de $S$.
Nous identifierons tacitement ces sites dans la suite.

\medskip\noindent
Soit $X_{Zar}$ le site associé à un schéma simplicial $X$.
Il existe un faisceau d'anneaux $\mathcal{O}$ sur $X_{Zar}$ dont la restriction
à $X_n$ est le faisceau structural $\mathcal{O}_{X_n}$. Cela résulte
du lemme \ref{spaces-simplicial-lemma-describe-sheaves-simplicial-site} ou du lemme
\ref{spaces-simplicial-lemma-describe-sheaves-simplicial-site-site}. Nous dirons que
{\it $\mathcal{O}$ est le faisceau structural du schéma simplicial $X$}.
À ce stade, tout ce qui a été développé pour les sites simpliciaux (annelés)
s'applique, voir les sections \ref{spaces-simplicial-section-simplicial-sites},
\ref{spaces-simplicial-section-augmentation-simplicial-sites},
\ref{spaces-simplicial-section-morphism-simplicial-sites},
\ref{spaces-simplicial-section-simplicial-sites-modules},
\ref{spaces-simplicial-section-cohomology-simplicial-sites},
\ref{spaces-simplicial-section-cohomology-augmentation-simplicial-sites},
\ref{spaces-simplicial-section-cohomology-simplicial-sites-modules},
\ref{spaces-simplicial-section-cohomology-augmentation-ringed-simplicial-sites},
\ref{spaces-simplicial-section-cartesian},
\ref{spaces-simplicial-section-glueing} et
\ref{spaces-simplicial-section-glueing-modules}.

\medskip\noindent
Soit $X$ un schéma simplicial muni de son faisceau structural $\mathcal{O}$.
Comme sur tout topos annelé, il existe une notion
de {\it $\mathcal{O}$-module quasi-cohérent sur $X_{Zar}$}; voir Modules
sur les sites, définition \ref{sites-modules-definition-site-local}.
Cependant, un $\mathcal{O}$-module quasi-cohérent sur $X_{Zar}$ n'est
autre qu'un $\mathcal{O}$-module cartésien $\mathcal{F}$ dont les restrictions
$\mathcal{F}_n$ sont quasi-cohérentes sur $X_n$; voir le
lemme \ref{spaces-simplicial-lemma-quasi-coherent-sheaf}.

\medskip\noindent
Soit $h : X \to Y$ un morphisme de schémas simpliciaux. Le
lemme \ref{spaces-simplicial-lemma-simplicial-space-site-functorial} ou
(la preuve du) lemme \ref{spaces-simplicial-lemma-morphism-simplicial-sites}
fournit un morphisme de sites $h_{Zar} : X_{Zar} \to Y_{Zar}$.
Rappelons que $h_{Zar}^{-1}$ et $h_{Zar, *}$ s'expriment simplement
en termes des composantes; voir le lemme
\ref{spaces-simplicial-lemma-describe-functoriality} ou le lemme
\ref{spaces-simplicial-lemma-morphism-simplicial-sites}.
Notons $\mathcal{O}_X$ et $\mathcal{O}_Y$ les faisceaux structuraux
respectifs de $X$ et de $Y$. On définit
$h_{Zar}^\sharp : h_{Zar, *}\mathcal{O}_X \to \mathcal{O}_Y$
comme le morphisme de faisceaux d'anneaux sur $Y_{Zar}$ donné par
$h_n^\sharp : h_{n, *}\mathcal{O}_{X_n} \to \mathcal{O}_{Y_n}$ sur $Y_n$.
On obtient un morphisme de sites annelés
$$
h_{Zar} : (X_{Zar}, \mathcal{O}_X) \longrightarrow (Y_{Zar}, \mathcal{O}_Y)
$$

\medskip\noindent
Soit $X$ un schéma simplicial muni de son faisceau structural $\mathcal{O}$.
Soit $S$ un schéma et $a_0 : X_0 \to S$ une augmentation de $X$.
Le
lemme \ref{spaces-simplicial-lemma-augmentation} ou le
lemme \ref{spaces-simplicial-lemma-augmentation-site}
fournit un morphisme correspondant de topoi $a : \Sh(X_{Zar}) \to \Sh(S)$.
Observons que $a^{-1}\mathcal{G}$ est le faisceau sur $X_{Zar}$ de composantes
$a_n^{-1}\mathcal{G}$. On peut donc utiliser les morphismes
$a_n^\sharp : a_n^{-1}\mathcal{O}_S \to \mathcal{O}_{X_n}$ pour définir
un morphisme $a^\sharp : a^{-1}\mathcal{O}_S \to \mathcal{O}$ ou, de manière équivalente,
par adjonction, un morphisme $a^\sharp : \mathcal{O}_S \to a_*\mathcal{O}$
(qui, comme d'habitude, porte le même nom). Cela nous place dans la situation
discutée dans la section
\ref{spaces-simplicial-section-cohomology-augmentation-ringed-simplicial-sites}.
On obtient donc un morphisme de topoi annelés
$$
a : (\Sh(X_{Zar}), \mathcal{O}) \longrightarrow (\Sh(S), \mathcal{O}_S)
$$

\medskip\noindent
Enfin, supposons donné un morphisme
$h : X \to Y$ entre des schémas simpliciaux $X$ et $Y$, munis de leurs faisceaux structuraux
$\mathcal{O}_X$, $\mathcal{O}_Y$, d'augmentations
$a_0 : X_0 \to X_{-1}$, $b_0 : Y_0 \to Y_{-1}$ et un morphisme
$h_{-1} : X_{-1} \to Y_{-1}$ tel que le diagramme
$$
\xymatrix{
X_0 \ar[r]_{h_0} \ar[d]_{a_0} & Y_0 \ar[d]^{b_0} \\
X_{-1} \ar[r]^{h_{-1}} & Y_{-1}
}
$$
soit commutatif. Les constructions décrites ci-dessus
fournissent alors le diagramme commutatif suivant de morphismes de topoi annelés :
$$
\xymatrix{
(\Sh(X_{Zar}), \mathcal{O}_X) \ar[r]_{h_{Zar}} \ar[d]_a &
(\Sh(Y_{Zar}), \mathcal{O}_Y) \ar[d]^b \\
(\Sh(X_{-1}), \mathcal{O}_{X_{-1}}) \ar[r]^{h_{-1}} &
(\Sh(Y_{-1}), \mathcal{O}_{Y_{-1}})
}
$$








\section{Descente en termes de schémas simpliciaux}
\label{spaces-simplicial-section-simplicial-descent}

\noindent
Les morphismes cartésiens sont définis comme suit.

\begin{definition}
\label{spaces-simplicial-definition-cartesian-morphism}
Soit $a : Y \to X$ un morphisme de schémas simpliciaux.
On dit que $a$ est {\it cartésien}, ou que {\it $Y$ est cartésien au-dessus de $X$},
si pour tout morphisme $\varphi : [n] \to [m]$ de $\Delta$ le diagramme
correspondant
$$
\xymatrix{
Y_m \ar[r]_a \ar[d]_{Y(\varphi)} & X_m \ar[d]^{X(\varphi)}\\
Y_n \ar[r]^{a} & X_n
}
$$
est un carré cartésien dans la catégorie des schémas.
\end{definition}

\noindent
Les morphismes cartésiens sont liés aux données de descente. Commençons par un lemme général
décrivant, pour un schéma simplicial fixé, la catégorie des schémas simpliciaux cartésiens
au-dessus de celui-ci. Dans ce lemme, $f^* : \Sch/X \to \Sch/Y$ désigne
le foncteur de changement de base associé à un morphisme de schémas $f : Y \to X$.

\begin{lemma}
\label{spaces-simplicial-lemma-characterize-cartesian-schemes}
Soit $X$ un schéma simplicial. La catégorie des schémas simpliciaux cartésiens
sur $X$ est équivalente à la catégorie des couples $(V, \varphi)$
où $V$ est un schéma sur $X_0$ et
$$
\varphi :
V \times_{X_0, d^1_1} X_1
\longrightarrow
X_1 \times_{d^1_0, X_0} V
$$
est un isomorphisme sur $X_1$ tel que
$(s_0^0)^*\varphi = \text{id}_V$ et tel que
$$
(d^2_1)^*\varphi = (d^2_0)^*\varphi \circ (d^2_2)^*\varphi
$$
comme morphismes de schémas sur $X_2$.
\end{lemma}

\begin{proof}
L'égalité affichée a un sens puisque
$d^1_1 \circ d^2_2 = d^1_1 \circ d^2_1$,
$d^1_1 \circ d^2_0 = d^1_0 \circ d^2_2$ et
$d^1_0 \circ d^2_0 = d^1_0 \circ d^2_1$ comme morphismes $X_2 \to X_0$, voir
Simplicial, remarque \ref{simplicial-remark-relations}; on peut donc
représenter ces morphismes comme suit
$$
\xymatrix{
&
X_2 \times_{d^1_1 \circ d^2_0, X_0} V
\ar[r]_-{(d^2_0)^*\varphi} &
X_2 \times_{d^1_0 \circ d^2_0, X_0} V
\ar@{=}[rd] & \\
X_2 \times_{d^1_0 \circ d^2_2, X_0} V
\ar@{=}[ru] & & &
X_2 \times_{d^1_0 \circ d^2_1, X_0} V \\
&
X_2 \times_{d^1_1 \circ d^2_2, X_0} V
\ar[lu]^{(d^2_2)^*\varphi} \ar@{=}[r] &
X_2 \times_{d^1_1 \circ d^2_1, X_0} V
\ar[ru]_{(d^2_1)^*\varphi}
}
$$
et la condition signifie que le diagramme est commutatif. Or,
à tout schéma simplicial cartésien $Y$ au-dessus de $X$, on peut associer
$V = Y_0$ et prendre pour $\varphi$ le composé
$$
V \times_{X_0, d^1_1} X_1 =
Y_0 \times_{X_0, d^1_1} X_1 = Y_1 =
X_1 \times_{X_0, d^1_0} Y_0 =
X_1 \times_{X_0, d^1_0} V
$$
des identifications données par la structure cartésienne. Pour montrer que ce foncteur
est une équivalence, construisons-en un quasi-inverse. La construction du
quasi-inverse est analogue à celle décrite dans
Descente, section \ref{descent-section-descent-modules}, à laquelle nous empruntons
les notations $\tau^n_i : [0] \to [n]$, $0 \mapsto i$ et
$\tau^n_{ij} : [1] \to [n]$, $0 \mapsto i$, $1 \mapsto j$.
Plus précisément, pour un couple $(V, \varphi)$
comme dans le lemme, posons $Y_n = X_n \times_{X(\tau^n_n), X_0} V$.
Pour $\beta : [n] \to [m]$, définissons alors
$V(\beta) : Y_m \to Y_n$ comme le changement de base par $X(\tau^m_{\beta(n)m})$
du morphisme $\varphi$ composé avec la projection
$X_m \times_{X(\beta), X_n} Y_n \to Y_n$. Cette définition a un sens, car
$$
X_m \times_{X(\tau^m_{\beta(n)m}), X_1} X_1 \times_{d^1_1, X_0} V
=
X_m \times_{X(\tau^m_m), X_0} V = Y_m
$$
et
$$
X_m \times_{X(\tau^m_{\beta(n)m}), X_1} X_1 \times_{d^1_0, X_0} V =
X_m \times_{X(\tau^m_{\beta(n)}), X_0} V =
X_m \times_{X(\beta), X_n} Y_n.
$$
Nous laissons au lecteur la vérification que la commutativité
du diagramme ci-dessus
entraîne la compatibilité des morphismes avec la composition, ainsi que la vérification
que les deux foncteurs sont quasi inverses l'un de l'autre.
\end{proof}

\begin{definition}
\label{spaces-simplicial-definition-fibre-products-simplicial-scheme}
Soit $f : X \to S$ un morphisme de schémas. Le {\it schéma simplicial
associé à $f$}, noté $(X/S)_\bullet$, est le foncteur
$\Delta^{opp} \to \Sch$, $[n] \mapsto X \times_S \ldots \times_S X$
décrit dans
Simplicial, exemple \ref{simplicial-example-fibre-products-simplicial-object}.
\end{definition}

\noindent
Ainsi, $(X/S)_n$ est le produit fibré de $(n + 1)$ copies de $X$ sur $S$.
Le morphisme $d^1_0 : X \times_S X \to X$ est la projection
$(x_0, x_1) \mapsto x_1$, et le morphisme $d^1_1$ est l'autre
projection. Le morphisme $s^0_0$ est le morphisme diagonal
$X \to X \times_S X$.

\begin{lemma}
\label{spaces-simplicial-lemma-cartesian-over}
Soit $f : X \to S$ un morphisme de schémas.
Soit $\pi : Y \to (X/S)_\bullet$ un morphisme cartésien
de schémas simpliciaux.
Posons $V = Y_0$ considéré comme un schéma sur $X$.
Les morphismes $d^1_0, d^1_1 : Y_1 \to Y_0$ et le morphisme
$\pi_1 : Y_1 \to X \times_S X$ induisent des isomorphismes
$$
\xymatrix{
V \times_S X & &
Y_1 \ar[ll]_-{(d^1_1, \text{pr}_1 \circ \pi_1)}
\ar[rr]^-{(\text{pr}_0 \circ \pi_1, d^1_0)} & &
X \times_S V.
}
$$
Notons $\varphi : V \times_S X \to X \times_S V$
l'isomorphisme ainsi obtenu.
Alors le couple $(V, \varphi)$ est une donnée de descente relative
à $X \to S$.
\end{lemma}

\begin{proof}
C'est un cas particulier d'une des assertions du
lemme \ref{spaces-simplicial-lemma-characterize-cartesian-schemes},
car l'équation affichée de ce lemme est
équivalente à la condition de cocycle de
Descente, définition \ref{descent-definition-descent-datum}.
\end{proof}

\begin{lemma}
\label{spaces-simplicial-lemma-cartesian-equivalent-descent-datum}
Soit $f : X \to S$ un morphisme de schémas. La construction
$$
\begin{matrix}
\text{catégorie des schémas cartésiens} \\
\text{au-dessus de } (X/S)_\bullet
\end{matrix}
\longrightarrow
\begin{matrix}
\text{catégorie des données de descente} \\
\text{relatives à } X/S
\end{matrix}
$$
du lemme \ref{spaces-simplicial-lemma-cartesian-over}
est une équivalence de catégories.
\end{lemma}

\begin{proof}
Le foncteur de gauche à droite est donné dans le lemme
\ref{spaces-simplicial-lemma-cartesian-over}.
Il s'agit donc d'un cas particulier du lemme
\ref{spaces-simplicial-lemma-characterize-cartesian-schemes}.
\end{proof}

\noindent
Nous pouvons reformuler le changement de base du
lemme \ref{descent-lemma-pullback} de Descente comme suit.
Supposons donnés un morphisme de schémas simpliciaux $f : X' \to X$ et un morphisme cartésien
de schémas simpliciaux $Y \to X$. Alors
le produit fibré (considéré comme une « image inverse »)
$$
f^*Y = Y \times_X X'
$$
est un schéma simplicial cartésien sur $X'$.
Supposons qu'on donne un diagramme commutatif de morphismes de schémas
$$
\xymatrix{
X' \ar[r]_f \ar[d] & X \ar[d] \\
S' \ar[r] & S.
}
$$
On obtient ainsi un morphisme de schémas simpliciaux
$$
f_\bullet : (X'/S')_\bullet \longrightarrow (X/S)_\bullet.
$$
Nous affirmons que le « changement de base » $f_\bullet^*$ associé au morphisme
$f_\bullet : (X'/S')_\bullet \to (X/S)_\bullet$ correspond, via
le lemme \ref{spaces-simplicial-lemma-cartesian-equivalent-descent-datum},
au changement de base défini, en termes de données de descente, dans
Descente, lemme
\ref{descent-lemma-pullback}.







\section{Modules quasi-cohérents sur les schémas simpliciaux}
\label{spaces-simplicial-section-modules-simplicial}

\begin{lemma}
\label{spaces-simplicial-lemma-pullback-cartesian-module}
Soit $f : V \to U$ un morphisme de schémas simpliciaux. Pour un
module quasi-cohérent $\mathcal{F}$ sur $U_{Zar}$, l'image inverse
$f^*\mathcal{F}$ est un module quasi-cohérent sur $V_{Zar}$.
\end{lemma}

\begin{proof}
Rappelons que $\mathcal{F}$ est cartésien avec $\mathcal{F}_n$
quasi-cohérent, voir le lemme \ref{spaces-simplicial-lemma-quasi-coherent-sheaf}.
D'après le lemme \ref{spaces-simplicial-lemma-describe-functoriality}, nous voyons que
$(f^*\mathcal{F})_n = f_n^*\mathcal{F}_n$ (quelques détails omis).
$(f^*\mathcal{F})_n$ est donc quasi-cohérent.
Le même fait et la propriété cartésienne pour $\mathcal{F}$
impliquent la propriété cartésienne pour $f^*\mathcal{F}$.
Ainsi $\mathcal{F}$ est quasi-cohérent par
le lemme \ref{spaces-simplicial-lemma-quasi-coherent-sheaf}, une fois encore.
\end{proof}

\begin{lemma}
\label{spaces-simplicial-lemma-pushforward-cartesian-module}
Soit $f : V \to U$ un morphisme cartésien de schémas simpliciaux.
Supposons que les morphismes $d^n_j : U_n \to U_{n - 1}$ soient
plats et que les morphismes $V_n \to U_n$ soient quasi-compacts et quasi-séparés.
Pour un module quasi-cohérent $\mathcal{G}$ sur $V_{Zar}$
l'image directe $f_*\mathcal{G}$ est un module quasi-cohérent sur $U_{Zar}$.
\end{lemma}

\begin{proof}
Si $\mathcal{F} = f_* \mathcal{G}$, alors
$\mathcal{F}_n = f_{n , *}\mathcal{G}_n$ par
le lemme \ref{spaces-simplicial-lemma-describe-functoriality}.
Les morphismes $\mathcal{F}(\varphi)$ sont définis à l'aide des morphismes de changement de base; voir
Cohomologie, section \ref{cohomology-section-base-change-map}.
Les faisceaux $\mathcal{F}_n$ sont quasi-cohérents par
Schémas, lemme \ref{schemes-lemma-push-forward-quasi-coherent}
et le fait que $\mathcal{G}_n$ est quasi-cohérent
d'après le lemme \ref{spaces-simplicial-lemma-quasi-coherent-sheaf}.
Les morphismes de changement de base le long des dégénérescences
$d^n_j$ sont des isomorphismes d'après Cohomologie des schémas, lemme
\ref{coherent-lemma-flat-base-change-cohomology}
et le fait que $\mathcal{G}$ est cartésien
d'après le lemme \ref{spaces-simplicial-lemma-quasi-coherent-sheaf}.
Ainsi, $\mathcal{F}$ est cartésien d'après le
lemme \ref{spaces-simplicial-lemma-check-cartesian-module}.
Par conséquent, $\mathcal{F}$ est quasi-cohérent d'après le
lemme \ref{spaces-simplicial-lemma-quasi-coherent-sheaf}.
\end{proof}

\begin{lemma}
\label{spaces-simplicial-lemma-adjoint-functors-cartesian-modules}
Soit $f : V \to U$ un morphisme cartésien de schémas
simpliciaux. Supposons que les morphismes $d^n_j : U_n \to U_{n - 1}$ soient
plats et que les morphismes $V_n \to U_n$ soient quasi-compacts et quasi-séparés.
Alors $f^*$ et $f_*$ forment une paire de foncteurs adjoints
entre les catégories de modules quasi-cohérents sur $U_{Zar}$ et $V_{Zar}$.
\end{lemma}

\begin{proof}
Nous avons vu dans les lemmes \ref{spaces-simplicial-lemma-pullback-cartesian-module} et
\ref{spaces-simplicial-lemma-pushforward-cartesian-module}
que l'assertion a un sens. La propriété d'adjonction résulte
du fait que chaque $f_n^*$ est adjoint à $f_{n, *}$.
\end{proof}

\begin{lemma}
\label{spaces-simplicial-lemma-cartesian-modules-with-section}
Soit $f : X \to S$ un morphisme de schémas qui admet une section
\footnote{En fait, il suffirait de supposer que $f$
admet fpqc-localement sur $S$ une section, puisque les modules quasi-cohérents
satisfont la descente d'après Descente,
section \ref{descent-section-fpqc-descent-quasi-coherent}.}.
Soit $(X/S)_\bullet$ le schéma simplicial
associé à $X \to S$, voir
définition \ref{spaces-simplicial-definition-fibre-products-simplicial-scheme}.
Alors l'image inverse définit une équivalence entre la catégorie des
$\mathcal{O}_S$-modules quasi-cohérents et la catégorie des modules quasi-cohérents
sur $((X/S)_\bullet)_{Zar}$.
\end{lemma}

\begin{proof}
Soit $\sigma : S \to X$ une section de $f$. Soit $(\mathcal{F}, \alpha)$
une paire comme dans le lemme \ref{spaces-simplicial-lemma-characterize-cartesian-modules}.
Posons $\mathcal{G} = \sigma^*\mathcal{F}$. Considérons le diagramme
$$
\xymatrix{
X \ar[r]_-{(\sigma \circ f, 1)} \ar[d]_f &
X \times_S X \ar[d]^{\text{pr}_0} \ar[r]_-{\text{pr}_1} & X \\
S \ar[r]^\sigma & X
}
$$
Notons que $\text{pr}_0 = d^1_1$ et $\text{pr}_1 = d^1_0$. Nous voyons donc que
$(\sigma \circ f, 1)^*\alpha$ définit un isomorphisme
$$
f^*\mathcal{G} = (\sigma \circ f, 1)^*\text{pr}_0^*\mathcal{F}
\longrightarrow
(\sigma \circ f, 1)^*\text{pr}_1^*\mathcal{F} = \mathcal{F}
$$
Nous omettons de vérifier que cet isomorphisme est compatible
avec $\alpha$ et l'isomorphisme canonique
$\text{pr}_0^*f^*\mathcal{G} \to \text{pr}_1^*f^*\mathcal{G}$.
\end{proof}






\section{Groupoïdes et schémas simpliciaux}
\label{spaces-simplicial-section-groupoids-simplicial}

\noindent
Étant donné un groupoïde en schémas, nous pouvons construire un schéma simplicial.
Il s'avère que la catégorie des faisceaux quasi-cohérents sur un groupoïde
est équivalente à la catégorie des faisceaux cartésiens quasi-cohérents
sur le schéma simplicial associé.

\begin{lemma}
\label{spaces-simplicial-lemma-groupoid-simplicial}
Soit $(U, R, s, t, c, e, i)$ un groupoïde en schémas sur $S$.
Il existe un schéma simplicial $X$ sur $S$
avec les propriétés suivantes
\begin{enumerate}
\item $X_0 = U$, $X_1 = R$, $X_2 = R \times_{s, U, t} R$,
\item $s_0^0 = e : X_0 \to X_1$,
\item $d^1_0 = s : X_1 \to X_0$, $d^1_1 = t : X_1 \to X_0$,
\item $s_0^1 = (e \circ t, 1) : X_1 \to X_2$,
$s_1^1 = (1, e \circ t) : X_1 \to X_2$,
\item $d^2_0 = \text{pr}_1 : X_2 \to X_1$,
$d^2_1 = c : X_2 \to X_1$,
$d^2_2 = \text{pr}_0$ et
\item $X = \text{cosk}_2 \text{sk}_2 X$.
\end{enumerate}
Pour tout $n$, on a $X_n = R \times_{s, U, t} \ldots \times_{s, U, t} R$
à $n$ facteurs. Le morphisme $d^n_j : X_n \to X_{n - 1}$ est donné sur les
foncteurs de points par
$$
(r_1, \ldots, r_n) \longmapsto (r_1, \ldots, c(r_j, r_{j + 1}), \ldots, r_n)
$$
pour $1 \leq j \leq n - 1$ alors que
$d^n_0(r_1, \ldots, r_n) = (r_2, \ldots, r_n)$ et
$d^n_n(r_1, \ldots, r_n) = (r_1, \ldots, r_{n - 1})$.
\end{lemma}

\begin{proof}
Il suffit de vérifier que les règles prescrites en (1), (2), (3), (4), (5)
définissent un schéma simplicial $2$-tronqué $U'$ sur $S$, puisque (6)
nous permet de définir $X = \text{cosk}_2 U'$, voir
Simplicial, lemme \ref{simplicial-lemma-existence-cosk}.
En utilisant l'approche du foncteur de points, il suffit de vérifier que
si $(\text{Ob}, \text{Arrows}, s, t, c, e, i)$ est un groupoïde, alors
$$
\xymatrix{
\text{Arrows} \times_{s, \text{Ob}, t} \text{Arrows}
\ar@<8ex>[d]^{\text{pr}_0}
\ar@<0ex>[d]_c
\ar@<-8ex>[d]_{\text{pr}_1}
\\
\text{Arrows}
\ar@<4ex>[d]^t
\ar@<-4ex>[d]_s
\ar@<4ex>[u]^{1, e}
\ar@<-4ex>[u]_{e, 1}
\\
\text{Ob}
\ar@<0ex>[u]_e
}
$$
est un ensemble simplicial $2$-tronqué. Nous omettons les détails.

\medskip\noindent
Enfin, la description de $X_n$ pour $n > 2$ suit par induction de
la description de $X_0$, $X_1$, $X_2$ et
Simplicial, remarque \ref{simplicial-remark-inductive-coskeleton}, et du
lemme \ref{simplicial-lemma-work-out}. On peut aussi montrer que
l'application de $\text{cosk}_2$ à l'ensemble simplicial $2$-tronqué ci-dessus
donne un ensemble simplicial dont le terme de degré $n$ est égal à
$\text{Arrows} \times_{s, \text{Ob}, t} \ldots \times_{s, \text{Ob}, t}
\text{Arrows}$ à $n$ facteurs et les morphismes de dégénérescence décrits dans le lemme.
Certains détails ont été omis.
\end{proof}

\begin{lemma}
\label{spaces-simplicial-lemma-quasi-coherent-groupoid-simplicial}
Soient $S$ un schéma et $(U, R, s, t, c)$ un groupoïde en schémas
sur $S$. Soit $X$ le schéma simplicial sur $S$ construit
dans le lemme \ref{spaces-simplicial-lemma-groupoid-simplicial}.
Alors la catégorie des modules quasi-cohérents sur $(U, R, s, t, c)$
est équivalente à la catégorie des modules quasi-cohérents sur $X_{Zar}$.
\end{lemma}

\begin{proof}
Cela ressort clairement des lemmes
\ref{spaces-simplicial-lemma-quasi-coherent-sheaf} et
\ref{spaces-simplicial-lemma-characterize-cartesian-modules}
et Groupoïdes, définition \ref{groupoids-definition-groupoid-module}.
\end{proof}

\noindent
Dans le lemme suivant, nous utiliserons la notion de morphisme cartésien
$V \to U$ de schémas simpliciaux, telle qu'elle est définie dans la définition
\ref{spaces-simplicial-definition-cartesian-morphism}.

\begin{lemma}
\label{spaces-simplicial-lemma-quasi-coherent-groupoid-R-cartesian}
Soit $(U, R, s, t, c)$ un groupoïde en schémas sur un schéma $S$.
Soit $X$ le schéma simplicial sur $S$ construit
dans le lemme \ref{spaces-simplicial-lemma-groupoid-simplicial}.
Soit $(R/U)_\bullet$ le schéma simplicial
associé à $s : R \to U$, voir
définition \ref{spaces-simplicial-definition-fibre-products-simplicial-scheme}.
Il existe un morphisme cartésien $t_\bullet : (R/U)_\bullet \to X$
de schémas simpliciaux, dont les morphismes en petits degrés sont donnés par
$$
\xymatrix{
R \times_{s, U, s} R \times_{s, U, s} R
\ar@<3ex>[r]_-{\text{pr}_{12}}
\ar@<0ex>[r]_-{\text{pr}_{02}}
\ar@<-3ex>[r]_-{\text{pr}_{01}}
\ar[dd]_{(r_0, r_1, r_2) \mapsto (r_0 \circ r_1^{-1}, r_1 \circ r_2^{-1})} &
R \times_{s, U, s} R
\ar@<1ex>[r]_-{\text{pr}_1} \ar@<-2ex>[r]_-{\text{pr}_0}
\ar[dd]_{(r_0, r_1) \mapsto r_0 \circ r_1^{-1}} &
R \ar[dd]^t
\\
\\
R \times_{s, U, t} R
\ar@<3ex>[r]_{\text{pr}_1}
\ar@<0ex>[r]_c
\ar@<-3ex>[r]_{\text{pr}_0} &
R \ar@<1ex>[r]_s \ar@<-2ex>[r]_t &
U
}
$$
\end{lemma}

\begin{proof}
Pour tout $n$, nous définissons $(R/U)_\bullet \to X_n$ par la règle
$$
(r_0, \ldots, r_n)
\longrightarrow
(r_0 \circ r_1^{-1}, \ldots, r_{n - 1} \circ r_n^{-1})
$$
La compatibilité avec les morphismes de dégénérescence résulte de la description de ceux-ci
dans le lemme \ref{spaces-simplicial-lemma-groupoid-simplicial}.
On omet de vérifier que les morphismes respectent les $s^n_j$.
Le lemme \ref{groupoids-lemma-diagram-pull} de Groupoïdes
(avec les rôles de $s$ et $t$ inversés)
montre que les deux carrés de droite sont cartésiens. Exactement de la même manière,
on montre que tous les autres carrés sont également cartésiens. Le
morphisme est donc cartésien.
\end{proof}




\section{Les données de descente définissent des relations d'équivalence}
\label{spaces-simplicial-section-equivalence-relation}

\noindent
Dans la section \ref{spaces-simplicial-section-simplicial-descent}, nous avons vu comment les données de descente relatives à
$X \to S$ peuvent se formuler en termes de schémas simpliciaux cartésiens
au-dessus de $(X/S)_\bullet$. Nous les relions ici aux relations d'équivalence
comme suit.

\begin{lemma}
\label{spaces-simplicial-lemma-equivalence-relation}
Soit $f : X \to S$ un morphisme de schémas.
Soit $\pi : Y \to (X/S)_\bullet$ un morphisme cartésien de schémas
simpliciaux; voir les définitions \ref{spaces-simplicial-definition-cartesian-morphism} et
\ref{spaces-simplicial-definition-fibre-products-simplicial-scheme}.
Alors le morphisme
$$
j = (d^1_1, d^1_0) : Y_1 \to Y_0 \times_S Y_0
$$
définit une relation d'équivalence sur $Y_0$ au-dessus de $S$;
voir Groupoïdes, définition \ref{groupoids-definition-equivalence-relation}.
\end{lemma}

\begin{proof}
Notons que $j$ est un monomorphisme. En effet, le composé
$Y_1 \to Y_0 \times_S Y_0 \to Y_0 \times_S X$
est un isomorphisme car $\pi$ est cartésien.

\medskip\noindent
Considérons le morphisme
$$
(d^2_2, d^2_0) : Y_2 \to Y_1 \times_{d^1_0, Y_0, d^1_1} Y_1.
$$
Ce morphisme est bien défini car $d_0 \circ d_2 = d_1 \circ d_0$;
voir Simplicial, remarque \ref{simplicial-remark-relations}.
Il s'agit en outre d'un morphisme au-dessus de $(X/S)_2$. C'est un isomorphisme
car $Y \to (X/S)_\bullet$ est cartésien. Notons par exemple que le membre
de droite est isomorphe à
$Y_0 \times_{\pi_0, X, \text{pr}_1} (X \times_S X \times_S X) =
X \times_S Y_0 \times_S X$
car $\pi$ est cartésien. Détails omis.

\medskip\noindent
Comme dans Groupoïdes, définition \ref{groupoids-definition-equivalence-relation},
on note $t = \text{pr}_0 \circ j = d^1_1$ et
$s = \text{pr}_1 \circ j = d^1_0$.
L'isomorphisme ci-dessus, combiné au morphisme
$d^2_1 : Y_2 \to Y_1$, nous donne un morphisme de composition
$$
c : Y_1 \times_{s, Y_0, t} Y_1 \longrightarrow Y_1
$$
sur $Y_0 \times_S Y_0$. Cela implique immédiatement
que pour tout schéma $T/S$ la relation
$Y_1(T) \subset Y_0(T) \times Y_0(T)$ est transitive.

\medskip\noindent
La réflexivité découle du fait que la restriction
du morphisme $j$ à la diagonale
$\Delta : X \to X \times_S X$ est un isomorphisme
(utilisez à nouveau la propriété cartésienne de $\pi$).

\medskip\noindent
Pour voir la symétrie, nous considérons le morphisme
$$
(d^2_2, d^2_1) : Y_2 \to Y_1 \times_{d^1_1, Y_0, d^1_1} Y_1.
$$
Ce morphisme est bien défini car $d_1 \circ d_2 = d_1 \circ d_1$;
voir Simplicial, remarque \ref{simplicial-remark-relations}.
C'est un isomorphisme
car $Y \to (X/S)_\bullet$ est cartésien.
Notons par exemple que le membre
de droite est isomorphe à
$Y_0 \times_{\pi_0, X, \text{pr}_0} (X \times_S X \times_S X) =
Y_0 \times_S X \times_S X$
car $\pi$ est cartésien. Détails omis.

\medskip\noindent
Soit $T/S$ un schéma. Convenons que $a \sim b$, pour $a, b \in Y_0(T)$,
signifie $(a, b) \in Y_1(T)$.
L'isomorphisme $(d^2_2, d^2_1)$ au-dessus de
implique que si $a \sim b$ et $a \sim c$, alors $b \sim c$.
Combiné avec la réflexivité, cela montre que $\sim$ est
une relation d'équivalence.
\end{proof}







\section{Un exemple}
\label{spaces-simplicial-section-example}

\noindent
Dans cette section, on montre que les réunions disjointes de spectres
d'anneaux artiniens peuvent être descendues le long d'un morphisme de schémas
quasi compact, surjectif et plat.

\begin{lemma}
\label{spaces-simplicial-lemma-equivalence-classes-points}
Soit $X \to S$ un morphisme de schémas. Supposons que $Y \to (X/S)_\bullet$
soit un morphisme cartésien de schémas simpliciaux. Pour un point $y \in Y_0$, posons
$$
T_y = \{y' \in Y_0 \mid \exists\ y_1 \in Y_1:
d^1_1(y_1) = y, d^1_0(y_1) = y'\}
$$
comme sous-ensemble de $Y_0$. Alors $y \in T_y$ et
$T_y \cap T_{y'} \not = \emptyset \Rightarrow T_y = T_{y'}$.
\end{lemma}

\begin{proof}
Il suffit de combiner le lemme \ref{spaces-simplicial-lemma-equivalence-relation} et le
lemme de Groupoïdes
\ref{groupoids-lemma-pre-equivalence-equivalence-relation-points}.
\end{proof}

\begin{lemma}
\label{spaces-simplicial-lemma-quasi-compact}
Soit $X \to S$ un morphisme de schémas.
Supposons que $Y \to (X/S)_\bullet$ soit un morphisme cartésien de schémas simpliciaux.
Soit $y \in Y_0$ un point. Si $X \to S$ est quasi-compact, alors
$$
T_y = \{y' \in Y_0 \mid \exists\ y_1 \in Y_1:
d^1_1(y_1) = y, d^1_0(y_1) = y'\}
$$
est un sous-ensemble quasi-compact de $Y_0$.
\end{lemma}

\begin{proof}
Soit $F_y$ la fibre schématique de $d^1_1 : Y_1 \to Y_0$
à $y$. On voit alors que $T_y$ est l'image du morphisme
$$
\xymatrix{
F_y \ar[r] \ar[d] &
Y_1 \ar[r]^{d^1_0} \ar[d]^{d^1_1} &
Y_0 \\
y \ar[r] &
Y_0 &
}
$$
Notons que $F_y$ est quasi-compact. Cela prouve le lemme.
\end{proof}

\begin{lemma}
\label{spaces-simplicial-lemma-descent-disjoint-union-Artinian-along-fields}
Soit $X \to S$ un morphisme plat surjectif quasi-compact.
Soit $(V, \varphi)$ une donnée de descente relative
à $X \to S$. Si $V$ est une réunion disjointe de spectres
d'anneaux artiniens, alors $(V, \varphi)$ est effective.
\end{lemma}

\begin{proof}
Soit $Y \to (X/S)_\bullet$ le morphisme cartésien de schémas simpliciaux
correspondant à $(V, \varphi)$ par
le lemme \ref{spaces-simplicial-lemma-cartesian-equivalent-descent-datum}.
Notons que $Y_0 = V$.
Écrivons $V = \coprod_{i \in I} \Spec(A_i)$, où chaque $A_i$ est local
artinien. Notons en outre $v_i \in V$ l'unique point fermé de
$\Spec(A_i)$ pour tout $i \in I$. Posons $i \sim j$ si et seulement si
$v_i \in T_{v_j}$, avec les notations du
lemme \ref{spaces-simplicial-lemma-equivalence-classes-points} ci-dessus.
D'après les lemmes \ref{spaces-simplicial-lemma-equivalence-classes-points} et \ref{spaces-simplicial-lemma-quasi-compact},
il s'agit d'une relation d'équivalence à classes d'équivalence finies.
Posons $\overline{I} = I/\sim$. On peut alors écrire
$V = \coprod_{\overline{i} \in \overline{I}} V_{\overline{i}}$
avec
$V_{\overline{i}} = \coprod_{i \in \overline{i}} \Spec(A_i)$.
Par construction, nous voyons que
$\varphi : V \times_S X \to X \times_S V$ envoie
les sous-espaces ouverts et fermés $V_{\overline{i}} \times_S X$
dans les sous-espaces ouverts et fermés $X \times_S V_{\overline{i}}$.
Autrement dit, nous obtenons des données de descente
$(V_{\overline{i}}, \varphi_{\overline{i}})$, et
$(V, \varphi)$ en est le coproduit dans la catégorie des
données de descente.
Puisque chacun des $V_{\overline{i}}$ est une union finie de spectres
d'anneaux locaux artiniens, le morphisme $V_{\overline{i}} \to X$
est affine, voir Morphismes, lemme \ref{morphisms-lemma-Artinian-affine}.
Puisque $\{X \to S\}$ est un recouvrement fpqc, toutes les
données de descente $(V_{\overline{i}}, \varphi_{\overline{i}})$ sont effectives
d'après Descente, lemme \ref{descent-lemma-affine}.
\end{proof}

\noindent
Certes, le lemme ci-dessus a un champ d'application très limité !









\section{Espaces algébriques simpliciaux}
\label{spaces-simplicial-section-simplicial-algebraic-spaces}

\noindent
Soit $S$ un schéma. Un {\it espace algébrique simplicial}
est un objet simplicial dans la catégorie des espaces algébriques sur $S$,
voir Simplicial, définition \ref{simplicial-definition-simplicial-object}.
Rappelons qu'un espace algébrique simplicial ressemble à
$$
\xymatrix{
X_2
\ar@<2ex>[r]
\ar@<0ex>[r]
\ar@<-2ex>[r]
&
X_1
\ar@<1ex>[r]
\ar@<-1ex>[r]
\ar@<1ex>[l]
\ar@<-1ex>[l]
&
X_0
\ar@<0ex>[l]
}
$$
Ici, il y a deux morphismes $d^1_0, d^1_1 : X_1 \to X_0$
et un seul morphisme $s^0_0 : X_0 \to X_1$, etc.
Ces morphismes satisfont notamment aux relations
$d^1_0 \circ s^0_0 = \text{id}_{X_0} = d^1_1 \circ s^0_0$, voir
Simplicial, lemme \ref{simplicial-lemma-characterize-simplicial-object}.
Il est utile de considérer $d^n_i : X_n \to X_{n - 1}$
comme la « projection qui oublie la $i$-ième coordonnée », et
$s^n_j : X_n \to X_{n + 1}$ comme le « morphisme diagonal qui répète
la $j$-ième coordonnée ».

\medskip\noindent
Un {\it morphisme d'espaces algébriques simpliciaux} $h : X \to Y$ est la même
chose qu'un morphisme d'objets simpliciaux dans
la catégorie des espaces algébriques sur $S$,
voir Simplicial, définition \ref{simplicial-definition-simplicial-object}.
Ainsi, $h$ est constitué de morphismes d'espaces algébriques $h_n : X_n \to Y_n$
tels que $h_{n - 1} \circ d^n_j = d^n_j \circ h_n$ et
$h_{n + 1} \circ s^n_j = s^n_j \circ h_n$ chaque fois que ces expressions sont définies.

\medskip\noindent
Une {\it augmentation} $a : X \to X_{-1}$
d'un espace algébrique simplicial $X$ est donnée par un morphisme
d'espaces algébriques $a_0 : X_0 \to X_{-1}$
tel que $a_0 \circ d^1_0 = a_0 \circ d^1_1$.
Voir Simplicial, section \ref{simplicial-section-augmentation}.
Dans cette situation, on note toujours $a_n : X_n \to X_{-1}$ le morphisme induit
pour $n \geq 0$.

\medskip\noindent
Soit $X$ un espace algébrique simplicial. Pour chaque $n$, nous avons le site
$X_{n, spaces, \etale}$ (Propriétés des espaces, définition
\ref{spaces-properties-definition-spaces-etale-site})
et pour chaque morphisme $\varphi : [m] \to [n]$ nous avons un morphisme de sites
$$
f_\varphi = X(\varphi)_{spaces, \etale} :
X_{n, spaces, \etale} \to X_{m, spaces, \etale},
$$
associé au morphisme des espaces algébriques
$X(\varphi) : X_n \to X_m$ (Propriétés des espaces, lemme
\ref{spaces-properties-lemma-functoriality-etale-site}).
Cela donne un objet simplicial dans la catégorie des sites.
Dans le lemme \ref{spaces-simplicial-lemma-simplicial-site-site}, nous avons construit un site
associé que nous notons $X_{spaces, \etale}$.
Un objet du site $X_{spaces, \etale}$ est
un espace algébrique $U$ \'etale sur $X_n$ pour un certain $n$;
et un morphisme $(\varphi, f) : U/X_n \to V/X_m$ est donné
par un morphisme $\varphi : [m] \to [n]$ dans $\Delta$ et un morphisme
$f : U \to V$ d'espaces algébriques tel que le diagramme
$$
\xymatrix{
U \ar[r]_f \ar[d] & V \ar[d] \\
X_n \ar[r]^{f_\varphi} & X_m
}
$$
est commutatif. Considérons les sous-catégories pleines
$$
X_{affine, \etale} \subset X_\etale \subset X_{spaces, \etale}
$$
dont les objets sont $U/X_n$ avec $U$ affine, respectivement un schéma.
Munies de leurs topologies naturelles
(voir
Propriétés des espaces, lemme \ref{spaces-properties-lemma-alternative},
définition \ref{spaces-properties-definition-etale-site} et
lemme \ref{spaces-properties-lemma-compare-etale-sites}),
ces catégories ont des foncteurs d'inclusion qui définissent des équivalences de topoi
$$
\Sh(X_{affine, \etale}) = \Sh(X_\etale) = \Sh(X_{spaces, \etale})
$$
Dans la suite, nous identifierons tacitement ces topos.
On dira que $X_\etale$ est le {\it petit site \'etale de $X$}
et son topos est le {\it petit topos \'etale de $X$}.

\medskip\noindent
Soit $X_\etale$ le petit site \'etale d'un espace algébrique simplicial $X$.
Il existe un faisceau d'anneaux $\mathcal{O}$ sur $X_\etale$ dont la restriction
à $X_n$ est le faisceau structural $\mathcal{O}_{X_n}$. Cela découle du lemme
\ref{spaces-simplicial-lemma-describe-sheaves-simplicial-site-site}. Nous dirons
{\it $\mathcal{O}$ est le faisceau structural de l'espace algébrique simplicial
$X$}.
À ce stade, tout ce qui a été développé pour les sites simpliciaux (annelés)
s'applique, voir les sections \ref{spaces-simplicial-section-simplicial-sites},
\ref{spaces-simplicial-section-augmentation-simplicial-sites},
\ref{spaces-simplicial-section-morphism-simplicial-sites},
\ref{spaces-simplicial-section-simplicial-sites-modules},
\ref{spaces-simplicial-section-cohomology-simplicial-sites},
\ref{spaces-simplicial-section-cohomology-augmentation-simplicial-sites},
\ref{spaces-simplicial-section-cohomology-simplicial-sites-modules},
\ref{spaces-simplicial-section-cohomology-augmentation-ringed-simplicial-sites},
\ref{spaces-simplicial-section-cartesian},
\ref{spaces-simplicial-section-glueing} et
\ref{spaces-simplicial-section-glueing-modules}.

\medskip\noindent
Soit $X$ un espace algébrique simplicial muni de son faisceau structural $\mathcal{O}$.
Comme sur tout topos annelé, il existe une notion
de {\it $\mathcal{O}$-module quasi-cohérent sur $X_\etale$}; voir Modules
sur les sites, définition \ref{sites-modules-definition-site-local}.
Cependant, un $\mathcal{O}$-module quasi-cohérent sur $X_\etale$ n'est
autre qu'un $\mathcal{O}$-module cartésien $\mathcal{F}$ dont les restrictions
$\mathcal{F}_n$ sont quasi-cohérentes sur $X_n$; voir le lemme
\ref{spaces-simplicial-lemma-quasi-coherent-sheaf}.

\medskip\noindent
Soit $h : X \to Y$ un morphisme d'espaces algébriques simpliciaux sur $S$.
D'après le lemme \ref{spaces-simplicial-lemma-morphism-simplicial-sites}, appliqué aux morphismes
des sites $(h_n)_{spaces, \etale} : X_{spaces, \etale} \to Y_{spaces, \etale}$
(Propriétés des espaces, lemme
\ref{spaces-properties-lemma-functoriality-etale-site})
on obtient un morphisme de petits topoi étales
$h_\etale : \Sh(X_\etale) \to \Sh(Y_\etale)$.
Rappelons que $h_\etale^{-1}$ et $h_{\etale, *}$ se décrivent simplement
en termes de leurs composantes; voir le
lemme \ref{spaces-simplicial-lemma-morphism-simplicial-sites}.
Notons $\mathcal{O}_X$, resp.\ $\mathcal{O}_Y$, le faisceau structural
de $X$, resp.\ de $Y$. On définit
$h_\etale^\sharp : h_{\etale, *}\mathcal{O}_X \to \mathcal{O}_Y$
comme le morphisme de faisceaux d'anneaux sur $Y_\etale$ donné par
$h_n^\sharp : h_{n, *}\mathcal{O}_{X_n} \to \mathcal{O}_{Y_n}$ sur $Y_n$.
On obtient un morphisme de topoi annelés
$$
h_\etale :
(\Sh(X_\etale), \mathcal{O}_X)
\longrightarrow
(\Sh(Y_\etale), \mathcal{O}_Y)
$$

\medskip\noindent
Soit $X$ un espace algébrique simplicial muni de son faisceau structural $\mathcal{O}$.
Soit $X_{-1}$ un espace algébrique sur $S$ et soit $a_0 : X_0 \to X_{-1}$
une augmentation de $X$. Par
le lemme \ref{spaces-simplicial-lemma-augmentation-site},
appliqué au morphisme des sites
$(a_0)_{spaces, \etale} :
X_{0, spaces, \etale} \to X_{-1, spaces, \etale}$
on obtient un morphisme correspondant de topoi
$a : \Sh(X_\etale) \to \Sh(X_{-1, \etale})$.
Observons que $a^{-1}\mathcal{G}$ est le faisceau sur
$X_\etale$ de composantes $a_n^{-1}\mathcal{G}$. On peut donc utiliser les morphismes
$a_n^\sharp : a_n^{-1}\mathcal{O}_{X_{-1}} \to \mathcal{O}_{X_n}$ pour définir
un morphisme $a^\sharp : a^{-1}\mathcal{O}_{X_{-1}} \to \mathcal{O}$ ou, de manière équivalente,
par adjonction, un morphisme $a^\sharp : \mathcal{O}_{X_{-1}} \to a_*\mathcal{O}$
(qui, comme d'habitude, porte le même nom). Cela nous place dans la situation
discutée dans la section
\ref{spaces-simplicial-section-cohomology-augmentation-ringed-simplicial-sites}.
On obtient donc un morphisme de topoi annelés
$$
a :
(\Sh(X_\etale), \mathcal{O})
\longrightarrow
(\Sh(X_{-1}), \mathcal{O}_{X_{-1}})
$$

\medskip\noindent
Faisons une dernière observation. Supposons donné un morphisme
$h : X \to Y$ entre des espaces algébriques simpliciaux $X$ et $Y$, munis de leurs faisceaux structuraux
$\mathcal{O}_X$, $\mathcal{O}_Y$, d'augmentations
$a_0 : X_0 \to X_{-1}$, $b_0 : Y_0 \to Y_{-1}$ et d'un morphisme
$h_{-1} : X_{-1} \to Y_{-1}$ tel que le diagramme
$$
\xymatrix{
X_0 \ar[r]_{h_0} \ar[d]_{a_0} & Y_0 \ar[d]^{b_0} \\
X_{-1} \ar[r]^{h_{-1}} & Y_{-1}
}
$$
soit commutatif. Les constructions décrites ci-dessus
fournissent alors le diagramme commutatif suivant de morphismes de topoi annelés :
$$
\xymatrix{
(\Sh(X_\etale), \mathcal{O}_X) \ar[r]_{h_\etale} \ar[d]_a &
(\Sh(Y_\etale), \mathcal{O}_Y) \ar[d]^b \\
(\Sh(X_{-1}), \mathcal{O}_{X_{-1}}) \ar[r]^{h_{-1}} &
(\Sh(Y_{-1}), \mathcal{O}_{Y_{-1}})
}
$$




\section{Hyperrecouvrements fppf d'espaces algébriques}
\label{spaces-simplicial-section-fppf-hypercovering}

\noindent
Cette section est l'analogue de la section \ref{spaces-simplicial-section-proper-hypercovering}
pour le cas des espaces algébriques et des hyperrecouvrements fppf.
Le lecteur qui le souhaite peut remplacer partout « espace algébrique »
par « schéma » et obtenir des résultats également valables.
Cela présente l'avantage de remplacer les références à
Compléments sur la cohomologie des espaces, section
\ref{spaces-more-cohomology-section-fppf-etale}
par des références à
Cohomologie étale, section \ref{etale-cohomology-section-fppf-etale}.

\medskip\noindent
Nous fixons un schéma de base $S$.
Soit $X$ un espace algébrique sur $S$ et soit $U$ un espace algébrique simplicial
sur $S$. Supposons que nous ayons une augmentation
$$
a : U \to X
$$
Voir la section \ref{spaces-simplicial-section-simplicial-algebraic-spaces}.
On dit que $U$ est un {\it hyperrecouvrement fppf} de $X$ si
\begin{enumerate}
\item $U_0 \to X$ est plat, localement de présentation finie, et surjectif,
\item $U_1 \to U_0 \times_X U_0$ est plat, localement de présentation finie, et
surjectif,
\item $U_{n + 1} \to (\text{cosk}_n\text{sk}_n U)_{n + 1}$
est plat, localement de présentation finie, et surjectif pour $n \geq 1$.
\end{enumerate}
La catégorie des espaces algébriques sur $S$ admet toutes les limites finies; les cosquelettes
utilisés dans la formulation ci-dessus existent donc.
$$
\fbox{Principe : les hyperrecouvrements fppf permettent de calculer la cohomologie étale.}
$$
L'idée essentielle de la démonstration consiste à comparer les topologies
fppf et \'etale sur la catégorie $\textit{Espaces}/S$.
La topologie fppf est en effet plus fine que la topologie \'etale, et l'on a :
(a) tout morphisme plat, localement de présentation finie et surjectif définit
un recouvrement fppf, et
(b) la cohomologie fppf des faisceaux provenant par image inverse du petit site \'etale
coïncide avec la cohomologie \'etale, comme on l'a vu dans
Compléments sur la cohomologie des espaces, section
\ref{spaces-more-cohomology-section-fppf-etale}.

\begin{lemma}
\label{spaces-simplicial-lemma-compare-simplicial-objects-fppf-etale}
Soit $S$ un schéma. Soit $X$ un espace algébrique sur $S$.
Soit $U$ un espace algébrique simplicial sur $S$.
Soit $a : U \to X$ une augmentation. Il existe un diagramme commutatif
$$
\xymatrix{
\Sh((\textit{Espaces}/U)_{fppf, total}) \ar[r]_-h \ar[d]_{a_{fppf}} &
\Sh(U_\etale) \ar[d]^a \\
\Sh((\textit{Espaces}/X)_{fppf}) \ar[r]^-{h_{-1}} &
\Sh(X_\etale)
}
$$
où la flèche verticale gauche est définie dans la section
\ref{spaces-simplicial-section-hypercovering}
et la flèche verticale droite est définie dans la section
\ref{spaces-simplicial-section-simplicial-algebraic-spaces}.
\end{lemma}

\begin{proof}
La notation $(\textit{Espaces}/U)_{fppf, total}$ indique que
nous utilisons la construction de
la section \ref{spaces-simplicial-section-hypercovering}
pour le site $(\textit{Espaces}/S)_{fppf}$ et l'objet simplicial
$U$ de ce site\footnote{On pourrait aussi
utiliser la topologie \'etale; on noterait alors
$(\textit{Espaces}/U)_{\etale, total}$.}.
Nous emploierons les sites $X_{spaces, \etale}$ et $U_{spaces, \etale}$
pour les topoi du membre de droite; cela est permis,
voir la discussion de la section \ref{spaces-simplicial-section-simplicial-algebraic-spaces}.

\medskip\noindent
Observons que $(\textit{Espaces}/U)_{fppf, total}$ et
$U_{spaces, \etale}$
relèvent tous deux du cas (A) de la situation \ref{spaces-simplicial-situation-simplicial-site}.
Cela résulte immédiatement de la construction de
$U_\etale$ dans la section \ref{spaces-simplicial-section-simplicial-algebraic-spaces}
et découle du lemme \ref{spaces-simplicial-lemma-sr-when-fibre-products}
pour $(\textit{Espaces}/U)_{fppf, total}$.
Considérons ensuite les foncteurs
$U_{n, spaces, \etale} \to (\textit{Espaces}/U_n)_{fppf}$, $U \mapsto U/U_n$
et
$X_{spaces, \etale} \to (\textit{Espaces}/X)_{fppf}$, $U \mapsto U/X$.
Nous avons vu que ces foncteurs définissent des morphismes de sites dans
Compléments sur la cohomologie des espaces, section
\ref{spaces-more-cohomology-section-fppf-etale}
où ils étaient notés $a_{U_n} = \epsilon_{U_n} \circ \pi_{u_n}$
et $a_X = \epsilon_X \circ \pi_X$.
On obtient ainsi un morphisme de sites simpliciaux compatible avec les augmentations
comme dans la remarque \ref{spaces-simplicial-remark-morphism-augmentation-simplicial-sites}
et on peut appliquer le lemme
\ref{spaces-simplicial-lemma-morphism-augmentation-simplicial-sites} pour conclure.
\end{proof}

\begin{lemma}
\label{spaces-simplicial-lemma-descent-sheaves-for-fppf-hypercovering}
Soit $S$ un schéma. Soit $X$ un espace algébrique sur $S$.
Soit $U$ un espace algébrique simplicial sur $S$. Soit $a : U \to X$
une augmentation. Si $a : U \to X$ est un hyperrecouvrement fppf de $X$,
alors
$$
a^{-1} : \Sh(X_\etale) \to \Sh(U_\etale)
\quad\text{et}\quad
a^{-1} : \textit{Ab}(X_\etale) \to \textit{Ab}(U_\etale)
$$
sont pleinement fidèles, d'image essentielle les faisceaux cartésiens et
de quasi-inverse $a_*$. Ici, $a : \Sh(U_\etale) \to \Sh(X_\etale)$
est comme dans la section \ref{spaces-simplicial-section-simplicial-algebraic-spaces}.
\end{lemma}

\begin{proof}
Nous allons démontrer l'énoncé pour les faisceaux d'ensembles. Il découlera
presque formellement de résultats déjà établis.
Considérons le diagramme du lemme
\ref{spaces-simplicial-lemma-compare-simplicial-objects-fppf-etale}.
Dans la preuve de ce lemme nous avons vu que
$h_{-1}$ est le morphisme $a_X$ de
Compléments sur la cohomologie des espaces, section
\ref{spaces-more-cohomology-section-fppf-etale}.
Il résulte donc de
Compléments sur la cohomologie des espaces, lemme
\ref{spaces-more-cohomology-lemma-comparison-fppf-etale}
que $(h_{-1})^{-1}$ est pleinement fidèle, de quasi-inverse $h_{-1, *}$.
Il en va de même pour les composantes $h_n$ de $h$.
D'après la description des foncteurs $h^{-1}$ et $h_*$ donnée dans le
lemme \ref{spaces-simplicial-lemma-morphism-simplicial-sites},
on conclut que $h^{-1}$ est pleinement fidèle, de quasi-inverse $h_*$.
Observons que $U$ est un hyperrecouvrement de $X$ dans $(\textit{Espaces}/S)_{fppf}$
au sens de la section \ref{spaces-simplicial-section-hypercovering}.
Par le lemme \ref{spaces-simplicial-lemma-hypercovering-X-simple-descent-sheaves}
on voit que $a_{fppf}^{-1}$ est pleinement fidèle, de quasi-inverse
$a_{fppf, *}$ et d'image essentielle les faisceaux cartésiens
sur $(\textit{Espaces}/U)_{fppf, total}$.
Une chasse formelle dans le diagramme montre alors que
$a^{-1}$ est pleinement fidèle.

\medskip\noindent
Enfin, supposons que $\mathcal{G}$ soit un faisceau cartésien sur $U_\etale$.
Alors $h^{-1}\mathcal{G}$ est un faisceau cartésien sur
$(\textit{Espaces}/U)_{fppf, total}$. Ainsi
$h^{-1}\mathcal{G} = a_{fppf}^{-1}\mathcal{H}$ pour un certain faisceau
$\mathcal{H}$ sur $(\textit{Espaces}/X)_{fppf}$.
En particulier, on obtient
$h_0^{-1}\mathcal{G}_0 = (a_{0, big, fppf})^{-1}\mathcal{H}$.
Comme $h_0 = a_{U_0}$ et que $U_0 \to X$ est
plat, localement de présentation finie et surjectif, le
lemme de Compléments sur la cohomologie des espaces
\ref{spaces-more-cohomology-lemma-descent-sheaf-fppf-etale} montre
qu'il existe un faisceau $\mathcal{F}$ sur $X_\etale$ et un isomorphisme
$\mathcal{H} = (h_{-1})^{-1}\mathcal{F}$.
Puisque $a_{fppf}^{-1}\mathcal{H} = h^{-1}\mathcal{G}$
on en déduit que $h^{-1}\mathcal{G} \cong h^{-1}a^{-1}\mathcal{F}$.
Par la pleine fidélité de $h^{-1}$, nous concluons que
$a^{-1}\mathcal{F} \cong \mathcal{G}$.

\medskip\noindent
Fixons un isomorphisme $\theta : a^{-1}\mathcal{F} \to \mathcal{G}$.
Pour terminer la preuve, il faut montrer $\mathcal{G} = a^{-1}a_*\mathcal{G}$
(afin d'établir que le quasi-inverse est donné par $a_*$; tout le reste
a été prouvé ci-dessus).
Puisque $a^{-1}$ est pleinement fidèle, on a $\text{id} \cong a_*a^{-1}$ d'après
Catégories, lemme \ref{categories-lemma-adjoint-fully-faithful}.
Ainsi $\mathcal{F} \cong a_*a^{-1}\mathcal{F}$ et
$a_*\theta : a_*a^{-1}\mathcal{F} \to a_*\mathcal{G}$
se combinent en un isomorphisme $\mathcal{F} \to a_*\mathcal{G}$.
En prenant l'image inverse par $a$ et en précomposant par $\theta^{-1}$,
on obtient l'isomorphisme souhaité.
\end{proof}

\begin{lemma}
\label{spaces-simplicial-lemma-cohomological-descent-for-fppf-hypercovering}
Soient $S$ un schéma et $X$ un espace algébrique sur $S$.
Soit $U$ un espace algébrique simplicial sur $S$. Soit $a : U \to X$
une augmentation. Si $a : U \to X$ est un hyperrecouvrement fppf de $X$,
alors pour $K \in D^+(X_\etale)$
$$
K \to Ra_*(a^{-1}K)
$$
est un isomorphisme. Ici, $a : \Sh(U_\etale) \to \Sh(X_\etale)$
est comme dans la section \ref{spaces-simplicial-section-simplicial-algebraic-spaces}.
\end{lemma}

\begin{proof}
Considérons le diagramme du lemme \ref{spaces-simplicial-lemma-compare-simplicial-objects-fppf-etale}.
Observons que $Rh_{n, *}h_n^{-1}$ est le foncteur identité
sur $D^+(U_{n, \etale})$ d'après
Compléments sur la cohomologie des espaces, lemme
\ref{spaces-more-cohomology-lemma-cohomological-descent-etale-fppf}.
Par conséquent, $Rh_*h^{-1}$ est le foncteur identité sur
$D^+(U_\etale)$ par
le lemme \ref{spaces-simplicial-lemma-direct-image-morphism-simplicial-sites}.
Nous avons
\begin{align*}
Ra_*(a^{-1}K)
& =
Ra_*Rh_*h^{-1}a^{-1}K \\
& =
Rh_{-1, *}Ra_{fppf, *}a_{fppf}^{-1}(h_{-1})^{-1}K \\
& =
Rh_{-1, *}(h_{-1})^{-1}K \\
& =
K
\end{align*}
La première égalité découle de la discussion ci-dessus; la deuxième,
de la commutativité du diagramme du lemme
\ref{spaces-simplicial-lemma-compare-simplicial-objects}; la troisième, du lemme
\ref{spaces-simplicial-lemma-hypercovering-X-simple-descent-bounded-abelian}
puisque $U$ est un hyperrecouvrement de $X$ dans $(\textit{Espaces}/S)_{fppf}$;
et la dernière, du lemme déjà utilisé de
Compléments sur la cohomologie des espaces
\ref{spaces-more-cohomology-lemma-cohomological-descent-etale-fppf}.
\end{proof}

\begin{lemma}
\label{spaces-simplicial-lemma-compute-via-fppf-hypercovering}
Soient $S$ un schéma et $X$ un espace algébrique sur $S$.
Soit $U$ un espace algébrique simplicial sur $S$. Soit $a : U \to X$
une augmentation. Si $a : U \to X$ est un hyperrecouvrement fppf de $X$, alors
$$
R\Gamma(X_\etale, K) = R\Gamma(U_\etale, a^{-1}K)
$$
pour $K \in D^+(X_\etale)$. Ici, $a : \Sh(U_\etale) \to \Sh(X_\etale)$
est comme dans la section \ref{spaces-simplicial-section-simplicial-algebraic-spaces}.
\end{lemma}

\begin{proof}
Cela découle du lemme
\ref{spaces-simplicial-lemma-cohomological-descent-for-fppf-hypercovering}
car $R\Gamma(U_\etale, -) = R\Gamma(X_\etale, -) \circ Ra_*$ par
Cohomologie sur les sites, remarque \ref{sites-cohomology-remark-before-Leray}.
\end{proof}

\begin{lemma}
\label{spaces-simplicial-lemma-fppf-hypercovering-equivalence-bounded}
Soient $S$ un schéma et $X$ un espace algébrique sur $S$.
Soit $U$ un espace algébrique simplicial sur $S$. Soit $a : U \to X$
une augmentation.
Soit $\mathcal{A} \subset \textit{Ab}(U_\etale)$
la sous-catégorie de Serre faible des faisceaux abéliens cartésiens.
Si $U$ est un hyperrecouvrement fppf de $X$, alors
le foncteur $a^{-1}$ définit une équivalence
$$
D^+(X_\etale) \longrightarrow D_\mathcal{A}^+(U_\etale)
$$
de quasi-inverse $Ra_*$. Ici, $a : \Sh(U_\etale) \to \Sh(X_\etale)$
est comme dans la section \ref{spaces-simplicial-section-simplicial-algebraic-spaces}.
\end{lemma}

\begin{proof}
Observons que $\mathcal{A}$ est une sous-catégorie de Serre faible d'après le
lemme \ref{spaces-simplicial-lemma-Serre-subcat-cartesian-modules}.
L'équivalence est une conséquence formelle
des résultats obtenus jusqu'à présent. On utilise les lemmes
\ref{spaces-simplicial-lemma-descent-sheaves-for-fppf-hypercovering} et
\ref{spaces-simplicial-lemma-cohomological-descent-for-fppf-hypercovering}, ainsi que le lemme de Cohomologie
sur les sites \ref{sites-cohomology-lemma-equivalence-bounded}.
\end{proof}

\begin{lemma}
\label{spaces-simplicial-lemma-spectral-sequence-fppf-hypercovering}
Soient $S$ un schéma et $X$ un espace algébrique sur $S$.
Soit $U$ un espace algébrique simplicial sur $S$. Soit $a : U \to X$
une augmentation. Soit $\mathcal{F}$ un faisceau abélien
sur $X_\etale$. Notons $\mathcal{F}_n$ son image inverse sur $U_{n, \etale}$.
Si $U$ est un hyperrecouvrement fppf de $X$, alors
il existe une suite spectrale canonique
$$
E_1^{p, q} = H^q_\etale(U_p, \mathcal{F}_p)
$$
convergeant vers $H^{p + q}_\etale(X, \mathcal{F})$.
\end{lemma}

\begin{proof}
Conséquence immédiate des lemmes \ref{spaces-simplicial-lemma-compute-via-fppf-hypercovering}
et \ref{spaces-simplicial-lemma-simplicial-sheaf-cohomology-site}.
\end{proof}








\section{Hyperrecouvrements fppf d'espaces algébriques : modules}
\label{spaces-simplicial-section-fppf-hypercovering-modules}

\noindent
Nous poursuivons l'étude de la descente (cohomologique) pour les hyperrecouvrements fppf,
entreprise dans la section \ref{spaces-simplicial-section-fppf-hypercovering},
en examinant ici le cas des faisceaux de modules.
Nous nous intéressons surtout aux modules quasi-cohérents,
pour lesquels une descente cohomologique non bornée est possible.

\begin{lemma}
\label{spaces-simplicial-lemma-compare-simplicial-objects-fppf-etale-modules}
Soit $S$ un schéma. Soit $X$ un espace algébrique sur $S$.
Soit $U$ un espace algébrique simplicial sur $S$.
Soit $a : U \to X$ une augmentation. Il existe un diagramme commutatif
$$
\xymatrix{
(\Sh((\textit{Espaces}/U)_{fppf, total}), \mathcal{O}_{big, total})
\ar[r]_-h \ar[d]_{a_{fppf}} &
(\Sh(U_\etale), \mathcal{O}_U) \ar[d]^a \\
(\Sh((\textit{Espaces}/X)_{fppf}), \mathcal{O}_{big}) \ar[r]^-{h_{-1}} &
(\Sh(X_\etale), \mathcal{O}_X)
}
$$
de topoi annelés, où la flèche verticale gauche est définie dans la section
\ref{spaces-simplicial-section-hypercovering-modules}
et la flèche verticale droite est définie dans la section
\ref{spaces-simplicial-section-simplicial-algebraic-spaces}.
\end{lemma}

\begin{proof}
Pour le diagramme sous-jacent de topoi, voir la discussion de
la preuve du lemme \ref{spaces-simplicial-lemma-compare-simplicial-objects-fppf-etale}.
Le faisceau $\mathcal{O}_U$ est le faisceau structural de l'espace algébrique simplicial
$U$, tel qu'il est défini dans la section
\ref{spaces-simplicial-section-simplicial-algebraic-spaces}.
Le faisceau $\mathcal{O}_X$ est le faisceau structural usuel de l'espace algébrique
$X$. Les faisceaux d'anneaux $\mathcal{O}_{big, total}$ et
$\mathcal{O}_{big}$ proviennent du faisceau structural sur
$(\textit{Espaces}/S)_{fppf}$ de la manière expliquée dans la section
\ref{spaces-simplicial-section-hypercovering-modules}
qui construit aussi $a_{fppf}$ comme morphisme de topoi annelés.
Les morphismes sur les composantes $h_n = a_{U_n}$ et $h_{-1} = a_X$
sont des morphismes de topoi annelés d'après
Compléments sur la cohomologie des espaces, section
\ref{spaces-more-cohomology-section-fppf-etale-modules}.
Enfin, puisque le foncteur continu
$u : U_{spaces, \etale} \to (\textit{Espaces}/U)_{fppf, total}$
utilisé pour définir $h$\footnote{Ce foncteur intervient dans la preuve du
lemme \ref{spaces-simplicial-lemma-compare-simplicial-objects-fppf-etale}
via une application du lemme
\ref{spaces-simplicial-lemma-morphism-augmentation-simplicial-sites}.}
est donné par $V/U_n \mapsto V/U_n$,
on voit que $h_*\mathcal{O}_{big, total} = \mathcal{O}_U$;
c'est ainsi que l'on munit $h$ de la structure d'un morphisme
de sites simpliciaux annelés comme dans
la remarque \ref{spaces-simplicial-remark-morphism-simplicial-sites-modules}.
On obtient alors un morphisme de topoi annelés $h$
par le lemme \ref{spaces-simplicial-lemma-morphism-simplicial-sites-modules}.
Notons que les morphismes $h_n$ coïncident bien
avec les morphismes $a_{U_n}$ décrits ci-dessus.
Nous omettons la vérification
que le diagramme est commutatif en tant que diagramme de topoi annelés
-- nous savons déjà qu'il est commutatif
en tant que diagramme de topoi).
\end{proof}

\begin{lemma}
\label{spaces-simplicial-lemma-descent-qcoh-for-fppf-hypercovering}
Soient $S$ un schéma et $X$ un espace algébrique sur $S$.
Soit $U$ un espace algébrique simplicial sur $S$, et soit $a : U \to X$
une augmentation. Si $a : U \to X$ est un hyperrecouvrement fppf de $X$,
alors
$$
a^* : \QCoh(\mathcal{O}_X) \to \QCoh(\mathcal{O}_U)
$$
est une équivalence, de quasi-inverse $a_*$.
Ici, $a : \Sh(U_\etale) \to \Sh(X_\etale)$
est comme dans la section \ref{spaces-simplicial-section-simplicial-algebraic-spaces}.
\end{lemma}

\begin{proof}
Considérons le diagramme du lemme
\ref{spaces-simplicial-lemma-compare-simplicial-objects-fppf-etale-modules}.
Dans la preuve de ce lemme nous avons vu que
$h_{-1}$ est le morphisme $a_X$ de
Compléments sur la cohomologie des espaces, section
\ref{spaces-more-cohomology-section-fppf-etale-modules}.
Il résulte donc de
Compléments sur la cohomologie des espaces, lemme
\ref{spaces-more-cohomology-lemma-review-quasi-coherent}
que
$$
(h_{-1})^* :
\QCoh(\mathcal{O}_X)
\longrightarrow
\QCoh(\mathcal{O}_{big})
$$
est une équivalence, de quasi-inverse $h_{-1, *}$.
Il en va de même pour les composantes $h_n$ de $h$.
Rappelons que $\QCoh(\mathcal{O}_U)$ et $\QCoh(\mathcal{O}_{big, total})$
sont formées des modules cartésiens dont les composantes sont quasi-cohérentes; voir le
lemme \ref{spaces-simplicial-lemma-quasi-coherent-sheaf}.
Puisque les foncteurs $h^*$ et $h_*$ du
lemme \ref{spaces-simplicial-lemma-morphism-simplicial-sites-modules}
coïncident avec les foncteurs $h_n^*$ et $h_{n, *}$ sur les composantes,
nous concluons que
$$
h^* :
\QCoh(\mathcal{O}_U)
\longrightarrow 
\QCoh(\mathcal{O}_{big, total})
$$
est une équivalence, de quasi-inverse $h_*$.
Observons que $U$ est un hyperrecouvrement de $X$ dans $(\textit{Espaces}/S)_{fppf}$
au sens de la section \ref{spaces-simplicial-section-hypercovering}.
D'après le lemme \ref{spaces-simplicial-lemma-hypercovering-X-simple-descent-modules},
$a_{fppf}^*$ est pleinement fidèle, de quasi-inverse
$a_{fppf, *}$, et son image essentielle est formée des
$\mathcal{O}_{fppf, total}$-modules cartésiens.
Ainsi, par la description de $\QCoh(\mathcal{O}_{big})$ et
$\QCoh(\mathcal{O}_{big, total})$ donnée par le lemme \ref{spaces-simplicial-lemma-quasi-coherent-sheaf},
nous obtenons une équivalence
$$
a_{fppf}^* :
\QCoh(\mathcal{O}_{big})
\longrightarrow
\QCoh(\mathcal{O}_{big, total})
$$
de quasi-inverse $a_{fppf, *}$.
Une chasse formelle dans le diagramme montre alors que
$a^*$ est pleinement fidèle sur $\QCoh(\mathcal{O}_X)$ et que son image est
contenue dans $\QCoh(\mathcal{O}_U)$.

\medskip\noindent
Enfin, supposons que $\mathcal{G}$ appartienne à $\QCoh(\mathcal{O}_U)$.
Alors $h^*\mathcal{G}$ appartient à $\QCoh(\mathcal{O}_{big, total})$.
Ainsi $h^*\mathcal{G} = a_{fppf}^*\mathcal{H}$ avec
$\mathcal{H} = a_{fppf, *}h^*\mathcal{G}$
dans $\QCoh(\mathcal{O}_{big})$ (voir ci-dessus).
On voit ensuite que $\mathcal{H} = (h_{-1})^*\mathcal{F}$
avec $\mathcal{F} = h_{-1, *}\mathcal{H}$ dans $\QCoh(\mathcal{O}_X)$.
En parcourant le diagramme, on en déduit que
$h^*\mathcal{G} \cong h^*a^*\mathcal{F}$.
Par la pleine fidélité de $h^*$, nous concluons que
$a^*\mathcal{F} \cong \mathcal{G}$.
Puisque $\mathcal{F} = h_{-1, *}a_{fppf, *}h^*\mathcal{G} =
a_*h_*h^*\mathcal{G} = a_*\mathcal{G}$, nous obtenons également
l'énoncé selon lequel le quasi-inverse est donné par $a_*$.
\end{proof}

\begin{lemma}
\label{spaces-simplicial-lemma-cohomological-descent-qcoh-for-fppf-hypercovering}
Soient $S$ un schéma et $X$ un espace algébrique sur $S$.
Soit $U$ un espace algébrique simplicial sur $S$. Soit $a : U \to X$
une augmentation. Si $a : U \to X$ est un hyperrecouvrement fppf de $X$,
alors, pour tout $\mathcal{F}$ qui est un $\mathcal{O}_X$-module quasi-cohérent,
le morphisme
$$
\mathcal{F} \to Ra_*(a^*\mathcal{F})
$$
est un isomorphisme. Ici, $a : \Sh(U_\etale) \to \Sh(X_\etale)$
est comme dans la section \ref{spaces-simplicial-section-simplicial-algebraic-spaces}.
\end{lemma}

\begin{proof}
Considérons le diagramme du lemme \ref{spaces-simplicial-lemma-compare-simplicial-objects-fppf-etale}.
Soit $\mathcal{F}_n = a_n^*\mathcal{F}$ la composante de degré $n$ de
$a^*\mathcal{F}$. C'est un $\mathcal{O}_{U_n}$-module quasi-cohérent.
Alors $\mathcal{F}_n = Rh_{n, *}h_n^*\mathcal{F}_n$
par Compléments sur la cohomologie des espaces, lemme
\ref{spaces-more-cohomology-lemma-cohomological-descent-etale-fppf-modules}.
Ainsi $a^*\mathcal{F} = Rh_*h^*a^*\mathcal{F}$ d'après le
lemme \ref{spaces-simplicial-lemma-direct-image-morphism-simplicial-sites-modules}.
Nous avons
\begin{align*}
Ra_*(a^*\mathcal{F})
& =
Ra_*Rh_*h^*a^*\mathcal{F} \\
& =
Rh_{-1, *}Ra_{fppf, *}a_{fppf}^*(h_{-1})^*\mathcal{F} \\
& =
Rh_{-1, *}(h_{-1})^*\mathcal{F} \\
& =
\mathcal{F}
\end{align*}
La première égalité découle de la discussion ci-dessus; la deuxième,
de la commutativité du diagramme du lemme
\ref{spaces-simplicial-lemma-compare-simplicial-objects}; la troisième, du lemme
\ref{spaces-simplicial-lemma-hypercovering-X-simple-descent-bounded-modules}
puisque $U$ est un hyperrecouvrement de $X$ dans $(\textit{Espaces}/S)_{fppf}$
et que $La_{fppf}^* = a_{fppf}^*$, car $a_{fppf}$ est plat
(c'est-à-dire que $a_{fppf}^{-1}\mathcal{O}_{big} = \mathcal{O}_{big, total}$;
voir la remarque \ref{spaces-simplicial-remark-augmentation-ringed}); enfin,
la dernière égalité découle du lemme déjà utilisé de
Compléments sur la cohomologie des espaces
\ref{spaces-more-cohomology-lemma-cohomological-descent-etale-fppf-modules}.
\end{proof}

\begin{lemma}
\label{spaces-simplicial-lemma-coh-descent-qcoh-unbounded-for-fppf-hypercovering}
Soient $S$ un schéma et $X$ un espace algébrique sur $S$.
Soit $U$ un espace algébrique simplicial sur $S$. Soit $a : U \to X$
une augmentation. Supposons que $a : U \to X$ soit un hyperrecouvrement fppf de $X$.
Alors $\QCoh(\mathcal{O}_U)$ est une sous-catégorie de Serre faible de
$\textit{Mod}(\mathcal{O}_U)$ et
$$
a^* : D_\QCoh(\mathcal{O}_X) \longrightarrow D_\QCoh(\mathcal{O}_U)
$$
est une équivalence de catégories, de quasi-inverse
$Ra_*$. Ici, $a : \Sh(U_\etale) \to \Sh(X_\etale)$
est comme dans la section \ref{spaces-simplicial-section-simplicial-algebraic-spaces}.
\end{lemma}

\begin{proof}
Observons d'abord que les morphismes $a_n : U_n \to X$ et $d^n_i : U_n \to U_{n - 1}$
sont plats, localement de présentation finie et surjectifs d'après
Hyperrecouvrements, remarque \ref{hypercovering-remark-P-covering}.

\medskip\noindent
Rappelons qu'un $\mathcal{O}_U$-module $\mathcal{F}$ est quasi-cohérent si et
seulement s'il est cartésien et si $\mathcal{F}_n$ est quasi-cohérent pour tout $n$.
Voir le lemme \ref{spaces-simplicial-lemma-quasi-coherent-sheaf}.
Par le lemme \ref{spaces-simplicial-lemma-Serre-subcat-cartesian-modules}
(et la platitude des morphismes $d^n_i : U_n \to U_{n - 1}$ établie ci-dessus)
les modules cartésiens forment une sous-catégorie de Serre faible de
$\textit{Mod}(\mathcal{O}_U)$. D'autre part,
$\QCoh(\mathcal{O}_{U_n}) \subset \textit{Mod}(\mathcal{O}_{U_n})$
est une sous-catégorie de Serre faible pour tout $n$
(Propriétés des espaces, lemme
\ref{spaces-properties-lemma-properties-quasi-coherent}).
Ces faits montrent ensemble que
$\QCoh(\mathcal{O}_U) \subset \textit{Mod}(\mathcal{O}_U)$
est une sous-catégorie de Serre faible.

\medskip\noindent
Pour achever la preuve, vérifions une à une les conditions (1) -- (5) du lemme
de Cohomologie sur les sites
\ref{sites-cohomology-lemma-equivalence-unbounded-one}.

\medskip\noindent
Pour (1). Comme $a_n$ est plat (voir ci-dessus), $a$ est plat
par le lemme \ref{spaces-simplicial-lemma-flat-augmentation-modules}.

\medskip\noindent
Pour (2). C'est exactement l'énoncé du
lemme \ref{spaces-simplicial-lemma-descent-qcoh-for-fppf-hypercovering}.

\medskip\noindent
Pour (3). C'est exactement l'énoncé du
lemme \ref{spaces-simplicial-lemma-cohomological-descent-qcoh-for-fppf-hypercovering}.

\medskip\noindent
Pour (4). Rappelons que l'on peut utiliser soit le site $U_\etale$, soit le site
$U_{spaces, \etale}$ pour définir le petit topos \'etale
$\Sh(U_\etale)$; voir la section \ref{spaces-simplicial-section-simplicial-algebraic-spaces}.
L'hypothèse de la situation de Cohomologie
sur les sites \ref{sites-cohomology-situation-olsson-laszlo}
est satisfaite pour le triplet
$(U_{spaces, \etale}, \mathcal{O}_U, \QCoh(\mathcal{O}_U))$
et, par le même raisonnement, pour le triplet
$(U_\etale, \mathcal{O}_U, \QCoh(\mathcal{O}_U))$.
Plus précisément, prenons pour
$$
\mathcal{B} \subset \Ob(U_\etale) \subset \Ob(U_{spaces, \etale})
$$
l'ensemble des objets affines. Pour $V/U_n \in \mathcal{B}$,
prenons $d_{V/U_n} = 0$ et pour $\text{Cov}_{V/U_n}$ l'ensemble des
recouvrements étales $\{V_i \to V\}$ tels que les $V_i$ soient affines.
On obtient alors l'annulation voulue, car pour
$\mathcal{F} \in \QCoh(\mathcal{O}_U)$
et tout $V/U_n \in \mathcal{B}$, nous avons
$$
H^p(V/U_n, \mathcal{F}) = H^p(V, \mathcal{F}_n)
$$
d'après le lemme \ref{spaces-simplicial-lemma-sanity-check-modules}. Dans le
membre de droite figure la cohomologie du faisceau quasi-cohérent
$\mathcal{F}_n$ sur $U_n$, évaluée sur l'objet affine $V$
de $U_{n, \etale}$. Ce groupe s'annule pour $p > 0$ d'après
Cohomologie des espaces, section
\ref{spaces-cohomology-section-higher-direct-image} et Cohomologie des schémas,
lemme
\ref{coherent-lemma-quasi-coherent-affine-cohomology-zero}.

\medskip\noindent
Pour (5). Il suffit de prendre pour $\mathcal{B} \subset \Ob(X_{spaces, \etale})$
l'ensemble des objets affines et d'utiliser les références données ci-dessus.
\end{proof}

\begin{lemma}
\label{spaces-simplicial-lemma-compute-via-fppf-hypercovering-modules}
Soit $S$ un schéma. Soit $X$ un espace algébrique sur $S$.
Soit $U$ un espace algébrique simplicial sur $S$. Soit $a : U \to X$
une augmentation. Si $a : U \to X$ est un hyperrecouvrement fppf de $X$, alors
$$
R\Gamma(X_\etale, K) = R\Gamma(U_\etale, a^*K)
$$
pour $K \in D_\QCoh(\mathcal{O}_X)$. Ici, $a : \Sh(U_\etale) \to \Sh(X_\etale)$
est comme dans la section \ref{spaces-simplicial-section-simplicial-algebraic-spaces}.
\end{lemma}

\begin{proof}
Cela découle du lemme
\ref{spaces-simplicial-lemma-coh-descent-qcoh-unbounded-for-fppf-hypercovering}
car $R\Gamma(U_\etale, -) = R\Gamma(X_\etale, -) \circ Ra_*$ par
Cohomologie sur les sites, remarque \ref{sites-cohomology-remark-before-Leray}.
\end{proof}

\begin{lemma}
\label{spaces-simplicial-lemma-spectral-sequence-fppf-hypercovering-modules}
Soient $S$ un schéma et $X$ un espace algébrique sur $S$.
Soit $U$ un espace algébrique simplicial sur $S$. Soit $a : U \to X$
une augmentation. Soit $\mathcal{F}$ un
$\mathcal{O}_X$-module quasi-cohérent. Notons $\mathcal{F}_n$ son image inverse sur
$U_{n, \etale}$. Si $U$ est un hyperrecouvrement fppf de $X$, alors
il existe une suite spectrale canonique
$$
E_1^{p, q} = H^q_\etale(U_p, \mathcal{F}_p)
$$
convergeant vers $H^{p + q}_\etale(X, \mathcal{F})$.
\end{lemma}

\begin{proof}
Conséquence immédiate de
les lemmes \ref{spaces-simplicial-lemma-compute-via-fppf-hypercovering-modules}
et \ref{spaces-simplicial-lemma-simplicial-module-cohomology-site}.
\end{proof}







\section{Descente fppf de complexes}
\label{spaces-simplicial-section-fppf-descent-derived}

\noindent
Dans cette section, on rassemble certains résultats établis plus haut
pour les recouvrements fppf d'espaces algébriques
et les catégories dérivées de modules quasi-cohérents.

\begin{lemma}
\label{spaces-simplicial-lemma-fppf-neg-ext-zero-hom}
Soit $X$ un espace algébrique sur un schéma $S$.
Soient $K, E \in D_\QCoh(\mathcal{O}_X)$.
Soit $a : U \to X$ un hyperrecouvrement fppf.
Supposons que, pour tout $n \geq 0$, on ait
$$
\Ext_{\mathcal{O}_{U_n}}^i(La_n^*K, La_n^*E) = 0
\text{ pour } i < 0
$$
Alors :
\begin{enumerate}
\item $\Ext_{\mathcal{O}_X}^i(K, E) = 0$ pour $i < 0$ et
\item il existe une suite exacte
$$
0
\to
\Hom_{\mathcal{O}_X}(K, E)
\to
\Hom_{\mathcal{O}_{U_0}}(La_0^*K, La_0^*E)
\to
\Hom_{\mathcal{O}_{U_1}}(La_1^*K, La_1^*E)
$$
\end{enumerate}
\end{lemma}

\begin{proof}
Posons $K_n = La_n^*K$ et $E_n = La_n^*E$. Ce sont alors les composantes des
systèmes simpliciaux de la catégorie dérivée de modules
(définition \ref{spaces-simplicial-definition-cartesian-derived-modules})
associés à $La^*K$ et $La^*E$
(lemme \ref{spaces-simplicial-lemma-cartesian-objects-derived-modules}),
où $a : U_\etale \to X_\etale$ est comme dans la section
\ref{spaces-simplicial-section-simplicial-algebraic-spaces}.
Démontrons d'abord (2). D'après le
lemme \ref{spaces-simplicial-lemma-coh-descent-qcoh-unbounded-for-fppf-hypercovering},
nous avons
$$
\Hom_{\mathcal{O}_X}(K, E) =
\Hom_{\mathcal{O}_U}(La^*K, La^*E)
$$
La suite prend donc la forme suivante :
$$
0
\to
\Hom_{\mathcal{O}_U}(La^*K, La^*E)
\to
\Hom_{\mathcal{O}_{U_0}}(K_0, E_0)
\to
\Hom_{\mathcal{O}_{U_1}}(K_1, E_1)
$$
La première flèche est injective par
le lemme \ref{spaces-simplicial-lemma-nullity-cartesian-modules-derived}.
L'image de cette flèche est le noyau de la seconde
par le lemme \ref{spaces-simplicial-lemma-hom-cartesian-modules-derived}.
Cela termine la preuve de (2).
La partie (1) résulte de la partie (2) appliquée à
$K[i]$ et $E$ pour $i > 0$.
\end{proof}

\begin{lemma}
\label{spaces-simplicial-lemma-fppf-glue-neg-ext-zero}
Soit $X$ un espace algébrique sur un schéma $S$.
Soit $a : U \to X$ un hyperrecouvrement fppf.
Supposons que soient donnés $K_0 \in D_\QCoh(U_0)$ et un isomorphisme
$$
\alpha :
L(f_{\delta_1^1})^*K_0
\longrightarrow
L(f_{\delta_0^1})^*K_0
$$
satisfaisant à la condition de cocycle sur $U_1$. Posons
$\tau^n_i : [0] \to [n]$, $0 \mapsto i$ et
posons $K_n = Lf_{\tau^n_n}^*K_0$.
Supposons $\Ext^i_{\mathcal{O}_{U_n}}(K_n, K_n) = 0$ pour $i < 0$.
Il existe alors un objet $K \in D_\QCoh(\mathcal{O}_X)$
et un isomorphisme $La_0^*K \to K$ compatible avec $\alpha$.
\end{lemma}

\begin{proof}
Les objets $K_n$ forment les termes d'un système simplicial de la
catégorie dérivée de modules, d'après le
lemme \ref{spaces-simplicial-lemma-construct-cartesian-objects-derived-modules}.
On obtient alors un objet
$K' \in D_\QCoh(\mathcal{O}_{U_\etale})$
tel que $(K_n, K_\varphi)$ soit le système déduit de $K'$; voir le lemme
\ref{spaces-simplicial-lemma-cartesian-module-derived-from-simplicial}.
Enfin, appliquons le lemme
\ref{spaces-simplicial-lemma-coh-descent-qcoh-unbounded-for-fppf-hypercovering}
pour voir que $K' = La^*K$ pour un certain $K \in D_\QCoh(\mathcal{O}_X)$
ce qui conclut.
\end{proof}











\section{Hyperrecouvrements propres d'espaces algébriques}
\label{spaces-simplicial-section-proper-hypercovering-spaces}

\noindent
Cette section est l'analogue de la section \ref{spaces-simplicial-section-proper-hypercovering}
pour le cas des espaces algébriques.
Le lecteur qui le souhaite peut remplacer partout « espace algébrique »
par « schéma » et obtenir des résultats également valables.
Cela présente l'avantage de remplacer les références à
Compléments sur la cohomologie des espaces, section
\ref{spaces-more-cohomology-section-ph-etale}
par des références à
Cohomologie étale, section \ref{etale-cohomology-section-ph-etale}.

\medskip\noindent
Nous fixons un schéma de base $S$.
Soit $X$ un espace algébrique sur $S$ et soit $U$ un espace algébrique simplicial
sur $S$. Supposons que nous ayons une augmentation
$$
a : U \to X
$$
Voir la section \ref{spaces-simplicial-section-simplicial-algebraic-spaces}.
On dit que $U$ est un {\it hyperrecouvrement propre} de $X$ si
\begin{enumerate}
\item $U_0 \to X$ est propre et surjectif,
\item $U_1 \to U_0 \times_X U_0$ est propre et surjectif,
\item $U_{n + 1} \to (\text{cosk}_n\text{sk}_n U)_{n + 1}$
est propre et surjectif pour $n \geq 1$.
\end{enumerate}
La catégorie des espaces algébriques sur $S$ admet toutes les limites finies; les cosquelettes
utilisés dans la formulation ci-dessus existent donc.
$$
\fbox{Principe : les hyperrecouvrements propres permettent de
calculer la cohomologie étale.}
$$
L'idée essentielle de la démonstration consiste à comparer les topologies
ph et \'etale sur la catégorie $\textit{Espaces}/S$.
La topologie ph est en effet plus fine que la topologie \'etale, et l'on a :
(a) tout morphisme propre surjectif définit un recouvrement ph, et
(b) la cohomologie ph des faisceaux provenant par image inverse du petit site étale
coïncide avec la cohomologie étale, comme on l'a vu dans
Compléments sur la cohomologie des espaces, section
\ref{spaces-more-cohomology-section-ph-etale}.

\medskip\noindent
Tous les résultats de cette section se généralisent au cas
où $U \to X$ n'est qu'un « hyperrecouvrement ph », c'est-à-dire un hyperrecouvrement
de $X$ dans le site $(\textit{Espaces}/S)_{ph}$
au sens de la section \ref{spaces-simplicial-section-hypercovering}. Si nous avons besoin de ce résultat,
nous en donnerons ici un énoncé précis et une démonstration.

\begin{lemma}
\label{spaces-simplicial-lemma-compare-simplicial-objects-ph-etale}
Soit $S$ un schéma. Soit $X$ un espace algébrique sur $S$.
Soit $U$ un espace algébrique simplicial sur $S$.
Soit $a : U \to X$ une augmentation. Il existe un diagramme commutatif
$$
\xymatrix{
\Sh((\textit{Espaces}/U)_{ph, total}) \ar[r]_-h \ar[d]_{a_{ph}} &
\Sh(U_\etale) \ar[d]^a \\
\Sh((\textit{Espaces}/X)_{ph}) \ar[r]^-{h_{-1}} &
\Sh(X_\etale)
}
$$
où la flèche verticale gauche est définie dans la section
\ref{spaces-simplicial-section-hypercovering}
et la flèche verticale droite est définie dans la section
\ref{spaces-simplicial-section-simplicial-algebraic-spaces}.
\end{lemma}

\begin{proof}
La notation $(\textit{Espaces}/U)_{ph, total}$ indique que
nous utilisons la construction de
la section \ref{spaces-simplicial-section-hypercovering}
pour le site $(\textit{Espaces}/S)_{ph}$ et l'objet simplicial
$U$ de ce site\footnote{Pour distinguer
$(\textit{Espaces}/U)_{fppf, total}$, défini à l'aide de la topologie fppf
dans la section \ref{spaces-simplicial-section-fppf-hypercovering}.}.
Nous emploierons les sites $X_{spaces, \etale}$ et $U_{spaces, \etale}$
pour les topoi du membre de droite; cela est permis,
voir la discussion de la section \ref{spaces-simplicial-section-simplicial-algebraic-spaces}.

\medskip\noindent
Observons que $(\textit{Espaces}/U)_{ph, total}$ et
$U_{spaces, \etale}$
relèvent tous deux du cas (A) de la situation \ref{spaces-simplicial-situation-simplicial-site}.
Cela résulte immédiatement de la construction de
$U_\etale$ dans la section \ref{spaces-simplicial-section-simplicial-algebraic-spaces}
et découle du lemme \ref{spaces-simplicial-lemma-sr-when-fibre-products}
pour $(\textit{Espaces}/U)_{ph, total}$.
Considérons ensuite les foncteurs
$U_{n, spaces, \etale} \to (\textit{Espaces}/U_n)_{ph}$, $U \mapsto U/U_n$
et
$X_{spaces, \etale} \to (\textit{Espaces}/X)_{ph}$, $U \mapsto U/X$.
Nous avons vu que ces foncteurs définissent des morphismes de sites dans
Compléments sur la cohomologie des espaces, section
\ref{spaces-more-cohomology-section-ph-etale}
où ils étaient notés $a_{U_n} = \epsilon_{U_n} \circ \pi_{u_n}$
et $a_X = \epsilon_X \circ \pi_X$.
On obtient ainsi un morphisme de sites simpliciaux compatible avec les augmentations
comme dans la remarque \ref{spaces-simplicial-remark-morphism-augmentation-simplicial-sites}
et on peut appliquer le lemme
\ref{spaces-simplicial-lemma-morphism-augmentation-simplicial-sites} pour conclure.
\end{proof}

\begin{lemma}
\label{spaces-simplicial-lemma-descent-sheaves-for-ph-hypercovering}
Soit $S$ un schéma. Soit $X$ un espace algébrique sur $S$.
Soit $U$ un espace algébrique simplicial sur $S$. Soit $a : U \to X$
une augmentation. Si $a : U \to X$ est un hyperrecouvrement propre de $X$,
alors
$$
a^{-1} : \Sh(X_\etale) \to \Sh(U_\etale)
\quad\text{et}\quad
a^{-1} : \textit{Ab}(X_\etale) \to \textit{Ab}(U_\etale)
$$
sont pleinement fidèles, d'image essentielle les faisceaux cartésiens et
de quasi-inverse $a_*$. Ici, $a : \Sh(U_\etale) \to \Sh(X_\etale)$
est comme dans la section \ref{spaces-simplicial-section-simplicial-algebraic-spaces}.
\end{lemma}

\begin{proof}
Nous allons démontrer l'énoncé pour les faisceaux d'ensembles. Il découlera
presque formellement de résultats déjà établis.
Considérons le diagramme du lemme
\ref{spaces-simplicial-lemma-compare-simplicial-objects-ph-etale}.
Dans la preuve de ce lemme nous avons vu que
$h_{-1}$ est le morphisme $a_X$ de
Compléments sur la cohomologie des espaces, section
\ref{spaces-more-cohomology-section-ph-etale}.
Il résulte donc de
Compléments sur la cohomologie des espaces, lemme
\ref{spaces-more-cohomology-lemma-comparison-ph-etale}
que $(h_{-1})^{-1}$ est pleinement fidèle, de quasi-inverse $h_{-1, *}$.
Il en va de même pour les composantes $h_n$ de $h$.
D'après la description des foncteurs $h^{-1}$ et $h_*$ donnée dans le
lemme \ref{spaces-simplicial-lemma-morphism-simplicial-sites},
on conclut que $h^{-1}$ est pleinement fidèle, de quasi-inverse $h_*$.
Observons que $U$ est un hyperrecouvrement de $X$ dans $(\textit{Espaces}/S)_{ph}$
au sens de la section \ref{spaces-simplicial-section-hypercovering}, puisqu'un morphisme propre
surjectif définit un recouvrement ph d'après Topologies sur les espaces, lemme
\ref{spaces-topologies-lemma-surjective-proper-ph}.
Par le lemme \ref{spaces-simplicial-lemma-hypercovering-X-simple-descent-sheaves}
on voit que $a_{ph}^{-1}$ est pleinement fidèle, de quasi-inverse
$a_{ph, *}$ et d'image essentielle les faisceaux cartésiens
sur $(\textit{Espaces}/U)_{ph, total}$.
Une chasse formelle dans le diagramme montre alors que
$a^{-1}$ est pleinement fidèle.

\medskip\noindent
Enfin, supposons que $\mathcal{G}$ soit un faisceau cartésien sur $U_\etale$.
Alors $h^{-1}\mathcal{G}$ est un faisceau cartésien sur
$(\textit{Espaces}/U)_{ph, total}$.
Ainsi $h^{-1}\mathcal{G} = a_{ph}^{-1}\mathcal{H}$ pour un certain faisceau
$\mathcal{H}$ sur $(\textit{Espaces}/X)_{ph}$.
Calculons, en employant des notations quelque peu pédantes :
\begin{align*}
(h_{-1})^{-1}(a_*\mathcal{G})
& =
(h_{-1})^{-1}
\text{Eq}(
\xymatrix{
a_{0, small, *}\mathcal{G}_0
\ar@<1ex>[r] \ar@<-1ex>[r] &
a_{1, small, *}\mathcal{G}_1
}
) \\
& =
\text{Eq}(
\xymatrix{
(h_{-1})^{-1}a_{0, small, *}\mathcal{G}_0
\ar@<1ex>[r] \ar@<-1ex>[r] &
(h_{-1})^{-1}a_{1, small, *}\mathcal{G}_1
}
) \\
& =
\text{Eq}(
\xymatrix{
a_{0, big, ph, *}h_0^{-1}\mathcal{G}_0
\ar@<1ex>[r] \ar@<-1ex>[r] &
a_{1, big, ph, *}h_1^{-1}\mathcal{G}_1
}
) \\
& =
\text{Eq}(
\xymatrix{
a_{0, big, ph, *}(a_{0, big, ph})^{-1}\mathcal{H}
\ar@<1ex>[r] \ar@<-1ex>[r] &
a_{1, big, ph, *}(a_{1, big, ph})^{-1}\mathcal{H}
}
) \\
& =
a_{ph, *}a_{ph}^{-1}\mathcal{H} \\
& =
\mathcal{H}
\end{align*}
La première égalité découle du lemme \ref{spaces-simplicial-lemma-augmentation-site};
la deuxième, de l'exactitude du foncteur $(h_{-1})^{-1}$;
la troisième, du
lemme de Compléments sur la cohomologie des espaces
\ref{spaces-more-cohomology-lemma-proper-push-pull-ph-etale}
(ici nous utilisons que $a_0 : U_0 \to X$ et $a_1: U_1 \to X$ sont propres);
la quatrième, de $a_{ph}^{-1}\mathcal{H} = h^{-1}\mathcal{G}$;
la cinquième, du lemme \ref{spaces-simplicial-lemma-augmentation-site}; enfin la sixième
a été établie ci-dessus. Puisque $a_{ph}^{-1}\mathcal{H} = h^{-1}\mathcal{G}$,
on en déduit que $h^{-1}\mathcal{G} \cong h^{-1}a^{-1}a_*\mathcal{G}$,
ce qui termine la preuve par la pleine fidélité de $h^{-1}$.
\end{proof}

\begin{lemma}
\label{spaces-simplicial-lemma-cohomological-descent-for-ph-hypercovering}
Soient $S$ un schéma et $X$ un espace algébrique sur $S$.
Soit $U$ un espace algébrique simplicial sur $S$. Soit $a : U \to X$
une augmentation. Si $a : U \to X$ est un hyperrecouvrement propre de $X$,
alors pour $K \in D^+(X_\etale)$
$$
K \to Ra_*(a^{-1}K)
$$
est un isomorphisme. Ici, $a : \Sh(U_\etale) \to \Sh(X_\etale)$
est comme dans la section \ref{spaces-simplicial-section-simplicial-algebraic-spaces}.
\end{lemma}

\begin{proof}
Considérons le diagramme du lemme \ref{spaces-simplicial-lemma-compare-simplicial-objects-ph-etale}.
Observons que $Rh_{n, *}h_n^{-1}$ est le foncteur identité
sur $D^+(U_{n, \etale})$ par
Compléments sur la cohomologie des espaces, lemme
\ref{spaces-more-cohomology-lemma-cohomological-descent-etale-ph}.
Par conséquent, $Rh_*h^{-1}$ est le foncteur identité sur
$D^+(U_\etale)$ par
le lemme \ref{spaces-simplicial-lemma-direct-image-morphism-simplicial-sites}.
Nous avons
\begin{align*}
Ra_*(a^{-1}K)
& =
Ra_*Rh_*h^{-1}a^{-1}K \\
& =
Rh_{-1, *}Ra_{ph, *}a_{ph}^{-1}(h_{-1})^{-1}K \\
& =
Rh_{-1, *}(h_{-1})^{-1}K \\
& =
K
\end{align*}
La première égalité découle de la discussion ci-dessus; la deuxième,
de la commutativité du diagramme du lemme
\ref{spaces-simplicial-lemma-compare-simplicial-objects}; la troisième, du lemme
le lemme \ref{spaces-simplicial-lemma-hypercovering-X-simple-descent-bounded-abelian}
puisque $U$ est un hyperrecouvrement de $X$ dans $(\textit{Espaces}/S)_{ph}$
d'après Topologies sur les espaces, lemme
\ref{spaces-topologies-lemma-surjective-proper-ph},
et la dernière, du lemme déjà utilisé de
Compléments sur la cohomologie des espaces, lemme
\ref{spaces-more-cohomology-lemma-cohomological-descent-etale-ph}.
\end{proof}

\begin{lemma}
\label{spaces-simplicial-lemma-compute-via-ph-hypercovering}
Soient $S$ un schéma et $X$ un espace algébrique sur $S$.
Soit $U$ un espace algébrique simplicial sur $S$. Soit $a : U \to X$
une augmentation. Si $a : U \to X$ est un hyperrecouvrement propre de $X$, alors
$$
R\Gamma(X_\etale, K) = R\Gamma(U_\etale, a^{-1}K)
$$
pour $K \in D^+(X_\etale)$. Ici, $a : \Sh(U_\etale) \to \Sh(X_\etale)$
est comme dans la section \ref{spaces-simplicial-section-simplicial-algebraic-spaces}.
\end{lemma}

\begin{proof}
Cela découle du lemme
\ref{spaces-simplicial-lemma-cohomological-descent-for-ph-hypercovering}
car $R\Gamma(U_\etale, -) = R\Gamma(X_\etale, -) \circ Ra_*$ par
Cohomologie sur les sites, remarque \ref{sites-cohomology-remark-before-Leray}.
\end{proof}

\begin{lemma}
\label{spaces-simplicial-lemma-ph-hypercovering-equivalence-bounded}
Soient $S$ un schéma et $X$ un espace algébrique sur $S$.
Soit $U$ un espace algébrique simplicial sur $S$. Soit $a : U \to X$
une augmentation.
Soit $\mathcal{A} \subset \textit{Ab}(U_\etale)$
la sous-catégorie de Serre faible des faisceaux abéliens cartésiens.
Si $U$ est un hyperrecouvrement propre de $X$, alors
le foncteur $a^{-1}$ définit une équivalence
$$
D^+(X_\etale) \longrightarrow D_\mathcal{A}^+(U_\etale)
$$
de quasi-inverse $Ra_*$. Ici, $a : \Sh(U_\etale) \to \Sh(X_\etale)$
est comme dans la section \ref{spaces-simplicial-section-simplicial-algebraic-spaces}.
\end{lemma}

\begin{proof}
Observons que $\mathcal{A}$ est une sous-catégorie de Serre faible d'après le
lemme \ref{spaces-simplicial-lemma-Serre-subcat-cartesian-modules}.
L'équivalence est une conséquence formelle
des résultats obtenus jusqu'à présent. On utilise les lemmes
\ref{spaces-simplicial-lemma-descent-sheaves-for-ph-hypercovering} et
\ref{spaces-simplicial-lemma-cohomological-descent-for-ph-hypercovering}, ainsi que le lemme de Cohomologie
sur les sites \ref{sites-cohomology-lemma-equivalence-bounded}.
\end{proof}

\begin{lemma}
\label{spaces-simplicial-lemma-spectral-sequence-ph-hypercovering}
Soit $S$ un schéma. Soit $X$ un espace algébrique sur $S$.
Soit $U$ un espace algébrique simplicial sur $S$. Soit $a : U \to X$
une augmentation. Soit $\mathcal{F}$ un faisceau abélien
sur $X_\etale$. Notons $\mathcal{F}_n$ son image inverse sur $U_{n, \etale}$.
Si $U$ est un hyperrecouvrement ph de $X$, alors
il existe une suite spectrale canonique
$$
E_1^{p, q} = H^q_\etale(U_p, \mathcal{F}_p)
$$
convergeant vers $H^{p + q}_\etale(X, \mathcal{F})$.
\end{lemma}

\begin{proof}
Conséquence immédiate des lemmes \ref{spaces-simplicial-lemma-compute-via-ph-hypercovering}
et \ref{spaces-simplicial-lemma-simplicial-sheaf-cohomology-site}.
\end{proof}














% OK, commencer ici.
%

\chapter{Dualité pour les espaces}



\phantomsection
\label{spaces-duality-section-phantom}


\section{Introduction}
\label{spaces-duality-section-introduction}

\noindent
Ce chapitre est l'analogue du chapitre correspondant pour les
schémas, voir Dualité pour les schémas, section \ref{duality-section-introduction}.
Le développement est semblable à celui des articles
\cite{Neeman-Grothendieck}, \cite{LN},
\cite{Lipman-notes} et \cite{Neeman-improvement}.




\section{Complexes dualisants sur les espaces algébriques}
\label{spaces-duality-section-dualizing-spaces}

\noindent
Soit $U$ un schéma localement noethérien. Soit $\mathcal{O}_\etale$
le faisceau structural de $U$ sur le petit site \'etale de $U$.
Nous dirons qu'un objet $K \in D_\QCoh(\mathcal{O}_\etale)$ est
un complexe dualisant sur $U$ si $K = \epsilon^*(\omega_U^\bullet)$
pour un complexe dualisant $\omega_U^\bullet$ au sens de
Dualité pour les schémas, section \ref{duality-section-dualizing-schemes}.
Ici, $\epsilon^* : D_\QCoh(\mathcal{O}_U) \to D_\QCoh(\mathcal{O}_\etale)$
est l'équivalence de Catégories dérivées des espaces, Lemme
\ref{spaces-perfect-lemma-derived-quasi-coherent-small-etale-site}.
La plupart des propriétés de $\omega_U^\bullet$ étudiées dans
Dualité pour les schémas, section \ref{duality-section-dualizing-schemes},
se transmettent à $K$ grâce à la discussion de
Catégories dérivées des espaces, sections
\ref{spaces-perfect-section-derived-quasi-coherent-etale} et
\ref{spaces-perfect-section-spell-out}.

\medskip\noindent
Nous définissons un complexe dualisant sur un espace algébrique localement noethérien
comme un complexe qui, localement pour la topologie \'etale, provient d'un complexe dualisant
sur le schéma correspondant.

\begin{lemma}
\label{spaces-duality-lemma-equivalent-definitions}
Soit $S$ un schéma. Soit $X$ un espace algébrique localement noethérien sur $S$.
Soit $K$ un objet de $D_\QCoh(\mathcal{O}_X)$. Les conditions suivantes sont équivalentes :
\begin{enumerate}
\item Pour tout morphisme \'etale $U \to X$, où $U$ est un schéma,
la restriction $K|_U$ est un complexe dualisant sur $U$ (au sens expliqué ci-dessus).
\item Il existe un morphisme \'etale surjectif $U \to X$, où $U$ est un
schéma, tel que $K|_U$ soit un complexe dualisant sur $U$.
\end{enumerate}
\end{lemma}

\begin{proof}
Supposons que $U \to X$ soit \'etale surjectif, où $U$ est un schéma.
Soit $V \to X$ un morphisme \'etale, où $V$ est un schéma.
Alors
$$
U \leftarrow U \times_X V \rightarrow V
$$
sont des morphismes \'etales de schémas, celui qui aboutit à $V$ étant surjectif.
Nous pouvons donc appliquer Dualité pour les schémas, Lemme \ref{duality-lemma-descent-ascent},
pour voir que, si $K|_U$ est un complexe dualisant sur $U$, alors
$K|_V$ est un complexe dualisant sur $V$.
\end{proof}

\begin{definition}
\label{spaces-duality-definition-dualizing-scheme}
Soit $S$ un schéma.
Soit $X$ un espace algébrique localement noethérien sur $S$.
Un objet $K$ de $D_\QCoh(\mathcal{O}_X)$ est appelé un
{\it complexe dualisant} si $K$ satisfait aux conditions équivalentes du
Lemme \ref{spaces-duality-lemma-equivalent-definitions}.
\end{definition}

\begin{lemma}
\label{spaces-duality-lemma-affine-duality}
Soit $A$ un anneau noethérien et soit $X = \Spec(A)$. Soit
$\mathcal{O}_\etale$ le faisceau structural de $X$ sur le
petit site \'etale de $X$. Soient $K, L$ des objets de $D(A)$.
Si $K \in D_{\textit{Coh}}(A)$ et si $L$ est de dimension injective
finie, alors
$$
\epsilon^*\widetilde{R\Hom_A(K, L)} =
R\SheafHom_{\mathcal{O}_\etale}(\epsilon^*\widetilde{K},
\epsilon^*\widetilde{L})
$$
dans $D(\mathcal{O}_\etale)$, où
$\epsilon : (X_\etale, \mathcal{O}_\etale) \to (X, \mathcal{O}_X)$
est comme dans Catégories dérivées des espaces, section
\ref{spaces-perfect-section-derived-quasi-coherent-etale}.
\end{lemma}

\begin{proof}
D'après Dualité pour les schémas, Lemme \ref{duality-lemma-affine-duality},
nous avons un isomorphisme canonique
$$
\widetilde{R\Hom_A(K, L)} =
R\SheafHom_{\mathcal{O}_X}(\widetilde{K}, \widetilde{L})
$$
dans $D(\mathcal{O}_X)$. Il existe un morphisme canonique
$$
\epsilon^*R\Hom_{\mathcal{O}_X}(\widetilde{K}, \widetilde{L})
\longrightarrow
R\SheafHom_{\mathcal{O}_\etale}(\epsilon^*\widetilde{K},
\epsilon^*\widetilde{L})
$$
dans $D(\mathcal{O}_\etale)$, voir Cohomologie sur les sites, Remarque
\ref{sites-cohomology-remark-prepare-fancy-base-change}.
Nous allons montrer que les membres de gauche et de droite de ce morphisme
ont des faisceaux de cohomologie isomorphes, mais nous omettrons de
vérifier que l'isomorphisme est donné par ce morphisme.

\medskip\noindent
Nous pouvons supposer que $L$ est donné par un complexe fini $I^\bullet$
de $A$-modules injectifs. Par récurrence sur la longueur de $I^\bullet$
et par compatibilité des constructions avec les triangles distingués,
nous nous ramenons au cas où $L = I[0]$, avec $I$ un $A$-module injectif.
Rappelons que les faisceaux de cohomologie de
$R\SheafHom_{\mathcal{O}_\etale}(\epsilon^*\widetilde{K},
\epsilon^*\widetilde{L}))$
sont les faisceaux associés au préfaisceau qui, à $U$ \'etale
sur $X$, associe le $i$-ème groupe d'Ext entre les restrictions de
$\epsilon^*\widetilde{K}$ et $\epsilon^*\widetilde{L}$
à $U_\etale$. Voir
Cohomologie sur les sites, Lemme
\ref{sites-cohomology-lemma-section-RHom-over-U}.
Si $U = \Spec(B)$ est affine, ce groupe d'Ext
est égal à $\text{Ext}^i_B(K \otimes_A B, L \otimes_A B)$
par l'équivalence de
Catégories dérivées des espaces, Lemme
\ref{spaces-perfect-lemma-derived-quasi-coherent-small-etale-site}, et
Catégories dérivées des schémas, Lemme
\ref{perfect-lemma-affine-compare-bounded}
(on utilise aussi les compatibilités détaillées dans
Catégories dérivées des espaces, Remarque
\ref{spaces-perfect-remark-match-total-direct-images}).
Puisque $A \to B$ est \'etale, nous voyons que
$I \otimes_A B$ est un $B$-module injectif
par Complexes dualisants, Lemme \ref{dualizing-lemma-injective-goes-up}.
Nous voyons donc que
\begin{align*}
\Ext^n_B(K \otimes_A B, I \otimes_A B)
& =
\Hom_B(H^{-n}(K \otimes_A B), I \otimes_A B) \\
& =
\Hom_{A_f}(H^{-n}(K) \otimes_A B, I \otimes_A B) \\
& =
\Hom_A(H^{-n}(K), I) \otimes_A B \\
& =
\text{Ext}^n_A(K, I) \otimes_A B
\end{align*}
L'avant-dernière égalité vaut parce que $H^{-n}(K)$ est un $A$-module fini, voir
Compléments d'algèbre, Lemme
\ref{more-algebra-lemma-pseudo-coherence-and-base-change-ext}.
Par conséquent, les faisceaux de cohomologie des membres de gauche et de droite
de l'égalité du lemme sont les mêmes.
\end{proof}

\begin{lemma}
\label{spaces-duality-lemma-dualizing-spaces}
Soit $S$ un schéma. Soit $X$ un espace algébrique localement noethérien sur $S$.
Soit $K$ un complexe dualisant sur $X$.
Alors $K$ est un objet de $D_{\textit{Coh}}(\mathcal{O}_X)$
et $D = R\SheafHom_{\mathcal{O}_X}(-, K)$ induit une anti-équivalence
$$
D :
D_{\textit{Coh}}(\mathcal{O}_X)
\longrightarrow
D_{\textit{Coh}}(\mathcal{O}_X)
$$
munie d'un isomorphisme canonique
$\text{id} \to D \circ D$. Si $X$ est quasi-compact, alors
$D$ échange $D^+_{\textit{Coh}}(\mathcal{O}_X)$ et
$D^-_{\textit{Coh}}(\mathcal{O}_X)$ et induit une équivalence
$D^b_{\textit{Coh}}(\mathcal{O}_X) \to D^b_{\textit{Coh}}(\mathcal{O}_X)$.
\end{lemma}

\begin{proof}
Soit $U \to X$ un morphisme \'etale, avec $U$ affine. Écrivons $U = \Spec(A)$ et
soit $\omega_A^\bullet$ un complexe dualisant pour $A$ correspondant à $K|_U$
comme dans le Lemme \ref{spaces-duality-lemma-equivalent-definitions} et
Dualité pour les schémas, Lemme \ref{duality-lemma-equivalent-definitions}.
D'après le Lemme \ref{spaces-duality-lemma-affine-duality}, le diagramme
$$
\xymatrix{
D_{\textit{Coh}}(A) \ar[r] \ar[d]_{R\Hom_A(-, \omega_A^\bullet)} &
D_{\textit{Coh}}(\mathcal{O}_\etale)
\ar[d]^{R\SheafHom_{\mathcal{O}_\etale}(-, K|_U)} \\
D_{\textit{Coh}}(A) \ar[r] &
D(\mathcal{O}_\etale)
}
$$
est commutatif, où $\mathcal{O}_\etale$ est le faisceau structural du
petit site \'etale de $U$. Comme la formation de $R\SheafHom$ commute
à la restriction, nous concluons que $D$ envoie
$D_{\textit{Coh}}(\mathcal{O}_X)$ dans
$D_{\textit{Coh}}(\mathcal{O}_X)$. De plus, le morphisme canonique
$$
L \longrightarrow
R\SheafHom_{\mathcal{O}_X}(R\SheafHom_{\mathcal{O}_X}(L, K), K)
$$
(Cohomologie sur les sites, Lemme \ref{sites-cohomology-lemma-internal-hom-evaluate})
est un isomorphisme pour tout $L$ de $D_{\textit{Coh}}(\mathcal{O}_X)$,
car il en est ainsi au-dessus de tout $U$ comme ci-dessus, d'après
Complexes dualisants, Lemme \ref{dualizing-lemma-dualizing}.
L'assertion sur les propriétés de bornitude du foncteur $D$
dans le cas quasi-compact résulte également des assertions correspondantes de
Complexes dualisants, Lemme \ref{dualizing-lemma-dualizing}.
\end{proof}

\noindent
Soit $(\mathcal{C}, \mathcal{O})$ un site annelé. Rappelons
qu'un objet $L$ de $D(\mathcal{O})$ est {\it inversible}
s'il est un objet inversible pour la structure monoïdale symétrique
sur $D(\mathcal{O}_X)$ donnée par le produit tensoriel dérivé. Dans
Cohomologie sur les sites, Lemme \ref{sites-cohomology-lemma-invertible-derived},
nous avons vu que cela signifie que $L$ est parfait et que, si $(\mathcal{C}, \mathcal{O})$
est un site localement annelé, alors, pour tout objet $U$ de $\mathcal{C}$,
il existe un recouvrement $\{U_i \to U\}$ de $U$ dans $\mathcal{C}$
tel que $L|_{U_i} \cong \mathcal{O}_{U_i}[-n_i]$
pour certains entiers $n_i$.

\medskip\noindent
Soit $S$ un schéma et soit $X$ un espace algébrique sur $S$.
Si $L$ dans $D(\mathcal{O}_X)$ est inversible, il existe une
décomposition en réunion disjointe $X = \coprod_{n \in \mathbf{Z}} X_n$
telle que $L|_{X_n}$ soit un module inversible placé en degré $n$.
En particulier, il s'ensuit que $L = \bigoplus H^n(L)[-n]$,
ce qui donne un complexe bien défini de $\mathcal{O}_X$-modules
(à différentielles nulles) représentant $L$.

\begin{lemma}
\label{spaces-duality-lemma-dualizing-unique-spaces}
Soit $S$ un schéma.
Soit $X$ un espace algébrique localement noethérien sur $S$.
Si $K$ et $K'$ sont des complexes dualisants sur $X$, alors $K'$
est isomorphe à $K \otimes_{\mathcal{O}_X}^\mathbf{L} L$
pour un objet inversible $L$ de $D(\mathcal{O}_X)$.
\end{lemma}

\begin{proof}
Posons
$$
L = R\SheafHom_{\mathcal{O}_X}(K, K')
$$
C'est un objet inversible de $D(\mathcal{O}_X)$, car cette assertion est vraie
localement sur les affines. Utiliser le Lemme \ref{spaces-duality-lemma-affine-duality} et
Complexes dualisants, Lemme
\ref{dualizing-lemma-dualizing-unique}, ainsi que sa démonstration.
Le morphisme d'évaluation $L \otimes_{\mathcal{O}_X}^\mathbf{L} K \to K'$
est un isomorphisme pour la même raison.
\end{proof}

\begin{lemma}
\label{spaces-duality-lemma-dimension-function-scheme}
Soit $S$ un schéma. Soit $X$ un espace algébrique localement noethérien
et quasi-séparé sur $S$.
Soit $\omega_X^\bullet$ un complexe dualisant sur $X$. Alors la fonction sur $X$
$|X| \to \mathbf{Z}$ définie par
$$
x \longmapsto \delta(x)\text{ tel que }
\omega_{X, \overline{x}}^\bullet[-\delta(x)]
\text{ soit un complexe dualisant normalisé sur }
\mathcal{O}_{X, \overline{x}}
$$
est une fonction de dimension sur $|X|$.
\end{lemma}

\begin{proof}
Soit $U$ un schéma et soit $U \to X$ un morphisme \'etale surjectif.
Soit $\omega_U^\bullet$ le complexe dualisant sur $U$ associé
à $\omega_X^\bullet|_U$.
Si $u \in U$ s'envoie sur $x \in |X|$, alors $\mathcal{O}_{X, \overline{x}}$
est l'hensélisé strict de $\mathcal{O}_{U, u}$. D'après
Complexes dualisants, Lemme \ref{dualizing-lemma-flat-unramified},
nous voyons que, si $\omega^\bullet$ est un complexe dualisant normalisé
pour $\mathcal{O}_{U, u}$, alors
$\omega^\bullet \otimes_{\mathcal{O}_{U, u}} \mathcal{O}_{X, \overline{x}}$
est un complexe dualisant normalisé pour $\mathcal{O}_{X, \overline{x}}$.
Nous voyons donc que la fonction de dimension $U \to \mathbf{Z}$ de
Dualité pour les schémas, Lemme \ref{duality-lemma-dimension-function-scheme},
pour le schéma $U$ et le complexe
$\omega_U^\bullet$ est égale au composé de $U \to |X|$ et de $\delta$.
En utilisant le fait que les spécialisations dans $|X|$ se relèvent en des spécialisations dans $U$
et que les spécialisations non triviales dans $U$ s'envoient sur des
spécialisations non triviales dans $X$
(Espaces raisonnables, Lemmes \ref{decent-spaces-lemma-decent-specialization} et
\ref{decent-spaces-lemma-decent-no-specializations-map-to-same-point}),
un argument topologique facile montre que $\delta$ est une fonction de dimension
sur $|X|$.
\end{proof}




 
\section{Adjoint à droite de l'image directe}
\label{spaces-duality-section-twisted-inverse-image}

\noindent
C'est l'analogue de Dualité pour les schémas, section
\ref{duality-section-twisted-inverse-image}.

\begin{lemma}
\label{spaces-duality-lemma-twisted-inverse-image}
\begin{reference}
C'est presque le même résultat que \cite[exemple 4.2]{Neeman-Grothendieck}.
\end{reference}
Soit $S$ un schéma.
Soit $f : X \to Y$ un morphisme entre espaces algébriques quasi-séparés et quasi-compacts
sur $S$. Le foncteur $Rf_* : D_\QCoh(X) \to D_\QCoh(Y)$
admet un adjoint à droite.
\end{lemma}

\begin{proof}
Nous allons démontrer l'existence d'un adjoint à droite en vérifiant les hypothèses de
Catégories dérivées, Proposition \ref{derived-proposition-brown}.
Tout d'abord, la catégorie $D_\QCoh(\mathcal{O}_X)$ admet les sommes directes, voir
Catégories dérivées des espaces, Lemme
\ref{spaces-perfect-lemma-quasi-coherence-direct-sums}.
La catégorie $D_\QCoh(\mathcal{O}_X)$ est compactement engendrée d'après
Catégories dérivées des espaces, Théorème
\ref{spaces-perfect-theorem-bondal-van-den-Bergh}.
Comme $X$ et $Y$ sont quasi-compacts et quasi-séparés, il en est de même de $f$, voir
Morphismes d'espaces, Lemmes
\ref{spaces-morphisms-lemma-compose-after-separated} et
\ref{spaces-morphisms-lemma-quasi-compact-permanence}.
Le foncteur $Rf_*$ commute donc aux sommes directes, voir
Catégories dérivées des espaces, Lemme
\ref{spaces-perfect-lemma-quasi-coherence-pushforward-direct-sums}.
Cela achève la démonstration.
\end{proof}

\begin{lemma}
\label{spaces-duality-lemma-twisted-inverse-image-bounded-below}
Notations et hypothèses comme dans le Lemme \ref{spaces-duality-lemma-twisted-inverse-image}.
Soit $a : D_\QCoh(\mathcal{O}_Y) \to D_\QCoh(\mathcal{O}_X)$ l'adjoint à droite
de $Rf_*$. Alors $a$ envoie
$D^+_\QCoh(\mathcal{O}_Y)$ dans $D^+_\QCoh(\mathcal{O}_X)$.
En fait, il existe un entier $N$ tel que
$H^i(K) = 0$ pour $i \leq c$ implique $H^i(a(K)) = 0$ pour $i \leq c - N$.
\end{lemma}

\begin{proof}
D'après Catégories dérivées des espaces, Lemme
\ref{spaces-perfect-lemma-quasi-coherence-direct-image},
le foncteur $Rf_*$ est de dimension cohomologique finie. Autrement dit,
il existe un entier $N$ tel que
$H^i(Rf_*L) = 0$ pour $i \geq N + c$ si $H^i(L) = 0$ pour $i \geq c$.
Soit $K \in D^+_\QCoh(\mathcal{O}_Y)$ tel que $H^i(K) = 0$ pour $i \leq c$.
Alors
$$
\Hom_{D(\mathcal{O}_X)}(\tau_{\leq c - N}a(K), a(K)) =
\Hom_{D(\mathcal{O}_Y)}(Rf_*\tau_{\leq c - N}a(K), K) = 0
$$
d'après ce qui précède. Cela implique manifestement que
$H^i(a(K)) = 0$ pour $i \leq c - N$.
\end{proof}

\noindent
Soit $S$ un schéma.
Soit $f : X \to Y$ un morphisme d'espaces algébriques quasi-séparés et quasi-compacts
sur $S$.
Notons $a$ l'adjoint à droite de
$Rf_* : D_\QCoh(\mathcal{O}_X) \to D_\QCoh(\mathcal{O}_Y)$. Pour tous
$K \in D_\QCoh(\mathcal{O}_Y)$ et $L \in D_\QCoh(\mathcal{O}_X)$,
nous obtenons un morphisme canonique
\begin{equation}
\label{spaces-duality-equation-sheafy-trace}
Rf_*R\SheafHom_{\mathcal{O}_X}(L, a(K))
\longrightarrow
R\SheafHom_{\mathcal{O}_Y}(Rf_*L, K)
\end{equation}
Ce morphisme est en effet défini comme le composé
$$
Rf_*R\SheafHom_{\mathcal{O}_X}(L, a(K)) \to
R\SheafHom_{\mathcal{O}_Y}(Rf_*L, Rf_*a(K)) \to
R\SheafHom_{\mathcal{O}_Y}(Rf_*L, K)
$$
où la première flèche est celle de
Cohomologie sur les sites, Remarque
\ref{sites-cohomology-remark-projection-formula-for-internal-hom},
et la seconde est la co-unité $Rf_*a(K) \to K$ de l'adjonction.

\begin{lemma}
\label{spaces-duality-lemma-iso-on-RSheafHom}
Soit $S$ un schéma.
Soit $f : X \to Y$ un morphisme d'espaces algébriques quasi-compacts et quasi-séparés
sur $S$.
Soit $a$ l'adjoint à droite de
$Rf_* : D_\QCoh(\mathcal{O}_X) \to D_\QCoh(\mathcal{O}_Y)$.
Soient $L \in D_\QCoh(\mathcal{O}_X)$ et
$K \in D_\QCoh(\mathcal{O}_Y)$.
Alors le morphisme (\ref{spaces-duality-equation-sheafy-trace})
$$
Rf_*R\SheafHom_{\mathcal{O}_X}(L, a(K))
\longrightarrow
R\SheafHom_{\mathcal{O}_Y}(Rf_*L, K)
$$
devient un isomorphisme après application du foncteur
$DQ_Y : D(\mathcal{O}_Y) \to D_\QCoh(\mathcal{O}_Y)$
décrit dans Catégories dérivées des espaces, section
\ref{spaces-perfect-section-better-coherator}.
\end{lemma}

\begin{proof}
L'énoncé a un sens puisque $DQ_Y$ existe d'après
Catégories dérivées des espaces, Lemme
\ref{spaces-perfect-lemma-better-coherator}.
Comme $DQ_Y$ est l'adjoint à droite du foncteur d'inclusion
$D_\QCoh(\mathcal{O}_Y) \to D(\mathcal{O}_Y)$,
pour démontrer le lemme, il suffit de montrer que, pour tout
$M \in D_\QCoh(\mathcal{O}_Y)$,
le morphisme (\ref{spaces-duality-equation-sheafy-trace}) induit une bijection
$$
\Hom_Y(M, Rf_*R\SheafHom_{\mathcal{O}_X}(L, a(K)))
\longrightarrow
\Hom_Y(M, R\SheafHom_{\mathcal{O}_Y}(Rf_*L, K))
$$
Pour le voir, nous utilisons la suite d'égalités suivante :
\begin{align*}
\Hom_Y(M, Rf_*R\SheafHom_{\mathcal{O}_X}(L, a(K)))
& =
\Hom_X(Lf^*M, R\SheafHom_{\mathcal{O}_X}(L, a(K))) \\
& =
\Hom_X(Lf^*M \otimes_{\mathcal{O}_X}^\mathbf{L} L, a(K)) \\
& =
\Hom_Y(Rf_*(Lf^*M \otimes_{\mathcal{O}_X}^\mathbf{L} L), K) \\
& =
\Hom_Y(M \otimes_{\mathcal{O}_Y}^\mathbf{L} Rf_*L, K) \\
& =
\Hom_Y(M, R\SheafHom_{\mathcal{O}_Y}(Rf_*L, K))
\end{align*}
La première égalité résulte de Cohomologie sur les sites, Lemme
\ref{sites-cohomology-lemma-adjoint}.
La deuxième résulte de Cohomologie sur les sites, Lemme
\ref{sites-cohomology-lemma-internal-hom}.
La troisième résulte de la construction de $a$.
La quatrième résulte de Catégories dérivées des espaces, Lemme
\ref{spaces-perfect-lemma-cohomology-base-change} (c'est l'étape importante).
La cinquième résulte de Cohomologie sur les sites, Lemme
\ref{sites-cohomology-lemma-internal-hom}.
\end{proof}

\begin{example}
\label{spaces-duality-example-iso-on-RSheafHom}
L'énoncé du Lemme \ref{spaces-duality-lemma-iso-on-RSheafHom} n'est pas vrai sans
appliquer le « cohérateur » $DQ_Y$. Voir
Dualité pour les schémas, Exemple \ref{duality-example-iso-on-RSheafHom}.
\end{example}

\begin{remark}
\label{spaces-duality-remark-iso-on-RSheafHom}
Dans la situation du Lemme \ref{spaces-duality-lemma-iso-on-RSheafHom}, nous avons
$$
DQ_Y(Rf_*R\SheafHom_{\mathcal{O}_X}(L, a(K))) =
Rf_* DQ_X(R\SheafHom_{\mathcal{O}_X}(L, a(K)))
$$
d'après Catégories dérivées des espaces, Lemme
\ref{spaces-perfect-lemma-pushforward-better-coherator}.
Ainsi, si $R\SheafHom_{\mathcal{O}_X}(L, a(K)) \in D_\QCoh(\mathcal{O}_X)$,
nous pouvons « supprimer » le $DQ_Y$ du membre de gauche du morphisme.
D'autre part, si nous savons que
$R\SheafHom_{\mathcal{O}_Y}(Rf_*L, K) \in D_\QCoh(\mathcal{O}_Y)$,
nous pouvons « supprimer » le $DQ_Y$ du membre de droite du morphisme.
Si ces deux conditions sont satisfaites, nous voyons que (\ref{spaces-duality-equation-sheafy-trace})
est un isomorphisme. En combinant ceci avec
Catégories dérivées des espaces, Lemme
\ref{spaces-perfect-lemma-quasi-coherence-internal-hom},
nous voyons que $Rf_*R\SheafHom_{\mathcal{O}_X}(L, a(K)) \to
R\SheafHom_{\mathcal{O}_Y}(Rf_*L, K)$ est un isomorphisme si
\begin{enumerate}
\item $L$ et $Rf_*L$ sont parfaits, ou
\item $K$ est borné inférieurement et $L$ et $Rf_*L$ sont pseudo-cohérents.
\end{enumerate}
Pour (2), nous utilisons le fait que $a(K)$ est borné inférieurement si $K$
est borné inférieurement, voir Lemme \ref{spaces-duality-lemma-twisted-inverse-image-bounded-below}.
\end{remark}

\begin{example}
\label{spaces-duality-example-iso-on-RSheafHom-noetherian}
Soit $S$ un schéma.
Soit $f : X \to Y$ un morphisme propre d'espaces algébriques noethériens
sur $S$, $L \in D^-_{\textit{Coh}}(X)$ et $K \in D^+_{\QCoh}(\mathcal{O}_Y)$.
Alors le morphisme $Rf_*R\SheafHom_{\mathcal{O}_X}(L, a(K)) \to
R\SheafHom_{\mathcal{O}_Y}(Rf_*L, K)$ est un isomorphisme.
En effet, les complexes $L$ et $Rf_*L$ sont pseudo-cohérents d'après
Catégories dérivées des espaces, Lemmes
\ref{spaces-perfect-lemma-identify-pseudo-coherent-noetherian} et
\ref{spaces-perfect-lemma-direct-image-coherent},
et la discussion de la Remarque \ref{spaces-duality-remark-iso-on-RSheafHom} s'applique.
\end{example}

\begin{lemma}
\label{spaces-duality-lemma-iso-global-hom}
Soit $S$ un schéma.
Soit $f : X \to Y$ un morphisme d'espaces algébriques quasi-séparés et quasi-compacts
sur $S$.
Pour tous $L \in D_\QCoh(\mathcal{O}_X)$ et $K \in D_\QCoh(\mathcal{O}_Y)$,
(\ref{spaces-duality-equation-sheafy-trace}) induit un isomorphisme
$R\Hom_X(L, a(K)) \to R\Hom_Y(Rf_*L, K)$ entre les Hom dérivés globaux.
\end{lemma}

\begin{proof}
Par construction (Cohomologie sur les sites, section
\ref{sites-cohomology-section-global-RHom}), les complexes
$$
R\Hom_X(L, a(K)) =
R\Gamma(X, R\SheafHom_{\mathcal{O}_X}(L, a(K))) =
R\Gamma(Y, Rf_*R\SheafHom_{\mathcal{O}_X}(L, a(K)))
$$
et
$$
R\Hom_Y(Rf_*L, K) = R\Gamma(Y, R\SheafHom_{\mathcal{O}_X}(Rf_*L, a(K)))
$$
Ainsi, le lemme est une conséquence du Lemme \ref{spaces-duality-lemma-iso-on-RSheafHom}.
En effet, un morphisme $E \to E'$ dans $D(\mathcal{O}_Y)$ qui induit
un isomorphisme $DQ_Y(E) \to DQ_Y(E')$ induit un quasi-isomorphisme
$R\Gamma(Y, E) \to R\Gamma(Y, E')$. Nous avons en effet
$H^i(Y, E) = \Ext^i_Y(\mathcal{O}_Y, E) = \Hom(\mathcal{O}_Y[-i], E) =
\Hom(\mathcal{O}_Y[-i], DQ_Y(E))$, car $\mathcal{O}_Y[-i]$
appartient à $D_\QCoh(\mathcal{O}_Y)$ et $DQ_Y$ est l'adjoint à droite
du foncteur d'inclusion $D_\QCoh(\mathcal{O}_Y) \to D(\mathcal{O}_Y)$.
\end{proof}







 
\section{Adjoint à droite de l'image directe et changement de base, I}
\label{spaces-duality-section-base-change-map}

\noindent
Définissons le morphisme de changement de base entre les adjoints à droite de l'image directe.
Soit $S$ un schéma. Considérons un diagramme cartésien
\begin{equation}
\label{spaces-duality-equation-base-change}
\vcenter{
\xymatrix{
X' \ar[r]_{g'} \ar[d]_{f'} & X \ar[d]^f \\
Y' \ar[r]^g & Y
}
}
\end{equation}
où $Y'$ et $X$ sont {\bf indépendants au sens de Tor} sur $Y$. Notons
$$
a  : D_\QCoh(\mathcal{O}_Y) \to D_\QCoh(\mathcal{O}_X)
\quad\text{et}\quad
a' : D_\QCoh(\mathcal{O}_{Y'}) \to D_\QCoh(\mathcal{O}_{X'})
$$
les adjoints à droite de $Rf_*$ et $Rf'_*$
(Lemme \ref{spaces-duality-lemma-twisted-inverse-image}).
Le morphisme de changement de base de
Cohomologie sur les sites, Remarque \ref{sites-cohomology-remark-base-change},
donne une transformation de foncteurs
$$
Lg^* \circ Rf_* \longrightarrow Rf'_* \circ L(g')^*
$$
sur les catégories dérivées de faisceaux à cohomologie quasi-cohérente.
Il donne donc une transformation entre les adjoints à droite en sens opposé
$$
a \circ Rg_* \longleftarrow Rg'_* \circ a'
$$

\begin{lemma}
\label{spaces-duality-lemma-flat-precompose-pus}
Dans le diagramme (\ref{spaces-duality-equation-base-change}), le morphisme
$a \circ Rg_* \leftarrow Rg'_* \circ a'$ est un isomorphisme.
\end{lemma}

\begin{proof}
Le morphisme de changement de base $Lg^* \circ Rf_* K \to Rf'_* \circ L(g')^*K$
est un isomorphisme pour tout $K$ de $D_\QCoh(\mathcal{O}_X)$ d'après
Catégories dérivées des espaces, Lemme
\ref{spaces-perfect-lemma-compare-base-change}
(on utilise ici l'hypothèse d'indépendance au sens de Tor).
La transformation correspondante entre foncteurs adjoints est donc
elle aussi un isomorphisme.
\end{proof}

\noindent
Nous pouvons alors considérer le
morphisme de foncteurs
$D_\QCoh(\mathcal{O}_Y) \to D_\QCoh(\mathcal{O}_{X'})$
donné par le composé
\begin{equation}
\label{spaces-duality-equation-base-change-map}
L(g')^* \circ a \to
L(g')^* \circ a \circ Rg_* \circ Lg^* \leftarrow
L(g')^* \circ Rg'_* \circ a' \circ Lg^* \to a' \circ Lg^*
\end{equation}
La première flèche provient du morphisme d'adjonction $\text{id} \to Rg_* Lg^*$,
et la dernière du morphisme d'adjonction $L(g')^*Rg'_* \to \text{id}$.
Nous avons besoin de l'hypothèse d'indépendance au sens de Tor pour inverser la flèche
du milieu, voir Lemme \ref{spaces-duality-lemma-flat-precompose-pus}.
On peut aussi considérer (\ref{spaces-duality-equation-base-change-map}), par
adjonction entre $L(g')^*$ et $R(g')_*$, comme une transformation naturelle
$$
a \to a \circ Rg_* \circ Lg^* \leftarrow Rg'_* \circ a' \circ Lg^*
$$
où, là encore, la seconde flèche est inversible. Si $M \in D_\QCoh(\mathcal{O}_X)$
et $K \in D_\QCoh(\mathcal{O}_Y)$,
alors ce morphisme est donné sur les foncteurs de Yoneda par
\begin{align*}
\Hom_X(M, a(K))
& =
\Hom_Y(Rf_*M, K) \\
& \to
\Hom_Y(Rf_*M, Rg_* Lg^*K) \\
& =
\Hom_{Y'}(Lg^*Rf_*M, Lg^*K) \\
& \leftarrow
\Hom_{Y'}(Rf'_* L(g')^*M, Lg^*K) \\
& =
\Hom_{X'}(L(g')^*M, a'(Lg^*K)) \\
& =
\Hom_X(M, Rg'_*a'(Lg^*K))
\end{align*}
(la flèche dirigée vers la gauche est inversible par le théorème de
changement de base donné dans
Catégories dérivées des espaces, Lemme
\ref{spaces-perfect-lemma-compare-base-change}),
ce qui rend les choses un peu plus explicites.

\medskip\noindent
Dans cette section, nous démontrons d'abord que le morphisme de changement de base satisfait
à certaines compatibilités naturelles relativement à l'empilement de carrés, comme dans
Cohomologie sur les sites, Remarques
\ref{sites-cohomology-remark-compose-base-change} et
\ref{sites-cohomology-remark-compose-base-change-horizontal},
pour le morphisme de changement de base usuel.
Nous suggérons au lecteur de sauter le reste de cette section en première lecture.

\begin{lemma}
\label{spaces-duality-lemma-compose-base-change-maps}
Soit $S$ un schéma. Considérons un diagramme commutatif
$$
\xymatrix{
X' \ar[r]_k \ar[d]_{f'} & X \ar[d]^f \\
Y' \ar[r]^l \ar[d]_{g'} & Y \ar[d]^g \\
Z' \ar[r]^m & Z
}
$$
d'espaces algébriques quasi-compacts et quasi-séparés sur $S$, dans lequel
les deux carrés sont cartésiens, et où $f$ et $l$
ainsi que $g$ et $m$ sont indépendants au sens de Tor.
Alors les morphismes (\ref{spaces-duality-equation-base-change-map})
des deux carrés se composent pour donner le morphisme de
changement de base du rectangle extérieur (voir la démonstration pour un énoncé précis).
\end{lemma}

\begin{proof}
Il résulte des hypothèses que $g \circ f$ et $m$ sont indépendants au sens de Tor
(nous omettons les détails), de sorte que l'énoncé a un sens.
Dans cette démonstration, nous écrivons $k^*$ à la place de $Lk^*$ et $f_*$ à la place
de $Rf_*$. Soient $a$, $b$ et $c$ les adjoints à droite du
Lemme \ref{spaces-duality-lemma-twisted-inverse-image}
pour $f$, $g$ et $g \circ f$, et de même pour les versions munies d'un prime.
La flèche correspondant au carré supérieur est le composé
$$
\gamma_{top} :
k^* \circ a \to k^* \circ a \circ l_* \circ l^*
\xleftarrow{\xi_{top}} k^* \circ k_* \circ a' \circ l^* \to a' \circ l^*
$$
où $\xi_{top} : k_* \circ a' \to a \circ l_*$
est un isomorphisme (et peut donc être inversé)
et est la flèche « duale » du morphisme de changement de base
$l^* \circ f_* \to f'_* \circ k^*$. Les flèches extérieures proviennent
des morphismes canoniques $1 \to l_* \circ l^*$ et $k^* \circ k_* \to 1$.
De même, pour le second carré, nous avons
$$
\gamma_{bot} :
l^* \circ b \to l^* \circ b \circ m_* \circ m^*
\xleftarrow{\xi_{bot}} l^* \circ l_* \circ b' \circ m^* \to b' \circ m^*
$$
Pour le rectangle extérieur, nous obtenons
$$
\gamma_{rect} :
k^* \circ c \to k^* \circ c \circ m_* \circ m^*
\xleftarrow{\xi_{rect}} k^* \circ k_* \circ c' \circ m^* \to c' \circ m^*
$$
Nous avons $(g \circ f)_* = g_* \circ f_*$ et donc
$c = a \circ b$, et de même $c' = a' \circ b'$.
L'assertion du lemme est que $\gamma_{rect}$
est égal au composé
$$
k^* \circ c = k^* \circ a \circ b \xrightarrow{\gamma_{top}}
a' \circ l^* \circ b \xrightarrow{\gamma_{bot}}
a' \circ b' \circ m^* = c' \circ m^*
$$
Pour le voir, considérons le diagramme suivant :
$$
\xymatrix{
& & k^* \circ a \circ b \ar[d] \ar[lldd] \\
& & k^* \circ a \circ l_* \circ l^* \circ b \ar[ld] \\
k^* \circ a \circ b \circ m_* \circ m^* \ar[r] &
k^* \circ a \circ l_* \circ l^* \circ b \circ m_* \circ m^* &
k^* \circ k_* \circ a' \circ l^* \circ b \ar[u]_{\xi_{top}} \ar[d] \ar[ld] \\
& k^*\circ k_* \circ a' \circ l^* \circ b \circ m_* \circ m^*
\ar[u]_{\xi_{top}} \ar[rd] &
a' \circ l^* \circ b \ar[d] \\
k^* \circ k_* \circ a' \circ b' \circ m^* \ar[uu]_{\xi_{rect}} \ar[ddrr] &
k^*\circ k_* \circ a' \circ l^* \circ l_* \circ b' \circ m^*
\ar[u]_{\xi_{bot}} \ar[l] \ar[dr] &
a' \circ l^* \circ b \circ m_* \circ m^* \\
& & a' \circ l^* \circ l_* \circ b' \circ m^* \ar[u]_{\xi_{bot}} \ar[d] \\
& & a' \circ b' \circ m^*
}
$$
En descendant du côté droit, nous obtenons le composé, et en descendant
du côté gauche, nous obtenons $\gamma_{rect}$.
Tous les quadrilatères du côté droit de ce diagramme sont commutatifs
d'après Catégories, Lemme \ref{categories-lemma-properties-2-cat-cats},
ou, plus simplement, d'après la discussion qui précède
Catégories, Définition \ref{categories-definition-horizontal-composition}.
Il suffit donc de montrer que le diagramme
$$
\xymatrix{
a \circ l_* \circ l^* \circ b \circ m_* &
a \circ b \circ m_* \ar[l] \\
k_* \circ a' \circ l^* \circ b \circ m_* \ar[u]_{\xi_{top}} & \\
k_* \circ a' \circ l^* \circ l_* \circ b' \ar[u]_{\xi_{bot}} \ar[r] &
k_* \circ a' \circ b' \ar[uu]_{\xi_{rect}}
}
$$
devient commutatif lorsque nous inversons les flèches $\xi_{top}$, $\xi_{bot}$
et $\xi_{rect}$ (noter que cela diffère du fait de demander que le
diagramme soit commutatif). Cependant, le diagramme
$$
\xymatrix{
& a \circ l_* \circ l^* \circ b \circ m_* \\
a \circ l_* \circ l^* \circ l_* \circ b'
\ar[ru]^{\xi_{bot}} & &
k_* \circ a' \circ l^* \circ b \circ m_* \ar[ul]_{\xi_{top}} \\
& k_* \circ a' \circ l^* \circ l_* \circ b'
\ar[ul]^{\xi_{top}} \ar[ur]_{\xi_{bot}}
}
$$
est commutatif d'après Catégories, Lemme \ref{categories-lemma-properties-2-cat-cats}.
Puisque les diagrammes
$$
\vcenter{
\xymatrix{
a \circ l_* \circ l^* \circ b \circ m_* & a \circ b \circ m \ar[l] \\
a \circ l_* \circ l^* \circ l_* \circ b' \ar[u] &
a \circ l_* \circ b' \ar[l] \ar[u]
}
}
\quad\text{et}\quad
\vcenter{
\xymatrix{
a \circ l_* \circ l^* \circ l_* \circ b' \ar[r] & a \circ l_* \circ b' \\
k_* \circ a' \circ l^* \circ l_* \circ b' \ar[u] \ar[r] &
k_* \circ a' \circ b' \ar[u]
}
}
$$
sont commutatifs (voir les références citées) et puisque le composé de
$l_* \to l_* \circ l^* \circ l_* \to l_*$ est l'identité,
nous trouvons qu'il suffit de démontrer que
$$
k \circ a' \circ b' \xrightarrow{\xi_{bot}} a \circ l_* \circ b
\xrightarrow{\xi_{top}} a \circ b \circ m_*
$$
est égal à $\xi_{rect}$ (via les identifications $a \circ b = c$
et $a' \circ b' = c'$). C'est l'assertion duale de
Cohomologie sur les sites, Remarque \ref{sites-cohomology-remark-compose-base-change},
ce qui achève la démonstration.
\end{proof}

\begin{lemma}
\label{spaces-duality-lemma-compose-base-change-maps-horizontal}
Soit $S$ un schéma. Considérons un diagramme commutatif
$$
\xymatrix{
X'' \ar[r]_{g'} \ar[d]_{f''} & X' \ar[r]_g \ar[d]_{f'} & X \ar[d]^f \\
Y'' \ar[r]^{h'} & Y' \ar[r]^h & Y
}
$$
d'espaces algébriques quasi-compacts et quasi-séparés sur $S$, dans lequel
les deux carrés sont cartésiens, et où $f$ et $h$
ainsi que $f'$ et $h'$ sont indépendants au sens de Tor.
Alors les morphismes (\ref{spaces-duality-equation-base-change-map})
des deux carrés se composent pour donner le morphisme de
changement de base du rectangle extérieur (voir la démonstration pour un énoncé précis).
\end{lemma}

\begin{proof}
Il résulte des hypothèses que $f$ et $h \circ h'$ sont indépendants au sens de Tor
(nous omettons les détails), de sorte que l'énoncé a un sens.
Dans cette démonstration, nous écrivons $g^*$ à la place de $Lg^*$ et $f_*$ à la place
de $Rf_*$. Soient $a$, $a'$ et $a''$ les adjoints à droite du
Lemme \ref{spaces-duality-lemma-twisted-inverse-image}
pour $f$, $f'$ et $f''$. La flèche correspondant au carré de droite
est le composé
$$
\gamma_{right} :
g^* \circ a \to g^* \circ a \circ h_* \circ h^*
\xleftarrow{\xi_{right}} g^* \circ g_* \circ a' \circ h^* \to a' \circ h^*
$$
où $\xi_{right} : g_* \circ a' \to a \circ h_*$
est un isomorphisme (et peut donc être inversé)
et est la flèche « duale » du morphisme de changement de base
$h^* \circ f_* \to f'_* \circ g^*$. Les flèches extérieures proviennent
des morphismes canoniques $1 \to h_* \circ h^*$ et $g^* \circ g_* \to 1$.
De même, pour le carré de gauche, nous avons
$$
\gamma_{left} :
(g')^* \circ a' \to (g')^* \circ a' \circ (h')_* \circ (h')^*
\xleftarrow{\xi_{left}}
(g')^* \circ (g')_* \circ a'' \circ (h')^* \to a'' \circ (h')^*
$$
Pour le rectangle extérieur, nous obtenons
$$
\gamma_{rect} :
k^* \circ a \to
k^* \circ a \circ m_* \circ m^* \xleftarrow{\xi_{rect}}
k^* \circ k_* \circ a'' \circ m^* \to
a'' \circ m^*
$$
où $k = g \circ g'$ et $m = h \circ h'$.
Nous avons $k^* = (g')^* \circ g^*$ et $m^* = (h')^* \circ h^*$.
L'assertion du lemme est que $\gamma_{rect}$
est égal au composé
$$
k^* \circ a =
(g')^* \circ g^* \circ a \xrightarrow{\gamma_{right}}
(g')^* \circ a' \circ h^* \xrightarrow{\gamma_{left}}
a'' \circ (h')^* \circ h^* = a'' \circ m^*
$$
Pour le voir, considérons le diagramme suivant :
$$
\xymatrix{
& (g')^* \circ g^* \circ a \ar[d] \ar[ddl] \\
& (g')^* \circ g^* \circ a \circ h_* \circ h^* \ar[ld] \\
(g')^* \circ g^* \circ a \circ h_* \circ (h')_* \circ (h')^* \circ h^* &
(g')^* \circ g^* \circ g_* \circ a' \circ h^*
\ar[u]_{\xi_{right}} \ar[d] \ar[ld] \\
(g')^* \circ g^* \circ g_* \circ a' \circ (h')_* \circ (h')^* \circ h^*
\ar[u]_{\xi_{right}} \ar[dr] &
(g')^* \circ a' \circ h^* \ar[d] \\
(g')^* \circ g^* \circ g_* \circ (g')_* \circ a'' \circ (h')^* \circ h^*
\ar[u]_{\xi_{left}} \ar[ddr] \ar[dr] &
(g')^* \circ a' \circ (h')_* \circ (h')^* \circ h^* \\
& (g')^*\circ (g')_* \circ a'' \circ (h')^* \circ h^*
\ar[u]_{\xi_{left}} \ar[d] \\
& a'' \circ (h')^* \circ h^*
}
$$
En descendant du côté droit, nous obtenons le composé, et en descendant
du côté gauche, nous obtenons $\gamma_{rect}$.
Tous les quadrilatères du côté droit de ce diagramme sont commutatifs
d'après Catégories, Lemme \ref{categories-lemma-properties-2-cat-cats},
ou, plus simplement, d'après la discussion qui précède
Catégories, Définition \ref{categories-definition-horizontal-composition}.
Nous voyons donc qu'il suffit de montrer que
$$
g_* \circ (g')_* \circ a'' \xrightarrow{\xi_{left}}
g_* \circ a' \circ (h')_* \xrightarrow{\xi_{right}}
a \circ h_* \circ (h')_*
$$
est égal à $\xi_{rect}$. C'est l'assertion duale de
Cohomologie, Remarque \ref{cohomology-remark-compose-base-change-horizontal},
ce qui achève la démonstration.
\end{proof}

\begin{remark}
\label{spaces-duality-remark-going-around}
Soit $S$ un schéma. Considérons un diagramme commutatif
$$
\xymatrix{
X'' \ar[r]_{k'} \ar[d]_{f''} & X' \ar[r]_k \ar[d]_{f'} & X \ar[d]^f \\
Y'' \ar[r]^{l'} \ar[d]_{g''} & Y' \ar[r]^l \ar[d]_{g'} & Y \ar[d]^g \\
Z'' \ar[r]^{m'} & Z' \ar[r]^m & Z
}
$$
d'espaces algébriques quasi-compacts et quasi-séparés sur $S$, dans lequel
tous les carrés sont cartésiens et où
$(f, l)$, $(g, m)$, $(f', l')$, $(g', m')$ sont des
paires de morphismes indépendants au sens de Tor.
Soient $a$, $a'$, $a''$, $b$, $b'$ et $b''$ les
adjoints à droite du Lemme \ref{spaces-duality-lemma-twisted-inverse-image}
pour $f$, $f'$, $f''$, $g$, $g'$ et $g''$.
Étiquetons les carrés du diagramme $A$, $B$, $C$, $D$
comme suit :
$$
\begin{matrix}
A & B \\
C & D
\end{matrix}
$$
Les morphismes (\ref{spaces-duality-equation-base-change-map})
des carrés sont alors les suivants (où nous utilisons $k^* = Lk^*$, etc.) :
$$
\begin{matrix}
\gamma_A : (k')^* \circ a' \to a'' \circ (l')^* &
\gamma_B : k^* \circ a \to a' \circ l^* \\
\gamma_C : (l')^* \circ b' \to b'' \circ (m')^* &
\gamma_D : l^* \circ b \to b' \circ m^*
\end{matrix}
$$
Pour les rectangles $2 \times 1$ et $1 \times 2$, nous avons quatre autres
morphismes de changement de base :
$$
\begin{matrix}
\gamma_{A + B} : (k \circ k')^* \circ a \to a'' \circ (l \circ l')^* \\
\gamma_{C + D} : (l \circ l')^* \circ b \to b'' \circ (m \circ m')^* \\
\gamma_{A + C} : (k')^* \circ (a' \circ b') \to (a'' \circ b'') \circ (m')^* \\
\gamma_{A + C} : k^* \circ (a \circ b) \to (a' \circ b') \circ m^*
\end{matrix}
$$
D'après le Lemme \ref{spaces-duality-lemma-compose-base-change-maps-horizontal}, nous avons
$$
\gamma_{A + B} = \gamma_A \circ \gamma_B, \quad
\gamma_{C + D} = \gamma_C \circ \gamma_D
$$
et d'après le Lemme \ref{spaces-duality-lemma-compose-base-change-maps}, nous avons
$$
\gamma_{A + C} = \gamma_C \circ \gamma_A, \quad
\gamma_{B + D} = \gamma_D \circ \gamma_B
$$
Il serait plus correct d'écrire ici
$\gamma_{A + B} = (\gamma_A \star \text{id}_{l^*}) \circ
(\text{id}_{(k')^*} \star \gamma_B)$ avec les notations de
Catégories, section \ref{categories-section-formal-cat-cat},
et de même pour les autres. Nous continuons toutefois à employer
l'abus de notation des démonstrations des
Lemmes \ref{spaces-duality-lemma-compose-base-change-maps} et
\ref{spaces-duality-lemma-compose-base-change-maps-horizontal},
qui consiste à omettre les produits $\star$ avec les identités, car on peut déterminer
ceux qu'il faut ajouter dès que la source et le but de la
transformation sont connus.
Cela étant dit, nous trouvons (a priori) deux transformations
$$
(k')^* \circ k^* \circ a \circ b
\longrightarrow
a'' \circ b'' \circ (m')^* \circ m^*
$$
à savoir
$$
\gamma_C \circ \gamma_A \circ \gamma_D \circ \gamma_B =
\gamma_{A + C} \circ \gamma_{B + D}
$$
et
$$
\gamma_C \circ \gamma_D \circ \gamma_A \circ \gamma_B =
\gamma_{C + D} \circ \gamma_{A + B}
$$
Le but de cette remarque est de signaler que ces transformations
sont égales. Pour le voir, il suffit en effet de montrer que
$$
\xymatrix{
(k')^* \circ a' \circ l^* \circ b \ar[r]_{\gamma_D} \ar[d]_{\gamma_A} &
(k')^* \circ a' \circ b' \circ m^* \ar[d]^{\gamma_A} \\
a'' \circ (l')^* \circ l^* \circ b \ar[r]^{\gamma_D} &
a'' \circ (l')^* \circ b' \circ m^*
}
$$
est commutatif. Cela résulte de
Catégories, Lemme \ref{categories-lemma-properties-2-cat-cats},
ou, plus simplement, de la discussion qui précède
Catégories, Définition \ref{categories-definition-horizontal-composition}.
\end{remark}







 
\section{Adjoint à droite de l'image directe et changement de base, II}
\label{spaces-duality-section-base-change-II}

\noindent
Dans cette section, nous démontrons que le morphisme de changement de base de la
section \ref{spaces-duality-section-base-change-map} est un isomorphisme
dans certains cas.

\begin{lemma}
\label{spaces-duality-lemma-more-base-change}
Dans le diagramme (\ref{spaces-duality-equation-base-change}), supposons en outre que
$g : Y' \to Y$ soit un morphisme de schémas affines et que $f : X \to Y$ soit propre.
Alors le morphisme de changement de base (\ref{spaces-duality-equation-base-change-map}) induit un
isomorphisme
$$
L(g')^*a(K) \longrightarrow a'(Lg^*K)
$$
dans les cas suivants :
\begin{enumerate}
\item pour tout $K \in D_\QCoh(\mathcal{O}_X)$ si $f$
est plat et de présentation finie,
\item pour tout $K \in D_\QCoh(\mathcal{O}_X)$ si $f$
est parfait et $Y$ noethérien,
\item pour $K \in D_\QCoh^+(\mathcal{O}_X)$ si $g$ est de dimension de Tor finie
et $Y$ est noethérien.
\end{enumerate}
\end{lemma}

\begin{proof}
Écrivons $Y = \Spec(A)$ et $Y' = \Spec(A')$. Comme changement de base d'un morphisme
affine, le morphisme $g'$ est affine. Soit $M$ un générateur parfait
de $D_\QCoh(\mathcal{O}_X)$, voir Catégories dérivées des espaces, Théorème
\ref{spaces-perfect-theorem-bondal-van-den-Bergh}. Alors $L(g')^*M$ est un
générateur de $D_\QCoh(\mathcal{O}_{X'})$, voir
Catégories dérivées des espaces, Remarque
\ref{spaces-perfect-remark-pullback-generator}.
Il suffit donc de montrer que (\ref{spaces-duality-equation-base-change-map})
induit un isomorphisme
\begin{equation}
\label{spaces-duality-equation-iso}
R\Hom_{X'}(L(g')^*M, L(g')^*a(K))
\longrightarrow
R\Hom_{X'}(L(g')^*M, a'(Lg^*K))
\end{equation}
de complexes de Hom globaux, voir
Cohomologie sur les sites, section \ref{sites-cohomology-section-global-RHom},
car cela impliquera que le cône de $L(g')^*a(K) \to a'(Lg^*K)$
est nul.
La démonstration est structurée comme suit : nous montrerons d'abord que
ces complexes de Hom sont isomorphes et, dans sa dernière partie,
que l'isomorphisme est induit par (\ref{spaces-duality-equation-iso}).

\medskip\noindent
Le membre de gauche. Comme $M$ est parfait, le morphisme canonique
$$
R\Hom_X(M, a(K)) \otimes^\mathbf{L}_A A'
\longrightarrow
R\Hom_{X'}(L(g')^*M, L(g')^*a(K))
$$
est un isomorphisme d'après Catégories dérivées des espaces, Lemme
\ref{spaces-perfect-lemma-affine-morphism-and-hom-out-of-perfect}.
En le combinant avec l'isomorphisme
$R\Hom_Y(Rf_*M, K) = R\Hom_X(M, a(K))$
du Lemme \ref{spaces-duality-lemma-iso-global-hom},
nous obtenons que le membre de gauche est égal à
$R\Hom_Y(Rf_*M, K) \otimes^\mathbf{L}_A A'$.

\medskip\noindent
Le membre de droite. Nous utilisons d'abord ici l'isomorphisme
$$
R\Hom_{X'}(L(g')^*M, a'(Lg^*K)) = R\Hom_{Y'}(Rf'_*L(g')^*M, Lg^*K)
$$
du Lemme \ref{spaces-duality-lemma-iso-global-hom}. Comme $f$ et $g$ sont
indépendants au sens de Tor, le morphisme de changement de base
$Lg^*Rf_*M \to Rf'_*L(g')^*M$ est un isomorphisme d'après
Catégories dérivées des espaces, Lemme
\ref{spaces-perfect-lemma-compare-base-change}.
Nous pouvons donc le réécrire sous la forme $R\Hom_{Y'}(Lg^*Rf_*M, Lg^*K)$.
Comme $Y$, $Y'$ sont affines et que $K$, $Rf_*M$ appartiennent à $D_\QCoh(\mathcal{O}_Y)$
(Catégories dérivées des espaces, Lemme
\ref{spaces-perfect-lemma-quasi-coherence-direct-image}),
nous avons un morphisme canonique
$$
\beta :
R\Hom_Y(Rf_*M, K) \otimes^\mathbf{L}_A A'
\longrightarrow
R\Hom_{Y'}(Lg^*Rf_*M, Lg^*K)
$$
dans $D(A')$. C'est la flèche de
Compléments d'algèbre, Équation (\ref{more-algebra-equation-base-change-RHom}),
où nous avons utilisé Catégories dérivées des schémas, Lemmes
\ref{perfect-lemma-affine-compare-bounded} et
\ref{perfect-lemma-quasi-coherence-internal-hom},
pour passer de la géométrie à l'algèbre et réciproquement.
\begin{enumerate}
\item Si $f$ est plat et de présentation finie, le complexe $Rf_*M$
est parfait sur $Y$ d'après Catégories dérivées des espaces, Lemme
\ref{spaces-perfect-lemma-flat-proper-perfect-direct-image-general},
et $\beta$ est un isomorphisme d'après
Compléments d'algèbre, Lemme \ref{more-algebra-lemma-base-change-RHom}, partie (1).
\item Si $f$ est parfait et $Y$ noethérien, le complexe $Rf_*M$
est parfait sur $Y$ d'après Compléments sur les morphismes d'espaces, Lemme
\ref{spaces-more-morphisms-lemma-perfect-proper-perfect-direct-image},
et $\beta$ est un isomorphisme comme précédemment.
\item Si $g$ est de dimension de Tor finie et $Y$ est noethérien,
le complexe $Rf_*M$ est pseudo-cohérent sur $Y$
(Catégories dérivées des espaces, Lemmes
\ref{spaces-perfect-lemma-direct-image-coherent} et
\ref{spaces-perfect-lemma-identify-pseudo-coherent-noetherian}),
et $\beta$ est un isomorphisme d'après
Compléments d'algèbre, Lemme \ref{more-algebra-lemma-base-change-RHom}, partie (4).
\end{enumerate}
Nous concluons que nous obtenons le même résultat que dans le paragraphe précédent.

\medskip\noindent
Dans la suite de la démonstration, nous montrons que les identifications des
membres de gauche et de droite de (\ref{spaces-duality-equation-iso})
données dans les deuxième et troisième paragraphes sont effectivement données par
(\ref{spaces-duality-equation-iso}). Pour que nos formules restent maniables,
nous noterons $(-, -)_X = R\Hom_X(-, -)$, écrirons $- \otimes A'$
à la place de $- \otimes_A^\mathbf{L} A'$, et abrégerons
$g^* = Lg^*$ et $f_* = Rf_*$. Considérons le diagramme
commutatif suivant :
$$
\xymatrix{
((g')^*M, (g')^*a(K))_{X'} \ar[d] &
(M, a(K))_X \otimes A' \ar[l]^-\alpha \ar[d] &
(f_*M, K)_Y \otimes A' \ar@{=}[l] \ar[d] \\
((g')^*M, (g')^*a(g_*g^*K))_{X'} &
(M, a(g_*g^*K))_X \otimes A' \ar[l]^-\alpha &
(f_*M, g_*g^*K)_Y \otimes A' \ar@{=}[l] \ar@/_4pc/[dd]_{\mu'} \\
((g')^*M, (g')^*g'_*a'(g^*K))_{X'} \ar[u] \ar[d] &
(M, g'_*a'(g^*K))_X \otimes A' \ar[u] \ar[l]^-\alpha \ar[ld]^\mu &
(f_*M, K) \otimes A' \ar[d]^\beta \\
((g')^*M, a'(g^*K))_{X'} &
(f'_*(g')^*M, g^*K)_{Y'} \ar@{=}[l] \ar[r] &
(g^*f_*M, g^*K)_{Y'}
}
$$
Les flèches étiquetées $\alpha$ sont les morphismes de
Catégories dérivées des espaces, Lemme
\ref{spaces-perfect-lemma-affine-morphism-and-hom-out-of-perfect},
pour le diagramme de sommets $X', X, Y', Y$.
La partie supérieure du diagramme est commutative, car les flèches horizontales sont
fonctorielles en leurs arguments.
Les flèches verticales du milieu proviennent de la transformation inversible
$g'_* \circ a' \to a \circ g_*$ du Lemme \ref{spaces-duality-lemma-flat-precompose-pus},
et le carré médian est donc commutatif.
Descendre le long du côté gauche donne (\ref{spaces-duality-equation-iso}).
Les flèches horizontales supérieures donnent les identifications utilisées dans le
deuxième paragraphe de la démonstration.
Les flèches horizontales inférieures, y compris $\beta$, donnent les identifications
utilisées dans le troisième paragraphe de la démonstration. Étant donnés $E \in D(A)$,
$E' \in D(A')$ et $c : E \to E'$ dans $D(A)$, nous noterons
$\mu_c : E \otimes A' \to E'$ le morphisme induit par $c$
et par l'adjonction entre restriction et changement de base ;
si $c$ est clair, nous écrivons $\mu = \mu_c$, c'est-à-dire que nous
omettons $c$ de la notation. Le morphisme $\mu$ du diagramme est de cette
forme, avec $c$ donné par l'identification
$(M, g'_*a(g^*K))_X = ((g')^*M, a'(g^*K))_{X'}$
; le triangle contenant $\mu$ est commutatif d'après
Catégories dérivées des espaces, Remarque
\ref{spaces-perfect-remark-multiplication-map}.

\medskip\noindent
Observons que
$$
\xymatrix{
(M, a(g_*g^*K))_X &
(f_*M, g_* g^*K)_Y \ar@{=}[l] &
(g^*f_*M, g^*K)_{Y'} \ar@{=}[l] \\
(M, g'_* a'(g^*K))_X \ar[u] &
((g')^*M, a'(g^*K))_{X'} \ar@{=}[l] &
(f'_*(g')^*M, g^*K)_{Y'} \ar@{=}[l] \ar[u]
}
$$
est commutatif par la définition même de la transformation
$g'_* \circ a' \to a \circ g_*$. Si $\mu'$ désigne, comme ci-dessus, le morphisme
correspondant à l'identification
$(f_*M, g_*g^*K)_X = (g^*f_*M, g^*K)_{Y'}$, alors
l'hexagone est lui aussi commutatif. Il suffit donc de montrer que
$\beta$ est égal au composé de
$(f_*M, K)_Y \otimes A' \to (f_*M, g_*g^*K)_X \otimes A'$
et de $\mu'$. Pour cela, il suffit de démontrer que les deux morphismes induits
$(f_*M, K)_Y \to (g^*f_*M, g^*K)_{Y'}$ sont égaux.
Autrement dit, il suffit de montrer que le diagramme
$$
\xymatrix{
R\Hom_A(E, K) \ar[rr]_{\text{induit par }\beta} \ar[rd] & &
R\Hom_{A'}(E \otimes_A^\mathbf{L} A', K \otimes_A^\mathbf{L} A') \\
& R\Hom_A(E, K \otimes_A^\mathbf{L} A') \ar[ru]
}
$$
est commutatif pour tous $E, K \in D(A)$. Puisque c'est ainsi que $\beta$ est construit dans
Compléments d'algèbre, section \ref{more-algebra-section-base-change-RHom},
la démonstration est achevée.
\end{proof}





 
\section{Adjoint à droite de l'image directe et morphismes trace}
\label{spaces-duality-section-trace}

\noindent
Soit $S$ un schéma.
Soit $f : X \to Y$ un morphisme d'espaces algébriques quasi-compacts et quasi-séparés
sur $S$.
Soit $a : D_\QCoh(\mathcal{O}_Y) \to D_\QCoh(\mathcal{O}_X)$
l'adjoint à droite du Lemme \ref{spaces-duality-lemma-twisted-inverse-image}. D'après
Catégories, section \ref{categories-section-adjoint}, nous obtenons une
transformation de foncteurs
$$
\text{Tr}_f : Rf_* \circ a \longrightarrow \text{id}
$$
Le morphisme correspondant $\text{Tr}_{f, K} : Rf_*a(K) \longrightarrow K$
pour $K \in D_\QCoh(\mathcal{O}_Y)$ est parfois appelé le {\it morphisme trace}.
C'est le morphisme ayant la propriété que la bijection
$$
\Hom_X(L, a(K)) \longrightarrow \Hom_Y(Rf_*L, K)
$$
pour $L \in D_\QCoh(\mathcal{O}_X)$ qui caractérise l'adjoint à droite
est donnée par
$$
\varphi \longmapsto \text{Tr}_{f, K} \circ Rf_*\varphi
$$
Le morphisme canonique (\ref{spaces-duality-equation-sheafy-trace})
$$
Rf_*R\SheafHom_{\mathcal{O}_X}(L, a(K))
\longrightarrow
R\SheafHom_{\mathcal{O}_Y}(Rf_*L, K)
$$
est obtenu par composition avec $\text{Tr}_{f, K}$.
Tout morphisme trace que nous considérerons dans cette section sera un
cas particulier de celui-ci. Avant d'en étudier quelques cas particuliers,
nous montrons que la formation du morphisme trace commute au changement de base.

\begin{lemma}[Morphisme trace et changement de base]
\label{spaces-duality-lemma-trace-map-and-base-change}
Supposons donné un diagramme (\ref{spaces-duality-equation-base-change}).
Alors les morphismes
$1 \star \text{Tr}_f : Lg^* \circ Rf_* \circ a \to Lg^*$ et
$\text{Tr}_{f'} \star 1 : Rf'_* \circ a' \circ Lg^* \to Lg^*$
coïncident via les morphismes de changement de base
$\beta : Lg^* \circ Rf_* \to Rf'_* \circ L(g')^*$
(Cohomologie sur les sites, Remarque \ref{sites-cohomology-remark-base-change})
et $\alpha : L(g')^* \circ a \to a' \circ Lg^*$
(\ref{spaces-duality-equation-base-change-map}).
Plus précisément, le diagramme
$$
\xymatrix{
Lg^* \circ Rf_* \circ a
\ar[d]_{\beta \star 1} \ar[r]_-{1 \star \text{Tr}_f} &
Lg^* \\
Rf'_* \circ L(g')^* \circ a \ar[r]^{1 \star \alpha} &
Rf'_* \circ a' \circ Lg^* \ar[u]_{\text{Tr}_{f'} \star 1}
}
$$
de transformations de foncteurs est commutatif.
\end{lemma}

\begin{proof}
Dans cette démonstration, nous écrivons $f_*$ pour $Rf_*$ et $g^*$ pour $Lg^*$, et nous
omettons les produits $\star$ avec les identités, puisqu'on peut déterminer ceux qu'il faut
ajouter dès que la source et le but de la transformation sont connus.
Rappelons que $\beta : g^* \circ f_* \to f'_* \circ (g')^*$ est un isomorphisme
et que $\alpha$ est défini à l'aide de
l'isomorphisme $\beta^\vee : g'_* \circ a' \to a \circ g_*$,
qui est l'adjoint de $\beta$ ; voir le Lemme \ref{spaces-duality-lemma-flat-precompose-pus}
et sa démonstration. Remarquons d'abord que la flèche horizontale supérieure
du diagramme du lemme est égale au composé
$$
g^* \circ f_* \circ a \to
g^* \circ f_* \circ a \circ g_* \circ g^* \to
g^* \circ g_* \circ g^* \to g^*
$$
où la première flèche est l'unité de $(g^*, g_*)$, la deuxième flèche
est $\text{Tr}_f$, et la troisième est la co-unité de $(g^*, g_*)$.
C'est une conséquence immédiate du fait que le composé
$g^* \to g^* \circ g_* \circ g^* \to g^*$ de l'unité et de la co-unité est l'identité.
Considérons le diagramme
$$
\xymatrix{
& g^* \circ f_* \circ a \ar[ld]_\beta \ar[d] \ar[r]_{\text{Tr}_f} & g^* \\
f'_* \circ (g')^* \circ a \ar[dr] &
g^* \circ f_* \circ a \circ g_* \circ g^* \ar[d]_\beta \ar[ru] &
g^* \circ f_* \circ g'_* \circ a' \circ g^* \ar[l]_{\beta^\vee} \ar[d]_\beta &
f'_* \circ a' \circ g^* \ar[lu]_{\text{Tr}_{f'}} \\
& f'_* \circ (g')^* \circ a \circ g_* \circ g^* &
f'_* \circ (g')^* \circ g'_* \circ a' \circ g^* \ar[ru] \ar[l]_{\beta^\vee}
}
$$
Dans ce diagramme, les deux carrés sont commutatifs d'après
Catégories, Lemme \ref{categories-lemma-properties-2-cat-cats},
ou, plus simplement, d'après la discussion qui précède
Catégories, Définition \ref{categories-definition-horizontal-composition}.
Le triangle est commutatif d'après la discussion ci-dessus. D'après
Catégories, Lemme
\ref{categories-lemma-transformation-between-functors-and-adjoints},
le carré
$$
\xymatrix{
g^* \circ f_* \circ g'_* \circ a' \ar[d]_{\beta^\vee} \ar[r]_-\beta &
f'_* \circ (g')^* \circ g'_* \circ a' \ar[d] \\
g^* \circ f_* \circ a \circ g_* \ar[r] &
\text{id}
}
$$
est commutatif, ce qui implique que le pentagone du grand diagramme est commutatif.
Comme $\beta$ et $\beta^\vee$ sont des isomorphismes, et comme le parcours
du bord extérieur du grand diagramme est égal à
$\text{Tr}_f \circ \alpha \circ \beta$ par définition, cela démontre le lemme.
\end{proof}

\noindent
Soit $S$ un schéma.
Soit $f : X \to Y$ un morphisme d'espaces algébriques quasi-compacts et quasi-séparés
sur $S$.
Soit $a : D_\QCoh(\mathcal{O}_Y) \to D_\QCoh(\mathcal{O}_X)$
l'adjoint à droite de $Rf_*$ comme dans le
Lemme \ref{spaces-duality-lemma-twisted-inverse-image}. D'après
Catégories, section \ref{categories-section-adjoint}, nous obtenons une
transformation de foncteurs
$$
\eta_f : \text{id} \to  a \circ Rf_*
$$
appelée l'unité de l'adjonction.

\begin{lemma}
\label{spaces-duality-lemma-unit-and-base-change}
Supposons donné un diagramme (\ref{spaces-duality-equation-base-change}). Alors les morphismes
$1 \star \eta_f : L(g')^* \to L(g')^* \circ a \circ Rf_*$ et
$\eta_{f'} \star 1 : L(g')^* \to a' \circ Rf'_* \circ L(g')^*$
coïncident via les morphismes de changement de base
$\beta : Lg^* \circ Rf_* \to Rf'_* \circ L(g')^*$
(Cohomologie sur les sites, Remarque \ref{sites-cohomology-remark-base-change})
et $\alpha : L(g')^* \circ a \to a' \circ Lg^*$
(\ref{spaces-duality-equation-base-change-map}).
Plus précisément, le diagramme
$$
\xymatrix{
L(g')^* \ar[r]_-{1 \star \eta_f} \ar[d]_{\eta_{f'} \star 1} &
L(g')^* \circ a \circ Rf_* \ar[d]^\alpha \\
a' \circ Rf'_* \circ L(g')^* &
a' \circ Lg^* \circ Rf_* \ar[l]_-\beta
}
$$
de transformations de foncteurs est commutatif.
\end{lemma}

\begin{proof}
Cette démonstration est duale de celle du Lemme \ref{spaces-duality-lemma-trace-map-and-base-change}.
Dans cette démonstration, nous écrivons $f_*$ pour $Rf_*$ et $g^*$ pour $Lg^*$, et nous
omettons les produits $\star$ avec les identités, puisqu'on peut déterminer ceux qu'il faut
ajouter dès que la source et le but de la transformation sont connus.
Rappelons que $\beta : g^* \circ f_* \to f'_* \circ (g')^*$ est un isomorphisme
et que $\alpha$ est défini à l'aide de
l'isomorphisme $\beta^\vee : g'_* \circ a' \to a \circ g_*$,
qui est l'adjoint de $\beta$ ; voir le Lemme \ref{spaces-duality-lemma-flat-precompose-pus}
et sa démonstration. Remarquons d'abord que la flèche verticale gauche
du diagramme du lemme est égale au composé
$$
(g')^* \to (g')^* \circ g'_* \circ (g')^* \to
(g')^* \circ g'_* \circ a' \circ f'_* \circ (g')^* \to
a' \circ f'_* \circ (g')^*
$$
où la première flèche est l'unité de $((g')^*, g'_*)$, la deuxième flèche
est $\eta_{f'}$, et la troisième est la co-unité de $((g')^*, g'_*)$.
C'est une conséquence immédiate du fait que le composé
$(g')^* \to (g')^* \circ (g')_* \circ (g')^* \to (g')^*$
de l'unité et de la co-unité est l'identité. Considérons le diagramme
$$
\xymatrix{
& (g')^* \circ a \circ f_* \ar[r] &
(g')^* \circ a \circ g_* \circ g^* \circ f_*
\ar[ld]_\beta \\
(g')^* \ar[ru]^{\eta_f} \ar[dd]_{\eta_{f'}} \ar[rd] &
(g')^* \circ a \circ g_* \circ f'_* \circ (g')^* &
(g')^* \circ g'_* \circ a' \circ g^* \circ f_*
\ar[u]_{\beta^\vee} \ar[ld]_\beta \ar[d] \\
& (g')^* \circ g'_* \circ a' \circ f'_* \circ (g')^*
\ar[ld] \ar[u]_{\beta^\vee} &
a' \circ g^* \circ f_* \ar[lld]^\beta \\
a' \circ f'_* \circ (g')^*
}
$$
Dans ce diagramme, les deux carrés sont commutatifs d'après
Catégories, Lemme \ref{categories-lemma-properties-2-cat-cats},
ou, plus simplement, d'après la discussion qui précède
Catégories, Définition \ref{categories-definition-horizontal-composition}.
Le triangle est commutatif d'après la discussion ci-dessus. D'après le dual de
Catégories, Lemme
\ref{categories-lemma-transformation-between-functors-and-adjoints},
le carré
$$
\xymatrix{
\text{id} \ar[r] \ar[d] &
g'_* \circ a' \circ g^* \circ f_* \ar[d]^\beta \\
g'_* \circ a' \circ g^* \circ f_* \ar[r]^{\beta^\vee} &
a \circ g_* \circ f'_* \circ (g')^*
}
$$
est commutatif, ce qui implique que le pentagone du grand diagramme est commutatif.
Comme $\beta$ et $\beta^\vee$ sont des isomorphismes, et comme le parcours
du bord extérieur du grand diagramme est égal à
$\beta \circ \alpha \circ \eta_f$ par définition, cela démontre le lemme.
\end{proof}





 
\section{Adjoint à droite de l'image directe et image réciproque}
\label{spaces-duality-section-compare-with-pullback}

\noindent
Soit $S$ un schéma.
Soit $f : X \to Y$ un morphisme d'espaces algébriques quasi-compacts et quasi-séparés
sur $S$.
Soit $a$ l'adjoint à droite de l'image directe, comme dans le
Lemme \ref{spaces-duality-lemma-twisted-inverse-image}. Pour $K, L \in D_\QCoh(\mathcal{O}_Y)$,
il existe un morphisme canonique
$$
Lf^*K \otimes^\mathbf{L}_{\mathcal{O}_X} a(L)
\longrightarrow
a(K \otimes_{\mathcal{O}_Y}^\mathbf{L} L)
$$
Ce morphisme est en effet adjoint d'un morphisme
$$
Rf_*(Lf^*K \otimes^\mathbf{L}_{\mathcal{O}_X} a(L)) =
K \otimes^\mathbf{L}_{\mathcal{O}_Y} Rf_*(a(L))
\longrightarrow
K \otimes^\mathbf{L}_{\mathcal{O}_Y} L
$$
(l'égalité résulte de Catégories dérivées des espaces, Lemme
\ref{spaces-perfect-lemma-cohomology-base-change}),
pour lequel nous utilisons le morphisme trace $Rf_*a(L) \to L$.
Lorsque $L = \mathcal{O}_Y$, nous obtenons un morphisme
\begin{equation}
\label{spaces-duality-equation-compare-with-pullback}
Lf^*K \otimes^\mathbf{L}_{\mathcal{O}_X} a(\mathcal{O}_Y) \longrightarrow a(K)
\end{equation}
fonctoriel en $K$ et compatible aux triangles distingués.

\begin{lemma}
\label{spaces-duality-lemma-compare-with-pullback-perfect}
Soit $S$ un schéma.
Soit $f : X \to Y$ un morphisme d'espaces algébriques quasi-compacts et quasi-séparés
sur $S$. Le morphisme
$Lf^*K \otimes^\mathbf{L}_{\mathcal{O}_X} a(L) \to
a(K \otimes_{\mathcal{O}_Y}^\mathbf{L} L)$
défini ci-dessus pour $K, L \in D_\QCoh(\mathcal{O}_Y)$
est un isomorphisme si $K$ est parfait. En particulier,
(\ref{spaces-duality-equation-compare-with-pullback}) est un isomorphisme si $K$ est parfait.
\end{lemma}

\begin{proof}
Soit $K^\vee$ le « dual » de $K$, voir
Cohomologie sur les sites, Lemme \ref{sites-cohomology-lemma-dual-perfect-complex}.
Pour $M \in D_\QCoh(\mathcal{O}_X)$, nous avons
\begin{align*}
\Hom_{D(\mathcal{O}_Y)}(Rf_*M, K \otimes^\mathbf{L}_{\mathcal{O}_Y} L)
& =
\Hom_{D(\mathcal{O}_Y)}(
Rf_*M \otimes^\mathbf{L}_{\mathcal{O}_Y} K^\vee, L) \\
& =
\Hom_{D(\mathcal{O}_X)}(
M \otimes^\mathbf{L}_{\mathcal{O}_X} Lf^*K^\vee, a(L)) \\
& =
\Hom_{D(\mathcal{O}_X)}(M,
Lf^*K \otimes^\mathbf{L}_{\mathcal{O}_X} a(L))
\end{align*}
La deuxième égalité résulte de la définition de $a$ et de la formule de projection
(Cohomologie sur les sites, Lemme
\ref{sites-cohomology-lemma-projection-formula}),
ou du résultat plus général Catégories dérivées des espaces, Lemme
\ref{spaces-perfect-lemma-cohomology-base-change}.
Le résultat découle donc du lemme de Yoneda.
\end{proof}

\begin{lemma}
\label{spaces-duality-lemma-restriction-compare-with-pullback}
Supposons donné un diagramme (\ref{spaces-duality-equation-base-change}).
Soit $K \in D_\QCoh(\mathcal{O}_Y)$. Le diagramme
$$
\xymatrix{
L(g')^*(Lf^*K \otimes^\mathbf{L}_{\mathcal{O}_X} a(\mathcal{O}_Y))
\ar[r] \ar[d] & L(g')^*a(K) \ar[d] \\
L(f')^*Lg^*K \otimes_{\mathcal{O}_{X'}}^\mathbf{L} a'(\mathcal{O}_{Y'})
\ar[r] & a'(Lg^*K)
}
$$
est commutatif, où les flèches horizontales sont les morphismes
(\ref{spaces-duality-equation-compare-with-pullback}) pour $K$ et $Lg^*K$,
et où les flèches verticales sont construites à l'aide de
Cohomologie sur les sites, Remarque \ref{sites-cohomology-remark-base-change}, et de
(\ref{spaces-duality-equation-base-change-map}).
\end{lemma}

\begin{proof}
Dans cette démonstration, nous écrirons $f_*$ pour $Rf_*$ et $f^*$ pour $Lf^*$, etc.,
et nous écrirons $\otimes$ pour $\otimes^\mathbf{L}_{\mathcal{O}_X}$, etc.
Écrivons (\ref{spaces-duality-equation-compare-with-pullback}) comme le composé
\begin{align*}
f^*K \otimes a(\mathcal{O}_Y)
& \to
a(f_*(f^*K \otimes a(\mathcal{O}_Y))) \\
& \leftarrow
a(K \otimes f_*a(\mathcal{O}_K)) \\
& \to
a(K \otimes \mathcal{O}_Y) \\
& \to
a(K)
\end{align*}
Ici, la première flèche est l'unité $\eta_f$, la deuxième est l'image par $a$
de Cohomologie sur les sites, Équation
(\ref{sites-cohomology-equation-projection-formula-map}), qui est un
isomorphisme d'après Catégories dérivées des espaces, Lemme
\ref{spaces-perfect-lemma-cohomology-base-change} ; la troisième flèche est
l'image par $a$ de $\text{id}_K \otimes \text{Tr}_f$, et la quatrième
est l'image par $a$ de l'isomorphisme $K \otimes \mathcal{O}_Y = K$.
La démonstration du lemme consiste à montrer que chacun de ces
morphismes donne un carré commutatif comme dans l'énoncé du lemme.
Pour $\eta_f$ et $\text{Tr}_f$, cela résulte des
Lemmes \ref{spaces-duality-lemma-unit-and-base-change} et
\ref{spaces-duality-lemma-trace-map-and-base-change}.
Pour la flèche qui utilise Cohomologie sur les sites, Équation
(\ref{sites-cohomology-equation-projection-formula-map}),
cela résulte de Cohomologie sur les sites, Remarque
\ref{sites-cohomology-remark-compatible-with-diagram}.
Pour le morphisme de multiplication, c'est clair. Cela achève la démonstration.
\end{proof}









 
\section{Adjoint à droite de l'image directe pour les morphismes propres et plats}
\label{spaces-duality-section-proper-flat}

\noindent
Pour les morphismes propres, plats et de présentation finie entre espaces algébriques
quasi-compacts et quasi-séparés, l'adjoint à droite de l'image directe
possède quelques propriétés remarquables.

\begin{lemma}
\label{spaces-duality-lemma-proper-flat}
Soit $S$ un schéma.
Soit $Y$ un espace algébrique quasi-compact et quasi-séparé sur $S$.
Soit $f : X \to Y$ un morphisme d'espaces algébriques propre, plat et
de présentation finie.
Soit $a$ l'adjoint à droite de
$Rf_* : D_\QCoh(\mathcal{O}_X) \to D_\QCoh(\mathcal{O}_Y)$ du
Lemme \ref{spaces-duality-lemma-twisted-inverse-image}. Alors $a$ commute aux sommes directes.
\end{lemma}

\begin{proof}
Soit $P$ un objet parfait de $D(\mathcal{O}_X)$. D'après
Catégories dérivées des espaces, Lemme
\ref{spaces-perfect-lemma-flat-proper-perfect-direct-image-general},
le complexe $Rf_*P$ est parfait sur $Y$.
Soit $K_i$ une famille d'objets de $D_\QCoh(\mathcal{O}_Y)$.
Alors
\begin{align*}
\Hom_{D(\mathcal{O}_X)}(P, a(\bigoplus K_i))
& =
\Hom_{D(\mathcal{O}_Y)}(Rf_*P, \bigoplus K_i) \\
& =
\bigoplus \Hom_{D(\mathcal{O}_Y)}(Rf_*P, K_i) \\
& =
\bigoplus \Hom_{D(\mathcal{O}_X)}(P, a(K_i))
\end{align*}
car un objet parfait est compact (Catégories dérivées des espaces,
Proposition \ref{spaces-perfect-proposition-compact-is-perfect}).
Comme $D_\QCoh(\mathcal{O}_X)$ possède un générateur parfait
(Catégories dérivées des espaces, Théorème
\ref{spaces-perfect-theorem-bondal-van-den-Bergh}),
nous concluons que le morphisme $\bigoplus a(K_i) \to a(\bigoplus K_i)$
est un isomorphisme, c'est-à-dire que $a$ commute aux sommes directes.
\end{proof}

\begin{lemma}
\label{spaces-duality-lemma-compare-with-pullback-flat-proper}
Soit $S$ un schéma.
Soit $Y$ un espace algébrique quasi-compact et quasi-séparé sur $S$.
Soit $f : X \to Y$ un morphisme d'espaces algébriques propre, plat et
de présentation finie.
Le morphisme (\ref{spaces-duality-equation-compare-with-pullback}) est un isomorphisme
pour tout objet $K$ de $D_\QCoh(\mathcal{O}_Y)$.
\end{lemma}

\begin{proof}
D'après le Lemme \ref{spaces-duality-lemma-proper-flat}, nous savons que $a$ commute
aux sommes directes. La collection des objets de
$D_\QCoh(\mathcal{O}_Y)$ pour lesquels (\ref{spaces-duality-equation-compare-with-pullback})
est un isomorphisme est donc une sous-catégorie strictement pleine, saturée et triangulée
de $D_\QCoh(\mathcal{O}_Y)$, qui est en outre
stable par sommes directes. Comme $D_\QCoh(\mathcal{O}_Y)$
est une catégorie de modules (Catégories dérivées des espaces, Théorème
\ref{spaces-perfect-theorem-DQCoh-is-Ddga}) engendrée par un unique
objet parfait (Catégories dérivées des espaces, Théorème
\ref{spaces-perfect-theorem-bondal-van-den-Bergh}),
nous pouvons raisonner comme dans
Compléments d'algèbre, Remarque \ref{more-algebra-remark-P-resolution},
pour voir qu'il suffit de démontrer que (\ref{spaces-duality-equation-compare-with-pullback})
est un isomorphisme pour un seul objet parfait.
Or le résultat vaut pour les objets parfaits, voir
Lemme \ref{spaces-duality-lemma-compare-with-pullback-perfect}.
\end{proof}

\begin{lemma}
\label{spaces-duality-lemma-properties-relative-dualizing}
Soit $Y$ un schéma affine. Soit $f : X \to Y$ un morphisme
d'espaces algébriques propre, plat et de présentation finie.
Soit $a$ l'adjoint à droite de
$Rf_* : D_\QCoh(\mathcal{O}_X) \to D_\QCoh(\mathcal{O}_Y)$ du
Lemme \ref{spaces-duality-lemma-twisted-inverse-image}.
Alors
\begin{enumerate}
\item $a(\mathcal{O}_Y)$ est un objet $Y$-parfait de $D(\mathcal{O}_X)$,
\item les faisceaux de cohomologie de $Rf_*a(\mathcal{O}_Y)$
sont nuls en degrés positifs,
\item $\mathcal{O}_X \to
R\SheafHom_{\mathcal{O}_X}(a(\mathcal{O}_Y), a(\mathcal{O}_Y))$
est un isomorphisme.
\end{enumerate}
\end{lemma}

\begin{proof}
Pour un objet parfait $E$ de $D(\mathcal{O}_X)$, nous avons
\begin{align*}
Rf_*(E \otimes_{\mathcal{O}_X}^\mathbf{L} \omega_{X/Y}^\bullet)
& =
Rf_*R\SheafHom_{\mathcal{O}_X}(E^\vee, \omega_{X/Y}^\bullet) \\
& =
R\SheafHom_{\mathcal{O}_Y}(Rf_*E^\vee, \mathcal{O}_Y) \\
& =
(Rf_*E^\vee)^\vee
\end{align*}
Pour la première égalité, voir
Cohomologie sur les sites, Lemme \ref{sites-cohomology-lemma-dual-perfect-complex}.
Pour la deuxième égalité, voir le Lemme \ref{spaces-duality-lemma-iso-on-RSheafHom},
la Remarque \ref{spaces-duality-remark-iso-on-RSheafHom}, et
Catégories dérivées des espaces, Lemme
\ref{spaces-perfect-lemma-flat-proper-perfect-direct-image-general}.
La troisième égalité est la définition du dual. En particulier,
ces références montrent aussi que le résultat est un objet parfait
de $D(\mathcal{O}_Y)$. Nous concluons que $\omega_{X/Y}^\bullet$
est $Y$-parfait d'après Compléments sur les morphismes d'espaces, Lemme
\ref{spaces-more-morphisms-lemma-characterize-relatively-perfect}.
Cela démontre (1).

\medskip\noindent
Soit $M$ un objet de $D_\QCoh(\mathcal{O}_Y)$. Alors
\begin{align*}
\Hom_Y(M, Rf_*a(\mathcal{O}_Y)) & =
\Hom_X(Lf^*M, a(\mathcal{O}_Y)) \\
& =
\Hom_Y(Rf_*Lf^*M, \mathcal{O}_Y) \\
& =
\Hom_Y(M \otimes_{\mathcal{O}_Y}^\mathbf{L} Rf_*\mathcal{O}_Y, \mathcal{O}_Y)
\end{align*}
La première égalité résulte de Cohomologie sur les sites, Lemme
\ref{sites-cohomology-lemma-adjoint}.
La deuxième égalité résulte de la construction de $a$.
La troisième égalité résulte de Catégories dérivées des espaces, Lemme
\ref{spaces-perfect-lemma-cohomology-base-change}.
Rappelons que $Rf_*\mathcal{O}_X$ est parfait, d'amplitude de Tor dans $[0, N]$
pour un certain $N$, voir
Catégories dérivées des espaces, Lemme
\ref{spaces-perfect-lemma-flat-proper-perfect-direct-image-general}.
Nous pouvons donc représenter $Rf_*\mathcal{O}_X$ par un complexe de
modules projectifs finis placé en degrés $[0, N]$
(en utilisant Compléments d'algèbre, Lemme \ref{more-algebra-lemma-perfect},
et le fait que $Y$ est affine).
Ainsi, si $M = \mathcal{O}_Y[-i]$ pour un certain $i > 0$, le dernier
groupe est nul. Comme $Y$ est affine, nous concluons que
$H^i(Rf_*a(\mathcal{O}_Y)) = 0$ pour $i > 0$.
Cela démontre (2).

\medskip\noindent
Soit $E$ un objet parfait de $D_\QCoh(\mathcal{O}_X)$. Alors
nous avons
\begin{align*}
\Hom_X(E, R\SheafHom_{\mathcal{O}_X}(a(\mathcal{O}_Y), a(\mathcal{O}_Y))
& =
\Hom_X(E \otimes_{\mathcal{O}_X}^\mathbf{L} a(\mathcal{O}_Y),
a(\mathcal{O}_Y)) \\
& =
\Hom_Y(Rf_*(E \otimes_{\mathcal{O}_X}^\mathbf{L} a(\mathcal{O}_Y)),
\mathcal{O}_Y) \\
& =
\Hom_Y(Rf_*(R\SheafHom_{\mathcal{O}_X}(E^\vee, a(\mathcal{O}_Y))),
\mathcal{O}_Y) \\
& =
\Hom_Y(R\SheafHom_{\mathcal{O}_Y}(Rf_*E^\vee, \mathcal{O}_Y),
\mathcal{O}_Y) \\
& =
R\Gamma(Y, Rf_*E^\vee) \\
& =
\Hom_X(E, \mathcal{O}_X)
\end{align*}
La première égalité résulte de Cohomologie sur les sites, Lemme
\ref{sites-cohomology-lemma-internal-hom}.
La deuxième égalité est la définition de $a$.
La troisième égalité provient de la construction du complexe parfait dual
$E^\vee$, voir Cohomologie sur les sites, Lemme
\ref{sites-cohomology-lemma-dual-perfect-complex}.
La quatrième égalité résulte de l'égalité
$Rf_*R\SheafHom_{\mathcal{O}_X}(E^\vee, \omega_{X/Y}^\bullet) =
R\SheafHom_{\mathcal{O}_Y}(Rf_*E^\vee, \mathcal{O}_Y)$
démontrée dans le premier paragraphe de la preuve.
La cinquième égalité résulte de la bidualité pour les complexes parfaits
(Cohomologie sur les sites, Lemme
\ref{sites-cohomology-lemma-dual-perfect-complex})
et du fait que $Rf_*E$ est parfait d'après
Catégories dérivées des espaces, Lemme
\ref{spaces-perfect-lemma-flat-proper-perfect-direct-image-general}.
La dernière égalité est la suite spectrale de Leray pour $f$.
Cette suite d'égalités montre essentiellement que (3)
vaut d'après le lemme de Yoneda. En effet, l'objet
$R\SheafHom(a(\mathcal{O}_Y), a(\mathcal{O}_Y))$
appartient à $D_\QCoh(\mathcal{O}_X)$ d'après Catégories dérivées des espaces, Lemme
\ref{spaces-perfect-lemma-quasi-coherence-internal-hom}.
En prenant $E = \mathcal{O}_X$ ci-dessus, nous obtenons un morphisme
$\alpha : \mathcal{O}_X \to
R\SheafHom_{\mathcal{O}_X}(a(\mathcal{O}_Y), a(\mathcal{O}_Y))$
correspondant à
$\text{id}_{\mathcal{O}_X} \in \Hom_X(\mathcal{O}_X, \mathcal{O}_X)$.
Comme tous les isomorphismes ci-dessus sont fonctoriels en $E$, nous
voyons que le cône de $\alpha$ est un objet $C$ de $D_\QCoh(\mathcal{O}_X)$
tel que $\Hom(E, C) = 0$ pour tout $E$ parfait.
Comme les objets parfaits engendrent
(Catégories dérivées des espaces, Théorème
\ref{spaces-perfect-theorem-bondal-van-den-Bergh}),
nous concluons que $\alpha$ est un isomorphisme.
\end{proof}









 
\section{Complexes dualisants relatifs pour les morphismes propres et plats}
\label{spaces-duality-section-relative-dualizing-proper-flat}

\noindent
Motivés par Dualité pour les schémas, sections
\ref{duality-section-proper-flat} et
\ref{duality-section-relative-dualizing-complexes},
ainsi que par les résultats de la
section \ref{spaces-duality-section-proper-flat},
nous posons la définition suivante.

\begin{definition}
\label{spaces-duality-definition-relative-dualizing-proper-flat}
Soit $S$ un schéma. Soit $f : X \to Y$ un morphisme propre et plat
d'espaces algébriques sur $S$, de présentation finie.
Un {\it complexe dualisant relatif} pour $X/Y$ est un couple
$(\omega_{X/Y}^\bullet, \tau)$ formé d'un
objet $Y$-parfait $\omega_{X/Y}^\bullet$ de $D(\mathcal{O}_X)$
et d'un morphisme
$$
\tau : Rf_*\omega_{X/Y}^\bullet \longrightarrow \mathcal{O}_Y
$$
tel que, pour tout carré cartésien
$$
\xymatrix{
X' \ar[r]_{g'} \ar[d]_{f'} & X \ar[d]^f \\
Y' \ar[r]^g & Y
}
$$
où $Y'$ est un schéma affine, le couple
$(L(g')^*\omega_{X/Y}^\bullet, Lg^*\tau)$
soit isomorphe au couple
$(a'(\mathcal{O}_{Y'}), \text{Tr}_{f', \mathcal{O}_{Y'}})$
étudié dans les sections
\ref{spaces-duality-section-twisted-inverse-image},
\ref{spaces-duality-section-base-change-map},
\ref{spaces-duality-section-base-change-II},
\ref{spaces-duality-section-trace},
\ref{spaces-duality-section-compare-with-pullback} et
\ref{spaces-duality-section-proper-flat}.
\end{definition}

\noindent
Il convient de faire ici plusieurs remarques.
\begin{enumerate}
\item Dans la Définition \ref{spaces-duality-definition-relative-dualizing-proper-flat},
on peut supprimer l'hypothèse que $\omega_{X/Y}^\bullet$ est $Y$-parfait.
En effet, en faisant parcourir à $Y'$ les membres d'un recouvrement \'etale de $Y$
par des affines, le Lemme \ref{spaces-duality-lemma-properties-relative-dualizing} montre
que les restrictions de $\omega_{X/Y}^\bullet$ aux membres d'un
recouvrement \'etale de $X$ sont $Y$-parfaites, ce qui implique que
$\omega_{X/Y}^\bullet$ est $Y$-parfait ; voir
Compléments sur les morphismes d'espaces, section
\ref{spaces-more-morphisms-section-relatively-perfect}.
\item Considérons un complexe dualisant relatif
$(\omega_{X/Y}^\bullet, \tau)$ et un carré cartésien comme dans la
Définition \ref{spaces-duality-definition-relative-dualizing-proper-flat}.
Nous interpréterons l'existence de l'isomorphisme
$(L(g')^*\omega_{X/Y}^\bullet, Lg^*\tau) \cong
(a'(\mathcal{O}_{Y'}), \text{Tr}_{f', \mathcal{O}_{Y'}})$
comme suit : elle affirme que, pour tout $M' \in D_\QCoh(\mathcal{O}_{X'})$,
le morphisme
$$
\Hom_{X'}(M', L(g')^*\omega_{X/Y}^\bullet)
\longrightarrow
\Hom_{Y'}(Rf'_*M', \mathcal{O}_{Y'}),\quad
\varphi' \longmapsto Lg^*\tau \circ Rf'_*\varphi'
$$
est un isomorphisme. Cela résulte de la définition de $a'$
et de la discussion de la section \ref{spaces-duality-section-trace}. En particulier,
le lemme de Yoneda garantit l'unicité de l'isomorphisme.
\item Si $Y$ est lui-même affine, alors un complexe dualisant relatif
$(\omega_{X/Y}^\bullet, \tau)$ existe et est canoniquement isomorphe
à $(a(\mathcal{O}_Y), \text{Tr}_{f, \mathcal{O}_Y})$, où
$a$ est l'adjoint à droite de $Rf_*$ comme dans le
Lemme \ref{spaces-duality-lemma-twisted-inverse-image}
et $\text{Tr}_f$ est comme dans la section \ref{spaces-duality-section-trace}.
En effet, un diagramme comme dans la définition donne
un isomorphisme $L(g')^*a(\mathcal{O}_Y) \to a'(\mathcal{O}_{Y'})$ d'après le
Lemme \ref{spaces-duality-lemma-more-base-change},
compatible aux morphismes trace d'après le
Lemme \ref{spaces-duality-lemma-trace-map-and-base-change}.
\end{enumerate}
Cela fournit exactement les données nécessaires pour recoller les
complexes dualisants relatifs donnés localement en un complexe global. Nous suggérons au lecteur
de sauter les démonstrations des lemmes suivants.

\begin{lemma}
\label{spaces-duality-lemma-relative-dualizing-RHom}
Soit $S$ un schéma. Soit $X \to Y$ un morphisme propre et plat
d'espaces algébriques, de présentation finie.
Si $(\omega_{X/Y}^\bullet, \tau)$ est un complexe dualisant relatif,
alors $\mathcal{O}_X \to
R\SheafHom_{\mathcal{O}_X}(\omega_{X/Y}^\bullet, \omega_{X/Y}^\bullet)$
est un isomorphisme et les faisceaux de cohomologie de $Rf_*\omega_{X/Y}^\bullet$
sont nuls en degrés positifs.
\end{lemma}

\begin{proof}
Il suffit de le démontrer après changement de base par un schéma affine \'etale
sur $Y$ ; cela résulte alors du
Lemme \ref{spaces-duality-lemma-properties-relative-dualizing}.
\end{proof}

\begin{lemma}
\label{spaces-duality-lemma-uniqueness-relative-dualizing}
Soit $S$ un schéma. Soit $X \to Y$ un morphisme propre et plat
d'espaces algébriques, de présentation finie.
Si $(\omega_j^\bullet, \tau_j)$, $j = 1, 2$,
sont deux complexes dualisants relatifs sur $X/Y$,
alors il existe un unique isomorphisme
$(\omega_1^\bullet, \tau_1) \to (\omega_2^\bullet, \tau_2)$.
\end{lemma}

\begin{proof}
Considérons $g : Y' \to Y$ \'etale, avec $Y'$ un schéma affine,
et notons $X' = Y' \times_Y X$ le changement de base.
D'après la Définition \ref{spaces-duality-definition-relative-dualizing-proper-flat}
et la discussion qui la suit, il existe un unique isomorphisme
$\iota : (\omega_1^\bullet|_{X'}, \tau_1|_{Y'}) \to
(\omega_2^\bullet|_{X'}, \tau_2|_{Y'})$. Si $Y'' \to Y'$
est un autre morphisme \'etale d'affines et $X'' = Y'' \times_Y X$,
alors $\iota|_{X''}$ est l'unique isomorphisme
$(\omega_1^\bullet|_{X''}, \tau_1|_{Y''}) \to
(\omega_2^\bullet|_{X''}, \tau_2|_{Y''})$ (par unicité).
De plus, nous avons
$$
\text{Ext}^p_{X'}(\omega_1^\bullet|_{X'}, \omega_2^\bullet|_{X'}) = 0,
\quad p < 0
$$
car
$\mathcal{O}_{X'} \cong
R\SheafHom_{\mathcal{O}_{X'}}(\omega_1^\bullet|_{X'}, \omega_1^\bullet|_{X'})
\cong
R\SheafHom_{\mathcal{O}_{X'}}(\omega_1^\bullet|_{X'}, \omega_2^\bullet|_{X'})$
d'après le Lemme \ref{spaces-duality-lemma-relative-dualizing-RHom}.

\medskip\noindent
Choisissons un hyperrecouvrement \'etale $b : V \to Y$ tel que chaque
$V_n = \coprod_{i \in I_n} Y_{n, i}$, avec $Y_{n, i}$ affine.
C'est possible d'après Hyperrecouvrements, Lemme
\ref{hypercovering-lemma-hypercovering-object} et
Remarque \ref{hypercovering-remark-take-unions-hypercovering-X}
(pour remplacer l'hyperrecouvrement produit
dans le lemme par celui qui possède des réunions disjointes en chaque degré).
Notons $X_{n, i} = Y_{n, i} \times_Y X$ et $U_n = V_n \times_Y X$,
de sorte que nous obtenons un hyperrecouvrement \'etale
$a : U \to X$ (Hyperrecouvrements, Lemme
\ref{hypercovering-lemma-hypercovering-morphism-sites})
avec $U_n = \coprod X_{n, i}$.
Les hypothèses d'Espaces simpliciaux, Lemme
\ref{spaces-simplicial-lemma-fppf-neg-ext-zero-hom},
sont satisfaites pour $a : U \to X$ et les complexes
$\omega_1^\bullet$ et $\omega_2^\bullet$.
Nous obtenons donc un unique morphisme
$\iota : \omega_1^\bullet \to \omega_2^\bullet$
dont la restriction à $X_{0, i}$ est l'unique
isomorphisme $(\omega_1^\bullet|_{X_{0, i}}, \tau_1|_{Y_{0, i}}) \to
(\omega_2^\bullet|_{X_{0, i}}, \tau_2|_{Y_{0, i}})$.
Il reste à vérifier que le diagramme
$$
\xymatrix{
Rf_*\omega_1^\bullet \ar[rd]_{\tau_1} \ar[rr]_{Rf_*\iota} & &
Rf_*\omega_1^\bullet \ar[ld]^{\tau_2} \\
& \mathcal{O}_Y
}
$$
est commutatif. Or nous savons que $Rf_*\omega_1^\bullet$ et
$Rf_*\omega_2^\bullet$ ont leurs faisceaux de cohomologie nuls en degrés
positifs (Lemme \ref{spaces-duality-lemma-relative-dualizing-RHom}) ;
cette commutativité peut donc se vérifier après
restriction aux affines $Y_{0, i}$, où elle vaut par construction.
\end{proof}

\begin{lemma}
\label{spaces-duality-lemma-covering-enough}
Soit $S$ un schéma. Soit $X \to Y$ un morphisme propre et plat
d'espaces algébriques, de présentation finie.
Soit $(\omega^\bullet, \tau)$ un couple formé
d'un objet $Y$-parfait de $D(\mathcal{O}_X)$ et d'un morphisme
$\tau : Rf_*\omega^\bullet \to \mathcal{O}_Y$.
Supposons donnés des diagrammes cartésiens
$$
\xymatrix{
X_i \ar[r]_{g_i'} \ar[d]_{f_i} & X \ar[d]^f \\
Y_i \ar[r]^{g_i} & Y
}
$$
avec $Y_i$ affine, tels que $\{g_i : Y_i \to Y\}$ soit un recouvrement \'etale,
et des isomorphismes de couples $(\omega^\bullet|_{X_i}, \tau|_{Y_i})
\to (a_i(\mathcal{O}_{Y_i}), \text{Tr}_{f_i, \mathcal{O}_{Y_i}})$
comme dans la Définition \ref{spaces-duality-definition-relative-dualizing-proper-flat}.
Alors $(\omega^\bullet, \tau)$ est un complexe dualisant relatif pour $X$ sur $Y$.
\end{lemma}

\begin{proof}
Soient $g : Y' \to Y$ et $X', f', g', a'$ comme dans la
Définition \ref{spaces-duality-definition-relative-dualizing-proper-flat}.
Posons $((\omega')^\bullet, \tau') = (L(g')^*\omega^\bullet, Lg^*\tau)$.
Nous pouvons trouver un recouvrement \'etale fini
$\{Y'_j \to Y'\}$ par des affines qui raffine $\{Y_i \times_Y Y' \to Y'\}$
(Topologies, Lemme \ref{topologies-lemma-etale-affine}).
Ainsi, pour chaque $j$, il existe un $i_j$ et un morphisme
$k_j : Y'_j \to Y_{i_j}$ sur $Y$. Considérons les produits fibrés
$$
\xymatrix{
X'_j \ar[r]_{h_j'} \ar[d]_{f'_j} &
X' \ar[d]^{f'} \\
Y'_j \ar[r]^{h_j} & Y'
}
$$
Notons $k'_j : X'_j \to X_{i_j}$ le morphisme induit (changement de base
de $k_j$ par $f_{i_j}$). En restreignant les isomorphismes donnés à $Y'_j$
par le morphisme $k'_j$, nous obtenons des isomorphismes de couples
$((\omega')^\bullet|_{X'_j}, \tau'|_{Y'_j})
\to (a_j(\mathcal{O}_{Y'_j}), \text{Tr}_{f'_j, \mathcal{O}_{Y'_j}})$.
Après avoir remplacé $f : X \to Y$ par $f' : X' \to Y'$, nous nous ramenons au
problème résolu dans le paragraphe suivant.

\medskip\noindent
Supposons $Y$ affine. Problème : montrer que $(\omega^\bullet, \tau)$ est
isomorphe à $(\omega_{X/Y}^\bullet, \text{Tr}) =
(a(\mathcal{O}_Y), \text{Tr}_{f, \mathcal{O}_Y})$.
Nous pouvons supposer que notre recouvrement $\{Y_i \to Y\}$ est donné par un unique
morphisme \'etale surjectif $\{g : Y' \to Y\}$ entre schémas affines.
En effet, nous pouvons d'abord remplacer $\{g_i: Y_i \to Y\}$ par un sous-recouvrement
fini, puis poser $g = \coprod g_i :  Y' = \coprod Y_i \to Y$ ;
nous omettons quelques détails. Posons $X' = Y' \times_Y X$, avec les morphismes
$f', g'$ de la Définition \ref{spaces-duality-definition-relative-dualizing-proper-flat}.
La seule donnée qui subsiste est alors un isomorphisme
$$
(\omega^\bullet|_{X'}, \tau|_{Y'}) \to
(a'(\mathcal{O}_{Y'}), \text{Tr}_{f', \mathcal{O}_{Y'}})
$$
Puisque $(\omega_{X/Y}^\bullet, \text{Tr})$ est un complexe dualisant relatif
(voir la discussion qui suit la
Définition \ref{spaces-duality-definition-relative-dualizing-proper-flat}),
il existe un unique isomorphisme
$$
(\omega_{X/Y}^\bullet|_{X'}, \text{Tr}|_{Y'}) \to
(a'(\mathcal{O}_{Y'}), \text{Tr}_{f', \mathcal{O}_{Y'}})
$$
L'unicité résulte par exemple du Lemme \ref{spaces-duality-lemma-uniqueness-relative-dualizing}.
En combinant les isomorphismes affichés, nous trouvons un isomorphisme
$$
\alpha :
(\omega^\bullet|_{X'}, \tau|_{Y'}) \to
(\omega_{X/Y}^\bullet|_{X'}, \text{Tr}|_{Y'})
$$
Posons $Y'' = Y' \times_Y Y'$ et $X'' = Y'' \times_Y X$ ; les deux
images réciproques de $\alpha$ sur $X''$ doivent être égales, encore par unicité.
Comme les auto-Ext négatifs de
$\omega_{X'/Y'}^\bullet$ sur $X'$ sont nuls (Lemme \ref{spaces-duality-lemma-relative-dualizing-RHom})
et que cette propriété subsiste après image réciproque par toute projection
$Y' \times_Y \ldots \times_Y Y' \to Y'$ (nous omettons un petit détail -- comparer
avec la démonstration du Lemme \ref{spaces-duality-lemma-uniqueness-relative-dualizing}),
nous trouvons que $\alpha$ descend en un isomorphisme
$\omega^\bullet \to \omega_{X/Y}^\bullet$
sur $X$ d'après Espaces simpliciaux, Lemme
\ref{spaces-simplicial-lemma-fppf-neg-ext-zero-hom}.
\end{proof}

\begin{lemma}
\label{spaces-duality-lemma-existence-relative-dualizing}
Soit $S$ un schéma. Soit $X \to Y$ un morphisme propre et plat
d'espaces algébriques, de présentation finie.
Il existe un complexe dualisant relatif $(\omega_{X/Y}^\bullet, \tau)$.
\end{lemma}

\begin{proof}
Choisissons un hyperrecouvrement \'etale $b : V \to Y$ tel que chaque
$V_n = \coprod_{i \in I_n} Y_{n, i}$, avec $Y_{n, i}$ affine.
C'est possible d'après Hyperrecouvrements, Lemme
\ref{hypercovering-lemma-hypercovering-object} et
Remarque \ref{hypercovering-remark-take-unions-hypercovering-X}
(pour remplacer l'hyperrecouvrement produit
dans le lemme par celui qui possède des réunions disjointes en chaque degré).
Notons $X_{n, i} = Y_{n, i} \times_Y X$ et $U_n = V_n \times_Y X$,
de sorte que nous obtenons un hyperrecouvrement \'etale
$a : U \to X$ (Hyperrecouvrements, Lemme
\ref{hypercovering-lemma-hypercovering-morphism-sites})
avec $U_n = \coprod X_{n, i}$.
Pour chaque $n, i$, il existe un complexe dualisant relatif
$(\omega_{n, i}^\bullet, \tau_{n, i})$ sur $X_{n, i}/Y_{n, i}$.
Voir la discussion qui suit la
Définition \ref{spaces-duality-definition-relative-dualizing-proper-flat}.
Pour $\varphi : [m] \to [n]$ et $i \in I_n$, considérons les morphismes
$g_{\varphi, i} : Y_{n, i} \to Y_{m, \alpha(\varphi)}$ et
$g'_{\varphi, i} : X_{n, i} \to X_{m, \alpha(\varphi)}$
qui font partie de la structure des hyperrecouvrements donnés
(Hyperrecouvrements, section \ref{hypercovering-section-hypercovering-sites}).
Nous avons alors un unique isomorphisme
$$
\iota_{n, i, \varphi} :
(L(g'_{n, i})^*\omega_{n, i}^\bullet, Lg_{n, i}^*\tau_{n, i})
\longrightarrow
(\omega_{m, \alpha(\varphi)(i)}^\bullet, \tau_{m, \alpha(\varphi)(i)})
$$
de couples ; voir la discussion qui suit la
Définition \ref{spaces-duality-definition-relative-dualizing-proper-flat}.
Observons que les auto-Ext négatifs de $\omega_{n, i}^\bullet$
sur $X_{n, i}$ sont nuls d'après le Lemme \ref{spaces-duality-lemma-relative-dualizing-RHom}.
Notons $(\omega_n^\bullet, \tau_n)$ le couple sur $U_n/V_n$ construit à l'aide
des couples $(\omega_{n, i}^\bullet, \tau_{n, i})$ pour $i \in I_n$.
Pour $\varphi : [m] \to [n]$ et $i \in I_n$, considérons les morphismes
$g_\varphi : V_n \to V_m$ et $g'_\varphi : U_n \to U_m$ qui font partie
de la structure des espaces algébriques simpliciaux $V$ et $U$.
Nous avons alors des isomorphismes uniques
$$
\iota_\varphi :
(L(g'_\varphi)^*\omega_n^\bullet, Lg_\varphi^*\tau_n)
\longrightarrow
(\omega_m^\bullet, \tau_m)
$$
de couples, construits à partir des isomorphismes sur les morceaux.
L'unicité garantit que ces isomorphismes satisfont
à la condition de transitivité formulée dans
Espaces simpliciaux, Définition
\ref{spaces-simplicial-definition-cartesian-derived-modules}.
Les hypothèses d'Espaces simpliciaux, Lemme
\ref{spaces-simplicial-lemma-fppf-glue-neg-ext-zero},
sont satisfaites pour $a : U \to X$, les complexes $\omega_n^\bullet$
et les isomorphismes $\iota_\varphi$\footnote{Ce lemme utilise seulement
$\omega_0^\bullet$ et les deux morphismes $\delta_1^1, \delta_0^1 : [1] \to [0]$.
Le lecteur peut sauter les premières lignes de la démonstration du
lemme cité,
car nous disposons ici déjà d'un système simplicial de la
catégorie dérivée des modules.}.
Nous obtenons ainsi un objet $\omega^\bullet$ de $D_\QCoh(\mathcal{O}_X)$
muni d'un isomorphisme
$\iota_0 : \omega^\bullet|_{U_0} \to \omega_0^\bullet$
compatible aux deux isomorphismes
$\iota_{\delta^1_0}$ et $\iota_{\delta^1_1}$.
Enfin, nous appliquons
Espaces simpliciaux, Lemme
\ref{spaces-simplicial-lemma-fppf-neg-ext-zero-hom},
pour obtenir un unique morphisme
$$
\tau : Rf_*\omega^\bullet \longrightarrow \mathcal{O}_Y
$$
dont la restriction à $V_0$ coïncide avec $\tau_0$ ; nous omettons quelques détails --
comparer par exemple avec la fin de la démonstration du
Lemme \ref{spaces-duality-lemma-uniqueness-relative-dualizing},
pour voir pourquoi nous avons
l'annulation requise des Ext négatifs.
D'après le Lemme \ref{spaces-duality-lemma-covering-enough},
le couple $(\omega^\bullet, \tau)$ est un complexe dualisant relatif,
ce qui achève la démonstration.
\end{proof}

\begin{lemma}
\label{spaces-duality-lemma-base-change-relative-dualizing}
Soit $S$ un schéma. Considérons un carré cartésien
$$
\xymatrix{
X' \ar[d]_{f'} \ar[r]_{g'} & X \ar[d]^f \\
Y' \ar[r]^g & Y
}
$$
d'espaces algébriques sur $S$. Supposons que $X \to Y$ soit propre, plat et
de présentation finie. Soit $(\omega_{X/Y}^\bullet, \tau)$ un
complexe dualisant relatif pour $f$. Alors
$(L(g')^*\omega_{X/Y}^\bullet, Lg^*\tau)$ est un complexe dualisant
relatif pour $f'$.
\end{lemma}

\begin{proof}
Observons que $L(g')^*\omega_{X/Y}^\bullet$ est $Y'$-parfait d'après
Compléments sur les morphismes d'espaces, Lemme
\ref{spaces-more-morphisms-lemma-base-change-relatively-perfect}.
L'autre condition de la
Définition \ref{spaces-duality-definition-relative-dualizing-proper-flat}
résulte de la transitivité des produits fibrés.
\end{proof}






 
\section{Comparaison avec le cas des schémas}
\label{spaces-duality-section-comparison}

\noindent
Il faudrait ajouter beaucoup plus de choses dans cette section.

\begin{lemma}
\label{spaces-duality-lemma-compare}
Soit $S$ un schéma. Soit $f : X \to Y$ un morphisme d'espaces algébriques
quasi-compacts et quasi-séparés sur $S$.
Supposons $X$ et $Y$ représentables, et soit $f_0 : X_0 \to Y_0$ un
morphisme de schémas représentant $f$ (notation peu élégante, mais temporaire).
Soit $a : D_\QCoh(\mathcal{O}_Y) \to D_\QCoh(\mathcal{O}_X)$
l'adjoint à droite de $Rf_*$ du Lemme \ref{spaces-duality-lemma-twisted-inverse-image}.
Soit $a_0 : D_\QCoh(\mathcal{O}_{Y_0}) \to D_\QCoh(\mathcal{O}_{X_0})$
l'adjoint à droite de $Rf_*$ de
Dualité pour les schémas, Lemme \ref{duality-lemma-twisted-inverse-image}.
Alors
$$
\xymatrix{
D_\QCoh(\mathcal{O}_{X_0})
\ar@{=}[rrrrrr]_{\text{Catégories dérivées des espaces, Lemme
\ref{spaces-perfect-lemma-derived-quasi-coherent-small-etale-site}}}
& & & & & &
D_\QCoh(\mathcal{O}_X) \\
D_\QCoh(\mathcal{O}_{Y_0}) \ar[u]^{a_0}
\ar@{=}[rrrrrr]^{\text{Catégories dérivées des espaces, Lemme
\ref{spaces-perfect-lemma-derived-quasi-coherent-small-etale-site}}}
& & & & & &
D_\QCoh(\mathcal{O}_Y) \ar[u]_a
}
$$
est commutatif.
\end{lemma}

\begin{proof}
Cela résulte de l'unicité des adjoints et des compatibilités de
Catégories dérivées des espaces, Remarque
\ref{spaces-perfect-remark-match-total-direct-images}.
\end{proof}













% OK, commencer ici.
%

\chapter{Espaces algébriques formels}



\phantomsection
\label{formal-spaces-section-phantom}


\section{Introduction}
\label{formal-spaces-section-introduction}

\noindent
Les schémas formels ont été introduits dans \cite{EGA}. Une version plus générale
des schémas formels a été introduite dans \cite{McQuillan}, et une autre
dans \cite{Yasuda}. Les espaces algébriques formels ont été introduits dans \cite{Kn}.
On trouvera des résultats connexes, ainsi que bien d'autres, dans
\cite{Abbes} et \cite{Fujiwara-Kato}.
Ce chapitre introduit la notion d'espace algébrique formel
que nous utiliserons. Notre définition est assez générale pour que la plupart
des classes de schémas ou d'espaces formels considérées dans la littérature en soient des sous-catégories pleines.

\medskip\noindent
Bien que nous comparions certaines de ces théories alternatives
à la nôtre, nous ne donnons pas toujours tous les détails lorsqu'ils ne sont pas nécessaires
au développement logique de la théorie.

\medskip\noindent
Outre l'introduction des espaces algébriques formels, nous établissons aussi quelques
propriétés très élémentaires et étudions plusieurs types de morphismes.










\section{Schémas formels à la EGA}
\label{formal-spaces-section-formal-schemes-EGA}

\noindent
Dans cette section, nous rappelons la construction des schémas formels de \cite{EGA}.
Cette notion, bien que très utile en géométrie algébrique,
n'est pas toujours celle qu'il convient de considérer. Il serait peut-être plus juste de dire
que la mise en place de la théorie comporte un certain nombre de choix, alors que, suivant
le but poursuivi, d'autres choix pourraient être préférables. On trouve en effet dans la littérature
de nombreuses théories voisines
adaptées au problème que les auteurs souhaitent étudier. Néanmoins, l'un
des grands avantages de la théorie esquissée ici est qu'elle permet de travailler
avec des objets géométriques bien déterminés.

\medskip\noindent
Avant de commencer, signalons une difficulté concernant la condition de faisceau
pour les faisceaux d'anneaux topologiques ou, plus généralement, d'espaces
topologiques. En effet, les grandes catégories suivantes :
\begin{enumerate}
\item la catégorie des espaces topologiques,
\item la catégorie des groupes topologiques,
\item la catégorie des anneaux topologiques,
\item la catégorie des modules topologiques sur un anneau topologique donné,
\end{enumerate}
munies de leurs foncteurs d'oubli naturels vers $\textit{Ensembles}$, ne sont pas
des exemples de types de structures algébriques au sens de
Faisceaux, section \ref{sheaves-section-algebraic-structures}.
Nous ne pouvons donc pas leur appliquer sans précaution les méthodes développées dans ce
chapitre. En revanche, chacune des catégories
énumérées ci-dessus admet des limites et des égalisateurs, et le foncteur d'oubli
vers les ensembles, les groupes, les anneaux ou les modules commute à ceux-ci
(voir Topologie, lemmes \ref{topology-lemma-limits},
\ref{topology-lemma-topological-group-limits},
\ref{topology-lemma-topological-ring-limits}, et
\ref{topology-lemma-topological-module-limits}).
Nous pouvons donc définir la notion de
faisceau comme dans Faisceaux, définition
\ref{sheaves-definition-sheaf-values-in-category},
et le préfaisceau sous-jacent d'ensembles, de groupes, d'anneaux ou de modules
est un faisceau. La différence essentielle est que, pour un recouvrement ouvert
$U = \bigcup_{i \in I} U_i$, le diagramme
$$
\xymatrix{
\mathcal{F}(U) \ar[r]
&
\prod\nolimits_{i\in I}
\mathcal{F}(U_i)
\ar@<1ex>[r] \ar@<-1ex>[r]
&
\prod\nolimits_{(i_0, i_1) \in I \times I}
\mathcal{F}(U_{i_0} \cap U_{i_1})
}
$$
doit être un diagramme égalisateur dans la catégorie des espaces
topologiques, des groupes topologiques, des anneaux topologiques ou des modules topologiques,
c'est-à-dire que la première application identifie
$\mathcal{F}(U)$ à un sous-espace de $\prod_{i \in I} \mathcal{F}(U_i)$,
ce dernier étant muni de la topologie produit.

\medskip\noindent
La fibre $\mathcal{F}_x$ d'un faisceau $\mathcal{F}$
d'espaces topologiques, de groupes topologiques, d'anneaux topologiques ou de
modules topologiques en un point $x \in X$ est définie comme la limite inductive sur les
voisinages ouverts
$$
\mathcal{F}_x = \colim_{x\in U} \mathcal{F}(U)
$$
dans la catégorie correspondante. Cela revient à prendre
la limite inductive au niveau des ensembles, des groupes, des anneaux ou des modules
(voir Topologie, lemmes \ref{topology-lemma-colimits},
\ref{topology-lemma-topological-group-colimits},
\ref{topology-lemma-topological-ring-colimits}, et
\ref{topology-lemma-topological-module-colimits}),
mais cette limite est munie d'une topologie. Attention :
la topologie obtenue dépend de la catégorie dans laquelle on travaille ; voir
Exemples, section \ref{examples-section-colimit-topology}.
On peut associer un faisceau aux préfaisceaux d'espaces topologiques,
de groupes topologiques, d'anneaux topologiques ou de modules topologiques,
et la prise des fibres commute à cette opération ; voir
la remarque \ref{formal-spaces-remark-sheafification-of-presheaves-in-top}.

\medskip\noindent
Soit $f : X \to Y$ une application continue d'espaces topologiques.
Il existe un foncteur $f_*$ de la catégorie des faisceaux d'espaces
topologiques, de groupes topologiques, d'anneaux topologiques ou de modules topologiques
vers la catégorie correspondante des faisceaux sur $Y$, défini en posant, comme d'habitude,
$f_*\mathcal{F}(V) = \mathcal{F}(f^{-1}V)$.
(Nous remettons à plus tard l'étude de l'image inverse dans ce cadre.)
Nous définissons la notion de $f$-application $\xi : \mathcal{G} \to \mathcal{F}$
entre un faisceau d'espaces topologiques $\mathcal{G}$ sur $Y$ et
un faisceau d'espaces topologiques $\mathcal{F}$ sur $X$ exactement de la
même manière que dans Faisceaux, définition \ref{sheaves-definition-f-map},
avec la condition supplémentaire que
$\xi_V : \mathcal{G}(V) \to \mathcal{F}(f^{-1}V)$ soit continue
pour tout ouvert $V \subset Y$. On a
$$
\{f\text{-applications de }\mathcal{G}\text{ vers }\mathcal{F}\} =
\Mor_{\Sh(Y, \textit{Top})}(\mathcal{G}, f_*\mathcal{F})
$$
comme dans Faisceaux, lemme \ref{sheaves-lemma-f-map}. Il en va de même pour les
faisceaux de groupes topologiques, d'anneaux topologiques et de modules topologiques. Enfin,
soit $\xi : \mathcal{G} \to \mathcal{F}$ une $f$-application comme ci-dessus.
Étant donné $x \in X$, d'image $y = f(x)$, il existe une application
continue
$$
\xi_x : \mathcal{G}_y \longrightarrow \mathcal{F}_x
$$
entre les fibres, définie exactement de la même manière que dans la discussion
qui suit Faisceaux, définition \ref{sheaves-definition-composition-f-maps}.

\medskip\noindent
La discussion précédente permet de définir une catégorie $LTRS$ des
« espaces localement annelés topologiques ». Un objet est un couple
$(X, \mathcal{O}_X)$ formé d'un espace topologique
$X$ et d'un faisceau d'anneaux topologiques $\mathcal{O}_X$ dont les fibres
$\mathcal{O}_{X, x}$ sont des anneaux locaux (lorsqu'on oublie la topologie).
Un morphisme $(X, \mathcal{O}_X) \to (Y, \mathcal{O}_Y)$ dans
$LTRS$ est un couple $(f, f^\sharp)$, où $f : X \to Y$ est une application
continue d'espaces topologiques et $f^\sharp : \mathcal{O}_Y \to \mathcal{O}_X$
est une $f$-application telle que, pour tout $x \in X$, l'application induite
$$
f^\sharp_x : \mathcal{O}_{Y, f(x)} \longrightarrow \mathcal{O}_{X, x}
$$
soit un homomorphisme local d'anneaux locaux (en oubliant les topologies).
La composition se définit exactement comme celle des
morphismes d'espaces localement annelés.

\medskip\noindent
Supposons maintenant que l'espace topologique $X$ possède une base formée
d'ouverts quasi-compacts. Étant donné un faisceau $\mathcal{F}$ d'ensembles, de groupes,
d'anneaux ou de modules sur un anneau, on peut munir $\mathcal{F}$
d'une structure de faisceau d'espaces topologiques, de groupes topologiques,
d'anneaux topologiques ou de modules topologiques.
En effet, si $U \subset X$ est un ouvert quasi-compact,
nous munissons $\mathcal{F}(U)$ de la topologie discrète. Si $U \subset X$
est quelconque, nous choisissons un recouvrement ouvert $U = \bigcup_{i \in I} U_i$
par des ouverts quasi-compacts et munissons $\mathcal{F}(U)$
de la topologie induite par $\prod_{i \in I} \mathcal{F}(U_i)$
(comme l'impose la discussion précédente).
Le lecteur vérifiera (nous omettons les détails) que l'on obtient ainsi un faisceau d'espaces topologiques,
de groupes topologiques, d'anneaux topologiques ou de modules topologiques.
Nous dirons qu'un faisceau d'espaces topologiques, de groupes topologiques,
d'anneaux topologiques ou de modules topologiques est
{\it pseudo-discret} si la topologie de $\mathcal{F}(U)$ est
discrète pour tout ouvert quasi-compact $U \subset X$. Alors
la construction précédente est adjointe au foncteur d'oubli
et induit une équivalence entre la catégorie des faisceaux
d'ensembles et celle des faisceaux pseudo-discrets d'espaces topologiques
(et de même pour les groupes, les anneaux et les modules).

\medskip\noindent
Grothendieck et Dieudonné définissent d'abord les schémas formels affines.
Ceux-ci correspondent aux anneaux topologiques admissibles $A$ ; voir
Compléments d'algèbre, définition \ref{more-algebra-definition-topological-ring}.
Plus précisément, étant donné $A$, on considère un système fondamental $I_\lambda$ d'idéaux
de définition de l'anneau $A$. (Dans tout anneau topologique admissible,
la famille de tous les idéaux de définition forme un système fondamental.)
Pour tout $\lambda$, on peut considérer le
schéma $\Spec(A/I_\lambda)$. Si $I_\lambda \subset I_\mu$, le
morphisme induit
$$
\Spec(A/I_\mu) \to \Spec(A/I_\lambda)
$$
est un épaississement, car $I_\mu^n \subset I_\lambda$ pour un certain $n$.
On peut aussi le voir en remarquant que chacune des
applications
$$
\Spec(A/I_\lambda) \to \Spec(A)
$$
est un homéomorphisme sur l'ensemble des idéaux premiers ouverts de $A$.
Ceci motive la définition
$$
\text{Spf}(A) = \{\text{idéaux premiers ouverts }\mathfrak p \subset A\}
$$
muni de la topologie induite par $\Spec(A)$. Pour tout $\lambda$,
nous pouvons considérer le faisceau structural $\mathcal{O}_{\Spec(A/I_\lambda)}$
comme un faisceau sur $\text{Spf}(A)$. Soit $\mathcal{O}_\lambda$ le
faisceau pseudo-discret d'anneaux topologiques correspondant ; voir ci-dessus.
On pose alors
$$
\mathcal{O}_{\text{Spf}(A)} = \lim \mathcal{O}_\lambda
$$
où la limite est prise dans la catégorie des faisceaux d'anneaux topologiques.
Le couple $(\text{Spf}(A), \mathcal{O}_{\text{Spf}(A)})$ est appelé le
{\it spectre formel} de $A$.

\medskip\noindent
Il convient alors de vérifier plusieurs points. Premièrement,
les fibres $\mathcal{O}_{\text{Spf}(A), x}$ sont des anneaux locaux
(lorsqu'on oublie la topologie). Deuxièmement, étant donné
$f \in A$, pour l'ouvert correspondant $D(f) \cap \text{Spf}(A)$,
on a
$$
\Gamma(D(f) \cap \text{Spf}(A), \mathcal{O}_{\text{Spf}(A)})
= A_{\{f\}} = \lim (A/I_\lambda)_f
$$
comme anneaux topologiques, où $I_\lambda$ est un système fondamental d'idéaux
de définition comme ci-dessus. De plus, l'anneau $A_{\{f\}}$ est lui aussi admissible, et
$(\text{Spf}(A_f), \mathcal{O}_{\text{Spf}(A_{\{f\}})})$
est isomorphe à
$(D(f) \cap \text{Spf}(A),
\mathcal{O}_{\text{Spf}(A)}|_{D(f) \cap \text{Spf}(A)})$.
Enfin, pour une paire d'anneaux topologiques admissibles $A, B$,
on a
\begin{equation}
\label{formal-spaces-equation-morphisms-affine-formal-schemes}
\Mor_{LTRS}((\text{Spf}(B), \mathcal{O}_{\text{Spf}(B)}),
(\text{Spf}(A), \mathcal{O}_{\text{Spf}(A)}))
= \Hom_{cont}(A, B)
\end{equation}
où $LTRS$ est la catégorie des « espaces localement annelés topologiques »
définie ci-dessus.

\medskip\noindent
Cela étant, on définit dans \cite{EGA} un {\it schéma formel} comme un couple
$(\mathfrak X, \mathcal{O}_\mathfrak X)$ où $\mathfrak X$
est un espace topologique et $\mathcal{O}_\mathfrak X$ un faisceau
d'anneaux topologiques tel que tout point possède un voisinage ouvert
isomorphe (dans $LTRS$) à un schéma formel affine.
Un {\it morphisme de schémas formels}
$f : (\mathfrak X, \mathcal{O}_\mathfrak X) \to
(\mathfrak Y, \mathcal{O}_\mathfrak Y)$
est un morphisme de la catégorie $LTRS$.

\medskip\noindent
Soit $A$ un anneau muni de la topologie discrète. Alors $A$ est
admissible et le schéma formel $\text{Spf}(A)$ est égal à
$\Spec(A)$. Le faisceau structural $\mathcal{O}_{\text{Spf}(A)}$
est le faisceau pseudo-discret d'anneaux topologiques associé
à $\mathcal{O}_{\Spec(A)}$ ; autrement dit, son faisceau
d'anneaux sous-jacent est égal à $\mathcal{O}_{\Spec(A)}$, et l'anneau
$\mathcal{O}_{\text{Spf}(A)}(U) = \mathcal{O}_{\Spec(A)}(U)$
sur un ouvert quasi-compact $U$ est muni de la topologie discrète,
mais pas en général. On peut donc associer à tout schéma affine
un schéma formel affine. De la même manière, on peut partir
d'un schéma quelconque $(X, \mathcal{O}_X)$ et lui associer
$(X, \mathcal{O}'_X)$, où $\mathcal{O}'_X$ est le
faisceau pseudo-discret d'anneaux topologiques dont le faisceau
d'anneaux sous-jacent est $\mathcal{O}_X$. Cette construction est
compatible aux morphismes et définit un foncteur
\begin{equation}
\label{formal-spaces-equation-compare-schemes-formal-schemes}
\textit{Schémas} \longrightarrow \textit{Schémas formels}
\end{equation}
Il résulte immédiatement de
(\ref{formal-spaces-equation-morphisms-affine-formal-schemes})
que ce foncteur est pleinement fidèle.

\medskip\noindent
Soit $\mathfrak X$ un schéma formel. Définissons la {\it taille}
du schéma formel par la formule
$\text{size}(\mathfrak X) = \max(\aleph_0, \kappa_1, \kappa_2)$,
où $\kappa_1$ est le cardinal de l'ensemble des ouverts formels affines de
$\mathfrak X$ et $\kappa_2$ le suprémum des cardinaux
des $\mathcal{O}_\mathfrak X(\mathfrak U)$, où
$\mathfrak U \subset \mathfrak X$ est un tel ouvert formel affine.

\begin{lemma}
\label{formal-spaces-lemma-fully-faithful}
Choisissons une catégorie de schémas $\Sch_\alpha$
comme dans Ensembles, lemme \ref{sets-lemma-construct-category}.
À tout schéma formel $\mathfrak X$, associons
$$
h_\mathfrak X : (\Sch_\alpha)^{opp} \longrightarrow \textit{Ensembles},\quad
h_\mathfrak X(S) = \Mor_{\textit{Schémas formels}}(S, \mathfrak X)
$$
son foncteur des points. Alors on a
$$
\Mor_{\textit{Schémas formels}}(\mathfrak X, \mathfrak Y) =
\Mor_{\textit{PSh}(\Sch_\alpha)}(h_\mathfrak X, h_\mathfrak Y)
$$
pourvu que la taille de $\mathfrak X$ ne soit pas trop grande.
\end{lemma}

\begin{proof}
Observons d'abord que $h_\mathfrak X$ vérifie la condition de faisceau pour
la topologie de Zariski, quel que soit le schéma formel $\mathfrak X$ (voir
Schémas, définition \ref{schemes-definition-representable-by-open-immersions}).
Cela découle du caractère local des morphismes dans la catégorie
des schémas formels. En outre, pour une immersion ouverte
$\mathfrak V \to \mathfrak W$ de schémas formels,
la transformation de foncteurs correspondante $h_\mathfrak V \to h_\mathfrak W$
est injective et représentable par des immersions ouvertes (voir
Schémas, définition \ref{schemes-definition-representable-by-open-immersions}).
Choisissons un recouvrement ouvert $\mathfrak X = \bigcup \mathfrak U_i$
d'un schéma formel par des schémas formels affines $\mathfrak U_i$.
La famille de foncteurs
$h_{\mathfrak U_i}$ recouvre alors $h_\mathfrak X$ (voir
Schémas, définition \ref{schemes-definition-representable-by-open-immersions}).
Enfin, remarquons que
$$
h_{\mathfrak U_i} \times_{h_\mathfrak X} h_{\mathfrak U_j} =
h_{\mathfrak U_i \cap \mathfrak U_j}
$$
Par conséquent, la donnée d'une application $h_\mathfrak X \to h_\mathfrak Y$
équivaut à celle d'une famille d'applications
$h_{\mathfrak U_i} \to h_\mathfrak Y$ qui coïncident sur les intersections.
Nous pouvons donc ramener la bijectivité (resp. l'injectivité) de l'application
du lemme à la bijectivité (resp. l'injectivité) pour les couples
$(\mathfrak U_i, \mathfrak Y)$
et à l'injectivité (resp. aucune condition)
pour $(\mathfrak U_i \cap \mathfrak U_j, \mathfrak Y)$.
On se ramène ainsi au cas où $\mathfrak X$ est un
schéma formel affine. Écrivons $\mathfrak X = \text{Spf}(A)$
pour un anneau topologique admissible $A$. Choisissons aussi un
système fondamental d'idéaux de définition $I_\lambda \subset A$.

\medskip\noindent
Nous pouvons aussi nous localiser sur $\mathfrak Y$.
Supposons en effet que $\mathfrak V \subset \mathfrak Y$ soit un
sous-schéma formel ouvert, et soit $\varphi : h_\mathfrak X \to h_\mathfrak Y$.
Alors
$$
h_\mathfrak V \times_{h_\mathfrak Y, \varphi} h_\mathfrak X \to h_\mathfrak X
$$
est représentable par des immersions ouvertes. Après changement de base à
$\Spec(A/I_\lambda)$ pour tout $\lambda$, on obtient un sous-schéma ouvert
$U_\lambda \subset \Spec(A/I_\lambda)$. Toutefois, si
$I_\lambda \subset I_\mu$, le morphisme $\Spec(A/I_\lambda) \to \Spec(A/I_\mu)$
transforme par image inverse $U_\mu$ en $U_\lambda$. Ceux-ci se recollent donc pour donner
un sous-schéma formel ouvert $\mathfrak U \subset \mathfrak X$.
Un argument immédiat (omis) montre que
$$
h_\mathfrak U = h_\mathfrak V \times_{h_\mathfrak Y} h_\mathfrak X
$$
On voit ainsi qu'un recouvrement ouvert
$\mathfrak Y = \bigcup \mathfrak V_j$ et une transformation
de foncteurs $\varphi :  h_\mathfrak X \to h_\mathfrak Y$
déterminent un recouvrement ouvert correspondant de $\mathfrak X$.
Comme $\mathfrak X$ est affine, on peut raffiner ce recouvrement en un
recouvrement ouvert fini
$\mathfrak X = \mathfrak U_1 \cup \ldots \cup \mathfrak U_n$
par des sous-schémas formels affines. Autrement dit, pour tout $i$, il
existe un $j$ et une application $\varphi_i : h_{\mathfrak U_i} \to h_{\mathfrak V_j}$
tels que le diagramme
$$
\xymatrix{
h_{\mathfrak U_i} \ar[r]_{\varphi_i} \ar[d] & h_{\mathfrak V_j} \ar[d] \\
h_{\mathfrak X} \ar[r]^\varphi & h_\mathfrak Y
}
$$
soit commutatif. Quelques arguments supplémentaires (que nous omettons) montrent alors
qu'il suffit d'établir la bijectivité du lemme lorsque
$\mathfrak X$ et $\mathfrak Y$ sont tous deux des schémas formels affines.

\medskip\noindent
Supposons que $\mathfrak X$ et $\mathfrak Y$ soient des schémas formels affines.
Écrivons $\mathfrak X = \text{Spf}(A)$ et $\mathfrak Y = \text{Spf}(B)$.
Soit $\varphi : h_\mathfrak X \to h_\mathfrak Y$ une transformation
de foncteurs. Soit $I_\lambda \subset A$ un système fondamental
d'idéaux de définition. Le morphisme d'inclusion canonique
$i_\lambda : \Spec(A/I_\lambda) \to \mathfrak X$ est envoyé sur un morphisme
$\varphi(i_\lambda) : \Spec(A/I_\lambda) \to \mathfrak Y$.
D'après (\ref{formal-spaces-equation-morphisms-affine-formal-schemes}), celui-ci correspond
à une application continue $\chi_\lambda : B \to A/I_\lambda$.
Comme $\varphi$ est une transformation de foncteurs, il s'ensuit
que, pour $I_\lambda \subset I_\mu$, le composé
$B \to A/I_\lambda \to A/I_\mu$ est égal à $\chi_\mu$.
Autrement dit, on obtient un homomorphisme d'anneaux
$$
\chi = \lim \chi_\lambda : B \longrightarrow \lim A/I_\lambda = A
$$
C'est un homomorphisme continu, car l'image réciproque
de $I_\lambda$ est ouverte pour tout $\lambda$ (puisque $A/I_\lambda$ est muni de la topologie
discrète et que $\chi_\lambda$ est continue). On obtient donc
un morphisme $\text{Spf}(\chi) : \mathfrak X \to \mathfrak Y$ par
(\ref{formal-spaces-equation-morphisms-affine-formal-schemes}).
Nous omettons la vérification que cette construction est l'inverse
de l'application du lemme dans ce cas.

\medskip\noindent
Remarques ensemblistes. Pour que cela fonctionne sur la catégorie donnée
de schémas $\Sch_\alpha$, il suffit de s'assurer que tous les
schémas utilisés dans la démonstration ci-dessus sont isomorphes à des objets de $\Sch_\alpha$.
En fait, une analyse attentive montre qu'il suffit que les
schémas $\Spec(A/I_\lambda)$ intervenant ci-dessus soient isomorphes à des
objets de $\Sch_\alpha$. Pour cela, il suffit certainement de supposer que
la taille de $\mathfrak X$ est au plus celle
d'un schéma appartenant à $\Sch_\alpha$.
\end{proof}

\begin{lemma}
\label{formal-spaces-lemma-formal-scheme-sheaf-fppf}
\begin{slogan}
Les schémas formels sont des faisceaux fpqc
\end{slogan}
Soit $\mathfrak X$ un schéma formel. Le foncteur des points
$h_\mathfrak X$ (voir le lemme \ref{formal-spaces-lemma-fully-faithful})
vérifie la condition de faisceau pour les recouvrements fpqc.
\end{lemma}

\begin{proof}
Topologies, lemme \ref{topologies-lemma-sheaf-property-fpqc}
nous ramène au cas d'un recouvrement de Zariski et d'un recouvrement
$\{\Spec(S) \to \Spec(R)\}$ où $R \to S$ est fidèlement plat.
Nous avons observé dans la démonstration du lemme \ref{formal-spaces-lemma-fully-faithful} 
que $h_\mathfrak X$ vérifie la condition de faisceau pour les recouvrements de Zariski.

\medskip\noindent
Supposons que $R \to S$ soit un homomorphisme d'anneaux fidèlement plat.
Notons $\pi : \Spec(S) \to \Spec(R)$ le
morphisme de schémas correspondant. Il est surjectif et plat.
Soit $f : \Spec(S) \to \mathfrak X$ un morphisme
tel que $f \circ \text{pr}_1 = f \circ \text{pr}_2$
comme applications $\Spec(S \otimes_R S) \to \mathfrak X$.
D'après Descente, lemme \ref{descent-lemma-equiv-fibre-product},
l'application $f$ entre les ensembles sous-jacents
est de la forme $f = g \circ \pi$ pour une certaine
application (ensembliste) $g : \Spec(R) \to \mathfrak X$.
D'après Morphismes, lemme \ref{morphisms-lemma-fpqc-quotient-topology},
et puisque $f$ est continue, on voit que $g$
est continue.

\medskip\noindent
Choisissons $y \in \Spec(R)$. Prenons un sous-schéma formel ouvert affine
$\mathfrak U \subset \mathfrak X$ contenant $g(y)$.
Écrivons $\mathfrak U = \text{Spf}(A)$ pour un anneau topologique
admissible $A$. D'après ce qui précède, on peut choisir $r \in R$ tel que
$y \in D(r) \subset g^{-1}(\mathfrak U)$.
La restriction de $f$ à $\pi^{-1}(D(r))$, à valeurs dans $\mathfrak U$,
correspond à une application continue d'anneaux $A \to S_r$ d'après
(\ref{formal-spaces-equation-morphisms-affine-formal-schemes}). Les deux homomorphismes d'anneaux induits
$A \to S_r \otimes_{R_r} S_r = (S \otimes_R S)_r$ sont égaux
par l'hypothèse sur $f$.
Remarquons que $R_r \to S_r$ est fidèlement plat.
D'après Descente, lemme \ref{descent-lemma-ff-exact}, l'égalisateur des
deux flèches $S_r \to S_r \otimes_{R_r} S_r$ est $R_r$.
On en déduit que $A \to S_r$ se factorise de manière unique par une application $A \to R_r$
qui est également continue, car elle a le même noyau (ouvert) que
l'application $A \to S_r$. Cette application donne à son tour un morphisme $D(r) \to \mathfrak U$ par
(\ref{formal-spaces-equation-morphisms-affine-formal-schemes}).

\medskip\noindent
Qu'avons-nous démontré jusqu'ici ? Pour tout $y \in \Spec(R)$,
il existe un ouvert affine standard
$y \in D(r) \subset \Spec(R)$ tel que le morphisme
$f|_{\pi^{-1}(D(r))} : \pi^{-1}(D(r)) \to \mathfrak X$ se factorise de manière unique
par un certain morphisme $D(r) \to \mathfrak X$. Nous omettons la
vérification que ces morphismes se recollent pour donner le
morphisme recherché $\Spec(R) \to \mathfrak X$.
\end{proof}

\begin{remark}[Variante de McQuillan]
\label{formal-spaces-remark-mcquillan}
Il existe une variante de la construction des schémas formels due à
McQuillan ; voir \cite{McQuillan}.
Il propose d'affaiblir légèrement la condition d'admissibilité.
Rappelons en effet qu'un anneau topologique admissible est un anneau topologique $A$
complet (et séparé selon nos conventions),
linéairement topologisé, qui possède un
idéal de définition, c'est-à-dire un
idéal ouvert $I$ tel que tout voisinage de $0$ contienne $I^n$
pour un certain $n \geq 1$.
McQuillan travaille avec ce que nous appellerons des anneaux topologiques
{\it faiblement admissibles}. Un anneau topologique faiblement admissible $A$ est un anneau topologique
complet (et séparé selon nos conventions),
linéairement topologisé, qui possède un
{\it idéal faible de définition}, c'est-à-dire un idéal ouvert $I$ tel que,
pour tout $f \in I$, on ait
$f^n \to 0$ lorsque $n \to \infty$. Comme dans le cas admissible,
si $I$ est un idéal faible de définition et si $J \subset A$ est un
idéal ouvert, alors $I \cap J$ est un idéal faible de définition.
Les idéaux faibles de définition forment donc un système fondamental de
voisinages ouverts de $0$, et
l'on peut suivre essentiellement la même démarche que ci-dessus
pour définir une catégorie plus grande de schémas formels fondée
sur cette notion. Les analogues des lemmes \ref{formal-spaces-lemma-fully-faithful} et
\ref{formal-spaces-lemma-formal-scheme-sheaf-fppf}
restent valables dans ce cadre (avec la même démonstration).
\end{remark}

\begin{remark}[Faisceau associé à un préfaisceau d'espaces topologiques]
\label{formal-spaces-remark-sheafification-of-presheaves-in-top}
\begin{reference}
\cite{Gray}
\end{reference}
Dans cette remarque, nous étudions brièvement la construction du faisceau associé à un préfaisceau
d'espaces topologiques. Les mêmes arguments s'appliquent exactement aux
préfaisceaux de groupes abéliens topologiques, d'anneaux topologiques et de
modules topologiques (sur un anneau topologique donné). Pour traiter
la question dans toute sa généralité, travaillons sur un site
$\mathcal{C}$. Le lecteur qui s'intéresse au cas des (pré)faisceaux
sur un espace topologique $X$ doit voir les objets de $\mathcal{C}$
comme les ouverts de $X$, les morphismes de $\mathcal{C}$ comme les inclusions
d'ouverts, et les recouvrements de $\mathcal{C}$ comme les recouvrements de $X$ ; voir
Sites, exemple \ref{sites-example-site-topological}.
Notons $\Sh(\mathcal{C}, \textit{Top})$ la catégorie des faisceaux
d'espaces topologiques sur $\mathcal{C}$ et
$\textit{PSh}(\mathcal{C}, \textit{Top})$ celle des préfaisceaux
d'espaces topologiques sur $\mathcal{C}$.
Soit $\mathcal{F}$ un préfaisceau d'espaces topologiques sur $\mathcal{C}$.
Le faisceau associé $\mathcal{F}^\#$ doit vérifier la formule
$$
\Mor_{\textit{PSh}(\mathcal{C}, \textit{Top})}(\mathcal{F}, \mathcal{G})
=
\Mor_{\Sh(\mathcal{C}, \textit{Top})}(\mathcal{F}^\#, \mathcal{G})
$$
fonctoriellement en $\mathcal{G}$ dans $\Sh(\mathcal{C}, \textit{Top})$.
Autrement dit, nous cherchons à construire l'adjoint à gauche
du foncteur d'inclusion
$\Sh(\mathcal{C}, \textit{Top}) \to \textit{PSh}(\mathcal{C}, \textit{Top})$.
Affirmons d'abord que $\Sh(\mathcal{C}, \textit{Top})$ admet des limites
et que le foncteur d'inclusion commute à celles-ci.
En effet, étant donnés une catégorie $\mathcal{I}$ et un foncteur
$i \mapsto \mathcal{G}_i$ à valeurs dans $\Sh(\mathcal{C}, \textit{Top})$,
nous définissons simplement
$$
(\lim \mathcal{G}_i)(U) = \lim \mathcal{G}_i(U)
$$
où la limite est prise dans la catégorie des espaces topologiques
(Topologie, lemme \ref{topology-lemma-limits}). Cela définit un faisceau,
car les limites commutent aux limites
(Catégories, lemme \ref{categories-lemma-colimits-commute}),
et en particulier aux produits et aux égalisateurs (qui sont les
opérations utilisées dans l'axiome des faisceaux). Enfin, un morphisme
de préfaisceaux $\mathcal{F} \to \lim \mathcal{G}_i$ est
manifestement la même chose qu'un système compatible de morphismes
$\mathcal{F} \to \mathcal{G}_i$. Autrement dit, l'objet
$\lim \mathcal{G}_i$ est la limite dans la catégorie
des préfaisceaux d'espaces topologiques, et a fortiori dans la
catégorie des faisceaux d'espaces topologiques.
Notre seconde affirmation est que tout morphisme de préfaisceaux
$\mathcal{F} \to \mathcal{G}$, où $\mathcal{G}$ est un objet de
$\Sh(\mathcal{C}, \textit{Top})$, se factorise par un sous-faisceau
$\mathcal{G}' \subset \mathcal{G}$ de taille bornée.
Nous définissons ici la {\it taille} $|\mathcal{H}|$
d'un faisceau d'espaces topologiques $\mathcal{H}$ comme le cardinal
$\sup_{U \in \Ob(\mathcal{C})} |\mathcal{H}(U)|$.
Pour démontrer l'affirmation, posons
$$
\mathcal{G}'(U) =
\left\{
\quad
s \in \mathcal{G}(U)
\quad \middle| \quad
\begin{matrix}
\text{il existe un recouvrement }\{U_i \to U\}_{i \in I} \\
\text{tel que }
s|_{U_i} \in \Im(\mathcal{F}(U_i) \to \mathcal{G}(U_i))
\end{matrix}
\quad
\right\}
$$
Nous munissons $\mathcal{G}'(U)$ de la topologie induite.
Alors $\mathcal{G}'$ est un faisceau d'espaces topologiques (détails omis),
et $\mathcal{G}' \to \mathcal{G}$ est un morphisme par lequel
l'application donnée $\mathcal{F} \to \mathcal{G}$ se factorise. De plus,
la taille de $\mathcal{G}'$ est bornée par un cardinal
$\kappa$ qui ne dépend que de $\mathcal{C}$ et du préfaisceau $\mathcal{F}$
(indication : utiliser le fait que, selon nos conventions, les recouvrements de $\mathcal{C}$
forment un ensemble). En rassemblant ces observations, on voit
que les hypothèses de Catégories, théorème
\ref{categories-theorem-adjoint-functor}
sont vérifiées, et l'on obtient le foncteur faisceau associé comme adjoint à gauche
du foncteur d'inclusion des faisceaux dans les préfaisceaux.
Enfin, soit $p$ un point du
site $\mathcal{C}$ défini par un foncteur $u : \mathcal{C} \to \textit{Ensembles}$ ;
voir Sites, définition \ref{sites-definition-point}.
Pour un espace topologique $M$, le préfaisceau défini par la règle
$$
U \mapsto \text{Map}(u(U), M) = \prod\nolimits_{x \in u(U)} M
$$
et muni de la topologie produit est un faisceau d'espaces topologiques.
Le même argument que celui donné dans la démonstration de
Sites, lemme \ref{sites-lemma-point-pushforward-sheaf}, montre donc que
$\mathcal{F}_p = \mathcal{F}^\#_p$ ; autrement dit, le foncteur faisceau associé
commute à la prise de la fibre en un point.
\end{remark}




\section{Conventions et notation}
\label{formal-spaces-section-conventions}

\noindent
Les conventions adoptées désormais seront semblables à celles de
Propriétés des espaces algébriques, section \ref{spaces-properties-section-conventions}.
Ainsi, nous supposerons toujours que tous les schémas appartiennent
à un grand site fppf $\Sch_{fppf}$. De plus, tout anneau $A$ considéré aura la
propriété que $\Spec(A)$ est (isomorphe à) un objet de ce grand site.
Pour les anneaux topologiques $A$, nous supposons seulement que tous leurs quotients discrets ont
cette propriété (mais nous faisons habituellement des hypothèses plus fortes ; comparer à la
remarque \ref{formal-spaces-remark-set-theoretic}).

\medskip\noindent
Soient $S$ un schéma et $X$ un « espace » sur $S$, c'est-à-dire un faisceau sur
$(\Sch/S)_{fppf}$. Dans ce chapitre, nous noterons $X \times_S X$ le
produit de $X$ par lui-même dans la catégorie des faisceaux sur $(\Sch/S)_{fppf}$,
et non $X \times X$. De plus, si $X$ et $Y$ sont des « espaces »,
nous dirons « soit $f : X \to Y$ un morphisme » pour indiquer que $f$ est une
transformation naturelle de foncteurs, c'est-à-dire une application de faisceaux sur
$(\Sch/S)_{fppf}$. De même, si $U$ est un schéma sur $S$ et si
$X$ est un « espace » sur $S$, nous dirons
« soit $f : U \to X$ un morphisme » ou
« soit $g : X \to U$ un morphisme » pour indiquer que $f$ ou $g$
est une application de faisceaux $h_U \to X$ ou $X \to h_U$, où $h_U$ est défini dans
Catégories, exemple \ref{categories-example-hom-functor}.






\section{Anneaux et modules topologiques}
\label{formal-spaces-section-topological-module}

\noindent
Cette section prolonge
Compléments d'algèbre, section \ref{more-algebra-section-topological-ring}.
Soient $R$ un anneau topologique et $M$ un $R$-module linéairement topologisé.
Lorsque nous dirons « {\it soit $M_\lambda$ un système fondamental de
sous-modules ouverts} », nous entendrons que tout $M_\lambda$ est un sous-module ouvert
et que tout voisinage de $0$ contient l'un des $M_\lambda$.
Autrement dit, cela signifie que les $M_\lambda$ forment un système fondamental
de voisinages de $0$ dans $M$, constitué de sous-modules.
De même, si $R$ est un anneau linéairement topologisé, nous dirons
« {\it soit $I_\lambda$ un système fondamental d'idéaux ouverts} »
pour signifier que les $I_\lambda$ forment un système fondamental
de voisinages de $0$ dans $R$, constitué d'idéaux.

\begin{example}
\label{formal-spaces-example-what-does-it-mean}
Soient $R$ un anneau linéairement topologisé et $M$ un
$R$-module linéairement topologisé. Soient $I_\lambda$ un système fondamental
d'idéaux ouverts de $R$ et $M_\mu$ un système fondamental de
sous-modules ouverts de $M$. La continuité de $+ : M \times M \to M$
est automatique, et la continuité de $R \times M \to M$ signifie que
$$
\forall f, x, \mu\ \exists \lambda, \nu,\ (f + I_\lambda)(x + M_\nu)
\subset fx + M_\mu
$$
Puisque $fM_\nu + I_\lambda M_\nu \subset M_\mu$ dès que
$M_\nu \subset M_\mu$, on voit que la condition équivaut à
$$
\forall x, \mu\ \exists \lambda\ I_\lambda x \subset M_\mu
$$
Cependant, il ne faut pas nécessairement qu'étant donné $\mu$, il existe $\lambda$
tel que $I_\lambda M \subset M_\mu$. Considérons par exemple
$R = k[[t]]$ muni de la topologie $t$-adique et
$M = \bigoplus_{n \in \mathbf{N}} R$, avec le système fondamental de
sous-modules ouverts donné par
$$
M_m = \bigoplus\nolimits_{n \in \mathbf{N}} t^{nm}R
$$
Comme tout $x \in M$ n'a qu'un nombre fini de coordonnées non nulles, on voit
qu'étant donnés $m$ et $x$, il existe $k$ tel que $t^k x \in M_m$.
Ainsi, $M$ est un $R$-module linéairement topologisé, mais il n'est pas vrai
qu'étant donné $m$, il existe $k$ tel que $t^kM \subset M_m$.
En revanche, si $R \to S$ est un homomorphisme continu d'anneaux
linéairement topologisés, l'énoncé correspondant est vrai :
pour tout idéal ouvert $J \subset S$, il existe un idéal ouvert
$I \subset R$ tel que $IS \subset J$ (comme le lecteur le déduira aisément
de la continuité de l'homomorphisme $R \to S$).
\end{example}

\begin{lemma}
\label{formal-spaces-lemma-closed}
Soit $R$ un anneau topologique. Soient $M$ un $R$-module linéairement
topologisé et $M_\lambda$, $\lambda \in \Lambda$, un système fondamental
de sous-modules ouverts. Soit $N \subset M$ un sous-module.
L'adhérence de $N$ est $\bigcap_{\lambda \in \Lambda} (N + M_\lambda)$.
\end{lemma}

\begin{proof}
Comme chaque $N + M_\lambda$ est ouvert, il est également fermé. Leur
intersection est donc fermée. Si $x \in M$ n'appartient pas à l'adhérence de $N$,
alors $(x + M_\lambda) \cap N = 0$ pour un certain $\lambda$. Ainsi,
$x \not \in N + M_\lambda$. Cela démontre le lemme.
\end{proof}

\noindent
Sauf mention contraire, nous munissons les sous-modules et les modules quotients
de la topologie induite. Soient $M$ un module linéairement topologisé
sur un anneau topologique $R$, et $0 \to N \to M \to Q \to 0$
une suite exacte courte de $R$-modules. Si les $M_\lambda$ forment un
système fondamental de sous-modules ouverts de $M$, alors
les $N \cap M_\lambda$ forment un système fondamental de sous-modules ouverts de $N$.
Si $\pi : M \to Q$ est l'application quotient, les $\pi(M_\lambda)$ forment un
système fondamental de sous-modules ouverts de $Q$. En particulier, ces topologies
induites sont des topologies linéaires.

\begin{lemma}
\label{formal-spaces-lemma-closure}
Soit $R$ un anneau topologique. Soit $M$ un $R$-module linéairement
topologisé. Soit $N \subset M$ un sous-module. Alors :
\begin{enumerate}
\item $0 \to N^\wedge \to M^\wedge \to (M/N)^\wedge$ est exacte, et
\item $N^\wedge$ est l'adhérence de l'image de $N \to M^\wedge$.
\end{enumerate}
\end{lemma}

\begin{proof}
Soient $M_\lambda$, $\lambda \in \Lambda$, un système fondamental de
sous-modules ouverts. Alors les $N \cap M_\lambda$ forment un système fondamental
de sous-modules ouverts de $N$, et les $M_\lambda + N/N$ un système fondamental
de sous-modules ouverts de $M/N$. On voit donc que (1) découle de
l'exactitude des suites
$$
0 \to N/N \cap M_\lambda \to M/M_\lambda \to M/(M_\lambda + N) \to 0
$$
et du fait que la prise de limites commute aux limites. La seconde
assertion en résulte, ainsi que du fait que $N \to N^\wedge$
est d'image dense et que le noyau de $M^\wedge \to (M/N)^\wedge$ est fermé.
\end{proof}

\begin{lemma}
\label{formal-spaces-lemma-quotient-by-closed}
Soit $R$ un anneau topologique. Soit $M$ un $R$-module complet et linéairement
topologisé. Soit $N \subset M$ un sous-module fermé. Si $M$ possède un
système fondamental dénombrable de voisinages de $0$, alors
$M/N$ est complet et l'application $M \to M/N$ est ouverte.
\end{lemma}

\begin{proof}
Soient $M_n$, $n \in \mathbf{N}$, un système fondamental de sous-modules ouverts de $M$.
Nous pouvons supposer $M_{n + 1} \subset M_n$
pour tout $n$. Les $(M_n + N)/N$ forment un système fondamental dans $M/N$.
Il faut donc montrer que $M/N = \lim M/(M_n + N)$. Considérons
les suites exactes courtes
$$
0 \to N/N \cap M_n \to M/M_n \to M/(M_n + N) \to 0
$$
Comme les applications de transition du système $\{N/N\cap M_n\}$ sont surjectives,
on voit que $M = \lim M/M_n$ (puisque $M$ est complet) se surjecte sur
$\lim M/(M_n + N)$ d'après
Algèbre, lemme \ref{algebra-lemma-ML-exact-sequence}.
Puisque $N$ est fermé, le noyau de $M \to \lim M/(M_n + N)$
est $N$ (voir le lemme \ref{formal-spaces-lemma-closed}). Enfin, $M \to M/N$
est ouverte par définition de la topologie quotient.
\end{proof}

\begin{lemma}
\label{formal-spaces-lemma-ses}
\begin{reference}
\cite[Theorem 8.1]{Ma}
\end{reference}
Soit $R$ un anneau topologique. Soit $M$ un $R$-module linéairement
topologisé. Soit $N \subset M$ un sous-module. Supposons que $M$ possède un
système fondamental dénombrable de voisinages de $0$. Alors :
\begin{enumerate}
\item $0 \to N^\wedge \to M^\wedge \to (M/N)^\wedge \to 0$ est exacte,
\item $N^\wedge$ est l'adhérence de l'image de $N \to M^\wedge$,
\item $M^\wedge \to (M/N)^\wedge$ est ouverte.
\end{enumerate}
\end{lemma}

\begin{proof}
La suite $0 \to N^\wedge \to M^\wedge \to (M/N)^\wedge$ est exacte,
et l'assertion (2) résulte du lemme \ref{formal-spaces-lemma-closure}.
On obtient ainsi une application canonique $c : M^\wedge/N^\wedge \to (M/N)^\wedge$.
Le module $M^\wedge/N^\wedge$ est complet, et
$M^\wedge \to M^\wedge/N^\wedge$ est ouverte d'après
le lemme \ref{formal-spaces-lemma-quotient-by-closed}.
La propriété universelle du complété donne une application canonique
$b : (M/N)^\wedge \to M^\wedge/N^\wedge$.
Alors $b$ et $c$ sont inverses l'une de l'autre, puisqu'ils le sont sur un sous-ensemble dense.
\end{proof}

\begin{lemma}
\label{formal-spaces-lemma-completion-adic-star}
Soit $R$ un anneau topologique. Soit $M$ un $R$-module topologique.
Soit $I \subset R$ un idéal de type fini. Supposons que $M$
possède un sous-module ouvert dont la topologie est $I$-adique. Alors $M^\wedge$
possède un sous-module ouvert dont la topologie est $I$-adique, et l'on a
$M^\wedge/I^n M^\wedge = M/I^nM$ pour tout $n \geq 1$.
\end{lemma}

\begin{proof}
Soit $M' \subset M$ un sous-module ouvert dont la topologie est $I$-adique.
Alors $\{I^nM'\}_{n \geq 1}$ est un système fondamental de sous-modules ouverts
de $M$. Ainsi, $M^\wedge = \lim M/I^nM'$ contient
$(M')^\wedge = \lim M'/I^nM'$ comme sous-module ouvert,
et la topologie sur $(M')^\wedge$ est
$I$-adique d'après Algèbre, lemme \ref{algebra-lemma-hathat-finitely-generated}.
Comme $I$ est de type fini, $I^n$ est de type fini ;
soit $f_1, \ldots, f_r$ un système de générateurs. Observons que la surjection
$(f_1, \ldots, f_r) : M^{\oplus r} \to I^n M$ est continue
et ouverte d'après la description précédente de la topologie sur $M$.
En appliquant le lemme \ref{formal-spaces-lemma-ses} à cette surjection et à la
suite exacte courte $0 \to I^nM \to M \to M/I^nM \to 0$,
on conclut que
$$
(f_1, \ldots, f_r) :
(M^\wedge)^{\oplus r} \longrightarrow M^\wedge
$$
se surjecte sur le noyau de la surjection $M^\wedge \to M/I^nM$.
Comme $f_1, \ldots, f_r$ engendrent $I^n$, cela conclut la démonstration.
\end{proof}

\begin{definition}
\label{formal-spaces-definition-toplogy-tensor-product}
\begin{reference}
Comparer à \cite[0, Section 7.7]{EGA}
\end{reference}
Soit $R$ un anneau topologique. Soient $M$ et $N$ des $R$-modules
linéairement topologisés. Le {\it produit tensoriel} de $M$ et $N$
est le produit tensoriel (usuel) $M \otimes_R N$, muni
de la topologie linéaire définie en déclarant que les
$$
\Im(M_\mu \otimes_R N + M \otimes_R N_\nu \longrightarrow M \otimes_R N)
$$
forment un système fondamental de sous-modules ouverts, lorsque
$M_\mu \subset M$ et $N_\nu \subset N$ parcourent des systèmes fondamentaux
de sous-modules ouverts de $M$ et de $N$.
Le {\it produit tensoriel complété}
$$
M \widehat{\otimes}_R N =
\lim M \otimes_R N/(M_\mu \otimes_R N + M \otimes_R N_\nu) =
\lim M/M_\mu \otimes_R N/N_\nu
$$
est le complété du produit tensoriel.
\end{definition}

\noindent
Observons que la topologie de $R$ n'intervient pas dans la construction
du produit tensoriel ni du produit tensoriel complété.
Si $R \to A$ et $R \to B$ sont des homomorphismes continus d'anneaux
linéairement topologisés, la construction précédente
fournit un produit tensoriel $A \otimes_R B$ et un produit
tensoriel complété $A \widehat{\otimes}_R B$.

\medskip\noindent
Consignons ici les notions introduites dans la remarque \ref{formal-spaces-remark-mcquillan}.

\begin{definition}
\label{formal-spaces-definition-weakly-admissible}
Soit $A$ un anneau linéairement topologisé.
\begin{enumerate}
\item Un élément $f \in A$ est dit {\it topologiquement nilpotent}
si $f^n \to 0$ lorsque $n \to \infty$.
\item Un {\it idéal faible de définition} de $A$ est un idéal ouvert
$I \subset A$ entièrement constitué d'éléments topologiquement nilpotents.
\item On dit que $A$ est {\it faiblement préadmissible} si $A$ possède un idéal faible
de définition.
\item On dit que $A$ est {\it faiblement admissible} si $A$ est faiblement préadmissible
et complet\footnote{Selon nos conventions, cela inclut la séparation.}.
\end{enumerate}
\end{definition}

\noindent
Étant donnés un idéal faible de définition $I$ dans un anneau linéairement topologisé
$A$ et un idéal ouvert $J$, l'intersection $I \cap J$ est un
idéal faible de définition. Par conséquent, s'il existe un idéal faible de définition,
il existe un système fondamental d'idéaux ouverts
constitué d'idéaux faibles de définition. En particulier,
si $A$ est un anneau topologique faiblement admissible, alors
$A = \lim A/I_\lambda$, où $\{I_\lambda\}$ est un système fondamental
d'idéaux faibles de définition.

\begin{lemma}
\label{formal-spaces-lemma-weakly-admissible-henselian}
Soit $A$ un anneau topologique faiblement admissible. Soit $I \subset A$
un idéal faible de définition. Alors $(A, I)$ est un couple hensélien.
\end{lemma}

\begin{proof}
Soit $A \to A'$ un homomorphisme d'anneaux étale et soit $\sigma : A' \to A/I$
un homomorphisme de $A$-algèbres. D'après Compléments d'algèbre, lemme
\ref{more-algebra-lemma-characterize-henselian-pair}, il suffit
de relever $\sigma$ en un homomorphisme de $A$-algèbres $A' \to A$.
Pour cela, comme $A$ est complet, il suffit de trouver,
pour tout idéal ouvert $J \subset I$, un unique homomorphisme de $A$-algèbres $A' \to A/J$
relevant $\sigma$. Comme $I$ est un idéal faible de définition,
l'idéal $I/J$ est localement nilpotent. On conclut par
Compléments d'algèbre, lemme \ref{more-algebra-lemma-locally-nilpotent-henselian}.
\end{proof}

\begin{lemma}
\label{formal-spaces-lemma-topologically-nilpotent}
Soit $B$ un anneau linéairement topologisé. L'ensemble des éléments topologiquement nilpotents
de $B$ est un idéal radical fermé de $B$.
Soit $\varphi : A \to B$ un homomorphisme continu d'anneaux linéairement topologisés.
\begin{enumerate}
\item Si $f \in A$ est topologiquement nilpotent, alors $\varphi(f)$ est
topologiquement nilpotent.
\item Si $I \subset A$ est constitué d'éléments topologiquement nilpotents,
alors l'adhérence de $\varphi(I)B$ est constituée d'éléments topologiquement
nilpotents.
\end{enumerate}
\end{lemma}

\begin{proof}
Soit $\mathfrak b \subset B$ l'ensemble des éléments topologiquement nilpotents.
Nous omettons la démonstration du fait que $\mathfrak b$ est un idéal radical
(c'est un bon exercice sur les définitions). Soit $g$ un élément de l'adhérence
de $\mathfrak b$. Nous voulons montrer que $g$ est topologiquement nilpotent.
Soit $J \subset B$ un idéal ouvert. Il faut montrer que
$g^e \in J$ pour un certain $e \geq 1$. On a $g \in \mathfrak b + J$
d'après le lemme \ref{formal-spaces-lemma-closed}. Ainsi, $g = f + h$
pour certains $f \in \mathfrak b$ et $h \in J$. Choisissons $m \geq 1$ tel que
$f^m \in J$. Alors $g^{m + 1} \in J$, comme voulu.

\medskip\noindent
Soit $\varphi : A \to B$ comme dans l'énoncé du lemme.
L'assertion (1) est claire, et l'assertion (2) résulte de celle-ci et
du fait que $\mathfrak b$ est un idéal fermé.
\end{proof}

\begin{lemma}
\label{formal-spaces-lemma-closure-image-ideal}
Soit $A \to B$ un homomorphisme continu d'anneaux linéairement topologisés.
Soit $I \subset A$ un idéal. L'adhérence de $IB$
est le noyau de $B \to B \widehat{\otimes}_A A/I$.
\end{lemma}

\begin{proof}
Soient les $J_\mu$ un système fondamental d'idéaux ouverts de $B$.
L'adhérence de $IB$ est $\bigcap (IB + J_\lambda)$ d'après le lemme \ref{formal-spaces-lemma-closed}.
Soient les $I_\mu$ un système fondamental d'idéaux ouverts de $A$.
Alors
$$
B \widehat{\otimes}_A A/I = \lim (B/J_\lambda \otimes_A A/(I_\mu + I)) =
\lim B/(J_\lambda + I_\mu B + I B)
$$
Comme $A \to B$ est continu, pour tout $\lambda$, il
existe $\mu$ tel que $I_\mu B \subset J_\lambda$ ; voir la discussion
de l'exemple \ref{formal-spaces-example-what-does-it-mean}. La limite peut donc
s'écrire $\lim B/(J_\lambda + IB)$, et le résultat est clair.
\end{proof}

\begin{lemma}
\label{formal-spaces-lemma-completed-tensor-product}
Soient $B \to A$ et $B \to C$ des homomorphismes continus d'anneaux
linéairement topologisés.
\begin{enumerate}
\item Si $A$ et $C$ sont faiblement préadmissibles, alors
$A \widehat{\otimes}_B C$ est faiblement admissible.
\item Si $A$ et $C$ sont préadmissibles, alors
$A \widehat{\otimes}_B C$ est admissible.
\item Si $A$ et $C$ possèdent un système fondamental dénombrable d'idéaux
ouverts, alors $A \widehat{\otimes}_B C$ possède un système fondamental dénombrable
d'idéaux ouverts.
\item Si $A$ et $C$ sont préadiques et possèdent des idéaux de définition
de type fini, alors $A \widehat{\otimes}_B C$ est adique et possède
un idéal de définition de type fini.
\item Si $A$ et $C$ sont des anneaux noethériens préadiques et si
$B/\mathfrak b \to A/\mathfrak a$ est de type fini,
où $\mathfrak a \subset A$ et $\mathfrak b \subset B$
sont les idéaux des éléments topologiquement nilpotents, alors
$A \widehat{\otimes}_B C$ est noethérien adique.
\end{enumerate}
\end{lemma}

\begin{proof}
Soient $I_\lambda \subset A$, $\lambda \in \Lambda$, et
$J_\mu \subset C$, $\mu \in M$,
des systèmes fondamentaux d'idéaux ouverts ; alors, par définition,
$$
A \widehat{\otimes}_B C =
\lim_{\lambda, \mu} A/I_\lambda \otimes_B C/J_\mu
$$
muni de la topologie limite. Un système fondamental d'idéaux ouverts
est donc donné par les noyaux $K_{\lambda, \mu}$ des applications
$A \widehat{\otimes}_B C \to A/I_\lambda \otimes_B C/J_\mu$.
Remarquons que $K_{\lambda, \mu}$ est l'adhérence de l'idéal
$I_\lambda(A \widehat{\otimes}_B C) + J_\mu(A \widehat{\otimes}_B C)$.
Enfin, nous avons un homomorphisme d'anneaux
$\tau : A \otimes_B C \to A \widehat{\otimes}_B C$ d'image dense.

\medskip\noindent
Démonstration de (1). Si $I_\lambda$ et $J_\mu$ sont constitués d'éléments topologiquement
nilpotents, il en va de même de $K_{\lambda, \mu}$ d'après
le lemme \ref{formal-spaces-lemma-topologically-nilpotent}. Ainsi, 
$A \widehat{\otimes}_B C$ est faiblement admissible par définition.

\medskip\noindent
Démonstration de (2). Supposons que, pour certains $\lambda_0$ et $\mu_0$, les idéaux
$I = I_{\lambda_0} \subset A$ et $J_{\mu_0} \subset C$ soient des idéaux de
définition. Alors, pour tout $\lambda$, il existe $n$ tel que
$I^n \subset I_\lambda$. Pour tout $\mu$, il existe $m$ tel que
$J^m \subset J_\mu$. Par conséquent,
$$
\left(I(A \widehat{\otimes}_B C) + J(A \widehat{\otimes}_B C)\right)^{n + m}
\subset
I_\lambda(A \widehat{\otimes}_B C) + J_\mu(A \widehat{\otimes}_B C)
$$
Il s'ensuit que l'idéal ouvert $K = K_{\lambda_0, \mu_0}$
vérifie $K^{n + m} \subset K_{\lambda, \mu}$. Ainsi, $K$
est un idéal de définition de $A \widehat{\otimes}_B C$,
et $A \widehat{\otimes}_B C$ est admissible par définition.

\medskip\noindent
Démonstration de (3). Si $\Lambda$ et $M$ sont dénombrables, il en va de même de
$\Lambda \times M$.

\medskip\noindent
Démonstration de (4). Supposons $\Lambda = \mathbf{N}$ et $M = \mathbf{N}$,
et soient $I \subset A$ et $J \subset C$ des idéaux de type fini
tels que $I_n = I^n$ et $J_n = J^n$. Alors
$$
I(A \widehat{\otimes}_B C) + J(A \widehat{\otimes}_B C)
$$
est un idéal de type fini, et l'on voit facilement que
$A \widehat{\otimes}_B C$ est le complété de
$A \otimes_B C$ pour cet idéal. L'assertion (4)
résulte donc de Algèbre, lemme \ref{algebra-lemma-hathat-finitely-generated}.

\medskip\noindent
Démonstration de (5).
Soit $\mathfrak c \subset C$ l'idéal des éléments topologiquement nilpotents.
Comme $A$ et $C$ sont noethériens adiques, on voit que
$\mathfrak a$ et $\mathfrak c$ sont des idéaux de définition (détails omis).
D'après (4), nous savons déjà que
$A \widehat{\otimes}_B C$ est adique et que
$\mathfrak a(A \widehat{\otimes}_B C) + \mathfrak c(A \widehat{\otimes}_B C)$
est un idéal de définition de type fini. Puisque
$$
A \widehat{\otimes}_B C /
\left(\mathfrak a(A \widehat{\otimes}_B C) +
\mathfrak c(A \widehat{\otimes}_B C)\right)
=
A/\mathfrak a \otimes_{B/\mathfrak b} C/\mathfrak c
$$
est noethérien, car c'est une algèbre de type fini sur l'anneau noethérien
$C/\mathfrak c$, on conclut par
Algèbre, lemme \ref{algebra-lemma-completion-Noetherian}.
\end{proof}









\section{Homomorphismes d'anneaux tendus}
\label{formal-spaces-section-taut-ring-maps}

\noindent
Il est commode de donner un nom à la propriété suivante
des homomorphismes continus entre anneaux linéairement topologisés.

\begin{definition}
\label{formal-spaces-definition-taut}
Soit $\varphi : A \to B$ un homomorphisme continu d'anneaux linéairement topologisés.
On dit que $\varphi$ est {\it tendu}\footnote{Cette terminologie n'est pas standard.
La définition s'étend aux modules : on dit qu'un
$A$-module linéairement topologisé $M$ est $A$-tendu si, pour tout idéal ouvert $I \subset A$, l'adhérence
de $IM$ dans $M$ est ouverte et si ces adhérences forment un système fondamental de
voisinages de $0$ dans $M$.}
si, pour tout idéal ouvert $I \subset A$, l'adhérence de l'idéal $\varphi(I)B$
est ouverte et si ces adhérences forment un système fondamental d'idéaux ouverts.
\end{definition}

\noindent
Si $\varphi : A \to B$ est un homomorphisme continu d'anneaux linéairement topologisés
et si les $I_\lambda$ forment un système fondamental d'idéaux ouverts de $A$, alors $\varphi$
est tendu si et seulement si
les adhérences des $I_\lambda B$ sont ouvertes et forment un système fondamental
d'idéaux ouverts de $A$.

\begin{lemma}
\label{formal-spaces-lemma-taut-weakly-admissible}
Soit $\varphi : A \to B$ un homomorphisme continu d'anneaux topologiques
faiblement admissibles. Les conditions suivantes sont équivalentes :
\begin{enumerate}
\item $\varphi$ est tendu ;
\item pour tout idéal faible de définition $I \subset A$, l'adhérence de
$\varphi(I)B$ est un idéal faible de définition de $B$, et ces idéaux forment un
système fondamental d'idéaux faibles de définition de $B$.
\end{enumerate}
\end{lemma}

\begin{proof}
Les remarques qui suivent la définition \ref{formal-spaces-definition-taut} montrent que
(2) implique (1). Réciproquement, supposons $\varphi$ tendu. Si $I \subset A$
est un idéal faible de définition, alors l'adhérence de $\varphi(I)B$
est ouverte par définition du caractère tendu et est constituée d'éléments topologiquement
nilpotents d'après le lemme \ref{formal-spaces-lemma-topologically-nilpotent}.
L'adhérence de $\varphi(I)B$ est donc un idéal faible de définition.
En outre, par définition du caractère tendu, ces idéaux forment un
système fondamental d'idéaux ouverts ; la condition (2) est donc vérifiée.
\end{proof}

\begin{lemma}
\label{formal-spaces-lemma-completion-taut}
Soit $A$ un anneau linéairement topologisé. L'homomorphisme $A \to A^\wedge$
de $A$ vers son complété est tendu.
\end{lemma}

\begin{proof}
Soient les $I_\lambda$ un système fondamental d'idéaux ouverts de $A$.
Rappelons que $A^\wedge = \lim A/I_\lambda$, muni de la topologie limite,
ce qui signifie que les noyaux $J_\lambda = \Ker(A^\wedge \to A/I_\lambda)$
forment un système fondamental d'idéaux ouverts de $A^\wedge$.
Puisque $J_\lambda$ est l'adhérence de $I_\lambda A^\wedge$
(comparer au lemme \ref{formal-spaces-lemma-closure-image-ideal}), on conclut.
\end{proof}

\begin{lemma}
\label{formal-spaces-lemma-composition-taut}
Soient $A \to B$ et $B \to C$ des homomorphismes continus d'anneaux
linéairement topologisés. Si $A \to B$ et $B \to C$ sont tendus, alors
$A \to C$ est tendu.
\end{lemma}

\begin{proof}
Omis. Indication : si $I \subset A$ est un idéal et si $J$ est l'adhérence
de $IB$, alors l'adhérence de $JC$ est égale à celle de $IC$.
\end{proof}

\begin{lemma}
\label{formal-spaces-lemma-permanence-taut}
Soient $A \to B$ et $B \to C$ des homomorphismes continus d'anneaux
linéairement topologisés. Si $A \to C$ est tendu, alors
$B \to C$ est tendu.
\end{lemma}

\begin{proof}
Soit $J \subset B$ un idéal ouvert d'image réciproque $I \subset A$.
Alors l'adhérence de $JC$ contient celle de $IC$. Cette
adhérence est donc ouverte, puisque $A \to C$ est tendu. Soient les $I_\lambda$ un système fondamental
d'idéaux ouverts de $A$. Soit $K_\lambda$ l'adhérence de
$I_\lambda C$. Puisque $A \to C$ est tendu, ils forment un système fondamental
d'idéaux ouverts de $C$. Notons $J_\lambda \subset B$ l'image réciproque
de $K_\lambda$. Alors l'adhérence de $J_\lambda C$ est $K_\lambda$.
On voit donc que les adhérences des idéaux $JC$, lorsque $J$ parcourt
les idéaux ouverts de $B$, forment un système fondamental d'idéaux ouverts de $C$.
\end{proof}

\begin{lemma}
\label{formal-spaces-lemma-base-change-taut}
Soient $A \to B$ et $A \to C$ des homomorphismes continus d'anneaux
linéairement topologisés. Si $A \to B$ est tendu, alors
$C \to B \widehat{\otimes}_A C$ est tendu.
\end{lemma}

\begin{proof}
Soit $K \subset C$ un idéal ouvert. Choisissons un idéal ouvert $I \subset A$
dont l'image dans $C$ est contenue dans $J$. Par hypothèse, l'adhérence $J$
de $IB$ est ouverte. Comme $A \to B$ est tendu, on voit que
$B \widehat{\otimes}_A C$ est la limite des anneaux $B/J \otimes_{A/I} C/K$
pour tous les choix de $K$ et de $I$, c'est-à-dire que
les idéaux $J(B \widehat{\otimes}_A C) + K(B \widehat{\otimes}_A C)$
forment un système fondamental d'idéaux ouverts. Or,
comme $B \to B \widehat{\otimes}_A C$ est continu, on voit
que $J$ s'envoie dans l'adhérence de $K(B \widehat{\otimes}_A C)$
(puisque $I$ s'envoie dans $K$). Cette adhérence est donc égale à
$J(B \widehat{\otimes}_A C) + K(B \widehat{\otimes}_A C)$,
ce qui achève la démonstration.
\end{proof}

\begin{lemma}
\label{formal-spaces-lemma-taut-ascent-countable}
Soit $\varphi : A \to B$ un homomorphisme continu d'anneaux
linéairement topologisés. Si $\varphi$ est tendu et si $A$
possède un système fondamental dénombrable d'idéaux ouverts, alors
$B$ possède un système fondamental dénombrable d'idéaux ouverts.
\end{lemma}

\begin{proof}
Cela résulte immédiatement des définitions.
\end{proof}

\begin{lemma}
\label{formal-spaces-lemma-taut-ascent-weakly-admissible}
Soit $\varphi : A \to B$ un homomorphisme continu d'anneaux
linéairement topologisés. Si $\varphi$ est tendu et si $A$
est faiblement préadmissible, alors $B$ est faiblement préadmissible.
\end{lemma}

\begin{proof}
Soit $I \subset A$ un idéal faible de définition. Alors l'adhérence
$J$ de $IB$ est ouverte et constituée d'éléments topologiquement nilpotents
d'après le lemme \ref{formal-spaces-lemma-topologically-nilpotent}. Par conséquent, $J$ est un idéal faible
de définition de $B$.
\end{proof}

\begin{lemma}
\label{formal-spaces-lemma-taut-ascent-admissible}
Soit $\varphi : A \to B$ un homomorphisme continu d'anneaux
linéairement topologisés. Si $\varphi$ est tendu et si $A$
est préadmissible, alors $B$ est préadmissible.
\end{lemma}

\begin{proof}
Soit $I \subset A$ un idéal de définition. Soient
les $I_\lambda \subset A$ un système fondamental d'idéaux ouverts.
Alors l'adhérence $J$ de $IB$ est ouverte, et les
adhérences $J_\lambda$ des $I_\lambda B$ sont ouvertes et forment un
système fondamental d'idéaux ouverts de $B$.
Pour tout $\lambda$, il existe $n$ tel que $I^n \subset I_\lambda$.
Observons que $J^n$ est contenu dans l'adhérence de $I^nB$.
Ainsi, $J^n \subset J_\lambda$, et l'on conclut que $J$ est un
idéal de définition.
\end{proof}

\begin{lemma}
\label{formal-spaces-lemma-dense-image-surjective}
Soit $\varphi : A \to B$ un homomorphisme continu d'anneaux
linéairement topologisés. Supposons que :
\begin{enumerate}
\item $\varphi$ soit tendu et d'image dense ;
\item $A$ soit complet et possède un système fondamental dénombrable
d'idéaux ouverts ; et
\item $B$ soit séparé.
\end{enumerate}
Alors $\varphi$ est surjectif et ouvert, $B$ est complet et $B = A/K$ pour
un certain idéal fermé $K \subset A$.
\end{lemma}

\begin{proof}
En combinant le lemme de l'application ouverte (Compléments d'algèbre, lemme
\ref{more-algebra-lemma-open-mapping})
et le caractère tendu de $\varphi$, 
on voit que l'application $\varphi$ est ouverte. Comme l'image de $\varphi$
est dense, $\varphi$ est surjectif. Le noyau $K$
de $\varphi$ est fermé, puisque $\varphi$ est continu.
Il s'ensuit que $B = A/K$ est complet ; voir par exemple
le lemme \ref{formal-spaces-lemma-quotient-by-closed}.
\end{proof}







\section{Homomorphismes d'anneaux adiques}
\label{formal-spaces-section-adic-ring-maps}

\noindent
Posons la définition suivante.

\begin{definition}
\label{formal-spaces-definition-adic-homomorphism}
Soient $A$ et $B$ des anneaux topologiques préadiques. Un homomorphisme d'anneaux
$\varphi : A \to B$ est dit {\it adique}\footnote{Cette terminologie n'est peut-être pas standard.}
s'il existe un idéal de définition $I \subset A$ tel que
la topologie de $B$ soit la topologie $I$-adique.
\end{definition}

\noindent
Si $\varphi : A \to B$ est un homomorphisme adique d'anneaux préadiques, alors
$\varphi$ est continu et la topologie de $B$ est la topologie $I$-adique
pour tout idéal de définition $I$ de $A$.

\begin{lemma}
\label{formal-spaces-lemma-composition-adic}
Soient $A \to B$ et $B \to C$ des homomorphismes continus d'anneaux
préadiques. Si $A \to B$ et $B \to C$ sont adiques, alors
$A \to C$ est adique.
\end{lemma}

\begin{proof}
Omis.
\end{proof}

\begin{lemma}
\label{formal-spaces-lemma-permanence-adic}
Soient $A \to B$ et $B \to C$ des homomorphismes continus d'anneaux
préadiques. Si $A \to C$ est adique, alors
$B \to C$ est adique.
\end{lemma}

\begin{proof}
Choisissons un idéal de définition $I$ de $A$. Comme $A \to C$ est adique,
on voit que $IC$ est un idéal de définition de $C$.
Comme $B \to C$ est continu, on peut trouver un
idéal de définition $J \subset B$ dont l'image est contenue dans $IC$.
Comme $A \to B$ est continu, l'image réciproque $I' \subset I$
de $J$ dans $I$ est aussi un idéal de définition de $A$.
Ainsi, $I'C \subset JC \subset IC$ est compris entre
deux idéaux de définition et est donc lui-même un idéal de définition.
\end{proof}

\begin{lemma}
\label{formal-spaces-lemma-adic-taut}
Soit $\varphi : A \to B$ un homomorphisme continu entre anneaux
topologiques préadiques. Si $\varphi$ est adique, alors $\varphi$ est tendu.
\end{lemma}

\begin{proof}
Cela résulte immédiatement des définitions.
\end{proof}

\noindent
Le lemme suivant affirme deux choses :
\begin{enumerate}
\item la propriété d'être adique monte le long des homomorphismes tendus
d'anneaux complets linéairement topologisés ; et
\item les propriétés « $\varphi$ est tendu » et « $\varphi$ est adique »
sont équivalentes pour les homomorphismes continus $\varphi : A \to B$ entre anneaux adiques
si $A$ possède un idéal de définition de type fini.
\end{enumerate}
En raison de (2), on peut dire que le « caractère tendu » généralise le « caractère adique »
aux homomorphismes continus entre anneaux linéairement topologisés quelconques.
Voir aussi la section \ref{formal-spaces-section-adic}.

\begin{lemma}
\label{formal-spaces-lemma-taut-is-adic}
Soit $\varphi : A \to B$ un homomorphisme continu d'anneaux linéairement topologisés.
Si $\varphi$ est tendu, si $A$ est préadique et possède un idéal de définition
de type fini, et si $B$ est complet, alors $B$ est adique et possède un idéal de définition
de type fini, et l'homomorphisme d'anneaux $\varphi$ est adique.
\end{lemma}

\begin{proof}
Choisissons un idéal de définition de type fini $I$ de $A$.
Soit $J_n$ l'adhérence de $\varphi(I^n)B$ dans $B$.
Comme $B$ est complet, on a $B = \lim B/J_n$.
Soit $B' = \lim B/I^nB$ le complété $I$-adique de $B$.
D'après Algèbre, lemme \ref{algebra-lemma-hathat-finitely-generated},
la topologie $I$-adique sur $B'$ est complète, et
$B'/I^nB' = B/I^nB$. Ainsi, l'homomorphisme d'anneaux $B' \to B$ est continu
et d'image dense, puisque $B' \to B/I^nB \to B/J_n$ est surjectif
pour tout $n$. Enfin, l'homomorphisme $B' \to B$ est tendu,
car $(I^nB')B = I^nB$ et $A \to B$ est tendu.
D'après le lemme \ref{formal-spaces-lemma-dense-image-surjective}, $B' \to B$ est ouvert
et surjectif. La topologie de $B$ est donc la topologie $I$-adique,
ce qui achève la démonstration.
\end{proof}








\section{Anneaux faiblement adiques}
\label{formal-spaces-section-weakly-adic}

\noindent
Nous conseillons au lecteur de sauter cette section. La notion suivante est une généralisation
naturelle des anneaux adiques.

\begin{definition}
\label{formal-spaces-definition-weakly-adic}
\begin{reference}
\cite[Definition 8.3.8]{Gabber-Ramero}
\end{reference}
Soit $A$ un anneau linéairement topologisé.
\begin{enumerate}
\item On dit que $A$ est {\it faiblement préadique}\footnote{Dans \cite{Gabber-Ramero}, les
auteurs disent que $A$ est {\it $c$-adique}.} s'il existe un idéal
$I \subset A$ tel que l'adhérence de $I^n$ soit ouverte pour tout $n \geq 0$
et que ces adhérences forment un système fondamental d'idéaux ouverts.
\item On dit que $A$ est {\it faiblement adique} si $A$ est faiblement préadique
et complet\footnote{Selon nos conventions, cela inclut la séparation.}.
\end{enumerate}
\end{definition}

\noindent
Pour les anneaux complets linéairement topologisés, on a les implications suivantes :
$$
\xymatrix{
\text{adique + noethérien} \ar@{=>}[d] \\
\text{adique + idéal de définition de type fini} \ar@{=>}[d] \\
\text{adique} \ar@{=>}[d] \\
\text{faiblement adique} \ar@{=>}[d] \\
\text{admissible + à base dénombrable} \ar@{=>}[d] \ar@{=>}[r] &
\text{admissible} \ar@{=>}[d] \\
\text{faiblement admissible + à base dénombrable} \ar@{=>}[r] &
\text{faiblement admissible}
}
$$
où « à base dénombrable » signifie que notre anneau topologique possède un système fondamental
dénombrable d'idéaux ouverts. On dispose d'un diagramme d'implications analogue
pour les anneaux linéairement topologisés non complets (c'est-à-dire en utilisant les notions
de préadique, faiblement préadique, préadmissible et faiblement préadmissible).
Contrairement à ce qui se passe pour les anneaux préadiques, le complété d'un anneau
faiblement préadique est faiblement adique, comme le montre le lemme suivant
qui caractérise les anneaux faiblement préadiques.

\begin{lemma}
\label{formal-spaces-lemma-weakly-adic}
Soit $A$ un anneau linéairement topologisé. Les conditions suivantes sont équivalentes :
\begin{enumerate}
\item $A$ est faiblement préadique ;
\item il existe un homomorphisme continu tendu d'anneaux $A' \to A$,
où $A'$ est un anneau topologique préadique ;
\item $A$ est préadmissible et il existe un idéal de définition $I$
tel que l'adhérence de $I^n$ soit ouverte pour tout $n \geq 1$ ; et
\item $A$ est préadmissible et, pour tout idéal de définition $I$,
l'adhérence de $I^n$ est ouverte pour tout $n \geq 1$.
\end{enumerate}
Le complété d'un anneau faiblement préadique est faiblement adique.
Si $A$ est faiblement adique, alors $A$ est admissible et possède un système fondamental
dénombrable d'idéaux ouverts.
\end{lemma}

\begin{proof}
Supposons (1). Choisissons un idéal $I$ tel que l'adhérence
de $I^n$ soit ouverte pour tout $n$ et que ces adhérences forment
un système fondamental d'idéaux ouverts. Notons $A' = A$ muni de
la topologie $I$-adique. Alors $A' \to A$ est tendu par définition,
et la condition (2) est vérifiée.

\medskip\noindent
Supposons (2). Soit $I' \subset A'$ un idéal de définition.
Notons $I$ l'adhérence de $I'A$. Le caractère tendu de $A' \to A$ signifie
que les adhérences $I_n$ des $(I')^nA$ sont ouvertes et forment un système fondamental
d'idéaux ouverts. Ainsi, $I = I_1$ est ouvert, et les adhérences de $I^n$
sont égales à $I_n$ ; elles sont donc ouvertes et forment un système fondamental
d'idéaux ouverts. En particulier, $I$ est un idéal de définition
tel que l'adhérence de $I^n$ soit ouverte pour tout $n$. La condition (3) est donc vérifiée.

\medskip\noindent
Si $I \subset A$ est comme dans (3), alors $I$ est un idéal comme dans
la définition \ref{formal-spaces-definition-weakly-adic}, et la condition (1) est vérifiée.
De plus, si $I' \subset A$ est un autre idéal de définition quelconque, alors
$I'$ est ouvert (voir Compléments d'algèbre, définition
\ref{more-algebra-definition-topological-ring})
et contient donc $I^n$ pour un certain $n \geq 1$.
Ainsi, $(I')^m$ contient $I^{nm}$ pour tout $m \geq 1$, et l'on
conclut que les adhérences de $(I')^m$ sont ouvertes pour tout $m$.
On voit ainsi que (3) implique (4).
L'implication (4) $\Rightarrow$ (3) est triviale.

\medskip\noindent
Soit $A$ faiblement préadique. Choisissons $A' \to A$ comme dans (2).
D'après les lemmes \ref{formal-spaces-lemma-completion-taut} et
\ref{formal-spaces-lemma-composition-taut},
le composé $A' \to A^\wedge$ est tendu. Ainsi, $A^\wedge$
est faiblement préadique par l'équivalence de (2) et (1).
Comme le complété d'un anneau linéairement topologisé $A$ est complet
(Compléments d'algèbre, section \ref{more-algebra-section-topological-ring}),
on voit que $A^\wedge$ est faiblement adique.

\medskip\noindent
Soit $A$ faiblement adique. Alors $A$ est complet
et préadmissible par (1) $\Rightarrow$ (3) ; par conséquent, $A$
est admissible. Bien entendu, par définition, $A$ possède un système fondamental
dénombrable d'idéaux ouverts.
\end{proof}

\noindent
Nous donnons deux critères qui garantissent qu'un anneau faiblement adique
est adique et possède un idéal de définition de type fini.

\begin{lemma}
\label{formal-spaces-lemma-weakly-adic-finite-generation-bis}
Soit $A$ un anneau complet linéairement topologisé. Soit $I \subset A$ un idéal
de type fini tel que l'adhérence de $I^n$ soit ouverte pour tout
$n \geq 0$ et que ces adhérences forment un système fondamental d'idéaux ouverts.
Alors $A$ est adique et possède un idéal de définition de type fini.
\end{lemma}

\begin{proof}
Notons $A'$ l'anneau $A$ muni de la topologie $I$-adique.
Les hypothèses nous disent que $A' \to A$ est tendu.
On conclut par le lemme \ref{formal-spaces-lemma-taut-is-adic} (plus précisément, ce
lemme nous dit également que $I$ est un idéal de définition).
\end{proof}

\begin{lemma}
\label{formal-spaces-lemma-weakly-adic-finite-generation}
Soit $A$ un anneau topologique faiblement adique. Soit $I$ un
idéal de définition tel que $I/I_2$ soit un module de type fini,
où $I_2$ est l'adhérence de $I^2$.
Alors $A$ est adique et possède un idéal de définition de type fini.
\end{lemma}

\begin{proof}
Nous utilisons la caractérisation du lemme \ref{formal-spaces-lemma-weakly-adic}
sans autre mention.
Choisissons $f_1, \ldots, f_r \in I$ dont les images engendrent $I/I_2$.
Posons $I' = (f_1, \ldots, f_r)$. On a $I' + I_2 = I$.
Alors $I_2$ est l'adhérence de $I^2 = (I' + I_2)^2 \subset I' + I_3$,
où $I_3$ est l'adhérence de $I^3$. Par conséquent, $I' + I_3 = I$.
En poursuivant ainsi, on voit que $I' + I_n = I$ pour tout $n \geq 2$,
où $I_n$ est l'adhérence de $I^n$.
Autrement dit, l'adhérence de $I'$ dans $A$ est $I$.
Par conséquent, l'adhérence de $(I')^n$ est $I_n$. Les adhérences de $(I')^n$
forment donc un système fondamental d'idéaux ouverts de $A$.
On conclut par le lemme \ref{formal-spaces-lemma-weakly-adic-finite-generation-bis}.
\end{proof}

\noindent
Une propriété essentielle de la condition « faiblement préadique » est qu'elle se transmet au but
par les homomorphismes d'anneaux tendus entre anneaux linéairement topologisés.

\begin{lemma}
\label{formal-spaces-lemma-taut-ascent-weakly-adic}
Soit $\varphi : A \to B$ un homomorphisme continu d'anneaux
linéairement topologisés. Si $\varphi$ est tendu et si $A$
est faiblement préadique, alors $B$ est faiblement préadique.
\end{lemma}

\begin{proof}
Soit $I \subset A$ un idéal tel que l'adhérence $I_n$ de $I^n$
soit ouverte et que ces adhérences définissent un système fondamental d'idéaux ouverts.
Alors l'adhérence de $I^nB$ est égale à celle de $I_nB$.
Comme $\varphi$ est tendu, ces adhérences sont ouvertes et forment un système fondamental
d'idéaux ouverts de $B$. Ainsi, $B$ est faiblement préadique.
\end{proof}

\begin{lemma}
\label{formal-spaces-lemma-completed-tensor-product-weakly-adic}
Soient $B \to A$ et $B \to C$ des homomorphismes continus d'anneaux
linéairement topologisés. Si $A$ et $C$ sont faiblement préadiques, alors
$A \widehat{\otimes}_B C$ est faiblement adique.
\end{lemma}

\begin{proof}
Nous utiliserons la caractérisation du lemme \ref{formal-spaces-lemma-weakly-adic}
sans autre mention. D'après le lemme \ref{formal-spaces-lemma-completed-tensor-product},
nous savons que $A \widehat{\otimes}_B C$ est admissible.
De plus, la démonstration de ce lemme montre que
l'adhérence $K \subset A \widehat{\otimes}_B C$
est un idéal de définition lorsque $I \subset A$ et $J \subset C$
le sont, cette adhérence étant celle de $I(A \widehat{\otimes}_B C) + J(A \widehat{\otimes}_B C)$.
Il suffit alors de montrer que
l'adhérence de $K^n$ est ouverte pour tout $n \geq 1$.
Puisque l'idéal $K^n$ contient
$I^n(A \widehat{\otimes}_B C) + J^n(A \widehat{\otimes}_B C)$,
puisque l'adhérence de $I^n$ dans $A$ est ouverte et
puisque l'adhérence de $J^n$ dans $C$ est ouverte, on
voit que l'adhérence de $K^n$ est ouverte dans $A \widehat{\otimes}_B C$.
\end{proof}






\section{Descente des propriétés}
\label{formal-spaces-section-taut-descent}

\noindent
Dans cette section, nous considérons la situation suivante :
\begin{enumerate}
\item $\varphi : A \to B$ est un homomorphisme continu d'anneaux
topologiques linéairement topologisés ;
\item $\varphi$ est tendu ; et
\item pour tout idéal ouvert $I \subset A$, si $J \subset B$
désigne l'adhérence de $IB$, alors l'homomorphisme $A/I \to B/J$
est fidèlement plat.
\end{enumerate}
Nous allons montrer que, dans cette situation, les propriétés de $B$ sont héritées
par $A$.

\begin{lemma}
\label{formal-spaces-lemma-taut-descent-countable}
Dans la situation ci-dessus, si $B$ possède un système fondamental dénombrable
d'idéaux ouverts, alors $A$ possède un système fondamental dénombrable d'idéaux ouverts.
\end{lemma}

\begin{proof}
Choisissons un système fondamental $B \supset J_1 \supset J_2 \supset \ldots$
d'idéaux ouverts. Puisque $\varphi$ est tendu, pour tout $n$, on peut trouver un
idéal ouvert $I_n$ tel que $J_n \supset I_nB$. Nous affirmons que les
$I_n$ forment un système fondamental d'idéaux ouverts de $A$.
En effet, supposons que $I \subset A$ soit ouvert. Comme $\varphi$ est tendu,
l'adhérence de $IB$ est ouverte et contient donc $J_n$ pour un certain $n$
assez grand. Ainsi, $I_nB \subset IB$. Soit $J$ l'adhérence de
$IB$ dans $B$. Puisque $A/I \to B/J$ est fidèlement plat, il est injectif.
Par conséquent, comme $I_n \to A/I \to B/J$ est nul puisque $I_nB \subset IB \subset J$,
on conclut que $I_n \to A/I$ est nul. Ainsi, $I_n \subset I$,
ce qui achève la démonstration.
\end{proof}

\begin{lemma}
\label{formal-spaces-lemma-taut-descent-weakly-admissible}
Dans la situation ci-dessus, si $B$ est faiblement préadmissible, alors
$A$ est faiblement préadmissible.
\end{lemma}

\begin{proof}
Soit $J \subset B$ un idéal faible de définition. Soit $I \subset A$
un idéal ouvert tel que $IB \subset J$. Pour montrer que $I$ est un idéal faible
de définition, il faut montrer que tout $f \in I$ est topologiquement nilpotent.
Soit $I' \subset A$ un idéal ouvert. Notons $J' \subset B$ l'adhérence
de $I'B$. Alors $A/I' \to B/J'$ est fidèlement plat, donc injectif.
Ainsi, pour montrer que $f^n \in I'$, il suffit de montrer
que $\varphi(f)^n \in J'$. C'est le cas pour $n \gg 0$, puisque
$\varphi(f) \in J$, que l'idéal $J$ est un idéal faible de définition de $B$
et que $J'$ est ouvert dans $B$.
\end{proof}

\begin{lemma}
\label{formal-spaces-lemma-taut-descent-admissible}
Dans la situation ci-dessus, si $B$ est préadmissible, alors $A$
est préadmissible.
\end{lemma}

\begin{proof}
Soit $J \subset B$ un idéal faible de définition. Soit $I \subset A$
un idéal ouvert tel que $IB \subset J$. Soit $I' \subset A$ un idéal ouvert.
Pour montrer que $I$ est un idéal de définition, il faut montrer que
$I^n \subset I'$ pour $n \gg 0$. Notons $J' \subset B$ l'adhérence
de $I'B$. Alors $A/I' \to B/J'$ est fidèlement plat, donc injectif.
Ainsi, pour montrer que $I^n \subset I'$, il suffit de montrer
que $\varphi(I)^n \subset J'$. C'est le cas pour $n \gg 0$, puisque
$\varphi(I) \subset J$, que l'idéal $J$ est un idéal de définition de $B$
et que $J'$ est ouvert dans $B$.
\end{proof}

\begin{lemma}
\label{formal-spaces-lemma-taut-descent-weakly-adic}
Dans la situation ci-dessus, si $B$ est faiblement préadique, alors $A$
est faiblement préadique.
\end{lemma}

\begin{proof}
Nous utiliserons sans autre mention la caractérisation des anneaux faiblement préadiques donnée dans
le lemme \ref{formal-spaces-lemma-weakly-adic}.
D'après le lemme \ref{formal-spaces-lemma-taut-descent-admissible},
l'anneau topologique $A$ est préadmissible.
Soit $I \subset A$ un idéal de définition.
Fixons $n \geq 1$. Pour démontrer le lemme, il faut
montrer que l'adhérence de $I^n$ est ouverte.
Soient les $I_\lambda \subset A$ un système fondamental d'idéaux ouverts.
Notons $J \subset B$, resp.\ $J_\lambda \subset B$,
l'adhérence de $IB$, resp.\ celle de $I_\lambda B$.
Comme $B$ est faiblement préadique, l'adhérence de $J^n$ est ouverte.
Il existe donc $\lambda$ tel que
$$
J_\lambda \subset \bigcap\nolimits_\mu (J^n + J_\mu)
$$
car le membre de droite est l'adhérence de $J^n$ d'après
le lemme \ref{formal-spaces-lemma-closed}.
Cela signifie que l'image de $J_\lambda$ dans $B/J_\mu$ est contenue dans
l'image de $J^n$ dans $B/J_\mu$. Observons que l'image de $J^n$ dans
$B/J_\mu$ est égale à celle de $I^nB$ dans $B/J_\mu$ (puisque tout élément
de $J$ est congru à un élément de $IB$ modulo $J_\mu$).
Comme $A/I_\mu \to B/J_\mu$ est fidèlement plat et que
$I_\lambda B \subset J_\lambda$,
on conclut que l'image de $I_\lambda$ dans $A/I_\mu$
est contenue dans l'image de $I^n$.
On en déduit que $I_\lambda$ est contenu dans l'adhérence
de $I^n$, ce qui achève la démonstration.
\end{proof}

\begin{lemma}
\label{formal-spaces-lemma-taut-descent-adic-star}
Dans la situation ci-dessus, si $B$ est adique et possède un idéal de définition
de type fini, et si $A$ est complet, alors $A$ est adique et
possède un idéal de définition de type fini.
\end{lemma}

\begin{proof}
Nous savons déjà que $A$ est faiblement adique, et a fortiori admissible, d'après
le lemme \ref{formal-spaces-lemma-taut-descent-weakly-adic}
(et le lemme \ref{formal-spaces-lemma-weakly-adic} pour voir que les anneaux adiques sont faiblement adiques).
Soit $I \subset A$ un idéal de définition.
Soit $J \subset B$ un idéal de définition de type fini.
Comme l'adhérence de $IB$ est ouverte, on peut trouver $n > 0$ tel que
$J^n$ soit contenu dans l'adhérence de $IB$. Après avoir remplacé $J$
par $J^n$, on peut donc supposer que $J$ est un idéal de définition
de type fini contenu dans l'adhérence de $IB$. D'après le lemme \ref{formal-spaces-lemma-closed},
cela implique en particulier que
$$
J \subset IB + J^2
$$
Considérons le $A$-module de type fini $M = (J + IB)/IB$.
L'équation affichée montre que $JM = M$.
D'après le lemme \ref{formal-spaces-lemma-weakly-admissible-henselian}, par exemple,
on voit que $J$ est contenu dans le radical de Jacobson de $B$.
Le lemme de Nakayama, plus précisément la partie (2) de
Algèbre, lemme \ref{algebra-lemma-NAK}, donne donc $M = 0$.
Ainsi, $J \subset IB$.

\medskip\noindent
Comme $J$ est de type fini, on peut trouver un idéal de type fini
$I' \subset I$ tel que $J \subset I'B$.
Comme $A \to B$ est continu, que $J \subset B$ est ouvert et que
$I$ est un idéal de définition, on peut trouver $n > 0$ tel que
$I^nB \subset J$. Soit $J_{n + 1} \subset B$ l'adhérence de
$I^{n + 1}B$.
On a
$$
I^n \cdot (B/J_{n + 1}) \subset J \cdot (B/J_{n + 1}) \subset
I' \cdot (B/J_{n + 1})
$$
Puisque $A/I^{n + 1} \to B/J_{n + 1}$ est fidèlement plat, cela
implique $I^n \cdot (A/I^{n + 1}) \subset I' \cdot (A/I^{n + 1})$,
ce qui signifie à son tour que
$$
I^n \subset I' + I^{n + 1}
$$
Cela implique $I^n \subset I'  + I^{n + k}$ pour tout $k \geq 1$,
ce qui implique à son tour $I^{nm} \subset (I')^m + I^{nm + k}$
pour tous $k, m \geq 1$. Il s'ensuit que l'adhérence de
$(I')^m$ contient $I^{nm}$. Comme l'adhérence de $I^{nm}$ est ouverte,
puisque $A$ est faiblement adique,
on conclut que l'adhérence de $(I')^m$ est ouverte pour tout $m$.
Comme ces adhérences forment un système fondamental d'idéaux ouverts
de $A$ (puisqu'il en est ainsi des adhérences de $I^n$),
on conclut par le lemme \ref{formal-spaces-lemma-weakly-adic-finite-generation-bis}.
\end{proof}









\section{Espaces algébriques formels affines}
\label{formal-spaces-section-affine-formal-algebraic-spaces}

\noindent
Dans cette section, nous introduisons les espaces algébriques formels affines.
Ceux-ci sont en fait les mêmes objets que ceux appelés schémas
formels affines dans \cite{BVGD}. Nous les appellerons cependant
espaces algébriques formels affines, afin d'éviter toute confusion avec
la notion de schéma formel affine définie dans \cite{EGA}.

\medskip\noindent
Rappelons qu'un épaississement de schémas est une immersion
fermée qui induit une surjection sur les espaces topologiques
sous-jacents ; voir Compléments sur les morphismes, définition
\ref{more-morphisms-definition-thickening}.

\begin{definition}
\label{formal-spaces-definition-affine-formal-algebraic-space}
Soit $S$ un schéma. On dit qu'un faisceau $X$ sur $(\Sch/S)_{fppf}$ est un
{\it espace algébrique formel affine} s'il existe :
\begin{enumerate}
\item un ensemble ordonné filtrant $\Lambda$ ;
\item un système $(X_\lambda, f_{\lambda \mu})$ indexé par $\Lambda$
dans $(\Sch/S)_{fppf}$ tel que :
\begin{enumerate}
\item chaque $X_\lambda$ soit affine ;
\item chaque $f_{\lambda \mu} : X_\lambda \to X_\mu$ soit un épaississement.
\end{enumerate}
\end{enumerate}
et tels que
$$
X \cong \colim_{\lambda \in \Lambda} X_\lambda
$$
comme faisceaux fppf, et que $X$ satisfasse une condition ensembliste
(voir la remarque \ref{formal-spaces-remark-set-theoretic}). Un
{\it morphisme d'espaces algébriques formels affines}
sur $S$ est une application de faisceaux.
\end{definition}

\noindent
Observons que le système $(X_\lambda, f_{\lambda \mu})$ ne fait pas
partie des données. Supposons que $U$ soit un schéma quasi-compact sur $S$.
Comme les applications de transition sont des monomorphismes, on a
$$
X(U) = \colim X_\lambda(U)
$$
d'après Sites, lemme \ref{sites-lemma-directed-colimits-sections}.
Ainsi, la construction du faisceau fppf inhérente à la limite inductive de la
définition se réduit à la construction du faisceau de Zariski, laquelle ne change
rien pour les schémas quasi-compacts.

\begin{lemma}
\label{formal-spaces-lemma-diagonal-affine-formal-algebraic-space}
Soit $S$ un schéma. Si $X$ est un espace algébrique formel affine sur
$S$, alors le morphisme diagonal $\Delta : X \to X \times_S X$
est représentable et est une immersion fermée.
\end{lemma}

\begin{proof}
Supposons donnés $U \to X$ et $V \to X$, où $U, V$ sont des schémas sur $S$.
Montrons que $U \times_X V$ est représentable. Écrivons $X = \colim X_\lambda$
comme dans la définition \ref{formal-spaces-definition-affine-formal-algebraic-space}.
La discussion précédente montre que, localement pour Zariski sur $U$ et $V$, les morphismes
se factorisent par un certain $X_\lambda$. Dans ce cas,
$U \times_X V = U \times_{X_\lambda} V$, qui est un schéma.
La diagonale est donc représentable ; voir
Espaces, lemme \ref{spaces-lemma-representable-diagonal}.
Étant donné $(a, b) : W \to X \times_S X$, où $W$ est un schéma sur $S$,
considérons l'application $X \times_{\Delta, X \times_S X, (a, b)} W \to W$.
Comme précédemment, localement sur $W$, les morphismes $a$ et $b$ se factorisent par
le schéma affine $X_\lambda$ pour un certain $\lambda$ ; on obtient alors
le morphisme
$X_\lambda
\times_{\Delta_\lambda, X_\lambda \times_S X_\lambda, (a, b)} W \to W$.
Il est obtenu par changement de base à partir de
$\Delta_\lambda : X_\lambda \to X_\lambda \times_S X_\lambda$,
qui est une immersion fermée puisque $X_\lambda \to S$ est séparé
(car $X_\lambda$ est affine).
Ainsi, $X \to X \times_S X$ est une immersion fermée.
\end{proof}

\noindent
Un morphisme de schémas $X \to X'$ est un épaississement s'il est une immersion fermée
et induit une surjection sur les ensembles de points sous-jacents ; voir
(Compléments sur les morphismes, définition
\ref{more-morphisms-definition-thickening}).
Par conséquent, la propriété d'être un épaississement est conservée par tout
changement de base et est locale pour la topologie fpqc sur le but ; voir
Espaces, section \ref{spaces-section-lists}.
Ainsi, Espaces, définition \ref{spaces-definition-relative-representable-property},
s'applique à la propriété « être un épaississement », et nous savons ce que signifie, pour une
transformation représentable $F \to G$ de
préfaisceaux sur $(\Sch/S)_{fppf}$, être un épaississement.
Observons que cela est compatible avec notre définition
(Compléments sur les morphismes d'espaces algébriques, définition
\ref{spaces-more-morphisms-definition-thickening})
des épaississements lorsque $F$ et $G$ sont des espaces algébriques.

\begin{lemma}
\label{formal-spaces-lemma-covering-by-thickenings}
Soient $X_\lambda, \lambda \in \Lambda$ et $X = \colim X_\lambda$
comme dans la définition \ref{formal-spaces-definition-affine-formal-algebraic-space}.
Alors $X_\lambda \to X$ est représentable et est un épaississement.
\end{lemma}

\begin{proof}
L'énoncé a un sens d'après la discussion de
Espaces, section \ref{spaces-section-representable} et
\ref{spaces-section-representable-properties}.
D'après le lemme \ref{formal-spaces-lemma-diagonal-affine-formal-algebraic-space},
les morphismes $X_\lambda \to X$ sont représentables.
Étant donné $U \to X$, où $U$ est un schéma,
la discussion qui suit
la définition \ref{formal-spaces-definition-affine-formal-algebraic-space}
montre que, localement pour Zariski sur $U$, le
morphisme se factorise par un certain $X_\mu$ avec $\lambda \leq \mu$.
Dans ce cas, $U \times_X X_\lambda = U \times_{X_\mu} X_\lambda$,
de sorte que $U \times_X X_\lambda \to U$ est obtenu par changement de base à partir de
l'épaississement $X_\lambda \to X_\mu$.
\end{proof}

\begin{lemma}
\label{formal-spaces-lemma-factor-through-thickening}
Soient $X_\lambda, \lambda \in \Lambda$ et $X = \colim X_\lambda$
comme dans la définition \ref{formal-spaces-definition-affine-formal-algebraic-space}.
Si $Y$ est un espace algébrique quasi-compact sur $S$, tout
morphisme $Y \to X$ se factorise par un $X_\lambda$.
\end{lemma}

\begin{proof}
Choisissons un schéma affine $V$ et un morphisme étale surjectif
$V \to Y$. Le composé $V \to Y \to X$ se factorise par
$X_\lambda$ pour un certain $\lambda$, d'après la discussion qui suit
la définition \ref{formal-spaces-definition-affine-formal-algebraic-space}.
Comme $V \to Y$ est une surjection de faisceaux, on conclut.
\end{proof}

\begin{lemma}
\label{formal-spaces-lemma-characterize-affine-formal-algebraic-space}
Soit $S$ un schéma. Soit $X$ un faisceau sur $(\Sch/S)_{fppf}$.
Alors $X$ est un espace algébrique formel affine si et seulement si
les conditions suivantes sont satisfaites :
\begin{enumerate}
\item tout morphisme $U \to X$, où $U$ est un schéma affine sur $S$,
se factorise par un morphisme $T \to X$ qui est représentable et est un
épaississement, avec $T$ un schéma affine sur $S$ ; et
\item une condition ensembliste comme dans la remarque \ref{formal-spaces-remark-set-theoretic}.
\end{enumerate}
\end{lemma}

\begin{proof}
Il résulte des lemmes \ref{formal-spaces-lemma-covering-by-thickenings} et
\ref{formal-spaces-lemma-factor-through-thickening} qu'un espace algébrique formel affine
satisfait (1) et (2). Pour démontrer la réciproque, nous pouvons
supposer que $X$ n'est pas vide.
Soit $\Lambda$ la catégorie des morphismes représentables $T \to X$ qui sont des
épaississements, où $T$ est un schéma affine sur $S$. Cette catégorie
est filtrante. Puisque $X$ n'est pas vide, $\Lambda$ contient au moins un
objet. Si $T \to X$ et $T' \to X$ appartiennent à $\Lambda$, on peut
factoriser $T \amalg T' \to X$ par $T'' \to X$ dans $\Lambda$. Entre
deux objets quelconques de $\Lambda$, il existe une flèche unique ou aucune. Ainsi,
$\Lambda$ est un ensemble ordonné filtrant et, par hypothèse,
$X = \colim_{T \to X\text{ dans }\Lambda} T$. Pour achever la démonstration,
il faut montrer que toute flèche $T \to T'$ de $\Lambda$ est un épaississement.
C'est vrai parce que $T' \to X$ est un monomorphisme de faisceaux, de sorte que
$T = T \times_{T'} T' = T \times_X T'$ ; par conséquent, le morphisme
$T \to T'$ est égal à la projection $T \times_X T' \to T'$, qui est
un épaississement puisque $T \to X$ en est un.
\end{proof}

\noindent
Pour un espace algébrique formel affine général $X$, rien ne garantit
que $X$ possède assez de fonctions pour séparer les points, par exemple.
Voir Exemples, section \ref{examples-section-affine-formal-algebraic-space}.
Pour caractériser ceux qui en possèdent assez, nous proposons le lemme suivant.

\begin{lemma}
\label{formal-spaces-lemma-mcquillan-affine-formal-algebraic-space}
Soit $S$ un schéma. Soit $X$ un faisceau fppf sur $(\Sch/S)_{fppf}$
qui satisfait la condition ensembliste de la
remarque \ref{formal-spaces-remark-set-theoretic}.
Les conditions suivantes sont équivalentes :
\begin{enumerate}
\item il existe un anneau topologique faiblement admissible $A$ sur $S$
(voir la remarque \ref{formal-spaces-remark-mcquillan}) tel que
$X = \colim_{I \subset A\text{ idéal faible de définition}} \Spec(A/I)$ ;
\item $X$ est un espace algébrique formel affine et
il existe une $S$-algèbre $A$ et une application $X \to \Spec(A)$
telles que, pour une immersion fermée $T \to X$ avec $T$ schéma affine,
le composé $T \to \Spec(A)$ soit une immersion fermée ;
\item $X$ est un espace algébrique formel affine et
il existe une $S$-algèbre $A$ et une application $X \to \Spec(A)$
telles que, pour une immersion fermée $T \to X$ avec $T$ schéma,
le composé $T \to \Spec(A)$ soit une immersion fermée ;
\item $X$ est un espace algébrique formel affine et, pour un choix de
$X = \colim X_\lambda$ comme dans
la définition \ref{formal-spaces-definition-affine-formal-algebraic-space},
les projections $\lim \Gamma(X_\lambda, \mathcal{O}_{X_\lambda})
\to \Gamma(X_\lambda, \mathcal{O}_{X_\lambda})$ sont surjectives ;
\item $X$ est un espace algébrique formel affine et, pour tout choix
de $X = \colim X_\lambda$ comme dans
la définition \ref{formal-spaces-definition-affine-formal-algebraic-space},
les projections $\lim \Gamma(X_\lambda, \mathcal{O}_{X_\lambda})
\to \Gamma(X_\lambda, \mathcal{O}_{X_\lambda})$ sont surjectives.
\end{enumerate}
De plus, l'anneau topologique faiblement admissible est
$A = \lim \Gamma(X_\lambda, \mathcal{O}_{X_\lambda})$,
muni de sa topologie limite, et les idéaux faibles de définition
classifient exactement les morphismes $T \to X$ qui sont représentables
et sont des épaississements.
\end{lemma}

\begin{proof}
Il est clair que (5) implique (4).

\medskip\noindent
Supposons (4) pour $X = \colim X_\lambda$ comme dans
la définition \ref{formal-spaces-definition-affine-formal-algebraic-space}.
Posons $A = \lim \Gamma(X_\lambda, \mathcal{O}_{X_\lambda})$.
Soit $T \to X$ une immersion fermée, avec $T$ un schéma
(remarquons que $T \to X$ est représentable d'après
le lemme \ref{formal-spaces-lemma-diagonal-affine-formal-algebraic-space}).
Puisque $X_\lambda \to X$ est un épaississement, il en va de même de
$X_\lambda \times_X T \to T$. D'autre part,
$X_\lambda \times_X T \to X_\lambda$ est une immersion fermée ;
$X_\lambda \times_X T$ est donc affine. Par conséquent, $T$ est affine
d'après Limits, Proposition \ref{limits-proposition-affine}.
Alors $T \to X$ se factorise par $X_\lambda$ pour un certain $\lambda$,
d'après le lemme \ref{formal-spaces-lemma-factor-through-thickening}.
Ainsi, $A \to \Gamma(X_\lambda, \mathcal{O}) \to \Gamma(T, \mathcal{O})$
est surjectif. La condition (3) est donc vérifiée.

\medskip\noindent
Il est clair que (3) implique (2).

\medskip\noindent
Supposons (2) pour $A$ et $X \to \Spec(A)$. Écrivons $X = \colim X_\lambda$
comme dans la définition \ref{formal-spaces-definition-affine-formal-algebraic-space}.
Alors $A_\lambda = \Gamma(X_\lambda, \mathcal{O})$ est un quotient
de $A$ par l'hypothèse (2). Ainsi, $A^\wedge = \lim A_\lambda$
est un anneau topologique complet ; voir la discussion de
Compléments d'algèbre, section \ref{more-algebra-section-topological-ring}.
Les applications $A^\wedge \to A_\lambda$ sont surjectives, puisque $A \to A_\lambda$ l'est.
Nous affirmons que, pour tout $\lambda$, le noyau $I_\lambda \subset A^\wedge$ de
$A^\wedge \to A_\lambda$ est un idéal faible de définition.
En effet, il est ouvert par définition de la topologie limite.
Si $f \in I_\lambda$, alors, pour tout $\mu \in \Lambda$,
l'image de $f$ dans $A_\mu$ est nulle dans tous les corps résiduels
des points de $X_\mu$. C'est donc un élément nilpotent
de $A_\mu$. Par conséquent, une puissance $f^n \in I_\mu$. Ainsi, $f^n \to 0$
lorsque $n \to 0$. L'anneau $A^\wedge$ est donc faiblement admissible.
Enfin, supposons que $I \subset A^\wedge$ soit un idéal faible
de définition. Alors $I \subset A^\wedge$ est ouvert ; il existe donc
$\lambda$ tel que $I \supset I_\lambda$. On obtient ainsi un morphisme
$\Spec(A^\wedge/I) \to \Spec(A_\lambda) \to X$.
Il s'ensuit que $X = \colim \Spec(A^\wedge/I)$, où
la limite inductive porte maintenant sur tous les idéaux faibles de définition.
La condition (1) est donc vérifiée.

\medskip\noindent
Supposons (1). Dans ce cas, il est clair que $X$ est un espace algébrique
formel affine. Soit $X = \colim X_\lambda$ une présentation quelconque comme dans
la définition \ref{formal-spaces-definition-affine-formal-algebraic-space}.
Pour tout $\lambda$, on peut trouver un idéal faible de définition
$I \subset A$ tel que $X_\lambda \to X$ se factorise par
$\Spec(A/I) \to X$ ; voir le lemme \ref{formal-spaces-lemma-factor-through-thickening}.
Alors $X_\lambda = \Spec(A/I_\lambda)$ avec $I \subset I_\lambda$.
Réciproquement, pour tout idéal faible de définition $I \subset A$,
le morphisme $\Spec(A/I) \to X$ se factorise par $X_\lambda$
pour un certain $\lambda$, c'est-à-dire que $I_\lambda \subset I$.
Il s'ensuit que chaque $I_\lambda$ est un idéal faible de définition
et qu'ils forment une partie cofinale de l'ensemble des idéaux faibles
de définition. Ainsi, $A = \lim A/I = \lim A/I_\lambda$,
et l'on voit que (5) est vraie et, de plus, que
$A = \lim \Gamma(X_\lambda, \mathcal{O}_{X_\lambda})$.
\end{proof}

\noindent
Ce lemme permet de poser la définition suivante.

\begin{definition}
\label{formal-spaces-definition-types-affine-formal-algebraic-space}
Soit $S$ un schéma. Soit $X$ un espace algébrique formel affine sur $S$.
On dit que $X$ est {\it de McQuillan} si $X$ satisfait les conditions équivalentes
du lemme \ref{formal-spaces-lemma-mcquillan-affine-formal-algebraic-space}. Soit $A$
l'anneau topologique faiblement admissible associé à $X$. On dit que :
\begin{enumerate}
\item $X$ est {\it classique} si $X$ est de McQuillan et si $A$ est admissible
(Compléments d'algèbre, définition \ref{more-algebra-definition-topological-ring}) ;
\item $X$ est {\it faiblement adique} si $X$ est de McQuillan et si $A$ est faiblement adique
(définition \ref{formal-spaces-definition-weakly-adic}) ;
\item $X$ est {\it adique} si $X$ est de McQuillan et si $A$ est adique
(Compléments d'algèbre, définition \ref{more-algebra-definition-topological-ring}) ;
\item $X$ est {\it adique*} si $X$ est de McQuillan, si $A$ est adique et si $A$
possède un idéal de définition de type fini ; et
\item $X$ est {\it noethérien} si $X$ est de McQuillan et si $A$ est
à la fois noethérien et adique.
\end{enumerate}
\end{definition}

\noindent
Dans \cite{Fujiwara-Kato}, les auteurs emploient la terminologie « de type idéal fini »
pour exprimer qu'un anneau topologique adique $A$ contient un idéal de définition
de type fini. Étant donné un espace algébrique formel affine $X$,
voici les implications entre les notions introduites dans la définition :
$$
\xymatrix{
X\text{ noethérien} \ar@{=>}[r] &
X\text{ adique*} \ar@{=>}[r] &
X\text{ adique} \ar@{=>}[lld] \\
X\text{ faiblement adique} \ar@{=>}[r] &
X\text{ classique} \ar@{=>}[r] &
X\text{ de McQuillan}
}
$$
Voir la discussion de la section \ref{formal-spaces-section-weakly-adic} et, pour un énoncé précis,
le lemme \ref{formal-spaces-lemma-implications-between-types}.

\begin{remark}
\label{formal-spaces-remark-compare-with-affine-formal-schemes}
Les espaces algébriques formels affines classiques correspondent aux
schémas formels affines considérés dans EGA (\cite{EGA}). Pour l'expliquer,
supposons que notre schéma de base soit $\Spec(\mathbf{Z})$. Soit
$\mathfrak X = \text{Spf}(A)$ un schéma formel affine.
Soit $h_\mathfrak X$ son foncteur des points comme dans
le lemme \ref{formal-spaces-lemma-fully-faithful}.
Alors $h_\mathfrak X = \colim h_{\Spec(A/I)}$, où la limite inductive
porte sur la famille des idéaux de définition de l'anneau topologique
admissible $A$. Cela résulte de
(\ref{formal-spaces-equation-morphisms-affine-formal-schemes})
après évaluation sur les schémas affines, et il suffit de vérifier l'égalité
sur les schémas affines puisque les deux membres sont des faisceaux fppf ; voir
le lemme \ref{formal-spaces-lemma-formal-scheme-sheaf-fppf}.
Ainsi, $h_\mathfrak X$ est un espace algébrique formel affine.
En fait, c'est un espace algébrique formel affine classique
d'après la définition \ref{formal-spaces-definition-types-affine-formal-algebraic-space}.
Le lemme \ref{formal-spaces-lemma-fully-faithful} nous dit donc que
la catégorie des schémas formels affines est équivalente à celle
des espaces algébriques formels affines classiques.
\end{remark}

\noindent
Le lien avec les schémas formels affines étant établi,
il est naturel de poser la définition suivante.

\begin{definition}
\label{formal-spaces-definition-affine-formal-spectrum}
Soit $S$ un schéma. Soit $A$ un anneau topologique faiblement admissible sur
$S$ ; voir la définition \ref{formal-spaces-definition-weakly-admissible}\footnote{Voir
Compléments d'algèbre, définition
\ref{more-algebra-definition-topological-ring}
pour le cas classique, et la remarque \ref{formal-spaces-remark-mcquillan}
pour une discussion des différences.}.
Le {\it spectre formel} de $A$ est l'espace algébrique formel affine
$$
\text{Spf}(A) = \colim \Spec(A/I)
$$
où la limite inductive est prise sur l'ensemble des idéaux faibles de définition de $A$
et dans la catégorie $\Sh((\Sch/S)_{fppf})$.
\end{definition}

\noindent
Un tel spectre formel est de McQuillan par construction et, réciproquement,
tout espace algébrique formel affine de McQuillan est isomorphe à un
spectre formel. Précisons que notre théorie contient des espaces algébriques
formels affines qui ne sont le spectre formel d'aucun
anneau topologique faiblement admissible.
Suivant \cite{Yasuda}, nous pourrions introduire les $S$-pro-anneaux
comme les pro-objets de la catégorie des $S$-algèbres ; voir Catégories,
remarque \ref{categories-remark-pro-category}. Tout
espace algébrique formel affine sur $S$ serait alors le spectre formel d'un tel
$S$-pro-anneau. Nous ne le ferons pas et travaillerons directement avec les
espaces algébriques formels affines correspondants.

\medskip\noindent
La construction du spectre formel est fonctorielle. Pour l'expliquer,
soit $\varphi : B \to A$ un homomorphisme continu d'anneaux topologiques
faiblement admissibles sur $S$. Alors
$$
\text{Spf}(\varphi) : \text{Spf}(B) \to \text{Spf}(A)
$$
est l'unique morphisme d'espaces algébriques formels affines tel que
les diagrammes
$$
\xymatrix{
\Spec(B/J) \ar[d] \ar[r] & \Spec(A/I) \ar[d] \\
\text{Spf}(B) \ar[r] & \text{Spf}(A)
}
$$
soient commutatifs pour tous les idéaux faibles de définition $I \subset A$ et $J \subset B$
tels que $\varphi(I) \subset J$. Comme la continuité de $\varphi$
implique que, pour tout idéal faible de définition $J \subset B$,
il existe un idéal faible de définition $I \subset A$ ayant la propriété
requise, on voit que les commutativités demandées déterminent
et définissent uniquement $\text{Spf}(\varphi)$.

\begin{lemma}
\label{formal-spaces-lemma-morphism-between-formal-spectra}
Soit $S$ un schéma. Soient $A$, $B$ des anneaux topologiques faiblement admissibles
sur $S$. Tout morphisme $f : \text{Spf}(B) \to \text{Spf}(A)$
d'espaces algébriques formels affines sur $S$
est égal à $\text{Spf}(f^\sharp)$ pour un unique homomorphisme continu
de $S$-algèbres $f^\sharp : A \to B$.
\end{lemma}

\begin{proof}
Soit $f : \text{Spf}(B) \to \text{Spf}(A)$ comme dans le lemme.
Soit $J \subset B$ un idéal faible de définition. D'après
le lemme \ref{formal-spaces-lemma-factor-through-thickening},
il existe un idéal faible de définition $I \subset A$ tel que
$\Spec(B/J) \to \text{Spf}(B) \to \text{Spf}(A)$
se factorise par $\Spec(A/I)$. D'après
Schémas, lemme \ref{schemes-lemma-morphism-into-affine},
on obtient un homomorphisme de $S$-algèbres $A/I \to B/J$.
Ces homomorphismes sont compatibles lorsque $J$ varie et définissent
l'homomorphisme $f^\sharp : A \to B$. Celui-ci est continu, car,
pour tout idéal faible de définition $J \subset B$, il existe un
idéal faible de définition $I \subset A$ tel que
$f^\sharp(I) \subset J$. L'égalité $f = \text{Spf}(f^\sharp)$
résulte du choix des homomorphismes d'anneaux $A/I \to B/J$ qui constituent $f^\sharp$.
\end{proof}

\begin{lemma}
\label{formal-spaces-lemma-presentation-representable}
Soit $S$ un schéma. Soit $f : X \to Y$ une application
de préfaisceaux sur $(\Sch/S)_{fppf}$. Si $X$ est un espace algébrique formel
affine et si $f$ est représentable par des espaces algébriques et localement quasi-fini,
alors $f$ est représentable (par des schémas).
\end{lemma}

\begin{proof}
Soient $T$ un schéma sur $S$ et $T \to Y$ une application. Il faut montrer que
l'espace algébrique $X \times_Y T$ est un schéma. Écrivons $X = \colim X_\lambda$
comme dans la définition
\ref{formal-spaces-definition-affine-formal-algebraic-space}.
Soit $W \subset X \times_Y T$
un sous-espace ouvert quasi-compact. La restriction de la projection
$X \times_Y T \to X$ à $W$ se factorise par $X_\lambda$ pour un certain $\lambda$.
Alors
$$
W \to X_\lambda \times_S T
$$
est un monomorphisme (donc séparé) et est localement quasi-fini (car
$W \to X \times_Y T \to T$ est localement quasi-fini par l'hypothèse
sur $X \to Y$ ; voir Morphismes d'espaces algébriques, lemme
\ref{spaces-morphisms-lemma-permanence-quasi-finite}).
Par conséquent, $W$ est un schéma d'après
Morphismes d'espaces algébriques, Proposition
\ref{spaces-morphisms-proposition-locally-quasi-finite-separated-over-scheme}.
Ainsi, $X \times_Y T$ est un schéma d'après
Propriétés des espaces algébriques, lemme \ref{spaces-properties-lemma-subscheme}.
\end{proof}






\section{Espaces algébriques formels affines à indexation dénombrable}
\sectionmark{Espaces formels affines à indexation dénombrable}
\label{formal-spaces-section-countably-indexed}

\noindent
Ce sont les espaces algébriques formels affines du lemme suivant.

\begin{lemma}
\label{formal-spaces-lemma-countable-affine-formal-algebraic-space}
Soient $S$ un schéma et $X$ un espace algébrique formel affine sur $S$.
Les assertions suivantes sont équivalentes :
\begin{enumerate}
\item il existe un système $X_1 \to X_2 \to X_3 \to \ldots$
d'épaississements de schémas affines sur $S$ tel que $X = \colim X_n$,
\item il existe une écriture $X = \colim X_\lambda$ comme dans la
Définition \ref{formal-spaces-definition-affine-formal-algebraic-space}
telle que $\Lambda$ soit dénombrable.
\end{enumerate}
\end{lemma}

\begin{proof}
Ceci résulte de l'observation qu'un ensemble ordonné filtrant dénombrable
possède une partie cofinale isomorphe à $(\mathbf{N}, \geq)$.
Voir la démonstration de Algèbre, lemme \ref{algebra-lemma-ML-limit-nonempty}.
\end{proof}

\begin{definition}
\label{formal-spaces-definition-countable}
Soient $S$ un schéma et $X$ un espace algébrique formel affine sur $S$.
On dit que $X$ est {\it à indexation dénombrable} si les conditions équivalentes du
Lemme \ref{formal-spaces-lemma-countable-affine-formal-algebraic-space} sont satisfaites.
\end{definition}

\noindent
Dans la terminologie de \cite{BVGD}, cela s'exprime en disant que
$X$ est un $\aleph_0$-ind-schéma.

\begin{lemma}
\label{formal-spaces-lemma-implications-between-types}
Soit $X$ un espace algébrique formel affine sur un schéma $S$.
\begin{enumerate}
\item Si $X$ est noethérien, alors $X$ est adique*.
\item Si $X$ est adique*, alors $X$ est adique.
\item Si $X$ est adique, alors $X$ est faiblement adique.
\item Si $X$ est faiblement adique, alors $X$ est classique.
\item Si $X$ est faiblement adique, alors $X$ est à indexation dénombrable.
\item Si $X$ est à indexation dénombrable, alors $X$ est de McQuillan.
\end{enumerate}
\end{lemma}

\begin{proof}
Les assertions (1), (2), (3) et (4) résultent de l'écriture $X = \text{Spf}(A)$,
où $A$ est un anneau faiblement admissible (donc complet) linéairement
topologisé, et des implications entre les différents
types de tels anneaux discutées dans la section \ref{formal-spaces-section-weakly-adic}.

\medskip\noindent
Démonstration de (5). Par définition, il existe un anneau topologique faiblement adique $A$
tel que $X = \colim \Spec(A/I)$, où la colimite porte sur les idéaux
de définition de $A$. Comme $A$ est faiblement adique, il existe en particulier
un système fondamental dénombrable d'idéaux ouverts $I_\lambda$, voir la
Définition \ref{formal-spaces-definition-weakly-adic}.
Alors $X = \colim \Spec(A/I_n)$ par définition de $\text{Spf}(A)$.
Ainsi, $X$ est à indexation dénombrable.

\medskip\noindent
Démonstration de (6). Écrivons $X = \colim X_n$ pour un système
$X_1 \to X_2 \to X_3 \to \ldots$ d'épaississements de schémas
affines sur $S$. Alors
$$
A = \lim \Gamma(X_n, \mathcal{O}_{X_n})
$$
l'homomorphisme canonique vers chacun des $\Gamma(X_n, \mathcal{O}_{X_n})$ est surjectif, car les applications de transition
sont surjectives puisque les morphismes $X_n \to X_{n + 1}$ sont des immersions
fermées. Par conséquent, $X$ est de McQuillan.
\end{proof}

\begin{lemma}
\label{formal-spaces-lemma-countably-indexed}
Soient $S$ un schéma et $X$ un préfaisceau sur $(\Sch/S)_{fppf}$.
Les assertions suivantes sont équivalentes :
\begin{enumerate}
\item $X$ est un espace algébrique formel affine à indexation dénombrable,
\item $X = \text{Spf}(A)$, où $A$ est une $S$-algèbre topologique faiblement
admissible qui possède un système fondamental dénombrable de voisinages de $0$,
\item $X = \text{Spf}(A)$, où $A$ est une $S$-algèbre topologique faiblement
admissible qui possède un système fondamental
$A \supset I_1 \supset I_2 \supset I_3 \supset \ldots$
d'idéaux faibles de définition,
\item $X = \text{Spf}(A)$, où $A$ est une $S$-algèbre topologique complète
ayant pour système fondamental de voisinages ouverts de $0$ une
suite dénombrable $A \supset I_1 \supset I_2 \supset I_3 \supset \ldots$
d'idéaux tels que $I_n/I_{n + 1}$ soit localement nilpotent, et
\item $X = \text{Spf}(A)$, où $A = \lim B/J_n$, muni de la topologie limite,
et où $B \supset J_1 \supset J_2 \supset J_3 \supset \ldots$ est une
suite d'idéaux d'une $S$-algèbre $B$ telle que $J_n/J_{n + 1}$ soit
localement nilpotent.
\end{enumerate}
\end{lemma}

\begin{proof}
Supposons (1). D'après le Lemme \ref{formal-spaces-lemma-implications-between-types},
on peut écrire $X = \text{Spf}(A)$, où $A$ est une $S$-algèbre topologique
faiblement admissible. Pour toute présentation $X = \colim X_n$ comme dans la
partie (1) du Lemme \ref{formal-spaces-lemma-countable-affine-formal-algebraic-space},
on voit que $A = \lim A_n$, avec $X_n = \Spec(A_n)$ et
$A_n = A/I_n$ pour un idéal faible de définition $I_n \subset A$.
Ceci résulte de la dernière assertion du
Lemme \ref{formal-spaces-lemma-mcquillan-affine-formal-algebraic-space},
qui implique en outre que $\{I_n\}$ est un système fondamental
de voisinages ouverts de $0$. Nous avons donc une suite
$$
A \supset I_1 \supset I_2 \supset I_3 \supset \ldots
$$
d'idéaux faibles de définition telle que $A = \lim A/I_n$. On voit ainsi
que la condition (1) implique chacune des conditions (2) -- (5).

\medskip\noindent
Supposons (5). Notons d'abord que la topologie limite sur
$A = \lim B/J_n$ est complète et linéairement topologisée, voir
Compléments d'algèbre, section \ref{more-algebra-section-topological-ring}.
Si $f \in A$ s'envoie sur zéro dans $B/J_1$, une certaine puissance s'envoie sur zéro
dans $B/J_2$, puisque son image dans $J_1/J_2$ est nilpotente, puis une nouvelle
puissance s'envoie sur zéro dans $J_2/J_3$, et ainsi de suite. On voit ainsi que
l'idéal ouvert $\Ker(A \to B/J_1)$ est un idéal faible de définition.
Ainsi, $A$ est faiblement admissible. On voit donc que (5) implique (2).

\medskip\noindent
Il est clair que (4) est un cas particulier de (5), en prenant $B = A$.
Il est clair que (3) est un cas particulier de (2).

\medskip\noindent
Supposons que $A$ soit comme en (2). Soit $E_n$ un système fondamental dénombrable
de voisinages de $0$ dans $A$. Comme $A$ est un anneau topologique faiblement
admissible, on peut trouver des idéaux ouverts $I_n \subset E_n$.
On peut aussi choisir un idéal faible de définition $J \subset A$.
Les $J \cap I_n$ forment alors un système fondamental d'idéaux faibles de définition
de $A$, et l'on obtient
$X = \text{Spf}(A) = \colim \Spec(A/(J \cap I_n))$,
ce qui montre que $X$ est un espace algébrique formel affine à indexation dénombrable.
\end{proof}

\begin{lemma}
\label{formal-spaces-lemma-characterize-noetherian-affine}
Soient $S$ un schéma et $X$ un espace algébrique formel affine.
Les assertions suivantes sont équivalentes :
\begin{enumerate}
\item $X$ est noethérien,
\item $X$ est adique* et, pour toute immersion fermée $T \to X$ où $T$ est un schéma,
$T$ est noethérien,
\item $X$ est adique* et, pour un certain choix de $X = \colim X_\lambda$ comme dans la
Définition \ref{formal-spaces-definition-affine-formal-algebraic-space},
les schémas $X_\lambda$ sont noethériens, et
\item $X$ est faiblement adique et, pour un certain choix de $X = \colim X_\lambda$
comme dans la Définition \ref{formal-spaces-definition-affine-formal-algebraic-space},
les schémas $X_\lambda$ sont noethériens.
\end{enumerate}
\end{lemma}

\begin{proof}
Supposons $X$ noethérien. Alors $X = \text{Spf}(A)$, où
$A$ est un anneau adique noethérien.
Soit $T \to X$ une immersion fermée, où $T$ est un schéma.
D'après le Lemme \ref{formal-spaces-lemma-mcquillan-affine-formal-algebraic-space},
on voit que $T$ est affine et que $T \to \Spec(A)$ est une immersion
fermée. Comme $A$ est noethérien, on voit que $T$ est noethérien.
On voit ainsi que (1) $\Rightarrow$ (2).

\medskip\noindent
Les implications (2) $\Rightarrow$ (3) et (2) $\Rightarrow$ (4) sont
immédiates (voir le Lemme \ref{formal-spaces-lemma-implications-between-types}).

\medskip\noindent
Pour démontrer (3) $\Rightarrow$ (1), écrivons $X = \text{Spf}(A)$
pour un anneau adique $A$ muni d'un idéal de définition $I$ de type fini.
On sait également que les anneaux $A/I_\lambda$ sont noethériens
pour un certain système fondamental d'idéaux ouverts $I_\lambda$.
Comme $I$ est ouvert, on peut trouver un $\lambda$ tel que $I_\lambda \subset I$.
Alors $A/I$ est noethérien, et l'on conclut que $A$ est noethérien par
Algèbre, lemme \ref{algebra-lemma-completion-Noetherian}.

\medskip\noindent
Pour démontrer (4) $\Rightarrow$ (3), écrivons $X = \text{Spf}(A)$
pour un anneau faiblement adique $A$. Alors $A$ est admissible et
possède un idéal de définition $I$, tandis que l'adhérence $I_2$ de $I^2$
est ouverte, voir le Lemme \ref{formal-spaces-lemma-weakly-adic}.
On sait également que les anneaux $A/I_\lambda$ sont noethériens
pour un certain système fondamental d'idéaux ouverts $I_\lambda$.
Choisissons un $\lambda$ tel que $I_\lambda \subset I_2$.
Alors $A/I_2$ est noethérien en tant que quotient de $A/I_\lambda$.
Par conséquent, $I/I_2$ est un $A$-module de type fini.
Ainsi, $A$ est un anneau adique possédant un idéal de définition
de type fini par le Lemme \ref{formal-spaces-lemma-weakly-adic-finite-generation}.
Donc $X$ est adique* et (3) est satisfaite.
\end{proof}





\section{Espaces algébriques formels}
\label{formal-spaces-section-formal-algebraic-spaces}

\noindent
Nous dérogeons ici à notre habitude d'introduire les notions nouvelles d'abord
pour les anneaux, puis pour les schémas et enfin pour les espaces algébriques, en
introduisant les espaces algébriques formels sans avoir d'abord introduit les
schémas formels. L'idée générale sera qu'un espace algébrique formel
est un faisceau pour la topologie fppf qui, localement pour la topologie étale, est un
schéma formel affine au sens de \cite{BVGD}.
On trouvera des résultats connexes dans \cite{Yasuda}.

\medskip\noindent
Dans la définition d'un espace algébrique formel, nous allons
emprunter de la terminologie à
Amorçage, sections
\ref{bootstrap-section-morphism-representable-by-spaces} et
\ref{bootstrap-section-representable-by-spaces-properties}.

\begin{definition}
\label{formal-spaces-definition-formal-algebraic-space}
Soit $S$ un schéma. On dit qu'un faisceau $X$ sur $(\Sch/S)_{fppf}$ est un
{\it espace algébrique formel} s'il existe une famille de morphismes
$\{X_i \to X\}_{i \in I}$ de faisceaux telle que
\begin{enumerate}
\item $X_i$ soit un espace algébrique formel affine,
\item $X_i \to X$ soit représentable par des espaces algébriques et étale,
\item $\coprod X_i \to X$ soit surjectif en tant que morphisme de faisceaux,
\end{enumerate}
et que $X$ satisfasse une condition ensembliste
(voir la Remarque \ref{formal-spaces-remark-set-theoretic}). Un
{\it morphisme d'espaces algébriques formels}
sur $S$ est un morphisme de faisceaux.
\end{definition}

\noindent
Discussion. Contrôle de cohérence : un espace algébrique formel affine est
un espace algébrique formel. Dans la situation de la définition,
les morphismes $X_i \to X$ sont représentables (par des schémas), voir le
Lemme \ref{formal-spaces-lemma-presentation-representable}.
D'après Amorçage, lemme
\ref{bootstrap-lemma-surjective-flat-locally-finite-presentation},
au lieu de demander que $\coprod X_i \to X$ soit
surjectif en tant que morphisme de faisceaux, on pourrait demander qu'il soit
surjectif (ce qui a un sens puisqu'il est représentable).

\medskip\noindent
Notre notion d'espace algébrique formel est {\bf très générale}.
En fait, même les espaces algébriques formels affines définis ci-dessus
sont des objets très pathologiques.

\begin{lemma}
\label{formal-spaces-lemma-diagonal-formal-algebraic-space}
Soit $S$ un schéma. Si $X$ est un espace algébrique formel sur
$S$, alors le morphisme diagonal $\Delta : X \to X \times_S X$
est représentable, est un monomorphisme, est localement quasi-fini,
est localement de type fini et est séparé.
\end{lemma}

\begin{proof}
Supposons donnés $U \to X$ et $V \to X$, avec $U, V$ des schémas sur $S$.
Alors $U \times_X V$ est un faisceau. Choisissons $\{X_i \to X\}$ comme dans la
Définition \ref{formal-spaces-definition-formal-algebraic-space}.
Pour tout $i$, le morphisme
$$
(U \times_X X_i) \times_{X_i} (V \times_X X_i)
= (U \times_X V) \times_X X_i \to U \times_X V
$$
est représentable et étale, car c'est un changement de base de $X_i \to X$,
et sa source est un schéma (utiliser les
Lemmes \ref{formal-spaces-lemma-diagonal-affine-formal-algebraic-space} et
\ref{formal-spaces-lemma-presentation-representable}). Ces morphismes forment une famille surjective ;
ainsi, $U \times_X V$ est un espace algébrique par
Amorçage, théorème \ref{bootstrap-theorem-final-bootstrap}.
Le morphisme $U \times_X V \to U \times_S V$ est un monomorphisme.
Il est également localement quasi-fini, car, en le précomposant par
le morphisme affiché ci-dessus, on obtient la composée
$$
(U \times_X X_i) \times_{X_i} (V \times_X X_i)
\to (U \times_X X_i) \times_S (V \times_X X_i)
\to U \times_S V
$$
qui est localement quasi-finie en tant que composée d'une immersion
fermée (Lemme \ref{formal-spaces-lemma-diagonal-affine-formal-algebraic-space})
et d'un morphisme étale, voir
Descente et espaces algébriques, lemme
\ref{spaces-descent-lemma-locally-quasi-finite-etale-local-source}.
On conclut donc que $U \times_X V$ est un schéma par
Morphismes d'espaces algébriques, Proposition
\ref{spaces-morphisms-proposition-locally-quasi-finite-separated-over-scheme}.
Ainsi, $\Delta$ est représentable, voir
Espaces, lemme \ref{spaces-lemma-representable-diagonal}.

\medskip\noindent
En fait, puisque nous avons montré ci-dessus que les morphismes de schémas
$U \times_X V \to U \times_S V$ sont toujours des monomorphismes et sont
localement quasi-finis, on conclut que $\Delta : X \to X \times_S X$
est un monomorphisme et est localement quasi-fini, voir
Espaces, lemme \ref{spaces-lemma-transformation-diagonal-properties}.
On peut ensuite appliquer le principe de
Espaces, lemme
\ref{spaces-lemma-representable-transformations-property-implication}
pour voir que $\Delta$ est séparé et localement de type fini.
En effet, un monomorphisme de schémas est séparé
(Schémas, lemme \ref{schemes-lemma-monomorphism-separated})
et un morphisme de schémas localement quasi-fini est
localement de type fini
(ceci résulte de la définition donnée dans
Morphismes, section \ref{morphisms-section-quasi-finite}).
\end{proof}

\begin{lemma}
\label{formal-spaces-lemma-space-to-formal-space}
Soit $S$ un schéma. Soit $f : X \to Y$ un morphisme d'un
espace algébrique sur $S$ vers un espace algébrique formel sur $S$.
Alors $f$ est représentable par des espaces algébriques.
\end{lemma}

\begin{proof}
Soit $Z \to Y$ un morphisme, où $Z$ est un schéma sur $S$.
Il faut montrer que $X \times_Y Z$ est un espace algébrique.
Choisissons un schéma $U$ et un morphisme étale surjectif $U \to X$.
Alors $U \times_Y Z \to X \times_Y Z$ est représentable, surjectif et étale
(Espaces, lemme
\ref{spaces-lemma-base-change-representable-transformations-property}),
et $U \times_Y Z$ est un schéma d'après le
Lemme \ref{formal-spaces-lemma-diagonal-formal-algebraic-space}.
Le résultat en découle par
Amorçage, théorème \ref{bootstrap-theorem-final-bootstrap}.
\end{proof}

\begin{remark}
\label{formal-spaces-remark-compare-with-formal-schemes}
À des questions ensemblistes près, la catégorie des schémas formels au sens d'EGA
(voir la section \ref{formal-spaces-section-formal-schemes-EGA}) est équivalente à une sous-catégorie
pleine de la catégorie des espaces algébriques formels. Pour l'expliquer,
supposons que notre schéma de base soit $\Spec(\mathbf{Z})$. D'après le
Lemme \ref{formal-spaces-lemma-formal-scheme-sheaf-fppf}, le foncteur des points
$h_\mathfrak X$ associé à un schéma formel $\mathfrak X$ est un faisceau
pour la topologie fppf. D'après le Lemme \ref{formal-spaces-lemma-fully-faithful},
le foncteur $\mathfrak X \mapsto h_\mathfrak X$ est un plongement pleinement fidèle
de la catégorie des schémas formels dans la catégorie des
faisceaux fppf. Étant donné un schéma formel $\mathfrak X$, choisissons un recouvrement ouvert
$\mathfrak X = \bigcup \mathfrak X_i$, où les $\mathfrak X_i$ sont des
schémas formels affines. Alors $h_{\mathfrak X_i}$
est un espace algébrique formel affine d'après la
Remarque \ref{formal-spaces-remark-compare-with-affine-formal-schemes}.
Les morphismes $h_{\mathfrak X_i} \to h_\mathfrak X$ sont représentables
et sont des immersions ouvertes. Ainsi, $\{h_{\mathfrak X_i} \to h_\mathfrak X\}$
est une famille comme dans la Définition \ref{formal-spaces-definition-formal-algebraic-space},
et l'on voit que $h_\mathfrak X$ est un espace algébrique formel.
\end{remark}

\begin{remark}
\label{formal-spaces-remark-set-theoretic}
Soient $S$ un schéma et $(\Sch/S)_{fppf}$ un grand site fppf comme dans
Topologies, définition \ref{topologies-definition-big-small-fppf}.
Comme condition ensembliste sur $X$ dans les
Définitions \ref{formal-spaces-definition-affine-formal-algebraic-space} et
\ref{formal-spaces-definition-formal-algebraic-space}, nous prenons la suivante :
il existe des objets $U, R$ de $(\Sch/S)_{fppf}$, un
morphisme $U \to X$ qui soit une surjection de faisceaux fppf, et
un morphisme $R \to U \times_X U$ qui soit une surjection de faisceaux fppf.
Autrement dit, nous demandons que notre faisceau soit le conoyau de
deux morphismes entre faisceaux représentables.
Voici quelques observations qui impliquent que cette notion se comporte
de façon raisonnable :
\begin{enumerate}
\item Supposons $X = \colim_{\lambda \in \Lambda} X_\lambda$
et que le système satisfasse les conditions (1) et (2) de la
Définition \ref{formal-spaces-definition-affine-formal-algebraic-space}. Alors
$U = \coprod_{\lambda \in \Lambda} X_\lambda \to X$ est une surjection
de faisceaux fppf. De plus, $U \times_X U$ est un sous-schéma fermé
de $U \times_S U$ d'après le Lemme \ref{formal-spaces-lemma-diagonal-affine-formal-algebraic-space}.
Ainsi, si $U$ est représentable par un objet de $(\Sch/S)_{fppf}$,
alors $U \times_S U$ l'est aussi (voir Ensembles, lemme \ref{sets-lemma-what-is-in-it}),
et la condition ensembliste est satisfaite. C'est toujours le cas
si $\Lambda$ est dénombrable, voir Ensembles, lemme \ref{sets-lemma-what-is-in-it}.
\item Contrôle de cohérence. Soit $\{X_i \to X\}_{i \in I}$ comme dans la
Définition \ref{formal-spaces-definition-formal-algebraic-space}
(avec la condition ensembliste formulée ci-dessus),
et supposons que chaque $X_i$ soit en fait un schéma affine.
Alors $X$ est un espace algébrique. En effet, si l'on choisit un
grand site fppf $(\Sch'/S)_{fppf}$ plus grand tel que $U' = \coprod X_i$
et $R' = \coprod X_i \times_X X_j$ soient représentables par des objets
de ce site, alors $X' = U'/R'$ sera un objet de la catégorie
des espaces algébriques pour ce choix. Une application de
Espaces, lemme \ref{spaces-lemma-fully-faithful} montre alors que
$X$ est un espace algébrique pour $(\Sch/S)_{fppf}$.
\item Soit $\{X_i \to X\}_{i \in I}$ une famille de morphismes de faisceaux
satisfaisant les conditions (1), (2), (3) de la
Définition \ref{formal-spaces-definition-formal-algebraic-space}.
Pour chaque $i$, on peut choisir $U_i \in \Ob((\Sch/S)_{fppf})$
et $U_i \to X_i$, qui soit une surjection de faisceaux.
Ainsi, si $I$ n'est pas trop grand (s'il est par exemple dénombrable),
$U = \coprod U_i \to X$ est une surjection de faisceaux et
$U$ est représentable par un objet de $(\Sch/S)_{fppf}$.
Pour obtenir un $R \in \Ob((\Sch/S)_{fppf})$ muni d'une surjection vers $U \times_X U$,
il suffit de supposer que la diagonale $\Delta : X \to X \times_S X$ n'est pas
trop pathologique ; c'est par exemple toujours possible si la diagonale de $X$ est
quasi-compacte, c'est-à-dire si $X$ est quasi-séparé.
\end{enumerate}
\end{remark}







\section{La réduction}
\label{formal-spaces-section-reduction}

\noindent
Tout espace algébrique formel possède un espace algébrique réduit
sous-jacent, comme le montre le lemme suivant.

\begin{lemma}
\label{formal-spaces-lemma-reduction-formal-algebraic-space}
Soit $S$ un schéma. Soit $X$ un espace algébrique formel sur $S$.
Il existe un espace algébrique réduit $X_{red}$ et un morphisme représentable
$X_{red} \to X$ qui est un épaississement. Tout morphisme $U \to X$,
où $U$ est un espace algébrique réduit, se factorise de manière unique par $X_{red}$.
\end{lemma}

\begin{proof}
Supposons d'abord que $X$ soit un espace algébrique formel affine.
Écrivons $X = \colim X_\lambda$ comme dans la
Définition \ref{formal-spaces-definition-affine-formal-algebraic-space}.
Puisque les morphismes de transition sont des épaississements, les schémas
affines $X_\lambda$ ont tous des réductions isomorphes $X_{red}$.
Le morphisme $X_{red} \to X$ est représentable et est un épaississement
par le Lemme \ref{formal-spaces-lemma-covering-by-thickenings} et par le fait que
les composés d'épaississements sont des épaississements. Nous omettons la
vérification de la propriété
universelle (utiliser Schémas, définition
\ref{schemes-definition-reduced-induced-scheme},
Schémas, lemme \ref{schemes-lemma-map-into-reduction},
Propriétés des espaces algébriques, définition
\ref{spaces-properties-definition-reduced-induced-space}, et
Propriétés des espaces algébriques, lemme \ref{spaces-properties-lemma-map-into-reduction}).

\medskip\noindent
Soient $X$ et $\{X_i \to X\}_{i \in I}$ comme dans la
Définition \ref{formal-spaces-definition-formal-algebraic-space}.
Pour tout $i$, soit $X_{i, red} \to X_i$ la réduction construite
ci-dessus. Pour $i, j \in I$, la projection
$X_{i, red} \times_X X_j \to X_{i, red}$ est un morphisme étale (par hypothèse)
de schémas (par le Lemme \ref{formal-spaces-lemma-presentation-representable}).
Ainsi $X_{i, red} \times_X X_j$ est réduit (voir
Descente, lemme \ref{descent-lemma-reduced-local-smooth}).
La projection $X_{i, red} \times_X X_j \to X_j$ se factorise donc
par $X_{j, red}$ en vertu de la propriété universelle. Nous en déduisons que
$$
R_{ij} = X_{i, red} \times_X X_j = X_{i, red} \times_X X_{j, red} =
X_i \times_X X_{j, red}
$$
car les morphismes $X_{i, red} \to X_i$ sont des injections de faisceaux.
Posons $U = \coprod X_{i, red}$, posons
$R = \coprod R_{ij}$, et notons $s, t : R \to U$ les deux
projections. En tant que faisceau, $R = U \times_X U$, et $s$ et $t$
sont étales. Alors $(t, s) : R \to U$ définit une relation d'équivalence
étale d'après les observations précédentes. Ainsi $X_{red} = U/R$ est un
espace algébrique par Espaces, théorème \ref{spaces-theorem-presentation}.
Par construction, le diagramme
$$
\xymatrix{
\coprod X_{i, red} \ar[r] \ar[d] & \coprod X_i \ar[d] \\
X_{red} \ar[r] & X
}
$$
est cartésien. Puisque la flèche verticale de droite est étale surjective
et que la flèche horizontale supérieure est représentable et est un épaississement,
nous concluons que $X_{red} \to X$ est représentable par
Amorçage, lemme \ref{bootstrap-lemma-after-fppf-sep-lqf}
(pour vérifier les hypothèses du lemme, utiliser qu'un morphisme étale
surjectif est surjectif, plat et localement de présentation
finie, et que les épaississements sont séparés et localement quasi-finis).
Nous pouvons ensuite appliquer Espaces, lemme
\ref{spaces-lemma-descent-representable-transformations-property}
pour conclure que $X_{red} \to X$ est un épaississement
(utiliser qu'être un épaississement équivaut à être
une immersion fermée surjective).

\medskip\noindent
Enfin, supposons que $U \to X$ soit un morphisme où
$U$ est un espace algébrique réduit sur $S$. Alors chaque $X_i \times_X U$
est étale sur $U$, donc réduit (d'après notre définition des
espaces algébriques réduits dans Propriétés des espaces algébriques, section
\ref{spaces-properties-section-types-properties}).
Alors $X_i \times_X U \to X_i$ se factorise par $X_{i, red}$.
Ainsi $U \to X$ se factorise par $X_{red}$, puisque
$\{X_i \times_X U \to U\}$ est un recouvrement étale.
\end{proof}

\begin{example}
\label{formal-spaces-example-reduction-affine-formal-spectrum}
Soit $A$ un anneau topologique faiblement admissible. Dans ce cas, on a
$$
\text{Spf}(A)_{red} = \Spec(A/\mathfrak a)
$$
où $\mathfrak a \subset A$ est l'idéal des éléments topologiquement
nilpotents. En effet, $\mathfrak a$ est un idéal radical
(Lemme \ref{formal-spaces-lemma-topologically-nilpotent})
qui est ouvert puisque $A$ est faiblement admissible.
\end{example}

\begin{lemma}
\label{formal-spaces-lemma-reduction-smooth}
Soit $S$ un schéma. Soit $f : X \to Y$ un morphisme d'espaces
algébriques formels sur $S$ qui est représentable par des
espaces algébriques et lisse (par exemple étale).
Alors $X_{red} = X \times_Y Y_{red}$.
\end{lemma}

\begin{proof}
(Le cas étale résulte directement de la construction de
l'espace algébrique réduit sous-jacent dans la démonstration du
Lemme \ref{formal-spaces-lemma-reduction-formal-algebraic-space}.)
Supposons $f$ lisse. Observons que $X \times_Y Y_{red} \to Y_{red}$
est un morphisme lisse d'espaces algébriques. Ainsi $X \times_Y Y_{red}$
est un espace algébrique réduit par Descente et espaces algébriques, lemme
\ref{spaces-descent-lemma-reduced-local-smooth}.
La propriété universelle de la réduction montre alors que le morphisme canonique
$X_{red} \to X \times_Y Y_{red}$ est un isomorphisme.
\end{proof}

\begin{lemma}
\label{formal-spaces-lemma-reduction-surjective}
Soit $S$ un schéma. Soit $f : X \to Y$ un morphisme d'espaces
algébriques formels sur $S$ qui est représentable par des
espaces algébriques. Alors $f$ est surjectif au sens de
Amorçage, définition \ref{bootstrap-definition-property-transformation}
si et seulement si $f_{red} : X_{red} \to Y_{red}$ est un
morphisme surjectif d'espaces algébriques.
\end{lemma}

\begin{proof}
Omis.
\end{proof}






\section{Limites inductives d'espaces algébriques le long d'épaississements}
\label{formal-spaces-section-global-colimits}

\noindent
Un type particulier d'espace algébrique formel est celui qui peut s'écrire
globalement comme une limite inductive cofiltrante d'espaces algébriques le long
d'épaississements, ainsi que dans le lemme suivant. Nous verrons plus loin
(à la section \ref{formal-spaces-section-quasi-compact-quasi-separated})
que tout espace algébrique formel quasi-compact et quasi-séparé
est une telle limite inductive globale.

\begin{lemma}
\label{formal-spaces-lemma-colimit-is-formal}
Soit $S$ un schéma. Supposons donnés un ensemble ordonné filtrant
$\Lambda$ et un système d'espaces algébriques $(X_\lambda, f_{\lambda \mu})$
sur $\Lambda$, où chaque $f_{\lambda \mu} : X_\lambda \to X_\mu$ est un
épaississement. Alors $X = \colim_{\lambda \in \Lambda} X_\lambda$
est un espace algébrique formel sur $S$.
\end{lemma}

\begin{proof}
Puisque nous prenons la limite inductive dans la catégorie des faisceaux fppf,
nous voyons que $X$ est un faisceau. Choisissons et fixons $\lambda \in \Lambda$. Choisissons un
recouvrement étale $\{X_{i, \lambda} \to X_\lambda\}$ où $X_i$ est un schéma
affine sur $S$, voir Propriétés des espaces algébriques, lemme
\ref{spaces-properties-lemma-cover-by-union-affines}.
Pour chaque $\mu \geq \lambda$, il existe un diagramme cartésien
$$
\xymatrix{
X_{i, \lambda} \ar[r] \ar[d] & X_{i, \mu} \ar[d] \\
X_\lambda \ar[r] & X_\mu
}
$$
dont les flèches verticales sont étales, voir
Compléments sur les morphismes d'espaces algébriques, théorème
\ref{spaces-more-morphisms-theorem-topological-invariance}
(on utilise aussi qu'un épaississement est une immersion fermée surjective qui
satisfait aux conditions du théorème). De plus, ces diagrammes sont
uniques à isomorphisme unique près, et donc
$X_{i, \mu} = X_\mu \times_{X_{\mu'}} X_{i, \mu'}$ pour
$\mu' \geq \mu$. Les morphismes $X_{i, \mu} \to X_{i, \mu'}$
sont des épaississements comme changements de base d'épaississements. Chaque $X_{i, \mu}$
est un schéma affine par Limits of Espaces, Proposition
\ref{spaces-limits-proposition-affine} et par le fait que
$X_{i, \lambda}$ est affine.
Posons $X_i = \colim_{\mu \geq \lambda} X_{i, \mu}$. Alors $X_i$ est
un espace algébrique formel affine. Le morphisme $X_i \to X$
est étale car, étant donnés un schéma affine $U$ et un morphisme $U \to X$,
celui-ci se factorise par $X_\mu$ pour un certain $\mu \geq \lambda$ (détails omis).
Nous voyons ainsi que $X$ est un espace algébrique formel.
\end{proof}

\noindent
Soit $S$ un schéma. Soit $X$ un espace algébrique formel sur $S$.
Comment démontrer ou vérifier que $X$ est une limite inductive globale comme dans le
Lemme \ref{formal-spaces-lemma-colimit-is-formal} ? Pour ce faire, nous cherchons des morphismes
$i : Z \to X$ où $Z$ est un espace algébrique sur $S$ et où $i$ est
surjectif et une immersion fermée ; autrement dit, $i$ est un épaississement.
Cela a un sens puisque $i$ est représentable par des espaces algébriques
(Lemme \ref{formal-spaces-lemma-space-to-formal-space}) et que nous pouvons appliquer
Amorçage, définition \ref{bootstrap-definition-property-transformation}
comme précédemment.

\begin{example}
\label{formal-spaces-example-david-hansen}
Soit $(A, \mathfrak m, \kappa)$ un anneau de valuation qui est
complet pour la topologie $(\pi)$-adique pour un certain $\pi \in \mathfrak m$ non nul.
Supposons aussi que $\mathfrak m$ ne soit pas de type fini.
Un exemple est $A = \mathcal{O}_{\mathbf{C}_p}$ et $\pi = p$,
où $\mathcal{O}_{\mathbf{C}_p}$ est l'anneau des entiers
du corps des nombres complexes $p$-adiques $\mathbf{C}_p$
(qui est le complété de la clôture algébrique de
$\mathbf{Q}_p$). Un autre exemple est
$$
A =
\left\{
\sum\nolimits_{\alpha \in \mathbf{Q},\ \alpha \geq 0} a_\alpha t^\alpha
\middle|
\begin{matrix}
a_\alpha \in \kappa \text{ et, pour tout }n\text{, il n'y a qu'un} \\
\text{nombre fini de coefficients non nuls }a_\alpha
\text{ tels que }\alpha \leq n
\end{matrix}
\right\}
$$
et $\pi = t$. Alors $X = \text{Spf}(A)$ est un espace algébrique formel
affine, et $\Spec(\kappa) \to X$ est un épaississement qui correspond
à l'idéal faible de définition $\mathfrak m \subset A$,
lequel n'est toutefois pas un idéal de définition.
\end{example}

\begin{remark}[Idéaux faibles de définition]
\label{formal-spaces-remark-weak-ideals-of-definition}
Soit $\mathfrak X$ un schéma formel au sens de McQuillan, voir
Remarque \ref{formal-spaces-remark-mcquillan}. Un {\it idéal faible de définition}
de $\mathfrak X$ est un faisceau d'idéaux
$\mathcal{I} \subset \mathcal{O}_\mathfrak X$ tel que,
pour tout sous-schéma formel ouvert affine $\mathfrak U \subset \mathfrak X$,
l'idéal
$\mathcal{I}(\mathfrak U) \subset \mathcal{O}_\mathfrak X(\mathfrak U)$
soit un idéal faible de définition de l'anneau topologique faiblement admissible
$\mathcal{O}_\mathfrak X(\mathfrak U)$.
Il suffit de vérifier la condition sur les membres d'un recouvrement ouvert affine.
Il existe une correspondance biunivoque
$$
\{\text{idéaux faibles de définition de }\mathfrak X\}
\leftrightarrow
\{\text{épaississements }i : Z \to h_\mathfrak X\text{ comme ci-dessus}\}
$$
Cette correspondance associe à $\mathcal{I}$ le schéma
$Z = (\mathfrak X, \mathcal{O}_\mathfrak X/\mathcal{I})$
muni du morphisme évident vers $\mathfrak X$.
Un {\it système fondamental d'idéaux faibles de définition}
est une famille d'idéaux faibles de définition
$\mathcal{I}_\lambda$ telle que, sur tout sous-schéma formel ouvert
affine $\mathfrak U \subset \mathfrak X$, les
idéaux
$$
I_\lambda = \mathcal{I}_\lambda(\mathfrak U) \subset
A = \Gamma(\mathfrak U, \mathcal{O}_\mathfrak X)
$$
forment un système fondamental d'idéaux faibles de définition de
l'anneau topologique faiblement admissible $A$. Il suffit de le vérifier
sur les membres d'un recouvrement ouvert affine. Nous concluons que
l'espace algébrique formel $h_\mathfrak X$ associé au
schéma formel de McQuillan $\mathfrak X$ est une limite inductive de schémas comme
dans le Lemme \ref{formal-spaces-lemma-colimit-is-formal} si et seulement s'il
existe un système fondamental d'idéaux faibles de définition
de $\mathfrak X$.
\end{remark}

\begin{remark}[Idéaux de définition]
\label{formal-spaces-remark-ideals-of-definition}
Soit $\mathfrak X$ un schéma formel \`a la EGA.
Un {\it idéal de définition} de $\mathfrak X$ est un faisceau d'idéaux
$\mathcal{I} \subset \mathcal{O}_\mathfrak X$ tel que,
pour tout sous-schéma formel ouvert affine $\mathfrak U \subset \mathfrak X$,
l'idéal
$\mathcal{I}(\mathfrak U) \subset \mathcal{O}_\mathfrak X(\mathfrak U)$
soit un idéal de définition de l'anneau topologique admissible
$\mathcal{O}_\mathfrak X(\mathfrak U)$.
Il suffit de vérifier la condition sur les membres d'un recouvrement ouvert affine.
Nous n'obtenons {\bf pas} la même correspondance entre idéaux de définition
et épaississements $Z \to h_\mathfrak X$ que dans la
Remarque \ref{formal-spaces-remark-weak-ideals-of-definition} ; un exemple
est donné dans l'Exemple \ref{formal-spaces-example-david-hansen}.
Un {\it système fondamental d'idéaux de définition}
est une famille d'idéaux de définition
$\mathcal{I}_\lambda$ telle que, sur tout sous-schéma formel ouvert
affine $\mathfrak U \subset \mathfrak X$, les
idéaux
$$
I_\lambda = \mathcal{I}_\lambda(\mathfrak U) \subset
A = \Gamma(\mathfrak U, \mathcal{O}_\mathfrak X)
$$
forment un système fondamental d'idéaux de définition de
l'anneau topologique admissible $A$. Il suffit de le vérifier
sur les membres d'un recouvrement ouvert affine. Supposons que $\mathfrak X$
soit quasi-compact et que $\{\mathcal{I}_\lambda\}_{\lambda \in \Lambda}$
soit un système fondamental d'idéaux faibles de définition.
Si $A$ est un anneau topologique admissible, alors tous les
idéaux ouverts suffisamment petits sont des idéaux de définition
(en effet, tout idéal ouvert contenu dans un idéal de définition
est un idéal de définition). Puisqu'il suffit donc de vérifier
la propriété sur les membres, en nombre fini, d'un recouvrement ouvert affine,
nous voyons que $\mathcal{I}_\lambda$ est un idéal de définition
pour $\lambda$ assez grand. D'après la discussion de la
Remarque \ref{formal-spaces-remark-weak-ideals-of-definition}, nous concluons que
l'espace algébrique formel $h_\mathfrak X$ associé au
schéma formel quasi-compact $\mathfrak X$ \`a la EGA
est une limite inductive de schémas comme dans le Lemme \ref{formal-spaces-lemma-colimit-is-formal}
si et seulement s'il existe un système fondamental d'idéaux de définition
de $\mathfrak X$.
\end{remark}






\section{Complété le long d'une partie fermée}
\label{formal-spaces-section-completion}

\noindent
Notre notion d'espace algébrique formel se prête bien à la formation du
complété le long d'une partie fermée.

\begin{lemma}
\label{formal-spaces-lemma-completion-affine-is-affine-formal-algebraic-space}
Soit $S$ un schéma. Soit $X$ un schéma affine sur $S$.
Soit $T \subset |X|$ une partie fermée. Alors le foncteur
$$
(\Sch/S)_{fppf} \longrightarrow \textit{Ensembles},\quad
U \longmapsto \{f : U \to X \mid f(|U|) \subset T\}
$$
est un espace algébrique formel affine de McQuillan.
\end{lemma}

\begin{proof}
Écrivons $X = \Spec(A)$ ; la partie $T$ correspond à l'idéal radical $I \subset A$.
Soit $U = \Spec(B)$ un schéma affine sur $S$ et soit
$f : U \to X$ un élément de $F(U)$. Alors $f$ correspond à un
homomorphisme d'anneaux $\varphi : A \to B$ tel que tout idéal premier de $B$ contienne
$\varphi(I) B$. Ainsi, tout élément de $\varphi(I)$ est nilpotent dans $B$ ; voir
Algèbre, lemme \ref{algebra-lemma-Zariski-topology}.
En posant $J = \Ker(\varphi)$, nous concluons que $I/J$ est un idéal localement
nilpotent de $A/J$. De manière équivalente, $V(J) = V(I) = T$.
Autrement dit, le foncteur du lemme est égal à
$\colim \Spec(A/J)$, où la limite inductive porte sur la
famille des idéaux $J$ tels que $V(J) = T$.
Notre foncteur est donc un espace algébrique formel affine. Il est de McQuillan
(définition \ref{formal-spaces-definition-types-affine-formal-algebraic-space})
car les applications $A \to A/J$ sont surjectives
et, par conséquent, $A^\wedge = \lim A/J \to A/J$ est surjective ; voir le
lemme \ref{formal-spaces-lemma-mcquillan-affine-formal-algebraic-space}.
\end{proof}

\begin{lemma}
\label{formal-spaces-lemma-completion-is-formal-algebraic-space}
Soit $S$ un schéma. Soit $X$ un espace algébrique sur $S$.
Soit $T \subset |X|$ une partie fermée. Alors le foncteur
$$
(\Sch/S)_{fppf} \longrightarrow \textit{Ensembles},\quad
U \longmapsto \{f : U \to X \mid f(|U|) \subset T\}
$$
est un espace algébrique formel.
\end{lemma}

\begin{proof}
Notons $F$ le foncteur. Soit $\{U_i \to U\}$ un recouvrement fppf.
Alors $\coprod |U_i| \to |U|$ est surjectif. Puisque $X$ est un
faisceau fppf, il s'ensuit que $F$ est un faisceau fppf.

\medskip\noindent
Soit $\{g_i : X_i \to X\}$ un recouvrement étale tel que $X_i$ soit affine
pour tout $i$ ; voir Propriétés des espaces algébriques, lemme
\ref{spaces-properties-lemma-cover-by-union-affines}.
Les morphismes $F \times_X X_i \to F$ sont étales
(voir Espaces, lemme
\ref{spaces-lemma-base-change-representable-transformations-property})
et l'application $\coprod F \times_X X_i \to F$ est une surjection de faisceaux.
Il suffit donc de démontrer que $F \times_X X_i$ est un espace algébrique
formel affine. Un point à valeurs dans $U$ de $F \times_X X_i$ est un
morphisme $U \to X_i$ dont l'image est contenue dans la partie fermée
$g_i^{-1}(T) \subset |X_i|$. L'assertion résulte donc du
lemme \ref{formal-spaces-lemma-completion-affine-is-affine-formal-algebraic-space}.
\end{proof}

\begin{definition}
\label{formal-spaces-definition-completion}
Soit $S$ un schéma. Soit $X$ un espace algébrique sur $S$.
Soit $T \subset |X|$ une partie fermée. L'espace algébrique formel
du Lemme \ref{formal-spaces-lemma-completion-is-formal-algebraic-space}
est appelé le {\it complété de $X$ le long de $T$}.
\end{definition}

\noindent
Dans \cite[chapitre I, section 10.8]{EGA}, la notation $X_{/T}$
sert à désigner le complété, et nous l'emploierons aussi
à l'occasion. Soit $f : X \to X'$ un
morphisme d'espaces algébriques sur un schéma $S$. Supposons
que $T \subset |X|$ et $T' \subset |X'|$ soient des parties fermées
telles que $|f|(T) \subset T'$. Il est alors clair que
$f$ définit un morphisme d'espaces algébriques formels
$$
X_{/T} \longrightarrow X'_{/T'}
$$
entre les complétés.

\begin{lemma}
\label{formal-spaces-lemma-map-completions-representable}
Soit $S$ un schéma. Soit $f : X' \to X$ un morphisme
d'espaces algébriques sur $S$. Soit $T \subset |X|$
une partie fermée et posons $T' = |f|^{-1}(T) \subset |X'|$.
Alors
$$
\xymatrix{
X'_{/T'} \ar[r] \ar[d] & X' \ar[d]^f \\
X_{/T} \ar[r] & X
}
$$
est un diagramme cartésien de faisceaux. En particulier, le morphisme
$X'_{/T'} \to X_{/T}$ est représentable par des espaces algébriques.
\end{lemma}

\begin{proof}
En effet, supposons que $Y \to X$ soit un morphisme d'un schéma vers $X$ tel
que $|Y|$ s'envoie dans $T$. Alors $Y \times_X X' \to X$ est un morphisme
d'espaces algébriques tel que $|Y \times_X X'|$ s'envoie dans $T'$. Le
foncteur $Y \times_{X_{/T}} X'_{/T'}$ est donc représenté par $Y \times_X X'$,
ce qui démontre le lemme.
\end{proof}

\begin{lemma}
\label{formal-spaces-lemma-reduction-completion}
Soit $S$ un schéma. Soit $X$ un espace algébrique sur $S$.
Soit $T \subset |X|$ une partie fermée. La réduction $(X_{/T})_{red}$
du complété $X_{/T}$ de $X$ le long de $T$ est
le sous-espace fermé induit réduit $Z$ de $X$ correspondant à $T$.
\end{lemma}

\begin{proof}
Cela résulte du Lemme \ref{formal-spaces-lemma-reduction-formal-algebraic-space},
Propriétés des espaces algébriques, définition
\ref{spaces-properties-definition-reduced-induced-space}
(qui utilise Propriétés des espaces algébriques, lemme
\ref{spaces-properties-lemma-reduced-closed-subspace} pour construire $Z$),
et de la définition de $X_{/T}$ que
$Z$ et $(X_{/T})_{red}$ sont des espaces algébriques réduits
caractérisés par la même propriété d'applications :
un morphisme $g : Y \to X$ dont la source est un espace algébrique réduit
se factorise par l'un ou l'autre si et seulement si $|Y|$ s'envoie dans $T \subset |X|$.
\end{proof}

\begin{lemma}
\label{formal-spaces-lemma-affine-formal-completion-types}
Soit $S$ un schéma. Soit $X = \Spec(A)$ un schéma affine sur $S$.
Soit $T \subset X$ une partie fermée. Soit $X_{/T}$ le
complété formel de $X$ le long de $T$.
\begin{enumerate}
\item Si $X \setminus T$ est quasi-compact, c'est-à-dire si $T$ est constructible,
alors $X_{/T}$ est adique*.
\item Si $T = V(I)$ pour un idéal de type fini $I \subset A$,
alors $X_{/T} = \text{Spf}(A^\wedge)$, où $A^\wedge$ est le
complété $I$-adique de $A$.
\item Si $X$ est noethérien, alors $X_{/T}$ est noethérien.
\end{enumerate}
\end{lemma}

\begin{proof}
Si (1) est vérifié, Algèbre, lemme \ref{algebra-lemma-qc-open} fournit
un idéal $I \subset A$ comme en (2). Si (3) est vérifié, nous pouvons
également trouver un idéal $I \subset A$ comme en (2). De plus, les complétés
d'anneaux noethériens sont noethériens d'après
Algèbre, lemme \ref{algebra-lemma-completion-Noetherian-Noetherian}.
Il suffit donc de démontrer (2).

\medskip\noindent
Preuve de (2).
Soit $I = (f_1, \ldots, f_r) \subset A$ un idéal définissant $T$.
Si $Z = \Spec(B)$ est un schéma affine et si $g : Z \to X$ est
un morphisme tel que $g(Z) \subset T$ (au sens ensembliste), alors
$g^\sharp(f_i)$ est nilpotent dans $B$ pour tout $i$. Ainsi,
$I^n$ s'envoie sur zéro dans $B$ pour un certain $n$. Nous voyons donc que
$X_{/T} = \colim \Spec(A/I^n) = \text{Spf}(A^\wedge)$.
\end{proof}

\noindent
Le lemme suivant est dû à Ofer Gabber.

\begin{lemma}
\label{formal-spaces-lemma-completion-countably-indexed}
\begin{reference}
Courriel d'Ofer Gabber du 11 septembre 2014.
\end{reference}
Soit $S$ un schéma. Soit $X = \Spec(A)$ un schéma affine sur $S$.
Soit $T \subset X$ un sous-schéma fermé.
\begin{enumerate}
\item Si le complété formel $X_{/T}$ est à indexation dénombrable
et s'il existe une famille dénombrable $f_1, f_2, f_3, \ldots \in A$ telle que
$T = V(f_1, f_2, f_3, \ldots)$, alors $X_{/T}$ est adique*.
\item La conclusion de (1) est fausse si l'on omet l'hypothèse selon laquelle
$T$ peut être défini par une famille dénombrable de fonctions sur $X$.
\end{enumerate}
\end{lemma}

\begin{proof}
L'hypothèse que $X_{/T}$ est à indexation dénombrable signifie qu'il existe une
suite d'idéaux
$$
A \supset J_1 \supset J_2 \supset J_3 \supset \ldots
$$
vérifiant $V(J_n) = T$ et telle que, pour tout idéal $J \subset A$ avec $V(J) = T$,
il existe un $n$ tel que $J \supset J_n$.

\medskip\noindent
Pour construire un exemple pour (2), soit $\omega_1$ le premier ordinal non
dénombrable. Soit $k$ un corps et soit
$A$ la $k$-algèbre engendrée par les $x_\alpha$, $\alpha \in \omega_1$,
et les $y_{\alpha \beta}$, où $\alpha \in \beta \in \omega_1$,
soumis aux relations $x_\alpha = y_{\alpha \beta} x_\beta$.
Soit $T = V(x_\alpha)$. Soit $J_n = (x_\alpha^n)$.
Si $J \subset A$ est un idéal tel que
$V(J) = T$, alors $x_\alpha^{n_\alpha} \in J$ pour un certain $n_\alpha \geq 1$.
L'un des ensembles $\{\alpha \mid n_\alpha = n\}$ doit être non borné dans
$\omega_1$. Les relations impliquent alors que $J_n \subset J$.

\medskip\noindent
Pour établir (2), il suffit maintenant de montrer que $A^\wedge = \lim A/J_n$
n'est complet pour la topologie définie par aucun idéal de type fini.
Pour $\gamma \in \omega_1$, soit $A_\gamma$ le quotient de $A$
par l'idéal engendré par les $x_\alpha$, $\alpha \in \gamma$, et les
$y_{\alpha \beta}$, $\alpha \in \gamma$. Comme $A/J_1$ est réduit,
tout élément topologiquement nilpotent $f$ de $\lim A/J_n$ appartient à
$J_1^\wedge = \lim J_1/J_n$. Cela signifie que $f$ est une série infinie
ne faisant intervenir qu'un nombre dénombrable de générateurs. Ainsi, $f$ s'annule dans
$A_\gamma^\wedge = \lim A_\gamma/J_nA_\gamma$ pour un certain $\gamma$.
Notons que $A^\wedge \to A_\gamma^\wedge$ est continu et ouvert d'après le
lemme \ref{formal-spaces-lemma-ses}.
Si la topologie de $A^\wedge$ était $I$-adique pour un idéal de type fini
$I \subset A^\wedge$, alors $I$ s'enverrait sur zéro dans un certain
$A_\gamma^\wedge$. Cela signifierait que $A_\gamma^\wedge$ est discret,
ce qui n'est pas le cas, puisqu'il existe une application surjective, continue et ouverte
(d'après le Lemme \ref{formal-spaces-lemma-ses})
$A_\gamma^\wedge \to k[[t]]$ donnée par
$x_\alpha \mapsto t$, $y_{\alpha \beta} \mapsto 1$ lorsque
$\gamma = \alpha$ ou $\gamma \in \alpha$.

\medskip\noindent
Avant de démontrer (1), établissons d'abord l'assertion suivante : si $I \subset A^\wedge$ est
un idéal de type fini dont l'adhérence $\bar I$ est ouverte, alors $I = \bar I$.
Puisque $V(J_n^2) = T$, il existe un $m$ tel que $J_n^2 \supset J_m$.
En passant à une sous-suite, nous pouvons donc supposer que $J_n^2 \supset J_{n + 1}$ pour tout
$n$. Posons $J_n^\wedge = \lim_{k \geq n} J_n/J_k \subset A^\wedge$.
Comme l'adhérence $\bar I = \bigcap (I + J_n^\wedge)$
(Lemme \ref{formal-spaces-lemma-closed}) est ouverte, il existe un $m$ tel que
$I + J_n^\wedge \supset J_m^\wedge$ pour tout $n \geq m$. Fixons un tel $m$.
Nous avons
$$
J_{n - 1}^\wedge I + J_{n + 1}^\wedge \supset
J_{n - 1}^\wedge (I + J_{n + 1}^\wedge) \supset
J_{n - 1}^\wedge J_m^\wedge
$$
pour tout $n \geq m + 1$. En effet, la première inclusion est triviale et la
seconde a été établie ci-dessus. Puisque $J_{n - 1}J_m \supset J_{n - 1}^2 \supset J_n$,
ces inclusions montrent que l'image de $J_n$ dans $A^\wedge$
est contenue dans l'idéal $J_{n - 1}^\wedge I + J_{n + 1}^\wedge$.
Comme cet idéal est ouvert, nous concluons que
$$
J_{n - 1}^\wedge I + J_{n + 1}^\wedge \supset J_n^\wedge.
$$
Écrivons $I = (g_1, \ldots, g_t)$. Prenons $f \in J_{m + 1}^\wedge$.
À l'aide de la dernière inclusion affichée, valable pour tout $n \geq m + 1$,
nous pouvons écrire, par récurrence sur $c \geq 0$,
$$
f = \sum f_{i, c} g_i \mod J_{m + 1+ c}^\wedge
$$
avec $f_{i, c} \in J_m^\wedge$ et
$f_{i, c} \equiv f_{i, c - 1} \bmod J_{m + c}^\wedge$.
Il s'ensuit que $IJ_m^\wedge \supset J_{m + 1}^\wedge$.
En combinant ceci avec $I + J_{m + 1}^\wedge \supset J_m^\wedge$,
nous concluons que $I$ est ouvert.

\medskip\noindent
Preuve de (1). Supposons $T = V(f_1, f_2, f_3, \ldots)$.
Soit $I_m \subset A^\wedge$ l'idéal engendré par $f_1, \ldots, f_m$.
Distinguons deux cas.

\medskip\noindent
Cas I : pour un certain $m$, l'adhérence de $I_m$ est ouverte.
Alors $I_m$ est ouvert d'après le résultat du paragraphe précédent.
Pour tout $n$, nous avons $(J_n)^2 \supset J_{n+1}$ par construction ; par conséquent,
l'adhérence de $(J_n^\wedge)^2$ contient $J_{n+1}^\wedge$
et est donc ouverte. En prenant $n$ assez grand, il s'ensuit que l'adhérence
du produit de deux idéaux ouverts quelconques de $A^\wedge$ est ouverte.
Démontrons que $I_m^k$ est ouvert pour $k \ge 1$, par récurrence sur $k$.
Le cas $k = 1$ est précisément notre hypothèse sur $m$ dans le cas I.
Pour $k > 1$, supposons $I_m^{k - 1}$ ouvert. Alors
$I_m^k = I_m^{k - 1} \cdot I_m$ est le produit de deux idéaux ouverts
et son adhérence est donc ouverte. Mais comme $I_m^k$
est de type fini, il s'ensuit que $I_m^k$
est ouvert d'après le paragraphe précédent (appliqué à $I = I_m^k$) ;
nous pouvons donc poursuivre la récurrence sur $k$.
Comme tout élément de $I_m$ est topologiquement nilpotent, nous concluons
que $I_m$ est un idéal de définition, ce qui prouve que $A^\wedge$
est adique et possède un idéal de définition de type fini, c'est-à-dire que
$X_{/T}$ est adique*.

\medskip\noindent
Cas II. Pour tout $m$, l'adhérence $\bar I_m$ de $I_m$ n'est pas ouverte.
La topologie de $A^\wedge/\bar I_m$ n'est alors pas discrète. Cela signifie
que nous pouvons choisir $\phi(m) \geq m$ tel que
$$
\Im(J_{\phi(m)} \to A/(f_1, \ldots, f_m)) \not =
\Im(J_{\phi(m) + 1} \to A/(f_1, \ldots, f_m))
$$
Pour le voir, nous avons utilisé l'égalité
$A^\wedge/(\bar I_m + J_n^\wedge) = A/((f_1, \ldots, f_m) + J_n)$.
Choisissons des exposants $e_i > 0$ tels que $f_i^{e_i} \in J_{\phi(m) + 1}$
pour $0 < m < i$. Posons $J = (f_1^{e_1}, f_2^{e_2}, f_3^{e_3}, \ldots)$.
Alors $V(J) = T$. Nous affirmons que $J \not \supset J_n$ pour tout $n$,
ce qui est une contradiction et montre que le cas II ne se produit pas.
En effet, l'image de $J$ dans $A/(f_1, \ldots, f_m)$ est contenue
dans l'image de $J_{\phi(m) + 1}$, laquelle est strictement contenue dans
l'image de $J_m$.
\end{proof}
















\section{Produits fibrés}
\label{formal-spaces-section-fibre-products}

\noindent
section obligée sur les produits fibrés d'espaces algébriques formels.

\begin{lemma}
\label{formal-spaces-lemma-etale-covering-by-formal-algebraic-spaces}
Soit $S$ un schéma. Soit $\{X_i \to X\}_{i \in I}$ une famille d'applications
de faisceaux sur $(\Sch/S)_{fppf}$. Supposons (a) que $X_i$ soit un
espace algébrique formel sur $S$, (b) que $X_i \to X$ soit représentable
par des espaces algébriques et étale, et (c) que $\coprod X_i \to X$
soit une surjection de faisceaux. Alors $X$ est un espace algébrique formel
sur $S$.
\end{lemma}

\begin{proof}
Pour chaque $i$, choisissons $\{X_{ij} \to X_i\}_{j \in J_i}$ comme dans la
définition \ref{formal-spaces-definition-formal-algebraic-space}.
Alors $\{X_{ij} \to X\}_{i \in I, j \in J_i}$ est une famille
comme dans la Définition \ref{formal-spaces-definition-formal-algebraic-space}
pour $X$.
\end{proof}

\begin{lemma}
\label{formal-spaces-lemma-fibre-products-general}
Soit $S$ un schéma. Soient $X, Y$ des espaces algébriques formels sur $S$
et soit $Z$ un faisceau dont la diagonale est représentable par des
espaces algébriques. Soient $X \to Z$ et $Y \to Z$ des applications de faisceaux.
Alors $X \times_Z Y$ est un espace algébrique formel.
\end{lemma}

\begin{proof}
Choisissons $\{X_i \to X\}$ et $\{Y_j \to Y\}$ comme dans la
définition \ref{formal-spaces-definition-formal-algebraic-space}.
Alors $\{X_i \times_Z Y_j \to X \times_Z Y\}$ est une famille
d'applications représentables par des espaces algébriques et étales.
Le Lemme \ref{formal-spaces-lemma-etale-covering-by-formal-algebraic-spaces}
montre donc qu'il suffit d'établir que $X \times_Z Y$ est un espace algébrique
formel lorsque $X$ et $Y$ sont des espaces algébriques formels affines.

\medskip\noindent
Supposons que $X$ et $Y$ soient des espaces algébriques formels affines.
Écrivons $X = \colim X_\lambda$ et $Y = \colim Y_\mu$ comme
dans la Définition \ref{formal-spaces-definition-affine-formal-algebraic-space}.
Alors $X \times_Z Y = \colim X_\lambda \times_Z Y_\mu$.
Chaque $X_\lambda \times_Z Y_\mu$ est un espace algébrique.
Pour $\lambda \leq \lambda'$ et $\mu \leq \mu'$, le morphisme
$$
X_\lambda \times_Z Y_\mu \to
X_\lambda \times_Z Y_{\mu'} \to
X_{\lambda'} \times_Z Y_{\mu'}
$$
est un épaississement, car c'est un composé de changements de base d'épaississements.
Nous concluons donc en appliquant le Lemme \ref{formal-spaces-lemma-colimit-is-formal}.
\end{proof}

\begin{lemma}
\label{formal-spaces-lemma-fibre-products}
Soit $S$ un schéma. La catégorie des espaces algébriques formels sur $S$
possède des produits fibrés.
\end{lemma}

\begin{proof}
C'est un cas particulier du Lemme \ref{formal-spaces-lemma-fibre-products-general},
car les espaces algébriques formels ont des diagonales représentables ; voir le
Lemme \ref{formal-spaces-lemma-diagonal-formal-algebraic-space}.
\end{proof}

\begin{lemma}
\label{formal-spaces-lemma-reduction-fibre-products}
Soit $S$ un schéma. Soient $X \to Z$ et $Y \to Z$ des morphismes
d'espaces algébriques formels sur $S$. Alors
$(X \times_Z Y)_{red} = (X_{red} \times_{Z_{red}} Y_{red})_{red}$.
\end{lemma}

\begin{proof}
Cela résulte de la propriété universelle de la réduction
du Lemme \ref{formal-spaces-lemma-reduction-formal-algebraic-space}.
\end{proof}

\noindent
Nous avons déjà démontré le lemme suivant (sans savoir que
les produits fibrés existent).

\begin{lemma}
\label{formal-spaces-lemma-diagonal-morphism-formal-algebraic-spaces}
Soit $S$ un schéma. Soit $f : X \to Y$ un morphisme d'espaces algébriques formels
sur $S$. Le morphisme diagonal $\Delta : X \to X \times_Y X$
est représentable (par des schémas), est un monomorphisme, est localement quasi-fini,
localement de type fini et séparé.
\end{lemma}

\begin{proof}
Soit $T$ un schéma et soit $T \to X \times_Y X$ un morphisme.
Alors
$$
T \times_{(X \times_Y X)} X = T \times_{(X \times_S X)} X
$$
Le résultat découle donc immédiatement du
Lemme \ref{formal-spaces-lemma-diagonal-formal-algebraic-space}.
\end{proof}






\section{Axiomes de séparation pour les espaces algébriques formels}
\label{formal-spaces-section-separation}

\noindent
Cette section porte sur les conditions de séparation « absolues » des espaces
algébriques formels. Nous étudierons plus loin les conditions de séparation
des morphismes d'espaces algébriques formels.

\begin{lemma}
\label{formal-spaces-lemma-characterize-quasi-separated}
Soit $S$ un schéma. Soit $X$ un espace algébrique formel sur $S$.
Les conditions suivantes sont équivalentes :
\begin{enumerate}
\item la réduction de $X$
(Lemme \ref{formal-spaces-lemma-reduction-formal-algebraic-space}) est un
espace algébrique quasi-séparé ;
\item pour $U \to X$, $V \to X$, où $U$, $V$ sont des schémas quasi-compacts,
le produit fibré $U \times_X V$ est quasi-compact ;
\item pour $U \to X$, $V \to X$, où $U$, $V$ sont affines,
le produit fibré $U \times_X V$ est quasi-compact.
\end{enumerate}
\end{lemma}

\begin{proof}
Observons que $U \times_X V$ est un schéma d'après le
Lemme \ref{formal-spaces-lemma-diagonal-formal-algebraic-space}.
Soient $U_{red}, V_{red}, X_{red}$ les réductions de $U, V, X$.
Alors
$$
U_{red} \times_{X_{red}} V_{red} = U_{red} \times_X V_{red} \to U \times_X V
$$
est un épaississement de schémas. L'équivalence de (1) et (2) en découle
clairement, compte tenu du lemme analogue pour les espaces algébriques ; voir
Propriétés des espaces algébriques, lemme
\ref{spaces-properties-lemma-characterize-quasi-separated}.
Nous omettons la démonstration de l'équivalence de (2) et (3).
\end{proof}

\begin{lemma}
\label{formal-spaces-lemma-characterize-separated}
Soit $S$ un schéma. Soit $X$ un espace algébrique formel sur $S$.
Les conditions suivantes sont équivalentes :
\begin{enumerate}
\item la réduction de $X$
(Lemme \ref{formal-spaces-lemma-reduction-formal-algebraic-space}) est un espace
algébrique séparé ;
\item pour $U \to X$, $V \to X$, où $U$, $V$ sont affines,
le produit fibré $U \times_X V$ est affine et l'application
$$
\mathcal{O}(U) \otimes_\mathbf{Z} \mathcal{O}(V)
\longrightarrow
\mathcal{O}(U \times_X V)
$$
est surjective.
\end{enumerate}
\end{lemma}

\begin{proof}
Si (2) est vérifiée, alors $X_{red}$ est un espace algébrique séparé : on applique
Propriétés des espaces algébriques, lemme
\ref{spaces-properties-lemma-characterize-quasi-separated}
aux morphismes $U \to X_{red}$ et $V \to X_{red}$,
avec $U, V$ affines, et on utilise que $U \times_{X_{red}} V = U \times_X V$.

\medskip\noindent
Supposons (1). Soient $U \to X$ et $V \to X$ comme en (2).
Observons que $U \times_X V$ est un schéma d'après le
Lemme \ref{formal-spaces-lemma-diagonal-formal-algebraic-space}.
Soient $U_{red}, V_{red}, X_{red}$ les réductions de $U, V, X$.
Alors
$$
U_{red} \times_{X_{red}} V_{red} = U_{red} \times_X V_{red} \to U \times_X V
$$
est un épaississement de schémas. Il s'ensuit que
$(U \times_X V)_{red} = (U_{red} \times_{X_{red}} V_{red})_{red}$.
En particulier, nous voyons que $(U \times_X V)_{red}$ est un schéma affine
et que l'application
$$
\mathcal{O}(U) \otimes_\mathbf{Z} \mathcal{O}(V)
\longrightarrow
\mathcal{O}((U \times_X V)_{red})
$$
est surjective ; voir Propriétés des espaces algébriques, lemme
\ref{spaces-properties-lemma-characterize-quasi-separated}.
Alors $U \times_X V$ est affine d'après
Limits of Espaces, Proposition \ref{spaces-limits-proposition-affine}.
D'autre part, le morphisme $U \times_X V \to U \times V$
de schémas affines est le composé
$$
U \times_X V = X \times_{(X \times_S X)} (U \times_S V)
\to U \times_S V \to U \times V
$$
Le premier morphisme est un monomorphisme localement de type fini
(lemme \ref{formal-spaces-lemma-diagonal-formal-algebraic-space}).
Le second morphisme est une immersion
(Schémas, lemme \ref{schemes-lemma-fibre-product-after-map}).
Le composé est donc un monomorphisme localement de type fini.
D'autre part, le composé est entier, car l'application entre les schémas
affines réduits sous-jacents est une immersion fermée
d'après ce qui précède, donc universellement fermée (utiliser
Morphismes, lemme \ref{morphisms-lemma-integral-universally-closed}).
Ainsi, l'homomorphisme d'anneaux
$$
\mathcal{O}(U) \otimes_\mathbf{Z} \mathcal{O}(V)
\longrightarrow
\mathcal{O}(U \times_X V)
$$
est un épimorphisme entier de type fini,
donc fini, donc surjectif (utiliser
Morphismes, lemme \ref{morphisms-lemma-finite-integral}
et
Algèbre, lemme \ref{algebra-lemma-finite-epimorphism-surjective}).
\end{proof}

\begin{definition}
\label{formal-spaces-definition-separated}
Soit $S$ un schéma. Soit $X$ un espace algébrique formel sur $S$.
On dit que
\begin{enumerate}
\item $X$ est {\it quasi-séparé} si les conditions équivalentes du
Lemme \ref{formal-spaces-lemma-characterize-quasi-separated} sont satisfaites ;
\item $X$ est {\it séparé} si les conditions équivalentes du
Lemme \ref{formal-spaces-lemma-characterize-separated} sont satisfaites.
\end{enumerate}
\end{definition}

\noindent
Le lemme suivant implique en particulier que le produit tensoriel
complété d'anneaux topologiques faiblement admissibles est
un anneau topologique faiblement admissible.

\begin{lemma}
\label{formal-spaces-lemma-fibre-product-affines-over-separated}
Soit $S$ un schéma. Soient $X \to Z$ et $Y \to Z$ des morphismes
d'espaces algébriques formels sur $S$. Supposons $Z$ séparé.
\begin{enumerate}
\item Si $X$ et $Y$ sont des espaces algébriques formels affines, alors
$X \times_Z Y$ l'est aussi ;
\item si $X$ et $Y$ sont des espaces algébriques formels affines de McQuillan, alors
$X \times_Z Y$ l'est aussi ;
\item si $X$, $Y$ et $Z$ sont des espaces algébriques formels affines de McQuillan
correspondant aux $S$-algèbres topologiques faiblement admissibles
$A$, $B$ et $C$, alors $X \times_Z Y$ correspond à
$A \widehat{\otimes}_C B$.
\end{enumerate}
\end{lemma}

\begin{proof}
Écrivons $X = \colim X_\lambda$ et $Y = \colim Y_\mu$ comme
dans la Définition \ref{formal-spaces-definition-affine-formal-algebraic-space}.
Alors $X \times_Z Y = \colim X_\lambda \times_Z Y_\mu$.
Puisque $Z$ est séparé, les produits fibrés sont affines ;
nous voyons donc que (1) est vérifiée. Supposons que $X$ et $Y$ correspondent
aux $S$-algèbres topologiques faiblement admissibles $A$ et $B$,
et que $X_\lambda = \Spec(A/I_\lambda)$ et $Y_\mu = \Spec(B/J_\mu)$.
Alors
$$
X_\lambda \times_Z Y_\mu \to
X_\lambda \times Y_\mu \to \Spec(A \otimes B)
$$
est une immersion fermée. L'une des conditions du
Lemme \ref{formal-spaces-lemma-mcquillan-affine-formal-algebraic-space}
est donc satisfaite et nous concluons que $X \times_Z Y$ est de McQuillan.
Si de plus $Z$ est de McQuillan et correspond à $C$, alors
$$
X_\lambda \times_Z Y_\mu = \Spec(A/I_\lambda \otimes_C B/J_\mu)
$$
nous voyons donc que l'anneau topologique faiblement admissible
correspondant à $X \times_Z Y$ est le produit tensoriel complété
(voir Définition \ref{formal-spaces-definition-toplogy-tensor-product}).
\end{proof}

\begin{lemma}
\label{formal-spaces-lemma-separated-from-separated}
Soit $S$ un schéma. Soit $X$ un espace algébrique formel sur $S$.
Soit $U \to X$ un morphisme, où $U$ est un espace algébrique
séparé sur $S$. Alors $U \to X$ est séparé.
\end{lemma}

\begin{proof}
L'énoncé a un sens, car $U \to X$ est représentable par des
espaces algébriques (Lemme \ref{formal-spaces-lemma-space-to-formal-space}).
Soit $T$ un schéma et soit $T \to X$ un morphisme. Nous devons montrer
que $U \times_X T \to T$ est séparé. Puisque $U \times_X T \to U \times_S T$
est un monomorphisme, il suffit de montrer que $U \times_S T \to T$
est séparé. Cela résulte du fait qu'il s'agit du changement de base de $U \to S$.
Nous avons utilisé dans l'argument précédent :
Morphismes d'espaces algébriques, lemmes
\ref{spaces-morphisms-lemma-base-change-separated},
\ref{spaces-morphisms-lemma-composition-separated},
\ref{spaces-morphisms-lemma-monomorphism-separated}, et
\ref{spaces-morphisms-lemma-separated-implies-morphism-separated}.
\end{proof}




\section{Espaces algébriques formels quasi-compacts}
\label{formal-spaces-section-quasi-compact}

\noindent
Voici la caractérisation des espaces algébriques formels
quasi-compacts.

\begin{lemma}
\label{formal-spaces-lemma-characterize-quasi-compact}
Soit $S$ un schéma. Soit $X$ un espace algébrique formel sur $S$.
Les conditions suivantes sont équivalentes :
\begin{enumerate}
\item la réduction de $X$
(Lemme \ref{formal-spaces-lemma-reduction-formal-algebraic-space}) est un espace
algébrique quasi-compact ;
\item on peut trouver $\{X_i \to X\}_{i \in I}$ comme dans la
Définition \ref{formal-spaces-definition-formal-algebraic-space}, avec $I$ fini ;
\item il existe un morphisme $Y \to X$ représentable par des espaces
algébriques, étale et surjectif, tel que
$Y$ soit un espace algébrique formel affine.
\end{enumerate}
\end{lemma}

\begin{proof}
Omis.
\end{proof}

\begin{definition}
\label{formal-spaces-definition-quasi-compact}
Soit $S$ un schéma. Soit $X$ un espace algébrique formel sur $S$.
On dit que $X$ est {\it quasi-compact} si les conditions équivalentes du
Lemme \ref{formal-spaces-lemma-characterize-quasi-compact} sont satisfaites.
\end{definition}

\begin{lemma}
\label{formal-spaces-lemma-characterize-quasi-compact-morphism}
Soit $S$ un schéma. Soit $f : X \to Y$ un morphisme d'espaces algébriques
formels sur $S$. Les conditions suivantes sont équivalentes :
\begin{enumerate}
\item l'application induite $f_{red} : X_{red} \to Y_{red}$ entre les réductions
(Lemme \ref{formal-spaces-lemma-reduction-formal-algebraic-space}) est un morphisme
quasi-compact d'espaces algébriques ;
\item pour tout schéma quasi-compact $T$ et tout morphisme $T \to Y$,
le produit fibré $X \times_Y T$ est un espace algébrique formel
quasi-compact ;
\item pour tout schéma affine $T$ et tout morphisme $T \to Y$,
le produit fibré $X \times_Y T$ est un espace algébrique formel
quasi-compact ;
\item il existe un recouvrement $\{Y_j \to Y\}$ comme dans la
Définition \ref{formal-spaces-definition-formal-algebraic-space}
tel que chaque $X \times_Y Y_j$ soit un espace algébrique formel quasi-compact.
\end{enumerate}
\end{lemma}

\begin{proof}
Omis.
\end{proof}

\begin{definition}
\label{formal-spaces-definition-quasi-compact-morphism}
Soit $S$ un schéma. Soit $f : X \to Y$ un morphisme
d'espaces algébriques formels sur $S$.
On dit que $f$ est {\it quasi-compact} si les conditions équivalentes du
Lemme \ref{formal-spaces-lemma-characterize-quasi-compact-morphism} sont satisfaites.
\end{definition}

\noindent
Cette notion coïncide avec celle qui existe déjà lorsque le morphisme
est représentable par des espaces algébriques (et, en particulier, lorsqu'il est
représentable).

\begin{lemma}
\label{formal-spaces-lemma-quasi-compact-representable}
Soit $S$ un schéma. Soit $f : X \to Y$ un morphisme d'espaces algébriques
formels sur $S$, représentable par des espaces algébriques.
Alors $f$ est quasi-compact au sens de la
Définition \ref{formal-spaces-definition-quasi-compact-morphism}
si et seulement si $f$ est quasi-compact au sens de
Amorçage, définition \ref{bootstrap-definition-property-transformation}.
\end{lemma}

\begin{proof}
Cela résulte immédiatement des définitions et du
Lemme \ref{formal-spaces-lemma-characterize-quasi-compact-morphism}.
\end{proof}





\section{Espaces algébriques formels quasi-compacts et quasi-séparés}
\label{formal-spaces-section-quasi-compact-quasi-separated}

\noindent
Le résultat suivant est dû à Yasuda ; voir
\cite[Proposition 3.32]{Yasuda}.

\begin{lemma}
\label{formal-spaces-lemma-structure-quasi-compact-quasi-separated}
\begin{reference}
\cite[Proposition 3.32]{Yasuda}
\end{reference}
Soit $S$ un schéma. Soit $X$ un espace algébrique formel quasi-compact et
quasi-séparé sur $S$. Alors $X = \colim X_\lambda$
pour un système d'espaces algébriques $(X_\lambda, f_{\lambda \mu})$
indexé par un ensemble ordonné filtrant $\Lambda$, où chaque
$f_{\lambda \mu} : X_\lambda \to X_\mu$ est un épaississement.
\end{lemma}

\begin{proof}
D'après le Lemme \ref{formal-spaces-lemma-characterize-quasi-compact}, nous pouvons choisir un
espace algébrique formel affine $Y$ et un morphisme étale surjectif
représentable $Y \to X$. Écrivons $Y = \colim Y_\lambda$ comme dans la
Définition \ref{formal-spaces-definition-affine-formal-algebraic-space}.

\medskip\noindent
Prenons $\lambda \in \Lambda$. Alors $Y_\lambda \times_X Y$ est un schéma d'après le
Lemme \ref{formal-spaces-lemma-presentation-representable}. La réduction
(lemme \ref{formal-spaces-lemma-reduction-formal-algebraic-space})
de $Y_\lambda \times_X Y$ est égale à la réduction de
$Y_{red} \times_{X_{red}} Y_{red}$, laquelle est quasi-compact puisque $X$
est quasi-séparé et que $Y_{red}$ est affine.
Par conséquent, $Y_\lambda \times_X Y$ est un schéma quasi-compact.
Il existe donc un $\mu \geq \lambda$ tel que
$\text{pr}_2 : Y_\lambda \times_X Y \to Y$ se factorise
par $Y_\mu$ ; voir le Lemme \ref{formal-spaces-lemma-factor-through-thickening}.
Soit $Z_\lambda$ l'image schématique du morphisme
$\text{pr}_2 : Y_\lambda \times_X Y \to Y_\mu$.
Celle-ci est indépendante du choix de $\mu$, et nous pouvons donc
considérer $Z_\lambda \subset Y$ comme l'image schématique
du morphisme $\text{pr}_2 : Y_\lambda \times_X Y \to Y$.
Observons que $Z_\lambda$ est aussi égale à l'image schématique
du morphisme $\text{pr}_1 : Y \times_X Y_\lambda \to Y$, puisque
celui-ci est isomorphe au morphisme ayant servi à définir $Z_\lambda$.
Nous affirmons que $Z_\lambda \times_X Y = Y \times_X Z_\lambda$ comme sous-foncteurs
de $Y \times_X Y$. En effet, puisque $Y \to X$ est étale, nous voyons que
$Z_\lambda \times_X Y$ est l'image schématique du morphisme
$$
\text{pr}_{13} = \text{pr}_1 \times \text{id}_Y :
Y \times_X Y_\lambda \times_X Y \longrightarrow Y \times_X Y
$$
d'après Morphismes d'espaces algébriques, lemme
\ref{spaces-morphisms-lemma-quasi-compact-scheme-theoretic-image}.
De même, $Y \times_X Z_\lambda$ est l'image schématique
du morphisme
$$
\text{pr}_{13} = \text{id}_Y \times \text{pr}_2 : 
Y \times_X Y_\lambda \times_X Y \longrightarrow Y \times_X Y
$$
L'affirmation en résulte. Alors
$R_\lambda = Z_\lambda \times_X Y = Y \times_X Z_\lambda$,
muni du morphisme $R_\lambda \to Z_\lambda \times_S Z_\lambda$,
définit une relation d'équivalence étale. Nous obtenons ainsi un espace
algébrique $X_\lambda = Z_\lambda/R_\lambda$. Par construction, le diagramme
$$
\xymatrix{
Z_\lambda \ar[r] \ar[d] & Y \ar[d] \\
X_\lambda \ar[r] & X
}
$$
est cartésien (car $X$ est le conoyau des deux projections
$R = Y \times_X Y \to Y$, car $Z_\lambda \subset Y$ est $R$-invariant,
et car $R_\lambda$ est la restriction de $R$ à $Z_\lambda$).
Ainsi, $X_\lambda \to X$ est représentable et est une immersion fermée ; voir
Espaces, lemme
\ref{spaces-lemma-morphism-sheaves-with-P-effective-descent-etale}.
D'autre part, puisque $Y_\lambda \subset Z_\lambda$, nous voyons que
$(X_\lambda)_{red} = X_{red}$ ; autrement dit, $X_\lambda \to X$
est un épaississement. Enfin, nous affirmons que
$$
X = \colim X_\lambda
$$
Nous avons $Y \times_X X_\lambda = Z_\lambda \supset Y_\lambda$. Tout
morphisme $T \to X$, où $T$ est un schéma sur $S$, se relève localement pour la topologie étale
en un morphisme vers $Y$, lequel se relève localement pour la topologie étale en un morphisme
vers un certain $Y_\lambda$. Ainsi, $T \to X$ se relève localement pour la topologie étale sur
$T$ en un morphisme vers $X_\lambda$. Ceci achève la démonstration.
\end{proof}

\begin{remark}
\label{formal-spaces-remark-structure-quasi-compact-quasi-separated}
Dans cette remarque, nous traduisons l'énoncé et la démonstration du
Lemme \ref{formal-spaces-lemma-structure-quasi-compact-quasi-separated}
dans le langage des schémas formels \`a la EGA.
La Remarque \ref{formal-spaces-remark-ideals-of-definition} montre
que le lemme se traduit comme suit :
\begin{itemize}
\item[(*)] Tout schéma formel quasi-compact et quasi-séparé
possède un système fondamental d'idéaux de définition.
\end{itemize}
Pour le démontrer, nous utilisons d'abord le principe de récurrence (reformulé pour les
schémas formels quasi-compacts et quasi-séparés) de
Cohomologie des schémas, lemme \ref{coherent-lemma-induction-principle}
pour nous ramener à la situation suivante :
$\mathfrak X = \mathfrak U \cup \mathfrak V$
où $\mathfrak U$, $\mathfrak V$ sont des sous-schémas formels ouverts,
où $\mathfrak V$ est affine, et où le résultat est vrai pour $\mathfrak U$,
$\mathfrak V$ et $\mathfrak U \cap \mathfrak V$. Choisissons des idéaux
de définition quelconques $\mathcal{I} \subset \mathcal{O}_\mathfrak U$
et $\mathcal{J} \subset \mathcal{O}_\mathfrak V$.
Par hypothèse, nous disposons d'un système fondamental d'idéaux
de définition sur $\mathfrak U$ et $\mathfrak V$ ; comme
$\mathfrak U \cap \mathfrak V$ est quasi-compact, nous pouvons trouver
des idéaux de définition $\mathcal{I}' \subset \mathcal{I}$
et $\mathcal{J}' \subset \mathcal{J}$
tels que
$$
\mathcal{I}'|_{\mathfrak U \cap \mathfrak V} \subset
\mathcal{J}|_{\mathfrak U \cap \mathfrak V}
\quad\text{et}\quad
\mathcal{J}'|_{\mathfrak U \cap \mathfrak V} \subset
\mathcal{I}|_{\mathfrak U \cap \mathfrak V}
$$
Soient $U \to U' \to \mathfrak U$ et $V \to V' \to \mathfrak V$ les
immersions fermées déterminées par les idéaux de définition
$\mathcal{I}' \subset \mathcal{I} \subset \mathcal{O}_\mathfrak U$
et
$\mathcal{J}' \subset \mathcal{J} \subset \mathcal{O}_\mathfrak V$.
Notons $\mathfrak U \cap V$ le sous-schéma ouvert de $V$ dont
l'espace topologique sous-jacent est celui de $\mathfrak U \cap \mathfrak V$.
Par notre choix de $\mathcal{I}'$, il existe une factorisation
$\mathfrak U \cap V \to U'$.
Nous définissons de même $U \cap \mathfrak V$, qui se factorise par $V'$.
Considérons alors
$$
Z_U = \text{image schématique de }
U \amalg (\mathfrak U \cap V) \longrightarrow U'
$$
et
$$
Z_V = \text{image schématique de }
(U \cap \mathfrak V) \amalg V \longrightarrow V'
$$
Comme la formation de l'image schématique d'un morphisme quasi-compact
commute à la restriction aux ouverts (Morphismes, lemme
\ref{morphisms-lemma-quasi-compact-scheme-theoretic-image})
on voit que $Z_U \cap \mathfrak V = \mathfrak U \cap Z_V$.
Ainsi, $Z_U$ et $Z_V$ se recollent en un schéma $Z$ muni
d'un morphisme $Z \to \mathfrak X$. Comme dans la discussion de la
remarque \ref{formal-spaces-remark-weak-ideals-of-definition}
nous voyons que $Z$ correspond à un idéal faible
de définition $\mathcal{I}_Z \subset \mathcal{O}_\mathfrak X$.
Notons que $Z_U \subset U'$ et que
$Z_V \subset V'$. La famille de tous les $\mathcal{I}_Z$
ainsi construits forme donc un système fondamental d'idéaux faibles
de définition. Une sous-famille fournit dès lors un système fondamental d'idéaux
de définition ; voir la Remarque \ref{formal-spaces-remark-ideals-of-definition}.
\end{remark}

\begin{lemma}
\label{formal-spaces-lemma-characterize-affine}
Soit $S$ un schéma. Soit $X$ un espace algébrique formel sur $S$.
Alors $X$ est un espace algébrique formel affine si et seulement si
sa réduction $X_{red}$ (Lemme \ref{formal-spaces-lemma-reduction-formal-algebraic-space})
est affine.
\end{lemma}

\begin{proof}
Par les Lemmes \ref{formal-spaces-lemma-characterize-quasi-separated} et
\ref{formal-spaces-lemma-characterize-quasi-compact} ainsi que par les
Définitions \ref{formal-spaces-definition-separated} et \ref{formal-spaces-definition-quasi-compact},
nous voyons que $X$ est quasi-compact et quasi-séparé.
D'après le lemme de Yasuda (Lemme \ref{formal-spaces-lemma-structure-quasi-compact-quasi-separated}),
nous pouvons écrire $X = \colim X_\lambda$ comme une limite inductive filtrante
d'épaississements d'espaces algébriques. Cependant, chaque $X_\lambda$
est affine d'après Limits of Espaces, lemme
\ref{spaces-limits-lemma-reduction-scheme}
car $(X_\lambda)_{red} = X_{red}$.
Ainsi, $X$ est un espace algébrique formel affine par définition.
\end{proof}






\section{Morphismes représentables par des espaces algébriques}
\label{formal-spaces-section-representable}

\noindent
Soit $f : X \to Y$ un morphisme d'espaces algébriques formels qui
est représentable par des espaces algébriques. Pour ce type de morphismes,
nous disposons d'une théorie abondante, grâce au travail accompli
dans les chapitres consacrés aux espaces algébriques.

\begin{lemma}
\label{formal-spaces-lemma-composition-representable}
Le composé de morphismes représentables par des espaces algébriques est
représentable par des espaces algébriques. Il en va de même pour les morphismes représentables
(par des schémas).
\end{lemma}

\begin{proof}
Voir Amorçage, lemme \ref{bootstrap-lemma-composition-transformation}.
\end{proof}

\begin{lemma}
\label{formal-spaces-lemma-base-change-representable}
Tout changement de base d'un morphisme représentable par des espaces algébriques est
représentable par des espaces algébriques. Il en va de même pour les morphismes représentables
(par des schémas).
\end{lemma}

\begin{proof}
Voir Amorçage, lemme \ref{bootstrap-lemma-base-change-transformation}.
\end{proof}

\begin{lemma}
\label{formal-spaces-lemma-permanence-representable}
Soit $S$ un schéma. Soient $f : X \to Y$ et $g : Y \to Z$ des morphismes
d'espaces algébriques formels sur $S$. Si $g \circ f : X \to Z$ est représentable
par des espaces algébriques, alors $f : X \to Y$ est représentable par des espaces algébriques.
\end{lemma}

\begin{proof}
Notons que la diagonale de $Y \to Z$ est représentable d'après le
Lemme \ref{formal-spaces-lemma-diagonal-morphism-formal-algebraic-spaces}.
Ainsi, $X \to Y$ est représentable par des espaces algébriques d'après
Amorçage, lemme \ref{bootstrap-lemma-representable-by-spaces-permanence}.
\end{proof}

\noindent
La propriété d'être représentable par des espaces algébriques est locale sur la
source et sur le but.

\begin{lemma}
\label{formal-spaces-lemma-representable-by-algebraic-spaces-local}
Soit $S$ un schéma. Soit $f : X \to Y$ un morphisme d'espaces algébriques
formels sur $S$. Les conditions suivantes sont équivalentes :
\begin{enumerate}
\item le morphisme $f$ est représentable par des espaces algébriques ;
\item il existe un diagramme commutatif
$$
\xymatrix{
U \ar[d] \ar[r] & V \ar[d] \\
X \ar[r] & Y
}
$$
où $U$, $V$ sont des espaces algébriques formels, les flèches verticales sont
représentables par des espaces algébriques, $U \to X$
est étale surjectif et $U \to V$ est représentable par des espaces algébriques ;
\item pour tout diagramme commutatif
$$
\xymatrix{
U \ar[d] \ar[r] & V \ar[d] \\
X \ar[r] & Y
}
$$
où $U$, $V$ sont des espaces algébriques formels et les flèches verticales sont
représentables par des espaces algébriques, le morphisme $U \to V$ est
représentable par des espaces algébriques ;
\item il existe un recouvrement $\{Y_j \to Y\}$ comme dans la
Définition \ref{formal-spaces-definition-formal-algebraic-space}
et, pour chaque $j$, un recouvrement $\{X_{ji} \to Y_j \times_Y X\}$ comme dans la
Définition \ref{formal-spaces-definition-formal-algebraic-space}, tels que
$X_{ji} \to Y_j$ soit représentable par des espaces algébriques pour tous $j$ et $i$ ;
\item il existe un recouvrement $\{X_i \to X\}$ comme dans la
Définition \ref{formal-spaces-definition-formal-algebraic-space}
et, pour chaque $i$, une factorisation $X_i \to Y_i \to Y$, où $Y_i$
est un espace algébrique formel affine, où $Y_i \to Y$ est représentable
par des espaces algébriques, et telle que $X_i \to Y_i$ soit représentable par des espaces
algébriques ;
\item en ajouter ici.
\end{enumerate}
\end{lemma}

\begin{proof}
Il est clair que (1) implique (2), car on peut prendre $U = X$ et $V = Y$.
Réciproquement, (2) implique (1) d'après
Amorçage, lemme \ref{bootstrap-lemma-representable-by-spaces-cover}
appliqué à $U \to X \to Y$.

\medskip\noindent
Supposons (1) vraie et considérons un diagramme comme en (3).
Alors $U \to Y$ est représentable par des espaces algébriques
(en tant que composé $U \to X \to Y$ ; voir
Amorçage, lemme \ref{bootstrap-lemma-composition-transformation})
et se factorise par $V$. Ainsi, $U \to V$ est représentable par des
espaces algébriques d'après le Lemme \ref{formal-spaces-lemma-permanence-representable}.

\medskip\noindent
Il est clair que (3) implique (2). Les conditions (1) -- (3) sont donc équivalentes.

\medskip\noindent
Observons que la condition (4) a un sens, car le produit fibré
$Y_j \times_Y X$ est un espace algébrique formel d'après le
Lemme \ref{formal-spaces-lemma-fibre-products}.
Il est clair que (4) implique (5).

\medskip\noindent
Supposons $X_i \to Y_i \to Y$ comme en (5). Posons alors
$V = \coprod Y_i$ et $U = \coprod X_i$ ; on voit que
(5) implique (2).

\medskip\noindent
Enfin, supposons (1) -- (3) vraies.
Nous pouvons donc choisir un recouvrement quelconque $\{Y_j \to Y\}$ comme dans la
Définition \ref{formal-spaces-definition-formal-algebraic-space}
et, pour chaque $j$, un recouvrement quelconque $\{X_{ji} \to Y_j \times_Y X\}$ comme dans la
Définition \ref{formal-spaces-definition-formal-algebraic-space}.
Alors $X_{ij} \to Y_j$ est représentable par des espaces algébriques d'après (3),
et nous voyons que (4) est vraie. Ceci achève la démonstration.
\end{proof}

\begin{lemma}
\label{formal-spaces-lemma-algebraic-space-over-affine-formal}
Soit $S$ un schéma. Soit $Y$ un espace algébrique formel affine sur $S$.
Soit $f : X \to Y$ une application de faisceaux sur $(\Sch/S)_{fppf}$ qui est
représentable par des espaces algébriques. Alors $X$ est un espace
algébrique formel.
\end{lemma}

\begin{proof}
Écrivons $Y = \colim Y_\lambda$ comme dans la
Définition \ref{formal-spaces-definition-affine-formal-algebraic-space}.
Pour chaque $\lambda$, le produit fibré
$X \times_Y Y_\lambda$ est un espace algébrique.
Ainsi, $X = \colim X \times_Y Y_\lambda$ est un espace
algébrique formel d'après le Lemme \ref{formal-spaces-lemma-colimit-is-formal}.
\end{proof}

\begin{lemma}
\label{formal-spaces-lemma-representable-by-algebraic-spaces}
Soit $S$ un schéma. Soit $Y$ un espace algébrique formel sur $S$.
Soit $f : X \to Y$ une application de faisceaux sur $(\Sch/S)_{fppf}$ qui est
représentable par des espaces algébriques. Alors $X$ est un espace
algébrique formel.
\end{lemma}

\begin{proof}
Soit $\{Y_i \to Y\}$ comme dans la
Définition \ref{formal-spaces-definition-formal-algebraic-space}.
Alors $X \times_Y Y_i \to X$ est une famille de morphismes
représentables par des espaces algébriques, étales et formant une famille
surjective. Il suffit donc de montrer que
$X \times_Y Y_i$ est un espace algébrique formel ; voir le
Lemme \ref{formal-spaces-lemma-etale-covering-by-formal-algebraic-spaces}.
Cela résulte du Lemme \ref{formal-spaces-lemma-algebraic-space-over-affine-formal}.
\end{proof}

\begin{lemma}
\label{formal-spaces-lemma-affine-representable-by-algebraic-spaces}
Soit $S$ un schéma. Soit $f : X \to Y$ un morphisme
d'espaces algébriques formels affines qui est représentable par des
espaces algébriques. Alors $f$ est représentable (par des schémas) et affine.
\end{lemma}

\begin{proof}
Nous allons montrer que $f$ est affine ; il s'ensuivra alors que
$f$ est représentable et affine d'après
Morphismes d'espaces algébriques, lemme \ref{spaces-morphisms-lemma-affine-local}.
Écrivons $Y = \colim Y_\mu$ et $X = \colim X_\lambda$ comme dans la
Définition \ref{formal-spaces-definition-affine-formal-algebraic-space}.
Soit $T \to Y$ un morphisme, où $T$ est un schéma
sur $S$. Nous devons montrer que $X \times_Y T \to T$ est affine ; voir
Amorçage, définition
\ref{bootstrap-definition-property-transformation}.
Pour ce faire, nous pouvons supposer $T$ affine et devons démontrer
que $X \times_Y T$ est affine. Dans ce cas, $T \to Y$ se factorise
par $Y_\mu \to Y$ pour un certain $\mu$ ; voir le
Lemme \ref{formal-spaces-lemma-factor-through-thickening}.
Puisque $f$ est quasi-compact, nous voyons que $X \times_Y T$ est
quasi-compact (Lemme \ref{formal-spaces-lemma-characterize-quasi-compact-morphism}).
Ainsi, $X \times_Y T \to X$ se factorise par $X_\lambda$ pour un certain
$\lambda$. De même, $X_\lambda \to Y$ se factorise par $Y_\mu$
après avoir augmenté $\mu$. Alors
$X \times_Y T = X_\lambda \times_{Y_\mu} T$.
Nous concluons, car les produits fibrés de schémas affines sont affines.
\end{proof}

\begin{lemma}
\label{formal-spaces-lemma-representable-affine}
Soit $S$ un schéma. Soit $\varphi : A \to B$ une application continue
d'anneaux topologiques faiblement admissibles sur $S$. Les conditions suivantes
sont équivalentes :
\begin{enumerate}
\item $\text{Spf}(\varphi) : \text{Spf}(B) \to \text{Spf}(A)$
est représentable par des espaces algébriques ;
\item $\text{Spf}(\varphi) : \text{Spf}(B) \to \text{Spf}(A)$
est représentable (par des schémas) ;
\item $\varphi$ est tendu ; voir la Définition \ref{formal-spaces-definition-taut}.
\end{enumerate}
\end{lemma}

\begin{proof}
Les assertions (1) et (2) sont équivalentes d'après le
lemme \ref{formal-spaces-lemma-affine-representable-by-algebraic-spaces}.

\medskip\noindent
Supposons que les conditions équivalentes (1) et (2) soient vérifiées.
Si $I \subset A$ est un idéal faible de définition, alors
$\Spec(A/I) \to \text{Spf}(A)$ est représentable et est un épaississement
(cela ressort de la construction du spectre formel,
mais résulte aussi du
lemme \ref{formal-spaces-lemma-mcquillan-affine-formal-algebraic-space}).
Alors $\Spec(A/I) \times_{\text{Spf}(A)} \text{Spf}(B) \to \text{Spf}(B)$
est représentable et est un épaississement en tant que changement de base.
Ainsi, d'après le
lemme \ref{formal-spaces-lemma-mcquillan-affine-formal-algebraic-space}
il existe un idéal faible de définition $J(I) \subset B$ tel que
$\Spec(A/I) \times_{\text{Spf}(A)} \text{Spf}(B) = \Spec(B/J(I))$
comme sous-foncteurs de $\text{Spf}(B)$. Nous obtenons un diagramme cartésien
$$
\xymatrix{
\Spec(B/J(I)) \ar[d] \ar[r] & \Spec(A/I) \ar[d] \\
\text{Spf}(B) \ar[r] & \text{Spf}(A)
}
$$
Le Lemme \ref{formal-spaces-lemma-fibre-product-affines-over-separated}
montre que $B/J(I) = B \widehat{\otimes}_A A/I$.
Il s'ensuit que $J(I)$ est l'adhérence de l'idéal $\varphi(I)B$ ; voir le
lemme \ref{formal-spaces-lemma-closure-image-ideal}.
Puisque $\text{Spf}(A) = \colim \Spec(A/I)$, où $I$ est comme ci-dessus,
nous obtenons $\text{Spf}(B) = \colim \Spec(B/J(I))$.
Ainsi, les idéaux $J(I)$ forment un système fondamental d'idéaux faibles
de définition (voir le
lemme \ref{formal-spaces-lemma-mcquillan-affine-formal-algebraic-space}).
La condition (3) est donc vérifiée.

\medskip\noindent
Supposons (3) vérifiée. Nous allons essentiellement reprendre en sens inverse les
arguments du paragraphe précédent.
Soit $I \subset A$ un idéal faible de définition.
D'après le Lemme \ref{formal-spaces-lemma-fibre-product-affines-over-separated},
nous obtenons un diagramme cartésien
$$
\xymatrix{
\text{Spf}(B \widehat{\otimes}_A A/I) \ar[d] \ar[r] & \Spec(A/I) \ar[d] \\
\text{Spf}(B) \ar[r] & \text{Spf}(A)
}
$$
Si $J(I)$ est l'adhérence de $IB$, alors $J(I)$ est ouvert dans $B$
par le caractère tendu de $\varphi$. Ainsi, si $J$ est ouvert dans $B$ et $J \subset J(B)$,
alors $B/J \otimes_A A/I = B/(IB + J) = B/J(I)$, car
$J(I) = \bigcap_{J \subset B\text{ ouvert}} (IB + J)$ d'après le Lemme \ref{formal-spaces-lemma-closed}.
Par conséquent, la limite définissant le produit tensoriel complété se stabilise et donne
$B \widehat{\otimes}_A A/I = B/J(I)$.
Ainsi, $\text{Spf}(B \widehat{\otimes}_A A/I) = \Spec(B/J(I))$.
Ceci prouve que $\text{Spf}(B) \times_{\text{Spf}(A)} \Spec(A/I)$
est représentable pour tout idéal faible de définition $I \subset A$.
Puisque tout morphisme $T \to \text{Spf}(A)$ avec $T$ quasi-compact
se factorise par $\Spec(A/I)$ pour un certain idéal faible de définition $I$
(lemme \ref{formal-spaces-lemma-factor-through-thickening})
nous concluons que $\text{Spf}(\varphi)$ est représentable, c'est-à-dire que
(2) est vérifiée. Ceci achève la démonstration.
\end{proof}

\begin{lemma}
\label{formal-spaces-lemma-property-goes-up-affine-morphism}
Soit $S$ un schéma. Soit $Y$ un espace algébrique formel affine.
Soit $f : X \to Y$ une application de faisceaux sur $(\Sch/S)_{fppf}$ qui
est représentable et affine. Alors
\begin{enumerate}
\item $X$ est un espace algébrique formel affine ;
\item si $Y$ est à indexation dénombrable, alors $X$ l'est aussi ;
\item si $Y$ est à indexation dénombrable et classique, alors $X$ est
à indexation dénombrable et classique ;
\item si $Y$ est faiblement adique, alors $X$ est faiblement adique ;
\item si $Y$ est adique*, alors $X$ est adique* ;
\item si $Y$ est noethérien et si $f$ est (localement) de type fini, alors
$X$ est noethérien.
\end{enumerate}
\end{lemma}

\begin{proof}
Preuve de (1). Écrivons $Y = \colim_{\lambda \in \Lambda} Y_\lambda$ comme dans la
Définition \ref{formal-spaces-definition-affine-formal-algebraic-space}.
Puisque $f$ est représentable et affine, les produits fibrés
$X_\lambda = Y_\lambda \times_Y X$ sont affines. De plus,
$X = \colim Y_\lambda \times_Y X$.
Ainsi, $X$ est un espace algébrique formel affine.

\medskip\noindent
Preuve de (2). Si $Y$ est à indexation dénombrable, alors, dans l'argument précédent,
nous pouvons supposer $\Lambda$ dénombrable.
Nous voyons aussitôt que $X$ est également à indexation dénombrable.

\medskip\noindent
Preuve de (3), (4) et (5). Dans chacun de ces cas, les hypothèses impliquent
que $Y$ est un espace algébrique formel affine à indexation dénombrable
(lemme \ref{formal-spaces-lemma-implications-between-types})
et donc que $X$ l'est aussi d'après (2). Nous pouvons ainsi écrire $X = \text{Spf}(A)$
et $Y = \text{Spf}(B)$ pour des $S$-algèbres topologiques faiblement
admissibles $A$ et $B$ ; voir le Lemme \ref{formal-spaces-lemma-countably-indexed}.
D'après le Lemme \ref{formal-spaces-lemma-morphism-between-formal-spectra},
le morphisme $f$ correspond à un homomorphisme continu de $S$-algèbres
$\varphi : B \to A$. Le
Lemme \ref{formal-spaces-lemma-representable-affine} montre que $\varphi$ est tendu. Nous concluons que
(3) résulte du Lemme \ref{formal-spaces-lemma-taut-ascent-admissible},
(4) résulte du Lemme \ref{formal-spaces-lemma-taut-ascent-weakly-adic}, et
(5) résulte du Lemme \ref{formal-spaces-lemma-taut-is-adic}.

\medskip\noindent
Preuve de (6). En combinant (3) avec le Lemme \ref{formal-spaces-lemma-implications-between-types},
nous voyons que $X$ est adique*. Nous pouvons donc utiliser le critère du
Lemme \ref{formal-spaces-lemma-characterize-noetherian-affine}.
Il nous dit d'abord que les schémas affines $Y_\lambda$ sont noethériens.
Puis $X_\lambda \to Y_\lambda$ est de type fini, donc $X_\lambda$
est également noethérien (Morphismes, lemme
\ref{morphisms-lemma-finite-type-noetherian}).
Le critère nous dit alors que $X$ est noethérien, ce qui achève la
démonstration.
\end{proof}

\begin{lemma}
\label{formal-spaces-lemma-property-goes-up-affine}
Soit $S$ un schéma. Soit $f : X \to Y$ un morphisme d'espaces algébriques
formels affines, représentable par des espaces algébriques. Alors
\begin{enumerate}
\item si $Y$ est à indexation dénombrable, alors $X$ l'est aussi ;
\item si $Y$ est à indexation dénombrable et classique, alors $X$ est à indexation
dénombrable et classique ;
\item si $Y$ est faiblement adique, alors $X$ est faiblement adique ;
\item si $Y$ est adique*, alors $X$ est adique* ;
\item si $Y$ est noethérien et si $f$ est (localement) de type fini, alors
$X$ est noethérien.
\end{enumerate}
\end{lemma}

\begin{proof}
Combiner les Lemmes \ref{formal-spaces-lemma-affine-representable-by-algebraic-spaces} et
\ref{formal-spaces-lemma-property-goes-up-affine-morphism}.
\end{proof}

\begin{example}
\label{formal-spaces-example-representable-morphism-from-completion}
Soit $B$ un anneau topologique faiblement admissible. Soit $B \to A$
un homomorphisme d'anneaux (sans topologie). Nous pouvons alors considérer
$$
A^\wedge = \lim A/JA
$$
où la limite porte sur tous les idéaux faibles de définition $J$ de $B$.
Alors $A^\wedge$ (muni de la topologie limite) est un anneau
complet linéairement topologisé. Le noyau (ouvert) $I$
de la surjection $A^\wedge \to A/JA$ est l'adhérence de $JA^\wedge$ ; voir le
Lemme \ref{formal-spaces-lemma-closed}. D'après le
lemme \ref{formal-spaces-lemma-topologically-nilpotent}
nous voyons que $I$ est formé d'éléments topologiquement nilpotents.
Ainsi, $I$ est un idéal faible de définition de $A^\wedge$, et nous concluons que
$A^\wedge$ est un anneau topologique faiblement admissible. Par conséquent,
$\varphi : B \to A^\wedge$ est une application tendue d'anneaux topologiques
faiblement admissibles et
$$
\text{Spf}(A^\wedge) \longrightarrow \text{Spf}(B)
$$
est un cas particulier du phénomène étudié dans le
Lemme \ref{formal-spaces-lemma-representable-affine}.
\end{example}

\begin{remark}[Avertissement]
\label{formal-spaces-remark-warning}
La discussion des Lemmes \ref{formal-spaces-lemma-representable-affine},
\ref{formal-spaces-lemma-property-goes-up-affine-morphism} et
\ref{formal-spaces-lemma-property-goes-up-affine}
est optimale aux deux sens suivants :
\begin{enumerate}
\item Si $A$ et $B$ sont des anneaux faiblement admissibles et si $\varphi : A \to B$
est une application continue, alors
$\text{Spf}(\varphi) : \text{Spf}(B) \to \text{Spf}(A)$ n'est en général
pas représentable ;
\item si $f : Y \to X$ est un morphisme représentable d'espaces algébriques
formels affines et si $X = \text{Spf}(A)$ est de McQuillan,
il n'en résulte pas que $Y$ soit de McQuillan.
\end{enumerate}
Un exemple pour (1) consiste à prendre $A = k$, un corps (muni de la topologie discrète),
et $B = k[[t]]$ muni de la topologie $t$-adique.
Un exemple pour (2) est donné dans
Exemples, section \ref{examples-section-affine-formal-algebraic-space}.
\end{remark}

\noindent
Malgré l'avertissement précédent, nous avons le résultat suivant.

\begin{lemma}
\label{formal-spaces-lemma-etale}
Soit $S$ un schéma. Soit $Y$ un espace algébrique formel affine de McQuillan
sur $S$, c'est-à-dire $Y = \text{Spf}(B)$ pour une $S$-algèbre topologique
faiblement admissible $B$. Alors il existe une équivalence de catégories entre
\begin{enumerate}
\item la catégorie des morphismes $f : X \to Y$
d'espaces algébriques formels affines qui sont représentables
par des espaces algébriques et étales ;
\item la catégorie des $B$-algèbres topologiques de la forme
$A^\wedge$, où $A$ est une $B$-algèbre étale et où
$A^\wedge = \lim A/JA$, avec $J \subset B$ parcourant les
idéaux faibles de définition de $B$.
\end{enumerate}
L'équivalence envoie $A^\wedge$ sur $X = \text{Spf}(A^\wedge)$.
En particulier, tout $X$ comme en (1) est de McQuillan.
\end{lemma}

\begin{proof}
Soit $A$ une $B$-algèbre étale. Alors $B/J \to A/JA$ est
étale pour tout idéal ouvert $J \subset B$. Par conséquent,
le morphisme $\text{Spf}(A^\wedge) \to Y$ est représentable et étale.
Le foncteur $\text{Spf}$ est pleinement fidèle d'après le
lemme \ref{formal-spaces-lemma-morphism-between-formal-spectra}.
Pour achever la démonstration, nous allons montrer au paragraphe suivant
que tout $X \to Y$ comme en (1) appartient à l'image essentielle.

\medskip\noindent
Choisissons un idéal faible de définition $J_0 \subset B$. Posons
$Y_0 = \Spec(B/J_0)$ et $X_0 = Y_0 \times_Y X$. Alors $X_0 \to Y_0$
est un morphisme étale de schémas affines (voir le
lemme \ref{formal-spaces-lemma-affine-representable-by-algebraic-spaces}).
Écrivons $X_0 = \Spec(A_0)$. D'après Algèbre, lemme \ref{algebra-lemma-lift-etale},
nous pouvons trouver un homomorphisme étale d'algèbres $B \to A$ tel que
$A_0 \cong A/J_0A$. Considérons un idéal de définition $J \subset J_0$.
Comme ci-dessus, nous pouvons écrire $\Spec(B/J) \times_Y X = \Spec(\bar A)$
pour un homomorphisme étale d'anneaux $B/J \to \bar A$. Alors
$B/J \to \bar A$ et $B/J \to A/JA$ sont tous deux des homomorphismes étales d'anneaux
relevant l'homomorphisme étale d'anneaux $B/J_0 \to A_0$. D'après
Compléments d'algèbre, lemme \ref{more-algebra-lemma-locally-nilpotent-henselian}
il existe un unique isomorphisme de $B/J$-algèbres
$\varphi_J : A/JA \to \bar A$ relevant l'identification modulo $J_0$.
Comme les applications $\varphi_J$ sont uniques, elles sont compatibles lorsque $J$ varie.
Ainsi,
$$
X = \colim \Spec(B/J) \times_Y X = \colim \Spec(A/JA) = \text{Spf}(A)
$$
et nous voyons que le lemme est démontré.
\end{proof}

\begin{lemma}
\label{formal-spaces-lemma-etale-surjective}
Avec les notations et hypothèses du Lemme \ref{formal-spaces-lemma-etale}, soit
$f : X \to Y$ correspondant à $B \to A^\wedge$. Les conditions suivantes sont équivalentes :
\begin{enumerate}
\item $f : X \to Y$ est surjectif ;
\item $B \to A$ est fidèlement plat ;
\item pour tout idéal faible de définition $J \subset B$,
l'homomorphisme d'anneaux $B/J \to A/JA$ est fidèlement plat ;
\item pour un certain idéal faible de définition $J \subset B$,
l'homomorphisme d'anneaux $B/J \to A/JA$ est fidèlement plat.
\end{enumerate}
\end{lemma}

\begin{proof}
Soit $J \subset B$ un idéal faible de définition. Comme tout élément de $J$
est topologiquement nilpotent, nous voyons que tout élément de $1 + J$ est
inversible. Il s'ensuit que $J$ est contenu dans le radical de Jacobson de $B$
(Algèbre, lemme \ref{algebra-lemma-contained-in-radical}).
Ainsi, un homomorphisme plat d'anneaux $B \to A$ est fidèlement plat si et seulement si
$B/J \to A/JA$ est fidèlement plat
(Algèbre, lemme \ref{algebra-lemma-ff-rings}).
Nous voyons ainsi que (2) -- (4) sont équivalentes.
Si (1) est vérifiée, alors, pour tout idéal faible de définition $J \subset B$,
le morphisme
$\Spec(A/JA) = \Spec(B/J) \times_Y X \to \Spec(B/J)$ est surjectif,
ce qui implique (3). Réciproquement, supposons (3).
Un morphisme $T \to Y$ avec $T$ quasi-compact
se factorise par $\Spec(B/J)$ pour un certain idéal de définition $J$ de $B$
(lemme \ref{formal-spaces-lemma-factor-through-thickening}).
Ainsi, $X \times_Y T = \Spec(A/JA) \times_{\Spec(B/J)} T \to T$
est surjectif en tant que changement de base du morphisme surjectif
$\Spec(A/JA) \to \Spec(B/J)$. La condition (1) est donc vérifiée.
\end{proof}





\section{Types d'espaces algébriques formels}
\label{formal-spaces-section-types}

\noindent
Dans cette section, nous définissons
les espaces algébriques formels
« localement noethériens »,
« localement adiques* »,
« localement faiblement adiques »,
« localement à indexation dénombrable et classiques », et
« localement à indexation dénombrable ».
Les types « localement adiques »,
« localement classiques » et
« localement de McQuillan »
manquent, car nous ne savons pas démontrer l'analogue des lemmes suivants
dans ces cas (il suffirait de démontrer l'analogue de ces
lemmes pour les recouvrements étales entre espaces algébriques formels affines).

\begin{lemma}
\label{formal-spaces-lemma-iff-countably-indexed}
Soit $S$ un schéma. Soit $X \to Y$ un morphisme d'espaces algébriques
formels affines qui est représentable par des espaces algébriques,
surjectif et plat. Alors $X$ est à indexation dénombrable si et seulement
si $Y$ est à indexation dénombrable.
\end{lemma}

\begin{proof}
Supposons $X$ à indexation dénombrable. Écrivons $X = \colim X_n$ comme dans le
Lemme \ref{formal-spaces-lemma-countable-affine-formal-algebraic-space}.
Écrivons $Y = \colim Y_\lambda$ comme dans la
Définition \ref{formal-spaces-definition-affine-formal-algebraic-space}.
Pour tout $n$, nous pouvons choisir un $\lambda_n$ tel que
$X_n \to Y$ se factorise par $Y_{\lambda_n}$ ; voir le
lemme \ref{formal-spaces-lemma-factor-through-thickening}.
D'autre part, pour tout $\lambda$, le schéma
$Y_\lambda \times_Y X$ est affine
(Lemme \ref{formal-spaces-lemma-affine-representable-by-algebraic-spaces})
et, par conséquent, $Y_\lambda \times_Y X \to X$ se factorise par
$X_n$ pour un certain $n$ (Lemme \ref{formal-spaces-lemma-factor-through-thickening}).
Considérons le diagramme
$$
\xymatrix{
Y_\lambda \times_Y X \ar[r] \ar[d] & X_n \ar[r] \ar[d] & X \ar[d] \\
Y_\lambda \ar@{..>}[r] \ar@/_1pc/[rr] & Y_{\lambda_n} \ar[r] & Y
}
$$
Si nous montrons que la flèche pointillée existe, nous concluons que
$Y = \colim Y_{\lambda_n}$ et que $Y$ est à indexation dénombrable. Pour cela, choisissons
un $\mu$ tel que $\mu \geq \lambda$ et $\mu \geq \lambda_n$.
Ainsi, $Y_\lambda \to Y$ et $Y_{\lambda_n} \to Y$ se factorisent tous deux
par $Y_\mu \to Y$.
Écrivons $Y_\mu = \Spec(B_\mu)$ ; le sous-schéma fermé $Y_\lambda$ correspond à
$J \subset B_\mu$, et le sous-schéma fermé $Y_{\lambda_n}$ correspond à
$J' \subset B_\mu$. Nous voulons montrer que $J' \subset J$.
Le diagramme précédent donne $J'A_\mu \subset JA_\mu$,
où $Y_\mu \times_Y X = \Spec(A_\mu)$.
Puisque $X \to Y$ est surjectif et plat, le morphisme
$Y_\lambda \times_Y X \to Y_\lambda$ est un morphisme fidèlement plat
de schémas affines ; ainsi, $B_\mu \to A_\mu$ est
fidèlement plat. Donc $J' \subset J$, comme voulu.

\medskip\noindent
Supposons $Y$ à indexation dénombrable. Alors $X$ est à indexation dénombrable
d'après le Lemme \ref{formal-spaces-lemma-property-goes-up-affine}.
\end{proof}

\begin{lemma}
\label{formal-spaces-lemma-iff-cic}
Soit $S$ un schéma. Soit $X \to Y$ un morphisme d'espaces algébriques
formels affines qui est représentable par des espaces algébriques,
surjectif et plat. Alors $X$ est à indexation dénombrable et classique
si et seulement si $Y$ est à indexation dénombrable et classique.
\end{lemma}

\begin{proof}
Nous avons déjà vu l'implication dans un sens au
Lemme \ref{formal-spaces-lemma-property-goes-up-affine}. Dans l'autre sens,
notons que, d'après le Lemme \ref{formal-spaces-lemma-iff-countably-indexed}, nous pouvons supposer
$X$ et $Y$ tous deux à indexation dénombrable. Ainsi, $X = \text{Spf}(A)$
et $Y = \text{Spf}(B)$ pour des $S$-algèbres topologiques faiblement
admissibles $A$ et $B$ ; voir le Lemme \ref{formal-spaces-lemma-countably-indexed}.
D'après le Lemme \ref{formal-spaces-lemma-morphism-between-formal-spectra},
le morphisme $X \to Y$ correspond à un homomorphisme continu de $S$-algèbres
$\varphi : B \to A$. Le
Lemme \ref{formal-spaces-lemma-representable-affine} montre que $\varphi$ est tendu.
Soit $J \subset B$ un idéal ouvert et soit $I \subset A$
l'adhérence de $JA$. D'après les
Lemmes \ref{formal-spaces-lemma-fibre-product-affines-over-separated} et
\ref{formal-spaces-lemma-closure-image-ideal}
nous voyons que $\Spec(B/J) \times_Y X = \Spec(A/I)$.
Ainsi, $B/J \to A/I$ est fidèlement plat (puisque $X \to Y$ est surjectif
et plat). Cela signifie que $\varphi : B \to A$ est comme dans la
section \ref{formal-spaces-section-taut-descent} (les rôles de $A$ et $B$ étant échangés).
Le lemme résulte donc du
Lemme \ref{formal-spaces-lemma-taut-descent-weakly-admissible}.
\end{proof}

\begin{lemma}
\label{formal-spaces-lemma-iff-weakly-adic}
Soit $S$ un schéma. Soit $X \to Y$ un morphisme d'espaces algébriques
formels affines qui est représentable par des espaces algébriques,
surjectif et plat. Alors $X$ est faiblement adique
si et seulement si $Y$ est faiblement adique.
\end{lemma}

\begin{proof}
La démonstration est exactement la même que celle du
Lemme \ref{formal-spaces-lemma-iff-cic}, sauf qu'à la fin on utilise le
Lemme \ref{formal-spaces-lemma-taut-descent-weakly-adic}.
\end{proof}

\begin{lemma}
\label{formal-spaces-lemma-iff-adic-star}
Soit $S$ un schéma. Soit $X \to Y$ un morphisme d'espaces algébriques
formels affines qui est représentable par des espaces algébriques,
surjectif et plat. Alors $X$ est adique* si et seulement si $Y$ est adique*.
\end{lemma}

\begin{proof}
La démonstration est exactement la même que celle du
Lemme \ref{formal-spaces-lemma-iff-cic}, sauf qu'à la fin on utilise le
Lemme \ref{formal-spaces-lemma-taut-descent-adic-star}.
\end{proof}

\begin{lemma}
\label{formal-spaces-lemma-iff-noetherian}
Soit $S$ un schéma. Soit $X \to Y$ un morphisme d'espaces algébriques
formels affines qui est représentable par des espaces algébriques,
surjectif, plat et (localement) de type fini. Alors $X$ est noethérien
si et seulement si $Y$ est noethérien.
\end{lemma}

\begin{proof}
Observons qu'un espace algébrique formel affine noethérien est adique* ; voir le
Lemme \ref{formal-spaces-lemma-implications-between-types}. Ainsi, d'après le
Lemme \ref{formal-spaces-lemma-iff-adic-star}, nous pouvons supposer que $X$ et $Y$
sont tous deux adiques*. Nous utiliserons le critère du
Lemme \ref{formal-spaces-lemma-characterize-noetherian-affine}
pour établir le lemme. Écrivons $Y = \colim Y_n$
comme dans le Lemme \ref{formal-spaces-lemma-countable-affine-formal-algebraic-space}.
Pour tout $n$, posons $X_n = Y_n \times_Y X$. Alors $X_n$ est un
schéma affine (Lemme \ref{formal-spaces-lemma-affine-representable-by-algebraic-spaces})
et $X = \colim X_n$. Chacun des morphismes $X_n \to Y_n$ est
fidèlement plat et de type fini. Le lemme résulte donc du
fait que, dans cette situation, $X_n$ est noethérien si et seulement si $Y_n$
est noethérien ; voir
Algèbre, lemme \ref{algebra-lemma-descent-Noetherian} (pour la descente)
et
Algèbre, lemme \ref{algebra-lemma-Noetherian-permanence} (pour la montée).
\end{proof}

\begin{lemma}
\label{formal-spaces-lemma-type-local}
Soit $S$ un schéma. Soit
$$
P \in
\left\{
\begin{matrix}
countably\ indexed,\\
\text{à indexation dénombrable et classique},\\
weakly\ adic,\ adic*,\ Noetherian
\end{matrix}
\right\}
$$
Soit $X$ un espace algébrique formel sur $S$.
Les conditions suivantes sont équivalentes :
\begin{enumerate}
\item si $Y$ est un espace algébrique formel affine et si
$f : Y \to X$ est représentable par des espaces algébriques et étale,
alors $Y$ possède la propriété $P$ ;
\item pour un certain $\{X_i \to X\}_{i \in I}$ comme dans la
Définition \ref{formal-spaces-definition-formal-algebraic-space},
chaque $X_i$ possède la propriété $P$.
\end{enumerate}
\end{lemma}

\begin{proof}
Il est clair que (1) implique (2). Supposons (2) et soit
$Y \to X$ comme en (1). Puisque les produits fibrés $X_i \times_X Y$
sont des espaces algébriques formels (Lemme \ref{formal-spaces-lemma-fibre-products-general}),
nous pouvons choisir des recouvrements $\{X_{ij} \to X_i \times_X Y\}$ comme dans la
Définition \ref{formal-spaces-definition-formal-algebraic-space}.
Puisque $Y$ est quasi-compact, il existe
$(i_1, j_1), \ldots, (i_n, j_n)$ tels que
$$
X_{i_1 j_1} \amalg \ldots \amalg X_{i_n j_n} \longrightarrow Y
$$
soit surjectif et étale. Alors $X_{i_kj_k} \to X_{i_k}$ est représentable
par des espaces algébriques et étale ; ainsi, $X_{i_kj_k}$ possède la propriété $P$ d'après le
Lemme \ref{formal-spaces-lemma-property-goes-up-affine}.
Alors $X_{i_1 j_1} \amalg \ldots \amalg X_{i_n j_n}$ est un
espace algébrique formel affine possédant la propriété $P$ (nous omettons un petit détail
concernant les réunions disjointes finies d'espaces algébriques formels affines).
Nous concluons donc en appliquant l'un des
lemmes \ref{formal-spaces-lemma-iff-countably-indexed},
\ref{formal-spaces-lemma-iff-cic},
\ref{formal-spaces-lemma-iff-weakly-adic},
\ref{formal-spaces-lemma-iff-adic-star}, et
\ref{formal-spaces-lemma-iff-noetherian}.
\end{proof}

\noindent
Le lemme précédent permet de poser la définition suivante.

\begin{definition}
\label{formal-spaces-definition-types-formal-algebraic-spaces}
Soit $S$ un schéma. Soit $X$ un espace algébrique formel sur $S$.
On dit que $X$ est
{\it localement à indexation dénombrable},
{\it localement à indexation dénombrable et classique},
{\it localement faiblement adique},
{\it localement adique*}, ou
{\it localement noethérien}
si les conditions équivalentes du Lemme \ref{formal-spaces-lemma-type-local}
sont satisfaites pour la propriété correspondante.
\end{definition}

\noindent
Le complété formel d'un espace algébrique localement noethérien
le long d'une partie fermée est un espace algébrique formel localement noethérien.

\begin{lemma}
\label{formal-spaces-lemma-formal-completion-types}
Soit $S$ un schéma. Soit $X$ un espace algébrique sur $S$.
Soit $T \subset |X|$ une partie fermée. Soit $X_{/T}$ le
complété formel de $X$ le long de $T$.
\begin{enumerate}
\item Si $X \setminus T \to X$ est quasi-compact,
alors $X_{/T}$ est localement adique* ;
\item si $X$ est localement noethérien, alors $X_{/T}$ est localement
noethérien.
\end{enumerate}
\end{lemma}

\begin{proof}
Choisissons un morphisme étale surjectif $U \to X$ avec $U = \coprod U_i$,
réunion disjointe de schémas affines ; voir Propriétés des espaces algébriques, lemme
\ref{spaces-properties-lemma-cover-by-union-affines}.
Soit $T_i \subset U_i$ l'image réciproque de $T$.
Nous avons $X_{/T} \times_X U_i = (U_i)_{/T_i}$
(Lemme \ref{formal-spaces-lemma-map-completions-representable}).
Ainsi, $\{(U_i)_{/T_i} \to X_{/T}\}$ est un recouvrement comme dans la
Définition \ref{formal-spaces-definition-formal-algebraic-space}.
De plus, si $X \setminus T \to X$ est quasi-compact, alors
$U_i \setminus T_i \to U_i$ l'est aussi, et si $X$ est localement noethérien,
$U_i$ l'est aussi. Le lemme résulte donc du cas affine, qui est le
Lemme \ref{formal-spaces-lemma-affine-formal-completion-types}.
\end{proof}

\begin{remark}[Avertissement]
\label{formal-spaces-remark-warning-completion}
Supposons $X = \Spec(A)$ et que $T \subset X$ soit le lieu des zéros d'un
idéal de type fini $I \subset A$. Soit $J = \sqrt{I}$ le
radical de $I$. Les définitions montrent alors que
$X_{/T} = \text{Spf}(A^\wedge)$, où $A^\wedge = \lim A/I^n$ est
le complété $I$-adique de $A$. D'autre part, l'application
$A^\wedge \to \lim A/J^n$ du complété $I$-adique vers
le complété $J$-adique peut ne pas être un isomorphisme d'anneaux.
À titre d'exemple, posons
$$
A = \bigcup\nolimits_{n \geq 1} \mathbf{C}[t^{1/n}]
$$
et $I = (t)$. Alors $J = \mathfrak m$ est l'idéal maximal
de l'anneau de valuation $A$ et $J^2 = J$. Ainsi, le complété $J$-adique
de $A$ est $\mathbf{C}$, tandis que le complété $I$-adique
est l'anneau de valuation décrit dans l'Exemple \ref{formal-spaces-example-david-hansen}
(mais il est notamment facile de voir que $A \subset A^\wedge$).
\end{remark}

\begin{lemma}
\label{formal-spaces-lemma-types-fibre-products}
Soit $S$ un schéma. Soient $X \to Y$ et $Z \to Y$ des
morphismes d'espaces algébriques formels sur $S$. Alors
\begin{enumerate}
\item Si $X$ et $Z$ sont localement à indexation dénombrable, alors $X \times_Y Z$
l'est aussi ;
\item si $X$ et $Z$ sont localement à indexation dénombrable et classiques,
alors $X \times_Y Z$ l'est aussi et est classique ;
\item si $X$ et $Z$ sont faiblement adiques, alors $X \times_Y Z$
est faiblement adique ;
\item si $X$ et $Z$ sont localement adiques*, alors $X \times_Y Z$ est
localement adique* ;
\item si $X$ et $Z$ sont localement noethériens et si $X_{red} \to Y_{red}$
est localement de type fini, alors $X \times_Y Z$ est localement noethérien.
\end{enumerate}
\end{lemma}

\begin{proof}
Choisissons un recouvrement $\{Y_j \to Y\}$ comme dans la
Définition \ref{formal-spaces-definition-formal-algebraic-space}.
Pour chaque $j$, choisissons un recouvrement $\{X_{ji} \to Y_j \times_Y X\}$
comme dans la Définition \ref{formal-spaces-definition-formal-algebraic-space}.
Pour chaque $j$, choisissons un recouvrement $\{Z_{jk} \to Y_j \times_Y Z\}$
comme dans la Définition \ref{formal-spaces-definition-formal-algebraic-space}.
Observons que $X_{ji} \times_{Y_j} Z_{jk}$ est un
espace algébrique formel affine d'après le
Lemme \ref{formal-spaces-lemma-fibre-product-affines-over-separated}.
Ainsi,
$$
\{X_{ji} \times_{Y_j} Z_{jk} \to X \times_Y Z\}
$$
est un recouvrement comme dans la Définition \ref{formal-spaces-definition-formal-algebraic-space}.
Il suffit donc de prouver (1), (2), (3) et (4) lorsque $X$, $Y$ et $Z$
sont des espaces algébriques formels affines.

\medskip\noindent
Supposons $X$ et $Z$ à indexation dénombrable. Écrivons $X = \colim X_n$ et
$Z = \colim Z_m$ comme dans le
Lemme \ref{formal-spaces-lemma-countable-affine-formal-algebraic-space}.
Écrivons $Y = \colim_{\lambda \in \Lambda} Y_\lambda$ comme dans la
Définition \ref{formal-spaces-definition-affine-formal-algebraic-space}.
Pour tous $n$ et $m$, nous pouvons trouver $\lambda_{n, m} \in \Lambda$
tel que $X_n \to Y$ et $Z_m \to Y$ se factorisent par $Y_{\lambda_{n, m}}$
(voir par exemple le Lemme \ref{formal-spaces-lemma-factor-through-thickening}).
Choisissons $\lambda_0 \in \Lambda$. Par récurrence sur $t \geq 1$,
choisissons un élément $\lambda_t \in \Lambda$ tel que
$\lambda_t \geq \lambda_{n, m}$ pour tous $1 \leq n, m \leq t$
et $\lambda_t \geq \lambda_{t - 1}$. Posons $Y' = \colim Y_{\lambda_t}$.
Alors $Y' \to Y$ est un monomorphisme tel que $X \to Y$ et $Z \to Y$
se factorisent par $Y'$. Nous pouvons donc remplacer $Y$ par $Y'$, c'est-à-dire
supposer $Y$ à indexation dénombrable.

\medskip\noindent
Supposons $X$, $Y$ et $Z$ à indexation dénombrable. D'après le
Lemme \ref{formal-spaces-lemma-countably-indexed}, nous pouvons écrire
$X = \text{Spf}(A)$, $Y = \text{Spf}(B)$, $Z = \text{Spf}(C)$
pour des anneaux topologiques faiblement admissibles $A$, $B$ et $C$.
Les morphismes $X \to Y$ et $Z \to Y$ sont donnés par des
homomorphismes continus d'anneaux $B \to A$ et $B \to C$ ; voir le
Lemme \ref{formal-spaces-lemma-morphism-between-formal-spectra}.
D'après le Lemme \ref{formal-spaces-lemma-fibre-product-affines-over-separated},
nous voyons que $X \times_Y Z = \text{Spf}(A \widehat{\otimes}_B C)$
et que $A \widehat{\otimes}_B C$ est un anneau topologique faiblement admissible.
En particulier, nous voyons que $X \times_Y Z$ est à indexation dénombrable
d'après la partie (3) du Lemme \ref{formal-spaces-lemma-completed-tensor-product}.
Ceci démontre (1).

\medskip\noindent
Preuve de (2). Dans ce cas, $X$ et $Z$ sont à indexation dénombrable ;
les arguments précédents montrent donc que $X \times_Y Z$ est le spectre formel
de $A \widehat{\otimes}_B C$, où $A$ et $C$ sont admissibles.
Alors $A \widehat{\otimes}_B C$ est admissible d'après la
partie (2) du Lemme \ref{formal-spaces-lemma-completed-tensor-product}.

\medskip\noindent
Preuve de (3). Comme précédemment, nous concluons que $X \times_Y Z$ est le spectre formel
de $A \widehat{\otimes}_B C$, où $A$ et $C$ sont faiblement adiques.
Alors $A \widehat{\otimes}_B C$ est faiblement adique d'après le
Lemme \ref{formal-spaces-lemma-completed-tensor-product-weakly-adic}.

\medskip\noindent
Preuve de (4). En raisonnant comme ci-dessus, cela résulte de la
partie (4) du Lemme \ref{formal-spaces-lemma-completed-tensor-product}.

\medskip\noindent
Preuve de (5). Pour déduire le cas (5) de la
partie (5) du Lemme \ref{formal-spaces-lemma-completed-tensor-product},
nous devons vérifier que les hypothèses coïncident.
Avec les notations du premier paragraphe de la démonstration,
si $X_{red} \to Y_{red}$ est localement de type fini, alors
$(X_{ji})_{red} \to (Y_j)_{red}$ est localement de type fini. Cela résulte de
Morphismes d'espaces algébriques, lemme \ref{spaces-morphisms-lemma-finite-type-local}
et du fait que, dans le diagramme commutatif
$$
\xymatrix{
(X_{ji})_{red} \ar[d] \ar[r] & (Y_j)_{red} \ar[d] \\
X_{red} \ar[r] & Y_{red}
}
$$
les morphismes verticaux sont étales. En effet, nous avons
$(X_{ji})_{red} = X_{ij} \times_X X_{red}$ et
$(Y_j)_{red} = Y_j \times_Y Y_{red}$
par le Lemme \ref{formal-spaces-lemma-reduction-smooth}. Comme ci-dessus, nous nous ramenons donc
au cas où $X$, $Y$, $Z$ sont des espaces algébriques formels affines,
où $X$, $Z$ sont noethériens et où $X_{red} \to Y_{red}$ est de type fini.
Ensuite, au deuxième paragraphe de la démonstration, nous avons remplacé $Y$ par $Y'$ ;
mais, par construction, $Y_{red} = Y'_{red}$, donc l'hypothèse de type fini
est préservée par ce remplacement. Nous voyons alors que
$X, Y, Z$ correspondent à $A, B, C$, et $X \times_Y Z$ à
$A \widehat{\otimes}_B C$, avec $A$, $C$ noethériens adiques.
Enfin, prendre la réduction revient à quotienter par
l'idéal des éléments topologiquement nilpotents
(exemple \ref{formal-spaces-example-reduction-affine-formal-spectrum})
et le fait que $X_{red} \to Y_{red}$ soit de type fini
signifie bien que $B/\mathfrak b \to A/\mathfrak a$ est de type fini.
Ceci achève la démonstration.
\end{proof}

\begin{lemma}
\label{formal-spaces-lemma-structure-locally-noetherian}
Soit $S$ un schéma. Soit $X$ un espace algébrique formel localement noethérien
sur $S$. Alors $X = \colim X_n$ pour un système $X_1 \to X_2 \to X_3 \to \ldots$
d'épaississements d'ordre fini d'espaces algébriques localement noethériens sur $S$,
où $X_1 = X_{red}$ et où $X_n$ est le $n$-ième voisinage infinitésimal de
$X_1$ dans $X_m$ pour tout $m \geq n$.
\end{lemma}

\begin{proof}
Nous ne faisons qu'esquisser la démonstration et omettons certains détails.
Posons $X_1 = X_{red}$. Définissons $X_n \subset X$ comme le sous-foncteur
caractérisé par la règle suivante : un morphisme $f : T \to X$, où $T$ est un schéma, se factorise
par $X_n$ si et seulement si la $n$-ième puissance du faisceau d'idéaux
de l'immersion fermée $X_1 \times_X T \to T$ est nulle. Alors $X_n \subset X$
est un sous-faisceau, car l'annulation des modules quasi-cohérents se vérifie
localement pour la topologie fppf. Nous affirmons que $X_n \to X$ est représentable par des schémas,
est une immersion fermée et que $X = \colim X_n$ (comme faisceaux fppf).
Pour le vérifier, nous pouvons travailler localement pour la topologie étale sur $X$. Nous pouvons donc supposer que
$X = \text{Spf}(A)$ est un espace algébrique formel affine noethérien.
Alors $X_1 = \Spec(A/\mathfrak a)$, où $\mathfrak a \subset A$
est l'idéal des éléments topologiquement nilpotents de l'anneau topologique
adique noethérien $A$. Alors $X_n = \Spec(A/\mathfrak a^n)$,
ce qui donne le résultat voulu.
\end{proof}









\section{Morphismes et homomorphismes continus d'anneaux}
\label{formal-spaces-section-morphisms-rings}

\noindent
Dans cette section, nous notons $\textit{WAdm}$ la catégorie des
anneaux topologiques faiblement admissibles et des homomorphismes continus d'anneaux.
Nous définissons les sous-catégories pleines
$$
\textit{WAdm} \supset
\textit{WAdm}^{count} \supset
\textit{WAdm}^{cic} \supset
\textit{WAdm}^{weakly\ adic} \supset
\textit{WAdm}^{adic*} \supset
\textit{WAdm}^{Noeth}
$$
dont les objets sont les suivants :
\begin{enumerate}
\item $\textit{WAdm}^{count}$ : les anneaux topologiques faiblement admissibles
$A$ qui possèdent un système fondamental dénombrable d'idéaux ouverts ;
\item $\textit{WAdm}^{cic}$ : les anneaux topologiques admissibles
$A$ qui possèdent un système fondamental dénombrable d'idéaux ouverts ;
\item $\textit{WAdm}^{weakly\ adic}$ : les anneaux topologiques faiblement adiques
(section \ref{formal-spaces-section-weakly-adic}) ;
\item $\textit{WAdm}^{adic*}$ : les anneaux topologiques adiques qui possèdent
un idéal de définition de type fini ;
\item $\textit{WAdm}^{Noeth}$ : les anneaux topologiques adiques qui
sont noethériens.
\end{enumerate}
Manifestement, les spectres formels de ces types d'anneaux sont les briques
élémentaires des espaces algébriques formels
localement à indexation dénombrable,
localement à indexation dénombrable et classiques,
localement faiblement adiques,
localement adiques* et
localement noethériens.

\medskip\noindent
Rappelons brièvement la relation entre les morphismes d'espaces algébriques
formels affines à indexation dénombrable et les
morphismes de $\textit{WAdm}^{count}$.
Soit $S$ un schéma. Soient $X$ et $Y$ des espaces algébriques formels
affines à indexation dénombrable. Écrivons $X = \text{Spf}(A)$
et $Y = \text{Spf}(B)$ pour des $S$-algèbres topologiques
$A$ et $B$ de $\textit{WAdm}^{count}$ ; voir le
Lemme \ref{formal-spaces-lemma-countably-indexed}.
D'après le Lemme \ref{formal-spaces-lemma-morphism-between-formal-spectra},
il existe une correspondance biunivoque entre les morphismes
$f : X \to Y$ et les homomorphismes continus
$$
\varphi : B \longrightarrow A
$$
de $S$-algèbres topologiques. Cette correspondance est donnée par
$f \mapsto f^\sharp$ et $\varphi \mapsto \text{Spf}(\varphi)$.

\medskip\noindent
Soit $S$ un schéma. Soit $f : X \to Y$ un morphisme d'espaces algébriques
formels localement à indexation dénombrable. Considérons un
diagramme commutatif
$$
\xymatrix{
U \ar[d] \ar[r] & V \ar[d] \\
X \ar[r] & Y
}
$$
où $U$ et $V$ sont des espaces algébriques formels affines et où $U \to X$ et $V \to Y$
sont représentables par des espaces algébriques et étales. D'après la
Définition \ref{formal-spaces-definition-types-formal-algebraic-spaces} (donc aussi le
Lemme \ref{formal-spaces-lemma-type-local}), $U$ et $V$ sont des espaces algébriques formels
affines à indexation dénombrable. La discussion du paragraphe
précédent montre que $U \to V$ est isomorphe à $\text{Spf}(\varphi)$
pour un certain homomorphisme continu
$$
\varphi : B \longrightarrow A
$$
de $S$-algèbres topologiques dans $\textit{WAdm}^{count}$.

\begin{lemma}
\label{formal-spaces-lemma-completion-in-sub}
Soit $A \in \Ob(\textit{WAdm})$. Soit $A \to A'$ un homomorphisme
d'anneaux (sans topologie). Soit $(A')^\wedge = \lim_{I \subset A\text{ w.i.d}} A'/IA'$
l'objet de $\textit{WAdm}$ construit dans
l'Exemple \ref{formal-spaces-example-representable-morphism-from-completion}.
\begin{enumerate}
\item Si $A$ appartient à $\textit{WAdm}^{count}$, il en est de même de $(A')^\wedge$.
\item Si $A$ appartient à $\textit{WAdm}^{cic}$, il en est de même de $(A')^\wedge$.
\item Si $A$ appartient à $\textit{WAdm}^{weakly\ adic}$, il en est de même de $(A')^\wedge$.
\item Si $A$ appartient à $\textit{WAdm}^{adic*}$, il en est de même de $(A')^\wedge$.
\item Si $A$ appartient à $\textit{WAdm}^{Noeth}$ et si $A'$ est noethérien, alors
$(A')^\wedge$ appartient à $\textit{WAdm}^{Noeth}$.
\end{enumerate}
\end{lemma}

\begin{proof}
Rappelons que $A \to (A')^\wedge$ est tendu ; voir la discussion de
l'Exemple \ref{formal-spaces-example-representable-morphism-from-completion}.
Les assertions (1), (2), (3) et (4) résultent donc des Lemmes
\ref{formal-spaces-lemma-taut-ascent-countable},
\ref{formal-spaces-lemma-taut-ascent-admissible},
\ref{formal-spaces-lemma-taut-ascent-weakly-adic}, et
\ref{formal-spaces-lemma-taut-is-adic}.
Enfin, supposons $A$ noethérien et adique.
D'après (4), $(A')^\wedge$ est adique.
Algèbre, lemme \ref{algebra-lemma-completion-Noetherian-Noetherian}
montre que $(A')^\wedge$ est noethérien. Ainsi, (5) est vérifié.
\end{proof}

\begin{situation}
\label{formal-spaces-situation-local-property}
Soit $P$ une propriété des morphismes de $\textit{WAdm}^{count}$.
Considérons les diagrammes commutatifs
\begin{equation}
\label{formal-spaces-equation-localize}
\vcenter{
\xymatrix{
A \ar[r] & (A')^\wedge \\
B \ar[r] \ar[u]^\varphi & (B')^\wedge \ar[u]_{\varphi'}
}
}
\end{equation}
qui satisfont aux conditions suivantes :
\begin{enumerate}
\item $A$ et $B$ sont des objets de $\textit{WAdm}^{count}$ ;
\item $A \to A'$ et $B \to B'$ sont des homomorphismes étales d'anneaux ;
\item $(A')^\wedge = \lim A'/IA'$, resp.\  $(B')^\wedge = \lim B'/JB'$
où $I \subset A$, resp.\ $J \subset B$,
parcourt les idéaux faibles de définition de $A$, resp.\ de $B$ ;
\item $\varphi : B \to A$ et $\varphi' : (B')^\wedge \to (A')^\wedge$
sont continus.
\end{enumerate}
D'après le Lemme \ref{formal-spaces-lemma-completion-in-sub}, les anneaux topologiques
$(A')^\wedge$ et $(B')^\wedge$ sont des objets de $\textit{WAdm}^{count}$.
Nous disons que $P$ est une {\it propriété locale} si les axiomes suivants sont vérifiés :
\begin{enumerate}
\item
\label{formal-spaces-item-axiom-1}
pour tout diagramme (\ref{formal-spaces-equation-localize}), on a
$P(\varphi) \Rightarrow P(\varphi')$,
\item
\label{formal-spaces-item-axiom-2}
pour tout diagramme (\ref{formal-spaces-equation-localize}) où $A \to A'$
est fidèlement plat, on a
$P(\varphi') \Rightarrow P(\varphi)$,
\item
\label{formal-spaces-item-axiom-3}
si $P(B \to A_i)$ pour $i = 1, \ldots, n$, alors
$P(B \to \prod_{i = 1, \ldots, n} A_i)$.
\end{enumerate}
L'axiome (\ref{formal-spaces-item-axiom-3})
a un sens puisque $\textit{WAdm}^{count}$ possède les produits finis.
\end{situation}

\begin{lemma}
\label{formal-spaces-lemma-property-defines-property-morphisms}
Soit $S$ un schéma. Soit $f : X \to Y$ un morphisme d'espaces algébriques formels
localement à indexation dénombrable sur $S$.
Soit $P$ une propriété locale des morphismes de $\textit{WAdm}^{count}$.
Les conditions suivantes sont équivalentes :
\begin{enumerate}
\item pour tout diagramme commutatif
$$
\xymatrix{
U \ar[d] \ar[r] & V \ar[d] \\
X \ar[r] & Y
}
$$
où $U$ et $V$ sont des espaces algébriques formels affines, où $U \to X$ et $V \to Y$
sont représentables par des espaces algébriques et étales, le morphisme $U \to V$
correspond à un morphisme de $\textit{WAdm}^{count}$ qui possède la propriété $P$ ;
\item il existe un recouvrement $\{Y_j \to Y\}$ comme dans la
Définition \ref{formal-spaces-definition-formal-algebraic-space} et, pour tout $j$,
un recouvrement $\{X_{ji} \to Y_j \times_Y X\}$ comme dans la
Définition \ref{formal-spaces-definition-formal-algebraic-space},
tel que chaque $X_{ji} \to Y_j$ corresponde
à un morphisme de $\textit{WAdm}^{count}$ qui possède la propriété $P$ ;
\item il existe un recouvrement $\{X_i \to X\}$ comme dans la
Définition \ref{formal-spaces-definition-formal-algebraic-space}
et, pour tout $i$, une factorisation $X_i \to Y_i \to Y$ où $Y_i$
est un espace algébrique formel affine, où $Y_i \to Y$ est représentable
par des espaces algébriques et étale, et où $X_i \to Y_i$ correspond
à un morphisme de $\textit{WAdm}^{count}$ qui possède la propriété $P$.
\end{enumerate}
\end{lemma}

\begin{proof}
Il est clair que (1) implique (2) et que (2) implique (3).
Supposons $\{X_i \to X\}$ et $X_i \to Y_i \to Y$ comme en (3),
et donnons-nous un diagramme comme en (1).
Puisque $Y_i \times_Y V$ est un espace algébrique formel
(Lemme \ref{formal-spaces-lemma-fibre-products-general}), nous pouvons choisir
des recouvrements $\{Y_{ij} \to Y_i \times_Y V\}$ comme dans la
Définition \ref{formal-spaces-definition-formal-algebraic-space}.
Pour tout $(i, j)$, nous pouvons de même choisir des recouvrements
$\{X_{ijk} \to Y_{ij} \times_{Y_i} X_i \times_X U\}$
comme dans la Définition \ref{formal-spaces-definition-formal-algebraic-space}.
Puisque $U$ est quasi-compact, nous pouvons choisir
$(i_1, j_1, k_1), \ldots, (i_n, j_n, k_n)$ tels que
$$
X_{i_1 j_1 k_1} \amalg \ldots \amalg X_{i_n j_n k_n} \longrightarrow U
$$
soit surjectif. Pour $s = 1, \ldots, n$, considérons le diagramme commutatif
$$
\xymatrix{
& & & X_{i_s j_s k_s} \ar[ld] \ar[d] \ar[rd] \\
X \ar[d] & X_{i_s} \ar[l] \ar[d] &
X_{i_s} \times_X U \ar[l] \ar[d] & Y_{i_s j_s} \ar[ld] \ar[rd] &
X_{i_s} \times_X U \ar[d] \ar[r] &
U \ar[d] \ar[r] & X \ar[d] \\
Y & Y_{i_s} \ar[l] &
Y_{i_s} \times_Y V \ar[l] & &
Y_{i_s} \times_Y V \ar[r] &
V \ar[r] & Y
}
$$
Convenons que $P$ est vérifiée pour un morphisme d'espaces algébriques formels
affines à indexation dénombrable si elle est vérifiée pour le
morphisme correspondant de $\textit{WAdm}^{count}$. Remarquons que les morphismes
$X_{i_s j_s k_s} \to X_{i_s}$, $Y_{i_s j_s} \to Y_{i_s}$
sont donnés par des complétés d'homomorphismes étales d'anneaux ; voir le Lemme \ref{formal-spaces-lemma-etale}.
Nous voyons donc que $P(X_{i_s} \to Y_{i_s})$ implique
$P(X_{i_s j_s k_s} \to Y_{i_s j_s})$ par l'axiome (\ref{formal-spaces-item-axiom-1}).
Remarquons que les morphismes $Y_{i_s j_s} \to V$ sont donnés par des complétés
d'homomorphismes étales d'anneaux (même lemme que précédemment).
En appliquant l'axiome (\ref{formal-spaces-item-axiom-2}) au diagramme
$$
\xymatrix{
X_{i_s j_s k_s} \ar@{=}[r] \ar[d] & X_{i_s j_s k_s} \ar[d] \\
Y_{i_s j_s} \ar[r] & V
}
$$
(cela est permis puisque les identités sont des homomorphismes fidèlement plats d'anneaux),
nous concluons que $P(X_{i_s j_s k_s} \to V)$ est vérifiée.
L'axiome (\ref{formal-spaces-item-axiom-3}) montre alors que
$P(\coprod_{s = 1, \ldots, n} X_{i_s j_s k_s} \to V)$ est vérifiée.
Comme le morphisme $\coprod X_{i_s j_s k_s} \to U$ est surjectif
par construction, le morphisme correspondant de $\textit{WAdm}^{count}$
est le complété d'un homomorphisme étale et fidèlement plat d'anneaux ; voir le
Lemme \ref{formal-spaces-lemma-etale-surjective}.
Une dernière application de l'axiome (\ref{formal-spaces-item-axiom-2})
(avec $B' = B$) implique que $P(U \to V)$ est vraie, comme voulu.
\end{proof}

\begin{remark}[Variante adique*]
\label{formal-spaces-remark-variant-adic-star}
Soit $P$ une propriété des morphismes de $\textit{WAdm}^{adic*}$.
Nous disons que $P$ est une {\it propriété locale} si les axiomes
(\ref{formal-spaces-item-axiom-1}), (\ref{formal-spaces-item-axiom-2}), (\ref{formal-spaces-item-axiom-3})
de la Situation \ref{formal-spaces-situation-local-property}
sont vérifiés pour les morphismes de $\textit{WAdm}^{adic*}$. De la même manière,
nous obtenons une variante du Lemme \ref{formal-spaces-lemma-property-defines-property-morphisms}
pour les morphismes entre espaces algébriques formels localement adiques* sur $S$.
\end{remark}

\begin{remark}[Variante noethérienne]
\label{formal-spaces-remark-variant-Noetherian}
Soit $P$ une propriété des morphismes de $\textit{WAdm}^{Noeth}$.
Nous disons que $P$ est une {\it propriété locale} si les axiomes
(\ref{formal-spaces-item-axiom-1}), (\ref{formal-spaces-item-axiom-2}), (\ref{formal-spaces-item-axiom-3}),
de la Situation \ref{formal-spaces-situation-local-property}
sont vérifiés pour les morphismes de $\textit{WAdm}^{Noeth}$. De la même manière,
nous obtenons une variante du Lemme \ref{formal-spaces-lemma-property-defines-property-morphisms}
pour les morphismes entre espaces algébriques formels localement noethériens sur $S$.
\end{remark}

\begin{situation}
\label{formal-spaces-situation-base-change-local-property}
Soit $P$ une propriété locale des morphismes de $\textit{WAdm}^{count}$ ; voir la
Situation \ref{formal-spaces-situation-local-property}. Nous disons que $P$ est {\it stable par
changement de base} si, étant donnés $B \to A$ et $B \to C$ dans $\textit{WAdm}^{count}$,
on a $P(B \to A) \Rightarrow P(C \to A \widehat{\otimes}_B C)$.
Cela a un sens puisque $A \widehat{\otimes}_B C$ est un objet de
$\textit{WAdm}^{count}$ d'après le Lemme \ref{formal-spaces-lemma-completed-tensor-product}.
\end{situation}

\begin{lemma}
\label{formal-spaces-lemma-base-change-property-morphisms}
Soit $S$ un schéma. Soit $P$ une propriété locale des morphismes de
$\textit{WAdm}^{count}$ qui est stable par changement de base.
Soient $f : X \to Y$ et $g : Z \to Y$ des morphismes d'espaces algébriques formels
localement à indexation dénombrable sur $S$. Si $f$ satisfait aux conditions équivalentes du
Lemme \ref{formal-spaces-lemma-property-defines-property-morphisms},
alors il en est de même de $\text{pr}_2 : X \times_Y Z \to Z$.
\end{lemma}

\begin{proof}
Choisissons un recouvrement $\{Y_j \to Y\}$ comme dans la
Définition \ref{formal-spaces-definition-formal-algebraic-space}.
Pour tout $j$, choisissons un recouvrement $\{X_{ji} \to Y_j \times_Y X\}$
comme dans la Définition \ref{formal-spaces-definition-formal-algebraic-space}.
Pour tout $j$, choisissons un recouvrement $\{Z_{jk} \to Y_j \times_Y Z\}$
comme dans la Définition \ref{formal-spaces-definition-formal-algebraic-space}.
Remarquons que $X_{ji} \times_{Y_j} Z_{jk}$ est un
espace algébrique formel affine à indexation dénombrable ; voir le
Lemme \ref{formal-spaces-lemma-types-fibre-products}.
Nous voyons alors que
$$
\{X_{ji} \times_{Y_j} Z_{jk} \to X \times_Y Z\}
$$
est un recouvrement comme dans la Définition \ref{formal-spaces-definition-formal-algebraic-space}.
De plus, les morphismes $X_{ji} \times_{Y_j} Z_{jk} \to Z$
se factorisent par $Z_{jk}$. Par hypothèse,
$X_{ji} \to Y_j$ correspond à un morphisme $B_j \to A_{ji}$ de
$\text{WAdm}^{count}$ qui possède la propriété $P$.
Les morphismes $Z_{jk} \to Y_j$ correspondent à des morphismes $B_j \to C_{jk}$ dans
$\text{WAdm}^{count}$. Puisque
$X_{ji} \times_{Y_j} Z_{jk} =
\text{Spf}(A_{ji} \widehat{\otimes}_{B_j} C_{jk})$
par le Lemme \ref{formal-spaces-lemma-fibre-product-affines-over-separated},
il suffit de montrer que
$C_{jk} \to A_{ji} \widehat{\otimes}_{B_j} C_{jk}$
possède la propriété $P$, ce que garantit précisément
la stabilité de $P$ par changement de base.
\end{proof}

\begin{remark}[Variante adique*]
\label{formal-spaces-remark-base-change-variant-adic-star}
Soit $P$ une propriété locale des morphismes de $\textit{WAdm}^{adic*}$ ; voir la
Remarque \ref{formal-spaces-remark-variant-adic-star}. Nous disons que $P$ est {\it stable par
changement de base} si, étant donnés $B \to A$ et $B \to C$ dans $\textit{WAdm}^{adic*}$,
on a $P(B \to A) \Rightarrow P(C \to A \widehat{\otimes}_B C)$.
Cela a un sens puisque $A \widehat{\otimes}_B C$ est un objet de
$\textit{WAdm}^{adic*}$ d'après le Lemme \ref{formal-spaces-lemma-completed-tensor-product}.
De la même manière, nous obtenons une variante du
Lemme \ref{formal-spaces-lemma-base-change-property-morphisms}
pour les morphismes entre espaces algébriques formels localement adiques* sur $S$.
\end{remark}

\begin{remark}[Variante noethérienne]
\label{formal-spaces-remark-base-change-variant-Noetherian}
Soit $P$ une propriété locale des morphismes de $\textit{WAdm}^{Noeth}$ ; voir la
Remarque \ref{formal-spaces-remark-variant-Noetherian}. Nous disons que $P$ est
{\it stable par changement de base} si, étant donnés $B \to A$ et $B \to C$
dans $\textit{WAdm}^{Noeth}$, la propriété $P(B \to A)$
implique à la fois que $A \widehat{\otimes}_B C$ est adique
noethérien\footnote{Voir le Lemme \ref{formal-spaces-lemma-completed-tensor-product}
pour un critère.} et que
$P(C \to A \widehat{\otimes}_B C)$.
De la même manière, nous obtenons une variante du
Lemme \ref{formal-spaces-lemma-base-change-property-morphisms}
pour les morphismes entre espaces algébriques formels localement noethériens sur $S$.
\end{remark}

\begin{remark}[Autre variante noethérienne]
\label{formal-spaces-remark-base-change-variant-variant-Noetherian}
Soient $P$ et $Q$ des propriétés locales des morphismes de
$\textit{WAdm}^{Noeth}$ ; voir la Remarque \ref{formal-spaces-remark-variant-Noetherian}.
Nous disons que {\it $P$ est stable par changement de base par $Q$}
si, étant donnés $B \to A$ et $B \to C$ dans $\textit{WAdm}^{Noeth}$
qui satisfont à $P(B \to A)$ et $Q(B \to C)$, alors
$A \widehat{\otimes}_B C$ est adique noethérien et
$P(C \to A \widehat{\otimes}_B C)$ est vérifiée.
En raisonnant exactement comme dans la démonstration du
Lemme \ref{formal-spaces-lemma-base-change-property-morphisms},
nous obtenons l'assertion suivante :
étant donnés des morphismes $f : X \to Y$ et $g : Y \to Z$ d'espaces algébriques
formels localement noethériens sur $S$
tels que
\begin{enumerate}
\item les conditions équivalentes du
Lemme \ref{formal-spaces-lemma-property-defines-property-morphisms}
soient vérifiées pour $f$ et $P$ ;
\item les conditions équivalentes du
Lemme \ref{formal-spaces-lemma-property-defines-property-morphisms}
soient vérifiées pour $g$ et $Q$ ;
\end{enumerate}
alors les conditions équivalentes du
Lemme \ref{formal-spaces-lemma-property-defines-property-morphisms}
sont vérifiées pour $\text{pr}_2 : X \times_Y Z \to Z$ et $P$.
\end{remark}

\begin{situation}
\label{formal-spaces-situation-composition-local-property}
Soit $P$ une propriété locale des morphismes de $\textit{WAdm}^{count}$ ; voir la
Situation \ref{formal-spaces-situation-local-property}. Nous disons que $P$ est {\it stable par
composition} si, étant donnés $B \to A$ et $C \to B$ dans $\textit{WAdm}^{count}$,
on a $P(B \to A) \wedge P(C \to B) \Rightarrow P(C \to A)$.
\end{situation}

\begin{lemma}
\label{formal-spaces-lemma-composition-property-morphisms}
Soit $S$ un schéma. Soit $P$ une propriété locale des morphismes de
$\textit{WAdm}^{count}$ qui est stable par composition.
Soient $f : X \to Y$ et $g : Y \to Z$ des morphismes d'espaces algébriques formels
localement à indexation dénombrable sur $S$. Si $f$ et $g$
satisfont aux conditions équivalentes du
Lemme \ref{formal-spaces-lemma-property-defines-property-morphisms},
alors il en est de même de $g \circ f : X \to Z$.
\end{lemma}

\begin{proof}
Choisissons un recouvrement $\{Z_k \to Z\}$ comme dans la
Définition \ref{formal-spaces-definition-formal-algebraic-space}.
Pour tout $k$, choisissons un recouvrement $\{Y_{kj} \to Z_k \times_Z Y\}$
comme dans la Définition \ref{formal-spaces-definition-formal-algebraic-space}.
Pour tous $k$ et $j$, choisissons un recouvrement $\{X_{kji} \to Y_{kj} \times_Y X\}$
comme dans la Définition \ref{formal-spaces-definition-formal-algebraic-space}.
Si $f$ et $g$
satisfont aux conditions équivalentes du
Lemme \ref{formal-spaces-lemma-property-defines-property-morphisms},
alors $X_{kji} \to Y_{jk}$ et $Y_{jk} \to Z_k$
correspondent aux flèches
$B_{kj} \to A_{kji}$ et $C_k \to B_{kj}$ de
$\text{WAdm}^{count}$ qui possèdent la propriété $P$.
Il en est donc de même de leurs composées, ce qui conclut.
\end{proof}

\begin{remark}[Variante adique*]
\label{formal-spaces-remark-composition-variant-adic-star}
Soit $P$ une propriété locale des morphismes de $\textit{WAdm}^{adic*}$ ; voir la
Remarque \ref{formal-spaces-remark-variant-adic-star}. Nous disons que $P$ est {\it stable par
composition} si, étant donnés $B \to A$ et $C \to B$ dans $\textit{WAdm}^{adic*}$,
on a $P(B \to A) \wedge P(C \to B) \Rightarrow P(C \to A)$.
De la même manière, nous obtenons une variante du
Lemme \ref{formal-spaces-lemma-composition-property-morphisms}
pour les morphismes entre espaces algébriques formels localement adiques* sur $S$.
\end{remark}

\begin{remark}[Variante noethérienne]
\label{formal-spaces-remark-composition-variant-Noetherian}
Soit $P$ une propriété locale des morphismes de $\textit{WAdm}^{Noeth}$ ; voir la
Remarque \ref{formal-spaces-remark-variant-Noetherian}. Nous disons que $P$ est
{\it stable par composition} si, étant donnés $B \to A$ et $C \to B$
dans $\textit{WAdm}^{Noeth}$, on a
$P(B \to A) \wedge P(C \to B) \Rightarrow P(C \to A)$.
De la même manière, nous obtenons une variante du
Lemme \ref{formal-spaces-lemma-composition-property-morphisms}
pour les morphismes entre espaces algébriques formels localement noethériens sur $S$.
\end{remark}

\begin{situation}
\label{formal-spaces-situation-permanence-local-property}
Soit $P$ une propriété locale des morphismes de $\textit{WAdm}^{count}$ ; voir la
Situation \ref{formal-spaces-situation-local-property}. Nous disons que $P$
{\it possède la propriété de simplification}
si, étant donnés $B \to A$ et $C \to B$ dans $\textit{WAdm}^{count}$,
on a $P(C \to B) \wedge P(C \to A) \Rightarrow P(B \to A)$.
\end{situation}

\begin{lemma}
\label{formal-spaces-lemma-permanence-property-morphisms}
Soit $S$ un schéma. Soit $P$ une propriété locale des morphismes de
$\textit{WAdm}^{count}$ qui possède la propriété de simplification.
Soient $f : X \to Y$ et $g : Y \to Z$ des morphismes d'espaces algébriques formels
localement à indexation dénombrable sur $S$. Si $g \circ f$ et $g$
satisfont aux conditions équivalentes du
Lemme \ref{formal-spaces-lemma-property-defines-property-morphisms},
alors il en est de même de $f : X \to Y$.
\end{lemma}

\begin{proof}
Choisissons un recouvrement $\{Z_k \to Z\}$ comme dans la
Définition \ref{formal-spaces-definition-formal-algebraic-space}.
Pour tout $k$, choisissons un recouvrement $\{Y_{kj} \to Z_k \times_Z Y\}$
comme dans la Définition \ref{formal-spaces-definition-formal-algebraic-space}.
Pour tous $k$ et $j$, choisissons un recouvrement $\{X_{kji} \to Y_{kj} \times_Y X\}$
comme dans la Définition \ref{formal-spaces-definition-formal-algebraic-space}.
Faisons correspondre à $X_{kji} \to Y_{jk}$ et $Y_{jk} \to Z_k$
les flèches
$B_{kj} \to A_{kji}$ et $C_k \to B_{kj}$ de
$\text{WAdm}^{count}$.
Si $g \circ f$ et $g$ satisfont aux conditions équivalentes du
Lemme \ref{formal-spaces-lemma-property-defines-property-morphisms},
alors $C_k \to B_{kj}$ et $C_k \to A_{kji}$ satisfont à $P$.
Ainsi, $B_{kj} \to A_{kji}$ y satisfait aussi, ce qui conclut.
\end{proof}

\begin{remark}[Variante adique*]
\label{formal-spaces-remark-permanence-variant-adic-star}
Soit $P$ une propriété locale des morphismes de $\textit{WAdm}^{adic*}$ ; voir la
Remarque \ref{formal-spaces-remark-variant-adic-star}. Nous disons que $P$ {\it possède la propriété de
simplification} si, étant donnés $B \to A$ et $C \to B$ dans $\textit{WAdm}^{adic*}$,
on a $P(C \to A) \wedge P(C \to B) \Rightarrow P(B \to A)$.
De la même manière, nous obtenons une variante du
Lemme \ref{formal-spaces-lemma-composition-property-morphisms}
pour les morphismes entre espaces algébriques formels localement adiques* sur $S$.
\end{remark}

\begin{remark}[Variante noethérienne]
\label{formal-spaces-remark-permanence-variant-Noetherian}
Soit $P$ une propriété locale des morphismes de $\textit{WAdm}^{Noeth}$ ; voir la
Remarque \ref{formal-spaces-remark-variant-Noetherian}. Nous disons que $P$
{\it possède la propriété de simplification} si, étant donnés $B \to A$ et $C \to B$
dans $\textit{WAdm}^{Noeth}$, on a
$P(C \to B) \wedge P(C \to A) \Rightarrow P(C \to B)$.
De la même manière, nous obtenons une variante du
Lemme \ref{formal-spaces-lemma-composition-property-morphisms}
pour les morphismes entre espaces algébriques formels localement noethériens sur $S$.
\end{remark}














\section{Homomorphismes d'anneaux tendus et représentabilité par des espaces algébriques}
\label{formal-spaces-section-taut}

\noindent
Dans cette section, nous montrons brièvement que les morphismes entre espaces algébriques
formels localement à indexation dénombrable correspondent, localement pour la topologie étale,
à des homomorphismes continus tendus entre anneaux topologiques faiblement
admissibles qui possèdent des systèmes fondamentaux dénombrables d'idéaux ouverts.
En fait, cela ressort assez clairement du
Lemme \ref{formal-spaces-lemma-representable-affine},
et nous invitons le lecteur à sauter cette section.

\begin{lemma}
\label{formal-spaces-lemma-representable-property-rings}
Soit $B \to A$ une flèche de $\textit{WAdm}^{count}$.
Les conditions suivantes sont équivalentes :
\begin{enumerate}
\item[(a)] $B \to A$ est tendu (Définition \ref{formal-spaces-definition-taut}) ;
\item[(b)] pour $B \supset J_1 \supset J_2 \supset J_3 \supset \ldots$
étant un système fondamental d'idéaux faibles de définition, il existe
un diagramme commutatif
$$
\xymatrix{
A \ar[r] & \ldots \ar[r] & A_3 \ar[r] & A_2 \ar[r] & A_1 \\
B \ar[r] \ar[u] & \ldots \ar[r] & B/J_3 \ar[r] \ar[u] &
B/J_2 \ar[r] \ar[u] & B/J_1 \ar[u]
}
$$
tel que $A_{n + 1}/J_nA_{n + 1} = A_n$ et $A = \lim A_n$
comme anneau topologique.
\end{enumerate}
De plus, ces conditions équivalentes définissent une propriété locale,
c'est-à-dire qu'elles satisfont aux axiomes (\ref{formal-spaces-item-axiom-1}), (\ref{formal-spaces-item-axiom-2}),
(\ref{formal-spaces-item-axiom-3}).
\end{lemma}

\begin{proof}
L'équivalence de (a) et (b) est immédiate. Nous donnerons ci-dessous une
démonstration algébrique des axiomes, mais il se trouve que nous les avons déjà
démontrés. En effet, d'après le Lemme \ref{formal-spaces-lemma-representable-affine},
les conditions équivalentes (a) et (b) signifient que le
morphisme correspondant d'espaces algébriques formels affines est représentable
par des espaces algébriques.
Puisque cette condition est « locale sur la source et le but pour la topologie étale » d'après le
Lemme \ref{formal-spaces-lemma-representable-by-algebraic-spaces-local},
nous obtenons immédiatement les axiomes (\ref{formal-spaces-item-axiom-1}), (\ref{formal-spaces-item-axiom-2}) et
(\ref{formal-spaces-item-axiom-3}).

\medskip\noindent
Démonstration algébrique directe de (\ref{formal-spaces-item-axiom-1}), (\ref{formal-spaces-item-axiom-2}),
(\ref{formal-spaces-item-axiom-3}). Donnons-nous un diagramme (\ref{formal-spaces-equation-localize}) comme dans la
Situation \ref{formal-spaces-situation-local-property}.
D'après l'Exemple \ref{formal-spaces-example-representable-morphism-from-completion},
les homomorphismes $A \to (A')^\wedge$ et $B \to (B')^\wedge$
satisfont à (a) et (b).

\medskip\noindent
Supposons (a) et (b) vérifiées pour $\varphi$. Soit $J \subset B$ un idéal faible
de définition. Alors l'adhérence de $JA$, resp.\ de $J(B')^\wedge$,
est un idéal faible de définition $I \subset A$, resp.\ $J' \subset (B')^\wedge$.
L'adhérence de $I(A')^\wedge$ est alors un idéal faible de définition
$I' \subset (A')^\wedge$. Un argument topologique montre que $I'$ est aussi
l'adhérence de $J(A')^\wedge$ et de $J'(A')^\wedge$.
Enfin, lorsque $J$ parcourt un système fondamental d'idéaux faibles de définition
de $B$, les idéaux $I$ et $I'$ en font autant dans $A$ et $(A')^\wedge$.
Il s'ensuit que (a) est vérifiée pour $\varphi'$. Cela démontre (\ref{formal-spaces-item-axiom-1}).

\medskip\noindent
Supposons $A \to A'$ fidèlement plat et (a) et (b) vérifiées pour $\varphi'$.
Soit $J \subset B$ un idéal faible de définition. En appliquant (a) et (b)
aux homomorphismes $B \to (B')^\wedge \to (A')^\wedge$, nous trouvons que
l'adhérence $I'$ de $J(A')^\wedge$ est un idéal faible de définition.
En particulier, $I'$ est ouvert, donc l'image réciproque de $I'$
dans $A$ est ouverte. Nous avons alors (explication ci-dessous)
\begin{align*}
A \cap I'
& =
A \cap \bigcap (J(A')^\wedge + \Ker((A')^\wedge \to A'/I_0A')) \\
& =
A \cap \bigcap \Ker((A')^\wedge \to A'/JA' + I_0 A') \\
& = \bigcap (JA + I_0)
\end{align*}
qui est l'adhérence de $JA$ d'après le Lemme \ref{formal-spaces-lemma-closed}.
Les intersections portent sur les idéaux faibles de définition $I_0 \subset A$.
La première égalité vient du fait qu'un système fondamental de voisinages de
$0$ dans $(A')^\wedge$ est formé des noyaux des homomorphismes $(A')^\wedge \to A'/I_0A'$.
La deuxième égalité est triviale. La troisième vient du fait que $A \to A'$
est fidèlement plat ; voir
Algèbre, lemme \ref{algebra-lemma-faithfully-flat-universally-injective}.
Ainsi, l'adhérence de $JA$ est ouverte. D'après le Lemme \ref{formal-spaces-lemma-topologically-nilpotent},
l'adhérence de $JA$
est un idéal faible de définition de $A$. Enfin, étant donné un idéal faible
de définition $I \subset A$, nous pouvons trouver $J$ tel que
$J(A')^\wedge$ soit contenu dans l'adhérence de $I(A')^\wedge$
d'après la propriété (a) pour $B \to (B')^\wedge$ et $\varphi'$.
Nous voyons donc que (a) est vérifiée pour $\varphi$. Cela démontre (\ref{formal-spaces-item-axiom-2}).

\medskip\noindent
Nous omettons la démonstration de (\ref{formal-spaces-item-axiom-3}).
\end{proof}

\begin{lemma}
\label{formal-spaces-lemma-representable-local-property}
Soit $S$ un schéma. Soit $f : X \to Y$ un morphisme d'espaces algébriques formels
localement à indexation dénombrable sur $S$.
Les conditions suivantes sont équivalentes :
\begin{enumerate}
\item pour tout diagramme commutatif
$$
\xymatrix{
U \ar[d] \ar[r] & V \ar[d] \\
X \ar[r] & Y
}
$$
où $U$ et $V$ sont des espaces algébriques formels affines, où $U \to X$ et $V \to Y$
sont représentables par des espaces algébriques et étales, le morphisme $U \to V$
correspond à un homomorphisme tendu $B \to A$ de $\textit{WAdm}^{count}$ ;
\item il existe un recouvrement $\{Y_j \to Y\}$ comme dans la
Définition \ref{formal-spaces-definition-formal-algebraic-space} et, pour tout $j$,
un recouvrement $\{X_{ji} \to Y_j \times_Y X\}$ comme dans la
Définition \ref{formal-spaces-definition-formal-algebraic-space},
tel que chaque $X_{ji} \to Y_j$ corresponde
à un homomorphisme tendu d'anneaux dans $\textit{WAdm}^{count}$ ;
\item il existe un recouvrement $\{X_i \to X\}$ comme dans la
Définition \ref{formal-spaces-definition-formal-algebraic-space}
et, pour tout $i$, une factorisation $X_i \to Y_i \to Y$ où $Y_i$
est un espace algébrique formel affine, où $Y_i \to Y$ est représentable
par des espaces algébriques et étale, et où $X_i \to Y_i$ correspond
à un homomorphisme tendu d'anneaux dans $\textit{WAdm}^{count}$ ;
\item $f$ est représentable par des espaces algébriques.
\end{enumerate}
\end{lemma}

\begin{proof}
Le fait qu'un homomorphisme de $\textit{WAdm}^{count}$ soit
« tendu » est une propriété locale d'après le
Lemme \ref{formal-spaces-lemma-representable-property-rings}.
Le Lemme \ref{formal-spaces-lemma-property-defines-property-morphisms}
nous dit donc exactement que (1), (2) et (3) sont équivalentes.
D'autre part, d'après le Lemme \ref{formal-spaces-lemma-representable-affine},
le caractère « tendu » des homomorphismes de $\textit{WAdm}^{count}$ correspond exactement à la
« représentabilité par des espaces algébriques » des morphismes correspondants
d'espaces algébriques formels affines à indexation dénombrable.
Ainsi, l'implication (1) $\Rightarrow$ (2) du
Lemme \ref{formal-spaces-lemma-representable-by-algebraic-spaces-local}
montre que (4) implique (1) du présent lemme.
De même, l'implication (4) $\Rightarrow$ (1) du
Lemme \ref{formal-spaces-lemma-representable-by-algebraic-spaces-local}
montre que (2) implique (4) du présent lemme.
\end{proof}









\section{Morphismes adiques}
\label{formal-spaces-section-adic}

\noindent
Cette section met en correspondance la notion parfois employée de
« morphisme adique » $f : X \to Y$ entre espaces algébriques formels localement adiques*
$X$ et $Y$, d'une part avec la représentabilité de $f$ par des
espaces algébriques et, d'autre part, avec notre notion d'
homomorphisme continu tendu d'anneaux. Rappelons d'abord que
le caractère tendu équivaut au caractère adique pour les anneaux adiques dont
l'idéal de définition est de type fini.

\begin{lemma}
\label{formal-spaces-lemma-adic-homomorphism}
Soient $A$ et $B$ des anneaux topologiques préadiques. Soit
$\varphi : A \to B$ un homomorphisme continu d'anneaux.
\begin{enumerate}
\item Si $\varphi$ est adique, alors $\varphi$ est tendu.
\item Si $B$ est complet, si $A$ possède un idéal de définition
de type fini et si $\varphi$ est tendu, alors $\varphi$ est adique.
\end{enumerate}
En particulier, les conditions « $\varphi$ est adique » et « $\varphi$ est tendu »
sont équivalentes dans la catégorie $\textit{WAdm}^{adic*}$.
\end{lemma}

\begin{proof}
La partie (1) est le Lemme \ref{formal-spaces-lemma-adic-taut}.
La partie (2) est le Lemme \ref{formal-spaces-lemma-taut-is-adic}.
La dernière assertion résulte de (1) et (2).
\end{proof}

\noindent
Soit $S$ un schéma. Soit $f : X \to Y$ un morphisme d'espaces algébriques
formels localement adiques* sur $S$.
D'après le Lemme \ref{formal-spaces-lemma-representable-local-property},
les conditions suivantes sont équivalentes :
\begin{enumerate}
\item $f$ est représentable par des espaces algébriques (autrement dit, les
conditions équivalentes du
Lemme \ref{formal-spaces-lemma-representable-by-algebraic-spaces-local} sont vérifiées) ;
\item pour tout diagramme commutatif
$$
\xymatrix{
U \ar[d] \ar[r] & V \ar[d] \\
X \ar[r] & Y
}
$$
où $U$ et $V$ sont des espaces algébriques formels affines, où $U \to X$ et $V \to Y$
sont représentables par des espaces algébriques et étales, le morphisme $U \to V$
correspond à un homomorphisme adique\footnote{De manière équivalente, tendu d'après le
Lemme \ref{formal-spaces-lemma-adic-homomorphism}.} dans $\textit{WAdm}^{adic*}$.
\end{enumerate}
Dans cette situation, nous dirons que {\it $f$ est un morphisme adique} (la définition
formelle figure ci-dessous). Cette notion ou terminologie ne sera définie et employée
que pour les morphismes entre espaces algébriques formels localement adiques*,
car, autrement, nous ne disposons pas de l'équivalence entre (1) et (2) ci-dessus.

\begin{definition}
\label{formal-spaces-definition-adic-morphism}
Soit $S$ un schéma. Soit $f : X \to Y$ un morphisme d'espaces algébriques formels
sur $S$. Supposons $X$ et $Y$ localement adiques*. Nous disons que $f$ est
un {\it morphisme adique} si $f$ est représentable par des espaces algébriques.
Voir la discussion précédente.
\end{definition}











\section{Morphismes de type fini}
\label{formal-spaces-section-finite-type}

\noindent
Compte tenu de l'organisation du projet Champs, la définition suivante
est véritablement celle qui convient, et les notions plus fortes devraient
porter un autre nom.

\begin{definition}
\label{formal-spaces-definition-finite-type}
Soit $S$ un schéma. Soit $f : Y \to X$ un morphisme d'espaces algébriques formels
sur $S$.
\begin{enumerate}
\item Nous disons que $f$ est {\it localement de type fini}
si $f$ est représentable par des espaces algébriques et est localement
de type fini au sens de
Amorçage, définition \ref{bootstrap-definition-property-transformation}.
\item Nous disons que $f$ est de {\it type fini} si $f$ est localement de type fini et
quasi-compact (Définition \ref{formal-spaces-definition-quasi-compact-morphism}).
\end{enumerate}
\end{definition}

\noindent
Nous étudierons la relation entre les morphismes de type fini de
certains espaces algébriques formels et les
homomorphismes continus d'anneaux $A \to B$ qui sont topologiquement de type fini
à la section \ref{formal-spaces-section-tft}.

\begin{lemma}
\label{formal-spaces-lemma-characterize-finite-type}
Soit $S$ un schéma. Soit $f : X \to Y$ un morphisme d'espaces algébriques formels
sur $S$. Les conditions suivantes sont équivalentes :
\begin{enumerate}
\item $f$ est de type fini ;
\item $f$ est représentable par des espaces algébriques et est de type fini au
sens de
Amorçage, définition \ref{bootstrap-definition-property-transformation}.
\end{enumerate}
\end{lemma}

\begin{proof}
Cela résulte de Amorçage, lemme
\ref{bootstrap-lemma-transformations-property-implication},
de l'implication « quasi-compact $+$ localement de type fini
$\Rightarrow$ de type fini » pour les morphismes d'espaces algébriques, et du
Lemme \ref{formal-spaces-lemma-quasi-compact-representable}.
\end{proof}

\begin{lemma}
\label{formal-spaces-lemma-composition-finite-type}
La composée de morphismes de type fini est de type fini.
Il en est de même de la propriété « localement de type fini ».
\end{lemma}

\begin{proof}
Voir Amorçage, lemme \ref{bootstrap-lemma-composition-transformation-property}
et utiliser Morphismes d'espaces algébriques, lemme
\ref{spaces-morphisms-lemma-composition-finite-type}.
\end{proof}

\begin{lemma}
\label{formal-spaces-lemma-base-change-finite-type}
Un changement de base d'un morphisme de type fini est de type fini.
Il en est de même de la propriété « localement de type fini ».
\end{lemma}

\begin{proof}
Voir Amorçage, lemme \ref{bootstrap-lemma-base-change-transformation-property}
et utiliser Morphismes d'espaces algébriques, lemme
\ref{spaces-morphisms-lemma-base-change-finite-type}.
\end{proof}

\begin{lemma}
\label{formal-spaces-lemma-permanence-finite-type}
Soit $S$ un schéma. Soient $f : X \to Y$ et $g : Y \to Z$ des morphismes
d'espaces algébriques formels sur $S$. Si $g \circ f : X \to Z$ est localement de
type fini, alors $f : X \to Y$ est localement de type fini.
\end{lemma}

\begin{proof}
Le Lemme \ref{formal-spaces-lemma-permanence-representable} montre que $f$ est
représentable par des espaces algébriques. Soit $T$ un schéma et soit
$T \to Z$ un morphisme. Nous pouvons alors appliquer
Morphismes d'espaces algébriques, lemme \ref{spaces-morphisms-lemma-permanence-finite-type}
aux morphismes
$T \times_Z X \to T \times_Z Y \to T$ d'espaces algébriques pour conclure.
\end{proof}

\noindent
La propriété d'être localement de type fini est locale sur la source et sur le but.

\begin{lemma}
\label{formal-spaces-lemma-finite-type-local}
Soit $S$ un schéma. Soit $f : X \to Y$ un morphisme d'espaces algébriques formels
sur $S$. Les conditions suivantes sont équivalentes :
\begin{enumerate}
\item le morphisme $f$ est localement de type fini ;
\item il existe un diagramme commutatif
$$
\xymatrix{
U \ar[d] \ar[r] & V \ar[d] \\
X \ar[r] & Y
}
$$
où $U$, $V$ sont des espaces algébriques formels, où les flèches verticales sont
représentables par des espaces algébriques et étales, où $U \to X$
est surjectif et où $U \to V$ est localement de type fini ;
\item pour tout diagramme commutatif
$$
\xymatrix{
U \ar[d] \ar[r] & V \ar[d] \\
X \ar[r] & Y
}
$$
où $U$, $V$ sont des espaces algébriques formels et où les flèches verticales
sont représentables par des espaces algébriques et étales, le morphisme
$U \to V$ est localement de type fini ;
\item il existe un recouvrement $\{Y_j \to Y\}$ comme dans la
Définition \ref{formal-spaces-definition-formal-algebraic-space}
et, pour tout $j$, un recouvrement $\{X_{ji} \to Y_j \times_Y X\}$ comme dans la
Définition \ref{formal-spaces-definition-formal-algebraic-space}, tels que
$X_{ji} \to Y_j$ soit localement de type fini pour tous $j$ et $i$ ;
\item il existe un recouvrement $\{X_i \to X\}$ comme dans la
Définition \ref{formal-spaces-definition-formal-algebraic-space}
et, pour tout $i$, une factorisation $X_i \to Y_i \to Y$ où $Y_i$
est un espace algébrique formel affine, où $Y_i \to Y$ est représentable
par des espaces algébriques et étale, telle que $X_i \to Y_i$ soit
localement de type fini ;
\item ajouter davantage ici.
\end{enumerate}
\end{lemma}

\begin{proof}
Dans chacun des cinq cas, le morphisme $f : X \to Y$ est représentable
par des espaces algébriques ; voir le
Lemme \ref{formal-spaces-lemma-representable-by-algebraic-spaces-local}.
Nous utiliserons ce fait ci-dessous sans autre mention.

\medskip\noindent
Il est clair que (1) implique (2), car nous pouvons prendre $U = X$ et $V = Y$.
Réciproquement, donnons-nous un diagramme comme en (2). Soit $T$ un schéma et
soit $T \to Y$ un morphisme. Nous pouvons alors considérer
$$
\xymatrix{
U \times_Y T \ar[d] \ar[r] & V \times_Y T \ar[d] \\
X \times_Y T \ar[r] & T
}
$$
Les flèches verticales sont étales et la flèche horizontale supérieure
est localement de type fini comme changement de base d'un tel morphisme.
Par conséquent, Morphismes d'espaces algébriques, lemme
\ref{spaces-morphisms-lemma-finite-type-local} montre que
$X \times_Y T \to T$ est localement de type fini. Autrement dit,
(1) est vérifiée.

\medskip\noindent
Supposons (1) vraie et considérons un diagramme comme en (3).
Alors $U \to Y$ est localement de type fini
(en tant que composée $U \to X \to Y$ ; voir
Amorçage, lemme \ref{bootstrap-lemma-composition-transformation-property}).
Soit $T$ un schéma et soit $T \to V$ un morphisme.
Alors la projection $T \times_V U \to T$ se factorise sous la forme
$$
T \times_V U = (T \times_Y U) \times_{(V \times_Y V)} V
\to T \times_Y U \to T
$$
La seconde flèche est localement de type fini (comme changement de base
de la composée $U \to X \to Y$), et
la première est le changement de base de la diagonale $V \to V \times_Y V$,
qui est localement de type fini d'après le
Lemme \ref{formal-spaces-lemma-diagonal-morphism-formal-algebraic-spaces}.

\medskip\noindent
Il est clair que (3) implique (2). Ainsi, (1) à (3) sont équivalentes.

\medskip\noindent
Remarquons que la condition de (4) a un sens, puisque le produit fibré
$Y_j \times_Y X$ est un espace algébrique formel d'après le
Lemme \ref{formal-spaces-lemma-fibre-products}.
Il est clair que (4) implique (5).

\medskip\noindent
Supposons $X_i \to Y_i \to Y$ comme en (5). Posons alors
$V = \coprod Y_i$ et $U = \coprod X_i$ pour voir que
(5) implique (2).

\medskip\noindent
Enfin, supposons (1) à (3) vraies.
Nous pouvons donc choisir un recouvrement quelconque $\{Y_j \to Y\}$ comme dans la
Définition \ref{formal-spaces-definition-formal-algebraic-space}
et, pour tout $j$, un recouvrement quelconque $\{X_{ji} \to Y_j \times_Y X\}$ comme dans la
Définition \ref{formal-spaces-definition-formal-algebraic-space}.
Alors $X_{ij} \to Y_j$ est localement de type fini d'après (3),
et nous voyons que (4) est vraie. Cela achève la démonstration.
\end{proof}

\begin{example}
\label{formal-spaces-example-finite-type-from-finite-type-ring-map}
Soit $S$ un schéma. Soit $A$ un anneau topologique faiblement admissible sur $S$.
Soit $A \to A'$ un homomorphisme d'anneaux de type fini. Alors
$$
(A')^\wedge = \lim_{I \subset A\ w.i.d.} A'/IA'
$$
est un anneau faiblement admissible, et le morphisme correspondant
$\text{Spf}((A')^\wedge) \to \text{Spf}(A)$ est représentable ;
voir l'Exemple \ref{formal-spaces-example-representable-morphism-from-completion}.
Si $T \to \text{Spf}(A)$ est un morphisme où $T$ est un schéma
quasi-compact, alors il se factorise par $\Spec(A/I)$ pour un certain idéal faible
de définition $I \subset A$ (Lemme \ref{formal-spaces-lemma-factor-through-thickening}).
Alors $T \times_{\text{Spf}(A)} \text{Spf}((A')^\wedge)$
est égal à $T \times_{\Spec(A/I)} \Spec(A'/IA')$, et
nous voyons que $\text{Spf}((A')^\wedge) \to \text{Spf}(A)$ est
de type fini.
\end{example}

\begin{lemma}
\label{formal-spaces-lemma-locally-finite-type-locally-noetherian}
Soit $S$ un schéma. Soit $f : X \to Y$ un morphisme d'espaces algébriques
formels sur $S$. Si $Y$ est localement noethérien et si
$f$ est localement de type fini, alors $X$ est localement noethérien.
\end{lemma}

\begin{proof}
Choisissons $\{Y_j \to Y\}$ et $\{X_{ij} \to Y_j \times_Y X\}$
comme dans le Lemme \ref{formal-spaces-lemma-finite-type-local}. Il résulte alors
du Lemme \ref{formal-spaces-lemma-property-goes-up-affine-morphism}
que chaque $X_{ij}$ est noethérien. Cela démontre le lemme.
\end{proof}

\begin{lemma}
\label{formal-spaces-lemma-fibre-product-Noetherian}
Soit $S$ un schéma. Soient $f : X \to Y$ et $Z \to Y$ des morphismes
d'espaces algébriques formels sur $S$. Si $Z$ est localement
noethérien et si $f$ est localement de type fini, alors
$Z \times_Y X$ est localement noethérien.
\end{lemma}

\begin{proof}
Le morphisme $Z \times_Y X \to Z$ est localement de type fini d'après le
Lemme \ref{formal-spaces-lemma-base-change-finite-type}. Le résultat découle donc
du Lemme \ref{formal-spaces-lemma-locally-finite-type-locally-noetherian}.
\end{proof}








\section{Morphismes surjectifs}
\label{formal-spaces-section-surjective}

\noindent
D'après le Lemme \ref{formal-spaces-lemma-reduction-surjective}, la définition suivante n'entre
pas en conflit avec les définitions existantes pour les morphismes d'espaces
algébriques ou pour les morphismes d'espaces algébriques formels représentables
par des espaces algébriques.

\begin{definition}
\label{formal-spaces-definition-surjective}
Soit $S$ un schéma. Un morphisme $f : X \to Y$ d'espaces algébriques formels
sur $S$ est dit {\it surjectif} s'il induit un morphisme surjectif
$X_{red} \to Y_{red}$ sur les espaces algébriques réduits sous-jacents.
\end{definition}

\begin{lemma}
\label{formal-spaces-lemma-composition-surjective}
La composée de deux morphismes surjectifs est un morphisme surjectif.
\end{lemma}

\begin{proof}
Démonstration omise.
\end{proof}

\begin{lemma}
\label{formal-spaces-lemma-base-change-surjective}
Un changement de base d'un morphisme surjectif est un morphisme surjectif.
\end{lemma}

\begin{proof}
Démonstration omise.
\end{proof}

\begin{lemma}
\label{formal-spaces-lemma-characterize-surjective}
Soit $S$ un schéma. Soit $f : X \to Y$ un morphisme d'espaces algébriques formels
sur $S$. Les conditions suivantes sont équivalentes :
\begin{enumerate}
\item $f$ est surjectif ;
\item pour tout schéma $T$ et tout morphisme $T \to Y$,
la projection $X \times_Y T \to T$ est un morphisme surjectif
d'espaces algébriques formels ;
\item pour tout schéma affine $T$ et tout morphisme $T \to Y$,
la projection $X \times_Y T \to T$ est un morphisme surjectif d'espaces algébriques
formels ;
\item il existe un recouvrement $\{Y_j \to Y\}$ comme dans la
Définition \ref{formal-spaces-definition-formal-algebraic-space},
tel que chaque $X \times_Y Y_j \to Y_j$ soit un morphisme surjectif
d'espaces algébriques formels ;
\item il existe un morphisme surjectif $Z \to Y$
d'espaces algébriques formels tel que $X \times_Y Z \to Z$ soit surjectif ;
\item ajouter davantage ici.
\end{enumerate}
\end{lemma}

\begin{proof}
Démonstration omise.
\end{proof}







\section{Monomorphismes}
\label{formal-spaces-section-monomorphisms}

\noindent
Voici la définition.

\begin{definition}
\label{formal-spaces-definition-monomorphism}
Soit $S$ un schéma.
Un morphisme d'espaces algébriques formels sur $S$ est appelé un
{\it monomorphisme} s'il est un morphisme injectif de faisceaux.
\end{definition}

\noindent
Voici un exemple. Soit $X$ un espace algébrique et soit
$T \subset |X|$ une partie fermée. Alors le morphisme
$X_{/T} \to X$ du complété formel de $X$ le long de $T$ vers $X$
est un monomorphisme. En particulier, les monomorphismes d'espaces algébriques
formels ne sont généralement pas représentables.

\begin{lemma}
\label{formal-spaces-lemma-composition-monomorphism}
La composée de deux monomorphismes est un monomorphisme.
\end{lemma}

\begin{proof}
Démonstration omise.
\end{proof}

\begin{lemma}
\label{formal-spaces-lemma-base-change-monomorphism}
Un changement de base d'un monomorphisme est un monomorphisme.
\end{lemma}

\begin{proof}
Démonstration omise.
\end{proof}

\begin{lemma}
\label{formal-spaces-lemma-characterize-monomorphisms}
Soit $S$ un schéma. Soit $f : X \to Y$ un morphisme d'espaces algébriques formels
sur $S$. Les conditions suivantes sont équivalentes :
\begin{enumerate}
\item $f$ est un monomorphisme ;
\item pour tout schéma $T$ et tout morphisme $T \to Y$,
la projection $X \times_Y T \to T$ est un monomorphisme
d'espaces algébriques formels ;
\item pour tout schéma affine $T$ et tout morphisme $T \to Y$,
la projection $X \times_Y T \to T$ est un monomorphisme d'espaces algébriques
formels ;
\item il existe un recouvrement $\{Y_j \to Y\}$ comme dans la
Définition \ref{formal-spaces-definition-formal-algebraic-space},
tel que chaque $X \times_Y Y_j \to Y_j$ soit un monomorphisme
d'espaces algébriques formels ;
\item il existe une famille de morphismes $\{Y_j \to Y\}$ telle
que $\coprod Y_j \to Y$ soit une surjection de faisceaux sur
$(\Sch/S)_{fppf}$ et que chaque $X \times_Y Y_j \to Y_j$ soit un
monomorphisme pour tout $j$ ;
\item il existe un morphisme $Z \to Y$ d'espaces algébriques formels
qui est représentable par des espaces algébriques, surjectif, plat et localement
de présentation finie, tel que $X \times_Y Z \to X$ soit un monomorphisme ;
\item ajouter davantage ici.
\end{enumerate}
\end{lemma}

\begin{proof}
Démonstration omise.
\end{proof}





\section{Immersions fermées}
\label{formal-spaces-section-closed-immersions}

\noindent
Voici la définition.

\begin{definition}
\label{formal-spaces-definition-closed-immersion}
Soit $S$ un schéma. Soit $f : Y \to X$ un morphisme d'espaces algébriques formels
sur $S$. Nous disons que $f$ est une {\it immersion fermée}
si $f$ est représentable par des espaces algébriques et est une immersion fermée
au sens de
Amorçage, définition \ref{bootstrap-definition-property-transformation}.
\end{definition}

\noindent
Lors de la lecture de cette section, veuillez sauter les lemmes préliminaires obligés.

\begin{lemma}
\label{formal-spaces-lemma-composition-closed-immersion}
La composée de deux immersions fermées est une immersion fermée.
\end{lemma}

\begin{proof}
Démonstration omise.
\end{proof}

\begin{lemma}
\label{formal-spaces-lemma-base-change-closed-immersion}
Un changement de base d'une immersion fermée est une immersion fermée.
\end{lemma}

\begin{proof}
Démonstration omise.
\end{proof}

\begin{lemma}
\label{formal-spaces-lemma-characterize-closed-immersion}
Soit $S$ un schéma. Soit $f : X \to Y$ un morphisme d'espaces algébriques formels
sur $S$. Les conditions suivantes sont équivalentes :
\begin{enumerate}
\item $f$ est une immersion fermée ;
\item pour tout schéma $T$ et tout morphisme $T \to Y$, la projection
$X \times_Y T \to T$ est une immersion fermée ;
\item pour tout schéma affine $T$ et tout morphisme $T \to Y$,
la projection $X \times_Y T \to T$ est une immersion fermée ;
\item il existe un recouvrement $\{Y_j \to Y\}$ comme dans la
Définition \ref{formal-spaces-definition-formal-algebraic-space},
tel que chaque $X \times_Y Y_j \to Y_j$ soit une immersion fermée ;
\item il existe un morphisme $Z \to Y$ d'espaces algébriques formels
qui est représentable par des espaces algébriques, surjectif, plat et localement
de présentation finie, tel que $X \times_Y Z \to X$ soit une
immersion fermée ;
\item ajouter davantage ici.
\end{enumerate}
\end{lemma}

\begin{proof}
Démonstration omise.
\end{proof}

\begin{lemma}
\label{formal-spaces-lemma-closed-immersion-into-McQuillan}
Soit $S$ un schéma. Soit $X$ un espace algébrique formel affine de McQuillan
sur $S$. Soit $f : Y \to X$ une immersion fermée d'espaces algébriques formels
sur $S$. Alors $Y$ est un espace algébrique formel affine de McQuillan et
$f$ correspond à un homomorphisme continu $A \to B$ de $S$-algèbres topologiques
faiblement admissibles qui est tendu, dont le noyau est fermé et dont
l'image est dense.
\end{lemma}

\begin{proof}
Écrivons $X = \text{Spf}(A)$, où $A$ est un anneau topologique faiblement admissible.
Soit $I_\lambda$ un système fondamental d'idéaux faiblement admissibles
de définition de $A$. Alors $Y \times_X \Spec(A/I_\lambda)$ est
un sous-schéma fermé de $\Spec(A/I_\lambda)$,
donc affine (Définition \ref{formal-spaces-definition-closed-immersion}).
Écrivons $Y \times_X \Spec(A/I_\lambda) = \Spec(B_\lambda)$.
L'homomorphisme d'anneaux $A/I_\lambda \to B_\lambda$
est surjectif. Par conséquent, les projections
$$
B = \lim B_\lambda \longrightarrow B_\lambda
$$
sont surjectives, puisque les composées $A \to B \to B_\lambda$ sont surjectives.
Il s'ensuit que $Y$ est de McQuillan d'après le
Lemme \ref{formal-spaces-lemma-mcquillan-affine-formal-algebraic-space}.
L'homomorphisme d'anneaux $A \to B$ est tendu d'après le Lemme \ref{formal-spaces-lemma-representable-affine}.
Le noyau est fermé parce que $B$ est complet et que $A \to B$ est
continu. Enfin, comme $A \to B_\lambda$ est surjectif pour tout $\lambda$,
nous voyons que l'image de $A$ dans $B$ est dense.
\end{proof}

\noindent
Malgré le résultat précédent, nous ne savons pas, en général, comment
se comportent les immersions fermées lorsque le but est un
espace algébrique formel affine de McQuillan ; voir la Remarque \ref{formal-spaces-remark-questions}.

\begin{example}
\label{formal-spaces-example-closed-immersion-from-quotient}
Soit $S$ un schéma. Soit $A$ un anneau topologique faiblement admissible sur $S$.
Soit $K \subset A$ un idéal fermé. En posant
$$
B = (A/K)^\wedge = \lim_{I \subset A\ w.i.d.} A/(I + K)
$$
le morphisme $\text{Spf}(B) \to \text{Spf}(A)$ est représentable ; voir
l'Exemple \ref{formal-spaces-example-representable-morphism-from-completion}.
Si $T \to \text{Spf}(A)$ est un morphisme où $T$ est un schéma
quasi-compact, alors il se factorise par $\Spec(A/I)$ pour un certain idéal faible
de définition $I \subset A$ (Lemme \ref{formal-spaces-lemma-factor-through-thickening}).
Alors $T \times_{\text{Spf}(A)} \text{Spf}(B)$
est égal à $T \times_{\Spec(A/I)} \Spec(A/(K + I))$, et
nous voyons que $\text{Spf}(B) \to \text{Spf}(A)$ est une immersion fermée.
Le noyau de $A \to B$ est $K$, puisque $K$ est fermé, mais
il faut prendre garde qu'en général l'homomorphisme d'anneaux $A \to B = (A/K)^\wedge$
n'est pas nécessairement surjectif.
\end{example}

\begin{lemma}
\label{formal-spaces-lemma-monomorphism-iso-over-red}
Soit $S$ un schéma. Soit $f : X \to Y$ un morphisme d'espaces algébriques
formels. Supposons que
\begin{enumerate}
\item $f$ soit représentable par des espaces algébriques ;
\item $f$ soit un monomorphisme ;
\item l'inclusion $Y_{red} \to Y$ se factorise par $f$ ;
\item $f$ soit localement de type fini ou $Y$ soit localement noethérien.
\end{enumerate}
Alors $f$ est une immersion fermée.
\end{lemma}

\begin{proof}
Les hypothèses (2) et (3) impliquent que
$X_{red} = X \times_Y Y_{red} = Y_{red}$.
Nous utiliserons ce fait sans autre mention.

\medskip\noindent
Si $Y' \to Y$ est un morphisme étale d'espaces algébriques formels sur $S$,
alors le changement de base $f' : X \times_Y Y' \to Y'$ satisfait aux conditions
(1) à (4). D'après le Lemme \ref{formal-spaces-lemma-characterize-closed-immersion},
nous pouvons donc supposer que $Y$ est un espace algébrique formel affine.

\medskip\noindent
Écrivons $Y = \colim_{\lambda \in \Lambda} Y_\lambda$ comme dans la
Définition \ref{formal-spaces-definition-affine-formal-algebraic-space}.
Alors $X_\lambda = X \times_Y Y_\lambda$ est un espace algébrique
muni d'un monomorphisme $f_\lambda : X_\lambda \to Y_\lambda$
qui induit un isomorphisme $X_{\lambda, red} \to Y_{\lambda, red}$.
Ainsi, $X_\lambda$ est un schéma affine d'après
Limits of Espaces, Proposition \ref{spaces-limits-proposition-affine}
(puisque $X_{\lambda, red} \to X_\lambda$ est surjectif et entier).
Pour achever la démonstration, il suffit de montrer que
$X_\lambda \to Y_\lambda$ est une immersion fermée,
ce que nous ferons au paragraphe suivant.

\medskip\noindent
Soit $X \to Y$ un monomorphisme de schémas affines tel que
$X_{red} = X \times_Y Y_{red} = Y_{red}$. En général, cela
n'implique pas que $X \to Y$ soit une immersion fermée ; voir
Exemples, section \ref{examples-section-epimorphism-not-surjective}.
Cependant, sous notre hypothèse (4), nous savons qu'au paragraphe précédent
soit $X_\lambda \to Y_\lambda$ est de type fini, soit $Y_\lambda$
est noethérien. Cela signifie que
$X \to Y$ correspond à un homomorphisme d'anneaux $R \to A$ tel que
$R/I \to A/IA$ soit un isomorphisme, où $I \subset R$ est le
nilradical (c'est-à-dire l'idéal localement nilpotent maximal de $R$),
et que soit $R \to A$ est de type fini, soit $R$ est noethérien.
Dans le premier cas, $R \to A$ est surjectif d'après
Algèbre, lemme \ref{algebra-lemma-surjective-mod-locally-nilpotent},
et, dans le second, $I$ est de type fini, donc
nilpotent ; ainsi $R \to A$ est surjectif par le lemme de Nakayama ; voir
Algèbre, lemme \ref{algebra-lemma-NAK}, partie (11).
\end{proof}








\section{Séries restreintes}
\label{formal-spaces-section-restricted-power-series}

\noindent
Soit $A$ un anneau topologique complet pour une
topologie linéaire (Compléments d'algèbre, définition
\ref{more-algebra-definition-topological-ring}).
Soit $I_\lambda$ un système fondamental d'idéaux ouverts.
Soit $r \geq 0$ un entier. Dans ce contexte, on note souvent
par
$$
A\{x_1, \ldots, x_r\} =
\lim_\lambda A/I_\lambda[x_1, \ldots, x_r] =
\lim_\lambda (A[x_1, \ldots, x_r]/I_\lambda A[x_1, \ldots, x_r])
$$
muni de la topologie limite. Autrement dit, il s'agit
du complété de l'anneau de polynômes relativement aux
idéaux $I_\lambda$. On peut considérer les éléments de $A\{x_1, \ldots, x_r\}$ comme des
séries
$$
f = \sum\nolimits_{E = (e_1, \ldots, e_r)} a_E x_1^{e_1} \ldots x_r^{e_r}
$$
en $x_1, \ldots, x_r$ à coefficients $a_E \in A$ qui tendent
vers zéro dans la topologie de $A$. Autrement dit, pour tout $\lambda$,
tous les $a_E$, sauf un nombre fini, appartiennent à $I_\lambda$.
Pour cette raison, les éléments de $A\{x_1, \ldots, x_r\}$ sont parfois
appelés des {\it séries restreintes}.
On note parfois cet anneau $A\langle x_1, \ldots, x_r\rangle$ ;
nous nous abstiendrons d'employer cette notation.

\begin{remark}[Propriété universelle des séries restreintes]
\label{formal-spaces-remark-universal-property}
\begin{reference}
\cite[Chapter 0, 7.5.3]{EGA}
\end{reference}
Soit $A \to C$ un homomorphisme continu d'anneaux complets munis d'une topologie linéaire.
Alors tout homomorphisme de $A$-algèbres $A[x_1, \ldots x_r] \to C$ se prolonge de manière unique en un
homomorphisme continu $A\{x_1, \ldots, x_r\} \to C$ sur les séries restreintes.
\end{remark}

\begin{remark}
\label{formal-spaces-remark-I-adic-completion-and-restricted-power-series}
Soit $A$ un anneau et soit $I \subset A$ un idéal. Si $A$ est complet pour la topologie
$I$-adique, alors le complété $I$-adique $A[x_1, \ldots, x_r]^\wedge$ de
$A[x_1, \ldots, x_r]$ est, comme anneau, l'anneau des séries restreintes sur $A$.
Cependant, il n'est pas clair que $A[x_1, \ldots, x_r]^\wedge$ soit complet pour la topologie
$I$-adique. Nous considérons la topologie sur $A\{x_1, \ldots, x_r\}$
comme la topologie limite (qui est toujours complète), tandis que nous considérons souvent
la topologie sur $A[x_1, \ldots, x_r]^\wedge$ comme la topologie $I$-adique
(qui n'est pas toujours complète). Si $I$ est de type fini, alors
$A\{x_1, \ldots, x_r\} = A[x_1, \ldots, x_r]^\wedge$ comme anneaux topologiques ;
voir Algèbre, lemme \ref{algebra-lemma-hathat-finitely-generated}.
\end{remark}





\section{Algèbres topologiquement de type fini}
\label{formal-spaces-section-tft}

\noindent
Voici notre définition. Celle-ci ne fait pas l'objet d'un consensus général.
De nombreux auteurs imposent des conditions supplémentaires, souvent parce qu'ils ne
s'intéressent qu'à certains types d'anneaux et non au cas le plus général.

\begin{definition}
\label{formal-spaces-definition-topologically-finite-type}
Soit $A \to B$ un homomorphisme continu d'anneaux topologiques
(Compléments d'algèbre, définition \ref{more-algebra-definition-topological-ring}).
Nous disons que $B$ est {\it topologiquement de type fini sur} $A$ s'il
existe un homomorphisme de $A$-algèbres $A[x_1, \ldots, x_n] \to B$ dont
l'image est dense dans $B$.
\end{definition}

\noindent
Si $A$ est un anneau complet muni d'une topologie linéaire, alors l'anneau des
séries restreintes $A\{x_1, \ldots, x_r\}$ est topologiquement de type
fini sur $A$. Si $k$ est un corps, alors l'anneau de séries formelles
$k[[x_1, \ldots, x_r]]$ est topologiquement de type fini sur $k$.

\medskip\noindent
Pour les homomorphismes continus tendus d'anneaux topologiques faiblement admissibles,
le fait d'être topologiquement de type fini correspond exactement aux morphismes
de type fini entre les espaces algébriques formels affines associés.

\begin{lemma}
\label{formal-spaces-lemma-topologically-finite-type-finite-type}
Soit $S$ un schéma. Soit $\varphi : A \to B$ un homomorphisme continu d'anneaux
topologiques faiblement admissibles sur $S$. Les conditions suivantes
sont équivalentes :
\begin{enumerate}
\item $\text{Spf}(\varphi) : Y = \text{Spf}(B) \to \text{Spf}(A) = X$
est de type fini ;
\item $\varphi$ est tendu et $B$ est topologiquement de type fini sur $A$.
\end{enumerate}
\end{lemma}

\begin{proof}
Nous pouvons utiliser le Lemme \ref{formal-spaces-lemma-representable-affine} pour relier le caractère tendu de
$\varphi$ à la représentabilité de $\text{Spf}(\varphi)$. Nous l'emploierons
ci-dessous sans autre mention. Il s'ensuit que $X = \colim \Spec(A/I)$
et $Y = \colim \Spec(B/J(I))$, où $I \subset A$ parcourt les idéaux faibles
de définition de $A$ et où $J(I)$ est l'adhérence de $IB$ dans $B$.

\medskip\noindent
Supposons (2).
Choisissons un homomorphisme d'anneaux $A[x_1, \ldots, x_r] \to B$ d'image dense.
Alors $A[x_1, \ldots, x_r] \to B \to B/J(I)$ a lui aussi une image dense,
ce qui signifie qu'il est surjectif. Par conséquent, $B/J(I)$ est de
type fini sur $A/I$. Soit $T \to X$ un morphisme où
$T$ est un schéma quasi-compact. Alors $T \to X$ se factorise par
$\Spec(A/I)$ pour un certain $I$ (Lemme \ref{formal-spaces-lemma-factor-through-thickening}).
Alors $T \times_X Y = T \times_{\Spec(A/I)} \Spec(B/J(I))$ ; voir la démonstration du
Lemme \ref{formal-spaces-lemma-representable-affine}.
Par conséquent, $T \times_Y X \to T$ est de type fini comme changement de base du
morphisme $\Spec(B/J(I)) \to \Spec(A/I)$ qui est de type
fini. Ainsi, (1) est vraie.

\medskip\noindent
Supposons (1). Choisissons un $I \subset A$ quelconque comme ci-dessus. Puisque
$\Spec(A/I) \times_X Y = \Spec(B/J(I))$, nous voyons que $A/I \to B/J(I)$
est de type fini. Choisissons $b_1, \ldots, b_r \in B$
dont les images engendrent $B/J(I)$ sur $A/I$. Nous affirmons que l'image
de l'homomorphisme d'anneaux $A[x_1, \ldots, x_r] \to B$ qui envoie $x_i$ sur $b_i$
est dense. Pour le montrer, soit $I' \subset I$ un second idéal faible
de définition. Alors nous avons
$$
B/(J(I') + IB) = B/J(I)
$$
parce que $J(I)$ est l'adhérence de $IB$ et que $J(I')$ est ouvert.
Nous pouvons donc appliquer Algèbre, lemme
\ref{algebra-lemma-surjective-mod-locally-nilpotent}
pour voir que $(A/I')[x_1, \ldots, x_r] \to B/J(I')$ est surjectif.
Ainsi, (2) est vraie, ce qui achève la démonstration.
\end{proof}

\noindent
Soit $A$ un anneau topologique complet pour une topologie
linéaire. Soit $(I_\lambda)$ un système fondamental d'idéaux ouverts.
Soit $\mathcal{C}$ la catégorie des systèmes projectifs $(B_\lambda)$ où
\begin{enumerate}
\item $B_\lambda$ est une $A/I_\lambda$-algèbre de type fini ;
\item $B_\mu \to B_\lambda$ est un homomorphisme de $A/I_\mu$-algèbres
qui induit un isomorphisme $B_\mu/I_\lambda B_\mu \to B_\lambda$.
\end{enumerate}
Les morphismes de $\mathcal{C}$ sont donnés par des systèmes compatibles d'homomorphismes.

\begin{lemma}
\label{formal-spaces-lemma-category-affine-over}
Soit $S$ un schéma. Soit $X$ un espace algébrique formel affine sur $S$.
Supposons $X$ de McQuillan et soit $A$ l'anneau topologique faiblement
admissible associé à $X$. Il existe alors une antiéquivalence de catégories
entre
\begin{enumerate}
\item la catégorie $\mathcal{C}$ introduite ci-dessus ;
\item la catégorie des morphismes $Y \to X$ de type fini entre
espaces algébriques formels affines.
\end{enumerate}
\end{lemma}

\begin{proof}
Soit $(I_\lambda)$ un système fondamental d'idéaux faiblement admissibles
de définition de $A$. Considérons $Y$ comme en (2). Alors
$Y \times_X \Spec(A/I_\lambda)$ est affine
(Définition \ref{formal-spaces-definition-finite-type}
et Lemme \ref{formal-spaces-lemma-affine-representable-by-algebraic-spaces}).
Écrivons $Y \times_X \Spec(A/I_\lambda) = \Spec(B_\lambda)$.
L'homomorphisme d'anneaux $A/I_\lambda \to B_\lambda$ est de type fini
parce que $\Spec(B_\lambda) \to \Spec(A/I_\lambda)$ est de
type fini
(d'après la Définition \ref{formal-spaces-definition-finite-type}).
Alors $(B_\lambda)$ est un objet de $\mathcal{C}$.

\medskip\noindent
Réciproquement, étant donné un objet $(B_\lambda)$ de $\mathcal{C}$, nous pouvons poser
$Y = \colim \Spec(B_\lambda)$. C'est un espace algébrique formel
affine. Nous affirmons que
$$
Y \times_X \Spec(A/I_\lambda) =
\left(\colim_\mu \Spec(B_\mu)\right) \times_X \Spec(A/I_\lambda) =
\Spec(B_\lambda)
$$
Pour le montrer, il suffit d'obtenir les mêmes valeurs en évaluant
sur un schéma quasi-compact $U$. Un morphisme
$U \to \left(\colim_\mu \Spec(B_\mu)\right) \times_X \Spec(A/I_\lambda)$
provient d'un morphisme
$U \to \Spec(B_\mu) \times_{\Spec(A/I_\mu)} \Spec(A/I_\lambda)$
pour un certain $\mu \geq \lambda$ (appliquer
deux fois le Lemme \ref{formal-spaces-lemma-factor-through-thickening}). Puisque
$\Spec(B_\mu) \times_{\Spec(A/I_\mu)} \Spec(A/I_\lambda) = \Spec(B_\lambda)$
d'après notre seconde hypothèse sur les objets de $\mathcal{C}$,
cela démontre l'assertion. Nous pouvons alors montrer que le morphisme
$Y \to X$ est de type fini. En effet, remarquons que,
pour tout morphisme $U \to X$ où $U$ est un schéma quasi-compact, nous obtenons
une factorisation $U \to \Spec(A/I_\lambda) \to X$ pour un certain $\lambda$
(voir le lemme cité ci-dessus). Par conséquent,
$$
Y \times_X U =
Y \times_X \Spec(A/I_\lambda)) \times_{\Spec(A/I_\lambda)} U =
\Spec(B_\lambda) \times_{\Spec(A/I_\lambda)} U
$$
est un schéma de type fini sur $U$, comme voulu. Ainsi, la construction
$(B_\lambda) \mapsto \colim \Spec(B_\lambda)$ fournit bien un foncteur
de la catégorie (1) vers la catégorie (2).

\medskip\noindent
Pour achever la démonstration, montrons que les constructions précédentes
définissent des foncteurs quasi-inverses entre les catégories (1) et (2).
Dans un sens, il faut montrer que
$$
\left(\colim_\mu \Spec(B_\mu)\right) \times_X \Spec(A/I_\lambda) =
\Spec(B_\lambda)
$$
pour tout objet $(B_\lambda)$ de la catégorie $\mathcal{C}$.
Nous l'avons démontré ci-dessus. Dans
l'autre sens, il faut montrer que
$$
Y = \colim (Y \times_X \Spec(A/I_\lambda))
$$
pour $Y$ dans la catégorie (2). Là encore, cela résulte de l'évaluation sur des
objets tests quasi-compacts et du fait que $X = \colim \Spec(A/I_\lambda)$.
\end{proof}

\begin{remark}
\label{formal-spaces-remark-questions}
Soit $A$ un anneau topologique faiblement admissible et soit $(I_\lambda)$
un système fondamental d'idéaux faibles de définition. Soit $X = \text{Spf}(A)$ ;
autrement dit, $X$ est un espace algébrique formel affine de McQuillan.
Soit $f : Y \to X$ un morphisme d'espaces algébriques formels affines.
En général, $Y$ n'est pas nécessairement de McQuillan. Plus précisément,
nous pouvons poser les questions suivantes :
\begin{enumerate}
\item Supposons que $f : Y \to X$ soit une immersion fermée. Alors
$Y$ est de McQuillan et $f$ correspond à un homomorphisme continu
$\varphi : A \to B$ d'anneaux topologiques faiblement admissibles
qui est tendu, dont le noyau $K \subset A$ est un idéal fermé et
dont l'image $\varphi(A)$ est dense dans $B$ ; voir le
Lemme \ref{formal-spaces-lemma-closed-immersion-into-McQuillan}.
Quelles conditions sur $A$ garantissent que $B = (A/K)^\wedge$ comme dans
l'Exemple \ref{formal-spaces-example-closed-immersion-from-quotient} ?
\item Quelles conditions sur $A$ garantissent que les immersions fermées
$f : Y \to X$ correspondent aux quotients $A/K$ de $A$ par des idéaux fermés,
autrement dit que l'homomorphisme continu correspondant $\varphi$ est surjectif
et ouvert ?
\item Supposons que $f : Y \to X$ soit de type fini. Nous obtenons alors
$Y = \colim \Spec(B_\lambda)$, où $(B_\lambda)$ est un objet de
$\mathcal{C}$ d'après le Lemme \ref{formal-spaces-lemma-category-affine-over}.
Dans ce cas, il existe bien un entier fixe $r$ tel
que $B_\lambda$ soit engendré par $r$ éléments sur $A/I_\lambda$ pour
tout $\lambda$ (l'argument est pour l'essentiel déjà donné dans la démonstration de
(1) $\Rightarrow$ (2) du
Lemme \ref{formal-spaces-lemma-topologically-finite-type-finite-type}).
Cependant, il n'est pas clair que les projections
$\lim B_\lambda \to B_\lambda$ soient surjectives ; autrement dit,
il n'est pas clair que $Y$ soit de McQuillan.
Existe-t-il un exemple où $Y$ n'est pas de McQuillan ?
\item Supposons que $f : Y \to X$ soit de type fini et que $Y$ soit de McQuillan.
Alors $f$ correspond à un homomorphisme continu $\varphi : A \to B$ d'anneaux topologiques
faiblement admissibles. En fait, $\varphi$ est tendu et
$B$ est topologiquement de type fini sur $A$ ; voir le
Lemme \ref{formal-spaces-lemma-topologically-finite-type-finite-type}.
Autrement dit, $f$ se factorise sous la forme
$$
Y \longrightarrow \mathbf{A}^r_X \longrightarrow X
$$
où la première flèche est une immersion fermée d'espaces algébriques formels
affines de McQuillan. Cependant, les questions (1) et
(2) se posent alors pour $Y \to \mathbf{A}^r_X$.
\end{enumerate}
Nous répondrons ci-dessous à ces questions lorsque
$X$ est à indexation dénombrable, c'est-à-dire lorsque $A$ possède un système fondamental
dénombrable d'idéaux ouverts. Si vous connaissez des réponses à ces questions
dans une plus grande généralité, ou des contre-exemples, veuillez écrire à
\href{mailto:stacks.project@gmail.com}{stacks.project@gmail.com}.
\end{remark}

\begin{lemma}
\label{formal-spaces-lemma-closed-immersion-into-countably-indexed}
Soit $S$ un schéma. Soit $X$ un espace algébrique formel affine à indexation
dénombrable sur $S$. Soit $f : Y \to X$ une immersion fermée d'espaces algébriques
formels sur $S$. Alors $Y$ est un espace algébrique formel affine à indexation dénombrable
et $f$ correspond à $A \to A/K$, où $A$ est un objet de
$\textit{WAdm}^{count}$
(section \ref{formal-spaces-section-morphisms-rings})
et où $K \subset A$ est un idéal fermé.
\end{lemma}

\begin{proof}
D'après le Lemme \ref{formal-spaces-lemma-countably-indexed},
nous avons $X = \text{Spf}(A)$, où $A$ est un objet de
$\textit{WAdm}^{count}$. Comme une immersion fermée est représentable
et affine, le Lemme \ref{formal-spaces-lemma-property-goes-up-affine-morphism} montre
que $Y$ est un espace algébrique formel affine à indexation dénombrable.
En appliquant de nouveau le Lemme \ref{formal-spaces-lemma-countably-indexed},
nous voyons donc que $Y = \text{Spf}(B)$, avec $B$ objet de
$\textit{WAdm}^{count}$. D'après le Lemme \ref{formal-spaces-lemma-closed-immersion-into-McQuillan},
nous concluons que $f$ est donné par un morphisme $A \to B$ de
$\textit{WAdm}^{count}$ qui est tendu et d'image dense.
Pour achever la démonstration, nous appliquons le
Lemme \ref{formal-spaces-lemma-dense-image-surjective}.
\end{proof}

\begin{lemma}
\label{formal-spaces-lemma-quotient-restricted-power-series}
Soit $B \to A$ une flèche de $\textit{WAdm}^{count}$ ; voir la
section \ref{formal-spaces-section-morphisms-rings}.
Les conditions suivantes sont équivalentes :
\begin{enumerate}
\item[(a)] $B \to A$ est tendu et $B/J \to A/I$ est de type fini pour
tout idéal faible de définition $J \subset B$, où l'idéal $I \subset A$ est
l'adhérence de $JA$ ;
\item[(b)] $B \to A$ est tendu et $B/J_\lambda \to A/I_\lambda$
est de type fini pour un système cofinal $(J_\lambda)$
d'idéaux faibles de définition de $B$, où
$I_\lambda \subset A$ est l'adhérence de $J_\lambda A$ ;
\item[(c)] $B \to A$ est tendu et $A$ est topologiquement de type
fini sur $B$ ;
\item[(d)] $A$ est isomorphe, comme $B$-algèbre topologique, à un quotient de
$B\{x_1, \ldots, x_n\}$ par un idéal fermé.
\end{enumerate}
De plus, ces conditions équivalentes définissent une propriété locale,
c'est-à-dire qu'elles satisfont aux
axiomes (\ref{formal-spaces-item-axiom-1}), (\ref{formal-spaces-item-axiom-2}), (\ref{formal-spaces-item-axiom-3}).
\end{lemma}

\begin{proof}
Les implications (a) $\Rightarrow$ (b), (c) $\Rightarrow$ (a),
(d) $\Rightarrow$ (c) résultent immédiatement des définitions.
Supposons (b) vérifiée et soient $J \subset B$ et $I \subset A$ comme en (a).
Choisissons un diagramme commutatif
$$
\xymatrix{
A \ar[r] & \ldots \ar[r] & A_3 \ar[r] & A_2 \ar[r] & A_1 \\
B \ar[r] \ar[u] & \ldots \ar[r] & B/J_3 \ar[r] \ar[u] &
B/J_2 \ar[r] \ar[u] & B/J_1 \ar[u]
}
$$
tel que $A_{n + 1}/J_nA_{n + 1} = A_n$ et $A = \lim A_n$, comme dans le
Lemme \ref{formal-spaces-lemma-representable-property-rings}.
Pour tout $m$, il existe un $\lambda$ tel que $J_\lambda \subset J_m$.
Comme $B/J_\lambda \to A/I_\lambda$ est de type fini, il s'ensuit
que $B/J_m \to A/I_m$ est de type fini.
Soient $\alpha_1, \ldots, \alpha_n \in A_1$ des générateurs de $A_1$ sur
$B/J_1$. Puisque $A$ est la limite dénombrable d'un système à morphismes de
transition surjectifs, nous pouvons trouver $a_1, \ldots, a_n \in A$ dont les images sont
$\alpha_1, \ldots, \alpha_n$ dans $A_1$. D'après la
Remarque \ref{formal-spaces-remark-universal-property}, nous obtenons un homomorphisme continu
$B\{x_1, \ldots, x_n\} \to A$ qui envoie $x_i$ sur $a_i$. Cet homomorphisme
induit des surjections $(B/J_m)[x_1, \ldots, x_n] \to A_m$ d'après
Algèbre, lemme \ref{algebra-lemma-surjective-mod-locally-nilpotent}.
Pour $m \geq 1$, nous obtenons une suite exacte courte
$$
0 \to K_m \to (B/J_m)[x_1, \ldots, x_n] \to A_m \to 0
$$
Les morphismes de transition induits $K_{m + 1} \to K_m$ sont surjectifs parce que
$A_{m + 1}/J_mA_{m + 1} = A_m$. Par conséquent, la limite projective de ces
suites exactes courtes est exacte ; voir
Algèbre, lemme \ref{algebra-lemma-ML-exact-sequence}.
Puisque $B\{x_1, \ldots, x_n\} = \lim (B/J_m)[x_1, \ldots, x_n]$
et $A = \lim A_m$,
nous concluons que $B\{x_1, \ldots, x_n\} \to A$ est surjectif et ouvert.
Comme $A$ est complet, le noyau est un idéal fermé. Nous voyons ainsi que
(a), (b), (c) et (d) sont équivalentes.

\medskip\noindent
Donnons-nous un diagramme (\ref{formal-spaces-equation-localize}) comme dans la
Situation \ref{formal-spaces-situation-local-property}.
D'après l'Exemple \ref{formal-spaces-example-finite-type-from-finite-type-ring-map},
les homomorphismes $A \to (A')^\wedge$ et $B \to (B')^\wedge$
satisfont à (a), (b), (c) et (d). De plus, d'après le
Lemme \ref{formal-spaces-lemma-representable-property-rings},
pour démontrer les axiomes (\ref{formal-spaces-item-axiom-1}) et (\ref{formal-spaces-item-axiom-2}),
nous pouvons supposer que $B \to A$ et $(B')^\wedge \to (A')^\wedge$
sont tous deux tendus. Choisissons un idéal faible de définition $J \subset B$. Soient
$J' \subset (B')^\wedge$, $I \subset A$, $I' \subset (A')^\wedge$
les adhérences de $J(B')^\wedge$, $JA$, $J(A')^\wedge$.
D'après ce qui précède, il suffit de considérer le
diagramme commutatif
$$
\xymatrix{
A/I \ar[r] & (A')^\wedge/I' \\
B/J \ar[r] \ar[u]^{\overline{\varphi}} &
(B')^\wedge/J' \ar[u]_{\overline{\varphi}'}
}
$$
et de montrer : (1) $\overline{\varphi}$ est de type fini
$\Rightarrow \overline{\varphi}'$ est de type fini ; (2) si
$A \to A'$ est fidèlement plat, alors
$\overline{\varphi}'$ est de type fini $\Rightarrow \overline{\varphi}$
est de type fini. Remarquons que $(B')^\wedge/J' = B'/JB'$ et
$(A')^\wedge/I' = A'/IA'$ par construction des topologies sur
$(B')^\wedge$ et $(A')^\wedge$. En particulier, les homomorphismes horizontaux
du diagramme sont étales. La partie (1) résulte alors de
Algèbre, lemme \ref{algebra-lemma-compose-finite-type},
et la partie (2) de
Descente, lemme \ref{descent-lemma-finite-type-local-source-fppf-algebra},
puisque l'homomorphisme d'anneaux $A/I \to (A')^\wedge/I' = A'/IA'$ est fidèlement plat
et étale.

\medskip\noindent
Nous omettons la démonstration de l'axiome (\ref{formal-spaces-item-axiom-3}).
\end{proof}

\begin{lemma}
\label{formal-spaces-lemma-quotient-restricted-power-series-admissible}
Dans le Lemme \ref{formal-spaces-lemma-quotient-restricted-power-series},
si $B$ est admissible (par exemple adique), alors les conditions équivalentes
(a) à (d) sont aussi équivalentes à la condition suivante :
\begin{enumerate}
\item[(e)] $B \to A$ est tendu et $B/J \to A/I$ est de type fini pour
un certain idéal de définition $J \subset B$, où $I \subset A$ est
l'adhérence de $JA$.
\end{enumerate}
\end{lemma}

\begin{proof}
Il suffit de montrer que (e) implique (a). Soit $J' \subset B$ un idéal faible
de définition et soit $I' \subset A$ l'adhérence de $J'A$. Nous devons
montrer que $B/J' \to A/I'$ est de type fini. Si l'assertion correspondante
est vraie pour le plus petit idéal faible de définition $J'' = J' \cap J$, alors elle
est vraie pour $J'$. Nous pouvons donc supposer $J' \subset J$. Comme $J$ est un idéal
de définition (et non seulement un idéal faible de définition), nous avons
$J^n \subset J'$ pour un certain $n \geq 1$. Nous pouvons donc considérer le
diagramme
$$
\xymatrix{
0 \ar[r] & I/I' \ar[r] & A/I' \ar[r] & A/I \ar[r] & 0 \\
0 \ar[r] & J/J' \ar[r] \ar[u] & B/J' \ar[r] \ar[u] & B/J \ar[r] \ar[u] & 0
}
$$
à lignes exactes. Puisque $I' \subset A$ est ouvert et que
$I$ est l'adhérence de $J A$, nous voyons que $I/I' = (J/J') \cdot A/I'$.
Comme $J/J'$ est un idéal nilpotent et que $B/J \to A/I$ est de type fini,
Algèbre, lemme \ref{algebra-lemma-finite-type-mod-nilpotent} montre
que $A/I'$ est de type fini sur $B/J'$, comme voulu.
\end{proof}

\begin{lemma}
\label{formal-spaces-lemma-representable-affine-finite-type}
Soit $S$ un schéma. Soit $f : X \to Y$ un morphisme d'espaces algébriques
formels affines. Supposons $Y$ à indexation dénombrable.
Les conditions suivantes sont équivalentes :
\begin{enumerate}
\item $f$ est localement de type fini ;
\item $f$ est de type fini ;
\item $f$ correspond à un morphisme $B \to A$ de $\textit{WAdm}^{count}$
(section \ref{formal-spaces-section-morphisms-rings})
qui satisfait aux conditions équivalentes du
Lemme \ref{formal-spaces-lemma-quotient-restricted-power-series}.
\end{enumerate}
\end{lemma}

\begin{proof}
Puisque $X$ et $Y$ sont affines, il est clair que les conditions (1)
et (2) sont équivalentes. Dans les cas (1) et (2), le morphisme $f$
est représentable par des espaces algébriques par définition, donc
affine d'après le Lemme \ref{formal-spaces-lemma-affine-representable-by-algebraic-spaces}.
Ainsi, si (1) ou (2) est vérifiée, nous voyons que
$X$ est à indexation dénombrable d'après le
Lemme \ref{formal-spaces-lemma-property-goes-up-affine-morphism}.
Écrivons $X = \text{Spf}(A)$ et $Y = \text{Spf}(B)$
pour des $S$-algèbres topologiques $A$ et $B$ de $\textit{WAdm}^{count}$ ; voir le
Lemme \ref{formal-spaces-lemma-countably-indexed}. D'après le
Lemme \ref{formal-spaces-lemma-morphism-between-formal-spectra},
nous voyons que $f$ correspond à un homomorphisme continu $B \to A$.
Le résultat découle alors du
Lemme \ref{formal-spaces-lemma-topologically-finite-type-finite-type}.
\end{proof}

\begin{lemma}
\label{formal-spaces-lemma-finite-type-local-property}
Soit $S$ un schéma. Soit $f : X \to Y$ un morphisme d'espaces algébriques formels
localement à indexation dénombrable sur $S$.
Les conditions suivantes sont équivalentes :
\begin{enumerate}
\item pour tout diagramme commutatif
$$
\xymatrix{
U \ar[d] \ar[r] & V \ar[d] \\
X \ar[r] & Y
}
$$
où $U$ et $V$ sont des espaces algébriques formels affines, où $U \to X$ et $V \to Y$
sont représentables par des espaces algébriques et étales, le morphisme $U \to V$
correspond à un morphisme de $\textit{WAdm}^{count}$ qui est
tendu et topologiquement de type fini ;
\item il existe un recouvrement $\{Y_j \to Y\}$ comme dans la
Définition \ref{formal-spaces-definition-formal-algebraic-space} et, pour tout $j$,
un recouvrement $\{X_{ji} \to Y_j \times_Y X\}$ comme dans la
Définition \ref{formal-spaces-definition-formal-algebraic-space},
tel que chaque $X_{ji} \to Y_j$ corresponde
à un morphisme de $\textit{WAdm}^{count}$ qui est
tendu et topologiquement de type fini ;
\item il existe un recouvrement $\{X_i \to X\}$ comme dans la
Définition \ref{formal-spaces-definition-formal-algebraic-space}
et, pour tout $i$, une factorisation $X_i \to Y_i \to Y$ où $Y_i$
est un espace algébrique formel affine, où $Y_i \to Y$ est représentable
par des espaces algébriques et étale, et où $X_i \to Y_i$ correspond
à un morphisme de $\textit{WAdm}^{count}$ qui est
tendu et topologiquement de type fini ;
\item $f$ est localement de type fini.
\end{enumerate}
\end{lemma}

\begin{proof}
D'après le Lemme \ref{formal-spaces-lemma-quotient-restricted-power-series},
la propriété
$P(\varphi)=$ « $\varphi$ est tendu et topologiquement de type fini »
est locale sur $\text{WAdm}^{count}$. Par conséquent, le
Lemme \ref{formal-spaces-lemma-property-defines-property-morphisms}
montre que les conditions (1), (2) et (3) sont équivalentes.
D'autre part, d'après le Lemme \ref{formal-spaces-lemma-representable-affine-finite-type},
la condition $P$ sur les morphismes de $\textit{WAdm}^{count}$
correspond exactement au fait que les morphismes d'espaces algébriques formels
affines à indexation dénombrable soient localement de type fini.
Ainsi, l'implication (1) $\Rightarrow$ (3) du
Lemme \ref{formal-spaces-lemma-finite-type-local}
montre que (4) implique (1) du présent lemme.
De même, l'implication (4) $\Rightarrow$ (1) du
Lemme \ref{formal-spaces-lemma-finite-type-local}
montre que (2) implique (4) du présent lemme.
\end{proof}







\section{Axiomes de séparation pour les morphismes}
\label{formal-spaces-section-separation-axioms}

\noindent
Cette section est l'analogue de
Morphismes d'espaces algébriques, section \ref{spaces-morphisms-section-separation-axioms}
pour les morphismes d'espaces algébriques formels.

\begin{definition}
\label{formal-spaces-definition-separated-morphism}
Soit $S$ un schéma.
Soit $f : X \to Y$ un morphisme d'espaces algébriques formels sur $S$.
Soit $\Delta_{X/Y} : X \to X \times_Y X$ le morphisme diagonal.
\begin{enumerate}
\item Nous disons que $f$ est {\it séparé} si $\Delta_{X/Y}$ est une immersion fermée.
\item Nous disons que $f$ est {\it quasi-séparé} si $\Delta_{X/Y}$ est quasi-compact.
\end{enumerate}
\end{definition}

\noindent
Puisque $\Delta_{X/Y}$ est représentable (par des schémas) d'après le
Lemme \ref{formal-spaces-lemma-diagonal-morphism-formal-algebraic-spaces},
nous pouvons tester ces propriétés en considérant les morphismes $T \to X \times_Y X$
issus de schémas affines $T$ et en vérifiant si
$$
E = T \times_{X \times_Y X} X \longrightarrow T
$$
est quasi-compact ou une immersion fermée ; voir le
Lemme \ref{formal-spaces-lemma-quasi-compact-representable} ou la
Définition \ref{formal-spaces-definition-closed-immersion}.
Remarquons que le schéma $E$ est l'égalisateur de deux morphismes
$a, b : T \to X$ qui coïncident comme morphismes vers $Y$,
et que $E \to T$ est un monomorphisme localement de type fini.

\begin{lemma}
\label{formal-spaces-lemma-base-change-separated}
Tous les axiomes de séparation énumérés dans la
Définition \ref{formal-spaces-definition-separated-morphism}
sont stables par changement de base.
\end{lemma}

\begin{proof}
Soient $f : X \to Y$ et $Y' \to Y$ des morphismes d'espaces algébriques formels.
Soit $f' : X' \to Y'$ le changement de base de $f$ par $Y' \to Y$. Alors
$\Delta_{X'/Y'}$ est le changement de base de $\Delta_{X/Y}$ par
le morphisme $X' \times_{Y'} X' \to X \times_Y X$. Chacune des propriétés
de la diagonale employées dans la Définition \ref{formal-spaces-definition-separated-morphism}
est stable par changement de base. Le lemme est donc vrai.
\end{proof}

\begin{lemma}
\label{formal-spaces-lemma-fibre-product-after-map}
Soit $S$ un schéma. Soient $f : X \to Z$, $g : Y \to Z$ et $Z \to T$
des morphismes d'espaces algébriques formels sur $S$. Considérons le
morphisme induit $i : X \times_Z Y \to X \times_T Y$. Alors
\begin{enumerate}
\item $i$ est représentable (par des schémas), localement de type fini,
localement quasi-fini, séparé et est un monomorphisme ;
\item si $Z \to T$ est séparé, alors $i$ est une immersion fermée ;
\item si $Z \to T$ est quasi-séparé, alors $i$ est quasi-compact.
\end{enumerate}
\end{lemma}

\begin{proof}
Par la théorie générale des catégories, le diagramme suivant
$$
\xymatrix{
X \times_Z Y \ar[r]_i \ar[d] & X \times_T Y \ar[d] \\
Z \ar[r]^-{\Delta_{Z/T}} \ar[r] & Z \times_T Z
}
$$
est un diagramme de produit fibré. Ainsi, $i$ est le changement de base du
morphisme diagonal $\Delta_{Z/T}$. Le lemme résulte donc
du Lemme \ref{formal-spaces-lemma-diagonal-morphism-formal-algebraic-spaces}.
\end{proof}

\begin{lemma}
\label{formal-spaces-lemma-composition-separated}
Tous les axiomes de séparation énumérés dans la
Définition \ref{formal-spaces-definition-separated-morphism}
sont stables par composition des morphismes.
\end{lemma}

\begin{proof}
Soient $f : X \to Y$ et $g : Y \to Z$ des morphismes d'espaces algébriques formels
auxquels s'applique l'axiome considéré.
La diagonale $\Delta_{X/Z}$ est la composée
$$
X \longrightarrow X \times_Y X \longrightarrow X \times_Z X.
$$
Notre axiome de séparation est défini en exigeant que la diagonale
possède une certaine propriété $\mathcal{P}$. D'après le
Lemme \ref{formal-spaces-lemma-fibre-product-after-map} ci-dessus, nous voyons que
la seconde flèche possède aussi cette propriété. Le lemme en résulte,
puisque la composée de morphismes (représentables) possédant la propriété
$\mathcal{P}$ possède encore la propriété $\mathcal{P}$.
\end{proof}

\begin{lemma}
\label{formal-spaces-lemma-separated-local}
Soit $S$ un schéma. Soit $f : X \to Y$ un morphisme d'espaces algébriques formels
sur $S$. Soit $\mathcal{P}$ l'un quelconque des axiomes de séparation de la
Définition \ref{formal-spaces-definition-separated-morphism}.
Les conditions suivantes sont équivalentes :
\begin{enumerate}
\item $f$ possède la propriété $\mathcal{P}$ ;
\item pour tout schéma $Z$ et tout morphisme $Z \to Y$, le
changement de base $Z \times_Y X \to Z$ de $f$ possède la propriété $\mathcal{P}$ ;
\item pour tout schéma affine $Z$ et tout morphisme $Z \to Y$, le
changement de base $Z \times_Y X \to Z$ de $f$ possède la propriété $\mathcal{P}$ ;
\item pour tout schéma affine $Z$ et tout morphisme $Z \to Y$, le
produit $Z \times_Y X$ possède la propriété $\mathcal{P}$ comme espace algébrique formel (voir la
Définition \ref{formal-spaces-definition-separated}) ;
\item il existe un recouvrement $\{Y_j \to Y\}$ comme dans la
Définition \ref{formal-spaces-definition-formal-algebraic-space},
tel que le changement de base $Y_j \times_Y X \to Y_j$ possède
la propriété $\mathcal{P}$ pour tout $j$.
\end{enumerate}
\end{lemma}

\begin{proof}
Nous utiliserons à plusieurs reprises le
Lemme \ref{formal-spaces-lemma-base-change-separated}
sans autre mention. En particulier, il est clair que
(1) implique (2) et que (2) implique (3).

\medskip\noindent
Supposons (3) et soit $Z \to Y$ un morphisme où $Z$ est un schéma affine.
Soient $U$, $V$ des schémas affines et soient $a : U \to Z \times_Y X$
et $b : V \to Z \times_Y X$ des morphismes. Alors
$$
U \times_{Z \times_Y X} V =
(Z \times_Y X) \times_{\Delta, (Z \times_Y X) \times_Z (Z \times_Y X)}
(U \times_Z V)
$$
et nous voyons qu'il est quasi-compact si $\mathcal{P} =$ « quasi-séparé »,
ou qu'il s'agit d'un schéma affine muni d'une immersion fermée dans
$U \times_Z V$ si $\mathcal{P} =$ « séparé ». Ainsi, (4) est vérifiée.

\medskip\noindent
Supposons (4) et soit $Z \to Y$ un morphisme où $Z$ est un schéma affine.
Soient $U$, $V$ des schémas affines et soient $a : U \to Z \times_Y X$
et $b : V \to Z \times_Y X$ des morphismes. En lisant l'argument précédent
à rebours, nous voyons que $U \times_{Z \times_Y X} V \to U \times_Z V$
est quasi-compact si $\mathcal{P} =$ « quasi-séparé », ou une immersion
fermée si $\mathcal{P} =$ « séparé ». Puisque nous pouvons choisir $U$ et
$V$ comme ci-dessus de façon que $U$ parcoure un
recouvrement étale de $Z \times_Y X$, nous trouvons
que les morphismes correspondants
$$
U \times_Z V \to (Z \times_Y X) \times_Z (Z \times_Y X)
$$
forment un recouvrement étale par des affines. Nous concluons donc que
$\Delta : (Z \times_Y X) \to (Z \times_Y X) \times_Z (Z \times_Y X)$
est quasi-compact, resp.\ une immersion fermée. Ainsi, (3) est vérifiée.

\medskip\noindent
Montrons que (3) implique (5). Supposons (3) et soit
$\{Y_j \to Y\}$ comme dans la
Définition \ref{formal-spaces-definition-formal-algebraic-space}.
Nous devons montrer que les morphismes
$$
\Delta_j :
Y_j \times_Y X
\longrightarrow
(Y_j \times_Y X) \times_{Y_j} (Y_j \times_Y X) =
Y_j \times_Y X \times_Y X
$$
possèdent la propriété correspondante (c'est-à-dire sont quasi-compacts ou des immersions fermées).
Écrivons $Y_j = \colim Y_{j, \lambda}$ comme dans la
Définition \ref{formal-spaces-definition-affine-formal-algebraic-space}.
En remplaçant $Y_j$ par $Y_{j, \lambda}$ dans la formule ci-dessus, nous obtenons la
propriété par notre hypothèse (3). Puisque la flèche affichée
est la colimite des flèches $\Delta_{j, \lambda}$ et que nous
pouvons vérifier si $\Delta_j$ possède la propriété correspondante
après changement de base par des schémas affines qui s'envoient dans
$Y_j \times_Y X \times_Y X$, nous concluons par le
Lemme \ref{formal-spaces-lemma-factor-through-thickening}.

\medskip\noindent
Montrons que (5) implique (1). Soit $\{Y_j \to Y\}$ comme en (5).
Nous avons alors le diagramme de produit fibré
$$
\xymatrix{
\coprod Y_j \times_Y X \ar[r] \ar[d] &
X \ar[d] \\
\coprod Y_j \times_Y X \times_Y X \ar[r] &
X \times_Y X
}
$$
Par hypothèse, la flèche verticale de gauche est quasi-compacte ou une immersion fermée.
Il résulte de
Espaces, lemme \ref{spaces-lemma-descent-representable-transformations-property}
que la flèche verticale de droite est, elle aussi, quasi-compacte ou une
immersion fermée.
\end{proof}







\section{Morphismes propres}
\label{formal-spaces-section-proper}

\noindent
Voici la définition que nous utiliserons.

\begin{definition}
\label{formal-spaces-definition-proper}
Soit $S$ un schéma. Soit $f : Y \to X$ un morphisme d'espaces algébriques formels
sur $S$. Nous disons que $f$ est {\it propre}
si $f$ est représentable par des espaces algébriques et est propre au sens de
Amorçage, définition \ref{bootstrap-definition-property-transformation}.
\end{definition}

\noindent
Il résulte des définitions qu'un morphisme propre est de type fini.

\begin{lemma}
\label{formal-spaces-lemma-proper-local}
Soit $S$ un schéma. Soit $f : X \to Y$ un morphisme d'espaces algébriques formels
sur $S$. Les conditions suivantes sont équivalentes :
\begin{enumerate}
\item $f$ est propre ;
\item pour tout schéma $Z$ et tout morphisme $Z \to Y$, le
changement de base $Z \times_Y X \to Z$ de $f$ est propre ;
\item pour tout schéma affine $Z$ et tout morphisme $Z \to Y$, le
changement de base $Z \times_Y X \to Z$ de $f$ est propre ;
\item pour tout schéma affine $Z$ et tout morphisme $Z \to Y$, le
produit $Z \times_Y X$ est un espace algébrique propre sur $Z$ ;
\item il existe un recouvrement $\{Y_j \to Y\}$ comme dans la
Définition \ref{formal-spaces-definition-formal-algebraic-space},
tel que le changement de base $Y_j \times_Y X \to Y_j$ soit propre pour tout $j$.
\end{enumerate}
\end{lemma}

\begin{proof}
Démonstration omise.
\end{proof}

\begin{lemma}
\label{formal-spaces-lemma-base-change-proper}
Les morphismes propres d'espaces algébriques formels sont préservés par changement de base.
\end{lemma}

\begin{proof}
C'est une conséquence immédiate du Lemme \ref{formal-spaces-lemma-proper-local}
et de la transitivité du changement de base.
\end{proof}







\section{Espaces algébriques formels et recouvrements fpqc}
\label{formal-spaces-section-fpqc}

\noindent
Cette section est l'analogue de Propriétés des espaces algébriques, section
\ref{spaces-properties-section-fpqc}. Veuillez lire cette section
au préalable.

\begin{lemma}
\label{formal-spaces-lemma-sheaf-fpqc}
\begin{slogan}
Les espaces algébriques formels sont des faisceaux fpqc
\end{slogan}
Soit $S$ un schéma. Soit $X$ un espace algébrique formel sur $S$. Alors
$X$ vérifie la propriété de faisceau pour la topologie fpqc.
\end{lemma}

\begin{proof}
La démonstration est {\bf identique} à celle de
Propriétés des espaces algébriques, Proposition
\ref{spaces-properties-proposition-sheaf-fpqc}.
Puisque $X$ est un faisceau pour la topologie de Zariski, il
suffit de montrer ce qui suit. Étant donné un morphisme plat surjectif
entre schémas affines $f : T' \to T$, on a :
$X(T)$ est l'égalisateur des deux applications $X(T') \to X(T' \times_T T')$.
Voir Topologies, lemme \ref{topologies-lemma-sheaf-property-fpqc}.

\medskip\noindent
Soient $a, b : T \to X$ deux morphismes tels que $a \circ f = b \circ f$.
Nous devons montrer que $a = b$. Considérons le produit fibré
$$
E = X \times_{\Delta_{X/S}, X \times_S X, (a, b)} T.
$$
Par le Lemme \ref{formal-spaces-lemma-diagonal-formal-algebraic-space},
le morphisme $\Delta_{X/S}$ est un monomorphisme représentable. Ainsi,
$E \to T$ est un monomorphisme de schémas. Notre hypothèse selon laquelle
$a \circ f = b \circ f$ implique que $T' \to T$ se factorise (uniquement) par $E$.
Considérons le diagramme commutatif
$$
\xymatrix{
T' \times_T E \ar[r] \ar[d] & E \ar[d] \\
T' \ar[r] \ar@/^5ex/[u] \ar[ru] & T
}
$$
Puisque la projection $T' \times_T E \to T'$ est un monomorphisme
muni d'une section, nous en déduisons qu'elle est un isomorphisme. Ainsi,
$E \to T$ est un isomorphisme d'après
Descente, lemme \ref{descent-lemma-descending-property-isomorphism}.
Cela signifie que $a = b$, comme souhaité.

\medskip\noindent
Ensuite, soit $c : T' \to X$ un morphisme tel que les deux composées
$T' \times_T T' \to T' \to X$ coïncident. Nous devons trouver un morphisme
$a : T \to X$ dont la composée avec $T' \to T$ est $c$. Choisissons un
schéma formel affine $U$ et un morphisme \'etale $U \to X$ tels que l'image
de $|U| \to |X_{red}|$ contienne l'image de $|c| : |T'| \to |X_{red}|$.
C'est possible d'après la
Définition \ref{formal-spaces-definition-formal-algebraic-space},
Propriétés des espaces algébriques, lemme \ref{spaces-properties-lemma-topology-points},
le fait qu'une réunion finie d'espaces algébriques formels affines est un
espace algébrique formel affine, et le fait que $|T'|$ est quasi-compact
(petit argument omis). Le morphisme $U \to X$ est représentable
par des schémas (Lemme \ref{formal-spaces-lemma-presentation-representable}) et
séparé (Lemme \ref{formal-spaces-lemma-separated-from-separated}). Ainsi,
$$
V = U \times_{X, c} T' \longrightarrow T'
$$
est un morphisme \'etale et séparé de schémas. Il est également surjectif
par notre choix de $U \to X$ (si l'on ne souhaite pas le démontrer, on peut
remplacer $U$ par une réunion disjointe d'espaces algébriques formels affines afin que
$U \to X$ soit surjectif ; tout le reste fonctionne également). Le fait que
$c \circ \text{pr}_0 = c \circ \text{pr}_1$ signifie que nous obtenons une
donnée de descente sur $V/T'/T$
(Descente, définition \ref{descent-definition-descent-datum})
car
\begin{align*}
V \times_{T'} (T' \times_T T')
& =
U \times_{X, c \circ \text{pr}_0} (T' \times_T T') \\
& =
(T' \times_T T') \times_{c \circ \text{pr}_1, X} U \\
& =
(T' \times_T T') \times_{T'} V
\end{align*}
Le morphisme $V \to T'$ est ind-quasi-affine d'après
Compléments sur les morphismes, lemme
\ref{more-morphisms-lemma-etale-separated-ind-quasi-affine}
(car les morphismes \'etales sont localement quasi-finis, voir
Morphismes, lemme \ref{morphisms-lemma-etale-locally-quasi-finite}).
D'après More on Groupoïdes, lemme \ref{more-groupoids-lemma-ind-quasi-affine},
la donnée de descente est effective. Soit $W \to T$ un morphisme
tel qu'il existe un isomorphisme $\alpha : T' \times_T W \to V$
compatible avec la donnée de descente sur $V$ et la donnée de descente canonique
sur $T' \times_T W$. Alors $W \to T$ est surjectif et \'etale
(Descente, lemmes \ref{descent-lemma-descending-property-surjective} et
\ref{descent-lemma-descending-property-etale}).
Considérons la composée
$$
b' : T' \times_T W \longrightarrow V = U \times_{X, c} T' \longrightarrow U
$$
Les deux composées
$b' \circ (\text{pr}_0, 1), 
b' \circ (\text{pr}_1, 1) :
(T' \times_T T') \times_T W \to T' \times_T W \to U$
coïncident par notre choix de $\alpha$ et la propriété correspondante de $c$
(calcul omis). Ainsi, $b'$ descend en un morphisme $b : W \to U$ d'après
Descente, lemme \ref{descent-lemma-fpqc-universal-effective-epimorphisms}.
Le diagramme
$$
\xymatrix{
T' \times_T W \ar[r] \ar[d] & W \ar[r]_b & U \ar[d] \\
T' \ar[rr]^c &  & X
}
$$
est commutatif. Cela signifie que nous avons démontré l'existence
de $a$ localement pour la topologie \'etale sur $T$, i.e., nous avons un $a' : W \to X$.
Toutefois, puisque nous avons démontré l'unicité
dans le premier paragraphe, cette solution locale pour la topologie \'etale
satisfait la condition de recollement, i.e., nous avons
$\text{pr}_0^*a' = \text{pr}_1^*a'$ dans $X(W \times_T W)$.
Puisque $X$ est un faisceau \'etale, il existe un unique $a \in X(T)$ qui se restreint
à $a'$ sur $W$.
\end{proof}





\section{Morphismes issus de schémas formels affines}
\label{formal-spaces-section-map-out-of}

\noindent
Nous démontrons quelques résultats qui seront utiles plus tard.
Dans l'article \cite{Bhatt-Algebraize}, le lecteur trouvera
des résultats très généraux de même nature.

\begin{lemma}
\label{formal-spaces-lemma-map-into-affine}
Soit $S$ un schéma. Soit $A$ une algèbre topologique faiblement admissible
sur $S$. Soit $X$ un schéma affine sur $S$. Alors
l'application naturelle
$$
\Mor_S(\Spec(A), X)
\longrightarrow
\Mor_S(\text{Spf}(A), X)
$$
est bijective.
\end{lemma}

\begin{proof}
Si $X$ est affine, disons $X = \Spec(B)$, il résulte du
Lemme \ref{formal-spaces-lemma-morphism-between-formal-spectra}
que les morphismes $\text{Spf}(A) \to \Spec(B)$ correspondent aux homomorphismes continus
d'algèbres sur $S$, $B \to A$, où $B$ porte la topologie discrète.
Ce sont simplement les homomorphismes d'algèbres sur $S$, qui correspondent aux morphismes
$\Spec(A) \to \Spec(B)$.
\end{proof}

\begin{lemma}
\label{formal-spaces-lemma-map-into-scheme}
Soit $S$ un schéma. Soit $A$ une algèbre topologique faiblement admissible
sur $S$ telle que $A/I$ soit un anneau local pour un certain idéal faible
de définition $I \subset A$. Soit $X$ un schéma sur $S$. Alors
l'application naturelle
$$
\Mor_S(\Spec(A), X)
\longrightarrow
\Mor_S(\text{Spf}(A), X)
$$
est bijective.
\end{lemma}

\begin{proof}
Soit $\varphi : \text{Spf}(A) \to X$ un morphisme. Puisque $\Spec(A/I)$
est local, nous voyons que $\varphi$ envoie $\Spec(A/I)$ dans un
ouvert affine $U \subset X$. Cela implique alors que $\Spec(A/J)$
s'envoie dans $U$ pour tout idéal de définition $J$. Nous pouvons donc
appliquer le Lemme \ref{formal-spaces-lemma-map-into-affine} pour voir que $\varphi$ provient
d'un morphisme $\Spec(A) \to X$. Cela démontre la surjectivité de l'application.
Nous omettons la démonstration de l'injectivité.
\end{proof}

\begin{lemma}
\label{formal-spaces-lemma-map-into-algebraic-space}
Soit $S$ un schéma. Soit $R$ une algèbre locale noethérienne complète sur $S$.
Soit $X$ un espace algébrique sur $S$. Alors l'application naturelle
$$
\Mor_S(\Spec(R), X)
\longrightarrow
\Mor_S(\text{Spf}(R), X)
$$
est bijective.
\end{lemma}

\begin{proof}
Soit $\mathfrak m$ l'idéal maximal de $R$. Nous devons montrer que
$$
\Mor_S(\Spec(R), X) \longrightarrow \lim \Mor_S(\Spec(R/\mathfrak m^n), X)
$$
est bijective pour $R$ comme ci-dessus.

\medskip\noindent
Injectivité : soient $x, x' : \Spec(R) \to X$
deux morphismes ayant la même image dans le membre de droite.
Considérons le produit fibré
$$
T = \Spec(R) \times_{(x, x'), X \times_S X, \Delta} X
$$
Alors $T$ est un schéma et $T \to \Spec(R)$ est localement de type fini,
est un monomorphisme, est séparé et est localement quasi-fini ; voir
Morphismes d'espaces algébriques, lemme \ref{spaces-morphisms-lemma-properties-diagonal}.
En particulier, $T$ est localement noethérien ; voir
Morphismes, lemme \ref{morphisms-lemma-finite-type-noetherian}.
Soit $t \in T$ l'unique point s'envoyant sur le point fermé de $\Spec(R)$,
qui existe puisque $x$ et $x'$ coïncident sur $R/\mathfrak m$. Alors
$R \to \mathcal{O}_{T, t}$ est un homomorphisme local d'anneaux noethériens tel que
$R/\mathfrak m^n \to \mathcal{O}_{T, t}/\mathfrak m^n\mathcal{O}_{T, t}$
est un isomorphisme pour tout $n$ (car $x$ et $x'$ coïncident sur
$\Spec(R/\mathfrak m^n)$ pour tout $n$). Puisque $\mathcal{O}_{T, t}$
s'injecte dans son complété (voir
Algèbre, lemme \ref{algebra-lemma-intersect-powers-ideal-module-zero})
nous concluons que $R = \mathcal{O}_{T, t}$. Ainsi, $x$ et $x'$ coïncident
sur $R$.

\medskip\noindent
Surjectivité : soit $(x_n)$ un élément du membre de droite.
Choisissons un schéma $U$ et un morphisme \'etale surjectif $U \to X$. 
Notons $x_0 : \Spec(k) \to X$ le morphisme induit sur le corps résiduel
$k = R/\mathfrak m$. Le morphisme de schémas
$U \times_{X, x_0} \Spec(k) \to \Spec(k)$ est \'etale et surjectif.
Ainsi, $U \times_{X, x_0} \Spec(k)$ est une réunion disjointe non vide de spectres
d'extensions finies séparables de $k$ ; voir
Morphismes, lemme \ref{morphisms-lemma-etale-over-field}.
Nous pouvons donc trouver une extension finie séparable $k'/k$
et un $k'$-point $u_0 : \Spec(k') \to U$ tels que
$$
\xymatrix{
\Spec(k') \ar[d] \ar[r]_-{u_0} & U \ar[d] \\
\Spec(k) \ar[r]^-{x_0} & X
}
$$
soit commutatif. Soit $R \subset R'$ une extension \'etale finie d'anneaux locaux
noethériens complets qui induit $k'/k$ sur les corps résiduels
(voir Algèbre, lemmes \ref{algebra-lemma-henselian-cat-finite-etale} et
\ref{algebra-lemma-complete-henselian}). Notons $x'_n$ la restriction
de $x_n$ à $\Spec(R'/\mathfrak m^nR')$. D'après
Compléments sur les morphismes d'espaces algébriques, lemme
\ref{spaces-more-morphisms-lemma-etale-formally-etale}
on peut trouver un élément
$(u'_n) \in \lim \Mor_S(\Spec(R'/\mathfrak m^nR'), U)$
qui s'envoie sur $(x'_n)$. D'après le Lemme \ref{formal-spaces-lemma-map-into-scheme},
la famille $(u'_n)$ provient d'un unique
morphisme $u' : \Spec(R') \to U$. Notons $x' : \Spec(R') \to X$ la
composée. Remarquons que $R' \otimes_R R'$ est un produit fini de spectres d'anneaux
locaux noethériens complets auxquels notre discussion s'applique.
Ainsi, le diagramme
$$
\xymatrix{
\Spec(R' \otimes_R R') \ar[r] \ar[d] & \Spec(R') \ar[d]^{x'} \\
\Spec(R') \ar[r]^{x'} & X
}
$$
est commutatif par l'injectivité démontrée ci-dessus et le fait que
$x'_n$ est la restriction de $x_n$, qui est défini sur $R/\mathfrak m^n$.
Puisque $\{\Spec(R') \to \Spec(R)\}$ est un recouvrement fppf, nous concluons
que $x'$ descend en un morphisme $x : \Spec(R) \to X$.
Nous omettons la démonstration du fait que $x_n$ est la restriction de $x$ à
$\Spec(R/\mathfrak m^n)$.
\end{proof}

\begin{lemma}
\label{formal-spaces-lemma-adic-into-completion}
Soit $S$ un schéma. Soit $X$ un espace algébrique sur $S$.
Soit $T \subset |X|$ un sous-ensemble fermé tel que
$X \setminus T \to X$ soit quasi-compact. Soit $R$ une algèbre locale
noethérienne complète sur $S$. Alors un morphisme adique $p : \text{Spf}(R) \to X_{/T}$
correspond à un unique morphisme $g : \Spec(R) \to X$ tel
que $g^{-1}(T) = \{\mathfrak m_R\}$. 
\end{lemma}

\begin{proof}
L'énoncé a un sens parce que $X_{/T}$ est adique* d'après le
Lemme \ref{formal-spaces-lemma-formal-completion-types} (et que nous sommes donc
autorisés à employer le terme adique pour les morphismes ; voir la
Définition \ref{formal-spaces-definition-adic-morphism}).
Soit $p$ donné. D'après le Lemme \ref{formal-spaces-lemma-map-into-algebraic-space},
nous obtenons un unique morphisme $g : \Spec(R) \to X$ correspondant à
la composée $\text{Spf}(R) \to X_{/T} \to X$.
Soit $Z \subset X$ le sous-espace fermé induit réduit de support
$T$. Le morphisme d'inclusion $Z \to X$ correspond à un morphisme
$Z \to X_{/T}$. Puisque $p$ est adique, il est représentable par des espaces
algébriques et nous obtenons
$$
\text{Spf}(R) \times_{X_{/T}} Z = \text{Spf}(R) \times_X Z
$$
est un espace algébrique muni d'une immersion fermée dans $\text{Spf}(R)$.
(L'égalité vaut parce que $X_{/T} \to X$ est un monomorphisme.)
Ainsi, ce produit fibré est égal à $\Spec(R/J)$ pour un certain
idéal $J \subset R$ qui contient $\mathfrak m_R^{n_0}$ pour un certain
$n_0 \geq 1$. Cela implique que $\Spec(R) \times_X Z$
est un sous-schéma fermé de $\Spec(R)$,
disons $\Spec(R) \times_X Z = \Spec(R/I)$, dont l'intersection avec
$\Spec(R/\mathfrak m_R^n)$ pour $n \geq n_0$ est égale à $\Spec(R/J)$.
En termes algébriques, cela signifie que
$I + \mathfrak m_R^n = J + \mathfrak m_R^n = J$ pour tout $n \geq n_0$.
D'après le théorème d'intersection de Krull,
cela implique que $I = J$, ce qui conclut.
\end{proof}








\section{Le petit site \'etale d'un espace algébrique formel}
\label{formal-spaces-section-etale-sites}

\noindent
La motivation de la définition suivante vient des schémas
formels classiques : l'espace topologique sous-jacent d'un schéma formel
$(\mathfrak X, \mathcal{O}_\mathfrak X)$
est l'espace topologique sous-jacent de la réduction $\mathfrak X_{red}$.

\medskip\noindent
Une remarque importante est la suivante. Supposons que $X$ soit un espace algébrique
de réduction $X_{red}$ (Propriétés des espaces algébriques, définition
\ref{spaces-properties-definition-reduced-induced-space}).
Alors nous avons
$$
X_{spaces, \etale} = X_{red, spaces, \etale},\quad
X_\etale = X_{red, \etale},\quad
X_{affine, \etale} = X_{red, affine, \etale}
$$
d'après Compléments sur les morphismes d'espaces algébriques, théorème
\ref{spaces-more-morphisms-theorem-topological-invariance} et
lemme \ref{spaces-more-morphisms-lemma-topological-invariance}.
Par conséquent, la définition suivante n'entre pas en conflit avec la notion
existante lorsque notre espace algébrique formel se trouve être
un espace algébrique.

\begin{definition}
\label{formal-spaces-definition-etale-sites}
Soit $S$ un schéma. Soit $X$ un espace algébrique formel de
réduction $X_{red}$ (Lemme \ref{formal-spaces-lemma-reduction-formal-algebraic-space}).
\begin{enumerate}
\item Le {\it petit site \'etale} $X_\etale$ de $X$ est
le site $X_{red, \etale}$ de Propriétés des espaces algébriques, définition
\ref{spaces-properties-definition-etale-site}.
\item Le site $X_{spaces, \etale}$ est le site
$X_{red, spaces, \etale}$ de Propriétés des espaces algébriques, définition
\ref{spaces-properties-definition-spaces-etale-site}.
\item Le site $X_{affine, \etale}$ est le site
$X_{red, affine, \etale}$ de Propriétés des espaces algébriques, lemme
\ref{spaces-properties-lemma-alternative}.
\end{enumerate}
\end{definition}

\noindent
Dans le Lemme \ref{formal-spaces-lemma-identify-spaces-etale}, nous verrons que
$X_{spaces, \etale}$ peut se décrire au moyen de
morphismes d'espaces algébriques formels qui sont représentables
par des espaces algébriques et \'etales.
D'après Propriétés des espaces algébriques, lemmes
\ref{spaces-properties-lemma-compare-etale-sites} et
\ref{spaces-properties-lemma-alternative}
nous avons les identifications
\begin{equation}
\label{formal-spaces-equation-etale-topos}
\Sh(X_\etale) = \Sh(X_{spaces, \etale}) = \Sh(X_{affine, \etale})
\end{equation}
Nous appellerons cela le {\it (petit) topos \'etale} de $X$.

\begin{lemma}
\label{formal-spaces-lemma-functoriality-etale-site}
Soit $S$ un schéma.
Soit $f : X \to Y$ un morphisme d'espaces algébriques formels sur $S$.
\begin{enumerate}
\item Il existe un foncteur continu
$Y_{spaces, \etale} \to X_{spaces, \etale}$
qui induit un morphisme de sites
$$
f_{spaces, \etale} : X_{spaces, \etale} \to Y_{spaces, \etale}.
$$
\item La règle $f \mapsto f_{spaces, \etale}$ est compatible avec les
compositions, autrement dit $(f \circ g)_{spaces, \etale}
= f_{spaces, \etale} \circ g_{spaces, \etale}$ (voir
Sites, définition \ref{sites-definition-composition-morphisms-sites}).
\item Le morphisme de topos associé à $f_{spaces, \etale}$
induit, via (\ref{formal-spaces-equation-etale-topos}), un morphisme de topos
$f_{small} : \Sh(X_\etale) \to \Sh(Y_\etale)$
dont la construction est compatible aux compositions.
\end{enumerate}
\end{lemma}

\begin{proof}
Le seul point est que $f$ induit un morphisme des réductions
$X_{red} \to Y_{red}$ d'après le Lemme \ref{formal-spaces-lemma-reduction-formal-algebraic-space}.
Ce lemme résulte donc immédiatement du lemme correspondant pour les
morphismes d'espaces algébriques (Propriétés des espaces algébriques,
lemme \ref{spaces-properties-lemma-functoriality-etale-site}).
\end{proof}

\noindent
Si le morphisme d'espaces algébriques formels $X \to Y$ est \'etale,
alors le morphisme de topos $\Sh(X_\etale) \to \Sh(Y_\etale)$
est une localisation. Voici un énoncé.

\begin{lemma}
\label{formal-spaces-lemma-etale-morphism-topoi}
Soient $S$ un schéma et $f : X \to Y$ un morphisme
d'espaces algébriques formels sur $S$. Supposons $f$ représentable
par des espaces algébriques et \'etale. Dans ce cas, il existe un
foncteur cocontinu $j : X_\etale \to Y_\etale$.
Le morphisme de topos $f_{small}$ est le
morphisme de topos associé à $j$ ; voir
Sites, lemme \ref{sites-lemma-cocontinuous-morphism-topoi}.
De plus, $j$ est également continu ; ainsi,
Sites, lemme \ref{sites-lemma-when-shriek} s'applique.
\end{lemma}

\begin{proof}
Cela résultera immédiatement du cas des espaces algébriques
(Propriétés des espaces algébriques, lemme
\ref{spaces-properties-lemma-etale-morphism-topoi})
si nous pouvons montrer que le morphisme induit $X_{red} \to Y_{red}$
est \'etale. Remarquons que $X \times_Y Y_{red}$ est un
espace algébrique, \'etale sur l'espace algébrique réduit $Y_{red}$,
et donc lui-même réduit (d'après notre définition des espaces algébriques
réduits dans Propriétés des espaces algébriques, section
\ref{spaces-properties-section-types-properties}).
Ainsi, $X_{red} = X \times_Y Y_{red}$, comme souhaité.
\end{proof}

\begin{lemma}
\label{formal-spaces-lemma-affine-identify-affine-etale}
Soit $S$ un schéma. Soit $X$ un espace algébrique formel affine sur $S$.
Alors $X_{affine, \etale}$ est équivalent à la catégorie dont les objets
sont les morphismes $\varphi : U \to X$ d'espaces algébriques formels tels que
\begin{enumerate}
\item $U$ est un espace algébrique formel affine,
\item $\varphi$ est représentable par des espaces algébriques et \'etale.
\end{enumerate}
\end{lemma}

\begin{proof}
Notons $\mathcal{C}$ la catégorie introduite dans le lemme.
Remarquons que, pour $\varphi : U \to X$ dans $\mathcal{C}$, le
morphisme $\varphi$ est représentable (par des schémas) et affine ; voir le
Lemme \ref{formal-spaces-lemma-affine-representable-by-algebraic-spaces}.
Rappelons que $X_{affine, \etale} = X_{red, affine, \etale}$.
Nous pouvons donc définir un foncteur
$$
\mathcal{C} \longrightarrow X_{affine, \etale},\quad
(U \to X) \longmapsto U \times_X X_{red}
$$
car $U \times_X X_{red}$ est un schéma affine.

\medskip\noindent
Pour achever la démonstration, nous allons construire un foncteur quasi-inverse.
À savoir, écrivons $X = \colim X_\lambda$ comme dans la
Définition \ref{formal-spaces-definition-affine-formal-algebraic-space}.
Pour chaque $\lambda$, l'inclusion $X_{red} \subset X_\lambda$
est un épaississement. Ainsi, pour tout $\lambda$, nous avons une
équivalence
$$
X_{red, affine, \etale} = X_{\lambda, affine, \etale}
$$
par exemple d'après
Compléments d'algèbre, lemme \ref{more-algebra-lemma-locally-nilpotent-henselian}.
Ainsi, si $U_{red} \to X_{red}$ est un morphisme \'etale avec
$U_{red}$ affine, nous obtenons un système de morphismes \'etales
$U_\lambda \to X_\lambda$ de schémas affines compatibles aux
morphismes de transition du système qui définit $X$. Nous pouvons donc prendre
$$
U = \colim U_\lambda
$$
comme espace algébrique formel affine sur $X$. La construction donne
$U \times_X X_\lambda = U_\lambda$. Cela montre que $U \to X$ est
représentable et \'etale. Nous omettons la vérification du fait que les constructions
sont mutuellement inverses.
\end{proof}

\begin{lemma}
\label{formal-spaces-lemma-affine-etale-mcquillan}
Soit $S$ un schéma. Soit $X$ un espace algébrique
formel affine sur $S$. Supposons que $X$ soit de McQuillan, i.e.,
égal à $\text{Spf}(A)$ pour une certaine
algèbre topologique faiblement admissible sur $S$, notée $A$.
Alors $(X_{affine, \etale})^{opp}$ est équivalent à
la catégorie suivante :
\begin{enumerate}
\item ses objets sont les $A$-algèbres de la forme
$B^\wedge = \lim B/JB$, où $A \to B$ est un homomorphisme d'anneaux \'etale
et $J$ parcourt les idéaux faibles de définition de $A$, et
\item ses morphismes sont les homomorphismes continus de $A$-algèbres.
\end{enumerate}
\end{lemma}

\begin{proof}
Combiner les Lemmes \ref{formal-spaces-lemma-affine-identify-affine-etale} et \ref{formal-spaces-lemma-etale}.
\end{proof}

\begin{lemma}
\label{formal-spaces-lemma-identify-spaces-etale}
Soit $S$ un schéma. Soit $X$ un espace algébrique formel sur $S$.
Alors $X_{spaces, \etale}$ est équivalent à la catégorie dont les objets
sont les morphismes $\varphi : U \to X$ d'espaces algébriques formels tels que
$\varphi$ soit représentable par des espaces algébriques et \'etale.
\end{lemma}

\begin{proof}
Notons $\mathcal{C}$ la catégorie introduite dans le lemme.
Rappelons que $X_{spaces, \etale} = X_{red, spaces, \etale}$.
Nous pouvons donc définir un foncteur
$$
\mathcal{C} \longrightarrow X_{spaces, \etale},\quad
(U \to X) \longmapsto U \times_X X_{red}
$$
car $U \times_X X_{red}$ est un espace algébrique \'etale sur $X_{red}$.

\medskip\noindent
Pour achever la démonstration, nous allons construire un foncteur quasi-inverse.
Choisissons un objet $\psi : V \to X_{red}$ de $X_{red, spaces, \etale}$.
Considérons le foncteur
$U_{V, \psi} : (\Sch/S)_{fppf} \to \textit{Ensembles}$ donné par
$$
U_{V, \psi}(T) = \{(a, b) \mid
a : T \to X,
\ b : T \times_{a, X} X_{red} \to V,
\ \psi \circ b = a|_{T \times_{a, X} X_{red}}\}
$$
Nous affirmons que la transformation $U_{V, \psi} \to X$, $(a, b) \mapsto a$
définit un objet de la catégorie $\mathcal{C}$.
Montrons d'abord que $U_{V, \psi}$ est un espace algébrique formel.
Remarquons que $U_{V, \psi}$ est un faisceau pour la topologie fppf (certains détails
sont omis). Supposons ensuite que $X_i \to X$ soit un recouvrement \'etale par des
espaces algébriques formels affines comme dans la
Définition \ref{formal-spaces-definition-formal-algebraic-space}.
Posons $V_i = V \times_{X_{red}} X_{i, red}$ et notons
$\psi_i : V_i \to X_{i, red}$ la projection. Alors 
nous avons
$$
U_{V, \psi} \times_X X_i = U_{V_i, \psi_i}
$$
par un argument formel, car $X_{i, red} = X_i \times_X X_{red}$
(puisque $X_i \to X$ est représentable par des espaces algébriques et \'etale).
Il suffit donc de montrer que $U_{V_i, \psi_i}$ est un
espace algébrique formel affine, car nous disposerons alors
d'un recouvrement $U_{V_i, \psi_i} \to U_{V, \psi}$ comme dans la
Définition \ref{formal-spaces-definition-formal-algebraic-space}.
D'autre part, nous avons vu dans la démonstration du
Lemme \ref{formal-spaces-lemma-etale-morphism-topoi}
que $\psi_i : V_i \to X_i$ est le changement de
base d'un morphisme représentable et \'etale
$U_i \to X_i$ d'espaces algébriques formels affines.
Il n'est alors pas difficile de voir que $U_i = U_{V_i, \psi_i}$,
comme souhaité.

\medskip\noindent
Nous omettons la vérification du fait que $U_{V, \psi} \to X$
est représentable par des espaces algébriques et \'etale.
Nous obtenons ainsi notre foncteur $(V, \psi) \mapsto (U_{V, \psi} \to X)$
dans l'autre sens.
Nous omettons la vérification du fait que les constructions
sont mutuellement inverses.
\end{proof}

\begin{lemma}
\label{formal-spaces-lemma-identify-affine-etale}
Soit $S$ un schéma. Soit $X$ un espace algébrique formel sur $S$.
Alors $X_{affine, \etale}$ est équivalent à la catégorie dont les objets
sont les morphismes $\varphi : U \to X$ d'espaces algébriques formels tels que
\begin{enumerate}
\item $U$ est un espace algébrique formel affine,
\item $\varphi$ est représentable par des espaces algébriques et \'etale.
\end{enumerate}
\end{lemma}

\begin{proof}
Cela résulte de la combinaison des Lemmes \ref{formal-spaces-lemma-identify-spaces-etale} et
\ref{formal-spaces-lemma-characterize-affine}.
\end{proof}






\section{Le faisceau structural}
\label{formal-spaces-section-structure-sheaf}


\noindent
Comme dans la section \ref{formal-spaces-section-etale-sites}, la discussion de la présente
section est motivée par les schémas formels classiques. Ceux-ci sont définis comme des couples
$(\mathfrak X, \mathcal{O}_\mathfrak X)$ où $\mathcal{O}_\mathfrak X$
est un faisceau d'anneaux topologiques.

\medskip\noindent
Soit $X$ un espace algébrique formel. Un {\it faisceau structural de $X$}
est un faisceau d'anneaux topologiques $\mathcal{O}_X$ sur le site \'etale
$X_\etale$ (défini dans la section \ref{formal-spaces-section-etale-sites}) tel que
$$
\mathcal{O}_X(U_{red}) = \lim \Gamma(U_\lambda, \mathcal{O}_{U_\lambda})
$$
comme anneaux topologiques chaque fois que
\begin{enumerate}
\item $\varphi : U \to X$ est un morphisme d'espaces algébriques formels,
\item $U$ est un espace algébrique formel affine,
\item $\varphi$ est représentable par des espaces algébriques et \'etale,
\item $U_{red} \to X_{red}$ est l'objet affine correspondant de
$X_\etale$, voir le
lemme \ref{formal-spaces-lemma-identify-affine-etale},
\item $U = \colim U_\lambda$ est une présentation de $U$ comme limite inductive au sens de la
définition \ref{formal-spaces-definition-affine-formal-algebraic-space}.
\end{enumerate}
Les faisceaux structuraux existent, mais peuvent se comporter de façon inattendue.

\begin{lemma}
\label{formal-spaces-lemma-structure-sheaf}
Tout espace algébrique formel possède un faisceau structural.
\end{lemma}

\begin{proof}
Soient $S$ un schéma et $X$ un espace algébrique formel sur $S$.
D'après (\ref{formal-spaces-equation-etale-topos}), il suffit de construire $\mathcal{O}_X$
comme faisceau d'anneaux topologiques sur $X_{affine, \etale}$.
Notons $\mathcal{C}$ la catégorie dont les objets
sont les morphismes $\varphi : U \to X$ d'espaces algébriques formels tels que
$U$ soit un espace algébrique formel affine et que
$\varphi$ soit représentable par des espaces algébriques et \'etale.
D'après le lemme \ref{formal-spaces-lemma-identify-affine-etale},
le foncteur $U \mapsto U_{red}$ est une équivalence de catégories
$\mathcal{C} \to X_{affine, \etale}$. La règle donnée
avant le lemme fournit donc déjà $\mathcal{O}_X$ comme préfaisceau
d'anneaux topologiques sur $X_{affine, \etale}$. Il suffit ainsi de vérifier
la condition de faisceau.

\medskip\noindent
Par définition de $X_{affine, \etale}$, un recouvrement correspond à une famille finie
$\{g_i : U_i \to U\}_{i = 1, \ldots, n}$ de morphismes de $\mathcal{C}$
telle que $\{U_{i, red} \to U_{red}\}$ soit un recouvrement \'etale.
Les morphismes $g_i$ sont représentables par des espaces algébriques
(lemme \ref{formal-spaces-lemma-permanence-representable}), donc affines
(lemme \ref{formal-spaces-lemma-affine-representable-by-algebraic-spaces}).
Ensuite, $g_i$ est \'etale (cela résulte formellement du
lemme \ref{spaces-properties-lemma-etale-permanence} de Propriétés des espaces algébriques,
puisque $U_i$ et $U$ sont \'etales sur $X$ au sens de la
section \ref{bootstrap-section-representable-by-spaces-properties} de Amorçage).
Enfin, écrivons $U = \colim U_\lambda$ comme dans la
définition \ref{formal-spaces-definition-affine-formal-algebraic-space}.

\medskip\noindent
Ces préparatifs étant faits, démontrons la propriété de
faisceau comme suit. Pour chaque $\lambda$, posons
$U_{i, \lambda} = U_i \times_U U_\lambda$ et
$U_{ij, \lambda} = (U_i \times_U U_j) \times_U U_\lambda$.
Selon ce qui précède, ce sont des schémas affines, $\{U_{i, \lambda} \to U_\lambda\}$
est un recouvrement \'etale, et
$U_{ij, \lambda} = U_{i, \lambda} \times_{U_\lambda} U_{j, \lambda}$.
De plus, $U_i = \colim U_{i, \lambda}$ et
$U_i \times_U U_j = \colim U_{ij, \lambda}$.
Pour chaque $\lambda$, on a une suite exacte
$$
0 \to
\Gamma(U_\lambda, \mathcal{O}_{U_\lambda}) \to
\prod\nolimits_i \Gamma(U_{i, \lambda}, \mathcal{O}_{U_{i, \lambda}}) \to
\prod\nolimits_{i, j} \Gamma(U_{ij, \lambda}, \mathcal{O}_{U_{ij, \lambda}})
$$
car le faisceau structural sur $U_\lambda$ vérifie la condition de faisceau
pour la topologie \'etale
(voir la proposition de Cohomologie \'etale
\ref{etale-cohomology-proposition-quasi-coherent-sheaf-fpqc}).
Comme les limites commutent aux limites, la limite projective de ces
suites exactes est une suite exacte
$$
0 \to
\lim \Gamma(U_\lambda, \mathcal{O}_{U_\lambda}) \to
\prod\nolimits_i \lim \Gamma(U_{i, \lambda}, \mathcal{O}_{U_{i, \lambda}}) \to
\prod\nolimits_{i, j} \lim
\Gamma(U_{ij, \lambda}, \mathcal{O}_{U_{ij, \lambda}})
$$
Cela signifie que
$$
0 \to
\mathcal{O}_X(U_{red}) \to
\prod\nolimits_i \mathcal{O}_X(U_{i, red}) \to
\prod\nolimits_{i, j} \mathcal{O}_X((U_i \times_U U_j)_{red})
$$
est exacte en tant que suite de groupes abéliens. Il reste à
montrer que la topologie limite sur $\mathcal{O}_X(U_{red})$ est induite
par la topologie limite sur
$\prod_{i = 1, \ldots, n} \mathcal{O}_X(U_{i, red})$.
Les noyaux des applications
$\mathcal{O}_X(U_{red}) \to \mathcal{O}_{U_\lambda}(U_\lambda)$
forment un système fondamental de sous-modules ouverts. Comme ils sont les images
réciproques des noyaux des applications
$\prod_i \mathcal{O}_X(U_{i, red}) \to
\prod_i \mathcal{O}_{U_{i, \lambda}}(U_{i, \lambda})$
et que ces noyaux forment un système fondamental de sous-modules ouverts
pour $\mathcal{O}_X(U_{red}) \to \mathcal{O}_{U_\lambda}(U_\lambda)$,
le résultat est acquis.
\end{proof}

\begin{remark}
\label{formal-spaces-remark-not-enough-sections}
Le faisceau structural ne possède pas toujours « assez de sections ».
Dans la section \ref{examples-section-affine-formal-algebraic-space} de Exemples,
nous avons vu qu'il existe des espaces algébriques formels affines qui
ne sont pas de McQuillan, et même des exemples dont les points ne sont pas
séparés par les fonctions régulières.
\end{remark}

\noindent
Dans le lemme suivant, nous montrons que le faisceau structural d'un schéma formel
affine à indexation dénombrable a une cohomologie supérieure nulle.
Sans l'hypothèse d'indexation dénombrable, ce résultat est sans doute faux en général.

\begin{lemma}
\label{formal-spaces-lemma-higher-vanishing-structure-sheaf}
Si $X$ est un espace algébrique formel affine à indexation dénombrable, alors
$H^n(X_\etale, \mathcal{O}_X) = 0$ pour $n > 0$.
\end{lemma}

\begin{proof}
Nous pouvons travailler avec $X_{affine, \etale}$, puisque celui-ci donne le même topos.
Nous allons appliquer le lemme
\ref{sites-cohomology-lemma-cech-vanish-collection}
de Cohomologie sur les sites pour établir cette annulation. Comme $X_{affine, \etale}$
possède les sommes disjointes finies, on est ramené au complexe de {\v C}ech
d'un recouvrement donné par une seule flèche $\{U_{red} \to V_{red}\}$
dans $X_{affine, \etale} = X_{red, affine, \etale}$
(voir Cohomologie \'etale, lemme \ref{etale-cohomology-lemma-cech-complex}).
Il faut donc montrer que
$$
0 \to \mathcal{O}_X(V_{red}) \to \mathcal{O}_X(U_{red}) \to
\mathcal{O}_X(U_{red} \times_{V_{red}} U_{red}) \to \ldots
$$
est exacte. Nous le ferons ci-dessous dans le cas $V_{red} = X_{red}$.
Le cas général se démontre exactement de la même manière.

\medskip\noindent
Rappelons que $X = \text{Spf}(A)$, où $A$ est un
anneau topologique faiblement admissible possédant un système
fondamental dénombrable d'idéaux faibles de définition.
Nous avons vu dans les lemmes
\ref{formal-spaces-lemma-affine-identify-affine-etale} et
\ref{formal-spaces-lemma-affine-etale-mcquillan}
que l'objet $U_{red}$ de $X_{affine, \etale}$
correspond à un morphisme $U \to X$ d'espaces algébriques formels
affines, représentable par des espaces algébriques et \'etale,
et que $U = \text{Spf}(B^\wedge)$, où $B$ est une $A$-algèbre \'etale.
La règle définissant le faisceau structural donne
$$
\mathcal{O}_X(U_{red}) = B^\wedge
$$
Rappelons que $B^\wedge = \lim B/JB$, où la limite porte sur les
idéaux faibles de définition $J \subset A$.
En développant les définitions, on obtient
$$
\mathcal{O}_X(U_{red} \times_{X_{red}} U_{red}) = (B \otimes_A B)^\wedge
$$
et ainsi de suite. Puisque $U \to X$ est un recouvrement, l'application $A \to B$
est fidèlement plate, voir le lemme \ref{formal-spaces-lemma-etale-surjective}.
Le complexe
$$
0 \to A \to B \to B \otimes_A B \to B \otimes_A B \otimes_A B \to \ldots
$$
est donc {\bf universellement} exact, voir le lemme \ref{descent-lemma-ff-exact} de Descente.
Notre but est de montrer que
$$
H^n(0 \to A^\wedge \to B^\wedge \to (B \otimes_A B)^\wedge
\to (B \otimes_A B \otimes_A B)^\wedge \to \ldots)
$$
est nul pour $n > 0$. Pour comprendre la situation, décomposons notre
complexe exact (avant passage au complété) en suites exactes courtes
$$
0 \to A \to B \to M_1 \to 0,\quad
0 \to M_i \to B^{\otimes_A i + 1} \to M_{i + 1} \to 0
$$
D'après ce qui précède, ce sont des suites exactes courtes universellement
exactes. Ainsi $JM_i = M_i \cap J(B^{\otimes_A i + 1})$ pour
tout idéal $J$ de $A$. En particulier, la
topologie de $M_i$ comme sous-module de $B^{\otimes_A i + 1}$
est identique à celle de $M_i$ comme module quotient de
$B^{\otimes_A i}$. Comme il existe dans $A$ un système fondamental dénombrable
d'idéaux faibles de définition, les suites
$$
0 \to A^\wedge \to B^\wedge \to M_1^\wedge \to 0,\quad
0 \to M_i^\wedge \to (B^{\otimes_A i + 1})^\wedge \to M_{i + 1}^\wedge \to 0
$$
restent exactes d'après le lemme \ref{formal-spaces-lemma-ses}. Cela démontre le lemme.
\end{proof}

\begin{remark}
\label{formal-spaces-remark-bad-quasi-coherent}
Même si le faisceau structural possède de bonnes propriétés, cela ne
signifie pas qu'il existe une bonne théorie des modules quasi-cohérents. Par exemple,
dans la section \ref{examples-section-nonabelian-QCoh} de Exemples,
nous avons vu que, pour presque tout espace algébrique formel affine noethérien,
la notion la plus naturelle de module quasi-cohérent conduit à une
catégorie de modules qui n'est pas abélienne.
\end{remark}





\section{Limites inductives d'espaces algébriques formels}
\label{formal-spaces-section-colimits-of-formal}

\noindent
Dans cette section, nous généralisons le résultat de la
section \ref{formal-spaces-section-global-colimits}
au cas des systèmes de morphismes d'espaces algébriques formels.
Remarquons que, dans les lemmes ci-dessous, la condition
« $f_{\lambda \mu} : X_\lambda \to X_\mu$ est une immersion fermée
induisant un isomorphisme $X_{\lambda, red} \to X_{\mu, red}$ »
peut se reformuler ainsi :
« $f_{\lambda \mu}$ est représentable et est un épaississement ».

\begin{lemma}
\label{formal-spaces-lemma-colimit-affine-formal}
Soit $S$ un schéma. Donnons-nous un ensemble ordonné filtrant
$\Lambda$ et un système d'espaces algébriques formels affines
$(X_\lambda, f_{\lambda \mu})$ indexé par $\Lambda$, où chaque
$f_{\lambda \mu} : X_\lambda \to X_\mu$ est une immersion fermée
induisant un isomorphisme $X_{\lambda, red} \to X_{\mu, red}$.
Alors $X = \colim_{\lambda \in \Lambda} X_\lambda$
est un espace algébrique formel affine sur $S$.
\end{lemma}

\begin{proof}
On peut écrire
$X_\lambda = \colim_{\omega \in \Omega_\lambda} X_{\lambda, \omega}$
comme limite inductive de schémas affines indexée par un ensemble ordonné filtrant
$\Omega_\lambda$, de telle sorte que les morphismes de transition
$X_{\lambda, \omega} \to X_{\lambda, \omega'}$
soient des épaississements. Pour tous $\lambda, \mu \in \Lambda$
et $\omega \in \Omega_\lambda$, avec $\mu \geq \lambda$,
il existe un $\omega' \in \Omega_\mu$ tel que le morphisme
$X_{\lambda, \omega} \to X_\mu$ se factorise par $X_{\mu, \omega'}$, voir le
lemme \ref{formal-spaces-lemma-factor-through-thickening}. Le
morphisme $X_{\lambda, \omega} \to X_{\mu, \omega'}$
est alors une immersion fermée induisant un isomorphisme sur les réductions,
donc un épaississement. Posons
$\Omega = \coprod_{\lambda \in \Lambda} \Omega_\lambda$
et convenons que $(\lambda, \omega) \leq (\mu, \omega')$ si et seulement si
$\lambda \leq \mu$ et si $X_{\lambda, \omega} \to X_\mu$ se factorise par
$X_{\mu, \omega'}$. Il résulte de ce qui précède que
$\Omega$ est un ensemble ordonné filtrant et que
$X = \colim_{\lambda \in \Lambda} X_\lambda =
\colim_{(\lambda, \omega) \in \Omega} X_{\lambda, \omega}$.
Ceci achève la démonstration.
\end{proof}

\begin{lemma}
\label{formal-spaces-lemma-colimit-formal-spaces-is-formal-space}
Soit $S$ un schéma. Donnons-nous un ensemble ordonné filtrant
$\Lambda$ et un système d'espaces algébriques formels
$(X_\lambda, f_{\lambda \mu})$ indexé par $\Lambda$, où chaque
$f_{\lambda \mu} : X_\lambda \to X_\mu$ est une immersion fermée
induisant un isomorphisme $X_{\lambda, red} \to X_{\mu, red}$.
Alors $X = \colim_{\lambda \in \Lambda} X_\lambda$
est un espace algébrique formel sur $S$.
\end{lemma}

\begin{proof}
Comme la limite inductive est prise dans la catégorie des faisceaux fppf,
$X$ est un faisceau. Choisissons et fixons $\lambda \in \Lambda$. Choisissons
un recouvrement $\{X_{i, \lambda} \to X_\lambda\}$ comme dans la définition
\ref{formal-spaces-definition-formal-algebraic-space}. En particulier,
En particulier, $\{X_{i, \lambda, red} \to X_{\lambda, red}\}$ est un
recouvrement \'etale par des schémas affines.
Pour chaque $\mu \geq \lambda$, il existe un diagramme cartésien
$$
\xymatrix{
X_{i, \lambda} \ar[r] \ar[d] & X_{i, \mu} \ar[d] \\
X_\lambda \ar[r] & X_\mu
}
$$
dont les flèches verticales sont \'etales. En effet, le morphisme \'etale
$X_{i, \lambda, red} \to X_{\lambda, red} = X_{\mu, red}$
correspond à un morphisme \'etale $X_{i, \mu} \to X_\mu$
d'espaces algébriques formels, où $X_{i, \mu}$ est un espace
algébrique formel affine, voir le
lemme \ref{formal-spaces-lemma-affine-identify-affine-etale}.
Le même lemme montre que le changement de base de $X_{i, \mu}$ à $X_\lambda$
coïncide avec $X_{i, \lambda}$. Il en résulte aussi que
$X_{i, \mu} = X_\mu \times_{X_{\mu'}} X_{i, \mu'}$ pour
$\mu' \geq \mu \geq \lambda$. Posons $X_i = \colim X_{i, \mu}$.
Alors $X_{i, \mu} = X_i \times_X X_\mu$ (comme foncteurs).
Comme tout morphisme $T \to X = \colim X_\mu$
issu d'un schéma affine (ou quasi-compact) $T$ se factorise par $X_\mu$
pour un certain $\mu$, on conclut que
$\colim X_{i, \mu} \to \colim X_\mu$ est \'etale.
Ainsi, s'il est possible de montrer que $\colim X_{i, \mu}$ est un espace
algébrique formel affine, le lemme en résultera.
Notons que les morphismes $X_{i, \mu} \to X_{i, \mu'}$
sont des immersions fermées, car ils s'obtiennent par changement de base de l'immersion fermée
$X_\mu \to X_{\mu'}$. Enfin, le morphisme
$X_{i, \mu, red} \to X_{i, \mu', red}$ est un isomorphisme,
puisque $X_{\mu, red} \to X_{\mu', red}$ en est un.
On est donc ramené au cas traité dans le
lemme \ref{formal-spaces-lemma-colimit-affine-formal}.
\end{proof}








\section{Recomplétion}
\label{formal-spaces-section-recompletion}

\noindent
Dans cette section, nous définissons le complété d'un espace algébrique
formel le long d'une partie fermée de sa réduction. C'est la généralisation
naturelle de la section \ref{formal-spaces-section-completion}.

\begin{lemma}
\label{formal-spaces-lemma-completion-affine-formal-is-affine-formal}
Soit $S$ un schéma. Soit $X$ un espace algébrique formel affine sur $S$.
Soit $T \subset |X_{red}|$ une partie fermée. Alors le foncteur
$$
X_{/T} : (\Sch/S)_{fppf} \longrightarrow \textit{Ensembles},\quad
U \longmapsto \{f : U \to X : f(|U|) \subset T\}
$$
est un espace algébrique formel affine.
\end{lemma}

\begin{proof}
Écrivons $X = \colim X_\lambda$ comme dans la
définition \ref{formal-spaces-definition-affine-formal-algebraic-space}.
Alors $X_{\lambda, red} = X_{red}$ et nous considérons
$T$ comme une partie fermée de $|X_\lambda| = |X_{\lambda, red}|$.
D'après le lemme \ref{formal-spaces-lemma-completion-affine-is-affine-formal-algebraic-space},
pour chaque $\lambda$, le complété
$(X_\lambda)_{/T}$ est un espace algébrique formel affine.
Les morphismes de transition $(X_\lambda)_{/T} \to (X_\mu)_{/T}$ sont des
immersions fermées, car ils s'obtiennent par changement de base des morphismes de transition
$X_\lambda \to X_\mu$, voir le lemme \ref{formal-spaces-lemma-map-completions-representable}.
De plus, les morphismes $((X_\lambda)_{/T})_{red} \to ((X_\mu)_/T)_{red}$
sont des isomorphismes d'après le lemme \ref{formal-spaces-lemma-reduction-completion}.
Comme $X_{/T} = \colim (X_\lambda)_{/T}$, on conclut
par le lemme \ref{formal-spaces-lemma-colimit-affine-formal}.
\end{proof}

\begin{lemma}
\label{formal-spaces-lemma-completion-fas-is-fas}
Soit $S$ un schéma. Soit $X$ un espace algébrique formel sur $S$.
Soit $T \subset |X_{red}|$ une partie fermée. Alors le foncteur
$$
X_{/T} : (\Sch/S)_{fppf} \longrightarrow \textit{Ensembles},\quad
U \longmapsto \{f : U \to X \mid f(|U|) \subset T\}
$$
est un espace algébrique formel.
\end{lemma}

\begin{proof}
Le foncteur $X_{/T}$ est un faisceau fppf, car si
$\{U_i \to U\}$ est un recouvrement fppf, alors
$\coprod |U_i| \to |U|$ est surjectif.

\medskip\noindent
Choisissons un recouvrement $\{g_i : X_i \to X\}_{i \in I}$ comme dans la définition
\ref{formal-spaces-definition-formal-algebraic-space}.
Les morphismes $X_i \times_X X_{/T} \to X_{/T}$ sont \'etales
(voir le lemme de Espaces
\ref{spaces-lemma-base-change-representable-transformations-property})
et l'application $\coprod X_i \times_X X_{/T} \to X_{/T}$ est une surjection de
faisceaux. Il suffit donc de montrer que $X_{/T} \times_X X_i$
est un espace algébrique formel affine. Un point à valeurs dans $U$ de
$X_i \times_X X_{/T}$ est un morphisme $U \to X_i$ dont l'image est
contenue dans la partie fermée
$|g_{i, red}|^{-1}(T) \subset |X_{i, red}|$. Cela résulte donc du
lemme \ref{formal-spaces-lemma-completion-affine-formal-is-affine-formal}.
\end{proof}

\begin{definition}
\label{formal-spaces-definition-completion-formal-algebraic-space}
Soit $S$ un schéma. Soit $X$ un espace algébrique formel sur $S$.
Soit $T \subset |X_{red}|$ une partie fermée. L'espace algébrique formel
$X_{/T}$ du lemme \ref{formal-spaces-lemma-completion-is-formal-algebraic-space}
est appelé le {\it complété de $X$ le long de $T$}.
\end{definition}

\noindent
Soit $f : X \to X'$ un morphisme d'espaces algébriques formels sur un schéma $S$.
Supposons que $T \subset |X_{red}|$ et $T' \subset |X'_{red}|$
soient des parties fermées telles que $|f_{red}|(T) \subset T'$. Il est alors clair que
$f$ définit un morphisme d'espaces algébriques formels
$$
X_{/T} \longrightarrow X'_{/T'}
$$
entre les complétés.

\begin{lemma}
\label{formal-spaces-lemma-map-recompletions}
Soit $S$ un schéma. Soit $f : X' \to X$ un morphisme
d'espaces algébriques formels sur $S$. Soit $T \subset |X_{red}|$
une partie fermée et soit $T' = |f_{red}|^{-1}(T) \subset |X'_{red}|$.
Alors
$$
\xymatrix{
X'_{/T'} \ar[r] \ar[d] & X' \ar[d]^f \\
X_{/T} \ar[r] & X
}
$$
est un diagramme cartésien d'espaces algébriques formels sur $S$.
\end{lemma}

\begin{proof}
Observons en effet que les flèches horizontales sont des monomorphismes
par construction. Il suffit donc de montrer qu'un morphisme
$g : U \to X'$ issu d'un schéma $U$ définit un point de $X'_{/T}$
si et seulement si $f \circ g$ définit un point de $X_{/T}$.
Autrement dit, il faut montrer que
$g(U)$ est contenu dans $T' \subset |X'_{red}|$
si et seulement si $(f \circ g)(U)$ est contenu dans $T \subset |X_{red}|$.
Cela résulte immédiatement du choix de $T'$ comme
image réciproque de $T$.
\end{proof}

\begin{lemma}
\label{formal-spaces-lemma-reduction-recompletion}
Soit $S$ un schéma. Soit $X$ un espace algébrique formel sur $S$.
Soit $T \subset |X_{red}|$ une partie fermée. La réduction
$(X_{/T})_{red}$ du complété $X_{/T}$ de $X$ le long de $T$ est
le sous-espace fermé induit réduit $Z$ de $X_{red}$ correspondant à $T$.
\end{lemma}

\begin{proof}
Il résulte du lemme \ref{formal-spaces-lemma-reduction-formal-algebraic-space},
de la définition de Propriétés des espaces algébriques
\ref{spaces-properties-definition-reduced-induced-space}
(qui utilise le lemme de Propriétés des espaces algébriques
\ref{spaces-properties-lemma-reduced-closed-subspace} pour construire $Z$),
et de la définition de $X_{/T}$ que
$Z$ et $(X_{/T})_{red}$ sont des espaces algébriques réduits
caractérisés par la même propriété universelle :
un morphisme $g : Y \to X$ dont la source est un espace algébrique réduit
se factorise par eux si et seulement si $|Y|$ s'envoie dans $T \subset |X|$.
\end{proof}

\begin{lemma}
\label{formal-spaces-lemma-recompletion-affine-types}
Soit $S$ un schéma. Soit $X$ un espace algébrique formel affine sur $S$.
Soit $T \subset X_{red}$ une partie fermée et soit $X_{/T}$
le complété formel de $X$ le long de $T$. Alors
\begin{enumerate}
\item $X_{/T}$ est un espace algébrique formel affine,
\item si $X$ est de McQuillan, alors $X_{/T}$ est de McQuillan,
\item si $|X_{red}| \setminus T$ est quasi-compact et si $X$
est à indexation dénombrable, alors $X_{/T}$ est à indexation dénombrable,
\item si $|X_{red}| \setminus T$ est quasi-compact et si $X$
est adique*, alors $X_{/T}$ est adique*,
\item si $X$ est noethérien, alors $X_{/T}$ est noethérien.
\end{enumerate}
\end{lemma}

\begin{proof}
L'assertion (1) est le lemme \ref{formal-spaces-lemma-completion-affine-formal-is-affine-formal}.
Si $X$ est de McQuillan, alors $X = \text{Spf}(A)$ pour un anneau topologique
faiblement admissible $A$. Alors $X_{/T} \to X \to \Spec(A)$ vérifie
la propriété (2) du lemme \ref{formal-spaces-lemma-mcquillan-affine-formal-algebraic-space} ;
par conséquent, $X_{/T}$ est de McQuillan, voir la
définition \ref{formal-spaces-definition-types-affine-formal-algebraic-space}.

\medskip\noindent
Supposons $X$ et $T$ comme dans (3).
Alors $X = \text{Spf}(A)$, où $A$ possède un système fondamental
$A \supset I_1 \supset I_2 \supset I_3 \supset \ldots$
constitué d'idéaux faibles de définition, voir le lemme \ref{formal-spaces-lemma-countably-indexed}.
D'après le lemme \ref{algebra-lemma-qc-open} de Algèbre,
on peut trouver un idéal de type fini
$\overline{J} = (\overline{f}_1, \ldots, \overline{f}_r) \subset A/I_1$
tel que $T$ soit défini par $\overline{J}$ dans $\Spec(A/I_1) = |X_{red}|$.
Choisissons des relèvements $f_i \in A$ de $\overline{f}_i$.
Si $Z = \Spec(B)$ est un schéma affine et si $g : Z \to X$ est
un morphisme tel que $g(Z) \subset T$ (au sens ensembliste), alors
$g^\sharp : A \to B$ se factorise par $A/I_n$ pour un certain $n$
et $g^\sharp(f_i)$ est nilpotent dans $B$ pour tout $i$. Ainsi,
$J_{m, n} = (f_1, \ldots, f_r)^m + I_n$ s'envoie sur zéro dans $B$
pour certains $n, m \geq 1$. Il en résulte que $X_{/T}$ est le spectre formel de
$\lim_{n, m} A/J_{m, n}$ et est donc à indexation dénombrable.
Cela démontre (3).

\medskip\noindent
Démonstration de (4). L'argument est le même que dans (3).
Toutefois, on peut ici choisir $I_n = I^n$ pour un idéal de type fini
$I \subset A$. Il est alors clair que $X_{/T}$ est le spectre formel de
$\lim A/J^n$, où $J = (f_1, \ldots, f_r) + I$. Quelques détails sont omis.

\medskip\noindent
Démonstration de (5). Dans ce cas, $X_{red}$ est le spectre d'un anneau noethérien
et l'hypothèse que $|X_{red}| \setminus T$ est quasi-compact
est donc satisfaite. Comme dans la démonstration de (4), on voit alors que
$X_{/T}$ est le spectre de $\lim A/J^n$, qui est un anneau topologique
adique noethérien, voir le
lemme \ref{algebra-lemma-completion-Noetherian-Noetherian} de Algèbre.
\end{proof}

\begin{lemma}
\label{formal-spaces-lemma-recompletion-types}
Soit $S$ un schéma. Soit $X$ un espace algébrique formel sur $S$.
Soit $T \subset X_{red}$ une partie fermée et soit $X_{/T}$
le complété formel de $X$ le long de $T$. Alors
\begin{enumerate}
\item si $X_{red} \setminus T \to X_{red}$ est quasi-compact et si $X$
est localement à indexation dénombrable, alors $X_{/T}$ l'est aussi,
\item si $X_{red} \setminus T \to X_{red}$ est quasi-compact et si $X$
est localement adique*, alors $X_{/T}$ l'est aussi, et
\item si $X$ est localement noethérien, alors $X_{/T}$ l'est aussi.
\end{enumerate}
\end{lemma}

\begin{proof}
Choisissons un recouvrement $\{X_i \to X\}$ comme dans la
définition \ref{formal-spaces-definition-formal-algebraic-space}.
Soit $T_i \subset X_{i, red}$ l'image réciproque de $T$.
On a $X_i \times_X X_{/T} = (X_i)_{/T_i}$
(lemme \ref{formal-spaces-lemma-map-recompletions}).
Ainsi, $\{(X_i)_{/T_i} \to X_{/T}\}$ est un recouvrement comme dans la
définition \ref{formal-spaces-definition-formal-algebraic-space}.
De plus, si $X_{red} \setminus T \to X_{red}$ est quasi-compact, alors
$X_{i, red} \setminus T_i \to X_{i, red}$ l'est aussi ;
et si $X$ est localement à indexation dénombrable, localement adique* ou
localement noethérien, alors $X_i$ est respectivement à indexation dénombrable, adique*
ou noethérien. Le lemme résulte donc du cas affine, à savoir le
lemme \ref{formal-spaces-lemma-recompletion-affine-types}.
\end{proof}








\section{Complété le long d'un sous-espace fermé}
\label{formal-spaces-section-completion-subspace}

\noindent
Cette section est l'analogue de la section \ref{formal-spaces-section-completion}
pour les complétés le long d'un sous-espace fermé.

\begin{definition}
\label{formal-spaces-definition-completion-subspace}
Soit $S$ un schéma. Soit $X$ un espace algébrique sur $S$.
Soit $Z \subset X$ un sous-espace fermé et notons $Z_n \subset X$
le $n$-ième voisinage infinitésimal. L'espace algébrique formel
$$
X^\wedge_Z = \colim Z_n
$$
(voir le lemme \ref{formal-spaces-lemma-colimit-formal-spaces-is-formal-space})
est appelé le {\it complété de $X$ le long de $Z$}.
\end{definition}

\noindent
Par exemple, si $X = \Spec(A)$ et si $Z$ est défini par l'idéal $I \subset A$,
alors
$$
X^\wedge_Z = \colim \Spec(A/I^n) = \text{Spf}(B)
$$
où $B = \lim A/I^n$ est muni de la topologie limite. Notons que $B$ est un
anneau topologique faiblement adique, mais n'est pas adique en général.

\medskip\noindent
Revenons au cas général. Si $T = |Z|$, il existe un morphisme canonique
$X^\wedge_Z \to X_{/T}$ qui compare les complétés le long de $Z$
et de $T$ (section \ref{formal-spaces-section-completion}) et qui n'est pas nécessairement un isomorphisme
(voir le lemme \ref{formal-spaces-lemma-when-isomorphism}).

\medskip\noindent
Soit $f : X \to X'$ un morphisme d'espaces algébriques sur un schéma $S$.
Supposons que $Z \subset X$ et $Z' \subset X'$ soient des sous-espaces fermés
tels que $f|_Z$ envoie $Z$ dans $Z'$ et induise un morphisme $Z \to Z'$.
Il est alors clair que $f$ définit un morphisme d'espaces algébriques formels
$$
X^\wedge_Z \longrightarrow (X')^\wedge_{Z'}
$$
entre les complétés.

\begin{lemma}
\label{formal-spaces-lemma-map-completions-subspaces-representable}
Soit $S$ un schéma. Soit $f : X' \to X$ un morphisme
d'espaces algébriques sur $S$. Soit $Z \subset X$
un sous-espace fermé et soit $Z' = f^{-1}(Z) = X' \times_X Z$.
Alors
$$
\xymatrix{
(X')^\wedge_{Z'} \ar[r] \ar[d] & X' \ar[d]^f \\
X^\wedge_Z \ar[r] & X
}
$$
est un diagramme cartésien de faisceaux. En particulier, le morphisme
$(X')^\wedge_{Z'} \to X^\wedge_Z$ est représentable par des espaces algébriques.
\end{lemma}

\begin{proof}
Supposons en effet que $Y \to X$ soit un morphisme d'un schéma dans $X$ tel
que $Y \to X$ se factorise par $Z$. Alors $Y \times_X X' \to X$
est un morphisme d'espaces algébriques tel que $Y \times_X X' \to X'$
se factorise par $Z'$. Puisque $Z'_n = X' \times_X Z_n$ pour tout $n \geq 1$,
il en va de même pour les voisinages infinitésimaux.
Le carré cartésien de foncteurs résulte donc
des formules $X^\wedge_Z = \colim Z_n$ et
$(X')^\wedge_{Z'} = \colim Z'_n$.
\end{proof}

\begin{lemma}
\label{formal-spaces-lemma-reduction-completion-subspace}
Soit $S$ un schéma. Soit $X$ un espace algébrique sur $S$.
Soit $Z \subset X$ un sous-espace fermé. La réduction $(X^\wedge_Z)_{red}$
du complété $X^\wedge_Z$ de $X$ le long de $Z$ est $Z_{red}$.
\end{lemma}

\begin{proof}
La démonstration est omise.
\end{proof}

\begin{lemma}
\label{formal-spaces-lemma-affine-formal-completion-subspace-types}
Soit $S$ un schéma. Soit $X = \Spec(A)$ un schéma affine sur $S$.
Soit $Z \subset X$ le sous-schéma fermé correspondant à l'idéal
$I \subset A$. Alors
\begin{enumerate}
\item L'espace algébrique formel affine $X^\wedge_Z$ est faiblement adique.
\item Si $I$ est de type fini, alors
$X^\wedge_Z = \text{Spf}(A^\wedge)$, où $A^\wedge$ est le
complété $I$-adique de $A$.
\item Si $Z \to X$ est de présentation finie, c'est-à-dire si $I$ est
de type fini, alors $X^\wedge_Z$ est adique*.
\item Si $X$ est noethérien, alors $X^\wedge_Z$ est noethérien.
\end{enumerate}
\end{lemma}

\begin{proof}
La démonstration est omise.
\end{proof}

\begin{lemma}
\label{formal-spaces-lemma-formal-completion-subspace-types}
Soit $S$ un schéma. Soit $X$ un espace algébrique sur $S$.
Soit $Z \subset X$ un sous-espace fermé. Soit $X^\wedge_Z$ le
complété formel de $X$ le long de $Z$.
\begin{enumerate}
\item L'espace algébrique formel $X^\wedge_Z$ est localement faiblement adique.
\item Si $Z \to X$ est de présentation finie,
alors $X^\wedge_Z$ est localement adique*.
\item Si $X$ est localement noethérien, alors $X_Z$ est localement
noethérien.
\end{enumerate}
\end{lemma}

\begin{proof}
La démonstration est omise.
\end{proof}

\begin{lemma}
\label{formal-spaces-lemma-when-isomorphism}
Soit $S$ un schéma. Soit $X$ un espace algébrique sur $S$.
Soit $Z \subset X$ un sous-espace fermé et posons $T = |Z|$. Le
morphisme canonique $c : X^\wedge_Z \to X_{/T}$ est un isomorphisme si $Z \to X$
est de présentation finie, mais pas en général.
\end{lemma}

\begin{proof}
Les deux constructions commutent à la localisation \'etale ; il suffit donc
de le démontrer lorsque $X$ est affine. Écrivons $X = \Spec(A)$ et supposons que $Z$ corresponde
à l'idéal $I \subset A$. Si $I$ est de type fini, alors
$X^\wedge_Z$ et $X_{/T}$ sont tous deux égaux à $\text{Spf}(A^\wedge)$, où $A^\wedge$
est le complété $I$-adique de $A$, voir les
lemmes \ref{formal-spaces-lemma-affine-formal-completion-types} et
\ref{formal-spaces-lemma-affine-formal-completion-subspace-types}.
Si $A = \mathbf{Z}[x_1, x_2, \ldots]$ et $I = (x_1, x_2, \ldots)$,
alors $\Spec(A/(x_n^n, n \geq 1)) \to X$ se factorise par $X_{/T}$,
mais pas par $X^\wedge_Z$ ; par conséquent, $c$ n'est pas un isomorphisme.
\end{proof}










% OK, start here.
%

\chapter{Algébrisation des espaces formels}



\phantomsection
\label{restricted-section-phantom}


\section{Introduction}
\label{restricted-section-introduction}

\noindent
Le but principal de ce chapitre est de démontrer le théorème d'Artin sur les
dilatations, voir le théorème \ref{restricted-theorem-dilatations} ; le résultat sur les
contractions sera étudié dans Axiomes d'Artin, section
\ref{artin-section-contractions}. Ces deux résultats utilisent des faits sur
les espaces algébriques formels ; c'est pourquoi, dans la partie centrale de
ce chapitre, nous poursuivons l'étude des espaces algébriques formels commencée
au chapitre précédent, voir Espaces formels, section
\ref{formal-spaces-section-introduction}. La première partie de ce chapitre est
consacrée à des préliminaires algébriques, portant principalement sur
l'algébrisation des algèbres rig-étales.

\medskip\noindent
Soient $A$ un anneau noethérien et $I \subset A$ un idéal. Dans la
première partie de ce chapitre (sections \ref{restricted-section-two-categories}
-- \ref{restricted-section-approximation-principal}), nous étudions la catégorie des
algèbres complètes pour la topologie $I$-adique, notées $B$, et
topologiquement de type fini sur l'anneau noethérien $A$. Nous montrons que
$B = A\{x_1, \ldots, x_n\}/J$ pour un certain idéal (fermé) $J$ de
l'anneau de séries formelles restreintes (où $A$ est muni de la topologie
$I$-adique). Nous montrons qu'il existe une bonne notion de complexe cotangent
naïf $\NL_{B/A}^\wedge$. Si une puissance de $I$ annule
$\NL_{B/A}^\wedge$, nous disons que $B$ est une algèbre rig-étale sur
$(A, I)$ ; il existe une notion analogue d'algèbre rig-lisse. Si $A$ est un
anneau G, nous montrons, à l'aide du théorème de Popescu, que toute algèbre
rig-lisse $B$ sur $(A, I)$ est le complété d'une $A$-algèbre de type fini ;
de manière informelle, nous disons que l'on peut « algébriser » $B$.
Cependant, le résultat principal de la première partie affirme que toute
algèbre rig-étale $B$ sur $(A, I)$ est algébrisable (sans supposer que $A$ est
un anneau G), voir le lemme \ref{restricted-lemma-approximate}. Pour des indications
bibliographiques sur ce type d'algébrisation, voir la remarque
\ref{restricted-remark-discussion}. Les références générales pour la première partie sont
\cite{EGA}, \cite{Abbes} et \cite{Fujiwara-Kato}.

\medskip\noindent
Dans la deuxième partie de ce chapitre (sections
\ref{restricted-section-finite-type-red} -- \ref{restricted-section-formal-modifications}), nous
étudions différents types de morphismes d'espaces algébriques formels à un
degré de généralité raisonnable (principalement pour les espaces algébriques
formels localement noethériens). La notion la plus intéressante est celle de
« modification formelle », introduite dans la dernière section. Nous vérifions
soigneusement que notre définition coïncide avec celle d'Artin dans
\cite{ArtinII}.

\medskip\noindent
Enfin, dans la troisième et dernière partie de ce chapitre (sections
\ref{restricted-section-completion-and-morphisms} -- \ref{restricted-section-modifications}), nous
démontrons le théorème principal et en donnons quelques applications. En fait,
nous déduisons le théorème d'Artin d'un résultat plus fort, à savoir le théorème
\ref{restricted-theorem-dilatations-general}. Très sommairement, ce théorème affirme que,
si $f : \mathfrak X \to \mathfrak X'$ est un morphisme rig-étale et si
$\mathfrak X'$ est le complété formel d'un espace algébrique localement
noethérien, alors $\mathfrak X$ l'est aussi. Dans les travaux d'Artin, le
morphisme $f$ est supposé propre et rig-surjectif.









\section{Deux catégories}
\label{restricted-section-two-categories}

\noindent
Soient $A$ un anneau et $I \subset A$ un idéal.
Dans cette section, ${}^\wedge$ désignera la complétion $I$-adique.
Posons $A_n = A/I^n$, de sorte que le complété $I$-adique de $A$ soit
$A^\wedge = \lim A_n$. Soit $\mathcal{C}$ la catégorie
\begin{equation}
\label{restricted-equation-C}
\mathcal{C} =
\left\{
\begin{matrix}
\text{des systèmes projectifs }\ldots \to B_3 \to B_2 \to B_1 \\
\text{où }B_n\text{ est une }A_n\text{-algèbre de type fini,}\\
B_{n + 1} \to B_n\text{ est un homomorphisme de }A_{n + 1}\text{-algèbres}\\
\text{qui induit }B_{n + 1}/I^nB_{n + 1} \cong B_n
\end{matrix}
\right\}
\end{equation}
Les morphismes de $\mathcal{C}$ sont donnés par des systèmes
d'homomorphismes. Soit $\mathcal{C}'$ la catégorie
\begin{equation}
\label{restricted-equation-C-prime}
\mathcal{C}' =
\left\{
\begin{matrix}
A\text{-algèbres }B\text{ qui sont complètes pour la topologie }I\text{-adique}\\
\text{et telles que }B/IB\text{ soit de type fini sur }A/I
\end{matrix}
\right\}
\end{equation}
Les morphismes de $\mathcal{C}'$ sont les homomorphismes de $A$-algèbres.
Il existe un foncteur
\begin{equation}
\label{restricted-equation-from-complete-to-systems}
\mathcal{C}' \longrightarrow \mathcal{C},\quad
B \longmapsto (B/I^nB)
\end{equation}
En effet, puisque $B/IB$ est de type fini sur $A/I$, les homomorphismes
d'anneaux $A_n = A/I^n \to B/I^nB$ sont de type fini d'après
Algèbre, lemme \ref{algebra-lemma-finite-type-mod-nilpotent}.

\begin{lemma}
\label{restricted-lemma-topologically-finite-type}
Soient $A$ un anneau et $I \subset A$ un idéal de type fini. Le foncteur
$$
\mathcal{C} \longrightarrow \mathcal{C}',\quad
(B_n) \longmapsto B = \lim B_n
$$
est un quasi-inverse de (\ref{restricted-equation-from-complete-to-systems}).
Les complétés $A[x_1, \ldots, x_r]^\wedge$ appartiennent à
$\mathcal{C}'$, et tout objet de $\mathcal{C}'$ est de la forme
$$
B = A[x_1, \ldots, x_r]^\wedge / J
$$
pour un certain idéal $J \subset A[x_1, \ldots, x_r]^\wedge$.
\end{lemma}

\begin{proof}
Soit $(B_n)$ un objet de $\mathcal{C}$. D'après
Algèbre, lemme \ref{algebra-lemma-limit-complete},
$B = \lim B_n$ est complète pour la topologie $I$-adique et
$B/I^nB = B_n$. Ainsi, $B$ est un objet de $\mathcal{C}'$ et l'on peut
retrouver l'objet $(B_n)$ en prenant les quotients. Réciproquement, si $B$ est
un objet de $\mathcal{C}'$, alors $B = \lim B/I^nB$ par hypothèse. Par
conséquent, $B \mapsto (B/I^nB)$ est un quasi-inverse du foncteur du lemme.

\medskip\noindent
Comme $A[x_1, \ldots, x_r]^\wedge = \lim A_n[x_1, \ldots, x_r]$,
c'est un objet de $\mathcal{C}'$ d'après la première assertion du lemme.
Enfin, soit $B$ un objet de $\mathcal{C}'$. Choisissons
$b_1, \ldots, b_r \in B$ dont les images dans $B/IB$ engendrent $B/IB$
comme algèbre sur $A/I$. Puisque $B$ est complète pour la topologie
$I$-adique, l'homomorphisme de $A$-algèbres
$A[x_1, \ldots, x_r] \to B$, $x_i \mapsto b_i$, se prolonge en un
homomorphisme de $A$-algèbres $A[x_1, \ldots, x_r]^\wedge \to B$.
Pour achever la démonstration, il faut montrer que cet homomorphisme est
surjectif. Cela résulte d'Algèbre, lemme
\ref{algebra-lemma-completion-generalities}, car notre homomorphisme
$A[x_1, \ldots, x_r] \to B$ est surjectif modulo $I$ et $B = B^\wedge$.
\end{proof}

\noindent
Signalons que, lorsque $A$ n'est pas noethérien, le quotient d'un objet de
$\mathcal{C}'$ peut ne pas être un objet de $\mathcal{C}'$. Voir Exemples,
lemme \ref{examples-lemma-noncomplete-quotient}. Nous montrons maintenant que
cela ne se produit pas lorsque $A$ est noethérien.

\begin{lemma}
\label{restricted-lemma-topologically-finite-type-Noetherian}
\begin{reference}
\cite[Proposition 7.5.5]{EGA1}
\end{reference}
Soient $A$ un anneau noethérien et $I \subset A$ un idéal. Alors :
\begin{enumerate}
\item tout objet de la catégorie $\mathcal{C}'$
(\ref{restricted-equation-C-prime}) est noethérien ;
\item si $B \in \Ob(\mathcal{C}')$ et si $J \subset B$ est un idéal,
alors $B/J$ est un objet de $\mathcal{C}'$ ;
\item pour toute $A$-algèbre de type fini $C$, le complété $I$-adique
$C^\wedge$ appartient à $\mathcal{C}'$ ;
\item en particulier, le complété $A[x_1, \ldots, x_r]^\wedge$
appartient à $\mathcal{C}'$.
\end{enumerate}
\end{lemma}

\begin{proof}
L'assertion (4) résulte d'Algèbre, lemme
\ref{algebra-lemma-completion-Noetherian-Noetherian}, puisque
$A[x_1, \ldots, x_r]$ est noethérien (Algèbre, lemme
\ref{algebra-lemma-Noetherian-permanence}). Pour démontrer (1), le lemme
\ref{restricted-lemma-topologically-finite-type} nous ramène au cas du complété de
l'anneau de polynômes, que nous venons de traiter. L'assertion (2) résulte
d'Algèbre, lemme \ref{algebra-lemma-completion-tensor}, qui affirme que tout
$B$-module fini est complet pour la topologie $IB$-adique. L'assertion (3) se
démontre de la même manière que l'assertion (4).
\end{proof}

\begin{remark}[Changement de base]
\label{restricted-remark-base-change}
Soient $\varphi : A_1 \to A_2$ un homomorphisme d'anneaux et
$I_i \subset A_i$ des idéaux tels que $\varphi(I_1^c) \subset I_2$ pour un
certain $c \geq 1$. Cela induit des homomorphismes d'anneaux
$A_{1, cn} = A_1/I_1^{cn} \to A_2/I_2^n = A_{2, n}$ pour tout $n \geq 1$.
Soit $\mathcal{C}_i$ la catégorie (\ref{restricted-equation-C}) associée à $(A_i, I_i)$.
Il existe un foncteur de changement de base
\begin{equation}
\label{restricted-equation-base-change-systems}
\mathcal{C}_1 \longrightarrow \mathcal{C}_2,\quad
(B_n) \longmapsto (B_{cn} \otimes_{A_{1, cn}} A_{2, n})
\end{equation}
Soit $\mathcal{C}_i'$ la catégorie (\ref{restricted-equation-C-prime}) associée à
$(A_i, I_i)$. Si $I_2$ est de type fini, il existe alors un foncteur de
changement de base
\begin{equation}
\label{restricted-equation-base-change-complete}
\mathcal{C}_1' \longrightarrow \mathcal{C}_2',\quad
B \longmapsto (B \otimes_{A_1} A_2)^\wedge
\end{equation}
car, dans ce cas, le complété est complet (Algèbre, lemme
\ref{algebra-lemma-hathat-finitely-generated}). Si $I_1$ et $I_2$ sont tous
deux de type fini, alors les deux foncteurs de changement de base coïncident
par l'intermédiaire des foncteurs
(\ref{restricted-equation-from-complete-to-systems}), qui sont des équivalences d'après
le lemme \ref{restricted-lemma-topologically-finite-type}.
\end{remark}

\begin{remark}[Changement de base par une immersion fermée]
\label{restricted-remark-take-bar}
Soient $A$ un anneau noethérien et $I \subset A$ un idéal. Soit
$\mathfrak a \subset A$ un idéal. Posons $\bar A = A/\mathfrak a$. Soit
$\bar I \subset \bar A$ un idéal tel que $I^c \bar A \subset \bar I$ et
$\bar I^d \subset I\bar A$ pour certains $c, d \geq 1$. Dans ce cas, le
foncteur de changement de base (\ref{restricted-equation-base-change-complete}) de
$(A, I)$ vers $(\bar A, \bar I)$ est donné par
$B \mapsto \bar B = B/\mathfrak aB$. En effet, on a
\begin{equation}
\label{restricted-equation-base-change-to-closed}
\bar B = (B \otimes_A \bar A)^\wedge = (B/\mathfrak a B)^\wedge =
B/\mathfrak a B
\end{equation}
la dernière égalité résultant du fait que tout $B$-module fini est complet
pour la topologie $I$-adique d'après Algèbre, lemme
\ref{algebra-lemma-completion-tensor}, et que, s'il est annulé par
$\mathfrak a$, il est aussi complet pour la topologie $\bar I$-adique d'après
Algèbre, lemme \ref{algebra-lemma-change-ideal-completion}.
\end{remark}







\section{Un complexe cotangent naïf}
\label{restricted-section-naive-cotangent-complex}

\noindent
Soient $A$ un anneau noethérien et $I \subset A$ un idéal. Soit $B$ une
$A$-algèbre complète pour la topologie $I$-adique et telle que
$A/I \to B/IB$ soit de type fini, c'est-à-dire un objet de
(\ref{restricted-equation-C-prime}). D'après le lemme
\ref{restricted-lemma-topologically-finite-type-Noetherian}, on peut écrire
$$
B = A[x_1, \ldots, x_r]^\wedge / J
$$
pour un certain idéal de type fini $J$. Pour un choix de présentation comme
ci-dessus, nous définissons dans ce cadre le {\it complexe cotangent naïf} par
la formule
\begin{equation}
\label{restricted-equation-NL}
\NL_{B/A}^\wedge = (J/J^2 \longrightarrow \bigoplus B\text{d}x_i)
\end{equation}
dont les termes sont placés en degrés $-1$ et $0$, l'homomorphisme envoyant
la classe résiduelle de $g \in J$ sur la différentielle
$\text{d}g = \sum (\partial g/\partial x_i) \text{d}x_i$. La dérivée
partielle est ici calculée en considérant $g$ comme une série formelle. Le
lemme suivant montre que $\NL_{B/A}^\wedge$ est bien défini à homotopie près.

\begin{lemma}
\label{restricted-lemma-NL-up-to-homotopy}
Soient $A$ un anneau noethérien et $I \subset A$ un idéal. Soit $B$ un objet
de (\ref{restricted-equation-C-prime}). Le complexe cotangent naïf
$\NL_{B/A}^\wedge$ est bien défini dans $K(B)$.
\end{lemma}

\begin{proof}
Le lemme signifie que, pour une seconde présentation
$B = A[y_1, \ldots, y_s]^\wedge / K$, les complexes de $B$-modules
$$
(J/J^2 \to B\text{d}x_i)
\quad\text{et}\quad
(K/K^2 \to \bigoplus B\text{d}y_j)
$$
sont homotopiquement équivalents. Pour le voir, on peut raisonner exactement
comme dans la démonstration d'Algèbre, lemme
\ref{algebra-lemma-NL-homotopy}.

\medskip\noindent
Étape 1. Si nous choisissons
$g_i(y_1, \ldots, y_s) \in A[y_1, \ldots, y_s]^\wedge$ ayant pour image
l'image de $x_i$ dans $B$, nous obtenons un unique homomorphisme continu de
$A$-algèbres
$$
A[x_1, \ldots, x_r]^\wedge \to A[y_1, \ldots, y_s]^\wedge,\quad
x_i \mapsto g_i(y_1, \ldots, y_s)
$$
compatible avec les surjections données sur $B$. Un tel homomorphisme est
appelé un morphisme de présentations. Il induit un homomorphisme de $J$ dans
$K$, donc un homomorphisme de $B$-modules $J/J^2 \to K/K^2$. En envoyant
$\text{d}x_i$ sur $\sum (\partial g_i/\partial y_j)\text{d}y_j$, nous
obtenons un morphisme de complexes
$$
(J/J^2 \to \bigoplus B\text{d}x_i)
\longrightarrow
(K/K^2 \to \bigoplus B\text{d}y_j)
$$
Bien entendu, on peut procéder de même en échangeant les rôles des deux
présentations, ce qui donne un morphisme de complexes en sens inverse.

\medskip\noindent
Étape 2. La construction précédente est compatible avec la composition des
morphismes de présentations. Pour achever la démonstration, il suffit donc de
montrer que, étant donné
$g_i(x_1, \ldots, x_r) \in A[x_1, \ldots, x_n]^\wedge$
ayant pour image celle de $x_i$ dans $B$, le morphisme de complexes induit
$$
(J/J^2 \to \bigoplus B\text{d}x_i)
\longrightarrow
(J/J^2 \to \bigoplus B\text{d}x_i)
$$
est homotope à l'identité. Pour le voir, considérons l'homomorphisme
$h : \bigoplus B \text{d}x_i \to J/J^2$ défini par
$\text{d}x_i \mapsto g_i(x_1, \ldots, x_n) - x_i$, puis calculons.
\end{proof}

\begin{lemma}
\label{restricted-lemma-NL-is-completion}
Soient $A$ un anneau noethérien et $I \subset A$ un idéal. Soit
$A \to B$ un homomorphisme d'anneaux de type fini. Choisissons une
présentation $\alpha : A[x_1, \ldots, x_n] \to B$. Alors
$\NL_{B^\wedge/A}^\wedge = \lim \NL(\alpha) \otimes_B B^\wedge$
comme complexes, et
$\NL_{B^\wedge/A}^\wedge = \NL_{B/A} \otimes_B^\mathbf{L} B^\wedge$
dans $D(B^\wedge)$.
\end{lemma}

\begin{proof}
L'énoncé a un sens puisque $B^\wedge$ est un objet de
(\ref{restricted-equation-C-prime}) d'après le lemme
\ref{restricted-lemma-topologically-finite-type-Noetherian}. Posons
$J = \Ker(\alpha)$. Le foncteur de complétion $I$-adique est exact sur les
modules finis sur $A[x_1, \ldots, x_n]$ et coïncide avec le foncteur
$M \mapsto M \otimes_{A[x_1, \ldots, x_n]} A[x_1, \ldots, x_n]^\wedge$ ;
voir Algèbre, lemmes \ref{algebra-lemma-completion-tensor} et
\ref{algebra-lemma-completion-flat}. En outre, les homomorphismes d'anneaux
$A[x_1, \ldots, x_n] \to A[x_1, \ldots, x_n]^\wedge$
et $B \to B^\wedge$ sont plats. Par conséquent,
$B^\wedge = A[x_1, \ldots, x_n]^\wedge / J^\wedge$ et
$$
(J/J^2) \otimes_B B^\wedge = (J/J^2)^\wedge = J^\wedge/(J^\wedge)^2
$$
Puisque $\NL(\alpha) = (J/J^2 \to \bigoplus B\text{d}x_i)$, voir Algèbre,
section \ref{algebra-section-netherlander}, nous en concluons que le complexe
$\NL_{B^\wedge/A}^\wedge$ est égal à
$\NL(\alpha) \otimes_B B^\wedge$. La dernière assertion résulte du fait que
$\NL_{B/A}$ est homotopiquement équivalent à $\NL(\alpha)$ et que
l'homomorphisme d'anneaux $B \to B^\wedge$ est plat (le changement de base
dérivé le long de $B \to B^\wedge$ n'est donc que le changement de base).
\end{proof}

\begin{lemma}
\label{restricted-lemma-NL-is-limit}
Soient $A$ un anneau noethérien et $I \subset A$ un idéal. Soit $B$ un objet
de (\ref{restricted-equation-C-prime}). Alors :
\begin{enumerate}
\item les pro-objets
$\{\NL_{B/A}^\wedge \otimes_B B/I^nB\}$ et $\{\NL_{B_n/A_n}\}$
sont strictement isomorphes dans $D(B)$ (voir la démonstration pour des
précisions) ;
\item $\NL_{B/A}^\wedge = R\lim \NL_{B_n/A_n}$ dans $D(B)$.
\end{enumerate}
Ici, $B_n$ et $A_n$ sont définis comme dans la section
\ref{restricted-section-two-categories}.
\end{lemma}

\begin{proof}
L'énoncé signifie ceci : pour tout $n$, nous disposons d'un complexe bien
défini $\NL_{B_n/A_n}$ de $B_n$-modules et de morphismes de transition
$\NL_{B_{n + 1}/A_{n + 1}} \to \NL_{B_n/A_n}$. Voir Algèbre, section
\ref{algebra-section-netherlander}. Nous pouvons donc considérer
$$
\ldots \to \NL_{B_3/A_3} \to \NL_{B_2/A_2} \to \NL_{B_1/A_1}
$$
comme un système projectif de complexes de $B$-modules et, a fortiori, comme
un système projectif dans $D(B)$. En outre, $R\lim \NL_{B_n/A_n}$ est une
limite homotopique de ce système projectif ; voir Catégories dérivées, section
\ref{derived-section-derived-limit}.

\medskip\noindent
Choisissons une présentation $B = A[x_1, \ldots, x_r]^\wedge / J$. Elle
définit des présentations
$$
B_n = B/I^nB = A_n[x_1, \ldots, x_r]/J_n
$$
où
$$
J_n = JA_n[x_1, \ldots, x_r] =
J/(J \cap I^nA[x_1, \ldots, x_r]^\wedge)
$$
Le complexe à deux termes
$J_n/J_n^2 \longrightarrow \bigoplus B_n \text{d}x_i$ représente
$\NL_{B_n/A_n}$ ; voir Algèbre, section
\ref{algebra-section-netherlander}. Le lemme d'Artin--Rees (Algèbre, lemme
\ref{algebra-lemma-Artin-Rees}), appliqué dans l'anneau noethérien
$A[x_1, \ldots, x_r]^\wedge$ (lemme
\ref{restricted-lemma-topologically-finite-type-Noetherian}), fournit un $c \geq 0$ tel
que l'on ait des surjections canoniques
$$
J/I^nJ \to J_n \to J/I^{n - c}J \to J_{n - c},\quad n \geq c
$$
pour tout $n \geq c$. On vérifie aussitôt que ces homomorphismes sont
compatibles aux différentielles, et l'on obtient des morphismes de complexes
$$
\NL_{B/A}^\wedge \otimes_B B/I^nB \to
\NL_{B_n/A_n} \to
\NL_{B/A}^\wedge \otimes_B B/I^{n - c}B \to
\NL_{B_{n - c}/A_{n - c}}
$$
compatibles aux morphismes de transition des systèmes projectifs
$\{\NL_{B/A}^\wedge \otimes_B B/I^nB\}$ et $\{\NL_{B_n/A_n}\}$. Cela
démontre l'assertion (1) du lemme.

\medskip\noindent
En vertu de l'assertion (1), et puisque des systèmes pro-isomorphes ont le même
$R\lim$, il suffit, pour démontrer (2), de montrer que
$\NL_{B/A}^\wedge$ est égal à
$R\lim \NL_{B/A}^\wedge \otimes_B B/I^nB$. Or $\NL_{B/A}^\wedge$ est un
complexe à deux termes $M^\bullet$ de $B$-modules finis, qui sont complets pour
la topologie $I$-adique, par exemple d'après Algèbre, lemme
\ref{algebra-lemma-completion-tensor}. Par conséquent,
$M^\bullet = \lim M^\bullet/I^nM^\bullet = R\lim M^\bullet/I^n M^\bullet$ ; voir
Compléments d'algèbre, lemme \ref{more-algebra-lemma-compute-Rlim-modules} et
remarque \ref{more-algebra-remark-how-unique}.
\end{proof}

\begin{lemma}
\label{restricted-lemma-NL-base-change}
Soit $(A_1, I_1) \to (A_2, I_2)$ comme dans la remarque
\ref{restricted-remark-base-change}, avec $A_1$ et $A_2$ noethériens. Soit $B_1$ un objet
de (\ref{restricted-equation-C-prime}) pour $(A_1, I_1)$. Soit $B_2$ le changement de
base de $B_1$. Il existe alors un morphisme canonique
$$
\NL_{B_1/A_1} \otimes_{B_2} B_1 \to \NL_{B_2/A_2}
$$
qui induit un isomorphisme sur $H^0$ et une surjection sur $H^{-1}$.
\end{lemma}

\begin{proof}
Choisissons une présentation
$B_1 = A_1[x_1, \ldots, x_r]^\wedge/J_1$. Puisque
$A_2/I_2^n[x_1, \ldots, x_r] =
A_1/I_1^{cn}[x_1, \ldots, x_r] \otimes_{A_1/I_1^{cn}} A_2/I_2^n$
on a
$$
A_2[x_1, \ldots, x_r]^\wedge =
(A_1[x_1, \ldots, x_r]^\wedge \otimes_{A_1} A_2)^\wedge
$$
où l'on prend la complétion $I_2$-adique des deux côtés (mais, bien entendu,
la complétion $I_1$-adique pour
$A_1[x_1, \ldots, x_r]^\wedge$). Posons
$J_2 = J_1 A_2[x_1, \ldots, x_r]^\wedge$. Un raisonnement analogue donne la
présentation
\begin{align*}
B_2
& =
(B_1 \otimes_{A_1} A_2)^\wedge \\
& =
\lim \frac{A_1/I_1^{cn}[x_1, \ldots, x_r]}{J_1(A_1/I_1^{cn}[x_1, \ldots, x_r])}
\otimes_{A_1/I_1^{cn}} A_2/I_2^n \\
& =
\lim \frac{A_2/I_2^n[x_1, \ldots, x_r]}{J_2(A_2/I_2^n[x_1, \ldots, x_r])} \\
& =
A_2[x_1, \ldots, x_r]^\wedge/J_2
\end{align*}
de $B_2$ sur $A_2$. On obtient par conséquent un diagramme commutatif
$$
\xymatrix{
\NL^\wedge_{B_1/A_1} : \ar[d] &
J_1/J_1^2 \ar[r]_-{\text{d}} \ar[d] & \bigoplus B_1\text{d}x_i \ar[d] \\
\NL^\wedge_{B_2/A_2} : &
J_2/J_2^2 \ar[r]^-{\text{d}} & \bigoplus B_2\text{d}x_i
}
$$
La flèche induite $J_1/J_1^2 \otimes_{B_1} B_2 \to J_2/J_2^2$ est surjective,
car $J_2$ est engendré par l'image de $J_1$. Cela détermine la flèche affichée
dans le lemme. Nous omettons de démontrer que cette flèche est bien définie à
homotopie près (c'est-à-dire indépendante, à homotopie près, du choix des
présentations). L'assertion relative au morphisme induit sur les modules de
cohomologie découle aisément de ce qui précède (détails omis).
\end{proof}

\begin{lemma}
\label{restricted-lemma-exact-sequence-NL}
Soient $A$ un anneau noethérien et $I \subset A$ un idéal. Soit $B \to C$ un
morphisme de (\ref{restricted-equation-C-prime}). Il existe alors une suite exacte
$$
\xymatrix{
C \otimes_B H^0(\NL_{B/A}^\wedge) \ar[r] &
H^0(\NL_{C/A}^\wedge) \ar[r] &
H^0(\NL_{C/B}^\wedge) \ar[r] & 0 \\
H^{-1}(\NL_{B/A}^\wedge \otimes_B C) \ar[r] &
H^{-1}(\NL_{C/A}^\wedge) \ar[r] &
H^{-1}(\NL_{C/B}^\wedge) \ar[llu]
}
$$
Voir la démonstration pour des précisions.
\end{lemma}

\begin{proof}
Notons que le produit tensoriel $\NL_{B/A}^\wedge \otimes_B C$ a un sens,
puisque $\NL_{B/A}^\wedge$ est bien défini à homotopie près d'après le lemme
\ref{restricted-lemma-NL-up-to-homotopy}. De plus, $(B, IB)$ est un couple où $B$ est un
anneau noethérien (lemme
\ref{restricted-lemma-topologically-finite-type-Noetherian}), et $C$ appartient à la
catégorie correspondante (\ref{restricted-equation-C-prime}). Tous les termes de la suite
à $6$ termes sont donc bien définis.

\medskip\noindent
Choisissons une présentation $B = A[x_1, \ldots, x_r]^\wedge/J$, puis une
présentation $C = B[y_1, \ldots, y_s]^\wedge/J'$. En combinant ces
présentations, on obtient une présentation
$$
C = A[x_1, \ldots, x_r, y_1, \ldots, y_s]^\wedge/K
$$
Le lecteur vérifie alors que l'on obtient un diagramme commutatif
$$
\xymatrix{
0 \ar[r] &
\bigoplus C \text{d}x_i \ar[r] &
\bigoplus C \text{d}x_i \oplus \bigoplus C \text{d}y_j \ar[r] &
\bigoplus C \text{d}y_j \ar[r] &
0 \\
&
J/J^2 \otimes_B C \ar[r] \ar[u] &
K/K^2 \ar[r] \ar[u] &
J'/(J')^2 \ar[r] \ar[u] &
0
}
$$
à lignes exactes. Notons que la flèche verticale de gauche est le produit
tensoriel de la flèche définissant $\NL_{B/A}^\wedge$ par $\text{id}_C$. Le
lemme résulte alors du lemme du serpent (Algèbre, lemme
\ref{algebra-lemma-snake}).
\end{proof}

\begin{lemma}
\label{restricted-lemma-transitive-lci-at-end}
Sous les hypothèses du lemme \ref{restricted-lemma-exact-sequence-NL}, supposons que
$B/I^nB \to C/I^nC$ soit un homomorphisme d'intersection complète locale pour
tout $n$. Alors
$H^{-1}(\NL_{B/A}^\wedge \otimes_B C) \to H^{-1}(\NL_{C/A}^\wedge)$
est injectif.
\end{lemma}

\begin{proof}
Pour tout $n \geq 1$, posons $A_n = A/I^n$, $B_n = B/I^nB$ et
$C_n = C/I^nC$. On a
\begin{align*}
H^{-1}(\NL_{B/A}^\wedge \otimes_B C)
& =
\lim H^{-1}(\NL_{B/A}^\wedge \otimes_B C_n) \\
& =
\lim H^{-1}(\NL_{B/A}^\wedge \otimes_B B_n \otimes_{B_n} C_n) \\
& =
\lim H^{-1}(\NL_{B_n/A_n} \otimes_{B_n} C_n)
\end{align*}
La première égalité résulte de Compléments d'algèbre, lemme
\ref{more-algebra-lemma-consequence-Artin-Rees}, et du fait que
$H^{-1}(\NL_{B/A}^\wedge \otimes_B C)$ est un $C$-module fini, donc complet
pour la topologie $I$-adique, par exemple d'après Algèbre, lemme
\ref{algebra-lemma-completion-tensor}. La deuxième égalité est immédiate. La
troisième résulte du lemme \ref{restricted-lemma-NL-is-limit}. Les homomorphismes
$H^{-1}(\NL_{B_n/A_n} \otimes_{B_n} C_n) \to H^{-1}(\NL_{C_n/A_n})$ sont
injectifs d'après Compléments d'algèbre, lemme
\ref{more-algebra-lemma-transitive-lci-at-end}. La démonstration s'achève car
nous avons aussi
$H^{-1}(\NL_{C/A}^\wedge) = \lim H^{-1}(\NL_{C_n/A_n})$
par le même raisonnement.
\end{proof}









\section{Algèbres rig-lisses}
\label{restricted-section-rig-smooth}

\noindent
Pour motiver la définition suivante, voir Compléments d'algèbre, remarque
\ref{more-algebra-remark-smoothness-ext-1-zero}.

\begin{definition}
\label{restricted-definition-rig-smooth-homomorphism}
Soient $A$ un anneau noethérien et $I \subset A$ un idéal. Soit $B$ un objet
de (\ref{restricted-equation-C-prime}). Nous disons que $B$ est {\it rig-lisse sur
$(A, I)$} s'il existe un entier $c \geq 0$ tel que $I^c$ annule
$\Ext^1_B(\NL_{B/A}^\wedge, N)$ pour tout $B$-module $N$.
\end{definition}

\noindent
Explicitons cette définition.

\begin{lemma}
\label{restricted-lemma-equivalent-with-artin-smooth}
Soient $A$ un anneau noethérien et $I \subset A$ un idéal. Soit $B$ un objet
de (\ref{restricted-equation-C-prime}). Écrivons
$B = A[x_1, \ldots, x_r]^\wedge/J$
(lemme \ref{restricted-lemma-topologically-finite-type-Noetherian}) et soit
$\NL_{B/A}^\wedge = (J/J^2 \to \bigoplus B\text{d}x_i)$ son complexe
cotangent naïf (\ref{restricted-equation-NL}). Les conditions suivantes sont équivalentes :
\begin{enumerate}
\item $B$ est rig-lisse sur $(A, I)$ ;
\item l'objet $\NL_{B/A}^\wedge$ de $D(B)$ satisfait les conditions
équivalentes (1) -- (6) de Compléments d'algèbre, lemme
\ref{more-algebra-lemma-ext-1-annihilated}, relativement à l'idéal $IB$ ;
\item il existe $c \geq 0$ tel que, pour tout $a \in I^c$, il existe un
homomorphisme $h : \bigoplus B\text{d}x_i \to J/J^2$ tel que
$a : J/J^2 \to J/J^2$ soit égal à $h \circ \text{d}$ ;
\item il existe $b_1, \ldots, b_s \in B$ tels que
$V(b_1, \ldots, b_s) \subset V(IB)$ et tels que, pour tout
$l = 1, \ldots, s$, il existe $m \geq 0$, $f_1, \ldots, f_m \in J$ et une
partie $T \subset \{1, \ldots, n\}$ vérifiant $|T| = m$ et les conditions
suivantes :
\begin{enumerate}
\item $\det_{i \in T, j \leq m}(\partial f_j/ \partial x_i)$
divise $b_l$ dans $B$ ;
\item $b_l J \subset (f_1, \ldots, f_m) + J^2$.
\end{enumerate}
\end{enumerate}
\end{lemma}

\begin{proof}
L'équivalence de (1), (2) et (3) résulte immédiatement de Compléments
d'algèbre, lemme \ref{more-algebra-lemma-ext-1-annihilated}.

\medskip\noindent
Supposons que $b_1, \ldots, b_s$ soient comme en (4). Puisque $B$ est
noethérien, l'inclusion $V(b_1, \ldots, b_s) \subset V(IB)$ implique
$I^cB \subset (b_1, \ldots, b_s)$ pour un certain $c \geq 0$ (par exemple
d'après Algèbre, lemme
\ref{algebra-lemma-Noetherian-power-ideal-kills-module}). Fixons
$1 \leq l \leq s$, ainsi que $m \geq 0$, $f_1, \ldots, f_m \in J$ et
$T \subset \{1, \ldots, n\}$ avec $|T| = m$, satisfaisant (4)(a) et (b).
Après inversion de $b_l$, on voit que
$$
\NL_{B/A}^\wedge \otimes_B B_{b_l} =
\left(
\bigoplus\nolimits_{j \leq m} B_{b_l} f_j
\longrightarrow
\bigoplus\nolimits_{i = 1, \ldots, n} B_{b_l} \text{d}x_i
\right)
$$
et que, de plus, la flèche est isomorphe à l'inclusion du facteur direct
$\bigoplus_{i \in T} B_{b_l} \text{d}x_i$. On en déduit que
$H^{-1}(\NL_{B/A}^\wedge)$ est de torsion annulée par une puissance de $b_l$
et que $H^0(\NL_{B/A}^\wedge)$ devient libre de type fini après inversion de
$b_l$. En combinant ceci avec l'inclusion
$I^cB \subset (b_1, \ldots, b_s)$, on voit que
$H^{-1}(\NL_{B/A}^\wedge)$ est de torsion annulée par une puissance de $IB$.
La condition (4) de Compléments d'algèbre, lemme
\ref{more-algebra-lemma-ext-1-annihilated}, est donc satisfaite. Ainsi, (4)
implique (2).

\medskip\noindent
Supposons satisfaites les conditions équivalentes (1), (2) et (3). Nous allons
démontrer (4), mais nous invitons vivement le lecteur à s'en convaincre par
lui-même. Le complexe $\NL_{B/A}^\wedge$ détermine un objet de
$D^b_{\textit{Coh}}(\Spec(B))$ dont la restriction à l'ouvert de Zariski
$U = \Spec(B) \setminus V(IB)$ est un module localement libre de type fini
$\mathcal{E}$ placé en degré $0$ (cela résulte, par exemple, de la quatrième
condition équivalente de Compléments d'algèbre, lemme
\ref{more-algebra-lemma-ext-1-annihilated}). Choisissons des générateurs
$f_1, \ldots, f_M$ de $J$. Ils déterminent une suite exacte
$$
\bigoplus\nolimits_{j = 1, \ldots, M} \mathcal{O}_U \cdot f_j \to
\bigoplus\nolimits_{i = 1, \ldots, n} \mathcal{O}_U \cdot \text{d}x_i \to
\mathcal{E} \to 0
$$
Soit $U = \bigcup_{l = 1, \ldots, s} U_l$ un recouvrement ouvert affine fini
tel que $\mathcal{E}|_{U_l}$ soit libre de rang $r_l = n - m_l$ pour un entier
$n \geq m_l \geq 0$. En remplaçant chaque $U_l$ par un recouvrement ouvert
affine, on peut supposer qu'il existe une partie
$T_l \subset \{1, \ldots, n\}$ telle que les éléments $\text{d}x_i$,
$i \in \{1, \ldots, n\} \setminus T_l$, aient pour images une base de
$\mathcal{E}|_{U_l}$. En répétant l'argument, on peut supposer qu'il existe une
partie $T'_l \subset \{1, \ldots, M\}$ de cardinal $m_l$ telle que les images
des $f_j$, $j \in T'_l$, forment une base du noyau de
$\mathcal{O}_{U_l} \cdot \text{d}x_i \to \mathcal{E}|_{U_l}$. Enfin, puisque
le recouvrement ouvert $U = \bigcup U_l$ peut être raffiné par un recouvrement
par des ouverts principaux (Algèbre, lemme
\ref{algebra-lemma-Zariski-topology}), on peut supposer que
$U_l = D(g_l)$ pour un certain $g_l \in B$. En particulier,
$V(g_1, \ldots, g_s) = V(IB)$. Un argument d'algèbre linéaire utilisant les
choix précédents montre que
$\det_{i \in T_l, j \in T'_l}(\partial f_j/ \partial x_i)$
a pour image un élément inversible de $B_{b_l}$. De même, l'annulation de la
cohomologie de $\NL_{B/A}^\wedge$ en degré $-1$ sur $U_l$ montre que
$J/J^2 + (f_j; j \in T')$ est annulé par une puissance de $b_l$. En remplaçant
chaque $g_l$ par une puissance convenable, on obtient les conditions (4)(a) et
(4)(b) du lemme. Certains détails sont omis.
\end{proof}

\begin{lemma}
\label{restricted-lemma-rig-smooth}
Soient $A$ un anneau noethérien et $I$ un idéal. Soit $B$ une $A$-algèbre de
type fini.
\begin{enumerate}
\item Si $\Spec(B) \to \Spec(A)$ est lisse au-dessus de
$\Spec(A) \setminus V(I)$, alors $B^\wedge$ est rig-lisse sur $(A, I)$.
\item Si $B^\wedge$ est rig-lisse sur $(A, I)$, alors il existe
$g \in 1 + IB$ tel que $\Spec(B_g)$ soit lisse au-dessus de
$\Spec(A) \setminus V(I)$.
\end{enumerate}
\end{lemma}

\begin{proof}
Nous utiliserons le lemme \ref{restricted-lemma-equivalent-with-artin-smooth} sans plus
de commentaires.

\medskip\noindent
Supposons (1). Rappelons que la formation de $\NL_{B/A}$ commute à la
localisation ; voir Algèbre, lemme \ref{algebra-lemma-localize-NL}. Par la
définition même des homomorphismes d'anneaux lisses (selon laquelle le complexe
cotangent naïf est quasi-isomorphe à un module projectif de type fini placé en
degré $0$), on voit donc que $\NL_{B/A}$ satisfait la quatrième condition
équivalente de Compléments d'algèbre, lemme
\ref{more-algebra-lemma-ext-1-annihilated}, relativement à l'idéal $IB$ (un
petit détail est omis). Puisque
$\NL_{B^\wedge/A}^\wedge = \NL_{B/A} \otimes_B B^\wedge$ d'après le lemme
\ref{restricted-lemma-NL-is-completion}, on en conclut que (2) est satisfaite d'après
Compléments d'algèbre, lemme
\ref{more-algebra-lemma-base-change-property-ext-1-annihilated}.

\medskip\noindent
Supposons (2). Choisissons une présentation
$B = A[x_1, \ldots, x_n]/J$, posons $N = J/J^2$ et considérons l'élément
$\xi \in \Ext^1_B(\NL_{B/A}, J/J^2)$ déterminé par l'identité de $J/J^2$.
En utilisant de nouveau l'égalité
$\NL_{B^\wedge/A}^\wedge = \NL_{B/A} \otimes_B B^\wedge$, on voit que notre
hypothèse implique que l'image
$$
\xi \otimes 1 \in
\Ext^1_{B^\wedge}(\NL_{B/A} \otimes_B B^\wedge, N \otimes_B B^\wedge) =
\Ext^1_{B^\wedge}(\NL_{B/A}, N) \otimes_B B^\wedge
$$
est annulée par $I^c$ pour un certain entier $c \geq 0$. L'égalité résulte,
par exemple, de Compléments d'algèbre, lemme
\ref{more-algebra-lemma-base-change-RHom} (mais elle se déduit aussi aisément
du lemme beaucoup plus simple de Compléments d'algèbre,
\ref{more-algebra-lemma-pseudo-coherence-and-base-change-ext}). Ainsi,
$M = I^cB\xi \subset \Ext^1_B(\NL_{B/A}, N)$ est un sous-module fini dont
l'image dans $\Ext^1_B(\NL_{B/A}, N) \otimes_B B^\wedge$ est nulle. Comme
$B \to B^\wedge$ est plat, cela signifie que $M \otimes_B B^\wedge$ est nul.
Le lemme de Nakayama (Algèbre, lemme \ref{algebra-lemma-NAK}) implique alors
que $M = I^cB\xi$ est annulé par un élément de la forme $g = 1 + x$ avec
$x \in IB$. Il s'ensuit que, pour tout $b \in I^cB$, il existe une flèche
pointillée $B$-linéaire rendant le diagramme commutatif
$$
\xymatrix{
J/J^2 \ar[r] \ar[d]^b & \bigoplus B\text{d}x_i \ar@{..>}[d]^h \\
J/J^2 \ar[r] & (J/J^2)_g
}
$$
Ainsi, $(\NL_{B/A})_{gb}$ est quasi-isomorphe à un module projectif de type
fini ; un petit détail est omis. Comme
$(\NL_{B/A})_{gb} = \NL_{B_{gb}/A}$ dans $D(B_{gb})$, cela montre que
$B_{gb}$ est lisse sur $\Spec(A)$. Puisque ceci vaut pour tout $b \in I^cB$,
on en conclut que $\Spec(B_g) \to \Spec(A)$ est lisse au-dessus de
$\Spec(A) \setminus V(I)$, comme souhaité.
\end{proof}

\begin{lemma}
\label{restricted-lemma-zero-ext-1-after-modding-out}
Soit $(A_1, I_1) \to (A_2, I_2)$ comme dans la remarque
\ref{restricted-remark-base-change}, avec $A_1$ et $A_2$ noethériens. Soit $B_1$ un objet
de (\ref{restricted-equation-C-prime}) pour $(A_1, I_1)$. Soit $B_2$ le changement de
base de $B_1$. Soient $f_1 \in B_1$ et son image $f_2 \in B_2$. Si
$\Ext^1_{B_1}(\NL_{B_1/A_1}^\wedge, N_1)$ est annulé par $f_1$ pour tout
$B_1$-module $N_1$, alors $\Ext^1_{B_2}(\NL_{B_2/A_2}^\wedge, N_2)$ est
annulé par $f_2$ pour tout $B_2$-module $N_2$.
\end{lemma}

\begin{proof}
D'après le lemme \ref{restricted-lemma-NL-base-change}, il existe un morphisme
$$
\NL_{B_1/A_1} \otimes_{B_2} B_1 \to \NL_{B_2/A_2}
$$
qui induit un isomorphisme sur $H^0$ et une surjection sur $H^{-1}$. Le
résultat découle donc des lemmes de Compléments d'algèbre
\ref{more-algebra-lemma-two-term-base-change},
\ref{more-algebra-lemma-base-change-property-ext-1-annihilated} et
\ref{more-algebra-lemma-surjection-property-ext-1-annihilated}
--- les deux derniers étant appliqués aux idéaux principaux
$(f_1) \subset B_1$ et $(f_2) \subset B_2$.
\end{proof}

\begin{lemma}
\label{restricted-lemma-base-change-rig-smooth-homomorphism}
Soit $A_1 \to A_2$ un homomorphisme d'anneaux noethériens. Soit
$I_i \subset A_i$ un idéal tel que $V(I_1A_2) = V(I_2)$. Soit $B_1$ un objet
de (\ref{restricted-equation-C-prime}) pour $(A_1, I_1)$. Soit $B_2$ le changement de
base de $B_1$, comme dans la remarque \ref{restricted-remark-base-change}. Si $B_1$ est
rig-lisse sur $(A_1, I_1)$, alors $B_2$ est rig-lisse sur $(A_2, I_2)$.
\end{lemma}

\begin{proof}
Cela résulte du lemme \ref{restricted-lemma-zero-ext-1-after-modding-out}, de la
définition \ref{restricted-definition-rig-smooth-homomorphism} et du fait que $I_2^c$ est
contenu dans $I_1A_2$ pour un certain $c \geq 0$, puisque $A_2$ est
noethérien.
\end{proof}













\section{Déformations d'homomorphismes d'anneaux}
\label{restricted-section-defos-ring-maps}

\noindent
Quelques résultats sur le relèvement des homomorphismes d'anneaux issus
d'algèbres rig-lisses.

\begin{remark}[Approximation linéaire]
\label{restricted-remark-linear-approximation}
Soient $A$ un anneau et $I \subset A$ un idéal de type fini. Soit $C$ une
algèbre complète pour la topologie $I$-adique sur $A$. Soit
$\psi : A[x_1, \ldots, x_r]^\wedge \to C$ un homomorphisme continu de
$A$-algèbres. Supposons donnés $\delta_i \in C$, $i = 1, \ldots, r$. On peut
alors considérer
$$
\psi' : A[x_1, \ldots, x_r]^\wedge \to C,\quad
x_i \longmapsto \psi(x_i) + \delta_i
$$
voir Espaces formels, remarque
\ref{formal-spaces-remark-universal-property}. On a alors
$$
\psi'(g) = \psi(g) + \sum \psi(\partial g/\partial x_i)\delta_i + \xi
$$
avec un terme d'erreur $\xi \in (\delta_i\delta_j)$. Cela se voit en écrivant
$g$ comme une série formelle et en procédant terme à terme. La convergence est
automatique puisque les coefficients de $g$ tendent vers zéro. Les détails
sont omis.
\end{remark}

\begin{remark}[Relèvement d'homomorphismes]
\label{restricted-remark-improve-homomorphism}
Soient $A$ un anneau noethérien et $I \subset A$ un idéal. Soit $B$ un objet
de (\ref{restricted-equation-C-prime}). Soit $C$ une algèbre complète pour la topologie
$I$-adique sur $A$. Soit $\psi_n : B \to C/I^nC$ un homomorphisme de
$A$-algèbres. L'obstruction au relèvement de $\psi_n$ en un homomorphisme de
$A$-algèbres à valeurs dans $C/I^{2n}C$ est un élément
$$
o(\psi_n) \in \Ext^1_B(\NL_{B/A}^\wedge, I^nC/I^{2n}C)
$$
comme nous allons l'expliquer. Choisissons une présentation
$B = A[x_1, \ldots, x_r]^\wedge/J$, puis un relèvement
$\psi : A[x_1, \ldots, x_r]^\wedge \to C$ de $\psi_n$. Puisque
$\psi(J) \subset I^nC$, on a $\psi(J^2) \subset I^{2n}C$, et donc un
homomorphisme $B$-linéaire
$$
o(\psi) :
J/J^2 \longrightarrow I^nC/I^{2n}C, \quad g \longmapsto \psi(g)
$$
qui se prolonge évidemment en un homomorphisme $C$-linéaire
$J/J^2 \otimes_B C \to I^nC/I^{2n}C$. Puisque
$\NL_{B/A}^\wedge = (J/J^2 \to \bigoplus B \text{d}x_i)$, on obtient
$o(\psi_n)$ comme image de $o(\psi)$ par l'identification
\begin{align*}
& \Ext^1_B(\NL_{B/A}^\wedge, I^nC/I^{2n}C) \\
& =
\Coker\left(\Hom_B(\bigoplus B\text{d}x_i, I^nC/I^{2n}C) \to
\Hom_B(J/J^2, I^nC/I^{2n}C)\right)
\end{align*}
Pour cette égalité, voir Compléments d'algèbre, lemme
\ref{more-algebra-lemma-map-out-of-almost-free}, assertion (1).

\medskip\noindent
Supposons que $o(\psi_n)$ ait une image nulle dans
$\Ext^1_B(\NL_{B/A}^\wedge, I^{n'}C/I^{2n'}C)$
pour un entier $n'$ tel que $n > n' > n/2$. Nous affirmons que cela permet de
trouver un homomorphisme de $A$-algèbres
$\psi'_{2n'} : B \to C/I^{2n'}C$ qui coïncide avec $\psi_n$ comme
homomorphisme à valeurs dans $C/I^{n'}C$. Le cas extrême $n' = n$ explique
pourquoi nous avons dit plus haut que $o(\psi_n)$ est l'obstruction au
relèvement de $\psi_n$ à $C/I^{2n}C$. Démontrons l'affirmation. L'hypothèse
selon laquelle l'image de $o(\psi_n)$ est nulle fournit un homomorphisme de
$B$-modules
$$
h : \bigoplus B\text{d}x_i \longrightarrow I^{n'}C/I^{2n'}C
$$
tel que $o(\psi)$ et $h \circ \text{d}$ coïncident comme homomorphismes à
valeurs dans $I^{n'}C/I^{2n'}C$. Écrivons
$h(\text{d}x_i) = \delta_i \bmod I^{2n'}C$ pour certains
$\delta_i \in I^{n'}C$. Considérons alors l'homomorphisme
$$
\psi' : A[x_1, \ldots, x_r]^\wedge \to C,\quad
x_i \longmapsto \psi(x_i) - \delta_i
$$
Un calcul sur les séries formelles montre que
$\psi'(J) \subset I^{2n'}C$. En effet, pour $g \in J$, on obtient
$$
\psi'(g) \equiv
\psi(g) - \sum \psi(\partial g/\partial x_i)\delta_i \equiv
o(\psi)(g) - (h \circ \text{d})(g) \equiv
0 \bmod I^{2n'}C
$$
Pour la première égalité, voir la remarque
\ref{restricted-remark-linear-approximation}. Par conséquent, $\psi'$ induit un
homomorphisme de $A$-algèbres $\psi'_{2n'} : B \to C/I^{2n'}C$, comme voulu.
\end{remark}

\begin{lemma}
\label{restricted-lemma-get-morphism-general-better}
Supposons données les données suivantes :
\begin{enumerate}
\item un entier $c \geq 0$ ;
\item un idéal $I$ d'un anneau noethérien $A$ ;
\item un objet $B$ de (\ref{restricted-equation-C-prime}) pour $(A, I)$ tel que $I^c$
annule $\Ext^1_B(\NL_{B/A}^\wedge, N)$ pour tout $B$-module $N$ ;
\item une algèbre noethérienne complète pour la topologie $I$-adique, munie
d'une structure de $A$-algèbre et notée $C$ ; notons
$d = d(\text{Gr}_I(C))$ et
$q_0 = q(\text{Gr}_I(C))$ les entiers obtenus dans Cohomologie locale,
section \ref{local-cohomology-section-uniform} ;
\item un entier $n \geq \max(q_0 + (d + 1)c, 2(d + 1)c + 1)$ ;
\item un homomorphisme de $A$-algèbres $\psi_n : B \to C/I^nC$.
\end{enumerate}
Il existe alors un homomorphisme $\varphi : B \to C$ de $A$-algèbres tel que
$\psi_n \bmod I^{n - (d + 1)c} = \varphi \bmod I^{n - (d + 1)c}$.
\end{lemma}

\begin{proof}
Considérons la classe d'obstruction
$$
o(\psi_n) \in \Ext^1_B(\NL_{B/A}^\wedge, I^nC/I^{2n}C)
$$
de la remarque \ref{restricted-remark-improve-homomorphism}. Pour tout $C/I^nC$-module
$N$, on a
\begin{align*}
\Ext^1_B(\NL_{B/A}^\wedge, N)
& =
\Ext^1_{C/I^nC}(\NL_{B/A}^\wedge \otimes_B^\mathbf{L} C/I^nC, N) \\
& =
\Ext^1_{C/I^nC}(\NL_{B/A}^\wedge \otimes_B C/I^nC, N)
\end{align*}
La première égalité résulte de Compléments d'algèbre, lemme
\ref{more-algebra-lemma-upgrade-adjoint-tensor-RHom}, et la seconde de
Compléments d'algèbre, lemme
\ref{more-algebra-lemma-two-term-base-change}. En particulier,
$\Ext^1_{C/I^nC}(\NL_{B/A}^\wedge \otimes_B C/I^nC, N)$ est annulé par
$I^cC$ pour tout $C/I^nC$-module $N$. On peut donc appliquer Cohomologie
locale, lemme \ref{local-cohomology-lemma-bound-two-term-complex}, pour voir
que l'image de $o(\psi_n)$ est nulle dans
$$
\Ext^1_{C/I^nC}(\NL_{B/A}^\wedge \otimes_B C/I^nC, I^{n'}C/I^{2n'}C) =
\Ext^1_B(\NL_{B/A}^\wedge, I^{n'}C/I^{2n'}C) =
$$
où $n' = n - (d + 1)c$. D'après la discussion de la remarque
\ref{restricted-remark-improve-homomorphism}, on obtient un homomorphisme
$$
\psi'_{2n'} : B \to C/I^{2n'}C
$$
qui coïncide avec $\psi_n$ modulo $I^{n'}$. Notons que $2n' > n$ puisque
$n \geq 2(d + 1)c + 1$.

\medskip\noindent
On peut répéter ce procédé. En partant de $n_0 = n$ et de
$\psi^0 = \psi_n$, on obtient une suite strictement croissante d'entiers
$$
n_0 < n_1 < n_2 < \ldots
$$
et des homomorphismes de $A$-algèbres $\psi^i : B \to C/I^{n_i}C$ tels que
$\psi^{i + 1}$ et $\psi^i$ coïncident modulo $I^{n_i - tc}$. Puisque $C$ est
complète pour la topologie $I$-adique, on peut prendre pour $\varphi$ la
limite des homomorphismes
$\psi^i \bmod I^{n_i - (d + 1)c} : B \to C/I^{n_i - (d + 1)c}C$
et le lemme en résulte.
\end{proof}

\noindent
Nous suggérons au lecteur de passer directement à la section suivante. En
effet, les deux lemmes qui suivent sont des conséquences du résultat précédent
si l'on y suppose l'algèbre $C$ noethérienne.

\begin{lemma}
\label{restricted-lemma-get-morphism-nonzerodivisor}
Soit $I = (a)$ un idéal principal d'un anneau noethérien $A$. Soit $B$ un
objet de (\ref{restricted-equation-C-prime}). Supposons donné un entier $c \geq 0$ tel
que $\Ext^1_B(\NL_{B/A}^\wedge, N)$ soit annulé par $a^c$ pour tout
$B$-module $N$. Soit $C$ une algèbre complète pour la topologie $I$-adique
sur $A$, telle que $a$ soit non-diviseur de zéro dans $C$. Soit $n > 2c$.
Pour tout homomorphisme de $A$-algèbres $\psi_n : B \to C/a^nC$, il existe
un homomorphisme de $A$-algèbres $\varphi : B \to C$ tel que
$\psi_n \bmod a^{n - c}C = \varphi \bmod a^{n - c}C$.
\end{lemma}

\begin{proof}
Considérons la classe d'obstruction
$$
o(\psi_n) \in \Ext^1_B(\NL_{B/A}^\wedge, a^nC/a^{2n}C)
$$
de la remarque \ref{restricted-remark-improve-homomorphism}. Puisque $a$ est
non-diviseur de zéro dans $C$, l'homomorphisme
$a^c : a^nC/a^{2n}C \to a^nC/a^{2n}C$ est isomorphe à l'homomorphisme
$a^nC/a^{2n}C \to a^{n - c}C/a^{2n - c}C$ dans la catégorie des $C$-modules.
Notre hypothèse sur $\NL_{B/A}^\wedge$ implique donc que l'image de la classe
$o(\psi_n)$ est nulle dans
$$
\Ext^1_B(\NL_{B/A}^\wedge, a^{n - c}C/a^{2n - c}C)
$$
et, a fortiori, dans
$$
\Ext^1_B(\NL_{B/A}^\wedge, a^{n - c}C/a^{2n - 2c}C)
$$
D'après la discussion de la remarque \ref{restricted-remark-improve-homomorphism}, on
obtient un homomorphisme
$$
\psi_{2n - 2c} : B \to C/a^{2n - 2c}C
$$
qui coïncide avec $\psi_n$ modulo $a^{n - c}C$. Notons que
$2n - 2c > n$ puisque $n > 2c$.

\medskip\noindent
On peut répéter ce procédé. En partant de $n_0 = n$ et de
$\psi^0 = \psi_n$, on obtient une suite strictement croissante d'entiers
$$
n_0 < n_1 < n_2 < \ldots
$$
et des homomorphismes de $A$-algèbres $\psi^i : B \to C/a^{n_i}C$ tels que
$\psi^{i + 1}$ et $\psi^i$ coïncident modulo $a^{n_i - c}C$. Puisque $C$ est
complète pour la topologie $I$-adique, on peut prendre pour $\varphi$ la
limite des homomorphismes
$\psi^i \bmod a^{n_i - c}C : B \to C/a^{n_i - c}C$
et le lemme en résulte.
\end{proof}

\begin{lemma}
\label{restricted-lemma-get-morphism-principal}
Soit $I = (a)$ un idéal principal d'un anneau noethérien $A$. Soit $B$ un
objet de (\ref{restricted-equation-C-prime}). Supposons donné un entier $c \geq 0$ tel
que $\Ext^1_B(\NL_{B/A}^\wedge, N)$ soit annulé par $a^c$ pour tout
$B$-module $N$. Soit $C$ une algèbre complète pour la topologie $I$-adique
sur $A$. Supposons donné un entier $d \geq 0$ tel que
$C[a^\infty] \cap a^dC = 0$. Soit $n > \max(2c, c + d)$. Pour tout
homomorphisme de $A$-algèbres $\psi_n : B \to C/a^nC$, il existe un
homomorphisme de $A$-algèbres $\varphi : B \to C$ tel que
$\psi_n \bmod a^{n - c} = \varphi \bmod a^{n - c}$.
\end{lemma}

\noindent
Si $C$ est noethérien, on a $C[a^\infty] = C[a^e]$ pour un certain
$e \geq 0$. D'après Artin--Rees (Algèbre, lemme
\ref{algebra-lemma-Artin-Rees}), il existe un entier $f$ tel que
$a^nC \cap C[a^\infty] \subset a^{n - f}C[a^\infty]$ pour tout $n \geq f$.
Alors $d = e + f$ est un entier comme dans le lemme. Cet argument s'applique
en particulier si $C$ est un objet de (\ref{restricted-equation-C-prime}), d'après le
lemme \ref{restricted-lemma-topologically-finite-type-Noetherian}.

\begin{proof}
Soit $C \to C'$ le quotient de $C$ par $C[a^\infty]$. La $A$-algèbre $C'$
est complète pour la topologie $I$-adique d'après Algèbre, lemme
\ref{algebra-lemma-quotient-complete}, et parce que
$\bigcap (C[a^\infty] + a^nC) = C[a^\infty]$ : pour $n \geq d$, la somme
$C[a^\infty] + a^nC$ est directe. Pour $m \geq d$, le
diagramme
$$
\xymatrix{
0 \ar[r] &
C[a^\infty] \ar[r] \ar[d] &
C \ar[r] \ar[d] & C' \ar[r] \ar[d] & 0 \\
0 \ar[r] &
C[a^\infty] \ar[r] &
C/a^m C \ar[r] & C'/a^m C' \ar[r] & 0
}
$$
a ses lignes exactes. Ainsi, $C$ est le produit fibré de $C'$ et de
$C/a^mC$ au-dessus de $C'/a^mC'$ pour tout $m \geq d$. D'après le lemme
\ref{restricted-lemma-get-morphism-nonzerodivisor}, on peut choisir un homomorphisme
$\varphi' : B \to C'$ tel que $\varphi'$ et $\psi_n$ coïncident comme
homomorphismes à valeurs dans $C'/a^{n - c}C'$. On obtient un homomorphisme
$(\varphi', \psi_n \bmod a^{n - c}C) : B \to
C' \times_{C'/a^{n - c}C'} C/a^{n - c}C$. Puisque $n - c \geq d$, cela
équivaut à un homomorphisme $\varphi : B \to C$. La démonstration est achevée.
\end{proof}










\section{Algébrisation des algèbres rig-lisses sur les anneaux G}
\label{restricted-section-over-G-ring}

\noindent
Si l'anneau de base $A$ est un anneau G noethérien, nous pouvons démontrer
\cite[III théorème 7]{Elkik} pour toute algèbre rig-lisse relativement à un
idéal quelconque $I \subset A$ (non nécessairement principal).

\begin{lemma}
\label{restricted-lemma-close-enough}
Soit $I$ un idéal d'un anneau noethérien $A$. Soit $r \geq 0$ et notons
$P = A[x_1, \ldots, x_r]$ le complété $I$-adique. Considérons une résolution
$$
P^{\oplus t} \xrightarrow{K} P^{\oplus m}
\xrightarrow{g_1, \ldots, g_m} P \to B \to 0
$$
d'un quotient de $P$. Supposons que $B$ soit rig-lisse sur $(A, I)$. Il existe
alors un entier $n$ tel que, pour tout complexe
$$
P^{\oplus t} \xrightarrow{K'} P^{\oplus m}
\xrightarrow{g'_1, \ldots, g'_m} P
$$
vérifiant $g_i - g'_i \in I^nP$ et
$K - K' \in I^n\text{Mat}(m \times t, P)$, il existe un isomorphisme
$B \to B'$ de $A$-algèbres, où $B' = P/(g'_1, \ldots, g'_m)$.
\end{lemma}

\begin{proof}
(A) D'après la définition \ref{restricted-definition-rig-smooth-homomorphism}, on peut
choisir $c \geq 0$ tel que $I^c$ annule
$\Ext^1_B(\NL_{B/A}^\wedge, N)$ pour tout $B$-module $N$.

\medskip\noindent
(B) D'après Compléments d'algèbre, lemmes
\ref{more-algebra-lemma-approximate-complex} et
\ref{more-algebra-lemma-approximate-complex-graded}, il existe une constante
$c_1 = c(g_1, \ldots, g_m, K)$ telle que, pour $n \geq c_1 + 1$, le complexe
$$
P^{\oplus t} \xrightarrow{K'} P^{\oplus m}
\xrightarrow{g'_1, \ldots, g'_m} P \to B' \to 0
$$
est exact et $\text{Gr}_I(B) \cong \text{Gr}_I(B')$.

\medskip\noindent
(C) Soient $d_0 = d(\text{Gr}_I(B))$ et $q_0 = q(\text{Gr}_I(B))$ les entiers
obtenus dans Cohomologie locale, section
\ref{local-cohomology-section-uniform}.

\medskip\noindent
Nous affirmons que
$n = \max(c_1 + 1, q_0 + (d_0 + 1)c, 2(d_0 + 1)c + 1)$ convient, où $c$ est
comme en (A), $c_1$ comme en (B), et $q_0, d_0$ comme en (C).

\medskip\noindent
Soient $g'_1, \ldots, g'_m$ et $K'$ comme dans le lemme. Puisque
$g_i = g'_i \in I^nP$, on obtient un homomorphisme canonique de $A$-algèbres
$$
\psi_n : B \longrightarrow B'/I^nB'
$$
qui induit un isomorphisme $B/I^nB \to B'/I^nB'$. Puisque
$\text{Gr}_I(B) \cong \text{Gr}_I(B')$, on a
$d_0 = d(\text{Gr}_I(B'))$ et $q_0 = q(\text{Gr}_I(B'))$ ; puisque
$n \geq \max(q_0 + (1 + d_0)c, 2(d_0 + 1)c + 1)$, on peut appliquer le lemme
\ref{restricted-lemma-get-morphism-general-better} pour trouver un homomorphisme de
$A$-algèbres
$$
\varphi : B \longrightarrow B'
$$
tel que
$\varphi \bmod I^{n - (d_0 + 1)c}B' = \psi_n \bmod I^{n - (d_0 + 1)c}B'$.
Puisque $n - (d_0 + 1)c > 0$, on voit que $\varphi$ est un homomorphisme de
$A$-algèbres qui, modulo $I$, induit l'isomorphisme $B/IB \to B'/IB'$ trouvé
ci-dessus. La fin de la démonstration établit que ces faits imposent à
$\varphi$ d'être un isomorphisme ; nous suggérons au lecteur d'en trouver sa
propre démonstration.

\medskip\noindent
En effet, il s'ensuit que $\varphi$ est surjectif, par exemple en appliquant
Algèbre, lemme \ref{algebra-lemma-completion-generalities}, assertion (1), et
en utilisant le fait que $B$ et $B'$ sont complets. Ainsi, $\varphi$ induit
une surjection $\text{Gr}_I(B) \to \text{Gr}_I(B')$, qui est nécessairement
un isomorphisme puisque sa source et son but sont des anneaux noethériens
isomorphes ; voir Algèbre, lemme
\ref{algebra-lemma-surjective-endo-noetherian-ring-is-iso} (on peut bien sûr
montrer que $\varphi$ induit l'isomorphisme obtenu ci-dessus, au prix d'un très
petit argument). Ainsi, $\varphi$ induit des homomorphismes injectifs
$I^eB/I^{e + 1}B \to I^eB'/I^{e + 1}B'$ pour tout $e \geq 0$. Cela implique
que $\varphi$ est injectif, car, pour tout $b \in B$, il existe
$e \geq 0$ tel que $b \in I^eB$, $b \not \in I^{e + 1}B$, d'après le
théorème d'intersection de Krull (Algèbre, lemme
\ref{algebra-lemma-intersect-powers-ideal-module-zero}). La démonstration est
achevée.
\end{proof}

\begin{lemma}
\label{restricted-lemma-algebraize-easy}
Soit $I$ un idéal d'un anneau noethérien $A$. Soit $C^h$ l'hensélisé d'une
$A$-algèbre de type fini $C$ relativement à l'idéal $IC$. Soit
$J \subset C^h$ un idéal. Il existe alors une $A$-algèbre de type fini $B$
telle que $B^\wedge \cong (C^h/J)^\wedge$.
\end{lemma}

\begin{proof}
D'après Compléments d'algèbre, lemme
\ref{more-algebra-lemma-henselization-Noetherian-pair}, l'anneau $C^h$ est
noethérien. Écrivons $J = (g_1, \ldots, g_m)$. L'anneau $C^h$ est la colimite
filtrante de $C$-algèbres étales $C'$ telles que
$C/IC \to C'/IC'$ soit un isomorphisme (voir la démonstration de Compléments
d'algèbre, lemme \ref{more-algebra-lemma-henselization}). Choisissons $C'$ tel
que $g_1, \ldots, g_m$ soient les images de
$g'_1, \ldots, g'_m \in C'$. En posant
$B = C'/(g'_1, \ldots, g'_m)$, on obtient une $A$-algèbre de type fini. Bien
entendu, $(C, IC)$ et $C', IC')$ ont les mêmes hensélisés et les mêmes
complétés. Il s'ensuit aisément que $B^\wedge = (C^h/J)^\wedge$.
\end{proof}

\begin{proposition}
\label{restricted-proposition-approximate}
Soit $I$ un idéal d'un anneau G noethérien $A$. Soit $B$ un objet de
(\ref{restricted-equation-C-prime}). Si $B$ est rig-lisse sur $(A, I)$, alors il existe
une $A$-algèbre de type fini $C$ et un isomorphisme
$B \cong C^\wedge$ de $A$-algèbres.
\end{proposition}

\begin{proof}
Choisissons une présentation $B = A[x_1, \ldots, x_r]^\wedge/J$. Posons
$P = A[x_1, \ldots, x_r]^\wedge$. Choisissons des générateurs
$g_1, \ldots, g_m \in J$, puis des générateurs $k_1, \ldots, k_t$ du module
des relations entre $g_1, \ldots, g_m$, c'est-à-dire tels que
$$
P^{\oplus t} \xrightarrow{k_1, \ldots, k_t}
P^{\oplus m} \xrightarrow{g_1, \ldots, g_m}
P \to B \to 0
$$
soit une résolution. Écrivons $k_i = (k_{i1}, \ldots, k_{im})$, de sorte que
l'on ait
\begin{equation}
\label{restricted-equation-relations-straight-up}
\sum\nolimits_j k_{ij}g_j = 0
\end{equation}
pour $i = 1, \ldots, t$. Notons $K = (k_{ij})$ la matrice $m \times t$ de
coefficients $k_{ij}$.

\medskip\noindent
Soit $A[x_1, \ldots, x_r]^h$ l'hensélisé du couple
$(A[x_1, \ldots, x_r], IA[x_1, \ldots, x_r])$ ; voir Compléments d'algèbre,
lemme \ref{more-algebra-lemma-henselization}. Nous pouvons, et allons,
considérer $A[x_1, \ldots, x_r]^h$ comme un sous-anneau de
$P = A[x_1, \ldots, x_r]^\wedge$ ; voir Compléments d'algèbre, lemme
\ref{more-algebra-lemma-henselization-Noetherian-pair}. Puisque $A$ est un
anneau G noethérien, $A[x_1, \ldots, x_r]$ l'est aussi ; voir Compléments
d'algèbre, proposition
\ref{more-algebra-proposition-finite-type-over-G-ring}. On dispose donc de
l'approximation pour l'homomorphisme
$A[x_1, \ldots, x_r]^h \to A[x_1, \ldots, x_r]^\wedge = P$
relativement à l'idéal engendré par $I$ ; voir Lissage des homomorphismes
d'anneaux, lemme \ref{smoothing-lemma-henselian-pair}. Choisissons un entier
$n$ suffisamment grand comme dans le lemme \ref{restricted-lemma-close-enough}. Par la
propriété d'approximation, on peut choisir
$g'_1, \ldots, g'_m \in A[x_1, \ldots, x_r]^h$
et une matrice
$K' = (k'_{ij}) \in \text{Mat}(m \times t, A[x_1, \ldots, x_r]^h)$
tels que $\sum\nolimits_j k'_{ij}g'_j = 0$ dans
$A[x_1, \ldots, x_r]^h$ et que $g_i - g'_i \in I^nP$ et
$K - K' \in I^n\text{Mat}(m \times t, P)$. Le choix de $n$ implique alors
qu'il existe un isomorphisme
$$
B \to P/(g'_1, \ldots, g'_m) =
\left(A[x_1, \ldots, x_r]^h/(g'_1, \ldots, g'_m)\right)^\wedge
$$
Le lemme \ref{restricted-lemma-algebraize-easy} achève la démonstration.
\end{proof}

\noindent
Le lemme suivant n'est pas vrai en général si $A$ n'est pas un anneau G, mais
seulement noethérien. En effet, si $(A, \mathfrak m)$ est local et si
$I = \mathfrak m$, le lemme équivaut à l'approximation d'Artin pour $A^h$
(comme dans Lissage des homomorphismes d'anneaux, théorème
\ref{smoothing-theorem-approximation-property}), laquelle n'est pas valable
pour tout anneau local noethérien.

\begin{lemma}
\label{restricted-lemma-fully-faithfulness}
Soit $A$ un anneau G noethérien. Soit $I \subset A$ un idéal. Soient $B, C$
des $A$-algèbres de type fini. Pour tout homomorphisme de $A$-algèbres
$\varphi : B^\wedge \to C^\wedge$ entre les complétés $I$-adiques et tout
$N \geq 1$, il existe :
\begin{enumerate}
\item un homomorphisme étale d'anneaux $C \to C'$ qui induit un isomorphisme
$C/IC \to C'/IC'$ ;
\item un homomorphisme de $A$-algèbres $\varphi : B \to C'$,
\end{enumerate}
tels que $\varphi$ et $\psi$ coïncident modulo $I^N$ à valeurs dans
$C^\wedge = (C')^\wedge$.
\end{lemma}

\begin{proof}
L'énoncé du lemme a un sens car $C \to C'$ est plat (Algèbre, lemme
\ref{algebra-lemma-etale}) et induit donc un isomorphisme
$C/I^nC \to C'/I^nC'$ pour tout $n$ (Compléments d'algèbre, lemme
\ref{more-algebra-lemma-neighbourhood-isomorphism}), donc un isomorphisme sur
les complétés. Soit $C^h$ l'hensélisé du couple $(C, IC)$ ; voir Compléments
d'algèbre, lemme \ref{more-algebra-lemma-henselization}. Alors $C^h$ est la
colimite filtrante des algèbres $C'$, et les homomorphismes
$C \to C' \to C^h$ induisent un isomorphisme sur les complétés (Compléments
d'algèbre, lemme
\ref{more-algebra-lemma-henselization-Noetherian-pair}). Il suffit donc de
montrer qu'il existe un homomorphisme de $A$-algèbres $B \to C^h$ congru à
$\psi$ modulo $I^N$. Écrivons
$B = A[x_1, \ldots, x_n]/(f_1, \ldots, f_m)$. L'homomorphisme d'anneaux
$\psi$ correspond à des éléments $\hat c_1, \ldots, \hat c_n \in C^\wedge$
tels que $f_j(\hat c_1, \ldots, \hat c_n) = 0$ pour
$j = 1, \ldots, m$. Comme $A$ est un anneau G noethérien, $C$ l'est aussi ;
voir Compléments d'algèbre, proposition
\ref{more-algebra-proposition-finite-type-over-G-ring}. Le lemme
\ref{smoothing-lemma-henselian-pair} de Lissage des homomorphismes d'anneaux
fournit donc des éléments $c_1, \ldots, c_n \in C^h$ tels que
$f_j(c_1, \ldots, c_n) = 0$ pour $j = 1, \ldots, m$ et tels que
$\hat c_i - c_i \in I^NC^h$. Cela détermine l'homomorphisme
$B \to C^h$ souhaité.
\end{proof}







\section{Algébrisation des algèbres rig-lisses}
\label{restricted-section-algebraization-rig-smooth}

\noindent
Il s'avère que, si l'algèbre rig-lisse possède une présentation particulière,
son algébrisation est immédiate. Voir aussi la remarque
\ref{restricted-remark-discussion} pour une discussion.

\begin{lemma}
\label{restricted-lemma-presentation-rig-smooth}
Soit $A$ un anneau. Soient $f_1, \ldots, f_m \in A[x_1, \ldots, x_n]$ et
posons $B = A[x_1, \ldots, x_n]/(f_1, \ldots, f_m)$. Supposons $m \leq n$
et posons $g = \det_{1 \leq i, j \leq m}(\partial f_j/\partial x_i)$. Alors :
\begin{enumerate}
\item $g$ annule $\Ext^1_B(\NL_{B/A}, N)$ pour tout $B$-module $N$ ;
\item si $n = m$, alors la multiplication par $g$ sur $\NL_{B/A}$ est $0$
dans $D(B)$.
\end{enumerate}
\end{lemma}

\begin{proof}
Soit $T$ la matrice $m \times m$ de coefficients
$\partial f_j/\partial x_i$ pour $1 \leq i, j \leq n$. Soit $K \in D(B)$
représenté par le complexe $T : B^{\oplus m} \to B^{\oplus m}$, dont les
termes sont placés en degrés $-1$ et $0$. D'après Compléments d'algèbre, lemme
\ref{more-algebra-lemma-silly}, l'homomorphisme $g : K \to K$ est nul dans
$D(B)$. Posons $J = (f_1, \ldots, f_m)$. Rappelons que $\NL_{B/A}$ est
homotopiquement équivalent à
$J/J^2 \to \bigoplus_{i = 1, \ldots, n} B\text{d}x_i$ ; voir Algèbre,
section \ref{algebra-section-netherlander}. Notons $L$ le complexe
$J/J^2 \to \bigoplus_{i = 1, \ldots, m} B\text{d}x_i$, de sorte que l'on ait
le morphisme quotient $\NL_{B/A} \to L$. On dispose aussi d'un morphisme
surjectif de complexes $K \to L$, obtenu en envoyant le $j$-ième vecteur de
base du terme $B^{\oplus m}$ en degré $-1$ sur la classe de $f_j$ dans
$J/J^2$. On a le diagramme
$$
\NL_{B/A} \to L \leftarrow K
$$
Le lemme \ref{more-algebra-lemma-two-term-surjection-map-zero} de Compléments
d'algèbre montre que la multiplication par $g$ sur $L$ est $0$ dans $D(B)$.
D'autre part, le triangle distingué
$B^{\oplus n - m}[0] \to \NL_{B/A} \to L$ montre que
$\Ext^1_B(L, N) \to \Ext^1_B(\NL_{B/A}, N)$ est surjectif pour tout
$B$-module $N$, et donc annulé par $g$. Cela démontre l'assertion (1). Si
$n = m$, alors $\NL_{B/A} = L$, et (2) est satisfaite.
\end{proof}

\begin{lemma}
\label{restricted-lemma-approximate-presentation-rig-smooth}
Soit $I$ un idéal d'un anneau noethérien $A$. Soit $B$ un objet de
(\ref{restricted-equation-C-prime}). Soit
$B = A[x_1, \ldots, x_r]^\wedge/J$ une présentation. Supposons qu'il existe
un élément $b \in B$, un entier $0 \leq m \leq r$ et des éléments
$f_1, \ldots, f_m \in J$ tels que :
\begin{enumerate}
\item $V(b) \subset V(IB)$ dans $\Spec(B)$ ;
\item l'image de
$\Delta = \det_{1 \leq i, j \leq m}(\partial f_j/\partial x_i)$
dans $B$ divise $b$ ;
\item $b J \subset (f_1, \ldots, f_m) + J^2$.
\end{enumerate}
Il existe alors une $A$-algèbre de type fini $C$ et un isomorphisme de
$A$-algèbres $B \cong C^\wedge$.
\end{lemma}

\begin{proof}
Les conditions impliquent que $B$ est rig-lisse sur $(A, I)$ ; voir le lemme
\ref{restricted-lemma-equivalent-with-artin-smooth}. Écrivons $b' \Delta = b$ dans $B$
pour un certain $b' \in B$. Écrivons $I = (a_1, \ldots, a_t)$. Puisque
$V(b) \subset V(IB)$, il existe un entier $c \geq 0$ tel que
$I^cB \subset bB$. Écrivons $bb_i = a_i^c$ dans $B$ pour certains $b_i \in B$.

\medskip\noindent
Choisissons un entier $n \gg 0$ (nous préciserons plus loin à quel point il
doit être grand). Choisissons des polynômes
$f'_1, \ldots, f'_m \in A[x_1, \ldots, x_r]$ tels que
$f_i - f'_i \in I^nA[x_1, \ldots, x_r]^\wedge$. Posons
$\Delta' = \det_{1 \leq i, j \leq m}(\partial f'_j/\partial x_i)$ et
considérons la $A$-algèbre de type fini
$$
C = A[x_1, \ldots, x_r, z_1, \ldots, z_t]/
(f'_1, \ldots, f'_m,
z_1\Delta' - a_1^c, \ldots, z_t\Delta' - a_t^c)
$$
Nous allons appliquer le lemme \ref{restricted-lemma-presentation-rig-smooth} à $C$.
Calculons
$$
\det\left(
\begin{matrix}
\text{matrice des dérivées partielles de} \\
f'_1, \ldots, f'_m, z_1\Delta' - a_1^c, \ldots, z_t\Delta' - a_t^c \\
\text{par rapport aux variables} \\
x_1, \ldots, x_m, z_1, \ldots, z_t
\end{matrix}
\right) =
(\Delta')^{t + 1}
$$
On voit donc que $\Ext^1_C(\NL_{C/A}, N)$ est annulé par
$(\Delta')^{t + 1}$ pour tout $C$-module $N$. Puisque $a_i^c$ est divisible
par $\Delta'$ dans $C$, on voit que $a_i^{(t + 1)c}$ annule aussi ces
$\Ext^1$. Ainsi, $I^{c_1}$ annule $\Ext^1_C(\NL_{C/A}, N)$ pour tout
$C$-module $N$, où $c_1 = 1 + t((t + 1)c - 1)$. La valeur exacte de $c_1$
importe peu pour la suite de l'argument ; ce qui compte, c'est qu'elle soit
indépendante de $n$.

\medskip\noindent
Puisque $\NL_{C^\wedge/A}^\wedge = \NL_{C/A} \otimes_C C^\wedge$ d'après le
lemme \ref{restricted-lemma-NL-is-completion}, on en conclut que la multiplication par
$I^{c_1}$ est également nulle sur
$\Ext^1_{C^\wedge}(\NL_{C^\wedge/A}^\wedge, N)$ pour tout $C^\wedge$-module
$N$ ; voir Compléments d'algèbre, lemmes
\ref{more-algebra-lemma-base-change-property-ext-1-annihilated} et
\ref{more-algebra-lemma-two-term-base-change}. En particulier, $C^\wedge$ est
rig-lisse sur $(A, I)$.

\medskip\noindent
Notons que nous avons un homomorphisme surjectif de $A$-algèbres
$$
\psi_n : C \longrightarrow B/I^nB
$$
qui envoie la classe de $x_i$ sur celle de $x_i$ et la classe de $z_i$ sur
celle de $b_ib'$. Cet homomorphisme est bien défini grâce aux choix de $b'$ et
de $b_i$ faits dans le premier paragraphe de la démonstration.

\medskip\noindent
Soient $d = d(\text{Gr}_I(B))$ et $q_0 = q(\text{Gr}_I(B))$ les entiers
obtenus dans Cohomologie locale, section
\ref{local-cohomology-section-uniform}. D'après le lemme
\ref{restricted-lemma-get-morphism-general-better}, si l'on prend
$n \geq \max(q_0 + (d + 1)c_1, 2(d + 1)c_1 + 1)$, on peut trouver un
homomorphisme $\varphi : C^\wedge \to B$ de $A$-algèbres, congru à $\psi_n$
modulo $I^{n - (d + 1)c_1}B$.

\medskip\noindent
Puisque $\varphi : C^\wedge \to B$ est surjectif modulo $I$, il est surjectif
(on peut, par exemple, utiliser Algèbre, lemme
\ref{algebra-lemma-completion-generalities}). Pour achever la démonstration, il
suffit de montrer que $\Ker(\varphi)/\Ker(\varphi)^2$ est annulé par une
puissance de $I$ ; voir Compléments d'algèbre, lemme
\ref{more-algebra-lemma-quotient-by-idempotent}.

\medskip\noindent
Puisque $\varphi$ est surjectif, on voit que les modules de cohomologie de
$\NL_{B/C^{\wedge}}^\wedge$ sont
$H^0(\NL_{B/C^{\wedge}}^\wedge) = 0$ et
$H^{-1}(\NL_{B/C^{\wedge}}^\wedge) = \Ker(\varphi)/\Ker(\varphi)^2$.
On a une suite exacte
$$
H^{-1}(\NL_{C^\wedge/A}^\wedge \otimes_{C^\wedge} B) \to
H^{-1}(\NL_{B/A}^\wedge) \to
H^{-1}(\NL_{B/C^{\wedge}}^\wedge) \to
H^0(\NL_{C^\wedge/A}^\wedge \otimes_{C^\wedge} B) \to
H^0(\NL_{B/A}^\wedge) \to 0
$$
d'après le lemme \ref{restricted-lemma-exact-sequence-NL}. Les deux premiers modules sont
annulés par une puissance de $I$, puisque $B$ et $C^\wedge$ sont rig-lisses
sur $(A, I)$. Il suffit donc de montrer que le noyau de l'homomorphisme
surjectif
$H^0(\NL_{C^\wedge/A}^\wedge \otimes_{C^\wedge} B) \to
H^0(\NL_{B/A}^\wedge)$ est annulé par une puissance de $I$. Pour cela, il
suffit de montrer qu'il est annulé par une puissance de $b$. Autrement dit,
il suffit de montrer que
$$
H^0(\NL_{C^\wedge/A}^\wedge) \otimes_{C^\wedge} B[1/b]
\longrightarrow
H^0(\NL_{B/A}^\wedge) \otimes_B B[1/b]
$$
est un isomorphisme. Or les deux membres sont des $B[1/b]$-modules libres de
rang $r - m$, de base $\text{d}x_{m + 1}, \ldots, \text{d}x_r$, ce qui achève
la démonstration.
\end{proof}

\begin{remark}
\label{restricted-remark-discussion}
Soit $I$ un idéal d'un anneau noethérien $A$. Soit $B$ un objet de
(\ref{restricted-equation-C-prime}) qui est rig-lisse sur $(A, I)$. Il est démontré dans
\cite[théorème 1.2]{gabber-zavyalov} que $B$ est isomorphe au complété $I$-adique
d'une $A$-algèbre de type fini. Ce résultat remplace les résultats partiels
suivants :
\begin{enumerate}
\item si $A$ est un anneau G, le résultat découle de la proposition
\ref{restricted-proposition-approximate} ;
\item si $B$ est rig-étale sur $(A, I)$, le résultat découle du lemme
\ref{restricted-lemma-approximate} ;
\item si $I$ est principal, le résultat découle de
\cite[III théorème 7]{Elkik}.
\end{enumerate}
\end{remark}







\section{Algèbres rig-étales}
\label{restricted-section-rig-etale}

\noindent
Au vu de notre définition des algèbres rig-lisses (définition
\ref{restricted-definition-rig-smooth-homomorphism}), la définition suivante ne devrait
pas surprendre.

\begin{definition}
\label{restricted-definition-rig-etale-homomorphism}
Soient $A$ un anneau noethérien et $I \subset A$ un idéal. Soit $B$ un objet
de (\ref{restricted-equation-C-prime}). Nous disons que $B$ est {\it rig-étale sur
$(A, I)$} s'il existe un entier $c \geq 0$ tel que, pour tout $a \in I^c$, la
multiplication par $a$ sur $\NL_{B/A}^\wedge$ soit nulle dans $D(B)$.
\end{definition}

\noindent
La condition (\ref{restricted-item-condition-artin}) du lemme suivant est l'une de celles
qu'emploie \cite{ArtinII} pour définir les modifications formelles. Nous
l'avons ajoutée à la liste afin de faciliter sa comparaison ultérieure avec
nos conditions.

\begin{lemma}
\label{restricted-lemma-equivalent-with-artin}
Soient $A$ un anneau noethérien et $I \subset A$ un idéal. Soit $B$ un objet
de (\ref{restricted-equation-C-prime}). Écrivons
$B = A[x_1, \ldots, x_r]^\wedge/J$
(lemme \ref{restricted-lemma-topologically-finite-type-Noetherian}) et soit
$\NL_{B/A}^\wedge = (J/J^2 \to \bigoplus B\text{d}x_i)$ son complexe
cotangent naïf (\ref{restricted-equation-NL}). Les conditions suivantes sont équivalentes :
\begin{enumerate}
\item $B$ est rig-étale sur $(A, I)$ ;
\item
\label{restricted-item-zero-on-NL}
il existe $c \geq 0$ tel que, pour tout $a \in I^c$, la multiplication par
$a$ sur $\NL_{B/A}^\wedge$ soit nulle dans $D(B)$ ;
\item
\label{restricted-item-zero-on-cohomology-NL}
il existe $c \geq 0$ tel que $H^i(\NL_{B/A}^\wedge)$, $i = -1, 0$, soit
annulé par $I^c$ ;
\item
\label{restricted-item-zero-on-cohomology-NL-truncations}
il existe $c \geq 0$ tel que $H^i(\NL_{B_n/A_n})$, $i = -1, 0$, soit annulé
par $I^c$ pour tout $n \geq 1$, où $A_n = A/I^n$ et $B_n = B/I^nB$ ;
\item
\label{restricted-item-condition-artin-pre-pre}
pour tout $a \in I$, il existe $c \geq 0$ tel que :
\begin{enumerate}
\item $a^c$ annule $H^0(\NL_{B/A}^\wedge)$ ;
\item il existe $f_1, \ldots, f_r \in J$ tels que
$a^c J \subset (f_1, \ldots, f_r) + J^2$ ;
\end{enumerate}
\item
\label{restricted-item-condition-artin-pre}
pour tout $a \in I$, il existe $f_1, \ldots, f_r \in J$ et $c \geq 0$ tels
que :
\begin{enumerate}
\item $\det_{1 \leq i, j \leq r}(\partial f_j/\partial x_i)$ divise $a^c$
dans $B$ ;
\item $a^c J \subset (f_1, \ldots, f_r) + J^2$ ;
\end{enumerate}
\item
\label{restricted-item-condition-artin}
si l'on choisit des générateurs $f_1, \ldots, f_t$ de $J$, on a :
\begin{enumerate}
\item l'idéal jacobien de $B$ sur $A$, c'est-à-dire l'idéal de $B$ engendré
par les mineurs $r \times r$ de la matrice
$(\partial f_j/\partial x_i)_{1 \leq i \leq r, 1 \leq j \leq t}$,
contient l'idéal $I^cB$ pour un certain $c$ ;
\item l'idéal de Cramer de $B$ sur $A$, c'est-à-dire l'idéal de $B$ engendré
par l'image dans $B$ du $r$-ième idéal de Fitting de $J$, considéré comme
$A[x_1, \ldots, x_r]^\wedge$-module, contient $I^cB$ pour un certain $c$.
\end{enumerate}
\end{enumerate}
\end{lemma}

\begin{proof}
L'équivalence de (1) et de (\ref{restricted-item-zero-on-NL}) reformule la définition
\ref{restricted-definition-rig-etale-homomorphism}.

\medskip\noindent
L'équivalence de (\ref{restricted-item-zero-on-NL}) et de
(\ref{restricted-item-zero-on-cohomology-NL}) résulte de Compléments d'algèbre, lemme
\ref{more-algebra-lemma-zero-in-derived}.

\medskip\noindent
L'équivalence de (\ref{restricted-item-zero-on-cohomology-NL}) et de
(\ref{restricted-item-zero-on-cohomology-NL-truncations}) résulte du fait que les systèmes
$\{\NL_{B_n/A_n}\}$ et $\NL_{B/A}^\wedge \otimes_B B_n$ sont strictement
isomorphes ; voir le lemme \ref{restricted-lemma-NL-is-limit}. Certains détails sont omis.

\medskip\noindent
Supposons (\ref{restricted-item-zero-on-NL}). Soit $a \in I$. Soit $c$ tel que la
multiplication par $a^c$ soit nulle sur $\NL_{B/A}^\wedge$. D'après
Compléments d'algèbre, lemme
\ref{more-algebra-lemma-map-out-of-almost-free}, assertion (1), il existe un
homomorphisme $\alpha : \bigoplus B\text{d}x_i \to J/J^2$ tel que
$\text{d} \circ \alpha$ et $\alpha \circ \text{d}$ soient tous deux la
multiplication par $a^c$. Soit $f_i \in J$ un élément dont la classe modulo
$J^2$ est égale à $\alpha(\text{d}x_i)$. Un calcul élémentaire montre que
(\ref{restricted-item-condition-artin-pre})(a), (b) sont satisfaites.

\medskip\noindent
Nous omettons de vérifier que (\ref{restricted-item-condition-artin-pre}) implique
(\ref{restricted-item-condition-artin-pre-pre}) ; il s'agit simplement d'un énoncé sur
les complexes à deux termes sur $B$ de la forme $M \to B^{\oplus r}$.

\medskip\noindent
Supposons (\ref{restricted-item-condition-artin-pre-pre}) satisfaite. Écrivons
$I = (a_1, \ldots, a_t)$. Soit $c_i \geq 0$ l'entier tel que
(\ref{restricted-item-condition-artin-pre-pre})(a), (b) soient satisfaites pour
$a_i^{c_i}$. Alors $I^{\sum c_i}$ annule $H^0(\NL_{B/A}^\wedge)$. Soient
$f_{i, 1}, \ldots, f_{i, r} \in J$ comme dans
(\ref{restricted-item-condition-artin-pre-pre})(b) pour $a_i$. Considérons le composé
$$
B^{\oplus r} \to J/J^2 \to \bigoplus B\text{d}x_i
$$
où le $j$-ième vecteur de base est envoyé sur la classe de $f_{i, j}$ dans
$J/J^2$. D'après (\ref{restricted-item-condition-artin-pre-pre})(a) et (b), le conoyau du
composé est annulé par $a_i^{2c_i}$. Cet homomorphisme est donc surjectif après
inversion de $a_i^{c_i}$, et par conséquent un isomorphisme (Algèbre, lemme
\ref{algebra-lemma-fun}). Ainsi, le noyau de
$B^{\oplus r} \to \bigoplus B\text{d}x_i$ est de torsion annulée par une
puissance de $a_i$, et donc
$H^{-1}(\NL_{B/A}^\wedge) = \Ker(J/J^2 \to \bigoplus B\text{d}x_i)$
est de torsion annulée par une puissance de $a_i$. Comme $B$ est noethérien
(lemme \ref{restricted-lemma-topologically-finite-type-Noetherian}), tous les modules, y
compris $H^{-1}(\NL_{B/A}^\wedge)$, sont finis. Ainsi, $a_i^{d_i}$ annule
$H^{-1}(\NL_{B/A}^\wedge)$ pour un certain $d_i \geq 0$. Il s'ensuit que
$I^{\sum d_i}$ annule $H^{-1}(\NL_{B/A}^\wedge)$, et la condition
(\ref{restricted-item-zero-on-cohomology-NL}) est satisfaite.

\medskip\noindent
Ainsi, les conditions
(\ref{restricted-item-zero-on-NL}),
(\ref{restricted-item-zero-on-cohomology-NL}),
(\ref{restricted-item-zero-on-cohomology-NL-truncations}),
(\ref{restricted-item-condition-artin-pre-pre}) et
(\ref{restricted-item-condition-artin-pre}) sont équivalentes. Il reste à montrer que ces
conditions équivalent à (\ref{restricted-item-condition-artin}). Notons que l'idéal de
Cramer $\text{Fit}_r(J) B$ est égal à $\text{Fit}_r(J/J^2)$, puisque
$J/J^2 = J \otimes_{A[x_1, \ldots, x_r]^\wedge} B$ ; voir Compléments
d'algèbre, lemme \ref{more-algebra-lemma-fitting-ideal-basics}, assertion (3).
Notons aussi que l'idéal jacobien n'est autre que
$\text{Fit}_0(H^0(\NL_{B/A}^\wedge))$. L'équivalence de
(\ref{restricted-item-zero-on-cohomology-NL}) et de (\ref{restricted-item-condition-artin}) est donc
une question purement algébrique, que nous traitons au paragraphe suivant.

\medskip\noindent
Soit $R$ un anneau noethérien et soit $I \subset R$ un idéal. Soit
$M \xrightarrow{d} R^{\oplus r}$ un complexe à deux termes. Il faut montrer
que les conditions suivantes sont équivalentes :
\begin{enumerate}
\item[(A)] la cohomologie de $M \to R^{\oplus r}$ est annulée par une
puissance de $I$ ;
\item[(B)] les idéaux $\text{Fit}_r(M)$ et $\text{Fit}_0(\text{Coker}(d))$
contiennent une puissance de $I$.
\end{enumerate}
Puisque $R$ est noethérien, on peut reformuler l'assertion (2) comme une
inclusion des sous-schémas fermés correspondants ; voir Algèbre, lemmes
\ref{algebra-lemma-Zariski-topology} et \ref{algebra-lemma-Noetherian-power}.
D'autre part, sur le complémentaire de $V(\text{Fit}_0(\Coker(d)))$, le
conoyau de $d$ est nul, et sur le complémentaire de $V(\text{Fit}_r(M))$, le
module $M$ est localement engendré par $r$ éléments ; voir Compléments
d'algèbre, lemme \ref{more-algebra-lemma-fitting-ideal-generate-locally}.
Ainsi, (B) équivaut à :
\begin{enumerate}
\item[(C)] hors de $V(I)$, le conoyau de $d$ est nul et le module $M$ est
localement engendré par $\leq r$ éléments.
\end{enumerate}
Cela équivaut évidemment à l'annulation de la cohomologie de
$M \to R^{\oplus r}$ sur $\Spec(R) \setminus V(I)$, laquelle équivaut à son
tour à (A). La démonstration est achevée.
\end{proof}

\begin{lemma}
\label{restricted-lemma-rig-etale-rig-smooth}
Soient $A$ un anneau noethérien et $I$ un idéal. Soit $B$ un objet de
(\ref{restricted-equation-C-prime}). Si $B$ est rig-étale sur $(A, I)$, alors $B$ est
rig-lisse sur $(A, I)$.
\end{lemma}

\begin{proof}
Cela résulte immédiatement des définitions
\ref{restricted-definition-rig-smooth-homomorphism} et
\ref{restricted-definition-rig-etale-homomorphism}.
\end{proof}

\begin{lemma}
\label{restricted-lemma-rig-etale}
Soient $A$ un anneau noethérien et $I$ un idéal. Soit $B$ une $A$-algèbre de
type fini.
\begin{enumerate}
\item Si $\Spec(B) \to \Spec(A)$ est étale au-dessus de
$\Spec(A) \setminus V(I)$, alors $B^\wedge$ satisfait les conditions
équivalentes du lemme \ref{restricted-lemma-equivalent-with-artin}.
\item Si $B^\wedge$ satisfait les conditions équivalentes du lemme
\ref{restricted-lemma-equivalent-with-artin}, alors il existe $g \in 1 + IB$ tel que
$\Spec(B_g)$ soit étale au-dessus de $\Spec(A) \setminus V(I)$.
\end{enumerate}
\end{lemma}

\begin{proof}
Supposons que $B^\wedge$ satisfasse les conditions équivalentes du lemme
\ref{restricted-lemma-equivalent-with-artin}. Le complexe cotangent naïf $\NL_{B/A}$ est
un complexe de $B$-modules de type fini ; par conséquent, $H^{-1}$ et $H^0$
sont des $B$-modules finis. Le foncteur de complétion est exact sur les
$B$-modules finis (Algèbre, lemme \ref{algebra-lemma-completion-flat}) et
$\NL_{B^\wedge/A}^\wedge$ est le complété du complexe $\NL_{B/A}$ (cela se
voit aisément en choisissant des présentations). L'hypothèse implique donc
qu'il existe $c \geq 0$ tel que $H^{-1}/I^nH^{-1}$ et $H^0/I^nH^0$ soient
annulés par $I^c$ pour tout $n$. D'après le lemme de Nakayama (Algèbre, lemme
\ref{algebra-lemma-NAK}), cela signifie que $I^cH^{-1}$ et $I^cH^0$ sont
annulés par un élément de la forme $g = 1 + x$ avec $x \in IB$. Après
inversion de $g$ (ce qui ne modifie pas les quotients $B/I^nB$), on voit que
la cohomologie de $\NL_{B/A}$ est annulée par $I^c$. Ainsi, $A \to B$ est
étale en tout idéal premier de $B$ qui n'est pas au-dessus de $V(I)$, par la
définition des homomorphismes d'anneaux étales ; voir Algèbre, définition
\ref{algebra-definition-etale}.

\medskip\noindent
Réciproquement, supposons que $\Spec(B) \to \Spec(A)$ soit étale au-dessus de
$\Spec(A) \setminus V(I)$. Alors, pour tout $a \in I$, il existe $c \geq 0$
tel que la multiplication par $a^c$ soit nulle sur $\NL_{B/A}$. Puisque
$\NL_{B^\wedge/A}^\wedge$ est le complété dérivé de $\NL_{B/A}$ (voir le lemme
\ref{restricted-lemma-NL-is-limit}), il s'ensuit que $B^\wedge$ satisfait les conditions
équivalentes du lemme \ref{restricted-lemma-equivalent-with-artin}.
\end{proof}

\begin{lemma}
\label{restricted-lemma-zero-after-modding-out}
Soit $(A_1, I_1) \to (A_2, I_2)$ comme dans la remarque
\ref{restricted-remark-base-change}, avec $A_1$ et $A_2$ noethériens. Soit $B_1$ un objet
de (\ref{restricted-equation-C-prime}) pour $(A_1, I_1)$. Soit $B_2$ le changement de
base de $B_1$. Si la multiplication par $f_1 \in B_1$ sur
$\NL^\wedge_{B_1/A_1}$ est nulle dans $D(B_1)$, alors la multiplication par
l'image $f_2 \in B_2$ sur $\NL^\wedge_{B_2/A_2}$ est nulle dans $D(B_2)$.
\end{lemma}

\begin{proof}
D'après le lemme \ref{restricted-lemma-NL-base-change}, il existe un morphisme
$$
\NL_{B_1/A_1} \otimes_{B_2} B_1 \to \NL_{B_2/A_2}
$$
qui induit un isomorphisme sur $H^0$ et une surjection sur $H^{-1}$. Le
résultat découle donc de Compléments d'algèbre, lemme
\ref{more-algebra-lemma-two-term-surjection-map-zero}.
\end{proof}

\begin{lemma}
\label{restricted-lemma-base-change-rig-etale-homomorphism}
Soit $A_1 \to A_2$ un homomorphisme d'anneaux noethériens. Soit
$I_i \subset A_i$ un idéal tel que $V(I_1A_2) = V(I_2)$. Soit $B_1$ un objet
de (\ref{restricted-equation-C-prime}) pour $(A_1, I_1)$. Soit $B_2$ le changement de
base de $B_1$, comme dans la remarque \ref{restricted-remark-base-change}. Si $B_1$ est
rig-étale sur $(A_1, I_1)$, alors $B_2$ est rig-étale sur $(A_2, I_2)$.
\end{lemma}

\begin{proof}
Cela résulte du lemme \ref{restricted-lemma-zero-after-modding-out}, de la définition
\ref{restricted-definition-rig-etale-homomorphism} et du fait que
$I_2^c \subset I_1A_2$ pour un certain $c \geq 0$, puisque $A_2$ est
noethérien.
\end{proof}

\begin{lemma}
\label{restricted-lemma-fully-faithful-etale-over-complement}
Soit $A$ un anneau noethérien. Soit $I \subset A$ un idéal. Soit $B$ une
$A$-algèbre de type fini telle que $\Spec(B) \to \Spec(A)$ soit étale
au-dessus de $\Spec(A) \setminus V(I)$. Soit $C$ une $A$-algèbre
noethérienne. Tout homomorphisme de $A$-algèbres
$B^\wedge \to C^\wedge$ entre les complétés $I$-adiques provient alors d'un
unique homomorphisme de $A$-algèbres
$$
B \longrightarrow C^h
$$
où $C^h$ est l'hensélisé du couple $(C, IC)$, comme dans Compléments
d'algèbre, lemme \ref{more-algebra-lemma-henselization}. De plus, tout
homomorphisme de $A$-algèbres $B \to C^h$ se factorise par une $C$-algèbre
étale $C'$ telle que $C/IC \to C'/IC'$ soit un isomorphisme.
\end{lemma}

\begin{proof}
L'unicité résulte du fait que $C^h$ est un sous-anneau de $C^\wedge$ ; voir,
par exemple, Compléments d'algèbre, lemme
\ref{more-algebra-lemma-henselization-Noetherian-pair}. La dernière assertion
résulte du fait que $C^h$ est la colimite filtrante de ces $C$-algèbres $C'$ ;
voir la démonstration de Compléments d'algèbre, lemme
\ref{more-algebra-lemma-henselization}. Cela étant établi, démontrons
l'existence.

\medskip\noindent
Soit $\varphi : B^\wedge \to C^\wedge$ l'homomorphisme donné. Il définit une
section
$$
\sigma : (B \otimes_A C)^\wedge \longrightarrow C^\wedge
$$
du complété de l'homomorphisme $C \to B \otimes_A C$. On peut remplacer
$(A, I, B, C, \varphi)$ par $(C, IC, B \otimes_A C, C, \sigma)$. On voit
ainsi que l'on peut supposer $A = C$.

\medskip\noindent
Existence dans le cas $A = C$. Dans ce cas, l'homomorphisme
$\varphi : B^\wedge \to A^\wedge$ est nécessairement surjectif. D'après les
lemmes \ref{restricted-lemma-rig-etale} et \ref{restricted-lemma-exact-sequence-NL}, les groupes de
cohomologie de
$\NL_{A^\wedge/\!_\varphi B^\wedge}^\wedge$
sont annulés par une puissance de $I$. Puisque $\varphi$ est surjectif, cela
implique que $\Ker(\varphi)/\Ker(\varphi)^2$ est annulé par une puissance de
$I$. Ainsi, $\varphi : B^\wedge \to A^\wedge$ est le complété d'une
$B$-algèbre de type fini $B \to D$ ; voir Compléments d'algèbre, lemme
\ref{more-algebra-lemma-quotient-by-idempotent}. Par conséquent, $A \to D$ est
un homomorphisme d'algèbres de type fini qui induit un isomorphisme
$A^\wedge \to D^\wedge$. D'après le lemme \ref{restricted-lemma-rig-etale}, on peut
remplacer $D$ par un localisé et supposer que $A \to D$ est étale hors de
$V(I)$. Puisque $A^\wedge \to D^\wedge$ est un isomorphisme, on voit que
$\Spec(D) \to \Spec(A)$ est aussi étale dans un voisinage de $V(ID)$ (par
exemple d'après Compléments sur les morphismes, lemme
\ref{more-morphisms-lemma-check-smoothness-on-infinitesimal-nbhds}). Ainsi,
$\Spec(D) \to \Spec(A)$ est étale. Par conséquent, $D$ s'envoie dans $A^h$,
et le lemme est démontré.
\end{proof}
















\section{Un argument de somme amalgamée}
\label{restricted-section-approximation}

\noindent
Le seul but de cette section est de démontrer le lemme suivant, qui jouera un
rôle essentiel dans l'algébrisation des algèbres rig-étales. Nous utiliserons
un peu de théorie des espaces algébriques pour le démontrer ; une version
antérieure de ce chapitre donnait une démonstration (beaucoup plus longue)
fondée sur l'algèbre et un peu de théorie des déformations, que le lecteur
intéressé trouvera dans l'historique du projet Champs.

\begin{lemma}
\label{restricted-lemma-lift-approximation}
Soient $A$ un anneau noethérien et $I \subset A$ un idéal. Soit
$J \subset A$ un idéal nilpotent. Considérons un diagramme commutatif
$$
\xymatrix{
C \ar[r] & C_0 \ar@{=}[r] & C/JC \\
& B_0 \ar[u] \\
A \ar[r] \ar[uu] & A_0 \ar[u] \ar@{=}[r] & A/J
}
$$
dont les flèches verticales sont de type fini et tel que :
\begin{enumerate}
\item $\Spec(C) \to \Spec(A)$ est étale au-dessus de
$\Spec(A) \setminus V(I)$ ;
\item $\Spec(B_0) \to \Spec(A_0)$ est étale au-dessus de
$\Spec(A_0) \setminus V(IA_0)$ ;
\item $B_0 \to C_0$ est étale et induit un isomorphisme
$B_0/IB_0 = C_0/IC_0$.
\end{enumerate}
On peut alors compléter le diagramme précédent en un diagramme commutatif
$$
\xymatrix{
C \ar[r] & C/JC \\
B \ar[u] \ar[r] & B_0 \ar[u] \\
A \ar[r] \ar[u] & A/J \ar[u]
}
$$
où $A \to B$ est de type fini, $B/JB = B_0$, $B \to C$ est étale et
$\Spec(B) \to \Spec(A)$ est étale au-dessus de
$\Spec(A) \setminus V(I)$.
\end{lemma}

\begin{proof}
Posons $X = \Spec(A)$, $X_0 = \Spec(A_0)$, $Y_0 = \Spec(B_0)$,
$Z = \Spec(C)$ et $Z_0 = \Spec(C_0)$. De plus, notons $U \subset X$,
$U_0 \subset X_0$, $V_0 \subset Y_0$, $W \subset Z$ et $W_0 \subset Z_0$
les complémentaires des lieux d'annulation de $I$. Le diagramme suivant aide
à visualiser la situation :
$$
\xymatrix{
Z \ar[dd] & Z_0 \ar[l] \ar[d] \\
& Y_0 \ar[d] \\
X & X_0 \ar[l]
}
\quad\quad\quad
\xymatrix{
W \ar[dd] & W_0 \ar[l] \ar[d] \\
& V_0 \ar[d] \\
U & U_0 \ar[l]
}
$$
Les conditions du lemme garantissent que
$$
\xymatrix{
W_0 \ar[r] \ar[d] & Z_0 \ar[d] \\
V_0 \ar[r] & Y_0
}
$$
est un carré distingué élémentaire ; voir Catégories dérivées d'espaces,
définition \ref{spaces-perfect-definition-elementary-distinguished-square}.
De plus, nous savons que $W_0 \to U_0$ et $V_0 \to U_0$ sont étales. Le
morphisme $X_0 \subset X$ est un épaississement d'ordre fini puisque $J$ est
supposé nilpotent. Par invariance topologique du site étale, on peut trouver un
unique morphisme étale de schémas $V \to X$ tel que
$V_0 = V \times_X X_0$, et relever le morphisme donné $W_0 \to V_0$ en un
unique morphisme $W \to V$ sur $X$. Voir Morphismes étales, théorème
\ref{etale-theorem-remarkable-equivalence}. Puisque $W_0 \to V_0$ est séparé,
le morphisme $W \to V$ est lui aussi séparé ; voir, par exemple, Compléments
sur les morphismes, lemme \ref{more-morphisms-lemma-deform-property}. D'après
Sommes amalgamées d'espaces, lemme
\ref{spaces-pushouts-lemma-construct-elementary-distinguished-square}, on peut
construire un carré distingué élémentaire
$$
\xymatrix{
W \ar[r] \ar[d] & Z \ar[d] \\
V \ar[r] & Y
}
$$
dans la catégorie des espaces algébriques sur $X$. Comme le changement de base
d'un carré distingué élémentaire est encore un carré distingué élémentaire
(Catégories dérivées d'espaces, lemme
\ref{spaces-perfect-lemma-make-more-elementary-distinguished-squares}), on voit
que
$$
\xymatrix{
W_0 \ar[r] \ar[d] & Z_0 \ar[d] \\
V_0 \ar[r] & Y \times_X X_0
}
$$
est un carré distingué élémentaire. Il existe donc un unique isomorphisme
$Y \times_X X_0 = Y_0$ compatible avec les deux carrés faisant intervenir ces
espaces, car les carrés distingués élémentaires sont des sommes amalgamées
(Sommes amalgamées d'espaces, lemme
\ref{spaces-pushouts-lemma-elementary-distinguished-square-pushout}). Il
s'ensuit que $Y$ est affine d'après Limites d'espaces, proposition
\ref{spaces-limits-proposition-affine}. Écrivons $Y = \Spec(B)$. Il est clair
que $B$ s'insère dans le diagramme voulu et satisfait toutes les propriétés
requises.
\end{proof}

















\section{Algébrisation des algèbres rig-étales}
\label{restricted-section-approximation-principal}

\noindent
Le but principal est de démontrer l'algébrisation des algèbres rig-étales sans
supposer que l'anneau noethérien de base $A$ soit un anneau G, et pour un idéal
arbitraire $I \subset A$ --- non nécessairement principal. Nous démontrons
d'abord le cas d'un idéal principal, puis utilisons le résultat de la section
\ref{restricted-section-approximation} pour achever la démonstration.

\begin{lemma}
\label{restricted-lemma-approximate-principal}
\begin{reference}
Cas rig-étale de \cite[III Theorem 7]{Elkik}
\end{reference}
Soit $A$ un anneau noethérien et soit $I = (a)$ un idéal principal. Soit $B$
un objet de (\ref{restricted-equation-C-prime}) rig-étale sur $(A, I)$. Il existe alors
une $A$-algèbre de type fini $C$ et un isomorphisme $B \cong C^\wedge$.
\end{lemma}

\begin{proof}
Choisissons une présentation $B = A[x_1, \ldots, x_r]^\wedge/J$. D'après le
lemme \ref{restricted-lemma-equivalent-with-artin}, assertion (6), on peut trouver
$c \geq 0$ et $f_1, \ldots, f_r \in J$ tels que
$\det_{1 \leq i, j \leq r}(\partial f_j/\partial x_i)$ divise $a^c$ dans $B$
et que $a^c J \subset (f_1, \ldots, f_r) + J^2$. Le lemme
\ref{restricted-lemma-approximate-presentation-rig-smooth} s'applique donc. Ceci achève la
démonstration, mais signalons que, dans ce cas, l'emploi du lemme
\ref{restricted-lemma-get-morphism-general-better} peut être remplacé par celui, beaucoup
plus simple, du lemme \ref{restricted-lemma-get-morphism-principal}.
\end{proof}

\begin{lemma}
\label{restricted-lemma-approximate}
Soit $A$ un anneau noethérien. Soit $I \subset A$ un idéal. Soit $B$ un objet
de (\ref{restricted-equation-C-prime}) rig-étale sur $(A, I)$. Il existe alors une
$A$-algèbre de type fini $C$ et un isomorphisme $B \cong C^\wedge$.
\end{lemma}

\begin{proof}
Nous démontrons ce lemme par récurrence sur le nombre de générateurs de $I$.
Écrivons $I = (a_1, \ldots, a_t)$. Si $t = 0$, alors $I = 0$ et il n'y a rien
à démontrer. Si $t = 1$, le lemme résulte du lemme
\ref{restricted-lemma-approximate-principal}. Supposons $t > 1$.

\medskip\noindent
Pour tout $m \geq 1$, posons $\bar A_m = A/(a_t^m)$. Considérons l'idéal
$\bar I_m = (\bar a_1, \ldots, \bar a_{t - 1})$ de $\bar A_m$. Notons que
$V(I \bar A_m) = V(\bar I_m)$. Soit $B_m = B/(a_t^m)$ le changement de base
de $B$ par l'homomorphisme $(A, I) \to (\bar A_m, \bar I_m)$ ; voir la
remarque \ref{restricted-remark-take-bar}. D'après le lemme
\ref{restricted-lemma-base-change-rig-etale-homomorphism}, $B_m$ est rig-étale sur
$(\bar A_m, \bar I_m)$.

\medskip\noindent
Par l'hypothèse de récurrence (sur $t$), on peut trouver une
$\bar A_m$-algèbre de type fini $C_m$ et un homomorphisme $C_m \to B_m$ qui
induit un isomorphisme $C_m^\wedge \cong B_m$, la complétion étant prise
relativement à $\bar I_m$. D'après le lemme \ref{restricted-lemma-rig-etale}, on peut
supposer que $\Spec(C_m) \to \Spec(\bar A_m)$ est étale au-dessus de
$\Spec(\bar A_m) \setminus V(\bar I_m)$.

\medskip\noindent
Nous affirmons que l'on peut choisir $A_m \to C_m \to B_m$ comme au paragraphe
précédent, de telle sorte qu'il existe en outre des isomorphismes
$C_m/(a_t^{m - 1}) \to C_{m - 1}$ compatibles avec la structure de
$A$-algèbre donnée et avec les homomorphismes vers
$B_{m - 1} = B_m/(a_t^{m - 1})$. Fixons d'abord un choix de
$A_1 \to C_1 \to B_1$. Supposons construits
$C_{m - 1} \to C_{m - 2} \to \ldots \to C_1$ avec les propriétés voulues.
Notons que $C_m/(a_t^{m - 1})$ est étale au-dessus de
$\Spec(\bar A_{m - 1}) \setminus V(\bar I_{m - 1})$. D'après le lemme
\ref{restricted-lemma-fully-faithful-etale-over-complement}, il existe une extension étale
$C_{m - 1} \to C'_{m - 1}$ induisant un isomorphisme modulo
$\bar I_{m - 1}$, et un homomorphisme de $\bar A_{m - 1}$-algèbres
$C_m/(a_t^{m - 1}) \to C'_{m - 1}$ induisant sur les complétés l'isomorphisme
$B_m/(a_t^{m - 1}) \to B_{m - 1}$. Notons que
$C_m/(a_t^{m - 1}) \to C'_{m - 1}$ est étale sur le complémentaire de
$V(\bar I_{m - 1})$ d'après Morphismes, lemme
\ref{morphisms-lemma-etale-permanence}, et qu'au-dessus de
$V(\bar I_{m - 1})$, il induit un isomorphisme sur les complétés et y est donc
également étale (par exemple d'après Compléments sur les morphismes, lemme
\ref{more-morphisms-lemma-check-smoothness-on-infinitesimal-nbhds}). Ainsi,
$C_m/(a_t^{m - 1}) \to C'_{m - 1}$ est étale. Par invariance topologique des
morphismes étales (Morphismes étales, théorème
\ref{etale-theorem-remarkable-equivalence}), il existe un homomorphisme
d'anneaux étale $C_m \to C'_m$ tel que
$C_m/(a_t^{m - 1}) \to C'_{m - 1}$ soit isomorphe à
$C_m/(a_t^{m - 1}) \to C'_m/(a_t^{m - 1})$. Notons que le complété
$\bar I_m$-adique de $C'_m$ est égal au complété $\bar I_m$-adique de $C_m$,
c'est-à-dire à $B_m$ (détails omis). Appliquons le lemme
\ref{restricted-lemma-lift-approximation} au diagramme
$$
\xymatrix{
 & C'_m \ar[r] & C'_m/(a_t^{m - 1}) \\
C''_m \ar@{..>}[ru] \ar@{..>}[rr] & & C_{m - 1} \ar[u] \\
 & \bar A_m \ar[r] \ar[uu] \ar@{..>}[lu] & \bar A_{m - 1} \ar[u]
}
$$
pour voir qu'il existe un « relèvement » $C''_m$ de $C_{m - 1}$ en une
algèbre sur $\bar A_m$ possédant toutes les propriétés voulues.

\medskip\noindent
Par construction, $(C_m)$ est un objet de la catégorie
(\ref{restricted-equation-C}) pour l'idéal principal $(a_t)$. Par conséquent, la limite
projective $B' = \lim C_m$ est une algèbre complète pour la topologie
$(a_t)$-adique sur $A$, telle que $B'/a_t B'$ soit de type fini sur $A/(a_t)$ ; voir le
lemme \ref{restricted-lemma-topologically-finite-type}. Par construction, on a
$C_m = B'/(a_t^m)$, le complété $I$-adique de $C_m$ est $B_m$, et le complété
$I$-adique de $B'$ est isomorphe à $B$. Pour tout $m$, le complexe
$\NL_{C_m/A_m}$ a, par construction, des modules de cohomologie finis à
support dans $V(IC_m) \subset \Spec(C_m)$. Ces modules sont donc complets pour
la topologie $I$-adique (puisqu'ils sont annulés par une puissance de $I$).
Comme $\NL^\wedge_{B_m/A_m}$ est le complété $I$-adique de
$\NL_{C_m/A_m}$, voir le lemme \ref{restricted-lemma-NL-is-completion}, l'homomorphisme
$\NL_{C_m/A_m} \to \NL_{B_m/A_m}$ induit un isomorphisme en cohomologie.
Puisque $B$ est rig-étale sur $A$, il existe donc une puissance fixe de $I$
qui annule la cohomologie de $\NL_{C_m/A_m}$ pour tout $m$ ; voir le lemme
\ref{restricted-lemma-equivalent-with-artin}. Il en résulte, par le même lemme, que $B'$
est rig-étale sur $(A, (a_t))$. Enfin, on peut appliquer le lemme
\ref{restricted-lemma-approximate-principal} à $B'$ sur $(A, (a_t))$, ce qui achève la
démonstration.
\end{proof}

\begin{lemma}
\label{restricted-lemma-approximate-by-etale-over-complement}
Soit $A$ un anneau noethérien. Soit $I \subset A$ un idéal. Soit $B$ une
algèbre complète pour la topologie $I$-adique sur $A$, telle que
$A/I \to B/IB$ soit de type fini. Les conditions équivalentes du lemme
\ref{restricted-lemma-equivalent-with-artin} équivalent encore à la suivante :
\begin{enumerate}
\item[(8)]
\label{restricted-item-algebraize}
il existe une $A$-algèbre de type fini $C$ telle que
$\Spec(C) \to \Spec(A)$ soit étale au-dessus de
$\Spec(A) \setminus V(I)$ et que $B \cong C^\wedge$.
\end{enumerate}
\end{lemma}

\begin{proof}
Il suffit de combiner les lemmes \ref{restricted-lemma-equivalent-with-artin},
\ref{restricted-lemma-approximate} et \ref{restricted-lemma-rig-etale}. Un petit détail est omis.
\end{proof}




\section{Morphismes de type fini}
\label{restricted-section-finite-type}

\noindent
Dans Espaces formels, section \ref{formal-spaces-section-finite-type}, nous
avons défini les morphismes de type fini d'espaces algébriques formels. Dans la
présente section, nous étudions les types correspondants d'homomorphismes
continus entre anneaux topologiques adiques possédant un idéal de définition
de type fini. Nous conseillons vivement au lecteur de sauter cette section.

\begin{lemma}
\label{restricted-lemma-finite-type}
Soient $A$ et $B$ des anneaux topologiques adiques possédant un idéal de
définition de type fini. Soit $\varphi : A \to B$ un homomorphisme continu
d'anneaux. Les conditions suivantes sont équivalentes :
\begin{enumerate}
\item $\varphi$ est adique et $B$ est topologiquement de type fini sur $A$ ;
\item $\varphi$ est tendu et $B$ est topologiquement de type fini sur $A$ ;
\item il existe un idéal de définition $I \subset A$ tel que la topologie de
$B$ soit la topologie $I$-adique, et il existe un idéal de définition
$I' \subset A$ tel que $A/I' \to B/I'B$ soit de type fini ;
\item pour tout idéal de définition $I \subset A$, la topologie de $B$ est la
topologie $I$-adique et $A/I \to B/IB$ est de type fini ;
\item il existe un idéal de définition $I \subset A$ tel que la topologie de
$B$ soit la topologie $I$-adique et que $B$ appartienne à la catégorie
(\ref{restricted-equation-C-prime}) ;
\item pour tout idéal de définition $I \subset A$, la topologie de $B$ est la
topologie $I$-adique et $B$ appartient à la catégorie
(\ref{restricted-equation-C-prime}) ;
\item $B$, comme $A$-algèbre topologique, est le quotient de
$A\{x_1, \ldots, x_r\}$ par un idéal fermé ;
\item $B$, comme $A$-algèbre topologique, est le quotient de
$A[x_1, \ldots, x_r]^\wedge$ par un idéal fermé, où
$A[x_1, \ldots, x_r]^\wedge$ est le complété de
$A[x_1, \ldots, x_r]$ relativement à un idéal de définition de $A$ ;
\item ajouter d'autres conditions ici.
\end{enumerate}
En outre, ces conditions équivalentes définissent une propriété locale des
morphismes de $\text{WAdm}^{adic*}$ au sens d'Espaces formels, remarque
\ref{formal-spaces-remark-variant-adic-star}.
\end{lemma}

\begin{proof}
Les homomorphismes d'anneaux tendus sont définis dans Espaces formels,
définition \ref{formal-spaces-definition-taut}. Les homomorphismes d'anneaux
adiques sont définis dans Espaces formels, définition
\ref{formal-spaces-definition-adic-homomorphism}. Le lemme résulte de la
combinaison des lemmes d'Espaces formels
\ref{formal-spaces-lemma-quotient-restricted-power-series},
\ref{formal-spaces-lemma-quotient-restricted-power-series-admissible} et
\ref{formal-spaces-lemma-adic-homomorphism}. Nous omettons les détails. Pour
être précis, il n'y a aucune différence entre les anneaux topologiques
$A[x_1, \ldots, x_n]^\wedge$ et $A\{x_1, \ldots, x_r\}$ ; voir Espaces
formels, remarque
\ref{formal-spaces-remark-I-adic-completion-and-restricted-power-series}.
\end{proof}

\begin{remark}
\label{restricted-remark-NL-well-defined-topological}
Soit $A \to B$ une flèche de $\text{WAdm}^{adic*}$, adique et
topologiquement de type fini (voir le lemme \ref{restricted-lemma-finite-type}). Écrivons
$B = A\{x_1, \ldots, x_r\}/J$. On peut alors poser\footnote{En fait, cette
construction vaut pour les flèches de $\text{WAdm}^{count}$ qui satisfont les
conditions équivalentes d'Espaces formels, lemme
\ref{formal-spaces-lemma-quotient-restricted-power-series}.}
$$
\NL_{B/A}^\wedge = \left(J/J^2 \longrightarrow \bigoplus B\text{d}x_i\right)
$$
Exactement comme dans la démonstration du lemme
\ref{restricted-lemma-NL-up-to-homotopy}, le lecteur peut montrer que ce complexe de
$B$-modules est bien défini à isomorphisme unique près dans la catégorie
homotopique $K(B)$. Si maintenant $A$ est noethérien et si $I \subset A$ est
un idéal de définition, cette construction redonne le complexe cotangent naïf
de $B$ sur $(A, I)$ défini par l'équation (\ref{restricted-equation-NL}) dans la section
\ref{restricted-section-naive-cotangent-complex}, simplement parce que
$A[x_1, \ldots, x_n]^\wedge$ coïncide avec $A\{x_1, \ldots, x_r\}$ d'après
Espaces formels, remarque
\ref{formal-spaces-remark-I-adic-completion-and-restricted-power-series}. En
particulier, toujours lorsque $A$ est un anneau topologique adique noethérien,
l'objet $\NL_{B/A}^\wedge$ est indépendant du choix de l'idéal de définition
$I \subset A$.
\end{remark}

\begin{lemma}
\label{restricted-lemma-base-change-finite-type}
Considérons la propriété $P$ des flèches de $\textit{WAdm}^{adic*}$ définie
dans le lemme \ref{restricted-lemma-finite-type}. Alors $P$ est stable par changement de
base au sens d'Espaces formels, remarque
\ref{formal-spaces-remark-base-change-variant-adic-star}.
\end{lemma}

\begin{proof}
L'énoncé a un sens d'après le lemme \ref{restricted-lemma-finite-type}. Pour le vérifier,
supposons donnés des morphismes $B \to A$ et $B \to C$ dans
$\textit{WAdm}^{adic*}$ et supposons que, comme $B$-algèbre topologique, on ait
$A = B\{x_1, \ldots, x_r\}/J$ pour un certain idéal fermé $J$. Alors
$A \widehat{\otimes}_B C$ est isomorphe au quotient
$C\{x_1, \ldots, x_r\}/J'$, où $J'$ est l'adhérence de
$JC\{x_1, \ldots, x_r\}$. Certains détails sont omis.
\end{proof}

\begin{lemma}
\label{restricted-lemma-composition-finite-type}
Considérons la propriété $P$ des flèches de $\textit{WAdm}^{adic*}$ définie
dans le lemme \ref{restricted-lemma-finite-type}. Alors $P$ est stable par composition au
sens d'Espaces formels, remarque
\ref{formal-spaces-remark-composition-variant-adic-star}.
\end{lemma}

\begin{proof}
L'énoncé a un sens d'après le lemme \ref{restricted-lemma-finite-type}. La manière la plus
simple de le démontrer consiste à vérifier (a) que les composés
d'homomorphismes d'anneaux adiques entre anneaux topologiques adiques sont
adiques et (b) que la composition des homomorphismes continus d'anneaux
préserve la propriété d'être topologiquement de type fini. Nous omettons les
détails.
\end{proof}

\noindent
Le lemme suivant affirme que les morphismes d'espaces algébriques formels
adiques* sont localement de type fini si et seulement si, localement pour la
topologie étale, ils sont donnés par les types d'homomorphismes d'anneaux
topologiques décrits dans le lemme \ref{restricted-lemma-finite-type}.

\begin{lemma}
\label{restricted-lemma-finite-type-morphisms}
Soit $S$ un schéma. Soit $f : X \to Y$ un morphisme d'espaces algébriques
formels localement adiques* sur $S$. Les conditions suivantes sont
équivalentes :
\begin{enumerate}
\item pour tout diagramme commutatif
$$
\xymatrix{
U \ar[d] \ar[r] & V \ar[d] \\
X \ar[r] & Y
}
$$
où $U$ et $V$ sont des espaces algébriques formels affines, où $U \to X$ et
$V \to Y$ sont représentables par des espaces algébriques et étales, le
morphisme $U \to V$ correspond à une flèche de $\textit{WAdm}^{adic*}$ qui est
adique et topologiquement de type fini ;
\item il existe un recouvrement $\{Y_j \to Y\}$ comme dans Espaces formels,
définition \ref{formal-spaces-definition-formal-algebraic-space}, et, pour tout
$j$, un recouvrement $\{X_{ji} \to Y_j \times_Y X\}$ comme dans Espaces
formels, définition \ref{formal-spaces-definition-formal-algebraic-space}, tel
que chaque $X_{ji} \to Y_j$ corresponde à une flèche de
$\textit{WAdm}^{adic*}$ adique et topologiquement de type fini ;
\item il existe un recouvrement $\{X_i \to X\}$ comme dans Espaces formels,
définition \ref{formal-spaces-definition-formal-algebraic-space}, et, pour tout
$i$, une factorisation $X_i \to Y_i \to Y$, où $Y_i$ est un espace algébrique
formel affine, $Y_i \to Y$ est représentable par des espaces algébriques et
étale, et $X_i \to Y_i$ correspond à une flèche de
$\textit{WAdm}^{adic*}$ adique et topologiquement de type fini ;
\item $f$ est localement de type fini.
\end{enumerate}
\end{lemma}

\begin{proof}
C'est une conséquence immédiate de l'équivalence de (1) et (2) dans le lemme
\ref{restricted-lemma-finite-type} et d'Espaces formels, lemme
\ref{formal-spaces-lemma-finite-type-local-property}.
\end{proof}










\section{Type fini sur les réductions}
\label{restricted-section-finite-type-red}

\noindent
Dans cette section, nous traitons brièvement des morphismes $X \to Y$
d'espaces algébriques formels localement à indexation dénombrable tels que
$X_{red} \to Y_{red}$ soit localement de type fini. Nous allons traduire
cette propriété en une condition algébrique. Pour comprendre cette condition,
il est utile de garder à l'esprit le fait suivant :
\begin{itemize}
\item Si $A$ est un anneau topologique faiblement admissible, alors l'ensemble
$\mathfrak a \subset A$ des éléments topologiquement nilpotents
est un idéal radical ouvert et
$\text{Spf}(A)_{red} = \Spec(A/\mathfrak a)$.
\end{itemize}
Voir Espaces formels, définition
\ref{formal-spaces-definition-weakly-admissible},
lemme \ref{formal-spaces-lemma-topologically-nilpotent}, et
exemple \ref{formal-spaces-example-reduction-affine-formal-spectrum}.

\begin{lemma}
\label{restricted-lemma-finite-type-red}
Pour une flèche $\varphi : A \to B$ de $\text{WAdm}^{count}$, considérons
la propriété $P(\varphi)=$``l'homomorphisme d'anneaux induit
$A/\mathfrak a \to B/\mathfrak b$ est de type fini'',
où $\mathfrak a \subset A$ et $\mathfrak b \subset B$ sont les idéaux
des éléments topologiquement nilpotents. Alors $P$ est une propriété locale
au sens d'Espaces formels, situation
\ref{formal-spaces-situation-local-property}.
\end{lemma}

\begin{proof}
Considérons un diagramme commutatif
$$
\xymatrix{
B \ar[r] & (B')^\wedge \\
A \ar[r] \ar[u]^\varphi & (A')^\wedge \ar[u]_{\varphi'}
}
$$
comme dans Espaces formels, situation
\ref{formal-spaces-situation-local-property}.
En appliquant $\text{Spf}$ à ce diagramme, on obtient
$$
\xymatrix{
\text{Spf}(B) \ar[d] &
\text{Spf}((B')^\wedge) \ar[l] \ar[d] \\
\text{Spf}(A) &
\text{Spf}((A')^\wedge) \ar[l]
}
$$
d'espaces algébriques formels affines dont les flèches horizontales sont
représentables par des espaces algébriques et étales d'après
Espaces formels, lemme \ref{formal-spaces-lemma-etale}.
On obtient donc un diagramme commutatif de schémas affines
$$
\xymatrix{
\text{Spf}(B)_{red} \ar[d]^f &
\text{Spf}((B')^\wedge)_{red} \ar[l]^g \ar[d]^{f'} \\
\text{Spf}(A)_{red} &
\text{Spf}((A')^\wedge)_{red} \ar[l]
}
$$
dont les flèches horizontales sont étales d'après
Espaces formels, lemme \ref{formal-spaces-lemma-reduction-smooth}.
D'après Espaces formels, exemple
\ref{formal-spaces-example-reduction-affine-formal-spectrum}, et
lemme \ref{formal-spaces-lemma-etale-surjective},
les conditions (1), (2) et (3) d'Espaces formels, situation
\ref{formal-spaces-situation-local-property}
se traduisent par les énoncés suivants :
\begin{enumerate}
\item si $f$ est localement de type fini, alors $f'$ est localement de type fini ;
\item si $f'$ est localement de type fini et si $g$ est surjectif, alors
$f$ est localement de type fini ;
\item si $T_i \to S$, $i = 1, \ldots, n$, sont localement de type fini, alors
$\coprod_{i = 1, \ldots, n} T_i \to S$ est localement de type fini.
\end{enumerate}
Les propriétés (1) et (2) résultent du fait que la propriété d'être localement
de type fini est locale sur la source et sur le but pour la topologie
étale ; voir la discussion dans Morphismes d'espaces, section
\ref{spaces-morphisms-section-finite-type}.
La propriété (3) est une conséquence immédiate de la définition.
\end{proof}

\begin{lemma}
\label{restricted-lemma-base-change-finite-type-red}
Considérons la propriété $P$ des flèches de $\textit{WAdm}^{count}$ définie
dans le lemme \ref{restricted-lemma-finite-type-red}. Alors $P$ est stable par changement
de base (Espaces formels, situation
\ref{formal-spaces-situation-base-change-local-property}).
\end{lemma}

\begin{proof}
L'énoncé a un sens d'après le lemme \ref{restricted-lemma-finite-type-red}.
Pour le démontrer, supposons donnés des homomorphismes $B \to A$ et $B \to C$
dans $\textit{WAdm}^{count}$ tels que $B/\mathfrak b \to A/\mathfrak a$
soit de type fini, où $\mathfrak b \subset B$ et $\mathfrak a \subset A$
sont les idéaux des éléments topologiquement nilpotents.
Puisque $A$ et $B$ sont faiblement admissibles, les idéaux
$\mathfrak a$ et $\mathfrak b$ sont ouverts.
Soit $\mathfrak c \subset C$ l'idéal (ouvert)
des éléments topologiquement nilpotents. Il existe alors une surjection
$A \widehat{\otimes}_B C \to
A/\mathfrak a \otimes_{B/\mathfrak b} C/\mathfrak c$
dont le noyau est un idéal faible de définition et est donc constitué
d'éléments topologiquement nilpotents
(comparer avec la démonstration d'Espaces formels,
lemme \ref{formal-spaces-lemma-completed-tensor-product}). Comme
$C/\mathfrak c \to A/\mathfrak a \otimes_{B/\mathfrak b} C/\mathfrak c$
est déjà de type fini en tant que changement de base de
$B/\mathfrak b \to A/\mathfrak a$, on conclut.
\end{proof}

\begin{lemma}
\label{restricted-lemma-composition-finite-type-red}
Considérons la propriété $P$ des flèches de $\textit{WAdm}^{count}$ définie
dans le lemme \ref{restricted-lemma-finite-type-red}. Alors $P$ est stable par composition
(Espaces formels, situation
\ref{formal-spaces-situation-composition-local-property}).
\end{lemma}

\begin{proof}
Omis. Indication : les composés d'homomorphismes d'anneaux de type fini sont
de type fini.
\end{proof}

\begin{lemma}
\label{restricted-lemma-finite-type-finite-type-red}
Soit $\varphi : A \to B$ une flèche de $\textit{WAdm}^{count}$.
Si $\varphi$ est tendu et topologiquement de type fini, alors $\varphi$
satisfait la condition définie dans le lemme \ref{restricted-lemma-finite-type-red}.
\end{lemma}

\begin{proof}
C'est une conséquence immédiate des définitions.
\end{proof}

\begin{lemma}
\label{restricted-lemma-Noetherian-finite-type-red}
Soit $\varphi : A \to B$ une flèche de $\textit{WAdm}^{Noeth}$
satisfaisant la condition définie dans le lemme \ref{restricted-lemma-finite-type-red}.
Alors $A \to B$ est topologiquement de type fini.
\end{lemma}

\begin{proof}
Soit $\mathfrak b \subset B$ l'idéal des éléments topologiquement nilpotents.
Choisissons $b_1, \ldots, b_r \in B$ dont les images engendrent
$B/\mathfrak b$ sur $A$.
Choisissons des générateurs $b_{r + 1}, \ldots, b_s$ de l'idéal
$\mathfrak b$. Nous affirmons que l'image de
$$
\varphi : A[x_1, \ldots, x_s] \longrightarrow B, \quad
x_i \longmapsto b_i
$$
est dense. En effet, si $b \in \mathfrak b^n$ pour un certain $n \geq 0$,
on peut écrire
$b = \sum b_E b_{r + 1}^{e_{r + 1}} \ldots b_s^{e_s}$, où les multi-indices
$E = (e_{r + 1}, \ldots, e_s)$ satisfont $|E| = \sum e_i = n$ et où $b_E \in B$.
On peut ensuite écrire $b_E = f_E(b_1, \ldots, b_r) + b'_E$
avec $b'_E \in \mathfrak b$ et $f_E \in A[x_1, \ldots, x_r]$.
En combinant ces expressions, on obtient $b \in \Im(\varphi) + \mathfrak b^{n + 1}$.
Par récurrence, on voit que $B = \Im(\varphi) + \mathfrak b^n$ pour tout
$n \geq 0$, ce qui donne le résultat puisque $\mathfrak b$ est un idéal
de définition de $B$.
\end{proof}

\begin{lemma}
\label{restricted-lemma-Noetherian-adic-finite-type-red}
Soit $\varphi : A \to B$ une flèche de $\textit{WAdm}^{Noeth}$.
Si $\varphi$ est adique, les conditions suivantes sont équivalentes :
\begin{enumerate}
\item $\varphi$ satisfait la condition définie dans le
lemme \ref{restricted-lemma-finite-type-red} ;
\item $\varphi$ satisfait la condition définie dans le
lemme \ref{restricted-lemma-finite-type}.
\end{enumerate}
\end{lemma}

\begin{proof}
Omis. Indication : pour démontrer (1) $\Rightarrow$ (2), utiliser le
lemme \ref{restricted-lemma-Noetherian-finite-type-red}.
\end{proof}

\begin{lemma}
\label{restricted-lemma-finite-type-red-morphisms}
Soit $S$ un schéma. Soit $f : X \to Y$ un morphisme d'espaces algébriques
formels localement à indexation dénombrable sur $S$.
Les conditions suivantes sont équivalentes :
\begin{enumerate}
\item pour tout diagramme commutatif
$$
\xymatrix{
U \ar[d] \ar[r] & V \ar[d] \\
X \ar[r] & Y
}
$$
où $U$ et $V$ sont des espaces algébriques formels affines, où $U \to X$ et
$V \to Y$ sont représentables par des espaces algébriques et étales, le
morphisme $U \to V$ correspond à une flèche de $\textit{WAdm}^{count}$
satisfaisant la propriété définie dans le lemme \ref{restricted-lemma-finite-type-red} ;
\item il existe un recouvrement $\{Y_j \to Y\}$ comme dans Espaces formels,
définition \ref{formal-spaces-definition-formal-algebraic-space},
et, pour tout $j$, un recouvrement $\{X_{ji} \to Y_j \times_Y X\}$ comme dans
Espaces formels, définition
\ref{formal-spaces-definition-formal-algebraic-space}, tel que chaque
$X_{ji} \to Y_j$ corresponde à une flèche de $\textit{WAdm}^{count}$
satisfaisant la propriété définie dans le lemme \ref{restricted-lemma-finite-type-red} ;
\item il existe un recouvrement $\{X_i \to X\}$ comme dans Espaces formels,
définition \ref{formal-spaces-definition-formal-algebraic-space}, et, pour tout
$i$, une factorisation $X_i \to Y_i \to Y$, où $Y_i$ est un espace algébrique
formel affine, où $Y_i \to Y$ est représentable par des espaces algébriques et
étale, et où $X_i \to Y_i$ correspond à une flèche de
$\textit{WAdm}^{count}$ satisfaisant la propriété définie dans le
lemme \ref{restricted-lemma-finite-type-red} ;
\item le morphisme $f_{red} : X_{red} \to Y_{red}$ est localement de type fini.
\end{enumerate}
\end{lemma}

\begin{proof}
L'équivalence de (1), (2) et (3) résulte du lemme
\ref{restricted-lemma-finite-type-red} et d'une application d'Espaces formels, lemme
\ref{formal-spaces-lemma-property-defines-property-morphisms}.
Soient $Y_j$ et $X_{ji}$ comme en (2). Alors :
\begin{itemize}
\item Les familles $\{Y_{j, red} \to Y_{red}\}$ et
$\{X_{ji, red} \to X_{red}\}$ sont des recouvrements étales par des schémas
affines. Cela résulte de la discussion dans la démonstration d'Espaces formels,
lemme
\ref{formal-spaces-lemma-reduction-formal-algebraic-space}
ou directement d'Espaces formels, lemme
\ref{formal-spaces-lemma-reduction-smooth}.
\item Si $X_{ji} \to Y_j$ correspond au morphisme
$B_j \to A_{ji}$ de $\textit{WAdm}^{count}$, alors
$X_{ji, red} \to Y_{j, red}$ correspond à l'homomorphisme d'anneaux
$B_j/\mathfrak b_j \to A_{ji}/\mathfrak a_{ji}$, où
$\mathfrak b_j$ et $\mathfrak a_{ji}$ sont les idéaux des éléments
topologiquement nilpotents. Cela résulte d'Espaces formels, exemple
\ref{formal-spaces-example-reduction-affine-formal-spectrum}.
Ainsi, $X_{ji, red} \to Y_{j, red}$ est localement de type fini si et seulement
si $B_j \to A_{ji}$ satisfait la propriété définie dans le
lemme \ref{restricted-lemma-finite-type-red}.
\end{itemize}
L'équivalence de (2) et (4) résulte de ces remarques, car la propriété d'être
localement de type fini est une propriété des morphismes d'espaces algébriques
qui est locale pour la topologie étale sur la source et sur le but ; voir la
discussion dans Morphismes d'espaces, section
\ref{spaces-morphisms-section-finite-type}.
\end{proof}


















\section{Morphismes plats}
\label{restricted-section-flat}

\noindent
Dans cette section, nous définissons les morphismes plats d'espaces
algébriques formels localement noethériens.

\begin{lemma}
\label{restricted-lemma-flat-axioms}
La propriété $P(\varphi)=$``$\varphi$ is flat'' des flèches de
$\textit{WAdm}^{Noeth}$ est une propriété locale au sens d'Espaces
formels, remarque \ref{formal-spaces-remark-variant-Noetherian}.
\end{lemma}

\begin{proof}
Rappelons ce que signifie l'énoncé. Tout d'abord,
$\textit{WAdm}^{Noeth}$ est la catégorie dont les objets sont les anneaux
topologiques noethériens adiques et dont les morphismes sont les
homomorphismes continus d'anneaux. Considérons un diagramme commutatif
$$
\xymatrix{
B \ar[r] & (B')^\wedge \\
A \ar[r] \ar[u]^\varphi & (A')^\wedge \ar[u]_{\varphi'}
}
$$
satisfaisant aux conditions suivantes :
$A$ et $B$ sont des anneaux topologiques noethériens adiques,
$A \to A'$ et $B \to B'$ sont des homomorphismes étales d'anneaux,
$(A')^\wedge = \lim A'/I^nA'$ pour un idéal de définition $I \subset A$,
$(B')^\wedge = \lim B'/J^nB'$ pour un idéal de définition $J \subset B$, et
$\varphi : A \to B$ et $\varphi' : (A')^\wedge \to (B')^\wedge$
sont continus. Notons que $(A')^\wedge$ et $(B')^\wedge$ sont des anneaux
topologiques noethériens adiques d'après Espaces formels, lemme
\ref{formal-spaces-lemma-completion-in-sub}. Nous devons montrer que
\begin{enumerate}
\item $\varphi$ est plat $\Rightarrow \varphi'$ est plat ;
\item si $B \to B'$ est fidèlement plat, alors $\varphi'$ est plat
$\Rightarrow \varphi$ est plat ; et
\item si $A \to B_i$ est plat pour $i = 1, \ldots, n$, alors
$A \to \prod_{i = 1, \ldots, n} B_i$ est plat.
\end{enumerate}
Nous utiliserons sans plus de précisions que les complétés d'anneaux
noethériens sont plats (Algèbre, lemme \ref{algebra-lemma-completion-flat}).
Comme $A \to A'$ et $B \to B'$ sont bien sûr plats, les flèches
horizontales du diagramme sont en particulier plates.

\medskip\noindent
Démonstration de (1). Si $\varphi$ est plat, alors la composée
$A \to (A')^\wedge \to (B')^\wedge$ est plate. Par suite,
$A' \to (B')^\wedge$ est plat d'après Compléments sur la platitude,
lemme \ref{flat-lemma-etale-flat-up-down}. Ainsi,
$(A')^\wedge \to (B')^\wedge$ est plat par application de Compléments
sur l'algèbre, lemme \ref{more-algebra-lemma-flat-after-completion}, avec
$R = A'$, l'idéal $I(A')$ et $M = (B')^\wedge = M^\wedge$.

\medskip\noindent
Démonstration de (2). Supposons $\varphi'$ plat et $B \to B'$ fidèlement
plat. La composée $A \to (A')^\wedge \to (B')^\wedge$ est alors plate.
D'autre part, $B \to (B')^\wedge$ est fidèlement plat d'après Espaces
formels, lemme \ref{formal-spaces-lemma-etale-surjective}. Il résulte donc
d'Algèbre, lemme \ref{algebra-lemma-flatness-descends-more-general}, que
$\varphi : A \to B$ est plat.

\medskip\noindent
Démonstration de (3). Omise.
\end{proof}

\begin{lemma}
\label{restricted-lemma-base-change-flat-continuous}
Notons $P$ la propriété des flèches de $\textit{WAdm}^{Noeth}$ définie
dans le lemme \ref{restricted-lemma-flat-axioms}. Notons $Q$ la propriété définie
dans le lemme \ref{restricted-lemma-finite-type-red}, considérée comme propriété des
flèches de $\textit{WAdm}^{Noeth}$. Notons $R$ la propriété définie dans
le lemme \ref{restricted-lemma-finite-type}, considérée comme propriété des flèches
de $\textit{WAdm}^{Noeth}$. Alors
\begin{enumerate}
\item $P$ est stable par changement de base par $Q$ (Espaces formels,
remarque \ref{formal-spaces-remark-base-change-variant-variant-Noetherian}) ;
et
\item $P + R$ est stable par changement de base (Espaces formels,
remarque \ref{formal-spaces-remark-base-change-variant-Noetherian}).
\end{enumerate}
\end{lemma}

\begin{proof}
L'énoncé a un sens puisque chacune des propriétés $P$, $Q$ et $R$ est
une propriété locale des morphismes de $\textit{WAdm}^{Noeth}$.
Soient $\varphi : B \to A$ et $\psi : B \to C$ des morphismes de
$\textit{WAdm}^{Noeth}$. Si $Q(\varphi)$ ou $Q(\psi)$ est satisfaite,
alors $A \widehat{\otimes}_B C$ est noethérien d'après Espaces formels,
lemme \ref{formal-spaces-lemma-completed-tensor-product}. Comme $R$
entraîne $Q$ (lemme \ref{restricted-lemma-finite-type-finite-type-red}), cela vaut
dans les deux cas (1) et (2). C'est le premier point à vérifier. Il reste
à montrer que $C \to A \widehat{\otimes}_B C$ est plat.

\medskip\noindent
Démonstration de (1). Fixons des idéaux de définition $I \subset A$ et
$J \subset B$. D'après le lemme \ref{restricted-lemma-Noetherian-finite-type-red},
l'homomorphisme d'anneaux $B \to C$ est topologiquement de type fini.
Ainsi, $B \to C/J^n$ est de type fini pour tout $n \geq 1$. L'anneau
$A \otimes_B C/J^n$ est donc noethérien (car il est de type fini sur
$A$ et $A$ est noethérien). Par conséquent, le complété $I$-adique
$A \widehat{\otimes}_B C/J^n$ de $A \otimes_B C/J^n$ est plat sur
$C/J^n$ : en effet, $C/J^n \to A \otimes_B C/J^n$ est plat comme
changement de base de $B \to A$, et
$A \otimes_B C/J^n \to A \widehat{\otimes}_B C/J^n$ est plat d'après
Algèbre, lemme \ref{algebra-lemma-completion-flat}. Remarquons que
$A \widehat{\otimes}_B C/J^n =
(A \widehat{\otimes}_B C)/J^n(A \widehat{\otimes}_B C)$ ; nous omettons
les détails. Nous en concluons que $M = A \widehat{\otimes}_B C$ est un
$C$-module complet pour la topologie $J$-adique, tel que $M/J^nM$ soit
plat sur $C/J^n$ pour tout $n \geq 1$. Il en résulte que $M$ est plat sur
$C$ d'après Compléments sur l'algèbre, lemme
\ref{more-algebra-lemma-limit-flat}.

\medskip\noindent
Démonstration de (2). Dans ce cas, $B \to A$ est adique, et l'on a donc
simplement $A \widehat{\otimes}_B C = \lim A \otimes_B C/J^n$.
Les anneaux $A \otimes_B C/J^n$ sont noethériens par application
d'Espaces formels, lemme \ref{formal-spaces-lemma-completed-tensor-product},
où l'on remplace $C$ par $C/J^n$. On conclut comme précédemment.
\end{proof}

\begin{lemma}
\label{restricted-lemma-composition-flat-continuous}
Notons $P$ la propriété des flèches de $\textit{WAdm}^{Noeth}$ définie
dans le lemme \ref{restricted-lemma-flat-axioms}. Alors $P$ est stable par composition
(Espaces formels, remarque
\ref{formal-spaces-remark-composition-variant-Noetherian}).
\end{lemma}

\begin{proof}
Cela résulte de ce que les composées d'homomorphismes plats d'anneaux
sont plates.
\end{proof}

\begin{definition}
\label{restricted-definition-flat}
Soit $S$ un schéma. Soit $f : X \to Y$ un morphisme d'espaces algébriques
formels localement noethériens sur $S$. On dit que $f$ est {\it plat} si,
pour tout diagramme commutatif
$$
\xymatrix{
U \ar[d] \ar[r] & V \ar[d] \\
X \ar[r] & Y
}
$$
où $U$ et $V$ sont des espaces algébriques formels affines et où $U \to X$
et $V \to Y$ sont représentables par des espaces algébriques et étales,
le morphisme $U \to V$ correspond à un homomorphisme plat d'anneaux
topologiques noethériens adiques.
\end{definition}

\noindent
Montrons que cette condition peut se vérifier localement pour la topologie
étale sur la source et sur le but.

\begin{lemma}
\label{restricted-lemma-flat-morphisms}
Soit $S$ un schéma. Soit $f : X \to Y$ un morphisme d'espaces algébriques
formels localement noethériens sur $S$. Les conditions suivantes sont
équivalentes :
\begin{enumerate}
\item $f$ est plat ;
\item pour tout diagramme commutatif
$$
\xymatrix{
U \ar[d] \ar[r] & V \ar[d] \\
X \ar[r] & Y
}
$$
où $U$ et $V$ sont des espaces algébriques formels affines et où $U \to X$
et $V \to Y$ sont représentables par des espaces algébriques et étales,
le morphisme $U \to V$ correspond à un homomorphisme plat dans
$\textit{WAdm}^{Noeth}$ ;
\item il existe un recouvrement $\{Y_j \to Y\}$ comme dans Espaces
formels, définition \ref{formal-spaces-definition-formal-algebraic-space},
et, pour chaque $j$, un recouvrement $\{X_{ji} \to Y_j \times_Y X\}$
comme dans Espaces formels, définition
\ref{formal-spaces-definition-formal-algebraic-space}, tel que chaque
$X_{ji} \to Y_j$ corresponde à un homomorphisme plat dans
$\textit{WAdm}^{Noeth}$ ; et
\item il existe un recouvrement $\{X_i \to X\}$ comme dans Espaces
formels, définition \ref{formal-spaces-definition-formal-algebraic-space},
et, pour chaque $i$, une factorisation $X_i \to Y_i \to Y$, où $Y_i$ est
un espace algébrique formel affine, $Y_i \to Y$ est représentable par des
espaces algébriques et étale, et $X_i \to Y_i$ correspond à un
homomorphisme plat dans $\textit{WAdm}^{Noeth}$.
\end{enumerate}
\end{lemma}

\begin{proof}
L'équivalence de (1) et (2) est la définition \ref{restricted-definition-flat}.
L'équivalence de (2), (3) et (4) résulte de ce que la platitude est une
propriété locale des flèches de $\text{WAdm}^{Noeth}$ d'après le lemme
\ref{restricted-lemma-flat-axioms}, puis de l'application de la variante d'Espaces
formels, lemme \ref{formal-spaces-lemma-property-defines-property-morphisms},
pour les morphismes entre espaces algébriques localement noethériens,
mentionnée dans Espaces formels, remarque
\ref{formal-spaces-remark-variant-Noetherian}.
\end{proof}

\begin{lemma}
\label{restricted-lemma-base-change-flat}
Soit $S$ un schéma. Soient $f : X \to Y$ et $g : Z \to Y$ des morphismes
d'espaces algébriques formels localement noethériens sur $S$.
\begin{enumerate}
\item Si $f$ est plat et si $g_{red} : Z_{red} \to Y_{red}$ est localement
de type fini, alors le changement de base $X \times_Y Z \to Z$ est plat.
\item Si $f$ est plat et localement de type fini, alors le changement de
base $X \times_Y Z \to Z$ est plat et localement de type fini.
\end{enumerate}
\end{lemma}

\begin{proof}
L'assertion (1) résulte de la combinaison d'Espaces formels, remarque
\ref{formal-spaces-remark-base-change-variant-variant-Noetherian}, du
lemme \ref{restricted-lemma-base-change-flat-continuous}, partie (1), du lemme
\ref{restricted-lemma-flat-morphisms} et du lemme
\ref{restricted-lemma-finite-type-red-morphisms}.

\medskip\noindent
L'assertion (2) résulte de la combinaison d'Espaces formels, remarque
\ref{formal-spaces-remark-base-change-variant-Noetherian}, du lemme
\ref{restricted-lemma-base-change-flat-continuous}, partie (2), du lemme
\ref{restricted-lemma-flat-morphisms} et du lemme \ref{restricted-lemma-finite-type-morphisms}.
\end{proof}

\begin{lemma}
\label{restricted-lemma-composition-flat}
Soit $S$ un schéma. Soient $f : X \to Y$ et $g : Y \to Z$ des morphismes
d'espaces algébriques formels localement noethériens sur $S$. Si $f$ et
$g$ sont plats, alors $g \circ f$ l'est aussi.
\end{lemma}

\begin{proof}
Combiner Espaces formels, remarque
\ref{formal-spaces-remark-composition-variant-Noetherian}, et le lemme
\ref{restricted-lemma-composition-flat-continuous}.
\end{proof}

\begin{lemma}
\label{restricted-lemma-representable-flat}
Soit $S$ un schéma. Soit $f : X \to Y$ un morphisme d'espaces algébriques
formels localement noethériens sur $S$. Si $f$ est représentable par des
espaces algébriques et plat au sens d'Amorçage, définition
\ref{bootstrap-definition-property-transformation}, alors $f$ est plat
au sens de la définition \ref{restricted-definition-flat}.
\end{lemma}

\begin{proof}
Il s'agit d'une vérification de cohérence dont la démonstration devrait
être triviale, mais ne l'est pas tout à fait. Nous invitons le lecteur à
passer cette démonstration. Supposons $f$ représentable par des espaces
algébriques et plat au sens d'Amorçage, définition
\ref{bootstrap-definition-property-transformation}. Considérons un
diagramme commutatif
$$
\xymatrix{
U \ar[d] \ar[r] & V \ar[d] \\
X \ar[r] & Y
}
$$
où $U$ et $V$ sont des espaces algébriques formels affines et où $U \to X$
et $V \to Y$ sont représentables par des espaces algébriques et étales.
Alors le morphisme $U \to V$ correspond à un homomorphisme tendu
$B \to A$ de $\textit{WAdm}^{Noeth}$ d'après Espaces formels, lemme
\ref{formal-spaces-lemma-representable-local-property}. Remarquons que
cela signifie que $B \to A$ est adique (Espaces formels, lemme
\ref{formal-spaces-lemma-adic-homomorphism}) et que, en particulier, pour
tout idéal de définition $J \subset B$, la topologie de $A$ est la
topologie $J$-adique et les diagrammes
$$
\xymatrix{
\Spec(A/J^nA) \ar[r] \ar[d] & \Spec(B/J^n) \ar[d] \\
U \ar[r] & V
}
$$
sont cartésiens.

\medskip\noindent
Soit $T \to V$ un morphisme, où $T$ est un schéma. Alors
\begin{align*}
X \times_Y T \to T\text{ est plat}
& \Rightarrow
U \times_Y T \to T\text{ est plat} \\
& \Rightarrow
U \times_V V \times_Y T \to T\text{ est plat} \\
& \Rightarrow
U \times_V V \times_Y T \to V \times_Y T\text{ est plat} \\
& \Rightarrow
U \times_V T \to T\text{ est plat}
\end{align*}
Le premier énoncé est l'hypothèse faite sur $f$. La première implication
vient de ce que $U \to X$ est étale, donc plat, et de ce que les composées
de morphismes plats d'espaces algébriques sont plates. La deuxième vient
de l'égalité $U \times_Y T = U \times_V V \times_Y T$. La troisième
résulte de Compléments sur la platitude, lemme
\ref{flat-lemma-etale-flat-up-down}. La quatrième vient de ce que l'on peut
effectuer le changement de base par le morphisme $T \to V \times_Y T$.
Nous en concluons que $U \to V$ est plat au sens d'Amorçage, définition
\ref{bootstrap-definition-property-transformation}. En termes de
l'homomorphisme continu d'anneaux $B \to A$, cela signifie que les
homomorphismes $B/J^n \to A/J^nA$ sont plats (voir le diagramme ci-dessus).

\medskip\noindent
Enfin, on peut conclure, par exemple à l'aide de Compléments sur l'algèbre,
lemme \ref{more-algebra-lemma-limit-flat}, que $B \to A$ est plat.
\end{proof}





\section{Points rig-fermés}
\label{restricted-section-rig-points}

\noindent
Nous développons juste assez de théorie pour pouvoir l'utiliser afin de tester
la rig-platitude dans une section ultérieure. Le lecteur trouvera davantage
de théorie dans \cite{BL-I}, où est notamment étudié le cas des
schémas formels localement noethériens.

\begin{lemma}
\label{restricted-lemma-rig-point}
Soit $A$ un anneau topologique adique noethérien. Soit
$\mathfrak q \subset A$ un idéal premier. Les conditions suivantes sont
équivalentes :
\begin{enumerate}
\item pour un idéal de définition $I \subset A$, on a
$I \not \subset \mathfrak q$ et $\mathfrak q$ est maximal
pour cette propriété ;
\item pour un idéal de définition $I \subset A$, l'idéal premier
$\mathfrak q$ définit un point fermé de $\Spec(A) \setminus V(I)$ ;
\item pour tout idéal de définition $I \subset A$, on a
$I \not \subset \mathfrak q$ et $\mathfrak q$ est maximal
pour cette propriété ;
\item pour tout idéal de définition $I \subset A$, l'idéal premier
$\mathfrak q$ définit un point fermé de $\Spec(A) \setminus V(I)$ ;
\item $\dim(A/\mathfrak q) = 1$ et, pour un idéal de définition
$I \subset A$, on a $I \not \subset \mathfrak q$ ;
\item $\dim(A/\mathfrak q) = 1$ et, pour tout idéal de définition
$I \subset A$, on a $I \not \subset \mathfrak q$ ;
\item $\dim(A/\mathfrak q) = 1$ et la topologie induite
sur $A/\mathfrak q$ est non triviale ;
\item $A/\mathfrak q$ est un anneau local noethérien complet intègre
de dimension $1$, dont l'idéal maximal est le radical de l'image de tout idéal
de définition de $A$ ; et
\item ajouter d'autres conditions ici.
\end{enumerate}
\end{lemma}

\begin{proof}
Il est clair que (1) et (2) sont équivalentes et, pour la même raison,
que (3) et (4) sont équivalentes.
Puisque $V(I)$ est indépendant du choix de l'idéal de définition
$I$ de $A$, on voit que (2) et (4) sont équivalentes.

\medskip\noindent
Supposons satisfaites les conditions équivalentes (1) -- (4).
Si $\dim(A/\mathfrak q) > 1$, on peut choisir un idéal maximal
$\mathfrak q \subset \mathfrak m \subset A$
tel que $\dim((A/\mathfrak q)_\mathfrak m) > 1$.
Alors $\Spec((A/\mathfrak q)_\mathfrak m) - V(I(A/\mathfrak q)_\mathfrak m)$
serait infini d'après Algèbre, Lemme
\ref{algebra-lemma-Noetherian-local-domain-dim-2-infinite-opens}.
Cela contredit le fait que $\mathfrak q$ est fermé dans
$\Spec(A) \setminus V(I)$.
On voit donc que (6) est satisfaite. Il est immédiat que (6) implique (5).

\medskip\noindent
Réciproquement, supposons (5) satisfaite. Soit $I \subset A$ un idéal de définition.
Puisque $A/\mathfrak q$ est complet
pour $I(A/\mathfrak q)$ (par exemple d'après
Algèbre, Lemme \ref{algebra-lemma-completion-tensor}),
on voit que tous les points fermés de $\Spec(A/\mathfrak q)$ sont
contenus dans $V(IA/\mathfrak q)$ d'après
Algèbre, Lemme \ref{algebra-lemma-radical-completion}.
Comme $\dim(A/\mathfrak q) = 1$ et $I \not \subset \mathfrak q$,
on en déduit deux faits : (a) $V(IA/\mathfrak q)$ doit contenir
tous les points distincts du point générique de $\Spec(A/\mathfrak q)$, et
(b) $V(IA/\mathfrak q)$ doit être un ensemble discret (fini).
D'après (a), $\mathfrak q$ est un point fermé de
$\Spec(A) \setminus V(I)$, et l'on conclut que (2) est satisfaite.

\medskip\noindent
Toujours sous l'hypothèse (5), on voit que l'espace discret fini
$V(IA/\mathfrak q)$ doit être un singleton d'après Compléments d'algèbre, Lemme
\ref{more-algebra-lemma-irreducible-henselian-pair-connected}
par exemple (et le fait que les couples complets sont henséliens ; voir
Compléments d'algèbre, Lemme \ref{more-algebra-lemma-complete-henselian}).
On voit donc que (8) est vraie.
Réciproquement, il est clair que (8) implique (5).

\medskip\noindent
À ce stade, nous savons que (1) -- (6) et (8) sont équivalentes.
Nous omettons la vérification qu'elles sont aussi équivalentes à (7).
\end{proof}

\noindent
Pour parler commodément de tels idéaux premiers, nous introduisons
la terminologie non standard suivante.

\begin{definition}
\label{restricted-definition-rig-closed}
Soit $A$ un anneau topologique adique noethérien. Soit
$\mathfrak q \subset A$ un idéal premier. Nous disons que
$\mathfrak q$ est {\it rig-fermé} si les conditions équivalentes
du Lemme \ref{restricted-lemma-rig-point} sont satisfaites.
\end{definition}

\noindent
Nous aurons besoin de quelques lemmes qui disent essentiellement qu'il existe
beaucoup d'idéaux premiers rig-fermés, même dans un cadre relatif.

\begin{lemma}
\label{restricted-lemma-rig-closed-point-relative-residue-field}
Soit $\varphi : A \to B$ dans $\textit{WAdm}^{Noeth}$.
Notons $\mathfrak a \subset A$ et $\mathfrak b \subset B$
les idéaux des éléments topologiquement nilpotents. Supposons que
$A/\mathfrak a \to B/\mathfrak b$ soit de type fini.
Soit $\mathfrak q \subset B$ rig-fermé.
Le corps résiduel $\kappa$ de l'anneau local $B/\mathfrak q$
est une $A/\mathfrak a$-algèbre de type fini.
\end{lemma}

\begin{proof}
Soit $\mathfrak q \subset \mathfrak m \subset B$ l'unique
idéal maximal contenant $\mathfrak q$.
Alors $\mathfrak b \subset \mathfrak m$. Par conséquent,
$A/\mathfrak a \to B/\mathfrak b \to B/\mathfrak m = \kappa$ est
de type fini.
\end{proof}

\begin{lemma}
\label{restricted-lemma-rig-closed-point-relative}
Soit $\varphi : A \to B$ une flèche de $\textit{WAdm}^{Noeth}$
qui est adique et topologiquement de type fini.
Soit $\mathfrak q \subset B$ rig-fermé.
Soit $\mathfrak p = \varphi^{-1}(\mathfrak q) \subset A$.
Soit $\mathfrak a \subset A$ l'idéal des éléments topologiquement
nilpotents. Les conditions suivantes sont équivalentes :
\begin{enumerate}
\item le corps résiduel $\kappa$ de $B/\mathfrak q$ est fini
sur $A/\mathfrak a$ ;
\item $\mathfrak p \subset A$ est rig-fermé ;
\item $A/\mathfrak p \subset B/\mathfrak q$ est une extension finie
d'anneaux.
\end{enumerate}
\end{lemma}

\begin{proof}
Supposons (1). Rappelons que $B/\mathfrak q$ est un anneau local
noethérien de dimension $1$, dont la topologie est la topologie adique
définie par l'idéal maximal. Puisque $\varphi$ est adique, on voit que
$A \to B/\mathfrak q$ est adique. Ainsi, $\varphi(\mathfrak a)$
est un idéal non nul de $B/\mathfrak q$. Par conséquent,
$B/\mathfrak q + \varphi(\mathfrak a)$
est de longueur finie. Ainsi, $B/\mathfrak q + \varphi(\mathfrak a)$
est fini comme $A/\mathfrak a$-module par hypothèse.
Il s'ensuit que $B/\mathfrak q$ est fini sur $A$ d'après
Algèbre, Lemme \ref{algebra-lemma-finite-over-complete-ring}.
Donc (3) est satisfaite.

\medskip\noindent
Supposons (3). Alors $\Spec(B/\mathfrak q) \to \Spec(A/\mathfrak p)$
est surjectif d'après
Algèbre, Lemme \ref{algebra-lemma-integral-overring-surjective}.
Cela implique (2).

\medskip\noindent
Supposons (2). Notons $\kappa'$ le corps résiduel de $A/\mathfrak p$.
D'après le Lemme \ref{restricted-lemma-rig-closed-point-relative-residue-field}
(et le Lemme \ref{restricted-lemma-finite-type-finite-type-red}),
l'extension $\kappa/\kappa'$ est une algèbre de type fini.
D'après le théorème des zéros de Hilbert (Algèbre, Lemme
\ref{algebra-lemma-field-finite-type-over-domain})
on voit que $\kappa/\kappa'$ est une extension finie.
On voit donc que (1) est satisfaite.
\end{proof}

\begin{lemma}
\label{restricted-lemma-rig-closed-jacobson}
Soit $\varphi : A \to B$ une flèche de $\textit{WAdm}^{Noeth}$
qui est adique et topologiquement de type fini.
Soit $\mathfrak q \subset B$ rig-fermé.
Si $A/I$ est de Jacobson pour un idéal de définition $I \subset A$, alors
$\mathfrak p = \varphi^{-1}(\mathfrak q) \subset A$
est rig-fermé.
\end{lemma}

\begin{proof}
D'après le Lemme \ref{restricted-lemma-rig-closed-point-relative-residue-field}
(combiné au Lemme \ref{restricted-lemma-finite-type-finite-type-red}),
le corps résiduel $\kappa$ de $B/\mathfrak q$ est de type fini sur
$A/\mathfrak a$. Puisque $A/\mathfrak a$ est de Jacobson, on
voit que $\kappa$ est fini sur $A/\mathfrak a$ d'après
Algèbre, Lemme \ref{algebra-lemma-silly-jacobson}.
On conclut par le Lemme \ref{restricted-lemma-rig-closed-point-relative}.
\end{proof}

\begin{lemma}
\label{restricted-lemma-rig-closed-point-in-image}
Soit $\varphi : A \to B$ une flèche de $\textit{WAdm}^{Noeth}$
qui est adique et topologiquement de type fini.
Soit $\mathfrak p \subset A$ rig-fermé.
Soient $\mathfrak a \subset A$ et $\mathfrak b \subset B$
les idéaux des éléments topologiquement nilpotents. Si $\varphi$
est plat, alors les conditions suivantes sont équivalentes :
\begin{enumerate}
\item l'idéal maximal de $A/\mathfrak p$ appartient à l'image de
$\Spec(B/\mathfrak b) \to \Spec(A/\mathfrak a)$ ;
\item il existe un idéal premier rig-fermé $\mathfrak q \subset B$
tel que $\mathfrak p = \varphi^{-1}(\mathfrak q)$.
\end{enumerate}
Dans ce cas, $\varphi$, $\mathfrak p$ et $\mathfrak q$
satisfont les conclusions du Lemme \ref{restricted-lemma-rig-closed-point-relative}.
\end{lemma}

\begin{proof}
L'implication (2) $\Rightarrow$ (1) est immédiate. Supposons (1).
Pour prouver l'existence de $\mathfrak q$,
on peut remplacer $A$ par $A/\mathfrak p$ et $B$ par $B/\mathfrak p B$
(certains détails sont omis). On peut donc supposer que $(A, \mathfrak m, \kappa)$
est un anneau local noethérien complet de dimension $1$,
que $\mathfrak m = \mathfrak a$ et que $\mathfrak p = (0)$. La condition (1)
dit simplement que $B_0 = B \otimes_A \kappa = B/\mathfrak m B = B/\mathfrak a B$
est non nul. Remarquons que $B_0$ est de type fini sur $\kappa$.
On peut donc raisonner par récurrence sur $\dim(B_0)$.
Si $\dim(B_0) = 0$, tout idéal premier minimal $\mathfrak q \subset B$
convient (la platitude de $A \to B$ assure que $\mathfrak q$ est
au-dessus de $\mathfrak p = (0)$).
Si $\dim(B_0) > 0$, on peut trouver un élément $b \in B$
dont l'image $b_0 \in B_0$ est un élément non diviseur de zéro
et non inversible ; voir Algèbre, Lemme
\ref{algebra-lemma-dim-not-zero-exists-nonzerodivisor-nonunit}.
D'après Algèbre, Lemme \ref{algebra-lemma-grothendieck},
l'anneau $B' = B/bB$ est plat sur $A$. Puisque
$B'_0 = B' \otimes_A \kappa = B_0/(b_0)$ est non nul,
on peut appliquer l'hypothèse de récurrence à $B'$ et conclure.
La dernière assertion du lemme résulte clairement du
Lemme \ref{restricted-lemma-rig-closed-point-relative}.
\end{proof}

\noindent
Introduisons une notation.

\begin{definition}
\label{restricted-definition-completed-principal-localization}
Soit $A$ un anneau topologique adique possédant un idéal de définition
de type fini. Soit $f \in A$. La {\it localisation principale complétée}
$A_{\{f\}}$ de $A$ est le complété de $A_f = A[1/f]$,
la localisation principale de $A$ en $f$, relativement à n'importe quel
idéal de définition de $A$.
\end{definition}

\noindent
Bien entendu, si $f$ est topologiquement nilpotent, alors $A_{\{f\}}$
est l'anneau nul.

\begin{lemma}
\label{restricted-lemma-rig-closed-point-in-localization}
Soit $A$ un anneau topologique adique noethérien.
Soit $\mathfrak p \subset A$ un idéal premier.
Soit $f \in A$ un élément dont l'image dans $A/\mathfrak p$ est inversible.
Alors
$$
\mathfrak p A_{\{f\}} =
\mathfrak p(A_f)^\wedge =
\mathfrak p \otimes_A (A_f)^\wedge =
(\mathfrak p_f)^\wedge
$$
est un idéal premier de quotient
$$
A/\mathfrak p = (A/\mathfrak p) \otimes_A (A_f)^\wedge =
(A_f)^\wedge / \mathfrak p (A_f)^\wedge = A_{\{f\}}/\mathfrak p A_{\{f\}}
$$
\end{lemma}

\begin{proof}
Puisque $A_f$ est noethérien, l'homomorphisme d'anneaux $A \to A_f \to (A_f)^\wedge$
est plat. Pour tout $A$-module fini $M$, on voit que
$M \otimes_A (A_f)^\wedge$ est le complété de $M_f$.
Si $f$ agit comme un élément inversible sur $M$, alors $M_f = M$ est déjà complet.
Voir la discussion dans Algèbre, section \ref{algebra-section-completion-noetherian}.
Ces observations donnent aisément les assertions.
\end{proof}

\begin{lemma}
\label{restricted-lemma-rig-closed-point-after-localization}
Soit $\varphi : A \to B$ une flèche de $\textit{WAdm}^{Noeth}$
qui est adique et topologiquement de type fini.
Soit $\mathfrak q \subset B$ rig-fermé.
Il existe un $f \in A$ dont l'image dans
$B/\mathfrak q$ est inversible et tel que l'on obtienne un diagramme
$$
\vcenter{
\xymatrix{
B \ar[r] &
B_{\{f\}} \\
A \ar[r] \ar[u]_\varphi &
A_{\{f\}} \ar[u]_{\varphi_{\{f\}}}
}
}
\quad\text{avec les idéaux premiers}\quad
\vcenter{
\xymatrix{
\mathfrak q \ar@{-}[r] \ar@{-}[d] &
\mathfrak q' \ar@{-}[d] \ar@{=}[r] &
\mathfrak q B_{\{f\}} \\
\mathfrak p \ar@{-}[r] &
\mathfrak p'
}
}
$$
tel que $\mathfrak p'$ soit rig-fermé, c'est-à-dire que
l'homomorphisme $A_{\{f\}} \to B_{\{f\}}$ et les idéaux premiers
$\mathfrak q'$ et $\mathfrak p'$ satisfassent
les conditions équivalentes du Lemme \ref{restricted-lemma-rig-closed-point-relative}.
\end{lemma}

\begin{proof}
Voir le Lemme \ref{restricted-lemma-rig-closed-point-in-localization}
pour la description de $\mathfrak q'$. La seule assertion du lemme
est que, pour un choix convenable de $f$, l'idéal premier $\mathfrak p'$
vérifie $\dim((A_f)^\wedge/\mathfrak p') = 1$.
D'après le Lemme \ref{restricted-lemma-rig-closed-point-relative}, cela signifie à son tour
simplement que le corps résiduel $\kappa$ de
$B/\mathfrak q = (B_f)^\wedge/\mathfrak q'$ est fini sur
$(A_f)^\wedge/\mathfrak a' = (A/\mathfrak a)_f$.
D'après le Lemme \ref{restricted-lemma-rig-closed-point-relative-residue-field},
on sait que $A/\mathfrak a \to \kappa$ est un homomorphisme d'algèbres
de type fini. Le théorème des zéros de Hilbert, sous la forme d'Algèbre,
Lemme \ref{algebra-lemma-field-finite-type-over-domain}, permet de trouver
un $f \in A$ dont l'image dans $\kappa$ est inversible et tel que
$\kappa$ soit fini sur $A_f$. Cela achève la démonstration.
\end{proof}

\begin{lemma}
\label{restricted-lemma-rig-closed-point-variables}
Soit $A$ un anneau topologique adique noethérien. Notons $A\{x_1, \ldots, x_n\}$
l'anneau des séries formelles restreintes sur $A$. Soit
$\mathfrak q \subset A\{x_1, \ldots, x_n\}$ un idéal premier.
Posons $\mathfrak q' = A[x_1, \ldots, x_n] \cap \mathfrak q$ et
$\mathfrak p = A \cap \mathfrak q$. Si $\mathfrak q$ et $\mathfrak p$
sont rig-fermés, alors l'homomorphisme
$$
A[x_1, \ldots, x_n]_{\mathfrak q'}
\to
A\{x_1, \ldots, x_n\}_\mathfrak q
$$
induit un isomorphisme sur les complétés relativement à leurs idéaux maximaux.
\end{lemma}

\begin{proof}
D'après le Lemme \ref{restricted-lemma-rig-closed-point-relative}, l'homomorphisme d'anneaux
$A/\mathfrak p \to A\{x_1, \ldots, x_n\}/\mathfrak q$ est fini.
Pour tout $m \geq 1$, le module $\mathfrak q^m/\mathfrak q^{m + 1}$
est fini sur $A$, puisqu'il s'agit d'un
$A\{x_1, \ldots, x_n\}/\mathfrak q$-module fini.
Ainsi, $A\{x_1, \ldots x_n\}/\mathfrak q^m$ est un $A$-module fini.
Par conséquent, $A[x_1, \ldots, x_n] \to A\{x_1, \ldots, x_n\}/\mathfrak q^m$
est surjectif (car son image est dense et est un $A$-sous-module).
On en déduit directement que
$A[x_1, \ldots, x_n]/(\mathfrak q')^m \to A\{x_1, \ldots, x_n\}/\mathfrak q^m$
est un isomorphisme pour tout $m$. Le lemme en résulte aisément.
Indication : choisissons un élément topologiquement nilpotent $g \in A$ qui
n'appartient pas à $\mathfrak p$. L'homomorphisme entre les complétés est alors
$$
\lim_m \left(A[x_1, \ldots, x_n]/(\mathfrak q')^m\right)_g
\longrightarrow
\left(A\{x_1, \ldots, x_n\}/\mathfrak q^m\right)_g
$$
Certains détails sont omis.
\end{proof}

\begin{lemma}
\label{restricted-lemma-rig-closed-point-etale}
Soit $\varphi : A \to B$ une flèche de $\textit{WAdm}^{Noeth}$.
Supposons que $\varphi$ soit adique, topologiquement de type fini et plat,
et que $A/I \to B/IB$ soit étale pour un (respectivement tout)
idéal de définition $I \subset A$. Soit $\mathfrak q \subset B$
rig-fermé tel que $\mathfrak p = A \cap \mathfrak q$
soit lui aussi rig-fermé. Alors
$\mathfrak p B_\mathfrak q = \mathfrak q B_\mathfrak q$.
\end{lemma}

\begin{proof}
Soit $\kappa$ le corps résiduel de l'anneau local complet de dimension $1$
$A/\mathfrak p$. Puisque $A/I \to B/IB$ est étale, on voit que
$B \otimes_A \kappa$ est un produit fini d'extensions finies séparables
de $\kappa$ ; voir
Algèbre, Lemme \ref{algebra-lemma-etale-over-field}.
L'une d'elles est le corps résiduel de $B/\mathfrak q$.
D'après Algèbre, Lemme \ref{algebra-lemma-finite-over-complete-ring}, on voit que
$B/\mathfrak p B$ est une $A/\mathfrak p$-algèbre finie.
Elle est également plate. En combinant ce qui précède,
on voit que $A/\mathfrak p  \to B /\mathfrak p B$
est fini étale ; voir
Algèbre, Lemme \ref{algebra-lemma-characterize-etale}.
Ainsi, $B/\mathfrak p B$ est réduit, ce qui implique l'assertion du
lemme (détails omis).
\end{proof}

\begin{lemma}
\label{restricted-lemma-fibre-regular}
Soit $A$ un anneau topologique adique noethérien.
Soit $\mathfrak p \subset A$ un idéal premier rig-fermé.
Pour tout $n \geq 1$, l'homomorphisme d'anneaux
$$
A/\mathfrak p
\longrightarrow
A\{x_1, \ldots, x_n\} \otimes_A A/\mathfrak p =
A/\mathfrak p\{x_1, \ldots, x_n\}
$$
est régulier. En particulier, l'algèbre
$A\{x_1, \ldots, x_n\} \otimes_A \kappa(\mathfrak p)$
est géométriquement régulière sur $\kappa(\mathfrak p)$.
\end{lemma}

\begin{proof}
Nous utiliserons quelques faits sur les homomorphismes réguliers d'anneaux,
que le lecteur trouvera dans Compléments d'algèbre, section \ref{more-algebra-section-regular}.
Puisque $A/\mathfrak p$ est un anneau local noethérien complet, il
est excellent (Compléments d'algèbre, Proposition
\ref{more-algebra-proposition-ubiquity-excellent}).
Ainsi, $A/\mathfrak p[x_1, \ldots, x_n]$ est excellent
(d'après la même référence). Par conséquent,
$A/\mathfrak p[x_1, \ldots, x_n] \to A/\mathfrak p\{x_1, \ldots, x_n\}$
est un homomorphisme régulier d'anneaux d'après
Compléments d'algèbre, Lemme
\ref{more-algebra-lemma-map-G-ring-to-completion-regular}.
Bien sûr, $A/\mathfrak p \to A/\mathfrak p[x_1, \ldots, x_n]$
est lisse, donc régulier. Comme la composée d'homomorphismes réguliers
d'anneaux est régulière, la démonstration est achevée.
\end{proof}







\section{Homomorphismes rig-plats}
\label{restricted-section-rig-flat-homomorphisms}

\noindent
Dans cette section, nous définissons les homomorphismes rig-plats
d'anneaux topologiques adiques noethériens.

\begin{lemma}
\label{restricted-lemma-naively-rig-flat-continuous}
Soit $\varphi : A \to B$ un morphisme dans $\textit{WAdm}^{adic*}$
(Espaces formels, section \ref{formal-spaces-section-morphisms-rings}).
Supposons que $\varphi$ soit adique. Les conditions suivantes sont équivalentes :
\begin{enumerate}
\item $B_f$ est plat sur $A$ pour tout
$f \in A$ topologiquement nilpotent ;
\item $B_g$ est plat sur $A$ pour tout
$g \in B$ topologiquement nilpotent ;
\item $B_\mathfrak q$ est plat sur $A$
pour tout idéal premier $\mathfrak q \subset B$ qui ne contient
aucun idéal de définition ;
\item $B_\mathfrak q$ est plat sur $A$ pour tout idéal premier
rig-fermé $\mathfrak q \subset B$ ; et
\item ajouter d'autres conditions ici.
\end{enumerate}
\end{lemma}

\begin{proof}
Cela résulte des définitions et d'Algèbre,
lemme \ref{algebra-lemma-flat-localization}.
\end{proof}

\begin{definition}
\label{restricted-definition-naively-rig-flat}
Soit $\varphi : A \to B$ un homomorphisme continu d'anneaux
topologiques adiques noethériens, c'est-à-dire que $\varphi$
est une flèche de $\textit{WAdm}^{Noeth}$. Nous disons que $\varphi$ est
{\it naïvement rig-plat} si $\varphi$ est adique, topologiquement
de type fini et satisfait les conditions équivalentes du
lemme \ref{restricted-lemma-naively-rig-flat-continuous}.
\end{definition}

\noindent
L'exemple ci-dessous montre que cette notion ne « se localise » pas.

\begin{example}
\label{restricted-example-recompletion-not-rig-flat}
D'après Exemples, lemme \ref{examples-lemma-nonreduced-recompletion},
il existe un anneau local noethérien intègre de dimension $2$, noté $(A, \mathfrak m)$,
complet relativement à un idéal principal $I = (a)$, et un élément
$f \in \mathfrak m$, $f \not \in I$, ayant la propriété suivante :
l'anneau $A_{\{f\}}[1/a]$ n'est pas réduit.
Ici, $A_{\{f\}}$ est le complété $I$-adique $(A_f)^\wedge$ de la localisation
principale $A_f$. Précisons que l'anneau $A_{\{f\}}[1/a]$ est non nul.
Soit $B = A_{\{f\}}/ \text{nil}(A_{\{f\}})$ le quotient
par son nilradical. Remarquons que $A \to B$ est adique et
topologiquement de type fini.
En fait, $B$ est un quotient de $A\{x\} = A[x]^\wedge$ par
l'homomorphisme qui envoie $x$ sur l'image de $1/f$ dans $B$.
Tout idéal premier $\mathfrak q$ de $B$ ne contenant pas $a$ doit
être au-dessus de $(0) \subset A$\footnote{En effet, on peut trouver
$\mathfrak q \subset \mathfrak q' \subset B$ avec $a \in \mathfrak q'$
parce que $B$ est $a$-adiquement complet. Alors $\mathfrak p' = A \cap \mathfrak q'$
contient $a$, mais non $f$, et est donc un idéal premier de hauteur $1$.
Par suite, $\mathfrak p = A \cap \mathfrak q$ doit être strictement
contenu dans $\mathfrak p'$ puisque $a \not \in \mathfrak p$. Comme
$\dim(A) = 2$, on voit que $\mathfrak p = (0)$.}. Ainsi, $B_\mathfrak q$
est plat sur $A$, car c'est un module sur le
corps des fractions de $A$. Par conséquent, $A \to B$ est naïvement rig-plat.
D'autre part, l'homomorphisme
$$
A_{\{f\}}
\longrightarrow
B_{\{f\}} = (B_f)^\wedge = B = A_{\{f\}} /\text{nil}(A_{\{f\}})
$$
n'est pas plat après inversion de $a$, car on obtient la surjection non triviale
$A_{\{f\}}[1/a] \to A_{\{f\}}[1/a]/\text{nil}(A_{\{f\}}[1/a])$.
Ainsi, $A_{\{f\}} \to B_{\{f\}}^\wedge$ n'est pas naïvement rig-plat !
\end{example}

\noindent
Il est facile de contourner ce problème en adoptant
la définition suivante.

\begin{definition}
\label{restricted-definition-rig-flat-continuous-homomorphism}
Soit $\varphi : A \to B$ un homomorphisme continu entre anneaux
topologiques adiques noethériens, c'est-à-dire que $\varphi$ est une flèche de
$\textit{WAdm}^{Noeth}$. Nous disons que $\varphi$ est {\it rig-plat} si $\varphi$
est adique, topologiquement de type fini et si, pour tout $f \in A$, l'homomorphisme induit
$$
A_{\{f\}} \longrightarrow B_{\{f\}}
$$
est naïvement rig-plat (définition \ref{restricted-definition-naively-rig-flat}).
\end{definition}

\noindent
En posant $f = 1$ dans la définition précédente, on voit que la rig-platitude
implique la rig-platitude naïve. L'exemple montre que la réciproque est fausse.
Cependant, dans de nombreuses situations, il n'est pas nécessaire de se soucier
de la différence entre la rig-platitude et sa version naïve, comme le montre le lemme suivant.

\begin{lemma}
\label{restricted-lemma-rig-flat-naive}
Soit $\varphi : A \to B$ une flèche de $\textit{WAdm}^{Noeth}$.
Si $A/I$ est de Jacobson pour un idéal de définition $I \subset A$
(ou, de manière équivalente, pour tout tel idéal) et si $\varphi$ est
naïvement rig-plat, alors $\varphi$ est rig-plat.
\end{lemma}

\begin{proof}
Supposons que $\varphi$ soit naïvement rig-plat. Énonçons d'abord quelques
conséquences immédiates des hypothèses. Soit donc $f \in A$.
Alors $A, B, A_{\{f\}}, B_{\{f\}}$
sont des anneaux topologiques adiques noethériens. Les homomorphismes
$A \to A_{\{f\}} \to B_{\{f\}}$ et $A \to B \to B_{\{f\}}$
sont adiques et topologiquement de type fini.
Les homomorphismes d'anneaux $A \to A_{\{f\}}$ et $B \to B_{\{f\}}$
sont plats, car ce sont les composées de $A \to A_f$ et $B \to B_f$
avec les homomorphismes de complétion, lesquels sont plats d'après
Algèbre, lemme \ref{algebra-lemma-completion-flat}.
Les quotients de chacun des anneaux
$A, B, A_{\{f\}}, B_{\{f\}}$ par $I$ sont de type fini
sur $A/I$, donc sont eux aussi de Jacobson
(Algèbre, proposition \ref{algebra-proposition-Jacobson-permanence}).

\medskip\noindent
Soit $\mathfrak q' \subset B_{\{f\}}$ rig-fermé.
Il suffit de prouver que $(B_{\{f\}})_{\mathfrak q'}$
est plat sur $A_{\{f\}}$ ; voir le lemme \ref{restricted-lemma-naively-rig-flat-continuous}.
D'après le lemme \ref{restricted-lemma-rig-closed-jacobson}, les idéaux premiers
$\mathfrak q \subset B$, $\mathfrak p' \subset A_{\{f\}}$
et $\mathfrak p \subset A$ situés sous $\mathfrak q'$ sont rig-fermés.
Nous allons appliquer
Algèbre, lemme \ref{algebra-lemma-yet-another-variant-local-criterion-flatness},
au diagramme
$$
\xymatrix{
B_\mathfrak q \ar[r] &
(B_{\{f\}})_{\mathfrak q'} \\
A_\mathfrak p \ar[u] \ar[r] &
(A_{\{f\}})_{\mathfrak p'} \ar[u]
}
$$
avec $M = B_\mathfrak q$.
La seule hypothèse qui n'a pas encore été vérifiée est que
$\mathfrak p$ engendre l'idéal maximal de
$(A_{\{f\}})_{\mathfrak p'}$. Cela résulte du
lemme \ref{restricted-lemma-rig-closed-point-in-localization} ;
on utilise ici le fait que $\mathfrak p$ et $\mathfrak p'$ sont rig-fermés
pour voir que l'image de $f$ dans $A/\mathfrak p$ est inversible
(c'est la seule étape de la démonstration qui échoue sans
l'hypothèse de Jacobson). En effet,
cela nous dit que $A/\mathfrak p \to A_{\{f\}}/\mathfrak p'$
est une inclusion finie d'anneaux locaux
(lemme \ref{restricted-lemma-rig-closed-point-relative})
et que l'image de $f$ dans le second est inversible.
\end{proof}

\begin{lemma}
\label{restricted-lemma-rig-flat-base-change}
Soient $\varphi : A \to B$ et $A \to C$ des flèches de $\textit{WAdm}^{Noeth}$.
Supposons que $\varphi$ soit rig-plat et que $A \to C$ soit adique et
topologiquement de type fini. Alors $C \to B \widehat{\otimes}_A C$ est rig-plat.
\end{lemma}

\begin{proof}
Supposons que $\varphi$ soit rig-plat. Soit $f \in C$ un élément.
Il faut montrer que
$C_{\{f\}} \to B \widehat{\otimes}_A C_{\{f\}}$
est naïvement rig-plat. Puisque l'on peut remplacer $C$ par $C_{\{f\}}$,
il suffit de montrer que $C \to B \widehat{\otimes}_A C$
est naïvement rig-plat.

\medskip\noindent
Si $A \to C$ est surjectif, ou plus généralement si $C$ est fini comme
$A$-module, alors $B \otimes_A C = B \widehat{\otimes}_A C$,
car tout module fini sur un anneau noethérien complet est complet ; voir
Algèbre, lemme \ref{algebra-lemma-completion-tensor}.
Par le changement de base usuel pour la platitude
(Algèbre, lemme \ref{algebra-lemma-flat-base-change}),
on voit que la rig-platitude naïve de $\varphi$ implique celle de
$C \to B \times_A C$ dans ce cas.

\medskip\noindent
Dans le cas général, on peut factoriser $A \to C$ sous la forme
$A \to A\{x_1, \ldots, x_n\} \to C$
où $A\{x_1, \ldots, x_n\}$ est l'anneau des séries formelles restreintes
et $A\{x_1, \ldots, x_n\} \to C$ est surjectif. Il
suffit donc de montrer que $C \to B \widehat{\otimes}_A B$ est naïvement
rig-plat lorsque $C = A\{x_1, \ldots, x_n\}$.
Puisque $A\{x_1, \ldots, x_n\} = A\{x_1, \ldots, x_{n - 1}\}\{x_n\}$,
une récurrence sur $n$ ramène au cas traité dans le paragraphe
suivant.

\medskip\noindent
Ici, $C = A\{x\}$. Remarquons que $B \widehat{\otimes}_A C = B\{x\}$.
Il faut montrer que $A\{x\} \to B\{x\}$ est naïvement rig-plat.
Soit $\mathfrak q \subset B\{x\}$ un idéal premier rig-fermé.
Il faut montrer que $B\{x\}_{\mathfrak q}$ est plat sur $A\{x\}$.
Posons $\mathfrak p = A \cap \mathfrak q$.
D'après le lemme \ref{restricted-lemma-rig-closed-point-after-localization},
on peut trouver un $f \in A$ tel que
l'image de $f$ dans $B\{x\}/\mathfrak q$ soit inversible et que
l'idéal premier induit $\mathfrak p'$ de $A_{\{f\}}$ soit rig-fermé.
Nous utiliserons ci-dessous que $A_{\{f\}}\{x\} = A\{x\}_{\{f\}}$, et de même
pour $B$ ; les détails sont omis. Considérons le diagramme
$$
\xymatrix{
(B\{x\})_{\mathfrak q} \ar[r] &
(B_{\{f\}}\{x\})_{\mathfrak q'} \\
A\{x\} \ar[r] \ar[u] &
A_{\{f\}}\{x\} \ar[u]
}
$$
Nous voulons montrer que la flèche verticale gauche est plate.
La flèche horizontale supérieure est fidèlement plate : c'est un
homomorphisme local d'anneaux locaux, et elle est plate puisque $B_{\{f\}}\{x\}$
est le complété d'une localisation de l'anneau noethérien
$B\{x\}$. De même, la flèche horizontale inférieure est plate.
Il suffit donc de prouver que la flèche verticale droite est plate.
Nous sommes ainsi ramenés au cas traité dans le paragraphe suivant.

\medskip\noindent
Ici, $C = A\{x\}$, et nous disposons d'un idéal premier rig-fermé
$\mathfrak q \subset B\{x\}$ tel que
$\mathfrak p = A \cap \mathfrak q$ soit lui aussi rig-fermé.
Il résulte alors du lemme \ref{restricted-lemma-rig-closed-point-relative}
que les idéaux premiers intermédiaires $B \cap \mathfrak q$ et
$A\{x\} \cap \mathfrak q$ sont eux aussi rig-fermés.
Considérons le diagramme
$$
\xymatrix{
(B[x])_{B[x] \cap \mathfrak q} \ar[r] &
(B\{x\})_{\mathfrak q} \\
(A[x])_{A[x] \cap \mathfrak q} \ar[r] \ar[u] &
(A\{x\})_{A\{x\} \cap \mathfrak q} \ar[u]
}
$$
formé d'homomorphismes locaux d'anneaux locaux noethériens.
D'après le lemme \ref{restricted-lemma-rig-closed-point-variables},
les flèches horizontales induisent des isomorphismes
sur les complétés. Nous savons déjà que la flèche verticale
gauche est plate (car $A \to B$ est naïvement rig-plat,
donc $A[x] \to B[x]$ est plat hors du lieu fermé
défini par un idéal de définition).
On conclut alors par
Compléments d'algèbre, lemme \ref{more-algebra-lemma-flat-completion}.
\end{proof}

\begin{lemma}
\label{restricted-lemma-rig-flat-local-etale}
Considérons un diagramme commutatif
$$
\xymatrix{
B \ar[r] & B' \\
A \ar[r] \ar[u]^\varphi & A' \ar[u]_{\varphi'}
}
$$
dans $\textit{WAdm}^{Noeth}$, dont toutes les flèches sont adiques et
topologiquement de type fini. Supposons que $A \to A'$ et $B \to B'$ soient plats.
Soit $I \subset A$ un idéal de définition.
Si $\varphi$ est rig-plat et si $A/I \to A'/IA'$
est étale, alors $\varphi'$ est rig-plat.
\end{lemma}

\begin{proof}
Pour $f \in A'$, les hypothèses du lemme restent vraies pour le diagramme
$$
\xymatrix{
B \ar[r] & (B')_{\{f\}} \\
A \ar[r] \ar[u]^\varphi & (A')_{\{f\}} \ar[u]
}
$$
Il suffit donc de prouver que $\varphi'$ est naïvement rig-plat.

\medskip\noindent
Prenons un idéal premier rig-fermé $\mathfrak q' \subset B'$.
Il faut montrer que $(B')_{\mathfrak q'}$ est plat sur $A'$.
On peut choisir un $f \in A$ dont l'image dans $B'/\mathfrak q'$ est inversible
et tel que l'idéal premier induit $\mathfrak p''$ de $A_{\{f\}}$ 
soit rig-fermé ; voir le lemme \ref{restricted-lemma-rig-closed-point-after-localization}.
Plus précisément, on a ici $\mathfrak q'' = \mathfrak q' B'_{\{f\}}$ et
$\mathfrak p'' = A_{\{f\}} \cap \mathfrak q''$.
Considérons le diagramme
$$
\xymatrix{
B'_{\mathfrak q'} \ar[r] &
(B'_{\{f\}})_{\mathfrak q''} \\
A \ar[r] \ar[u] &
A_{\{f\}} \ar[u]
}
$$
Nous voulons montrer que la flèche verticale gauche est plate.
La flèche horizontale supérieure est fidèlement plate : c'est un
homomorphisme local d'anneaux locaux, et elle est plate puisque $B'_{\{f\}}$
est le complété d'une localisation de l'anneau noethérien
$B'_f$. De même, la flèche horizontale inférieure est plate.
Il suffit donc de prouver que la flèche verticale droite est plate.
Enfin, toutes les hypothèses du lemme restent vraies pour le diagramme
$$
\xymatrix{
B_{\{f\}} \ar[r] &
B'_{\{f\}} \\
A_{\{f\}} \ar[r] \ar[u] &
A'_{\{f\}} \ar[u]
}
$$
Cela nous ramène au cas traité dans le paragraphe suivant.

\medskip\noindent
Prenons un idéal premier rig-fermé $\mathfrak q' \subset B'$
et supposons que $\mathfrak p = A \cap \mathfrak q'$ soit lui aussi rig-fermé.
Alors les idéaux premiers $\mathfrak q = B \cap \mathfrak q'$
et $\mathfrak p' = A' \cap \mathfrak q'$ sont eux aussi rig-fermés ; voir
le lemme \ref{restricted-lemma-rig-closed-point-relative}.
Nous allons appliquer
Algèbre, lemme \ref{algebra-lemma-yet-another-variant-local-criterion-flatness},
au diagramme
$$
\xymatrix{
B_\mathfrak q \ar[r] &
B'_{\mathfrak q'} \\
A_\mathfrak p \ar[u] \ar[r] &
A'_{\mathfrak p'} \ar[u]
}
$$
avec $M = B_\mathfrak q$. La seule hypothèse qui n'a pas encore été vérifiée
est que $\mathfrak p$ engendre l'idéal maximal de
$A'_{\mathfrak p'}$. Cela résulte du
lemme \ref{restricted-lemma-rig-closed-point-etale}.
\end{proof}

\begin{lemma}
\label{restricted-lemma-rig-flat-local-down}
Considérons un diagramme commutatif
$$
\xymatrix{
B \ar[r] & B' \\
A \ar[r] \ar[u]^\varphi & A' \ar[u]_{\varphi'}
}
$$
dans $\textit{WAdm}^{Noeth}$, dont toutes les flèches sont adiques et
topologiquement de type fini. Supposons $A \to A'$ plat et $B \to B'$ fidèlement plat.
Si $\varphi'$ est rig-plat, alors $\varphi$ est rig-plat.
\end{lemma}

\begin{proof}
Pour $f \in A$, les hypothèses du lemme restent vraies pour le diagramme
$$
\xymatrix{
B_{\{f\}} \ar[r] & (B')_{\{f\}} \\
A_{\{f\}} \ar[r] \ar[u]^\varphi & (A')_{\{f\}} \ar[u]
}
$$
(Pour vérifier la condition de fidèle platitude : la fidèle platitude
de $B \to B'$ équivaut à la platitude de $B \to B'$ et à la
surjectivité de $\Spec(B'/IB') \to \Spec(B/IB)$
pour un idéal de définition $I \subset A$.)
Il suffit donc de prouver que $\varphi$ est naïvement rig-plat.
Or nous savons que $\varphi'$ est naïvement rig-plat et
que $\Spec(B') \to \Spec(B)$ est surjectif. Le
résultat en découle immédiatement.
\end{proof}

\noindent
Enfin, nous pouvons montrer que la rig-platitude est une propriété locale.

\begin{lemma}
\label{restricted-lemma-rig-flat-axioms}
La propriété $P(\varphi)=$ « $\varphi$ est rig-plat » sur les flèches
de $\textit{WAdm}^{Noeth}$ est une propriété locale au sens d'
Espaces formels, remarque \ref{formal-spaces-remark-variant-adic-star}.
\end{lemma}

\begin{proof}
Rappelons ce que signifie l'énoncé. Tout d'abord, 
$\textit{WAdm}^{Noeth}$ est la catégorie dont les objets sont les
anneaux topologiques adiques noethériens et dont les morphismes sont les
homomorphismes continus d'anneaux. Considérons un diagramme commutatif
$$
\xymatrix{
B \ar[r] & (B')^\wedge \\
A \ar[r] \ar[u]^\varphi & (A')^\wedge \ar[u]_{\varphi'}
}
$$
satisfaisant les conditions suivantes :
$A$ et $B$ sont des anneaux topologiques adiques noethériens,
$A \to A'$ et $B \to B'$ sont des homomorphismes étales d'anneaux,
$(A')^\wedge = \lim A'/I^nA'$ pour un idéal de définition $I \subset A$,
$(B')^\wedge = \lim B'/J^nB'$ pour un idéal de définition $J \subset B$, et
$\varphi : A \to B$ et $\varphi' : (A')^\wedge \to (B')^\wedge$
sont continus. Remarquons que $(A')^\wedge$ et $(B')^\wedge$ sont des
anneaux topologiques adiques noethériens d'après
Espaces formels, lemme \ref{formal-spaces-lemma-completion-in-sub}.
Il faut montrer :
\begin{enumerate}
\item $\varphi$ est rig-plat $\Rightarrow \varphi'$ est rig-plat ;
\item si $B \to B'$ est fidèlement plat, alors $\varphi'$ est rig-plat
$\Rightarrow \varphi$ est rig-plat ; et
\item si $A \to B_i$ est rig-plat pour $i = 1, \ldots, n$, alors
$A \to \prod_{i = 1, \ldots, n} B_i$ est rig-plat.
\end{enumerate}
Les propriétés d'être adique et topologiquement de type fini satisfont
les conditions (1), (2) et (3) ; voir le lemme \ref{restricted-lemma-finite-type}.
Ainsi, lorsque nous vérifions (1), (2) et (3) pour la propriété
« rig-plat », nous pouvons déjà supposer que tous nos homomorphismes d'anneaux
sont adiques et topologiquement de type fini. Alors (1) et (2) résultent
des lemmes \ref{restricted-lemma-rig-flat-local-etale} et
\ref{restricted-lemma-rig-flat-local-down}.
Nous omettons la démonstration immédiate de (3).
\end{proof}

\begin{lemma}
\label{restricted-lemma-composition-rig-flat-continuous}
La propriété $P(\varphi)=$ « $\varphi$ est rig-plat »
sur les flèches de $\textit{WAdm}^{Noeth}$ est stable par composition
au sens d'Espaces formels, remarque
\ref{formal-spaces-remark-composition-variant-Noetherian}.
\end{lemma}

\begin{proof}
L'énoncé a un sens d'après le lemme \ref{restricted-lemma-rig-flat-axioms}.
Pour voir qu'il est vrai, supposons donnés des morphismes rig-plats
$A \to B$ et $B \to C$ dans $\textit{WAdm}^{Noeth}$.
Alors $A \to C$ est adique et topologiquement de type fini
d'après le lemme \ref{restricted-lemma-composition-finite-type}.
Pour achever la démonstration, il faut montrer que, pour tout $f \in A$,
l'homomorphisme $A_{\{f\}} \to C_{\{f\}}$ est naïvement rig-plat.
Puisque $A_{\{f\}} \to B_{\{f\}}$ et $B_{\{f\}} \to C_{\{f\}}$
sont naïvement rig-plats, il suffit de montrer qu'une
composée d'homomorphismes naïvement rig-plats est naïvement rig-plate.
C'est une conséquence d'Algèbre, lemme \ref{algebra-lemma-composition-flat}.
\end{proof}











\section{Morphismes rig-plats}
\label{restricted-section-rig-flat-morphisms}

\noindent
Dans cette section, nous utilisons le travail accompli dans la
section \ref{restricted-section-rig-flat-homomorphisms}
pour définir les morphismes rig-plats d'espaces algébriques localement noethériens.

\begin{definition}
\label{restricted-definition-rig-flat}
Soit $S$ un schéma. Soit $f : X \to Y$ un morphisme d'espaces algébriques
formels localement noethériens sur $S$. Nous disons que $f$ est
{\it rig-plat} si, pour tout diagramme commutatif
$$
\xymatrix{
U \ar[d] \ar[r] & V \ar[d] \\
X \ar[r] & Y
}
$$
où $U$ et $V$ sont des espaces algébriques formels affines, où $U \to X$ et $V \to Y$
sont représentables par des espaces algébriques et étales, le morphisme $U \to V$
correspond à un homomorphisme rig-plat d'anneaux topologiques adiques noethériens.
\end{definition}

\noindent
Montrons que cette condition peut être vérifiée localement pour la topologie
étale sur la source et sur le but.

\begin{lemma}
\label{restricted-lemma-rig-flat-morphisms}
Soit $S$ un schéma. Soit $f : X \to Y$ un morphisme d'espaces
algébriques formels localement noethériens sur $S$.
Les conditions suivantes sont équivalentes :
\begin{enumerate}
\item $f$ est rig-plat ;
\item pour tout diagramme commutatif
$$
\xymatrix{
U \ar[d] \ar[r] & V \ar[d] \\
X \ar[r] & Y
}
$$
où $U$ et $V$ sont des espaces algébriques formels affines, où $U \to X$ et $V \to Y$
sont représentables par des espaces algébriques et étales, le morphisme $U \to V$
correspond à un homomorphisme rig-plat dans $\textit{WAdm}^{Noeth}$ ;
\item il existe un recouvrement $\{Y_j \to Y\}$ comme dans
Espaces formels,
définition \ref{formal-spaces-definition-formal-algebraic-space},
et, pour tout $j$,
un recouvrement $\{X_{ji} \to Y_j \times_Y X\}$ comme dans
Espaces formels,
définition \ref{formal-spaces-definition-formal-algebraic-space},
tel que chaque $X_{ji} \to Y_j$ corresponde
à un homomorphisme rig-plat dans $\textit{WAdm}^{Noeth}$ ; et
\item il existe un recouvrement $\{X_i \to X\}$ comme dans
Espaces formels,
définition \ref{formal-spaces-definition-formal-algebraic-space},
et, pour tout $i$, une factorisation $X_i \to Y_i \to Y$, où $Y_i$
est un espace algébrique formel affine, $Y_i \to Y$ est représentable
par des espaces algébriques et étale, et $X_i \to Y_i$ correspond
à un homomorphisme rig-plat dans $\textit{WAdm}^{Noeth}$.
\end{enumerate}
\end{lemma}

\begin{proof}
L'équivalence de (1) et (2) est la définition \ref{restricted-definition-rig-flat}.
L'équivalence de (2), (3) et (4) résulte du fait que
la rig-platitude est une propriété locale des flèches de
$\text{WAdm}^{Noeth}$ d'après le lemme \ref{restricted-lemma-rig-flat-axioms},
et de l'application de la variante d'
Espaces formels, lemme
\ref{formal-spaces-lemma-property-defines-property-morphisms}
pour les morphismes entre espaces algébriques localement noethériens
mentionnée dans
Espaces formels, remarque
\ref{formal-spaces-remark-variant-Noetherian}.
\end{proof}

\begin{lemma}
\label{restricted-lemma-base-change-rig-flat}
Soit $S$ un schéma. Soient $f : X \to Y$ et $g : Z \to Y$
des morphismes d'espaces algébriques formels localement noethériens sur $S$.
Si $f$ est rig-plat et $g$ est localement de type fini, alors le changement de base
$X \times_Y Z \to Z$ est rig-plat.
\end{lemma}

\begin{proof}
Cela résulte d'Espaces formels, remarque
\ref{formal-spaces-remark-base-change-variant-variant-Noetherian}
et de la discussion d'Espaces formels, section
\ref{formal-spaces-section-adic},
ainsi que du
lemme \ref{restricted-lemma-rig-flat-base-change}.
\end{proof}

\begin{lemma}
\label{restricted-lemma-composition-rig-flat}
Soit $S$ un schéma. Soient $f : X \to Y$ et $g : Y \to Z$
des morphismes d'espaces algébriques formels localement noethériens sur $S$.
Si $f$ et $g$ sont rig-plats, alors $g \circ f$ l'est aussi.
\end{lemma}

\begin{proof}
Cela résulte d'Espaces formels, remarque
\ref{formal-spaces-remark-composition-variant-Noetherian}
et du lemme \ref{restricted-lemma-composition-rig-flat-continuous}.
\end{proof}










\section{Homomorphismes rig-lisses}
\label{restricted-section-rig-smooth-homomorphisms}

\noindent
Dans cette section, nous démontrons quelques propriétés des
homomorphismes rig-lisses d'anneaux topologiques
adiques noethériens nécessaires pour introduire les
morphismes rig-lisses d'espaces algébriques formels
localement noethériens.

\begin{lemma}
\label{restricted-lemma-rig-smooth-continuous}
Soit $A \to B$ un morphisme dans $\textit{WAdm}^{Noeth}$
(Espaces formels, section \ref{formal-spaces-section-morphisms-rings}).
Les conditions suivantes sont équivalentes :
\begin{enumerate}
\item[(a)] $A \to B$ satisfait les conditions équivalentes du
lemme \ref{restricted-lemma-finite-type} et il existe un idéal de définition
$I \subset B$ tel que $B$ soit rig-lisse sur $(A, I)$ ; et
\item[(b)] $A \to B$ satisfait les conditions équivalentes du
lemme \ref{restricted-lemma-finite-type} et, pour tout idéal de définition
$I \subset A$, l'algèbre $B$ est rig-lisse sur $(A, I)$.
\end{enumerate}
\end{lemma}

\begin{proof}
Soient $I$ et $I'$ des idéaux de définition de $A$. Il existe alors un
entier $c \geq 0$ tel que $I^c \subset I'$ et $(I')^c \subset I$. Par conséquent,
$B$ est rig-lisse sur $(A, I)$ si et seulement si
$B$ est rig-lisse sur $(A, I')$. Cela résulte de la
définition \ref{restricted-definition-rig-smooth-homomorphism},
des inclusions $I^c \subset I'$ et $(I')^c \subset I$, ainsi que
du fait que le complexe cotangent naïf $\NL_{B/A}^\wedge$
est indépendant du choix de l'idéal de définition de $A$ d'après la
remarque \ref{restricted-remark-NL-well-defined-topological}.
\end{proof}

\begin{definition}
\label{restricted-definition-rig-smooth-continuous-homomorphism}
Soit $\varphi : A \to B$ un homomorphisme continu d'anneaux
topologiques adiques noethériens, c'est-à-dire que $\varphi$
est une flèche de $\textit{WAdm}^{Noeth}$. Nous disons que
$\varphi$ est {\it rig-lisse} si les conditions équivalentes
du lemme \ref{restricted-lemma-rig-smooth-continuous} sont satisfaites.
\end{definition}

\noindent
Cela définit une propriété locale.

\begin{lemma}
\label{restricted-lemma-rig-smooth-axioms}
La propriété $P(\varphi)=$ « $\varphi$ est rig-lisse » sur les flèches
de $\textit{WAdm}^{Noeth}$ est une propriété locale au sens d'
Espaces formels, remarque \ref{formal-spaces-remark-variant-Noetherian}.
\end{lemma}

\begin{proof}
Rappelons ce que signifie l'énoncé. Tout d'abord, 
$\textit{WAdm}^{Noeth}$ est la catégorie dont les objets sont les
anneaux topologiques adiques noethériens et dont les morphismes sont les
homomorphismes continus d'anneaux. Considérons un diagramme commutatif
$$
\xymatrix{
B \ar[r] & (B')^\wedge \\
A \ar[r] \ar[u]^\varphi & (A')^\wedge \ar[u]_{\varphi'}
}
$$
satisfaisant les conditions suivantes :
$A$ et $B$ sont des anneaux topologiques adiques noethériens,
$A \to A'$ et $B \to B'$ sont des homomorphismes étales d'anneaux,
$(A')^\wedge = \lim A'/I^nA'$ pour un idéal de définition $I \subset A$,
$(B')^\wedge = \lim B'/J^nB'$ pour un idéal de définition $J \subset B$, et
$\varphi : A \to B$ et $\varphi' : (A')^\wedge \to (B')^\wedge$
sont continus. Remarquons que $(A')^\wedge$ et $(B')^\wedge$ sont des
anneaux topologiques adiques noethériens d'après
Espaces formels, lemme \ref{formal-spaces-lemma-completion-in-sub}.
Il faut montrer :
\begin{enumerate}
\item $\varphi$ est rig-lisse $\Rightarrow \varphi'$ est rig-lisse ;
\item si $B \to B'$ est fidèlement plat, alors $\varphi'$ est rig-lisse
$\Rightarrow \varphi$ est rig-lisse ; et
\item si $A \to B_i$ est rig-lisse pour $i = 1, \ldots, n$, alors
$A \to \prod_{i = 1, \ldots, n} B_i$ est rig-lisse.
\end{enumerate}
Les conditions équivalentes du lemme \ref{restricted-lemma-finite-type} satisfont
les conditions (1), (2) et (3).
Ainsi, lors de la vérification de (1), (2) et (3) pour la propriété
« rig-lisse », nous pouvons déjà supposer dans chaque cas que nos
homomorphismes d'anneaux satisfont les conditions équivalentes du
lemme \ref{restricted-lemma-finite-type}.

\medskip\noindent
Choisissons un idéal de définition $I \subset A$. D'après les remarques précédentes,
la topologie de chaque anneau du diagramme est la topologie $I$-adique,
et $B$, $(A')^\wedge$ et $(B')^\wedge$ appartiennent à la catégorie
(\ref{restricted-equation-C-prime}) pour $(A, I)$.
Puisque $A \to A'$ et $B \to B'$ sont étales, les complexes
$\NL_{A'/A}$ et $\NL_{B'/B}$ sont nuls ; par conséquent,
$\NL_{(A')^\wedge/A}^\wedge$ et $\NL_{(B')^\wedge/B}^\wedge$
sont nuls d'après le lemme \ref{restricted-lemma-NL-is-completion}.
En appliquant le lemme \ref{restricted-lemma-exact-sequence-NL} à
$A \to (A')^\wedge \to (B')^\wedge$, on obtient les isomorphismes
$$
H^i(\NL_{(B')^\wedge/(A')^\wedge}^\wedge) \to H^i(\NL_{(B')^\wedge/A}^\wedge)
$$
Ainsi, $\NL_{(B')^\wedge/A}^\wedge \to \NL_{(B')^\wedge/(A')^\wedge}$
est un quasi-isomorphisme. Les homomorphismes d'anneaux $B/I^nB \to B'/I^nB'$ sont étales,
donc sont des intersections complètes locales
(Algèbre, lemme \ref{algebra-lemma-etale-standard-smooth}).
On peut donc appliquer les
lemmes \ref{restricted-lemma-exact-sequence-NL} et
\ref{restricted-lemma-transitive-lci-at-end} à
$A \to B \to (B')^\wedge$, et l'on obtient les isomorphismes
$$
H^i(\NL_{B/A}^\wedge \otimes_B (B')^\wedge) \to
H^i(\NL_{(B')^\wedge/A}^\wedge)
$$
On en conclut que
$\NL_{B/A}^\wedge \otimes_B (B')^\wedge \to \NL_{(B')^\wedge/A}^\wedge$
est un quasi-isomorphisme. En combinant ces deux observations, on obtient
$$
\NL_{(B')^\wedge/(A')^\wedge}^\wedge \cong
\NL_{B/A}^\wedge \otimes_B (B')^\wedge
$$
dans $D((B')^\wedge)$.
Ces préparatifs étant achevés, nous pouvons commencer la démonstration proprement dite.

\medskip\noindent
Preuve de (1). Supposons que $\varphi$ soit rig-lisse. Il existe alors un $c \geq 0$
tel que $\Ext^1_B(\NL_{B/A}^\wedge, N)$ soit annulé par $I^c$
pour tout $B$-module $N$. D'après
Compléments d'algèbre, lemmes \ref{more-algebra-lemma-two-term-base-change} et
\ref{more-algebra-lemma-base-change-property-ext-1-annihilated}
cette propriété est préservée par le changement de base $B \to (B')^\wedge$.
Ainsi, $\Ext^1_{(B')^\wedge}(\NL_{(B')^\wedge/(A')^\wedge}^\wedge, N)$
est annulé par $I^c(A')^\wedge$ pour tout $(B')^\wedge$-module $N$,
ce qui montre que $\varphi'$ est rig-lisse.
Ceci prouve (1).

\medskip\noindent
Pour prouver (2), supposons que $B \to B'$ soit fidèlement plat et que $\varphi'$
soit rig-lisse. Il existe alors un $c \geq 0$ tel que
$\Ext^1_{(B')^\wedge}(\NL_{(B')^\wedge/(A')^\wedge}^\wedge, N')$
soit annulé par $I^c(B')^\wedge$ pour tout $(B')^\wedge$-module $N'$.
La composée $B \to B' \to (B')^\wedge$ est plate
(Algèbre, lemme \ref{algebra-lemma-completion-flat}) ;
ainsi, pour tout $B$-module $N$, on a
$$
\Ext^1_B(\NL_{B/A}^\wedge, N) \otimes_B (B')^\wedge =
\Ext^1_{(B')^\wedge}(\NL_{B/A}^\wedge \otimes_B (B')^\wedge,
N \otimes_B (B')^\wedge)
$$
d'après Compléments d'algèbre, lemme \ref{more-algebra-lemma-base-change-RHom}, partie (3)
(quelques détails mineurs sont omis). On voit donc que ce module est annulé
par $I^c$. Or $B \to (B')^\wedge$ est en fait fidèlement plat,
car $B \to B'$ est fidèlement plat par hypothèse (Espaces formels, lemme
\ref{formal-spaces-lemma-etale-surjective}).
On en conclut que $\Ext^1_B(\NL_{B/A}^\wedge, N)$ est annulé par $I^c$.
Ainsi, $\varphi$ est rig-lisse. Ceci prouve (2).

\medskip\noindent
Pour prouver (3), posons $B = \prod_{i = 1, \ldots, n} B_i$.
Il suffit d'observer que $\NL_{B/A}^\wedge$ est la somme
directe des complexes $\NL_{B_i/A}^\wedge$, considérés comme complexes
de $B$-modules.
\end{proof}

\begin{lemma}
\label{restricted-lemma-base-change-rig-smooth-continuous}
Considérons les propriétés $P(\varphi)=$ « $\varphi$ est rig-lisse »
et $Q(\varphi)$= « $\varphi$ est adique » sur les flèches de $\textit{WAdm}^{Noeth}$.
Alors $P$ est stable par changement de base par $Q$ au sens d'
Espaces formels, remarque
\ref{formal-spaces-remark-base-change-variant-variant-Noetherian}.
\end{lemma}

\begin{proof}
L'énoncé a un sens d'après le lemme \ref{restricted-lemma-rig-smooth-continuous}.
Pour voir qu'il est vrai, supposons donnés des morphismes
$B \to A$ et $B \to C$ dans $\textit{WAdm}^{Noeth}$,
et supposons que $B \to A$ soit rig-lisse et que $B \to C$ soit adique
(Espaces formels, définition
\ref{formal-spaces-definition-adic-homomorphism}).
On peut alors choisir un idéal de définition $I \subset B$
tel que la topologie de $A$ et de $C$ soit la topologie $I$-adique.
Dans cette situation, il résulte immédiatement du
lemme \ref{restricted-lemma-base-change-rig-smooth-homomorphism} que
$A \widehat{\otimes}_B C$ est rig-lisse sur $(C, IC)$.
\end{proof}

\begin{lemma}
\label{restricted-lemma-composition-rig-smooth-continuous}
La propriété $P(\varphi)=$ « $\varphi$ est rig-lisse »
sur les flèches de $\textit{WAdm}^{Noeth}$ est stable par composition
au sens d'Espaces formels, remarque
\ref{formal-spaces-remark-composition-variant-Noetherian}.
\end{lemma}

\begin{proof}
Nous recommandons vivement au lecteur de chercher sa propre démonstration
et de ne pas lire celle qui suit. L'énoncé a un sens d'après le
lemme \ref{restricted-lemma-rig-smooth-continuous}.
Pour voir qu'il est vrai, supposons donnés des morphismes rig-lisses
$A \to B$ et $B \to C$ dans $\textit{WAdm}^{Noeth}$.
On peut alors choisir un idéal de définition $I \subset A$
tel que la topologie de $C$ et de $B$ soit la topologie $I$-adique.
D'après le lemme \ref{restricted-lemma-exact-sequence-NL}, on obtient une suite exacte
$$
\xymatrix{
C \otimes_B H^0(\NL_{B/A}^\wedge) \ar[r] &
H^0(\NL_{C/A}^\wedge) \ar[r] &
H^0(\NL_{C/B}^\wedge) \ar[r] & 0 \\
H^{-1}(\NL_{B/A}^\wedge \otimes_B C) \ar[r] &
H^{-1}(\NL_{C/A}^\wedge) \ar[r] &
H^{-1}(\NL_{C/B}^\wedge) \ar[llu]
}
$$
Remarquons que $H^{-1}(\NL_{B/A}^\wedge \otimes_B C)$
et $H^{-1}(\NL_{C/B}^\wedge)$ sont annulés par
une puissance de $I$ ; cela résulte de la partie (2) du
lemme \ref{restricted-lemma-equivalent-with-artin-smooth},
combinée aux lemmes de Compléments d'algèbre
\ref{more-algebra-lemma-two-term-base-change} et
\ref{more-algebra-lemma-base-change-property-ext-1-annihilated}
(pour traiter le changement de base $B \to C$).
Ainsi, $H^{-1}(\NL_{C/A}^\wedge)$ est annulé par une puissance de $I$.
Ensuite, d'après la caractérisation des algèbres rig-lisses donnée dans la
partie (2) du lemme \ref{restricted-lemma-equivalent-with-artin-smooth},
qui renvoie elle-même à la partie (5) de
Compléments d'algèbre, lemme \ref{more-algebra-lemma-ext-1-annihilated},
on peut choisir $f_1, \ldots, f_s \in IB$ et $g_1, \ldots, g_t \in IC$
tels que $V(f_1, \ldots, f_s) = V(IB)$ et
$V(g_1, \ldots, g_t) = V(IC)$, et tels que
$H^0(\NL_{B/A}^\wedge)_{f_i}$ soit un $B_{f_i}$-module projectif fini et
$H^0(\NL_{C/B}^\wedge)_{g_j}$ un $C_{g_j}$-module projectif fini.
Comme les cohomologies en degré $-1$ s'annulent après localisation en
$f_ig_j$, on obtient une suite exacte courte
$$
0 \to
(C \otimes_B H^0(\NL_{B/A}^\wedge))_{f_ig_j} \to
H^0(\NL_{C/A}^\wedge)_{f_ig_j} \to
H^0(\NL_{C/B}^\wedge)_{f_ig_j} \to 0
$$
et l'on conclut que $H^0(\NL_{C/A}^\wedge)_{f_ig_j}$ est un
$C_{f_ig_j}$-module projectif fini, puisqu'il est une extension des deux autres.
Ainsi, par le critère de la partie (2) du
lemme \ref{restricted-lemma-equivalent-with-artin-smooth},
et donc par celui de la partie (4) de
Compléments d'algèbre, lemme \ref{more-algebra-lemma-ext-1-annihilated},
on conclut que $C$ est rig-lisse sur $(A, I)$.
\end{proof}

\noindent
Le lemme suivant peut s'interpréter en disant qu'un homomorphisme rig-lisse
est « rig-syntomique » ou « rig-plat$+$rig-intersection complète locale ».

\begin{lemma}
\label{restricted-lemma-rig-smooth-rig-flat}
Soit $\varphi : A \to B$ une flèche de $\textit{WAdm}^{Noeth}$.
Si $\varphi$ est rig-lisse, alors $\varphi$ est rig-plat et,
pour toute présentation $B = A\{x_1, \ldots, x_n\}/J$
et tout idéal premier $J \subset \mathfrak q \subset A\{x_1, \ldots, x_n\}$
ne contenant pas un idéal de définition, l'idéal
$J_\mathfrak q \subset A\{x_1, \ldots, x_n\}_\mathfrak q$
est engendré par une suite régulière.
\end{lemma}

\begin{proof}
Soit $f \in A$. Pour prouver que $\varphi$ est rig-plat, il faut montrer
que $\varphi_{\{f\}} : A_{\{f\}} \to B_{\{f\}}$ est naïvement rig-plat.
Or, soit en considérant $\varphi_{\{f\}}$ comme un changement de base de $\varphi$
et en utilisant le lemme \ref{restricted-lemma-base-change-rig-smooth-continuous},
soit en utilisant le fait que la rig-lissité
est une propriété locale (lemme \ref{restricted-lemma-rig-smooth-axioms}), on voit que
$\varphi_{\{f\}}$ est rig-lisse. Il suffit donc de montrer
que $\varphi$ est naïvement rig-plat.

\medskip\noindent
Choisissons une présentation $B = A\{x_1, \ldots, x_n\}/J$.
Pour vérifier la seconde partie du lemme, il suffit
de vérifier que $J_\mathfrak q \subset A\{x_1, \ldots, x_n\}_\mathfrak q$
est engendré par une suite régulière pour $J \subset \mathfrak q$,
l'idéal $\mathfrak q$ étant maximal parmi ceux qui ne contiennent pas
d'idéal de définition ; voir
Algèbre, lemme \ref{algebra-lemma-regular-sequence-in-neighbourhood}
(qui montre que l'ensemble des idéaux premiers de $V(J)$ en lesquels
$J$ est engendré par une suite régulière est ouvert).
Autrement dit, on peut supposer que $\mathfrak q$ est rig-fermé
dans $A\{x_1, \ldots, x_n\}$. Pour vérifier que
$B$ est naïvement rig-plat, il suffit également de
vérifier que les localisés correspondants $B_\mathfrak q$
sont plats sur $A$.

\medskip\noindent
Soit $\mathfrak q \subset A\{x_1, \ldots, x_n\}$ rig-fermé avec
$J \subset \mathfrak q$. D'après le
lemme \ref{restricted-lemma-rig-closed-point-after-localization},
on peut choisir un $f \in A$ dont l'image dans
$A\{x_1, \ldots, x_n\}/\mathfrak q$ est inversible et tel
que l'idéal premier induit $\mathfrak p'$ de $A_{\{f\}}$ soit rig-fermé.
Nous utiliserons ci-dessous que
$A_{\{f\}}\{x_1, \ldots, x_n\} = A\{x_1, \ldots, x_n\}_{\{f\}}$ ;
les détails sont omis. Considérons le diagramme
$$
\xymatrix{
A\{x_1, \ldots, x_n\}_{\mathfrak q} / J_\mathfrak q \ar[r] &
A_{\{f\}}\{x_1, \ldots, x_n\}_{\mathfrak q'}/
J A_{\{f\}}\{x_1, \ldots, x_n\}_{\mathfrak q'} \\
A\{x_1, \ldots, x_n\}_{\mathfrak q} \ar[r] \ar[u] &
A_{\{f\}}\{x_1, \ldots, x_n\}_{\mathfrak q'} \ar[u] \\
A \ar[r] \ar[u] &
A_{\{f\}} \ar[u]
}
$$
La flèche horizontale médiane est fidèlement plate : c'est un
homomorphisme local d'anneaux locaux, et elle est plate puisque $A_{\{f\}}\{x_1, \ldots, x_n\}$
est le complété d'une localisation de l'anneau noethérien
$A\{x_1, \ldots, x_n\}$. De même, la flèche horizontale inférieure est plate.
Ainsi, pour montrer que $J_\mathfrak q$ est engendré par une suite régulière
et que $A \to A\{x_1, \ldots, x_n\}_{\mathfrak q} / J_\mathfrak q$
est plat, il suffit de prouver les mêmes assertions pour
$J A_{\{f\}}\{x_1, \ldots, x_n\}_{\mathfrak q'}$ et
$A_{\{f\}} \to A_{\{f\}}\{x_1, \ldots, x_n\}_{\mathfrak q'}/
J A_{\{f\}}\{x_1, \ldots, x_n\}_{\mathfrak q'}$.
Voir Algèbre, lemme \ref{algebra-lemma-flat-increases-depth}, ou
Compléments d'algèbre, lemme \ref{more-algebra-lemma-flat-descent-regular-ideal},
pour l'assertion sur les suites régulières. Enfin, nous avons déjà vu que
$A_{\{f\}} \to B_{\{f\}}$ est rig-lisse.
Cela nous ramène au cas traité dans le paragraphe suivant.

\medskip\noindent
Soit $\mathfrak q \subset A\{x_1, \ldots, x_n\}$ rig-fermé avec
$J \subset \mathfrak q$, et supposons en outre que $\mathfrak p = A \cap \mathfrak q$
soit lui aussi rig-fermé. D'après la caractérisation des algèbres rig-lisses
donnée dans le lemme \ref{restricted-lemma-equivalent-with-artin-smooth},
après réordonnement des variables $x_1, \ldots, x_n$,
on peut trouver $m \geq 0$ et $f_1, \ldots, f_m \in J$ tels que
\begin{enumerate}
\item $J_\mathfrak q$ est engendré par $f_1, \ldots, f_m$ ; et
\item $\det_{1 \leq i, j \leq m}(\partial f_j/ \partial x_i)$
a une image inversible dans $A\{x_1, \ldots, x_n\}_\mathfrak q$.
\end{enumerate}
D'après le lemme \ref{restricted-lemma-fibre-regular}, l'anneau fibre
$$
F = A\{x_1, \ldots, x_n\} \otimes_A \kappa(\mathfrak p)
$$
est régulier. Remarquons que les $A$-dérivations $\partial / \partial x_i$
se prolongent (de manière unique) en des dérivations $D_i : F \to F$. D'après
Compléments d'algèbre, lemme \ref{more-algebra-lemma-quotient-sequence-regular},
on voit que les images de $f_1, \ldots, f_m$ forment une suite régulière dans
$F_\mathfrak q$. La platitude de $A \to A\{x_1, \ldots, x_n\}$
et Algèbre, lemme \ref{algebra-lemma-grothendieck-regular-sequence},
montrent alors que les images de $f_1, \ldots, f_m$ forment une suite régulière dans
$A\{x_1, \ldots, x_m\}_\mathfrak q$, et que le quotient par
ces éléments est plat sur $A$. Cela achève la démonstration.
\end{proof}

\begin{lemma}
\label{restricted-lemma-exact-sequence-NL-rig-smooth}
Soient $A \to B \to C$ des flèches de $\textit{WAdm}^{Noeth}$
qui sont adiques et topologiquement de type fini. Si $B \to C$
est rig-lisse, alors le noyau de l'homomorphisme
$$
H^{-1}(\NL_{B/A}^\wedge \otimes_B C) \to H^{-1}(\NL_{C/A}^\wedge)
$$
(voir le lemme \ref{restricted-lemma-exact-sequence-NL})
est annulé par un idéal de définition.
\end{lemma}

\begin{proof}
Soit $\overline{\mathfrak q} \subset C$ un idéal premier qui ne contient
aucun idéal de définition. Puisque les modules considérés sont finis,
il suffit de montrer que
$$
H^{-1}(\NL_{B/A}^\wedge \otimes_B C)_{\overline{\mathfrak q}} \to
H^{-1}(\NL_{C/A}^\wedge)_{\overline{\mathfrak q}}
$$
est injectif. Comme dans la démonstration du lemme \ref{restricted-lemma-exact-sequence-NL},
choisissons des présentations $B = A\{x_1, \ldots, x_r\}/J$,
$C = B\{y_1, \ldots, y_s\}/J'$ et
$C = A\{x_1, \ldots, x_r, y_1, \ldots, y_s\}/K$.
En examinant le diagramme de la démonstration du lemme \ref{restricted-lemma-exact-sequence-NL},
on voit qu'il suffit de montrer que $J/J^2 \otimes_B C \to K/K^2$
est injectif après localisation en l'idéal premier
$\mathfrak q \subset A\{x_1, \ldots, x_r, y_1, \ldots, y_s\}$
correspondant à $\overline{\mathfrak q}$. Comparer avec
Compléments d'algèbre, lemme \ref{more-algebra-lemma-transitive-lci-at-end},
et sa démonstration. Cela revient à demander que $J/KJ \to K/K^2$
soit injectif après localisation en $\mathfrak q$. De manière équivalente,
il faut montrer que $J_\mathfrak q \cap K^2_\mathfrak q = (KJ)_\mathfrak q$.
D'après le lemme \ref{restricted-lemma-rig-smooth-rig-flat},
on sait que $(K/J)_\mathfrak q = J'_\mathfrak q$
est engendré par une suite régulière.
La propriété d'intersection voulue résulte donc de
Compléments d'algèbre, lemme
\ref{more-algebra-lemma-conormal-sequence-H1-regular-ideal}
(et du fait qu'un idéal engendré par une suite régulière
est $H_1$-régulier ; voir
Compléments d'algèbre, section \ref{more-algebra-section-ideals}).
\end{proof}









\section{Morphismes rig-lisses}
\label{restricted-section-rig-smooth-morphisms}

\noindent
Dans cette section, nous utilisons le travail accompli dans la
section \ref{restricted-section-rig-smooth-homomorphisms}
pour définir les morphismes rig-lisses d'espaces algébriques localement noethériens.

\begin{definition}
\label{restricted-definition-rig-smooth}
Soit $S$ un schéma. Soit $f : X \to Y$ un morphisme d'espaces algébriques
formels localement noethériens sur $S$. Nous disons que $f$ est
{\it rig-lisse} si, pour tout diagramme commutatif
$$
\xymatrix{
U \ar[d] \ar[r] & V \ar[d] \\
X \ar[r] & Y
}
$$
où $U$ et $V$ sont des espaces algébriques formels affines, où $U \to X$ et $V \to Y$
sont représentables par des espaces algébriques et étales, le morphisme $U \to V$
correspond à un homomorphisme rig-lisse d'anneaux topologiques adiques noethériens.
\end{definition}

\noindent
Montrons que cette condition peut être vérifiée localement pour la topologie
étale sur la source et sur le but.

\begin{lemma}
\label{restricted-lemma-rig-smooth-morphisms}
Soit $S$ un schéma. Soit $f : X \to Y$ un morphisme d'espaces
algébriques formels localement noethériens sur $S$.
Les conditions suivantes sont équivalentes :
\begin{enumerate}
\item $f$ est rig-lisse ;
\item pour tout diagramme commutatif
$$
\xymatrix{
U \ar[d] \ar[r] & V \ar[d] \\
X \ar[r] & Y
}
$$
où $U$ et $V$ sont des espaces algébriques formels affines, où $U \to X$ et $V \to Y$
sont représentables par des espaces algébriques et étales, le morphisme $U \to V$
correspond à un homomorphisme rig-lisse dans $\textit{WAdm}^{Noeth}$ ;
\item il existe un recouvrement $\{Y_j \to Y\}$ comme dans
Espaces formels,
définition \ref{formal-spaces-definition-formal-algebraic-space},
et, pour tout $j$,
un recouvrement $\{X_{ji} \to Y_j \times_Y X\}$ comme dans
Espaces formels,
définition \ref{formal-spaces-definition-formal-algebraic-space},
tel que chaque $X_{ji} \to Y_j$ corresponde
à un homomorphisme rig-lisse dans $\textit{WAdm}^{Noeth}$ ; et
\item il existe un recouvrement $\{X_i \to X\}$ comme dans
Espaces formels,
définition \ref{formal-spaces-definition-formal-algebraic-space},
et, pour tout $i$, une factorisation $X_i \to Y_i \to Y$, où $Y_i$
est un espace algébrique formel affine, $Y_i \to Y$ est représentable
par des espaces algébriques et étale, et $X_i \to Y_i$ correspond
à un homomorphisme rig-lisse dans $\textit{WAdm}^{Noeth}$.
\end{enumerate}
\end{lemma}

\begin{proof}
L'équivalence de (1) et (2) est la définition \ref{restricted-definition-rig-smooth}.
L'équivalence de (2), (3) et (4) résulte du fait que
la rig-lissité est une propriété locale des flèches de
$\text{WAdm}^{Noeth}$ d'après le lemme \ref{restricted-lemma-rig-smooth-axioms},
et de l'application de la variante d'
Espaces formels, lemme
\ref{formal-spaces-lemma-property-defines-property-morphisms}
pour les morphismes entre espaces algébriques localement noethériens
mentionnée dans
Espaces formels, remarque
\ref{formal-spaces-remark-variant-Noetherian}.
\end{proof}

\begin{lemma}
\label{restricted-lemma-base-change-rig-smooth}
Soit $S$ un schéma. Soient $f : X \to Y$ et $g : Z \to Y$
des morphismes d'espaces algébriques formels localement noethériens sur $S$.
Si $f$ est rig-lisse et $g$ est adique, alors le changement de base
$X \times_Y Z \to Z$ est rig-lisse.
\end{lemma}

\begin{proof}
Cela résulte d'Espaces formels, remarque
\ref{formal-spaces-remark-base-change-variant-variant-Noetherian}
et de la discussion d'Espaces formels, section
\ref{formal-spaces-section-adic},
ainsi que du lemme \ref{restricted-lemma-base-change-rig-smooth-continuous}.
\end{proof}

\begin{lemma}
\label{restricted-lemma-composition-rig-smooth}
Soit $S$ un schéma. Soient $f : X \to Y$ et $g : Y \to Z$
des morphismes d'espaces algébriques formels localement noethériens sur $S$.
Si $f$ et $g$ sont rig-lisses, alors $g \circ f$ l'est aussi.
\end{lemma}

\begin{proof}
Cela résulte d'Espaces formels, remarque
\ref{formal-spaces-remark-composition-variant-Noetherian}
et du lemme \ref{restricted-lemma-composition-rig-smooth-continuous}.
\end{proof}

\begin{lemma}
\label{restricted-lemma-rig-smooth-rig-flat-morphism}
Soit $S$ un schéma. Soit $f : X \to Y$ un morphisme d'espaces
algébriques formels localement noethériens sur $S$.
Si $f$ est rig-lisse, alors $f$ est rig-plat.
\end{lemma}

\begin{proof}
Cela résulte immédiatement du lemme \ref{restricted-lemma-rig-smooth-rig-flat}
et des définitions.
\end{proof}









\section{Homomorphismes rig-étales}
\label{restricted-section-rig-etale-homomorphisms}

\noindent
Dans cette section, nous démontrons quelques propriétés des homomorphismes
rig-étales d'anneaux topologiques adiques noethériens nécessaires pour
introduire les morphismes rig-étales d'espaces algébriques localement noethériens.

\begin{lemma}
\label{restricted-lemma-rig-etale-continuous}
Soit $A \to B$ un morphisme dans $\textit{WAdm}^{Noeth}$
(Espaces formels, section \ref{formal-spaces-section-morphisms-rings}).
Les conditions suivantes sont équivalentes :
\begin{enumerate}
\item[(a)] $A \to B$ satisfait les conditions équivalentes du
lemme \ref{restricted-lemma-finite-type} et il existe un idéal de définition
$I \subset B$ tel que $B$ soit rig-étale sur $(A, I)$ ; et
\item[(b)] $A \to B$ satisfait les conditions équivalentes du
lemme \ref{restricted-lemma-finite-type} et, pour tout idéal de définition
$I \subset A$, l'algèbre $B$ est rig-étale sur $(A, I)$.
\end{enumerate}
\end{lemma}

\begin{proof}
Soient $I$ et $I'$ des idéaux de définition de $A$. Il existe alors un
entier $c \geq 0$ tel que $I^c \subset I'$ et $(I')^c \subset I$. Par conséquent,
$B$ est rig-étale sur $(A, I)$ si et seulement si
$B$ est rig-étale sur $(A, I')$. Cela résulte de la
définition \ref{restricted-definition-rig-etale-homomorphism},
des inclusions $I^c \subset I'$ et $(I')^c \subset I$, ainsi que
du fait que le complexe cotangent naïf $\NL_{B/A}^\wedge$
est indépendant du choix de l'idéal de définition de $A$ d'après la
remarque \ref{restricted-remark-NL-well-defined-topological}.
\end{proof}

\begin{definition}
\label{restricted-definition-rig-etale-continuous-homomorphism}
Soit $\varphi : A \to B$ un homomorphisme continu d'anneaux
topologiques adiques noethériens, c'est-à-dire que $\varphi$
est une flèche de $\textit{WAdm}^{Noeth}$. Nous disons que
$\varphi$ est {\it rig-étale} si les conditions équivalentes
du lemme \ref{restricted-lemma-rig-etale-continuous} sont satisfaites.
\end{definition}

\noindent
Cela définit une propriété locale.

\begin{lemma}
\label{restricted-lemma-rig-etale-axioms}
La propriété $P(\varphi)=$ « $\varphi$ est rig-étale » sur les flèches
de $\textit{WAdm}^{Noeth}$ est une propriété locale au sens d'
Espaces formels, remarque \ref{formal-spaces-remark-variant-Noetherian}.
\end{lemma}

\begin{proof}
Cette démonstration est exactement la même que celle du
lemme \ref{restricted-lemma-rig-smooth-axioms}.
Rappelons ce que signifie l'énoncé. Tout d'abord, 
$\textit{WAdm}^{Noeth}$ est la catégorie dont les objets sont les
anneaux topologiques adiques noethériens et dont les morphismes sont les
homomorphismes continus d'anneaux. Considérons un diagramme commutatif
$$
\xymatrix{
B \ar[r] & (B')^\wedge \\
A \ar[r] \ar[u]^\varphi & (A')^\wedge \ar[u]_{\varphi'}
}
$$
satisfaisant les conditions suivantes :
$A$ et $B$ sont des anneaux topologiques adiques noethériens,
$A \to A'$ et $B \to B'$ sont des homomorphismes étales d'anneaux,
$(A')^\wedge = \lim A'/I^nA'$ pour un idéal de définition $I \subset A$,
$(B')^\wedge = \lim B'/J^nB'$ pour un idéal de définition $J \subset B$, et
$\varphi : A \to B$ et $\varphi' : (A')^\wedge \to (B')^\wedge$
sont continus. Remarquons que $(A')^\wedge$ et $(B')^\wedge$ sont des
anneaux topologiques adiques noethériens d'après
Espaces formels, lemme \ref{formal-spaces-lemma-completion-in-sub}.
Il faut montrer :
\begin{enumerate}
\item $\varphi$ est rig-étale $\Rightarrow \varphi'$ est rig-étale ;
\item si $B \to B'$ est fidèlement plat, alors $\varphi'$ est rig-étale
$\Rightarrow \varphi$ est rig-étale ; et
\item si $A \to B_i$ est rig-étale pour $i = 1, \ldots, n$, alors
$A \to \prod_{i = 1, \ldots, n} B_i$ est rig-étale.
\end{enumerate}
Les conditions équivalentes du lemme \ref{restricted-lemma-finite-type} satisfont
les conditions (1), (2) et (3).
Ainsi, lors de la vérification de (1), (2) et (3) pour la propriété
« rig-étale », nous pouvons déjà supposer dans chaque cas que nos
homomorphismes d'anneaux satisfont les conditions équivalentes du
lemme \ref{restricted-lemma-finite-type}.

\medskip\noindent
Choisissons un idéal de définition $I \subset A$. D'après les remarques précédentes,
la topologie de chaque anneau du diagramme est la topologie $I$-adique,
et $B$, $(A')^\wedge$ et $(B')^\wedge$ appartiennent à la catégorie
(\ref{restricted-equation-C-prime}) pour $(A, I)$.
Puisque $A \to A'$ et $B \to B'$ sont étales, les complexes
$\NL_{A'/A}$ et $\NL_{B'/B}$ sont nuls ; par conséquent,
$\NL_{(A')^\wedge/A}^\wedge$ et $\NL_{(B')^\wedge/B}^\wedge$
sont nuls d'après le lemme \ref{restricted-lemma-NL-is-completion}.
En appliquant le lemme \ref{restricted-lemma-exact-sequence-NL} à
$A \to (A')^\wedge \to (B')^\wedge$, on obtient les isomorphismes
$$
H^i(\NL_{(B')^\wedge/(A')^\wedge}^\wedge) \to H^i(\NL_{(B')^\wedge/A}^\wedge)
$$
Ainsi, $\NL_{(B')^\wedge/A}^\wedge \to \NL_{(B')^\wedge/(A')^\wedge}$
est un quasi-isomorphisme. Les homomorphismes d'anneaux $B/I^nB \to B'/I^nB'$ sont étales,
donc sont des intersections complètes locales
(Algèbre, lemme \ref{algebra-lemma-etale-standard-smooth}).
On peut donc appliquer les
lemmes \ref{restricted-lemma-exact-sequence-NL} et
\ref{restricted-lemma-transitive-lci-at-end} à
$A \to B \to (B')^\wedge$, et l'on obtient les isomorphismes
$$
H^i(\NL_{B/A}^\wedge \otimes_B (B')^\wedge) \to
H^i(\NL_{(B')^\wedge/A}^\wedge)
$$
On en conclut que
$\NL_{B/A}^\wedge \otimes_B (B')^\wedge \to \NL_{(B')^\wedge/A}^\wedge$
est un quasi-isomorphisme. En combinant ces deux observations, on obtient
$$
\NL_{(B')^\wedge/(A')^\wedge}^\wedge \cong
\NL_{B/A}^\wedge \otimes_B (B')^\wedge
$$
dans $D((B')^\wedge)$.
Ces préparatifs étant achevés, nous pouvons commencer la démonstration proprement dite.

\medskip\noindent
Preuve de (1). Supposons que $\varphi$ soit rig-étale. Il existe alors un $c \geq 0$
tel que la multiplication par $a \in I^c$ soit nulle sur $\NL_{B/A}^\wedge$
dans $D(B)$. Cette propriété est préservée par le changement de base
$B \to (B')^\wedge$ ; voir
Compléments d'algèbre, lemme \ref{more-algebra-lemma-two-term-base-change}.
L'isomorphisme précédent montre que $\varphi'$ est rig-étale.
Ceci prouve (1).

\medskip\noindent
Pour prouver (2), supposons que $B \to B'$ soit fidèlement plat et que $\varphi'$
soit rig-étale. Il existe alors un $c \geq 0$ tel que
la multiplication par $a \in I^c$ soit nulle sur
$\NL_{(B')^\wedge/(A')^\wedge}^\wedge$ dans $D((B')^\wedge)$.
D'après l'isomorphisme précédent, $a^c$ annule les
modules de cohomologie de
$\NL_{B/A}^\wedge \otimes_B (B')^\wedge$.
La composée $B \to (B')^\wedge$ est fidèlement plate,
car $B \to B'$ est fidèlement plat par hypothèse ; voir
Espaces formels, lemme \ref{formal-spaces-lemma-etale-surjective}.
Ainsi, les modules de cohomologie de $\NL_{B/A}^\wedge$ sont annulés
par $I^c$. Il résulte du lemme \ref{restricted-lemma-equivalent-with-artin}
que $\varphi$ est rig-étale.
Ceci prouve (2).

\medskip\noindent
Pour prouver (3), posons $B = \prod_{i = 1, \ldots, n} B_i$.
Il suffit d'observer que $\NL_{B/A}^\wedge$ est la somme
directe des complexes $\NL_{B_i/A}^\wedge$, considérés comme complexes
de $B$-modules.
\end{proof}

\begin{lemma}
\label{restricted-lemma-base-change-rig-etale-continuous}
Considérons les propriétés $P(\varphi)=$ « $\varphi$ est rig-étale »
et $Q(\varphi)$= « $\varphi$ est adique » sur les flèches de $\textit{WAdm}^{Noeth}$.
Alors $P$ est stable par changement de base par $Q$ au sens d'
Espaces formels, remarque
\ref{formal-spaces-remark-base-change-variant-variant-Noetherian}.
\end{lemma}

\begin{proof}
L'énoncé a un sens d'après le lemme \ref{restricted-lemma-rig-etale-continuous}.
Pour voir qu'il est vrai, supposons donnés des morphismes
$B \to A$ et $B \to C$ dans $\textit{WAdm}^{Noeth}$,
et supposons que $B \to A$ soit rig-étale et que $B \to C$ soit adique
(Espaces formels, définition
\ref{formal-spaces-definition-adic-homomorphism}).
On peut alors choisir un idéal de définition $I \subset B$
tel que la topologie de $A$ et de $C$ soit la topologie $I$-adique.
Dans cette situation, il résulte immédiatement que
$A \widehat{\otimes}_B C$ est rig-étale sur $(C, IC)$ d'après le
lemme \ref{restricted-lemma-base-change-rig-etale-homomorphism}.
\end{proof}

\begin{lemma}
\label{restricted-lemma-composition-rig-etale-continuous}
La propriété $P(\varphi)=$ « $\varphi$ est rig-étale »
sur les flèches de $\textit{WAdm}^{Noeth}$ est stable par composition
au sens d'Espaces formels, remarque
\ref{formal-spaces-remark-composition-variant-Noetherian}.
\end{lemma}

\begin{proof}
L'énoncé a un sens d'après le
lemme \ref{restricted-lemma-rig-etale-continuous}.
Pour voir qu'il est vrai, supposons donnés des morphismes rig-étales
$A \to B$ et $B \to C$ dans $\textit{WAdm}^{Noeth}$.
On peut alors choisir un idéal de définition $I \subset A$
tel que la topologie de $C$ et de $B$ soit la topologie $I$-adique.
D'après le lemme \ref{restricted-lemma-exact-sequence-NL}, on obtient une suite exacte
$$
\xymatrix{
C \otimes_B H^0(\NL_{B/A}^\wedge) \ar[r] &
H^0(\NL_{C/A}^\wedge) \ar[r] &
H^0(\NL_{C/B}^\wedge) \ar[r] & 0 \\
H^{-1}(\NL_{B/A}^\wedge \otimes_B C) \ar[r] &
H^{-1}(\NL_{C/A}^\wedge) \ar[r] &
H^{-1}(\NL_{C/B}^\wedge) \ar[llu]
}
$$
Il existe un $c \geq 0$ tel que, pour tout $a \in I$, 
la multiplication par $a^c$ soit nulle sur $\NL_{B/A}^\wedge$ dans $D(B)$ et sur
$\NL_{C/B}^\wedge$ dans $D(C)$. Alors, bien sûr,
la multiplication par $a^c$ est aussi nulle sur $\NL_{B/A}^\wedge \otimes_B C$
dans $D(C)$. Ainsi,
$H^0(\NL_{B/A}^\wedge) \otimes_A C$, 
$H^0(\NL_{C/B}^\wedge)$, 
$H^{-1}(\NL_{B/A}^\wedge \otimes_B C)$, et
$H^{-1}(\NL_{C/B}^\wedge)$
sont annulés par $a^c$. La suite exacte montre que
la multiplication par $a^{2c}$ est nulle sur
$H^0(\NL_{C/A}^\wedge)$ et $H^{-1}(\NL_{C/A}^\wedge)$.
Il résulte du lemme \ref{restricted-lemma-equivalent-with-artin}
que $C$ est rig-étale sur $(A, I)$, comme voulu.
\end{proof}

\begin{lemma}
\label{restricted-lemma-permanence-rig-etale-continuous}
La propriété $P(\varphi)=$ « $\varphi$ est rig-étale »
sur les flèches de $\textit{WAdm}^{Noeth}$ possède la propriété de simplification
au sens d'Espaces formels, remarque
\ref{formal-spaces-remark-permanence-variant-Noetherian}.
\end{lemma}

\begin{proof}
L'énoncé a un sens d'après le
lemme \ref{restricted-lemma-rig-etale-continuous}.
Pour voir qu'il est vrai, supposons donnés des homomorphismes
$A \to B$ et $B \to C$ dans $\textit{WAdm}^{Noeth}$,
avec $A \to C$ et $A \to B$ rig-étales.
Il faut montrer que $B \to C$ est rig-étale.
On peut alors choisir un idéal de définition $I \subset A$
tel que la topologie de $C$ et de $B$ soit la topologie $I$-adique.
D'après le lemme \ref{restricted-lemma-exact-sequence-NL}, on obtient une suite exacte
$$
\xymatrix{
C \otimes_B H^0(\NL_{B/A}^\wedge) \ar[r] &
H^0(\NL_{C/A}^\wedge) \ar[r] &
H^0(\NL_{C/B}^\wedge) \ar[r] & 0 \\
H^{-1}(\NL_{B/A}^\wedge \otimes_B C) \ar[r] &
H^{-1}(\NL_{C/A}^\wedge) \ar[r] &
H^{-1}(\NL_{C/B}^\wedge) \ar[llu]
}
$$
Il existe un $c \geq 0$ tel que, pour tout $a \in I$, 
la multiplication par $a^c$ soit nulle sur $\NL_{B/A}^\wedge$ dans $D(B)$ et sur
$\NL_{C/A}^\wedge$ dans $D(C)$. Ainsi,
$H^0(\NL_{B/A}^\wedge) \otimes_A C$,
$H^0(\NL_{C/A}^\wedge)$, et
$H^{-1}(\NL_{C/A}^\wedge)$
sont annulés par $a^c$. La suite exacte montre que
la multiplication par $a^{2c}$ est nulle sur
$H^0(\NL_{C/B}^\wedge)$ et $H^{-1}(\NL_{C/B}^\wedge)$.
Il résulte du lemme \ref{restricted-lemma-equivalent-with-artin}
que $C$ est rig-étale sur $(B, IB)$, comme voulu.
\end{proof}














\section{Morphismes rig-étales}
\label{restricted-section-rig-etale-morphisms}

\noindent
Dans cette section, nous utilisons le travail accompli dans la
section \ref{restricted-section-rig-etale-homomorphisms}
pour définir les morphismes rig-étales d'espaces algébriques localement noethériens.

\begin{definition}
\label{restricted-definition-rig-etale}
Soit $S$ un schéma. Soit $f : X \to Y$ un morphisme d'espaces algébriques
formels localement noethériens sur $S$. Nous disons que $f$ est
{\it rig-étale} si, pour tout diagramme commutatif
$$
\xymatrix{
U \ar[d] \ar[r] & V \ar[d] \\
X \ar[r] & Y
}
$$
où $U$ et $V$ sont des espaces algébriques formels affines, où $U \to X$ et $V \to Y$
sont représentables par des espaces algébriques et étales, le morphisme $U \to V$
correspond à un homomorphisme rig-étale d'anneaux topologiques adiques noethériens.
\end{definition}

\noindent
Montrons que cette condition peut être vérifiée localement pour la topologie
étale sur la source et sur le but.

\begin{lemma}
\label{restricted-lemma-rig-etale-morphisms}
Soit $S$ un schéma. Soit $f : X \to Y$ un morphisme d'espaces
algébriques formels localement noethériens sur $S$.
Les conditions suivantes sont équivalentes :
\begin{enumerate}
\item $f$ est rig-étale ;
\item pour tout diagramme commutatif
$$
\xymatrix{
U \ar[d] \ar[r] & V \ar[d] \\
X \ar[r] & Y
}
$$
où $U$ et $V$ sont des espaces algébriques formels affines, où $U \to X$ et $V \to Y$
sont représentables par des espaces algébriques et étales, le morphisme $U \to V$
correspond à un homomorphisme rig-étale dans $\textit{WAdm}^{Noeth}$ ;
\item il existe un recouvrement $\{Y_j \to Y\}$ comme dans
Espaces formels,
définition \ref{formal-spaces-definition-formal-algebraic-space},
et, pour tout $j$,
un recouvrement $\{X_{ji} \to Y_j \times_Y X\}$ comme dans
Espaces formels,
définition \ref{formal-spaces-definition-formal-algebraic-space},
tel que chaque $X_{ji} \to Y_j$ corresponde
à un homomorphisme rig-étale dans $\textit{WAdm}^{Noeth}$ ; et
\item il existe un recouvrement $\{X_i \to X\}$ comme dans
Espaces formels,
définition \ref{formal-spaces-definition-formal-algebraic-space},
et, pour tout $i$, une factorisation $X_i \to Y_i \to Y$, où $Y_i$
est un espace algébrique formel affine, $Y_i \to Y$ est représentable
par des espaces algébriques et étale, et $X_i \to Y_i$ correspond
à un homomorphisme rig-étale dans $\textit{WAdm}^{Noeth}$.
\end{enumerate}
\end{lemma}

\begin{proof}
L'équivalence de (1) et (2) est la définition \ref{restricted-definition-rig-etale}.
L'équivalence de (2), (3) et (4) résulte du fait que
la rig-étalité est une propriété locale des flèches de
$\text{WAdm}^{Noeth}$ d'après le lemme \ref{restricted-lemma-rig-etale-axioms},
et de l'application de la variante d'
Espaces formels, lemme
\ref{formal-spaces-lemma-property-defines-property-morphisms}
pour les morphismes entre espaces algébriques localement noethériens,
mentionnée dans Espaces formels, remarque
\ref{formal-spaces-remark-variant-Noetherian}.
\end{proof}

\noindent
Précisons qu'un morphisme rig-étale est localement de type fini.

\begin{lemma}
\label{restricted-lemma-rig-etale-finite-type}
Un morphisme rig-étale d'espaces algébriques formels localement noethériens
est localement de type fini.
\end{lemma}

\begin{proof}
La propriété $P$ du lemme \ref{restricted-lemma-rig-etale-axioms}
implique les conditions équivalentes (a), (b), (c) et (d) d'
Espaces formels, lemme
\ref{formal-spaces-lemma-quotient-restricted-power-series}.
L'assertion résulte donc d'Espaces formels,
lemme \ref{formal-spaces-lemma-finite-type-local-property}.
\end{proof}

\begin{lemma}
\label{restricted-lemma-rig-etale-rig-smooth-morphism}
Un morphisme rig-étale d'espaces algébriques formels localement noethériens
est rig-lisse.
\end{lemma}

\begin{proof}
Cela résulte des définitions et du
lemme \ref{restricted-lemma-rig-etale-rig-smooth}.
\end{proof}

\begin{lemma}
\label{restricted-lemma-base-change-rig-etale}
Soit $S$ un schéma. Soient $f : X \to Y$ et $g : Z \to Y$
des morphismes d'espaces algébriques formels localement noethériens sur $S$.
Si $f$ est rig-étale et $g$ est adique, alors le changement de base
$X \times_Y Z \to Z$ est rig-étale.
\end{lemma}

\begin{proof}
Cela résulte d'Espaces formels, remarque
\ref{formal-spaces-remark-base-change-variant-variant-Noetherian}
et de la discussion d'Espaces formels, section
\ref{formal-spaces-section-adic},
ainsi que du lemme \ref{restricted-lemma-base-change-rig-etale-continuous}.
\end{proof}

\begin{lemma}
\label{restricted-lemma-composition-rig-etale}
Soit $S$ un schéma. Soient $f : X \to Y$ et $g : Y \to Z$
des morphismes d'espaces algébriques formels localement noethériens sur $S$.
Si $f$ et $g$ sont rig-étales, alors $g \circ f$ l'est aussi.
\end{lemma}

\begin{proof}
Cela résulte d'Espaces formels, remarque
\ref{formal-spaces-remark-composition-variant-Noetherian}
et du lemme \ref{restricted-lemma-composition-rig-etale-continuous}.
\end{proof}

\begin{lemma}
\label{restricted-lemma-rig-etale-permanence}
Soit $S$ un schéma. Soient $f : X \to Y$ et $g : Y \to Z$
des morphismes d'espaces algébriques formels localement noethériens sur $S$.
Si $g \circ f$ et $g$ sont rig-étales, alors $f$ l'est aussi.
\end{lemma}

\begin{proof}
Cela résulte d'Espaces formels, remarque
\ref{formal-spaces-remark-permanence-variant-Noetherian}
et du lemme \ref{restricted-lemma-permanence-rig-etale-continuous}.
\end{proof}

\begin{lemma}
\label{restricted-lemma-rig-etale-alternative-permanence}
Soit $S$ un schéma. Soient $f : X \to Y$ et $g : Y \to Z$
des morphismes d'espaces algébriques formels localement noethériens sur $S$.
Si $g \circ f$ est rig-étale et $g$ est un monomorphisme adique, alors
$f$ est rig-étale.
\end{lemma}

\begin{proof}
Utiliser le lemme \ref{restricted-lemma-base-change-rig-etale} et le fait que
$f$ est le changement de base de $g \circ f$ par $g$.
\end{proof}

\begin{lemma}
\label{restricted-lemma-closed-immersion-rig-smooth}
Soit $S$ un schéma. Soit $f : X \to Y$ un morphisme d'espaces algébriques
formels. Supposons que $X$ et $Y$ soient localement noethériens et que $f$ soit une
immersion fermée. Les conditions suivantes sont équivalentes :
\begin{enumerate}
\item $f$ est rig-lisse ;
\item $f$ est rig-étale ;
\item pour tout espace algébrique formel affine $V$ et tout morphisme
$V \to Y$ représentable par des espaces algébriques et étale,
le morphisme $X \times_Y V \to V$ correspond à un morphisme surjectif
$B \to A$ dans $\textit{WAdm}^{Noeth}$ dont le noyau $J$ possède la propriété
suivante : $I(J/J^2) = 0$ pour un idéal de définition $I$ de $B$.
\end{enumerate}
\end{lemma}

\begin{proof}
Remarquons que, étant donnés $V$ et $V \to Y$ comme en (2), sans autre
hypothèse sur $f$, le morphisme $X \times_Y V \to V$
correspond à un morphisme surjectif $B \to A$ dans $\textit{WAdm}^{Noeth}$
d'après Espaces formels, lemme
\ref{formal-spaces-lemma-closed-immersion-into-countably-indexed}.

\medskip\noindent
On a (2) $\Rightarrow$ (1) d'après le
lemme \ref{restricted-lemma-rig-etale-rig-smooth-morphism}.

\medskip\noindent
Démonstration de (3) $\Rightarrow$ (2). Supposons (3). D'après le
lemme \ref{restricted-lemma-rig-etale-morphisms},
il suffit de montrer que les homomorphismes d'anneaux
$B \to A$ qui apparaissent en (3) sont rig-étales au sens de la
définition \ref{restricted-definition-rig-etale-continuous-homomorphism}.
Soit $I$ comme en (3). Le complexe cotangent naïf
$\NL_{A/B}^\wedge$ de $A$ sur $(B, I)$ est le complexe de $A$-modules
obtenu en plaçant $J/J^2$ en degré $-1$. Par conséquent, $A$ est
rig-étale sur $(B, I)$ d'après la
définition \ref{restricted-definition-rig-etale-homomorphism}.

\medskip\noindent
Supposons (1), et soient $V$ et $B \to A$ comme en (3).
D'après la définition \ref{restricted-definition-rig-smooth}, on voit que
$B \to A$ est rig-lisse. Choisissons un idéal de définition quelconque $I \subset B$.
Alors $A$ est rig-lisse sur $(B, I)$.
Comme ci-dessus, le complexe $\NL_{A/B}^\wedge$ est 
obtenu en plaçant $J/J^2$ en degré $-1$.
Ainsi, d'après le lemme \ref{restricted-lemma-equivalent-with-artin-smooth},
on voit que $J/J^2$ est annulé par
une puissance $I^n$ pour un $n \geq 1$. Puisque $B$ est adique, on voit
que $I^n$ est un idéal de définition de $B$, ce qui achève
la démonstration.
\end{proof}














\section{Morphismes rig-surjectifs}
\label{restricted-section-rig-surjective}

\noindent
Pour les morphismes localement de type fini entre espaces algébriques formels
localement noethériens, on peut utiliser une définition empruntée à \cite{ArtinII}.
Voir la remarque \ref{restricted-remark-rig-surjective-more-general} pour une discussion
du traitement des cas plus généraux.

\begin{definition}
\label{restricted-definition-rig-surjective}
Soit $S$ un schéma. Soit $f : X \to Y$ un morphisme d'espaces
algébriques formels sur $S$. Supposons que $X$ et $Y$ soient localement
noethériens et que $f$ soit localement de type fini. Nous disons que
$f$ est {\it rig-surjectif} si, pour tout diagramme à flèches pleines
$$
\xymatrix{
\text{Spf}(R') \ar@{..>}[r] \ar@{..>}[d] & X \ar[d]^f \\
\text{Spf}(R) \ar[r]^-p & Y
}
$$
où $R$ est un anneau de valuation discrète complet et où
$p$ est un morphisme adique, il existe une
extension d'anneaux de valuation discrète complets $R \subset R'$
et un morphisme $\text{Spf}(R') \to X$ rendant commutatif le diagramme affiché.
\end{definition}

\noindent
Nous verrons dans les lemmes ci-dessous que cette notion se comporte raisonnablement
dans le contexte des espaces algébriques formels localement noethériens et des morphismes
localement de type fini.
Dans la remarque suivante, nous discutons plusieurs manières de modifier cette définition
pour une classe plus large de morphismes d'espaces algébriques formels.

\begin{remark}
\label{restricted-remark-rig-surjective-more-general}
La condition formulée dans la définition \ref{restricted-definition-rig-surjective}
ne convient pas même aux morphismes de type fini
d'espaces algébriques formels localement adiques*.
Par exemple, si $A = (\bigcup_{n \geq 1} k[t^{1/n}])^\wedge$,
où la complétion est la complétion $t$-adique, alors
il n'existe aucun morphisme adique $\text{Spf}(R) \to \text{Spf}(A)$
où $R$ soit un anneau de valuation discrète complet.
Ainsi, tout morphisme $X \to \text{Spf}(A)$ serait rig-surjectif ;
mais, puisque $A$ est intègre et que $t \in A$ est non nul, nous voulons
considérer que $A$ possède au moins un « point rigide », et nous ne voulons
pas autoriser $X = \emptyset$. Pour traiter ce
cas particulier, on peut considérer des morphismes adiques
$$
\text{Spf}(R) \longrightarrow Y
$$
où $R$ est un anneau de valuation complet relativement à un idéal
principal $J$ dont le radical est $\mathfrak m_R = \sqrt{J}$.
Dans ce cas, le groupe de valeurs de $R$ peut se plonger dans
$(\mathbf{R}, +)$, et l'on obtient le point de vue utilisé par
Berkovich pour définir un espace analytique associé à $Y$ ; voir
\cite{Berkovich}. Une autre approche est défendue par Huber. Dans sa théorie,
on abandonne l'hypothèse que $\Spec(R/J)$ est un singleton ; voir
\cite{Huber-continuous-valuations}.
\end{remark}

\begin{lemma}
\label{restricted-lemma-composition-rig-surjective}
\begin{slogan}
La rig-surjectivité des morphismes localement de type fini est préservée par
composition
\end{slogan}
Soit $S$ un schéma. Soient $f : X \to Y$ et $g : Y \to Z$ des morphismes d'espaces
algébriques formels sur $S$. Supposons $X$, $Y$ et $Z$ localement noethériens,
et $f$ et $g$ localement de type fini. Si $f$ et $g$ sont rig-surjectifs,
alors $g \circ f$ l'est aussi.
\end{lemma}

\begin{proof}
Cela résulte directement des définitions
(et d'Espaces formels, lemme \ref{formal-spaces-lemma-composition-finite-type}).
\end{proof}

\begin{lemma}
\label{restricted-lemma-base-change-rig-surjective}
Soit $S$ un schéma. Soient $f : X \to Y$ et $Z \to Y$ des morphismes
d'espaces algébriques formels sur $S$. Supposons $X$, $Y$ et $Z$ localement
noethériens, et $f$ et $g$ localement de type fini. Si $f$ est
rig-surjectif, alors le changement de base $Z \times_Y X \to Z$ l'est aussi.
\end{lemma}

\begin{proof}
Cela résulte directement des définitions (et d'
Espaces formels, lemmes \ref{formal-spaces-lemma-fibre-product-Noetherian} et
\ref{formal-spaces-lemma-base-change-finite-type}).
\end{proof}

\begin{lemma}
\label{restricted-lemma-rig-surjective-alternative-permanence}
Soit $S$ un schéma. Soient $f : X \to Y$ et $g : Y \to Z$
des morphismes localement de type fini d'espaces algébriques formels
localement noethériens sur $S$. Si $g \circ f$ est rig-surjectif
et si $g$ est un monomorphisme, alors $f$ est rig-surjectif.
\end{lemma}

\begin{proof}
Utiliser le lemme \ref{restricted-lemma-base-change-rig-surjective} et le fait que
$f$ est le changement de base de $g \circ f$ par $g$.
\end{proof}

\begin{lemma}
\label{restricted-lemma-permanence-rig-surjective}
Soit $S$ un schéma. Soient $f : X \to Y$ et $g : Y \to Z$ des morphismes d'espaces
algébriques formels sur $S$. Supposons $X$, $Y$ et $Z$ localement noethériens,
et $f$ et $g$ localement de type fini. Si $g \circ f : X \to Z$
est rig-surjectif, alors $g : Y \to Z$ l'est aussi.
\end{lemma}

\begin{proof}
Cela résulte immédiatement de la définition.
\end{proof}

\begin{lemma}
\label{restricted-lemma-etale-covering-rig-surjective}
Soit $S$ un schéma. Soit $f : X \to Y$ un morphisme d'espaces algébriques
formels localement noethériens, représentable par des espaces algébriques, étale
et surjectif. Alors $f$ est rig-surjectif.
\end{lemma}

\begin{proof}
Soit $p : \text{Spf}(R) \to Y$ un morphisme adique, où $R$ est un anneau
de valuation discrète complet. Soit $Z = \text{Spf}(R) \times_Y X$. Alors
$Z \to \text{Spf}(R)$ est représentable par des espaces algébriques, étale et
surjectif. Ainsi, $Z$ est non vide. Choisissons un espace algébrique formel affine
non vide $V$ et un morphisme étale $V \to Z$ (ce qui est possible par nos définitions).
Alors $V \to \text{Spf}(R)$ correspond à $R \to A^\wedge$, où
$R \to A$ est un homomorphisme étale d'anneaux ; voir Espaces formels, lemme
\ref{formal-spaces-lemma-etale}. Puisque $A^\wedge \not = 0$
(puisque $V \not = \emptyset$), on peut trouver un idéal maximal $\mathfrak m$
de $A$ au-dessus de $\mathfrak m_R$. Alors $A_\mathfrak m$ est un anneau
de valuation discrète (Compléments d'algèbre, lemme
\ref{more-algebra-lemma-Dedekind-etale-extension}).
Alors $R' = A_\mathfrak m^\wedge$ est un anneau de valuation discrète complet
(Compléments d'algèbre, lemme \ref{more-algebra-lemma-completion-dvr}).
En appliquant Espaces formels, lemme
\ref{formal-spaces-lemma-morphism-between-formal-spectra}.
on trouve le morphisme voulu $\text{Spf}(R') \to V \to Z \to X$.
\end{proof}

\noindent
Il résulte des lemmes précédents que l'on peut vérifier si
$f : X \to Y$ est rig-surjectif localement pour la topologie étale sur $Y$.

\begin{lemma}
\label{restricted-lemma-upshot}
Soit $S$ un schéma. Soit $f : X \to Y$ un morphisme d'espaces algébriques
formels localement noethériens qui est localement de type fini.
Soit $\{g_i : Y_i \to Y\}$ une famille de morphismes d'espaces algébriques
formels représentables par des espaces algébriques et
étales, telle que $\coprod g_i$ soit surjectif.
Alors $f$ est rig-surjectif si et seulement si chaque
$f_i : X \times_Y Y_i \to Y_i$ est rig-surjectif.
\end{lemma}

\begin{proof}
En effet, si $f$ est rig-surjectif, tout changement de base l'est aussi
(lemme \ref{restricted-lemma-base-change-rig-surjective}).
Réciproquement, si tous les $f_i$ sont rig-surjectifs, alors
$\coprod f_i : \coprod X \times_Y Y_i \to \coprod Y_i$ l'est aussi.
D'après le lemme \ref{restricted-lemma-etale-covering-rig-surjective},
le morphisme $\coprod g_i : \coprod Y_i \to Y$ est rig-surjectif.
Ainsi, $\coprod X \times_Y Y_i \to Y$ est rig-surjectif
(lemme \ref{restricted-lemma-composition-rig-surjective}).
Puisque ce morphisme se factorise par $X \to Y$, on voit que $X \to Y$
est rig-surjectif d'après le lemme \ref{restricted-lemma-permanence-rig-surjective}.
\end{proof}

\begin{lemma}
\label{restricted-lemma-faithfully-flat-rig-surjective}
Soit $A$ un anneau noethérien complet relativement à un idéal $I$.
Soit $B$ une algèbre $I$-adiquement complète sur $A$.
Si $A/I^n \to B/I^nB$ est de type fini et plat pour tout $n$, et
fidèlement plat pour $n = 1$, alors $\text{Spf}(B) \to \text{Spf}(A)$
est rig-surjectif.
\end{lemma}

\begin{proof}
Nous utiliserons sans autre mention le fait que les morphismes entre spectres formels
sont donnés par les homomorphismes continus entre les anneaux topologiques correspondants ; voir
Espaces formels, lemme \ref{formal-spaces-lemma-morphism-between-formal-spectra}.
Soit $\varphi : A \to R$ un homomorphisme continu vers un anneau de
valuation discrète complet $A$. Cela implique $\varphi(I) \subset \mathfrak m_R$.
D'autre part, puisque nous devons seulement produire le relèvement
$\varphi' : B' \to R'$ dans le cas où $\varphi$ correspond à un morphisme
adique, on peut supposer que $\varphi(I) \not = 0$. On peut donc considérer
le changement de base $C = B \widehat{\otimes}_A R$ ; voir par exemple la
remarque \ref{restricted-remark-base-change}.
Alors $C$ est une algèbre $\mathfrak m_R$-adiquement complète sur $R$
telle que $C/\mathfrak m_R^n C$ soit de type fini et plate sur
$R/\mathfrak m_R^n$, et que $C/\mathfrak m_R C$ soit non nulle.
Choisissons un idéal maximal $\mathfrak m \subset C$ au-dessus de
$\mathfrak m_R$. Par platitude (qui implique la propriété de descente), on voit que
$\Spec(C_\mathfrak m) \setminus V(\mathfrak m_R C_\mathfrak m)$
est un ouvert non vide. On peut donc choisir un idéal premier
$\mathfrak q \subset \mathfrak m$ tel que $\mathfrak q$ définisse un point fermé de
$\Spec(C_\mathfrak m) \setminus \{\mathfrak m\}$ et tel que
$\mathfrak q \not \in V(IC_\mathfrak m)$ ; voir
Propriétés, lemme \ref{properties-lemma-complement-closed-point-Jacobson}.
Alors $C/\mathfrak q$ est un anneau local intègre de dimension $1$, et l'on peut trouver
$C/\mathfrak q \subset R'$, où $R'$ est un anneau de valuation discrète
(Algèbre, lemme \ref{algebra-lemma-exists-dvr}).
Par construction, $\mathfrak m_R R' \subset \mathfrak m_{R'}$,
et l'on voit que $C \to R'$ se prolonge en un homomorphisme continu
$C \to (R')^\wedge$ (en fait, dans la situation présente, on peut choisir $R'$ tel que
$R' = (R')^\wedge$, mais cela ne nous est pas nécessaire).
Puisque le complété d'un anneau de valuation discrète est un anneau de
valuation discrète, on voit que l'hypothèse donne un diagramme commutatif
d'anneaux
$$
\xymatrix{
(R')^\wedge & C \ar[l] & B \ar[l] \\
R \ar[u] & R \ar[l] \ar[u] & A \ar[l] \ar[u]
}
$$
qui fournit le relèvement voulu.
\end{proof}

\begin{lemma}
\label{restricted-lemma-flat-rig-surjective}
Soit $A$ un anneau noethérien complet relativement à un idéal $I$.
Soit $B$ une algèbre $I$-adiquement complète sur $A$. Supposons que
\begin{enumerate}
\item la $I$-torsion de $A$ soit $0$ ;
\item $A/I^n \to B/I^nB$ soit plat et de type fini pour tout $n$.
\end{enumerate}
Alors $\text{Spf}(B) \to \text{Spf}(A)$ est rig-surjectif si et seulement
si $A/I \to B/IB$ est fidèlement plat.
\end{lemma}

\begin{proof}
La fidèle platitude implique la rig-surjectivité d'après le
lemme \ref{restricted-lemma-faithfully-flat-rig-surjective}.
Pour prouver la réciproque, nous utiliserons sans autre mention que
l'annulation de la $I$-torsion équivaut à celle de la torsion par une puissance de $I$
(Compléments d'algèbre, lemme \ref{more-algebra-lemma-torsion-free}).
Nous utiliserons également sans autre mention le fait que les morphismes entre
spectres formels sont donnés par des homomorphismes continus entre les anneaux
topologiques correspondants ; voir
Espaces formels, lemme \ref{formal-spaces-lemma-morphism-between-formal-spectra}.

\medskip\noindent
Supposons que $\text{Spf}(B) \to \text{Spf}(A)$ soit rig-surjectif.
Choisissons un idéal maximal $I \subset \mathfrak m \subset A$.
L'ouvert $U = \Spec(A_\mathfrak m) \setminus V(I_\mathfrak m)$
de $\Spec(A_\mathfrak m)$ est non vide, car la $I_\mathfrak m$-torsion de
$A_\mathfrak m$ est nulle
(utiliser Algèbre, lemme \ref{algebra-lemma-Noetherian-power-ideal-kills-module}).
On peut donc trouver un idéal premier $\mathfrak q \subset A_\mathfrak m$ qui définit
un point de $U$ (c'est-à-dire $IA_\mathfrak m \not \subset \mathfrak q$)
et qui correspond à un point fermé
de $\Spec(A_\mathfrak m) \setminus \{\mathfrak m\}$ ; voir
Propriétés, lemme \ref{properties-lemma-complement-closed-point-Jacobson}.
Alors $A_\mathfrak m/\mathfrak q$ est un anneau local intègre de dimension $1$.
On peut donc trouver un homomorphisme local injectif d'anneaux locaux
$A_\mathfrak m/\mathfrak q \subset R$, où $R$ est un anneau de valuation discrète
(Algèbre, lemme \ref{algebra-lemma-exists-dvr}).
Par construction, $IR \subset \mathfrak m_R$, et l'on voit que
$A \to R$ se prolonge en un homomorphisme continu $A \to R^\wedge$.
Puisque le complété d'un anneau de valuation discrète est un anneau de
valuation discrète, on voit que l'hypothèse donne un diagramme commutatif
d'anneaux
$$
\xymatrix{
R' & B \ar[l] \\
R^\wedge \ar[u] & A \ar[l] \ar[u]
}
$$
On trouve ainsi un idéal premier de $B$ au-dessus de $\mathfrak m$. Il s'ensuit
que $\Spec(B/IB) \to \Spec(A/I)$ est surjectif, d'où la fidèle platitude de
$A/I \to B/IB$
(Algèbre, lemme \ref{algebra-lemma-ff-rings}).
\end{proof}

\begin{lemma}
\label{restricted-lemma-monomorphism-rig-surjective}
Soit $S$ un schéma. Soit $f : X \to Y$ un morphisme d'espaces algébriques
formels. Supposons $X$ et $Y$ localement noethériens, $f$ localement de type
fini, et $f$ un monomorphisme. Alors $f$ est rig-surjectif si et seulement si
tout morphisme adique $\text{Spf}(R) \to Y$, où $R$ est un anneau de valuation
discrète complet, se factorise par $X$.
\end{lemma}

\begin{proof}
Un sens est immédiat. Pour l'autre, supposons que $\text{Spf}(R) \to Y$
soit un morphisme adique tel qu'il existe une extension d'anneaux de
valuation discrète complets $R \subset R'$ avec
$\text{Spf}(R') \to \text{Spf}(R) \to X$ se factorisant par $Y$. Alors
$\Spec(R'/\mathfrak m_R^n R') \to \Spec(R/\mathfrak m_R^n)$ est surjectif
et plat ; les morphismes $\Spec(R/\mathfrak m_R^n) \to X$ se factorisent donc
par $X$, puisque $X$ satisfait la condition de faisceau pour les recouvrements fpqc ; voir
Espaces formels, lemme \ref{formal-spaces-lemma-sheaf-fpqc}.
Autrement dit, $\text{Spf}(R) \to Y$ se factorise par $X$.
\end{proof}

\begin{lemma}
\label{restricted-lemma-closed-immersion-rig-surjective}
Soit $S$ un schéma. Soit $f : X \to Y$ un morphisme d'espaces algébriques
formels. Supposons que $X$ et $Y$ soient localement noethériens et que $f$ soit une
immersion fermée. Les conditions suivantes sont équivalentes :
\begin{enumerate}
\item $f$ est rig-surjectif ; et
\item pour tout espace algébrique formel affine $V$ et tout morphisme
$V \to Y$ représentable par des espaces algébriques et étale,
le morphisme $X \times_Y V \to V$ correspond à un morphisme surjectif
$B \to A$ dans $\textit{WAdm}^{Noeth}$ dont le noyau $J$ possède la propriété
suivante : $IJ^n = 0$ pour un idéal de définition $I$ de $B$
et un $n \geq 1$.
\end{enumerate}
\end{lemma}

\begin{proof}
Remarquons que, étant donnés $V$ et $V \to Y$ comme en (2), sans autre
hypothèse sur $f$, le morphisme $X \times_Y V \to V$
correspond à un morphisme surjectif $B \to A$ dans $\textit{WAdm}^{Noeth}$
d'après Espaces formels, lemme
\ref{formal-spaces-lemma-closed-immersion-into-countably-indexed}.

\medskip\noindent
Supposons (1). D'après le lemme \ref{restricted-lemma-base-change-rig-surjective}, on voit que
$\text{Spf}(A) \to \text{Spf}(B)$ est rig-surjectif.
Soit $I \subset B$ un idéal de définition. Puisque $B$ est adique,
$I^m \subset B$ est un idéal de définition pour tout $m \geq 1$.
Si $I^m J^n \not = 0$ pour tous $n, m \geq 1$, alors
$IJ$ n'est pas nilpotent, donc $V(IJ) \not = \Spec(B)$.
On peut donc trouver un idéal premier $\mathfrak p \subset B$
avec $\mathfrak p \not \in V(I) \cup V(J)$.
Remarquons que $I(B/\mathfrak p) \not = B/\mathfrak p$ ;
on peut donc trouver un idéal maximal
$\mathfrak p + I \subset \mathfrak m \subset B$.
D'après Algèbre, lemme \ref{algebra-lemma-exists-dvr},
on peut trouver un anneau de valuation discrète $R$
et un homomorphisme local injectif d'anneaux $(B/\mathfrak p)_\mathfrak m \to R$.
Il est clair que l'homomorphisme d'anneaux $B \to R$ ne peut se factoriser par $A = B/J$.
D'après le lemme \ref{restricted-lemma-monomorphism-rig-surjective},
cela contredit la rig-surjectivité de $\text{Spf}(A) \to \text{Spf}(B)$.
Ainsi, pour certains $n, m$, on a bien $I^n J^m = 0$,
ce qui montre (2).

\medskip\noindent
Supposons (2). D'après le lemme \ref{restricted-lemma-upshot}, il suffit de montrer
que $\text{Spf}(A) \to \text{Spf}(B)$ est rig-surjectif.
Choisissons un idéal de définition $I \subset B$ et un entier $n$
tels que $I J^n = 0$.
Considérons un homomorphisme d'anneaux $B \to R$, où $R$ est un anneau de
valuation discrète et où l'image de $I$ est non nulle. Puisque $R$ est intègre,
on en déduit que l'image de $J$ dans $R$ est nulle. Ainsi, $B \to R$
se factorise par la surjection $B \to A$, et l'on conclut par la
définition des morphismes rig-surjectifs.
\end{proof}

\begin{lemma}
\label{restricted-lemma-closed-immersion-rig-smooth-rig-surjective}
Soit $S$ un schéma. Soit $f : X \to Y$ un morphisme d'espaces algébriques
formels. Supposons que $X$ et $Y$ soient localement noethériens et que $f$ soit une
immersion fermée. Les conditions suivantes sont équivalentes :
\begin{enumerate}
\item $f$ est rig-lisse et rig-surjectif ;
\item $f$ est rig-étale et rig-surjectif ; et
\item pour tout espace algébrique formel affine $V$ et tout morphisme
$V \to Y$ représentable par des espaces algébriques et étale,
le morphisme $X \times_Y V \to V$ correspond à un morphisme surjectif
$B \to A$ dans $\textit{WAdm}^{Noeth}$ dont le noyau $J$ possède la propriété
suivante : $IJ = 0$ pour un idéal de définition $I$ de $B$.
\end{enumerate}
\end{lemma}

\begin{proof}
Soient $I$ et $J$ des idéaux d'un anneau $B$ tels que $IJ^n = 0$ et
$I(J/J^2) = 0$. Alors $I^nJ = 0$ (démonstration omise).
Le lemme résulte donc d'une combinaison immédiate des
lemmes \ref{restricted-lemma-closed-immersion-rig-smooth} et
\ref{restricted-lemma-closed-immersion-rig-surjective}.
\end{proof}

\begin{lemma}
\label{restricted-lemma-rig-etale-descent}
Soit $S$ un schéma. Soient $f : X \to Y$ et $g : Y \to Z$
des morphismes d'espaces algébriques formels localement noethériens sur $S$.
Supposons que :
\begin{enumerate}
\item $g$ soit localement de type fini ;
\item $f$ soit rig-lisse (respectivement rig-étale) et rig-surjectif ;
\item $g \circ f$ soit rig-lisse (respectivement rig-étale).
\end{enumerate}
Alors $g$ est rig-lisse (respectivement rig-étale).
\end{lemma}

\begin{proof}
Nous le démontrerons dans le cas rig-lisse et indiquerons à la fin
les modifications nécessaires pour traiter le cas rig-étale.
Considérons un diagramme commutatif
$$
\xymatrix{
X \times_Y V \ar[r] \ar[d] &
V \ar[d] \ar[r] &
W \ar[d] \\
X \ar[r] &
Y \ar[r] &
Z
}
$$
où $V$ et $W$ sont des espaces algébriques formels affines, où $V \to Y$ et $W \to Z$
sont représentables par des espaces algébriques et étales. Il faut montrer que
$V \to W$ correspond à un homomorphisme rig-lisse d'anneaux topologiques
adiques noethériens ; voir la définition \ref{restricted-definition-rig-smooth}.
On peut écrire $V = \text{Spf}(B)$ et $W = \text{Spf}(C)$, et dire que
$V \to W$ correspond à un homomorphisme adique d'anneaux $C \to B$ qui est
topologiquement de type fini ; voir le lemme \ref{restricted-lemma-finite-type-morphisms}.

\medskip\noindent
Nous utiliserons ci-dessous sans autre mention que $X \times_Y V \to V$
est rig-lisse et rig-surjectif ; voir les
lemmes \ref{restricted-lemma-base-change-rig-smooth} et
\ref{restricted-lemma-base-change-rig-surjective}.
En outre, la composée $X \times_Y V \to V \to W$ est rig-lisse
puisque $g \circ f$ est rig-lisse.

\medskip\noindent
Soit $I \subset C$ un idéal de définition. Le module
Supposons, pour obtenir une contradiction, que $C \to B$ ne soit pas rig-lisse.
Cela signifie qu'il existe un idéal premier $\mathfrak q \subset B$
ne contenant pas $IB$ tel que soit $H^{-1}(\NL_{B/C}^\wedge)_\mathfrak p$
est non nul, soit $H^0(\NL_{B/C}^\wedge)_\mathfrak p$ n'est pas un
$B_\mathfrak q$-module libre de type fini. Voir le
lemme \ref{restricted-lemma-equivalent-with-artin-smooth} ; certains détails sont
omis. On peut choisir un idéal maximal
$IB + \mathfrak q \subset \mathfrak m$. D'après
Algèbre, lemme \ref{algebra-lemma-exists-dvr},
on peut trouver un anneau de valuation discrète complet $R$ et un
homomorphisme local injectif d'anneaux $(B/\mathfrak q)_\mathfrak m \to R$.

\medskip\noindent
Après avoir remplacé $R$ par une extension, on peut supposer donné un relèvement
$\text{Spf}(R) \to X \times_Y V$ du morphisme adique
$\text{Spf}(R) \to V = \text{Spf}(B)$. Choisissons un recouvrement étale
$\{\text{Spf}(A_i) \to X \times_Y V\}$ comme dans
Espaces formels, définition
\ref{formal-spaces-definition-formal-algebraic-space}.
D'après le lemme \ref{restricted-lemma-etale-covering-rig-surjective},
on peut supposer que $\text{Spf}(R) \to X \times_Y V$ se relève en un
morphisme $\text{Spf}(R) \to \text{Spf}(A_i)$ pour un certain $i$
(cela peut nécessiter de remplacer $R$ par une autre extension).
Posons $A = A_i$. Considérons les homomorphismes d'anneaux
$$
C \to B \to A \to R
$$
Soit $\mathfrak p \subset A$ le noyau de l'homomorphisme
$A \to R$, et remarquons que $\mathfrak p$ est au-dessus de $\mathfrak q$.
Nous savons que $C \to A$ et $B \to A$ sont rig-lisses.
En particulier, l'homomorphisme d'anneaux $B_\mathfrak q \to A_\mathfrak p$
est plat d'après le lemme \ref{restricted-lemma-rig-smooth-rig-flat}.
Considérons la suite exacte associée
$$
\xymatrix{
&
H^0(\NL_{B/C}^\wedge) \otimes_B A_\mathfrak p \ar[r] &
H^0(\NL_{A/C}^\wedge)_\mathfrak p \ar[r] &
H^0(\NL_{A/B}^\wedge)_\mathfrak p \ar[r] & 0 \\
0 \ar[r] &
H^{-1}(\NL_{B/C}^\wedge \otimes_B A)_\mathfrak p \ar[r] &
H^{-1}(\NL_{A/C}^\wedge)_\mathfrak p \ar[r] &
H^{-1}(\NL_{A/B}^\wedge)_\mathfrak p \ar[llu]
}
$$
des lemmes \ref{restricted-lemma-exact-sequence-NL} et
\ref{restricted-lemma-exact-sequence-NL-rig-smooth}.
La rig-lissité de $C \to A$ et de $B \to A$ permet de conclure que
$H^{-1}(\NL_{B/C}^\wedge \otimes_B A)_\mathfrak p = 0$
et que $H^0(\NL_{B/C}^\wedge) \otimes_B A_\mathfrak p$
est libre de type fini, comme noyau d'une surjection entre
$A_\mathfrak p$-modules libres de type fini. Puisque $B_\mathfrak q \to A_\mathfrak p$
est plat, donc fidèlement plat, cela implique que
$H^{-1}(\NL_{B/C}^\wedge)_\mathfrak q = 0$
et que $H^0(\NL_{B/C}^\wedge)_\mathfrak q$
est libre de type fini, ce qui est la contradiction recherchée.

\medskip\noindent
Dans le cas rig-étale, on raisonne exactement de la même manière,
mais la conclusion obtenue est que
$H^{-1}(\NL_{B/C}^\wedge)_\mathfrak q$
et $H^0(\NL_{B/C}^\wedge)_\mathfrak q$ sont tous deux nuls.
\end{proof}







\section{Espaces algébriques formels sur des anneaux de valuation discrète complets}
\label{restricted-section-over-cdrv}

\noindent
Dans cette section, nous emploierons la terminologie suivante :
si $A$ est un anneau topologique faiblement admissible, nous
dirons que « {\it $X$ est un espace algébrique formel sur $A$} » lorsque $X$
est un espace algébrique formel muni d'un
morphisme $p : X \to \text{Spf}(A)$ d'espaces algébriques formels.
Dans cette situation, nous appellerons $p$ le {\it morphisme structural}.

\begin{lemma}
\label{restricted-lemma-flat-locus}
Soit $X$ un espace algébrique formel localement noethérien sur
un anneau de valuation discrète complet $A$.
Il existe alors une immersion fermée $X' \to X$
d'espaces algébriques formels telle que $X'$ soit plat sur $A$
et que tout morphisme $Y \to X$ d'espaces algébriques formels localement
noethériens, avec $Y$ plat sur $A$, se factorise par $X'$.
\end{lemma}

\begin{proof}
Soit $\pi \in A$ une uniformisante. Rappelons qu'un $A$-module
est plat si et seulement si sa torsion annulée par une puissance de $\pi$ est $0$.

\medskip\noindent
Supposons d'abord que $X$ soit un espace algébrique formel affine.
Alors $X = \text{Spf}(B)$, où $B$ est une $A$-algèbre noethérienne adique.
Dans ce cas, posons $X' = \text{Spf}(B')$, où
$B' = B/\pi\text{-power torsion}$. Il est clair que $X'$ est plat
sur $A$ et que $X' \to X$ est une immersion fermée.
Soit $g : Y \to X$ un morphisme d'espaces algébriques formels localement
noethériens, avec $Y$ plat sur $A$. Choisissons un recouvrement
$\{Y_j \to Y\}$ comme dans Espaces formels, définition
\ref{formal-spaces-definition-formal-algebraic-space}.
Alors $Y_j = \text{Spf}(C_j)$ avec $C_j$ plat sur $A$.
Par conséquent, le morphisme $Y_j \to X$, qui correspond à un homomorphisme
continu de $R$-algèbres $B \to C_j$, se factorise par $X'$, car
$B \to C_j$ annule manifestement la torsion tuée par une puissance de $\pi$.
Comme $\{Y_j \to Y\}$ est un recouvrement et que $X' \to X$
est un monomorphisme, nous concluons que $g$ se factorise par $X'$.

\medskip\noindent
Soient $X$ et $\{X_i \to X\}_{i \in I}$ comme dans
Espaces formels, définition
\ref{formal-spaces-definition-formal-algebraic-space}.
Pour tout $i$, soit $X'_i \to X_i$ la partie plate construite
ci-dessus. Pour $i, j \in I$, la projection
$X'_i \times_X X_j \to X'_i$ est un morphisme étale (par hypothèse)
de schémas (d'après Espaces formels, lemme
\ref{formal-spaces-lemma-presentation-representable}).
Ainsi, $X'_i \times_X X_j$ est plat sur $A$, puisque les morphismes
étales représentables par des espaces algébriques sont plats
(lemme \ref{restricted-lemma-representable-flat}).
La projection $X'_i \times_X X_j \to X_j$ se factorise donc
par $X'_j$ en vertu de la propriété universelle. Nous en concluons que
$$
R_{ij} = X'_i \times_X X_j = X'_i \times_X X'_j = X_i \times_X X'_j
$$
car les morphismes $X'_i \to X_i$ sont des injections de faisceaux.
Posons $U = \coprod X'_i$ et
$R = \coprod R_{ij}$, et notons $s, t : R \to U$ les deux
projections. En tant que faisceau, $R = U \times_X U$, et $s$ et $t$
sont étales. D'après les observations précédentes, $(t, s) : R \to U$
définit alors une relation d'équivalence étale. Ainsi, $X' = U/R$ est un
espace algébrique d'après Espaces, théorème \ref{spaces-theorem-presentation}.
Par construction, le diagramme
$$
\xymatrix{
\coprod X'_i \ar[r] \ar[d] & \coprod X_i \ar[d] \\
X' \ar[r] & X
}
$$
est cartésien. Comme la flèche verticale de droite est étale surjective
et que la flèche horizontale supérieure est représentable et est une immersion
fermée, nous concluons que $X' \to X$ est représentable d'après
Amorçage, lemme \ref{bootstrap-lemma-after-fppf-sep-lqf}.
Nous pouvons alors appliquer Espaces, lemme
\ref{spaces-lemma-descent-representable-transformations-property}
pour conclure que $X' \to X$ est une immersion fermée.

\medskip\noindent
Enfin, supposons que $Y \to X$ soit un morphisme, où
$Y$ est un espace algébrique formel localement noethérien plat sur $A$.
Alors chaque $X_i \times_X Y$ est étale sur $Y$ et
donc plat sur $A$ (voir ci-dessus).
Le morphisme $X_i \times_X Y \to X_i$ se factorise donc par $X'_i$.
Par conséquent, $Y \to X$ se factorise par $X'$, puisque
$\{X_i \times_X Y \to Y\}$ est un recouvrement étale.
\end{proof}

\begin{lemma}
\label{restricted-lemma-flat-and-diagonal-rig-surjective}
Soit $X$ un espace algébrique formel localement noethérien qui est
localement de type fini sur un anneau de valuation discrète complet $A$.
Soit $X' \subset X$ comme dans le lemme \ref{restricted-lemma-flat-locus}.
Si $X \to X \times_{\text{Spf}(A)} X$ est rig-étale et rig-surjectif,
alors $X' = \text{Spf}(A)$ ou $X' = \emptyset$.
\end{lemma}

\begin{proof}
(Parenthèse : la diagonale est toujours localement de type fini d'après
Espaces formels, lemme
\ref{formal-spaces-lemma-diagonal-morphism-formal-algebraic-spaces}
et $X \times_{\text{Spf}(A)} X$ est localement noethérien d'après
Espaces formels, lemmes
\ref{formal-spaces-lemma-base-change-finite-type} et
\ref{formal-spaces-lemma-locally-finite-type-locally-noetherian}.
Il est donc légitime d'imposer ces conditions au morphisme diagonal.)
Le diagramme
$$
\xymatrix{
X' \ar[r] \ar[d] & X' \times_{\text{Spf}(A)} X' \ar[d] \\
X \ar[r] & X \times_{\text{Spf}(A)} X
}
$$
est cartésien. Par conséquent, $X' \to X' \times_{\text{Spf}(A)} X'$
est rig-étale et rig-surjectif d'après le
lemme \ref{restricted-lemma-base-change-rig-surjective}.
Choisissons un espace algébrique formel affine $U$ et un morphisme
$U \to X'$ qui soit étale et représentable par des espaces algébriques.
Alors $U = \text{Spf}(B)$, où $B$ est un anneau topologique noethérien adique
qui est une $A$-algèbre plate, dont la topologie est la topologie
$\pi$-adique, où $\pi \in A$ est une uniformisante, et tel que
$A/\pi^n A \to B/\pi^n B$ soit de type fini pour tout $n$.
Pour usage ultérieur, remarquons que cela implique en particulier que, si
$B \not = 0$, alors l'application $\text{Spf}(B) \to \text{Spf}(A)$ est une
surjection de faisceaux (rappelons que nous employons, comme toujours,
la topologie fppf). En répétant l'argument précédent, nous voyons que
$$
W = U \times_{X'} U =
X' \times_{X' \times_{\text{Spf}(A)} X'} (U \times_{\text{Spf}(A)} U)
\longrightarrow
U \times_{\text{Spf}(A)} U
$$
est une immersion fermée, rig-étale et rig-surjective. Nous avons
$U \times_{\text{Spf}(A)} U = \text{Spf}(B \widehat{\otimes}_A B)$
d'après Espaces formels, lemme
\ref{formal-spaces-lemma-fibre-product-affines-over-separated}.
Alors $B \widehat{\otimes}_A B$ est une $A$-algèbre plate,
car c'est le complété $\pi$-adique de la $A$-algèbre plate $B \otimes_A B$.
Ainsi, $W = U \times_{\text{Spf}(A)} U$ d'après le
lemme \ref{restricted-lemma-closed-immersion-rig-smooth-rig-surjective}.
Autrement dit, nous avons $U \times_{X'} U = U \times_{\text{Spf}(A)} U$,
ce qui signifie à son tour que l'image de $U \to X'$ (comme application
de faisceaux) s'envoie injectivement dans $\text{Spf}(A)$.
Choisissons un recouvrement $\{U_i \to X'\}$ comme dans Espaces formels, définition
\ref{formal-spaces-definition-formal-algebraic-space}.
En particulier, $\coprod U_i \to X'$ est une surjection de faisceaux.
En appliquant ce qui précède à $U_i \coprod U_j \to X'$ (et en utilisant
le fait que $U_i \amalg U_j$ est lui aussi un espace algébrique formel
affine), nous voyons que $X' \to \text{Spf}(A)$ est une injection
de faisceaux fppf. Comme $X'$ est plat sur $A$, soit $X'$ est vide
(si $U_i$ est vide pour tout $i$), soit l'application est un isomorphisme
(si $U_i$ est non vide pour un certain $i$, auquel cas nous avons vu que
$U_i \to \text{Spf}(A)$ est une surjection de faisceaux),
ce qui achève la démonstration.
\end{proof}

\begin{lemma}
\label{restricted-lemma-rig-monomorphism-rig-surjective}
Soit $S$ un schéma. Soit $f : X \to Y$ un morphisme d'espaces algébriques
formels. Supposons que
\begin{enumerate}
\item $X$ et $Y$ soient localement noethériens,
\item $f$ soit localement de type fini,
\item $\Delta_f : X \to X \times_Y X$ soit rig-étale et rig-surjectif.
\end{enumerate}
Alors $f$ est rig-surjectif si et seulement si tout morphisme adique
$\text{Spf}(R) \to Y$, où $R$ est un anneau de valuation
discrète complet, se relève en un morphisme $\text{Spf}(R) \to X$.
\end{lemma}

\begin{proof}
Une implication est triviale. Pour la réciproque, supposons que
$\text{Spf}(R) \to Y$ soit un morphisme adique tel qu'il existe une extension
d'anneaux de valuation discrète complets $R \subset R'$ pour laquelle
$\text{Spf}(R') \to \text{Spf}(R) \to X$ se factorise par $Y$.
Considérons le diagramme de produit fibré
$$
\xymatrix{
\text{Spf}(R') \ar[r] \ar[rd] &
\text{Spf}(R) \times_Y X \ar[r] \ar[d]^p &
X \ar[d]^f \\
&
\text{Spf}(R) \ar[r] &
Y
}
$$
Le morphisme $p$ est localement de type fini en tant que changement de base de
$f$, voir Espaces formels, lemme
\ref{formal-spaces-lemma-base-change-finite-type}.
Le morphisme diagonal $\Delta_p$ est le changement de base de
$\Delta_f$ ; il est donc rig-étale et rig-surjectif.
D'après le lemme \ref{restricted-lemma-flat-and-diagonal-rig-surjective},
le lieu plat de $\text{Spf}(R) \times_Y X$ sur $R$
est soit $\emptyset$, soit égal à $\text{Spf}(R)$.
Toutefois, puisque $\text{Spf}(R')$ se factorise par ce lieu,
nous concluons que celui-ci n'est pas vide et obtenons donc
un morphisme $\text{Spf}(R) \to \text{Spf}(R) \times_Y X \to X$,
comme souhaité.
\end{proof}










\section{Le foncteur de complétion}
\label{restricted-section-completion-functor}

\noindent
Dans cette section, nous considérons la situation suivante.
Fixons d'abord un schéma de base $S$. Tous les anneaux, anneaux topologiques,
schémas, espaces algébriques et espaces algébriques formels, ainsi que
leurs morphismes, seront sur $S$. Fixons ensuite un
espace algébrique $X$ et une partie fermée $T \subset |X|$.
Notons $U \subset X$ le sous-espace ouvert tel que $|U| = |X| \setminus T$.
On peut représenter la situation ainsi :
$$
U \to X \quad |X| = |U| \amalg T
$$
Dans cette situation, étant donné un espace algébrique $X'$ sur $X$, c'est-à-dire
un espace algébrique $X'$ muni d'un morphisme $f : X' \to X$, nous
notons $T' \subset |X'|$ l'image réciproque de $T$ et désignons par
$U' \subset X'$ le sous-espace ouvert tel que $|U'| = |X'| \setminus T'$.
On peut représenter la situation ainsi :
$$
U' = f^{-1}U
\quad\quad
\vcenter{
\xymatrix{
U' \ar[d] \ar[r] & X' \ar[d]_f \\
U \ar[r] & X
}
}
\quad\quad
\vcenter{
\xymatrix{
|U'| \ar[r] \ar[d] & |X'| \ar[d]^{|f|} & T' \ar[l] \ar[d] \\
|U| \ar[r] & |X| & T \ar[l]
}
}
\quad\quad
T' = |f|^{-1}T
$$
Nous relierons les propriétés de $f$ à celles du morphisme induit
$$
f_{/T} : X'_{/T'} \longrightarrow X_{/T}
$$
de complétés formels. Comme l'indique la formule affichée, nous
noterons ce morphisme $f_{/T}$. Nous avons déjà vu que
$f_{/T}$ est représentable par des espaces algébriques dans
Espaces formels, lemme
\ref{formal-spaces-lemma-map-completions-representable}.
En fait, comme le montre la démonstration de ce lemme, le diagramme
$$
\xymatrix{
X'_{/T'} \ar[d]_{f_{/T}} \ar[r] & X' \ar[d]^f \\
X_{/T} \ar[r] & X
}
$$
est cartésien. Il convient de garder ce fait à l'esprit en lisant les lemmes
énoncés et démontrés ci-dessous.

\begin{lemma}
\label{restricted-lemma-map-completions-finite-type}
Dans la situation ci-dessus, si $f$ est localement de type fini, alors
$f_{/T}$ est localement de type fini.
\end{lemma}

\begin{proof}
(Les morphismes de type fini d'espaces algébriques formels sont étudiés dans
Espaces formels, section \ref{formal-spaces-section-finite-type}.)
En effet, supposons que $Z \to X$ soit un morphisme d'un schéma vers $X$ tel
que $|Z|$ s'envoie dans $T$. Le carré cartésien ci-dessus montre que
$Z \times_X X'$ est un espace algébrique qui représente
$Z \times_{X_{/T}} X'_{/T'}$. Comme $Z \times_X X' \to Z$
est localement de type fini d'après Morphismes d'espaces, lemme
\ref{spaces-morphisms-lemma-base-change-finite-type}, cela permet de conclure.
\end{proof}

\begin{lemma}
\label{restricted-lemma-map-completions-etale}
Dans la situation ci-dessus, si $f$ est étale, alors $f_{/T}$ est étale.
\end{lemma}

\begin{proof}
Le même argument que dans la démonstration du
lemme \ref{restricted-lemma-map-completions-finite-type} ramène l'assertion à
Morphismes d'espaces, lemme \ref{spaces-morphisms-lemma-base-change-etale}.
\end{proof}

\begin{lemma}
\label{restricted-lemma-closed-immersion-gives-closed-immersion}
Dans la situation ci-dessus, si $f$ est une immersion fermée, alors
$f_{/T}$ est une immersion fermée.
\end{lemma}

\begin{proof}
(Les immersions fermées d'espaces algébriques formels sont étudiées dans
Espaces formels, section
\ref{formal-spaces-section-closed-immersions}.)
Le même argument que dans la démonstration du
lemme \ref{restricted-lemma-map-completions-finite-type} ramène l'assertion à
Espaces, lemme
\ref{spaces-lemma-base-change-immersions}.
\end{proof}

\begin{lemma}
\label{restricted-lemma-proper-gives-proper}
Dans la situation ci-dessus, si $f$ est propre, alors $f_{/T}$ est propre.
\end{lemma}

\begin{proof}
(Les morphismes propres d'espaces algébriques formels sont étudiés dans
Espaces formels, section \ref{formal-spaces-section-proper}.)
Le même argument que dans la démonstration du
lemme \ref{restricted-lemma-map-completions-finite-type} ramène l'assertion à
Morphismes d'espaces, lemme
\ref{spaces-morphisms-lemma-base-change-proper}.
\end{proof}

\begin{lemma}
\label{restricted-lemma-quasi-compact-gives-quasi-compact}
Dans la situation ci-dessus, si $f$ est quasi-compact, alors
$f_{/T}$ est quasi-compact.
\end{lemma}

\begin{proof}
(Les morphismes quasi-compacts d'espaces algébriques formels sont étudiés dans
Espaces formels, section \ref{formal-spaces-section-quasi-compact}.)
Il faut montrer que $(X'_{/T'})_{red} \to (X_{/T})_{red}$ est un morphisme
quasi-compact d'espaces algébriques. D'après
Espaces formels, lemme \ref{formal-spaces-lemma-reduction-completion},
il s'agit du morphisme $Z' \to Z$, où $Z' \subset X'$, resp.\ $Z \subset X$,
est la structure d'espace algébrique réduit induite sur $T'$, resp.\ $T$.
Il s'ensuit que $Z' \to f^{-1}Z = Z \times_X X'$ est un épaississement
(une immersion fermée induisant un isomorphisme sur les espaces topologiques
sous-jacents). Comme $Z \times_X X' \to Z$ est quasi-compact en tant que
changement de base de $f$ (Morphismes d'espaces, lemme
\ref{spaces-morphisms-lemma-base-change-quasi-compact})
nous concluons que $Z' \to Z$ l'est aussi d'après
Compléments sur les morphismes d'espaces, lemme
\ref{spaces-more-morphisms-lemma-thicken-property-morphisms}.
\end{proof}

\begin{remark}
\label{restricted-remark-diagonal-gives-diagonal}
Dans la situation ci-dessus, considérons les morphismes diagonaux
$\Delta_f : X' \to X' \times_X X'$ et
$\Delta_{f_{/T}} : X'_{/T'} \to X'_{/T'} \times_{X_{/T}} X'_{/T'}$.
On voit facilement que
$$
X'_{/T'} \times_{X_{/T}} X'_{/T'} = (X' \times_X X')_{/T''}
$$
comme sous-foncteurs de $X' \times_X X'$, où $T'' \subset |X' \times_X X'|$
est l'image réciproque de $T$. Nous voyons donc que
$\Delta_{f_{/T}} = (\Delta_f)_{/T''}$. Nous utiliserons ce fait ci-dessous
pour montrer que les propriétés de $\Delta_f$ sont héritées par $\Delta_{f_{/T}}$.
\end{remark}

\begin{lemma}
\label{restricted-lemma-quasi-separated-gives-quasi-separated}
Dans la situation ci-dessus, si $f$ est (quasi-)séparé, alors
$f_{/T}$ l'est aussi.
\end{lemma}

\begin{proof}
(Les conditions de séparation des morphismes d'espaces algébriques formels
sont étudiées dans Espaces formels, section
\ref{formal-spaces-section-separation-axioms}.)
Il faut montrer que si $\Delta_f$ est quasi-compact, resp.\ une immersion
fermée, alors il en va de même pour $\Delta_{f_{/T}}$. Cela résulte de la
discussion de la remarque \ref{restricted-remark-diagonal-gives-diagonal} et des
lemmes \ref{restricted-lemma-quasi-compact-gives-quasi-compact} et
\ref{restricted-lemma-closed-immersion-gives-closed-immersion}.
\end{proof}

\begin{lemma}
\label{restricted-lemma-smooth-gives-rig-smooth}
Dans la situation ci-dessus, si $X$ est localement noethérien,
si $f$ est localement de type fini et si $U' \to U$ est lisse, alors
$f_{/T}$ est rig-lisse.
\end{lemma}

\begin{proof}
La stratégie de la démonstration est la suivante : se ramener au cas où $X$ et
$X'$ sont affines, traduire le cas affine en algèbre, puis appliquer le
lemme \ref{restricted-lemma-rig-smooth}. Nous invitons le lecteur à omettre les détails.

\medskip\noindent
Choisissons un morphisme étale surjectif $W \to X$, où $W = \coprod W_i$
est une réunion disjointe de schémas affines, voir Propriétés des espaces, lemme
\ref{spaces-properties-lemma-cover-by-union-affines}.
Pour chaque $i$, choisissons un morphisme étale surjectif
$W'_i \to W_i \times_X X'$, où $W'_i = \coprod W'_{ij}$ est une réunion
disjointe de schémas affines. En particulier,
$\coprod W'_{ij} \to X'$ est surjectif et étale.
Notons $f_{ij} : W_{ij} \to W_i$ le morphisme donné.
Notons $T_i \subset W_i$ et $T'_{ij} \subset W_{ij}$ les images
réciproques de $T$. Comme la complétion le long de l'image réciproque de
$T$ produit des diagrammes cartésiens (voir ci-dessus), nous avons
$(W_i)_{/T_i} = W_i \times_X X_{/T}$ et, de même,
$(W'_{ij})_{/T'_{ij}} = W'_{ij} \times_{X'} X'_{/T'}$.
Rappelons en outre que $(W_i)_{/T_i}$ et $(W'_{ij})_{/T'_{ij}}$
sont des espaces algébriques formels affines.
Ainsi, $\{W'_{ij})_{/T'_{ij}} \to X'_{/T'}\}$ est un recouvrement
comme dans Espaces formels, définition
\ref{formal-spaces-definition-formal-algebraic-space}.
D'après le lemme \ref{restricted-lemma-rig-smooth-morphisms},
il suffit de démontrer que
$$
(W'_{ij})_{/T'_{ij}} \longrightarrow (W_i)_{/T_i}
$$
est rig-lisse. Observons que $W'_{ij} \to W_i$
est localement de type fini et induit un morphisme lisse
$W'_{ij} \setminus T'_{ij} \to W_i \setminus T_i$
(car c'est vrai pour $f$ et que ces propriétés des morphismes sont
locales pour la topologie étale sur la source et le but).
Observons que $W_i$ est localement noethérien (car $X$ est localement
noethérien et cette propriété est locale pour la topologie étale sur
l'espace algébrique). Il suffit donc de démontrer le lemme lorsque
$X$ et $X'$ sont des schémas affines.

\medskip\noindent
Supposons que $X = \Spec(A)$ et $X' = \Spec(A')$ soient des schémas affines.
Comme $X$ est noethérien, nous voyons que $A$ est noethérien.
Le morphisme $f$ est donné par un homomorphisme d'anneaux
$A \to A'$ de type fini. Soit $I \subset A$ un idéal définissant $T$.
Alors $IA'$ définit $T'$. De plus, $\Spec(A') \to \Spec(A)$ est lisse sur
$\Spec(A) \setminus T$. Soient $A^\wedge$ et $(A')^\wedge$ les complétés
$I$-adiques. Nous avons $X_{/T} = \text{Spf}(A^\wedge)$ et
$X'_{/T'} = \text{Spf}((A')^\wedge)$, voir la démonstration d'
Espaces formels, lemme
\ref{formal-spaces-lemma-formal-completion-types}.
D'après le lemme \ref{restricted-lemma-rig-smooth}, nous voyons que
$(A')^\wedge$ est rig-lisse sur $(A. I)$,
ce qui signifie à son tour que $A^\wedge \to (A')^\wedge$ est rig-lisse ;
cela implique enfin que $X'_{/T'} \to X_{/T}$ est rig-lisse d'après le
lemme \ref{restricted-lemma-rig-smooth-morphisms}.
\end{proof}

\begin{lemma}
\label{restricted-lemma-etale-gives-rig-etale}
Dans la situation ci-dessus, si $X$ est localement noethérien,
si $f$ est localement de type fini et si $U' \to U$ est étale, alors
$f_{/T}$ est rig-étale.
\end{lemma}

\begin{proof}
La démonstration est exactement la même que celle du
lemme \ref{restricted-lemma-smooth-gives-rig-smooth}, à ceci près que les
lemmes \ref{restricted-lemma-rig-smooth} et \ref{restricted-lemma-rig-smooth-morphisms}
sont remplacés par les
lemmes \ref{restricted-lemma-rig-etale} et \ref{restricted-lemma-rig-etale-morphisms}.
\end{proof}

\begin{lemma}
\label{restricted-lemma-completion-proper-surjective-rig-surjective}
Dans la situation ci-dessus, si $X$ est localement noethérien,
si $f$ est propre et si $U' \to U$ est surjectif, alors $f_{/T}$ est rig-surjectif.
\end{lemma}

\begin{proof}
(L'énoncé a un sens d'après le
lemme \ref{restricted-lemma-map-completions-finite-type} et
Espaces formels, lemme \ref{formal-spaces-lemma-formal-completion-types}.)
Soit $R$ un anneau de valuation discrète complet de corps des fractions $K$.
Soit $p : \text{Spf}(R) \to X_{/T}$ un morphisme adique d'espaces
algébriques formels. D'après Espaces formels, lemme
\ref{formal-spaces-lemma-adic-into-completion}
le composé $\text{Spf}(R) \to X_{/T} \to X$
correspond à un morphisme $q : \Spec(R) \to X$
qui envoie $\Spec(K)$ dans $U$. Comme $U' \to U$ est propre et surjectif,
nous voyons que $\Spec(K) \times_U U'$ est non vide et propre sur $K$.
Nous pouvons donc choisir une extension de corps $K'/K$ et un diagramme
commutatif
$$
\xymatrix{
\Spec(K') \ar[r] \ar[d] & U' \ar[r] \ar[d] & X' \ar[d] \\
\Spec(K) \ar[r] & U \ar[r] & X
}
$$
Soit $R' \subset K'$ un anneau de valuation discrète dominant $R$ et de
corps des fractions $K'$, voir Algèbre, lemme \ref{algebra-lemma-exists-dvr}.
Comme $\Spec(K) \to X$ se prolonge en $\Spec(R) \to X$, le critère valuatif
de propreté (Morphismes d'espaces, lemme
\ref{spaces-morphisms-lemma-characterize-proper}) montre que nous pouvons
prolonger notre $K'$-point de $U'$ en un morphisme
$\Spec(R') \to X'$ au-dessus de $\Spec(R) \to X$.
Il s'ensuit que l'image réciproque de $T'$ dans $\Spec(R')$ est le
point fermé, et nous obtenons un morphisme adique
$\text{Spf}((R')^\wedge) \to X'_{/T'}$ relevant $p$,
comme souhaité (notons que $(R')^\wedge$ est un anneau de valuation discrète
complet d'après Compléments d'algèbre, lemme
\ref{more-algebra-lemma-completion-dvr}).
\end{proof}

\begin{lemma}
\label{restricted-lemma-separated-mono-open-diagonal-rig-surjective}
Dans la situation ci-dessus, si $X$ est localement noethérien,
si $f$ est séparé et localement de type fini et si $U' \to U$ est
un monomorphisme, alors $\Delta_{f_{/T}}$ est rig-surjectif.
\end{lemma}

\begin{proof}
La diagonale $\Delta_f : X' \to X' \times_X X'$ est une immersion
fermée et la restriction $U' \to U' \times_U U'$ de $\Delta_f$
est surjective. Le lemme résulte donc de la discussion de la
remarque \ref{restricted-remark-diagonal-gives-diagonal} et du
lemme \ref{restricted-lemma-completion-proper-surjective-rig-surjective}.
\end{proof}












\section{Modifications formelles}
\label{restricted-section-formal-modifications}

\noindent
Dans cette section, nous définissons et étudions la notion de modification
formelle d'Artin pour les espaces algébriques formels localement noethériens.
Commençons par la définition.

\begin{definition}
\label{restricted-definition-formal-modification}
Soit $S$ un schéma. Soit $f : X \to Y$ un morphisme d'espaces
algébriques formels localement noethériens sur $S$. Nous disons que $f$
est une {\it modification formelle} si
\begin{enumerate}
\item $f$ est un morphisme propre (Espaces formels, définition
\ref{formal-spaces-definition-proper}),
\item $f$ est rig-étale,
\item $f$ est rig-surjectif,
\item $\Delta_f : X \to X \times_Y X$ est rig-surjectif.
\end{enumerate}
\end{definition}

\noindent
Un exemple typique est donné au
lemme \ref{restricted-lemma-modification-gives-formal-modification},
et nous montrerons en fait plus loin que toute modification formelle est
« formellement locale » de ce type, voir le
lemme \ref{restricted-lemma-formal-modifications-locally-algebraic}.
Comparons ces conditions à celles de l'article d'Artin.

\begin{remark}
\label{restricted-remark-compare-formal-modification-artin}
Dans \cite[définition 1.7]{ArtinII}, une modification formelle est définie
comme un morphisme propre $f : X \to Y$ d'espaces algébriques formels
localement noethériens satisfaisant aux trois conditions suivantes\footnote{Nous
ne traduirons pas entièrement ces conditions dans le langage développé dans
le projet Champs. Nous espérons néanmoins que cette discussion sera utile au
lecteur.}
\begin{enumerate}
\item[(\romannumeral1)] les idéaux de Cramer et jacobien de
$f$ contiennent chacun un idéal de définition de $X$,
\item[(\romannumeral2)] l'idéal définissant l'application
diagonale $\Delta : X \to X \times_Y X$
est annulé par un idéal de définition de $X \times_Y X$, et
\item[(\romannumeral3)] tout morphisme adique $\text{Spf}(R) \to Y$
se relève en $\text{Spf}(R) \to X$ dès que $R$ est un
anneau de valuation discrète complet.
\end{enumerate}
Comparons ces conditions à notre liste ci-dessus.

\medskip\noindent
Pour (\romannumeral1). La propriété (\romannumeral1) coïncide avec notre
condition que $f$ soit un morphisme rig-étale : cela résulte du
lemme \ref{restricted-lemma-equivalent-with-artin}, partie (\ref{restricted-item-condition-artin}).

\medskip\noindent
Pour (\romannumeral2). Supposons que $f$ soit rig-étale. Alors
$\Delta_f : X \to X \times_Y X$ est rig-étale en tant que morphisme
entre espaces algébriques formels localement noethériens qui sont
rig-étales sur $X$ (par $\text{id}_X$ pour le premier et par
$\text{pr}_1$ pour le second).
Voir les lemmes \ref{restricted-lemma-base-change-rig-etale} et
\ref{restricted-lemma-rig-etale-permanence}.
La propriété (\romannumeral2) coïncide donc avec notre condition
que $\Delta_f$ soit rig-surjectif d'après le
lemme \ref{restricted-lemma-closed-immersion-rig-smooth-rig-surjective}.

\medskip\noindent
Pour (\romannumeral3). La propriété (\romannumeral3) ne coïncide pas tout à
fait avec notre notion de morphisme rig-surjectif, car Artin exige que tout
morphisme adique $\text{Spf}(R) \to Y$ se relève en un morphisme vers $X$,
tandis que notre notion de rig-surjectivité affirme seulement l'existence
d'un relèvement après remplacement de $R$ par une extension.
Toutefois, puisque nous savons déjà, par (\romannumeral1) et
(\romannumeral2), que $\Delta_f$ est rig-étale et rig-surjectif,
ces conditions sont équivalentes d'après le
lemme \ref{restricted-lemma-rig-monomorphism-rig-surjective}.
\end{remark}

\begin{lemma}
\label{restricted-lemma-modification-gives-formal-modification}
Soient $S$, $f : X' \to X$, $T \subset |X|$, $U \subset X$,
$T' \subset |X'|$ et $U' \subset X'$ comme dans la
section \ref{restricted-section-completion-functor}.
Si $X$ est localement noethérien, si $f$ est propre et si $U' \to U$ est un
isomorphisme, alors $f_{/T} : X'_{/T'} \to X_{/T}$ est une modification formelle.
\end{lemma}

\begin{proof}
D'après Espaces formels, lemme \ref{formal-spaces-lemma-formal-completion-types},
la source et le but de la flèche sont des espaces algébriques formels
localement noethériens.
Les autres conditions résultent des
lemmes \ref{restricted-lemma-proper-gives-proper},
\ref{restricted-lemma-etale-gives-rig-etale},
\ref{restricted-lemma-completion-proper-surjective-rig-surjective} et
\ref{restricted-lemma-separated-mono-open-diagonal-rig-surjective}.
\end{proof}

\begin{lemma}
\label{restricted-lemma-base-change-formal-modification}
Soit $S$ un schéma. Soit $f : X \to Y$ un morphisme d'espaces algébriques
formels localement noethériens sur $S$ qui est une modification formelle.
Alors, pour tout morphisme adique $Y' \to Y$ d'espaces algébriques formels
localement noethériens, le changement de base $f' : X \times_Y Y' \to Y'$
est une modification formelle.
\end{lemma}

\begin{proof}
Le morphisme $f'$ est propre d'après
Espaces formels, lemme \ref{formal-spaces-lemma-base-change-proper}.
Le morphisme $f'$ est rig-étale d'après le
lemme \ref{restricted-lemma-base-change-rig-etale}.
Le morphisme $f'$ est donc rig-surjectif d'après le
lemme \ref{restricted-lemma-base-change-rig-surjective}. Posons $X' = X \times_ Y'$.
Le morphisme $\Delta_{f'}$ est le changement de base de
$\Delta_f$ par le morphisme adique $X' \times_{Y'} X' \to X \times_Y X$.
Ainsi, $\Delta_{f'}$ est rig-surjectif d'après le
lemme \ref{restricted-lemma-base-change-rig-surjective}.
\end{proof}









\section{Complétions et morphismes, I}
\label{restricted-section-completion-and-morphisms}

\noindent
Dans cette section, nous établissons quelques résultats préliminaires sur les
complétions que nous utiliserons dans la démonstration du
théorème \ref{restricted-theorem-dilatations-general}.
Bien que les lemmes énoncés et démontrés ici ne soient pas triviaux
(certains reposent sur notre étude de l'algébrisation des algèbres rig-étales),
nous conseillons néanmoins au lecteur d'omettre cette section en première lecture.

\begin{lemma}
\label{restricted-lemma-algebraize-rig-etale-affine}
Soit $T \subset X$ une partie fermée d'un schéma affine noethérien $X$.
Soit $W$ un espace algébrique formel affine noethérien.
Soit $g : W \to X_{/T}$ un morphisme rig-étale. Il existe alors
un schéma affine $X'$ et un morphisme de type fini $f : X' \to X$,
étale sur $X \setminus T$, ainsi qu'un isomorphisme
$X'_{/f^{-1}T} \cong W$ compatible avec $f_{/T}$ et $g$.
De plus, si $W \to X_{/T}$ est étale, alors $X' \to X$ est étale.
\end{lemma}

\begin{proof}
L'existence de $X'$ est une reformulation du
lemme \ref{restricted-lemma-approximate-by-etale-over-complement}.
La dernière assertion résulte de
Compléments sur les morphismes, lemme
\ref{more-morphisms-lemma-check-smoothness-on-infinitesimal-nbhds}.
\end{proof}

\begin{lemma}
\label{restricted-lemma-algebraize-morphism-rig-etale}
Supposons donnés
\begin{enumerate}
\item des schémas affines noethériens $X$, $X'$ et $Y$,
\item une partie fermée $T \subset |X|$,
\item un morphisme $f : X' \to X$ localement de type fini
et étale sur $X \setminus T$,
\item un morphisme $h : Y \to X$,
\item un morphisme $\alpha : Y_{/T} \to X'_{/T}$ sur $X_{/T}$
(voir la démonstration pour la notation).
\end{enumerate}
Il existe alors un morphisme étale $b : Y' \to Y$ de schémas affines
qui induit un isomorphisme $b_{/T} : Y'_{/T} \to Y_{/T}$
et un morphisme $a : Y' \to X'$ sur $X$
tels que $\alpha = a_{/T} \circ b_{/T}^{-1}$.
\end{lemma}

\begin{proof}
Dans l'énoncé, la notation utilisant l'indice ${}_{/T}$ renvoie à la
construction qui, à un morphisme de schémas $g : V \to X$,
associe le morphisme $g_{/T} : V_{/g^{-1}T} \to X_{/T}$ d'espaces
algébriques formels ; c'est un foncteur de la catégorie des schémas sur $X$
vers la catégorie des espaces algébriques formels sur $X_{/T}$, voir la
section \ref{restricted-section-completion-functor}.
Cela étant dit, le lemme n'est qu'une reformulation du
lemme \ref{restricted-lemma-fully-faithful-etale-over-complement}.
\end{proof}

\begin{lemma}
\label{restricted-lemma-factor}
Soit $S$ un schéma. Soient $f : X \to Y$ et $g : Z \to Y$ des morphismes
d'espaces algébriques. Soit $T \subset |X|$ une partie fermée.
Supposons que
\begin{enumerate}
\item $X$ soit localement noethérien,
\item $g$ soit un monomorphisme localement de type fini,
\item $f|_{X \setminus T} : X \setminus T \to Y$ se factorise par $g$, et
\item $f_{/T} : X_{/T} \to Y$ se factorise par $g$,
\end{enumerate}
alors $f$ se factorise par $g$.
\end{lemma}

\begin{proof}
Considérons le produit fibré $E = X \times_Y Z \to X$.
Par hypothèse, l'immersion ouverte $X \setminus T \to X$
se factorise par $E$, et tout morphisme $\varphi : X' \to X$ tel que
$|\varphi|(|X'|) \subset T$ se factorise lui aussi par $E$, voir
Espaces formels, section \ref{formal-spaces-section-completion}.
D'après Compléments sur les morphismes d'espaces, lemme
\ref{spaces-more-morphisms-lemma-check-smoothness-on-infinitesimal-nbhds}
cela implique que $E \to X$ est étale en tout point de $E$
qui s'envoie sur un point de $T$. Ainsi, $E \to X$ est un
monomorphisme étale, donc une immersion ouverte
(Morphismes d'espaces, lemme
\ref{spaces-morphisms-lemma-etale-universally-injective-open}).
Il s'ensuit que $E = X$, puisque nos hypothèses impliquent que $|X| = |E|$.
\end{proof}

\begin{lemma}
\label{restricted-lemma-faithful-general}
Soit $S$ un schéma. Soient $X$ et $W$ des espaces algébriques sur $S$, avec
$X$ localement noethérien. Soit $T \subset |X|$ une partie fermée.
Soient $a, b : X \to W$ des morphismes d'espaces algébriques sur $S$ tels
que $a|_{X \setminus T} = b|_{X \setminus T}$ et que
$a_{/T} = b_{/T}$ comme morphismes $X_{/T} \to W$. Alors $a = b$.
\end{lemma}

\begin{proof}
Soit $E$ l'égalisateur de $a$ et $b$. Alors $E$ est un espace algébrique
et $E \to X$ est localement de type fini et est un monomorphisme, voir
Morphismes d'espaces, lemme \ref{spaces-morphisms-lemma-properties-diagonal}.
Nos hypothèses permettent d'appliquer le lemme \ref{restricted-lemma-factor} aux deux
morphismes $f = \text{id} : X \to X$ et $g : E \to X$, ainsi qu'à la partie
fermée $T$ de $|X|$.
\end{proof}

\begin{lemma}
\label{restricted-lemma-faithful}
Soit $S$ un schéma. Soient $X$ et $Y$ des espaces algébriques localement
noethériens sur $S$. Soient $T \subset |X|$ et $T' \subset |Y|$ des parties
fermées. Soient $a, b : X \to Y$ des morphismes d'espaces algébriques sur
$S$ tels que $a|_{X \setminus T} = b|_{X \setminus T}$, que
$|a|(T) \subset T'$ et $|b|(T) \subset T'$, et que
$a_{/T} = b_{/T}$ comme morphismes $X_{/T} \to Y_{/T'}$.
Alors $a = b$.
\end{lemma}

\begin{proof}
C'est une conséquence du lemme plus général \ref{restricted-lemma-faithful-general}.
\end{proof}

\begin{lemma}
\label{restricted-lemma-equivalence-relation}
Soit $S$ un schéma. Soit $X$ un espace algébrique localement noethérien
sur $S$. Soit $T \subset |X|$ une partie fermée.
Soient $s, t : R \to U$ deux morphismes d'espaces algébriques sur $X$.
Supposons que
\begin{enumerate}
\item $R$ et $U$ soient localement de type fini sur $X$,
\item le changement de base de $s$ et $t$ à $X \setminus T$
soit une relation d'équivalence étale, et
\item le complété formel
$(t_{/T}, s_{/T}) : R_{/T} \to U_{/T} \times_{X_{/T}} U_{/T}$
soit lui aussi une relation d'équivalence (voir la démonstration pour la notation).
\end{enumerate}
Alors $(t, s) : R \to U \times_X U$ est une relation d'équivalence étale.
\end{lemma}

\begin{proof}
Dans l'énoncé, la notation utilisant l'indice ${}_{/T}$ renvoie à la
construction qui, à un morphisme $f : X' \to X$ d'espaces algébriques,
associe le morphisme $f_{/T} : X'_{/f^{-1}T} \to X_{/T}$ d'espaces
algébriques formels, voir la section \ref{restricted-section-completion-functor}.
Les morphismes $s, t : R \to U$ sont étales sur $X \setminus T$
par hypothèse. Comme les complétés formels des applications
$s, t : R \to U$ sont étales, nous voyons que $s$ et $t$ sont étales,
par exemple d'après Compléments sur les morphismes, lemme
\ref{more-morphisms-lemma-check-smoothness-on-infinitesimal-nbhds}.
En appliquant le lemme \ref{restricted-lemma-factor} aux morphismes
$\text{id} : R \times_{U \times_X U} R \to R \times_{U \times_X U} R$
et $\Delta : R \to R \times_{U \times_X U} R$, nous concluons que
$(t, s)$ est un monomorphisme. En l'appliquant de nouveau à
$(t \circ \text{pr}_0, s \circ \text{pr}_1) :
R \times_{s, U, t} R \to U \times_X U$ et $(t, s) : R \to U \times_X U$,
nous constatons que la « transitivité » est satisfaite. Nous omettons la
démonstration des deux autres axiomes d'une relation d'équivalence.
\end{proof}

\begin{lemma}
\label{restricted-lemma-smash-away-from-T}
Soit $S$ un schéma. Soit $X$ un espace algébrique localement noethérien sur
$S$, et soit $T \subset |X|$ une partie fermée. Soit $f : X' \to X$ un
morphisme d'espaces algébriques localement de type fini et étale hors de $T$.
Il existe une factorisation
$$
X' \longrightarrow X'' \longrightarrow X
$$
de $f$ ayant les propriétés suivantes :
$X'' \to X$ est localement de type fini,
$X'' \to X$ est un isomorphisme sur $X \setminus T$, et
$X'_{/T} \to X''_{/T}$ est un isomorphisme (voir la démonstration pour la notation).
\end{lemma}

\begin{proof}
Dans l'énoncé, la notation utilisant l'indice ${}_{/T}$ renvoie à la
construction qui, à un morphisme $f : X' \to X$ d'espaces algébriques,
associe le morphisme $f_{/T} : X'_{/f^{-1}T} \to X_{/T}$ d'espaces
algébriques formels, voir la section \ref{restricted-section-completion-functor}.
Nous utiliserons également les notations $U \subset X$ et $U' \subset X'$
pour les sous-espaces ouverts tels que $|U| = |X| \setminus T$ et
$U' = |X'| \setminus f^{-1}T$, introduits à la
section \ref{restricted-section-completion-functor}.

\medskip\noindent
Après avoir remplacé $X'$ par $X' \amalg U$, nous pouvons supposer, et
supposons désormais, que l'image de $X' \to X$ contient $U$.
Soit
$$
R = X' \amalg_{U'} (U' \times_U U')
$$
la somme amalgamée de $U' \to X'$ et du morphisme diagonal
$U' \to U' \times_U U' = U' \times_X U'$. Comme $U' \to X$ est étale,
cette diagonale est une immersion ouverte, et nous voyons que $R$ est un
espace algébrique (cela résulte par exemple d'
Espaces, lemme \ref{spaces-lemma-glueing-algebraic-spaces}).
Les deux projections $U' \times_U U' \to U'$ se prolongent à $R$
et nous obtenons deux morphismes étales $s, t : R \to X'$.
En vérifiant séparément sur chaque morceau, nous constatons que $R$
est une relation d'équivalence étale sur $X'$. Posons $X'' = X'/R$,
qui est un espace algébrique d'après
Amorçage, théorème \ref{bootstrap-theorem-final-bootstrap}.
Par construction, nous avons la factorisation de l'énoncé, et
le morphisme $X'' \to X$ est localement de type fini (car cela
se vérifie localement pour la topologie étale, c'est-à-dire sur $X'$).
Comme $U' \to U$ est un morphisme étale surjectif
et que $s^{-1}(U') = t^{-1}(U') = U' \times_U U'$,
nous voyons que $U'' = U \times_X X'' \to U$ est un isomorphisme.
Enfin, il faut montrer que le morphisme $X' \to X''$ induit un isomorphisme
$X'_{/T} \to X''_{/T}$. Pour cela, notons que le complété formel de $R$
le long de l'image réciproque de $T$ est égal au complété formel de
$X'$ le long de l'image réciproque de $T$, par notre choix de $R$ ! D'après
notre construction du complété formel dans
Espaces formels, section \ref{formal-spaces-section-completion},
nous avons $X''_{/T} = (X'_{/T}) / (R_{/T})$ comme faisceaux. Comme
$X'_{/T} = R_{/T}$, nous concluons que $X'_{/T} = X''_{/T}$,
ce qui achève la démonstration.
\end{proof}








\section{Recollement rigide de morphismes}
\label{restricted-section-glue-formal-to-actual}

\noindent
Soient $X$ et $W$ des espaces algébriques, avec $X$ noethérien. Soit
$Z \subset X$ un sous-espace fermé de complémentaire ouvert $U$.
La proposition ci-dessous affirme, en gros, que
$$
\{\text{morphismes }X \to W\} =
\{\text{morphismes compatibles }U \to W\text{ et }X_{/Z} \to W\}
$$
où la compatibilité de $a : U \to W$ et $b : X_{/Z} \to W$
signifie que $a$ et $b$ définissent le même « morphisme d'espaces rigides ».
Introduire la catégorie des « espaces rigides » demande beaucoup de travail,
mais nous n'avons pas à le faire pour énoncer précisément ce que signifie
la condition dans ce cas.

\begin{proposition}
\label{restricted-proposition-glue-modification}
Soit $S$ un schéma. Soit $X$ un espace algébrique localement noethérien sur
$S$. Soit $T \subset |X|$ une partie fermée de sous-espace ouvert
complémentaire $U \subset X$. Soit $f : X' \to X$ un morphisme propre
d'espaces algébriques tel que $f^{-1}(U) \to U$ soit un isomorphisme.
Pour tout espace algébrique $W$ sur $S$, l'application
$$
\Mor_S(X, W) \longrightarrow
\Mor_S(X', W) \times_{\Mor_S(X'_{/T}, W)} \Mor_S(X_{/T}, W)
$$
est bijective.
\end{proposition}

\begin{proof}
Soient $w' : X' \to W$ et $\hat w : X_{/T} \to W$ des morphismes
qui déterminent le même morphisme $X'_{/T} \to W$ par composition
avec $X'_{/T} \to X$ et $X'_{/T} \to X_{/T}$. Il faut montrer
qu'il existe un unique morphisme $w : X \to W$ dont la composition
avec $X' \to X$ et $X_{/T} \to X$ redonne $w'$ et $\hat w$.
L'unicité résulte immédiatement du lemme \ref{restricted-lemma-faithful-general}.

\medskip\noindent
Les hypothèses sur $T$ et $f$ sont préservées par changement de base
par tout morphisme étale $X_1 \to X$ d'espaces algébriques.
Comme les espaces algébriques formels sont des faisceaux pour la
topologie étale et que l'unicité est déjà acquise,
il suffit de démontrer l'existence après avoir remplacé $X$ par
les membres d'un recouvrement étale. Nous pouvons donc supposer que $X$
est un schéma affine noethérien.

\medskip\noindent
Supposons que $X$ soit un schéma affine noethérien. Nous allons construire
le morphisme $w : X \to W$ à l'aide des résultats de
Sommes amalgamées d'espaces, section
\ref{spaces-pushouts-section-coequalizer-glue}.
Il peut être utile de lire une partie de cette section avant de poursuivre
la lecture de la démonstration.

\medskip\noindent
Posons $X'' = X' \times_X X'$ et considérons les deux morphismes
$a = w' \circ \text{pr}_1 : X'' \to W$ et
$b = w' \circ \text{pr}_2 : X'' \to W$.
Nous voyons alors que $a$ et $b$ coïncident sur l'ouvert $U$
et que $a_{/T}$ et $b_{a/T}$ coïncident (car ils sont
tous deux égaux au composé $X''_{/T} \to X_{/T} \to W$,
où la seconde flèche est $\hat w$).
Le lemme \ref{restricted-lemma-faithful-general} donne donc
$a = b$.

\medskip\noindent
Notons $Z \subset X$ la structure de sous-schéma fermé réduit induite
sur $T$. Pour $n \geq 1$, notons $Z_n \subset X$ le $n$-ième voisinage
infinitésimal de $Z$. Posons $w_n = \hat w|_{Z_n} : Z_n \to W$,
de sorte que $\hat w = \colim w_n$ sur $X_{/T} = \colim Z_n$.
Posons $Y_n = X' \amalg Z_n$. Considérons les deux projections
$$
s_n, t_n : R_n = Y_n \times_X Y_n \longrightarrow Y_n
$$
Soit $Y_n \to X_n \to X$ le coégalisateur de $s_n$ et $t_n$ comme dans
Sommes amalgamées d'espaces, section
\ref{spaces-pushouts-section-coequalizer-glue}
(en particulier, ce coégalisateur existe et possède de bonnes propriétés,
voir Sommes amalgamées d'espaces, lemme
\ref{spaces-pushouts-lemma-coequalizer}).
D'après l'égalité $a = b$ du paragraphe précédent et la coïncidence
de $w'$ et $\hat w$ sur $X'_{/T}$, nous voyons que le morphisme
$$
w' \amalg w_n : Y_n \longrightarrow W
$$
égalise les morphismes $s_n$ et $t_n$. Ainsi, pour tout $n \geq 1$,
il existe un morphisme $w^n : X_n \to W$ compatible avec $w'$ et $w_n$.
De plus, pour $m \geq 1$, le composé
$$
X_n \to X_{n + m} \xrightarrow{w^{n + m}} W
$$
est égal à $w^n$ par construction (car l'assertion correspondante vaut
pour $w' \amalg w_{n + m}$ et $w' \amalg w_n$). D'après
Sommes amalgamées d'espaces, lemme
\ref{spaces-pushouts-lemma-essentially-constant} et de
la remarque \ref{spaces-pushouts-remark-essentially-constant},
le système d'espaces algébriques $X_n$ est essentiellement constant
de valeur $X$, ce qui permet de conclure.
\end{proof}














\section{Algébrisation des morphismes rig-étales}
\label{restricted-section-algebraization}

\noindent
Dans cette section, nous démontrons une généralisation du résultat sur les
dilatations de l'article d'Artin \cite{ArtinII}.

\medskip\noindent
Les notations de cette section seront celles de la
section \ref{restricted-section-completion-functor}, à ceci près que nos espaces
algébriques et espaces algébriques formels seront localement noethériens.

\medskip\noindent
Ainsi, fixons d'abord un schéma de base $S$. Tous les anneaux, anneaux
topologiques, schémas, espaces algébriques et espaces algébriques formels,
ainsi que leurs morphismes, seront sur $S$. Fixons ensuite un espace
algébrique localement noethérien $X$ et une partie fermée $T \subset |X|$.
Notons $U \subset X$ le sous-espace ouvert tel que $|U| = |X| \setminus T$.
On peut représenter la situation ainsi :
$$
U \to X \quad |X| = |U| \amalg T
$$
Étant donné un morphisme d'espaces algébriques $f : X' \to X$, nous
emploierons les notations $U' = f^{-1}U$, $T' = |f|^{-1}(T)$ et
$f_{/T} : X'_{/T'} \to X_{/T}$ comme dans la
section \ref{restricted-section-completion-functor}. Nous écrirons parfois
$X'_{/T}$ au lieu de $X'_{/T'}$ et, plus généralement, pour un morphisme
$a : X' \to X''$ d'espaces algébriques sur $X$, nous noterons
$a_{/T} : X'_{/T} \to X''_{/T}$ le morphisme induit d'espaces algébriques
formels obtenu en complétant le morphisme $a$ le long des images réciproques
de $T$ dans $X'$ et $X''$.

\medskip\noindent
Dans cette situation, nous considérerons le foncteur
\begin{equation}
\label{restricted-equation-completion-functor}
\left\{
\begin{matrix}
\text{morphismes d'espaces algébriques}\\
f : X' \to X\text{ qui sont localement}\\
\text{de type fini et tels que}\\
U' \to U\text{ est un isomorphisme}
\end{matrix}
\right\}
\longrightarrow
\left\{
\begin{matrix}
\text{morphismes }g : W \to X_{/T}\\
\text{d'espaces algébriques formels}\\
\text{avec }W\text{ localement noethérien}\\
\text{et }g\text{ rig-étale}
\end{matrix}
\right\}
\end{equation}
qui envoie $f : X' \to X$ sur $f_{/T} : X'_{/T'} \to X_{/T}$.
Cette définition a un sens, car $f_{/T}$ est rig-étale d'après le
lemme \ref{restricted-lemma-etale-gives-rig-etale}.

\begin{lemma}
\label{restricted-lemma-functoriality-completion-functor}
Dans la situation ci-dessus, soit $X_1 \to X$ un morphisme d'espaces
algébriques, avec $X_1$ localement noethérien. Notons
$T_1 \subset |X_1|$ l'image réciproque de $T$ et $U_1 \subset X_1$
l'image réciproque de $U$. Nous notons
\begin{enumerate}
\item $\mathcal{C}_{X, T}$ la catégorie dont les objets sont les
morphismes d'espaces algébriques $f : X' \to X$ qui sont localement
de type fini et tels que $U' = f^{-1}U \to U$ soit un isomorphisme,
\item $\mathcal{C}_{X_1, T_1}$ la catégorie dont les objets sont les
morphismes d'espaces algébriques $f_1 : X_1' \to X_1$ qui sont localement
de type fini et tels que $f_1^{-1}U_1 \to U_1$ soit un isomorphisme,
\item $\mathcal{C}_{X_{/T}}$ la catégorie dont les objets sont les
morphismes $g : W \to X_{/T}$ d'espaces algébriques formels,
avec $W$ localement noethérien et $g$ rig-étale,
\item $\mathcal{C}_{X_{1, /T_1}}$ la catégorie dont les objets sont les
morphismes $g_1 : W_1 \to X_{1, /T_1}$ d'espaces algébriques formels,
avec $W_1$ localement noethérien et $g_1$ rig-étale.
\end{enumerate}
Alors le diagramme
$$
\xymatrix{
\mathcal{C}_{X, T} \ar[d] \ar[r] &
\mathcal{C}_{X_{/T}} \ar[d] \\
\mathcal{C}_{X_1, T_1} \ar[r] &
\mathcal{C}_{X_{1, /T_1}}
}
$$
est commutatif, où les flèches horizontales sont données par
(\ref{restricted-equation-completion-functor})
et les flèches verticales par changement de base le long de
$X_1 \to X$ et de $X_{1, /T_1} \to X_{/T}$.
\end{lemma}

\begin{proof}
Cela résulte immédiatement du fait que le foncteur de complétion
$(h : Y \to X) \mapsto Y_{/T} = Y_{/|h|^{-1}T}$
sur la catégorie des espaces algébriques sur $X$
commute aux produits fibrés.
\end{proof}

\begin{lemma}
\label{restricted-lemma-completion-functor-fully-faithful}
Dans la situation ci-dessus, soit $f : X' \to X$ un morphisme d'espaces
algébriques localement de type fini et qui est un isomorphisme sur $U$.
Soit $g : Y \to X$ un morphisme, avec $Y$ localement noethérien. Alors la
complétion définit une bijection
$$
\Mor_X(Y, X') \longrightarrow \Mor_{X_{/T}}(Y_{/T}, X'_{/T})
$$
En particulier, le foncteur (\ref{restricted-equation-completion-functor}) est
pleinement fidèle.
\end{lemma}

\begin{proof}
Soient $a, b : Y \to X'$ des morphismes sur $X$ tels que
$a_{/T} = b_{/T}$. Nous voyons alors que $a$ et $b$ coïncident sur le
sous-espace ouvert $g^{-1}U$ et après complétion le long de $g^{-1}T$.
Ainsi, $a = b$ d'après le lemme \ref{restricted-lemma-faithful}.
Autrement dit, l'application de complétion est toujours injective.

\medskip\noindent
Soit $\alpha : Y_{/T} \to X'_{/T}$ un morphisme d'espaces algébriques
formels sur $X_{/T}$. Il faut montrer qu'il existe un morphisme
$a : Y \to X'$ sur $X$ tel que $\alpha = a_{/T}$. La démonstration procède
par une réduction standard, mais laborieuse, au cas affine, puis applique le
lemme \ref{restricted-lemma-algebraize-morphism-rig-etale}.

\medskip\noindent
Soit $\{h_i : Y_i \to Y\}$ un recouvrement étale d'espaces algébriques.
Si, pour tout $i$, nous pouvons trouver un morphisme $a_i : Y_i \to X'$
sur $X$ dont le complété $(a_i)_{/T} : (Y_i)_{/T} \to X'_{/T}$ est égal à
$\alpha \circ (h_i)_{/T}$, alors nous obtenons un morphisme $a : Y \to X'$
tel que $\alpha = a_{/T}$. En effet, observons d'abord que
$(a_i)_{/T} \circ \text{pr}_1 = (a_j)_{/T} \circ \text{pr}_2$
comme morphismes $(Y_i \times_Y Y_j)_{/T} \to X'_{/T}$, par coïncidence
avec $\alpha$ (on utilise ici le fait que la complétion ${}_{/T}$ commute
aux produits fibrés). L'injectivité déjà démontrée montre alors que
$a_i \circ \text{pr}_1 = a_j \circ \text{pr}_2$ comme morphismes
comme morphismes $Y_i \times_Y Y_j \to X'$. Comme $X'$ est un faisceau
fppf, cela signifie que la famille de morphismes $a_i$ descend en un
morphisme $a : Y \to X'$. Nous avons $\alpha = a_{/T}$, car
$\{(a_i)_{/T} : (Y_i)_{/T} \to X'_{/T}\}$ est un recouvrement étale.

\medskip\noindent
D'après le paragraphe précédent, pour démontrer l'existence, nous pouvons
supposer que $Y$ est affine et que $g : Y \to X$ se factorise par
$g_1 : Y \to X_1$ et par un morphisme étale $X_1 \to X$, avec $X_1$ affine.
Nous pouvons alors considérer $T_1 \subset |X_1|$, image réciproque de $T$,
et poser $X'_1 = X' \times_X X_1$, avec projection $f_1 : X'_1 \to X_1$, et
$$
\alpha_1 = (\alpha, (g_1)_{/T_1}) :
Y_{/T_1} = Y_{/T}
\longrightarrow
X'_{/T} \times_{X_{/T}} (X_1)_{/T_1} = (X'_1)_{/T_1}
$$
Nous concluons qu'il suffit de démontrer l'existence pour $\alpha_1$
sur $X_1$ ; autrement dit, nous pouvons remplacer $X, T, X', Y, f, g, \alpha$
par $X_1, T_1, X'_1, Y, g_1, \alpha_1$.
Nous sommes ainsi ramenés au cas décrit dans le paragraphe suivant.

\medskip\noindent
Supposons que $Y$ et $X$ soient affines. Rappelons que $(Y_{/T})_{red}$
est un schéma affine (isomorphe à la structure de schéma réduit induite
sur $g^{-1}T \subset Y$, voir Espaces formels, lemme
\ref{formal-spaces-lemma-reduction-completion}).
Ainsi, $\alpha_{red} : (Y_{/T})_{red} \to (X'_{/T})_{red}$
a une image quasi-compacte $E$ dans $f^{-1}T$ (c'est l'espace topologique
sous-jacent à $(X'_{/T})_{red}$ d'après le même lemme).
Nous pouvons donc trouver un schéma affine $V$ et un morphisme étale
$h : V \to X'$ tel que l'image de $h$ contienne $E$.
Choisissons un diagramme cartésien sur les flèches pleines
$$
\xymatrix{
Y'_{/T} \ar@{..>}[rd] \ar@{..>}[r] &
W \ar[d] \ar[r] & V_{/T} \ar[d]^{h_{/T}} \\
& Y_{/T} \ar[r]^\alpha & X'_{/T}
}
$$
Par construction, le morphisme $W \to Y_{/T}$ est représentable par des
espaces algébriques, étale et surjectif
(la surjectivité se voit en considérant les réductions, voir
Espaces formels, lemme
\ref{formal-spaces-lemma-reduction-surjective}).
D'après le lemme \ref{restricted-lemma-algebraize-rig-etale-affine},
nous pouvons écrire $W = Y'_{/T}$, où $Y' \to Y$ est étale et $Y'$ affine.
Cela fournit les flèches pointillées du diagramme.
Comme $W \to Y_{/T}$ est surjectif, nous voyons que
l'image de $Y' \to Y$ contient $g^{-1}T$.
Ainsi, $\{Y' \to Y, Y \setminus g^{-1}T \to Y\}$
est un recouvrement étale. Comme $f$ est un isomorphisme sur $U$, nous avons
un unique morphisme $Y \setminus g^{-1}T \to X'$ sur $X$
qui coïncide avec $\alpha$ sur les complétés
(car le complété de $Y \setminus g^{-1}T$ est vide).
Il suffit donc de démontrer l'existence pour $Y'$,
ce qui nous ramène au cas étudié dans le paragraphe suivant.

\medskip\noindent
D'après le paragraphe précédent, nous pouvons supposer que $Y$ est affine
et que $\alpha$ se factorise en $Y_{/T} \to V_{/T} \to X'_{/T}$,
où $V$ est un schéma affine étale sur $X'$. Nous pouvons encore remplacer
$Y$ par les membres d'un recouvrement étale affine.
D'après le lemme \ref{restricted-lemma-algebraize-morphism-rig-etale},
nous pouvons trouver un morphisme étale $b : Y' \to Y$ de schémas affines
qui induit un isomorphisme $b_{/T} : Y'_{/T} \to Y_{/T}$
et un morphisme $c : Y' \to V$ tel que $c_{/T} \circ b_{/T}^{-1}$
soit le morphisme donné $Y_{/T} \to V_{/T}$. Si nous posons
$a' : Y' \to X'$ égal au composé de $c$ et $V \to X'$,
nous obtenons $a'_{/T} = \alpha \circ b_{/T}$ ; autrement dit, nous avons
l'existence pour $Y'$ et $\alpha \circ b_{/T}$.
Nous concluons en considérant une fois encore le recouvrement étale
$\{Y' \to Y, Y \setminus g^{-1}T \to Y\}$.
\end{proof}

\begin{lemma}
\label{restricted-lemma-dilatations-affine}
Dans la situation ci-dessus, supposons que $X$ soit affine. Alors le foncteur
(\ref{restricted-equation-completion-functor}) est une équivalence.
\end{lemma}

\noindent
Avant de démontrer ce lemme, examinons un exemple. Supposons que
$S = \Spec(k)$, $X = \mathbf{A}^1_k$ et $T = \{0\}$. Alors
$X_{/T} = \text{Spf}(k[[x]])$. Soit $W = \text{Spf}(k[[x]] \times k[[x]])$.
Le morphisme correspondant $f : X' \to X$ est alors la droite affine à
origine dédoublée envoyée sur la droite affine
(Schémas, exemple \ref{schemes-example-affine-space-zero-doubled}).
De plus, c'est le résultat de la construction du
lemme \ref{restricted-lemma-smash-away-from-T}
appliquée à $X \amalg X$ sur $X$.

\begin{proof}
Nous savons déjà que le foncteur est pleinement fidèle, voir le
lemme \ref{restricted-lemma-completion-functor-fully-faithful}.
Surjectivité essentielle. Soit $g : W \to X_{/T}$ un morphisme
d'espaces algébriques formels, avec $W$ localement noethérien
et $g$ rig-étale. Nous allons démontrer en plusieurs étapes que $W$
appartient à l'image essentielle.

\medskip\noindent
Étape 1 : $W$ est un espace algébrique formel affine. Nous pouvons alors
trouver $U \to X$ de type fini et étale sur $X \setminus T$ tel que
$U_{/T}$ soit isomorphe à $W$, voir le
lemme \ref{restricted-lemma-algebraize-rig-etale-affine}.
Le lemme \ref{restricted-lemma-smash-away-from-T} montre donc que $W$
appartient à l'image essentielle.

\medskip\noindent
Étape 2 : $W$ est séparé. Choisissons $\{W_i \to W\}$ comme dans
Espaces formels, définition \ref{formal-spaces-definition-formal-algebraic-space}.
D'après l'étape 1, les espaces algébriques formels $W_i$ et
$W_i \times_W W_j$ appartiennent à l'image essentielle.
Écrivons $W_i = (X'_i)_{/T}$ et $W_i \times_W W_j = (X'_{ij})_{/T}$.
Par pleine fidélité, nous obtenons des morphismes
$t_{ij} : X'_{ij} \to X'_i$ et $s_{ij} : X'_{ij} \to X'_j$ correspondant
aux projections $W_i \times_W W_j \to W_i$ et $W_i \times_W W_j \to W_j$.
Considérons les données
$$
R = \coprod X'_{ij},\quad
V = \coprod X'_i,\quad
s = \coprod s_{ij},\quad
t = \coprod t_{ij}
$$
(Nous ne pouvons pas employer la lettre $U$, car elle l'est déjà.)
En appliquant le lemme \ref{restricted-lemma-equivalence-relation}, nous constatons que
$(t, s) : R \to V \times_X V$ définit une relation d'équivalence étale
sur $V$ au-dessus de $X$. Nous pouvons donc prendre le quotient
$X' = V/R$, qui est un espace algébrique, voir
Amorçage, théorème \ref{bootstrap-theorem-final-bootstrap}.
Comme la complétion commute aux produits fibrés et à la formation des
faisceaux quotients, nous obtenons
$X'_{/T} \cong W$ comme espaces algébriques formels sur $X_{/T}$.

\medskip\noindent
Étape 3 : $W$ est quelconque. Choisissons $\{W_i \to W\}$ comme dans
Espaces formels, définition \ref{formal-spaces-definition-formal-algebraic-space}.
Les espaces algébriques formels $W_i$ et $W_i \times_W W_j$ sont séparés.
L'étape 2 montre donc que les espaces algébriques formels $W_i$ et
$W_i \times_W W_j$ appartiennent à l'image essentielle. Le même argument
que dans le paragraphe précédent montre alors que $W$ appartient lui aussi
à l'image essentielle. Cela achève la démonstration.
\end{proof}

\begin{theorem}
\label{restricted-theorem-dilatations-general}
Soit $S$ un schéma. Soit $X$ un espace algébrique localement noethérien sur
$S$. Soit $T \subset |X|$ une partie fermée. Soit $U \subset X$ le
sous-espace ouvert tel que $|U| = |X| \setminus T$. Le foncteur de complétion
(\ref{restricted-equation-completion-functor})
$$
\left\{
\begin{matrix}
\text{morphismes d'espaces algébriques}\\
f : X' \to X\text{ qui sont localement}\\
\text{de type fini et tels que}\\
f^{-1}U \to U\text{ est un isomorphisme}
\end{matrix}
\right\}
\longrightarrow
\left\{
\begin{matrix}
\text{morphismes }g : W \to X_{/T}\\
\text{d'espaces algébriques formels}\\
\text{avec }W\text{ localement noethérien}\\
\text{et }g\text{ rig-étale}
\end{matrix}
\right\}
$$
qui envoie $f : X' \to X$ sur $f_{/T} : X'_{/T'} \to X_{/T}$ est une équivalence.
\end{theorem}

\begin{proof}
Le foncteur est pleinement fidèle d'après le
lemme \ref{restricted-lemma-completion-functor-fully-faithful}.
Soit $g : W \to X_{/T}$ un morphisme d'espaces algébriques formels,
avec $W$ localement noethérien et $g$ rig-étale.
Pour achever la démonstration, montrons que $W$ appartient à l'image essentielle.

\medskip\noindent
Choisissons un recouvrement étale $\{X_i \to X\}$ tel que $X_i$ soit affine
pour tout $i$. Notons $U_i \subset X_i$ l'image réciproque de $U$ et
$T_i \subset X_i$ l'image réciproque de $T$.
Rappelons que $(X_i)_{/T_i} = (X_i)_{/T} = (X_i \times_X X)_{/T}$ et
$W_i = X_i \times_X W = (X_i)_{/T} \times_{X_{/T}} W$, voir
le lemme \ref{restricted-lemma-functoriality-completion-functor}.
Observons que nous obtenons des isomorphismes
$$
\alpha_{ij} :
W_i \times_{X_{/T}} (X_j)_{/T}
\longrightarrow 
(X_i)_{/T} \times_{X_{/T}} W_j
$$
qui satisfont à une condition de cocycle convenable.
En appliquant le lemme \ref{restricted-lemma-dilatations-affine} à
$X_i, T_i, U_i, W_i \to (X_i)_{/T}$
nous obtenons un morphisme $X'_i \to X_i$ d'espaces algébriques,
localement de type fini et qui est un isomorphisme sur $U_i$,
ainsi qu'un isomorphisme $\beta_i : (X'_i)_{/T} \cong W_i$ sur $(X_i)_{/T}$.
Par pleine fidélité, nous obtenons un isomorphisme
$$
a_{ij} : X'_i \times_X X_j \longrightarrow X_i \times_X X'_j
$$
sur $X_i \times_X X_j$ tel que
$\alpha_{ij} = \beta_j|_{X_i \times_X X_j}
\circ (a_{ij})_{/T} \circ \beta_i^{-1}|_{X_i \times_X X_j}$.
En utilisant de nouveau la pleine fidélité (cette fois sur
$X_i \times_X X_j \times_X X_k$),
nous voyons que ces morphismes $a_{ij}$ satisfont à la même
condition de cocycle que les $\alpha_{ij}$.
Autrement dit, nous obtenons une donnée de descente
(comme dans Descente sur les espaces, définition
\ref{spaces-descent-definition-descent-datum-for-family-of-morphisms})
$(X'_i, a_{ij})$ relative à la famille $\{X_i \to X\}$.
D'après Amorçage, lemme \ref{bootstrap-lemma-descend-algebraic-space},
cette donnée de descente est effective. Nous obtenons donc un morphisme
$f : X' \to X$ d'espaces algébriques et des isomorphismes
$h_i : X' \times_X X_i \to X'_i$ sur $X_i$ tels que
$a_{ij} = h_j|_{X_i \times_X X_j} \circ h_i^{-1}|_{X_i \times_X X_j}$.
Le lecteur vérifiera que les isomorphismes qui en résultent
$$
(X' \times_X X_i)_{/T}
\xrightarrow{\beta_i \circ (h_i)_{/T}}
W_i
$$
sur $X_i$ se recollent en un isomorphisme $X'_{/T} \to W$
sur $X_{/T}$ ; certains détails sont omis.
\end{proof}









\section{Complétions et morphismes, II}
\label{restricted-section-completion-and-morphisms-bis}

\noindent
Pour obtenir le théorème d'Artin sur les dilatations, nous devons faire
correspondre les modifications formelles aux modifications algébriques dans
l'équivalence donnée par le théorème \ref{restricted-theorem-dilatations-general}.
Nous invitons le lecteur à omettre cette section.

\begin{lemma}
\label{restricted-lemma-output-quasi-compact}
Sous les hypothèses et avec les notations du
théorème \ref{restricted-theorem-dilatations-general}, soit $f : X' \to X$
correspondant à $g : W \to X_{/T}$.
Alors $f$ est quasi-compact si et seulement si $g$ est quasi-compact.
\end{lemma}

\begin{proof}
Si $f$ est quasi-compact, alors $g$ est quasi-compact d'après le
lemme \ref{restricted-lemma-quasi-compact-gives-quasi-compact}.
Réciproquement, supposons que $g$ soit quasi-compact.
Choisissons un recouvrement étale $\{X_i \to X\}$ avec $X_i$ affine.
Il suffit de démontrer que le changement de base $X' \times_X X_i \to X_i$
est quasi-compact, voir
Morphismes d'espaces, lemme \ref{spaces-morphisms-lemma-quasi-compact-local}.
D'après Espaces formels, lemme
\ref{formal-spaces-lemma-characterize-quasi-compact-morphism}
les changements de base $W_i \times_{X_{/T}} (X_i)_{/T}  \to (X_i)_{/T}$
sont quasi-compacts.
Le lemme \ref{restricted-lemma-functoriality-completion-functor}
nous ramène au cas décrit dans le paragraphe suivant.

\medskip\noindent
Supposons que $X$ soit affine et que $g : W \to X_{/T}$ soit quasi-compact.
Il faut montrer que $X'$ est quasi-compact.
Soit $V \to X'$ un morphisme étale surjectif, où
$V = \coprod_{j \in J} V_j$ est une réunion disjointe de schémas affines.
Alors $V_{/T} \to X'_{/T} = W$ est un morphisme étale surjectif.
Comme $W$ est quasi-compact, nous pouvons trouver une partie finie
$J' \subset J$ telle que $\coprod_{j \in J'} (V_j)_{/T} \to W$ soit surjective.
Il s'ensuit que
$$
U \amalg \coprod\nolimits_{j \in J'} V_j \longrightarrow X'
$$
est surjectif (et donc que $X'$ est quasi-compact).
En effet, nous avons $|X'| = |U| \amalg |W_{red}|$, puisque $X'_{/T} = W$.
\end{proof}

\begin{lemma}
\label{restricted-lemma-output-quasi-separated}
Sous les hypothèses et avec les notations du
théorème \ref{restricted-theorem-dilatations-general}, soit $f : X' \to X$
correspondant à $g : W \to X_{/T}$.
Alors $f$ est quasi-séparé si et seulement si $g$ l'est.
\end{lemma}

\begin{proof}
Si $f$ est quasi-séparé, alors $g$ est quasi-séparé d'après le
lemme \ref{restricted-lemma-quasi-separated-gives-quasi-separated}.
Réciproquement, supposons que $g$ soit quasi-séparé. Il faut montrer
que $f$ est quasi-séparé. Exactement comme dans la démonstration du
lemme \ref{restricted-lemma-output-quasi-compact}, nous pouvons le vérifier
sur les membres d'un recouvrement étale de $X$ par des schémas affines,
en utilisant Morphismes d'espaces, lemme
\ref{spaces-morphisms-lemma-separated-local}
et Espaces formels, lemme
\ref{formal-spaces-lemma-separated-local}.
Nous pouvons donc supposer, et supposons désormais, que $X$ est affine.

\medskip\noindent
Soit $V \to X'$ un morphisme étale surjectif, où
$V = \coprod_{j \in J} V_j$ est une réunion disjointe de schémas affines.
Pour montrer que $X'$ est quasi-séparé, il suffit de montrer que
$V_j \times_{X'} V_{j'}$ est quasi-compact pour tous $j, j' \in J$.
Comme $W$ est quasi-séparé, les produits fibrés
$(V_j \times_Y V_{j'})_{/T} = (V_j)_{/T} \times_{X'_{/T}} (V_{j'})_{/T}$
sont quasi-compacts pour tous $j, j' \in J$. Comme $X$ est affine
noethérien et que $U' \to U$ est un isomorphisme, nous voyons que
$$
(V_j \times_{X'} V_{j'}) \times_X U =
(V_j \times_X V_{j'}) \times_X U
$$
est quasi-compact. Nous concluons donc grâce à l'égalité
$$
|V_j \times_{X'} V_{j'}| =
|(V_j \times_{X'} V_{j'}) \times_X U| \amalg
|(V_j \times_{X'} V_{j'})_{/T, red}|
$$
et au fait qu'un espace algébrique formel est quasi-compact si et seulement
si l'espace algébrique réduit qui lui est associé l'est.
\end{proof}

\begin{lemma}
\label{restricted-lemma-output-separated}
Sous les hypothèses et avec les notations du
théorème \ref{restricted-theorem-dilatations-general}, soit $f : X' \to X$
correspondant à $g : W \to X_{/T}$.
Alors $f$ est séparé $\Leftrightarrow$ $g$ est séparé et
$\Delta_g : W \to W \times_{X_{/T}} W$ est rig-surjectif.
\end{lemma}

\begin{proof}
Si $f$ est séparé, alors $g$ est séparé et $\Delta_g$
est rig-surjectif d'après les
lemmes \ref{restricted-lemma-quasi-separated-gives-quasi-separated} et
\ref{restricted-lemma-separated-mono-open-diagonal-rig-surjective}.
Supposons que $g$ soit séparé et que $\Delta_g$ soit rig-surjectif.
Exactement comme dans la démonstration du
lemme \ref{restricted-lemma-output-quasi-compact},
nous pouvons le vérifier sur les membres d'un recouvrement étale de $X$
par des schémas affines, en utilisant
Morphismes d'espaces, lemme \ref{spaces-morphisms-lemma-base-change-separated}
(la propriété d'être séparé pour les morphismes d'espaces algébriques est
locale sur la base), Espaces formels, lemme
\ref{formal-spaces-lemma-base-change-separated}
(la propriété d'être séparé pour les morphismes d'espaces algébriques
formels est préservée par changement de base), et le
lemme \ref{restricted-lemma-base-change-rig-surjective} (la rig-surjectivité
est préservée par changement de base).
Nous pouvons donc supposer, et supposons désormais, que $X$ est affine.
De plus, nous savons déjà que $f : X' \to X$ est quasi-séparé d'après le
lemme \ref{restricted-lemma-output-quasi-separated}.

\medskip\noindent
D'après Cohomologie des espaces, lemme
\ref{spaces-cohomology-lemma-check-separated-dvr} et
remarque \ref{spaces-cohomology-remark-variant},
il suffit de montrer que, pour tout diagramme commutatif
$$
\xymatrix{
\Spec(K) \ar[r] \ar[d] & X' \ar[d] \\
\Spec(R) \ar[r]^p \ar@{-->}[ru] & X' \times_X X'
}
$$
où $R$ est un anneau de valuation discrète complet de corps des fractions
$K$, il existe une flèche pointillée rendant le diagramme commutatif (cela
donnera la partie unicité du critère valuatif). Soit
$h : \Spec(R) \to X$ le composé de $p$ avec le morphisme
$Y \times_X Y \to X$. Il y a trois cas :
Cas I : $h(\Spec(R)) \subset U$. Ce cas est trivial,
car $U' = X' \times_X U \to U$ est un isomorphisme.
Cas II : $h$ envoie $\Spec(R)$ dans $T$. Ce cas résulte
de notre hypothèse que $g : W \to X_{/T}$ est séparé. En effet,
si $Z$ désigne la structure de sous-espace fermé réduit induite
sur $T$, alors $h$ se factorise par $Z$ et
$$
W \times_{X_{/T}} Z = X' \times_X Z \longrightarrow Z
$$
est séparé par hypothèse (voir par exemple
Espaces formels, lemme \ref{formal-spaces-lemma-separated-local}),
ce qui donne la propriété de relèvement d'après
Cohomologie des espaces, lemme \ref{spaces-cohomology-lemma-check-separated-dvr},
appliqué à la flèche affichée. Cas III : $h(\Spec(K))$ n'appartient pas à
$T$, mais $h$ envoie le point fermé de $\Spec(R)$ dans $T$. Dans ce cas,
le morphisme correspondant
$$
p_{/T} :
\text{Spf}(R)
\longrightarrow
(X' \times_X X')_{/T} =
W \times_{X_{/T}} W
$$
est un morphisme adique (d'après
Espaces formels, lemme
\ref{formal-spaces-lemma-map-completions-representable} et
définition \ref{formal-spaces-definition-adic-morphism}).
Notre hypothèse que
$\Delta_g : W \to W \times_{X_{/T}} W$ est rig-surjectif implique donc que
nous pouvons relever $p_{/T}$ en un morphisme
$\text{Spf}(R) \to W = X'_{/T}$, voir le
lemme \ref{restricted-lemma-monomorphism-rig-surjective}.
En algébrisant le composé $\text{Spf}(R) \to X'$ à l'aide d'
Espaces formels, lemme \ref{formal-spaces-lemma-map-into-algebraic-space},
nous obtenons un morphisme $\Spec(R) \to X'$ relevant $p$, comme souhaité.
\end{proof}

\begin{lemma}
\label{restricted-lemma-output-proper}
Sous les hypothèses et avec les notations du
théorème \ref{restricted-theorem-dilatations-general}, soit $f : X' \to X$
correspondant à $g : W \to X_{/T}$.
Alors $f$ est propre si et seulement si $g$ est une modification formelle
(définition \ref{restricted-definition-formal-modification}).
\end{lemma}

\begin{proof}
Si $f$ est propre, alors $g$ est une modification formelle d'après le
lemme \ref{restricted-lemma-modification-gives-formal-modification}.
Supposons que $g$ soit une modification formelle. D'après les
lemmes \ref{restricted-lemma-output-quasi-compact} et \ref{restricted-lemma-output-separated},
nous voyons que $f$ est quasi-compact et séparé.

\medskip\noindent
D'après Cohomologie des espaces, lemme
\ref{spaces-cohomology-lemma-check-proper-dvr} et
remarque \ref{spaces-cohomology-remark-variant},
il suffit de montrer que, pour tout diagramme commutatif
$$
\xymatrix{
\Spec(K) \ar[r] \ar[d] & X' \ar[d]^f \\
\Spec(R) \ar[r]^p \ar@{-->}[ru] & X
}
$$
où $R$ est un anneau de valuation discrète complet de corps des fractions
$K$, il existe une flèche pointillée rendant le diagramme commutatif.
Il y a trois cas :
Cas I : $p(\Spec(R)) \subset U$. Ce cas est trivial,
car $U' \to U$ est un isomorphisme.
Cas II : $p$ envoie $\Spec(R)$ dans $T$. Ce cas résulte
de notre hypothèse que $g : W \to X_{/T}$ est propre. En effet,
si $Z$ désigne la structure de sous-espace fermé réduit induite
sur $T$, alors $p$ se factorise par $Z$ et
$$
W \times_{X_{/T}} Z = X' \times_X Z \longrightarrow Z
$$
est propre par hypothèse, ce qui donne la propriété de relèvement d'après
Cohomologie des espaces, lemme \ref{spaces-cohomology-lemma-check-proper-dvr},
appliqué à la flèche affichée. Cas III : $p(\Spec(K))$ n'appartient pas à
$T$, mais $p$ envoie le point fermé de $\Spec(R)$ dans $T$. Dans ce cas,
le morphisme correspondant
$$
p_{/T} : \text{Spf}(R) \longrightarrow X'_{/T} = W
$$
est un morphisme adique (d'après
Espaces formels, lemme
\ref{formal-spaces-lemma-map-completions-representable} et
définition \ref{formal-spaces-definition-adic-morphism}).
Notre hypothèse que $g : W \to X_{/T}$ est rig-surjectif
implique donc que nous pouvons relever $g_{/T}$ en un morphisme
$\text{Spf}(R') \to W = X'_{/T}$
pour une certaine extension d'anneaux de valuation discrète complets
$R \subset R'$. En algébrisant le composé $\text{Spf}(R') \to X'$ à l'aide d'
Espaces formels, lemme \ref{formal-spaces-lemma-map-into-algebraic-space},
nous obtenons un morphisme $\Spec(R') \to X'$ relevant $p$, comme souhaité.
\end{proof}

\begin{lemma}
\label{restricted-lemma-output-etale}
Sous les hypothèses et avec les notations du
théorème \ref{restricted-theorem-dilatations-general}, soit $f : X' \to X$
correspondant à $g : W \to X_{/T}$.
Alors $f$ est étale si et seulement si $g$ est étale.
\end{lemma}

\begin{proof}
Si $f$ est étale, alors $g$ est étale d'après le
lemme \ref{restricted-lemma-map-completions-etale}.
Réciproquement, supposons que $g$ soit étale.
Comme $f$ est un isomorphisme sur $U$, nous voyons que $f$ est étale
sur $U$. Il suffit donc de démontrer que $f$ est étale
en tout point de $X'$ situé au-dessus de $T$. Notons $Z \subset X$
le sous-espace fermé réduit dont l'espace topologique sous-jacent
est $|Z| = T \subset |X|$, voir
Propriétés des espaces, définition
\ref{spaces-properties-definition-reduced-induced-space}.
Si $Z_n \subset X$ désigne le $n$-ième voisinage infinitésimal,
nous avons $X_{/T} = \colim Z_n$. Puisque $X'_{/T} = W \to X_{/T}$,
nous concluons que $f^{-1}(Z_n) = X' \times_X Z_n \to Z_n$
est étale par l'étalité supposée de $g$.
D'après Compléments sur les morphismes d'espaces, lemme
\ref{spaces-more-morphisms-lemma-check-smoothness-on-infinitesimal-nbhds}
nous concluons que $f$ est étale aux points situés au-dessus de $T$.
\end{proof}





\section{Théorème d'Artin sur les dilatations}
\label{restricted-section-dilatations}

\noindent
Dans cette section, nous employons une fonte différente pour les espaces
algébriques formels afin de souligner la similitude des énoncés avec les
énoncés correspondants de \cite{ArtinII}. Voici le premier théorème principal
de ce chapitre.

\begin{theorem}
\label{restricted-theorem-dilatations}
\begin{reference}
\cite[Theorem 3.2]{ArtinII}
\end{reference}
Soit $S$ un schéma. Soit $X$ un espace algébrique localement noethérien sur
$S$. Soit $T \subset |X|$ une partie fermée. Soit
$\mathfrak X = X_{/T}$
le complété formel de $X$ le long de $T$. Soit
$$
\mathfrak f : \mathfrak X' \to \mathfrak X
$$
une modification formelle (définition \ref{restricted-definition-formal-modification}).
Alors il existe un unique morphisme propre $f : X' \to X$ qui est un
isomorphisme sur le complémentaire de $T$ dans $X$ et dont le complété
$f_{/T}$ redonne $\mathfrak f$.
\end{theorem}

\begin{proof}
Cela résulte du théorème \ref{restricted-theorem-dilatations-general}
et du lemme \ref{restricted-lemma-output-proper}.
\end{proof}

\noindent
Voici la caractérisation des modifications formelles
annoncée dans la section \ref{restricted-section-formal-modifications}.

\begin{lemma}
\label{restricted-lemma-formal-modifications-locally-algebraic}
Soit $S$ un schéma. Soit $\mathfrak X' \to \mathfrak X$
une modification formelle (définition \ref{restricted-definition-formal-modification})
d'espaces algébriques formels localement noethériens sur $S$. Étant donnés
\begin{enumerate}
\item un anneau topologique noethérien adique quelconque $A$,
\item un morphisme adique quelconque $\text{Spf}(A) \longrightarrow \mathfrak X$,
\end{enumerate}
il existe un morphisme propre $X \to \Spec(A)$ d'espaces algébriques
et un isomorphisme
$$
\text{Spf}(A) \times_{\mathfrak X} \mathfrak X'
\longrightarrow
X_{/Z}
$$
sur $\text{Spf}(A)$, entre le changement de base de $\mathfrak X$
et le complété formel de $X$ le long de la « fibre fermée »
$Z = X \times_{\Spec(A)} \text{Spf}(A)_{red}$ de $X$ sur $A$.
\end{lemma}

\begin{proof}
Le morphisme $\text{Spf}(A) \times_{\mathfrak X} \mathfrak X'
\to \text{Spf}(A)$ est une modification formelle d'après le
lemme \ref{restricted-lemma-base-change-formal-modification}.
L'assertion résulte donc du théorème \ref{restricted-theorem-dilatations}.
\end{proof}







\section{Application aux modifications}
\label{restricted-section-modifications}

\noindent
Soit $A$ un anneau noethérien et soit $I \subset A$ un idéal. Posons
$X = \Spec(A)$ et $U = X \setminus V(I)$. Dans cette section,
nous considérerons la catégorie
\begin{equation}
\label{restricted-equation-modification}
\left\{
f : X' \longrightarrow X
\quad \middle| \quad
\begin{matrix}
X'\text{ est un espace algébrique}\\
f\text{ est localement de type fini}\\
f^{-1}(U) \to U\text{ est un isomorphisme}
\end{matrix}
\right\}
\end{equation}
Un morphisme de $X'/X$ vers $X''/X$ sera un morphisme d'espaces algébriques
$X' \to X''$ sur $X$.

\medskip\noindent
Soit $A \to B$ un homomorphisme d'anneaux noethériens et soit $J \subset B$
un idéal tel que $J = \sqrt{I B}$. Le changement de base le long du
morphisme $\Spec(B) \to \Spec(A)$ fournit alors un foncteur
de la catégorie (\ref{restricted-equation-modification}) pour $A$
vers la catégorie (\ref{restricted-equation-modification}) pour $B$.

\begin{lemma}
\label{restricted-lemma-Noetherian-local-ring}
Soit $A \to B$ un homomorphisme d'anneaux noethériens induisant un
isomorphisme sur les complétés $I$-adiques pour un certain idéal $I \subset A$
(par exemple, si $B$ est le complété $I$-adique de $A$).
Alors le changement de base définit une équivalence entre la
catégorie (\ref{restricted-equation-modification}) pour $(A, I)$
et la catégorie (\ref{restricted-equation-modification}) pour $(B, IB)$.
\end{lemma}

\begin{proof}
Posons $X = \Spec(A)$ et $T = V(I)$.
Posons $X_1 = \Spec(B)$ et $T_1 = V(IB)$.
D'après le théorème \ref{restricted-theorem-dilatations-general} (en fait, seul le cas
affine traité dans le lemme \ref{restricted-lemma-dilatations-affine} est nécessaire),
la catégorie (\ref{restricted-equation-modification}) pour $X$ et $T$
est équivalente à la catégorie des morphismes rig-étales
$W \to X_{/T}$ d'espaces algébriques formels localement noethériens.
De même, la catégorie (\ref{restricted-equation-modification})
pour $X_1$ et $T_1$ est équivalente à la catégorie des morphismes rig-étales
$W_1 \to X_{1, /T_1}$ d'espaces algébriques formels localement
noethériens. Comme $X_{/T} = \text{Spf}(A^\wedge)$
et $X_{1, /T_1} = \text{Spf}(B^\wedge)$ (Espaces formels, lemme
\ref{formal-spaces-lemma-affine-formal-completion-types}), nous voyons que
ces catégories sont équivalentes en vertu de notre hypothèse selon laquelle
$A^\wedge \to B^\wedge$ est un isomorphisme. Nous omettons la vérification
que cette équivalence est donnée par changement de base.
\end{proof}

\begin{lemma}
\label{restricted-lemma-Noetherian-local-ring-properties}
Notations et hypothèses comme dans le lemme \ref{restricted-lemma-Noetherian-local-ring}.
Soit $f : X' \to \Spec(A)$ correspondant à $g : Y' \to \Spec(B)$
par l'équivalence. Alors $f$ est quasi-compact, quasi-séparé, séparé,
propre, fini, et ajouter d'autres conditions ici si et seulement si $g$ l'est.
\end{lemma}

\begin{proof}
On peut le déduire pour les propriétés
quasi-compact, quasi-séparé, séparé et propre
en utilisant les lemmes \ref{restricted-lemma-output-quasi-compact},
\ref{restricted-lemma-output-quasi-separated},
\ref{restricted-lemma-output-separated},
\ref{restricted-lemma-output-quasi-separated} et
\ref{restricted-lemma-output-proper},
pour traduire la propriété correspondante en une propriété du complété
formel, puis en utilisant l'argument de la démonstration du
lemme \ref{restricted-lemma-Noetherian-local-ring}.
Il existe toutefois l'argument direct suivant, par descente fpqc.
On peut d'abord se ramener à démontrer le lemme pour $A \to A^\wedge$
et $B \to B^\wedge$, puisque $A^\wedge  \to B^\wedge$ est un isomorphisme.
Notons ensuite que $\{U \to \Spec(A), \Spec(A^\wedge) \to \Spec(A)\}$ est un
recouvrement fpqc, avec $U = \Spec(A) \setminus V(I)$ comme précédemment.
Le changement de base de $f$ par $U \to \Spec(A)$ est $\text{id}_U$
par définition de notre catégorie (\ref{restricted-equation-modification}).
Soit $P$ une propriété des morphismes d'espaces algébriques qui est locale
pour la topologie fpqc sur la base (Descente sur les espaces, définition
\ref{spaces-descent-definition-property-morphisms-local})
et telle que $P$ soit satisfaite par les morphismes identiques.
Nous voyons alors que $P$ est satisfaite par $f$ si et seulement si $P$ est
satisfaite par $g$. Cela s'applique lorsque $P$ est l'une des propriétés
quasi-compact, quasi-séparé, séparé, propre et fini, d'après
Descente sur les espaces, lemmes
\ref{spaces-descent-lemma-descending-property-quasi-compact},
\ref{spaces-descent-lemma-descending-property-quasi-separated},
\ref{spaces-descent-lemma-descending-property-separated},
\ref{spaces-descent-lemma-descending-property-proper} et
\ref{spaces-descent-lemma-descending-property-finite}.
\end{proof}

\begin{lemma}
\label{restricted-lemma-equivalence-to-completion}
Soit $A \to B$ un homomorphisme local d'anneaux locaux noethériens tel que
\begin{enumerate}
\item $A \to B$ soit plat,
\item $\mathfrak m_B = \mathfrak m_A B$, et
\item $\kappa(\mathfrak m_A) = \kappa(\mathfrak m_B)$
\end{enumerate}
Alors le foncteur de changement de base de la catégorie
(\ref{restricted-equation-modification}) pour $(A, \mathfrak m_A)$ vers la catégorie
(\ref{restricted-equation-modification}) pour $(B, \mathfrak m_B)$
est une équivalence.
\end{lemma}

\begin{proof}
Les conditions signifient que $A \to B$ induit un isomorphisme sur les
complétés, voir
Compléments d'algèbre, lemme \ref{more-algebra-lemma-flat-unramified}.
Ce lemme est donc un cas particulier du
lemme \ref{restricted-lemma-Noetherian-local-ring}.
\end{proof}

\begin{lemma}
\label{restricted-lemma-dominate-by-admissible-blowup}
Soit $(A, \mathfrak m, \kappa)$ un anneau local noethérien.
Soit $f : X \to S$ un objet de (\ref{restricted-equation-modification})
tel que $f$ soit propre.
Alors il existe un éclatement $U$-admissible $S' \to S$
qui domine $X$.
\end{lemma}

\begin{proof}
C'est un cas particulier de Compléments sur les morphismes d'espaces,
lemme \ref{spaces-more-morphisms-lemma-dominate-modification-by-blowup}.
\end{proof}








% OK, start here.
%

\chapter{Retour sur la résolution des surfaces}



\phantomsection
\label{spaces-resolve-section-phantom}


\section{Introduction}
\label{spaces-resolve-section-introduction}

\noindent
Ce chapitre traite de la résolution des singularités des
espaces algébriques noethériens de dimension $2$.
Dans un chapitre antérieur, nous avons déjà étudié, à la suite de
Lipman \cite{Lipman}, la résolution des surfaces dans le cas des schémas.
Voir Résolution des surfaces, section
\ref{resolve-section-introduction}.
La plupart des résultats de ce chapitre se déduisent immédiatement
des résultats relatifs aux schémas.

\medskip\noindent
Sauf mention expresse du contraire, tous les objets géométriques
de ce chapitre seront des espaces algébriques. Ainsi, lorsque nous dirons
``soit $f : X \to Y$ une modification'', cela signifiera que
$f$ est un morphisme au sens de Espaces au-dessus de corps, définition
\ref{spaces-over-fields-definition-modification}.
Il en ira de même pour les morphismes propres, et ainsi de suite.










\section{Modifications}
\label{spaces-resolve-section-modifications}

\noindent
Soit $(A, \mathfrak m, \kappa)$ un anneau local noethérien. Posons
$S = \Spec(A)$ et $U = S \setminus \{\mathfrak m\}$. Dans cette section,
nous considérerons la catégorie
\begin{equation}
\label{spaces-resolve-equation-modification}
\left\{
f : X \longrightarrow S
\quad \middle| \quad
\begin{matrix}
X\text{ est un espace algébrique}\\
f\text{ est un morphisme propre}\\
f^{-1}(U) \to U\text{ est un isomorphisme}
\end{matrix}
\right\}
\end{equation}
Un morphisme de $X/S$ vers $X'/S$ sera un morphisme d'espaces algébriques
$X \to X'$ compatible aux morphismes structuraux au-dessus de $S$. Dans
Algébrisation des espaces formels, section
\ref{restricted-section-modifications}, nous avons vu que cette catégorie
ne dépend que du complété de $A$ et démontré quelques propriétés élémentaires
de ses objets. Dans cette section, nous étudions plus particulièrement les cas
où $\dim(A) \leq 2$ ou bien où la fibre fermée est de dimension au plus $1$.

\begin{lemma}
\label{spaces-resolve-lemma-modification}
Soit $(A, \mathfrak m, \kappa)$ un domaine local noethérien
de dimension $2$ tel que $U = \Spec(A) \setminus \{\mathfrak m\}$
soit un schéma normal. Alors toute modification $f : X \to \Spec(A)$
est un morphisme comme dans (\ref{spaces-resolve-equation-modification}).
\end{lemma}

\begin{proof}
Soit $f : X \to S$ une modification. Il nous faut montrer que
$f^{-1}(U) \to U$ est un isomorphisme. Comme tout point fermé $u$ de $U$
est de codimension $1$, cela résulte de
Espaces au-dessus de corps, lemme
\ref{spaces-over-fields-lemma-modification-normal-iso-over-codimension-1}.
\end{proof}

\begin{lemma}
\label{spaces-resolve-lemma-closed-immersion-on-fibre}
Soit $(A, \mathfrak m, \kappa)$ un anneau local noethérien.
Soit $g : X \to Y$ un morphisme de la catégorie
(\ref{spaces-resolve-equation-modification}). Si le morphisme induit
$X_\kappa \to Y_\kappa$ entre les fibres spéciales est une immersion fermée,
alors $g$ est une immersion fermée.
\end{lemma}

\begin{proof}
C'est un cas particulier de
Compléments sur les morphismes d'espaces, lemme
\ref{spaces-more-morphisms-lemma-where-closed-immersion}.
\end{proof}

\begin{lemma}
\label{spaces-resolve-lemma-dimension-special-fibre}
Soit $(A, \mathfrak m, \kappa)$ un domaine local noethérien
de dimension $\geq 1$.
Soit $f : X \to \Spec(A)$ un morphisme d'espaces algébriques.
Supposons satisfaite au moins l'une des conditions suivantes :
\begin{enumerate}
\item $f$ est une modification (Espaces au-dessus de corps, définition
\ref{spaces-over-fields-definition-modification}),
\item $f$ est une altération (Espaces au-dessus de corps, définition
\ref{spaces-over-fields-definition-alteration}),
\item $f$ est localement de type fini et quasi-séparé, $X$ est intègre,
et il existe exactement un point de $|X|$ dont l'image est le point générique
de $\Spec(A)$,
\item $f$ est localement de type fini, $X$ est raisonnable, et les points
de $|X|$ dont l'image est le point générique de $\Spec(A)$ sont
les points génériques des composantes irréductibles de $|X|$,
\item ajouter d'autres cas ici.
\end{enumerate}
Alors $\dim(X_\kappa) \leq \dim(A) - 1$.
\end{lemma}

\begin{proof}
Les cas (1), (2) et (3) sont des cas particuliers de (4). Choisissons un
schéma affine $U = \Spec(B)$ et un morphisme étale $U \to X$.
L'homomorphisme d'anneaux $A \to B$ est de type fini. Il nous faut montrer que
$\dim(U_\kappa) \leq \dim(A) - 1$. Puisque $X$ est raisonnable, les points
génériques des composantes irréductibles de $U$ sont les points situés
au-dessus des points génériques des composantes irréductibles de $|X|$ ; voir
Espaces raisonnables, lemme
\ref{decent-spaces-lemma-decent-generic-points}.
La fibre de $\Spec(B) \to \Spec(A)$ au-dessus de $(0)$ est donc l'ensemble
(fini) des idéaux premiers minimaux $\mathfrak q_1, \ldots, \mathfrak q_r$
de $B$. Ainsi, $A_f \to B_f$ est fini pour un certain $f \in A$ non nul
(Algèbre, lemme \ref{algebra-lemma-generically-finite}).
Il s'ensuit que $\kappa(\mathfrak q_i)$ est une extension finie du
corps des fractions de $A$.
Soit $\mathfrak q \subset B$ un idéal premier au-dessus de $\mathfrak m$. Alors
$$
\dim(B_\mathfrak q) = \max \dim((B/\mathfrak q_i)_{\mathfrak q})
\leq \dim(A)
$$
l'inégalité résultant de la formule de dimension appliquée à
$A \subset B/\mathfrak q_i$ ; voir Algèbre, lemme
\ref{algebra-lemma-dimension-formula}.
Toutefois, la dimension de $B_\mathfrak q/\mathfrak m B_\mathfrak q$
(qui est l'anneau local de $U_\kappa$ au point correspondant)
est inférieure d'au moins une unité, car les idéaux premiers minimaux
$\mathfrak q_i$ n'appartiennent pas à $V(\mathfrak m)$. On conclut par
Propriétés, lemme \ref{properties-lemma-dimension}.
\end{proof}

\begin{lemma}
\label{spaces-resolve-lemma-modification-of-dim-2-is-projective-over-complete}
Si $(A, \mathfrak m, \kappa)$ est un domaine local noethérien complet
de dimension $2$, alors toute modification de $\Spec(A)$ est projective sur $A$.
\end{lemma}

\begin{proof}
Par Compléments sur les morphismes d'espaces, lemme
\ref{spaces-more-morphisms-lemma-projective-over-complete}
il suffit de montrer que la fibre spéciale de toute modification
$X$ de $\Spec(A)$ est de dimension $\leq 1$.
Cela résulte du lemme \ref{spaces-resolve-lemma-dimension-special-fibre}.
\end{proof}






\section{Stratégie}
\label{spaces-resolve-section-strategy}

\noindent
Soit $S$ un schéma. Soit $X$ un espace algébrique raisonnable sur $S$.
Soient $x_1, \ldots, x_n \in |X|$ des points fermés deux à deux distincts.
Pour chaque $i$, choisissons un voisinage étale élémentaire
$(U_i, u_i) \to (X, x_i)$ comme dans Espaces raisonnables, lemme
\ref{decent-spaces-lemma-decent-space-elementary-etale-neighbourhood}.
Cela signifie que $U_i$ est un schéma affine, que $U_i \to X$ est étale,
que $u_i$ est l'unique point de $U_i$ au-dessus de $x_i$, et que
$\Spec(\kappa(u_i)) \to X$ est un monomorphisme représentant $x_i$.
Quitte à restreindre $U_i$, nous pouvons et voulons supposer que, si
$j \not = i$, aucun point de $U_i$ ne s'envoie sur $x_j$.
Remarquons que $u_i \in U_i$ est un point fermé.

\medskip\noindent
Notons $\mathcal{C}_{X, \{x_1, \ldots, x_n\}}$ la catégorie des
morphismes d'espaces algébriques $f : Y \to X$ qui induisent un isomorphisme
$f^{-1}(X \setminus \{x_1, \ldots, x_n\}) \to
X \setminus \{x_1, \ldots, x_n\}$.
Pour chaque $i$, notons $\mathcal{C}_{U_i, u_i}$ la catégorie des
morphismes d'espaces algébriques $g_i : Y_i \to U_i$ qui induisent un
isomorphisme $g_i^{-1}(U_i \setminus \{u_i\}) \to U_i \setminus \{u_i\}$.
Le changement de base définit un foncteur
\begin{equation}
\label{spaces-resolve-equation-equivalence}
F :
\mathcal{C}_{X, \{x_1, \ldots, x_n\}}
\longrightarrow
\mathcal{C}_{U_1, u_1} \times \ldots \times \mathcal{C}_{U_n, u_n}
\end{equation}
Le lemme suivant permettra de ramener au cas des schémas au moins certains
des problèmes de ce chapitre.

\begin{lemma}
\label{spaces-resolve-lemma-equivalence}
Le foncteur $F$ de (\ref{spaces-resolve-equation-equivalence}) est une équivalence.
\end{lemma}

\begin{proof}
Pour $n = 1$, c'est Limites d'espaces, lemme
\ref{spaces-limits-lemma-excision-modifications}.
Pour $n > 1$, le lemme se démontre exactement de la même manière, ou s'en
déduit. Supposons par exemple que les $g_i : Y_i \to U_i$ soient des objets
de $\mathcal{C}_{U_i, u_i}$. Le cas $n = 1$ fournit alors des
$f'_i : Y'_i \to X$ qui sont des isomorphismes au-dessus de
$X \setminus \{x_i\}$ et dont le changement de base à $U_i$ est $f_i$.
On peut alors poser
$$
f : Y = Y'_1 \times_X \ldots \times_X Y'_n \to X
$$
C'est un objet de $\mathcal{C}_{X, \{x_1, \ldots, x_n\}}$ dont le
changement de base par $U_i \to X$ redonne $g_i$. Le foncteur est donc
essentiellement surjectif. Nous omettons la démonstration de sa
pleine fidélité.
\end{proof}

\begin{lemma}
\label{spaces-resolve-lemma-equivalence-properties}
Soient $X, x_i, U_i \to X, u_i$ comme dans
(\ref{spaces-resolve-equation-equivalence}). Si $f : Y \to X$ correspond aux
$g_i : Y_i \to U_i$ par $F$, alors $f$ est quasi-compact, quasi-séparé,
séparé, localement de présentation finie, de présentation finie, localement
de type fini, de type fini, propre, entier ou fini si et seulement si tous les
$g_i$ le sont, pour $i = 1, \ldots, n$.
\end{lemma}

\begin{proof}
Cela résulte de Limites d'espaces, lemme
\ref{spaces-limits-lemma-excision-modifications-properties}.
\end{proof}

\begin{lemma}
\label{spaces-resolve-lemma-equivalence-fibre}
Soient $X, x_i, U_i \to X, u_i$ comme dans
(\ref{spaces-resolve-equation-equivalence}). Si $f : Y \to X$ correspond aux
$g_i : Y_i \to U_i$ par $F$, alors
$Y_{x_i} \cong (Y_i)_{u_i}$ comme espaces algébriques.
\end{lemma}

\begin{proof}
C'est clair, puisque $u_i \to x_i$ est un isomorphisme.
\end{proof}







\section{Domination par des transformations quadratiques}
\label{spaces-resolve-section-quadratic-spaces}

\noindent
Nous ne définissons l'éclatement d'un espace en un point que si $X$ est
raisonnable.

\begin{definition}
\label{spaces-resolve-definition-blowup-at-point}
Soit $S$ un schéma. Soit $X$ un espace algébrique raisonnable sur $S$.
Soit $x \in |X|$ un point fermé. D'après Espaces raisonnables, lemme
\ref{decent-spaces-lemma-decent-space-closed-point}, on peut représenter $x$
par une immersion fermée $i : \Spec(k) \to X$.
Par {\it éclatement $X' \to X$ de $X$ en $x$}, on entend l'éclatement de $X$
le long du sous-espace fermé $Z = i(\Spec(k)) \subset X$.
\end{definition}

\noindent
À ce degré de généralité, l'éclatement de $X$ en $x$ n'est pas nécessairement
propre. Toutefois, si $X$ est localement noethérien, il résulte de
Diviseurs sur les espaces, lemme
\ref{spaces-divisors-lemma-blowing-up-projective}, que cet éclatement est
propre. Rappelons qu'un espace algébrique localement noethérien est noethérien
si et seulement s'il est quasi-compact et quasi-séparé. De plus, pour un
espace algébrique localement noethérien, être quasi-séparé équivaut à être
raisonnable (Espaces raisonnables, lemme
\ref{decent-spaces-lemma-locally-Noetherian-decent-quasi-separated}).

\begin{lemma}
\label{spaces-resolve-lemma-equivalence-sequence-blowups}
Soient $X, x_i, U_i \to X, u_i$ comme dans
(\ref{spaces-resolve-equation-equivalence}), et supposons que $f : Y \to X$ corresponde aux
$g_i : Y_i \to U_i$ par $F$. Il existe alors une factorisation
$$
Y = Z_m \to Z_{m - 1} \to \ldots \to Z_1 \to Z_0 = X
$$
de $f$ dans laquelle $Z_{j + 1} \to Z_j$ est l'éclatement de $Z_j$ en un
point fermé $z_j$ situé au-dessus de $\{x_1, \ldots, x_n\}$ si et seulement
si, pour chaque $i$, il existe une factorisation
$$
Y_i = Z_{i, m_i} \to Z_{i, m_i - 1} \to \ldots \to Z_{i, 1} \to Z_{i, 0} = U_i
$$
de $g_i$ dans laquelle $Z_{i, j + 1} \to Z_{i, j}$ est l'éclatement de
$Z_{i, j}$ en un point fermé $z_{i, j}$ situé au-dessus de $u_i$.
\end{lemma}

\begin{proof}
Un éclatement est un morphisme représentable. Dans l'un et l'autre cas,
on voit donc par récurrence que $Z_j \to X$ ou $Z_{i, j} \to U_i$ est
représentable. Par conséquent, chaque $Z_j$ ou $Z_{i, j}$ est un espace
algébrique raisonnable d'après Espaces raisonnables, lemme
\ref{decent-spaces-lemma-representable-named-properties}.
Cela montre que les assertions ont un sens (puisque l'éclatement n'est défini
que pour les espaces raisonnables).
Pour démontrer l'équivalence, partons d'une suite d'éclatements
$Z_m \to Z_{m - 1} \to \ldots \to Z_1 \to Z_0 = X$.
Le premier morphisme $Z_1 \to X$ est obtenu en éclatant l'un des $x_i$,
disons $x_1$. En appliquant $F$ à $Z_1 \to X$, on obtient l'éclatement
$Z_{1, 1} \to U_1$ en $u_1$, précisément l'éclatement en $u_1$, tandis que
$Z_{i, 0} = U_i$ pour $i > 1$.
À l'étape suivante, soit on éclate l'un des $x_i$, avec $i \geq 2$, sur
$Z_1$, soit on choisit un point fermé $z_1$ de la fibre de $Z_1 \to X$
au-dessus de $x_1$. Dans le premier cas, la marche à suivre est claire ;
dans le second, on utilise l'isomorphisme
$(Z_1)_{x_1} \cong (Z_{1, 1})_{u_1}$
(lemme \ref{spaces-resolve-lemma-equivalence-fibre}) pour obtenir le point fermé
$z_{1, 1} \in Z_{1, 1}$ correspondant à $z_1$. On définit alors
$Z_{1, 2} \to Z_{1, 1}$ comme l'éclatement en $z_{1, 1}$. En poursuivant
ainsi, on construit les factorisations de chacun des $g_i$.

\medskip\noindent
Réciproquement, étant données des suites d'éclatements
$Z_{i, m_i} \to Z_{i, m_i - 1} \to \ldots \to Z_{i, 1} \to Z_{i, 0} = U_i$
on construit exactement de la même manière la suite d'éclatements de $X$.
\end{proof}

\begin{lemma}
\label{spaces-resolve-lemma-make-ideal-principal}
Soit $S$ un schéma. Soit $X$ un espace algébrique noethérien sur $S$.
Soit $T \subset |X|$ un ensemble fini de points fermés $x$ tels que
(1) $X$ soit régulier en $x$ et (2) l'anneau local de $X$ en $x$ soit de
dimension $2$. Soit $\mathcal{I} \subset \mathcal{O}_X$ un faisceau
quasi-cohérent d'idéaux tel que $\mathcal{O}_X/\mathcal{I}$ soit supporté
par $T$. Il existe alors une suite
$$
X_m \to X_{m - 1} \to \ldots \to X_1 \to X_0 = X
$$
où $X_{j + 1} \to X_j$ est l'éclatement de $X_j$ en un point fermé $x_j$
situé au-dessus d'un point de $T$, telle que
$\mathcal{I}\mathcal{O}_{X_n}$ soit un faisceau inversible d'idéaux.
\end{lemma}

\begin{proof}
Écrivons $T = \{x_1, \ldots, x_r\}$. Choisissons des voisinages étales
élémentaires $(U_i, u_i) \to (X, x_i)$ comme dans la section
\ref{spaces-resolve-section-strategy}. Pour chaque $i$, la restriction
$\mathcal{I}_i = \mathcal{I}|_{U_i} \subset \mathcal{O}_{U_i}$
est un faisceau quasi-cohérent d'idéaux supporté en $u_i$.
L'anneau local de $U_i$ en $u_i$ est régulier et de dimension $2$.
Nous pouvons donc appliquer Résolution des surfaces, lemme
\ref{resolve-lemma-make-ideal-principal}, et obtenir une suite
$$
X_{i, m_i} \to X_{i, m_i - 1} \to \ldots \to X_1 \to X_{i, 0} = U_i
$$
d'éclatements en des points fermés situés au-dessus de $u_i$ telle que
$\mathcal{I}_i \mathcal{O}_{X_{i, m_i}}$ soit inversible.
Le lemme \ref{spaces-resolve-lemma-equivalence-sequence-blowups} fournit une suite
d'éclatements
$$
X_m \to X_{m - 1} \to \ldots \to X_1 \to X_0 = X
$$
comme dans l'énoncé du lemme, dont le changement de base aux $U_i$ redonne
les suites données. Il en résulte que $\mathcal{I}\mathcal{O}_{X_n}$ est un
faisceau inversible d'idéaux. En effet, nous le savons au-dessus de
$X \setminus \{x_1, \ldots, x_n\}$, et cela est vrai par construction dans
un voisinage étale de la fibre de chacun des $x_i$.
\end{proof}

\begin{lemma}
\label{spaces-resolve-lemma-dominate-by-blowing-up-in-points}
Soit $S$ un schéma. Soit $X$ un espace algébrique noethérien sur $S$.
Soit $T \subset |X|$ un ensemble fini de points fermés $x$ tels que
(1) $X$ soit régulier en $x$ et (2) l'anneau local de $X$ en $x$ soit de
dimension $2$. Soit $f : Y \to X$ un morphisme propre d'espaces algébriques
qui est un isomorphisme au-dessus de $U = X \setminus T$. Il existe alors
une suite
$$
X_n \to X_{n - 1} \to \ldots \to X_1 \to X_0 = X
$$
où $X_{i + 1} \to X_i$ est l'éclatement de $X_i$ en un point fermé $x_i$
situé au-dessus d'un point de $T$, ainsi qu'une factorisation
$X_n \to Y \to X$ de la composée.
\end{lemma}

\begin{proof}
D'après Compléments sur les morphismes d'espaces, lemme
\ref{spaces-more-morphisms-lemma-dominate-modification-by-blowup},
il existe un éclatement $U$-admissible $X' \to X$ qui domine $Y \to X$.
On peut donc supposer qu'il existe un faisceau d'idéaux
$\mathcal{I} \subset \mathcal{O}_X$ tel que
$\mathcal{O}_X/\mathcal{I}$ soit supporté par $T$ et que $Y$ soit
l'éclatement de $X$ le long de $\mathcal{I}$.
D'après le lemme \ref{spaces-resolve-lemma-make-ideal-principal}, il existe une suite
$$
X_n \to X_{n - 1} \to \ldots \to X_1 \to X_0 = X
$$
où $X_{i + 1} \to X_i$ est l'éclatement de $X_i$ en un point fermé $x_i$
situé au-dessus d'un point de $T$, telle que
$\mathcal{I}\mathcal{O}_{X_n}$ soit un faisceau inversible d'idéaux.
La propriété universelle de l'éclatement
(Diviseurs sur les espaces, lemme
\ref{spaces-divisors-lemma-universal-property-blowing-up})
fournit la factorisation voulue.
\end{proof}





\section{Domination par des éclatements normalisés}
\label{spaces-resolve-section-normalized-blowups}

\noindent
Dans cette section, nous montrons qu'une modification d'une surface peut être
dominée par une suite d'éclatements normalisés en des points.

\begin{definition}
\label{spaces-resolve-definition-normalized-blowup}
Soit $S$ un schéma. Soit $X$ un espace algébrique raisonnable sur $S$
satisfaisant aux conditions équivalentes de Morphismes d'espaces, lemme
\ref{spaces-morphisms-lemma-prepare-normalization}.
Soit $x \in |X|$ un point fermé. L'{\it éclatement normalisé de $X$ en $x$}
est la composée $X'' \to X' \to X$, où $X' \to X$ est l'éclatement de $X$
en $x$ (définition \ref{spaces-resolve-definition-blowup-at-point}) et où
$X'' \to X'$ est la normalisation de $X'$.
\end{definition}

\noindent
La normalisation $X'' \to X'$ est bien définie comme espace algébrique,
car $X'$ satisfait aux conditions équivalentes de
Morphismes d'espaces, lemme
\ref{spaces-morphisms-lemma-prepare-normalization}, d'après
Diviseurs sur les espaces, lemme
\ref{spaces-divisors-lemma-blowup-finite-nr-irreducibles}.
Voir Morphismes d'espaces, définition
\ref{spaces-morphisms-definition-normalization}
pour la définition de la normalisation.

\medskip\noindent
En général, l'éclatement normalisé n'est pas nécessairement propre, même
lorsque $X$ est noethérien. Rappelons qu'un espace algébrique est de Nagata
s'il admet un recouvrement étale par des espaces affines qui sont les spectres
d'anneaux de Nagata (Propriétés des espaces, définition
\ref{spaces-properties-definition-type-property} et
remarque \ref{spaces-properties-remark-list-properties-local-etale-topology},
et Propriétés, définition \ref{properties-definition-nagata}).

\begin{lemma}
\label{spaces-resolve-lemma-Nagata-normalized-blowup}
Dans la définition \ref{spaces-resolve-definition-normalized-blowup}, si $X$ est de Nagata,
alors l'éclatement normalisé de $X$ en $x$ est un espace algébrique normal
de Nagata, propre sur $X$.
\end{lemma}

\begin{proof}
Le morphisme d'éclatement $X' \to X$ est propre
(puisque $X$ est localement noethérien, nous pouvons appliquer
Diviseurs sur les espaces, lemme
\ref{spaces-divisors-lemma-blowing-up-projective}).
Ainsi, $X'$ est de Nagata (Morphismes d'espaces, lemme
\ref{spaces-morphisms-lemma-finite-type-nagata}).
La normalisation $X'' \to X'$ est donc finie (Morphismes d'espaces, lemme
\ref{spaces-morphisms-lemma-nagata-normalization}), et l'on en conclut que
$X'' \to X$ est également propre (Morphismes d'espaces, lemmes
\ref{spaces-morphisms-lemma-finite-proper} et
\ref{spaces-morphisms-lemma-composition-proper}).
Il s'ensuit que l'éclatement normalisé est un espace algébrique normal
(Morphismes d'espaces, lemme
\ref{spaces-morphisms-lemma-normalization-normal})
et de Nagata.
\end{proof}

\noindent
Voici, pour les éclatements normalisés, l'analogue du lemme
\ref{spaces-resolve-lemma-equivalence-sequence-blowups}.

\begin{lemma}
\label{spaces-resolve-lemma-equivalence-sequence-normalized-blowups}
Soient $X, x_i, U_i \to X, u_i$ comme dans
(\ref{spaces-resolve-equation-equivalence}), et supposons que $f : Y \to X$ corresponde aux
$g_i : Y_i \to U_i$ par $F$. Supposons que $X$ satisfasse aux conditions
équivalentes de Morphismes d'espaces, lemme
\ref{spaces-morphisms-lemma-prepare-normalization}.
Il existe alors une factorisation
$$
Y = Z_m \to Z_{m - 1} \to \ldots \to Z_1 \to Z_0 = X
$$
de $f$ dans laquelle $Z_{j + 1} \to Z_j$ est l'éclatement normalisé de $Z_j$
en un point fermé $z_j$ situé au-dessus de $\{x_1, \ldots, x_n\}$ si et
seulement si, pour chaque $i$, il existe une factorisation
$$
Y_i = Z_{i, m_i} \to Z_{i, m_i - 1} \to \ldots \to Z_{i, 1} \to Z_{i, 0} = U_i
$$
de $g_i$ dans laquelle $Z_{i, j + 1} \to Z_{i, j}$ est l'éclatement
normalisé de $Z_{i, j}$ en un point fermé $z_{i, j}$ situé au-dessus de $u_i$.
\end{lemma}

\begin{proof}
Cela résulte exactement du même argument que dans la démonstration du lemme
\ref{spaces-resolve-lemma-equivalence-sequence-blowups}.
\end{proof}

\noindent
Un espace algébrique de Nagata est localement noethérien.

\begin{lemma}
\label{spaces-resolve-lemma-dominate-by-normalized-blowing-up}
Soit $S$ un schéma. Soit $X$ un espace algébrique noethérien de Nagata sur
$S$, avec $\dim(X) = 2$. Soit $f : Y \to X$ un morphisme propre birationnel.
Il existe alors un diagramme commutatif
$$
\xymatrix{
X_n \ar[r] \ar[d] &
X_{n - 1} \ar[r] &
\ldots \ar[r] &
X_1 \ar[r] &
X_0 \ar[d] \\
Y \ar[rrrr]  & & & & X
}
$$
où $X_0 \to X$ est la normalisation et où $X_{i + 1} \to X_i$ est
l'éclatement normalisé de $X_i$ en un point fermé.
\end{lemma}

\begin{proof}
Bien que l'on puisse démontrer ce lemme directement pour les espaces
algébriques, nous poursuivrons la méthode employée ci-dessus pour nous ramener
au cas des schémas.

\medskip\noindent
Nous utiliserons le fait que les espaces algébriques noethériens sont
quasi-séparés et que leurs points possèdent donc des corps résiduels bien
définis (par exemple d'après Espaces raisonnables, lemme
\ref{decent-spaces-lemma-decent-space-elementary-etale-neighbourhood}).
Nous utiliserons sans autre mention les résultats de Morphismes d'espaces,
sections
\ref{spaces-morphisms-section-nagata},
\ref{spaces-morphisms-section-dimension-formula} et
\ref{spaces-morphisms-section-normalization}.
Nous pouvons remplacer $Y$ par sa normalisation. Soit $X_0 \to X$ la
normalisation. Le morphisme $Y \to X$ se factorise par $X_0$.
Nous pouvons donc supposer $X$ et $Y$ normaux.

\medskip\noindent
Supposons $X$ et $Y$ normaux. Le morphisme $f : Y \to X$ est un isomorphisme
au-dessus d'un ouvert qui contient tout point de codimension $0$ ou $1$ de
$Y$, ainsi que tout point de $Y$ au-dessus duquel la fibre est finie ; voir
Espaces au-dessus de corps, lemme
\ref{spaces-over-fields-lemma-modification-normal-iso-over-codimension-1}.
On voit donc qu'il existe un ensemble fini de points fermés $T \subset |X|$
tel que $f$ soit un isomorphisme au-dessus de $X \setminus T$.
D'après Compléments sur les morphismes d'espaces, lemme
\ref{spaces-more-morphisms-lemma-dominate-modification-by-blowup}
il existe un éclatement $X \setminus T$-admissible $Y' \to X$ qui domine $Y$.
Après avoir remplacé $Y$ par la normalisation de $Y'$, on voit que l'on peut
supposer le morphisme $Y \to X$ représentable.

\medskip\noindent
Écrivons $T = \{x_1, \ldots, x_r\}$. Choisissons des voisinages étales
élémentaires $(U_i, u_i) \to (X, x_i)$ comme dans la section
\ref{spaces-resolve-section-strategy}. Pour chaque $i$, le morphisme
$Y_i = Y \times_X U_i \to U_i$ est un morphisme propre birationnel qui est un
isomorphisme au-dessus de $U_i \setminus \{u_i\}$. Nous pouvons donc appliquer
Résolution des surfaces, lemme
\ref{resolve-lemma-dominate-by-normalized-blowing-up}
et obtenir une suite
$$
X_{i, m_i} \to X_{i, m_i - 1} \to \ldots \to X_1 \to X_{i, 0} = U_i
$$
d'éclatements normalisés en des points fermés situés au-dessus de $u_i$ telle
que $X_{i, m_i}$ domine $Y_i$.
Le lemme \ref{spaces-resolve-lemma-equivalence-sequence-normalized-blowups} fournit une suite
d'éclatements normalisés
$$
X_m \to X_{m - 1} \to \ldots \to X_1 \to X_0 = X
$$
comme dans l'énoncé du lemme, dont le changement de base à nos $U_i$ redonne
les suites données. Il résulte alors de l'équivalence de catégories du lemme
\ref{spaces-resolve-lemma-equivalence} que $X_m$ domine $Y$.
\end{proof}















\section{Changement de base au complété}
\label{spaces-resolve-section-aux}

\noindent
Le lemme simple qui suit sera un outil utile dans la suite.

\begin{lemma}
\label{spaces-resolve-lemma-iso-completions}
Soit $(A, \mathfrak m, \kappa)$ un anneau local dont l'idéal maximal
$\mathfrak m$ est de type fini. Soit $X$ un espace algébrique raisonnable
sur $A$. Posons $Y = X \times_{\Spec(A)} \Spec(A^\wedge)$, où $A^\wedge$
est le complété $\mathfrak m$-adique de $A$.
Soit $q \in |Y|$ un point d'image $p \in |X|$, situé au-dessus du point fermé
de $\Spec(A)$. L'homomorphisme d'anneaux locaux henséliens
$\mathcal{O}_{X, p}^h \to \mathcal{O}_{Y, q}^h$ induit un isomorphisme entre
leurs complétés.
\end{lemma}

\begin{proof}
Cela résulte immédiatement du cas des schémas : choisissons un voisinage
étale élémentaire $(U, u) \to (X, p)$ comme dans Espaces raisonnables, lemme
\ref{decent-spaces-lemma-decent-space-elementary-etale-neighbourhood},
et posons $V = U \times_X Y = U \times_{\Spec(A)} \Spec(A^\wedge)$ ainsi que
$v = (u, q)$.
Le cas des schémas est Résolution des surfaces, lemme
\ref{resolve-lemma-iso-completions}.
\end{proof}

\begin{lemma}
\label{spaces-resolve-lemma-port-regularity-to-completion}
Soit $(A, \mathfrak m, \kappa)$ un anneau local noethérien.
Soit $X \to \Spec(A)$ un morphisme localement de type fini, où $X$ est un
espace algébrique raisonnable. Posons
$Y = X \times_{\Spec(A)} \Spec(A^\wedge)$. Soit $y \in |Y|$
d'image $x \in |X|$. Alors :
\begin{enumerate}
\item si $\mathcal{O}_{Y, y}^h$ est régulier, alors
$\mathcal{O}_{X, x}^h$ est régulier ;
\item si $y$ appartient à la fibre fermée, alors
$\mathcal{O}_{Y, y}^h$ est régulier
$\Leftrightarrow \mathcal{O}_{X, x}^h$ est régulier ;
\item si $X$ est propre sur $A$, alors $X$ est régulier si et seulement si
$Y$ est régulier.
\end{enumerate}
\end{lemma}

\begin{proof}
Par localisation étale, les deux premières assertions résultent immédiatement
de l'analogue de ce lemme pour les schémas ; voir Résolution des surfaces,
lemme \ref{resolve-lemma-port-regularity-to-completion}.
Pour (3), puisque $Y \to X$ est surjectif (car $A \to A^\wedge$ est
fidèlement plat), la régularité de $Y$ entraîne celle de $X$ d'après (1).
Réciproquement, si $X$ est régulier, alors les anneaux locaux henséliens de
$Y$ sont réguliers en tous les points de la fibre spéciale.
Soit $y \in |Y|$ un point quelconque.
Dans le cas propre, l'application $|Y| \to |\Spec(A^\wedge)|$ est fermée ;
on peut donc trouver une spécialisation $y \leadsto y_0$, où $y_0$ appartient
à la fibre fermée. Choisissons un voisinage étale élémentaire
$(V, v_0) \to (Y, y_0)$ comme dans Espaces raisonnables, lemme
\ref{decent-spaces-lemma-decent-space-elementary-etale-neighbourhood}.
Comme $Y$ est raisonnable, on peut relever $y \leadsto y_0$ en une
spécialisation $v \leadsto v_0$ dans $V$ (Espaces raisonnables, lemme
\ref{decent-spaces-lemma-decent-specialization}).
On en conclut que $\mathcal{O}_{V, v}$ est un localisé de
$\mathcal{O}_{V, v_0}$, donc qu'il est régulier, ce qui achève la démonstration.
\end{proof}

\begin{lemma}
\label{spaces-resolve-lemma-formally-unramified}
Soit $(A, \mathfrak m)$ un anneau local noethérien. Soit $X$ un espace
algébrique sur $A$. Supposons que :
\begin{enumerate}
\item $A$ est analytiquement non ramifié
(Algèbre, définition \ref{algebra-definition-analytically-unramified}) ;
\item $X$ est localement de type fini sur $A$ ;
\item $X \to \Spec(A)$ est étale en tout point de codimension $0$ de $X$.
\end{enumerate}
Alors la normalisation de $X$ est finie sur $X$.
\end{lemma}

\begin{proof}
Choisissons un schéma $U$ et un morphisme étale surjectif $U \to X$.
Alors $U \to \Spec(A)$ satisfait aux hypothèses, donc aux conclusions, de
Résolution des surfaces, lemme \ref{resolve-lemma-formally-unramified}.
\end{proof}











\section{Propriétés qui s'en déduisent}
\label{spaces-resolve-section-existence-gives}

\noindent
Dans cette section, nous montrons que, pour un espace algébrique noethérien
intègre, l'existence d'une altération régulière entraîne de nombreuses
conséquences. Les lecteurs que n'intéressent pas les ``mauvais'' espaces
algébriques noethériens peuvent omettre cette section.

\begin{lemma}
\label{spaces-resolve-lemma-regular-alteration-implies}
Soit $S$ un schéma. Soit $Y$ un espace algébrique noethérien intègre sur $S$.
Supposons qu'il existe une altération $f : X \to Y$, avec $X$ régulier.
Alors la normalisation $Y^\nu \to Y$ est finie et $Y$ possède un ouvert dense
qui est régulier.
\end{lemma}

\begin{proof}
Par localisation étale, il suffit de le démontrer lorsque
$Y = \Spec(A)$, où $A$ est un domaine noethérien.
Soit $B$ la clôture intégrale de $A$ dans son corps des fractions.
Posons $C = \Gamma(X, \mathcal{O}_X)$. D'après
Cohomologie des espaces, lemme
\ref{spaces-cohomology-lemma-proper-pushforward-coherent}
nous voyons que $C$ est un $A$-module fini.
Comme $X$ est normal (Propriétés des espaces, lemme
\ref{spaces-properties-lemma-regular-normal})
nous voyons que $C$ est un domaine normal (Espaces au-dessus de corps, lemme
\ref{spaces-over-fields-lemma-normal-integral-sections}).
Ainsi $B \subset C$, et l'on conclut que $B$ est fini sur $A$ puisque $A$ est
noethérien.

\medskip\noindent
Il existe un ouvert non vide $V \subset Y$ tel que $f^{-1}V \to V$ soit fini ;
voir Espaces au-dessus de corps, définition
\ref{spaces-over-fields-definition-alteration}.
Quitte à restreindre $V$, on peut supposer que $f^{-1}V \to V$ est plat
(Morphismes d'espaces, proposition
\ref{spaces-morphisms-proposition-generic-flatness-reduced}).
Ainsi, $f^{-1}V \to V$ est fidèlement plat. Dès lors, $V$ est régulier d'après
Algèbre, lemme \ref{algebra-lemma-descent-regular}.
\end{proof}

\begin{lemma}
\label{spaces-resolve-lemma-regular-alteration-implies-local}
Soit $(A, \mathfrak m, \kappa)$ un domaine local noethérien.
Supposons qu'il existe une altération $f : X \to \Spec(A)$ avec $X$ régulier.
Alors :
\begin{enumerate}
\item il existe un $f \in A$ non nul tel que $A_f$ soit régulier ;
\item la clôture intégrale $B$ de $A$ dans son corps des fractions est finie
sur $A$ ;
\item le complété $\mathfrak m$-adique de $B$ est un anneau normal, c'est-à-dire
que les complétés de $B$ en ses idéaux maximaux sont des domaines normaux ;
\item la fibre formelle générique de $A$ est régulière.
\end{enumerate}
\end{lemma}

\begin{proof}
Les assertions (1) et (2) résultent du lemme
\ref{spaces-resolve-lemma-regular-alteration-implies}.
Il nous faut reprendre une partie de la démonstration de ce lemme afin de
fixer les notations pour la preuve de (3). Posons
$C = \Gamma(X, \mathcal{O}_X)$. D'après Cohomologie des espaces, lemme
\ref{spaces-cohomology-lemma-proper-pushforward-coherent}
nous voyons que $C$ est un $A$-module fini.
Comme $X$ est normal (Propriétés des espaces, lemme
\ref{spaces-properties-lemma-regular-normal})
nous voyons que $C$ est un domaine normal (Espaces au-dessus de corps, lemme
\ref{spaces-over-fields-lemma-normal-integral-sections}).
Ainsi $B \subset C$, et l'on conclut que $B$ est fini sur $A$ puisque $A$ est
noethérien. D'après Résolution des surfaces, lemme
\ref{resolve-lemma-algebra-helper}, il suffit, pour prouver (3), de montrer que
le complété $\mathfrak m$-adique $C^\wedge$ est normal.

\medskip\noindent
D'après Algèbre, lemme \ref{algebra-lemma-completion-finite-extension},
le complété $C^\wedge$ est le produit des complétés de $C$ en les idéaux
premiers de $C$ qui sont au-dessus de $\mathfrak m$. Ceux-ci sont en nombre
fini et ce sont
les idéaux maximaux $\mathfrak m_1, \ldots, \mathfrak m_r$ de $C$.
(Le résultat correspondant pour $B$ explique la dernière précision de
l'énoncé.) En remplaçant $A$ par $C_{\mathfrak m_i}$ et $X$ par
$X_i = X \times_{\Spec(C)} \Spec(C_{\mathfrak m_i})$, on se ramène donc au
cas du paragraphe suivant. (Notons que
$\Gamma(X_i, \mathcal{O}) = C_{\mathfrak m_i}$ d'après Cohomologie des espaces,
lemme \ref{spaces-cohomology-lemma-flat-base-change-cohomology}.)

\medskip\noindent
Ici, $A$ est un domaine local noethérien normal et $f : X \to \Spec(A)$ est
une altération régulière telle que $\Gamma(X, \mathcal{O}_X) = A$.
Il nous faut montrer que le complété $A^\wedge$ de $A$ est un domaine normal.
D'après le lemme \ref{spaces-resolve-lemma-port-regularity-to-completion},
$Y = X \times_{\Spec(A)} \Spec(A^\wedge)$ est régulier.
Or $\Gamma(Y, \mathcal{O}_Y) = A^\wedge$ d'après Cohomologie des espaces,
lemme \ref{spaces-cohomology-lemma-flat-base-change-cohomology}.
On conclut comme précédemment que $A^\wedge$ est normal.
En effet, $Y$ est normal d'après Propriétés des espaces, lemme
\ref{spaces-properties-lemma-regular-normal}.
Il est connexe parce que $\Gamma(Y, \mathcal{O}_Y) = A^\wedge$ est local.
Ainsi, $Y$ est normal et intègre (pour un espace algébrique séparé, être
connexe et normal entraîne être intègre). Par conséquent,
$\Gamma(Y, \mathcal{O}_Y) = A^\wedge$ est un domaine normal d'après
Espaces au-dessus de corps, lemme
\ref{spaces-over-fields-lemma-normal-integral-sections}.
Ceci démontre (3).

\medskip\noindent
Démontrons (4). Notons $\eta \in \Spec(A)$ le point générique et indiquons
par un indice $\eta$ le changement de base à $\eta$.
Comme $f$ est une altération, le schéma $X_\eta$ est fini et fidèlement plat
sur $\eta$. Puisque $Y = X \times_{\Spec(A)} \Spec(A^\wedge)$ est régulier
d'après le lemme \ref{spaces-resolve-lemma-port-regularity-to-completion}, on voit que
$Y_\eta$ est régulier (comme limite d'ouverts de $Y$).
Le morphisme $Y_\eta \to \Spec(A^\wedge \otimes_A \kappa(\eta))$ est alors
fini et fidèlement plat sur la fibre formelle générique. On conclut par
Algèbre, lemme \ref{algebra-lemma-descent-regular}.
\end{proof}











\section{Résolution}
\label{spaces-resolve-section-resolution}

\noindent
Voici une définition.

\begin{definition}
\label{spaces-resolve-definition-resolution}
Soit $S$ un schéma. Soit $Y$ un espace algébrique noethérien intègre sur $S$.
Une {\it résolution des singularités} de $X$ est une modification
$f : X \to Y$ telle que $X$ soit régulier.
\end{definition}

\noindent
Dans le cas des surfaces, nous souhaitons parfois disposer d'un peu plus
d'information.

\begin{definition}
\label{spaces-resolve-definition-resolution-surface}
Soit $S$ un schéma. Soit $Y$ un espace algébrique noethérien intègre de
dimension $2$ sur $S$. On dit que $Y$ admet une
{\it résolution des singularités par éclatements normalisés}
s'il existe une suite
$$
Y_n \to X_{n - 1} \to \ldots \to Y_1 \to Y_0 \to Y
$$
où :
\begin{enumerate}
\item $Y_i$ est propre sur $Y$ pour $i = 0, \ldots, n$ ;
\item $Y_0 \to Y$ est la normalisation ;
\item $Y_i \to Y_{i - 1}$ est un éclatement normalisé pour
$i = 1, \ldots, n$ ;
\item $Y_n$ est régulier.
\end{enumerate}
\end{definition}

\noindent
Remarquons que la condition (1) implique que la normalisation $Y_0$ de $Y$
est finie sur $Y$ et que les normalisations intervenant dans les éclatements
normalisés sont elles aussi finies.
Nous arrivons enfin au théorème principal de ce chapitre.

\begin{theorem}
\label{spaces-resolve-theorem-resolve}
Soit $S$ un schéma. Soit $Y$ un espace algébrique noethérien intègre de
dimension deux sur $S$. Les conditions suivantes sont équivalentes :
\begin{enumerate}
\item il existe une altération $X \to Y$ avec $X$ régulier ;
\item il existe une résolution des singularités de $Y$ ;
\item $Y$ admet une résolution des singularités par éclatements normalisés ;
\item la normalisation $Y^\nu \to Y$ est finie, $Y^\nu$ possède un nombre
fini de points singuliers $y_1, \ldots, y_m \in |Y|$, et, pour chaque $i$,
le complété de l'anneau local hensélien $\mathcal{O}_{Y^\nu, y_i}^h$ est
normal.
\end{enumerate}
\end{theorem}

\begin{proof}
Les implications (3) $\Rightarrow$ (2) $\Rightarrow$ (1) sont immédiates.

\medskip\noindent
Soit $X \to Y$ une altération avec $X$ régulier. Alors $Y^\nu \to Y$ est
fini d'après le lemme \ref{spaces-resolve-lemma-regular-alteration-implies}.
Considérons la factorisation $f : X \to Y^\nu$ de
Morphismes d'espaces, lemme \ref{spaces-morphisms-lemma-normalization-normal}.
D'après Espaces au-dessus de corps, lemme
\ref{spaces-over-fields-lemma-finite-in-codim-1}, le morphisme $f$ est fini
au-dessus d'un ouvert $V \subset Y^\nu$ contenant tout point de codimension
$\leq 1$ de $Y^\nu$. Alors $f$ est plat au-dessus de $V$ d'après
Algèbre, lemme \ref{algebra-lemma-CM-over-regular-flat}, et parce qu'un anneau
local normal de dimension $\leq 2$ est de Cohen--Macaulay par le critère de
Serre (Algèbre, lemme \ref{algebra-lemma-criterion-normal}).
L'ouvert $V$ est donc régulier d'après Algèbre, lemme
\ref{algebra-lemma-descent-regular}.
Comme $Y^\nu$ est noethérien, on conclut que
$Y^\nu \setminus V = \{y_1, \ldots, y_m\}$ est fini.
Pour chaque $i$, soit $\mathcal{O}_{Y^\nu, y_i}^h$ l'anneau local hensélien.
Alors $X \times_Y \Spec(\mathcal{O}_{Y^\nu, y_i}^h)$ est une altération
régulière de $\Spec(\mathcal{O}_{Y^\nu, y_i}^h)$
(nous omettons quelques détails).
D'après le lemme \ref{spaces-resolve-lemma-regular-alteration-implies-local}, le complété de
$\mathcal{O}_{Y^\nu, y_i}^h$ est normal. On obtient ainsi
(1) $\Rightarrow$ (4).

\medskip\noindent
Supposons (4). Il nous faut démontrer (3). Nous pouvons immédiatement
remplacer $Y$ par sa normalisation. Soient
$y_1, \ldots, y_m \in |Y|$ les points singuliers. Choisissons des voisinages
étales élémentaires $(V_i, v_i) \to (Y, y_i)$ comme dans la section
\ref{spaces-resolve-section-strategy}. Pour chaque $i$, l'anneau local hensélien
$\mathcal{O}_{Y^\nu, y_i}^h$ est le hensélisé de
$\mathcal{O}_{V_i, v_i}$. Ces anneaux ont donc des complétés isomorphes.
Le résultat pour les schémas (Résolution des surfaces, théorème
\ref{resolve-theorem-resolve}) montre alors qu'il existe des suites finies
d'éclatements normalisés
$$
X_{i, n_i} \to X_{i, n_i - 1} \to \ldots \to V_i
$$
dont les centres sont uniquement des points situés au-dessus de $v_i$, telles
que $X_{i, n_i}$ soit régulier. D'après le lemme
\ref{spaces-resolve-lemma-equivalence-sequence-normalized-blowups}, il existe une suite
d'éclatements normalisés
$$
X_n \to X_{n - 1} \to \ldots \to X_1 \to Y
$$
et, bien entendu, $X_n$ est lui aussi régulier (il suffit de considérer les
anneaux locaux). Cela achève la démonstration.
\end{proof}















\section{Exemples}
\label{spaces-resolve-section-examples}

\noindent
Voici quelques exemples liés aux résultats précédents de ce chapitre.

\begin{example}
\label{spaces-resolve-example-factorial}
\begin{reference}
\cite[4(c)]{Samuel-UFD}
\end{reference}
Soit $k$ un corps. L'anneau $A = k[x, y, z]/(x^r + y^s + z^t)$ est factoriel
lorsque $r, s, t$ sont des entiers deux à deux premiers entre eux.
En effet, puisque $x^r + y^s + z^t$ est irréductible, $A$ est un domaine.
L'élément $z$ est premier, c'est-à-dire qu'il engendre un idéal premier de
$A$. D'autre part, si $t = 1 + ers$ pour un certain $e$, alors
$$
A[1/z] \cong k[x', y', 1/z]
$$
où $x' = x/z^{es}$, $y' = y/z^{er}$ et $z = (x')^r + (y')^s$.
Ainsi, $A[1/z]$ est un localisé d'un anneau de polynômes, et il est donc
factoriel. Un argument de Nagata montre alors que $A$ est factoriel.
Voir Algèbre, lemme \ref{algebra-lemma-invert-prime-elements}.
Un argument analogue s'applique lorsque $t$ n'est pas congru à $1$ modulo
$rs$.
\end{example}

\begin{example}
\label{spaces-resolve-example-completion-not-factorial}
\begin{reference}
Voir \cite{Brieskorn} et \cite{Lipman-rational} pour la non-annulation, en
général, des groupes de Picard locaux.
\end{reference}
L'anneau $A = \mathbf{C}[[x, y, z]]/(x^r + y^s + z^t)$ n'est pas factoriel
lorsque $1 < r < s < t$ sont des entiers deux à deux premiers entre eux qui
ne sont pas égaux à $2, 3, 5$. Considérons par exemple le cas particulier
$A = \mathbf{C}[[x, y, z]]/(x^2 + y^5 + z^7)$.
Considérons les homomorphismes
$$
\psi_\zeta : \mathbf{C}[[x, y, z]]/(x^2 + y^5 + z^7) \to \mathbf{C}[[t]]
$$
donnés par
$$
x \mapsto t^7,\quad
y \mapsto t^3,\quad
z \mapsto -\zeta t^2(1 + t)^{1/7}
$$
où $\zeta$ est une racine $7$-ième de l'unité. Le noyau
$\mathfrak p_\zeta$ de $\psi_\zeta$ est un idéal premier de hauteur un ; si
$A$ était factoriel, il serait donc principal, engendré, disons, par
$f_\zeta \in \mathbf{C}[[x, y, z]]$.
Notons que $V(x^3 - y^7) = \bigcup V(\mathfrak p_\zeta)$ et que
$A/(x^3 - y^7)$ est réduit hors du point fermé. Par conséquent, toujours sous
l'hypothèse que $A$ est factoriel, on obtiendrait
$$
\prod\nolimits_\zeta f_\zeta = u(x^3 - y^7) + a(x^2 + y^5 + z^7)
\quad\text{in}\quad
\mathbf{C}[[x, y, z]]
$$
pour une unité $u \in \mathbf{C}[[x, y, z]]$ et un élément
$a \in \mathbf{C}[[x, y, z]]$. Après multiplication par une constante, on
peut supposer que $u(0, 0, 0) = 1$. Notons que le membre de gauche s'annule à
l'ordre $7$. Ainsi $a = - x \bmod \mathfrak m^2$. Mais il apparaît alors dans
le membre de droite un terme $xy^5$ qui n'apparaît pas dans celui de gauche.
C'est une contradiction.
\end{example}

\begin{example}
\label{spaces-resolve-example-not-blow-up}
Il existe un anneau local noethérien excellent de dimension $2$ et une
modification $X \to S = \Spec(A)$ qui n'est pas un schéma.
Esquissons une construction. Soit $X$ une surface normale sur $\mathbf{C}$
ayant un unique point singulier $x \in X$. Supposons qu'il existe une
résolution $\pi : X' \to X$ telle que la fibre exceptionnelle
$C = \pi^{-1}(x)_{red}$ soit une courbe projective lisse. Supposons en outre
qu'il existe un point $c \in C$ tel que, si $\mathcal{O}_C(nc)$ appartient à
l'image de $\Pic(X') \to \Pic(C)$, alors $n = 0$.
Soit ensuite $X'' \to X'$ l'éclatement au point non singulier $c$.
Notons $C' \subset X''$ la transformée stricte de $C$ et $E \subset X''$ la
fibre exceptionnelle. D'après les résultats d'Artin
(\cite{ArtinII} ; on peut utiliser, par exemple, \cite{Mumford-topology} pour
voir que le fibré normal de $C'$ est négatif), on peut contracter la courbe
$C'$ dans $X''$ et obtenir un espace algébrique $X'''$. Voici le diagramme :
$$
\xymatrix{
& X'' \ar[ld] \ar[rd] \\
X' \ar[rd] &  & X''' \ar[ld] \\
& X
}
$$
Nous affirmons que $X'''$ n'est pas un schéma. Cela fournit notre exemple,
car $X'''$ est un schéma si et seulement si le changement de base de $X'''$
à $A = \mathcal{O}_{X, x}$ est un schéma (nous omettons les détails).
Si $X'''$ était un schéma, l'image de $C'$ dans $X'''$ posséderait un voisinage
affine. Le complémentaire de ce voisinage serait un diviseur de Cartier
effectif sur $X'''$ (car $X'''$ est non singulier à l'exception de $1$ point).
Ce diviseur de Cartier effectif correspondrait à un diviseur de Cartier
effectif sur $X''$ qui rencontre $E$ et évite $C'$. En prenant son image dans
$X'$, on obtiendrait un diviseur de Cartier effectif qui rencontre $C$, au sens
ensembliste, en $c$. C'est impossible, puisque, par hypothèse, aucun multiple
de $c$ n'est la restriction d'un diviseur de Cartier.

\medskip\noindent
Pour terminer, il nous faut trouver une telle surface singulière $X$. On peut
simplement prendre pour $X$ la surface affine définie par
$$
x^3 + y^3 + z^3 + x^4 + y^4 + z^4 = 0
$$
dans $\mathbf{A}^3_\mathbf{C} = \Spec(\mathbf{C}[x, y, z])$, le point
singulier étant $(0, 0, 0)$. Alors $(0, 0, 0)$ est l'unique point singulier.
En éclatant $X$ le long de l'idéal maximal correspondant à $(0, 0, 0)$, on
obtient trois cartes, chacune isomorphe à la surface affine lisse
$$
1 + s^3 + t^3 + x(1 + s^4 + t^4) = 0
$$
qui est non singulière et dont le diviseur exceptionnel $C$ est défini par
$x = 0$. Le lecteur reconnaîtra en $C$ une courbe elliptique. Enfin, la
surface $X$ est rationnelle, comme le montre la projection depuis $(0, 0, 0)$,
ou encore parce que l'équation de l'éclatement permet de résoudre en $x$.
Le groupe de Picard d'une surface rationnelle non singulière est dénombrable,
tandis que celui d'une courbe elliptique sur les nombres complexes ne l'est
pas. On peut donc choisir un point fermé $c$ comme indiqué.
\end{example}













% OK, start here.
%

\chapter{Théorie des déformations formelles}



\phantomsection
\label{formal-defos-section-phantom}




\section{Introduction}
\label{formal-defos-section-introduction}

\noindent
Ce chapitre développe la théorie des déformations formelles sous une forme
applicable plus loin dans le projet Champs, en suivant de près Rim
\cite[Exposee VI]{SGA7-I} et Schlessinger \cite{Sch}. Nous recommandons vivement
au lecteur qui découvre ce sujet de lire d'abord l'article de Schlessinger :
il est suffisamment général pour la plupart des applications, et les résultats
de Schlessinger sont effectivement utilisés dans la plupart des articles
qui font appel à ce type de théorie des déformations formelles.

\medskip\noindent
Soit $\Lambda$ un anneau local noethérien complet de corps résiduel $k$,
et notons $\mathcal{C}_\Lambda$ la catégorie des $\Lambda$-algèbres
artiniennes locales de corps résiduel $k$. Étant donné un foncteur
$F : \mathcal{C}_\Lambda \to \textit{Ensembles}$ tel que $F(k)$
soit un ensemble à un élément, l'article de Schlessinger introduit des
conditions (H1)--(H4) telles que :
\begin{enumerate}
\item $F$ admet une « enveloppe » si et seulement si (H1)--(H3) sont satisfaites.
\item $F$ est pro-représentable si et seulement si (H1)--(H4) sont satisfaites.
\end{enumerate}
Le but de ce chapitre est de généraliser ces résultats de deux manières,
exactement comme dans l'article de Rim :
\begin{enumerate}
\item[(A)] Le foncteur $F$ est remplacé par une catégorie $\mathcal{F}$
cofibrée en groupoïdes sur $\mathcal{C}_\Lambda$ ; voir la
section \ref{formal-defos-section-CLambda}.
\item[(B)] On prend pour $\Lambda$ un anneau noethérien et pour
$\Lambda \to k$ un homomorphisme fini d'anneaux à valeurs dans un corps.
La catégorie $\mathcal{C}_\Lambda$ est celle des $\Lambda$-algèbres
artiniennes locales $A$ munies d'une identification donnée
$A/\mathfrak m_A = k$.
\end{enumerate}
L'analogue de la condition selon laquelle $F(k)$ est un ensemble à un élément
est que $\mathcal{F}(k)$ soit le groupoïde trivial. Si $\mathcal{F}$ satisfait
cette condition, on dit que c'est une {\it catégorie de prédéformation}, mais
on ne fait pas cette hypothèse en général. L'article de Rim
\cite[Exposee VI]{SGA7-I} est la source originale des résultats de ce document.
Mentionnons aussi l'utile article \cite{Talpo-Vistoli}, qui traite de la théorie
des déformations avec des groupoïdes, mais dans un cadre moins général qu'ici.

\medskip\noindent
Le « complété » $\widehat{\mathcal{C}}_\Lambda$ de la catégorie
$\mathcal{C}_\Lambda$ joue un rôle important. Un objet de
$\widehat{\mathcal{C}}_\Lambda$ est une $\Lambda$-algèbre locale
noethérienne complète $R$ dont le corps résiduel est identifié à $k$ ; voir la
section \ref{formal-defos-section-category-completion-CLambda}.
D'une part, $\mathcal{C}_\Lambda \subset \widehat{\mathcal{C}}_\Lambda$
est une sous-catégorie strictement pleine et, d'autre part,
$\widehat{\mathcal{C}}_\Lambda$ est une sous-catégorie pleine de la catégorie
des pro-objets de $\mathcal{C}_\Lambda$. Un foncteur
$\mathcal{C}_\Lambda \to \textit{Ensembles}$ est {\it pro-représentable}
s'il est isomorphe à la restriction d'un foncteur représentable
$\underline{R} = \Mor_{\widehat{\mathcal{C}}_\Lambda}(R, -)$
à $\mathcal{C}_\Lambda$, où
$R \in \Ob(\widehat{\mathcal{C}}_\Lambda)$.

\medskip\noindent
Les {\it catégories cofibrées en groupoïdes} sont duales des catégories
fibrées en groupoïdes ; nous les introduisons à la section
\ref{formal-defos-section-preliminary}. Un morphisme {\it lisse} de catégories cofibrées
en groupoïdes sur $\mathcal{C}_\Lambda$ est un morphisme qui satisfait le
critère infinitésimal de relèvement pour les objets ; voir la
section \ref{formal-defos-section-smooth-morphisms}.
Cette définition est analogue à celle d'un homomorphisme d'anneaux
formellement lisse ; voir Algèbre, définition
\ref{algebra-definition-formally-smooth}. Elle est exactement duale de la
notion de Critères de représentabilité, section
\ref{criteria-section-formally-smooth}. Cette notion est importante, car nous
voulons finalement démontrer que certaines catégories cofibrées en groupoïdes
admettent une présentation lisse pro-représentable, à l'instar de la
caractérisation des champs algébriques dans Champs algébriques, sections
\ref{algebraic-section-stack-to-presentation} et
\ref{algebraic-section-smooth-groupoid-gives-algebraic-stack}.
Un {\it objet formel versel} d'une catégorie $\mathcal{F}$ cofibrée en
groupoïdes sur $\mathcal{C}_\Lambda$ est un objet
$\xi \in \widehat{\mathcal{F}}(R)$ du complété tel que le morphisme
associé
$\underline{\xi} : \underline{R}|_{\mathcal{C}_\Lambda} \to \mathcal{F}$
soit lisse.

\medskip\noindent
Dans la section \ref{formal-defos-section-schlessinger-conditions}, nous définissons sur
$\mathcal{F}$ les conditions (S1) et (S2), qui généralisent les conditions
(H1) et (H2) de Schlessinger. L'analogue de la condition (H3) de
Schlessinger---à savoir que l'espace tangent de $\mathcal{F}$ est de dimension
finie---ne reçoit pas de nom. Une étape essentielle du développement de la
théorie consiste à établir l'existence d'objets formels versels pour les
catégories de prédéformation qui satisfont (S1), (S2) et (H3) ; voir le
lemme \ref{formal-defos-lemma-versal-object-existence}.
La notion d'{\it enveloppe} d'un foncteur, au sens de Schlessinger,
$F : \mathcal{C}_\Lambda \to \textit{Ensembles}$
est, dans notre terminologie, un objet formel versel
$\xi \in \widehat{F}(R)$ tel que l'application induite sur les espaces tangents
$d\underline{\xi} : T\underline{R}|_{\mathcal{C}_\Lambda} \to TF$
soit un isomorphisme. Dans la littérature, une enveloppe est souvent appelée
objet « miniversel ». Nous n'employons pas ce terme, pour la raison suivante.
Il peut arriver qu'un foncteur possède un objet formel versel sans posséder
d'enveloppe. De plus, nous montrons à la section
\ref{formal-defos-section-minimal-versal} que, si une catégorie de prédéformation possède
un objet formel versel, elle en possède toujours un qui est {\it minimal}
(au sens de la définition \ref{formal-defos-definition-minimal-versal}) et qui est unique
à isomorphisme près ; voir le lemme \ref{formal-defos-lemma-minimal-versal}.
Il peut cependant arriver que l'objet formel versel minimal n'induise pas
un isomorphisme sur les espaces tangents ! (Voir les exemples
\ref{formal-defos-example-do-not-get-S2} et
\ref{formal-defos-example-smooth-continued}.)

\medskip\noindent
Compte tenu des différences signalées ci-dessus, le théorème
\ref{formal-defos-theorem-miniversal-object-existence} est la généralisation directe de
(1) : il redonne le résultat de Schlessinger lorsque $\mathcal{F}$ est un
foncteur et, sous les conditions (S1) et (S2), il caractérise les objets
formels versels minimaux au moyen de l'application
$d\underline{\xi} : T\underline{R}|_{\mathcal{C}_\Lambda} \to TF$
entre espaces tangents.

\medskip\noindent
À la section \ref{formal-defos-section-RS-condition}, nous définissons sur $\mathcal{F}$
la condition (RS) de Rim, qui généralise la condition (H4) de Schlessinger.
Une {\it catégorie de déformation} est, par définition, une catégorie de
prédéformation satisfaisant (RS). Les analogues des foncteurs
pro-représentables sont les catégories cofibrées en groupoïdes sur
$\mathcal{C}_\Lambda$ qui admettent une {\it présentation par un groupoïde
en foncteurs lisse et pro-représentable} sur $\mathcal{C}_\Lambda$ ; voir les
définitions \ref{formal-defos-definition-groupoid-in-functors},
\ref{formal-defos-definition-prorepresentable-groupoid-in-functors}, et
\ref{formal-defos-definition-smooth-groupoid-in-functors}.
Cette notion de présentation tient compte de la structure de groupoïde des
fibres de $\mathcal{F}$. Dans le théorème
\ref{formal-defos-theorem-presentation-deformation-groupoid}, nous démontrons que
$\mathcal{F}$ admet une présentation par un groupoïde en foncteurs lisse et
pro-représentable si et seulement si $\mathcal{F}$ possède un espace tangent et un espace
d'automorphismes infinitésimaux sont de dimension finie. C'est la
généralisation de (2) ci-dessus : elle se réduit au résultat de Schlessinger
lorsque $\mathcal{F}$ est un foncteur. Enfin, la section
\ref{formal-defos-section-minimality} explique comment utiliser les objets formels versels
minimaux pour construire une présentation minimale (unique à isomorphisme
près) par un groupoïde en foncteurs lisse et pro-représentable.

\medskip\noindent
Nous obtenons également l'explication conceptuelle suivante des conditions
de Schlessinger. Si une catégorie de prédéformation $\mathcal{F}$ satisfait
(RS), alors le foncteur associé des classes d'isomorphie
$\overline{\mathcal{F}}: \mathcal{C}_\Lambda \to \textit{Ensembles}$
satisfait (H1) et (H2) (lemmes \ref{formal-defos-lemma-RS-implies-S1-S2} et
\ref{formal-defos-lemma-S1-S2-associated-functor}).
Réciproquement, si un foncteur $F : \mathcal{C}_\Lambda \to \textit{Ensembles}$
se présente naturellement comme le foncteur des classes d'isomorphie d'une
catégorie $\mathcal{F}$ cofibrée en groupoïdes, on constate en pratique qu'un
argument établissant que $F$ satisfait (H1) et (H2) établit également que
$\mathcal{F}$ satisfait (RS). Des exemples sont étudiés dans Problèmes de
déformation, section
\ref{examples-defos-section-introduction}.
De plus, si $\mathcal{F}$ satisfait (RS), la condition (H4) pour
$\overline{\mathcal{F}}$ admet une interprétation simple en termes de
prolongement des automorphismes des objets de $\mathcal{F}$
(lemme \ref{formal-defos-lemma-RS-associated-functor}). Ces observations suggèrent de
considérer (RS) comme la condition fondamentale de recollement de la théorie
des déformations.




\section{Notations et conventions}
\label{formal-defos-section-notations-conventions}

\noindent
Les anneaux sont commutatifs et possèdent un élément unité $1$. L'idéal maximal d'un anneau local
$A$ est noté $\mathfrak{m}_A$. L'ensemble des entiers strictement positifs est
noté $\mathbf{N} = \{1, 2, 3, \ldots\}$. Si $U$ est un objet d'une catégorie
$\mathcal{C}$, nous notons $\underline{U}$ le foncteur
$\Mor_\mathcal{C}(U, -): \mathcal{C} \to \textit{Ensembles}$ ; voir les remarques
\ref{formal-defos-remarks-cofibered-groupoids} (\ref{formal-defos-item-definition-yoneda}). Attention :
cette convention peut entrer en conflit avec la notation d'autres chapitres,
où nous employons parfois $\underline{U}$ pour désigner
$h_U(-) = \Mor_\mathcal{C}(-, U)$.

\medskip\noindent
Dans tout ce chapitre, $\Lambda$ est un anneau noethérien et
$\Lambda \to k$ est un homomorphisme fini de $\Lambda$ vers un corps.
Le noyau de cet homomorphisme est noté $\mathfrak m_\Lambda$ et son image
$k' \subset k$. On verra que $\mathfrak m_\Lambda$ est un idéal maximal,
que $k' = \Lambda/\mathfrak m_\Lambda$ est un corps et que l'extension
$k/k'$ est finie. Voir la discussion autour de
(\ref{formal-defos-equation-k-prime}).


\section{La catégorie de base}
\label{formal-defos-section-CLambda}

\noindent
Motivation. Une application importante de la théorie des déformations formelles
concerne les critères de représentabilité par des espaces algébriques.
Supposons donnés une base localement noethérienne $S$ et un foncteur
$F : (\Sch/S)_{fppf}^{opp} \to \textit{Ensembles}$.
Soit $k$ un corps de type fini sur $S$, c'est-à-dire que l'on se donne un
morphisme de type fini $\Spec(k) \to S$.
L'un des critères d'Artin exige que, pour tout élément
$x \in F(\Spec(k))$, le foncteur de prédéformation associé au triplet
$(S, k, x)$ soit pro-représentable. D'après Morphismes, lemme
\ref{morphisms-lemma-point-finite-type}, dire que $k$ est de type fini sur
$S$ signifie qu'il existe un ouvert affine $\Spec(\Lambda) \subset S$ tel que
$k$ soit une $\Lambda$-algèbre finie. Cela explique le choix, dans tout ce
chapitre, de la catégorie de base suivante.

\begin{definition}
\label{formal-defos-definition-CLambda}
Soit $\Lambda$ un anneau noethérien et soit $\Lambda \to k$ un homomorphisme
fini d'anneaux, où $k$ est un corps. On définit {\it $\mathcal{C}_\Lambda$}
comme la catégorie dont
\begin{enumerate}
\item les objets sont les couples $(A, \varphi)$, où $A$ est une
$\Lambda$-algèbre artinienne locale et où
$\varphi : A/\mathfrak m_A \to k$ est un isomorphisme de
$\Lambda$-algèbres ;
\item les morphismes $f : (B, \psi) \to (A, \varphi)$ sont les
homomorphismes locaux de $\Lambda$-algèbres tels que
$\varphi \circ (f \bmod \mathfrak m) = \psi$.
\end{enumerate}
On dit que l'on est dans le {\it cas classique} si $\Lambda$ est un anneau
local noethérien complet et si $k$ est son corps résiduel.
\end{definition}

\noindent
Remarquons que, si $\Lambda \to k$ est surjectif et si $A$ est une
$\Lambda$-algèbre artinienne locale, l'identification $\varphi$, lorsqu'elle
existe, est unique. De plus, dans ce cas, tout homomorphisme de
$\Lambda$-algèbres $A \to B$ est compatible avec les identifications. Ainsi,
$\mathcal{C}_\Lambda$ n'est alors que la catégorie des $\Lambda$-algèbres
artiniennes locales dont le corps résiduel « est » $k$. Par abus de notation,
dans le cas général, les objets de $\mathcal{C}_\Lambda$ seront aussi désignés
simplement par $A$. Nous écrirons en outre souvent
$A/\mathfrak m = k$, autrement dit nous feindrons que tous les anneaux de
$\mathcal{C}_\Lambda$ ont pour corps résiduel $k$ (comme tous les
homomorphismes d'anneaux de $\mathcal{C}_\Lambda$ sont compatibles avec les
identifications données, cela ne devrait jamais poser de difficulté).
Dans toute la suite du chapitre, l'anneau de base $\Lambda$ et le corps $k$
sont fixés. La catégorie $\mathcal{C}_\Lambda$ sera la catégorie de base des
catégories cofibrées considérées ci-dessous.

\begin{definition}
\label{formal-defos-definition-small-extension}
Soit $f: B \to A$ un homomorphisme d'anneaux dans $\mathcal{C}_\Lambda$.
On dit que $f$ est une {\it petite extension} s'il est surjectif et si
$\Ker(f)$ est un idéal principal non nul, annulé par $\mathfrak{m}_B$.
\end{definition}

\noindent
Le lemme suivant permet souvent de ramener les arguments portant sur des
homomorphismes d'anneaux surjectifs dans $\mathcal{C}_\Lambda$ au cas des
petites extensions.

\begin{lemma}
\label{formal-defos-lemma-factor-small-extension}
Soit $f: B \to A$ un homomorphisme d'anneaux surjectif dans
$\mathcal{C}_\Lambda$. Alors $f$ se décompose en un composé de petites
extensions.
\end{lemma}

\begin{proof}
Soit $I$ le noyau de $f$. L'idéal maximal $\mathfrak{m}_B$ est nilpotent,
car $B$ est artinien ; écrivons $\mathfrak{m}_B^n = 0$. On obtient donc une
décomposition
$$
B = B/I\mathfrak{m}_B^{n-1} \to B/I\mathfrak{m}_B^{n-2} \to
\ldots \to B/I \cong A
$$
de $f$ en un composé d'homomorphismes surjectifs dont les noyaux sont annulés
par l'idéal maximal. Il suffit donc de démontrer le lemme lorsque $f$ est lui-même
un tel homomorphisme, c'est-à-dire lorsque $I$ est annulé par
$\mathfrak{m}_B$. Dans ce cas, $I$ est un espace vectoriel sur $k$, de
dimension finie ; voir Algèbre, lemme
\ref{algebra-lemma-artinian-finite-length}. En choisissant une base
$x_1, \ldots, x_n$ de l'espace vectoriel $I$ sur $k$, on obtient une
décomposition
$$
B \to B/(x_1) \to \ldots \to  B/(x_1, \ldots, x_n) \cong  A
$$
de $f$ en un composé de petites extensions.
\end{proof}

\noindent
Le lemme suivant affirme que l'on peut calculer la longueur d'un module sur
une $\Lambda$-algèbre locale de corps résiduel $k$ en fonction de sa longueur
sur $\Lambda$. Pour expliquer la notation de l'énoncé, soit $k' \subset k$
l'image de l'homomorphisme fini fixé $\Lambda \to k$. Remarquons que
$k' \subset k$ est une extension finie d'anneaux. Par conséquent, $k'$ est un
corps et $k/k'$ une extension finie de corps ; voir Algèbre, lemme
\ref{algebra-lemma-integral-under-field}. De plus, puisque
$\Lambda \to k'$ est surjectif, son noyau est un idéal maximal
$\mathfrak m_\Lambda$. Ainsi
\begin{equation}
\label{formal-defos-equation-k-prime}
[k : k'] = [k : \Lambda/\mathfrak m_\Lambda] < \infty
\end{equation}
et, dans le cas classique, on a $k = k'$. La notation
$k' = \Lambda/\mathfrak m_\Lambda$ sera fixée dans tout ce chapitre.

\begin{lemma}
\label{formal-defos-lemma-length}
Soit $A$ une $\Lambda$-algèbre locale de corps résiduel $k$.
Soit $M$ un $A$-module. Alors
$[k : k'] \text{length}_A(M) = \text{length}_\Lambda(M)$.
Dans le cas classique, on a
$\text{length}_A(M) = \text{length}_\Lambda(M)$.
\end{lemma}

\begin{proof}
Si $M$ est un $A$-module simple, alors $M \cong k$ comme $A$-module ; voir
Algèbre, lemme \ref{algebra-lemma-characterize-length-1}.
Dans ce cas, $\text{length}_A(M) = 1$ et
$\text{length}_\Lambda(M) = [k' : k]$ ; voir
Algèbre, lemme \ref{algebra-lemma-dimension-is-length}.
Si $\text{length}_A(M)$ est finie, le résultat s'obtient en choisissant une
filtration de $M$ par des sous-$A$-modules à quotients simples, puis en
utilisant l'additivité ; voir Algèbre, lemme
\ref{algebra-lemma-length-additive}. Si $\text{length}_A(M)$ est infinie, le
résultat découle de l'inégalité évidente
$\text{length}_A(M) \leq \text{length}_\Lambda(M)$.
\end{proof}

\begin{lemma}
\label{formal-defos-lemma-surjective}
Soit $A \to B$ un homomorphisme d'anneaux dans $\mathcal{C}_\Lambda$.
Les conditions suivantes sont équivalentes :
\begin{enumerate}
\item $f$ est surjectif ;
\item $\mathfrak m_A/\mathfrak m_A^2 \to \mathfrak m_B/\mathfrak m_B^2$
est surjectif ;
\item $\mathfrak m_A/(\mathfrak m_\Lambda A + \mathfrak m_A^2)
\to \mathfrak m_B/(\mathfrak m_\Lambda B + \mathfrak m_B^2)$ est surjectif.
\end{enumerate}
\end{lemma}

\begin{proof}
Pour tout homomorphisme d'anneaux $f : A \to B$ dans
$\mathcal{C}_\Lambda$, on a $f(\mathfrak m_A) \subset \mathfrak m_B$ ; cela
résulte par exemple du fait que $\mathfrak m_A$ et $\mathfrak m_B$ sont les
ensembles des éléments nilpotents de $A$ et de $B$. Supposons $f$ surjectif.
Soit $y \in \mathfrak m_B$. Choisissons $x \in A$ tel que $f(x) = y$.
Comme $f$ induit un isomorphisme
$A/\mathfrak m_A \to B/\mathfrak m_B$, on a $x \in \mathfrak m_A$.
L'homomorphisme induit
$\mathfrak m_A/\mathfrak m_A^2 \to \mathfrak m_B/\mathfrak m_B^2$
est donc surjectif. Ainsi, (1) implique (2).

\medskip\noindent
Il est clair que (2) implique (3). L'homomorphisme $A \to B$ donne le
diagramme commutatif canonique
$$
\xymatrix{
\mathfrak m_\Lambda/\mathfrak m_\Lambda^2 \otimes_{k'} k \ar[r] \ar[d] &
\mathfrak m_A/\mathfrak m_A^2 \ar[r] \ar[d] &
\mathfrak m_A/(\mathfrak m_\Lambda A + \mathfrak m_A^2) \ar[r] \ar[d] & 0 \\
\mathfrak m_\Lambda/\mathfrak m_\Lambda^2 \otimes_{k'} k \ar[r] &
\mathfrak m_B/\mathfrak m_B^2 \ar[r] &
\mathfrak m_B/(\mathfrak m_\Lambda B + \mathfrak m_B^2) \ar[r] & 0
}
$$
dont les lignes sont exactes. Ainsi, si (3) est satisfaite, (2) l'est aussi.

\medskip\noindent
Supposons (2). Pour montrer que $A \to B$ est surjectif, il suffit, d'après
le lemme de Nakayama (Algèbre, lemme \ref{algebra-lemma-NAK}), de montrer que
$A/\mathfrak m_A \to B/\mathfrak m_AB$ est surjectif. (Remarquons que
$\mathfrak m_A$ est un idéal nilpotent.) Comme
$k = A/\mathfrak m_A = B/\mathfrak m_B$, il suffit de montrer que
$\mathfrak m_AB \to \mathfrak m_B$ est surjectif. En appliquant une nouvelle
fois le lemme de Nakayama, on voit qu'il suffit que
$\mathfrak m_AB/\mathfrak m_A\mathfrak m_B \to \mathfrak m_B/\mathfrak m_B^2$
soit surjectif, ce qui est précisément l'hypothèse.
\end{proof}

\noindent
Si $A \to B$ est un homomorphisme d'anneaux dans $\mathcal{C}_\Lambda$, alors
l'homomorphisme
$\mathfrak m_A/(\mathfrak m_\Lambda A + \mathfrak m_A^2)
\to \mathfrak m_B/(\mathfrak m_\Lambda B + \mathfrak m_B^2)$
est l'homomorphisme induit sur les espaces cotangents relatifs. En voici la
définition formelle.

\begin{definition}
\label{formal-defos-definition-tangent-space-ring}
Soit $R \to S$ un homomorphisme local d'anneaux locaux. L'{\it espace
cotangent relatif}\footnote{Attention : nous verrons plus loin que, dans
notre cadre général, l'espace tangent d'un objet
$A \in \mathcal{C}_\Lambda$ sur $\Lambda$ ne doit pas être défini simplement
comme le dual $k$-linéaire de l'espace cotangent relatif. En réalité, la bonne
définition de l'espace cotangent relatif est
$\Omega_{S/R} \otimes_S S/\mathfrak m_S$.} de $R$ sur $S$ est l'espace
vectoriel sur $S/\mathfrak m_S$
$\mathfrak m_S/(\mathfrak m_R S + \mathfrak m_S^2)$.
\end{definition}

\noindent
Si $f_1: A_1 \to A$ et $f_2: A_2 \to A$ sont deux homomorphismes d'anneaux,
le produit fibré $A_1 \times_A A_2$ est le sous-anneau de $A_1 \times A_2$
formé des éléments dont les deux projections dans $A$ sont égales. Dans tout
ce chapitre, nous considérerons des conditions faisant intervenir un tel
produit fibré lorsque $f_1$ et $f_2$ appartiennent à
$\mathcal{C}_\Lambda$. Le produit fibré n'est pas toujours un objet de
$\mathcal{C}_\Lambda$.

\begin{example}
\label{formal-defos-example-fibre-product}
Soient $p$ un nombre premier et $n \in \mathbf{N}$.
Posons $\Lambda = \mathbf{F}_p(t_1, t_2, \ldots, t_n)$ et
$k = \mathbf{F}_p(x_1, \ldots, x_n)$, l'homomorphisme $\Lambda \to k$ étant
donné par $t_i \mapsto x_i^p$. Posons $A = k[\epsilon] = k[x]/(x^2)$.
Alors $A$ est un objet de $\mathcal{C}_\Lambda$. Supposons que
$D : k \to k$ soit une dérivation de $k$ sur $\Lambda$, par exemple
$D = \partial/\partial x_i$. Alors l'homomorphisme
$$
f_D : k \longrightarrow k[\epsilon], \quad
a \mapsto a + D(a)\epsilon
$$
est un morphisme de $\mathcal{C}_\Lambda$. Posons $A_1 = A_2 = k$, puis
$f_1 = f_{\partial/\partial x_1}$ et $f_2(a) = a$. Alors
$A_1 \times_A A_2 = \{a \in k \mid \partial/\partial x_1(a) = 0\}$
n'a pas pour image tout $k$. Le produit fibré n'est donc pas un objet de
$\mathcal{C}_\Lambda$.
\end{example}

\noindent
On verra que ce problème ne peut se produire que si l'extension de corps
résiduels $k/k'$ (\ref{formal-defos-equation-k-prime}) est inséparable et si ni $f_1$ ni
$f_2$ n'est surjectif.

\begin{lemma}
\label{formal-defos-lemma-fiber-product-CLambda}
Soient $f_1 : A_1 \to A$ et $f_2 : A_2 \to A$ des homomorphismes d'anneaux
dans $\mathcal{C}_\Lambda$. Alors :
\begin{enumerate}
\item Si $f_1$ ou $f_2$ est surjectif, alors
$A_1 \times_A A_2$ appartient à $\mathcal{C}_\Lambda$.
\item Si $f_2$ est une petite extension, il en est de même de
$A_1 \times_A A_2 \to A_1$.
\item Si l'extension de corps $k/k'$ est séparable, alors
$A_1 \times_A A_2$ appartient à $\mathcal{C}_\Lambda$.
\end{enumerate}
\end{lemma}

\begin{proof}
L'anneau $A_1 \times_A A_2$ est une $\Lambda$-algèbre par l'homomorphisme
$\Lambda \to A_1 \times_A A_2$ induit par les homomorphismes
$\Lambda \to A_1$ et $\Lambda \to A_2$. C'est un anneau local dont l'unique
idéal maximal est
$$
\mathfrak m_{A_1} \times_{\mathfrak m_A} \mathfrak m_{A_2} =
\Ker(A_1 \times_A A_2 \longrightarrow k)
$$
Un anneau est artinien si et seulement s'il est de longueur finie comme module
sur lui-même ; voir Algèbre, lemme
\ref{algebra-lemma-artinian-finite-length}. Puisque $A_1$ et $A_2$ sont
artiniens, le lemme \ref{formal-defos-lemma-length} montre que
$\text{length}_\Lambda(A_1)$ et $\text{length}_\Lambda(A_2)$, donc aussi
$\text{length}_\Lambda(A_1 \times A_2)$, sont toutes finies. Comme
$A_1 \times_A A_2 \subset A_1 \times A_2$ est un sous-$\Lambda$-module, il
vient que
$\text{length}_{A_1 \times_A A_2}(A_1 \times_A A_2) \leq
\text{length}_\Lambda(A_1 \times_A A_2)$ est finie. Ainsi, $A_1
\times_A A_2$ est artinien. Le seul obstacle à ce que
$A_1 \times_A A_2$ soit un objet de $\mathcal{C}_\Lambda$ est donc la
possibilité que son corps résiduel s'envoie dans un sous-corps propre de $k$
par l'homomorphisme ci-dessus
$A_1 \times_A A_2 \to A \to A/\mathfrak m_A = k$.

\medskip\noindent
Démonstration de (1). Si $f_2$ est surjectif, la projection
$A_1 \times_A A_2 \to A_1$ est surjective. Par suite, le composé
$A_1 \times_A A_2 \to A_1 \to A_1/\mathfrak m_{A_1} = k$ est surjectif, et
l'on conclut que $A_1 \times_A A_2$ est un objet de $\mathcal{C}_\Lambda$.

\medskip\noindent
Démonstration de (2). Si $f_2$ est une petite extension, alors
$A_2 \to A$ et $A_1 \times_A A_2  \to A_1$ sont tous deux surjectifs et ont
le même noyau. Le noyau de $A_1 \times_A A_2  \to A_1$ est donc un espace
vectoriel de dimension $1$ sur $k$, et l'on voit que
$A_1 \times_A A_2  \to A_1$ est une petite extension.

\medskip\noindent
Démonstration de (3). Choisissons $\overline{x} \in k$ tel que
$k = k'(\overline{x})$ (voir Corps, lemme
\ref{fields-lemma-primitive-element}). Soit $P'(T) \in k'[T]$ le polynôme
minimal de $\overline{x}$ sur $k'$. Puisque $k/k'$ est séparable, on a
$\text{d}P/\text{d}T(\overline{x}) \not = 0$.
Choisissons un polynôme unitaire $P \in \Lambda[T]$ dont l'image est $P'$ par
l'homomorphisme surjectif $\Lambda[T] \to k'[T]$. Comme
$A, A_1, A_2$ sont henséliens, voir Algèbre, lemme
\ref{algebra-lemma-local-dimension-zero-henselian}, il existe
$x, x_1, x_2 \in A, A_1, A_2$ tels que $P(x) = 0, P(x_1) = 0,
P(x_2) = 0$ et tels que les images de $x, x_1, x_2$ dans $k$ soient toutes
$\overline{x}$.
Alors $(x_1, x_2) \in A_1 \times_A A_2$, car $x_1, x_2$ s'envoient sur
$x \in A$ par unicité ; voir Algèbre, lemme
\ref{algebra-lemma-uniqueness}. Ainsi, le corps résiduel de
$A_1 \times_A A_2$ contient un générateur de $k$ sur $k'$, ce qui conclut.
\end{proof}

\noindent
Définissons maintenant les surjections essentielles dans
$\mathcal{C}_\Lambda$. Une condition nécessaire et suffisante pour qu'une
surjection de $\mathcal{C}_\Lambda$ soit essentielle est donnée au lemme
\ref{formal-defos-lemma-essential-surjection}.

\begin{definition}
\label{formal-defos-definition-essential-surjection}
Soit $f: B \to A$ un homomorphisme d'anneaux dans $\mathcal{C}_\Lambda$.
On dit que $f$ est une {\it surjection essentielle} s'il possède les
propriétés suivantes :
\begin{enumerate}
\item $f$ est surjectif.
\item Si $g: C \to B$ est un homomorphisme d'anneaux dans
$\mathcal{C}_\Lambda$ tel que $f \circ g$ soit surjectif, alors $g$ est
surjectif.
\end{enumerate}
\end{definition}

\noindent
À l'aide du lemme \ref{formal-defos-lemma-surjective}, on peut caractériser comme suit les
surjections essentielles dans $\mathcal{C}_\Lambda$.

\begin{lemma}
\label{formal-defos-lemma-essential-surjection-mod-squares}
Soit $f: B \to A$ un homomorphisme d'anneaux dans $\mathcal{C}_\Lambda$.
Les conditions suivantes sont équivalentes :
\begin{enumerate}
\item $f$ est une surjection essentielle ;
\item l'homomorphisme $B/\mathfrak m_B^2 \to A/\mathfrak m_A^2$ est une
surjection essentielle ;
\item l'homomorphisme
$B/(\mathfrak m_\Lambda B + \mathfrak m_B^2) \to
A/(\mathfrak m_\Lambda A + \mathfrak m_A^2)$ est une surjection essentielle.
\end{enumerate}
\end{lemma}

\begin{proof}
Supposons (3). Soit $C \to B$ un homomorphisme d'anneaux dans
$\mathcal{C}_\Lambda$ tel que $C \to A$ soit surjectif. Alors
$C \to A/(\mathfrak m_\Lambda A + \mathfrak m_A^2)$ est également surjectif.
L'hypothèse entraîne que
$C \to B/(\mathfrak m_\Lambda B + \mathfrak m_B^2)$ est surjectif. Par deux
applications du lemme \ref{formal-defos-lemma-surjective}, $C \to B$ est donc surjectif.

\medskip\noindent
Supposons (1). Soit $C \to B/(\mathfrak m_\Lambda B + \mathfrak m_B^2)$ un
morphisme de $\mathcal{C}_\Lambda$ tel que
$C \to A/(\mathfrak m_\Lambda A + \mathfrak m_A^2)$ soit surjectif. Posons
$C' = C \times_{B/(\mathfrak m_\Lambda B + \mathfrak m_B^2)} B$
qui est un objet de $\mathcal{C}_\Lambda$ d'après le lemme
\ref{formal-defos-lemma-fiber-product-CLambda}. Remarquons que
$C' \to A/(\mathfrak m_\Lambda A + \mathfrak m_A^2)$ est encore surjectif ;
le lemme \ref{formal-defos-lemma-surjective} montre donc que $C' \to A$ est surjectif.
Notre hypothèse implique alors que $C' \to B$ est surjectif. Il en résulte que
$C' \to B/(\mathfrak m_\Lambda B + \mathfrak m_B^2)$ est surjectif, ce qui,
par construction de $C'$, implique que
$C \to B/(\mathfrak m_\Lambda B + \mathfrak m_B^2)$ est surjectif.

\medskip\noindent
Dans le premier paragraphe, nous avons démontré (3) $\Rightarrow$ (1), et dans
le second, (1) $\Rightarrow$ (3). L'équivalence de (2) et (3) est un cas
particulier de l'équivalence de (1) et (3), ce qui conclut.
\end{proof}

\noindent
Pour analyser plus avant les surjections essentielles dans
$\mathcal{C}_\Lambda$, introduisons quelques notations. Supposons que $A$ soit
un objet de $\mathcal{C}_\Lambda$, ou plus généralement une
$\Lambda$-algèbre quelconque munie d'une surjection de $\Lambda$-algèbres
$A \to k$. On dispose d'une suite exacte canonique
\begin{equation}
\label{formal-defos-equation-sequence}
\mathfrak m_A/\mathfrak m_A^2 \xrightarrow{\text{d}_A}
\Omega_{A/\Lambda} \otimes_A k \to
\Omega_{k/\Lambda} \to 0
\end{equation}
voir Algèbre, lemme \ref{algebra-lemma-differential-seq}.
Remarquons que $\Omega_{k/\Lambda} = \Omega_{k/k'}$, avec $k'$ comme dans
(\ref{formal-defos-equation-k-prime}). Soit $H_1(L_{k/\Lambda})$ le premier module
d'homologie du complexe cotangent naïf de $k$ sur $\Lambda$ ; voir Algèbre,
définition \ref{algebra-definition-naive-cotangent-complex}. On peut alors
prolonger (\ref{formal-defos-equation-sequence}) en la suite exacte
\begin{equation}
\label{formal-defos-equation-sequence-extended}
H_1(L_{k/\Lambda}) \to
\mathfrak m_A/\mathfrak m_A^2 \xrightarrow{\text{d}_A}
\Omega_{A/\Lambda} \otimes_A k \to
\Omega_{k/\Lambda} \to 0,
\end{equation}
voir Algèbre, lemme \ref{algebra-lemma-exact-sequence-NL}.
Si $B \to A$ est un homomorphisme d'anneaux dans $\mathcal{C}_\Lambda$, ou,
plus généralement, un homomorphisme de $\Lambda$-algèbres munies de
surjections de $\Lambda$-algèbres sur $k$, on obtient un diagramme commutatif
\begin{equation}
\label{formal-defos-equation-sequence-functorial}
\vcenter{
\xymatrix{
H_1(L_{k/\Lambda}) \ar[r] \ar@{=}[d] &
\mathfrak m_B/\mathfrak m_B^2 \ar[r]_{\text{d}_B} \ar[d] &
\Omega_{B/\Lambda} \otimes_B k \ar[r] \ar[d] &
\Omega_{k/\Lambda} \ar[r] \ar@{=}[d] & 0 \\
H_1(L_{k/\Lambda}) \ar[r] &
\mathfrak m_A/\mathfrak m_A^2 \ar[r]^{\text{d}_A} &
\Omega_{A/\Lambda} \otimes_A k \ar[r] &
\Omega_{k/\Lambda} \ar[r] & 0
}
}
\end{equation}
dont les lignes sont exactes.

\begin{lemma}
\label{formal-defos-lemma-H1-separable-case}
Il existe un homomorphisme canonique
$$
\mathfrak m_\Lambda/\mathfrak m_\Lambda^2 \longrightarrow H_1(L_{k/\Lambda}).
$$
Si $k' \subset k$ est séparable (par exemple si la caractéristique de $k$ est
nulle), cet homomorphisme induit un isomorphisme
$\mathfrak m_\Lambda/\mathfrak m_\Lambda^2 \otimes_{k'} k = H_1(L_{k/\Lambda})$.
Si $k = k'$ (par exemple dans le cas classique), alors
$\mathfrak m_\Lambda/\mathfrak m_\Lambda^2 = H_1(L_{k/\Lambda})$.
Le composé
$$
\mathfrak m_\Lambda/\mathfrak m_\Lambda^2 \longrightarrow
H_1(L_{k/\Lambda}) \longrightarrow \mathfrak m_A/\mathfrak m_A^2
$$
provient de l'homomorphisme canonique $\mathfrak m_\Lambda \to \mathfrak m_A$.
\end{lemma}

\begin{proof}
On a $H_1(L_{k'/\Lambda}) = \mathfrak m_\Lambda/\mathfrak m_\Lambda^2$, car
$\Lambda \to k'$ est surjectif de noyau $\mathfrak m_\Lambda$.
L'homomorphisme provient de la fonctorialité du complexe cotangent naïf.
Si $k' \subset k$ est séparable, alors $k' \to k$ est un homomorphisme
d'anneaux \'etale ; voir Algèbre, lemme
\ref{algebra-lemma-etale-over-field}. Son complexe cotangent naïf a donc des
groupes d'homologie triviaux ; voir Algèbre, définition
\ref{algebra-definition-etale}. Appliqué aux homomorphismes d'anneaux
$\Lambda \to k' \to k$, le lemme
\ref{algebra-lemma-exact-sequence-NL} d'Algèbre entraîne alors que
$\mathfrak m_\Lambda/\mathfrak m_\Lambda^2 \otimes_{k'} k = H_1(L_{k/\Lambda})$.
Nous omettons la démonstration de la dernière assertion.
\end{proof}

\begin{lemma}
\label{formal-defos-lemma-essential-surjection}
Soit $f: B \to A$ un homomorphisme d'anneaux dans $\mathcal{C}_\Lambda$.
On conserve les notations de (\ref{formal-defos-equation-sequence-functorial}).
\begin{enumerate}
\item Les conditions équivalentes du lemme \ref{formal-defos-lemma-surjective} qui
caractérisent la surjectivité de $f$ sont aussi équivalentes aux suivantes :
\begin{enumerate}
\item $\Im(\text{d}_B) \to \Im(\text{d}_A)$ est surjectif ;
\item l'homomorphisme $\Omega_{B/\Lambda} \otimes_B k \to
\Omega_{A/\Lambda} \otimes_A k$ est surjectif.
\end{enumerate}
\item Les conditions suivantes sont équivalentes :
\begin{enumerate}
\item $f$ est une surjection essentielle (voir le lemme
\ref{formal-defos-lemma-essential-surjection-mod-squares}) ;
\item l'homomorphisme $\Im(\text{d}_B) \to \Im(\text{d}_A)$ est un
isomorphisme ;
\item l'homomorphisme $\Omega_{B/\Lambda} \otimes_B k \to
\Omega_{A/\Lambda} \otimes_A k$ est un isomorphisme.
\end{enumerate}
\item Si $k/k'$ est séparable, alors $f$ est une surjection essentielle si et
seulement si l'homomorphisme
$\mathfrak m_B/(\mathfrak m_\Lambda B + \mathfrak m_B^2) \to
\mathfrak m_A/(\mathfrak m_\Lambda A + \mathfrak m_A^2)$
est un isomorphisme.
\item Si $f$ est une petite extension, alors $f$ n'est pas essentielle si et
seulement si $f$ admet une section $s: A \to B$ dans $\mathcal{C}_\Lambda$
telle que $f \circ s = \text{id}_A$.
\end{enumerate}
\end{lemma}

\begin{proof}
Démonstration de (1). Il résulte de (\ref{formal-defos-equation-sequence-functorial}) que
(1)(a) et (1)(b) sont équivalentes. En outre, si $A \to B$ est surjectif,
alors (1)(a) et (1)(b) sont satisfaites. Supposons (1)(a). Puisque le noyau de
$\text{d}_A$ est l'image de $H_1(L_{k/\Lambda})$, lequel s'envoie aussi dans
$\mathfrak m_B/\mathfrak m_B^2$, on en déduit que
$\mathfrak m_B/\mathfrak m_B^2 \to \mathfrak m_A/\mathfrak m_A^2$
est surjectif. Ainsi, $B \to A$ est surjectif d'après le lemme
\ref{formal-defos-lemma-surjective}, ce qui achève la démonstration de (1).

\medskip\noindent
Démonstration de (2). L'équivalence de (2)(b) et (2)(c) résulte immédiatement
de (\ref{formal-defos-equation-sequence-functorial}).

\medskip\noindent
Supposons (2)(b). Soit $g : C \to B$ un homomorphisme d'anneaux dans
$\mathcal{C}_\Lambda$ tel que $f \circ g$ soit surjectif. On en déduit que
$\mathfrak m_C/\mathfrak m_C^2 \to \mathfrak m_A/\mathfrak m_A^2$
est surjectif d'après le lemme \ref{formal-defos-lemma-surjective}. Par conséquent,
$\Im(\text{d}_C) \to \Im(\text{d}_A)$ est surjectif et, par hypothèse,
$\Im(\text{d}_C) \to \Im(\text{d}_B)$ l'est aussi. Il résulte de (1) que
$C \to B$ est surjectif.

\medskip\noindent
Supposons (2)(a). Alors $f$ est surjectif et l'on voit que
$\Omega_{B/\Lambda} \otimes_B k \to \Omega_{A/\Lambda} \otimes_A k$
est surjectif. Soit $K$ son noyau. Remarquons que
$K = \text{d}_B(\Ker(\mathfrak m_B/\mathfrak m_B^2 \to
\mathfrak m_A/\mathfrak m_A^2))$ d'après
(\ref{formal-defos-equation-sequence-functorial}). Choisissons une décomposition
$$
\Omega_{B/\Lambda} \otimes_B k =
\Omega_{A/\Lambda} \otimes_A k \oplus K
$$
d'espaces vectoriels sur $k$. L'homomorphisme
$\text{d} : B \to \Omega_{B/\Lambda}$ induit, par projection sur $K$, un
homomorphisme $D : B \to K$. Posons
$C = \{b \in B \mid D(b) = 0\}$. La règle de Leibniz montre que c'est une
sous-$\Lambda$-algèbre de $B$. Soit $\overline{x} \in k$. Choisissons
$x \in B$ d'image $\overline{x}$. Si $D(x) \not = 0$, il existe un élément
$y \in \mathfrak m_B$ tel que $D(y) = D(x)$. Ainsi, $x - y \in C$ est un
élément d'image $\overline{x}$. Par conséquent, $C \to k$ est surjectif et
$C$ est un objet de $\mathcal{C}_\Lambda$. De même, choisissons
$\omega \in \Im(\text{d}_A)$. D'après (1), il existe
$x \in \mathfrak m_B$ tel que $\text{d}_B(x)$ s'envoie sur $\omega$. Si
$D(x) \not = 0$, il existe un élément $y \in \mathfrak m_B$ dont l'image dans
$\mathfrak m_A/\mathfrak m_A^2$ est nulle et tel que $D(y) = D(x)$.
Alors $z = x - y$ est un élément de $\mathfrak m_C$ dont l'image
$\text{d}_C(z) \in \Omega_{C/k} \otimes_C k$ s'envoie sur $\omega$.
Ainsi, $\Im(\text{d}_C) \to \Im(\text{d}_A)$ est surjectif. On déduit de
(1) que $C \to A$ est surjectif, puis de l'hypothèse que $C \to B$ est
surjectif. Par suite, $D = 0$, c'est-à-dire $K = 0$, c'est-à-dire que (2)(c)
est satisfaite. Cela achève la démonstration de (2).

\medskip\noindent
Démonstration de (3). Si $k'/k$ est séparable, alors
$H_1(L_{k/\Lambda}) =
\mathfrak m_\Lambda/\mathfrak m_\Lambda^2 \otimes_{k'} k$ ; voir le lemme
\ref{formal-defos-lemma-H1-separable-case}. Par conséquent, $\Im(\text{d}_A) =
\mathfrak m_A/(\mathfrak m_\Lambda A + \mathfrak m_A^2)$
et il en va de même pour $B$. Ainsi, (3) résulte de (2).

\medskip\noindent
Démonstration de (4). Une section $s$ de $f$ n'est pas surjective (par
définition, une petite extension a un noyau non trivial) ; ainsi, $f$ n'est pas
une surjection essentielle. Réciproquement, supposons que $f$ soit une petite
extension, mais non une surjection essentielle. Choisissons un homomorphisme
d'anneaux $C \to B$ dans $\mathcal{C}_\Lambda$, non surjectif, tel que
$C \to A$ soit surjectif. Soit $C' \subset B$ l'image de $C \to B$. Alors
$C' \not = B$, mais l'application de $C'$ vers $A$ est surjective. Puisque
$f : B \to A$ est une petite extension,
$\text{length}_C(B) = \text{length}_C(A) + 1$. Ainsi,
$\text{length}_C(C') \leq \text{length}_C(A)$, car $C'$ est un sous-anneau
propre de $B$. Or $C' \to A$ est surjectif ; on a donc nécessairement
$\text{length}_C(C') = \text{length}_C(A)$, et $C' \to A$ est un
isomorphisme qui fournit la section cherchée.
\end{proof}

\begin{example}
\label{formal-defos-example-essential-surjection}
Soit $\Lambda = k[[x]]$ l'anneau des séries formelles en $1$ variable sur $k$.
Posons $A = k$ et $B = \Lambda/(x^2)$. Alors $B \to A$ est une surjection
essentielle d'après le lemme \ref{formal-defos-lemma-essential-surjection}, car c'est une
petite extension et l'homomorphisme $B \to A$ n'admet pas d'inverse à droite
(dans la catégorie $\mathcal{C}_\Lambda$). Mais l'homomorphisme
$$
k \cong \mathfrak m_B/\mathfrak m_B^2
\longrightarrow
\mathfrak m_A/\mathfrak m_A^2 = 0
$$
n'est pas un isomorphisme. Ainsi, dans le lemme
\ref{formal-defos-lemma-essential-surjection} (3), il est nécessaire de considérer
l'homomorphisme entre espaces cotangents relatifs
$\mathfrak m_B/(\mathfrak m_\Lambda B + \mathfrak m_B^2) \to
\mathfrak m_A/(\mathfrak m_\Lambda A + \mathfrak m_A^2)$.
\end{example}







\section{La catégorie de base complétée}
\label{formal-defos-section-category-completion-CLambda}

\noindent
Le « complété » suivant de la catégorie $\mathcal{C}_\Lambda$ servira de
catégorie de base au complété d'une catégorie cofibrée en groupoïdes sur
$\mathcal{C}_\Lambda$ (section \ref{formal-defos-section-formal-objects}).

\begin{definition}
\label{formal-defos-definition-completion-CLambda}
Soit $\Lambda$ un anneau noethérien et soit $\Lambda \to k$ un homomorphisme
fini d'anneaux, où $k$ est un corps. On définit
{\it $\widehat{\mathcal{C}}_\Lambda$} comme la catégorie dont
\begin{enumerate}
\item les objets sont les couples $(R, \varphi)$, où $R$ est une
$\Lambda$-algèbre locale noethérienne complète et où
$\varphi : R/\mathfrak m_R \to k$ est un isomorphisme de
$\Lambda$-algèbres ;
\item les morphismes $f : (S, \psi) \to (R, \varphi)$ sont les
homomorphismes locaux de $\Lambda$-algèbres tels que
$\varphi \circ (f \bmod \mathfrak m) = \psi$.
\end{enumerate}
\end{definition}

\noindent
Comme dans la discussion qui suit la définition \ref{formal-defos-definition-CLambda}, nous
désignerons en général un objet de $\widehat{\mathcal{C}}_\Lambda$ simplement
par $R$, l'identification
$R/\mathfrak m_R = k$ étant sous-entendue. Nous étudions dans cette section
quelques propriétés élémentaires des objets et des morphismes de la catégorie
$\widehat{\mathcal{C}}_\Lambda$, parallèlement à l'étude de la catégorie
$\mathcal{C}_\Lambda$ menée dans la section précédente.

\medskip\noindent
Observons d'abord que tout objet $A \in \mathcal{C}_\Lambda$ est un objet de
$\widehat{\mathcal{C}}_\Lambda$, puisqu'un anneau local artinien est toujours
noethérien et complet pour la topologie définie par son idéal maximal (qui est
un idéal nilpotent). De plus, il résulte clairement des définitions que
$\mathcal{C}_\Lambda \subset \widehat{\mathcal{C}}_\Lambda$
est la sous-catégorie strictement pleine formée de tous les anneaux artiniens.
Réciproquement, tout objet de $\widehat{\mathcal{C}}_\Lambda$ est en fait une
limite d'objets de $\mathcal{C}_\Lambda$.

\medskip\noindent
Supposons que $R$ soit un objet de $\widehat{\mathcal{C}}_\Lambda$.
Considérons les anneaux $R_n = R/\mathfrak m_R^n$, pour
$n \in \mathbf{N}$. Ce sont des anneaux locaux noethériens qui possèdent un
unique idéal premier, lequel est nilpotent ; ils sont donc artiniens, voir
Algèbre, proposition \ref{algebra-proposition-dimension-zero-ring}.
Les homomorphismes d'anneaux
$$
\ldots \to R_{n + 1} \to R_n \to \ldots \to R_2 \to R_1 = k
$$
sont tous surjectifs. Par définition, la complétude de $R$ signifie que
$R = \lim R_n$. Si $f : R \to S$ est un homomorphisme d'anneaux dans
$\widehat{\mathcal{C}}_\Lambda$, on obtient un système d'homomorphismes
$f_n : R_n \to S_n$ dont la limite est l'homomorphisme donné.

\begin{lemma}
\label{formal-defos-lemma-surjective-cotangent-space}
Soit $f: R \to S$ un homomorphisme d'anneaux dans
$\widehat{\mathcal{C}}_\Lambda$. Les conditions suivantes sont équivalentes :
\begin{enumerate}
\item $f$ est surjectif ;
\item l'homomorphisme
$\mathfrak m_R/\mathfrak m_R^2 \to \mathfrak m_S/\mathfrak m_S^2$
est surjectif ;
\item l'homomorphisme
$\mathfrak m_R/(\mathfrak m_\Lambda R + \mathfrak m_R^2) \to
\mathfrak m_S/(\mathfrak m_\Lambda S + \mathfrak m_S^2)$
est surjectif.
\end{enumerate}
\end{lemma}

\begin{proof}
Remarquons que, pour $n \geq 2$, on a l'égalité d'espaces cotangents relatifs
$$
\mathfrak m_R/(\mathfrak m_\Lambda R + \mathfrak m_R^2)
=
\mathfrak m_{R_n}/(\mathfrak m_\Lambda R_n + \mathfrak m_{R_n}^2)
$$
et de même pour $S$. Le lemme \ref{formal-defos-lemma-surjective} montre donc que
$R_n \to S_n$ est surjectif pour tout $n$.
Soit maintenant $K_n$ le noyau de $R_n \to S_n$. Les suites
$$
0 \to K_n \to R_n \to S_n \to 0
$$
forment une suite exacte de systèmes projectifs dirigés. Le système $(K_n)$
vérifie la condition de Mittag-Leffler, puisque chaque $K_n$ est artinien.
D'après Algèbre, lemme \ref{algebra-lemma-ML-exact-sequence}, le passage à la
limite préserve donc l'exactitude. Ainsi,
$\lim R_n \to \lim S_n$ est surjectif, c'est-à-dire que $f$ est surjectif.
\end{proof}

\begin{lemma}
\label{formal-defos-lemma-CLambdahat-pushouts}
La catégorie $\widehat{\mathcal{C}}_\Lambda$ admet des sommes amalgamées.
\end{lemma}

\begin{proof}
Soient $R \to S_1$ et $R \to S_2$ des morphismes de
$\widehat{\mathcal{C}}_\Lambda$. Considérons l'anneau
$C = S_1 \otimes_R S_2$. Cet anneau possède l'idéal maximal de type fini
$\mathfrak m = \mathfrak m_{S_1} \otimes S_2 +
S_1 \otimes \mathfrak m_{S_2}$, de corps résiduel $k$.
Soit $C^\wedge$ le complété de $C$ pour la topologie $\mathfrak m$-adique.
Alors $C^\wedge$ est un anneau noethérien complet pour la topologie définie
par l'idéal maximal $\mathfrak m^\wedge = \mathfrak mC^\wedge$, dont le corps
résiduel est identifié à $k$ ; voir Algèbre, lemme
\ref{algebra-lemma-completion-Noetherian}. Ainsi, $C^\wedge$ est un objet de
$\widehat{\mathcal{C}}_\Lambda$. Les homomorphismes
$S_1 \to C^\wedge$ et $S_2 \to C^\wedge$ font alors de $C^\wedge$ une somme
amalgamée au-dessus de $R$ dans $\widehat{\mathcal{C}}_\Lambda$ (détails omis).
\end{proof}

\noindent
Nous n'aurons pas besoin du lemme suivant.

\begin{lemma}
\label{formal-defos-lemma-CLambdahat-coproducts}
La catégorie $\widehat{\mathcal{C}}_\Lambda$ admet les coproduits de deux
objets.
\end{lemma}

\begin{proof}
Soient $R$ et $S$ des objets de $\widehat{\mathcal{C}}_\Lambda$.
Considérons l'anneau $C = R \otimes_\Lambda S$. Il existe un homomorphisme
surjectif canonique $C \to R \otimes_\Lambda S \to k \otimes_\Lambda k \to k$,
le dernier homomorphisme étant celui de multiplication. Le noyau de
$C \to k$ est un idéal maximal $\mathfrak m$. Remarquons que $\mathfrak m$
est engendré par $\mathfrak m_R C$, $\mathfrak m_S C$ et par un nombre fini
d'éléments de $C$ dont les images engendrent le noyau de
$k \otimes_\Lambda k \to k$. Ainsi, $\mathfrak m$ est un idéal de type fini.
Soit $C^\wedge$ le complété de $C$ pour la topologie $\mathfrak m$-adique.
Alors $C^\wedge$ est un anneau noethérien complet pour la topologie définie
par l'idéal maximal $\mathfrak m^\wedge = \mathfrak mC^\wedge$, de corps
résiduel $k$ ; voir Algèbre, lemme
\ref{algebra-lemma-completion-Noetherian}. Ainsi, $C^\wedge$ est un objet de
$\widehat{\mathcal{C}}_\Lambda$. Les homomorphismes $R \to C^\wedge$ et
$S \to C^\wedge$ font alors de $C^\wedge$ un coproduit dans
$\widehat{\mathcal{C}}_\Lambda$ (détails omis).
\end{proof}

\noindent
Dans une catégorie, un coproduit vide est un objet initial. Dans le cas
classique, $\widehat{\mathcal{C}}_\Lambda$ possède un objet initial, à savoir
$\Lambda$ lui-même. Plus généralement, si $k' = k$, alors le complété
$\Lambda^\wedge$ de $\Lambda$ pour la topologie $\mathfrak m_\Lambda$-adique
est un objet initial. Plus généralement encore, si $k' \subset k$ est
séparable, alors $\widehat{\mathcal{C}}_\Lambda$ possède aussi un objet
initial. En effet, choisissons un polynôme unitaire $P \in \Lambda[T]$ tel que
$k \cong k'[T]/(P')$, où $p' \in k'[T]$ est l'image de $P$. Alors
$R = \Lambda^\wedge[T]/(P)$ est un objet initial ; voir la démonstration du
lemme \ref{formal-defos-lemma-fiber-product-CLambda}.

\medskip\noindent
Si $R$ est un objet initial comme ci-dessus, alors on a
$\mathcal{C}_\Lambda = \mathcal{C}_R$ et
$\widehat{\mathcal{C}}_\Lambda = \widehat{\mathcal{C}}_R$, ce qui ramène en
fait toute la discussion de ce chapitre au cas classique. En revanche, si
$k' \subset k$ est inséparable, il n'existe pas d'objet initial.

\begin{lemma}
\label{formal-defos-lemma-derivations-finite}
Soit $S$ un objet de $\widehat{\mathcal{C}}_\Lambda$.
Alors $\dim_k \text{Der}_\Lambda(S, k) < \infty$.
\end{lemma}

\begin{proof}
Choisissons $x_1, \ldots, x_n \in \mathfrak m_S$ dont les images forment une
base sur $k$ de l'espace cotangent relatif
$\mathfrak m_S/(\mathfrak m_\Lambda S + \mathfrak m_S^2)$.
Choisissons $y_1, \ldots, y_m \in S$ dont les images dans $k$ engendrent $k$
sur $k'$. Nous affirmons que $\dim_k \text{Der}_\Lambda(S, k) \leq n + m$.
Pour le voir, il suffit de démontrer que, si $D(x_i) = 0$ et
$D(y_j) = 0$, alors $D = 0$. Soit $a \in S$. Il existe un polynôme
$P = \sum \lambda_J y^J$, avec $\lambda_J \in \Lambda$, dont l'image dans
$k$ est égale à celle de $a$ dans $k$. On voit alors que
$D(a - P) = D(a) - D(P) = D(a)$ grâce à l'hypothèse
$D(y_j) = 0$ pour tout $j$. On peut donc supposer
$a \in \mathfrak m_S$. Écrivons $a = \sum a_i x_i$, avec $a_i \in S$.
La règle de Leibniz donne
$$
D(a) = \sum x_iD(a_i) + \sum a_iD(x_i) = \sum x_iD(a_i)
$$
car nous avons supposé $D(x_i) = 0$. On a $\sum x_iD(a_i) = 0$, puisque la
multiplication par $x_i$ est nulle sur $k$.
\end{proof}

\begin{lemma}
\label{formal-defos-lemma-derivations-surjective}
Soit $f : R \to S$ un morphisme de $\widehat{\mathcal{C}}_\Lambda$.
Si $\text{Der}_\Lambda(S, k) \to \text{Der}_\Lambda(R, k)$ est injectif,
alors $f$ est surjectif.
\end{lemma}

\begin{proof}
Si $f$ n'est pas surjectif, alors
$\mathfrak m_S/(\mathfrak m_R S + \mathfrak m_S^2)$ est non nul d'après le
lemme \ref{formal-defos-lemma-surjective-cotangent-space}. Par conséquent,
$Q = S/(f(R) + \mathfrak m_R S + \mathfrak m_S^2)$ est lui aussi non nul.
Remarquons que $Q$ est un espace vectoriel sur $k = R/\mathfrak m_R$ par
$f$. Munissons $Q$ d'une structure de $S$-module par
$S \to k$. L'application quotient $D : S \to Q$ est une $R$-dérivation : si
$a_1, a_2 \in S$, on peut écrire $a_1 = f(b_1) + a_1'$ et
$a_2 = f(b_2) + a_2'$, avec $b_1, b_2 \in R$ et
$a_1', a_2' \in \mathfrak m_S$. Alors
Les éléments $b_i$ et $a_i$ ont la même image dans $k$ pour $i = 1, 2$, et
\begin{align*}
a_1a_2 & = (f(b_1) + a_1')(f(b_2) + a_2') \\
& = f(b_1)a_2' + f(b_2)a_1' \\
& = f(b_1)(f(b_2) + a_2') + f(b_2)(f(b_1) + a_1') \\
& = f(b_1)a_2 + f(b_2)a_1
\end{align*}
dans $Q$, ce qui démontre la règle de Leibniz. Ainsi, $D : S \to Q$ est une
$\Lambda$-dérivation dont le composé avec $R \to S$ est nul. Puisque
$Q \not = 0$, il existe aussi des dérivations $D : S \to k$ dont le composé
avec $R \to S$ est nul, c'est-à-dire que
$\text{Der}_\Lambda(S, k) \to \text{Der}_\Lambda(R, k)$ n'est pas injectif.
\end{proof}

\begin{lemma}
\label{formal-defos-lemma-m-adic-topology}
Soit $R$ un objet de $\widehat{\mathcal{C}}_\Lambda$. Soit $(J_n)$ une suite
décroissante d'idéaux telle que $\mathfrak m_R^n \subset J_n$.
Posons $J = \bigcap J_n$. Alors la suite $(J_n/J)$ définit la topologie
$\mathfrak m_{R/J}$-adique sur $R/J$.
\end{lemma}

\begin{proof}
Il est clair que $\mathfrak m_{R/J}^n \subset J_n/J$. Il suffit donc de
montrer que, pour tout $n$, il existe un $N$ tel que
$J_N/J \subset \mathfrak m_{R/J}^n$. Cela équivaut à
$J_N \subset \mathfrak m_R^n + J$. Pour tout $n$, l'anneau
$R/\mathfrak m_R^n$ est artinien ; il existe donc un $N_n$ tel que
$$
J_{N_n} + \mathfrak m_R^n = J_{N_n + 1} + \mathfrak m_R^n = \ldots
$$
Posons $E_n = (J_{N_n} + \mathfrak m_R^n)/\mathfrak m_R^n$.
Posons $E = \lim E_n \subset \lim R/\mathfrak m_R^n = R$.
Remarquons que $E \subset J$ : en effet, pour tout $f \in E$ et tout $m$, on
a $f \in J_m + \mathfrak m_R^n$ pour tout $n \gg 0$, donc $f \in J_m$ par le
théorème d'intersection de Krull ; voir Algèbre, lemme
\ref{algebra-lemma-intersect-powers-ideal-module-zero}. Comme les
homomorphismes de transition $E_n \to E_{n - 1}$ sont tous surjectifs, on voit
que $J$ se surjecte sur $E_n$. Ainsi, $N = N_n$ convient.
\end{proof}

\begin{lemma}
\label{formal-defos-lemma-limit-artinian}
Soit $\ldots \to A_3 \to A_2 \to A_1$ une suite d'homomorphismes d'anneaux
surjectifs dans $\mathcal{C}_\Lambda$. Si
$\dim_k (\mathfrak m_{A_n}/\mathfrak m_{A_n}^2)$ est bornée, alors
$S = \lim A_n$ est un objet de $\widehat{\mathcal{C}}_\Lambda$, et les idéaux
$I_n = \Ker(S \to A_n)$ définissent la topologie $\mathfrak m_S$-adique sur
$S$.
\end{lemma}

\begin{proof}
Nous utiliserons librement le fait que les homomorphismes $S \to A_n$ sont
surjectifs pour tout $n$. Remarquons que les homomorphismes
$\mathfrak m_{A_{n + 1}}/\mathfrak m_{A_{n + 1}}^2 \to
\mathfrak m_{A_n}/\mathfrak m_{A_n}^2$ sont surjectifs ; voir le lemme
\ref{formal-defos-lemma-surjective-cotangent-space}. Par conséquent, pour $n$ assez grand,
la dimension $\dim_k (\mathfrak m_{A_n}/\mathfrak m_{A_n}^2)$ se stabilise à
un entier, disons $r$. Il existe donc
$x_1, \ldots, x_r \in \mathfrak m_S$ dont les images dans $A_n$ engendrent
$\mathfrak m_{A_n}$. Choisissons en outre $y_1, \ldots, y_t \in S$ dont les
images dans $k$ engendrent $k$ sur $\Lambda$. On obtient alors un homomorphisme
d'anneaux
$P = \Lambda[z_1, \ldots, z_{r + t}] \to S$, $z_i \mapsto x_i$ et
$z_{r + j} \mapsto y_j$ tel que le composé $P \to S \to A_n$ soit surjectif
pour tout $n$. Soit $\mathfrak m \subset P$ le noyau de $P \to k$.
Soit $R = P^\wedge$ le complété $\mathfrak m$-adique de $P$ ; c'est un objet
de $\widehat{\mathcal{C}}_\Lambda$. Comme on dispose encore du système
compatible d'homomorphismes surjectifs $R \to A_n$, on obtient un
homomorphisme $R \to S$. Posons $J_n = \Ker(R \to A_n)$.
Posons $J = \bigcap J_n$. Le lemme \ref{formal-defos-lemma-m-adic-topology} montre que
$R/J = \lim R/J_n = \lim A_n = S$ et que les idéaux $J_n/J = I_n$ définissent
la topologie $\mathfrak m$-adique. (Remarquons que, pour tout $n$, on a
$\mathfrak m_R^{N_n} \subset J_n$ pour un certain $N_n$, sans avoir
nécessairement $N_n = n$ ; il peut donc être nécessaire de renuméroter les
idéaux $J_n$ avant d'appliquer le lemme.)
\end{proof}

\begin{lemma}
\label{formal-defos-lemma-power-series}
Soient $R', R \in \Ob(\widehat{\mathcal{C}}_\Lambda)$. Supposons que
$R = R' \oplus I$ pour un idéal $I$ de $R$. Soient
$x_1, \ldots, x_r \in I$ des éléments dont les images forment une base de
$I/\mathfrak m_R I$. Posons $S = R'[[X_1, \ldots, X_r]]$ et considérons
l'homomorphisme de $R'$-algèbres $S \to R$ qui envoie $X_i$ sur $x_i$.
Supposons que, pour tout $n \gg 0$, l'homomorphisme
$S/\mathfrak m_S^n \to R/\mathfrak m_R^n$ admette un inverse à gauche dans
$\mathcal{C}_\Lambda$. Alors $S \to R$ est un isomorphisme.
\end{lemma}

\begin{proof}
Puisque $R = R' \oplus I$, on a
$$
\mathfrak m_R/\mathfrak m_R^2 =
\mathfrak m_{R'}/\mathfrak m_{R'}^2 \oplus I/\mathfrak m_RI
$$
et, de même,
$$
\mathfrak m_S/\mathfrak m_S^2 =
\mathfrak m_{R'}/\mathfrak m_{R'}^2 \oplus \bigoplus kX_i
$$
Par conséquent, pour $n > 1$, l'homomorphisme
$S/\mathfrak m_S^n \to R/\mathfrak m_R^n$ induit un isomorphisme sur les
espaces cotangents. Ainsi, un inverse à gauche
$h_n : R/\mathfrak m_R^n \to S/\mathfrak m_S^n$ est surjectif d'après le
lemme \ref{formal-defos-lemma-surjective-cotangent-space}. Comme $h_n$ est injectif en tant
qu'inverse à gauche, c'est un isomorphisme. Les surjections canoniques
$S/\mathfrak m_S^n \to R/\mathfrak m_R^n$ sont donc toutes des isomorphismes,
ce qui conclut.
\end{proof}




\section{Catégories cofibrées en groupoïdes}
\label{formal-defos-section-preliminary}

\noindent
Pour développer la théorie, nous travaillons avec des catégories {\it cofibrées}
en groupoïdes. Nous supposons connues la définition et les propriétés
fondamentales des catégories {\it fibrées} en groupoïdes, voir
Catégories, section \ref{categories-section-fibred-groupoids}.

\begin{definition}
\label{formal-defos-definition-category-cofibred-groupoids}
Soit $\mathcal{C}$ une catégorie. Une {\it catégorie cofibrée en groupoïdes
au-dessus de $\mathcal{C}$} est une catégorie $\mathcal{F}$ munie d'un foncteur
$p: \mathcal{F} \to \mathcal{C}$ tel que $\mathcal{F}^{opp}$ soit une catégorie
fibrée en groupoïdes au-dessus de $\mathcal{C}^{opp}$ au moyen de
$p^{opp}: \mathcal{F}^{opp} \to \mathcal{C}^{opp}$.
\end{definition}

\noindent
Plus explicitement, le foncteur $p: \mathcal{F} \to \mathcal{C}$ est cofibré
en groupoïdes si les deux conditions suivantes sont satisfaites :
\begin{enumerate}
\item Pour tout morphisme $f: U \to V$ de $\mathcal{C}$ et tout objet
$x$ situé au-dessus de $U$, il existe un morphisme $x \to y$ de $\mathcal{F}$
situé au-dessus de $f$.
\item Pour toute paire de morphismes $a: x \to y$ et $b: x \to z$
de $\mathcal{F}$ et tout morphisme $f: p(y) \to p(z)$ tel que $p(b) = f
\circ p(a)$, il existe un unique morphisme $c: y \to z$ de $\mathcal
F$ situé au-dessus de $f$ tel que $b = c \circ a$.
\end{enumerate}

\begin{remarks}
\label{formal-defos-remarks-cofibered-groupoids}
Tout ce qui concerne les catégories fibrées en groupoïdes se transpose
directement au cadre cofibré. Les remarques suivantes ont pour objet de fixer
les notations. Soit $\mathcal{C}$ une catégorie.
\begin{enumerate}
\item Nous omettons souvent le foncteur $p: \mathcal{F} \to \mathcal{C}$ des
notations.
\item La catégorie fibre au-dessus d'un objet $U$ de $\mathcal{C}$ est notée
$\mathcal{F}(U)$. Ses objets sont ceux de $\mathcal{F}$ situés au-dessus de $U$
et ses morphismes sont ceux de $\mathcal{F}$ situés au-dessus de $\text{id}_U$.
Si $x, y$ sont des objets de $\mathcal{F}(U)$, nous écrivons parfois
$\Mor_U(x, y)$ pour $\Mor_{\mathcal{F}(U)}(x, y)$.
\item Les catégories fibres $\mathcal{F}(U)$ sont des groupoïdes, voir
Catégories, lemme \ref{categories-lemma-fibred-groupoids}.
Ainsi, tous les morphismes de $\mathcal{F}(U)$ sont des isomorphismes.
Nous écrivons parfois $\text{Aut}_U(x)$ pour $\Mor_{\mathcal{F}(U)}(x, x)$.
\item
\label{formal-defos-item-pushforward}
Soit $\mathcal{F}$ une catégorie cofibrée en groupoïdes au-dessus de
$\mathcal{C}$, soit $f: U \to V$ un morphisme de $\mathcal{C}$ et
soit $x \in \Ob(\mathcal{F}(U))$.
Une {\it image directe} de $x$ par $f$ est un morphisme
$x \to y$ de $\mathcal{F}$ situé au-dessus de $f$. Une image directe
est déterminée à un isomorphisme unique près (voir la discussion qui suit
Catégories, définition \ref{categories-definition-cartesian-over-C}).
Nous écrivons parfois $x \to f_*x$ pour désigner « l'image directe » de $x$
par $f$.
\item Un {\it choix d'images directes pour $\mathcal{F}$} consiste à choisir
une image directe de $x$ par $f$ pour toute paire $(x, f)$ comme ci-dessus. Nous
pouvons effectuer un tel choix pour $\mathcal{F}$ à l'aide de l'axiome du choix.
\item Soit $\mathcal{F}$ une catégorie cofibrée en groupoïdes au-dessus de
$\mathcal{C}$. Étant donné un choix d'images directes pour $\mathcal{F}$, il
existe un pseudo-foncteur associé $\mathcal{C} \to \textit{Groupoïdes}$.
Nous n'utiliserons jamais cette construction et n'en donnons donc pas les détails.
\item
\label{formal-defos-item-cofibered-morphism}
Un morphisme de catégories cofibrées en groupoïdes au-dessus de $\mathcal{C}$
est un foncteur commutant aux projections sur $\mathcal{C}$. Si $\mathcal{F}$
et $\mathcal{F}'$ sont des catégories cofibrées en groupoïdes au-dessus de
$\mathcal{C}$, nous notons les morphismes de $\mathcal{F}$ vers $\mathcal{F}'$
par $\Mor_\mathcal{C}(\mathcal{F}, \mathcal{F}')$.
\item
\label{formal-defos-item-definition-cofibered-groupoids-2-category}
Les catégories cofibrées en groupoïdes forment une $(2, 1)$-catégorie
$\text{Cof}(\mathcal{C})$. Ses 1-morphismes sont ceux décrits en
(\ref{formal-defos-item-cofibered-morphism}). Si $p : \mathcal{F} \to C$ et
$p': \mathcal{F}' \to \mathcal{C}$ sont les projections de catégories
cofibrées en groupoïdes et si $\varphi, \psi : \mathcal{F} \to \mathcal{F}'$
sont des $1$-morphismes, un 2-morphisme $t : \varphi \to \psi$ est un morphisme
de foncteurs tel que $p'(t_x) = \text{id}_{p(x)}$ pour tout $x \in \Ob(\mathcal{F})$.
\item
\label{formal-defos-item-construction-associated-cofibered-groupoid}
Soit $F : \mathcal{C} \to \textit{Groupoïdes}$ un foncteur. Il
existe une catégorie cofibrée en groupoïdes $\mathcal{F} \to \mathcal{C}$
associée à $F$ de la façon suivante. Un objet de $\mathcal{F}$ est une paire $(U, x)$
où $U \in \Ob(\mathcal{C})$ et $x \in \Ob(F(U))$. Un
morphisme $(U, x) \to (V, y)$ est une paire $(f, a)$ où
$f \in \Mor_\mathcal{C}(U, V)$ et
$a \in \Mor_{F(V)}(F(f)(x), y)$.
Le foncteur $\mathcal{F} \to \mathcal{C}$ envoie $(U, x)$ sur $U$. Voir
Catégories, section \ref{categories-section-presheaves-groupoids}.
\item
\label{formal-defos-item-associated-functor-isomorphism-classes}
Soit $\mathcal{F}$ cofibrée en groupoïdes au-dessus de $\mathcal{C}$.
Pour $U \in \Ob(\mathcal{C})$, posons $\overline{\mathcal{F}}(U)$ égal à
l'ensemble des classes d'isomorphie des objets de la catégorie $\mathcal{F}(U)$.
Si $f : U \to V$ est un morphisme de $\mathcal{C}$, nous obtenons alors une
application d'ensembles $\overline{\mathcal{F}}(U) \to \overline{\mathcal{F}}(V)$ en
envoyant la classe d'isomorphie de $x$ sur celle d'une image directe
$f_*x$ de $x$ ; voir (\ref{formal-defos-item-pushforward}). Alors
$\overline{\mathcal{F}} : \mathcal{C} \to \textit{Ensembles}$ est un
foncteur. De même, si $\varphi : \mathcal{F} \to \mathcal{G}$ est un
morphisme de catégories cofibrées, nous notons
$\overline{\varphi}: \overline{\mathcal{F}} \to  \overline{\mathcal{G}}$
le morphisme de foncteurs associé.
\item
\label{formal-defos-item-convention-cofibered-sets}
Soit $F: \mathcal{C} \to \textit{Ensembles}$ un foncteur. Nous pouvons regarder un
ensemble comme une catégorie discrète, c'est-à-dire comme un groupoïde dont les
seuls morphismes sont les identités. La construction
(\ref{formal-defos-item-construction-associated-cofibered-groupoid}) associe alors à $F$ une
catégorie cofibrée en ensembles. On obtient ainsi un foncteur pleinement fidèle
de la catégorie des foncteurs $\mathcal{C} \to \textit{Ensembles}$ vers la catégorie
des catégories cofibrées en groupoïdes au-dessus de $\mathcal{C}$.
Nous identifions la catégorie des foncteurs à son image par ce foncteur.
Ainsi, si $F : \mathcal{C} \to \textit{Ensembles}$ est un foncteur, nous notons aussi
par $F$ la catégorie cofibrée en ensembles associée ; et si
$\varphi : F \to G$ est un morphisme de foncteurs, nous notons encore $\varphi$
le morphisme correspondant de catégories cofibrées en ensembles, et réciproquement.
Voir Catégories, section \ref{categories-section-fibred-in-sets}.
\item
\label{formal-defos-item-definition-yoneda}
Soit $U$ un objet de $\mathcal{C}$. Nous notons $\underline{U}$ le
foncteur
$\Mor_\mathcal{C}(U, -): \mathcal{C} \to
\textit{Ensembles}$. Cela définit un foncteur pleinement fidèle de $\mathcal
C^{opp}$ vers la catégorie des foncteurs $\mathcal{C} \to
\textit{Ensembles}$. Ainsi, si $f : U \to V$ est un morphisme, nous pouvons
encore noter $f$ le morphisme induit $\underline{V}
\to \underline{U}$, et réciproquement.
\item
\label{formal-defos-item-fibre-product}
Produits fibrés de catégories cofibrées en groupoïdes : si $\mathcal{F}
\to \mathcal{H}$ et $\mathcal{G} \to \mathcal{H}$ sont des morphismes
de catégories cofibrées en groupoïdes au-dessus de $\mathcal{C}_\Lambda$, une
construction de leur 2-produit fibré est donnée par celle de leur
2-produit fibré comme catégories au-dessus de $\mathcal{C}_\Lambda$, décrite dans
Catégories, lemme \ref{categories-lemma-2-product-categories-over-C}.
\item
\label{formal-defos-item-product}
Produits de catégories cofibrées en groupoïdes : si $\mathcal{F}$ et
$\mathcal{G}$ sont des catégories cofibrées en groupoïdes au-dessus de
$\mathcal{C}_\Lambda$, leur produit est défini comme le $2$-produit fibré
$\mathcal{F} \times_{\mathcal{C}_\Lambda} \mathcal{G}$ décrit dans
Catégories, lemme \ref{categories-lemma-2-product-categories-over-C}.
\item
\label{formal-defos-item-definition-restricting-base-category}
Restriction de la catégorie de base : soit $p : \mathcal{F} \to \mathcal{C}$
une catégorie cofibrée en groupoïdes, et soit $\mathcal{C}'$ une sous-catégorie
pleine de $\mathcal{C}$. La restriction $\mathcal{F}|_{\mathcal{C}'}$
est la sous-catégorie pleine de $\mathcal{F}$ dont les objets sont situés
au-dessus des objets de $\mathcal{C}'$. C'est une catégorie cofibrée en
groupoïdes au moyen du foncteur
$p|_{\mathcal{C}'}: \mathcal{F}|_{\mathcal{C}'} \to \mathcal{C}'$.
\end{enumerate}
\end{remarks}









\section{Foncteurs pro-représentables et catégories de prédéformation}
\sectionmark{Foncteurs pro-représentables et prédéformations}
\label{formal-defos-section-cofibered-groupoids}

\noindent
Notre objectif fondamental est de comprendre les catégories cofibrées en
groupoïdes au-dessus de $\mathcal{C}_\Lambda$ et de
$\widehat{\mathcal{C}}_\Lambda$. Comme $\mathcal{C}_\Lambda$ est une
sous-catégorie pleine de $\widehat{\mathcal{C}}_\Lambda$, nous pouvons
restreindre les catégories cofibrées en groupoïdes au-dessus de
$\widehat{\mathcal{C}}_\Lambda$ à $\mathcal{C}_\Lambda$ ; voir les remarques
\ref{formal-defos-remarks-cofibered-groupoids}
(\ref{formal-defos-item-definition-restricting-base-category}).
Nous pouvons en particulier le faire pour les foncteurs, et notamment pour
les foncteurs représentables. Les foncteurs sur $\mathcal{C}_\Lambda$
ainsi obtenus sont appelés foncteurs pro-représentables.

\begin{definition}
\label{formal-defos-definition-prorepresentable}
Soit $F : \mathcal{C}_\Lambda \to \textit{Ensembles}$ un foncteur.
On dit que $F$ est {\it pro-représentable} s'il existe un isomorphisme
de foncteurs $F \cong \underline{R}|_{\mathcal{C}_\Lambda}$ pour un certain
$R \in \Ob(\widehat{\mathcal{C}}_\Lambda)$.
\end{definition}

\noindent
Remarquons que si $F : \mathcal{C}_\Lambda \to \textit{Ensembles}$ est
pro-représentable par $R \in \Ob(\widehat{\mathcal{C}}_\Lambda)$, alors
$$
F(k) = \Mor_{\widehat{\mathcal{C}}_\Lambda}(R, k) = \{*\}
$$
est réduit à un élément. Les catégories cofibrées en groupoïdes au-dessus de
$\mathcal{C}_\Lambda$ qui interviennent en théorie des déformations
satisferont souvent une condition analogue.

\begin{definition}
\label{formal-defos-definition-predeformation-category}
Une {\it catégorie de prédéformation} $\mathcal{F}$ est une catégorie
cofibrée en groupoïdes au-dessus de $\mathcal{C}_\Lambda$ telle que
$\mathcal{F}(k)$ soit équivalente à une catégorie possédant un seul objet et
un seul morphisme, c'est-à-dire que $\mathcal{F}(k)$ contient au moins un
objet et qu'il existe un unique morphisme entre deux objets quelconques. Un
{\it morphisme de catégories de prédéformation} est un morphisme de catégories
cofibrées en groupoïdes au-dessus de
$\mathcal{C}_\Lambda$.
\end{definition}

\noindent
Une catégorie de prédéformation possède la propriété suivante.
Soit $x_0 \in \Ob(\mathcal{F}(k))$. Tout objet de $\mathcal{F}$ est alors
muni d'un unique morphisme vers $x_0$. En effet, si $x$ est un objet de
$\mathcal{F}$ au-dessus de $A$, nous pouvons choisir une image directe
$x \to q_*x$, où $q : A \to k$ est le morphisme quotient. Il existe un unique
isomorphisme $q_*x \to x_0$, et la composée $x \to q_*x \to x_0$ est le
morphisme recherché.

\begin{remark}
\label{formal-defos-remark-predeformation-functor}
On dit qu'un foncteur $F: \mathcal{C}_\Lambda \to \textit{Ensembles}$ est un
{\it foncteur de prédéformation} si l'ensemble cofibré associé est une
catégorie de prédéformation, c'est-à-dire\ si $F(k)$ est un ensemble à un
seul élément. Ainsi,
si $\mathcal{F}$ est une catégorie de prédéformation, alors
$\overline{\mathcal{F}}$ est un foncteur de prédéformation.
\end{remark}

\begin{remark}
\label{formal-defos-remark-localize-cofibered-groupoid}
Soit $p : \mathcal{F} \to \mathcal{C}_\Lambda$ une catégorie cofibrée en
groupoïdes, et soit $x \in \Ob(\mathcal{F}(k))$. On note
$\mathcal{F}_x$ la catégorie des objets au-dessus de $x$.
Un objet de $\mathcal{F}_x$ est une flèche $y \to x$.
Un morphisme $(y \to x) \to (z \to x)$ dans $\mathcal{F}_x$ est un diagramme
commutatif
$$
\xymatrix{
y \ar[rr] \ar[dr] & & z \ar[dl] \\
& x &
}
$$
Il existe un foncteur d'oubli $\mathcal{F}_x \to \mathcal{F}$. On définit le
foncteur $p_x : \mathcal{F}_x \to \mathcal{C}_\Lambda$ comme la composée
$\mathcal{F}_x \to \mathcal{F} \xrightarrow{p} \mathcal{C}_\Lambda$.
Alors $p_x : \mathcal{F}_x \to \mathcal{C}_\Lambda$ est une catégorie de
prédéformation (démonstration omise). On peut ainsi passer d'une catégorie
cofibrée en groupoïdes arbitraire au-dessus de $\mathcal{C}_\Lambda$ à une
catégorie de prédéformation en tout $x \in \Ob(\mathcal{F}(k))$.
\end{remark}






\section{Objets formels et catégories complétées}
\label{formal-defos-section-formal-objects}

\noindent
Dans cette section, nous expliquons comment passer des catégories cofibrées
en groupoïdes au-dessus de $\mathcal{C}_\Lambda$ aux catégories cofibrées en
groupoïdes au-dessus de $\widehat{\mathcal{C}}_\Lambda$, et réciproquement.

\begin{definition}
\label{formal-defos-definition-formal-objects}
Soit $\mathcal{F}$ une catégorie cofibrée en groupoïdes au-dessus de
$\mathcal{C}_\Lambda$. La {\it catégorie $\widehat{\mathcal{F}}$ des objets
formels de $\mathcal{F}$} est la catégorie dont les objets et les morphismes
sont les suivants.
\begin{enumerate}
\item Un {\it objet formel $\xi = (R, \xi_n, f_n)$ de $\mathcal{F}$} est
constitué d'un objet $R$ de $\widehat{\mathcal{C}}_\Lambda$, d'une famille,
indexée par $n \in \mathbf{N}$, d'objets $\xi_n$ de
$\mathcal{F}(R/\mathfrak m_R^n)$, et de morphismes
$f_n : \xi_{n + 1} \to \xi_n$ situés au-dessus de la projection
$R/\mathfrak m_R^{n + 1} \to R/\mathfrak m_R^n$.
\item Soient $\xi = (R, \xi_n, f_n)$ et $\eta = (S, \eta_n, g_n)$ des
objets formels de $\mathcal{F}$. Un {\it morphisme d'objets formels
$a : \xi \to \eta$} est constitué d'un morphisme $a_0 : R \to S$ dans
$\widehat{\mathcal{C}}_\Lambda$ et d'une famille de morphismes
$a_n : \xi_n \to \eta_n$ de $\mathcal{F}$ situés au-dessus de
$R/\mathfrak m_R^n \to S/\mathfrak m_S^n$,
de telle sorte que, pour tout $n$, le diagramme
$$
\xymatrix{
\xi_{n + 1} \ar[r]_{f_n} \ar[d]_{a_{n + 1}} & \xi_n \ar[d]^{a_n} \\
\eta_{n + 1} \ar[r]^{g_n} & \eta_n
}
$$
soit commutatif.
\end{enumerate}
\end{definition}

\noindent
La catégorie des objets formels est munie d'un foncteur $\widehat{p}:
\widehat{\mathcal{F}} \to \widehat{\mathcal{C}}_\Lambda$ qui envoie un objet
$(R, \xi_n, f_n)$ sur $R$ et un morphisme
$(R, \xi_n, f_n) \to (S, \eta_n, g_n)$ sur le morphisme $R \to S$.

\begin{lemma}
\label{formal-defos-lemma-completion-cofibred}
Soit $p : \mathcal{F} \to \mathcal{C}_\Lambda$ une catégorie cofibrée en
groupoïdes. Alors
$\widehat{p} : \widehat{\mathcal{F}} \to \widehat{\mathcal{C}}_\Lambda$
est une catégorie cofibrée en groupoïdes.
\end{lemma}

\begin{proof}
Soit $R \to S$ un morphisme d'anneaux dans
$\widehat{\mathcal{C}}_\Lambda$. Soit $(R, \xi_n, f_n)$ un objet de
$\widehat{\mathcal{F}}$. Pour tout $n$, choisissons une image directe
$\xi_n \to \eta_n$ de $\xi_n$ le long de
$R/\mathfrak m_R^n \to S/\mathfrak m_S^n$. Pour tout $n$, il existe un unique
morphisme $g_n : \eta_{n + 1} \to \eta_n$ dans $\mathcal{F}$, situé au-dessus
de $S/\mathfrak m_S^{n + 1} \to S/\mathfrak m_S^n$, tel que
$$
\xymatrix{
\xi_{n + 1} \ar[d] \ar[r]_{f_n} & \xi_n \ar[d] \\
\eta_{n + 1} \ar[r]^{g_n} & \eta_n
}
$$
soit commutatif (par le premier axiome d'une catégorie cofibrée en
groupoïdes). Nous obtenons donc un morphisme
$(R, \xi_n, f_n) \to (S, \eta_n, g_n)$ situé au-dessus de $R \to S$ ; le
premier axiome d'une catégorie cofibrée en groupoïdes est ainsi vérifié pour
$\widehat{\mathcal{F}}$. Pour vérifier le second axiome, supposons donnés des
morphismes
$a : (R, \xi_n, f_n) \to (S, \eta_n, g_n)$ et
$b : (R, \xi_n, f_n) \to (T, \theta_n, h_n)$ dans $\widehat{\mathcal{F}}$
et un morphisme $c_0 : S \to T$ dans $\widehat{\mathcal{C}}_\Lambda$ tel que
$c_0 \circ a_0 = b_0$. Le second axiome d'une catégorie cofibrée en groupoïdes
appliqué à $\mathcal{F}$ fournit d'uniques morphismes
$c_n : \eta_n \to \theta_n$, situées au-dessus de
$S/\mathfrak m_S^n \to T/\mathfrak m_T^n$, telles que
$c_n \circ a_n = b_n$. En posant $c = (c_n)_{n \geq 0}$, on obtient le
morphisme recherché $c : (S, \eta_n, g_n) \to (T, \theta_n, h_n)$ dans
$\widehat{\mathcal{F}}$ (nous omettons la vérification de l'égalité
$h_n \circ c_{n + 1} = c_n \circ g_n$).
\end{proof}

\begin{definition}
\label{formal-defos-definition-completion}
Soit $p : \mathcal{F} \to \mathcal{C}_\Lambda$ une catégorie cofibrée en
groupoïdes. La catégorie cofibrée en groupoïdes
$\widehat{p} : \widehat{\mathcal  F} \to \widehat{\mathcal{C}}_\Lambda$
est appelée le {\it complété de $\mathcal{F}$}.
\end{definition}

\noindent
Si $\mathcal{F}$ est une catégorie cofibrée en groupoïdes au-dessus de $\mathcal
C_\Lambda$, nous avons défini $\widehat{\mathcal{F}}(R)$, pour $R \in
\Ob(\widehat{\mathcal{C}}_\Lambda)$, à l'aide de la filtration de $R$
par les puissances de son idéal maximal. Supposons maintenant que
$\mathcal{I} = (I_n)$ soit une filtration de $R$ par des idéaux qui induit la
topologie $\mathfrak{m}_R$-adique. On définit
$\widehat{\mathcal{F}}_\mathcal{I}(R)$ comme la catégorie ayant les objets et
les morphismes suivants :
\begin{enumerate}
\item Un objet est une famille $(\xi_n, f_n)_{n \in \mathbf{N}}$ d'objets
$\xi_n$ de $\mathcal{F}(R/I_n)$ et de morphismes
$f_n : \xi_{n + 1} \to \xi_n$ situés au-dessus des projections
$R/I_{n + 1} \to R/I_n$.
\item Un morphisme $a : (\xi_n, f_n) \to (\eta_n, g_n)$ est une famille de
morphismes $a_n : \xi_n \to \eta_n$ dans $\mathcal{F}(R/I_n)$ telle que,
pour tout $n$, le diagramme
$$
\xymatrix{
\xi_{n + 1} \ar[r]^{f_n} \ar[d]_{a_{n + 1}} & \xi_n \ar[d]^{a_n} \\
\eta_{n + 1} \ar[r]^{g_n} & \eta_n
}
$$
soit commutatif.
\end{enumerate}

\begin{lemma}
\label{formal-defos-lemma-formal-objects-different-filtration}
Dans la situation ci-dessus, $\widehat{\mathcal{F}}_\mathcal{I}(R)$ est
équivalente à la catégorie $\widehat{\mathcal{F}}(R)$.
\end{lemma}

\begin{proof}
On peut définir une équivalence
$\widehat{\mathcal{F}}_\mathcal{I}(R) \to \widehat{\mathcal{F}}(R)$ de la
manière suivante. Pour tout $n$, soit $m(n)$ le plus petit $m$ tel que
$I_m \subset \mathfrak m_R^n$. Étant donné un objet $(\xi_n, f_n)$ de
$\widehat{\mathcal{F}}_\mathcal{I}(R)$, soit $\eta_n$ l'image directe de
$\xi_{m(n)}$ le long de $R/I_{m(n)} \to R/\mathfrak m_R^n$. Soit
$g_n : \eta_{n + 1} \to \eta_n$ l'unique morphisme de $\mathcal{F}$ situé
au-dessus de $R/\mathfrak m_R^{n + 1} \to R/\mathfrak m_R^n$ tel que
$$
\xymatrix{
\xi_{m(n + 1)} \ar[rrr]_{f_{m(n)} \circ \ldots \circ f_{m(n + 1) - 1}} \ar[d]
& & & \xi_{m(n)} \ar[d] \\
\eta_{n + 1} \ar[rrr]^{g_n} & & & \eta_n
}
$$
soit commutatif (l'existence et l'unicité sont assurées par les axiomes d'une
catégorie cofibrée). Le foncteur
$\widehat{\mathcal{F}}_\mathcal{I}(R) \to \widehat{\mathcal{F}}(R)$
envoie $(\xi_n, f_n)$ sur $(R, \eta_n, g_n)$. Nous omettons la vérification
qu'il s'agit bien d'une équivalence de catégories.
\end{proof}

\begin{remark}
\label{formal-defos-remark-different-sequence-ideals}
Soit $p : \mathcal{F} \to \mathcal{C}_\Lambda$ une catégorie cofibrée en
groupoïdes. Supposons que, pour tout
$R \in \Ob(\widehat{\mathcal{C}}_\Lambda)$, on se soit donné une filtration
$\mathcal{I}_R$ de $R$ par des idéaux. Si $\mathcal{I}_R$ induit la topologie
$\mathfrak m_R$-adique sur $R$ pour tout $R$, on peut définir une catégorie
$\widehat{\mathcal{F}}_\mathcal{I}$ en calquant la définition de
$\widehat{\mathcal{F}}$. Cette catégorie est munie d'un morphisme
$\widehat{p}_\mathcal{I} : \widehat{\mathcal{F}}_\mathcal{I} \to
\widehat{\mathcal{C}}_\Lambda$ qui en fait une catégorie cofibrée en
groupoïdes et tel que $\widehat{\mathcal{F}}_\mathcal{I}(R)$ soit isomorphe à
$\widehat{\mathcal{F}}_{\mathcal{I}_R}(R)$, telle qu'elle a été définie
ci-dessus. Les catégories cofibrées en groupoïdes
$\widehat{\mathcal{F}}_\mathcal{I}$ et $\widehat{\mathcal{F}}$ sont
équivalentes : au-dessus d'un objet
$R \in \Ob(\widehat{\mathcal{C}}_\Lambda)$, on utilise l'équivalence du
lemme \ref{formal-defos-lemma-formal-objects-different-filtration}.
\end{remark}

\begin{remark}
\label{formal-defos-remark-completion-functor}
Soit $F: \mathcal{C}_\Lambda \to \textit{Ensembles}$ un foncteur. En identifiant
les foncteurs aux ensembles cofibrés, le complété de $F$ est le foncteur
$\widehat{F} : \widehat{\mathcal{C}}_\Lambda \to \textit{Ensembles}$
donné par $\widehat{F}(S) = \lim F(S/\mathfrak{m}_S^{n})$. Cela concorde avec
la définition de l'article de Schlessinger \cite{Sch}.
\end{remark}

\begin{remark}
\label{formal-defos-remark-restrict-completion}
Soit $\mathcal{F}$ une catégorie cofibrée en groupoïdes au-dessus de
$\mathcal{C}_\Lambda$. Nous affirmons qu'il existe une équivalence canonique
$$
can :
\widehat{\mathcal{F}}|_{\mathcal{C}_\Lambda}
\longrightarrow
\mathcal{F}.
$$
En effet, soient $A \in \Ob(\mathcal{C}_\Lambda)$ et
$(A, \xi_n, f_n)$ un objet de
$\widehat{\mathcal{F}}|_{\mathcal{C}_\Lambda}(A)$. Comme $A$ est artinien,
il existe un plus petit $m \in \mathbf{N}$ tel que
$\mathfrak m_A^m = 0$. Alors $can$ envoie $(A, \xi_n, f_n)$ sur $\xi_m$.
Ce foncteur est une équivalence de catégories cofibrées en groupoïdes en
vertu de Catégories, lemme
\ref{categories-lemma-equivalence-fibred-categories}, car il induit une
équivalence sur toutes les catégories fibres d'après le lemme
\ref{formal-defos-lemma-formal-objects-different-filtration} et parce que la topologie
$\mathfrak m_A$-adique d'un anneau local artinien $A$ provient de l'idéal
nul. Nous identifierons fréquemment $\mathcal{F}$ à une sous-catégorie pleine
de $\widehat{\mathcal{F}}$ au moyen d'un quasi-inverse du foncteur $can$.
\end{remark}

\begin{remark}
\label{formal-defos-remark-completion-morphism}
Soit $\varphi : \mathcal{F} \to \mathcal{G}$ un morphisme de catégories
cofibrées en groupoïdes au-dessus de $\mathcal{C}_\Lambda$. Il en résulte un
morphisme
$\widehat{\varphi}: \widehat{\mathcal{F}} \to \widehat{\mathcal{G}}$
de catégories cofibrées en groupoïdes au-dessus de
$\widehat{\mathcal{C}}_\Lambda$. Il envoie un objet
$\xi = (R, \xi_n, f_n)$ de $\widehat{\mathcal{F}}$ sur
$(R, \varphi(\xi_n), \varphi(f_n))$, et un morphisme
$(a_0 : R \to S, a_n : \xi_n \to \eta_n)$ entre des objets $\xi$ et $\eta$
de $\widehat{\mathcal{F}}$ sur
$(a_0 : R \to S, \varphi(a_n) : \varphi(\xi_n) \to \varphi(\eta_n))$.
Enfin, si $t : \varphi \to \varphi'$ est un $2$-morphisme entre des
$1$-morphismes $\varphi, \varphi': \mathcal{F} \to \mathcal{G}$ de
catégories cofibrées en groupoïdes, on obtient un $2$-morphisme
$\widehat{t} : \widehat{\varphi} \to \widehat{\varphi}'$. Plus précisément,
pour $\xi = (R, \xi_n, f_n)$ comme ci-dessus, on pose
$\widehat{t}_\xi = (t_{\varphi(\xi_n)})$. La complétion définit donc un
foncteur entre $2$-catégories
$$
\widehat{~} :
\text{Cof}(\mathcal{C}_\Lambda)
\longrightarrow
\text{Cof}(\widehat{\mathcal{C}}_\Lambda)
$$
de la $2$-catégorie des catégories cofibrées en groupoïdes au-dessus de
$\mathcal{C}_\Lambda$ vers la $2$-catégorie des catégories cofibrées en
groupoïdes au-dessus de $\widehat{\mathcal{C}}_\Lambda$.
\end{remark}

\begin{remark}
\label{formal-defos-remark-completion-restriction-adjoint}
Nous affirmons que le foncteur de complétion de la remarque
\ref{formal-defos-remark-completion-morphism} et le foncteur de restriction
$|_{\mathcal{C}_\Lambda} : \text{Cof}(\widehat{\mathcal{C}}_\Lambda)
\to \text{Cof}(\mathcal{C}_\Lambda)$ des remarques
\ref{formal-defos-remarks-cofibered-groupoids}
(\ref{formal-defos-item-definition-restricting-base-category})
ont une relation de « 2-adjonction » au sens précis suivant. Soient
$\mathcal{F} \in \Ob(\text{Cof}(\mathcal{C}_\Lambda))$ et
$\mathcal{G} \in \Ob(\text{Cof}(\widehat{\mathcal{C}}_\Lambda))$.
Il existe alors une équivalence de catégories
$$
\Phi :
\Mor_{\mathcal{C}_\Lambda}(
\mathcal{G}|_{\mathcal{C}_\Lambda}, \mathcal{F})
\longrightarrow
\Mor_{\widehat{\mathcal{C}}_\Lambda}(\mathcal{G}, \widehat{\mathcal{F}})
$$
Pour décrire cette équivalence, définissons des morphismes canoniques
$\mathcal{G} \to \widehat{\mathcal{G}|_{\mathcal{C}_\Lambda}}$ et
$\widehat{\mathcal{F}}|_{\mathcal{C}_\Lambda} \to \mathcal{F}$ comme suit :
\begin{enumerate}
\item Soient $R \in \Ob(\widehat{\mathcal{C}}_\Lambda))$ et $\xi$ un objet
de la catégorie fibre $\mathcal{G}(R)$. Choisissons une image directe
$\xi \to \xi_n$ de $\xi$ vers $R/\mathfrak m_R^n$ pour chaque
$n \in \mathbf{N}$, et soit $f_n : \xi_{n + 1} \to \xi_n$ le morphisme
induit. Alors $\mathcal{G} \to \widehat{\mathcal{G}|_{\mathcal{C}_\Lambda}}$
envoie $\xi$ sur
$(R, \xi_n, f_n)$.
\item Il s'agit de l'équivalence
$can : \widehat{\mathcal{F}}|_{\mathcal{C}_\Lambda} \to \mathcal{F}$
de la remarque \ref{formal-defos-remark-restrict-completion}.
\end{enumerate}
Cela étant, l'équivalence
$\Phi : \Mor_{\mathcal{C}_\Lambda}(
\mathcal{G}|_{\mathcal{C}_\Lambda}, \mathcal{F})  \to
\Mor_{\widehat{\mathcal{C}}_\Lambda}(\mathcal{G},
\widehat{\mathcal{F}})$
envoie un morphisme
$\varphi : \mathcal{G}|_{\mathcal{C}_\Lambda} \to \mathcal{F}$
sur
$$
\mathcal{G} \to \widehat{\mathcal{G}|_{\mathcal{C}_\Lambda}}
\xrightarrow{\widehat{\varphi}} \widehat{\mathcal{F}}
$$
Il existe un quasi-inverse
$\Psi :
\Mor_{\widehat{\mathcal{C}}_\Lambda}(
\mathcal{G}, \widehat{\mathcal{F}}) \to
\Mor_{\mathcal{C}_\Lambda}(
\mathcal{G}|_{\mathcal{C}_\Lambda}, \mathcal{F})$
de $\Phi$ qui envoie $\psi : \mathcal{G} \to \widehat{\mathcal{F}}$ sur
$$
\mathcal{G}|_{\mathcal{C}_\Lambda} \xrightarrow{\psi|_{\mathcal{C}_\Lambda}}
\widehat{\mathcal{F}}|_{\mathcal{C}_\Lambda} \to \mathcal{F}.
$$
Nous omettons de vérifier que $\Phi$ et $\Psi$ sont quasi-inverses. Nous
n'abordons pas non plus la fonctorialité de $\Phi$ (car elle nous conduirait
sur le terrain des 3-catégories, que nous voulons éviter à tout prix).
\end{remark}

\begin{remark}
\label{formal-defos-remark-completion-restriction-cofset-adjoint}
Pour une catégorie $\mathcal{C}$, on note $\text{CofSet}(\mathcal{C})$ la
catégorie des ensembles cofibrés au-dessus de $\mathcal{C}$. C'est une
$1$-catégorie isomorphe à la catégorie des foncteurs
$\mathcal{C} \to \textit{Ensembles}$. Voir les remarques
\ref{formal-defos-remarks-cofibered-groupoids}
(\ref{formal-defos-item-convention-cofibered-sets}).
Les foncteurs de complétion et de restriction induisent des foncteurs
$\widehat{~} : \text{CofSet}(\mathcal{C}_\Lambda) \to
\text{CofSet}(\widehat{\mathcal{C}}_\Lambda)$ et
$|_{\mathcal{C}_\Lambda} : \text{CofSet}(\widehat{\mathcal{C}}_\Lambda) \to
\text{CofSet}(\mathcal{C}_\Lambda)$, que l'on désigne par les mêmes symboles.
Considérées comme foncteurs entre catégories d'ensembles cofibrés, la
complétion et la restriction sont adjointes au sens 1-catégorique usuel :
la même construction que dans la remarque
\ref{formal-defos-remark-completion-restriction-adjoint} définit une bijection fonctorielle
$$
\Mor_{\mathcal{C}_\Lambda}(G|_{\mathcal{C}_\Lambda}, F)
\longrightarrow
\Mor_{\widehat{\mathcal{C}}_\Lambda}(G, \widehat{F})
$$
pour $F \in \Ob(\text{CofSet}(\mathcal{C}_\Lambda))$ et
$G \in \Ob(\text{CofSet}(\widehat{\mathcal{C}}_\Lambda))$.
Le morphisme $\widehat{F}|_{\mathcal{C}_\Lambda} \to F$ est, là encore, un
isomorphisme.
\end{remark}

\begin{remark}
\label{formal-defos-remark-restrict-complete-continuous-functor}
Soit $G : \widehat{\mathcal{C}}_\Lambda \to \textit{Ensembles}$ un foncteur qui
commute aux limites. Alors le morphisme
$G \to \widehat{G|_{\mathcal{C}_\Lambda}}$ décrite dans la remarque
\ref{formal-defos-remark-completion-restriction-adjoint} est un isomorphisme. En effet, si
$S$ est un objet de $\widehat{\mathcal{C}}_\Lambda$, on dispose de bijections
canoniques
$$
\widehat{G|_{\mathcal{C}_\Lambda}}(S) =
\lim_n G(S/\mathfrak{m}_S^n) =
G(\lim_n S/\mathfrak{m}_S^n) = G(S).
$$
En particulier, si $R$ est un objet de $\widehat{\mathcal{C}}_\Lambda$, alors
$\underline{R} = \widehat{\underline{R}|_{\mathcal{C}_\Lambda}}$, car le
foncteur représentable $\underline{R}$ commute aux limites par définition des
limites.
\end{remark}

\begin{remark}
\label{formal-defos-remark-formal-objects-yoneda}
Soit $R$ un objet de $\widehat{\mathcal{C}}_\Lambda$. Il définit un foncteur
$\underline{R}: \widehat{\mathcal{C}}_\Lambda \to \textit{Ensembles}$
comme décrit dans les remarques \ref{formal-defos-remarks-cofibered-groupoids}
(\ref{formal-defos-item-definition-yoneda}). Comme d'habitude, on identifie ce foncteur à
l'ensemble cofibré associé. Si $\mathcal{F}$ est une catégorie cofibrée
au-dessus de $\mathcal{C}_\Lambda$, il existe une équivalence de catégories
\begin{equation}
\label{formal-defos-equation-formal-objects-maps}
\Mor_{\mathcal{C}_\Lambda}(
\underline{R}|_{\mathcal{C}_\Lambda}, \mathcal{F})
\longrightarrow
\widehat{\mathcal{F}}(R).
\end{equation}
Elle est donnée par la composée
$$
\Mor_{\mathcal{C}_\Lambda}(
\underline{R}|_{\mathcal{C}_\Lambda}, \mathcal{F})
\xrightarrow{\Phi}
\Mor_{\widehat{\mathcal{C}}_\Lambda}(
\underline{R}, \widehat{\mathcal{F}})
\xrightarrow{\sim}
\widehat{\mathcal{F}}(R)
$$
où $\Phi$ est celle de la remarque
\ref{formal-defos-remark-completion-restriction-adjoint}, et où la seconde équivalence
provient du lemme de 2-Yoneda (l'analogue cofibré de Catégories, lemme
\ref{categories-lemma-yoneda-2category}). Explicitement, l'équivalence envoie
un morphisme
$\varphi : \underline{R}|_{\mathcal{C}_\Lambda} \to \mathcal{F}$
sur l'objet formel
$(R, \varphi(R \to R/\mathfrak{m}_R^n), \varphi(f_n))$ dans
$\widehat{\mathcal{F}}(R)$, où
$f_n : R/\mathfrak m_R^{n + 1} \to R/\mathfrak m_R^n$ est la projection.

\medskip\noindent
Supposons qu'un choix d'images directes ait été fait pour $\mathcal{F}$.
Pour tout $\xi \in \Ob(\widehat{\mathcal{F}}(R))$, nous construisons
explicitement un morphisme
$\underline{\xi} : \underline{R}|_{\mathcal{C}_\Lambda} \to \mathcal{F}$
qui s'envoie sur $\xi$ par (\ref{formal-defos-equation-formal-objects-maps}). Écrivons
$\xi = (R, \xi_n, f_n)$. Se donner un objet $\alpha$ de
$\underline{R}|_{\mathcal{C}_\Lambda}$ revient à se donner un morphisme
$\alpha : R \to A$ de $\widehat{\mathcal{C}}_\Lambda$, avec $A$ artinien.
Soit $m \in \mathbf{N}$ le plus petit entier tel que
$\mathfrak m_A^m = 0$. Alors $\alpha$ se factorise par un unique
$\alpha_m : R/\mathfrak m_R^m \to A$, et l'on peut poser
$\underline{\xi}(\alpha) = \alpha_{m, *}\xi_m$. Nous omettons la description
de $\underline{\xi}$ sur les morphismes, ainsi que la démonstration du fait
que $\underline{\xi}$ s'envoie sur $\xi$ par
(\ref{formal-defos-equation-formal-objects-maps}).

\medskip\noindent
Supposons qu'un choix d'images directes ait été fait pour
$\widehat{\mathcal{F}}$. Dans ce cas, la démonstration de Catégories, lemme
\ref{categories-lemma-yoneda-2category}, fournit un quasi-inverse explicite
$$
\iota :
\widehat{\mathcal{F}}(R) \longrightarrow
\Mor_{\widehat{\mathcal{C}}_\Lambda}(
\underline{R}, \widehat{\mathcal{F}})
$$
de l'équivalence de 2-Yoneda, qui envoie $\xi$ sur le morphisme
$\iota(\xi) : \underline{R} \to \widehat{\mathcal{F}}$ envoyant
$f \in \underline{R}(S) = \Mor_{\mathcal{C}_\Lambda}(R, S)$
sur $f_*\xi$. Un quasi-inverse de (\ref{formal-defos-equation-formal-objects-maps}) est
donc
$$
\widehat{\mathcal{F}}(R)
\xrightarrow{\iota}
\Mor_{\widehat{\mathcal{C}}_\Lambda}(
\underline{R}, \widehat{\mathcal{F}})
\xrightarrow{\Psi}
\Mor_{\mathcal{C}_\Lambda}(
\underline{R}|_{\mathcal{C}_\Lambda}, \mathcal{F})
$$
où $\Psi$ est celle de la remarque
\ref{formal-defos-remark-completion-restriction-adjoint}. Pour
$\xi \in \Ob(\widehat{\mathcal{F}}(R))$, on a
$\Psi(\iota(\xi)) \cong \underline{\xi}$, où $\underline{\xi}$ est comme au
paragraphe précédent, puisque tous deux s'envoient sur $\xi$ par
l'équivalence de catégories (\ref{formal-defos-equation-formal-objects-maps}). En utilisant
$\underline{R} = \widehat{\underline{R}|_{\mathcal{C}_\Lambda}}$ (voir la
remarque \ref{formal-defos-remark-restrict-complete-continuous-functor}) et en développant
les définitions de $\Phi$ et $\Psi$, on conclut que $\iota(\xi)$ est isomorphe
au complété de $\underline{\xi}$.
\end{remark}

\begin{remark}
\label{formal-defos-remark-formal-objects-yoneda-map}
Soit $\mathcal{F}$ une catégorie cofibrée en groupoïdes au-dessus de
$\mathcal{C}_\Lambda$. Soient $\xi = (R, \xi_n, f_n)$ et
$\eta = (S, \eta_n, g_n)$ des objets formels de $\mathcal{F}$. Soit
$a = (a_n) : \xi \to \eta$ un morphisme d'objets formels, c'est-à-dire un
morphisme de $\widehat{\mathcal{F}}$. Soit
$f = \widehat{p}(a) = a_0 : R \to S$ la projection de $a$ dans
$\widehat{\mathcal{C}}_\Lambda$. On obtient alors un diagramme
$2$-commutatif
$$
\xymatrix{
\underline{R}|_{\mathcal{C}_\Lambda} \ar[rd]_{\underline{\xi}} & &
\underline{S}|_{\mathcal{C}_\Lambda} \ar[ll]^f \ar[ld]^{\underline{\eta}} \\
& \mathcal{F}
}
$$
où $\underline{\xi}$ et $\underline{\eta}$ sont les morphismes construits
dans la remarque \ref{formal-defos-remark-formal-objects-yoneda}. Pour le voir, soit
$\alpha : S \to A$ un objet de
$\underline{S}|_{\mathcal{C}_\Lambda}$ (voir loc.\ cit.). Soit
$m \in \mathbf{N}$ le plus petit entier tel que $\mathfrak m_A^m = 0$. On
obtient un diagramme commutatif
$$
\xymatrix{
R \ar[d]^f \ar[r] & R/\mathfrak m_R^m \ar[d]_{f_m} \ar[rd]^{\beta_m} \\
S \ar[r] & S/\mathfrak m_S^m \ar[r]^{\alpha_m} & A
}
$$
tel que la composée des flèches du bas soit $\alpha$. Alors
$\underline{\eta}(\alpha) = \alpha_{m, *}\eta_m$ et
$\underline{\xi}(\alpha \circ f) = \beta_{m, *}\xi_m$. Comme le morphisme
$a_m : \xi_m \to \eta_m$ est situé au-dessus de $f_m$, on obtient un
morphisme canonique
$$
\underline{\xi}(\alpha \circ f) = \beta_{m, *}\xi_m
\longrightarrow
\underline{\eta}(\alpha) = \alpha_{m, *}\eta_m
$$
situé au-dessus de $\text{id}_A$ et tel que
$$
\xymatrix{
\xi_m \ar[r] \ar[d]^{a_m} &  \beta_{m, *}\xi_m \ar[d] \\
\eta_m \ar[r] &  \alpha_{m, *}\eta_m
}
$$
soit commutatif en vertu des axiomes d'une catégorie cofibrée en groupoïdes.
Ceci définit une transformation de foncteurs
$\underline{\xi} \circ f \to \underline{\eta}$ qui témoigne de la
2-commutativité du premier diagramme de cette remarque.
\end{remark}

\begin{remark}
\label{formal-defos-remark-spell-out-formal-object}
Selon la remarque \ref{formal-defos-remark-formal-objects-yoneda}, se donner un objet
formel $\xi$ de $\mathcal{F}$ équivaut à se donner un foncteur
pro-représentable $U : \mathcal{C}_\Lambda \to \textit{Ensembles}$ et un morphisme
$U \to \mathcal{F}$.
\end{remark}




\section{Morphismes lisses}
\label{formal-defos-section-smooth-morphisms}

\noindent
Dans cette section, nous étudions les morphismes lisses de catégories
cofibrées en groupoïdes au-dessus de $\mathcal{C}_\Lambda$.

\begin{definition}
\label{formal-defos-definition-smooth-morphism}
Soit $\varphi : \mathcal{F} \to \mathcal{G}$ un morphisme de catégories
cofibrées en groupoïdes au-dessus de $\mathcal{C}_\Lambda$. On dit que
$\varphi$ est {\it lisse} s'il satisfait la condition suivante. Soient
$B \to A$ un morphisme surjectif d'anneaux dans $\mathcal{C}_\Lambda$,
$y \in
\Ob(\mathcal{G}(B)), x \in \Ob(\mathcal{F}(A))$, et $y
\to \varphi(x)$ un morphisme situé au-dessus de $B \to A$. Alors il
existe $x' \in \Ob(\mathcal{F}(B))$, un morphisme $x' \to x$ situé au-dessus
de $B \to A$, et un morphisme $\varphi(x') \to y$ situé au-dessus de
$\text{id}: B \to B$, tels que le diagramme
$$
\xymatrix{
\varphi(x') \ar[r] \ar[dr] & y \ar[d] \\
& \varphi(x)
}
$$
soit commutatif.
\end{definition}

\begin{lemma}
\label{formal-defos-lemma-smoothness-small-extensions}
Soit $\varphi : \mathcal{F} \to \mathcal{G}$ un morphisme de catégories
cofibrées en groupoïdes au-dessus de $\mathcal{C}_\Lambda$. Alors $\varphi$
est lisse dès que la condition de la définition
\ref{formal-defos-definition-smooth-morphism} est supposée satisfaite seulement pour les
petites extensions $B \to A$.
\end{lemma}

\begin{proof}
Soient $B \to A$ un morphisme surjectif d'anneaux dans
$\mathcal{C}_\Lambda$, $y \in \Ob(\mathcal{G}(B))$,
$x \in \Ob(\mathcal{F}(A))$, et $y \to \varphi(x)$ un morphisme situé
au-dessus de $B \to A$. D'après le lemme
\ref{formal-defos-lemma-factor-small-extension}, on peut factoriser $B \to A$ en petites
extensions $B = B_n \to B_{n-1} \to \ldots \to B_0 = A$. Raisonnons par
récurrence sur $n$. Si $n = 1$, le résultat découle de l'hypothèse. Si
$n > 1$, notons $f : B = B_n \to B_{n - 1}$ et
$g : B_{n - 1} \to B_0 = A$. Choisissons une image directe $y \to f_* y$
de $y$ le long de $f$, de sorte que le morphisme $y \to \varphi(x)$ se
factorise en $y \to f_* y \to \varphi(x)$. Par l'hypothèse de récurrence, on
peut trouver $x_{n - 1} \to x$ situé au-dessus de
$g : B_{n - 1} \to A$ et $a : \varphi(x_{n - 1}) \to f_*y$ situé au-dessus
de $\text{id} : B_{n - 1} \to B_{n - 1}$ tels que
$$
\xymatrix{
\varphi(x_{n - 1}) \ar[r]_-a \ar[dr] & f_*y \ar[d] \\
& \varphi(x)
}
$$
soit commutatif. On peut appliquer l'hypothèse à la composée
$y \to \varphi(x_{n - 1})$ de $y \to f_*y$ et de
$a^{-1} : f_*y \to \varphi(x_{n - 1})$. On obtient
$x_n \to x_{n - 1}$ situé au-dessus de $B_n \to B_{n - 1}$ et
$\varphi(x_n) \to y$ situé au-dessus de $\text{id} : B_n \to B_n$, de sorte
que le diagramme
$$
\xymatrix{
\varphi(x_n) \ar[r] \ar[d] & y \ar[d] \\
\varphi(x_{n - 1}) \ar[r]^-a \ar[dr] & f_*y \ar[d] \\
& \varphi(x)
}
$$
soit commutatif. La composée $x_n \to x_{n - 1} \to x$ et le morphisme
$\varphi(x_n) \to y$ sont alors les morphismes exigés par la définition de la
lissité.
\end{proof}

\begin{remark}
\label{formal-defos-remark-smoothness-2-categorical}
Soit $\varphi : \mathcal{F} \to \mathcal{G}$ un morphisme de catégories
cofibrées en groupoïdes au-dessus de $\mathcal{C}_\Lambda$. Soit $B \to A$
un morphisme d'anneaux dans $\mathcal{C}_\Lambda$. Des choix d'images directes
le long de $B
\to A$ pour les objets des catégories fibres $\mathcal{F}(B)$
et $\mathcal{G}(B)$ déterminent des foncteurs
$\mathcal{F}(B) \to \mathcal{F}(A)$ et
$\mathcal{G}(B) \to \mathcal{G}(A)$ qui s'insèrent dans un diagramme
$2$-commutatif
$$
\xymatrix{
\mathcal{F}(B) \ar[r]^{\varphi} \ar[d] & \mathcal{G}(B) \ar[d] \\
\mathcal{F}(A) \ar[r]^{\varphi}        & \mathcal{G}(A) .
}
$$
Il en résulte un foncteur induit $\mathcal{F}(B) \to \mathcal{F}(A)
\times_{\mathcal{G}(A)} \mathcal{G}(B)$. En développant les définitions, on
voit que $\varphi : \mathcal{F} \to \mathcal{G}$ est lisse si et seulement
si ce foncteur induit est essentiellement surjectif chaque fois que
$B \to A$ est surjectif (ou, de manière équivalente d'après le lemme
\ref{formal-defos-lemma-smoothness-small-extensions}, chaque fois que $B \to A$ est une
petite extension).
\end{remark}

\begin{remark}
\label{formal-defos-remark-compare-smooth-schlessinger}
La caractérisation des morphismes lisses donnée dans la remarque
\ref{formal-defos-remark-smoothness-2-categorical} est analogue à la notion de morphisme
lisse de foncteurs due à Schlessinger ; cf.\ \cite[définition 2.2.]{Sch}. En
fait, lorsque $\mathcal{F}$ et $\mathcal{G}$ sont cofibrées en ensembles,
notre notion équivaut à celle de Schlessinger. Soient en effet, dans ce cas,
$F, G : \mathcal{C}_\Lambda \to \textit{Ensembles}$ les foncteurs correspondants ;
voir les remarques \ref{formal-defos-remarks-cofibered-groupoids}
(\ref{formal-defos-item-convention-cofibered-sets}).
Alors $F \to G$ est lisse si et seulement si, pour toute surjection d'anneaux
$B \to A$ dans $\mathcal{C}_\Lambda$, l'application
$F(B) \to F(A) \times_{G(A)} G(B)$ est surjective.
\end{remark}

\begin{remark}
\label{formal-defos-remark-smooth-to-iso-classes}
Soit $\mathcal{F}$ une catégorie cofibrée en groupoïdes au-dessus de
$\mathcal{C}_\Lambda$. Alors le morphisme
$\mathcal{F} \to \overline{\mathcal{F}}$ est lisse. Supposons en effet que
$f : B \to A$ soit un morphisme d'anneaux dans $\mathcal{C}_\Lambda$. Soient
$x \in \Ob(\mathcal{F}(A))$ et
$\overline{y} \in \overline{\mathcal{F}}(B)$
la classe d'isomorphie de $y \in \Ob(\mathcal{F}(B))$, telle que
$\overline{f_*y} = \overline{x}$. Il suffit alors de prendre $x' = y$, le
morphisme induit $x' = y \to x$ au-dessus de $B \to A$, et l'égalité
$\overline{x'} = \overline{y}$ pour résoudre le problème posé dans la
définition \ref{formal-defos-definition-smooth-morphism}.
\end{remark}

\noindent
Si $R \to S$ est un morphisme d'anneaux de
$\widehat{\mathcal{C}}_\Lambda$, il induit un morphisme
$\underline{S} \to \underline{R}$ entre les foncteurs
$\underline{S}, \underline{R}: \widehat{\mathcal{C}}_\Lambda
\to \textit{Ensembles}$. Dans cette situation, la lissité de la restriction
$\underline{S}|_{\mathcal{C}_\Lambda} \to
\underline{R}|_{\mathcal{C}_\Lambda}$ est une notion familière :

\begin{lemma}
\label{formal-defos-lemma-smooth-morphism-power-series}
Soit $R \to S$ un morphisme d'anneaux dans
$\widehat{\mathcal{C}}_\Lambda$. Alors le morphisme induit
$\underline{S}|_{\mathcal{C}_\Lambda} \to \underline{R}|_{\mathcal{C}_\Lambda}$
est lisse si et seulement si $S$ est un anneau de séries formelles sur $R$.
\end{lemma}

\begin{proof}
Supposons que $S$ soit un anneau de séries formelles sur $R$, disons
$S = R[[x_1, \ldots, x_n]]$. La lissité de
$\underline{S}|_{\mathcal{C}_\Lambda} \to \underline{R}|_{\mathcal{C}_\Lambda}$
signifie ce qui suit (voir la remarque
\ref{formal-defos-remark-compare-smooth-schlessinger}). Étant donnés un morphisme
surjectif d'anneaux $B \to A$ dans $\mathcal{C}_\Lambda$, un morphisme
d'anneaux $R \to B$ et un morphisme d'anneaux $S \to A$ tels que le diagramme
en traits pleins
$$
\xymatrix{
S \ar[r] \ar@{..>}[rd] & A \\
R \ar[u] \ar[r] & B \ar[u]
}
$$
soit commutatif, il existe une flèche en pointillés qui rend le diagramme
commutatif. (Remarquer la similitude avec Algèbre, définition
\ref{algebra-definition-formally-smooth}.) Pour construire la flèche en
pointillés, choisissons des éléments $b_i \in B$ dont les images dans $A$
sont égales aux images de $x_i$ dans $A$. Remarquons que
$b_i \in \mathfrak m_B$, puisque $x_i$ s'envoie sur un élément de
$\mathfrak m_A$. Il existe donc un unique morphisme de $R$-algèbres
$R[[x_1, \ldots, x_n]] \to B$ qui envoie $x_i$ sur $b_i$ et qui peut servir
de flèche en pointillés.

\medskip\noindent
Réciproquement, supposons que
$\underline{S}|_{\mathcal{C}_\Lambda} \to \underline{R}|_{\mathcal{C}_\Lambda}$
soit lisse. Soient $x_1, \ldots, x_n \in S$ des éléments dont les images
forment une base de l'espace cotangent relatif
$\mathfrak m_S/(\mathfrak m_R S + \mathfrak m_S^2)$ de $S$ sur $R$.
Posons $T = R[[X_1, \ldots, X_n]]$. Remarquons que l'on a à la fois
$$
S/(\mathfrak m_R S + \mathfrak m_S^2) \cong
R/\mathfrak m_R[x_1, \ldots, x_n]/(x_ix_j)
$$
et
$$
T/(\mathfrak m_R T + \mathfrak m_T^2) \cong
R/\mathfrak m_R[X_1, \ldots, X_n]/(X_iX_j).
$$
Soit
$S/(\mathfrak m_R S + \mathfrak m_S^2) \to
T/(\mathfrak m_R T + \mathfrak m_T^2)$
l'isomorphisme de $R$-algèbres locales qui envoie la classe de $x_i$ sur la
classe de $X_i$. Soit
$f_1 : S \to T/(\mathfrak m_R T + \mathfrak m_T^2)$ la composée
$S \to S/(\mathfrak m_R S + \mathfrak m_S^2)
\to T/(\mathfrak m_R T + \mathfrak m_T^2)$.
L'hypothèse selon laquelle
$\underline{S}|_{\mathcal{C}_\Lambda} \to \underline{R}|_{\mathcal{C}_\Lambda}$
est lisse signifie que l'on peut relever $f_1$ en un morphisme
$f_2 : S \to T/\mathfrak{m}_T^2$, puis en un morphisme
$f_3 : S \to T/\mathfrak{m}_T^3$, et ainsi de suite pour tout $n \geq 1$.
On obtient donc un morphisme induit $f : S \to T = \lim T/\mathfrak m_T^n$
de $R$-algèbres locales. Par notre choix de $f_1$, le morphisme $f$ induit
un isomorphisme
$\mathfrak m_S/(\mathfrak m_R S + \mathfrak m_S^2) \to
\mathfrak m_T/(\mathfrak m_R T + \mathfrak m_T^2)$
d'espaces cotangents relatifs. Le lemme
\ref{formal-defos-lemma-surjective-cotangent-space} montre donc que $f$ est surjective
(où l'on considère $f$ comme un morphisme dans
$\widehat{\mathcal{C}}_R$). Choisissons des antécédents $y_i \in S$ des
$X_i \in T$ par $f$. Comme $T$ est un anneau de séries formelles sur $R$, il
existe un homomorphisme de $R$-algèbres locales $s : T \to S$ envoyant $X_i$
sur $y_i$. Par construction, $f \circ s = \text{id}$ ; donc $s$ est
injectif. Mais $s$ induit un isomorphisme sur les espaces cotangents relatifs,
puisque $f$ en induit un ; il est donc également surjectif, de nouveau par le
lemme \ref{formal-defos-lemma-surjective-cotangent-space}. Ainsi $s$ et $f$ sont des
isomorphismes.
\end{proof}

\noindent
Les morphismes lisses possèdent les propriétés fonctorielles suivantes.

\begin{lemma}
\label{formal-defos-lemma-smooth-properties}
Soient $\varphi : \mathcal{F} \to \mathcal{G}$ et $\psi : \mathcal{G}
\to \mathcal{H}$ des morphismes de catégories cofibrées en groupoïdes
au-dessus de $\mathcal{C}_\Lambda$.
\begin{enumerate}
\item Si $\varphi$ et $\psi$ sont lisses, alors $\psi \circ \varphi$ est
lisse.
\item Si $\varphi$ est essentiellement surjectif et si
$\psi \circ \varphi$ est lisse, alors $\psi$ est lisse.
\item Si $\mathcal{G}' \to \mathcal{G}$ est un morphisme de catégories
cofibrées en groupoïdes et si $\varphi$ est lisse, alors
$\mathcal{F} \times_\mathcal{G} \mathcal{G}' \to \mathcal{G}'$ est lisse.
\end{enumerate}
\end{lemma}

\begin{proof}
Les assertions (1) et (2) résultent immédiatement des définitions. Nous
omettons la démonstration de (3). Indications : utiliser la formulation de la
lissité donnée dans la remarque \ref{formal-defos-remark-smoothness-2-categorical} ainsi
que le fait que $\mathcal{F} \times_\mathcal{G} \mathcal{G}'$ est le
$2$-produit fibré ; voir les remarques \ref{formal-defos-remarks-cofibered-groupoids}
(\ref{formal-defos-item-fibre-product}).
\end{proof}

\begin{lemma}
\label{formal-defos-lemma-smooth-morphism-essentially-surjective}
Soit $\varphi : \mathcal{F} \to \mathcal{G}$ un morphisme lisse de
catégories cofibrées en groupoïdes au-dessus de $\mathcal{C}_\Lambda$.
Supposons que $\varphi : \mathcal{F}(k) \to \mathcal{G}(k)$ soit
essentiellement surjectif. Alors $\varphi : \mathcal{F} \to \mathcal{G}$ et
$\widehat{\varphi} : \widehat{\mathcal{F}} \to \widehat{\mathcal{G}}$
sont essentiellement surjectifs.
\end{lemma}

\begin{proof}
Soit $y$ un objet de $\mathcal{G}$ situé au-dessus de
$A \in \Ob(\mathcal{C}_\Lambda)$. Soit $y \to y_0$ une image directe de $y$
le long de $A \to k$. L'hypothèse de surjectivité essentielle de
$\varphi : \mathcal{F}(k) \to \mathcal{G}(k)$ donne un objet $x_0$ de
$\mathcal{F}$ situé au-dessus de $k$ et un isomorphisme
$y_0 \to \varphi(x_0)$. La lissité de $\varphi$ implique qu'il existe un
objet $x$ de $\mathcal{F}$ au-dessus de $A$ dont l'image $\varphi(x)$ est
isomorphe à $y$. Ainsi, $\varphi : \mathcal{F} \to \mathcal{G}$ est
essentiellement surjectif.

\medskip\noindent
Soit $\eta = (R, \eta_n, g_n)$ un objet de $\widehat{\mathcal{G}}$.
Construisons un objet $\xi$ de $\widehat{\mathcal{F}}$ muni d'un isomorphisme
$\eta \to \varphi(\xi)$. Par l'hypothèse de surjectivité essentielle de
$\varphi : \mathcal{F}(k) \to \mathcal{G}(k)$, il existe un morphisme
$\eta_1 \to \varphi(\xi_1)$ dans $\mathcal{G}(k)$ pour un certain
$\xi_1 \in \Ob(\mathcal{F}(k))$. Le morphisme
$\eta_2 \xrightarrow{g_1} \eta_1 \to \varphi(\xi_1)$
est situé au-dessus du morphisme surjectif d'anneaux
$R/\mathfrak m_R^2 \to k$ ; la lissité de $\varphi$ donne donc un objet
$\xi_2 \in \Ob(\mathcal{F}(R/\mathfrak m_R^2))$, un morphisme
$f_1: \xi_2 \to \xi_1$ situé au-dessus de $R/\mathfrak m_R^2 \to k$, et un
morphisme $\eta_2 \to \varphi(\xi_2)$ tels que
$$
\xymatrix{
\varphi(\xi_2)  \ar[r]^{\varphi(f_1)} &  \varphi(\xi_{1})   \\
\eta_2   \ar[u] \ar[r]^{g_1}  & \eta_1  \ar[u] \\
}
$$
soit commutatif. En poursuivant ainsi, on construit un objet
$\xi = (R, \xi_n, f_n)$ de $\widehat{\mathcal{F}}$ et un morphisme
$\eta \to \varphi(\xi) = (R, \varphi(\xi_n), \varphi(f_n))$ dans
$\widehat{\mathcal{G}}(R)$.
\end{proof}

\noindent
Nous chercherons plus loin à construire des morphismes lisses de foncteurs
pro-représentables vers des catégories de prédéformation $\mathcal{F}$.
D'après la discussion de la remarque \ref{formal-defos-remark-formal-objects-yoneda}, ces
morphismes correspondent à certains objets formels de $\mathcal{F}$ : ce
sont, plus précisément, les objets formels dits versels de $\mathcal{F}$.

\begin{definition}
\label{formal-defos-definition-versal}
Soit $\mathcal{F}$ une catégorie cofibrée en groupoïdes. Soit $\xi$ un objet
formel de $\mathcal{F}$ situé au-dessus de
$R \in \Ob(\widehat{\mathcal{C}}_\Lambda)$. On dit que $\xi$ est {\it versel}
si le morphisme correspondant
$\underline{\xi}: \underline{R}|_{\mathcal{C}_\Lambda} \to \mathcal{F}$
de la remarque \ref{formal-defos-remark-formal-objects-yoneda} est lisse.
\end{definition}

\begin{remark}
\label{formal-defos-remark-versal-object}
Soit $\mathcal{F}$ une catégorie cofibrée en groupoïdes au-dessus de $\mathcal
C_\Lambda$, et soit $\xi$ un objet formel de $\mathcal{F}$. Il
résulte de la définition de la lissité que la versalité de $\xi$ équivaut à la
condition suivante. Si
$$
\xymatrix{
& y \ar[d] \\
\xi \ar[r] & x
}
$$
est un diagramme dans $\widehat{\mathcal{F}}$ tel que $y \to x$ soit situé
au-dessus d'un morphisme surjectif $B \to A$ d'anneaux artiniens (on peut
supposer qu'il s'agit d'une petite extension), alors il existe un morphisme
$\xi \to y$ tel que
$$
\xymatrix{
& y \ar[d] \\
\xi \ar[r] \ar[ur] & x
}
$$
soit commutatif. En particulier, la condition pour $\xi$ d'être versel ne
dépend pas des choix d'images directes effectués lors de la construction de
$\underline{\xi} : \underline{R}|_{\mathcal{C}_\Lambda} \to \mathcal{F}$
dans la remarque \ref{formal-defos-remark-formal-objects-yoneda}.
\end{remark}

\begin{lemma}
\label{formal-defos-lemma-versal-object-quasi-initial}
Soit $\mathcal{F}$ une catégorie de prédéformation. Soit $\xi$ un objet
formel versel de $\mathcal{F}$. Pour tout objet formel $\eta$ de
$\widehat{\mathcal{F}}$, il existe un morphisme $\xi \to \eta$.
\end{lemma}

\begin{proof}
Par hypothèse, le morphisme
$\underline{\xi} : \underline{R}|_{\mathcal{C}_\Lambda} \to \mathcal{F}$
est lisse. Alors
$\iota(\xi) : \underline{R} \to \widehat{\mathcal{F}}$
est le complété de $\underline{\xi}$ ; voir la remarque
\ref{formal-defos-remark-formal-objects-yoneda}. D'après le lemme
\ref{formal-defos-lemma-smooth-morphism-essentially-surjective}, il existe un objet $f$ de
$\underline{R}$ tel que $\iota(\xi)(f) = \eta$. Ainsi, $f$ est un morphisme
d'anneaux $f : R \to S$ dans $\widehat{\mathcal{C}}_\Lambda$. L'égalité
$\iota(\xi)(f) = \eta$ signifie que $f_*\xi \cong \eta$, ce qui signifie
exactement qu'il existe un morphisme $\xi \to \eta$ situé au-dessus de $f$.
\end{proof}





\section{Catégories lisses ou non obstruées}
\label{formal-defos-section-smooth}

\noindent
Soit $p : \mathcal{F} \to \mathcal{C}_\Lambda$ une catégorie cofibrée en
groupoïdes. On peut considérer $\mathcal{C}_\Lambda$ comme une catégorie
cofibrée en groupoïdes au-dessus de $\mathcal{C}_\Lambda$ au moyen du foncteur
identité. De cette manière,
$p : \mathcal{F} \longrightarrow \mathcal{C}_\Lambda$
devient un morphisme de catégories cofibrées en groupoïdes au-dessus de
$\mathcal{C}_\Lambda$.

\begin{definition}
\label{formal-defos-definition-cofibered-groupoid-projection-smooth}
Soit $p : \mathcal{F} \to \mathcal{C}_\Lambda$ une catégorie cofibrée en
groupoïdes. On dit que $\mathcal{F}$ est {\it lisse} ou {\it non obstruée} si
son morphisme structural $p$ est lisse au sens de la définition
\ref{formal-defos-definition-smooth-morphism}.
\end{definition}

\noindent
C'est la notion « absolue » de lissité pour une catégorie cofibrée en
groupoïdes au-dessus de $\mathcal{C}_\Lambda$, bien qu'il soit plus correct de
dire que $\mathcal{F}$ est lisse sur $\Lambda$. Il faut manier avec prudence
l'expression « {\it $\mathcal{F}$ est non obstruée} » : il peut arriver que
$\mathcal{F}$ possède une théorie de l'obstruction dont les espaces
d'obstructions ne sont pas nuls, alors même que $\mathcal{F}$ est lisse.

\begin{remark}
\label{formal-defos-remark-smooth-on-top}
Supposons que $\mathcal{F}$ soit une catégorie de prédéformation admettant un
morphisme lisse $\varphi : \mathcal U \to \mathcal{F}$ depuis une catégorie
de prédéformation $\mathcal{U}$. D'après le lemme
\ref{formal-defos-lemma-smooth-morphism-essentially-surjective}, $\varphi$ est
essentiellement surjectif ; le lemme \ref{formal-defos-lemma-smooth-properties} montre donc
que $p: \mathcal{F} \to \mathcal{C}_\Lambda$ est lisse si et seulement si la
composée $\mathcal U \xrightarrow{\varphi} \mathcal{F} \xrightarrow{p}
\mathcal{C}_\Lambda$ est lisse, c'est-à-dire\ que $\mathcal{F}$ est lisse si
et seulement si $\mathcal{U}$ est lisse.
\end{remark}

\begin{lemma}
\label{formal-defos-lemma-smooth}
Soit $R \in \Ob(\widehat{\mathcal{C}}_\Lambda)$. Les conditions suivantes
sont équivalentes :
\begin{enumerate}
\item $\underline{R}|_{\mathcal{C}_\Lambda}$ est lisse ;
\item $\Lambda \to R$ est formellement lisse pour la topologie
$\mathfrak m_R$-adique ;
\item $\Lambda \to R$ est plat et $R \otimes_\Lambda k'$ est
géométriquement régulier sur $k'$ ;
\item $\Lambda \to R$ est plat et $k' \to R \otimes_\Lambda k'$ est
formellement lisse pour la topologie $\mathfrak m_R$-adique.
\end{enumerate}
Dans le cas classique, elles équivalent également à :
\begin{enumerate}
\item[(5)] $R$ est isomorphe à $\Lambda[[x_1, \ldots, x_n]]$ pour un certain
$n$.
\end{enumerate}
\end{lemma}

\begin{proof}
La lissité de
$p : \underline{R}|_{\mathcal{C}_\Lambda} \to \mathcal{C}_\Lambda$
signifie qu'étant donnés $B \to A$ surjectif dans $\mathcal{C}_\Lambda$ et
$R \to A$, on peut trouver la flèche en pointillés dans le diagramme
$$
\xymatrix{
R \ar[r] \ar@{-->}[rd] & A \\
\Lambda \ar[r] \ar[u] & B \ar[u]
}
$$
Cela est certainement vrai si $\Lambda \to R$ est formellement lisse pour la
topologie $\mathfrak m_R$-adique ; voir Compléments d'algèbre, définitions
\ref{more-algebra-definition-formally-smooth-adic} et
\ref{more-algebra-definition-formally-smooth}. Réciproquement, si cela est
vrai, Compléments d'algèbre, lemme \ref{more-algebra-lemma-fs-local}, montre
que $\Lambda \to R$ est formellement lisse pour la topologie
$\mathfrak m_R$-adique. Ainsi, (1) et (2) sont équivalentes.

\medskip\noindent
L'équivalence de (2), (3) et (4) est Compléments d'algèbre, proposition
\ref{more-algebra-proposition-fs-flat-fibre-fs}. L'équivalence avec (5)
résulte, par exemple, du lemme \ref{formal-defos-lemma-smooth-morphism-power-series} et du
fait que, dans le cas classique, $\mathcal{C}_\Lambda$ s'identifie à
$\underline{\Lambda}|_{\mathcal{C}_\Lambda}$.
\end{proof}

\begin{lemma}
\label{formal-defos-lemma-smooth-power-series-classical}
Soit $\mathcal{F}$ une catégorie de prédéformation. Soit $\xi$ un objet
formel versel de $\mathcal{F}$ situé au-dessus de
$R \in \Ob(\widehat{\mathcal{C}}_\Lambda)$. Les conditions suivantes sont
équivalentes :
\begin{enumerate}
\item $\mathcal{F}$ est non obstruée ;
\item $\Lambda \to R$ est formellement lisse pour la topologie
$\mathfrak m_R$-adique.
\end{enumerate}
Dans le cas classique, elles équivalent également à :
\begin{enumerate}
\item[(3)] $R \cong \Lambda[[x_1, \ldots, x_n]]$ pour un certain $n$.\
\end{enumerate}
\end{lemma}

\begin{proof}
Si (1) est vérifiée, c'est-à-dire si $\mathcal{F}$ est non obstruée, alors la
composée
$$
\underline{R}|_{\mathcal{C}_\Lambda}
\xrightarrow{\underline{\xi}}
\mathcal{F}
\to
\mathcal{C}_\Lambda
$$
est lisse ; voir le lemme \ref{formal-defos-lemma-smooth-properties}. Le lemme
\ref{formal-defos-lemma-smooth} donne donc (2). Réciproquement, si (2) est vérifiée, la
composée est lisse ; de plus, la première flèche est essentiellement
surjective d'après le lemme \ref{formal-defos-lemma-versal-object-quasi-initial}. Le lemme
\ref{formal-defos-lemma-smooth-properties} montre donc que la seconde flèche est lisse, ce
qui signifie par définition que $\mathcal{F}$ est non obstruée. Dans le cas
classique, l'équivalence avec (3) résulte du lemme \ref{formal-defos-lemma-smooth}.
\end{proof}

\begin{lemma}
\label{formal-defos-lemma-exists-smooth}
Il existe $R \in \Ob(\widehat{\mathcal{C}}_\Lambda)$ tel que les conditions
équivalentes du lemme \ref{formal-defos-lemma-smooth} soient satisfaites et que, de plus,
$H_1(L_{k/\Lambda}) = \mathfrak m_R/\mathfrak m_R^2$ et
$\Omega_{R/\Lambda} \otimes_R k = \Omega_{k/\Lambda}$.
\end{lemma}

\begin{proof}
Dans le cas classique, on choisit $R = \Lambda$. Plus généralement, si
l'extension de corps résiduels $k/k'$ est séparable, il existe une unique
extension finie étale $\Lambda^\wedge \to R$ (Algèbre, lemmes
\ref{algebra-lemma-complete-henselian} et
\ref{algebra-lemma-henselian-cat-finite-etale}) du complété
$\Lambda^\wedge$ de $\Lambda$ qui induit l'extension $k/k'$ sur les corps
résiduels.

\medskip\noindent
Dans le cas général, procédons comme suit. Choisissons une $\Lambda$-algèbre
lisse $P$ et une surjection de $\Lambda$-algèbres $P \to k$. (On peut, par
exemple, prendre pour $P$ une algèbre de polynômes.) Notons $\mathfrak m_P$ le
noyau de $P \to k$. La suite de Jacobi--Zariski, voir
(\ref{formal-defos-equation-sequence-extended}) et Algèbre, lemme
\ref{algebra-lemma-exact-sequence-NL}, est une suite exacte
$$
0 \to H_1(\NL_{k/\Lambda}) \to
\mathfrak m_P/\mathfrak m_P^2 \to
\Omega_{P/\Lambda} \otimes_P k \to
\Omega_{k/\Lambda} \to 0
$$
Le $0$ de gauche provient de ce que $P/k$ est lisse :
$\NL_{P/\Lambda}$ est alors quasi-isomorphe à un module projectif de type fini placé
en degré $0$, et donc $H_1(\NL_{P/\Lambda} \otimes_P k) = 0$. Supposons que
$f \in \mathfrak m_P$ s'envoie sur un élément non nul de
$\Omega_{P/\Lambda} \otimes_P k$. En posant $P' = P/(f)$, on dispose d'une
surjection de $\Lambda$-algèbres $P' \to k$. Remarquons que $P'$ est lisse en
$\mathfrak m_{P'}$ ; cela résulte de Compléments sur les morphismes, lemme
\ref{more-morphisms-lemma-slice-smooth-given-element}. Après avoir remplacé
$P$ par une localisation principale de $P'$, on voit donc que
$\dim(\mathfrak m_P/\mathfrak m_P^2)$ diminue. En répétant l'opération un
nombre fini de fois, on peut supposer que l'application
$\mathfrak m_P/\mathfrak m_P^2 \to
\Omega_{P/\Lambda} \otimes_P k$ est nulle, de sorte que la suite exacte se
décompose en les deux isomorphismes
$H_1(L_{k/\Lambda}) = \mathfrak m_P/\mathfrak m_P^2$ et
$\Omega_{P/\Lambda} \otimes_P k = \Omega_{k/\Lambda}$.

\medskip\noindent
Soit $R$ le complété $\mathfrak m_P$-adique de $P$. Alors $R$ est un objet de
$\widehat{\mathcal{C}}_\Lambda$ : c'est un anneau local noethérien complet
(voir Algèbre, lemme
\ref{algebra-lemma-completion-Noetherian-Noetherian}) dont le corps résiduel
s'identifie à $k$. Nous affirmons que $R$ convient.

\medskip\noindent
Remarquons d'abord que le morphisme $P \to R$ induit les isomorphismes
$\mathfrak m_P/\mathfrak m_P^2 = \mathfrak m_R/\mathfrak m_R^2$ et
$\Omega_{P/\Lambda} \otimes_P k = \Omega_{R/\Lambda} \otimes_R k$. En effet,
$\mathfrak m_P/\mathfrak m_P^2$ et
$\Omega_{P/\Lambda} \otimes_P k$ ne dépendent que de la $\Lambda$-algèbre
$P/\mathfrak m_P^2$ ; voir Algèbre, lemme
\ref{algebra-lemma-differential-mod-power-ideal}. Il en va de même pour $R$,
et l'on a $P/\mathfrak m_P^2 = R/\mathfrak m_R^2$. La fonctorialité de la
suite de Jacobi--Zariski (\ref{formal-defos-equation-sequence-functorial}) permet d'en
déduire $H_1(L_{k/\Lambda}) = \mathfrak m_R/\mathfrak m_R^2$ et
$\Omega_{R/\Lambda} \otimes_R k = \Omega_{k/\Lambda}$, puisque ces égalités
sont vraies pour $P$.

\medskip\noindent
Enfin, puisque $\Lambda \to P$ est lisse, Algèbre, proposition
\ref{algebra-proposition-smooth-formally-smooth}, montre que
$\Lambda \to P$ est formellement lisse. Alors $\Lambda \to P$ est
formellement lisse pour
la topologie $\mathfrak m_P$-adique d'après Compléments d'algèbre, lemme
\ref{more-algebra-lemma-formally-smooth}. Le complété $R$ hérite de cette
propriété par Compléments d'algèbre, lemme
\ref{more-algebra-lemma-formally-smooth-completion}, ce qui achève la
démonstration. En fait, on peut montrer que, chaque fois que
$\underline{R}|_{\mathcal{C}_\Lambda}$ est lisse, $R$ est isomorphe au
complété d'une algèbre lisse sur $\Lambda$, mais nous n'utiliserons pas ce
fait.
\end{proof}

\begin{example}
\label{formal-defos-example-smooth-explicit}
Voici un exemple plus explicite d'un $R$ comme dans le lemme
\ref{formal-defos-lemma-exists-smooth}. Soient $p$ un nombre premier et
$n \in \mathbf{N}$. Posons
$\Lambda = \mathbf{F}_p(t_1, t_2, \ldots, t_n)$ et
$k = \mathbf{F}_p(x_1, \ldots, x_n)$, l'application $\Lambda \to k$ étant
donnée par $t_i \mapsto x_i^p$. On peut alors prendre
$$
R = \Lambda[x_1, \ldots, x_n]^\wedge_{(x_1^p - t_1, \ldots, x_n^p - t_n)}
$$
On ne peut pas faire « mieux » dans cet exemple, c'est-à-dire que l'on ne
peut pas approcher $\mathcal{C}_\Lambda$ par un objet lisse plus petit de
$\widehat{\mathcal{C}}_\Lambda$ (on peut montrer que la dimension de $R$ doit
être au moins $n$, puisque l'application
$\Omega_{R/\Lambda} \otimes_R k \to \Omega_{k/\Lambda}$ est surjective). Nous
reviendrons plus loin sur ce phénomène de manière plus détaillée.
\end{example}







\section{Conditions de Schlessinger}
\label{formal-defos-section-schlessinger-conditions}

\noindent
Dans la suite, nous considérerons souvent des produits fibrés
$A_1 \times_A A_2$ d'anneaux de la catégorie $\mathcal{C}_\Lambda$. Nous
avons vu dans l'exemple \ref{formal-defos-example-fibre-product} qu'un tel produit fibré
n'est pas nécessairement un objet de $\mathcal{C}_\Lambda$. Toutefois, dans
presque tous les cas ci-dessous, l'un des deux morphismes $A_i \to A$ est
surjectif, et $A_1 \times_A A_2$ est alors un objet de
$\mathcal{C}_\Lambda$ d'après le lemme \ref{formal-defos-lemma-fiber-product-CLambda}.
Nous utiliserons ce résultat sans autre mention.

\medskip\noindent
On note $k[\epsilon]$ l'algèbre des nombres duaux sur $k$. Plus généralement,
pour un $k$-espace vectoriel $V$, on note $k[V]$ la $k$-algèbre dont l'espace
vectoriel sous-jacent est $k \oplus V$ et dont la multiplication est donnée
par $(a, v) \cdot (a', v') = (aa', av' + a'v)$. Lorsque $V = k$, $k[V]$ est
l'algèbre des nombres duaux sur $k$. Pour tout $k$-espace vectoriel $V$ de
dimension finie, l'anneau $k[V]$ appartient à $\mathcal{C}_\Lambda$.

\begin{definition}
\label{formal-defos-definition-S1-S2}
Soit $\mathcal{F}$ une catégorie cofibrée en groupoïdes au-dessus de $\mathcal
C_\Lambda$. On définit les {\it conditions (S1) et (S2)} pour
$\mathcal{F}$ comme suit :
\begin{enumerate}
\item[(S1)] Tout diagramme de $\mathcal{F}$
$$
\vcenter{
\xymatrix{
           & x_2 \ar[d] \\
x_1 \ar[r] & x
}
}
\quad\text{au-dessus de}\quad
\vcenter{
\xymatrix{
           & A_2 \ar[d] \\
A_1 \ar[r] & A
}
}
$$
dans $\mathcal{C}_\Lambda$, avec $A_2 \to A$ surjectif, peut être complété en
un diagramme commutatif
$$
\vcenter{
\xymatrix{
y \ar[r] \ar[d] & x_2 \ar[d] \\
x_1 \ar[r]      & x
}
}
\quad\text{au-dessus de}\quad
\vcenter{
\xymatrix{
A_1 \times_A A_2 \ar[r] \ar[d] & A_2 \ar[d] \\
A_1 \ar[r]      & A.
}
}
$$
\item[(S2)]
La condition (S1) est satisfaite pour les diagrammes de $\mathcal{F}$ situés
au-dessus d'un diagramme de $\mathcal{C}_\Lambda$ de la forme
$$
\xymatrix{
          & k[\epsilon] \ar[d] \\
A  \ar[r] & k.
}
$$
De plus, si l'on a deux diagrammes commutatifs dans $\mathcal{F}$
$$
\vcenter{
\xymatrix{
y \ar[r]_c \ar[d]_a & x_\epsilon \ar[d]^e \\
x \ar[r]^d          & x_0
}
}
\quad\text{et}\quad
\vcenter{
\xymatrix{
y' \ar[r]_{c'} \ar[d]_{a'} & x_\epsilon \ar[d]^e \\
x \ar[r]^d                 & x_0
}
}
\quad\text{au-dessus de}\quad
\vcenter{
\xymatrix{
A \times_k k[\epsilon] \ar[r] \ar[d] & k[\epsilon] \ar[d] \\
A  \ar[r] & k
}
}
$$
alors il existe un morphisme $b : y \to y'$ dans
$\mathcal{F}(A \times_k k[\epsilon])$ tel que $a = a' \circ b$.
\end{enumerate}
\end{definition}

\noindent
On peut expliquer en partie le sens des conditions (S1) et (S2) à l'aide des
catégories fibres. Supposons que $f_1 : A_1 \to A$ et $f_2 : A_2 \to A$
soient des morphismes d'anneaux dans $\mathcal{C}_\Lambda$, avec $f_2$
surjectif. Notons $p_i : A_1 \times_A A_2 \to A_i$ les projections.
Supposons fixé un choix d'images directes pour $\mathcal{F}$. Le diagramme
commutatif d'anneaux se traduit alors par un diagramme $2$-commutatif
$$
\xymatrix{
\mathcal{F}(A_1 \times_A A_2) \ar[r]_-{p_{2, *}} \ar[d]_{p_{1, *}} &
\mathcal{F}(A_2) \ar[d]^{f_{2, *}} \\
\mathcal{F}(A_1) \ar[r]^{f_{1, *}} & \mathcal{F}(A)
}
$$
de catégories fibres, d'où un foncteur
\begin{equation}
\label{formal-defos-equation-compare}
\mathcal{F}(A_1 \times_A A_2) \to
\mathcal{F}(A_1) \times_{\mathcal{F}(A)} \mathcal{F}(A_2)
\end{equation}
à valeurs dans le $2$-produit fibré des catégories. La condition (S1) exige
que ce foncteur soit essentiellement surjectif. La première partie de la
condition (S2) exige que ce foncteur soit essentiellement surjectif lorsque
$f_2$ est le morphisme $k[\epsilon] \to k$. En outre, dans ce cas, la seconde
partie de (S2) implique que deux objets qui deviennent isomorphes dans la
cible sont isomorphes dans la source (mais elle n'est {\it pas} équivalente à
cette assertion). L'avantage de formuler les conditions comme dans la
définition est qu'aucun choix n'est nécessaire.

\begin{lemma}
\label{formal-defos-lemma-S1-small-extensions}
Soit $\mathcal{F}$ une catégorie cofibrée en groupoïdes au-dessus de $\mathcal
C_\Lambda$. Alors $\mathcal{F}$ satisfait (S1) dès que la condition
(S1) est supposée satisfaite seulement lorsque $A_2 \to A$ est une petite
extension.
\end{lemma}

\begin{proof}
Démonstration omise. Indications : appliquer le lemme
\ref{formal-defos-lemma-factor-small-extension} et procéder par récurrence comme dans la
démonstration du lemme \ref{formal-defos-lemma-smoothness-small-extensions}.
\end{proof}

\begin{remark}
\label{formal-defos-remark-compare-S1-S2-schlessinger}
Lorsque $\mathcal{F}$ est cofibrée en ensembles, les conditions (S1) et (S2)
sont exactement les conditions (H1) et (H2) de l'article de Schlessinger
\cite{Sch}. En effet, pour un foncteur $F: \mathcal{C}_\Lambda \to
\textit{Ensembles}$, les conditions (S1) et (S2) s'énoncent ainsi :
\begin{enumerate}
\item [(S1)] Si $A_1 \to A$ et $A_2 \to A$ sont des morphismes dans
$\mathcal{C}_\Lambda$, avec $A_2 \to A$ surjectif, alors l'application
induite $F(A_1 \times_A A_2) \to F(A_1) \times_{F(A)} F(A_2)$ est surjective.
\item [(S2)] Si $A \to k$ est un morphisme dans $\mathcal{C}_\Lambda$,
alors l'application induite
$F(A \times_k k[\epsilon]) \to F(A) \times_{F(k)} F(k[\epsilon])$
est bijective.
\end{enumerate}
L'injectivité de l'application
$F(A \times_k k[\epsilon]) \to F(A) \times_{F(k)} F(k[\epsilon])$
résulte de la seconde partie de la condition (S2) et du fait que les
morphismes sont des identités.
\end{remark}

\begin{lemma}
\label{formal-defos-lemma-S2-extensions}
Soit $\mathcal{F}$ une catégorie cofibrée en groupoïdes au-dessus de
$\mathcal{C}_\Lambda$. Si $\mathcal{F}$ satisfait (S2), alors la condition
(S2) reste valable lorsque $k[\epsilon]$ est remplacé par $k[V]$, pour tout
$k$-espace vectoriel $V$ de dimension finie.
\end{lemma}

\begin{proof}
Lorsque $\mathcal{F}$ est cofibrée en ensembles, c'est-à-dire correspond à un
foncteur $F : \mathcal{C}_\Lambda \to \textit{Ensembles}$, cela résulte de la
description de (S2) pour $F$ donnée dans la remarque
\ref{formal-defos-remark-compare-S1-S2-schlessinger} et du fait que
$k[V] \cong k[\epsilon] \times_k \ldots \times_k k[\epsilon]$
avec $\dim_k V$ facteurs. C'est le cas des foncteurs que nous utiliserons dans
la suite du chapitre.

\medskip\noindent
Démontrons le cas général par récurrence sur $\dim(V)$. Si $\dim(V) = 1$,
alors $k[V] \cong k[\epsilon]$ et le résultat est vrai par hypothèse. Si
$\dim(V) > 1$, écrivons $V = V' \oplus k\epsilon$. Choisissons un diagramme
$$
\vcenter{
\xymatrix{
& x_V \ar[d] \\
x \ar[r] & x_0
}
}
\quad\text{au-dessus de}\quad
\vcenter{
\xymatrix{
& k[V] \ar[d] \\
A \ar[r] & k
}
}
$$
Choisissons un morphisme $x_V \to x_{V'}$ situé au-dessus de
$k[V] \to k[V']$ et un morphisme $x_V \to x_\epsilon$ situé au-dessus de
$k[V] \to k[\epsilon]$. Remarquons que le morphisme $x_V \to x_0$ se
factorise à la fois en $x_V \to x_{V'} \to x_0$ et en
$x_V \to x_\epsilon \to x_0$. Par l'hypothèse de récurrence, on peut trouver
un diagramme
$$
\vcenter{
\xymatrix{
y' \ar[d] \ar[r] & x_{V'} \ar[d] \\
x \ar[r] & x_0
}
}
\quad\text{au-dessus de}\quad
\vcenter{
\xymatrix{
A \times_k k[V'] \ar[d] \ar[r] & k[V'] \ar[d] \\
A \ar[r] & k
}
}
$$
Cela nous donne un diagramme commutatif
$$
\vcenter{
\xymatrix{
& x_\epsilon \ar[d] \\
y' \ar[r] & x_0
}
}
\quad\text{au-dessus de}\quad
\vcenter{
\xymatrix{
& k[\epsilon] \ar[d] \\
A \times_k k[V'] \ar[r] & k
}
}
$$
La condition (S2) fournit donc un diagramme commutatif
$$
\vcenter{
\xymatrix{
y \ar[d] \ar[r] & x_\epsilon \ar[d] \\
y' \ar[r] & x_0
}
}
\quad\text{au-dessus de}\quad
\vcenter{
\xymatrix{
(A \times_k k[V']) \times_k k[\epsilon] \ar[d] \ar[r] & k[\epsilon] \ar[d] \\
A \times_k k[V'] \ar[r] & k
}
}
$$
Remarquons que
$(A \times_k k[V']) \times_k k[\epsilon] = A \times_k k[V' \oplus k\epsilon]
= A \times_k k[V]$. Nous affirmons que $y$ s'insère dans le diagramme
commutatif voulu. Pour le voir, soit $y \to y_V$ un morphisme situé au-dessus
de $A \times_k k[V] \to k[V]$. Les morphismes
$y \to y' \to x_{V'}$ et $y \to x_\epsilon$ se factorisent par le morphisme
$y \to y_V$ (d'après les axiomes des catégories cofibrées en groupoïdes).
Ainsi, $y_V$ et $x_V$ s'insèrent tous deux dans des diagrammes commutatifs
$$
\vcenter{
\xymatrix{
y_V \ar[r] \ar[d] & x_\epsilon \ar[d] \\
x_{V'} \ar[r]          & x_0
}
}
\quad\text{et}\quad
\vcenter{
\xymatrix{
x_V \ar[r] \ar[d] & x_\epsilon \ar[d] \\
x_{V'} \ar[r]                 & x_0
}
}
$$
La seconde partie de (S2) donne donc un isomorphisme $y_V \to x_V$
compatible avec $y_V \to x_{V'}$ et $x_V \to x_{V'}$, et en particulier
avec les morphismes vers $x_0$. La composée $y \to y_V \to x_V$ s'insère
alors dans le diagramme commutatif requis
$$
\vcenter{
\xymatrix{
y \ar[r] \ar[d] & x_V \ar[d] \\
x \ar[r] & x_0
}
}
\quad\text{au-dessus de}\quad
\vcenter{
\xymatrix{
A \times_k k[V] \ar[d] \ar[r] & k[V] \ar[d] \\
A \ar[r] & k
}
}
$$
On voit ainsi que la première partie de $(S2)$ reste vraie lorsque
$k[\epsilon]$ est remplacé par $k[V]$.

\medskip\noindent
Pour démontrer la seconde partie, supposons donnés deux diagrammes
commutatifs
$$
\vcenter{
\xymatrix{
y \ar[r] \ar[d] & x_V \ar[d] \\
x \ar[r] & x_0
}
}
\quad\text{et}\quad
\vcenter{
\xymatrix{
y' \ar[r] \ar[d] & x_V \ar[d] \\
x \ar[r] & x_0
}
}
\quad\text{au-dessus de}\quad
\vcenter{
\xymatrix{
A \times_k k[V] \ar[d] \ar[r] & k[V] \ar[d] \\
A \ar[r] & k
}
}
$$
Nous utiliserons les morphismes $x_V \to x_{V'} \to x_0$ et
$x_V \to x_\epsilon \to x_0$ introduits au premier paragraphe de la
démonstration. Choisissons des morphismes $y \to y_{V'}$ et
$y' \to y'_{V'}$ situés au-dessus de
$A \times_k k[V] \to A \times_k k[V']$. Les axiomes d'une catégorie cofibrée
impliquent que l'on peut trouver des diagrammes commutatifs
$$
\vcenter{
\xymatrix{
y_{V'} \ar[r] \ar[d] & x_{V'} \ar[d] \\
x \ar[r] & x_0
}
}
\quad\text{et}\quad
\vcenter{
\xymatrix{
y'_{V'} \ar[r] \ar[d] & x_{V'} \ar[d] \\
x \ar[r] & x_0
}
}
\quad\text{au-dessus de}\quad
\vcenter{
\xymatrix{
A \times_k k[V'] \ar[d] \ar[r] & k[V'] \ar[d] \\
A \ar[r] & k
}
}
$$
Par l'hypothèse de récurrence, on obtient un isomorphisme
$b : y_{V'} \to y'_{V'}$ compatible avec les morphismes $y_{V'} \to x$ et
$y'_{V'} \to x$, donc en particulier avec les morphismes vers $x_0$. On a
alors des diagrammes commutatifs
$$
\vcenter{
\xymatrix{
y \ar[r] \ar[d] & x_\epsilon \ar[d] \\
y'_{V'} \ar[r] & x_0
}
}
\quad\text{et}\quad
\vcenter{
\xymatrix{
y' \ar[r] \ar[d] & x_\epsilon \ar[d] \\
y'_{V'} \ar[r] & x_0
}
}
\quad\text{au-dessus de}\quad
\vcenter{
\xymatrix{
A \times_k k[\epsilon] \ar[d] \ar[r] & k[\epsilon] \ar[d] \\
A \ar[r] & k
}
}
$$
où le morphisme $y \to y'_{V'}$ est la composée
$y \to y_{V'} \xrightarrow{b} y'_{V'}$, et où les morphismes
$y \to x_\epsilon$ et $y' \to x_\epsilon$ sont les composées des morphismes
$y \to x_V$ et $y' \to x_V$ avec le morphisme $x_V \to x_\epsilon$. La
seconde partie de (S2) assure alors l'existence d'un isomorphisme $y \to y'$
compatible avec les morphismes vers $y'_{V'}$, donc en particulier avec
les morphismes vers $x$ (puisque $b$ était compatible avec les morphismes
vers $x$).
\end{proof}

\begin{lemma}
\label{formal-defos-lemma-S1-S2-associated-functor}
Soit $\mathcal{F}$ une catégorie cofibrée en groupoïdes au-dessus de
$\mathcal{C}_\Lambda$.
\begin{enumerate}
\item Si $\mathcal{F}$ satisfait (S1), alors $\overline{\mathcal{F}}$ la
satisfait aussi.
\item Si $\mathcal{F}$ satisfait (S2), alors $\overline{\mathcal{F}}$ la
satisfait aussi, pourvu que l'une au moins des conditions suivantes soit
vérifiée :
\begin{enumerate}
\item $\mathcal{F}$ est une catégorie de prédéformation ;
\item la catégorie $\mathcal{F}(k)$ est un ensemble ou un ensoïde ;
\item pour tout morphisme $x_\epsilon \to x_0$ de $\mathcal{F}$ situé
au-dessus de $k[\epsilon] \to k$, l'application d'image directe
$\text{Aut}_{k[\epsilon]}(x_\epsilon) \to \text{Aut}_k(x_0)$
est surjective.
\end{enumerate}
\end{enumerate}
\end{lemma}

\begin{proof}
Supposons que $\mathcal{F}$ satisfasse (S1). Soient des morphismes d'anneaux
$f_i : A_i \to A$ dans $\mathcal{C}_\Lambda$, avec $f_2$ surjectif, et des
$x_i \in \mathcal{F}(A_i)$ tels que les images directes
$f_{1, *}(x_1)$ et $f_{2, *}(x_2)$ soient isomorphes. On peut alors choisir un
objet $x$ de $\mathcal{F}$ au-dessus de $A$ isomorphe à chacune de ces images,
et l'on obtient un diagramme comme dans (S1). On trouve donc un objet $y$ de
$\mathcal{F}$ au-dessus de $A_1 \times_A A_2$ dont l'image directe vers
$A_1$, respectivement $A_2$, est isomorphe à $x_1$, respectivement $x_2$.
Ainsi, (S1) est vérifiée pour $\overline{\mathcal{F}}$.

\medskip\noindent
Supposons que $\mathcal{F}$ satisfasse (S2). La première partie de (S2) pour
$\overline{\mathcal{F}}$ résulte de l'argument précédent. La seconde partie
de (S2) pour $\overline{\mathcal{F}}$ signifie que l'application
$$
\overline{\mathcal{F}}(A \times_k k[\epsilon]) \to
\overline{\mathcal{F}}(A)
\times_{\overline{\mathcal{F}}(k)} \overline{\mathcal{F}}(k[\epsilon])
$$
est injective pour tout anneau $A$ de $\mathcal{C}_\Lambda$. Supposons que
$y, y' \in \mathcal{F}(A \times_k k[\epsilon])$. En utilisant les axiomes
des catégories cofibrées, on peut choisir des diagrammes commutatifs
$$
\vcenter{
\xymatrix{
y \ar[r]_c \ar[d]_a & x_\epsilon \ar[d]^e \\
x \ar[r]^d          & x_0
}
}
\quad\text{et}\quad
\vcenter{
\xymatrix{
y' \ar[r]_{c'} \ar[d]_{a'} & x'_\epsilon \ar[d]^{e'} \\
x' \ar[r]^{d'}                 & x'_0
}
}
\quad\text{au-dessus de}\quad
\vcenter{
\xymatrix{
A \times_k k[\epsilon] \ar[d] \ar[r] & k[\epsilon] \ar[d] \\
A \ar[r] & k
}
}
$$
Supposons qu'il existe des isomorphismes
$\alpha : x \to x'$ dans $\mathcal{F}(A)$ et
$\beta : x_\epsilon \to x'_\epsilon$ dans $\mathcal{F}(k[\epsilon])$. Cela
signifie également qu'il existe un isomorphisme $\gamma : x_0 \to x'_0$
compatible avec $\alpha$. Pour démontrer (S2) pour
$\overline{\mathcal{F}}$, il faut montrer qu'il existe un isomorphisme
$y \to y'$ dans $\mathcal{F}(A \times_k k[\epsilon])$. D'après (S2) pour
$\mathcal{F}$, un tel morphisme existera si l'on peut choisir les
isomorphismes $\alpha$, $\beta$ et $\gamma$ de telle sorte que
$$
\xymatrix{
x \ar[d]^\alpha \ar[r] & x_0 \ar[d]^\gamma &
x_\epsilon \ar[d]^\beta \ar[l]^e \\
x' \ar[r] & x'_0 & x'_\epsilon \ar[l]_{e'}
}
$$
soit commutatif (on peut alors remplacer $x$ par $x'$ et $x_\epsilon$ par
$x'_\epsilon$ dans le diagramme affiché précédent). Le carré de gauche est
commutatif par notre choix de $\gamma$. On peut factoriser $e' \circ \beta$
en $\gamma' \circ e$ pour un second morphisme
$\gamma' : x_0 \to x'_0$. Il reste à savoir si l'on peut faire en sorte que
$\gamma = \gamma'$. C'est clair si $\mathcal{F}(k)$ est un ensemble ou un
ensoïde. De plus, si
$\text{Aut}_{k[\epsilon]}(x_\epsilon) \to \text{Aut}_k(x_0)$
est surjective, on peut modifier le choix de $\beta$ en le précomposant avec
un automorphisme de $x_\epsilon$ dont l'image est
$\gamma^{-1} \circ \gamma'$, ce qui donne le résultat.
\end{proof}

\begin{lemma}
\label{formal-defos-lemma-S1-S2-localize}
Soit $\mathcal{F}$ une catégorie cofibrée en groupoïdes au-dessus de
$\mathcal{C}_\Lambda$. Soit $x_0 \in \Ob(\mathcal{F}(k))$. Soit
$\mathcal{F}_{x_0}$ la catégorie cofibrée en groupoïdes au-dessus de
$\mathcal{C}_\Lambda$ construite dans la remarque
\ref{formal-defos-remark-localize-cofibered-groupoid}.
\begin{enumerate}
\item Si $\mathcal{F}$ satisfait (S1), alors $\mathcal{F}_{x_0}$ la satisfait
aussi.
\item Si $\mathcal{F}$ satisfait (S2), alors $\mathcal{F}_{x_0}$ la satisfait
aussi.
\end{enumerate}
\end{lemma}

\begin{proof}
Tout diagramme de $\mathcal{F}_{x_0}$ comme dans la définition
\ref{formal-defos-definition-S1-S2} donne un diagramme dans $\mathcal{F}$, et le résultat
fourni par la condition (S1) ou (S2) pour ce diagramme dans $\mathcal{F}$ peut
également être considéré comme un résultat dans $\mathcal{F}_{x_0}$.
\end{proof}

\begin{lemma}
\label{formal-defos-lemma-lifting-section}
Soit $p: \mathcal{F} \to \mathcal{C}_\Lambda$ une catégorie cofibrée en
groupoïdes. Considérons un diagramme de $\mathcal{F}$
$$
\vcenter{
\xymatrix{
y \ar[r] \ar[d]_a & x_\epsilon \ar[d]_e \\
x \ar[r]^d        & x_0
}
}
\quad\text{au-dessus de}\quad
\vcenter{
\xymatrix{
A \times_k k[\epsilon] \ar[r] \ar[d] & k[\epsilon] \ar[d] \\
A \ar[r] & k.
}
}
$$
dans $\mathcal{C}_\Lambda$. Supposons que $\mathcal{F}$ satisfasse (S2).
Alors il existe un morphisme $s : x \to y$ tel que
$a \circ s = \text{id}_x$ si et seulement s'il existe un morphisme
$s_\epsilon : x \to x_\epsilon$ tel que $e \circ s_\epsilon = d$.
\end{lemma}

\begin{proof}
Le sens direct est clair. Réciproquement, supposons qu'il existe un morphisme
$s_\epsilon : x \to x_\epsilon$ tel que $e \circ s_\epsilon = d$.
Remarquons que $p(s_\epsilon) : A \to k[\epsilon]$ est un morphisme d'anneaux
compatible avec le morphisme $A \to k$. On obtient donc
$$
\sigma = (\text{id}_A, p(s_\epsilon)) : A \to A \times_k k[\epsilon].
$$
Choisissons une image directe $x \to \sigma_*x$. Par construction, on peut
factoriser $s_\epsilon$ en $x \to \sigma_*x \to x_\epsilon$. De plus, comme
$\sigma$ est une section de $A \times_k k[\epsilon] \to A$, on obtient un
morphisme $\sigma_*x \to x$ tel que $x \to \sigma_*x \to x$ soit
$\text{id}_x$. L'égalité $e \circ s_\epsilon = d$ montre que le diagramme
$$
\xymatrix{
\sigma_*x \ar[r] \ar[d] & x_\epsilon \ar[d]_e \\
x \ar[r]^d        & x_0
}
$$
est commutatif. La condition (S2) fournit donc un morphisme
$\sigma_*x \to y$ tel que $\sigma_*x \to y \to x$ soit le morphisme donné
$\sigma_*x \to x$. Il suffit maintenant de prendre pour solution la composée
$a : x \to y$ égale à $x \to \sigma_*x \to y$.
\end{proof}

\begin{lemma}
\label{formal-defos-lemma-lifting-along-small-extension}
Considérons un diagramme commutatif dans une catégorie de prédéformation
$\mathcal{F}$
$$
\vcenter{
\xymatrix{
y \ar[r] \ar[d] & x_2 \ar[d]^{a_2} \\
x_1 \ar[r]^{a_1}        & x
}
}
\quad\text{au-dessus de}
\vcenter{
\xymatrix{
A_1 \times_A A_2 \ar[r] \ar[d] & A_2 \ar[d]^{f_2} \\
A_1 \ar[r]^{f_1} & A
}
}
$$
dans $\mathcal{C}_\Lambda$, où $f_2 : A_2 \to A$ est une petite extension.
Supposons qu'il existe un morphisme $h : A_1 \to A_2$ tel que
$f_2 = f_1 \circ h$. Soit $I = \Ker(f_2)$. Considérons le morphisme d'anneaux
$$
g : A_1 \times_A A_2 \longrightarrow k[I] = k \oplus I, \quad
(u, v) \longmapsto \overline{u} \oplus (v - h(u))
$$
Choisissons une image directe $y \to g_*y$. Supposons que $\mathcal{F}$
satisfasse (S2). S'il existe un morphisme $x_1 \to g_*y$, alors il existe un
morphisme $b: x_1 \to x_2$ tel que $a_1 =  a_2 \circ b$.
\end{lemma}

\begin{proof}
Remarquons que
$\text{id}_{A_1} \times g : A_1 \times_A A_2 \to A_1 \times_k k[I]$
est un isomorphisme et que $k[I] \cong k[\epsilon]$. On a donc un diagramme
$$
\vcenter{
\xymatrix{
y \ar[r] \ar[d] & g_*y \ar[d] \\
x_1 \ar[r]        & x_0
}
}
\quad\text{au-dessus de}\quad
\vcenter{
\xymatrix{
A_1 \times_k k[\epsilon] \ar[r] \ar[d] & k[\epsilon] \ar[d] \\
A_1 \ar[r] & k.
}
}
$$
où $x_0$ est un objet de $\mathcal{F}$ situé au-dessus de $k$ (tout objet de
$\mathcal{F}$ possède un unique morphisme vers $x_0$ ; voir la discussion qui
suit la définition \ref{formal-defos-definition-predeformation-category}). Si l'on dispose
d'un morphisme $x_1 \to g_*y$, le lemme \ref{formal-defos-lemma-lifting-section} fournit
une section $s : x_1 \to y$ du morphisme $y \to x_1$. En la composant
avec le morphisme $y \to x_2$, on obtient $b : x_1 \to x_2$, qui vérifie
$a_1 =  a_2 \circ b$ puisque le diagramme du lemme est commutatif et que $s$
est une section.
\end{proof}





\section{Espaces tangents des foncteurs}
\label{formal-defos-section-tangent-spaces-functors}

\noindent
Soit $R$ un anneau. On note $\text{Mod}_R$ la catégorie des $R$-modules et
$\text{Mod}^{fg}_R$ la catégorie des $R$-modules de type fini.

\begin{definition}
\label{formal-defos-definition-linear}
Soit $L: \text{Mod}^{fg}_R \to \text{Mod}_R$, respectivement
$L: \text{Mod}_R \to \text{Mod}_R$, un foncteur. On dit que $L$ est
{\it $R$-linéaire} si, pour toute paire d'objets $M, N$ de
$\text{Mod}^{fg}_R$, respectivement de $\text{Mod}_R$, l'application
$$
L : \Hom_R(M, N) \longrightarrow \Hom_R(L(M), L(N))
$$
est un morphisme de $R$-modules.
\end{definition}

\begin{remark}
\label{formal-defos-remark-linear-enriched-over-modules}
On peut définir la notion de $R$-linéarité pour tout foncteur entre catégories
enrichies sur $\text{Mod}_R$. Nous avons formulé la définition spécifiquement
pour les foncteurs $L: \text{Mod}^{fg}_R \to \text{Mod}_R$ et
$L: \text{Mod}_R \to \text{Mod}_R$, car ce sont les cas dont nous avons eu
besoin jusqu'ici.
\end{remark}

\begin{remark}
\label{formal-defos-remark-linear-functor}
Si $L: \text{Mod}^{fg}_R \to \text{Mod}_R$ est un foncteur $R$-linéaire,
alors $L$ préserve les produits finis et envoie le module nul sur le module
nul ; voir Homologie, lemme \ref{homology-lemma-additive-additive}.
Réciproquement, si un foncteur $\text{Mod}^{fg}_R \to \textit{Ensembles}$ préserve
les produits finis et envoie le module nul sur un ensemble à un seul élément, il admet un
unique relèvement en un foncteur $R$-linéaire ; voir le lemme
\ref{formal-defos-lemma-linear-functor}.
\end{remark}

\begin{lemma}
\label{formal-defos-lemma-linear-functor}
Soit $L: \text{Mod}^{fg}_R \to \textit{Ensembles}$, respectivement
$L: \text{Mod}_R \to \textit{Ensembles}$, un foncteur. Supposons que $L(0)$ soit
un ensemble à un seul élément et que $L$ préserve les produits finis. Il existe alors un
unique foncteur $R$-linéaire
$\widetilde{L} : \text{Mod}^{fg}_R \to \text{Mod}_R$,
respectivement $\widetilde{L} : \text{Mod}_R \to \text{Mod}_R$, tel que
$$
\vcenter{
\xymatrix{
& \text{Mod}_R \ar[dr]^{\text{forget}} &   \\
\text{Mod}^{fg}_R  \ar[ur]^{\widetilde{L}} \ar[rr]^{L} &  &
\textit{Ensembles}
}
}
\quad\text{resp.}\quad
\vcenter{
\xymatrix{
& \text{Mod}_R \ar[dr]^{\text{forget}} &   \\
\text{Mod}_R  \ar[ur]^{\widetilde{L}} \ar[rr]^{L} &  &
\textit{Ensembles}
}
}
$$
soit commutatif.
\end{lemma}

\begin{proof}
Nous ne le démontrons que dans le cas
$L: \text{Mod}^{fg}_R \to \textit{Ensembles}$. Soit $M$ un $R$-module de type
fini. On définit $\widetilde{L}(M)$ comme l'ensemble $L(M)$ muni de la
structure de $R$-module suivante.

\medskip\noindent
Multiplication : si $r \in R$, la multiplication par $r$ sur $L(M)$ est, par
définition, l'application $L(M) \to L(M)$ induite par l'homothétie
$r \cdot : M \to M$.

\medskip\noindent
Addition : l'application somme $M \times M \to M: (m_1, m_2) \mapsto m_1 + m_2$
induit une application
$L(M \times M) \to L(M)$. Par hypothèse, $L(M \times M)$ est canoniquement
isomorphe à $L(M) \times L(M)$. L'addition sur $L(M)$ est définie par
l'application $L(M) \times L(M) \cong L(M \times M) \to L(M)$.

\medskip\noindent
Zéro : il existe un unique morphisme $0 \to M$. L'élément nul de $L(M)$
est l'image de $L(0) \to L(M)$.

\medskip\noindent
Nous omettons de vérifier que cela définit un $R$-module
$\widetilde{L}(M)$, et que cette structure est l'unique pour laquelle la
construction soit $R$-linéaire et fonctorielle en $M$.
\end{proof}

\begin{lemma}
\label{formal-defos-lemma-morphism-linear-functors}
Soient $L_1, L_2: \text{Mod}^{fg}_R \to \textit{Ensembles}$ des foncteurs qui
envoient $0$ sur un ensemble à un seul élément et préservent les produits finis. Soit
$t : L_1 \to L_2$ un morphisme de foncteurs. Alors $t$ induit un morphisme
$\widetilde{t} : \widetilde{L}_1 \to \widetilde{L}_2$ entre les foncteurs
fournis par le lemme \ref{formal-defos-lemma-linear-functor}, donné simplement par
$\widetilde{t}_M = t_M: \widetilde{L}_1(M) \to \widetilde{L}_2(M)$ pour tout
$M \in \Ob(\text{Mod}^{fg}_R)$. Autrement dit,
$t_M: \widetilde{L}_1(M) \to \widetilde{L}_2(M)$ est un morphisme de
$R$-modules.
\end{lemma}

\begin{proof}
Omis.
\end{proof}

\noindent
Dans le cas où $R = K$ est un corps, un foncteur $K$-linéaire
$L : \text{Mod}^{fg}_K \to \text{Mod}_K$ est déterminé par sa valeur $L(K)$.

\begin{lemma}
\label{formal-defos-lemma-linear-functor-over-field}
Soit $K$ un corps. Soit $L: \text{Mod}^{fg}_K \to
\text{Mod}_K$ un foncteur
$K$-linéaire. Alors $L$ est isomorphe au foncteur
$L(K) \otimes_K - : \text{Mod}^{fg}_K \to
\text{Mod}_K$.
\end{lemma}

\begin{proof}
Pour $V \in \Ob(\text{Mod}^{fg}_K)$, l'isomorphisme
$L(K) \otimes_K V \to L(V)$ est donné sur les tenseurs purs par
$x \otimes v \mapsto L(f_v)(x)$, où $f_v: K \to V$ est l'application
$K$-linéaire envoyant $1 \mapsto v$. Lorsque $V = K$, c'est l'isomorphisme
$L(K) \otimes_K K \to L(K)$ donné par la multiplication, c'est-à-dire par
l'action de $K$. Pour $V$
quelconque, c'est un isomorphisme en vertu du cas $V = K$ et du fait que $L$
commute aux produits finis (remarque \ref{formal-defos-remark-linear-functor}).
\end{proof}

\noindent
Pour un anneau $R$ et un $R$-module $M$, soit $R[M]$ la $R$-algèbre dont le
$R$-module sous-jacent est $R \oplus M$ et dont la multiplication est donnée
par $(r, m) \cdot (r', m') = (rr', rm' + r'm)$. Lorsque $M = R$, c'est
l'algèbre des nombres duaux sur $R$, que l'on note $R[\epsilon]$.

\medskip\noindent
Soit maintenant $S$ un anneau, et supposons que $R$ soit une $S$-algèbre.
L'affectation $M \mapsto R[M]$ détermine alors un foncteur
$\text{Mod}_R \to S\text{-Alg}/R$, où $S\text{-Alg}/R$ désigne la catégorie
des $S$-algèbres au-dessus de $R$. Remarquons que $S\text{-Alg}/R$ admet les
produits finis : si $A_1 \to R$ et $A_2 \to R$ sont deux objets, alors
$A_1 \times_R A_2$ est un produit.

\begin{lemma}
\label{formal-defos-lemma-preserves-products}
Soit $R$ une $S$-algèbre. Alors le foncteur
$\text{Mod}_R \to S\text{-Alg}/R$ décrit ci-dessus préserve les produits
finis.
\end{lemma}

\begin{proof}
Il s'agit simplement de l'assertion suivante : si $M$ et $N$ sont des
$R$-modules, le morphisme $R[M \times N] \to R[M] \times_R R[N]$ est un
isomorphisme dans $S\text{-Alg}/R$.
\end{proof}

\begin{lemma}
\label{formal-defos-lemma-tangent-space-functor}
Soit $R$ une $S$-algèbre, et soit $\mathcal{C}$ une sous-catégorie strictement
pleine de $S\text{-Alg}/R$ contenant $R[M]$ pour tout
$M \in \Ob(\text{Mod}^{fg}_R)$. Soit
$F: \mathcal{C} \to \textit{Ensembles}$ un foncteur. Supposons que $F(R)$ soit un
ensemble à un seul élément et que, pour tous $M, N \in
\Ob(\text{Mod}^{fg}_R)$, l'application
induite
$$
F(R[M] \times_R R[N]) \to F(R[M]) \times F(R[N])
$$
soit une bijection. Alors $F(R[M])$ possède une structure naturelle de
$R$-module pour tout $M
\in \Ob(\text{Mod}^{fg}_R)$.
\end{lemma}

\begin{proof}
Remarquons que $R \cong R[0]$ et
$R[M] \times_R R[N] \cong R[M \times N]$ ; ainsi, $R$ et
$R[M] \times_R R[N]$ sont des objets de $\mathcal{C}$ par nos hypothèses sur
$\mathcal{C}$. Les conditions imposées à $F$ ont donc un sens. Le foncteur
$\text{Mod}_R \to S\text{-Alg}/R$ du lemme \ref{formal-defos-lemma-preserves-products} se
restreint à un foncteur $\text{Mod}^{fg}_R \to \mathcal{C}$ par l'hypothèse
sur $\mathcal{C}$. Soit $L$ la composée
$\text{Mod}^{fg}_R \to \mathcal{C} \to \textit{Ensembles}$, c'est-à-dire
$L(M) = F(R[M])$. Alors $L$ préserve les produits finis d'après le lemme
\ref{formal-defos-lemma-preserves-products} et l'hypothèse sur $F$. Le lemme
\ref{formal-defos-lemma-linear-functor} montre donc que $L(M) = F(R[M])$ possède une
structure naturelle de $R$-module pour tout
$M \in \Ob(\text{Mod}^{fg}_R)$.
\end{proof}

\begin{definition}
\label{formal-defos-definition-tangent-space-over-R}
Soit $\mathcal{C}$ une catégorie comme dans le lemme
\ref{formal-defos-lemma-tangent-space-functor}. Soit
$F : \mathcal{C} \to \textit{Ensembles}$ un foncteur tel que $F(R)$ soit un
ensemble à un seul élément. L'{\it espace tangent $TF$ de $F$} est $F(R[\epsilon])$.
\end{definition}

\noindent
Lorsque $F : \mathcal{C} \to \textit{Ensembles}$ satisfait les hypothèses du
lemme \ref{formal-defos-lemma-tangent-space-functor}, l'espace tangent $TF$ possède une
structure naturelle de $R$-module.

\begin{example}
\label{formal-defos-example-tangent-space-functor}
Puisque $\mathcal{C}_\Lambda$ contient tous les $k[V]$ pour les espaces
vectoriels $V$ de dimension finie, la définition
\ref{formal-defos-definition-tangent-space-over-R} s'applique avec $S = \Lambda$, $R = k$,
$\mathcal{C} = \mathcal{C}_\Lambda$ et
$F : \mathcal{C}_\Lambda \to \textit{Ensembles}$ un foncteur de prédéformation.
L'espace tangent est $TF = F(k[\epsilon])$.
\end{example}

\begin{example}
\label{formal-defos-example-tangent-space-prorepresentable-functor}
Calculons l'espace tangent de l'exemple
\ref{formal-defos-example-tangent-space-functor} lorsque
$F : \mathcal{C}_\Lambda \to \textit{Ensembles}$ est un foncteur
pro-représentable, disons $F = \underline{S}|_{\mathcal{C}_\Lambda}$ pour
$S \in \Ob(\widehat{\mathcal{C}}_\Lambda)$. Alors $F$ commute aux limites
arbitraires et satisfait donc les hypothèses du lemme
\ref{formal-defos-lemma-tangent-space-functor}. On calcule
$$
TF = F(k[\epsilon]) = \Mor_{\mathcal{C}_\Lambda}(S, k[\epsilon])
= \text{Der}_\Lambda(S, k)
$$
et, plus généralement, pour un $k$-espace vectoriel $V$ de dimension finie,
on a
$$
F(k[V]) = \Mor_{\mathcal{C}_\Lambda}(S, k[V]) = \text{Der}_\Lambda(S, V).
$$
Explicitement, un morphisme de $\Lambda$-algèbres $f : S \to k[V]$
compatible avec les augmentations $q : S \to k$ et $k[V] \to k$ correspond
à la dérivation $D$ définie par $s \mapsto f(s) - q(s)$. Réciproquement, une
$\Lambda$-dérivation $D : S \to V$ correspond au morphisme
$f : S \to k[V]$ de $\mathcal{C}_\Lambda$ défini par
$f(s) = q(s) + D(s)$. Ces identifications étant fonctorielles, les structures
de $k$-espace vectoriel de $TF$ et de $\text{Der}_\Lambda(S, k)$ se
correspondent (voir le lemme \ref{formal-defos-lemma-morphism-linear-functors}). Il résulte
du lemme \ref{formal-defos-lemma-derivations-finite} que $\dim_k TF$ est finie.
\end{example}

\begin{example}
\label{formal-defos-example-tangent-space-classical-prorepresentable-functor}
Le calcul de l'exemple
\ref{formal-defos-example-tangent-space-prorepresentable-functor} se simplifie dans le cas
classique. Dans ce cas, l'espace tangent du foncteur
$F = \underline{S}|_{\mathcal{C}_\Lambda}$ est simplement l'espace cotangent
relatif de $S$ sur $\Lambda$, soit $TF = T_{S/\Lambda}$. Cela reste en fait
valable, plus généralement, lorsque l'extension de corps $k/k'$ est
séparable. Voir Exercices, exercice
\ref{exercises-exercise-tangent-space-Zariski}.
\end{example}

\begin{lemma}
\label{formal-defos-lemma-morphism-tangent-spaces}
Soient $F, G: \mathcal{C} \to \textit{Ensembles}$ des foncteurs satisfaisant les
hypothèses du lemme \ref{formal-defos-lemma-tangent-space-functor}. Soit $t : F \to G$ un
morphisme de foncteurs. Pour tout $M \in \Ob(\text{Mod}^{fg}_R)$,
l'application $t_{R[M]}: F(R[M]) \to G(R[M])$ est un morphisme de
$R$-modules, où $F(R[M])$ et $G(R[M])$ sont munis des structures de
$R$-modules du lemme \ref{formal-defos-lemma-tangent-space-functor}. En particulier,
$t_{R[\epsilon]} : TF \to TG$ est un morphisme de $R$-modules.
\end{lemma}

\begin{proof}
Cela résulte du lemme \ref{formal-defos-lemma-morphism-linear-functors}.
\end{proof}

\begin{example}
\label{formal-defos-example-tangent-space-map-prorepresentable-functor}
Supposons que $f : R \to S$ soit un morphisme d'anneaux dans
$\widehat{\mathcal{C}}_\Lambda$. Posons
$F = \underline{R}|_{\mathcal{C}_\Lambda}$ et
$G = \underline{S}|_{\mathcal{C}_\Lambda}$. Le morphisme d'anneaux $f$ induit
une transformation de foncteurs $G \to F$. Le lemme
\ref{formal-defos-lemma-morphism-tangent-spaces} donne une application $k$-linéaire
$TG \to TF$. Il s'agit de l'application
$$
TG = \text{Der}_\Lambda(S, k) \longrightarrow \text{Der}_\Lambda(R, k) = TF
$$
comme il résulte des identifications canoniques
$F(k[V]) = \text{Der}_\Lambda(R, V)$ et
$G(k[V]) = \text{Der}_\Lambda(S, V)$ de l'exemple
\ref{formal-defos-example-tangent-space-prorepresentable-functor}, ainsi que de la règle de
calcul de l'application sur les espaces tangents.
\end{example}

\begin{lemma}
\label{formal-defos-lemma-tangent-space-tensor}
Soit $F: \mathcal{C} \to \textit{Ensembles}$ un foncteur satisfaisant les
hypothèses du lemme \ref{formal-defos-lemma-tangent-space-functor}. Supposons que $R = K$
soit un corps. Alors $F(K[V]) \cong TF \otimes_K V$ pour tout $K$-espace
vectoriel $V$ de dimension finie.
\end{lemma}

\begin{proof}
Cela résulte du lemme \ref{formal-defos-lemma-linear-functor-over-field}.
\end{proof}






\section{Espaces tangents des catégories de prédéformation}
\label{formal-defos-section-tangent-spaces}

\noindent
Nous définirons les espaces tangents des foncteurs de prédéformation à l'aide
de la définition générale \ref{formal-defos-definition-tangent-space-over-R}. Nous l'avons
explicitée dans l'exemple \ref{formal-defos-example-tangent-space-functor}. Elle s'applique
aux catégories de prédéformation en considérant le foncteur associé des
classes d'isomorphie.

\begin{definition}
\label{formal-defos-definition-tangent-space}
Soit $\mathcal{F}$ une catégorie de prédéformation. L'{\it espace tangent
$T \mathcal{F}$ de $\mathcal{F}$} est l'ensemble
$\overline{\mathcal{F}}(k[\epsilon])$ des classes d'isomorphie d'objets de la
catégorie fibre $\mathcal
F(k[\epsilon])$.
\end{definition}

\noindent
Ainsi, $T \mathcal{F}$ n'est autre que l'espace tangent du foncteur associé
$\overline{\mathcal{F}}: \mathcal{C}_\Lambda \to \textit{Ensembles}$. Il possède
une structure naturelle d'espace vectoriel lorsque $\mathcal{F}$ satisfait
(S2), ou plus généralement dès que $\overline{\mathcal{F}}$ la satisfait.

\begin{lemma}
\label{formal-defos-lemma-tangent-space-vector-space}
Soit $\mathcal{F}$ une catégorie de prédéformation telle que
$\overline{\mathcal{F}}$ satisfasse (S2)\footnote{Par exemple, si
$\mathcal{F}$ satisfait (S2), voir le lemme
\ref{formal-defos-lemma-S1-S2-associated-functor}.}. Alors $T \mathcal{F}$ possède une
structure naturelle de $k$-espace vectoriel. Pour tout espace vectoriel $V$
de dimension finie, on a
$\overline{\mathcal{F}}(k[V]) = T\mathcal{F} \otimes_k V$
fonctoriellement en $V$.
\end{lemma}

\begin{proof}
Posons
$F = \overline{\mathcal{F}} : \mathcal{C}_\Lambda \to \textit{Ensembles}$.
C'est un foncteur de prédéformation et $F$ satisfait (S2). D'après le lemme
\ref{formal-defos-lemma-S2-extensions} (et la traduction de la remarque
\ref{formal-defos-remark-compare-S1-S2-schlessinger}), on voit que
$$
F(A \times_k k[V]) \longrightarrow F(A) \times F(k[V])
$$
est une bijection pour tout espace vectoriel $V$ de dimension finie et tout
$A \in \Ob(\mathcal{C}_\Lambda)$. En particulier, si $A = k[W]$, on voit que
$F(k[W] \times_k k[V]) = F(k[W]) \times F(k[V])$. Autrement dit, les
hypothèses du lemme \ref{formal-defos-lemma-tangent-space-functor} sont satisfaites, et
$TF = T \mathcal{F}$ possède une structure naturelle de $k$-espace
vectoriel. La dernière assertion résulte du lemme
\ref{formal-defos-lemma-tangent-space-tensor}.
\end{proof}

\noindent
Un morphisme de catégories de prédéformation induit une application sur les
espaces tangents.

\begin{definition}
\label{formal-defos-definition-differential}
Soit $\varphi : \mathcal{F} \to \mathcal{G}$ un morphisme de catégories de
prédéformation. La {\it différentielle
$d \varphi : T \mathcal{F} \to T \mathcal{G}$ de $\varphi$} est l'application
obtenue en évaluant le morphisme de foncteurs
$\overline{\varphi}: \overline{\mathcal{F}} \to \overline{\mathcal{G}}$ en
$A = k[\epsilon]$.
\end{definition}

\begin{lemma}
\label{formal-defos-lemma-k-linear-differential}
Soit $\varphi : \mathcal{F} \to \mathcal{G}$ un morphisme de catégories de
prédéformation. Supposons que $\overline{\mathcal{F}}$ et
$\overline{\mathcal{G}}$ satisfassent toutes deux (S2). Alors
$d \varphi : T \mathcal{F} \to T \mathcal{G}$ est $k$-linéaire.
\end{lemma}

\begin{proof}
Dans la démonstration du lemme \ref{formal-defos-lemma-tangent-space-vector-space}, nous
avons vu que $\overline{\mathcal{F}}$ et $\overline{\mathcal{G}}$ satisfont
les hypothèses du lemme \ref{formal-defos-lemma-tangent-space-functor}. Le résultat découle
donc du lemme \ref{formal-defos-lemma-morphism-tangent-spaces}.
\end{proof}

\begin{remark}
\label{formal-defos-remark-tangent-space-cofibered-groupoid}
On peut étendre les notions d'espace tangent et de différentielle aux
catégories cofibrées en groupoïdes arbitraires comme suit. Soit $\mathcal{F}$
une catégorie cofibrée en groupoïdes au-dessus de $\mathcal{C}_\Lambda$, et
soit $x \in \Ob(\mathcal{F}(k))$. Comme dans la remarque
\ref{formal-defos-remark-localize-cofibered-groupoid}, on obtient une catégorie de
prédéformation $\mathcal{F}_x$. On définit
$$
T_x\mathcal{F} = T\mathcal{F}_x
$$
comme l'{\it espace tangent de $\mathcal{F}$ en $x$}. Si
$\varphi : \mathcal{F} \to \mathcal{G}$ est un morphisme de catégories
cofibrées en groupoïdes au-dessus de $\mathcal{C}_\Lambda$ et si
$x \in \Ob(\mathcal{F}(k))$, il existe un morphisme induit
$\varphi_x: \mathcal{F}_x \to \mathcal{G}_{\varphi(x)}$. On définit la
{\it différentielle
$d_x \varphi : T_x \mathcal{F} \to T_{\varphi(x)} \mathcal{G}$
de $\varphi$ en $x$} comme l'application
$d \varphi_x: T \mathcal{F}_x \to T \mathcal{G}_{\varphi(x)}$.
Si $\mathcal{F}$ et $\mathcal{G}$ satisfont toutes deux (S2), alors tous ces
espaces tangents possèdent une structure naturelle de $k$-espace vectoriel et
toutes les différentielles
$d_x \varphi : T_x \mathcal{F} \to T_{\varphi(x)} \mathcal{G}$
sont $k$-linéaires (utiliser les lemmes \ref{formal-defos-lemma-S1-S2-localize} et
\ref{formal-defos-lemma-k-linear-differential}).
\end{remark}

\noindent
Les observations suivantes sont sans intérêt dans le cas classique ou lorsque
$k/k'$ est une extension de corps séparable, car
$\text{Der}_\Lambda(k, k)$ et $\text{Der}_\Lambda(k, V)$ sont alors nuls. Il
existe une identification canonique
$$
\Mor_{\mathcal{C}_\Lambda}(k, k[\epsilon]) =
\text{Der}_\Lambda(k, k).
$$
En effet, pour $D \in \text{Der}_\Lambda(k, k)$, soit
$f_D : k \to k[\epsilon]$ le morphisme $a \mapsto a + D(a)\epsilon$. Plus
généralement, pour un espace vectoriel $V$ de dimension finie sur $k$, on a
$$
\Mor_{\mathcal{C}_\Lambda}(k, k[V]) =
\text{Der}_\Lambda(k, V)
$$
et nous utiliserons la même notation $f_D$ pour le morphisme associé à la
dérivation $D$. On a également
$$
\Mor_{\mathcal{C}_\Lambda}(k[W], k[V]) =
\Hom_k(V, W) \oplus \text{Der}_\Lambda(k, V)
$$
où $(\varphi, D)$ correspond au morphisme
$f_{\varphi, D} : a + w \mapsto a + \varphi(w) + D(a)$. Nous noterons parfois
$f_{1, D} : a + v \to a + v + D(a)$ l'automorphisme de $k[V]$ déterminé par
la dérivation $D : k \to V$. Remarquons que
$f_{1, D} \circ f_{1, D'} = f_{1, D + D'}$.

\medskip\noindent
Soit $\mathcal{F}$ une catégorie de prédéformation au-dessus de
$\mathcal{C}_\Lambda$. Soit $x_0 \in \Ob(\mathcal{F}(k))$. D'après ce qui
précède, il existe une application canonique
$$
\gamma_V :
\text{Der}_\Lambda(k, V)
\longrightarrow
\overline{\mathcal{F}}(k[V])
$$
définie par $D \mapsto f_{D, *}(x_0)$. Il existe en outre une action
$$
a_V : \text{Der}_\Lambda(k, V) \times \overline{\mathcal{F}}(k[V])
\longrightarrow
\overline{\mathcal{F}}(k[V])
$$
définie par $(D, x) \mapsto f_{1, D, *}(x)$. Ces deux applications sont
compatibles, c'est-à-dire que
$f_{1, D, *}f_{D', *}x_0 = f_{D + D', *}x_0$, comme le montre le calcul de
leurs composées. Remarquons que les applications $\gamma_V$ et $a_V$ ne
dépendent pas du choix de $x_0$, puisque $x_0$ est unique à isomorphisme près.

\begin{lemma}
\label{formal-defos-lemma-action-linear}
Soit $\mathcal{F}$ une catégorie de prédéformation au-dessus de
$\mathcal{C}_\Lambda$. Si $\overline{\mathcal{F}}$ satisfait (S2), alors les
applications $\gamma_V$ sont $k$-linéaires et l'on a
$a_V(D, x) = x + \gamma_V(D)$.
\end{lemma}

\begin{proof}
Dans la démonstration du lemme \ref{formal-defos-lemma-tangent-space-vector-space}, nous
avons vu que le foncteur $V \mapsto \overline{\mathcal{F}}(k[V])$ envoie $0$
sur un ensemble à un seul élément et les produits sur des produits. Il en va de même du
foncteur $V \mapsto \text{Der}_\Lambda(k, V)$. Le lemme
\ref{formal-defos-lemma-morphism-linear-functors} montre donc que $\gamma_V$ est linéaire.
Soit $D : k \to V$ une $\Lambda$-dérivation. Posons
$D_1 : k \to V^{\oplus 2}$ égal à $a \mapsto (D(a), 0)$. Alors
$$
\xymatrix{
k[V \times V] \ar[r]_{+} \ar[d]^{f_{1, D_1}} & k[V] \ar[d]^{f_{1, D}} \\
k[V \times V] \ar[r]^{+} & k[V]
}
$$
est commutatif. En développant les définitions et en utilisant
$\overline{F}(V \times V) = \overline{F}(V) \times \overline{F}(V)$, cela
signifie que $a_D(x_1) + x_2 = a_D(x_1 + x_2)$ pour tous
$x_1, x_2 \in \overline{F}(V)$. Il suffit donc de montrer que
$a_V(D, 0) = 0 + \gamma_V(D)$, où $0 \in \overline{F}(V)$ est le vecteur
nul. Par définition, celui-ci est l'élément $f_{0, *}(x_0)$. Comme
$f_D = f_{1, D} \circ f_0$, le résultat voulu s'ensuit.
\end{proof}

\noindent
Deux cas particuliers des constructions précédentes sont l'application
\begin{equation}
\label{formal-defos-equation-map}
\gamma : \text{Der}_\Lambda(k, k) \longrightarrow T\mathcal{F}
\end{equation}
ainsi que l'action
\begin{equation}
\label{formal-defos-equation-action}
a : \text{Der}_\Lambda(k, k) \times T\mathcal{F} \longrightarrow T\mathcal{F}
\end{equation}
définies pour toute catégorie de prédéformation $\mathcal{F}$. Remarquons que,
si $\varphi : \mathcal{F} \to \mathcal{G}$ est un morphisme de catégories de
prédéformation, on obtient les diagrammes commutatifs
$$
\vcenter{
\xymatrix{
\text{Der}_\Lambda(k, k) \ar[r]_-\gamma \ar[rd]_\gamma &
T\mathcal{F} \ar[d]_{d\varphi} \\
& T\mathcal{G}
}
}
\quad\text{et}\quad
\vcenter{
\xymatrix{
\text{Der}_\Lambda(k, k) \times T\mathcal{F} \ar[r]_-a
\ar[d]_{1 \times d\varphi} &
T\mathcal{F} \ar[d]^{d\varphi} \\
\text{Der}_\Lambda(k, k) \times T\mathcal{G} \ar[r]^-a &
T\mathcal{G}
}
}
$$










\section{Objets formels versels}
\label{formal-defos-section-versal-objects}

\noindent
L'existence d'un objet formel versel impose à $\mathcal{F}$ de posséder la
propriété (S1).

\begin{lemma}
\label{formal-defos-lemma-versal-object-S1}
Soit $\mathcal{F}$ une catégorie de prédéformation. Supposons que
$\mathcal{F}$ possède un objet formel versel. Alors $\mathcal{F}$ satisfait
(S1).
\end{lemma}

\begin{proof}
Soit $\xi$ un objet formel versel de $\mathcal{F}$. Soit
$$
\xymatrix{
           & x_2 \ar[d] \\
x_1 \ar[r] & x
}
$$
un diagramme dans $\mathcal{F}$ tel que $x_2 \to x$ soit situé au-dessus d'un
morphisme surjectif d'anneaux. Comme le morphisme naturel
$\widehat{\mathcal{F}}|_{\mathcal{C}_\Lambda} \xrightarrow{\sim} \mathcal{F}$
est une équivalence (voir la remarque \ref{formal-defos-remark-restrict-completion}), on
peut aussi considérer ce diagramme comme un diagramme dans
$\widehat{\mathcal{F}}$. D'après le lemme
\ref{formal-defos-lemma-versal-object-quasi-initial}, il existe un morphisme
$\xi \to x_1$ ; la remarque \ref{formal-defos-remark-versal-object} fournit donc également
un morphisme $\xi \to x_2$ qui rend le diagramme
$$
\xymatrix{
\xi \ar[r] \ar[d]          & x_2 \ar[d] \\
x_1 \ar[r] & x
}
$$
commutatif. Si $x_1 \to x$ et $x_2 \to x$ sont situés au-dessus des
morphismes d'anneaux $A_1 \to A$ et $A_2 \to A$, l'image directe de $\xi$
vers $A_1 \times_A A_2$ fournit un objet $y$ comme l'exige (S1).
\end{proof}

\noindent
Lorsque notre catégorie cofibrée satisfait (S1) et (S2), on peut caractériser
les objets formels versels de la manière suivante.

\begin{lemma}
\label{formal-defos-lemma-versal-criterion}
Soit $\mathcal{F}$ une catégorie de prédéformation satisfaisant (S1) et (S2).
Soit $\xi$ un objet formel de $\mathcal{F}$ correspondant à
$\underline{\xi} : \underline{R}|_{\mathcal{C}_\Lambda} \to \mathcal{F}$ ;
voir la remarque \ref{formal-defos-remark-formal-objects-yoneda}. Alors $\xi$ est versel
si et seulement si les deux conditions suivantes sont satisfaites :
\begin{enumerate}
\item l'application
$d\underline{\xi} : T\underline{R}|_{\mathcal{C}_\Lambda} \to T\mathcal{F}$
sur les espaces tangents est surjective ;
\item étant donné un diagramme dans $\widehat{\mathcal{F}}$
$$
\vcenter{
\xymatrix{
            &  y \ar[d] \\
\xi \ar[r]  &  x
}
}
\quad\text{au-dessus de}\quad
\vcenter{
\xymatrix{
         &   B  \ar[d]^{f} \\
R \ar[r] &   A
}
}
$$
au-dessus d'un diagramme de $\widehat{\mathcal{C}}_\Lambda$ où $B \to A$ est
une petite extension d'anneaux artiniens, il existe un morphisme d'anneaux
$R \to B$ tel que
$$
\xymatrix{
         &   B  \ar[d]^{f} \\
R \ar[ur] \ar[r] &   A
}
$$
soit commutatif.
\end{enumerate}
\end{lemma}

\begin{proof}
Si $\xi$ est versel, (1) découle du lemme
\ref{formal-defos-lemma-smooth-morphism-essentially-surjective} et (2) de la remarque
\ref{formal-defos-remark-versal-object}. Supposons (1) et (2) satisfaites. D'après la
remarque \ref{formal-defos-remark-versal-object}, il faut montrer que, pour un diagramme de
$\widehat{\mathcal{F}}$ comme dans (2), il existe $\xi \to y$ tel que
$$
\xymatrix{
            &  y \ar[d] \\
\xi \ar[ur] \ar[r]  &  x
}
$$
soit commutatif. Soit $b : R \to B$ le morphisme fourni par (2). Posons
$y' = b_*\xi$ et choisissons une factorisation $\xi \to y' \to x$, située
au-dessus de $R \to B \to A$, du morphisme donné $\xi \to x$. La condition
(S1) fournit un diagramme commutatif
$$
\vcenter{
\xymatrix{
z  \ar[r] \ar[d]          &  y \ar[d] \\
y' \ar[r]  &  x
}
}
\quad\text{au-dessus de}\quad
\vcenter{
\xymatrix{
B \times_A B \ar[d] \ar[r] &   B  \ar[d]^{f} \\
B \ar[r]^{f} &   A .
}
}
$$
Posons $I = \Ker(f)$. Soit $\overline{g} : B \times_A B \to k[I]$ le
morphisme d'anneaux $(u, v) \mapsto \overline{u} \oplus (v - u)$ ; cf.\ le
lemme \ref{formal-defos-lemma-lifting-along-small-extension}. Par (1), il existe un
morphisme $\xi \to \overline{g}_*z$ situé au-dessus d'un morphisme d'anneaux
$i : R \to k[\epsilon]$. Choisissons un quotient artinien
$b_1 : R \to B_1$ tel que $b : R \to B$ et $i : R \to k[\epsilon]$ se
factorisent tous deux par $R \to B_1$, donnant ainsi
$h : B_1 \to B$ et $i' : B_1 \to k[\epsilon]$. Choisissons une image directe
$y_1 = b_{1, *}\xi$, une factorisation $\xi \to y_1 \to y'$ située au-dessus
de $R \to B_1 \to B$ de $\xi \to y'$, et une factorisation
$\xi \to y_1 \to \overline{g}_*z$ située au-dessus de
$R \to B_1 \to k[\epsilon]$ de $\xi \to \overline{g}_*z$. Une nouvelle
application de (S1) donne
$$
\vcenter{
\xymatrix{
z_1  \ar[r] \ar[d] & z \ar[r] \ar[d] & y \ar[d] \\
y_1 \ar[r] & y' \ar[r] &  x
}
}
\quad\text{au-dessus de}\quad
\vcenter{
\xymatrix{
B_1 \times_A B \ar[d] \ar[r] & B \times_A B \ar[r] \ar[d] & B \ar[d]^{f} \\
B_1 \ar[r]  & B \ar[r] &  A .
}
}
$$
Remarquons que le morphisme $g : B_1 \times_A B \to k[I]$ du lemme
\ref{formal-defos-lemma-lifting-along-small-extension} (définie à l'aide de $h$) est la
composée de $B_1 \times_A B \to B \times_A B$ et du morphisme
$\overline{g}$ ci-dessus. Par construction, il existe un morphisme
$y_1 \to g_*z_1 \cong \overline{g}_*z$ ! Le lemme
\ref{formal-defos-lemma-lifting-along-small-extension} s'applique donc (aux rectangles
extérieurs des diagrammes précédents) et fournit un morphisme $y_1 \to y$ ;
en le précomposant avec $\xi \to y_1$, on obtient le morphisme recherché
$\xi \to y$.
\end{proof}

\noindent
Si $\mathcal{F}$ possède la propriété (S1), alors le « plus grand quotient sur
lequel un relèvement existe » existe. Voici un énoncé précis.

\begin{lemma}
\label{formal-defos-lemma-largest-closed-where-lift}
Soit $\mathcal{F}$ une catégorie cofibrée en groupoïdes au-dessus de
$\mathcal{C}_\Lambda$ qui satisfait (S1). Soit $B \to A$ une surjection dans
$\mathcal{C}_\Lambda$, de noyau $I$ annulé par $\mathfrak m_B$. Soit
$x \in \mathcal{F}(A)$. L'ensemble des idéaux
$$
\mathcal{J} = \{ J \subset I \mid
\text{il existe un }y \to x\text{ au-dessus de }B/J \to A\}
$$
possède un plus petit élément.
\end{lemma}

\begin{proof}
Remarquons que $\mathcal{J}$ n'est pas vide, puisque
$I \in \mathcal{J}$. De plus, si $J \in \mathcal{J}$ et
$J \subset J' \subset I$, alors $J' \in \mathcal{J}$, car on peut prendre
l'image directe de l'objet $y$ en un objet $y'$ au-dessus de $B/J'$. Soient
$J$ et $K$ des éléments de l'ensemble affiché. Nous affirmons que
$J \cap K \in \mathcal{J}$, ce qui démontrera le lemme. Comme $I$ est un
$k$-espace vectoriel, on peut trouver un idéal $J \subset J' \subset I$ tel
que $J \cap K = J' \cap K$ et $J' + K = I$. D'après ce qui précède, on peut
remplacer $J$ par $J'$ et supposer $J + K = I$. Dans ce cas,
$$
A/(J \cap K) = A/J \times_{A/I} A/K.
$$
L'existence d'un élément $z \in \mathcal{F}(A/(J \cap K))$ qui s'envoie sur
$x$ résulte donc, par (S1), de l'existence supposée des éléments au-dessus de
$A/J$ et de $A/K$.
\end{proof}

\noindent
Nous améliorerons plus loin le résultat suivant.

\begin{lemma}
\label{formal-defos-lemma-versal-object-existence}
Soit $\mathcal{F}$ une catégorie cofibrée en groupoïdes au-dessus de
$\mathcal{C}_\Lambda$. Supposons les conditions suivantes satisfaites :
\begin{enumerate}
\item $\mathcal{F}$ est une catégorie de prédéformation.
\item $\mathcal{F}$ satisfait (S1).
\item $\mathcal{F}$ satisfait (S2).
\item $\dim_k T\mathcal{F}$ est finie.
\end{enumerate}
Alors $\mathcal{F}$ possède un objet formel versel.
\end{lemma}

\begin{proof}
Supposons (1), (2), (3) et (4) satisfaites. Choisissons un objet
$R \in \Ob(\widehat{\mathcal{C}}_\Lambda)$ tel que
$\underline{R}|_{\mathcal{C}_\Lambda}$ soit lisse. Voir le lemme
\ref{formal-defos-lemma-exists-smooth}. Soit $r = \dim_k T\mathcal{F}$ et posons
$S = R[[X_1, \ldots, X_r]]$.

\medskip\noindent
Nous allons construire par récurrence, pour $n \geq 2$, des paires
$(J_n, f_{n - 1} : \xi_n \to \xi_{n - 1})$, où les $J_n \subset S$ forment
une suite décroissante d'idéaux et où
$f_{n - 1} : \xi_n \to \xi_{n - 1}$ est un morphisme de $\mathcal{F}$ situé
au-dessus de la projection $S/J_n \to S/J_{n - 1}$.

\medskip\noindent
Étape 1. Soit $J_1 = \mathfrak m_S$. Soit $\xi_1$ l'unique objet de
$\mathcal{F}$ au-dessus de $k = S/J_1 = S/\mathfrak m_S$, à isomorphisme
unique près.

\medskip\noindent
Étape 2. Soit $J_2 = \mathfrak m_S^2 + \mathfrak{m}_R S$. Alors
$S/J_2 = k[V]$, où $V = kX_1 \oplus \ldots \oplus kX_r$. La condition (S2)
pour $\overline{\mathcal{F}}$ donne une bijection
$$
\overline{\mathcal{F}}(S/J_2)
\longrightarrow
T\mathcal{F} \otimes_k V,
$$
voir les lemmes \ref{formal-defos-lemma-S1-S2-associated-functor} et
\ref{formal-defos-lemma-tangent-space-vector-space}. Choisissons une base
$\theta_1, \ldots, \theta_r$ de $T\mathcal{F}$ et posons
$\xi_2 = \sum \theta_i \otimes X_i \in \Ob(\mathcal{F}(S/J_2))$. Ce choix a
pour effet que
$$
d\xi_2 :
\Mor_{\mathcal{C}_\Lambda}(S/J_2, k[\epsilon])
\longrightarrow
T\mathcal{F}
$$
est surjective. Soit $f_1 : \xi_2 \to \xi_1$ l'unique morphisme.

\medskip\noindent
Étape de récurrence. Supposons construite
$(J_n, f_{n - 1} : \xi_n \to \xi_{n - 1})$ pour un certain $n \geq 2$. Il
existe un plus petit élément $J_{n + 1}$ de l'ensemble des idéaux
$J \subset S$ satisfaisant : (a) $\mathfrak m_S J_n \subset J \subset J_n$
et (b) il existe un morphisme $\xi_{n + 1} \to \xi_n$ situé au-dessus de
$S/J \to S/J_n$ ; voir le lemme \ref{formal-defos-lemma-largest-closed-where-lift}. Soit
$f_n : \xi_{n + 1} \to \xi_n$ un morphisme quelconque de $\mathcal{F}$ situé
au-dessus de $S/J_{n + 1} \to S/J_n$.

\medskip\noindent
Posons $J = \bigcap J_n$, $\overline{S} = S/J$ et
$\overline{J}_n = J_n/J$. D'après le lemme \ref{formal-defos-lemma-m-adic-topology}, la
suite d'idéaux $(\overline{J}_n)$ induit la topologie
$\mathfrak m_{\overline{S}}$-adique sur $\overline{S}$. Comme
$(\xi_n, f_n)$ est un objet de
$\widehat{\mathcal{F}}_\mathcal{I}(\overline{S})$, où $\mathcal{I}$ est la
filtration $(\overline{J}_n)$ de $\overline{S}$, on voit que
$(\xi_n, f_n)$ induit un objet $\xi$ de
$\widehat{\mathcal{F}}(\overline{S})$ ; voir le lemme
\ref{formal-defos-lemma-formal-objects-different-filtration}.

\medskip\noindent
Montrons que $\xi$ est versel. Pour la versalité, il suffit de vérifier les
conditions (1) et (2) du lemme \ref{formal-defos-lemma-versal-criterion}. La condition (1)
résulte de notre choix de $\xi_2$ à l'étape 2. Considérons un diagramme dans
$\widehat{\mathcal{F}}$
$$
\vcenter{
\xymatrix{
            &  y \ar[d] \\
\xi \ar[r]  &  x
}
}
\quad\text{au-dessus de}\quad
\vcenter{
\xymatrix{
         &   B  \ar[d]^{f} \\
\overline{S} \ar[r] &   A
}
}
$$
au-dessus d'un diagramme de $\widehat{\mathcal{C}}_\Lambda$, où
$f: B \to A$ est une petite extension d'anneaux artiniens. Il faut montrer
qu'il existe un morphisme $\overline{S} \to B$ s'insérant dans le
diagramme de droite. Choisissons $n$ tel que $\overline{S} \to A$ se
factorise par $\overline{S} \to S/J_n$. C'est possible puisque la suite
$(\overline{J}_n)$ induit la topologie $\mathfrak m_{\overline{S}}$-adique,
comme nous l'avons vu. L'image directe de $\xi$ le long de
$\overline{S} \to S/J_n$ est $\xi_n$. On peut factoriser $\xi \to x$ en
$\xi \to \xi_n \to x$ ; on obtient donc un diagramme dans $\mathcal{F}$
$$
\vcenter{
\xymatrix{
            &  y \ar[d] \\
\xi_n \ar[r]  &  x
}
}
\quad\text{au-dessus de}\quad
\vcenter{
\xymatrix{
         &   B  \ar[d]^{f} \\
S/J_n \ar[r] &   A .
}
}
$$
Pour vérifier la condition (2) du lemme \ref{formal-defos-lemma-versal-criterion}, il
suffit de compléter le diagramme
$$
\xymatrix{
S/J_{n + 1} \ar[d] \ar@{-->}[r] & B \ar[d]^{f} \\
S/J_n   \ar[r] & A
}
$$
ou, de manière équivalente, le diagramme
$$
\xymatrix{
  &  S/J_n \times_A B \ar[d]^{p_1} \\
S/J_{n + 1} \ar@{-->}[ur] \ar[r] &  S/J_n.
}
$$
Si $p_1$ admet une section, la conclusion est acquise. Sinon, d'après le lemme
\ref{formal-defos-lemma-fiber-product-CLambda} (2), $p_1$ est une petite extension ; le
lemme \ref{formal-defos-lemma-essential-surjection} (4) montre donc que $p_1$ est une
surjection essentielle. Rappelons que $S = R[[X_1, \ldots, X_r]]$ et que
nous avons choisi $R$ de sorte que
$\underline{R}|_{\mathcal{C}_\Lambda}$ soit lisse. Il existe donc un
morphisme $h : R \to B$ relevant le morphisme
$R \to S \to S/J_n \to A$. Par la propriété universelle d'un anneau de
séries formelles, il existe un morphisme de $R$-algèbres
$h : S = R[[X_1, \ldots, X_r]] \to B$ relevant le morphisme donné
$S \to S/J_n \to A$. Celui-ci induit un morphisme
$g: S \to S/J_n \times_A B$ qui rend commutatif le carré en traits pleins du
diagramme
$$
\xymatrix{
S \ar[d] \ar[r]_-g  &  S/J_n \times_A B \ar[d]^{p_1} \\
S/J_{n + 1} \ar@{-->}[ur] \ar[r] &  S/J_n
}
$$
Alors $g$ est une surjection, puisque $p_1$ est une surjection essentielle.
Nous affirmons que l'idéal $K = \Ker(g)$ de $S$ satisfait les conditions (a)
et (b) de la construction de $J_{n + 1}$ à l'étape de récurrence précédente.
En effet, $K \subset J_n$ est clair et
$\mathfrak m_SJ_n \subset K$, puisque $p_1$ est une petite extension ; cela
démontre (a). En appliquant (S1) à
$$
\xymatrix{
            &  y \ar[d] \\
\xi_n \ar[r]  &  x,
}
$$
on obtient un relèvement de $\xi_n$ à
$S/K \cong S/J_n \times_A B$, donc (b) est satisfaite. Puisque $J_{n + 1}$
était le plus petit idéal possédant les propriétés (a) et (b), on a
$J_{n + 1} \subset K$. Le morphisme recherché
$S/J_{n+1} \to S/K \cong S/J_n \times_A B$ existe donc.
\end{proof}

\begin{remark}
\label{formal-defos-remark-compare-schlessinger-H3-pre}
Soit $F : \mathcal{C}_\Lambda \to \textit{Ensembles}$ un foncteur de
prédéformation satisfaisant (S1) et (S2). La condition
$\dim_k TF < \infty$ est précisément la condition (H3) de l'article de
Schlessinger. Rappelons que (S1) et (S2) correspondent aux conditions (H1) et
(H2) ; voir la remarque \ref{formal-defos-remark-compare-S1-S2-schlessinger}. Le lemme
\ref{formal-defos-lemma-versal-object-existence} nous dit donc que
$$
(H1) + (H2) + (H3)
\Rightarrow
\text{ there exists a versal formal object}
$$
pour les foncteurs de prédéformation. Nous établirons le lien avec les
enveloppes dans la remarque \ref{formal-defos-remark-compare-schlessinger-H3}.
\end{remark}





\section{Objets formels versels minimaux}
\label{formal-defos-section-minimal-versal}

\noindent
Nous effectuons quelques préparatifs afin de comprendre la (non-)unicité des
objets formels versels. Il s'avère que, si une catégorie de prédéformation
possède un objet formel versel, elle possède un objet formel versel minimal,
et que deux tels objets sont isomorphes. En outre, tous les objets formels
versels sont « plus ou moins » les mêmes, quitte à remplacer l'anneau de base
par une extension en séries formelles.

\medskip\noindent
Soit $\mathcal{F}$ une catégorie cofibrée en groupoïdes au-dessus de
$\mathcal{C}_\Lambda$. Pour tout objet $x$ de $\mathcal{F}$ situé au-dessus
de $A \in \Ob(\mathcal{C}_\Lambda)$, considérons la catégorie
$\mathcal{S}_x$ dont les objets sont
$$
\Ob(\mathcal{S}_x) =
\{x' \to x \mid x' \to x\text{ lies over }A' \subset A\}
$$
et dont les morphismes sont les morphismes au-dessus de $x$. Pour tout
$y \to x$ dans $\mathcal{F}$ situé au-dessus de $f : B \to A$ dans
$\mathcal{C}_\Lambda$, il existe un foncteur
$f_* : \mathcal{S}_y \to \mathcal{S}_x$ défini comme suit. Étant donné
$y' \to y$ situé au-dessus de $B' \subset B$, posons $A' = f(B')$ et soit
$y' \to x'$ au-dessus de $B' \to f(B')$ l'image directe de $y'$. Les axiomes
d'une catégorie cofibrée en groupoïdes donnent un unique morphisme
$x' \to x$ situé au-dessus de $f(B') \to A$ tel que
$$
\xymatrix{
y' \ar[d] \ar[r] & x' \ar[d] \\
y \ar[r] & x
}
$$
soit commutatif. Alors $x' \to x$ est un objet de $\mathcal{S}_x$. On dit
qu'un objet $x' \to x$ de $\mathcal{S}_x$ est {\it minimal} si tout morphisme
$(x'_1 \to x) \to (x' \to x)$ dans $\mathcal{S}_x$ est un isomorphisme,
c'est-à-dire si $x'$ et $x'_1$ sont définis sur le même sous-anneau de $A$.
Comme $A$ est de longueur finie comme $\Lambda$-module, les objets minimaux
existent toujours.

\begin{lemma}
\label{formal-defos-lemma-smallest-where-descends}
Soit $\mathcal{F}$ une catégorie cofibrée en groupoïdes au-dessus de
$\mathcal{C}_\Lambda$ qui satisfait (S1).
\begin{enumerate}
\item Pour $y \to x$ dans $\mathcal{F}$, l'image d'un objet minimal de
$\mathcal{S}_y$ est un objet minimal de $\mathcal{S}_x$.
\item Pour $y \to x$ dans $\mathcal{F}$ situé au-dessus d'une surjection
$f : B \to A$ dans $\mathcal{C}_\Lambda$, tout objet minimal de
$\mathcal{S}_x$ est l'image d'un objet minimal de $\mathcal{S}_y$.
\end{enumerate}
\end{lemma}

\begin{proof}
Démonstration de (1). Supposons $y \to x$ situé au-dessus de $f : B \to A$.
Soit $y' \to y$, situé au-dessus de $B' \subset B$, un objet minimal de
$\mathcal{S}_y$. Soit
$$
\vcenter{
\xymatrix{
y' \ar[d] \ar[r] & x' \ar[d] \\
y \ar[r] & x
}
}
\quad\text{au-dessus de}\quad
\vcenter{
\xymatrix{
B' \ar[d] \ar[r] & f(B') \ar[d] \\
B \ar[r] & A
}
}
$$
le diagramme de la construction de $f_*$ ci-dessus. Supposons que
$(x'' \to x) \to (x' \to x)$ soit un morphisme de $\mathcal{S}_x$, avec
$x'' \to x'$ situé au-dessus de $A'' \subset f(B')$. Par (S1), il existe
$y'' \to y'$ situé au-dessus de $B' \times_{f(B')} A'' \to B'$. La
minimalité de $y' \to y$ implique que
$B' \times_{f(B')} A'' \to B'$ est un isomorphisme, donc que
$A'' = f(B')$ ; ainsi, $x' \to x$ est minimal.

\medskip\noindent
Démonstration de (2). Supposons $f : B \to A$ surjective et $y \to x$ situé
au-dessus de $f$. Soit $x' \to x$ un objet minimal de $\mathcal{S}_x$ situé
au-dessus de $A' \subset A$. Par (S1), il existe $y' \to y$ situé au-dessus
de $B' = f^{-1}(A') = B \times_A A' \to B$ dont l'image dans
$\mathcal{S}_x$ est $x' \to x$. Ainsi, $f_*(y' \to y) = x' \to x$.
Choisissons un morphisme $(y'' \to y) \to (y' \to y)$ dans
$\mathcal{S}_y$ tel que $y'' \to y$ soit minimal (c'est possible d'après la
remarque sur les longueurs qui précède le lemme). Alors $f_*(y'' \to y)$ est
un objet de $\mathcal{S}_x$ qui s'envoie sur $x' \to x$ (par fonctorialité de
$f_*$), donc est isomorphe à $x' \to x$ par minimalité de $x' \to x$.
\end{proof}

\begin{lemma}
\label{formal-defos-lemma-smallest-where-descends-versal}
Soit $\mathcal{F}$ une catégorie cofibrée en groupoïdes au-dessus de
$\mathcal{C}_\Lambda$ qui satisfait (S1). Soit $\xi$ un objet formel versel de
$\mathcal{F}$ situé au-dessus de $R$. Il existe un morphisme $\xi' \to \xi$
situé au-dessus de $R' \subset R$ ayant les propriétés de minimalité
suivantes :
\begin{enumerate}
\item pour tout $f : R \to A$, avec $A \in \Ob(\mathcal{C}_\Lambda)$, les
images directes
$$
\vcenter{
\xymatrix{
\xi' \ar[d] \ar[r] & x' \ar[d] \\
\xi \ar[r] & x
}
}
\quad\text{au-dessus de}\quad
\vcenter{
\xymatrix{
R' \ar[d] \ar[r] & f(R') \ar[d] \\
R \ar[r] & A
}
}
$$
produisent un objet minimal $x' \to x$ de $\mathcal{S}_x$ ;
\item pour tout morphisme d'objets formels $\xi'' \to \xi'$, le morphisme
correspondant $R'' \to R'$ est surjectif.
\end{enumerate}
\end{lemma}

\begin{proof}
Écrivons $\xi = (R, \xi_n, f_n)$. Posons $R'_1 = k$ et
$\xi'_1 = \xi_1$. Supposons construits des objets minimaux
$\xi'_m \to \xi_m$ de $\mathcal{S}_{\xi_m}$ situés au-dessus de
$R'_m \subset R/\mathfrak m_R^m$ pour $m \leq n$, ainsi que des morphismes
$f'_m : \xi'_{m + 1} \to \xi'_m$ compatibles avec $f_m$
pour $m \leq n - 1$. D'après
le lemme \ref{formal-defos-lemma-smallest-where-descends} (2), il existe un objet minimal
$\xi'_{n + 1} \to \xi_{n + 1}$ situé au-dessus de
$R'_{n + 1} \subset R/\mathfrak m_R^{n + 1}$ dont l'image est
$\xi'_n \to \xi_n$ au-dessus de $R'_n \subset R/\mathfrak m_R^n$. Cela donne
le diagramme commutatif
$$
\xymatrix{
\xi'_{n + 1} \ar[r]_{f'_n} \ar[d] & \xi'_n \ar[d] \\
\xi_{n + 1} \ar[r]^{f_n} & \xi_n
}
$$
par construction. De plus, le morphisme d'anneaux
$R'_{n + 1} \to R'_n$ est surjectif. Posons $R' = \lim_n R'_n$. Alors
$R' \to R$ est injectif.

\medskip\noindent
Toutefois, il n'est pas clair a priori que $R'$ soit noethérien. Pour le
montrer, on utilise la versalité de $\xi$. Celle-ci implique l'existence d'un
morphisme $\xi \to \xi'_n$ dans $\widehat{\mathcal{F}}$ ; voir le lemme
\ref{formal-defos-lemma-versal-object-quasi-initial}. Le morphisme correspondant
$R \to R'_n$ est nécessairement surjective (puisque
$\xi'_n \to \xi_n$ est minimal dans $\mathcal{S}_{\xi_n}$). Les dimensions
des espaces cotangents sont donc bornées, et le lemme
\ref{formal-defos-lemma-limit-artinian} implique que $R'$ est noethérien, c'est-à-dire un
objet de $\widehat{\mathcal{C}}_\Lambda$. D'après le lemme
\ref{formal-defos-lemma-formal-objects-different-filtration} (et le résultat sur les
filtrations du lemme \ref{formal-defos-lemma-limit-artinian}), la suite d'éléments
$\xi'_n$ définit un objet formel $\xi'$ au-dessus de $R'$, et l'on dispose
d'un morphisme $\xi' \to \xi$.

\medskip\noindent
Par construction, (1) est vérifiée pour $R \to R/\mathfrak m_R^n$, pour tout
$n$. Comme tout morphisme $R \to A$ de (1) se factorise par
$R \to R/\mathfrak m_R^n \to A$, la propriété (1) pour $x' \to x$ au-dessus
de $f(R) \subset A$ résulte, par le lemme
\ref{formal-defos-lemma-smallest-where-descends} (1), de la minimalité de
$\xi'_n \to \xi_n$ au-dessus de $R'_n \to R/\mathfrak m_R^n$.

\medskip\noindent
Si $R'' \to R'$ comme dans (2) n'était pas surjectif, alors
$R'' \to R' \to R'_n$ ne serait pas surjectif pour un certain $n$, et
$\xi'_n \to \xi_n$ ne serait pas minimal, contradiction. Cela démontre (2).
\end{proof}

\begin{lemma}
\label{formal-defos-lemma-descends-versal}
Soit $\mathcal{F}$ une catégorie cofibrée en groupoïdes au-dessus de
$\mathcal{C}_\Lambda$ qui satisfait (S1). Soit $\xi$ un objet formel versel de
$\mathcal{F}$ situé au-dessus de $R$. Soit $\xi' \to \xi$ un morphisme
d'objets formels situé au-dessus de $R' \subset R$, construit comme dans le
lemme \ref{formal-defos-lemma-smallest-where-descends-versal}. Alors
$$
R \cong R'[[x_1, \ldots, x_r]]
$$
est un anneau de séries formelles sur $R'$. De plus, $\xi'$ est lui aussi un
objet formel versel.
\end{lemma}

\begin{proof}
D'après le lemme \ref{formal-defos-lemma-versal-object-quasi-initial}, il existe un
morphisme $\xi \to \xi'$. Selon le lemme
\ref{formal-defos-lemma-smallest-where-descends-versal}, le morphisme correspondant
$f : R \to R'$ induit une surjection $f|_{R'} : R' \to R'$. C'est un
isomorphisme d'après Algèbre, lemme
\ref{algebra-lemma-surjective-endo-noetherian-ring-is-iso}. Ainsi,
$I = \Ker(f)$ est un idéal de $R$ tel que $R = R' \oplus I$. Soient
$x_1, \ldots, x_n \in I$ des éléments qui forment une base de
$I/\mathfrak m_RI$. Considérons le morphisme
$S = R'[[X_1, \ldots, X_r]] \to R$ qui envoie $X_i$ sur $x_i$. Pour tout
$n \geq 1$, on obtient une surjection de $R'$-algèbres artiniennes
$B = S/\mathfrak m_S^n \to R/\mathfrak m_R^n = A$. Notons
$y \in \Ob(\mathcal{F}(B)$, respectivement $x \in \Ob(\mathcal{F}(A))$,
l'image directe de $\xi'$ le long de $R' \to S \to B$, respectivement de
$R' \to S \to A$. Remarquons que $x$ est aussi l'image directe de $\xi$ le
long de $R \to A$, puisque $\xi$ est l'image directe de $\xi'$ le long de
$R' \to R$. On dispose donc d'un diagramme en traits pleins
$$
\vcenter{
\xymatrix{
& y \ar[d] \\
\xi \ar[r] \ar@{..>}[ru] & x
}
}
\quad\text{au-dessus de}\quad
\vcenter{
\xymatrix{
& S/\mathfrak m_S^n \ar[d] \\
R \ar[r] \ar@{..>}[ru] & R/\mathfrak m_R^n
}
}
$$
Comme $\xi$ est versel, la remarque \ref{formal-defos-remark-versal-object} fournit les
flèches en pointillés qui complètent ces diagrammes. En particulier, les
morphismes $S/\mathfrak m_S^n \to R/\mathfrak m_R^n$ admettent des sections
$h_n : R/\mathfrak m_R^n \to S/\mathfrak m_S^n$. Il résulte du lemme
\ref{formal-defos-lemma-power-series} que $S \to R$ est un isomorphisme.

\medskip\noindent
Puisque $\xi$ est une image directe de $\xi'$ le long de $R' \to R$, la
remarque \ref{formal-defos-remark-formal-objects-yoneda-map} fournit un diagramme
commutatif
$$
\xymatrix{
\underline{R}|_{\mathcal{C}_\Lambda} \ar[rr] \ar[rd]_{\underline{\xi}} & &
\underline{R'}|_{\mathcal{C}_\Lambda} \ar[ld]^{\underline{\xi'}} \\
& \mathcal{F}
}
$$
Comme $R' \to R$ admet un inverse à gauche (à savoir $R \to R/I = R'$),
$\underline{R}|_{\mathcal{C}_\Lambda} \to
\underline{R'}|_{\mathcal{C}_\Lambda}$ est essentiellement surjectif. Le
lemme \ref{formal-defos-lemma-smooth-properties} montre donc que $\underline{\xi'}$ est
lisse, c'est-à-dire que $\xi'$ est un objet formel versel.
\end{proof}

\noindent
Les lemmes précédents motivent la définition suivante.

\begin{definition}
\label{formal-defos-definition-minimal-versal}
Soit $\mathcal{F}$ une catégorie de prédéformation. On dit qu'un objet formel
versel $\xi$ de $\mathcal{F}$ est {\it minimal}\footnote{Cette terminologie
n'est peut-être pas standard. Beaucoup d'auteurs relient cette notion à des
propriétés des espaces tangents. Nous établirons ce lien dans la section
\ref{formal-defos-section-miniversal-objects-existence}.} si, pour tout morphisme d'objets
formels $\xi' \to \xi$, le morphisme sous-jacent sur les anneaux est
surjective. Un objet formel versel minimal est parfois appelé {\it miniversel}.
\end{definition}

\noindent
Les résultats de cette section montrent que cette définition est raisonnable.
Tout d'abord, l'existence d'un objet formel versel implique que
$\mathcal{F}$ satisfait (S1). Les lemmes précédents montrent ensuite qu'il
existe un objet formel versel minimal. Enfin, deux objets formels versels
minimaux quelconques sont isomorphes. Voici un résumé de nos résultats (avec
des démonstrations détaillées).

\begin{lemma}
\label{formal-defos-lemma-minimal-versal}
Soit $\mathcal{F}$ une catégorie de prédéformation qui possède un objet
formel versel. Alors :
\begin{enumerate}
\item $\mathcal{F}$ possède un objet formel versel minimal ;
\item les objets versels minimaux sont uniques à isomorphisme près ;
\item tout objet versel est l'image directe d'un objet versel minimal le long
d'une extension en séries formelles.
\end{enumerate}
\end{lemma}

\begin{proof}
Supposons que $\mathcal{F}$ possède un objet formel versel $\xi$ au-dessus de
$R$. Elle satisfait alors (S1) ; voir le lemme
\ref{formal-defos-lemma-versal-object-S1}. Soit $\xi' \to \xi$ au-dessus de
$R' \subset R$ l'un des morphismes construits dans le lemme
\ref{formal-defos-lemma-smallest-where-descends-versal}. D'après le lemme
\ref{formal-defos-lemma-descends-versal}, $\xi'$ est versel ; il est donc, par
construction, un objet formel versel minimal. Cela démontre (1). De plus,
$R \cong R'[[x_1, \ldots, x_n]]$, ce qui démontre (3).

\medskip\noindent
Supposons que $\xi_i/R_i$ soient deux objets formels versels minimaux. Le
lemme \ref{formal-defos-lemma-versal-object-quasi-initial} fournit des morphismes
$\xi_1 \to \xi_2$ et $\xi_2 \to \xi_1$. Les morphismes d'anneaux
correspondants $f : R_1 \to R_2$ et $g : R_2 \to R_1$ sont surjectifs par
minimalité. Les composées $g \circ f : R_1 \to R_1$ et
$f \circ g : R_2 \to R_2$ sont donc des isomorphismes d'après Algèbre, lemme
\ref{algebra-lemma-surjective-endo-noetherian-ring-is-iso}. Ainsi, $f$ et $g$
sont des isomorphismes, et par conséquent les morphismes
$\xi_1 \to \xi_2$ et $\xi_2 \to \xi_1$ sont des isomorphismes (car
$\widehat{\mathcal{F}}$ est cofibrée en groupoïdes d'après le lemme
\ref{formal-defos-lemma-completion-cofibred}). Cela démontre (2) et achève la démonstration
du lemme.
\end{proof}








\section{Objets formels miniversels et espaces tangents}
\label{formal-defos-section-miniversal-objects-existence}

\noindent
La notion générale de minimalité introduite dans la définition
\ref{formal-defos-definition-minimal-versal} peut parfois se déduire du comportement sur
les espaces tangents. Soit $\xi$ un objet formel de la catégorie de
prédéformation $\mathcal{F}$, et soit
$\underline{\xi} : \underline{R}|_{\mathcal{C}_\Lambda} \to \mathcal{F}$
le morphisme correspondant. On peut alors considérer la condition
\begin{equation}
\label{formal-defos-equation-bijective}
d\underline{\xi} : \text{Der}_\Lambda(R, k) \to T\mathcal{F} 
\text{ is bijective}
\end{equation}
et la condition
\begin{equation}
\label{formal-defos-equation-bijective-orbits}
d\underline{\xi} : \text{Der}_\Lambda(R, k) \to T\mathcal{F} 
\text{ is bijective on }\text{Der}_\Lambda(k, k)\text{-orbits.}
\end{equation}
Nous utilisons ici l'identification
$T\underline{R}|_{\mathcal{C}_\Lambda} = \text{Der}_\Lambda(R, k)$ de
l'exemple \ref{formal-defos-example-tangent-space-prorepresentable-functor}, ainsi que
l'action (\ref{formal-defos-equation-action}) des dérivations sur les espaces tangents. Si
$k' \subset k$ est séparable, alors $\text{Der}_\Lambda(k, k) = 0$ et les deux
conditions sont équivalentes. Il s'avère qu'en présence de la condition (S2),
un objet formel versel est minimal si et seulement si $\underline{\xi}$
satisfait (\ref{formal-defos-equation-bijective-orbits}). De plus, si $\underline{\xi}$
satisfait (\ref{formal-defos-equation-bijective}), alors $\mathcal{F}$ satisfait (S2).

\begin{lemma}
\label{formal-defos-lemma-miniversal-object-unique}
Soit $\mathcal{F}$ une catégorie de prédéformation. Soit $\xi$ un objet
formel versel de $\mathcal{F}$ satisfaisant
(\ref{formal-defos-equation-bijective-orbits}). Alors $\xi$ est un objet formel versel
minimal. En particulier, de tels $\xi$ sont uniques à isomorphisme près.
\end{lemma}

\begin{proof}
Si $\xi$ n'est pas minimal, il existe un morphisme $\xi' \to \xi$ situé
au-dessus de $R' \to R$ tel que $R = R'[[x_1, \ldots, x_n]]$ avec $n > 0$ ;
voir le lemme \ref{formal-defos-lemma-minimal-versal}. Ainsi, $d\underline{\xi}$ se
factorise en
$$
\text{Der}_\Lambda(R, k) \to
\text{Der}_\Lambda(R', k) \to T\mathcal{F}
$$
et l'on voit que (\ref{formal-defos-equation-bijective-orbits}) ne peut être satisfaite,
car $D : f \mapsto \partial/\partial x_1(f) \bmod \mathfrak m_R$ est un
élément du noyau de la première flèche qui n'appartient pas à l'image de
$\text{Der}_\Lambda(k, k) \to \text{Der}_\Lambda(R, k)$.
\end{proof}

\begin{lemma}
\label{formal-defos-lemma-miniversal-object-existence-1}
Soit $\mathcal{F}$ une catégorie de prédéformation. Soit $\xi$ un objet
formel versel de $\mathcal{F}$ satisfaisant (\ref{formal-defos-equation-bijective}). Alors :
\begin{enumerate}
\item $\mathcal{F}$ satisfait (S1).
\item $\mathcal{F}$ satisfait (S2).
\item $\dim_k T\mathcal{F}$ est finie.
\end{enumerate}
\end{lemma}

\begin{proof}
La condition (S1) résulte du lemme \ref{formal-defos-lemma-versal-object-S1}. La première
partie de (S2) est satisfaite puisque (S1) l'est. Soient
$$
\vcenter{
\xymatrix{
y \ar[r]_c \ar[d]_a & x_\epsilon \ar[d]^e \\
x \ar[r]^d          & x_0
}
}
\quad\text{et}\quad
\vcenter{
\xymatrix{
y' \ar[r]_{c'} \ar[d]_{a'} & x_\epsilon \ar[d]^e \\
x \ar[r]^d                 & x_0
}
}
\quad\text{au-dessus de}\quad
\vcenter{
\xymatrix{
A \times_k k[\epsilon] \ar[r] \ar[d] & k[\epsilon] \ar[d] \\
A  \ar[r] & k
}
}
$$
des diagrammes comme dans la seconde partie de (S2). Comme ci-dessus, on peut
trouver des morphismes $b : \xi \to y$ et $b' : \xi \to y'$ tels que
$$
\xymatrix{
\xi \ar[r]^{b'} \ar[d]_b          & y' \ar[d]^{a'} \\
y \ar[r]^{a} & x
}
$$
soit commutatif. Notons $p : \mathcal{F} \to \mathcal{C}_\Lambda$ le
morphisme structural. Écrivons $\widehat{p}(\xi) = R$, c'est-à-dire que $\xi$
est situé au-dessus de $R \in \Ob(\widehat{\mathcal{C}}_\Lambda)$. L'image
directe de $\xi$ par $p(c) \circ p(b)$ est $x_\epsilon$, et l'image directe de
$\xi$ par $p(c') \circ p(b')$ est $x_\epsilon$. Puisque $\xi$ satisfait
(\ref{formal-defos-equation-bijective}), on a
$p(c) \circ p(b) = p(c') \circ p(b')$ comme morphismes
$R \to k[\epsilon]$. Par conséquent, $p(b) = p(b')$ comme morphismes
$R \to A \times_k k[\epsilon]$. Ainsi, $y$ et $y'$ sont isomorphes à l'image
directe de $\xi$ le long de ce morphisme, et l'on obtient l'unique
morphisme recherché $y \to y'$ au-dessus de $A \times_k k[\epsilon]$,
compatible avec $b$ et $b'$.

\medskip\noindent
Enfin, l'exemple \ref{formal-defos-example-tangent-space-prorepresentable-functor} montre
que $\dim_k T\mathcal{F} = \dim_k T\underline{R}|_{\mathcal{C}_\Lambda}$
est finie.
\end{proof}

\begin{example}
\label{formal-defos-example-do-not-get-S2}
Il existe des catégories de prédéformation qui possèdent un objet formel
versel satisfaisant (\ref{formal-defos-equation-bijective-orbits}) mais ne satisfont pas
(S2). Un exemple immédiat consiste à prendre
$F = \underline{k[\epsilon]}/G$, où
$G \subset \text{Aut}_{\mathcal{C}_\Lambda}(k[\epsilon])$ est un sous-groupe
fini non trivial. En effet, le morphisme $\underline{k[\epsilon]} \to F$ est
lisse, mais l'espace tangent de $F$ ne possède pas de structure naturelle de
$k$-espace vectoriel (puisqu'il est le quotient d'un $k$-espace vectoriel par
un groupe fini).
\end{example}

\begin{lemma}
\label{formal-defos-lemma-construct-bijective-orbits}
Soit $\mathcal{F}$ une catégorie de prédéformation satisfaisant (S2) et
possédant un objet formel versel. Alors son objet formel versel minimal
satisfait (\ref{formal-defos-equation-bijective-orbits}).
\end{lemma}

\begin{proof}
Soit $\xi$ un objet formel versel minimal pour $\mathcal{F}$ ; voir le lemme
\ref{formal-defos-lemma-minimal-versal}. Supposons $\xi$ situé au-dessus de
$R \in \Ob(\widehat{\mathcal{C}}_\Lambda)$. Pour interpréter
(\ref{formal-defos-equation-bijective-orbits}), remarquons que $T\mathcal{F}$ possède une
structure naturelle de $k$-espace vectoriel (voir le lemme
\ref{formal-defos-lemma-tangent-space-vector-space}), que
$d\underline{\xi} : \text{Der}_\Lambda(R, k) \to T\mathcal{F}$ est linéaire
(voir le lemme \ref{formal-defos-lemma-k-linear-differential}), et que l'action de
$\text{Der}_\Lambda(k, k)$ est donnée par l'addition (voir le lemme
\ref{formal-defos-lemma-action-linear}). Considérons le diagramme
$$
\xymatrix{
& \Hom_k(\mathfrak m_R/\mathfrak m_R^2, k) \\
K \ar[r] & \text{Der}_\Lambda(R, k) \ar[r]^{d\underline{\xi}} \ar[u] &
T\mathcal{F} \\
& \text{Der}_\Lambda(k, k) \ar[u] \ar[ru]
}
$$
L'espace vectoriel $K$ est le noyau de $d\underline{\xi}$. Remarquons que la
colonne médiane est exacte en son milieu, puisqu'elle est duale de la suite
(\ref{formal-defos-equation-sequence}). Si (\ref{formal-defos-equation-bijective-orbits}) n'est pas
satisfaite, on peut trouver un élément non nul $D \in K$ dont l'image dans
$\Hom_k(\mathfrak m_R/\mathfrak m_R^2, k)$ n'est pas nulle. Il existe donc
$t \in \mathfrak m_R$ tel que $D(t) = 1$. Posons
$R' = \{a \in R \mid D(a) = 0\}$. Comme $D$ est une dérivation, c'est un
sous-anneau de $R$. Puisque $D(t) = 1$, le morphisme $R' \to k$ est
surjective (comparer avec la démonstration du lemme
\ref{formal-defos-lemma-essential-surjection}). Remarquons que
$\mathfrak m_{R'} = \Ker(D : \mathfrak m_R \to k)$ est un idéal de $R$ et
que $\mathfrak m_R^2 \subset \mathfrak m_{R'}$. Ainsi,
$$
\mathfrak m_R/\mathfrak m_R^2 =
\mathfrak m_{R'}/\mathfrak m_R^2 + k\overline{t}
$$
ce qui implique que le morphisme
$$
R'/\mathfrak m_R^2 \times_k k[\epsilon] \to R/\mathfrak m_R^2
$$
qui envoie $\epsilon$ sur $\overline{t}$ est un isomorphisme. En particulier,
il existe un morphisme $R/\mathfrak m_R^2 \to R'/\mathfrak m_R^2$.

\medskip\noindent
Soit $\xi \to y$ un morphisme situé au-dessus de
$R \to R/\mathfrak m_R^2$. Soit $y \to x$ un morphisme situé au-dessus de
$R/\mathfrak m_R^2 \to R'/\mathfrak m_R^2$. Soit $y \to x_\epsilon$ un
morphisme situé au-dessus de $R/\mathfrak m_R^2 \to k[\epsilon]$. Soit $x_0$
l'unique objet de $\mathcal{F}$ au-dessus de $k$, à isomorphisme unique près.
Les axiomes d'une catégorie cofibrée en groupoïdes donnent un diagramme
commutatif
$$
\vcenter{
\xymatrix{
y \ar[r] \ar[d] & x_\epsilon \ar[d] \\
x \ar[r]   & x_0
}
}
\quad\text{au-dessus de}\quad
\vcenter{
\xymatrix{
R'/\mathfrak m_R^2 \times_k k[\epsilon] \ar[r] \ar[d] & k[\epsilon] \ar[d] \\
R'/\mathfrak m_R^2 \ar[r] & k.
}
}
$$
Comme $D \in K$, on voit que $x_\epsilon$ est isomorphe à
$0 \in \mathcal{F}(k[\epsilon])$, c'est-à-dire que $x_\epsilon$ est l'image
directe de $x_0$ par $k \to k[\epsilon], a \mapsto a$. Le lemme
\ref{formal-defos-lemma-lifting-section} montre donc qu'il existe un morphisme $x \to y$.
Comme
$\text{length}_\Lambda(R'/\mathfrak m_R^2) <
\text{length}_\Lambda(R/\mathfrak m_R^2)$
le morphisme d'anneaux correspondant
$R'/\mathfrak m_R^2 \to R/\mathfrak m_R^2$ n'est pas surjectif. Cela
contredit la minimalité de $\xi/R$ ; voir la partie (1) du lemme
\ref{formal-defos-lemma-smallest-where-descends-versal}. Cette contradiction montre qu'un
tel $D$ ne peut exister, d'où le résultat.
\end{proof}

\begin{theorem}
\label{formal-defos-theorem-miniversal-object-existence}
Soit $\mathcal{F}$ une catégorie de prédéformation. Considérons les conditions
suivantes :
\begin{enumerate}
\item $\mathcal{F}$ possède un objet formel versel minimal satisfaisant
(\ref{formal-defos-equation-bijective}) ;
\item $\mathcal{F}$ possède un objet formel versel minimal satisfaisant
(\ref{formal-defos-equation-bijective-orbits}) ;
\item les conditions suivantes sont satisfaites :
\begin{enumerate}
\item $\mathcal{F}$ satisfait (S1).
\item $\mathcal{F}$ satisfait (S2).
\item $\dim_k T\mathcal{F}$ est finie.
\end{enumerate}
\end{enumerate}
On a toujours
$$
(1) \Rightarrow (3) \Rightarrow (2).
$$
Si $k' \subset k$ est séparable, les trois conditions sont équivalentes.
\end{theorem}

\begin{proof}
Le lemme \ref{formal-defos-lemma-miniversal-object-existence-1} montre que
(1) $\Rightarrow$ (3). Les lemmes \ref{formal-defos-lemma-versal-object-existence} et
\ref{formal-defos-lemma-construct-bijective-orbits} montrent que
(3) $\Rightarrow$ (2). Si $k' \subset k$ est séparable, alors
$\text{Der}_\Lambda(k, k) = 0$, et l'on voit que
(\ref{formal-defos-equation-bijective}) $=$ (\ref{formal-defos-equation-bijective-orbits}),
c'est-à-dire que (1) et (2) sont la même condition.

\medskip\noindent
Une autre démonstration de (3) $\Rightarrow$ (1) dans le cas classique
consiste à compléter légèrement la démonstration du lemme
\ref{formal-defos-lemma-versal-object-existence} pour voir que l'on peut construire
directement, dans ce cas, un objet versel qui satisfait
(\ref{formal-defos-equation-bijective}). Cela évite d'utiliser le lemme
\ref{formal-defos-lemma-versal-object-existence} dans le cas classique. Nous omettons les
détails.
\end{proof}

\begin{remark}
\label{formal-defos-remark-compare-schlessinger-H3}
Soit $F : \mathcal{C}_\Lambda \to \textit{Ensembles}$ un foncteur de
prédéformation satisfaisant (S1), (S2) et $\dim_k TF < \infty$. Rappelons que
ces conditions correspondent aux conditions (H1), (H2) et (H3) de l'article
de Schlessinger ; voir la remarque
\ref{formal-defos-remark-compare-schlessinger-H3-pre}. Dans le cas classique (ou si
$k' \subset k$ est séparable), nous introduisons à la suite de Schlessinger
la notion d'enveloppe : une {\it enveloppe} est un objet formel versel
$\xi \in \widehat{F}(R)$ tel que
$d\underline{\xi} : T\underline{R}|_{\mathcal{C}_\Lambda} \to TF$
soit un isomorphisme, c'est-à-dire que (\ref{formal-defos-equation-bijective}) soit
satisfaite. Ainsi, le théorème \ref{formal-defos-theorem-miniversal-object-existence} nous
dit que
$$
(H1) + (H2) + (H3)
\Rightarrow
\text{ there exists a hull}
$$
dans le cas classique. Autrement dit, notre théorème retrouve le théorème de
Schlessinger sur l'existence des enveloppes.
\end{remark}

\begin{remark}
\label{formal-defos-remark-compose-minimal-into-iso-classes}
Soit $\mathcal{F}$ une catégorie de prédéformation. Rappelons que
$\mathcal{F} \to \overline{\mathcal{F}}$ est lisse ; voir la remarque
\ref{formal-defos-remark-smooth-to-iso-classes}. Par conséquent, si
$\xi \in \widehat{\mathcal{F}}(R)$ est un objet formel versel, la composée
$$
\underline{R}|_{\mathcal{C}_\Lambda} \longrightarrow
\mathcal{F} \longrightarrow \overline{\mathcal{F}}
$$
est lisse (lemme \ref{formal-defos-lemma-smooth-properties}), et l'on conclut que l'image
$\overline{\xi}$ de $\xi$ dans $\overline{\mathcal{F}}$ est un objet formel
versel. Si (\ref{formal-defos-equation-bijective}) est satisfaite, alors
$\overline{\xi}$ induit un isomorphisme
$T\underline{R}|_{\mathcal{C}_\Lambda} \to T\overline{\mathcal{F}}$
car $\mathcal{F} \to \overline{\mathcal{F}}$ identifie les espaces tangents.
Dans ce cas, $\overline{\xi}$ est donc une enveloppe de
$\overline{\mathcal{F}}$ ; voir la remarque
\ref{formal-defos-remark-compare-schlessinger-H3}. D'après le théorème
\ref{formal-defos-theorem-miniversal-object-existence}, on peut toujours trouver un tel
$\xi$ si $k' \subset k$ est séparable et si $\mathcal{F}$ est une catégorie
de prédéformation satisfaisant (S1), (S2) et
$\dim_k T\mathcal{F} < \infty$.
\end{remark}

\begin{example}
\label{formal-defos-example-smooth-continued}
Dans le lemme \ref{formal-defos-lemma-exists-smooth}, nous avons construit des objets
$R \in \widehat{\mathcal{C}}_\Lambda$ tels que
$\underline{R}|_{\mathcal{C}_\Lambda}$ soit lisse et que
$$
H_1(L_{k/\Lambda}) = \mathfrak m_R/\mathfrak m_R^2
\quad\text{et}\quad
\Omega_{R/\Lambda} \otimes_R k = \Omega_{k/\Lambda}
$$
Réinterprétons ceci à l'aide du théorème précédent. Considérons
$\mathcal{F} = \mathcal{C}_\Lambda$ comme une catégorie cofibrée en
groupoïdes au-dessus d'elle-même (au moyen du foncteur identité). Alors
$\mathcal{F}$ est une catégorie de prédéformation, satisfait (S1) et (S2), et
l'on a $T\mathcal{F} = 0$. Ainsi, $\mathcal{F}$ satisfait la condition (3) du
théorème \ref{formal-defos-theorem-miniversal-object-existence}. Le théorème implique que
(2) est satisfaite, c'est-à-dire que l'on peut trouver un objet formel versel
minimal $\xi \in \widehat{\mathcal{F}}(S)$, situé au-dessus d'un certain
$S \in \widehat{\mathcal{C}}_\Lambda$, et satisfaisant
(\ref{formal-defos-equation-bijective-orbits}). Le lemme \ref{formal-defos-lemma-smooth} montre que
$\Lambda \to S$ est formellement lisse pour la topologie
$\mathfrak m_S$-adique (car
$\underline{\xi} : \underline{R}|_{\mathcal{C}_\Lambda} \to
\mathcal{F} = \mathcal{C}_\Lambda$ est lisse). La condition
(\ref{formal-defos-equation-bijective-orbits}) dit maintenant que
$\text{Der}_\Lambda(S, k) \to 0$ est bijective sur les orbites de
$\text{Der}_\Lambda(k, k)$. Cela signifie que l'injection
$\text{Der}_\Lambda(k, k) \to \text{Der}_\Lambda(S, k)$ est également
surjective. Autrement dit,
$\Omega_{S/\Lambda} \otimes_S k = \Omega_{k/\Lambda}$. Puisque
$\Lambda \to S$ est formellement lisse pour la topologie
$\mathfrak m_S$-adique, on peut appliquer le lemme
\ref{more-algebra-lemma-formally-smooth-JZ} de Compléments d'algèbre ; la suite
exacte (\ref{formal-defos-equation-sequence-extended}) se traduit alors par les deux identifications
$$
H_1(L_{k/\Lambda}) = \mathfrak m_S/\mathfrak m_S^2
\quad\text{et}\quad
\Omega_{S/\Lambda} \otimes_S k = \Omega_{k/\Lambda}
$$
En remontant l'argument, on voit que le $R$ construit dans le lemme
\ref{formal-defos-lemma-exists-smooth} porte un objet versel minimal. Par l'unicité des
objets versels minimaux (lemme \ref{formal-defos-lemma-minimal-versal}), on conclut aussi
que $R \cong S$, c'est-à-dire que les deux constructions donnent le même
résultat.
\end{example}







\section{Conditions de Rim--Schlessinger et catégories de déformation}
\sectionmark{Conditions de Rim--Schlessinger}
\label{formal-defos-section-RS-condition}

\noindent
Il existe une propriété très naturelle des catégories cofibrées en groupoïdes
au-dessus de $\mathcal{C}_\Lambda$, qui est facile à vérifier en pratique et
qui entraîne les propriétés (S1) et (S2) de Schlessinger introduites plus haut.

\begin{definition}
\label{formal-defos-definition-RS}
Soit $\mathcal{F}$ une catégorie cofibrée en groupoïdes au-dessus de $\mathcal
C_\Lambda$. On dit que $\mathcal{F}$ satisfait à la {\it condition (RS)} si,
pour tout diagramme de $\mathcal{F}$
$$
\vcenter{
\xymatrix{
           & x_2 \ar[d] \\
x_1 \ar[r] & x
}
}
\quad\text{au-dessus de}\quad
\vcenter{
\xymatrix{
           & A_2 \ar[d] \\
A_1 \ar[r] & A
}
}
$$
dans $\mathcal{C}_\Lambda$, où $A_2 \to A$ est surjectif, il existe un produit
fibré $x_1 \times_x x_2$ dans $\mathcal{F}$ tel que le diagramme
$$
\vcenter{
\xymatrix{
x_1 \times_x x_2 \ar[r] \ar[d] & x_2 \ar[d] \\
x_1 \ar[r]      & x
}
}
\quad\text{lies over}\quad
\vcenter{
\xymatrix{
A_1 \times_A A_2 \ar[r] \ar[d] & A_2 \ar[d] \\
A_1 \ar[r]      & A.
}
}
$$
\end{definition}

\begin{lemma}
\label{formal-defos-lemma-RS-fiber-square}
Soit $\mathcal{F}$ une catégorie cofibrée en groupoïdes au-dessus de
$\mathcal{C}_\Lambda$ satisfaisant à (RS). Étant donné un diagramme commutatif
de $\mathcal{F}$
$$
\vcenter{
\xymatrix{
y \ar[r] \ar[d] & x_2 \ar[d]   \\
x_1 \ar[r]      & x
}
}
\quad\text{au-dessus de}\quad
\vcenter{
\xymatrix{
A_1 \times_A A_2 \ar[r] \ar[d] & A_2 \ar[d] \\
A_1 \ar[r]      & A.
}
}
$$
où $A_2 \to A$ est surjectif, ce diagramme est cartésien.
\end{lemma}

\begin{proof}
Puisque $\mathcal{F}$ satisfait à (RS), il existe un diagramme de produit fibré
$$
\vcenter{
\xymatrix{
x_1 \times_x x_2 \ar[r] \ar[d] & x_2 \ar[d] \\
x_1 \ar[r]      & x
}
}
\quad\text{au-dessus de}\quad
\vcenter{
\xymatrix{
A_1 \times_A A_2 \ar[r] \ar[d] & A_2 \ar[d] \\
A_1 \ar[r]      & A.
}
}
$$
Le morphisme induit $y \to x_1 \times_x x_2$ est situé au-dessus de
$\text{id} : A_1 \times_A A_1 \to A_1 \times_A A_1$ ; c'est donc un
isomorphisme.
\end{proof}

\begin{lemma}
\label{formal-defos-lemma-RS-small-extension}
Soit $\mathcal{F}$ une catégorie cofibrée en groupoïdes au-dessus de $\mathcal
C_\Lambda$. Alors $\mathcal{F}$ satisfait à (RS) s'il suffit de supposer que
la condition de la définition \ref{formal-defos-definition-RS} est vérifiée lorsque
$A_2 \to A$ est une petite extension.
\end{lemma}

\begin{proof}
Appliquer le lemme \ref{formal-defos-lemma-factor-small-extension}. La démonstration est
analogue à celle du lemme \ref{formal-defos-lemma-smoothness-small-extensions}.
\end{proof}

\begin{lemma}
\label{formal-defos-lemma-RS-2-categorical}
Soit $\mathcal{F}$ une catégorie cofibrée en groupoïdes au-dessus de
$\mathcal{C}_\Lambda$. Les assertions suivantes sont équivalentes :
\begin{enumerate}
\item $\mathcal{F}$ satisfait à (RS) ;
\item le foncteur
$\mathcal{F}(A_1 \times_A A_2) \to
\mathcal{F}(A_1) \times_{\mathcal{F}(A)} \mathcal{F}(A_2)$
de (\ref{formal-defos-equation-compare}) est une équivalence de catégories dès que
$A_2 \to A$ est surjectif ;
\item il en est de même qu'en (2) dès que $A_2 \to A$ est une petite extension.
\end{enumerate}
\end{lemma}

\begin{proof}
Supposons (1). Le lemme \ref{formal-defos-lemma-RS-fiber-square} montre que tout objet de
$\mathcal{F}(A_1 \times_A A_2)$ est de la forme $x_1 \times_x x_2$. De plus,
$$
\Mor_{A_1 \times_A A_2}(x_1 \times_x x_2, y_1 \times_y y_2) =
\Mor_{A_1}(x_1, y_1) \times_{\Mor_A(x, y)}
\Mor_{A_2}(x_2, y_2).
$$
Ainsi, $\mathcal{F}(A_1 \times_A A_2)$ est un $2$-produit fibré de
$\mathcal{F}(A_1)$ et $\mathcal{F}(A_2)$ au-dessus de $\mathcal{F}(A)$ d'après
Catégories, remarque \ref{categories-remark-other-description-2-fibre-product}.
Autrement dit, (2) est vérifiée.

\medskip\noindent
L'implication (2) $\Rightarrow$ (3) est immédiate.

\medskip\noindent
Supposons (3). Soient $q_1 : A_1 \to A$ et $q_2 : A_2 \to A$, avec $q_2$
une petite extension. Nous utiliserons la description du $2$-produit fibré
$\mathcal{F}(A_1) \times_{\mathcal{F}(A)} \mathcal{F}(A_2)$ donnée dans
Catégories, remarque \ref{categories-remark-other-description-2-fibre-product}.
Soit donc $y \in \mathcal{F}(A_1 \times_A A_2)$ correspondant à
$(x_1, x_2, x, a_1 : x_1 \to x, a_2 : x_2 \to x)$.
Soit $z$ un objet de $\mathcal{F}$ au-dessus de $C$. Alors
\begin{align*}
\Mor_\mathcal{F}(z, y) & =
\{(f, \alpha) \mid f : C \to A_1 \times_A A_2,
\alpha : f_*z \to y\} \\
& = \{(f_1, f_2, \alpha_1, \alpha_2) \mid
f_i : C \to A_i, \ \alpha_i : f_{i, *}z \to x_i, \\
& \quad\quad\quad\quad
q_1 \circ f_1 = q_2 \circ f_2, \ q_{1, *} \alpha_1 = q_{2, *}\alpha_2\} \\
& =
\Mor_\mathcal{F}(z, x_1) \times_{\Mor_\mathcal{F}(z, x)}
\Mor_\mathcal{F}(z, x_2)
\end{align*}
d'où $y$ est un produit fibré de $x_1$ et $x_2$ au-dessus de $x$. Ainsi,
$\mathcal{F}$ satisfait à (RS) lorsque $A_2 \to A$ est une petite extension.
Le lemme \ref{formal-defos-lemma-RS-small-extension} montre donc que (RS) est vérifiée.
\end{proof}

\begin{remark}
\label{formal-defos-remark-compare-schlessinger-H4}
Lorsque $\mathcal{F}$ est cofibrée en ensembles, la condition (RS) est
exactement la condition (H4) de l'article de Schlessinger
\cite[théorème 2.11]{Sch}. Plus précisément, pour un foncteur
$F: \mathcal{C}_\Lambda \to \textit{Ensembles}$, la condition (RS) affirme ceci :
si $A_1 \to A$ et $A_2 \to A$ sont des morphismes de
$\mathcal{C}_\Lambda$, avec $A_2 \to A$ surjectif, alors l'application induite
$F(A_1 \times_A A_2) \to F(A_1) \times_{F(A)} F(A_2)$ est bijectif.
\end{remark}

\begin{lemma}
\label{formal-defos-lemma-RS-implies-S1-S2}
Soit $\mathcal{F}$ une catégorie cofibrée en groupoïdes au-dessus de
$\mathcal{C}_\Lambda$. La condition (RS) pour $\mathcal{F}$ entraîne à la
fois (S1) et (S2) pour $\mathcal{F}$.
\end{lemma}

\begin{proof}
La reformulation du lemme \ref{formal-defos-lemma-RS-2-categorical} et l'explication de
(S1) qui suit la définition \ref{formal-defos-definition-S1-S2} montrent immédiatement que
(RS) entraîne (S1). Cela démontre la première partie de (S2). La seconde partie
de (S2) résulte du lemme \ref{formal-defos-lemma-RS-fiber-square}, qui donne
$y = x_1 \times_{d, x_0, e} x_2 = y'$ lorsque $y, y'$ sont comme dans la
seconde partie de la définition de (S2), dans la définition
\ref{formal-defos-definition-S1-S2}. (En fait, le morphisme $y \to y'$ est compatible à la
fois avec $a, a'$ et avec $c, c'$ !)
\end{proof}

\noindent
Le lemme suivant est l'analogue du lemme
\ref{formal-defos-lemma-S1-S2-associated-functor}. Rappelons que, si $\mathcal{F}$ est une
catégorie cofibrée en groupoïdes au-dessus de $\mathcal{C}_\Lambda$ et si $x$
est un objet de $\mathcal{F}$ au-dessus de $A$, on note
$\text{Aut}_A(x) = \Mor_A(x, x) = \Mor_{\mathcal{F}(A)}(x, x)$.
Si $x' \to x$ est un morphisme de $\mathcal{F}$ au-dessus de $A' \to A$,
il existe un homomorphisme de groupes bien défini
$\text{Aut}_{A'}(x') \to \text{Aut}_A(x)$.

\begin{lemma}
\label{formal-defos-lemma-RS-associated-functor}
Soit $\mathcal{F}$ une catégorie cofibrée en groupoïdes au-dessus de
$\mathcal{C}_\Lambda$ satisfaisant à (RS). Les conditions suivantes sont
équivalentes :
\begin{enumerate}
\item $\overline{\mathcal{F}}$ satisfait à (RS).
\item Soient $f_1: A_1 \to A$ et $f_2: A_2 \to A$ des morphismes
d'anneaux de $\mathcal{C}_\Lambda$, avec $f_2$ surjectif. L'application induite
entre ensembles de classes d'isomorphie
$$
\overline{\mathcal{F}(A_1) \times_{\mathcal{F}(A)} \mathcal{F}(A_2)}
\to \overline{\mathcal{F}}(A_1) \times_{\overline{\mathcal{F}}(A)}
\overline{\mathcal{F}}(A_2)
$$
est injective.
\item Pour tout morphisme $x' \to x$ de $\mathcal{F}$ au-dessus d'un
morphisme d'anneaux surjectif $A' \to A$, l'homomorphisme
$\text{Aut}_{A'}(x') \to \text{Aut}_A(x)$ est surjective.
\item Pour tout morphisme $x' \to x$ de $\mathcal{F}$ au-dessus d'une petite
extension $A' \to A$, l'homomorphisme
$\text{Aut}_{A'}(x') \to \text{Aut}_A(x)$ est surjective.
\end{enumerate}
\end{lemma}

\begin{proof}
Nous montrons que (1) est équivalente à (2), puis que (2) est équivalente à
(3). L'équivalence de (3) et (4) résulte du lemme
\ref{formal-defos-lemma-factor-small-extension}.

\medskip\noindent
Soient $f_1: A_1 \to A$ et $f_2: A_2 \to A$ des morphismes d'anneaux de
$\mathcal{C}_\Lambda$, avec $f_2$ surjectif. D'après la remarque
\ref{formal-defos-remark-compare-schlessinger-H4}, $\overline{\mathcal{F}}$ satisfait à
(RS) si et seulement si l'application
$$
\overline{\mathcal{F}}(A_1 \times_A A_2) \to \overline{\mathcal
F}(A_1) \times_{\overline{\mathcal{F}}(A)} \overline{\mathcal{F}}(A_2)
$$
est bijective pour tous $f_1, f_2$ de cette sorte. Cette application est au
moins surjective, car telle est la condition (S1), et
$\overline{\mathcal{F}}$ satisfait à (S1) d'après les lemmes
\ref{formal-defos-lemma-RS-implies-S1-S2} et \ref{formal-defos-lemma-S1-S2-associated-functor}. De plus,
elle se factorise sous la forme
$$
\overline{\mathcal{F}}(A_1 \times_A A_2)
\longrightarrow
\overline{\mathcal{F}(A_1) \times_{\mathcal{F}(A)} \mathcal{F}(A_2)}
\longrightarrow
\overline{\mathcal{F}}(A_1) \times_{\overline{\mathcal{F}}(A)}
\overline{\mathcal{F}}(A_2),
$$
où la première flèche est une bijection puisque
$$
\mathcal{F}(A_1 \times_A A_2)
\longrightarrow
\mathcal{F}(A_1) \times_{\mathcal{F}(A)} \mathcal{F}(A_2)
$$
est une équivalence par (RS) pour $\mathcal{F}$. Ainsi, (1) est équivalente à
(2).

\medskip\noindent
Supposons (2). Soit $x' \to x$ un morphisme de $\mathcal{F}$ au-dessus d'un
morphisme d'anneaux surjectif $f: A' \to A$. Soit
$a \in \text{Aut}_A(x)$. Les objets
$$
(x', x', a : x \to x), \ (x', x', \text{id} : x \to x)
$$
de
$\mathcal{F}(A') \times_{\mathcal{F}(A)} \mathcal{F}(A')$
ont la même image dans
$\overline{\mathcal{F}}(A') \times_{\overline{\mathcal{F}}(A)}
\overline{\mathcal{F}}(A')$. D'après (2), il existe des morphismes
$b_1, b_2 : x' \to x'$ tels que
$$
\xymatrix{
x \ar[r]_a \ar[d]_{f_*b_1} & x \ar[d]^{f_*b_2} \\
x \ar[r]^{\text{id}} & x
}
$$
soit commutatif. Ainsi, l'image de
$b_2^{-1} \circ b_1 \in \text{Aut}_{A'}(x')$ est
$a \in \text{Aut}_A(x)$. L'assertion (3) est donc vérifiée.

\medskip\noindent
Supposons (3). Supposons que
$$
(x_1, x_2, a : (f_1)_*x_1 \to (f_2)_*x_2),
\ (x'_1, x'_2, a' : (f_1)_*x'_1 \to (f_2)_*x'_2)
$$
soient des objets de
$\mathcal{F}(A_1) \times_{\mathcal{F}(A)} \mathcal{F}(A_2)$
ayant la même image dans
$\overline{\mathcal{F}}(A_1) \times_{\overline{\mathcal{F}}(A)}
\overline{\mathcal{F}}(A_2)$. Il existe alors des morphismes $b_1: x_1 \to
x'_1$ dans $\mathcal{F}(A_1)$ et $b_2: x_2 \to x'_2$ dans $\mathcal
F(A_2)$. D'après (3), on peut modifier $b_2$ par un automorphisme de $x_2$
au-dessus de $A_2$ de telle sorte que le diagramme
$$
\xymatrix{
(f_1)_*x_1 \ar[r]_a \ar[d]_{(f_1)_*b_1} & (f_2)_*x_2 \ar[d]^{(f_2)_*b_2} \\
(f_1)_*x'_1 \ar[r]^{a'} & (f_2)_*x'_2.
}
$$
soit commutatif. Cela démontre $(x_1, x_2, a) \cong (x'_1, x'_2, a')$ dans
$\overline{\mathcal{F}(A_1) \times_{\mathcal{F}(A)} \mathcal{F}(A_2)}$.
Ainsi, (2) est vérifiée.
\end{proof}

\noindent
Définissons enfin la notion de catégorie de déformation.

\begin{definition}
\label{formal-defos-definition-deformation-category}
Une {\it catégorie de déformation} est une catégorie de prédéformation
$\mathcal{F}$ satisfaisant à (RS). Un morphisme de catégories de déformation
est un morphisme de catégories au-dessus de $\mathcal{C}_\Lambda$.
\end{definition}

\begin{remark}
\label{formal-defos-remark-deformation-functor}
On dit qu'un foncteur $F: \mathcal{C}_\Lambda \to \textit{Ensembles}$ est un
{\it foncteur de déformation} si l'ensemble cofibré associé est une catégorie
de déformation, c'est-à-dire si $F(k)$ est un ensemble à un seul élément et si $F$
satisfait à (RS). Si $\mathcal{F}$ est une catégorie de déformation, alors
$\overline{\mathcal{F}}$ est un foncteur de prédéformation, mais pas
nécessairement un foncteur de déformation, comme le montre le lemme
\ref{formal-defos-lemma-RS-associated-functor}.
\end{remark}

\begin{example}
\label{formal-defos-example-prorepresentable-deformation-functor}
Un foncteur pro-représentable $F$ est un foncteur de déformation. En effet,
supposons $R \in \Ob(\widehat{\mathcal{C}}_\Lambda)$ et
$F(A) = \Mor_{\widehat{\mathcal{C}}_\Lambda}(R, A)$. Il existe un unique
morphisme $R \to k$, de sorte que $F(k)$ est un ensemble à un seul élément. Comme
$$
\Hom_\Lambda(R, A_1 \times_A A_2) =
\Hom_\Lambda(R, A_1) \times_{\Hom_\Lambda(R, A)}
\Hom_\Lambda(R, A_2)
$$
il en est de même pour les morphismes de $\widehat{\mathcal{C}}_\Lambda$, et
l'on voit que $F$ satisfait à (RS).
\end{example}

\noindent
Voici l'une de nos remarques typiques sur le passage d'une catégorie cofibrée
en groupoïdes à la catégorie de prédéformation en un point au-dessus de $k$ :
ce passage préserve (RS).

\begin{lemma}
\label{formal-defos-lemma-localize-RS}
Soit $\mathcal{F}$ une catégorie cofibrée en groupoïdes au-dessus de
$\mathcal{C}_\Lambda$. Soit $x_0 \in \Ob(\mathcal{F}(k))$. Soit
$\mathcal{F}_{x_0}$ la catégorie cofibrée en groupoïdes au-dessus de
$\mathcal{C}_\Lambda$ construite dans la remarque
\ref{formal-defos-remark-localize-cofibered-groupoid}. Si $\mathcal{F}$ satisfait à (RS),
alors $\mathcal{F}_{x_0}$ y satisfait aussi. En particulier,
$\mathcal{F}_{x_0}$ est une catégorie de déformation.
\end{lemma}

\begin{proof}
Tout diagramme comme dans la définition \ref{formal-defos-definition-RS} au sein de
$\mathcal{F}_{x_0}$ donne lieu à un diagramme dans $\mathcal{F}$, et le
résultat de (RS) pour ce diagramme dans $\mathcal{F}$ peut également être
considéré comme un résultat pour $\mathcal{F}_{x_0}$.
\end{proof}

\noindent
Le lemme suivant est l'analogue du fait que les $2$-produits fibrés de champs
algébriques sont des champs algébriques.

\begin{lemma}
\label{formal-defos-lemma-RS-fiber-product-morphisms}
Supposons que
$$
\xymatrix{
\mathcal{H} \times_\mathcal{F} \mathcal{G} \ar[r] \ar[d] &
\mathcal{G} \ar[d]^g \\
\mathcal{H} \ar[r]^f & \mathcal{F}
}
$$
soit un $2$-produit fibré de catégories cofibrées en groupoïdes au-dessus de
$\mathcal{C}_\Lambda$. Si $\mathcal{F}, \mathcal{G}, \mathcal{H}$ satisfont
toutes à (RS), alors $\mathcal{H} \times_\mathcal{F} \mathcal{G}$ satisfait à
(RS).
\end{lemma}

\begin{proof}
Si $A$ est un objet de $\mathcal{C}_\Lambda$, un objet de la catégorie fibre
de $\mathcal{H} \times_\mathcal{F} \mathcal{G}$ au-dessus de $A$ est un
triplet $(u, v, a)$, où $u \in \mathcal{H}(A)$, $v \in \mathcal{G}(A)$ et
$a : f(u) \to g(v)$ est un morphisme de $\mathcal{F}(A)$. Considérons un
diagramme de $\mathcal{H} \times_\mathcal{F} \mathcal{G}$
$$
\vcenter{
\xymatrix{
           & (u_2, v_2, a_2) \ar[d] \\
(u_1, v_1, a_1) \ar[r] & (u, v, a)
}
}
\quad\text{au-dessus de}\quad
\vcenter{
\xymatrix{
           & A_2 \ar[d] \\
A_1 \ar[r] & A
}
}
$$
dans $\mathcal{C}_\Lambda$, avec $A_2 \to A$ surjectif. Puisque
$\mathcal{H}$ et $\mathcal{G}$ satisfont à (RS), il existe des produits fibrés
$u_1 \times_u u_2$ et $v_1 \times_v v_2$ au-dessus de
$A_1 \times_A A_2$. Puisque $\mathcal{F}$ satisfait à (RS), le lemme
\ref{formal-defos-lemma-RS-fiber-square} montre que
$$
\vcenter{
\xymatrix{
f(u_1 \times_u u_2) \ar[r] \ar[d] & f(u_2) \ar[d] \\
f(u_1) \ar[r] & f(u)
}
}
\quad\text{et}\quad
\vcenter{
\xymatrix{
g(v_1 \times_v v_2) \ar[r] \ar[d] & g(v_2) \ar[d] \\
g(v_1) \ar[r] & g(v)
}
}
$$
sont tous deux des carrés cartésiens dans $\mathcal{F}$. On peut donc voir
$a_1 \times_a a_2$ comme un morphisme de $f(u_1 \times_u u_2)$ vers
$g(v_1 \times_v v_2)$ au-dessus de $A_1 \times_A A_2$. Il s'ensuit que
$$
\xymatrix{
 (u_1 \times_u u_2, v_1 \times_v v_2, a_{1} \times_a a_2) \ar[d] \ar[r] &
(u_2, v_2, a_2) \ar[d] \\
(u_1, v_1, a_1) \ar[r] & (u, v, a)
}
$$
est un carré cartésien dans $\mathcal{H} \times_\mathcal{F} \mathcal{G}$,
comme voulu.
\end{proof}


















\section{Relèvements d'objets}
\label{formal-defos-section-lifts}

\noindent
Le contenu de cette section est que, pour une catégorie de déformation,
l'ensemble des relèvements le long d'une petite extension est un espace
principal homogène sous l'espace tangent.

\begin{definition}
\label{formal-defos-definition-lifts}
Soit $\mathcal{F}$ une catégorie cofibrée en groupoïdes au-dessus de
$\mathcal{C}_\Lambda$. Soit $f: A' \to A$ un morphisme de
$\mathcal{C}_\Lambda$. Soit $x \in \mathcal{F}(A)$. La catégorie
$\textit{Lift}(x, f)$ des relèvements de $x$ le long de $f$ est la catégorie
dont les objets et les morphismes sont les suivants.
\begin{enumerate}
\item Objets : un {\it relèvement de $x$ le long de $f$} est un morphisme
$x' \to x$ au-dessus de $f$.
\item Morphismes : un {\it morphisme de relèvements} de
$a_1 : x'_1 \to x$ vers $a_2 : x'_2 \to x$ est un morphisme
$b : x'_1 \to x'_2$ dans $\mathcal{F}(A')$ tel que $a_2 = a_1 \circ b$.
\end{enumerate}
L'ensemble $\text{Lift}(x, f)$ des relèvements de $x$ le long de $f$ est
l'ensemble des classes d'isomorphie de $\textit{Lift}(x, f)$.
\end{definition}

\begin{remark}
\label{formal-defos-remark-omit-arrow}
Lorsque le morphisme $f: A' \to A$ ressort clairement du contexte, on peut
écrire $\textit{Lift}(x, A')$ et $\text{Lift}(x, A')$ au lieu de
$\textit{Lift}(x, f)$ et $\text{Lift}(x, f)$.
\end{remark}

\begin{remark}
\label{formal-defos-remark-tangent-space-lifting}
Soit $\mathcal{F}$ une catégorie cofibrée en groupoïdes au-dessus de
$\mathcal{C}_\Lambda$. Soit $x_0 \in \Ob(\mathcal{F}(k))$. Soit $V$ un espace
vectoriel de dimension finie. Alors $\text{Lift}(x_0, k[V])$ est l'ensemble
des classes d'isomorphie de $\mathcal{F}_{x_0}(k[V])$, où
$\mathcal{F}_{x_0}$ est la catégorie de prédéformation des objets de
$\mathcal{F}$ au-dessus de $x_0$ ; voir la remarque
\ref{formal-defos-remark-localize-cofibered-groupoid}. Par conséquent, si $\mathcal{F}$
satisfait à (S2), il en est de même de $\mathcal{F}_{x_0}$ (voir le lemme
\ref{formal-defos-lemma-S1-S2-localize}), et le lemme
\ref{formal-defos-lemma-tangent-space-vector-space} donne
$$
\text{Lift}(x_0, k[V]) = T\mathcal{F}_{x_0} \otimes_k V
$$
comme espaces vectoriels sur $k$.
\end{remark}

\begin{remark}
\label{formal-defos-remark-lift-bijections}
Soit $\mathcal{F}$ une catégorie cofibrée en groupoïdes au-dessus de $\mathcal
C_\Lambda$ satisfaisant à (RS). Soit
$$
\xymatrix{
A_1 \times_A A_2 \ar[r] \ar[d] & A_2 \ar[d] \\
A_1 \ar[r] & A
}
$$
un carré cartésien dans $\mathcal{C}_\Lambda$, tel que $A_1 \to A$ ou
$A_2 \to A$ soit surjectif. Soit $x \in \Ob(\mathcal{F}(A))$. Étant donnés
des relèvements $x_1 \to x$ et $x_2 \to x$ de $x$ à $A_1$ et $A_2$, (RS)
fournit un relèvement $x_1 \times_x x_2 \to x$ de $x$ à
$A_1 \times_A A_2$. Réciproquement, d'après le lemme
\ref{formal-defos-lemma-RS-fiber-square}, tout relèvement de $x$ à
$A_1 \times_A A_2$ est de cette forme. On obtient donc une bijection
$$
\text{Lift}(x, A_1) \times \text{Lift}(x, A_2)
\longrightarrow
\text{Lift}(x, A_1 \times_A A_2).
$$
De même, si $x_1 \to x$ est un relèvement fixé de $x$ à $A_1$, il existe une
bijection
$$
\text{Lift}(x_1, A_1 \times_A A_2)
\longrightarrow
\text{Lift}(x, A_2).
$$
Soit maintenant
$$
\xymatrix{
A_1' \times_A A_2 \ar[r] \ar[d] & A_1 \times_A A_2 \ar[r] \ar[d] & A_2
\ar[d] \\
A_1' \ar[r] & A_1 \ar[r] & A
}
$$
une composée de carrés cartésiens dans $\mathcal{C}_\Lambda$, où
$A'_1 \to A_1$ et $A_1 \to A$ sont tous deux surjectifs. Soit $x_1 \to x$ un
morphisme au-dessus de $A_1 \to A$. D'après ce qui précède, on a alors les
bijections
\begin{align*}
\text{Lift}(x_1, A_1' \times_A A_2)
& = \text{Lift}(x_1, A_1') \times \text{Lift}(x_1, A_1 \times_A A_2) \\
& = \text{Lift}(x_1, A_1') \times \text{Lift}(x, A_2).
\end{align*}
\end{remark}

\begin{lemma}
\label{formal-defos-lemma-free-transitive-action}
Soit $\mathcal{F}$ une catégorie de déformation. Soit $A' \to A$ un
morphisme d'anneaux surjectif de $\mathcal{C}_\Lambda$, dont le noyau $I$
est annulé par $\mathfrak m_{A'}$. Soit $x \in \Ob(\mathcal{F}(A))$. Si
$\text{Lift}(x, A')$ est non vide, alors $T\mathcal{F} \otimes_k I$ agit sur
$\text{Lift}(x, A')$ et en fait un espace principal homogène.
\end{lemma}

\begin{proof}
Considérons le morphisme d'anneaux $g : A' \times_A A' \to k[I]$ défini
par la règle $g(a_1, a_2) = \overline{a_1} \oplus a_2 - a_1$ (comparer au
lemme \ref{formal-defos-lemma-lifting-along-small-extension}). Il existe un isomorphisme
$$
A' \times_A A' \xrightarrow{\sim} A' \times_k k[I]
$$
donné par $(a_1, a_2) \mapsto (a_1, g(a_1, a_2))$. Cet isomorphisme est
compatible avec les projections sur le premier facteur $A'$, et donc avec celles de
$A' \times_A A'$ et $A' \times_k k[I]$ vers $A$. Il existe ainsi une bijection
\begin{equation}
\label{formal-defos-equation-one}
\text{Lift}(x, A' \times_A A')
\longrightarrow
\text{Lift}(x, A' \times_k k[I])
\end{equation}
D'après la remarque \ref{formal-defos-remark-lift-bijections}, il existe une bijection
\begin{equation}
\label{formal-defos-equation-two}
\text{Lift}(x, A') \times \text{Lift}(x, A')
\longrightarrow
\text{Lift}(x, A' \times_A A')
\end{equation}
On a un diagramme commutatif
$$
\xymatrix{
A' \times_k k[I] \ar[r] \ar[d] & A \times_k k[I] \ar[r] \ar[d] & k[I] \ar[d] \\
A' \ar[r] & A \ar[r] & k.
}
$$
Ainsi, si l'on choisit une image directe $x \to x_0$ de $x$ le long de
$A \to k$, la fin de la remarque \ref{formal-defos-remark-lift-bijections} fournit une
bijection
\begin{equation}
\label{formal-defos-equation-three}
\text{Lift}(x, A' \times_k k[I])
\longrightarrow
\text{Lift}(x, A') \times \text{Lift}(x_0, k[I])
\end{equation}
En composant (\ref{formal-defos-equation-two}), (\ref{formal-defos-equation-one}) et
(\ref{formal-defos-equation-three}), on obtient une bijection
$$
\Phi :
\text{Lift}(x, A') \times \text{Lift}(x, A')
\longrightarrow
\text{Lift}(x, A') \times \text{Lift}(x_0, k[I]).
$$
Cette bijection est compatible avec les projections sur les premiers facteurs. La remarque
\ref{formal-defos-remark-tangent-space-lifting} donne
$\text{Lift}(x_0, k[I]) = T\mathcal{F} \otimes_k I$. Si $\text{pr}_2$ est la
seconde projection de $\text{Lift}(x, A') \times \text{Lift}(x, A')$, on
obtient une application
$$
a = \text{pr}_2 \circ \Phi^{-1} :
\text{Lift}(x, A') \times (T\mathcal{F} \otimes_k I)
\longrightarrow
\text{Lift}(x, A').
$$
En explicitant ce qui précède, on voit que $a(x' \to x, \theta)$ est l'unique
relèvement $x'' \to x$ tel que $g_*(x', x'') = \theta$ dans
$\text{Lift}(x_0, k[I]) = T\mathcal{F} \otimes_k I$. Pour voir qu'il s'agit
d'une action de $T\mathcal{F} \otimes_k I$ sur $\text{Lift}(x, A')$, il faut
montrer ceci : si $x', x'', x'''$ sont des relèvements de $x$ et si
$g_*(x', x'') = \theta$, $g_*(x'', x''') = \theta'$, alors
$g_*(x', x''') = \theta + \theta'$. Cela résulte du diagramme commutatif
$$
\xymatrix{
A' \times_A A' \times_A A'
\ar[rrrrr]_-{(a_1, a_2, a_3) \mapsto (g(a_1, a_2), g(a_2, a_3))}
\ar[rrrrrd]_{(a_1, a_2, a_3) \mapsto g(a_1, a_3)} & & & & &
k[I] \times_k k[I] = k[I \times I] \ar[d]^{+} \\
& & & & & k[I]
}
$$
L'action est libre et transitive puisque $\Phi$ est bijective.
\end{proof}

\begin{remark}
\label{formal-defos-remark-free-transitive-action-functorial}
L'action du lemme \ref{formal-defos-lemma-free-transitive-action} est fonctorielle. Soit
$\varphi : \mathcal{F} \to \mathcal{G}$ un morphisme de catégories de
déformation. Soit $A' \to A$ un morphisme d'anneaux surjectif dont le
noyau $I$ est annulé par $\mathfrak m_{A'}$. Soit
$x \in \Ob(\mathcal{F}(A))$. Dans cette situation, $\varphi$ induit les
flèches verticales du diagramme commutatif suivant
$$
\xymatrix{
\text{Lift}(x, A') \times (T\mathcal{F} \otimes_k I)
\ar[d]_{(\varphi, d\varphi \otimes \text{id}_I)} \ar[r] &
\text{Lift}(x, A') \ar[d]^\varphi \\
\text{Lift}(\varphi(x), A') \times (T\mathcal{G} \otimes_k I) \ar[r] &
\text{Lift}(\varphi(x), A')
}
$$
La commutativité résulte du fait que chacune des applications
(\ref{formal-defos-equation-two}), (\ref{formal-defos-equation-one}) et (\ref{formal-defos-equation-three}) de la
démonstration du lemme \ref{formal-defos-lemma-free-transitive-action} donne lieu à un
diagramme commutatif analogue.
\end{remark}






\section{Théorème de Schlessinger sur les foncteurs pro-représentables}
\label{formal-defos-section-schlessingers-theorem}

\noindent
Nous en déduisons le théorème de Schlessinger qui caractérise les foncteurs
pro-représentables sur $\mathcal{C}_\Lambda$.

\begin{lemma}
\label{formal-defos-lemma-minimal-smooth-morphism-functors}
Soient $F, G: \mathcal{C}_\Lambda \to \textit{Ensembles}$ des foncteurs de
déformation. Soit $\varphi : F \to G$ un morphisme lisse qui induit un
isomorphisme $d\varphi : TF \to TG$ d'espaces tangents. Alors $\varphi$ est un
isomorphisme.
\end{lemma}

\begin{proof}
Nous montrons que $F(A) \to G(A)$ est une bijection pour tout $A \in
\Ob(\mathcal{C}_\Lambda)$, par récurrence sur
$\text{length}_A(A)$. Pour $A = k$, l'assertion résulte de
l'hypothèse que $F$ et $G$ sont des foncteurs de déformation. Supposons
l'assertion vraie pour les anneaux de longueur strictement inférieure à $n$ et
soit $A'$ un anneau de longueur $n$. Choisissons une petite extension
$f : A' \to A$. On a un diagramme commutatif
$$
\xymatrix{
F(A') \ar[r] \ar[d]_{F(f)} & G(A') \ar[d]^{G(f)} \\
F(A) \ar[r]^{\sim} & G(A)
}
$$
où l'application $F(A) \to G(A)$ est une bijection. Par la lissité de $F
\to G$, $F(A') \to G(A')$ est surjective (lemme
\ref{formal-defos-lemma-smooth-morphism-essentially-surjective}). Il suffit donc de vérifier
la bijectivité sur les fibres $F(f)^{-1}(x) \to
G(f)^{-1}(\varphi(x))$, pour les $x \in F(A)$ tels que $F(f)^{-1}(x)$ soit
non vide. Ces fibres sont précisément $\text{Lift}(x, A')$ et
$\text{Lift}(\varphi(x), A')$, et l'hypothèse fournit un isomorphisme
$d\varphi \otimes \text{id} :
TF \otimes_k \Ker(f)  \to TG \otimes_k \Ker(f)$.
Ainsi, d'après le lemme \ref{formal-defos-lemma-free-transitive-action} et la remarque
\ref{formal-defos-remark-free-transitive-action-functorial}, pour tout $x \in F(A)$ tel que
$F(f)^{-1}(x)$ soit non vide, l'application
$F(f)^{-1}(x) \to G(f)^{-1}(\varphi(x))$ est équivariante pour les actions
libres et transitives de
$TF \otimes_k \Ker(f)$. Elle est donc bijective.
\end{proof}

\noindent
Notons que, lorsque l'extension $k' \subset k$ est séparable, la condition (c)
du théorème ci-dessous est vide.

\begin{theorem}
\label{formal-defos-theorem-Schlessinger-prorepresentability}
Soit $F: \mathcal{C}_\Lambda \to \textit{Ensembles}$ un foncteur. Alors $F$ est
pro-représentable si et seulement si (a) $F$ est un foncteur de déformation,
(b) $\dim_k TF$ est fini, et (c) $\gamma : \text{Der}_\Lambda(k, k) \to TF$
est injective.
\end{theorem}

\begin{proof}
Supposons $F$ pro-représenté par
$R \in \widehat{\mathcal{C}}_\Lambda$. L'exemple
\ref{formal-defos-example-prorepresentable-deformation-functor} montre que $F$ est un
foncteur de déformation. L'exemple
\ref{formal-defos-example-tangent-space-prorepresentable-functor} montre que
$\dim_k TF$ est fini. Enfin, l'exemple
\ref{formal-defos-example-tangent-space-map-prorepresentable-functor} identifie
$\text{Der}_\Lambda(k, k) \to TF$ à
$\text{Der}_\Lambda(k, k) \to \text{Der}_\Lambda(R, k)$, qui est injectif
puisque $R \to k$ est surjectif.

\medskip\noindent
Réciproquement, supposons (a), (b) et (c) vérifiées. Le lemme
\ref{formal-defos-lemma-RS-implies-S1-S2} montre que (S1) et (S2) sont vérifiées. Le théorème
\ref{formal-defos-theorem-miniversal-object-existence} fournit donc un objet formel versel
minimal $\xi$ de $F$ tel que (\ref{formal-defos-equation-bijective-orbits}) soit vérifiée.
Disons que $\xi$ est situé au-dessus de $R$. L'application
$$
d\underline{\xi} : \text{Der}_\Lambda(R, k) \to T\mathcal{F}
$$
est bijective sur les orbites de $\text{Der}_\Lambda(k, k)$. Puisque l'action
de $\text{Der}_\Lambda(k, k)$ sur le membre de gauche est libre d'après (c) et
le lemme \ref{formal-defos-lemma-action-linear}, cette application est bijective. Le lemme
\ref{formal-defos-lemma-minimal-smooth-morphism-functors} montre donc que
$\underline{\xi}$ est un isomorphisme.
\end{proof}




\section{Automorphismes infinitésimaux}
\label{formal-defos-section-infinitesimal-automorphisms}

\noindent
Soit $\mathcal{F}$ une catégorie cofibrée en groupoïdes au-dessus de
$\mathcal{C}_\Lambda$. Étant donné un morphisme $x' \to x$ de $\mathcal{F}$
au-dessus de $A' \to A$, on a un homomorphisme de groupes induit
$$
\text{Aut}_{A'}(x') \to \text{Aut}_A(x).
$$
Le lemme \ref{formal-defos-lemma-RS-associated-functor} affirme que le conoyau de cet
homomorphisme détermine si la condition (RS) sur $\mathcal{F}$ passe à
$\overline{\mathcal{F}}$. Dans cette section, nous étudions le noyau de cet
homomorphisme. Nous verrons qu'il mesure lui aussi l'écart entre $\mathcal{F}$
et $\overline{\mathcal{F}}$.

\begin{definition}
\label{formal-defos-definition-relative-infinitesimal-auts}
Soit $\mathcal{F}$ une catégorie cofibrée en groupoïdes au-dessus de $\mathcal
C_\Lambda$. Soit $x' \to x$ un morphisme de $\mathcal{F}$ au-dessus de
$A' \to A$. Le noyau
$$
\text{Inf}(x'/x) = \Ker(\text{Aut}_{A'}(x') \to \text{Aut}_A(x))
$$
est le {\it groupe des automorphismes infinitésimaux de $x'$ au-dessus de $x$}.
\end{definition}

\begin{definition}
\label{formal-defos-definition-infinitesimal-auts}
Soit $\mathcal{F}$ une catégorie cofibrée en groupoïdes au-dessus de $\mathcal
C_\Lambda$. Soit $x_0 \in \Ob(\mathcal{F}(k))$. Supposons choisie une image
directe $x_0 \to x_0'$ de $x_0$ le long du morphisme
$k \to k[\epsilon], a \mapsto a$. Il existe alors un unique morphisme
$x'_0 \to x_0$ tel que $x_0 \to x_0' \to x_0$ soit l'identité de $x_0$.
Alors
$$
\text{Inf}_{x_0}(\mathcal F) = \text{Inf}(x'_0/x_0)
$$
est le {\it groupe des automorphismes infinitésimaux de $x_0$}.
\end{definition}

\begin{remark}
\label{formal-defos-remark-choice-pushforward-immaterial-infinitesimal-aut}
À isomorphisme canonique près, $\text{Inf}_{x_0}(\mathcal{F})$ ne dépend pas
du choix de l'image directe $x_0 \to x_0'$, car deux images directes
quelconques sont canoniquement isomorphes. De plus, si
$y_0 \in \mathcal{F}(k)$ et $x_0 \cong y_0$ dans $\mathcal{F}(k)$, alors
$\text{Inf}_{x_0}(\mathcal{F}) \cong \text{Inf}_{y_0}(\mathcal{F})$,
où l'isomorphisme ne dépend que du choix d'un isomorphisme $x_0 \to y_0$. En
particulier, $\text{Aut}_k(x_0)$ agit sur $\text{Inf}_{x_0}(\mathcal{F})$.
\end{remark}

\begin{remark}
\label{formal-defos-remark-trivial-aut-point}
Supposons que $\mathcal{F}$ soit une catégorie de prédéformation. Alors
\begin{enumerate}
\item pour $x_0 \in \Ob(\mathcal{F}(k))$, le groupe d'automorphismes
$\text{Aut}_k(x_0)$ est trivial et, par conséquent,
$\text{Inf}_{x_0}(\mathcal{F}) = \text{Aut}_{k[\epsilon]}(x'_0)$ ;
\item pour $x_0, y_0 \in \Ob(\mathcal{F}(k))$, il existe un unique
isomorphisme $x_0 \to y_0$, et donc une identification canonique
$\text{Inf}_{x_0}(\mathcal{F}) = \text{Inf}_{y_0}(\mathcal{F})$.
\end{enumerate}
Puisque $\mathcal{F}(k)$ est non vide, en choisissant
$x_0 \in \Ob(\mathcal{F}(k))$ et en posant
$$
\text{Inf}(\mathcal{F}) = \text{Inf}_{x_0}(\mathcal{F})
$$
on obtient un {\it groupe des automorphismes infinitésimaux de $\mathcal{F}$}
bien défini. Avec cette notation, on a
$\text{Inf}(\mathcal{F}_{x_0}) = \text{Inf}_{x_0}(\mathcal{F})$.
Comparer à l'égalité
$T\mathcal{F}_{x_0} = T_{x_0}\mathcal{F}$
de la remarque \ref{formal-defos-remark-tangent-space-cofibered-groupoid}.
\end{remark}

\noindent
Nous verrons que $\text{Inf}_{x_0}(\mathcal{F})$ est muni d'une structure
naturelle d'espace vectoriel sur $k$ lorsque $\mathcal{F}$ satisfait à (RS).
Nous verrons en même temps que, si $\mathcal{F}$ satisfait à (RS), les
automorphismes infinitésimaux
$\text{Inf}(x'/x)$ d'un morphisme $x' \to x$ au-dessus d'une petite extension
sont gouvernés par $\text{Inf}_{x_0}(\mathcal{F})$, où $x_0$ est une image
directe de $x$ dans $\mathcal{F}(k)$. À cette fin, nous introduisons comme suit
le foncteur des automorphismes de tout objet $x \in \Ob(\mathcal{F})$.

\begin{definition}
\label{formal-defos-definition-automorphism-functor}
Soit $p : \mathcal{F} \to \mathcal{C}$ une catégorie cofibrée en groupoïdes
au-dessus d'une catégorie de base arbitraire $\mathcal{C}$. Supposons effectué
un choix d'images directes. Soient $x \in \Ob(\mathcal{F})$ et $U = p(x)$.
Notons $U/\mathcal{C}$ la catégorie des objets sous $U$. Le {\it foncteur des
automorphismes de $x$} est le foncteur
$\mathit{Aut}(x) : U/\mathcal{C} \to \textit{Ensembles}$ qui envoie un objet
$f : U \to V$ sur $\text{Aut}_V(f_*x)$ et un morphisme
$$
\xymatrix{
V' \ar[rr] &                    & V\\
          & U \ar[ul]^{f'}  \ar[ur]_f &
}
$$
sur l'homomorphisme de groupes
$\text{Aut}_{V'}(f'_*x) \to \text{Aut}_V(f_*x)$
provenant de l'unique morphisme $f'_*x \to f_*x$ situé au-dessus de
$V' \to V$ et compatible avec $x \to f'_*x$ et $x \to f_*x$.
\end{definition}

\noindent
Nous nous intéresserons aux foncteurs des automorphismes d'objets d'une
catégorie cofibrée en groupoïdes $\mathcal{F}$ au-dessus de
$\mathcal{C}_\Lambda$. Si $A \in \Ob(\mathcal{C}_\Lambda)$, la catégorie
$A/\mathcal{C}_\Lambda$ n'est autre que la catégorie $\mathcal{C}_A$,
c'est-à-dire la catégorie définie dans la section \ref{formal-defos-section-CLambda}, où
l'on prend $\Lambda = A$ et $k = A/\mathfrak m_A$. Par conséquent, le foncteur
des automorphismes d'un objet $x \in \Ob(\mathcal{F}(A))$ est un foncteur
$\mathit{Aut}(x) : \mathcal{C}_A \to \textit{Ensembles}$.

\medskip\noindent
Le lemme suivant pourrait se déduire du lemme
\ref{formal-defos-lemma-RS-fiber-product-morphisms} en considérant l'« inertie » d'une
catégorie cofibrée en groupoïdes ; voir par exemple Champs, section
\ref{stacks-section-the-inertia-stack}, et Catégories, section
\ref{categories-section-inertia}. Il est toutefois plus facile de le vérifier
directement.

\begin{lemma}
\label{formal-defos-lemma-Aut-functor-RS}
Soit $\mathcal{F}$ une catégorie cofibrée en groupoïdes au-dessus de
$\mathcal{C}_\Lambda$ satisfaisant à (RS). Soit
$x \in \Ob(\mathcal{F}(A))$. Alors
$\mathit{Aut}(x): \mathcal{C}_A \to \textit{Ensembles}$ satisfait à (RS).
\end{lemma}

\begin{proof}
Le fait que $\mathit{Aut}(x)$ satisfait à (RS) résulte de la pleine fidélité du
foncteur
$\mathcal{F}(A_1 \times_A A_2) \to
\mathcal{F}(A_1) \times_{\mathcal{F}(A)} \mathcal{F}(A_2)$ du lemme
\ref{formal-defos-lemma-RS-2-categorical}.
\end{proof}

\begin{lemma}
\label{formal-defos-lemma-Aut-functor-tangent-space}
Soit $\mathcal{F}$ une catégorie cofibrée en groupoïdes au-dessus de
$\mathcal{C}_\Lambda$ satisfaisant à (RS). Soit
$x \in \Ob(\mathcal{F}(A))$. Soit $x_0$ une image directe de $x$ dans
$\mathcal{F}(k)$.
\begin{enumerate}
\item $T_{\text{id}_{x_0}} \mathit{Aut}(x)$ possède une structure naturelle
d'espace vectoriel sur $k$ telle que l'addition coïncide avec la composition
dans $T_{\text{id}_{x_0}} \mathit{Aut}(x)$. En particulier, la composition
dans $T_{\text{id}_{x_0}} \mathit{Aut}(x)$ est commutative.
\item Il existe un isomorphisme canonique
$T_{\text{id}_{x_0}} \mathit{Aut}(x) \to
T_{\text{id}_{x_0}} \mathit{Aut}(x_0)$
d'espaces vectoriels sur $k$.
\end{enumerate}
\end{lemma}

\begin{proof}
Appliquons la remarque \ref{formal-defos-remark-localize-cofibered-groupoid} au foncteur
$\mathit{Aut}(x) : \mathcal{C}_A \to \textit{Ensembles}$ et à l'élément
$\text{id}_{x_0} \in \mathit{Aut}(x)(k)$. On obtient un foncteur de
prédéformation $F = \mathit{Aut}(x)_{\text{id}_{x_0}}$. D'après les lemmes
\ref{formal-defos-lemma-Aut-functor-RS} et \ref{formal-defos-lemma-localize-RS}, $F$ est un foncteur de
déformation. Par définition,
$T_{\text{id}_{x_0}} \mathit{Aut}(x) = TF = F(k[\epsilon])$
possède la structure naturelle d'espace vectoriel sur $k$ décrite dans le
lemme \ref{formal-defos-lemma-tangent-space-functor}.

\medskip\noindent
L'addition est définie comme la composée
$$
F(k[\epsilon]) \times F(k[\epsilon]) \longrightarrow
F(k[\epsilon] \times_k k[\epsilon]) \longrightarrow
F(k[\epsilon])
$$
où la première application est l'inverse de la bijection garantie par (RS) et
la seconde est induite par le morphisme de $k$-algèbres
$k[\epsilon] \times_k k[\epsilon] \to k[\epsilon]$
qui envoie $(\epsilon, 0)$ et $(0, \epsilon)$ sur $\epsilon$. Si $A \to B$
est un morphisme d'anneaux de $\mathcal{C}_\Lambda$, alors
$F(A) \to F(B)$ est un homomorphisme de groupes, où
$F(A) = \mathit{Aut}(x)_{\text{id}_{x_0}}(A)$ et
$F(B) = \mathit{Aut}(x)_{\text{id}_{x_0}}(B)$ sont des groupes pour la
composition. On en déduit que
$+ : F(k[\epsilon]) \times F(k[\epsilon])\to F(k[\epsilon])$
est un homomorphisme de groupes, $F(k[\epsilon])$ étant considéré comme groupe pour la
composition. En prenant $\text{id} \in F(k[\epsilon])$ pour élément neutre, on
voit que $+(v, \text{id}) =
+(\text{id}, v) = v$ pour tout
$v \in F(k[\epsilon])$, car $(\text{id}, v)$ est l'image directe de $v$ le
long du morphisme d'anneaux
$k[\epsilon] \to k[\epsilon] \times_k k[\epsilon]$ défini par
$\epsilon \mapsto (\epsilon, 0)$. En général, si $G$ est un groupe de loi
$\circ$ et si $+ : G \times G \to G$ est un homomorphisme de groupes tel que
$+(g, 1) = +(1, g) = g$, où $1$ est l'élément neutre de $G$, alors
$+ = \circ$. Cela montre que l'addition de la structure d'espace vectoriel sur
$k$ de $F(k[\epsilon])$ coïncide avec la composition.

\medskip\noindent
Enfin, (2) s'obtient en explicitant les définitions. En effet,
$T_{\text{id}_{x_0}} \mathit{Aut}(x)$ est l'ensemble des automorphismes
$\alpha$ de l'image directe de $x$ le long de $A \to k \to k[\epsilon]$ qui
sont triviaux modulo $\epsilon$. D'autre part,
$T_{\text{id}_{x_0}} \mathit{Aut}(x_0)$ est l'ensemble des automorphismes de
l'image directe de $x_0$ le long de $k \to k[\epsilon]$ qui sont triviaux
modulo $\epsilon$. Puisque $x_0$ est l'image directe de $x$ le long de
$A \to k$, le résultat est clair.
\end{proof}

\begin{remark}
\label{formal-defos-remark-infaut-lifting-equalities}
Signalons quelques relations élémentaires entre les groupes d'automorphismes
infinitésimaux, les relèvements et les espaces tangents aux foncteurs des
automorphismes. Soit $\mathcal{F}$ une catégorie cofibrée en groupoïdes
au-dessus de $\mathcal{C}_\Lambda$. Soit $x' \to x$ un morphisme situé au-dessus
d'un morphisme d'anneaux $A' \to A$. Les définitions donnent alors l'égalité
$$
\text{Inf}(x'/x) = \text{Lift}(\text{id}_x, A')
$$
où il s'agit des relèvements de $\text{id}_x$ comme objet de
$\mathit{Aut}(x')$. Si $x_0 \in \Ob(\mathcal{F}(k))$ et si $x'_0$ est son
image directe dans $\mathcal{F}(k[\epsilon])$, alors, en appliquant ce qui précède
à $x'_0 \to x_0$, on obtient
$$
\text{Inf}_{x_0}(\mathcal{F}) =
\text{Lift}(\text{id}_{x_0}, k[\epsilon]) =
T_{\text{id}_{x_0}} \mathit{Aut}(x_0),
$$
la dernière égalité résultant directement des définitions.
\end{remark}

\begin{lemma}
\label{formal-defos-lemma-infaut-vector-space}
Soit $\mathcal{F}$ une catégorie cofibrée en groupoïdes au-dessus de
$\mathcal{C}_\Lambda$ satisfaisant à (RS). Soit
$x_0 \in \Ob(\mathcal{F}(k))$. Alors, comme ensemble,
$\text{Inf}_{x_0}(\mathcal{F})$ est égal à
$T_{\text{id}_{x_0}} \mathit{Aut}(x_0)$ ; il possède donc une structure
naturelle d'espace vectoriel sur $k$ telle que l'addition coïncide avec la
composition des automorphismes.
\end{lemma}

\begin{proof}
L'égalité d'ensembles est celle obtenue à la fin de la remarque
\ref{formal-defos-remark-infaut-lifting-equalities}, et l'assertion sur la structure
d'espace vectoriel résulte du lemme \ref{formal-defos-lemma-Aut-functor-tangent-space}.
\end{proof}

\begin{lemma}
\label{formal-defos-lemma-k-linear-infaut}
Soit $\varphi : \mathcal{F} \to \mathcal{G}$ un morphisme de catégories
cofibrées en groupoïdes au-dessus de $\mathcal{C}_\Lambda$ satisfaisant à
(RS). Soit $x_0 \in \Ob(\mathcal{F}(k))$. Alors $\varphi$ induit une
application $k$-linéaire
$\text{Inf}_{x_0}(\mathcal{F}) \to \text{Inf}_{\varphi(x_0)}(\mathcal{G})$.
\end{lemma}

\begin{proof}
Il est clair que $\varphi$ induit un morphisme de foncteurs
$\mathit{Aut}(x_0) \to \mathit{Aut}(\varphi(x_0))$
qui envoie l'identité sur l'identité. Le résultat découle donc de celui sur les
espaces tangents ; voir le lemme \ref{formal-defos-lemma-k-linear-differential}.
\end{proof}

\begin{lemma}
\label{formal-defos-lemma-lifted-automorphisms-torsor}
Soit $\mathcal{F}$ une catégorie cofibrée en groupoïdes au-dessus de
$\mathcal{C}_\Lambda$ satisfaisant à (RS). Soit $x' \to x$ un morphisme
situé au-dessus d'un morphisme d'anneaux surjectif $A' \to A$, de noyau $I$
annulé par $\mathfrak m_{A'}$. Soit $x_0$ une image directe de $x$ dans
$\mathcal{F}(k)$. Alors $\text{Inf}(x'/x)$ est un espace principal homogène
sous $T_{\text{id}_{x_0}} \mathit{Aut}(x') \otimes_k I
= \text{Inf}_{x_0}(\mathcal{F}) \otimes_k I$.
\end{lemma}

\begin{proof}
C'est simplement l'analogue du lemme \ref{formal-defos-lemma-free-transitive-action} dans
le cadre des faisceaux d'automorphismes. Plus précisément, appliquons la
remarque \ref{formal-defos-remark-localize-cofibered-groupoid} au foncteur
$\mathit{Aut}(x') : \mathcal{C}_{A'} \to \textit{Ensembles}$ et à l'élément
$\text{id}_{x_0} \in \mathit{Aut}(x)(k)$. On obtient un foncteur de
prédéformation $F = \mathit{Aut}(x')_{\text{id}_{x_0}}$. D'après les lemmes
\ref{formal-defos-lemma-Aut-functor-RS} et \ref{formal-defos-lemma-localize-RS}, $F$ est un foncteur de
déformation. Le lemme \ref{formal-defos-lemma-free-transitive-action} donne donc une action
libre et transitive de $TF \otimes_k I$ sur
$\text{Lift}(\text{id}_x, A')$. Comme
$\text{Lift}(\text{id}_x, A')$ est un groupe, cet ensemble est toujours non
vide. Notons que l'on a les égalités d'espaces vectoriels
$$
TF = T_{\text{id}_{x_0}} \mathit{Aut}(x') \otimes_k I =
\text{Inf}_{x_0}(\mathcal{F}) \otimes_k I
$$
d'après le lemme \ref{formal-defos-lemma-Aut-functor-tangent-space}. L'égalité
$\text{Inf}(x'/x) = \text{Lift}(\text{id}_x, A')$ de la remarque
\ref{formal-defos-remark-infaut-lifting-equalities} achève la démonstration.
\end{proof}

\begin{lemma}
\label{formal-defos-lemma-infaut-trivial}
Soit $\mathcal{F}$ une catégorie cofibrée en groupoïdes au-dessus de
$\mathcal{C}_\Lambda$ satisfaisant à (RS). Soit $x' \to x$ un morphisme de
$\mathcal{F}$ situé au-dessus d'un morphisme d'anneaux surjectif. Soit $x_0$ une
image directe de $x$ dans $\mathcal{F}(k)$. Si
$\text{Inf}_{x_0}(\mathcal{F}) = 0$, alors $\text{Inf}(x'/x) = 0$.
\end{lemma}

\begin{proof}
Cela résulte des lemmes \ref{formal-defos-lemma-factor-small-extension} et
\ref{formal-defos-lemma-lifted-automorphisms-torsor}.
\end{proof}

\begin{lemma}
\label{formal-defos-lemma-infdef-trivial}
Soit $\mathcal{F}$ une catégorie cofibrée en groupoïdes au-dessus de
$\mathcal{C}_\Lambda$ satisfaisant à (RS). Soit
$x_0 \in \Ob(\mathcal{F}(k))$. Alors $\text{Inf}_{x_0}(\mathcal{F}) = 0$ si
et seulement si le morphisme naturel
$\mathcal{F}_{x_0} \to \overline{\mathcal{F}_{x_0}}$ de catégories cofibrées
en groupoïdes est une équivalence.
\end{lemma}

\begin{proof}
Le morphisme $\mathcal{F}_{x_0} \to \overline{\mathcal{F}_{x_0}}$ est une
équivalence si et seulement si $\mathcal{F}_{x_0}$ est fibrée en ensoïdes ;
cf.\ Catégories, section \ref{categories-section-fibred-in-setoids} (un
ensoïde est, par définition, un groupoïde dans lequel l'identité est le seul
automorphisme de chaque objet). Montrons que
$\text{Inf}_{x_0}(\mathcal{F}) = 0$ si et seulement si cette condition est
vérifiée pour $\mathcal{F}_{x_0}$. Il est évident que, si
$\mathcal{F}_{x_0}$ est fibrée en ensoïdes, alors
$\text{Inf}_{x_0}(\mathcal{F}) = 0$. Réciproquement, supposons
$\text{Inf}_{x_0}(\mathcal{F}) = 0$. Soit $A$ un objet de
$\mathcal{C}_\Lambda$. D'après le lemme \ref{formal-defos-lemma-infaut-trivial},
$\text{Inf}(x/x_0) = 0$ pour tout objet $x \to x_0$ de
$\mathcal{F}_{x_0}(A)$. Puisque, par définition, $\text{Inf}(x/x_0)$ est le
groupe des automorphismes de $x \to x_0$ dans $\mathcal{F}_{x_0}(A)$, cela
montre que $\mathcal{F}_{x_0}(A)$ est un ensoïde.
\end{proof}








\section{Applications}
\label{formal-defos-section-applications}

\noindent
Rassemblons quelques résultats sur les catégories de déformation dont nous
nous servirons plus loin.

\begin{lemma}
\label{formal-defos-lemma-deformation-categories-fiber-product-morphisms}
Soient $f : \mathcal{H} \to \mathcal{F}$ et
$g : \mathcal{G} \to \mathcal{F}$ des $1$-morphismes entre catégories de
déformation. Alors
\begin{enumerate}
\item $\mathcal{W} = \mathcal{H} \times_\mathcal{F} \mathcal{G}$ est une
catégorie de déformation ;
\item on a une suite exacte à $6$ termes d'espaces vectoriels
$$
0 \to \text{Inf}(\mathcal{W})
\to \text{Inf}(\mathcal{H}) \oplus \text{Inf}(\mathcal{G})
\to \text{Inf}(\mathcal{F}) \to
T\mathcal{W} \to T\mathcal{H} \oplus T\mathcal{G} \to T\mathcal{F}
$$
\end{enumerate}
\end{lemma}

\begin{proof}
La partie (1) résulte du lemme \ref{formal-defos-lemma-RS-fiber-product-morphisms} et du fait
que $\mathcal{W}(k)$ est le produit fibré de deux ensoïdes ayant une unique
classe d'isomorphie au-dessus d'un ensoïde ayant une unique classe
d'isomorphie.

\medskip\noindent
Partie (2). Soit $w_0 \in \Ob(\mathcal{W}(k))$, et soient $x_0, y_0, z_0$
l'image de $w_0$ dans $\mathcal{F}, \mathcal{H}, \mathcal{G}$. Alors
$\text{Inf}(\mathcal{W}) = \text{Inf}_{w_0}(\mathcal{W})$, et de même pour
$\mathcal{H}$, $\mathcal{G}$ et $\mathcal{F}$ ; voir la remarque
\ref{formal-defos-remark-trivial-aut-point}. Les lemmes \ref{formal-defos-lemma-k-linear-differential}
et \ref{formal-defos-lemma-k-linear-infaut} fournissent toutes les applications linéaires,
sauf le « morphisme bord »
$\delta : \text{Inf}_{x_0}(\mathcal{F}) \to T\mathcal{W}$. Nous insérerons
plus tard les signes convenables.

\medskip\noindent
Construction de $\delta$. Choisissons une image directe $w_0 \to w'_0$ le
long de $k \to k[\epsilon]$. Notons $x'_0, y'_0, z'_0$ les images de $w'_0$
dans $\mathcal{F}, \mathcal{H}, \mathcal{G}$. On obtient notamment des
isomorphismes $b' : f(y'_0) \to x'_0$ et $c' : x'_0 \to g(z'_0)$. Notons
$b : f(y_0) \to x_0$ et $c : x_0 \to g(z_0)$ les images directes le long de
$k[\epsilon] \to k$. Cela signifie que
$w'_0 = (k[\epsilon], y'_0, z'_0, c' \circ b')$ et
$w_0 = (k, y_0, z_0, c \circ b)$, eu égard à la forme explicite du produit
fibré de catégories ; voir les remarques \ref{formal-defos-remarks-cofibered-groupoids}
(\ref{formal-defos-item-fibre-product}). Étant donné $\alpha : x'_0 \to x'_0$, posons
$\delta(\alpha) = (k[\epsilon], y'_0, z'_0, c' \circ \alpha \circ b')$
qui est bien un objet de $\mathcal{W}$ au-dessus de $k[\epsilon]$ et qui est
muni d'un morphisme
$(k[\epsilon], y'_0, z'_0, c' \circ \alpha \circ b') \to w_0$ au-dessus de
$k[\epsilon] \to k$, puisque l'image directe de $\alpha$ est l'identité
au-dessus de $k$. Plus généralement, pour tout espace vectoriel sur $k$, disons
$V$, on peut définir une application
$$
\text{Lift}(\text{id}_{x_0}, k[V])
\longrightarrow
\text{Lift}(w_0, k[V])
$$
au moyen d'exactement les mêmes formules. Cette construction est fonctorielle
en l'espace vectoriel $V$ (nous omettons les détails). Le lemme
\ref{formal-defos-lemma-morphism-linear-functors} montre donc que $\delta$ est $k$-linéaire.

\medskip\noindent
Ces applications étant construites, il est immédiat de montrer que la suite
est exacte. L'injectivité de la première application résulte du fait que
$f \times g : \mathcal{W} \to \mathcal{H} \times \mathcal{G}$ est fidèle. Si
$(\beta, \gamma) \in
\text{Inf}_{y_0}(\mathcal{H}) \oplus \text{Inf}_{z_0}(\mathcal{G})$
ont la même image dans $\text{Inf}_{x_0}(\mathcal{F})$, alors
$(\beta, \gamma)$ définit un automorphisme de
$w'_0 = (k[\epsilon], y'_0, z'_0, c' \circ b')$, d'où l'exactitude au
deuxième terme. Si $\alpha$ ci-dessus donne la déformation triviale
$(k[\epsilon], y'_0, z'_0, c' \circ \alpha \circ b')$
de $w_0$, alors l'isomorphisme
$w'_0 = (k[\epsilon], y'_0, z'_0, c' \circ b') \to
(k[\epsilon], y'_0, z'_0, c' \circ \alpha \circ b')$
produit une paire $(\beta, \gamma)$ qui est un antécédent de $\alpha$. Si
$w = (k[\epsilon], y, z, \phi)$ est une déformation de $w_0$ telle que
$y'_0 \cong y$ et $z \cong z'_0$, alors le morphisme
$$
f(y'_0) \to f(y) \xrightarrow{\phi} g(z) \to g(z'_0)
$$
définit un automorphisme $\alpha$ qui a pour image $w$ par $\delta$.
Enfin, si $y$ et $z$ sont des déformations de $y_0$ et $z_0$ et s'il existe un isomorphisme
$\phi : f(y) \to g(z)$ de déformations de
$f(y_0) = x_0 = g(z_0)$, on obtient un antécédent
$w = (k[\epsilon], y, z, \phi)$ de $(x, y)$ dans $T\mathcal{W}$. Cela achève
la démonstration.
\end{proof}

\begin{lemma}
\label{formal-defos-lemma-map-fibre-products-smooth}
Soient $\mathcal{H}_1 \to \mathcal{G}$, $\mathcal{H}_2 \to \mathcal{G}$ et
$\mathcal{G} \to \mathcal{F}$ des morphismes de catégories cofibrées en
groupoïdes au-dessus de $\mathcal{C}_\Lambda$. Supposons que
\begin{enumerate}
\item $\mathcal{F}$ et $\mathcal{G}$ soient des catégories de déformation ;
\item $T\mathcal{G} \to T\mathcal{F}$ soit injectif ;
\item $\text{Inf}(\mathcal{G}) \to \text{Inf}(\mathcal{F})$ soit surjectif.
\end{enumerate}
Alors $\mathcal{H}_1 \times_\mathcal{G} \mathcal{H}_2 \to
\mathcal{H}_1 \times_\mathcal{F} \mathcal{H}_2$ est lisse.
\end{lemma}

\begin{proof}
Notons $p_i : \mathcal{H}_i \to \mathcal{G}$ et
$q : \mathcal{G} \to \mathcal{F}$ les morphismes donnés. Soit $A' \to A$ une
petite extension de $\mathcal{C}_\Lambda$. Un objet de
$\mathcal{H}_1 \times_\mathcal{F} \mathcal{H}_2$ au-dessus de $A'$ est un
triplet $(x'_1, x'_2, a')$, où $x'_i$ est un objet de $\mathcal{H}_i$
au-dessus de $A'$ et où $a' : q(p_1(x'_1)) \to q(p_2(x'_2))$ est un morphisme
de la catégorie fibre de $\mathcal{F}$ au-dessus de $A'$. L'image directe le
long de $A' \to A$ donne $(x_1, x_2, a)$. Un relèvement fixé de cet objet en
$\mathcal{H}_1 \times_\mathcal{G} \mathcal{H}_2$ au-dessus de $A$
consiste en la donnée d'un morphisme $b : p_1(x_1) \to p_2(x_2)$ situé au-dessus de $A$ tel que
$q(b) = a$. Il faut donc montrer que l'on peut relever $b$ en un morphisme
$b' : p_1(x'_1) \to p_2(x'_2)$ dont l'image par $q$ est $a'$.

\medskip\noindent
Observons que l'on peut considérer
$$
p_1(x'_1) \to p_1(x_1) \xrightarrow{b} p_2(x_2)
\quad\text{et}\quad
p_2(x'_2) \to p_2(x_2)
$$
comme deux objets de $\textit{Lift}(p_2(x_2), A' \to A)$. Le foncteur $q$
envoie ces objets sur les deux objets
$$
q(p_1(x'_1)) \to q(p_1(x_1)) \xrightarrow{b} q(p_2(x_2))
\quad\text{et}\quad
q(p_2(x'_2)) \to q(p_2(x_2))
$$
de $\textit{Lift}(q(p_2(x_2)), A' \to A)$, lesquels sont isomorphes au moyen
du morphisme $a' : q(p_1(x'_1)) \to q(p_2(x'_2))$. D'autre part, le foncteur
$$
q : \textit{Lift}(p_2(x_2), A' \to A) \to
\textit{Lift}(q(p_2(x_2)), A' \to A)
$$
induit une injection sur les classes d'isomorphie d'après le lemme
\ref{formal-defos-lemma-free-transitive-action} et notre hypothèse sur les espaces tangents.
Il existe donc un morphisme $b' : p_1(x_1') \to p_2(x'_2)$ dont l'image
directe sur $A$ est $b$. Il peut toutefois être nécessaire d'ajuster le choix
de $b'$ afin d'obtenir $q(b') = a'$. Il suffit pour cela de voir que
$q : \text{Inf}(p_2(x'_2)/p_2(x_2)) \to \text{Inf}(q(p_2(x'_2))/q(p_2(x_2)))$
est surjectif. Cela résulte de notre hypothèse sur les automorphismes
infinitésimaux et du lemme \ref{formal-defos-lemma-lifted-automorphisms-torsor}.
\end{proof}

\begin{lemma}
\label{formal-defos-lemma-easy-check-smooth}
Soit $f : \mathcal{F} \to \mathcal{G}$ un morphisme de catégories de
déformation. Soit $x_0 \in \Ob(\mathcal{F}(k))$, d'image
$y_0 \in \Ob(\mathcal{G}(k))$. Si
\begin{enumerate}
\item l'application $T\mathcal{F} \to T\mathcal{G}$ est surjective ;
\item pour toute petite extension $A' \to A$ de $\mathcal{C}_\Lambda$ et tout
$x \in \mathcal{F}(A)$ d'image $y \in \mathcal{G}(A)$, s'il existe un
relèvement de $y$ à $A'$, alors il existe un relèvement de $x$ à $A'$,
\end{enumerate}
alors $\mathcal{F} \to \mathcal{G}$ est lisse (et réciproquement).
\end{lemma}

\begin{proof}
Soit $A' \to A$ une petite extension. Soit $x \in \mathcal{F}(A)$. Soit
$y' \to f(x)$ un morphisme de $\mathcal{G}$ au-dessus de $A' \to A$.
Considérons le foncteur $\text{Lift}(A', x) \to \text{Lift}(A', f(x))$ induit
par $f$. Il faut montrer qu'il existe un objet $x' \to x$ de
$\text{Lift}(A', x)$ dont l'image est $y' \to f(x)$ ; voir le lemme
\ref{formal-defos-lemma-smoothness-small-extensions}. La condition (2) montre que
$\text{Lift}(A', x)$ n'est pas la catégorie vide. La condition (1) et le lemme
\ref{formal-defos-lemma-free-transitive-action} montrent alors que l'application sur les
classes d'isomorphie est surjective, comme voulu.
\end{proof}

\begin{lemma}
\label{formal-defos-lemma-map-between-smooth}
Soient $\mathcal{F} \to \mathcal{G} \to \mathcal{H}$ des morphismes de
catégories cofibrées en groupoïdes au-dessus de $\mathcal{C}_\Lambda$. Si
\begin{enumerate}
\item $\mathcal{F}$ et $\mathcal{G}$ sont des catégories de déformation ;
\item l'application $T\mathcal{F} \to T\mathcal{G}$ est surjective ;
\item $\mathcal{F} \to \mathcal{H}$ est lisse.
\end{enumerate}
Alors $\mathcal{F} \to \mathcal{G}$ est lisse.
\end{lemma}

\begin{proof}
Soit $A' \to A$ une petite extension de $\mathcal{C}_\Lambda$, et soit
$x \in \mathcal{F}(A)$ d'image $y \in \mathcal{G}(A)$. Supposons qu'il existe
un relèvement $y' \in \mathcal{G}(A')$. D'après le lemme
\ref{formal-defos-lemma-easy-check-smooth}, il suffit de vérifier que $x$ possède aussi un
relèvement. Prenons l'image $z' \in \mathcal{H}(A')$ de $y'$. Puisque
$\mathcal{F} \to \mathcal{H}$ est lisse, il existe
$x' \in \mathcal{F}(A')$ ayant pour images à la fois
$x \in \mathcal{F}(A)$ et $z' \in \mathcal{H}(A')$ ; voir la définition
\ref{formal-defos-definition-smooth-morphism}. Cela achève la démonstration.
\end{proof}





\section{Groupoïdes en foncteurs sur une catégorie arbitraire}
\label{formal-defos-section-groupoids-arbitrary}

\noindent
Commençons par des généralités sur les groupoïdes en foncteurs sur une
catégorie arbitraire. Dans la section suivante, nous passerons à la catégorie
$\mathcal{C}_\Lambda$. Par souci de clarté, nous appellerons parfois
« catégorie groupoïde » un groupoïde ordinaire, c'est-à-dire une catégorie
dont tous les morphismes sont des isomorphismes.

\begin{definition}
\label{formal-defos-definition-groupoid-in-functors}
Soit $\mathcal{C}$ une catégorie. La {\it catégorie des groupoïdes en
foncteurs sur $\mathcal{C}$} est la catégorie dont les objets et les morphismes
sont les suivants.
\begin{enumerate}
\item Objets : un {\it groupoïde en foncteurs sur $\mathcal{C}$} est un
quintuplet
$(U, R, s, t, c)$, où $U, R : \mathcal{C} \to \textit{Ensembles}$ sont des
foncteurs, et $s, t : R \to U$ et $c : R \times_{s, U, t} R \to R$ sont
des morphismes ayant la propriété suivante : pour tout objet $T$ de
$\mathcal{C}$, le quintuplet
$$
(U(T), R(T), s, t, c)
$$
est une catégorie groupoïde.
\item Morphismes : un {\it morphisme
$(U, R, s, t, c) \to (U', R', s', t', c')$ de groupoïdes en foncteurs sur
$\mathcal{C}$} consiste en des morphismes $U \to U'$ et $R \to R'$ ayant la
propriété suivante : pour tout objet $T$ de $\mathcal{C}$, les applications
induites $U(T) \to U'(T)$ et $R(T) \to R'(T)$ définissent un foncteur entre
catégories groupoïdes
$$
(U(T), R(T), s, t, c) \to (U'(T), R'(T), s', t', c').
$$
\end{enumerate}
\end{definition}

\begin{remark}
\label{formal-defos-remark-confusion-groupoids-in-functors}
Un groupoïde en foncteurs sur $\mathcal{C}$ équivaut à la donnée d'un foncteur
$\mathcal{C} \to \textit{Groupoïdes}$, et un morphisme de groupoïdes en
foncteurs sur $\mathcal{C}$ équivaut à un morphisme des foncteurs
correspondants $\mathcal{C} \to \textit{Groupoïdes}$ (où l'on considère
$\textit{Groupoïdes}$ comme une 1-catégorie). Pour nos
besoins, il est toutefois plus commode d'employer la terminologie des
groupoïdes en foncteurs. En effet, considérer un groupoïde en foncteurs comme
le foncteur correspondant $\mathcal{C} \to \textit{Groupoïdes}$, ou, de manière
équivalente, comme la catégorie cofibrée en groupoïdes associée à ce foncteur,
peut prêter à confusion (remarque
\ref{formal-defos-remark-smooth-groupoid-in-functors-warning}).
\end{remark}

\begin{remark}
\label{formal-defos-remark-identity-inverse}
Soit $(U, R, s, t, c)$ un groupoïde en foncteurs sur une catégorie
$\mathcal{C}$. Il existe d'uniques morphismes $e : U \to R$ et $i : R \to R$
tels que, pour tout objet $T$ de $\mathcal{C}$, $e: U(T) \to R(T)$ envoie
$x \in U(T)$ sur la flèche identique de $x$, et $i: R(T) \to R(T)$ envoie
$a \in U(T)$ sur l'inverse de $a$ dans la catégorie groupoïde
$(U(T), R(T), s, t, c)$. Nous appellerons parfois $s$, $t$, $c$, $e$ et $i$
la « source », le « but », la « composition », l'« identité » et l'« inverse ».
\end{remark}

\begin{definition}
\label{formal-defos-definition-representable}
Soit $\mathcal{C}$ une catégorie. Un groupoïde en foncteurs sur $\mathcal{C}$
est {\it représentable} s'il est isomorphe à un groupoïde de la forme
$(\underline{U}, \underline{R}, s, t, c)$, où $U$ et $R$ sont des objets de
$\mathcal{C}$ et où la somme amalgamée $R \amalg_{s, U, t} R$ existe.
\end{definition}

\begin{remark}
\label{formal-defos-remark-reason-existence-coproduct}
Ainsi, un groupoïde représentable en foncteurs sur $\mathcal{C}$ est donné par
des objets $U$ et $R$ de $\mathcal{C}$, ainsi que par des morphismes
$s, t : U \to R$ et $c : R \to R \amalg_{s, U, t} R$, tels que
$(\underline{U}, \underline{R}, s, t, c)$ satisfasse à la condition de la
définition \ref{formal-defos-definition-groupoid-in-functors}. On exige l'existence de la
somme amalgamée $R \amalg_{s, U, t} R$ afin que le morphisme de composition
$c$ soit défini au niveau des morphismes de $\mathcal{C}$. Cette exigence sera
toujours satisfaite plus bas lorsque nous considérerons des groupoïdes
représentables en foncteurs sur $\widehat{\mathcal{C}}_\Lambda$, car, d'après
le lemme \ref{formal-defos-lemma-CLambdahat-pushouts}, la catégorie
$\widehat{\mathcal{C}}_\Lambda$ admet des sommes amalgamées.
\end{remark}

\begin{remark}
\label{formal-defos-remark-simplify-terminology}
Nous dirons « {\it soit $(\underline{U}, \underline{R}, s, t, c)$ un
groupoïde en foncteurs sur $\mathcal{C}$} » pour signifier que nous disposons
d'un groupoïde représentable en foncteurs. Cela signifie donc que $U$ et $R$
sont des objets de $\mathcal{C}$, qu'il existe des morphismes
$s, t : U \to R$, que la somme amalgamée $R \amalg_{s, U, t} R$ existe, qu'il
existe un morphisme $c : R \to R \amalg_{s, U, t} R$, et que
$(\underline{U}, \underline{R}, s, t, c)$ est un groupoïde en foncteurs sur
$\mathcal{C}$.
\end{remark}

\noindent
Introduisons une notation pour la restriction des groupoïdes en foncteurs. Elle
sera utile plus bas lorsque nous restreindrons de $\widehat{\mathcal
C}_\Lambda$ à $\mathcal{C}_\Lambda$.

\begin{definition}
\label{formal-defos-definition-restricting-groupoids-in-functors}
Soit $(U, R, s, t, c)$ un groupoïde en foncteurs sur une catégorie
$\mathcal{C}$. Soit $\mathcal{C}'$ une sous-catégorie de $\mathcal{C}$. La
{\it restriction $(U, R, s, t, c)|_{\mathcal{C}'}$ de $(U, R, s, t, c)$ à
$\mathcal{C}'$} est le groupoïde en foncteurs sur $\mathcal{C}'$ donné par
$(U|_{\mathcal{C}'}, R|_{\mathcal
C'}, s|_{\mathcal{C}'}, t|_{\mathcal{C}'}, c|_{\mathcal{C}'})$.
\end{definition}

\begin{remark}
\label{formal-defos-remark-notation-restriction}
Dans la situation de la définition
\ref{formal-defos-definition-restricting-groupoids-in-functors}, nous notons souvent
$s|_{\mathcal{C}'}, t|_{\mathcal{C}'}, c|_{\mathcal{C}'}$ simplement $s, t, c$.
\end{remark}

\begin{definition}
\label{formal-defos-definition-quotient}
Soit $(U, R, s, t, c)$ un groupoïde en foncteurs sur une catégorie
$\mathcal{C}$.
\begin{enumerate}
\item La correspondance $T \mapsto  (U(T), R(T), s, t, c)$ définit un foncteur
$\mathcal{C} \to \textit{Groupoïdes}$. La {\it catégorie quotient cofibrée en
groupoïdes $[U/R] \to \mathcal{C}$} est la catégorie cofibrée en groupoïdes
au-dessus de $\mathcal{C}$ associée à ce foncteur (comme dans les remarques
\ref{formal-defos-remarks-cofibered-groupoids}
(\ref{formal-defos-item-construction-associated-cofibered-groupoid})).
\item Le {\it morphisme quotient $U \to [U/R]$} est le morphisme de catégories
cofibrées en groupoïdes au-dessus de $\mathcal{C}$ défini par les règles
suivantes :
\begin{enumerate}
\item $x \in U(T)$ est envoyé sur l'objet $(T, x) \in \Ob([U/R](T))$ ;
\item $x \in U(T)$ et $f : T \to T'$ donnent lieu au morphisme
$(f, \text{id}_{U(f)(x)}): (T, x) \to (T, U(f)(x))$ au-dessus de
$f : T \to T'$.
\end{enumerate}
\end{enumerate}
\end{definition}





\section{Groupoïdes en foncteurs sur la catégorie de base}
\label{formal-defos-section-prorepresentable-groupoids-in-functors}

\noindent
Dans cette section, nous étudions les groupoïdes en foncteurs sur
$\mathcal{C}_\Lambda$. Notre but ultime est de montrer que les groupoïdes en
foncteurs sur $\mathcal{C}_\Lambda$ qui sont pro-représentables fournissent des
« présentations » de catégories de déformation qui se comportent bien, de la
même manière que les groupoïdes lisses en espaces algébriques fournissent des
présentations de champs algébriques ; cf.\ Champs algébriques, section
\ref{algebraic-section-stack-to-presentation}.

\begin{definition}
\label{formal-defos-definition-prorepresentable-groupoid-in-functors}
Un groupoïde en foncteurs sur $\mathcal{C}_\Lambda$ est
{\it pro-représentable} s'il est isomorphe à
$(\underline{R_0}, \underline{R_1}, s, t, c)|_{\mathcal{C}_\Lambda}$
pour un certain groupoïde en foncteurs représentable
$(\underline{R_0}, \underline{R_1}, s, t, c)$ sur la catégorie
$\widehat{\mathcal{C}}_\Lambda$.
\end{definition}

\noindent
Soit $(U, R, s, t, c)$ un groupoïde en foncteurs sur
$\mathcal{C}_\Lambda$. En passant aux complétés, on obtient un quintuplet
$(\widehat{U}, \widehat{R}, \widehat{s}, \widehat{t}, \widehat{c})$. D'après
la remarque \ref{formal-defos-remark-completion-restriction-cofset-adjoint}, la complétion,
vue comme foncteur sur $\text{CofSet}(\mathcal{C}_\Lambda)$, est un adjoint à
droite ; elle commute donc aux limites. En particulier, il existe un
isomorphisme canonique
$$
\widehat{R \times_{s, U, t} R}
\longrightarrow
\widehat{R} \times_{\widehat{s}, \widehat{U}, \widehat{t}} \widehat{R},
$$
de sorte que $\widehat{c}$ peut être considéré comme un foncteur
$\widehat{R} \times_{\widehat{s}, \widehat{U}, \widehat{t}} \widehat{R} \to
\widehat{R}$. Alors
$(\widehat{U}, \widehat{R}, \widehat{s}, \widehat{t}, \widehat{c})$
est un groupoïde en foncteurs sur $\widehat{\mathcal{C}}_\Lambda$, dont les
morphismes identité et inverse sont les complétés de ceux de
$(U, R, s, t, c)$.

\begin{definition}
\label{formal-defos-definition-completion-groupoid-in-functors}
Soit $(U, R, s, t, c)$ un groupoïde en foncteurs sur
$\mathcal{C}_\Lambda$. Le {\it complété $(U, R, s, t, c)^{\wedge}$ de
$(U, R, s, t, c)$} est le groupoïde en foncteurs
$(\widehat{U}, \widehat{R}, \widehat{s}, \widehat{t}, \widehat{c})$ sur
$\widehat{\mathcal{C}}_\Lambda$ décrit ci-dessus.
\end{definition}

\begin{remark}
\label{formal-defos-remark-groupoid-in-functors-complete-restrict}
Soit $(U, R, s, t, c)$ un groupoïde en foncteurs sur
$\mathcal{C}_\Lambda$. Il existe alors un isomorphisme canonique
$(U, R, s, t, c)^{\wedge}|_{\mathcal{C}_\Lambda} \cong (U, R, s, t, c)$ ; voir
la remarque \ref{formal-defos-remark-restrict-completion}. D'autre part, soit
$(U, R, s, t, c)$ un groupoïde en foncteurs sur
$\widehat{\mathcal{C}}_\Lambda$ tel que
$U, R : \widehat{\mathcal{C}}_\Lambda \to \textit{Ensembles}$
commutent tous deux aux limites, par exemple si $U, R$ sont représentables.
Il existe alors un isomorphisme canonique
$((U, R, s, t, c)|_{\mathcal{C}_\Lambda})^{\wedge} \cong (U, R, s, t, c)$.
Cela résulte de la remarque \ref{formal-defos-remark-restrict-complete-continuous-functor}.
\end{remark}

\begin{lemma}
\label{formal-defos-lemma-groupoid-in-functors-prorep-equivalences}
Soit $(U, R, s, t, c)$ un groupoïde en foncteurs sur
$\mathcal{C}_\Lambda$.
\begin{enumerate}
\item $(U, R, s, t, c)$ est pro-représentable si et seulement si son complété
est représentable comme groupoïde en foncteurs sur
$\widehat{\mathcal{C}}_\Lambda$.
\item $(U, R, s, t, c)$ est pro-représentable si et seulement si $U$ et $R$
sont pro-représentables.
\end{enumerate}
\end{lemma}

\begin{proof}
La partie (1) résulte de la remarque
\ref{formal-defos-remark-groupoid-in-functors-complete-restrict}. Pour (2), l'implication
directe découle immédiatement de la définition d'un groupoïde en foncteurs
pro-représentable. Réciproquement, supposons $U$ et $R$
pro-représentables, et écrivons $U \cong \underline{R_0}|_{\mathcal{C}_\Lambda}$ et
$R \cong \underline{R_1}|_{\mathcal{C}_\Lambda}$ pour des objets $R_0$ et
$R_1$ de $\widehat{\mathcal{C}}_\Lambda$. Puisque
$\underline{R_0} \cong \widehat{\underline{R_0}|_{\mathcal{C}_\Lambda}}$ et
$\underline{R_1} \cong \widehat{\underline{R_1}|_{\mathcal{C}_\Lambda}}$
d'après la remarque \ref{formal-defos-remark-restrict-complete-continuous-functor}, on voit
que le complété $(U, R, s, t, c)^\wedge$ est un groupoïde en foncteurs de la
forme
$(\underline{R_0}, \underline{R_1}, \widehat{s}, \widehat{t}, \widehat{c})$.
Par le lemme \ref{formal-defos-lemma-CLambdahat-pushouts}, la somme amalgamée
$\underline{R_1} \times_{\widehat{s}, \underline{R_1}, \widehat{t}}
\underline{R_1}$ existe. Ainsi,
$(\underline{R_0}, \underline{R_1}, \widehat{s}, \widehat{t}, \widehat{c})$
est un groupoïde représentable en foncteurs sur
$\widehat{\mathcal{C}}_\Lambda$. Enfin, la restriction
$(\underline{R_0}, \underline{R_1}, s, t, c)|_{\mathcal{C}_\Lambda}$
redonne $(U, R, s, t, c)$ d'après la remarque
\ref{formal-defos-remark-groupoid-in-functors-complete-restrict} ; ainsi,
$(U, R, s, t, c)$ est pro-représentable par définition.
\end{proof}






\section{Groupoïdes en foncteurs lisses sur la catégorie de base}
\label{formal-defos-section-smooth-minimal-groupoids-in-functors}

\noindent
La notion de lissité des groupoïdes en foncteurs sur $\mathcal{C}_\Lambda$ est
définie comme suit.

\begin{definition}
\label{formal-defos-definition-smooth-groupoid-in-functors}
Soit $(U, R, s, t, c)$ un groupoïde en foncteurs sur
$\mathcal{C}_\Lambda$. On dit que $(U, R, s, t, c)$ est {\it lisse} si
$s, t: R \to U$ sont lisses.
\end{definition}

\begin{remark}
\label{formal-defos-remark-smooth-groupoid-in-functors-warning}
Notons que cette terminologie peut prêter à confusion : si
$(U, R, s, t, c)$ est un groupoïde en foncteurs lisse, alors le quotient
$[U/R]$ n'est pas nécessairement une catégorie cofibrée en groupoïdes lisse au
sens de la définition \ref{formal-defos-definition-cofibered-groupoid-projection-smooth}.
Cependant, la lissité de $(U, R, s, t, c)$ entraîne la lissité du morphisme
quotient $U \to [U/R]$ (et lui est en fait équivalente), comme nous le verrons
dans le lemme \ref{formal-defos-lemma-smooth-quotient-morphism}.
\end{remark}

\begin{remark}
\label{formal-defos-remark-smooth-power-series-prorepresentable-smooth-groupoid-in-functors}
Soit $(\underline{R_0}, \underline{R_1}, s, t, c)|_{\mathcal{C}_\Lambda}$ un
groupoïde en foncteurs pro-représentable sur $\mathcal{C}_\Lambda$. Alors
$(\underline{R_0}, \underline{R_1}, s, t, c)|_{\mathcal{C}_\Lambda}$ est
lisse si et seulement si $R_1$ est un anneau de séries formelles sur $R_0$
pour chacune des structures définies par $s$ et par $t$. Cela résulte du lemme
\ref{formal-defos-lemma-smooth-morphism-power-series}.
\end{remark}

\begin{lemma}
\label{formal-defos-lemma-smooth-quotient-morphism}
Soit $(U, R, s, t, c)$ un groupoïde en foncteurs sur
$\mathcal{C}_\Lambda$. Les assertions suivantes sont équivalentes :
\begin{enumerate}
\item Le groupoïde en foncteurs $(U, R, s, t, c)$ est lisse.
\item Le morphisme $s : R \to U$ est lisse.
\item Le morphisme $t : R \to U$ est lisse.
\item Le morphisme quotient $U \to [U/R]$ est lisse.
\end{enumerate}
\end{lemma}

\begin{proof}
L'assertion (2) est équivalente à (3), puisque l'inverse $i: R \to R$ de
$(U, R, s, t, c)$ est un isomorphisme et que $t = s \circ i$. Par définition,
(1) équivaut à la conjonction de (2) et (3), donc à chacune d'elles séparément.

\medskip\noindent
Montrons enfin que (2) est équivalente à (4). En explicitant les définitions :
\begin{enumerate}
\item[(2)] La lissité de $s: R \to U$ équivaut à la condition suivante : si
$f: B \to A$ est un morphisme d'anneaux surjectif dans
$\mathcal{C}_\Lambda$, si $a \in R(A)$ et si $y \in U(B)$ vérifie
$s(a) = U(f)(y)$, alors il existe $a' \in R(B)$ tel que
$R(f)(a') = a$ et $s(a') = y$.
\item[(4)] La lissité de $U \to [U/R]$ équivaut à la condition suivante : si
$f: B \to A$ est un morphisme d'anneaux surjectif dans
$\mathcal{C}_\Lambda$ et si $(f, a) : (B, y) \to (A, x)$ est un morphisme de
$[U/R]$, alors il existe $x' \in U(B)$ et $b \in R(B)$, avec
$s(b) = x'$, $t(b) = y$, tels que $c(a, R(f)(b)) = e(x)$. Ici,
$e : U \to R$ désigne l'identité, et la notation $(f, a)$ est celle des
remarques \ref{formal-defos-remarks-cofibered-groupoids}
(\ref{formal-defos-item-construction-associated-cofibered-groupoid}) ; en particulier,
$a \in R(A)$, avec $s(a) = U(f)(y)$ et $t(a) = x$.
\end{enumerate}
Supposons (4) vérifiée et soient $f, a, y$ comme en (2). Posons $x = t(a)$,
de sorte que l'on ait un morphisme $(f, a): (B, y) \to (A, x)$. Alors (4)
fournit $x'$ et $b$, et $a' = i(b)$ satisfait aux exigences de (2).
Réciproquement, supposons (2) vérifiée et soit
$(f, a): (B, y) \to (A, x)$ comme en (4). Alors (2) fournit
$a' \in R(B)$, et $x' = t(a')$ et $b = i(a')$ satisfont aux exigences de (4).
\end{proof}





\section{Catégories de déformation comme quotients de groupoïdes en foncteurs}
\label{formal-defos-section-deformation-categories-as-quotients}

\noindent
Nous étudions des conditions sur un groupoïde en foncteurs sur
$\mathcal{C}_\Lambda$ qui garantissent que son quotient est une catégorie de
déformation, et nous calculons l'espace tangent et l'espace des automorphismes
infinitésimaux d'un tel quotient.

\begin{lemma}
\label{formal-defos-lemma-smooth-RS-groupoid-in-functors-quotient}
Soit $(U, R, s, t, c)$ un groupoïde en foncteurs lisse sur
$\mathcal{C}_\Lambda$. Supposons que $U$ et $R$ satisfassent à (RS). Alors
$[U/R]$ satisfait à (RS).
\end{lemma}

\begin{proof}
Soit
$$
\xymatrix{
                           &     (A_2, x_2) \ar[d]^{(f_2, a_2)} \\
(A_1, x_1) \ar[r]^{(f_1, a_1)} &     (A, x)
}
$$
un diagramme de $[U/R]$ tel que $f_2: A_2 \to A$ soit surjectif. La notation
est celle des remarques \ref{formal-defos-remarks-cofibered-groupoids}
(\ref{formal-defos-item-construction-associated-cofibered-groupoid}). Ainsi,
$f_1: A_1 \to A, f_2: A_2 \to A$ sont des morphismes de
$\mathcal{C}_\Lambda$, $x \in U(A)$, $x_1 \in U(A_1)$, $x_2 \in U(A_2)$, et
$a_1, a_2 \in R(A)$, avec $s(a_1) = U(f_1)(x_1)$, $t(a_1) = x$ et
$s(a_2) = U(f_2)(x_2)$, $t(a_2) = x$. Construisons comme suit un produit fibré
au-dessus de $A_1 \times_A A_2$ pour ce diagramme de $[U/R]$.

\medskip\noindent
Soit $a = c(i(a_1), a_2)$, où $i: R \to R$ est le morphisme inverse. Alors
$a \in R(A)$, $x_2 \in U(A_2)$ et $s(a) = U(f_2)(x_2)$. On a donc un élément
$(a, x_2) \in R(A) \times_{s, U(A), U(f_2)} U(A_2)$. Par la lissité de
$s : R \to U$, il existe un élément $\widetilde{a} \in R(A_2)$ tel que
$R(f_2)(\widetilde{a}) = a$ et $s(\widetilde{a}) = x_2$. En particulier,
$U(f_2)(t(\widetilde{a})) = t(a) = U(f_1)(x_1)$. Ainsi, $x_1$ et
$t(\widetilde{a})$ définissent un élément
$$
(x_1, t(\widetilde{a})) \in U(A_1) \times_{U(A)} U(A_2).
$$
Puisque $U$ satisfait à (RS), on a une identification
$U(A_1) \times_{U(A)} U(A_2) = U(A_1 \times_A A_2)$. Notons
$x_1 \times t(\widetilde{a}) \in U(A_1 \times_A A_2)$ l'élément correspondant
à $(x_1, t(\widetilde{a})) \in U(A_1) \times_{U(A)} U(A_2)$. Soient
$p_1, p_2$ les projections de $A_1 \times_A A_2$. Nous affirmons que
$$
\xymatrix{
(A_1 \times_A A_2, x_1 \times t(\widetilde{a})) \ar[d]_{(p_1, e(x_1))}
\ar[rr]_-{(p_2, i(\widetilde{a}))} & & (A_2, x_2) \ar[d]^{(f_2, a_2)} \\
(A_1, x_1) \ar[rr]^{(f_1, a_1)} & & (A, x)
}
$$
est un carré cartésien dans $[U/R]$. (Notons que $e: U \to R$ désigne
l'identité.)

\medskip\noindent
Le diagramme est commutatif parce que
$c(a_2, R(f_2)(i(\widetilde{a}))) = c(a_2, i(a)) = a_1$.
Pour vérifier qu'il est cartésien, soit
$$
\xymatrix{
(B, z) \ar[d]_{(g_1, b_1)} \ar[rr]_{(g_2, b_2)} & & (A_2, x_2)
\ar[d]^{(f_2, a_2)} \\
(A_1, x_1) \ar[rr]^{(f_1, a_1)} & & (A, x)
}
$$
un diagramme commutatif de $[U/R]$. Nous allons montrer qu'il existe un unique
morphisme
$(g, b) : (B, z) \to (A_1 \times_A A_2, x_1 \times t(\widetilde{a}))$
compatible aux morphismes vers $(A_1, x_1)$ et $(A_2, x_2)$. Nécessairement,
$g = (g_1, g_2) : B \to A_1 \times_A A_2$. Puisque, par hypothèse, $R$
satisfait à (RS), on a une identification
$R(A_1 \times_A A_2) = R(A_1) \times_{R(A)} R(A_2)$. On peut donc écrire
$b = (b'_1, b'_2)$ pour certains $b'_1 \in R(A_1)$, $b'_2 \in R(A_2)$ ayant
la même image dans $R(A)$. Alors
$((g_1, g_2), (b'_1, b'_2)) : (B, z) \to
(A_1 \times_A A_2, x_1 \times t(\widetilde{a}))$
est compatible avec les projections si et seulement si $b'_1 = b_1$ et
$b'_2 = c(\widetilde{a}, b_2)$, ce qui démontre l'existence et l'unicité.
\end{proof}

\begin{lemma}
\label{formal-defos-lemma-deformation-groupoid-quotient}
Soit $(U, R, s, t, c)$ un groupoïde en foncteurs lisse sur
$\mathcal{C}_\Lambda$. Supposons que $U$ et $R$ soient des foncteurs de
déformation. Alors :
\begin{enumerate}
\item Le quotient $[U/R]$ est une catégorie de déformation.
\item L'espace tangent de $[U/R]$ est
$$
T[U/R] = \Coker(ds-dt: TR \to TU).
$$
\item L'espace des automorphismes infinitésimaux de $[U/R]$ est
$$
\text{Inf}([U/R]) =
\Ker(ds \oplus dt : TR \to TU \oplus TU).
$$
\end{enumerate}
\end{lemma}

\begin{proof}
Puisque $U$ et $R$ sont des foncteurs de déformation, $[U/R]$ est une
catégorie de prédéformation. Comme (RS) est vérifiée par définition pour les
foncteurs de déformation, le lemme
\ref{formal-defos-lemma-smooth-RS-groupoid-in-functors-quotient} montre que (RS) est
vérifiée pour [U/R]. Ainsi, $[U/R]$ est une catégorie de déformation. Les
assertions (2) et (3) résultent directement des définitions.
\end{proof}







\section{Présentations des catégories cofibrées en groupoïdes}
\label{formal-defos-section-presentation-categories-cofibred-in-groupoids}

\noindent
Une présentation se définit comme suit.

\begin{definition}
\label{formal-defos-definition-presentation}
Soit $\mathcal{F}$ une catégorie cofibrée en groupoïdes au-dessus d'une
catégorie $\mathcal{C}$. Soit $(U, R, s, t, c)$ un groupoïde en foncteurs sur
$\mathcal{C}$. Une {\it présentation de $\mathcal{F}$ par
$(U, R, s, t, c)$} est une équivalence
$\varphi : [U/R] \to \mathcal{F}$ de catégories cofibrées en groupoïdes
au-dessus de $\mathcal{C}$.
\end{definition}

\noindent
Les deux lemmes généraux suivants serviront à construire des présentations.

\begin{lemma}
\label{formal-defos-lemma-presentation-construction}
Soit $\mathcal{F}$ une catégorie cofibrée en groupoïdes au-dessus d'une
catégorie $\mathcal{C}$. Soit $U : \mathcal{C} \to \textit{Ensembles}$ un
foncteur. Soit $f : U \to \mathcal{F}$ un morphisme de catégories cofibrées
en groupoïdes au-dessus de $\mathcal{C}$. Définissons $R, s, t, c$ comme suit :
\begin{enumerate}
\item $R : \mathcal{C} \to \textit{Ensembles}$ est le foncteur
$U \times_{f, \mathcal{F}, f} U$.
\item $t, s : R \to U$ sont respectivement la première et la seconde
projection.
\item $c : R \times_{s, U, t} R \to R$ est le morphisme donné par la
projection sur le premier et le dernier facteur de
$U \times_{f, \mathcal{F}, f} U \times_{f, \mathcal{F}, f} U$
sous l'isomorphisme canonique
$R \times_{s, U, t} R \to
U \times_{f, \mathcal{F}, f} U \times_{f, \mathcal{F}, f} U$.
\end{enumerate}
Alors $(U, R, s, t, c)$ est un groupoïde en foncteurs sur $\mathcal{C}$.
\end{lemma}

\begin{proof}
La démonstration est omise.
\end{proof}

\begin{lemma}
\label{formal-defos-lemma-presentation-morphism}
Soit $\mathcal{F}$ une catégorie cofibrée en groupoïdes au-dessus d'une
catégorie $\mathcal{C}$. Soit $U : \mathcal{C} \to \textit{Ensembles}$ un
foncteur. Soit $f : U \to \mathcal{F}$ un morphisme de catégories cofibrées
en groupoïdes au-dessus de $\mathcal{C}$. Soit $(U, R, s, t, c)$ le groupoïde
en foncteurs sur $\mathcal{C}$ construit à partir de
$f : U \to \mathcal{F}$ dans le lemme
\ref{formal-defos-lemma-presentation-construction}. Il existe alors un morphisme naturel
$[f] : [U/R] \to \mathcal{F}$ tel que :
\begin{enumerate}
\item $[f]: [U/R] \to \mathcal{F}$ est pleinement fidèle.
\item $[f]: [U/R] \to \mathcal{F}$ est une équivalence si et seulement si
$f : U \to \mathcal{F}$ est essentiellement surjectif.
\end{enumerate}
\end{lemma}

\begin{proof}
La démonstration est omise.
\end{proof}







\section{Présentations des catégories de déformation}
\label{formal-defos-section-presentation-deformation-categories}

\noindent
D'après le lemme suivant, un morphisme lisse d'un foncteur de prédéformation
vers une catégorie de prédéformation $\mathcal{F}$ donne lieu à une
présentation de $\mathcal{F}$ par un groupoïde en foncteurs lisse.

\begin{lemma}
\label{formal-defos-lemma-smooth-groupoid-in-functors-construction}
Soit $\mathcal{F}$ une catégorie cofibrée en groupoïdes au-dessus de
$\mathcal{C}_\Lambda$. Soit $U : \mathcal{C}_\Lambda \to \textit{Ensembles}$ un
foncteur. Soit $f : U \to \mathcal{F}$ un morphisme lisse de catégories
cofibrées en groupoïdes. Alors :
\begin{enumerate}
\item Si $(U, R, s, t, c)$ est le groupoïde en foncteurs sur
$\mathcal{C}_\Lambda$ construit à partir de $f : U \to \mathcal{F}$ dans le
lemme \ref{formal-defos-lemma-presentation-construction}, alors $(U, R, s, t, c)$ est
lisse.
\item Si $f : U(k) \to \mathcal{F}(k)$ est essentiellement surjectif, alors le
morphisme $[f] : [U/R] \to \mathcal{F}$ du lemme
\ref{formal-defos-lemma-presentation-morphism} est une équivalence.
\end{enumerate}
\end{lemma}

\begin{proof}
La construction du lemme \ref{formal-defos-lemma-presentation-construction} donne un
diagramme commutatif
$$
\xymatrix{
R = U \times_{f, \mathcal{F}, f} U \ar[r]_-s \ar[d]_t & U
\ar[d]^f \\
U \ar[r]^f & \mathcal{F}
}
$$
où $t, s$ sont respectivement la première et la seconde projection. Les morphismes $t, s$
sont donc lisses d'après le lemme \ref{formal-defos-lemma-smooth-properties}. Ainsi, (1)
est vérifiée.

\medskip\noindent
Si l'hypothèse de (2) est vérifiée, le lemme
\ref{formal-defos-lemma-smooth-morphism-essentially-surjective} montre que le morphisme
$f : U \to \mathcal{F}$ est essentiellement surjectif. Le lemme
\ref{formal-defos-lemma-presentation-morphism} montre donc que le morphisme
$[f] : [U/R] \to \mathcal{F}$ est une équivalence.
\end{proof}

\begin{lemma}
\label{formal-defos-lemma-deformation-functor-diagonal}
Soit $\mathcal{F}$ une catégorie de déformation. Soit
$U : \mathcal{C}_\Lambda \to \textit{Ensembles}$ un foncteur de déformation. Soit
$f: U \to \mathcal{F}$ un morphisme de catégories cofibrées en groupoïdes.
Alors $U \times_{f, \mathcal{F}, f} U$ est un foncteur de déformation dont
l'espace tangent s'insère dans une suite exacte d'espaces vectoriels sur $k$
$$
0 \to \text{Inf}(\mathcal{F}) \to
T(U \times_{f, \mathcal{F}, f} U) \to TU \oplus TU
$$
\end{lemma}

\begin{proof}
Cela résulte du lemme
\ref{formal-defos-lemma-deformation-categories-fiber-product-morphisms} et du fait que
$\text{Inf}(U) = (0)$.
\end{proof}

\begin{lemma}
\label{formal-defos-lemma-prorepresentable-groupoid-in-functors-construction}
Soit $\mathcal{F}$ une catégorie de déformation. Soit
$U : \mathcal{C}_\Lambda \to \textit{Ensembles}$ un foncteur pro-représentable.
Soit $f : U \to \mathcal{F}$ un morphisme de catégories cofibrées en
groupoïdes. Soit $(U, R, s, t, c)$ le groupoïde en foncteurs sur
$\mathcal{C}_\Lambda$ construit à partir de $f : U \to \mathcal{F}$ dans le
lemme \ref{formal-defos-lemma-presentation-construction}. Si
$\dim_k \text{Inf}(\mathcal{F}) < \infty$, alors $(U, R, s, t, c)$ est
pro-représentable.
\end{lemma}

\begin{proof}
Notons que $U$ est un foncteur de déformation d'après l'exemple
\ref{formal-defos-example-prorepresentable-deformation-functor}. Le lemme
\ref{formal-defos-lemma-deformation-functor-diagonal} montre que
$R = U \times_{f, \mathcal{F}, f} U$ est un foncteur de déformation dont
l'espace tangent $TR = T(U \times_{f, \mathcal{F}, f} U)$ s'insère dans une
suite exacte $0 \to \text{Inf}(\mathcal{F}) \to TR \to TU \oplus TU$. Le
premier espace est de dimension finie par hypothèse, et $TU$ est de dimension
finie d'après l'exemple
\ref{formal-defos-example-tangent-space-prorepresentable-functor} ; ainsi,
$\dim TR < \infty$. L'application
$\gamma : \text{Der}_\Lambda(k, k) \to TR$ (voir (\ref{formal-defos-equation-map})) est
injective, car son composé avec $TR \to TU$ est injectif, d'après le théorème
\ref{formal-defos-theorem-Schlessinger-prorepresentability} appliqué au foncteur
pro-représentable $U$. Le théorème
\ref{formal-defos-theorem-Schlessinger-prorepresentability} montre donc que $R$ est
pro-représentable. Il résulte alors du lemme
\ref{formal-defos-lemma-groupoid-in-functors-prorep-equivalences} que
$(U, R, s, t, c)$ est pro-représentable.
\end{proof}

\begin{theorem}
\label{formal-defos-theorem-presentation-deformation-groupoid}
Soit $\mathcal{F}$ une catégorie cofibrée en groupoïdes au-dessus de
$\mathcal{C}_\Lambda$. Alors $\mathcal{F}$ admet une présentation par un
groupoïde en foncteurs lisse et pro-représentable sur $\mathcal{C}_\Lambda$ si et
seulement si les conditions suivantes sont vérifiées :
\begin{enumerate}
\item $\mathcal{F}$ est une catégorie de déformation.
\item $\dim_k T\mathcal{F}$ est finie.
\item $\dim_k \text{Inf}(\mathcal{F})$ est finie.
\end{enumerate}
\end{theorem}

\begin{proof}
Rappelons qu'un foncteur pro-représentable est un foncteur de déformation ;
voir l'exemple \ref{formal-defos-example-prorepresentable-deformation-functor}. Ainsi, si
$\mathcal{F}$ est équivalente à un groupoïde en foncteurs lisse et
pro-représentable, les conditions (1), (2) et (3) résultent des assertions
(1), (2) et (3) du lemme \ref{formal-defos-lemma-deformation-groupoid-quotient}.

\medskip\noindent
Réciproquement, supposons les conditions (1), (2) et (3) vérifiées. La
condition (1) entraîne que (S1) et (S2) sont vérifiées ; voir le lemme
\ref{formal-defos-lemma-RS-implies-S1-S2}. D'après le lemme
\ref{formal-defos-lemma-versal-object-existence}, il existe un objet formel versel $\xi$.
En posant $U = \underline{R}|_{\mathcal{C}_\Lambda}$, le morphisme associé
$\underline{\xi} : U \to \mathcal{F}$ est lisse (c'est la définition d'un
objet formel versel). Soit $(U, R, s, t, c)$ le groupoïde en foncteurs construit
dans le lemme \ref{formal-defos-lemma-presentation-construction} à partir du morphisme
$\underline{\xi}$. Le lemme
\ref{formal-defos-lemma-smooth-groupoid-in-functors-construction} montre que
$(U, R, s, t, c)$ est un groupoïde en foncteurs lisse et que
$[U/R] \to \mathcal{F}$ est une équivalence. Le lemme
\ref{formal-defos-lemma-prorepresentable-groupoid-in-functors-construction} montre que
$(U, R, s, t, c)$ est pro-représentable. Ainsi,
$[U/R] \to \mathcal{F}$ est la présentation cherchée de $\mathcal{F}$.
\end{proof}








\section{Remarques sur la minimalité}
\label{formal-defos-section-minimality}

\noindent
Le théorème principal de ce chapitre est le théorème
\ref{formal-defos-theorem-presentation-deformation-groupoid} ci-dessus. Il décrit
complètement les catégories cofibrées en groupoïdes au-dessus de
$\mathcal{C}_\Lambda$ qui admettent une présentation par un groupoïde en foncteurs
lisse et pro-représentable. Dans cette section, nous expliquons brièvement
comment la minimalité étudiée dans les sections \ref{formal-defos-section-minimal-versal} et
\ref{formal-defos-section-miniversal-objects-existence} permet d'obtenir une présentation
lisse et pro-représentable « minimale ».

\begin{definition}
\label{formal-defos-definition-minimal-groupoid-in-functors}
Soit $(U, R, s, t, c)$ un groupoïde en foncteurs lisse et pro-représentable sur
$\mathcal{C}_\Lambda$.
\begin{enumerate}
\item On dit que $(U, R, s, t, c)$ est {\it normalisé} si le groupoïde
$(U(k[\epsilon]), R(k[\epsilon]), s, t, c)$ est totalement discontinu,
c'est-à-dire s'il n'existe aucun morphisme entre deux objets distincts.
\item On dit que $(U, R, s, t, c)$ est {\it minimal} si $U \to [U/R]$ est
donné par un objet formel versel minimal de $[U/R]$.
\end{enumerate}
\end{definition}

\noindent
La différence entre les deux notions est liée à celle entre les conditions
(\ref{formal-defos-equation-bijective}) et (\ref{formal-defos-equation-bijective-orbits}), et elle
disparaît lorsque $k' \subset k$ est séparable. En outre, un groupoïde en foncteurs
lisse, pro-représentable et normalisé est minimal, comme le montre le lemme
suivant. Voici un énoncé précis.

\begin{lemma}
\label{formal-defos-lemma-characterize-minimal-groupoid-in-functors}
Soit $(U, R, s, t, c)$ un groupoïde en foncteurs lisse et pro-représentable sur
$\mathcal{C}_\Lambda$.
\begin{enumerate}
\item $(U, R, s, t, c)$ est normalisé si et seulement si le morphisme
$U \to [U/R]$ induit un isomorphisme sur les espaces tangents ;
\item $(U, R, s, t, c)$ est minimal si et seulement si le noyau de
$TU \to T[U/R]$ est contenu dans l'image de
$\text{Der}_\Lambda(k, k) \to TU$.
\end{enumerate}
\end{lemma}

\begin{proof}
La partie (1) résulte immédiatement des définitions. Pour voir (2), posons
$\mathcal{F} = [U/R]$. Puisque $\mathcal{F}$ possède une présentation, c'est
une catégorie de déformation ; voir le théorème
\ref{formal-defos-theorem-presentation-deformation-groupoid}. Elle satisfait en particulier
à (RS), (S1) et (S2) ; voir le lemme \ref{formal-defos-lemma-RS-implies-S1-S2}. Rappelons
que les objets formels versels minimaux sont uniques à isomorphisme près ; voir
le lemme \ref{formal-defos-lemma-minimal-versal}. D'après le théorème
\ref{formal-defos-theorem-miniversal-object-existence}, un objet versel minimal induit un
morphisme
$\underline{\xi} : \underline{R}|_{\mathcal{C}_\Lambda} \to \mathcal{F}$
qui satisfait à (\ref{formal-defos-equation-bijective-orbits}). Puisque
$U \cong \underline{R}|_{\mathcal{C}_\Lambda}$ au-dessus de $\mathcal{F}$, on
voit que $TU \to T\mathcal{F} = T[U/R]$ satisfait à la propriété énoncée dans
le lemme.
\end{proof}

\noindent
Le quotient d'un groupoïde en foncteurs minimal et pro-représentable sur $\mathcal
C_\Lambda$ n'admet pas d'autoéquivalence qui ne soit un automorphisme. Pour le
démontrer, commençons par le lemme suivant.

\begin{lemma}
\label{formal-defos-lemma-surjective-morphism-prorepresentable-functor}
Soit $U: \mathcal{C}_\Lambda \to \textit{Ensembles}$ un foncteur
pro-représentable. Soit $\varphi : U \to U$ un morphisme tel que
$d\varphi : TU \to TU$ soit un isomorphisme. Alors $\varphi$ est un
isomorphisme.
\end{lemma}

\begin{proof}
Si $U \cong \underline{R}|_{\mathcal{C}_\Lambda}$ pour un certain
$R \in \Ob(\widehat{\mathcal{C}}_\Lambda)$, la complétion de $\varphi$ donne un
morphisme $\underline{R} \to \underline{R}$. Si $f: R \to R$ est le morphisme
correspondant dans $\widehat{\mathcal{C}}_\Lambda$, alors $f$ induit un
isomorphisme $\text{Der}_\Lambda(R, k) \to \text{Der}_\Lambda(R, k)$ ; voir
l'exemple \ref{formal-defos-example-tangent-space-map-prorepresentable-functor}. En
particulier, $f$ est surjectif d'après le lemme
\ref{formal-defos-lemma-derivations-surjective}. Or un endomorphisme surjectif d'un anneau
noethérien est un isomorphisme (voir Algèbre, lemme
\ref{algebra-lemma-surjective-endo-noetherian-ring-is-iso}) ; on conclut que
$f$, puis $\underline{R}
\to \underline{R}$, puis $\varphi : U \to U$ sont des isomorphismes.
\end{proof}

\begin{lemma}
\label{formal-defos-lemma-minimal-prorepresentable-groupoid-autoequivalence}
Soit $(U, R, s, t, c)$ un groupoïde en foncteurs minimal, lisse et
pro-représentable sur $\mathcal{C}_\Lambda$. Si $\varphi : [U/R] \to [U/R]$ est une
équivalence de catégories cofibrées en groupoïdes, alors $\varphi$ est un
isomorphisme.
\end{lemma}

\begin{proof}
Un morphisme $\varphi : [U/R] \to [U/R]$ équivaut à un morphisme
$\varphi : (U, R, s, t, c) \to (U, R, s, t, c)$ de groupoïdes en foncteurs
au-dessus de $\mathcal{C}_\Lambda$, au sens de la définition
\ref{formal-defos-definition-groupoid-in-functors}. Notons $\phi : U \to U$ et
$\psi : R \to R$ les morphismes correspondants. Puisque le diagramme
$$
\xymatrix{
& \text{Der}_\Lambda(k, k) \ar[dr]_\gamma \ar[dl]^\gamma \\
TU \ar[rr]_{d\phi} \ar[d] & & TU \ar[d]  \\
T[U/R] \ar[rr]^{d\varphi} & & T[U/R]
}
$$
est commutatif, que $d\varphi$ est bijectif, et que nous disposons de la
caractérisation de la minimalité du lemme
\ref{formal-defos-lemma-characterize-minimal-groupoid-in-functors}, on conclut que $d\phi$
est injectif (et donc bijectif pour des raisons de dimension). Par le lemme
\ref{formal-defos-lemma-surjective-morphism-prorepresentable-functor},
$\phi : U \to U$ est donc un isomorphisme. Un argument analogue, faisant
intervenir la suite exacte
$$
0 \to \text{Inf}([U/R]) \to TR \to TU \oplus TU
$$
du lemme \ref{formal-defos-lemma-deformation-functor-diagonal}, montre que
$\psi : R \to R$ est un isomorphisme. Cela résulte aussi du fait que
$R = U \times_{[U/R]} U$ et que $\varphi$ et $\phi$ sont des isomorphismes.
\end{proof}

\begin{lemma}
\label{formal-defos-lemma-minimal-prorepresentable-groupoid-equivalence}
Soient $(U, R, s, t, c)$ et $(U', R', s', t', c')$ des groupoïdes en foncteurs
minimaux, lisses et pro-représentables sur $\mathcal{C}_\Lambda$. Si
$\varphi : [U/R] \to [U'/R']$ est une équivalence de catégories cofibrées en
groupoïdes, alors $\varphi$ est un isomorphisme.
\end{lemma}

\begin{proof}
Soit $\psi : [U'/R'] \to [U/R]$ un quasi-inverse de $\varphi$. Alors
$\psi \circ \varphi$ et $\varphi \circ \psi$ sont des isomorphismes d'après
le lemme \ref{formal-defos-lemma-minimal-prorepresentable-groupoid-autoequivalence} ; ainsi,
$\varphi$ et $\psi$ sont des isomorphismes.
\end{proof}

\noindent
Le lemme suivant résume une partie de ce que nous avons vu plus haut dans ce
chapitre.

\begin{lemma}
\label{formal-defos-lemma-minimal-groupoid-in-functors-construction}
Soit $\mathcal{F}$ une catégorie de déformation telle que
$\dim_k T\mathcal{F} <\infty$ et
$\dim_k \text{Inf}(\mathcal{F}) < \infty$. Il existe alors un objet formel
versel minimal $\xi$ de $\mathcal{F}$. Supposons $\xi$ au-dessus de
$R \in \Ob(\widehat{\mathcal{C}}_\Lambda)$. Soit
$U = \underline{R}|_{\mathcal{C}_\Lambda}$. Soit
$f = \underline{\xi} : U \to \mathcal{F}$ le morphisme associé. Soit
$(U, R, s, t, c)$ le groupoïde en foncteurs sur $\mathcal{C}_\Lambda$
construit à partir de $f : U \to \mathcal{F}$ dans le lemme
\ref{formal-defos-lemma-presentation-construction}. Alors $(U, R, s, t, c)$ est un
groupoïde en foncteurs minimal, lisse et pro-représentable sur
$\mathcal{C}_\Lambda$, et il existe une équivalence
$[U/R] \to \mathcal{F}$.
\end{lemma}

\begin{proof}
Comme $\mathcal{F}$ est une catégorie de déformation, elle satisfait à (S1)
et (S2) ; voir le lemme \ref{formal-defos-lemma-RS-implies-S1-S2}. D'après le lemme
\ref{formal-defos-lemma-versal-object-existence}, il existe un objet formel versel. D'après
le lemme \ref{formal-defos-lemma-minimal-versal}, il existe un objet formel versel minimal
$\xi/R$ comme dans l'énoncé. En posant
$U = \underline{R}|_{\mathcal{C}_\Lambda}$, le morphisme associé
$\underline{\xi} : U \to \mathcal{F}$ est lisse (c'est la définition d'un
objet formel versel). Soit $(U, R, s, t, c)$ le groupoïde en foncteurs construit
dans le lemme \ref{formal-defos-lemma-presentation-construction} à partir du morphisme
$\underline{\xi}$. D'après le lemme
\ref{formal-defos-lemma-smooth-groupoid-in-functors-construction},
$(U, R, s, t, c)$ est un groupoïde en foncteurs lisse et
$[U/R] \to \mathcal{F}$ est une équivalence. D'après le lemme
\ref{formal-defos-lemma-prorepresentable-groupoid-in-functors-construction},
$(U, R, s, t, c)$ est pro-représentable. Enfin, $(U, R, s, t, c)$ est minimal,
car $U \to [U/R] = \mathcal{F}$ correspond à l'objet formel versel minimal
$\xi$.
\end{proof}

\noindent
Les présentations par des groupoïdes en foncteurs minimaux et pro-représentables
satisfont à la propriété d'unicité suivante.

\begin{lemma}
\label{formal-defos-lemma-minimal-presentations-equivalent}
Soit $\mathcal{F}$ une catégorie cofibrée en groupoïdes au-dessus de
$\mathcal{C}_\Lambda$. Supposons qu'il existe des présentations de
$\mathcal{F}$ par des groupoïdes en foncteurs minimaux, lisses et
pro-représentables $(U, R, s, t, c)$ et $(U', R', s', t', c')$. Alors
$(U, R, s, t, c)$ et $(U', R', s', t', c')$ sont isomorphes.
\end{lemma}

\begin{proof}
Cela résulte du lemme
\ref{formal-defos-lemma-minimal-prorepresentable-groupoid-equivalence} et de l'observation
qu'un morphisme $[U/R] \to [U'/R']$ équivaut à un morphisme de groupoïdes en
foncteurs (par notre construction explicite de $[U/R]$ dans la définition
\ref{formal-defos-definition-quotient}).
\end{proof}

\noindent
En résumé, nous avons démontré le théorème suivant.

\begin{theorem}
\label{formal-defos-theorem-minimal-smooth-prorepresentable-presentations}
Soit $\mathcal{F}$ une catégorie cofibrée en groupoïdes au-dessus de
$\mathcal{C}_\Lambda$. Considérons les conditions suivantes :
\begin{enumerate}
\item $\mathcal{F}$ admet une présentation par un groupoïde en foncteurs normalisé,
lisse et pro-représentable sur $\mathcal{C}_\Lambda$ ;
\item $\mathcal{F}$ admet une présentation par un groupoïde en foncteurs lisse
et pro-représentable sur $\mathcal{C}_\Lambda$ ;
\item $\mathcal{F}$ admet une présentation par un groupoïde en foncteurs minimal,
lisse et pro-représentable sur $\mathcal{C}_\Lambda$ ;
\item $\mathcal{F}$ satisfait aux conditions suivantes :
\begin{enumerate}
\item $\mathcal{F}$ est une catégorie de déformation.
\item $\dim_k T\mathcal{F}$ est finie.
\item $\dim_k \text{Inf}(\mathcal{F})$ est finie.
\end{enumerate}
\end{enumerate}
Alors (2), (3) et (4) sont équivalentes et sont entraînées par (1). Si
$k' \subset k$ est séparable, (1), (2), (3) et (4) sont toutes équivalentes.
De plus, les groupoïdes en foncteurs minimaux, lisses et pro-représentables qui
fournissent une présentation de $\mathcal{F}$ sont uniques à isomorphisme près.
\end{theorem}

\begin{proof}
Le lemme \ref{formal-defos-lemma-characterize-minimal-groupoid-in-functors} montre que (1)
entraîne (3), et qu'elle lui est équivalente si $k' \subset k$ est séparable.
Il est clair que (3) entraîne (2). Le théorème
\ref{formal-defos-theorem-presentation-deformation-groupoid} montre que (2) entraîne (4).
Le lemme \ref{formal-defos-lemma-minimal-groupoid-in-functors-construction} montre que (4)
entraîne (3). Cela démontre toutes les implications. La dernière assertion
d'unicité résulte du lemme \ref{formal-defos-lemma-minimal-presentations-equivalent}.
\end{proof}








\section{Unicité des anneaux versels}
\label{formal-defos-section-uniqueness}

\noindent
Étant donnés $R, S$ dans $\widehat{\mathcal{C}}_\Lambda$, on dit que des
morphismes $f, g : R \to S$ sont {\it formellement homotopes} s'il existe un
$r \geq 0$ et des morphismes $h : R \to R[[t_1, \ldots, t_r]]$ et
$k : R[[t_1, \ldots, t_r]] \to S$ dans $\widehat{\mathcal{C}}_\Lambda$ tels
que, pour tout $a \in R$, on ait
\begin{enumerate}
\item $h(a) \bmod (t_1, \ldots, t_r) = a$,
\item $f(a) = k(a)$,
\item $g(a) = k(h(a))$.
\end{enumerate}
On dit que $(r, h, k)$ est une {\it homotopie formelle} entre $f$ et $g$.

\begin{lemma}
\label{formal-defos-lemma-homotopy}
La relation « être formellement homotopes » est une relation d'équivalence sur
les ensembles de morphismes de $\widehat{\mathcal{C}}_\Lambda$.
\end{lemma}

\begin{proof}
Supposons donnés $r \geq 1$ et deux morphismes
$h_1, h_2 : R \to R[[t_1, \ldots, t_r]]$ tels que
$h_1(a) \bmod (t_1, \ldots, t_r) =  h_2(a) \bmod (t_1, \ldots, t_r) = a$
pour tout $a \in R$, ainsi qu'un morphisme
$k : R[[t_1, \ldots, t_r]] \to S$. Nous affirmons que $k \circ h_1$ est
formellement homotope à $k \circ h_2$. La symétrie inhérente à cette assertion
montrera que notre notion est symétrique. En effet, le morphisme
$$
\Psi :
R[[t_1, \ldots, t_r]]
\longrightarrow
R[[t_1, \ldots, t_r]],\quad
\sum a_I t^I \longmapsto \sum h_1(a_I)t^I
$$
est un isomorphisme. Posons $h(a) = \Psi^{-1}(h_2(a))$ pour $a \in R$ et
$k' = k \circ \Psi$ ; alors $(r, h, k')$ est une homotopie formelle entre
$k \circ h_1$ et $k \circ h_2$, ce qui démontre l'assertion.

\medskip\noindent
Supposons donnés trois morphismes $f_1, f_2, f_3 : R \to S$ comme ci-dessus,
une homotopie formelle $(r_1, h_1, k_1)$ entre $f_1$ et $f_2$, et une
homotopie formelle $(r_2, h_2, k_2)$ entre $f_3$ et $f_2$ (!). Après
renumérotation des coordonnées, on peut supposer
$h_2 : R \to R[[t_{r_1 + 1}, \ldots, t_{r_1 + r_2}]]$
et
$k_2 : R[[t_{r_1 + 1}, \ldots, t_{r_1 + r_2}]] \to S$.
En choisissant un isomorphisme convenable
$$
R[[t_1, \ldots, t_{r_1 + r_2}]]
\longrightarrow
R[[t_{r_1 + 1}, \ldots, t_{r_1 + r_2}]]
\widehat{\otimes}_{h_2, R, h_1}
R[[t_1, \ldots, t_{r_1}]]
$$
on peut insérer ces morphismes dans un diagramme commutatif
$$
\xymatrix{
R \ar[r]_{h_1} \ar[d]^{h_2} &
R[[t_1, \ldots, t_{r_1}]] \ar[d]^{h_2'} \\
R[[t_{r_1 + 1}, \ldots, t_{r_1 + r_2}]] \ar[r]^{h_1'} &
R[[t_1, \ldots, t_{r_1 + r_2}]]
}
$$
avec $h_2'(t_i) = t_i$ pour $1 \leq i \leq r_1$ et
$h_1'(t_i) = t_i$ pour $r_1 + 1 \leq i \leq r_2$. Nous omettons certains
détails. Puisque ce diagramme présente une somme amalgamée dans la catégorie
$\widehat{\mathcal{C}}_\Lambda$ (voir la démonstration du lemme
\ref{formal-defos-lemma-CLambdahat-pushouts}) et que
$k_1 \circ h_1 = f_2 = k_2 \circ h_2$, on conclut qu'il existe un morphisme
$$
k : R[[t_1, \ldots, t_{r_1 + r_2}]] \to S
$$
avec $k_1 = k \circ h_2'$ et $k_2 = k \circ h_1'$. Notons
$h = h_1' \circ h_2 = h_2' \circ h_1$. Alors
\begin{enumerate}
\item $k(h_1'(a)) = k_2(a) = f_3(a)$ ;
\item $k(h_2'(a)) = k_1(a) = f_1(a)$.
\end{enumerate}
L'assertion du premier paragraphe de la démonstration montre donc que $f_1$ et
$f_3$ sont formellement homotopes.
\end{proof}

\begin{lemma}
\label{formal-defos-lemma-composition-homotopic}
Dans la catégorie $\widehat{\mathcal{C}}_\Lambda$, si
$f_1, f_2 : R \to S$ sont formellement homotopes et si $g : S \to S'$ est un
morphisme, alors $g \circ f_1$ et $g \circ f_2$ sont formellement homotopes.
\end{lemma}

\begin{proof}
En effet, si $(r, h, k)$ est une homotopie formelle entre $f_1$ et $f_2$,
alors $(r, h, g \circ k)$ est une homotopie formelle entre $g \circ f_1$ et
$g \circ f_2$.
\end{proof}

\begin{lemma}
\label{formal-defos-lemma-versal-unique-up-to-homotopy}
Soit $\mathcal{F}$ une catégorie de déformation au-dessus de
$\mathcal{C}_\Lambda$ telle que $\dim_k T\mathcal{F} < \infty$ et
$\dim_k \text{Inf}(\mathcal{F}) < \infty$. Soit $\xi$ un objet formel versel
au-dessus de $R$. Soit $\eta$ un objet formel au-dessus de $S$. Alors deux
morphismes quelconques
$$
f, g : R \to S
$$
tels que $f_*\xi \cong \eta \cong g_*\xi$ sont formellement homotopes.
\end{lemma}

\begin{proof}
D'après le théorème \ref{formal-defos-theorem-presentation-deformation-groupoid} et sa
démonstration, $\mathcal{F}$ admet une présentation par le groupoïde en foncteurs
suivant :
$$
(\underline{R}, \underline{R_1}, s, t, c, e, i)|_{\mathcal{C}_\Lambda}
$$
Il est lisse et pro-représentable sur $\mathcal{C}_\lambda$ et présente $\mathcal{F}$. Alors les
morphismes $s : R \to R_1$ et $t : R \to R_1$ sont des homomorphismes
d'anneaux formellement lisses, et $e : R_1 \to R$ est une section. En
particulier, on peut choisir un isomorphisme
$R_1 = R[[t_1, \ldots, t_r]]$ pour un certain $r \geq 0$, tel que $s$ soit
l'inclusion $R \subset R[[t_1, \ldots, t_r]]$ et que $t$ corresponde à un
morphisme $h : R \to R[[t_1, \ldots, t_r]]$ vérifiant
$h(a) \bmod (t_1, \ldots, t_r) = a$ pour tout $a \in R$. L'existence de
l'isomorphisme $\alpha : f_*\xi \to g_*\xi$ signifie exactement qu'il existe
un morphisme $k : R_1 \to S$ tel que $f = k \circ s$ et $g = k \circ t$.
Cela signifie exactement que $(r, h, k)$ est une homotopie formelle entre $f$
et $g$.
\end{proof}

\begin{lemma}
\label{formal-defos-lemma-homotopic-minimal-prime}
Dans la catégorie $\widehat{\mathcal{C}}_\Lambda$, si
$f_1, f_2 : R \to S$ sont formellement homotopes et si
$\mathfrak p \subset R$ est un idéal premier minimal, alors
$f_1(\mathfrak p)S = f_2(\mathfrak p)S$ comme idéaux.
\end{lemma}

\begin{proof}
Supposons que $(r, h, k)$ soit une homotopie formelle entre $f_1$ et $f_2$.
Nous affirmons que
$\mathfrak pR[[t_1, \ldots, t_r]] = h(\mathfrak p)R[[t_1, \ldots, t_r]]$.
L'assertion implique le lemme en composant encore avec $k$. Pour démontrer
l'assertion, observons que l'application
$\mathfrak p \mapsto \mathfrak pR[[t_1, \ldots, t_r]]$
est une bijection de l'ensemble des idéaux premiers minimaux de $R$ sur celui
des idéaux premiers minimaux de $R[[t_1, \ldots, t_r]]$. Enfin,
$h(\mathfrak p)R[[t_1, \ldots, t_r]]$ est un idéal premier minimal puisque $h$ est
plat ; il est donc de la forme $\mathfrak q R[[t_1, \ldots, t_r]]$ pour un
certain idéal premier minimal $\mathfrak q \subset R$, d'après ce qui précède.
Mais, puisque $h \bmod (t_1, \ldots, t_r) = \text{id}_R$ par définition d'une
homotopie formelle, on conclut que $\mathfrak q = \mathfrak p$, comme voulu.
\end{proof}











\section{Changement du corps résiduel}
\label{formal-defos-section-change-of-field}

\noindent
Dans cette section, nous examinons rapidement ce qui se passe lorsque l'on
remplace le corps résiduel $k$ par une extension finie. Soit $\Lambda$ un
anneau noethérien et soit $\Lambda \to k$ un homomorphisme fini, où $k$ est un
corps. Dans tout ce chapitre, nous avons noté $\mathcal{C}_\Lambda$ la
catégorie des $\Lambda$-algèbres locales artiniennes dont le corps résiduel est
identifié à $k$ ; voir la définition \ref{formal-defos-definition-CLambda}. Cependant,
puisque nous étudierons dans cette section ce qui arrive lorsque l'on change
$k$, nous utiliserons plutôt la notation $\mathcal{C}_{\Lambda, k}$ afin de
faire apparaître la dépendance en $k$.

\begin{situation}
\label{formal-defos-situation-change-of-fields}
Soit $\Lambda$ un anneau noethérien, et soient $\Lambda \to k \to l$ des
homomorphismes d'anneaux finis, où $k$ et $l$ sont des corps. Ainsi, $l/k$ est
une extension finie de corps. Un objet typique de
$\mathcal{C}_{\Lambda, l}$ sera noté $B$, et un objet typique de
$\mathcal{C}_{\Lambda, k}$ sera noté $A$. Définissons
\begin{equation}
\label{formal-defos-equation-comparison}
\mathcal{C}_{\Lambda, l} \longrightarrow \mathcal{C}_{\Lambda, k},
\quad
B \longmapsto B \times_l k
\end{equation}
Étant donnée une catégorie cofibrée en groupoïdes
$p : \mathcal{F} \to \mathcal{C}_{\Lambda, k}$, on obtient une catégorie
cofibrée en groupoïdes associée
$$
p_{l/k} : \mathcal{F}_{l/k} \longrightarrow \mathcal{C}_{\Lambda, l}
$$
en posant $\mathcal{F}_{l/k}(B) = \mathcal{F}(B \times_l k)$.
\end{situation}

\noindent
Le foncteur (\ref{formal-defos-equation-comparison}) est bien défini : puisque
$B \times_l k \subset B$, on a
\begin{align*}
[k : k']\ \text{length}_{B \times_l k}(B \times_l k) & =
\text{length}_\Lambda(B \times_l k) \\
& \leq \text{length}_\Lambda(B) \\
& = [l : k']\ \text{length}_B(B) < \infty
\end{align*}
(voir le lemme \ref{formal-defos-lemma-length}) ; par conséquent, $B \times_l k$ est
artinien (voir Algèbre, lemme
\ref{algebra-lemma-artinian-finite-length}). Ainsi, $B \times_l k$ est un
anneau local artinien de corps résiduel $k$. Notons que
(\ref{formal-defos-equation-comparison}) commute aux produits fibrés
$$
(B_1 \times_B B_2) \times_l k =
(B_1 \times_l k) \times_{(B \times_l k)} (B_2 \times_l k)
$$
et transforme les homomorphismes d'anneaux surjectifs en homomorphismes
d'anneaux surjectifs. Nous employons la notation « lourde »
$\mathcal{F}_{l/k}$ afin d'éviter toute confusion avec la construction de la
remarque \ref{formal-defos-remark-localize-cofibered-groupoid}. Voici quelques observations
élémentaires.

\begin{lemma}
\label{formal-defos-lemma-elementary-properties-change-of-field}
Avec les notations et les hypothèses de la situation
\ref{formal-defos-situation-change-of-fields}.
\begin{enumerate}
\item On a $\overline{\mathcal{F}_{l/k}} = (\overline{\mathcal{F}})_{l/k}$.
\item Si $\mathcal{F}$ est une catégorie de prédéformation, alors
$\mathcal{F}_{l/k}$ est une catégorie de prédéformation.
\item Si $\mathcal{F}$ satisfait à (S1), alors $\mathcal{F}_{l/k}$ satisfait à
(S1).
\item Si $\mathcal{F}$ satisfait à (S2), alors $\mathcal{F}_{l/k}$ satisfait à
(S2).
\item Si $\mathcal{F}$ satisfait à (RS), alors $\mathcal{F}_{l/k}$ satisfait à
(RS).
\end{enumerate}
\end{lemma}

\begin{proof}
La partie (1) résulte immédiatement des définitions.

\medskip\noindent
Puisque $\mathcal{F}_{l/k}(l) = \mathcal{F}(k)$, la partie (2) résulte de la
définition ; voir la définition \ref{formal-defos-definition-predeformation-category}.

\medskip\noindent
La partie (3) résulte du fait que le foncteur (\ref{formal-defos-equation-comparison})
commute aux produits fibrés et transforme les morphismes surjectifs en
morphismes surjectifs ; voir la définition \ref{formal-defos-definition-S1-S2}.

\medskip\noindent
Partie (4). Pour le voir, considérons un diagramme
$$
\xymatrix{
          & l[\epsilon] \ar[d] \\
B  \ar[r] & l
}
$$
de $\mathcal{C}_{\Lambda, l}$ comme dans la définition
\ref{formal-defos-definition-S1-S2}. En appliquant le foncteur
(\ref{formal-defos-equation-comparison}), on obtient
$$
\xymatrix{
          & k[l\epsilon] \ar[d] \\
B \times_l k  \ar[r] & k
}
$$
où $l\epsilon$ désigne l'espace vectoriel de dimension finie sur $k$
$l\epsilon \subset l[\epsilon]$. D'après le lemme
\ref{formal-defos-lemma-S2-extensions}, la condition (S2) est aussi satisfaite par $\mathcal{F}$
pour ce diagramme. Ainsi, (S2) est satisfaite par $\mathcal{F}_{l/k}$.

\medskip\noindent
La partie (5) résulte de la caractérisation de (RS) donnée dans le lemme
\ref{formal-defos-lemma-RS-2-categorical}, partie (2), et du fait que
(\ref{formal-defos-equation-comparison}) commute aux produits fibrés.
\end{proof}

\noindent
Le lemme suivant s'applique notamment lorsque $\mathcal{F}$ satisfait à (S2)
et est une catégorie de prédéformation ; voir le lemme
\ref{formal-defos-lemma-S1-S2-associated-functor}.

\begin{lemma}
\label{formal-defos-lemma-tangent-space-change-of-field}
Avec les notations et les hypothèses de la situation
\ref{formal-defos-situation-change-of-fields}. Supposons que $\mathcal{F}$ soit une
catégorie de prédéformation et que $\overline{\mathcal{F}}$ satisfasse à (S2).
Il existe alors un isomorphisme canonique $l$-linéaire
$$
T\mathcal{F} \otimes_k l \longrightarrow T\mathcal{F}_{l/k}
$$
entre les espaces tangents.
\end{lemma}

\begin{proof}
D'après le lemme \ref{formal-defos-lemma-elementary-properties-change-of-field}, on peut
remplacer $\mathcal{F}$ par $\overline{\mathcal{F}}$. De plus,
$T\mathcal{F}$ et $T\mathcal{F}_{l/k}$ possèdent respectivement une structure
canonique d'espace vectoriel sur $k$ et sur $l$ ; voir le lemme
\ref{formal-defos-lemma-tangent-space-vector-space}. Alors
$$
T\mathcal{F}_{l/k} = \mathcal{F}_{l/k}(l[\epsilon])
= \mathcal{F}(k[l\epsilon]) = T\mathcal{F} \otimes_k l
$$
où la dernière égalité résulte du lemme
\ref{formal-defos-lemma-tangent-space-vector-space}. Plus généralement, étant donné un
espace vectoriel de dimension finie sur $l$, noté $V$, on a
$$
\mathcal{F}_{l/k}(l[V]) = \mathcal{F}(k[V_k]) = T\mathcal{F} \otimes_k V_k
$$
où $V_k$ désigne $V$ considéré comme espace vectoriel sur $k$. On en déduit
que les foncteurs $V \mapsto \mathcal{F}_{l/k}(l[V])$ et
$V \mapsto T\mathcal{F} \otimes_k V_k$ sont canoniquement identifiés comme
foncteurs à valeurs dans la catégorie des ensembles. Le lemme
\ref{formal-defos-lemma-linear-functor} montre qu'il existe au plus une façon de rendre l'un
ou l'autre foncteur $l$-linéaire. Les isomorphismes sont donc compatibles avec les
structures d'espaces vectoriels sur $l$, ce qui conclut.
\end{proof}

\begin{lemma}
\label{formal-defos-lemma-inf-aut-change-of-field}
Avec les notations et les hypothèses de la situation
\ref{formal-defos-situation-change-of-fields}. Supposons que $\mathcal{F}$ soit une
catégorie de déformation. Il existe alors un isomorphisme canonique
$l$-linéaire
$$
\text{Inf}(\mathcal{F}) \otimes_k l
\longrightarrow
\text{Inf}(\mathcal{F}_{l/k})
$$
entre les espaces d'automorphismes infinitésimaux.
\end{lemma}

\begin{proof}
Soit $x_0 \in \Ob(\mathcal{F}(k))$, et notons $x_{l, 0}$ l'objet correspondant
de $\mathcal{F}_{l/k}$ au-dessus de $l$. Rappelons que
$\text{Inf}(\mathcal{F}) = \text{Inf}_{x_0}(\mathcal{F})$ et
$\text{Inf}(\mathcal{F}_{l/k}) = \text{Inf}_{x_{l, 0}}(\mathcal{F}_{l/k})$ ;
voir la remarque \ref{formal-defos-remark-trivial-aut-point}. Rappelons que la structure
d'espace vectoriel sur $\text{Inf}_{x_0}(\mathcal{F})$ provient de son
identification à l'espace tangent du foncteur $\mathit{Aut}(x_0)$ défini sur la
catégorie $\mathcal{C}_{k, k}$ des $k$-algèbres locales artiniennes de corps
résiduel $k$. De même, $\text{Inf}_{x_{l, 0}}(\mathcal{F}_{l/k})$ est l'espace
tangent de $\mathit{Aut}(x_{l, 0})$, défini sur la catégorie
$\mathcal{C}_{l, l}$ des $l$-algèbres locales artiniennes de corps résiduel
$l$. En explicitant les définitions, on voit que $\mathit{Aut}(x_{l, 0})$ est
la restriction de $\mathit{Aut}(x_0)_{l/k}$ (qui est défini sur
$\mathcal{C}_{k, l}$) à $\mathcal{C}_{l, l}$. Puisqu'il n'y a aucune
différence entre l'espace tangent de $\mathit{Aut}(x_0)_{l/k}$ considéré comme
foncteur sur $\mathcal{C}_{k, l}$ ou sur $\mathcal{C}_{l, l}$, le lemme résulte
du lemme \ref{formal-defos-lemma-tangent-space-change-of-field} et du fait que
$\mathit{Aut}(x_0)$ satisfait à (RS) d'après le lemme
\ref{formal-defos-lemma-Aut-functor-RS} (d'où (S2) par le lemme
\ref{formal-defos-lemma-RS-implies-S1-S2}).
\end{proof}

\begin{lemma}
\label{formal-defos-lemma-change-of-fields-smooth}
Avec les notations et les hypothèses de la situation
\ref{formal-defos-situation-change-of-fields}. Si $\mathcal{F} \to \mathcal{G}$ est un
morphisme lisse de catégories cofibrées en groupoïdes au-dessus de
$\mathcal{C}_{\Lambda, k}$, alors
$\mathcal{F}_{l/k} \to \mathcal{G}_{l/k}$ est un morphisme lisse de catégories
cofibrées en groupoïdes au-dessus de $\mathcal{C}_{\Lambda, l}$.
\end{lemma}

\begin{proof}
Cela résulte immédiatement des définitions et du fait que
(\ref{formal-defos-equation-comparison}) préserve les surjections.
\end{proof}

\noindent
On pourrait dire bien davantage sur la relation entre $\mathcal{F}$ et
$\mathcal{F}_{l/k}$ (notamment sur la relation entre les déformations
verselles), et nous ajouterons ici ces compléments au besoin.

\begin{lemma}
\label{formal-defos-lemma-change-of-field-versal-ring}
Avec les notations et les hypothèses de la situation
\ref{formal-defos-situation-change-of-fields}. Soit $\xi$ un objet formel versel de
$\mathcal{F}$ au-dessus de
$R \in \Ob(\widehat{\mathcal{C}}_{\Lambda, k})$. Il existe alors
\begin{enumerate}
\item un $S \in \Ob(\widehat{\mathcal{C}}_{\Lambda, l})$ et un homomorphisme
local de $\Lambda$-algèbres $R \to S$, formellement lisse pour la topologie
$\mathfrak m_S$-adique et induisant sur les corps résiduels l'extension donnée
$l/k$ ;
\item un objet formel versel de $\mathcal{F}_{l/k}$ au-dessus de $S$.
\end{enumerate}
\end{lemma}

\begin{proof}
Construction de $S$. Choisissons une surjection
$R[x_1, \ldots, x_n] \to l$ de $R$-algèbres. Son noyau est un idéal maximal
$\mathfrak m$.
Prenons pour $S$ le complété $\mathfrak m$-adique de l'anneau noethérien
$R[x_1, \ldots, x_n]$. Alors $S$ appartient à
$\widehat{\mathcal{C}}_{\Lambda, l}$ d'après Algèbre, lemme
\ref{algebra-lemma-completion-Noetherian-Noetherian}. Le morphisme $R \to S$
est formellement lisse pour la topologie $\mathfrak m_S$-adique d'après
Compléments d'algèbre, lemmes \ref{more-algebra-lemma-formally-smooth} et
\ref{more-algebra-lemma-formally-smooth-completion}, et parce que
$R \to R[x_1, \ldots, x_n]$ est formellement lisse. (Comparer avec la
démonstration du lemme \ref{formal-defos-lemma-exists-smooth}.)

\medskip\noindent
Puisque $\xi$ est versel, la transformation
$\underline{\xi} : \underline{R}|_{\mathcal{C}_{\Lambda, k}} \to \mathcal{F}$
est lisse. D'après le lemme \ref{formal-defos-lemma-change-of-fields-smooth}, le morphisme
induit
$$
(\underline{R}|_{\mathcal{C}_{\Lambda, k}})_{l/k}
\longrightarrow
\mathcal{F}_{l/k}
$$
est lisse. Il suffit donc de construire un morphisme lisse
$\underline{S}|_{\mathcal{C}_{\Lambda, l}} \to
(\underline{R}|_{\mathcal{C}_{\Lambda, k}})_{l/k}$. Donner un tel morphisme
revient à donner, pour tout objet $B$ de $\mathcal{C}_{\Lambda, l}$, une
application d'ensembles
$$
\Mor_{\widehat{\mathcal{C}}_{\Lambda, l}}(S, B)
\longrightarrow
\Mor_{\widehat{\mathcal{C}}_{\Lambda, k}}(R, B \times_l k)
$$
fonctorielle en $B$. Un élément $\varphi : S \to B$ du membre de gauche est
envoyé sur le composé $R \to S \to B$, dont l'image est contenue dans la
sous-$\Lambda$-algèbre $B \times_l k$. La lissité du morphisme signifie
qu'étant donnés une surjection $B' \to B$ et un diagramme commutatif
$$
\xymatrix{
S \ar[r] & B \ar@{=}[r] & B \\
R \ar[u] \ar[r] & B' \times_l k \ar[u] \ar[r] & B' \ar[u]
}
$$
il faut trouver un homomorphisme d'anneaux $S \to B'$ s'insérant dans le
rectangle extérieur. L'existence de ce morphisme est garantie par le fait que
nous avons choisi $R \to S$ formellement lisse pour la topologie $\mathfrak m_S$-adique ; voir
Compléments d'algèbre, lemme \ref{more-algebra-lemma-lift-continuous}.
\end{proof}
























\bibliography{my}
\bibliographystyle{amsalpha}
\end{document}
